% Preamble for the swarm-transcribed Weber volumes.
% Deliberately close to weber_v1_ge.tex so the verified p1-148 region and the swarm
% region can live in one document.
\documentclass[11pt]{article}
\usepackage[ngerman,english]{babel}
\usepackage[a4paper,margin=0.9in]{geometry}
\usepackage[T1]{fontenc}
\usepackage[utf8]{inputenc}
\usepackage{lmodern}
\usepackage[leqno]{amsmath}
\usepackage{amssymb,mathtools,mathrsfs}
\usepackage{microtype}
\usepackage{fancyhdr}
\usepackage{array,longtable,enumitem,multicol}
\usepackage{tikz}
\usepackage{graphicx}
\graphicspath{{../figures/b1/}{../figures/b2/}{../figures/b3/}}
\setlength{\parindent}{1.5em}
\setlength{\parskip}{0.18em}
\emergencystretch=6em
\setlength{\headheight}{14pt}
\pagestyle{fancy}
\fancyhf{}
\cfoot{\thepage}

% --- Weber conventions -------------------------------------------------------
% Sperrung (letter-spacing) is Weber's emphasis. There is no bold in his body text.
\providecommand{\gesperrt}[1]{\textls[130]{#1}}
% Printed-page marker: a no-op in the PDF. The machine-readable anchor is the
% "% !PRINTED n" comment in the source, which the collector and Codex both read.
\providecommand{\weberpage}[1]{}
% Displayed equation numbered in the LEFT margin, as Weber prints them.
\providecommand{\weq}[2]{\begin{equation}\tag{#1}#2\end{equation}}
% Figure placeholder. A later pass replaces it with TikZ or a cleaned scan crop.
% A marked gap is honest; a fabricated diagram is not.
\providecommand{\weberfigure}[2]{%
  \begin{center}\fbox{\parbox{0.78\linewidth}{\centering\small [Figur, S.~#1: #2]}}\end{center}}
% Source-faithful figure crop.  The asset search path works both from the live
% swarm_build directory and from a self-contained release bundle with assets/.
\providecommand{\weberfigureimage}[2]{%
  \begin{center}\includegraphics[width=#2\linewidth]{#1}\end{center}}

% Text symbols agents reach for that are not always defined. Providing them in the preamble
% fixes every occurrence at once and leaves the generated fragments untouched, which is the
% right trade whenever the fix is "define the macro" rather than "change the text".
\usepackage{textcomp}
% amsmath caps matrix environments at 10 columns. Weber prints permutation matrices far wider
% than that -- Bd III has an 11-column pmatrix of the w_nu -- and the failure is the opaque
% "Extra alignment tab has been changed to \cr".
\setcounter{MaxMatrixCols}{30}
\providecommand{\textquotedblbase}{\glqq}
\providecommand{\textquotedblleft}{``}
\providecommand{\textquotedblright}{''}
\providecommand{\textquotesinglbase}{\glq}

% --- Greek capitals are SLANTED in this fount -------------------------------------------
% RULING 2026-08-07, after the same observation was reported from five separate pages
% (printed 271, 275, 285, 333 and the nu/v cousin on 172).
%
% Verified at 8x zoom on printed 333 (b1 idx 358): every capital Delta on that page --
% "f(x) - f(alpha) = Delta", "Delta_alpha = f(alpha+delta) - f(alpha)" -- is the INCLINED
% form, which LaTeX calls \varDelta. \Delta, LaTeX's default, is upright and is not what
% Vieweg set. The same inclination is visible on capital Phi (printed 171, 282).
%
% This is a property of the TYPEFACE, not of any occurrence -- the same kind of fact as the
% varkappa ruling -- so it holds wherever the letter appears, and the right fix is one line
% here rather than several thousand edits across the fragments. \let leaves every fragment
% untouched, applies to the whole corpus at once, and is reversible by deleting these lines.
%
% Delta and Phi are the two verified by eye. The rest follow from its being a single fount;
% if any capital is later found upright on the page, delete its line and the fragments still
% say exactly what they said before.
\let\Gamma\varGamma
\let\Delta\varDelta
\let\Theta\varTheta
\let\Lambda\varLambda
\let\Xi\varXi
\let\uprightPi\Pi
\let\Pi\varPi
\let\Sigma\varSigma
\let\Upsilon\varUpsilon
\let\Phi\varPhi
\let\Psi\varPsi
\let\Omega\varOmega
% Lower-case rho is the OPEN-TAILED form in this fount. Verified at 10x on printed 419
% ("wenn also überhaupt ein \varrho vorhanden ist"): the tail curves down and left rather than
% descending straight, which is \varrho, not \rho. Same kind of fact as the capitals above, and
% the same one-line treatment. (kappa is already handled: the fragments write \varkappa
% explicitly, on the ruling recorded in resolve_class.py.)
\let\rho\varrho

% --- Fractions written in running text -------------------------------------------------
% Gauss's German prose says things like "von 6\tfrac{1}{2} Jahren" and "\tfrac{1}{2} Grad":
% a fraction as an ordinary word, not inside math. Bare \tfrac there is
%     ! Missing $ inserted.
% and the error points at the sentence, not at the macro.
%
% Wrapping the three fraction macros in \ensuremath fixes every occurrence at once and leaves
% the fragments untouched -- the right trade whenever the fix is "define the macro" rather than
% "change the text", as with the Greek capitals above. Inside math \ensuremath is a no-op, so
% the thousands of legitimate in-math uses are unaffected.
\let\swarmfrac\frac
\let\swarmtfrac\tfrac
\let\swarmdfrac\dfrac
\renewcommand{\frac}[2]{\ensuremath{\swarmfrac{#1}{#2}}}
\renewcommand{\tfrac}[2]{\ensuremath{\swarmtfrac{#1}{#2}}}
\renewcommand{\dfrac}[2]{\ensuremath{\swarmdfrac{#1}{#2}}}

% --- Unicode vulgar fractions -----------------------------------------------------------
% Agents sometimes type the single Unicode glyph the page shows -- "erstere ⅔" -- rather than a
% macro. utf8 inputenc has no default binding for U+2153..U+215E and stops the build with
%     ! LaTeX Error: Unicode character ⅔ (U+2154)
% Declaring them costs nothing, keeps the fragment reading exactly as the agent saw the page,
% and is the same "define it in the preamble" trade as the fraction macros above. Typing the
% glyph is not wrong; it is arguably closer to the page than \tfrac would be.
\DeclareUnicodeCharacter{2153}{\ensuremath{\swarmtfrac{1}{3}}}
\DeclareUnicodeCharacter{2154}{\ensuremath{\swarmtfrac{2}{3}}}
\DeclareUnicodeCharacter{2155}{\ensuremath{\swarmtfrac{1}{5}}}
\DeclareUnicodeCharacter{2156}{\ensuremath{\swarmtfrac{2}{5}}}
\DeclareUnicodeCharacter{2157}{\ensuremath{\swarmtfrac{3}{5}}}
\DeclareUnicodeCharacter{2158}{\ensuremath{\swarmtfrac{4}{5}}}
\DeclareUnicodeCharacter{2159}{\ensuremath{\swarmtfrac{1}{6}}}
\DeclareUnicodeCharacter{215A}{\ensuremath{\swarmtfrac{5}{6}}}
\DeclareUnicodeCharacter{215B}{\ensuremath{\swarmtfrac{1}{8}}}
\DeclareUnicodeCharacter{215C}{\ensuremath{\swarmtfrac{3}{8}}}
\DeclareUnicodeCharacter{215D}{\ensuremath{\swarmtfrac{5}{8}}}
\DeclareUnicodeCharacter{215E}{\ensuremath{\swarmtfrac{7}{8}}}
% Prime and double prime: arcminutes and arcseconds. Gauss's geodesy and astronomy volumes are
% full of "43° 37' 6\"741", and agents type the true Unicode prime characters rather than the
% ASCII apostrophe -- which is right, since that is what the page prints.
\DeclareUnicodeCharacter{2032}{\ensuremath{{}^{\prime}}}
\DeclareUnicodeCharacter{2033}{\ensuremath{{}^{\prime\prime}}}
\DeclareUnicodeCharacter{2034}{\ensuremath{{}^{\prime\prime\prime}}}
\DeclareUnicodeCharacter{00BD}{\ensuremath{\swarmtfrac{1}{2}}}
\DeclareUnicodeCharacter{00BC}{\ensuremath{\swarmtfrac{1}{4}}}
\DeclareUnicodeCharacter{00BE}{\ensuremath{\swarmtfrac{3}{4}}}

% Macros agents invent for a text fraction. \textfrac is not a real LaTeX command, but the
% intent is unambiguous and providing it is cheaper and safer than rewriting the fragment --
% the alternative is editing a table cell by hand and risking its alignment.
\providecommand{\textfrac}[2]{\ensuremath{\swarmtfrac{#1}{#2}}}

% Common shorthands inherited from the volume I file, so fragments that use them compile.
\providecommand{\eps}{\varepsilon}
\providecommand{\ii}{\mathrm{i}}
\providecommand{\tg}{\operatorname{tg}}
% Gauss writes the trigonometric functions in the period's Latin abbreviations -- "tang" for
% tangent, "cotang" for cotangent -- rather than Weber's "tg"/"cotg". Same function, different
% century's shorthand; both are reproduced as the author sets them.
\providecommand{\tang}{\operatorname{tang}}
\providecommand{\cotang}{\operatorname{cotang}}
\providecommand{\arctang}{\operatorname{arc\,tang}}
\providecommand{\ctg}{\operatorname{cotg}}
\providecommand{\cotg}{\operatorname{cotg}}
\providecommand{\sn}{\operatorname{sn}}
\providecommand{\cn}{\operatorname{cn}}
\providecommand{\dn}{\operatorname{dn}}
\providecommand{\am}{\operatorname{am}}
\providecommand{\dd}{\,d}
\providecommand{\tht}{\vartheta}
\providecommand{\vark}{\varkappa}
\providecommand{\Norm}{\operatorname{N}}
\providecommand{\ord}{\operatorname{ord}}
\providecommand{\ind}{\operatorname{ind}}
\providecommand{\abs}[1]{\left|#1\right|}
\providecommand{\sect}[2]{\section*{\S.\,#1. #2}}
\providecommand{\frp}{\mathfrak p}
\providecommand{\frP}{\mathfrak P}
\providecommand{\frG}{\mathfrak G}
\providecommand{\frA}{\mathfrak A}
\providecommand{\frB}{\mathfrak B}

% --- reader additions ---------------------------------------------------------
\usepackage{url}
\usepackage[unicode,bookmarksnumbered=false,bookmarksopen=true,bookmarksopenlevel=1,colorlinks=true,
  linkcolor=black,urlcolor=blue]{hyperref}
\usepackage{bookmark}
\renewcommand{\headrulewidth}{0pt}
\renewcommand{\weberpage}[1]{\phantomsection\label{wp:#1}}
\newcommand{\inhaltkopf}[1]{\begin{center}\large\textbf{#1}\end{center}\medskip}
\newcommand{\inhaltbuch}[2]{\par\medskip\begin{center}\if\relax\detokenize{#1}\relax\else#1\\[2pt]\fi{\large\textbf{#2}}\end{center}}
\newcommand{\inhaltabschnitt}[2]{\par\smallskip\begin{center}\small\if\relax\detokenize{#1}\relax\else#1\\\fi\textbf{#2}\end{center}}
\newcommand{\inhaltpar}[3]{\par\noindent\hangindent=3.4em\hbox to 3.4em{\S~#1.\hfil}#2\dotfill\,\hyperref[wp:#3]{#3}\par}
\newcommand{\inhaltzeile}[2]{\par\noindent #1\dotfill\,\hyperref[wp:#2]{#2}\par}

\begin{document}
\selectlanguage{english}
\pagenumbering{roman}
\begin{center}
{\Large Heinrich Weber}\\[6pt]
{\LARGE\textbf{Lehrbuch der Algebra}}\\[4pt]
{\large (Textbook of Algebra)}\\[8pt]
{\large Volume II}\\[4pt]
Braunschweig: Friedrich Vieweg und Sohn, 1896\\[24pt]
{\large English translation}\\[4pt]
with the page numbers of the original
\end{center}
\clearpage
\noindent Translated from the scan \url{https://archive.org/details/lehrbuchderalgeb02weberich} (Internet Archive).
The German text was transcribed from the scan and translated into English by AI models (Anthropic Claude and
OpenAI Codex), and both were checked against the scan by AI. Weber's text is in the public domain; the
transcription, translation and typesetting are dedicated to the public domain under CC0~1.0 (\url{https://creativecommons.org/publicdomain/zero/1.0/}).
Source files: \url{https://github.com/KokunoYumeto/weber-lehrbuch-der-algebra}.
\clearpage

% ===== front matter, scan leaf 6 =====
\begin{center}
\vspace*{5cm}
{\Large LEHRBUCH}\\[1cm]
{\small DER}\\[1.2cm]
{\Huge \gesperrt{ALGEBRA.}}\\[0.8cm]
{\large [Textbook of Algebra]}\\[1.2cm]
\rule{3cm}{0.4pt}
\end{center}

\clearpage
% ===== front matter, scan leaf 8 =====
\begin{center}
\vspace*{2cm}
{\LARGE LEHRBUCH}\\[1.2cm]
{\small DER}\\[1.2cm]
{\Huge \gesperrt{ALGEBRA}.}\\[0.5cm]
{\large [Textbook of Algebra]}\\[1.5cm]
{\large \textsc{By}}\\[1.2cm]
{\LARGE HEINRICH WEBER,}\\[0.2cm]
{\footnotesize PROFESSOR OF MATHEMATICS AT THE UNIVERSITY OF STRASSBURG.}\\[1.2cm]
\rule{2.5cm}{0.4pt}\\[1.2cm]
{\large \gesperrt{IN TWO VOLUMES}.}\\[0.8cm]
\rule{2cm}{0.4pt}\\[0.6cm]
{\large \textbf{SECOND VOLUME.}}\\[3.5cm]
\rule{\textwidth}{1pt}\\[0.8cm]
{\Large BRAUNSCHWEIG,}\\[0.3cm]
{\small PRINTED AND PUBLISHED BY FRIEDRICH VIEWEG AND SON.}\\[0.3cm]
{\Large 1896.}
\end{center}

\clearpage
% ===== front matter, scan leaf 9 =====
\begin{center}
\vspace*{0.35\textheight}
\rule{0.5\textwidth}{0.4pt}\\[1ex]
\gesperrt{All rights reserved.}\\[0.5ex]
\rule{0.5\textwidth}{0.4pt}
\end{center}

\clearpage
% ===== front matter, scan leaf 10 =====
\begin{center}
{\large PREFACE TO THE SECOND VOLUME.}

\vspace{8pt}
\rule{2.5cm}{0.4pt}
\end{center}

\vspace{6pt}

In accordance with the intention announced in the preface to the first volume, I am today able to hand over the second volume of my Lehrbuch der Algebra to the public. The plan set out there has been carried out in its essential points. In the applications I have endeavoured to choose such problems as have already gained an independent interest in other fields, geometry or the theory of functions, and which at the same time illustrate the main points of the algebraic theory from as many sides as possible.

The application of the theory of algebraic numbers has been carried as far as the theory of the cyclotomic numbers. If life and capacity for work hold out, I hope to present, in a continuation of my work, the further applications to the field of elliptic functions, which are contained only in part in my book ``Elliptische Functionen und algebraische Zahlen'' [Elliptic Functions and Algebraic Numbers].

During the working out and the printing of the second volume, too, I have had the help and advice of the friends whom I already named in the preface to the first edition. But the first volume has also already won many a new friend, who has furthered my work with hints and advice. To all of them I express my thanks here, and I add the request that they may continue to keep their interest in the work.

\bigskip

\gesperrt{Strassburg}, July 1896.

\begin{flushright}
\textbf{\gesperrt{The Author}.}\hspace*{2em}
\end{flushright}

\clearpage
% ===== front matter, scan leaf 12 =====
\inhaltkopf{CONTENTS OF THE SECOND VOLUME.}

\inhaltbuch{First Book.}{Groups.}

\inhaltabschnitt{First Section.}{General Group Theory.}
\inhaltpar{1}{Definition of Groups}{3}
\inhaltpar{2}{The Divisors of Finite Groups}{7}
\inhaltpar{3}{Normal Divisors of a Group}{10}
\inhaltpar{4}{Composition of the Parts}{12}
\inhaltpar{5}{Multiple Isomorphism}{15}
\inhaltpar{6}{The Composition Series and the Theorem of C. \gesperrt{Jordan}}{17}
\inhaltpar{7}{Further Theorems on the Composition Series}{24}
\inhaltpar{8}{Metacyclic Groups}{27}

\inhaltabschnitt{Second Section.}{Abelian Groups.}
\inhaltpar{9}{Representation of Abelian Groups by a Basis}{32}
\inhaltpar{10}{The Invariants of the Abelian Groups}{39}
\inhaltpar{11}{Group Characters}{43}
\inhaltpar{12}{Divisors of an Abelian Group. Reciprocal Groups}{48}
\inhaltpar{13}{The Genera in an Abelian Group}{52}
\inhaltpar{14}{Indices with Respect to an Odd Prime Power as Modulus}{54}
\inhaltpar{15}{Indices for a Power of 2 as Modulus}{58}
\inhaltpar{16}{The Groups of Number Classes with Respect to a Composite Modulus}{60}

\inhaltabschnitt{Third Section.}{The Group of the Cyclotomic Fields.}
\inhaltpar{17}{The Resolvents of the Theory of Cyclotomy}{63}
\inhaltpar{18}{Cyclotomic Fields}{67}
\inhaltpar{19}{Primary and Non-primary Divisors of the Group $\mathfrak{N}$}{72}
\inhaltpar{20}{The Cyclotomic Periods}{74}
\inhaltpar{21}{Cyclotomic Fields with a Given Group}{79}
\inhaltpar{22}{Determination of the Group $\mathfrak{A}$}{92}

\clearpage
% ===== front matter, scan leaf 13 =====
% !PRINTED VIII
\weberpage{VIII}
\inhaltabschnitt{Fourth Section.}{Cubic and Biquadratic Abelian Fields.}
\inhaltpar{23}{Cubic cyclotomic fields}{94}
\inhaltpar{24}{Biquadratic cyclotomic fields}{101}
\inhaltpar{25}{Cubic Abelian equations}{107}
\inhaltpar{26}{Biquadratic Abelian equations}{110}
\inhaltabschnitt{Fifth Section.}{Constitution of General Groups.}
\inhaltpar{27}{Construction of groups according to Cayley}{114}
\inhaltpar{28}{Relation of general groups to permutation groups}{117}
\inhaltpar{29}{The first Sylow theorem}{121}
\inhaltpar{30}{The second Sylow theorem}{125}
\inhaltpar{31}{Groups of degree $p^\alpha$}{127}
\inhaltpar{32}{Theorem of Frobenius}{129}
\inhaltpar{33}{Groups of degree $p^\alpha q$}{134}
\inhaltpar{34}{Simple groups}{136}
\inhaltpar{35}{Groups of degree $p q$}{140}
\inhaltpar{36}{Bounds for the index of a divisor of the symmetric permutation group}{143}
\inhaltbuch{Second Book.}{Linear Groups.}
\inhaltabschnitt{Sixth Section.}{Groups of Linear Substitutions.}
\inhaltpar{37}{Linear substitutions and their composition}{151}
\inhaltpar{38}{Substitution of the ratios}{158}
\inhaltpar{39}{Permutations as linear substitutions}{159}
\inhaltpar{40}{The invariants of finite groups of linear substitutions}{161}
\inhaltpar{41}{The theorem of Hilbert}{165}
\inhaltpar{42}{Finiteness of the system of invariants of a finite linear substitution group}{168}
\inhaltpar{43}{The form problem}{171}
\inhaltpar{44}{Klein's extension of the fundamental algebraic problem}{176}
\inhaltpar{45}{Influence of relative invariants}{177}
\inhaltpar{46}{The extended concept of invariant}{179}
\inhaltpar{47}{Normal forms}{181}
\inhaltabschnitt{Seventh Section.}{Groups of Binary Linear Substitutions.}
\inhaltpar{48}{Ternary orthogonal substitutions}{184}
\inhaltpar{49}{Linear fractional substitutions}{189}

\clearpage
% ===== front matter, scan leaf 14 =====
% !PRINTED IX
\weberpage{IX}
\inhaltpar{50}{Conditions of reality}{193}
\inhaltpar{51}{Finite groups of linear fractional substitutions. Poles of the groups}{195}
\inhaltpar{52}{The different kinds of possible groups}{199}
\inhaltpar{53}{Transformation of the substitutions of $G$ to simple forms}{203}
\inhaltpar{54}{The fundamental forms}{205}

\inhaltabschnitt{Eighth Section.}{The Polyhedral Groups.}
\inhaltpar{55}{The cyclic groups and the dihedral groups}{209}
\inhaltpar{56}{The tetrahedral group}{212}
\inhaltpar{57}{The octahedral group}{216}
\inhaltpar{58}{The icosahedral group}{220}
\inhaltpar{59}{The divisors of the icosahedral group}{227}
\inhaltpar{60}{The fundamental forms of the icosahedral group}{230}
\inhaltpar{61}{The invariants of the icosahedron}{232}
\inhaltpar{62}{Polyhedral groups of the second kind. Crystallographic groups}{234}

\inhaltabschnitt{Ninth Section.}{Congruence Groups.}
\inhaltpar{63}{Function congruences}{242}
\inhaltpar{64}{Congruence fields}{245}
\inhaltpar{65}{Congruence groups in the field $\mathfrak{E}$}{250}
\inhaltpar{66}{Simplicity of the Group $E$}{254}
\inhaltpar{67}{Congruence fields of the second degree}{259}
\inhaltpar{68}{The real linear congruence group $L_p$}{261}
\inhaltpar{69}{Imaginary form of the group $L_p$}{266}
\inhaltpar{70}{Divisors of the group $L_p$ whose degree is divisible by $p$}{271}
\inhaltpar{71}{Divisors of the group $L_p$ whose degree is not divisible by $p$}{273}
\inhaltpar{72}{Constitution of the group $L_7$ of degree 168}{282}

\inhaltbuch{Third Book.}{Applications of Group Theory.}

\inhaltabschnitt{Tenth Section.}{General Theory of Metacyclic Equations.}
\inhaltpar{73}{The resolvents of the composition series}{289}
\inhaltpar{74}{Metacyclic equations}{292}
\inhaltpar{75}{Metacyclic equations whose degree is a power of a prime}{297}
\inhaltpar{76}{Representation of the Abelian group $Q$}{298}
\inhaltpar{77}{Analytic Representation of the Permutations}{299}
\inhaltpar{78}{Representation of the metacyclic group $P$}{301}
\inhaltpar{79}{Ternary linear congruence group for the modulus 2}{306}
\inhaltpar{80}{Resolvents of the equation of the 8\textsuperscript{th} degree}{308}

\clearpage
% ===== front matter, scan leaf 15 =====
% !PRINTED X
\weberpage{X}
\inhaltpar{81}{The triple systems of the resolvents}{310}
\inhaltpar{82}{Application to equations of the 8\textsuperscript{th} degree}{313}
\inhaltpar{83}{Metacyclic equations of the 8\textsuperscript{th} degree}{314}
\inhaltpar{84}{Biquadratic equations}{319}
\inhaltabschnitt{Eleventh Section.}{The Points of Inflection of a Curve of the Third Order.}
\inhaltpar{85}{Ternary forms and algebraic curves}{322}
\inhaltpar{86}{Singular points. Points of inflection. Double tangents}{324}
\inhaltpar{87}{Fundamental covariants of a ternary form}{327}
\inhaltpar{88}{The \gesperrt{Hessian} curve}{329}
\inhaltpar{89}{Inflexion points of a curve of the third order}{331}
\inhaltpar{90}{Transformation of the cubic form to the canonical form}{333}
\inhaltpar{91}{The invariants of the curve of the third order and the biquadratic equation}{336}
\inhaltpar{92}{Triple equations}{342}
\inhaltpar{93}{The group of the triple equations}{344}
\inhaltpar{94}{Reality conditions of the triple equations}{348}
\inhaltabschnitt{Twelfth Section.}{The Double Tangents of a Curve of the Fourth Order.}
\inhaltpar{95}{Number of the double tangents of a curve of the fourth order}{351}
\inhaltpar{96}{The \gesperrt{Steiner} complexes}{357}
\inhaltpar{97}{Pairs and triples of complexes}{363}
\inhaltpar{98}{The \gesperrt{Aronhold} systems of seven}{367}
\inhaltpar{99}{The \gesperrt{Hesse-Cayley} algorithm for the designation of the double tangents}{369}
\inhaltpar{100}{Rational determination of the curve from a complete system of seven}{374}
\inhaltpar{101}{The \gesperrt{Galois} group of the problem of the double tangents}{380}
\inhaltpar{102}{Representation of the group}{384}
\inhaltpar{103}{Simplicity of the group of the problem of the double tangents}{387}
\inhaltpar{104}{Reality of the double tangents}{391}
\inhaltpar{105}{Proof of the existence of the four cases}{398}
\inhaltabschnitt{Thirteenth Section.}{General Theory of the Equation of the Fifth Degree.}
\inhaltpar{106}{Statement of the problem}{403}
\inhaltpar{107}{\gesperrt{Lüroth}'s theorem}{404}
\inhaltpar{108}{Resolvents with one parameter}{407}
\inhaltpar{109}{Group of the resolvents with one parameter}{410}
\inhaltpar{110}{The icosahedral equation}{415}
\inhaltpar{111}{The resolvents of the icosahedral equation}{419}
\inhaltpar{112}{The principal resolvent of the fifth degree}{422}

\clearpage
% ===== front matter, scan leaf 16 =====
% !PRINTED XI
\weberpage{XI}
\inhaltpar{113}{Resolvents of the sixth degree}{426}
\inhaltpar{114}{On the solution of the icosahedral equation by transcendental functions}{429}
\inhaltabschnitt{Fourteenth Section.}{Groups of Linear Ternary Substitutions.}
\inhaltpar{115}{Ternary linear substitution group of the 168\textsuperscript{th} degree}{433}
\inhaltpar{116}{Poles and axes of the ternary groups}{438}
\inhaltpar{117}{Application to the group $G_{168}$. Sevenfold poles}{442}
\inhaltpar{118}{The principal axes}{444}
\inhaltpar{119}{The threefold and sixfold poles}{448}
\inhaltpar{120}{The configuration of the group $G_{168}$}{451}
\inhaltpar{121}{Invariant curves of the group $G_{168}$}{453}
\inhaltpar{122}{The first invariant of the group $G_{168}$ and the fundamental curve}{454}
\inhaltpar{123}{The higher invariants}{459}
\inhaltpar{124}{The complete system of invariants}{461}
\inhaltabschnitt{Fifteenth Section.}{The Form Problem of the Group $G_{168}$ and the Theory of Equations of the Seventh Degree.}
\inhaltpar{125}{The resolvents of the form problem}{466}
\inhaltpar{126}{Reduction of the general resolvent of the seventh degree to the special one}{471}
\inhaltpar{127}{Permutation group of seven digits of degree 168}{473}
\inhaltpar{128}{Equations of the seventh degree with a group of the 168\textsuperscript{th} degree}{476}
\inhaltpar{129}{Contragredient groups}{478}
\inhaltpar{130}{Solution of the equation of the seventh degree with the group $P_{168}$ by means of the form problem of the group $G_{168}$}{481}
\inhaltpar{131}{Possibility of determining the functions $X_1, X_2, X_3$}{484}
\inhaltbuch{Fourth Book.}{Algebraic Numbers.}
\inhaltabschnitt{Sixteenth Section.}{Numbers and Functionals of an Algebraic Field.}
\inhaltpar{132}{Definition of algebraic numbers}{489}
\inhaltpar{133}{Algebraic integers}{490}
\inhaltpar{134}{Algebraic fields}{493}
\inhaltpar{135}{Integral functions in an algebraic field}{495}
\inhaltpar{136}{The functionals of an algebraic field $\Omega$ and the extended field $\overline{\Omega}$}{499}
\inhaltpar{137}{Integral functionals}{504}
\inhaltpar{138}{Divisibility. Associated functionals. Units}{509}
\inhaltpar{139}{Greatest common divisor}{512}

\clearpage
% ===== front matter, scan leaf 17 =====
% !PRINTED XII
\weberpage{XII}

\inhaltpar{140}{Prime Functionals in the Field $\Omega$}{515}
\inhaltpar{141}{Decomposition of the Integral and Fractional Functionals into Prime Factors}{516}
\inhaltpar{142}{Integral Functions in the Field $\overline{\Omega}$}{520}
\inhaltpar{143}{The Prime Factors of the Numbers of the Field $\Omega$}{523}

\inhaltabschnitt{Seventeenth Section.}{Theory of Algebraic Fields.}

\inhaltpar{144}{Basis of an Algebraic Number Field. Discriminants}{527}
\inhaltpar{145}{The Minimal Basis and the Field Discriminant}{529}
\inhaltpar{146}{The Bases of the Functionals}{532}
\inhaltpar{147}{The Absolute Norms of the Functionals}{536}
\inhaltpar{148}{Complete Residue System with Respect to a Modulus}{539}
\inhaltpar{149}{Congruences}{541}
\inhaltpar{150}{\gesperrt{Fermat}'s Theorem}{544}
\inhaltpar{151}{The \gesperrt{Dedekind} Ideals}{547}
\inhaltpar{152}{Equivalence}{552}
\inhaltpar{153}{The Class Number of the Field $\Omega$}{554}
\inhaltpar{154}{The Group of Ideal Classes}{557}

\inhaltabschnitt{Eighteenth Section.}{Discriminants.}

\inhaltpar{155}{Minimum of a Quadratic Form}{559}
\inhaltpar{156}{Application to Algebraic Fields}{566}
\inhaltpar{157}{Prime Factors of the Natural Prime Numbers}{569}
\inhaltpar{158}{\gesperrt{Dedekind}'s Theorem on the Field Discriminant}{578}

\inhaltabschnitt{Nineteenth Section.}{Relation of a Field to Its Divisors.}

\inhaltpar{159}{Relative Norms}{583}
\inhaltpar{160}{Prime Ideals in the Relatively Normal Field}{588}
\inhaltpar{161}{Primitive Roots of the Prime Ideals}{591}
\inhaltpar{162}{The Partial Fundamental Ideal}{593}
\inhaltpar{163}{The Ideals in the Divisors of the Field $\Omega$}{597}

\inhaltabschnitt{Twentieth Section.}{Quadratic Fields.}

\inhaltpar{164}{Basis of a Quadratic Field}{601}
\inhaltpar{165}{The Prime Ideals in $\Omega$}{603}
\inhaltpar{166}{Functionals in the Quadratic Field and Quadratic Irrational Numbers}{607}
\inhaltpar{167}{Equivalent Functionals and Equivalent Numbers}{610}
\inhaltpar{168}{Composition of the Quadratic Irrational Numbers}{613}

\clearpage
% ===== front matter, scan leaf 18 =====
% !PRINTED XIII
\weberpage{XIII}

\inhaltabschnitt{Twenty-first Section.}{Cyclotomic Fields.}
\inhaltpar{169}{Decomposition of the prime number $q$ in the cyclotomic field $\Omega_{q^\varkappa}$}{617}
\inhaltpar{170}{Minimal basis of the field $\Omega_m$}{620}
\inhaltpar{171}{The prime ideals in the field $\Omega_m$}{622}
\inhaltpar{172}{The conjugate prime ideals}{623}
\inhaltpar{173}{Representation of the prime factors of $p$}{628}
\inhaltpar{174}{\gesperrt{Kummer}'s theorem}{632}
\inhaltpar{175}{The roots of unity in the field $\Omega_m$}{638}
\inhaltpar{176}{The real field $H_m$ contained in $\Omega_m$}{641}
\inhaltpar{177}{The prime ideals of the field $H_m$}{643}
\inhaltpar{178}{The units of the field $H_m$}{645}

\inhaltabschnitt{Twenty-second Section.}{Abelian Fields and Cyclotomic Fields.}
\inhaltpar{179}{Simple \gesperrt{Abel}ian fields}{648}
\inhaltpar{180}{The resolvents}{651}
\inhaltpar{181}{Preparation for the Proof}{654}
\inhaltpar{182}{Proof of the first lemma for an odd $m$}{659}
\inhaltpar{183}{Proof of the second lemma for an odd $m$}{666}
\inhaltpar{184}{Preliminary remarks on the case of an even $m$}{667}

\inhaltabschnitt{Twenty-third Section.}{Class Number.}
\inhaltpar{185}{\gesperrt{Dirichlet}'s theorem on the units}{670}
\inhaltpar{186}{Systems of independent units and systems of exponents of the units}{677}
\inhaltpar{187}{Fundamental systems of units}{681}
\inhaltpar{188}{Reduced numbers}{684}
\inhaltpar{189}{Bounds for the number of integers of the field $\Omega$ divisible by an ideal}{686}
\inhaltpar{190}{Determination of the volume}{690}
\inhaltpar{191}{Theorems from the theory of series}{694}
\inhaltpar{192}{Application to the determination of the class number}{701}

\inhaltabschnitt{Twenty-fourth Section.}{Class Number of the Cyclotomic Fields.}
\inhaltpar{193}{Representation of the class number in the cyclotomic field $\Omega_m$}{705}
\inhaltpar{194}{Determination of the sums $X$}{708}
\inhaltpar{195}{On the class number in the real field contained in $\Omega_m$}{711}
\inhaltpar{196}{Class number in the field of the eighth roots of unity}{717}
\inhaltpar{197}{Recursive computation of the class number in the field $\Omega_m$ when $m$ is a power of 2}{719}
\inhaltpar{198}{The class number factor $A$}{722}
\inhaltpar{199}{The class number factor $B$}{726}

\clearpage
% ===== front matter, scan leaf 19 =====
% !PRINTED XIV
\weberpage{XIV}
\inhaltpar{200}{Normal units in $H_m$}{728}
\inhaltpar{201}{Fundamental system of units of the field $H_m$}{735}
\inhaltpar{202}{Positive units}{742}
\inhaltabschnitt{Twenty-fifth Section.}{Transcendental Numbers.}
\inhaltpar{203}{Countable sets}{745}
\inhaltpar{204}{Uncountable sets}{748}
\inhaltpar{205}{Transcendence of the number $e$}{751}
\inhaltpar{206}{Transcendence of the number $\pi$}{756}
\inhaltpar{207}{The general theorem of \gesperrt{Lindemann} on the exponential function}{760}
\inhaltbuch{}{Addenda.}
\inhaltzeile{I. Irreducibility of the cyclotomic equation}{769}
\inhaltzeile{II. The irreducibility of the cyclotomic equation and the theorem on the primes contained in a linear form}{777}
\inhaltzeile{Index to Volumes I and II}{785}
\inhaltzeile{Corrections}{795}

\clearpage
\pagenumbering{arabic}

% ===== scan leaf 22, printed 1 =====
% !PRINTED 1
\weberpage{1}

\begin{center}
\gesperrt{FIRST BOOK.}

\rule{2cm}{0.4pt}

\vspace{2em}
\gesperrt{GROUPS.}

\rule{4cm}{0.4pt}
\end{center}

\vspace{4em}
%% [running head, not body text -- commented out 2026-08-07] Page furniture (running
%% head, folio, signature) is excluded from this edition by standing policy.
%% \noindent{\small Weber, Algebra.\ II. \hfill 1}


% ===== scan leaf 23, printed 2 =====
\weberpage{2}% page number not printed on this page
% !PRINTED none
% Blank page (verso of the Book-I / "GRUPPEN." half-title, page 1). No printed
% text is visible in any of the three bands beyond scan noise/specks.


% ===== scan leaf 24, printed 3 =====
\weberpage{3}% page number not printed on this page
\pdfbookmark[0]{First Book. Groups.}{bm1}
\pdfbookmark[1]{First Section. General Group Theory.}{bm2}
\pdfbookmark[2]{§ 1. Definition of Groups}{bm3}
% !PRINTED none

\begin{center}
\gesperrt{First Chapter.}

\vspace{1em}
\gesperrt{General Group Theory.}

\rule{2cm}{0.4pt}

\vspace{1em}
\S.\,1.

\gesperrt{Definition of Groups.}
\end{center}

\noindent In the first volume, in connection with the permutations, we became acquainted with the concept of a group, and we made important algebraic applications of it. It must now be our next task to grasp this concept, so exceedingly important in the whole of recent mathematics, in a more general way, and to become acquainted with the laws that govern it. We place the following definition at the head:

A system $P$ of things (elements) of any kind whatever becomes a group when the following conditions are fulfilled:

\noindent 1.\ \gesperrt{A rule is given, by which from a first and a second element of the system a wholly determinate third element of the same system is derived.}

One writes symbolically, if $a$ is the first, $b$ the second, $c$ the third element:
\[
ab = c, \quad c = ab,
\]
and calls $c$ \gesperrt{compounded} from $a$ and $b$, and $a$ and $b$ the \gesperrt{components} of $c$.

In this composition the commutative law is not in general presupposed, that is, $ab$ may be different from $ba$; on the other hand,

\noindent 2.\ \gesperrt{the associative law is presupposed,} that is, if $a, b, c$ are any three elements of $P$, then
\[
(ab)c = a(bc),
\]
%CONT>>


% ===== scan leaf 25, printed 4 =====
% !PRINTED 4
\weberpage{4}

%% [running head, not body text -- commented out 2026-08-07] Page furniture (running
%% head, folio, signature) is excluded from this edition by standing policy.
%% \noindent{\small 4 \hfill Erster Abschnitt. \hfill \S.\,1.}

%<<CONT
and from this it follows, by the mode of inference of complete induction, that one always arrives at the same result when, in an arbitrary sequence of elements from $P$ in finite number, $a, b, c, d \ldots$, one first compounds two adjacent elements, then again two adjacent ones, and so on, until the whole sequence is reduced to a single element, which is denoted by $abcd \ldots$. Thus, for example:
\[
\begin{aligned}
abcd &= (ab)cd = [(ab)c]d = (ab)(cd)\\
&= a(bc)d = [a(bc)]d = a[(bc)d]\\
&= ab(cd) = (ab)(cd) = a[b(c.d)]^{1)}. % [sic: the source prints a dot between c and d]
\end{aligned}
\]

\noindent 3.\ \gesperrt{It is presupposed. that, if $ab = ab'$ or $ab = a'b$, then necessarily $b = b'$ in the first case, and $a = a'$ in the second.} % [sic: the source prints a full stop after vorausgesetzt]

If $P$ comprises a finite number of elements, the group is called a \gesperrt{finite one}, and the number of its elements its \gesperrt{degree}.

For finite groups there results from 1., 2., 3. the \gesperrt{consequence:}

\noindent 4.\ \gesperrt{If, of the three elements $a, b, c$ from $P$, any two are given, then the third can always be determined, and in only one way, so that}
\[
ab = c
\]
\gesperrt{holds.}

For if $a, b$ are the given elements, the assertion 4. coincides with 1. But if $a$ and $c$ are given, let the second component $b$ in the compositum $ab$ run through the whole system $P$, whose degree may be $= n$. Then by 1. and 3. one obtains in $ab$ nothing but distinct elements of $P$, and since their number is $= n$, all the elements of $P$, hence also $c$, must occur among them. One reasons in the same way when $b$ and $c$ are given, by letting $a$ run through the whole system $P$.

For infinite groups one can no longer reason in this way. \gesperrt{For infinite groups, therefore, the property 4. is further taken up into the definition as a requirement.}

\noindent\rule{3cm}{0.4pt}

\noindent{\footnotesize $^{1)}$ Cf.\ \gesperrt{Dirichlet-Dedekind}, Vorlesungen \"uber Zahlentheorie, \S.\,2.}


% ===== scan leaf 26, printed 5 =====
% !PRINTED 5
\weberpage{5}

In the permutation groups considered in the first volume one will easily recognise the marks of the general concept of a group.

We now draw from this definition first some quite general consequences.

By 4. there is for every given $b$ \gesperrt{an} element $e$ in $P$ that satisfies the condition
\weq{1}{eb = b}
and this $e$ is independent of $b$; for from (1) it follows, for every $c$,
\[
ebc = bc,
\]
and $bc$ can by 4. denote every element in $P$. Likewise there is an element $e'$ that for every $b$ satisfies the condition
\weq{2}{be' = b}
This element $e'$, however, is not different from $e$; for if we set $b = e'$ in (1) and $b = e$ in (2), it follows that
\[
e e' = e', \, . \, e e' = e, % [sic: the source prints a full stop between the two identities]
\]
hence
\[
e = e'.
\]

The element $e$ changes nothing when it is compounded with any elements of $P$, and is called the \gesperrt{unit of the group}.

In many cases it can, without misunderstanding, be denoted simply by ``1''.

To every element $a$ there is by 4. a definite element $a^{-1}$ that satisfies the condition
\weq{3}{a^{-1} a = e}
From (3), (1) and (2) it follows that
\[
a^{-1} a a^{-1} = e a^{-1} = a^{-1} = a^{-1} e,
\]
and consequently, by 3.,
\weq{4}{a a^{-1} = e.}

The two elements $a$, $a^{-1}$ are said to be \gesperrt{opposite} or \gesperrt{reciprocal} to one another.

In special cases the commutative law can also hold in the composition of the elements of a group $P$; that is, for any two elements $a$, $b$ of the group
\[
a b = b a
\]
can hold.


% ===== scan leaf 27, printed 6 =====
% !PRINTED 6
\weberpage{6}

\gesperrt{5. Groups which have this property are called commutative groups or also Abelian groups.}

If the elements of two groups
\[
a, b, c, d \ldots
\]
and
\[
a', b', c', d' \ldots
\]
correspond to one another in a mutually one-to-one manner, so that whenever $ab = c$, then also $a'b' = c'$, the groups are called \gesperrt{isomorphic}, and the evident theorem holds that two groups isomorphic with a third are isomorphic with one another. Accordingly one may gather all groups that are isomorphic with one another into a class of groups, which is itself again a group, whose elements are the generic concepts obtained when one gathers the corresponding elements of the individual isomorphic groups into one general concept. The individual mutually isomorphic groups are then to be regarded as different representatives of one generic concept.

The properties of groups that can come into consideration are of different kinds. They may, namely, either attach to the particular groups and be derived from the nature of the elements of which the group consists, or else from the nature of the law of composition. Or they may attach to the groups as such, and must then admit of being derived solely from the definition of the concept of a group. The latter properties are common to all isomorphic groups and may be designated as invariant properties of the group. If a comparison is permitted, one might recall the distinction between the metric and the projective properties in geometry.

To the first kind of properties, which are derived from the particular nature of the elements, belong, e.\,g., in the case of permutation groups, the properties of transitivity and intransitivity, of primitivity and imprimitivity; to the invariant properties belong permutability or non-permutability, the degree, the divisors and their index, the normal divisors. It is with these invariant properties, for which it is immaterial from which representative of a class of isomorphic groups they are derived, that we have first to occupy ourselves.


% ===== scan leaf 28, printed 7 =====
% !PRINTED 7
\weberpage{7}
\pdfbookmark[2]{§ 2. The Divisors of Finite Groups}{bm4}

\begin{center}
\S.\,2.

\gesperrt{The Divisors of Finite Groups.}
\end{center}

It is, as we have already seen in the first volume, a particularly important question whether a part of the elements of a group is itself again a group. For this it is necessary and sufficient that any two elements of this part yield under composition an element belonging again to the same part. This smaller group is then called a \gesperrt{divisor} or \gesperrt{factor}\footnote{Also called a subgroup.} of the whole group. The unit element ``1'' forms in every group a group by itself, and is therefore a divisor in every group.

We now restrict ourselves, as almost always in what follows, to \gesperrt{finite groups}, and shall assume that
\[
P = 1, a_1, a_2 \ldots a_{n-1}
\]
is a group of degree $n$ with the elements $a_0 = 1$, $a_1$, $a_2 \ldots a_{n-1}$, and that
\[
Q = 1, a_1, a_2 \ldots a_{\nu-1}
\]
is a divisor of $P$ of degree $\nu$. Every such divisor $Q$ must have the fundamental properties of a group, and therefore certainly contains the unit element ``1'', and with each of its elements $a_1$ at the same time the opposite $a_1^{-1}$. If $\nu < n$, we call $Q$ a \gesperrt{proper divisor} of $P$. If then $b$ is an element of $P$ not contained in $Q$, the elements
\[
Q_1 = b, a_1 b, a_2 b \ldots a_{\nu-1} b
\]
are all different from one another (\S.\,1, 3.) and all not contained in $Q$; for if, say, $a_1 b$ were contained in $Q$, then, since $Q$ is a group, $a_1^{-1} a_1 b = b$ would also have to be contained in $Q$, contrary to the hypothesis. If the whole group $P$ is not yet exhausted by $Q$ and $Q_1$, we take one of the remaining elements, which we may denote by $c$, and form
\[
Q_2 = c, a_1 c, a_2 c \ldots a_{\nu-1} c,
\]
and easily convince ourselves that the elements of $Q_2$ are all different not only from one another, but also from the elements of $Q$ and
%CONT>>


% ===== scan leaf 29, printed 8 =====
% !PRINTED 8
\weberpage{8}
%<<CONT
$Q_1$ are distinct. For if we had, say, $a_1 c = a_2 b$, it would follow that $c = a_1^{-1} a_2 b$; but $a_1^{-1} a_2$ is contained in $Q$, hence identical with one of the elements $1, a_1, a_2 \ldots a_{\nu-1}$, and $c$ would therefore, contrary to hypothesis, be contained in $Q_1$. We continue in this way, forming the systems $Q, Q_1, Q_2 \ldots Q_{\mu-1}$, of which the last shall be
\[
Q_{\mu-1} = g, a_1 g, a_2 g \ldots a_{\nu-1} g
\]
But then $P$ must be exhausted by the systems $Q, Q_1, Q_2 \ldots Q_{\mu-1}$, and since these systems contain nothing but distinct elements, it follows that
\[
n = \nu \mu,
\]
or the theorem:

\noindent 1.\ \gesperrt{The degree of a group is divisible by the degree of each of its divisors. The quotient $\mu = n : \nu$ shall be called the index of the divisor $Q$.}

The systems $Q_1, Q_2 \ldots Q_{\mu-1}$ we call, as in the special case of the permutation groups, the \gesperrt{subsidiary groups} belonging to $Q$, and denote them by
\[
Q_1 = Qb, \quad Q_2 = Qc \ldots Q_{\mu-1} = Qg,
\]
so that we can also write symbolically
\weq{1}{P = Q + Qb + Qc + \cdots + Qg}
At times we shall also designate the whole system $Q, Q_1 \ldots Q_{\mu-1}$ ($Q$ itself included) as a system of subsidiary groups, and accordingly call the representation (1) the \gesperrt{decomposition of $P$ into a system of subsidiary groups}.

From the mode of formation of the subsidiary groups it follows further that, if $b$, $c$ are any two elements from $P$, the two systems $Qb$ and $Qc$ are either entirely identical or contain no common element. For if $a_1 b = a_2 c$, where $a_1$, $a_2$ are two elements from $Q$, then for every other element $a$ from $Q$ we also have
\[
a a_1 b = a a_2 c,
\]
and when $a$ runs through the whole group $Q$, then, as already remarked in \S.\,1, each of the two $a a_1$ and $a a_2$ also runs through the whole group $Q$; consequently $Qb$ and $Qc$ are identical.

As in Bd.\,I, \S.\,154, 2. we can also decompose $P$ into subsidiary groups in the following way:
\[
P = Q + b^{-1} Q + c^{-1} Q + \cdots + g^{-1} Q.
\]


% ===== scan leaf 30, printed 9 =====
% !PRINTED 9
\weberpage{9}

The subsidiary groups do not have the characteristics of a group. For in
order that two elements from $Q_1$ should, on composition, again yield an
element of $Q_1$, we should have to have, say,
\[
a_1 b a_2 b = a_3 b,
\]
hence $a_1 b a_2 = a_3$, $b = a_1^{-1} a_3 a_2^{-1}$. Thus $b$ would,
contrary to the hypothesis, be contained in $Q$.

The name subsidiary group is therefore to be understood only in an
improper sense.

Two divisors $Q$, $Q'$ of $P$ always have the element 1 in common. But
they may also have still other elements in common, and these common
elements form a group, which we shall denote by $R$. For if the elements
$a$ and $b$ belong both to $Q$ and to $Q'$, then on account of the group
character of $Q$ and $Q'$ the same holds of the compositum $ab$; hence $R$
is a group. This common group we call the \gesperrt{greatest common divisor} of $Q$
and $Q'$, or also, with a geometrical allusion, the \gesperrt{intersection of} $Q$
and $Q'$. It follows likewise that the elements which are common to any
number of divisors of $P$ form a group, which we likewise designate as the
greatest common divisor or the intersection of all these groups.

In every group we can form divisors by the following procedure.

One divisor is always the unit element by itself.

If we denote the repeated composition of an element with itself by powers
(with $a^0$ the unit element), then the sequence of elements
\[
1,\ a,\ a^2,\ a^3 \ldots,
\]
which all belong to the group $P$, cannot contain nothing but distinct
elements. If, then, the first element that recurs for the second time is
\[
a^\mu = a^{\mu+\alpha},
\]
it follows that $a^\alpha = 1$ must hold, and that $\alpha$ is the least
positive number that satisfies this condition. Then, if $m = q\alpha$ is
any multiple of $\alpha$, $a^m = 1$, and conversely: whenever $a^m = 1$,
$m$ must be divisible by $\alpha$; for otherwise one could set
$m = q\alpha + \alpha'$, where $\alpha'$ is positive and less than
$\alpha$, and we should have $a^{\alpha'} = 1$, contrary to the
hypothesis,
%CONT>>


% ===== scan leaf 31, printed 10 =====
% !PRINTED 10
\weberpage{10}
\pdfbookmark[2]{§ 3. Normal Divisors of a Group}{bm5}

%<<CONT
that $\alpha$ be the smallest positive exponent for which $a^\alpha = 1$.
We have, always and only then,
\[
a^\mu = a^{\mu'},
\]
when $\mu \equiv \mu'\ (\mathrm{mod}\ \alpha)$, and here the exponents may
also be negative, if $a^{-r}$ denotes the $r$th power of the element
$a^{-1}$, or the element opposite to $a^r$.

The sequence
\[
A = 1,\ a,\ a^2 \ldots a^{\alpha-1}
\]
then contains $\alpha$ elements of $P$ different from one another, among
which all powers of $a$ are represented. But the elements $A$ evidently
form a group of degree $\alpha$, because the exponents simply add under
composition. This group is a divisor of $P$, and hence $\alpha$ is
\gesperrt{a divisor of $n$}. Hence for every element $a^n = 1$.

\gesperrt{The group $A$ is called the period of the element $a$, and $\alpha$ is
also called the degree of the element $a$.}

\begin{center}
\S.\,3.

\gesperrt{Normal divisors of a group.}
\end{center}

If $P$ is, as above, a group and $Q$ a divisor of $P$, and further $b$ an
element of $P$ not contained in $Q$, then the subsidiary group $Qb$ is not
a group. On the other hand the system $b^{-1}Qb$, or written out at
length
\[
1,\ b^{-1}a_1 b,\ b^{-1}a_2 b \ldots b^{-1}a_{\nu-1} b
\]
certainly forms a group, because
\[
b^{-1}a_1 b\ .\ b^{-1}a_2 b = b^{-1}a_1 a_2 b
\]
and this group is isomorphic with $Q$. The group $b^{-1}Qb$ is called
\gesperrt{the group transformed by $b$}. If $b$ itself belongs to $Q$, then
$b^{-1}Qb$ is identical with $Q$. If one takes for $b$ the various
elements of $P$, one obtains a whole family of such groups, which we call
the \gesperrt{divisors of $P$ conjugate to $Q$}, or also briefly \gesperrt{conjugate groups}.

It may happen that all the conjugate divisors are identical with one
another. We have already seen in the case of the permutation groups that
this case is of particular importance, and we therefore now introduce in
general the following definition:
%CONT>>


% ===== scan leaf 32, printed 11 =====
% !PRINTED 11
\weberpage{11}

%<<CONT
\gesperrt{If $Q$ is a divisor of $P$ which is identical with all of its conjugate
divisors, then $Q$ is called a normal divisor of
$P$}\footnote{Also called a distinguished or invariant subgroup.}.

The group formed from the single element 1 is a normal divisor of
every group. We obtain, further, a normal divisor in the greatest
common divisor $R$ of all the divisors conjugate to any divisor
$Q$ of $P$.

For it has already been proved above that $R$, as the intersection of
several divisors of $P$, is a group.

If now
\weq{1}{Q,\ Q',\ Q'' \ldots}
are the divisors of $P$ conjugate to $Q$, and $b$ is any element of $P$,
then the system of groups
\weq{2}{b^{-1}Qb,\ b^{-1}Q'b,\ b^{-1}Q''b \ldots}
is not different from the system (1). But if $R$ is the intersection of the
groups (1), then $b^{-1}Rb$ is the intersection of (2), and consequently
$R$ is identical with $b^{-1}Rb$, i.e. $R$ is a normal divisor of $P$.

If $N$ is a normal divisor of any group $P$, and $b$ an arbitrary
element in $P$, then
\weq{3}{b^{-1}Nb = N \quad \text{or} \quad Nb = bN.}

If $P$ is of degree $n$, $N$ of degree $\nu$ and of index $\mu$, so that
$n = \mu\nu$, then we can choose the $\mu$ elements $1, b_1 \ldots b_{\mu-1}$ in
such a way that the decomposition of $P$ into the subsidiary groups
\[
\begin{aligned}
P &= N + Nb_1 + Nb_2 + \cdots + Nb_{\mu-1} \\
  &= N + b_1 N + b_2 N + \cdots + b_{\mu-1} N
\end{aligned}
\]
results. Thus
\weq{4}{N,\ N_1 = Nb_1 = b_1 N,\ N_2 = Nb_2 = b_2 N \ldots,\quad N_{\mu-1} = Nb_{\mu-1} = b_{\mu-1} N}
is the system of the subsidiary groups.

If $a$ is any element in $N$, then $Na$ is identical with $N$; hence
$Nab$ is also identical with $Nb$, and since every element $c$ of $P$
occurs in $N$ or in one of its subsidiary groups, and is therefore
contained in the form $ab_k$, $Nc = cN$ must agree with one of the
systems (4).


% ===== scan leaf 33, printed 12 =====
% !PRINTED 12
\weberpage{12}
\pdfbookmark[2]{§ 4. Composition of the Parts}{bm6}

We mention further the following theorem, whose correctness follows immediately from
the isomorphism of transformed groups:

\gesperrt{If $Q$ is a divisor of $P$, and $R$ a divisor of $Q$, then, if $Q'$
and $R'$ arise from $Q$ and $R$ by transformation by the same element of $P$,
$R'$ is also a divisor of $Q'$; and if $R$ is a normal divisor of $Q$,
then $R'$ is also a normal divisor of $Q'$.}

\begin{center}
\S.\,4.

\gesperrt{Composition of the Parts.}
\end{center}

If $A$ and $B$ are any two sequences of elements from a group $P$
(groups or not), then by the symbolic product $AB$ we shall understand
the totality of all elements that one obtains when each element $a$ of $A$
is compounded with each element $b$ of $B$, according to the rule holding
in $P$, into $ab$. This manner of composition of
$A$ and $B$ into $AB$ we shall call the \gesperrt{composition of the parts} (of $P$).
We distinguish here between a \gesperrt{part} and a \gesperrt{divisor} of
$P$, so that a divisor is always to be a group, which is not necessary
for a part.

One can likewise compound three and more parts $A, B, C \ldots$ of $P$,
and the associative law holds for the composition of the parts, as
follows immediately from the fact that this law holds in $P$. Accordingly the
meaning of a compositum $ABC \ldots$ is uniquely determined.

In the composition $AB$ one of the parts $A$, $B$ may also consist of a
single element, and then the notation $AB$ agrees with the notation
used above for the subsidiary groups.

For the composition of the parts the following theorems hold:

\noindent 1.\ \gesperrt{The necessary and sufficient condition
for a part $A$ of $P$ to be a group consists in the
equation}
\weq{1}{AA = A.}
For this equation says nothing else than that, if $a$ and $a'$ are contained in
$A$, then $aa'$ also belongs to $A$, and hence that $A$ is a group.


% ===== scan leaf 34, printed 13 =====
% !PRINTED 13
\weberpage{13}

%% [running head, not body text -- commented out 2026-08-07] Page furniture (running
%% head, folio, signature) is excluded from this edition by standing policy.
%% \gesperrt{\S.\,4.} \hfill \gesperrt{Group of the subsidiary groups.} \hfill 13

\gesperrt{2. By the two equations}

\weq{2}{AA = A,}
\weq{3}{AB = BA,}

\gesperrt{in which $A$ may be a definite part, $B$ any arbitrary part of $P$, it is asserted that $A$ is a normal divisor of $P$.}

For by (2), $A$ is a divisor of $P$, and by (3), for every element $b$ of $P$
\[
b^{-1}Ab = A,
\]
i.\,e. $A$ is a normal divisor.

\gesperrt{3. If $A$ and $B$ are two divisors of $P$ (groups), then either of the two equations}

\weq{4}{AB = A, \quad BA = A,}

\gesperrt{is the necessary and sufficient condition for $B$ to be a divisor of $A$.}

For if $B$ is a group and a divisor of the group $A$, and if $a$ is contained in $A$, and $b$ in $B$ and hence also in $A$, then $ab$ is also contained in $A$. Thus $A$ contains every element of $AB$.

But since, secondly, $B$ as a group also contains the element 1, $AB$ also contains every element of $A$, and hence $AB$ is identical with $A$.

In the same way one sees that $BA$ is identical with $A$. Conversely, it follows from either of the equations (4), since $A$ as a group contains the unit element, that every element of $B$ is contained in $A$, and hence that $B$ is a divisor of $A$.

\gesperrt{4. Under the composition of the parts, the system of the subsidiary groups belonging to a normal divisor of $P$ itself forms a group, in which the normal divisor $N$ forms the unit, and}
\[
Nb, \quad Nb^{-1}
\]
\gesperrt{are opposite elements.}

For if $N$ is a normal divisor of $P$ and
\weq{5}{N, \; N_1 = Nb_1, \; N_2 = Nb_2, \dots}
is the system of the subsidiary groups, then
\[
N_h N_k = N b_h N b_k.
\]


% ===== scan leaf 35, printed 14 =====
% !PRINTED 14
\weberpage{14}

%% [running head, not body text -- commented out 2026-08-07] Page furniture (running
%% head, folio, signature) is excluded from this edition by standing policy.
%% \gesperrt{Erster Abschnitt.} \hfill 14 \hfill \gesperrt{\S.\,4.}

But on account of the property of normal divisors, $N$ is permutable with every $b$, hence $b_h N = N b_h$, and consequently
\[
N_h N_k = N N b_h b_k.
\]

Because $b_h b_k$ is contained in $P$, $N b_h b_k$ is identical with one of the subsidiary groups (5), and since $NN = N$ may be set (by 1.), it follows, if we set $N b_h b_k = N_l$,
\weq{6}{N_h N_k = N_l.}

With this, the property \S.\,1, 1. of groups is established for the composition of the subsidiary groups. That \S.\,1, 2. also holds, namely the associative law, we have already pointed out. It therefore remains to prove \S.\,1, 3., namely that, if
\weq{7}{N_h N_i = N_k N_l}
is set, and $N_h = N_k$, then also $N_i = N_l$, and if $N_i = N_l$, then also $N_h = N_k$ must hold.

By the representation (5), however, we can represent (7) in the twofold manner:
\weq{8}{\begin{aligned}
N b_h b_i &= N b_k b_l \\
b_h b_i N &= b_k b_l N.
\end{aligned}}

If now $N_i = N_l$, we may also assume $b_i = b_l$, and obtain from (8), by composition with $b_i^{-1}$,
\[
N b_h = N b_k;
\]
and if $N b_h = N b_k$, we assume $b_h = b_k$ and obtain from (8), by composition with $b_h^{-1}$,
\[
b_i N = b_l N,
\]
hence also $N_i = N_l$.

It is thus proved that the system of subsidiary groups becomes a group under the composition of the parts; that in this group the normal divisor $N$ is the unit element follows immediately from 1., because by it
\[
N N_k = N N b_k = N_k
\]
holds. And because $N b_k N b_k^{-1} = N b_k b_k^{-1} = N$, $N b_k$ and $N b_k^{-1}$ are opposite elements.

The group of the quantities $N_k$, whose existence is hereby proved, we call the \gesperrt{group of the subsidiary groups}, or the \gesperrt{group complementary to $N$} with respect to $P$. We
%CONT>>


% ===== scan leaf 36, printed 15 =====
% !PRINTED 15
\weberpage{15}
\pdfbookmark[2]{§ 5. Multiple Isomorphism}{bm7}

%% [running head, not body text -- commented out 2026-08-07] Page furniture (running
%% head, folio, signature) is excluded from this edition by standing policy.
%% \gesperrt{\S.\,5.} \hfill \gesperrt{Mehrstufiger Isomorphismus.} \hfill 15

%<<CONT
denote it, following the example of Hölder\footnote{Mathematische Annalen, Bd.\,34. The symbol recalls a quotient, with which the complementary group does indeed have a certain analogy.}, by $P/N$. The degree of the complementary group is equal to the index of the group $N$.

At the close of these considerations we mention one further theorem on normal divisors.

\noindent 5.\ \gesperrt{If $N$ is a divisor of $P$, and $N'$ a divisor of $N$ and at the same time a normal divisor of $P$, then $N'$ is also a normal divisor of $N$.}

This is self-evident, because for every element $b$ of $P$ we have $b^{-1}N'b = N'$, and because every element of $N$ is also an element of $P$.

\gesperrt{But one may not conclude conversely that, if $N$ is a normal divisor of $P$ and $N'$ a normal divisor of $N$, then $N'$ must also be a normal divisor of $P$.}

\begin{center}
\S.\,5.

\gesperrt{Multiple Isomorphism.}
\end{center}

If, as above, $N$ is a normal divisor of $P$ of degree $\nu$ and index $\mu$, so that $n = \mu \nu$, then the group $P/N$ complementary to $P$ is of degree $\mu$. We will denote this group, or one isomorphic with it, by $Q$, and its elements, which correspond to the subsidiary groups $N, Nb_1, Nb_2 \dots$, by $A, B, C \dots$ To each of these elements there correspond $\nu$ elements of the group $P$, say

\noindent to the element $A$ the elements $a, a_1 \dots a_{\nu-1}$

\noindent to the element $B$ the elements $b, b_1 \dots b_{\nu-1}$

\noindent $. \; . \; . \; . \; . \; . \; . \; . \; . \; . \; . \; . \; . \; . \; . \; . \; . \; . \; . \; . \; .$

This correspondence is then of such a kind that every composed element $ab$ corresponds to the composed element $AB$. For the elements $A$ and $B$ correspond to the subsidiary groups $Na$ and $Nb$, and, because $N$ is a normal divisor, we have
\[
Na \, Nb = Nab.
\]

Thus in the composition of the $A, B \dots$ on the one hand and of the $a, b \dots$ on the other, the same law holds here as in
%CONT>>


% ===== scan leaf 37, printed 16 =====
% !PRINTED 16
\weberpage{16}

%% [running head, not body text -- commented out 2026-08-07] Page furniture (running
%% head, folio, signature) is excluded from this edition by standing policy.
%% 16 \hfill \gesperrt{Erster Abschnitt.} \hfill \gesperrt{\S.\,5.}

%<<CONT
the isomorphic groups (\S.\,1), only with the difference that to each element $A$ there corresponds not merely one element $a$, but several. This fact gives occasion to extend the concept of isomorphism, as is done by the following definition:

\gesperrt{A group $P$ with the elements $a, b, c \dots$ is called (multiply) isomorphic with a group $Q$ with the elements $A, B, C \dots$ if the two groups can be so related to one another that to each of the elements $A, B, C \dots$ there correspond one or more of the elements $a, b, c \dots$, and indeed in such a way that each of the elements of $Q$ corresponds to one and only one of the elements of $P$, and that, when $a$ and $A$, $b$ and $B$ correspond to one another, $ab$ corresponds to the element $AB$.}
% [sic: Weber prints ``elements of Q''; the uniqueness condition requires P.]

Thus it can first be proved that to each of the elements $A, B, C \dots$ there corresponds an equal number of elements $a, b, c \dots$, and that therefore \gesperrt{the degree of $Q$ is a divisor of the degree of $P$}. For let $A$ be the unit element in $Q$, to which the elements $a, a_1 \dots a_{\nu-1}$ in $P$ may correspond. These latter elements must then form a group of degree $\nu$, which we shall denote by $N$, because by the definition of isomorphism the element $a a_1$ must correspond to the element $AA = A$. If then $b$ is an element of $P$ corresponding to the element $B$, then again all the elements $Nb$ must correspond to the same element $B$ of $Q$. For every $ab$ must correspond to the element $AB = B$. If $b_1, b$ are two elements corresponding to the element $B$, then $b_1 b^{-1} = a$ corresponds to the element $A$, and is therefore contained in $N$; therefore $b_1 = ab$ is contained in $Nb$. By $Nb$, therefore, all the elements that correspond to the element $B$ are exhausted, and to each element of $Q$ there correspond $\nu$ elements of $P$. The isomorphism is called \gesperrt{$\nu$-fold}. But every element $b^{-1}ab$ likewise corresponds to the unit element $A$, and therefore $b^{-1}Nb = N$, i.\,e. $N$ is a normal divisor of $P$, and $Q$ is simply isomorphic with $P/N$.

We see, then, that there is no other multiple isomorphism than that between the complementary group of a normal divisor and the whole group. But one could extend the concept of isomorphism still further, by calling two groups which are $\mu$- and $\nu$-fold
%CONT>>


% ===== scan leaf 38, printed 17 =====
% !PRINTED 17
\weberpage{17}
\pdfbookmark[2]{§ 6. The Composition Series and the Theorem of C. Jordan}{bm8}
%<<CONT
isomorphic, $\mu$-$\nu$-fold isomorphic to one another. By far
the most important is the simple isomorphism, which we shall
therefore designate as isomorphism pure and simple, while a
multiple isomorphism shall always be made recognisable by an
added word $^{1)}$.

\begin{center}
\S.\,6.

\gesperrt{The composition series and the theorem of C. Jordan.}
\end{center}

A group $P$ is called \gesperrt{simple} if it has no normal
divisor besides itself and unity, \gesperrt{composite} if further
normal divisors are present.

A normal divisor which has no other normal divisor \gesperrt{above it},
which therefore is not part of another proper normal divisor
of $P$, is called \gesperrt{a greatest normal divisor}
of $P$ $^{2)}$. If $P_1$ is a greatest normal divisor of $P$, $P_2$ a
greatest normal divisor of $P_1$, $P_3$ a greatest normal divisor
of $P_2$, and so on, then the series of groups

\weq{1}{P,\, P_1,\, P_2,\, P_3 \ldots 1,}

which necessarily ends with the group 1 formed from the unit
element alone, is called a \gesperrt{composition series} of the
group $P$.

We shall denote the degrees of the groups (1) by $n, n_1, n_2, n_3 \ldots 1$,
and the quotients $n : n_1$, $n_1 : n_2$, $n_2 : n_3 \ldots$, that is,
the indices of $P_1$ with respect to $P$, of $P_2$ with respect to $P_1$, and so on,
by $\nu_1, \nu_2, \nu_3 \ldots$ The series of integers

\weq{2}{\nu_1,\, \nu_2,\, \nu_3 \ldots}

we call the \gesperrt{index series} of the group $P$.

For a given group $P$ the greatest normal divisors $P_1, P_2, P_3 \ldots$
can in general be selected in several different
ways, and accordingly the degree numbers
$n_1, n_2, n_3 \ldots$ and the indices $\nu_1, \nu_2, \nu_3 \ldots$ may also turn out
%CONT>>

$^{1)}$ Cf. C.\,\gesperrt{Jordan}, Traité des substitutions. \gesperrt{Netto}, Substitutionentheorie
. The designation ``$\mu$-fold'' is due to \gesperrt{Netto}. According to \gesperrt{Jordan}
the simple isomorphism is called ``holohedral'', the multiply
``merohedral'' isomorphism. The simple isomorphism is sometimes
 also designated as equivalence.

$^{2)}$ Distinguished maximal subgroup, according to another designation.

%% [page furniture, not body text -- commented out 2026-08-10] Page furniture (running
%% head, folio, signature) is excluded from this edition by standing policy.
%% Weber, Algebra. II. \hfill 2


% ===== scan leaf 39, printed 18 =====
% !PRINTED 18
\weberpage{18}
%<<CONT
turn out. But the following beautiful and important theorem
of C.\,\gesperrt{Jordan} $^{1)}$ holds:

\gesperrt{I. However the composition series of a group
$P$ may be chosen, the index series is, apart
from the arrangement, always the same.}

The proof rests on the consideration of the complementary
groups introduced in the preceding section.

Let $Q$ and $Q_1$ be normal divisors of $P$, and at the same time $Q_1$
a divisor (hence, by \S.\,4, 5. a normal divisor) of $Q$.

Then the following theorem holds:

\gesperrt{1. The group $Q/Q_1$ is a normal divisor of the
group $P/Q_1$, and the index of $Q/Q_1$ with respect
to $P/Q_1$ is equal to the index of $Q$ with respect
to $P$.}

To see its correctness, one need only consider more
closely the mode of formation of the individual groups.

Let $n, q, q_1$ be the degrees of $P, Q, Q_1$ and

\weq{3}{n = \nu q,\quad q = \nu_1 q_1,\quad n = \mu q_1,\quad \mu = \nu \nu_1.}

One now decomposes $Q$ into the subsidiary groups with respect to $Q_1$, i.e. one
sets, if $b_1, b_2 \ldots b_{\nu_1 - 1}$ are suitably determined elements in $Q$,

\weq{4}{Q = Q_1 + Q_1 b_1 + \cdots Q_1 b_{\nu_1 - 1}.}

The elements of the group $Q/Q_1$ are

\weq{5}{Q_1,\, Q_1 b_1,\, \ldots\, Q_1 b_{\nu_1 - 1}.}

If we decompose $P$ into the subsidiary groups with respect to $Q_1$, then
the elements (5) certainly occur among them, and consequently $Q/Q_1$
is a divisor of $P/Q_1$.

It therefore remains to show that this divisor is a
normal divisor.

But if we denote by $a$ any element of $P$, then
all elements of $P/Q_1$ are contained in the form $Q_1 a$, and
we therefore have only to show that, if $b$ is any element
in $Q$, every element

\weq{6}{Q_1 a^{-1} Q_1 b Q_1 a}

is contained in $Q/Q_1$, i.e. in the form $Q_1 b$.

$^{1)}$ C.\,\gesperrt{Jordan}, Traité des substitutions, Paris 1870. \gesperrt{Netto}, Substitutionentheorie,
Leipzig 1882. We follow here the proof which
\gesperrt{Hölder} has given following \gesperrt{Netto} (Math. Annalen, Bd. 34).


% ===== scan leaf 40, printed 19 =====
% !PRINTED 19
\weberpage{19}
But this follows immediately from

\weq{7}{Q_1 a^{-1} Q_1 b Q_1 a = Q_1 a^{-1} b a,}

because together with $b$, $a^{-1} b a$ also occurs in $Q$ (because $Q$ is a
normal divisor of $P$). And that the determination of the
indices is also correct is shown by the enumeration. For the degrees of
$P/Q_1$ and $Q/Q_1$ are $\nu \nu_1$ and $\nu_1$, hence the index is $\nu$.

Alongside this there is the following theorem:

\gesperrt{2. If $Q_1$ is a normal divisor of $P$ of degree $q_1$,
and the group $P/Q_1$ itself has a divisor $M$
of degree $m$, then $P$ has a divisor $Q$ of
degree $q = m q_1$. At the same time $Q_1$ is a normal divisor
 of $Q$, and $M = Q/Q_1$. If $M$ is a normal divisor
 of $P/Q_1$, then $Q$ is also a normal divisor
of $P$.}

Let $\mu$ be the index of $Q_1$ with respect to $P$, and let $P$ be
decomposed into the subsidiary groups:

\weq{8}{P = Q_1 + Q_1 e_1 + Q_1 e_2 \cdots + Q_1 e_{\mu - 1},}

so that $e_1, e_2 \ldots e_{\mu - 1}$ are suitably chosen elements in $P$ and

\weq{9}{Q_1,\, Q_1 e_1,\, Q_1 e_2 \ldots Q_1 e_{\mu - 1}}

are the elements of $P/Q_1$. Now if a divisor $M$ is contained in
the group $P/Q_1$, it must consist of a part
of the elements (9), say of

\weq{10}{Q_1,\, Q_1 b_1,\, Q_1 b_2 \ldots Q_1 b_{\nu_1 - 1},}

and if this system is to form a group, then for
any two elements $b_i, b_k$ the product

\weq{11}{Q_1 b_i Q_1 b_k = Q_1 b_i b_k}

must again be contained in (10), hence say $= Q_1 b_h$. If, then,
the system $1, b_1, b_2 \ldots b_{\nu_1 - 1}$ is denoted by $B$, the relation
 (11) can be written, by the method of the composition of the parts,
thus:

\weq{12}{Q_1 B Q_1 B = Q_1 B,}

and so, by \S.\,4, 1., it says that the elements of $P$ contained in

\weq{13}{Q = Q_1 B}

form a group which contains the group
$Q_1$ as a divisor and is itself a divisor of $P$. That
$Q_1$ is a normal divisor of $Q$ follows from \S.\,4, 5., because $Q_1$, as a normal
%CONT>>


% ===== scan leaf 41, printed 20 =====
% !PRINTED 20
\weberpage{20}
%<<CONT
divisor of $P$ is presupposed, and from the representation (10)
of the group $M$ it follows that $M = Q/Q_1$.

It remains to show that if $M$ is a normal divisor of
$P/Q_1$, then $Q$ is also a normal divisor of $P$. This follows thus:

If $e$ is any element in $P$ and $Q_1 b_k$ an element in $M$,
then from the assumption that $M$ is a normal divisor of $P/Q_1$ it follows that:

\[ Q_1 e^{-1} Q_1 b_k Q_1 e = Q_1 b_h, \]

which we can also express, according to the composition of the parts, by the formula

\weq{14}{Q_1 e^{-1} Q_1 B Q_1 e = Q_1 B}

Now since $Q_1$ is a normal divisor of $P$,
$Q_1$ can be interchanged with each of the other factors, so
that from (14) follows:

\[ e^{-1} Q_1 B e = Q_1 B, \]

or

\weq{15}{e^{-1} Q e = Q,}

whereby it is expressed that $Q$ is a normal divisor of $P$.

From Theorems 1. and 2. results the corollary:

\gesperrt{3. A normal divisor $Q$ of $P$ is a greatest
normal divisor if and only if the complementary
 group $P/Q$ is simple.}

If $Q$ and $Q'$ are two normal divisors of $P$, then the
symbolic product $Q Q'$ is also a normal divisor.

For, since $Q$ is a normal divisor, we have

\[ Q Q' = Q' Q, \]

hence

\[ Q Q' Q Q' = Q Q Q' Q' = Q Q', \]

and this, by \S.\,4, 1., is the criterion that $Q Q'$ is a group.
That it is a normal divisor of $P$ then follows
likewise, by \S.\,4, 2., from

\[ Q Q' B = Q B Q' = B Q Q', \]

where $B$ is any part of $P$.

Now if $Q$ is a \gesperrt{greatest normal divisor} of $P$, and $Q'$ is
different from $Q$, then $Q Q'$ is a normal divisor of $P$ which
contains both $Q$ and $Q'$ within itself and hence is different from $Q$.
Consequently, since above $Q$ there is no normal divisor of $P$
except $P$ itself,

\weq{16}{Q Q' = P,}

and likewise we must also have

\weq{17}{Q' Q = P}
%CONT>>


% ===== scan leaf 42, printed 21 =====
% !PRINTED 21
\weberpage{21}
%<<CONT
be. These theorems hold if only one of the two groups $Q$, $Q'$
is a greatest normal divisor of $P$. We now assume that they
both have this property, and understand by $R$ the greatest
common divisor of $Q$ and $Q'$. This group $R$ is then a normal
divisor of $P$ as well as of $Q$ and $Q'$. For if $a$ is any
element in $P$, and $c$ is contained in $Q$ as well as in $Q'$,
then $a^{-1}ca$ is also in $Q$ and in $Q'$, and hence also
contained in $R$. Therefore $a^{-1}Ra = R$, and $R$ is a normal
divisor of $P$ and consequently a normal divisor of $Q$ and of $Q'$
(\S.\,4, 5.). In order to decompose $Q$ into the subsidiary groups
of $R$, we choose a suitable system of elements $1, b_1, b_2, b_3 \ldots$
in $Q$, so that $R, Rb_1, Rb_2, Rb_3 \ldots$ all become different
from one another, and set
\[
Q = R + Rb_1 + Rb_2 + Rb_3 + \ldots,
\]
which we can also, if we set
\weq{18}{B = 1, b_1, b_2, b_3 \ldots}
denote briefly by
\weq{19}{Q = RB}
Now by (13) we have $Q'Q = P$, hence also
\weq{20}{P = Q'RB = Q'B}
or
\[
P = Q' + Q'b_1 + Q'b_2 + Q'b_3 \ldots
\]

The subsidiary groups
\weq{21}{Q',\ Q'b_1,\ Q'b_2,\ Q'b_3 \ldots}
are, however, all different from one another. For if we had, say,
\[
Q'b_2 = Q'b_1,
\]
it would follow, since $Q'$ contains the unit, that there is an element
$e$ in $Q'$ such that
\[
b_2 = eb_1
\]
would hold. This element $e$, however, is equal to $b_2b_1^{-1}$, and
hence, since the $b$ are contained in $Q$, is likewise contained in $Q$
and therefore also in $R$. But then we should also have $Rb_2 = Rb_1$,
which is contrary to the assumption.

Accordingly we can form the elements of the group complementary to $R$
with respect to $Q$ and of the group complementary to $Q'$ with respect
to $P$. We obtain these two groups
\weq{22}{\begin{aligned}
Q/R &= R,\ Rb_1,\ Rb_2,\ Rb_3 \ldots \\
P/Q' &= Q',\ Q'b_1,\ Q'b_2,\ Q'b_3 \ldots
\end{aligned}}


% ===== scan leaf 43, printed 22 =====
% !PRINTED 22
\weberpage{22}
From this, however, one easily recognises that these groups are not
only of the same degree, but that they are \gesperrt{isomorphic}. For we have
simultaneously
\[
Rb_1\,Rb_2 = Rb_1b_2
\]
and
\[
Q'b_1\,Q'b_2 = Q'b_1b_2.
\]
In the same way one can also conclude that the two groups
\[
Q'/R,\quad P/Q
\]
are isomorphic.

But the groups $P/Q$ and $P/Q'$ are, since $Q$, $Q'$ are greatest
normal divisors, simple by 3), and consequently the groups $Q'/R$
and $Q/R$ isomorphic with them are also simple, and consequently we
have the theorem:

\gesperrt{4. If $Q$ and $Q'$ are two different greatest normal
divisors of $P$, and if $R$ is the intersection of $Q$ and $Q'$,
then $R$ is a greatest normal divisor both of $Q$ and of $Q'$, and
the index of $R$ with respect to $Q$ is equal to the index of $Q'$
with respect to $P$.}

With this we have gained the foundation for demonstrating very
simply the theorem of the constancy of the index series.

For a better survey we write the various composition series to be
drawn into comparison in such a way that we place the index of each
term with respect to the preceding one under the term.

Accordingly, let any two given composition series of $P$, with the
associated indices, be the following:
\weq{23}{\begin{aligned}
&P,\ Q,\ Q_1,\ Q_2.\ Q_3 \ldots \\ % [sic: the source prints a full stop after Q_2.]
&\hphantom{P,\ }\nu,\ \nu_1,\ \nu_2,\ \nu_3 \ldots
\end{aligned}}
\weq{24}{\begin{aligned}
&P,\ Q',\ Q_1',\ Q_2',\ Q_3' \ldots \\
&\hphantom{P,\ }\nu',\ \nu_1',\ \nu_2',\ \nu_3' \ldots
\end{aligned}}
Now let $R$ be the intersection of $Q$ and $Q'$, and let $\mu$ and
$\mu'$ be the indices of $R$ with respect to $Q$ and $Q'$. On
account of the isomorphism of the groups $P/Q$ and $Q'/R$ we then
have $\mu' = \nu$, and for the corresponding reason $\mu = \nu'$. We
form a composition series of $R$ with the associated index series
\weq{25}{\begin{aligned}
&R,\ R_1,\ R_2,\ R_3 \ldots \\
&\hphantom{R,\ }\mu_1,\ \mu_2,\ \mu_3 \ldots,
\end{aligned}}
%CONT>>


% ===== scan leaf 44, printed 23 =====
% !PRINTED 23
\weberpage{23}
%<<CONT
and since $R$ is now a greatest normal divisor of $Q$ and of $Q'$,
we can form from this two new composition series of $P$,
namely
\weq{26}{\begin{aligned}
&P,\ Q,\ R,\ R_1,\ R_2,\ R_3 \ldots \\
&\hphantom{P,\ }\nu,\ \nu',\ \mu_1,\ \mu_2,\ \mu_3 \ldots
\end{aligned}}
\weq{27}{\begin{aligned}
&P,\ Q',\ R,\ R_1,\ R_2,\ R_3 \ldots \\
&\hphantom{P,\ }\nu',\ \nu,\ \mu_1,\ \mu_2,\ \mu_3 \ldots
\end{aligned}}
The two composition series (26) and (27) of $P$ therefore have
the same index series.

If we now assume that the theorem to be proved has already been
established as correct for groups whose degree is $\leqq \frac{1}{2}n$
(where $n$ denotes the degree of $P$), it follows that all index
series of $Q$, whose degree is after all a proper divisor of $n$
and hence at most $= \frac{1}{2}n$, agree with one another, and hence
that the series
\[
\nu',\ \mu_1,\ \mu_2,\ \mu_3 \ldots
\]
\[
\nu_1,\ \nu_2,\ \nu_3,\ \nu_4 \ldots,
\]
agree, apart from the arrangement. In the same way the index
series of $Q'$
\[
\nu,\ \mu_1,\ \mu_2,\ \mu_3 \ldots
\]
\[
\nu_1',\ \nu_2',\ \nu_3',\ \nu_4' \ldots
\]
agree. Hence the two index series of $P$
\[
\nu,\ \nu_1,\ \nu_2,\ \nu_3,\ \nu_4 \ldots
\]
\[
\nu',\ \nu_1',\ \nu_2',\ \nu_3',\ \nu_4' \ldots
\]
also agree with one another, since the first agrees with $\nu, \nu', \mu_1, \mu_2, \mu_3
\ldots$, and the second with $\nu', \nu, \mu_1, \mu_2, \mu_3 \ldots$

For groups of degree 2, however, which have only the single normal
divisor 1, the theorem is evident, and it is therefore proved in
general by complete induction.

If in two composition series of $P$, which together with their
index series we shall now represent thus
\weq{28}{\begin{aligned}
&P,\ P_1,\ P_2,\ P_3 \ldots P_{\mu-1},\ 1 \\
&\hphantom{P,\ }j_1,\ j_2,\ j_3 \ldots j_{\mu-1},\ j_\mu
\end{aligned}}
\weq{29}{\begin{aligned}
&P,\ P_1',\ P_2',\ P_3' \ldots P_{\mu-1}',\ 1 \\
&\hphantom{P,\ }j_1',\ j_2',\ j_3' \ldots j_{\mu-1}',\ j_\mu'
\end{aligned}}
a common term $P_r = P_s'$ occurs, then the theorem on the
constancy of the index series holds also for the group $P_r =
P_s'$, and from this it follows first that we must have $s = r$, and
%CONT>>


% ===== scan leaf 45, printed 24 =====
% !PRINTED 24
\weberpage{24}
\pdfbookmark[2]{§ 7. Further Theorems on the Composition Series}{bm9}
%<<CONT
that the index series $j_{r+1}, j_{r+2} \ldots j_\mu$ and $j_{r+1}',
j_{r+2}' \ldots j_\mu'$ agree, apart from the order.
Consequently the preceding parts of the index series
\[
j_1, j_2 \ldots j_r \quad \text{and} \quad j_1', j_2' \ldots j_r'
\]
must also agree.

\begin{center}
\S.\,7.

\gesperrt{Further Theorems on the Composition Series.}
\end{center}

Of the composition series many further theorems hold, of
which we here derive the most important.

\gesperrt{II. If $P_\nu$ is any normal divisor of $P$, then there
is a composition series of $P$ in which $P_\nu$ occurs.}

If $P_\nu$ is a greatest normal divisor of $P$, we can set $P_\nu =
P_1$ and continue the composition series from there on arbitrarily.
But if $P_\nu$ is not a greatest normal divisor of $P$, then we
seek a greatest normal divisor above $P_\nu$, which we denote by
$P_1$. Then $P_\nu$ is a normal divisor of $P_1$. If it is a
greatest one, we set $P_\nu = P_2$ and continue the composition
series from there arbitrarily. But if $P_\nu$ is still not a greatest
normal divisor of $P_1$, we seek a greatest normal divisor of
$P_1$ above $P_\nu$ and proceed in this way until we finally
arrive at $P_\nu$, which must necessarily occur after a finite number
of steps. If we then append a composition series of $P_\nu$, we
have a composition series of $P$ in which $P_\nu$ occurs, as
theorem II. requires.

In a composition series of a group $P$, by virtue of the
definition, every term is a normal divisor of every preceding one.
The first term $P_1$ and the last term 1 are at the same time
normal divisors of $P$ itself. For the intermediate $P_2, P_3
\ldots$ this will in general not be the case. In the composition
series there are thus certain distinguished terms which are at the
same time normal divisors of $P$ itself, and which play a special
role\footnote{They form the so-called principal series of $P$
(\gesperrt{Netto}).}. On this we derive in what follows one more
important theorem.

Suppose $Q$ is a term of a composition series of $P$ which is
at the same time a normal divisor of $P$, while the term following upon $Q$,
%CONT>>


% ===== scan leaf 46, printed 25 =====
% !PRINTED 25
\weberpage{25}

%<<CONT
% the German page break falls inside "das auf $Q$ | folgende Glied"; the English absorbs
% "following" onto the previous page, so this page opens directly with $Q_1$
$Q_1$, is not a normal divisor of $P$. Then among the groups $a^{-1} Q_1 a$ conjugate with respect to $P$ there are several distinct ones, which we shall denote by

\weq{1}{Q_1,\ Q_1',\ Q_1'',\ Q_1''' \cdots}

\gesperrt{The greatest common divisor $R$ of all these groups is again a normal divisor of $P$.}

Now all the groups (1) are greatest normal divisors of $Q$, as follows by the concluding theorem of \S.\,3 from the fact that $Q, Q_1$ are isomorphic with $a^{-1} Q a, a^{-1} Q_1 a$, while $Q$, as a normal divisor of $P$, is identical with $a^{-1} Q a$. We can therefore let any one of the groups (1), which all have the same index $\nu$, follow $Q$ in the composition series.

If now $Q_2$ is the greatest common divisor of $Q_1, Q_1'$ (and hence distinct from $Q_1$), then by the theorem \S.\,6, 4. $Q_2$ is a greatest normal divisor of $Q_1$ and $Q_1'$, whose index is likewise $\nu$. We can therefore continue the composition series from $Q$ on in the following two ways:

\weq{2}{
\begin{array}{ccc@{\qquad}ccc}
Q, & Q_1, & Q_2, & Q & Q_1' & Q_2 \\
& \nu & \nu & & \nu & \nu
\end{array}
}

We now wish to generalise the law expressed herein and to prove it by complete induction.

If $Q_2$ is not yet equal to $R$, we adjoin to $Q_1, Q_1'$ a third group $Q_1''$ from the sequence (1), so that the intersection $Q_3$ of $Q_1, Q_1', Q_1''$ becomes a proper divisor of $Q_2$, and we continue in this way until we have arrived at the greatest common divisor $R$ of all the groups (1).

Accordingly, let:

\weq{3}{
\begin{array}{l}
Q_2 \text{ the intersection of } Q_1,\ Q_1' \\
Q_3 \text{ \ ,,\qquad ,,\qquad ,,\ } Q_1,\ Q_1',\ Q_1'' \\
\cdot\ \cdot\ \cdot\ \cdot\ \cdot\ \cdot\ \cdot\ \cdot\ \cdot\ \cdot\ \cdot\ \cdot\ \cdot \\
Q_r \text{ \ ,,\qquad ,,\qquad ,,\ } Q_1,\ Q_1',\ Q_1'' \cdots Q_1^{(r-1)}
\end{array}
}

and let the arrangement of the groups $Q_1, Q_1' \cdots Q_1^{(r-1)}$ be so chosen that, of the groups $Q_1, Q_2 \cdots Q_r$, each is a proper divisor of the one preceding it. This arrangement can be continued so long as $Q_r$ is not equal to $R$.

We wish to prove that we can choose the composition series of $P$ in such a way that in it

\weq{4}{
\begin{array}{cccc}
Q, & Q_1, & Q_2 & \cdots\ Q_r\\
& \nu' & \nu & \nu % [sic: the source prints \nu' although the text says all indices are \nu.]
\end{array}
}
%CONT>>


% ===== scan leaf 47, printed 26 =====
% !PRINTED 26
\weberpage{26}

%<<CONT
occurs, and that the indices in this part of the series are all
equal to $\nu$.

The assertion is true for $r = 2$, and we prove it in general
by complete induction, under the assumption that it has already
been proved true for $r$.

If $Q_r$ is not equal to $R$, we choose $Q_1^{(r)}$ so that $Q_{r+1}$ is the
intersection of $Q_1, Q_1' \cdots Q_1^{(r-1)}, Q_1^{(r)}$, and moreover denote
the intersection of $Q_1, Q_1' \cdots Q_1^{(r-2)}, Q_1^{(r)}$ by $Q_r'$. Then
$Q_r'$ is distinct from $Q_r$, and $Q_{r+1}$ is the intersection of
$Q_r$ and $Q_r'$. Now by the assumption (4) we have two
corresponding pieces of the composition series with the constant
index series $\nu$

\weq{5}{
\begin{array}{l}
Q,\ Q_1 \cdots Q_{r-1}\ Q_r,\\
Q,\ Q_1 \cdots Q_{r-1}\ Q_r';
\end{array}
}

for the second of these series is formed exactly like the first
according to the prescription (3), only that $Q_1^{(r)}$ has taken the
place of $Q_1^{(r-1)}$. By the theorem of \S.\,6, 2., therefore, $Q_{r+1}$ is
the greatest normal divisor both of $Q_r$ and of $Q_r'$, and indeed
of index $\nu$, and it follows therefore, as was to be proved, that
the composition series can be continued thus:

\weq{6}{Q,\ Q_1 \cdots Q_{r-1},\ Q_r,\ Q_{r+1},}

and that the newly added index is likewise $\nu$.

This we can now continue until we have reached the group $R$,
and we obtain a piece of the composition series

\weq{7}{
\begin{array}{ccc}
Q,\ Q_1,\ Q_2 & \cdots & R,\\
\nu\quad\ \nu & & \nu
\end{array}
}

whose last term $R$ is again a normal divisor of $P$ and in
which all the indices agree.

If we continue the composition series onward from $R$, the
index may change; but from there on it can again be kept
equal, until one once more arrives at a normal divisor of $P$.

With this the following theorem is proved:

III.\ \gesperrt{One can so arrange the composition series of $P$ with
its index series}

\weq{8}{
\begin{array}{c}
P,\ P_1,\ P_2,\ P_3 \cdots\\
j_1,\ j_2,\ j_3 \cdots
\end{array}
}

\gesperrt{that, whenever a change occurs in the index
series, that is, whenever $j_{r+1}$ is different from $j_r$,
$P_r$ is a normal divisor of $P$.}


% ===== scan leaf 48, printed 27 =====
% !PRINTED 27
\weberpage{27}
\pdfbookmark[2]{§ 8. Metacyclic Groups}{bm10}

\begin{center}
\S.\,8.

\gesperrt{Metacyclic Groups.}
\end{center}

We defined metacyclic groups in \S.\,176 of the first volume, in connection with the permutation groups, and there learned their significance for the solution of equations. The concept, however, is independent of the particular meaning of the group elements, and we therefore now define generally:

\gesperrt{A group whose index series consists entirely of primes shall be called a metacyclic group.}

For these groups there holds a theorem important for the applications to algebra, but which is independent of these applications and which we therefore derive here. The theorem reads:

IV.\ \gesperrt{A metacyclic group has a commutative normal
divisor different from the unit
group}\footnote{C. J o r d a n, Traité des substitutions, Livre IV.}.

If $P$ is a metacyclic group and

\weq{1}{
\begin{array}{c}
P,\ P_1,\ P_2 \cdots P_{\mu-1},\ 1\\
j_1,\ j_2 \cdots j_{\mu-1},\ j_\mu
\end{array}
}

is the composition series, whose index series $j_1, j_2 \cdots$ consists entirely of primes, then $P_{\mu-1}$ is in any case a commutative group; for if $\pi$ is any element of $P_{\mu-1}$ different from 1, then the powers

\weq{2}{1,\ \pi,\ \pi^2 \cdots \pi^{j_\mu - 1}}

(because $j_\mu$ is a prime) are all different from one another, and therefore make up the whole group $P_{\mu-1}$; but the cyclic group (2) is certainly commutative, because for any two exponents $h, k$ we always have

\[
\pi^h \pi^k = \pi^k \pi^h = \pi^{h+k}
\]

At the same time, by the definition of the composition series. % [sic] the full stop after "Compositionsreihe" is as printed; a comma is evidently intended
$P_{\mu-1}$ is a normal divisor of $P_{\mu-2}$.

We accordingly now assume that we have a commutative normal divisor $Q_\nu$, different from the unit, of some group $P_\nu$ of the sequence (1), and suppose moreover,
%CONT>>


% ===== scan leaf 49, printed 28 =====
% !PRINTED 28
\weberpage{28}

%<<CONT
if there are several divisors of $P_\nu$ of the same character as $Q_\nu$,
that one of the lowest possible degree is taken for $Q_\nu$.

If we then show how one can derive from $Q_\nu$ a commutative
normal divisor $Q_{\nu-1}$ of $P_{\nu-1}$ which at the same time
contains $Q_\nu$, then we can continue this procedure until
we have arrived at a commutative normal divisor $Q$ of $P$ itself,
as our theorem requires.

If we choose any element $\gamma$ in the group $P_{\nu-1}$
which is not contained in $P_\nu$, then, since $P_\nu$ is a normal divisor
of $P_{\nu-1}$,

\weq{3}{\gamma P_\nu = P_\nu \gamma,}

and if therefore $h$ is the lowest exponent for which $\gamma^h$ is
contained in $P_\nu$, then $h$ is greater than 1 and

\weq{4}{\gamma^h P_\nu = P_\nu,}

and the elements contained in the system of subsidiary groups

\weq{5}{P_\nu,\ \gamma P_\nu,\ \gamma^2 P_\nu \ldots \gamma^{h-1} P_\nu}

which are all different from one another,
form a group whose degree, if the degree of $P_\nu$ is denoted by $p_\nu$,
is equal to $h p_\nu$. But the group (5) is also
a divisor of $P_{\nu-1}$, and consequently $h p_\nu$ must be a divisor of
$p_{\nu-1} = j_\nu p_\nu$, and $h$ is a divisor of $j_\nu$. But since $j_\nu$
is a prime, $h = j_\nu$ must hold, and by (5) the
whole group $P_{\nu-1}$ is exhausted:

\weq{6}{P_{\nu-1} = P_\nu + \gamma P_\nu + \gamma^2 P_\nu + \cdots \gamma^{j_\nu - 1} P_\nu.}

We now consider the system of transformed groups

\weq{7}{\gamma^{-1} Q_\nu \gamma,}

in which $\gamma$ runs through all the elements of $P_{\nu-1}$.

These groups are all isomorphic with one another and are
therefore, just like $Q_\nu$, commutative. They are all contained in $P_\nu$,
and since $\gamma^{-1} P_\nu \gamma = P_\nu$, they are, by the concluding theorem of \S.\,3,
normal divisors of $P_\nu$.

If all the groups (7) are identical with one another, then
$Q_\nu$ is a normal divisor of $P_{\nu-1}$, and we attain our purpose
of determining $Q_{\nu-1}$ simply by taking for $Q_{\nu-1}$ either $Q_\nu$ itself,
or, in case a commutative normal divisor of lower
degree of $P_{\nu-1}$ should still exist, this latter.


% ===== scan leaf 50, printed 29 =====
% !PRINTED 29
\weberpage{29}
In the other case we will denote the mutually distinct ones among the groups (7) by
\weq{8}{Q_\nu, \; Q'_\nu, \; Q''_\nu \; \ldots}
and from them derive a group
\weq{9}{R = (Q_\nu, \; Q'_\nu, \; Q''_\nu \; \ldots)}
which shall be defined by this, that it is the smallest group that contains each of the groups (8) as a divisor, and which we may denote as the \gesperrt{least common multiple} of the groups $Q_\nu, Q'_\nu, Q''_\nu \ldots$.

That there is one and only one such group $R$, and that every group that is divisible by $Q_\nu, Q'_\nu, Q''_\nu \ldots$ must also be divisible by $R$, is easy to see. For, first, there are finite groups, e.g. the group $P_\nu$, that contain all the groups (8) as divisors, and secondly, if two groups $R, R'$ contain all these groups, the same holds also of the intersection of $R$ and $R'$. If therefore $R$ is of the lowest possible degree, this intersection must be identical with $R$, and $R$ is thus a \gesperrt{divisor of} $R'$. If $R'$ is also of the lowest possible degree, then $R'$ is not different from $R$.

Of this group $R$ we now prove:

a) that it is a normal divisor of $P_{\nu-1}$,

b) that it is a commutative group.

When we have proved both of these, we are at our goal; for then, if there is not also a commutative normal divisor of $P_{\nu-1}$ of \gesperrt{lower degree}, we can choose $R$ itself for $Q_{\nu-1}$.

But the first is clear without further ado. For if $\gamma$ denotes any element of $P_{\nu-1}$, then
\[ \gamma^{-1} R \gamma = (\gamma^{-1} Q_\nu \gamma, \; \gamma^{-1} Q'_\nu \gamma, \; \gamma^{-1} Q''_\nu \gamma \; \ldots) \]
is the least common multiple of the groups
\[ \gamma^{-1} Q_\nu \gamma, \; \gamma^{-1} Q'_\nu \gamma, \; \gamma^{-1} Q''_\nu \gamma \; \ldots \]
But these groups, by the definition (7), are in their totality not different from the groups
\[ Q_\nu, \; Q'_\nu, \; Q''_\nu \; \ldots \]
and consequently
\[ \gamma^{-1} R \gamma = R, \]
that is, $R$ is a normal divisor of $P_{\nu-1}$.


% ===== scan leaf 51, printed 30 =====
% !PRINTED 30
\weberpage{30}
But in order to prove that $R$ is a commutative group, we must first set out two auxiliary considerations.

1) It is to be proved that any two of the groups $Q_\nu, Q'_\nu, Q''_\nu \ldots$ have no common divisor other than the unit.

For this we need only prove that $Q_\nu, Q'_\nu$ have this property. For if
\[ Q'_\nu = \gamma'^{-1} Q_\nu \gamma', \quad Q''_\nu = \gamma''^{-1} Q_\nu \gamma'' \]
have the greatest common divisor $T$, then $Q_\nu$ and $Q'''_\nu = \gamma' Q''_\nu \gamma'^{-1}$ have the greatest common divisor $\gamma' T \gamma'^{-1}$.

If we assume, then, that $T$ is the greatest common divisor of $Q_\nu$ and $Q'_\nu$, then $T$ is certainly a commutative group, because $Q_\nu, Q'_\nu$ are such. But $T$ is also a normal divisor of $P_\nu$. For $Q_\nu$ and $Q'_\nu$ are such divisors; and if therefore $\tau$ is an element of $T$ and $\pi$ an element of $P_\nu$, then $\pi^{-1}\tau\pi$ is contained both in $Q_\nu$ and in $Q'_\nu$, and hence also in $T$. Consequently $\pi^{-1} T \pi = T$.

Now we have assumed, however, that $Q_\nu$ is a commutative normal divisor of $P_\nu$ above the unit of the lowest possible degree; further, $Q_\nu$ and $Q'_\nu$ were different from one another. Hence $T$ can only be the unit group itself, q. e. d.

2) To form the group $R$, we may proceed in this way: we compound the elements of $Q_\nu, Q'_\nu, Q''_\nu \ldots$ with one another, in arbitrary order and with repetition, so long as new elements still arise. For all the compositions so formed constitute a group $C$, which is contained in $R$, because $R$ contains the individual elements of $Q_\nu, Q'_\nu, Q''_\nu \ldots$. On the other hand, the groups $Q_\nu, Q'_\nu, Q''_\nu \ldots$ are also contained in $C$, and consequently $R$ is contained in $C$ and $C$ is identical with $R$.

In order, then, finally to prove that $R$ is a commutative group, it only remains for us to show that, if $\alpha, \beta$ are any two elements of $Q_\nu, Q'_\nu, Q''_\nu \ldots$, the permutability
\weq{10}{\alpha\beta = \beta\alpha}
holds.

Now if $\alpha, \beta$ both belong to the same group $Q_\nu$, then (10) is a consequence of our assumption that $Q_\nu$ is commutative; and it is likewise so if both occur in one of the other groups $Q'_\nu, Q''_\nu \ldots$.


% ===== scan leaf 52, printed 31 =====
% !PRINTED 31
\weberpage{31}
If however $\alpha$ and $\beta$ are contained in two different groups, say in $Q_\nu$ and $Q'_\nu$, we reason as follows: since $\beta$ is contained in $P_\nu$ and $Q_\nu$ is a normal divisor of $P_\nu$,
\weq{11}{\beta^{-1}\alpha\beta = \alpha'}
is likewise contained in $Q_\nu$. From this it follows:
\weq{12}{\beta\alpha'\alpha^{-1} = \alpha\beta\alpha^{-1} = \beta',}
and since $\alpha$ is contained in $P_\nu$ and $Q'_\nu$ is a normal divisor of $P_\nu$, $\beta'$ is likewise contained in $Q'_\nu$. But from (12) it follows that
\[ \alpha'\alpha^{-1} = \beta^{-1}\beta' = \delta, \]
and accordingly the element $\delta$ is contained both in $Q_\nu$ and in $Q'_\nu$. By what was proved under 1), it can therefore only be the unit element, i.e. we have
\[ \alpha' = \alpha, \quad \beta' = \beta, \]
and accordingly it follows from (11) that $\alpha\beta = \beta\alpha$ must hold, as was to be proved.

\begin{center}
\rule{3cm}{0.4pt}
\end{center}
% [The source prints this plain horizontal rule below the paragraph, marking the close of \S.\,8; the remainder of the page is blank.]


% ===== scan leaf 53, printed 32 =====
\weberpage{32}% page number not printed on this page
\pdfbookmark[1]{Second Section. Abelian Groups.}{bm11}
\pdfbookmark[2]{§ 9. Representation of Abelian Groups by a Basis}{bm12}
% !PRINTED none
% [part-title / section-opening page: no running-head page number visible in any of the three bands; centred "Zweiter Abschnitt." / "Abel'sche Gruppen." title with rule, then \S.\,9 heading, then body text begins]
\begin{center}
\gesperrt{Second Chapter.}

\gesperrt{Abelian Groups.}

\rule{0.3\linewidth}{0.4pt}

\S.\,9.

\gesperrt{Representation of Abelian groups by a basis.}
\end{center}

We have already, in \S.\,1, called those groups in which the commutative law holds for composition commutative or \gesperrt{Abelian} groups. In these groups, which are of particular importance in the applications, much simpler laws prevail than in the general groups, as we shall now see. For the composition of the elements in these groups the same rules hold as in the multiplication of numbers, and we accordingly call the composita of elements of such a group, as in multiplication, products and powers.

Here too we restrict ourselves to the consideration of finite groups.

From the definition of Abelian groups it follows that the commutative law also holds for the composition of the parts of such a group (\S.\,4). Every divisor of an Abelian group is itself an Abelian group, and is at the same time a normal divisor, so that in these groups we may omit the addition ``Normal''.

A group which consists only of the repetitions of one element, such as
\[ 1, \; A, \; A^2 \; \ldots \; A^{a-1}, \]
we have already earlier (Bd. I, \S.\,148) called a cyclic group, and we retain this designation here. If $a$ is the smallest positive number for which $A^a = 1$, and thus $a$ is the degree of the group, then $a$ is also called the degree of the element $A$. If $m$ is any exponent for which $A^m = 1$, then $m$ is a
%CONT>>


% ===== scan leaf 54, printed 33 =====
% !PRINTED 33
\weberpage{33}
%<<CONT
multiple of $a$. The element 1 has degree 1. All cyclic groups are commutative.

Now for every Abelian group $S$ there holds the fundamental theorem that a system of elements $A, B, C \ldots$ of degrees $a, b, c \ldots$ can always be selected in such a way that in the form
\[
\Theta = A^\alpha B^\beta C^\gamma \ldots
\]
every element $\Theta$ of $S$, and each only once, can be represented, when $\alpha$ runs through a complete system of residues to the modulus $a$, $\beta$ a complete system of residues to the modulus $b$, $\gamma$ a complete system of residues to the modulus $c$, etc.

A system of elements that is suited to such a representation is called a \gesperrt{basis of the group}, and we can accordingly state the theorem to be proved thus:

\gesperrt{I. Every Abelian group can be represented by a basis.}

The proof of the theorem rests on the following sequence of inferences:

\gesperrt{1. If $n$ is the degree of the group $S$, if $A_1, A_2 \ldots A_n$ are its elements and $a_1, a_2 \ldots a_n$ their degrees, and if further $\alpha_1, \alpha_2 \ldots \alpha_n$ run through complete systems of residues to the moduli $a_1, a_2 \ldots a_n$, then in the form}
\weq{1}{\Theta = A_1^{\alpha_1} A_2^{\alpha_2} \ldots A_n^{\alpha_n}}
\gesperrt{every element of $S$ is represented, and each equally often.}

That one obtains every element at all in the form (1) is seen immediately; for one obtains, e.\,g., $A_1$ by setting $\alpha_1 = 1, \alpha_2 = 0 \ldots \alpha_n = 0$.

It remains to be proved that one obtains every element equally often.

But if
\weq{2}{1 = A_1^{h_1} A_2^{h_2} \ldots A_n^{h_n}}
is a representation of the element 1, then $\Theta$ does not change when $\alpha_1, \alpha_2 \ldots \alpha_n$ in (1) are replaced by $\alpha_1 + h_1, \alpha_2 + h_2 \ldots \alpha_n + h_n$. Thus every element $\Theta$ is represented by (1) at least as often as the element 1 is by (2).

If, on the other hand, the two sequences of exponents $\alpha_1, \alpha_2 \ldots \alpha_n$ and $\alpha_1', \alpha_2' \ldots \alpha_n'$ give two representations of the element $\Theta$, then one has
%CONT>>


% ===== scan leaf 55, printed 34 =====
% !PRINTED 34
\weberpage{34}
%<<CONT
\weq{3}{1 = A_1^{\alpha_1' - \alpha_1} A_2^{\alpha_2' - \alpha_2} \ldots A_n^{\alpha_n' - \alpha_n},}
thus a representation of the element 1. Hence by (2) we may set:
\[
\alpha_1' = \alpha_1 + h_1,\ \alpha_2' = \alpha_2 + h_2 \ldots \alpha_n' = \alpha_n + h_n.
\]
Consequently there can also be no more distinct representations of the element $\Theta$ than the number of representations of the element 1 amounts to. We denote the number of these representations by $h$. In all, however, the number of the various possible determinations of the $\alpha$ in (1) is equal to the product $a_1 a_2 \ldots a_n$, and since $n$ is the number of elements of $S$, it follows that
\weq{4}{n h = a_1 a_2 \ldots a_n.}
From the formula (4) results the theorem:

\gesperrt{2. If $r$ is a prime dividing the degree $n$ of the group, then there is in $S$ an element of degree $r$.}

For from (4) one sees first of all that one of the degrees $a_1, a_2, \ldots a_n$ is divisible by $r$. If then, say, $a_1 = rk$, then $A_1^k$ is an element of degree $r$.

\gesperrt{3. If $A, B \ldots$ are any elements whatever in $S$ of the degrees $a, b \ldots$, then $AB \ldots$ is also an element in $S$, and the degree of this product is a divisor of the least common multiple $m$ of $a, b \ldots$}

For we have
\[
(AB \ldots)^m = A^m B^m \ldots = 1,
\]
and consequently $m$ is a multiple of the degree of $AB \ldots$

\gesperrt{4. If the degree $n$ of the group $S$ is decomposed into two factors $n = ab$ such that $a$ and $b$ are relatively prime, then there are in $S$ exactly $a$ elements $A$ whose degree is a divisor of $a$, and $b$ elements $B$ whose degree is a divisor of $b$, and in the form}
\weq{5}{\Theta = AB}
\gesperrt{all the elements of $S$ are contained, and each only once.}

To see the correctness of this theorem, we make the following consideration. The totality $\mathfrak{A}$ of the elements $A$ whose
%CONT>>


% ===== scan leaf 56, printed 35 =====
% !PRINTED 35
\weberpage{35}
%<<CONT
degree is a divisor of $a$, is a group contained in $S$; for if $A, A'$ are two such elements, then by 3. the degree of $AA'$ is a divisor of $a$, and $AA'$ is therefore also contained in $\mathfrak{A}$.

Likewise the totality $\mathfrak{B}$ of the elements $B$ whose degree is a divisor of $b$ is a group. The element 1 occurs in $\mathfrak{A}$ as well as in $\mathfrak{B}$; apart from this the two contain no common element.

The degree $a'$ of $\mathfrak{A}$ is relatively prime to $b$. For if $r$ is a prime dividing $a'$, then by 2. there is in $\mathfrak{A}$ an element of degree $r$. But since the degree of every element of $\mathfrak{A}$ is a divisor of $a$, $r$ too is a divisor of $a$ and not of $b$.

In the same way one proves that the degree $b'$ of $\mathfrak{B}$ is relatively prime to $a$.

Now determine (by Bd.\,I, \S.\,118) two integers $x$ and $y$ so that
\weq{6}{a x + b y = 1}
and take any element $\Theta$ of $S$. Then
\weq{7}{\Theta = \Theta^{ax} \Theta^{by},}
and now, since $(\Theta^{by})^a = 1$, $\Theta^{by}$ is contained in $\mathfrak{A}$, and for the same reason $\Theta^{ax}$ in $\mathfrak{B}$. Hence it follows from (7):
\weq{8}{\Theta = AB.}
Accordingly every element $\Theta$ is contained in the form $AB$. But such a representation is also possible in only one way. For if $A', B'$ are two elements from $\mathfrak{A}$ and $\mathfrak{B}$, and if
\[
AB = A'B',
\]
then, on raising to the power $by = 1 - ax$, it follows that $A = A'$ and consequently also $B = B'$. But the number of the distinct products of the form $AB$ is $= a'b'$, and therefore one has
\[
n = ab = a'b',
\]
and since $a$ is relatively prime to $b'$ and $b$ relatively prime to $a'$:
\[
a' = a,\quad b' = b.
\]
With this the theorem 4. is proved in all its parts.

If now the two groups $\mathfrak{A}, \mathfrak{B}$ are represented by bases, then it follows from formula (8) that $S$ too is represented by a basis, and the basis of $S$ contains the elements of the bases of $\mathfrak{A}$ and $\mathfrak{B}$ and no others.


% ===== scan leaf 57, printed 36 =====
% !PRINTED 36
\weberpage{36}
If $a$ and $b$ are further decomposable into factors which are relatively prime to one another, then we can proceed with the groups $\mathfrak{A}$ and $\mathfrak{B}$ in the same way again, and we thus arrive at the result that our Theorem I. is proved in general, provided we can still establish it for groups whose degree is a power of a prime.

So let the degree $n$ of the group $S$ now be a power of a prime $p$
\weq{9}{n = p^k.}
The degrees of all elements of $S$, which are of course divisors of $n$, must then likewise be powers of $p$. It is to be shown that such a group is representable by a basis.

Choose in $S$ an element $A$ of the highest possible degree $a$. Then $a$ is a power of $p$, and the degrees of all the other elements are divisors of $a$, so that for every element $\Theta$ of $S$
\weq{10}{\Theta^a = 1}
holds. The elements
\weq{11}{1, A, A^2 \ldots A^{a-1}}
are all different from one another, and their totality is a divisor of $S$. If the group $S$ is thereby exhausted, that is, if every element is of the form $A^\alpha$, then $S$ is represented by a one-term basis. But if $S$ is not yet exhausted by (11), then for every element $\Theta$ of $S$ there will nevertheless be a certain exponent $h$ such that $\Theta^h$ is contained in the sequence (11). This will certainly occur when $h$ is the degree of $\Theta$, that is, $\Theta^h = 1$. Among these numbers $h$ one will be the least positive, which we shall denote by $b$. Thus for every element $\Theta$ there is a certain least positive number $b$ such that
\weq{12}{\Theta^b = A^\lambda}
is contained in the sequence (11).

This number $b$ is a divisor of $a$, hence also a power of $p$, and at the same time a divisor of $\lambda$. For if we set $a = qb + b'$, where $0 \leqq b' < b$, it follows from (10) and (12) that
\[
\Theta^a = A^{\lambda q} \Theta^{b'} = 1,
\]
or $\Theta^{b'} = A^{-\lambda q}$, and so if $b'$ is not zero, there is, contrary to the assumption, a number still smaller than $b$, namely $b'$, which satisfies the requirement (12).


% ===== scan leaf 58, printed 37 =====
% !PRINTED 37
\weberpage{37}
From (12) it follows further that
\[
1 = \Theta^a = A^{\frac{\lambda a}{b}};
\]
hence $\lambda a : b$ must be a multiple of $a$, and consequently $\lambda$ a multiple of $b$. If we therefore set
\weq{13}{B = \Theta A^{-\frac{\lambda}{b}}}
then $B^b = 1$, and at the same time $B^b$ is the lowest power of $B$ that becomes equal to a power of $A$, because, if $B^{b'}$ is a power of $A$, the same, by (13), holds also of $\Theta^{b'}$. We assume the element $\Theta$ to be so chosen that $b$ becomes as great as possible. Then for every other element $\Theta_1$ of $S$ as well, $\Theta_1^b$ is always a power of $A$ whose exponent is divisible by $b$.

If now
\[
\begin{gathered}
\alpha = 0, 1, 2 \ldots a-1 \\
\beta = 0, 1, 2 \ldots b-1,
\end{gathered}
\]
then the elements
\weq{14}{A^\alpha B^\beta}
are, $ab$ in number, all different from one another. They form a divisor $S'$ contained in $S$, which is represented by the basis $A, B$. If $S'$ is identical with $S$, we have reached our goal.

But if $S$ is not yet exhausted by (14), we continue in the same manner.

By an application of complete induction we shall draw the conclusion at once in full generality.

Let a divisor $S_{\nu-1}$ of $S$ have been found, represented by a basis in the form
\weq{15}{S_{\nu-1} = A_1^{\alpha_1} A_2^{\alpha_2} \ldots A_{\nu-1}^{\alpha_{\nu-1}},}
about which we make the following assumptions:

1) Let the degrees of $A_1, A_2 \ldots A_{\nu-1}$ be $a_1, a_2 \ldots a_{\nu-1}$, and let
\[
a_1 \geqq a_2 \geqq \ldots \geqq a_{\nu-1}.
\]

2) For every element $\Theta$ contained in $S$ let
\weq{16}{\Theta^{a_{\nu-1}} = A_1^{\lambda_1} A_2^{\lambda_2} \ldots A_{\nu-2}^{\lambda_{\nu-2}},}
that is, contained in $S_{\nu-2}$, and let the exponents $\lambda_1, \lambda_2 \ldots \lambda_{\nu-2}$ be divisible by $a_{\nu-1}$.

The group $S'$ derived above satisfies these requirements if $\nu = 3$ and $a_{\nu-1} = b$ is set, and it is now to be shown
%CONT>>


% ===== scan leaf 59, printed 38 =====
% !PRINTED 38
\weberpage{38}
%<<CONT
how, if $S_{\nu-1}$ does not yet exhaust the whole group $S$, one can derive from it a more comprehensive group $S_\nu$ of the same kind.

For every element $\Theta$ of $S$ there will be a certain lowest positive exponent $a_\nu$ for which $\Theta^{a_\nu}$ is contained in $S_{\nu-1}$, and this $a_\nu$ is, by (16), equal to or smaller than $a_{\nu-1}$, and is moreover a power of $p$, since it must be a divisor of the degree of $\Theta$. We choose $\Theta$ so that the exponent $a_\nu$ becomes as large as possible. If $\Theta_1$ is any other element in $S$, and $\Theta_1^h$ the lowest power of $\Theta_1$ that is contained in $S_{\nu-1}$, then $h$ too is a power of $p$ and equal to or smaller than $a_\nu$. Hence $\Theta_1^{a_\nu}$ is contained in $S_{\nu-1}$.

But if we set
\weq{17}{\Theta_1^{a_\nu} = A_1^{\lambda_1} A_2^{\lambda_2} \ldots A_{\nu-1}^{\lambda_{\nu-1}},}
then all the exponents $\lambda$ are divisible by $a_\nu$. For we have
\[
\Theta_1^{a_{\nu-1}} = A_1^{\frac{\lambda_1 a_{\nu-1}}{a_\nu}} A_2^{\frac{\lambda_2 a_{\nu-1}}{a_\nu}} \ldots A_{\nu-1}^{\frac{\lambda_{\nu-1} a_{\nu-1}}{a_\nu}},
\]
and by hypothesis 2) the exponents of $A_1, A_2 \ldots A_{\nu-1}$ on the right-hand side of this formula must be integers divisible by $a_{\nu-1}$. Consequently $\lambda_1, \lambda_2 \ldots \lambda_{\nu-1}$ are divisible by $a_\nu$. We can therefore, returning to the special $\Theta$ considered above and understanding by $h_1, h_2 \ldots h_{\nu-1}$ integers, set
\[
\Theta^{a_\nu} = A_1^{h_1 a_\nu} A_2^{h_2 a_\nu} \ldots A_{\nu-1}^{h_{\nu-1} a_\nu}
\]
and if we then take
\[
A_\nu = \Theta A_1^{-h_1} A_2^{-h_2} \ldots A_{\nu-1}^{-h_{\nu-1}}
\]
then
\weq{18}{A_\nu^{a_\nu} = 1}
is the lowest power of $A_\nu$ that is contained in $S_{\nu-1}$ (since otherwise a still lower power of $\Theta$ would also be contained in $S_{\nu-1}$).

Then
\weq{19}{A_1^{\alpha_1} A_2^{\alpha_2} \ldots A_\nu^{\alpha_\nu}, \quad \alpha_i = 0, 1, 2 \ldots a_i - 1}
is $= 1$ only if $\alpha_1, \alpha_2 \ldots \alpha_\nu$ are divisible by $a_1, a_2 \ldots a_\nu$ and consequently are $= 0$, and accordingly the elements represented in (19) are all different from one another. But these elements form a group $S_\nu$ with the basis $A_1, A_2 \ldots A_\nu$, which satisfies requirement 1).


% ===== scan leaf 60, printed 39 =====
% !PRINTED 39
\weberpage{39}
\pdfbookmark[2]{§ 10. The Invariants of the Abelian Groups}{bm13}
At the same time, as formula (17) shows, if $\Theta_1$ is an arbitrary element in $S$, then $\Theta_1^{a_\nu}$ is contained in $S_{\nu-1}$, and the exponents $\lambda_1, \lambda_2 \ldots \lambda_{\nu-1}$ are divisible by $a_\nu$. Hence requirement 2) is also satisfied.

With this, then, our Theorem I. is proved. At the same time we see from this derivation \gesperrt{that for an arbitrary Abelian group $S$ the elements of the basis can always be taken in such a way that their degrees are prime powers.}

From the representability by a basis it also follows that every Abelian group is metacyclic; for if the elements of such a group $S$ are represented in the form:
\[
A = A_1^{\alpha_1} A_2^{\alpha_2} \ldots A_\nu^{\alpha_\nu},
\]
and if $p$ is a prime which divides the degree $a_1$ of $A_1$, then the elements
\[
A' = A_1^{\alpha_1 p} A_2^{\alpha_2} \ldots A_\nu^{\alpha_\nu},
\]
when $\alpha_1$ runs through a complete system of residues to the modulus $a_1 : p$, form a divisor $S'$ of $S$ of index $p$, which is representable by the basis $A_1^p, A_2 \ldots A_\nu$. From $S'$ one can again find in the same way a divisor of prime index, and so forth. Hence $S$ is metacyclic (\S.\,8).

\begin{center}
\S.\,10.

\gesperrt{The Invariants of the Abelian Groups.}
\end{center}

The proof which we communicated in the previous section for the possibility of representing an Abelian group by a basis contains at the same time a way of finding such a basis, and indeed one in which the degrees of the basis elements are prime powers. Nevertheless it can happen that one and the same group can be represented by bases in different ways.

If we consider, for example, two elements $A, B$ whose degrees $a, b$ are relatively prime, then by these, as elements of a two-term basis, a group
\weq{1}{A^\alpha B^\beta \qquad \begin{aligned}
\alpha &= 0, 1, 2 \ldots a-1 \\
\beta &= 0, 1, 2 \ldots b-1
\end{aligned}}
is represented. But the same group can also be represented by the one-term basis $AB$ in the form
\weq{2}{(AB)^s \qquad s = 0, 1, 2 \ldots ab-1.}


% ===== scan leaf 61, printed 40 =====
% !PRINTED 40
\weberpage{40}
For if $\alpha$ and $\beta$ are given arbitrarily, one can determine $s$ with respect to the modulus $ab$ so that
\[
s \equiv \alpha \ (\mathrm{mod}\ a), \quad s \equiv \beta \ (\mathrm{mod}\ b),
\]
whereby (2) becomes identical with (1). Nevertheless, in a certain sense the constitution of the basis is completely determined by the nature of the groups, according to the following theorem:

\gesperrt{II. If
\[
A_1, A_2 \ldots A_\nu
\]
with the degrees
\[
a_1, a_2 \ldots a_\nu
\]
and
\[
B_1, B_2 \ldots B_\mu
\]
with the degrees
\[
b_1, b_2 \ldots b_\mu
\]
are two bases of an Abelian group $S$ of degree $n$, if $p$ is a prime dividing $n$ and
\weq{1}{a_1 = p_1 a_1', \quad a_2 = p_2 a_2', \quad \ldots \quad a_\nu = p_\nu a_\nu',}
\weq{2}{b_1 = p_1' b_1', \quad b_2 = p_2' b_2', \quad \ldots \quad b_\mu = p_\mu' b_\mu',}
where $p_1, p_2 \ldots p_\nu, p_1', p_2' \ldots p_\mu'$ are the highest powers of $p$ that divide $a_1, a_2 \ldots a_\nu, b_1, b_2 \ldots b_\mu$, then all the prime powers $p_1, p_2 \ldots p_\nu$ that are greater than $1$ also occur among the $p_1', p_2' \ldots p_\mu'$, and conversely.}

To prove this theorem, we take the elements of the two bases $A$ and $B$ to be so ordered that
\weq{3}{\begin{aligned}
p_1 &\geqq p_2 \geqq \ldots \geqq p_\nu \\
p_1' &\geqq p_2' \geqq \ldots \geqq p_\mu'.
\end{aligned}}
The integers $a_1', a_2' \ldots a_\nu', b_1', b_2' \ldots b_\mu'$ occurring in (1) and (2) are by their definition not divisible by $p$, and if we therefore denote by $m$ the least common multiple of $a_1', a_2' \ldots a_\nu'$, then $m$ too is not divisible by $p$. But if
\weq{4}{\Theta = A_1^{\alpha_1} A_2^{\alpha_2} \ldots A_\nu^{\alpha_\nu}}
is an arbitrary element of $S$, it follows that
\[
\Theta^{p_1 m} = 1
\]
must hold.


% ===== scan leaf 62, printed 41 =====
% !PRINTED 41
\weberpage{41}

If we substitute $B_1, B_2 \ldots B_\mu$ for $\Theta$ here, it follows that $p_1 m$
is divisible by each of the numbers $b_1, b_2 \ldots b_\mu$. Hence $p_1$ is
divisible by $p_1'$ and $m$ by the least common
multiple of $b_1', b_2' \ldots b_\mu'$. In this argument we may now
interchange $A$ with $B$ throughout. Hence $p_1'$ must also be divisible
by $p_1$, and therefore
\weq{5}{p_1 = p_1',}
and moreover it follows that $m$ is also the least common
 multiple of $b_1', b_2' \ldots b_\mu'$.

In order to prove our theorem in general from this, we assume
that it has been proved for some $s$ that
\weq{6}{p_1 = p_1', p_2 = p_2', \ldots p_{s-1} = p_{s-1}'}
must hold. By (4), for every element $\Theta$:
\weq{7}{\Theta^{p_s m} = A_1^{\alpha_1 p_s m} A_2^{\alpha_2 p_s m} \ldots A_{s-1}^{\alpha_{s-1} p_s m},}
in which $m$, as above, is the least common multiple of
$a_1', a_2' \ldots a_\nu'$ and at the same time of $b_1', b_2' \ldots b_\mu'$, and hence is not
divisible by $p$.

We now determine the number of mutually distinct
 elements contained in the form (7).

If we let $\alpha_1$ run through the sequence of numbers $0, 1, 2 \ldots \dfrac{p_1}{p_s} - 1$,
the elements
\weq{8}{1,\ A_1^{p_s m},\ A_1^{2 p_s m} \ldots A_1^{\left(\frac{p_1}{p_s}-1\right) p_s m}}
are all distinct from one another, while $A_1^{\frac{p_1}{p_s} p_s m}$ becomes $= 1$
again, so that all the other powers $A_1^{\alpha_1 p_s m}$ are found again in the
sequence (8). We conclude in the same way with regard to the remaining
factors of (7), from which the exact number of the mutually
 distinct elements contained in $\Theta^{p_s m}$
\weq{9}{z = \frac{p_1}{p_s} \frac{p_2}{p_s} \ldots \frac{p_{s-1}}{p_s}}
is found.

Of the two numbers $p_s$, $p_s'$, if they are not equal,
one will be the larger. We shall therefore assume that
\weq{10}{p_s' \geqq p_s.}

We now express $\Theta$ in terms of the base $B$, and set
\weq{11}{\Theta^{p_s m} = B_1^{\beta_1 p_s m} B_2^{\beta_2 p_s m} \ldots B_\mu^{\beta_\mu p_s m},}
%CONT>>


% ===== scan leaf 63, printed 42 =====
% !PRINTED 42
\weberpage{42}

%<<CONT
and we now count again how many distinct elements are contained
in this form. The two values for $z$ must then agree.

If we count, as above in (7), the number of the distinct
elements contained in the form
\weq{12}{B_1^{\beta_1 p_s m} B_2^{\beta_2 p_s m} \ldots B_{s-1}^{\beta_{s-1} p_s m}}
then by the assumption (6) this too yields the value $z$. Hence
in the form
\weq{13}{B_s^{\beta_s p_s m} \ldots B_\mu^{\beta_\mu p_s m}}
only the single element $1$ can be contained, from which it is to be
concluded that $p_s$ must be divisible by $p_s'$. But this is compatible
with (10) only under the supposition that
\weq{14}{p_s = p_s' \text{ holds.}}

With this our Theorem II. is proved.

The prime powers contained in the degree numbers of a basis
of $S$ are thus entirely independent of the particular choice of the
basis, and we therefore call them the \gesperrt{invariants of the
group}. The product of all the invariants is equal to the degree $n$
of the group.

By \S.\,9 we can determine a basis for $S$ in which the
degrees of the elements are all prime powers. By
Theorem II. these degrees are the invariants of the group,
and are thus completely determined by the group.

That the invariants also completely determine the nature
of the group follows from the theorem:

III. \gesperrt{Two groups with the same invariants are
isomorphic, and isomorphic groups have the same
 invariants.}

For if two groups $S$ and $S'$ have basis elements agreeing
in number and in degree, if say
\[ \Theta = A_1^{\alpha_1} A_2^{\alpha_2} \ldots A_\nu^{\alpha_\nu} \]
are the elements of $S$, and
\[ \Theta' = B_1^{\alpha_1} B_2^{\alpha_2} \ldots B_\nu^{\alpha_\nu} \]
the elements of $S'$, where in both cases the $\alpha_1, \alpha_2 \ldots \alpha_\nu$ run through
systems of residues with respect to the moduli $a_1, a_2 \ldots a_\nu$, then we need
only let $\Theta$ and $\Theta'$ correspond to one another when
%CONT>>


% ===== scan leaf 64, printed 43 =====
% !PRINTED 43
\weberpage{43}
\pdfbookmark[2]{§ 11. Group Characters}{bm14}

%<<CONT
the exponents $\alpha$ are the same in both; then the two groups
are related to each other isomorphically.

If therefore two groups have the same invariants, we can
make these invariants the degrees of the basis elements in both
groups, and then obtain this case.

If conversely two groups are isomorphic, then the elements
of the other group corresponding to the basis elements of the
one form a basis of the latter, and hence the invariants
must also be the same in both\textsuperscript{1}).

\begin{center}
§. 11.

\gesperrt{Group Characters}.
\end{center}

It is a principal problem in the theory of \gesperrt{Abel}ian
groups to find all divisors of an \gesperrt{Abel}ian group. The
solution of this problem is essentially facilitated by the
introduction of the concept of group characters, which is also
of importance in various other respects.

Let $S$ be an \gesperrt{Abel}ian group of degree $n$, whose elements
shall be denoted by $A$, and let
\weq{1}{A_1, A_2 \ldots A_\nu}
be the elements of a basis of $S$ of degrees $a_1, a_2 \ldots a_\nu$.
Every element $A$ is therefore, and indeed in only one way,
representable in the form
\weq{2}{A = A_1^{\alpha_1} A_2^{\alpha_2} \ldots A_\nu^{\alpha_\nu}}

\textsuperscript{1}) On the theory of \gesperrt{Abel}ian groups compare:

\gesperrt{Gauss}, D\'emonstration de quelques th\'eor\`emes concernant les
périodes des classes des formes binaires du second degré. Werke,
Bd. II, S. 266.

\gesperrt{Schering}, Die Fundamentalclassen der zusammensetzbaren
arithmetischen Formen, Göttinger Abhandlungen, Bd. 14.

\gesperrt{Frobenius} and \gesperrt{Stickelberger}, Ueber Gruppen vertauschbarer
 Elemente. Crelle's Journal, Bd. 86, S. 217.

\gesperrt{Weber}, ``Theorie der \gesperrt{Abel}'schen Zahlkörper'', Acta mathematica,
Bd. 8 u. 9.

The concept of the invariants was first introduced in the memoir
by \gesperrt{Frobenius} and \gesperrt{Stickelberger}, but defined somewhat differently than
here. This older definition, which proves less suitable in its application to cyclotomic
numbers, was followed by the author in the
memoir cited in the Acta mathematica.
%CONT>>


% ===== scan leaf 65, printed 44 =====
% !PRINTED 44
\weberpage{44}

%<<CONT
representable, so that
\weq{3}{0 \leqq \alpha_1 < a_1,\ 0 \leqq \alpha_2 < a_2, \ldots 0 \leqq \alpha_\nu < a_\nu.}

The exponents $\alpha_1, \alpha_2 \ldots \alpha_\nu$, or any numbers congruent to them
with respect to the moduli $a_1, a_2 \ldots a_\nu$, are called the \gesperrt{indices}
of the element $A$, and the following theorem holds:

1. \gesperrt{One obtains the indices of a compositum $A A'$
by adding the corresponding indices of the two
factors.}

If we assign to each element $A$ in any manner a \gesperrt{numerical value},
then we may call this assignment a function
of $A$. Such a function $\chi(A)$ shall now be called a
\gesperrt{group character} if for every two elements
$A$, $A'$ of $S$ it satisfies the condition:
\weq{4}{\chi(AA') = \chi(A)\chi(A'),}
and consequently also the more general one
\[ \chi(AA'A''\ldots) = \chi(A)\chi(A')\chi(A'')\ldots \]
Two such functions $\chi$ and $\chi_1$ are regarded as different
if there is at least one element $A$ for which $\chi(A)$
is different from $\chi_1(A)$. We shall now determine these
characters more closely.

If we set $A' = 1$, then from (4) there results:
\weq{5}{\chi(A) = \chi(A)\chi(1),}
and if we exclude once and for all the case that all
$\chi(A) = 0$:
\weq{6}{\chi(1) = 1,}
hence the first theorem:

2. \gesperrt{Every character has the value 1 for the unit
element.}

If we further set $A' = A$, then from (4) it follows:
\[ \chi(A^2) = [\chi(A)]^2, \]
and from this by complete induction for every exponent $h$:
\weq{7}{\chi(A^h) = \chi(A)^h.}
If we therefore denote by $a$ the degree of the element $A$, then
it follows, when one sets $h = a$ in (7) and uses (6):
\weq{8}{\chi(A)^a = 1,}
hence:
%CONT>>


% ===== scan leaf 66, printed 45 =====
% !PRINTED 45
\weberpage{45}
%<<CONT

3. \gesperrt{The values of a character $\chi(A)$ are roots of unity whose degree is a divisor of the degree of $A$.}

After this we can easily determine all the characters. For if we set
\weq{9}{
\chi(A_1) = \omega_1,\ \chi(A_2) = \omega_2 \ldots \chi(A_\nu) = \omega_\nu,
}
then
\weq{10}{\omega_1^{a_1} = 1,\ \omega_2^{a_2} = 1, \ldots \omega_\nu^{a_\nu} = 1,}
and by (2) and (4)
\weq{11}{\chi(A) = \omega_1^{\alpha_1}\omega_2^{\alpha_2}\ldots\omega_\nu^{\alpha_\nu}.}

Conversely, if $\omega_1, \omega_2 \ldots \omega_\nu$ denotes any solution of the equations (10), then a function of $A$ is determined by (11) which, by 1., satisfies the condition (4) and is therefore a character of the group.

Each of the equations (10) has $a_1, a_2 \ldots a_\nu$ roots, and if we combine each with each, there result
\[
a_1 a_2 \ldots a_\nu = n
\]
combinations. All these combinations lead to distinct characters $\chi(A)$. For if we denote by $\omega_1', \omega_2' \ldots \omega_\nu'$ a second combination of roots of the equation (10), and if for all elements $A$
\[
\omega_1^{\alpha_1}\omega_2^{\alpha_2}\ldots\omega_\nu^{\alpha_\nu} = \omega_1'^{\alpha_1}\omega_2'^{\alpha_2}\ldots\omega_\nu'^{\alpha_\nu},
\]
then, if one sets $\alpha_1 = 1, \alpha_2 = 0 \ldots \alpha_\nu = 0$, it follows that $\omega_1 = \omega_1'$, and likewise $\omega_2 = \omega_2' \ldots \omega_\nu = \omega_\nu'$. From this follows the theorem:

4. \gesperrt{There are $n$ and no more distinct characters of an Abelian group of degree $n$, all of which are represented by the formula (11).}

If we understand by $\varepsilon_1, \varepsilon_2 \ldots \varepsilon_\nu$ a system of primitive roots of the equations (10), we can set
\weq{12}{\omega_1 = \varepsilon_1^{\beta_1},\ \omega_2 = \varepsilon_2^{\beta_2}, \ldots \omega_\nu = \varepsilon_\nu^{\beta_\nu}}
and can take the $\beta_1, \beta_2 \ldots \beta_\nu$, just like the $\alpha_1, \alpha_2 \ldots \alpha_\nu$, each from a complete system of residues with respect to the moduli $a_1, a_2 \ldots a_\nu$. Then, with $\varepsilon_1, \varepsilon_2 \ldots \varepsilon_\nu$ held fixed, we obtain all the $n$ group characters in the form
\weq{13}{\chi(A) = \varepsilon_1^{\alpha_1\beta_1}\varepsilon_2^{\alpha_2\beta_2}\ldots\varepsilon_\nu^{\alpha_\nu\beta_\nu}.}


% ===== scan leaf 67, printed 46 =====
% !PRINTED 46
\weberpage{46}

Among these characters there occurs also the one which is obtained by setting $\beta_1 = 0, \beta_2 = 0 \ldots \beta_\nu = 0$, and which has the value 1 for all elements $A$. This we shall call the \gesperrt{unit character}.

The $n$ characters of $S$ can now themselves in turn be combined into an \gesperrt{Abel}ian group, and indeed into a group isomorphic with $S$.

For by (13) each of the $n$ characters is determined by a system of numbers $\beta_1, \beta_2 \ldots \beta_\nu$ taken with respect to the modulus system $a_1, a_2 \ldots a_\nu$, and this system of numbers $\beta$ is at the same time the system of the indices of a definite element $B$ of $S$, namely of
\weq{14}{B = A_1^{\beta_1}A_2^{\beta_2}\ldots A_\nu^{\beta_\nu},}
so that to each element $B$ of $S$ there corresponds a definite character, which we denote by $\chi_B(A)$ or, omitting $A$, by $\chi_B$. This assignment of the characters $\chi$ to the elements $B$ will change, however, if another basis is taken as the foundation, or if the roots of unity $\varepsilon_1, \varepsilon_2 \ldots \varepsilon_\nu$ are chosen differently.

If $B'$ is a second element of $S$ with the indices $\beta_1', \beta_2' \ldots \beta_\nu'$, we obtain in the same way the character $\chi_{B'}$, and if we now understand by $\chi_{BB'}$ the character which for every $A$ is determined by the formula
\weq{15}{\chi_{BB'}(A) = \varepsilon_1^{\alpha_1(\beta_1+\beta_1')}\varepsilon_2^{\alpha_2(\beta_2+\beta_2')}\ldots\varepsilon_\nu^{\alpha_\nu(\beta_\nu+\beta_\nu')}}
then the characters $\chi$ are thereby joined into a group isomorphic with $S$. The unit element in this \gesperrt{group of characters} is the unit character.

Accordingly the group of the characters can also be represented by a basis, which one obtains by setting
\weq{16}{\chi_1(A) = \varepsilon_1^{\alpha_1},\ \chi_2(A) = \varepsilon_2^{\alpha_2} \ldots \chi_\nu(A) = \varepsilon_\nu^{\alpha_\nu}}
Then every character is contained in the form:
\weq{17}{\chi_B = \chi_1^{\beta_1}\chi_2^{\beta_2}\ldots\chi_\nu^{\beta_\nu},}
and, if $B, B'$ are two elements in $S$, we have:
\weq{18}{\chi_B\chi_{B'} = \chi_{BB'}.}

If $\chi_1, \chi_2$ are any two of the characters and $\chi_1\chi_2$ the one composed from the two, then by (15), for every element $A$,
\weq{19}{\chi_1(A)\chi_2(A) = \chi_1\chi_2(A).}


% ===== scan leaf 68, printed 47 =====
% !PRINTED 47
\weberpage{47}

We shall express the result obtained also as a theorem:

5. \gesperrt{The characters of a group $S$ can be combined into a group isomorphic with $S$.}

Let us also mention the theorem which follows immediately from the definition of $\chi_B$ \gesperrt{through} the formula (13):
\weq{20}{\chi_B(A) = \chi_A(B).}

Finally the following theorem also holds:

6. \gesperrt{If $A$ runs through all the elements of the group $S$, then for a fixed character $\chi$:}
\weq{21}{\overset{A}{\Sigma}\,\chi(A) = n \text{ or } = 0,}
\gesperrt{according as $\chi$ is the unit character or not. Likewise, if the element $A$ is held fixed and $\chi$ runs through the sequence of characters:}
\weq{22}{\overset{\chi}{\Sigma}\,\chi(A) = n \text{ or } = 0,}
\gesperrt{according as $A$ is the unit element or not.}

For the case that $\chi$ or $A$ are the unit elements of their groups, the formulae (21) and (22) are evident, since then each of the $n$ terms of the sum has the value 1. But if $\chi$ is not the unit character, then there is an element $B$ in $S$ such that $\chi(B)$ is not $= 1$. If we then set
\[
\overset{A}{\Sigma}\,\chi(A) = \sigma,
\]
it follows by multiplication with $\chi(B)$ that:
\[
\overset{A}{\Sigma}\,\chi(AB) = \sigma\chi(B).
\]
But since $AB$ runs through the whole group $S$ together with $A$, we have also $\overset{A}{\Sigma}\,\chi(AB) = \sigma$, hence $\sigma[1 - \chi(B)] = 0$ or $\sigma = 0$, q.\,e.\,d.

In the same way one proves the formula (22), which moreover also follows immediately from (21) by (20).


% ===== scan leaf 69, printed 48 =====
% !PRINTED 48
\weberpage{48}
\pdfbookmark[2]{§ 12. Divisors of an Abelian Group. Reciprocal Groups}{bm15}

\begin{center}
\S.\,12.

\gesperrt{Divisors of an Abelian Group}.

\gesperrt{Reciprocal Groups}.
\end{center}

A divisor $T$ of an \gesperrt{Abel}ian group $S$ is, as already remarked above, always a normal divisor, and hence by \S.\,4 there is a group complementary to $T$,
\[
S/T,
\]
whose elements are the subsidiary groups of $T$, namely
\weq{1}{T,\ TA',\ TA''\ldots,}
where $A', A'' \ldots$ denote certain elements of $S$. The number of the elements (1), that is, the degree of the group $S/T$, is, if $t$ is the degree of $T$, equal to the index of the divisor $T$ of $S$, hence equal to
\[
\frac{n}{t} = j.
\]

But this group $S/T$ is itself again an \gesperrt{Abel}ian one; for we have
\[
TA'\,TA'' = TA'A'',\ \ TA''\,TA' = TA''A',
\]
so that the two are equal to each other.

We have now to determine the characters of the group $S/T$. We denote the elements of this group by
\weq{2}{T,\ T_1,\ T_2 \ldots T_{j-1},}
and one of their characters by $\xi(T_i)$.

From this function $\xi(T)$ we can now derive a function $\xi(A)$ of the elements of $S$ by setting
\weq{3}{\xi(A_i) = \xi(T_i)}
whenever $A_i$ is any element occurring in the subsidiary group $T_i$. For the elements of the group $T$ itself we then have $\xi(A) = 1$.

But this function $\xi(A_i)$ is contained among the characters of $S$. For if $A_i$ occurs in $T_i$ and $A_k$ in $T_k$, then $A_iA_k$ occurs in $T_iT_k$, and consequently
\weq{4}{\xi(A_i)\xi(A_k) = \xi(T_i)\xi(T_k) = \xi(T_iT_k) = \xi(A_iA_k).}
But this, by \S.\,11,(4), is the definition of a character of $S$.


% ===== scan leaf 70, printed 49 =====
% !PRINTED 49
\weberpage{49}

If conversely one of the characters $\chi = \xi$ of the group $S$
has the same value for all elements of the group $T$, then this
can only be the value 1, since the unit element of $S$ certainly
occurs in $T$ (\S.\,11, 2.). If then $A_i$ is any element
of $T_i$, and $A$ runs through the whole group $T$, then
$AA_i$ runs through the subsidiary group $T_i$. But then
$\xi(A) = 1$ and consequently
\weq{5}{\xi(AA_i) = \xi(A_i),}
i.\,e. $\xi(A)$ has one and the same value for all elements of a
subsidiary group $T_i$, and $\xi(A)$ can therefore be regarded as a
function of $T_i$ and denoted by $\xi(T_i)$.

Since now, if $A_i$ occurs in $T_i$ and $A_k$ in $T_k$, $A_iA_k$
is contained in $T_iT_k$, it follows that
\weq{6}{\xi(T_i)\xi(T_k) = \xi(T_iT_k),}
i.\,e. $\xi(T_i)$ is contained among the characters of the group
$S/T$. There are therefore exactly $j$ such characters $\xi(A)$,
which are defined by the property that they have the value 1 for
every element in $T$. These $j$ characters $\xi$ form a group under
the composition of the characters; for if $\xi_1(A) = 1$,
$\xi_2(A) = 1$, then also $\xi_1\xi_2(A) = 1$ [\S.\,11, (18)].
Since by (4) the functions $\xi$ can also be regarded as the
characters of the group $S/T$, by \S.\,11, 5. the group of the
$\xi$ is isomorphic with the group $S/T$.

With this, then, the following theorem is proved:

7.\ \gesperrt{If an Abelian group $S$ has a divisor $T$ of
degree $t$ and of index $j$, then among the characters of
$S$ there are exactly $j$ and no more that have the value 1
for all elements of $T$, while for every element of $S$ not
contained in $T$ at least one of them is different from 1. These
characters form a group isomorphic with $S/T$.}

We shall call these characters \gesperrt{belonging to the group $T$}.

Since every character $\chi_B$ corresponds to a definite element $B$ of $S$,
it follows also that, when $\chi_B$ runs through the group of the
characters belonging to $T$, $B$, by the formula \S.\,11, (18),
will run through a group which is of degree $j$ and isomorphic with the
group of the $\chi_B$ and hence also with the group $S/T$. This
%CONT>>


% ===== scan leaf 71, printed 50 =====
% !PRINTED 50
\weberpage{50}
%<<CONT
group of the $B$, which is thus also a divisor of $S$, we shall
denote by $U$ and call the \gesperrt{group reciprocal to $T$}. Here too,
however, it is to be remarked that the concept of the reciprocal
group depends in general on the choice of the basis and of the
roots of unity $\varepsilon$, and thus does not belong to the
groups $S$ and $T$ as such.

The group reciprocal to $T'$ is characterised by the following
theorem:
% [sic: Weber prints T' here; the newly defined group is reciprocal to T.]

8.\ \gesperrt{If one lets $A$ run through the elements of a
divisor $T$ of $S$ and seeks all elements $B$ of $S$ which for
every $A$ satisfy the condition}
\weq{7}{\chi_B(A) = 1}
\gesperrt{then $B$ runs through the elements of the group $U$
reciprocal to $T$, whose degree is equal to the index of $T$.}

By the theorem of \S.\,11, (20), $T'$ is the reciprocal group to
$U$, and the relation of these two groups is therefore mutual.
% [sic: Weber prints T' here; reciprocity requires T.]

Of these reciprocal groups there holds further the theorem:

9.\ \gesperrt{If $T'$ is a divisor of $T$, then conversely the
reciprocal group $U$ of $T'$ is a divisor of the group $U'$
reciprocal to $T''$.}
% [sic: Weber prints T' and T'' here; the theorem requires T and T', respectively.]

For every element $A'$ of $T''$ is at the same time contained in
$T'$, and consequently, if $B$ runs through the group $U$,
% [sic: Weber prints T'' and T' here; the proof requires T' and T, respectively.]
$\chi_B(A') = 1$. Hence $B$ must occur in the group $U'$.

If from the totality of the characters $\chi$ one selects an
arbitrary number $\xi_1, \xi_2 \dots \xi_\mu$, no matter whether
they form a group or not, then all elements $A$ which satisfy
the $\mu$ conditions
\weq{8}{\xi_1(A) = 1,\ \xi_2(A) = 1, \dots \xi_\mu(A) = 1}
form a group, because from $\xi_i(A) = 1$, $\xi_i(A') = 1$ it
follows that $\xi_i(AA') = 1$ also. If the
$\xi_1, \xi_2 \dots \xi_\mu$ are not groups, then from the
equations (8) there follow so many further
$\xi_{\mu+1}(A) = 1 \dots$ that the
$\xi_1, \xi_2 \dots \xi_\mu$ are completed to a group.

10.\ \gesperrt{From Theorem 7. it follows that in this way all
possible divisors of $S$ can be obtained.}


% ===== scan leaf 72, printed 51 =====
% !PRINTED 51
\weberpage{51}

It is perhaps desirable to illustrate these theorems by a simple
example. Let the degree of the group $S$ be the square of a
prime $n = p^2$. Then $S$ has either only the one
invariant $p^2$, and is then cyclic:
\[
\text{a)}\quad S = 1, A, A^2 \dots A^{p^2-1},
\]
or $S$ has two invariants, which are both equal to $p$; then
$S$ consists of the elements:
\[
\text{b)}\quad S = A_1^{\alpha_1}A_2^{\alpha_2}, \quad \alpha_1,
\alpha_2 = 0, 1 \dots p-1.
\]
In case a) we obtain the $p^2$ characters
\[
\chi_B(A^\alpha) = \varepsilon^{\beta\alpha},
\]
if $\varepsilon$ is a primitive root of the equation
$\varepsilon^{pp} = 1$.
% [sic: the source prints \varepsilon^{pp}; the cyclic group order and surrounding p^2 formulas require \varepsilon^{p^2}.]
If, in order to form the divisors of $S$ according to Theorem 10.,
we take one of the characters $\chi_B$ in which
$\beta$ is not divisible by $p$, then $\chi_B(A^\alpha)$
can be $= 1$ only when $\alpha$ is divisible by $p^2$,
that is, $= 0$; i.\,e., we obtain only the divisor of $S$ consisting
of the unit element. But if we assume
$\beta = p\beta'$ to be divisible by $p$, but not by $p^2$, then
$\chi_{p\beta'}(A^\alpha)$ is $= 1$ then and only then when
$\alpha = p\alpha'$ is divisible by $p$. We thus obtain the
divisor
\[
T = 1, A^p, A^{2p} \dots A^{(p-1)p},
\]
and the divisor $U$ reciprocal to $T$ is here identical with $T$.

In case b) we must, in order to form the characters, assume two
$p$th roots of unity $\varepsilon^{\beta_1},
\varepsilon^{\beta_2}$, and obtain, if
\[
B = A_1^{\beta_1}A_2^{\beta_2}
\]
is set,
\[
\chi_B(A_1^{\alpha_1}A_2^{\alpha_2}) = \varepsilon^{\alpha_1\beta_1 + \alpha_2\beta_2}.
\]
If we set two of these characters $= 1$, that is
\[
\alpha_1\beta_1 + \alpha_2\beta_2 \equiv 0, \quad
\alpha_1\beta_1' + \alpha_2\beta_2' \equiv 0,
\]
then it follows, if the determinant $\beta_1\beta_2' - \beta_2\beta_1'$
is not divisible by $p$, that $\alpha_1$ and $\alpha_2$ must be
divisible by $p$. But if the determinant is divisible by $p$,
then the one of these two congruences is a consequence of the
other. In the first case, therefore, we obtain a group consisting
only of the unit element. We therefore obtain all divisors
$T_{\beta_1,\beta_2}$ distinct from 1 and $S$ when, for a fixed
$\beta_1, \beta_2$, we make the $\alpha_1, \alpha_2$ satisfy in all
possible ways the congruence
\weq{9}{\alpha_1\beta_1 + \alpha_2\beta_2 \equiv 0\ (\mathrm{mod}\ p)}
%CONT>>


% ===== scan leaf 73, printed 52 =====
% !PRINTED 52
\weberpage{52}
\pdfbookmark[2]{§ 13. The Genera in an Abelian Group}{bm16}
%<<CONT
accordingly. If $\alpha_1, \alpha_2$ is a solution of this
congruence in which the two numbers are not both zero, then we obtain
all the solutions if we let the factor $h$ in $h\alpha_1, h\alpha_2$
run through the sequence of numbers $0, 1 \dots p-1$, and the
group $T$ therefore becomes, if $\alpha_1, \alpha_2$ is a fixed
pair of values satisfying the condition (9),
\[
(A_1^{\alpha_1}A_2^{\alpha_2})^h, \quad h = 0, 1, \dots p-1.
\]

The reciprocal group $U$ is obtained by seeking all the values of
the $\beta$ that satisfy the condition (9), and one thus finds
the group $U$ in the form
\[
(A_1^{\beta_1}A_2^{\beta_2})^h, \quad h = 0, 1, \dots p-1,
\]
in which $\alpha_1, \alpha_2, \beta_1, \beta_2$ are now four fixed
numbers satisfying the condition (9). If one sets $\beta_1 = 0$,
$\beta_2 = 1$, $\alpha_1 = 1$, $\alpha_2 = 0$, one obtains the
special case of the two reciprocal groups
\[
A_1^h, \quad A_2^h, \quad h = 0, 1, \dots p-1.
\]

\begin{center}
\S.\,13.

\gesperrt{The Genera in an Abelian Group.}
\end{center}

An element of an \gesperrt{Abel}ian group $S$ that is identical with its
opposite, whose second power is therefore equal to the unit
element, is called a \gesperrt{two-sided
element}\footnote{Cf. Bd.~I, S.~390, Anmerkung.}. The unit
element therefore always belongs to the two-sided ones.

If we represent the elements of $S$ by a basis
$A_1, A_2 \dots A_\nu$, whose elements have the degrees
$a_1, a_2 \dots a_\nu$, then
\weq{1}{A = A_1^{\alpha_1}A_2^{\alpha_2} \dots A_\nu^{\alpha_\nu}}
will be a two-sided element if and only if
\weq{2}{2\alpha_1 \equiv 0\ (\mathrm{mod}\ a_1), \ 2\alpha_2 \equiv 0\ (\mathrm{mod}\ a_2), \dots 2\alpha_\nu \equiv 0\ (\mathrm{mod}\ a_\nu).}
Now if among the invariants of the group a power of 2 occurs
$\lambda$ times, then we can assume the basis of $S$ so
ordered that $a_1, a_2 \dots a_\lambda$ are even and
$a_{\lambda+1} \dots a_\nu$ odd numbers. Then the solutions
of (2) come out as:
%CONT>>


% ===== scan leaf 74, printed 53 =====
% !PRINTED 53
\weberpage{53}
%<<CONT

\[
\alpha_1 = \frac{\eta_1 a_1}{2},\ \alpha_2 = \frac{\eta_2 a_2}{2} \cdots \alpha_\lambda = \frac{\eta_\lambda a_\lambda}{2},\ \alpha_{\lambda+1}=0 \ldots \alpha_\nu = 0,
\]

in which $\eta_1, \eta_2 \ldots \eta_\lambda$ can be equal to 0 or equal to 1, and one obtains all two-sided elements in the form

\[
A_1^{\frac{\eta_1 a_1}{2}} A_2^{\frac{\eta_2 a_2}{2}} \ldots A_\lambda^{\frac{\eta_\lambda a_\lambda}{2}},
\]

and their number, since all combinations of 0 and 1 are admissible for the exponents $\eta$, is $= 2^\lambda$.

If $\chi$ is an arbitrary character of $S$, then, if $A$ is a two-sided element, $\chi(A) = \pm 1$, since

\[
\chi(A)^2 = \chi(1) = 1
\]

holds.

Similarly, we call a character two-sided if it is a two-sided element in the group of characters. Since the group of characters is isomorphic with the group $S$, there are just as many two-sided characters as two-sided elements, namely $2^\lambda$. By \S.\,11, (17) they are represented by:

\[
\chi_1^{\frac{\eta_1 a_1}{2}} \chi_2^{\frac{\eta_2 a_2}{2}} \ldots \chi_\lambda^{\frac{\eta_\lambda a_\lambda}{2}}.
\]

\gesperrt{The two-sided characters have the value $\pm 1$ for every element.}

The two-sided elements form by themselves a group $T$ of degree $2^\lambda$. Likewise the two-sided characters form a group isomorphic with it, and one obtains the group $U$ reciprocal to $T$ by seeking out all elements $A$ for which all two-sided characters have the value $+1$. This group, which by \S.\,12, 7. is a divisor of $S$ of index $2^\lambda$, depends, by the last definition, neither on the choice of the basis nor on the roots of unity representing the characters, and is completely determined by the nature of the group $S$. If we denote it by $G$, then the system of subsidiary groups

\weq{3}{G,\ G_1,\ G_2 \ldots G_{2^\lambda - 1}}

is characterised by the fact that for all elements of one of these systems the two-sided characters yield one and the same system of values. The systems $G, G_1, G_2 \ldots$ are called the \gesperrt{genera} (Genera) contained in $S$. The group $G$ in particular is called the \gesperrt{principal genus}.


% ===== scan leaf 75, printed 54 =====
% !PRINTED 54
\weberpage{54}
\pdfbookmark[2]{§ 14. Indices with Respect to an Odd Prime Power as Modulus}{bm17}

\gesperrt{The number of genera is thus as great as the number of two-sided elements}\footnote{\gesperrt{Gauss} first introduced these concepts into the theory of quadratic forms, and used the expression ``Genera''. Disqu. arithm. art. 228 f.}.

\begin{center}
\S.\,14.
\end{center}

\begin{center}
\gesperrt{Indices for an odd prime power}\\
\gesperrt{as modulus.}
\end{center}

The most important example of a commutative group is offered by the natural numbers, when they are combined with one another by ordinary multiplication.

In order to derive from them a finite group, we take an arbitrary positive integer $m$ as modulus and decompose $m$ into its prime factors

\weq{1}{m = 2^\lambda q_1^{\varkappa_1} q_2^{\varkappa_2} \ldots,}

in which $q_1, q_2 \ldots$ denote distinct odd primes, $\lambda, \varkappa_1, \varkappa_2 \ldots$ positive exponents, $\lambda$ possibly also zero, namely when $m$ is odd.

Now all numbers, positive as well as negative, that are congruent to one another with respect to the modulus $m$ are united into a \gesperrt{number class}, and each of these number classes is represented by a representative, say by the residue of the division, hence by one of the numbers $0, 1, 2 \ldots m-1$.

If $a$ and $a'$ are two numbers of such a class, and $b$ and $b'$ two numbers of a second class, then the products $ab$ and $a'b'$ also belong to the same class. For if $a \equiv a'$, $b \equiv b'$, then also $ab \equiv a'b'\ (\mathrm{mod}\ m)$. By multiplication, therefore, not only the numbers but also the classes are compounded.

Nevertheless these number classes in their totality still form no group; for from a congruence

\[
ab \equiv ac\ (\mathrm{mod}\ m)
\]

it follows necessarily that $b \equiv c$ only when $a$ is relatively prime to $m$. Thus the requirement \S.\,1, 3. is here in general not fulfilled.

The greatest common divisor which a number $a$ has with $m$ is, for the entire class represented by $a$, the
%CONT>>


% ===== scan leaf 76, printed 55 =====
% !PRINTED 55
\weberpage{55}
%<<CONT
same, and may be called the divisor of the class. If $d, d'$ are the divisors of two classes $a, a'$, then the greatest common divisor of $m$ and $dd'$ is the divisor of the class $aa'$. From this it follows that the number classes whose numbers are relatively prime to $m$ always yield such number classes again under composition, and for these the requirement of \S.\,1, 3. is then also fulfilled.

1.\ \gesperrt{The number classes whose individuals are relatively prime to the modulus thus form, under composition by multiplication, an Abelian group.}

This group is the object of our considerations, in which our chief goal shall be the determination of a basis.

We shall denote by $N$ every number class that contains only numbers coprime to $m$, and the group of the number classes $N$, whose existence we have now demonstrated, by $\mathfrak{N}$. By $n$ we shall denote every number coprime to $m$.

The degree of this group is as great as the number of the relative primes to $m$ that are at the same time positive and not greater than $m$, and this number we have determined in \S.\,132 of the first volume. It is

\weq{2}{\mu = \varphi(m) = 2^{\lambda-1} q_1^{\varkappa_1 - 1}(q_1 - 1)\, q_2^{\varkappa_2 - 1}(q_2 - 1) \ldots}

or

\[
= q_1^{\varkappa_1 - 1}(q_1 - 1)\, q_2^{\varkappa_2 - 1}(q_2 - 1) \ldots,
\]

when $\lambda = 0$.

The degree of an element $N$ of this group is the least positive exponent $e$ to which one must raise a representative $n$ of $N$ in order that $n^e$ become congruent to unity with respect to the modulus $m$. Since every $e$ must be a divisor of the degree of the group $\varphi(m)$, there follows the \gesperrt{generalised Fermat's theorem}:

\weq{3}{n^{\varphi(m)} \equiv 1\ (\mathrm{mod}\ m),}

a formula which holds for every number $n$ coprime to $m$.

If now $q$ is one of the odd prime factors of $m$, and $q^\varkappa$ the highest power of $q$ dividing $m$, then we take a primitive root $g$ of $q$, which, if $\varkappa$ is greater than 1, we choose so that $g^q - g$ is not divisible by $q^2$ (Bd. I, \S.\,184). For the sake of brevity we shall call a number $g$ satisfying this last condition a \gesperrt{primitive root of $q^2$}.


% ===== scan leaf 77, printed 56 =====
% !PRINTED 56
\weberpage{56}

If now $\varkappa = 1$, then by the definition of $g$ the numbers

\[
1,\ g,\ g^2 \ldots g^{q-2}
\]

are all different from one another with respect to the modulus $q$, while $g^{q-1}$ is again congruent to 1.

But if $\varkappa > 1$, then:

\weq{4}{g^{q-1} = 1 + hq,}

and by our assumption about $g$, $h$ is not divisible by $q$. If we raise the equation (4) to the power $q$, and apply the binomial theorem on the right, then it follows that

\[
g^{q(q-1)} = 1 + hq^2 + h^2\, \frac{q^3(q-1)}{2} + \cdots,
\]

hence

\weq{5}{g^{q(q-1)} = 1 + h_1 q^2,}

wherein

\[
h_1 = h + h^2\, \frac{q(q-1)}{2} + \cdots
\]

is not divisible by $q$. (This would no longer be correct for $q = 2$, and for that reason the prime 2 requires a different treatment.)

If we raise (5) once more to the $q$th power, it follows that

\[
g^{q^2(q-1)} = 1 + h_2 q^3,
\]

and we can continue in this way, and obtain for every arbitrary positive exponent $\lambda$

\weq{6}{g^{q^{\lambda-1}(q-1)} = 1 + hq^\lambda,}

wherein $h$ is an integer not divisible by $q$. If we set for brevity:

\weq{7}{q^{\varkappa-1}(q-1) = \varphi(q^\varkappa) = c,}

then we have

\weq{8}{g^c \equiv 1\ (\mathrm{mod}\ q^\varkappa),}

and it remains to be shown that $c$ is the least positive exponent for which the congruence (8) is satisfied. Let us assume that $e$ is this least exponent, so that

\weq{9}{g^e \equiv 1\ (\mathrm{mod}\ q^\varkappa),}

then $e$ must be a divisor of $c$, because otherwise the congruence (8) would also be satisfied if $c$ were replaced by the residue of the division of $c$ by $e$, which is smaller than $e$.

On the other hand, $e$ must be divisible by $q-1$, because the congruence $g^e \equiv 1\ (\mathrm{mod}\ q)$ contained in (9) holds only for such exponents as are divisible by $q-1$.


% ===== scan leaf 78, printed 57 =====
% !PRINTED 57
\weberpage{57}

Thus $e$ is of the form $q^{\lambda-1}(q-1)$ with a positive $\lambda$. But that $\lambda$ cannot be smaller than $\varkappa$ follows from (6), from which it is to be seen that $q^\lambda$ is the highest power of $q$ that divides the difference

\[
g^{q^{\lambda-1}(q-1)} - 1
\]

Thus $c$ is the least positive number for which the congruence (8) is satisfied, and this is equivalent to the statement that the numbers

\weq{10}{1, g, g^2 \ldots g^{c-1}}

are all incongruent with respect to the modulus $q^\varkappa$; $g$ can therefore also be designated as a \gesperrt{primitive root} of $q^\varkappa$.

The number of terms of this sequence (10) is equal to $c$. But just as great is also the number of numbers $n$ incongruent with respect to the modulus $q^\varkappa$ and not divisible by $q$, and with this is proved:

\noindent\textbf{2.} \gesperrt{If} $q$ \gesperrt{is an odd prime,} $n$ \gesperrt{any number not divisible by} $q$, $g$ \gesperrt{a primitive root of} $q^2$ \gesperrt{and} $\varkappa$ \gesperrt{a positive exponent, then one and only one number} $\gamma$ \gesperrt{can be determined with respect to the modulus} $c$ \gesperrt{by which the congruence}

\[
n \equiv g^\gamma \pmod{q^\varkappa}
\]

\gesperrt{is satisfied.}

The number $c$ is always divisible by 2, and we therefore further emphasise the particular theorem that

\weq{11}{g^{\frac12 c} \equiv -1 \pmod{q^\varkappa}}

holds. For in the product, divisible by $q^\varkappa$,

\[
g^c - 1 = (g^{\frac12 c} - 1)(g^{\frac12 c} + 1)
\]

the first factor is not divisible by $q^\varkappa$, and consequently the second factor must be divisible by $q$. But then it follows that the first factor also cannot be divisible by $q$, since otherwise the sum of the two factors, and hence $g$ itself, would also have to be divisible by $q$. Consequently the second factor must be divisible by $q^\varkappa$.

Here, then, there is complete analogy with the theorems on primitive roots and index systems which we came to know in \S.\,136 of the first volume, and one can therefore also here designate $\gamma$ as the index of $n$ for the modulus $q^\varkappa$.


% ===== scan leaf 79, printed 58 =====
% !PRINTED 58
\weberpage{58}
\pdfbookmark[2]{§ 15. Indices for a Power of 2 as Modulus}{bm18}

\begin{center}
\S.\,15.

\gesperrt{Indices for a Power of 2 as Modulus}.
\end{center}

A similar theorem must now be established for the powers $2^\lambda$ of the prime 2, which, as already remarked above, behaves differently.

If $\lambda = 1$, then every number not divisible by 2 is $n \equiv 1 \pmod{2}$. If $\lambda = 2$, then $-1$ is to be regarded as a primitive root of 4, for every odd number $n$ satisfies one of the congruences

\[
n \equiv (-1)^\alpha \pmod{4},
\]

in which $\alpha = 0$ or $= 1$.

But already for the modulus 8 no primitive root exists any longer, that is, no number by whose powers all classes of odd numbers with respect to the modulus 8 can be represented. For if $g$ is any odd number, then among the powers of $g$ at most the two $1, g$ are distinct with respect to the modulus 8, because always $g^2 \equiv 1 \pmod{8}$. If therefore one takes $g$ not congruent to 1 and not congruent to $-1 \pmod{8}$, thus $g = 3$ or $= 5$, then every odd number is congruent to one of the four numbers

\[
(-1)^\alpha g^\beta \qquad \alpha = 0, 1, \quad \beta = 0, 1
\]

with respect to the modulus 8.

The number of classes of odd numbers with respect to the modulus $2^\lambda$ is $2^{\lambda-1}$. Now, however, for every odd number $g$, provided $\lambda > 2$,

\weq{1}{g^{2^{\lambda-2}} \equiv 1 \pmod{2^\lambda}.}

For if we assume (1) to be correct and accordingly set

\[
g^{2^{\lambda-2}} = 1 + h\, 2^\lambda,
\]

and raise to the square, it follows that

\[
g^{2^{\lambda-1}} \equiv 1 \pmod{2^{\lambda+1}}.
\]

If therefore the formula (1) is correct for $\lambda$, then it is so also for $\lambda + 1$, and since it holds for $\lambda = 3$, it holds generally. From this it follows that among the powers of an odd number $g$ at most $2^{\lambda-2}$ distinct ones with respect to the modulus $2^\lambda$ can occur.

On the other hand, however, it follows easily for $g = 5$ that $5^{2^{\lambda-2}}$ is the lowest power that becomes congruent to unity with respect to the modulus $2^\lambda$. For if $5^e$ is the lowest power that is congruent to
%CONT>>


% ===== scan leaf 80, printed 59 =====
% !PRINTED 59
\weberpage{59}

%<<CONT
unity, then $e$ is by (1) a divisor of $2^{\lambda-2}$, hence a power of 2, and if $e < 2^{\lambda-2}$, then we should have to have

\[
5^{2^{\lambda-3}} \equiv 1 \pmod{2^\lambda}
\]

But this is not possible, for, for every $\lambda$ greater than 2, the formula

\[
5^{2^{\lambda-3}} = 1 + 2^{\lambda-1} h\ ^{1)}
\]

holds with odd $h$ --- a formula which, just like the formula (1), can be proved by complete induction. Hence, if $\lambda \geqq 3$,
the $2^{\lambda-2}$ powers

\weq{2}{1, 5, 5^2 \ldots 5^{2^{\lambda-2}-1}}

are all different from one another with respect to the modulus $2^\lambda$. Now a relation of the form $5^h \equiv -5^\varkappa$ is impossible for the modulus 4, and therefore all the more for every higher power of 2, and consequently the quantities

\weq{3}{-1, -5, -5^2 \ldots -5^{2^{\lambda-2}-1}}

are, with respect to the modulus $2^\lambda$, different both among themselves and from the quantities (2); and since their total number amounts to $2^{\lambda-1}$, every odd number is congruent, with respect to the modulus $2^\lambda$, to one and only one of the quantities (2), (3).

If, then, we set

\weq{4}{a = 2, \quad b = 2^{\lambda-2},}

we obtain the following theorem for $\lambda \geqq 3$:

\noindent\textbf{3.} \gesperrt{For every odd number} $n$ \gesperrt{one can assign, completely determined with respect to the pair of moduli} $a, b$\gesperrt{, a pair of numbers} $\alpha, \beta$\gesperrt{, so that}

\[
n \equiv (-1)^\alpha 5^\beta \pmod{2^\lambda}
\]

\gesperrt{holds.}

The case $\lambda = 2$ can be subsumed under this, because then $b = 1$ and $\beta = 0$ can be set. In order to include the case $\lambda = 1$ under it as well, which however offers no particular interest, one would have to set $a = 1, b = 1$.

\noindent\rule{4cm}{0.4pt}

\noindent $^{1)}$ If one wished to take the number 3 as base in place of the number 5, this formula would not yet be valid for $\lambda = 3$, though it would be for every larger $\lambda$, and hence 3 could serve as base just as well as 5. The reason for preferring the base 5 lies in the fact that all powers of this base are $\equiv 1 \pmod{4}$.


% ===== scan leaf 81, printed 60 =====
% !PRINTED 60
\weberpage{60}
\pdfbookmark[2]{§ 16. The Groups of Number Classes with Respect to a Composite Modulus}{bm19}

\begin{center}
\S.\,16.

\gesperrt{The Groups of Number Classes with Respect to a Composite Modulus}.
\end{center}

From the combination of the two theorems 2. and 3. of the two preceding sections there now results the following result, by which the problem that we posed at the beginning of \S.\,14 is completely solved:

\noindent\textbf{4.} \gesperrt{If the modulus}

\[
m = 2^\lambda q_1^{\varkappa_1} q_2^{\varkappa_2} \ldots
\]

\gesperrt{is, and} $g_1, g_2 \ldots$ \gesperrt{are primitive roots of the squares of the primes} $q_1, q_2 \ldots$\gesperrt{, and if further}

\[
a = 2, \quad b = \tfrac12 \varphi(2^\lambda), \quad c_1 = \varphi(q_1^{\varkappa_1}), \quad c_2 = \varphi(q_2^{\varkappa_2}) \ldots
\]

\gesperrt{is set, then one can, for every number} $n$ \gesperrt{that is relatively prime to} $m$\gesperrt{, determine uniquely a system of numbers} $\alpha, \beta, \gamma_1, \gamma_2 \ldots$ \gesperrt{with respect to the moduli} $a, b, c_1, c_2 \ldots$\gesperrt{, such that the equations}

\weq{1}{
\begin{aligned}
n &\equiv (-1)^\alpha 5^\beta \pmod{2^\lambda} \\
&\equiv g_1^{\gamma_1} \pmod{q_1^{\varkappa_1}} \\
&\equiv g_2^{\gamma_2} \pmod{q_2^{\varkappa_2}} \\
&\; \cdots\cdots\cdots\cdots\cdots\cdots
\end{aligned}
}

\gesperrt{are satisfied.}

\gesperrt{The numbers $\alpha, \beta, \gamma_1, \gamma_2 \ldots$ are called the system of indices of $n$ for the modulus $m$.}

Here it is at first assumed that $\lambda \geqq 3$. But the theorem holds also for the remaining values of $\lambda$, if

\[
\text{for } \lambda = 0, 1, \quad a = 1, \quad b = 1
\]
\[
\text{\textquotedbl{}} \quad \lambda = 2, \qquad\quad a = 2, \quad b = 1
\]

is set. For $\lambda = 0, 1$ one is to set the indices $\alpha$ and $\beta = 0$, or else omit them entirely; for $\lambda = 2$, $\beta$ drops out, while $\alpha$ can be equal to 0 or equal to 1.

Now, in order to represent the group $\mathfrak{N}$ of the number classes $N$ with respect to the modulus $m$ by a basis, one determines the number classes

\weq{2}{A, B, C_1, C_2 \ldots,}

or at least representatives of these classes, which is always possible, from the congruences
%CONT>>


% ===== scan leaf 82, printed 61 =====
% !PRINTED 61
\weberpage{61}
%<<CONT
\[
\begin{aligned}
A &\equiv -1 \pmod{2^{\lambda}} &&\equiv 1 \pmod{q_1^{\varkappa_1}} &&\equiv 1 \pmod{q_2^{\varkappa_2}} \ldots\\
B &\equiv \phantom{-}5\;'' &&\equiv 1\;'' &&\equiv 1\;'' \ldots\\
C_1 &\equiv \phantom{-}1\;'' &&\equiv g_1\;'' &&\equiv 1\;'' \ldots\\
C_2 &\equiv \phantom{-}1\;'' &&\equiv 1\;'' &&\equiv g_2\;'' \ldots
\end{aligned}
\]

\noindent
$\ldots\ldots\ldots\ldots\ldots\ldots\ldots\ldots\ldots\ldots\ldots\ldots\ldots\ldots\ldots\ldots\ldots\ldots\ldots\ldots\ldots\ldots$

and by 4.\ one then obtains, for every number $n$ of our group,

\weq{3}{n \equiv A^{\alpha} B^{\beta} C_1^{\gamma_1} C_2^{\gamma_2} \ldots \pmod{m}.}

The $A, B, C_1, C_2 \ldots$ are thus the elements of a basis of the
group $\mathfrak{N}$, of degrees $a, b, c_1, c_2 \ldots$

If $m$ is odd or divisible only by the first power of 2,
the two elements $A, B$ drop out of the basis. If $m$ is divisible
by 4, but not by 8, then $B$ drops out and the element $A$ of the
second degree remains.

\gesperrt{5. If one decomposes the numbers $a, b, c_1, c_2 \ldots$ into powers
of primes different from one another, one obtains the invariants of the group $\mathfrak{N}$.}

In order to comprehend the various cases together, we denote
the elements of the basis (2) by

\weq{4}{C_{-1}, C_0, C_1 \ldots C_\mu.}

Here $C_{-1}, C_0$ are to stand for $A$ and $B$, and are accordingly
to be set equal to 1 when $\lambda = 0$ or $\lambda = 1$. When $\lambda = 2$,
we are to set $B = C_{-1} = 1$, $A = C_0$, and when $\lambda > 1$, $A = C_{-1}$,
$B = C_0$.
% [sic: Weber prints \lambda > 1; the preceding cases require \lambda > 2 here.]

Let the degrees of these elements be denoted by

\weq{5}{c_{-1}, c_0, c_1 \ldots c_\mu}

and the indices of a number from $\mathfrak{N}$, which are to be taken with
the quantities (5) as moduli, by

\weq{6}{\nu_{-1}, \nu_0, \nu_1 \ldots \nu_\mu.}

If $\lambda = 0$ or 1, then $\nu_{-1}$ and $\nu_0$ have only the value 0;
if $\lambda = 2$, then $\nu_{-1}$ has the value 0 and $\nu_0$ the value 0 or 1.
If $\lambda \geqq 3$, then $\nu_{-1}$ has one of the two values 0, 1, and $\nu_0$
one of the values $0, 1, 2, \ldots c_0 - 1$; $\mu$ is always equal to the
number of odd primes, different from one another,
that divide $m$. Leaving $\nu_{-1}, c_{-1}$ aside, we shall
call $\nu_0, \nu_1 \ldots \nu_\mu$ the \gesperrt{indices} of $n$ corresponding to the primes $2, q_1 \ldots q_\mu$,
and $c_0, c_1 \ldots c_\mu$ the \gesperrt{index moduli} corresponding
to the same primes. These index moduli are
\[
c_0 = \varphi(2^{\lambda-1}), \quad c_1 = \varphi(q_1^{\varkappa_1}), \ldots c_\mu = \varphi(q_\mu^{\varkappa_\mu}),
\]
%CONT>>


% ===== scan leaf 83, printed 62 =====
% !PRINTED 62
\weberpage{62}
%<<CONT
and only when $\lambda = 2$ is $c_0 = 2$, and not $= 1$, to be set.

If we then denote by $\varepsilon_{-1}, \varepsilon_0, \varepsilon_1 \ldots \varepsilon_\mu$ a system of primitive
roots of unity of the degrees $c_{-1}, c_0, c_1 \ldots c_\mu$, and by
$\beta_{-1}, \beta_0, \beta_1 \ldots \beta_\mu$ the indices of a number $b$, then we obtain the
characters of the group $\mathfrak{N}$ in the form

\weq{7}{\chi_b(n) = \varepsilon_{-1}^{\beta_{-1}\nu_{-1}} \varepsilon_0^{\beta_0\nu_0} \varepsilon_1^{\beta_1\nu_1} \ldots \varepsilon_\mu^{\beta_\mu\nu_\mu}.}

Here $\varepsilon_{-1} = 1$ when $\lambda = 0, 1, 2$; in the other
cases $\varepsilon_{-1} = -1$, and $\varepsilon_1, \ldots \varepsilon_\mu$ are primitive roots of unity
 of the degrees
\[
c_1 = \varphi(q_1^{\varkappa_1}), \ldots c_\mu = \varphi(q_\mu^{\varkappa_\mu}).
\]

\begin{center}
\rule{3cm}{0.4pt}
\end{center}


% ===== scan leaf 84, printed 63 =====
\weberpage{63}% page number not printed on this page
\pdfbookmark[1]{Third Section. The Group of the Cyclotomic Fields.}{bm20}
\pdfbookmark[2]{§ 17. The Resolvents of the Theory of Cyclotomy}{bm21}
% !PRINTED none
\begin{center}
\gesperrt{Third Chapter.}

\gesperrt{The Group of the Cyclotomic Fields.}

\rule{3cm}{0.4pt}

\S.\,17.

\gesperrt{The Resolvents of the Theory of the Division of the Circle.}
\end{center}

Of the theorems on \gesperrt{Abelian} groups we make an
application to the theory of the division of the circle for the case in which the
degree of the roots of unity is not a prime, but a
higher power of a prime. If $q$ is an odd prime
and $m = q^{\varkappa}$, $\varkappa > 1$, we take a primitive congruence root
 $g$ of $m$ and set, for every number $n$ not divisible
 by $q$,
\weq{1}{n \equiv g^{\nu} \pmod{m}.}
Then $\nu$ is the \gesperrt{index} of $n$. If $n$ runs through the group $\mathfrak{N}$
of the number classes not divisible by $q$ with respect to the modulus $m$,
of degree
\weq{2}{c = \varphi(q^{\varkappa}) = q^{\varkappa - 1}(q - 1),}
then $\nu$ runs through a complete system of residues with respect to the modulus $c$. We
now take a primitive $m$th root of unity $r$ and a primitive
 $c$th root of unity $\varepsilon$ and form the \gesperrt{Lagrange}
\gesperrt{resolvents}
\weq{3}{(\varepsilon^{\beta}, r) = \sum^{n} \varepsilon^{\beta \nu} r^n,}
in which $\beta$ is an arbitrary exponent. The question at issue is
when such a resolvent can vanish. If
$\varkappa = 1$, so that $m$ is a prime, it vanishes, as we saw in
Bd.\ I, \S.\,169, for no value of $\beta$. In the general
 case of an arbitrary $\varkappa$, let $n'$ denote an arbitrary number
from $\mathfrak{N}$. Then $n n'$ runs through the whole group $\mathfrak{N}$ together with $n$.
One can therefore replace $n$ by $n n'$ in (3), if at the same time
%CONT>>


% ===== scan leaf 85, printed 64 =====
% !PRINTED 64
\weberpage{64}
%<<CONT
$\nu$ is replaced by $\nu + \nu'$, if $\nu'$ denotes the index of $n'$.
By this there follows from (3):
\[
\varepsilon^{-\beta \nu'} (\varepsilon^{\beta}, r) = \sum^{n} \varepsilon^{\beta \nu} r^{n n'},
\]
and if one therefore multiplies by $r^{n'}$ and takes the sum over $n'$:
\weq{4}{(\varepsilon^{-\beta}, r)(\varepsilon^{\beta}, r) = \sum^{n, n'} \varepsilon^{\beta \nu} r^{n'(n+1)}.}
Thus first of all the sum
\weq{5}{\sigma = \sum^{n'} r^{n'(n+1)}}
is to be determined. But the sum $\sigma$ always vanishes when
$n + 1$ is not divisible by $q^{\varkappa - 1}$ (by Bd.\ I, \S.\,133); for
then $r^{n+1}$ is a root of unity whose degree is a higher than
the first power of $q$. Hence $\sigma$ can be different from zero
only when $n$ has the form
\weq{6}{n \equiv -1 + t q^{\varkappa - 1} \pmod{m},}
and then we shall denote its value by $\sigma_t$. We obtain
all values of $n$ with respect to the modulus $m$ that are contained
in the form (6) if we let $t$ run through a complete system of residues
with respect to the modulus $q$, so for instance the values
\weq{7}{t = 0, 1, 2 \ldots q - 1}
may be set.
The sum $\sigma$ consists of $\varphi(m)$ terms. If $t = 0$, then
each of these terms becomes $= 1$ and it follows that
\weq{8}{\sigma_0 = \varphi(m) = q^{\varkappa} - q^{\varkappa - 1};}
for every other value of $t$, $r^{n+1}$ is a primitive $q$th root of unity
, and in $\sigma_t$ every such root of unity then occurs
$\varphi(m) : (q - 1) = q^{\varkappa - 1}$ times. But the sum of all primitive
$q$th roots of unity is equal to $-1$, and hence it follows for
$t = 1, 2 \ldots q - 1$ that
\weq{9}{\sigma_t = -q^{\varkappa - 1}.}
With this the sum $\sigma$ is determined.

But in order to derive from this the value of the sum in (4),
it is still necessary to ascertain the index $\nu$ of the numbers $n$ of the form (6).

For these numbers, however,
\[
g^{\nu} \equiv -1 \pmod{q^{\varkappa - 1}},
\]
and since by \gesperrt{Fermat's} theorem [\S.\,14, (11)]
\weq{10}{g^{1/2 \varphi(q^{\varkappa})} \equiv -1 \pmod{q^{\varkappa}}}
%CONT>>


% ===== scan leaf 86, printed 65 =====
% !PRINTED 65
\weberpage{65}
%<<CONT
it follows that
\[
g^\nu \equiv g^{\frac{1}{2}\varphi(q^\varkappa)} \ (\mathrm{mod}\ q^{\varkappa-1});
\]
hence $\nu \equiv \tfrac{1}{2}\varphi(q^\varkappa)\ \bigl[\mathrm{mod}\ \varphi(q^{\varkappa-1})\bigr]$, since the index of a number for the modulus $q^{\varkappa-1}$ is completely determined with respect to the modulus $\varphi(q^{\varkappa-1})$. Thus we have
\weq{11}{\nu \equiv \tfrac{1}{2}\varphi(q^\varkappa) + \tau\varphi(q^{\varkappa-1})\ \bigl[\mathrm{mod}\ \varphi(q^\varkappa)\bigr],}
and in this, since $\varphi(q^\varkappa) = q\,\varphi(q^{\varkappa-1})$, $\tau$ now runs, together with $t$, through a complete system of residues with respect to the modulus $q$. To the value $t=0$ corresponds the value $\tau=0$, on account of (10).

But we have
\[
\varepsilon^{\frac{1}{2}\varphi(q^\varkappa)} = -1, \qquad \varepsilon^{\varphi(q^{\varkappa-1})} = \varrho,
\]
where $\varrho$ is a primitive $q$th root of unity, and hence by (11)
\[
\varepsilon^\nu = -\varrho^\tau,
\]
and accordingly there results from (8) and (9):
\[
(\varepsilon^{-\beta}, r)(\varepsilon^{\beta}, r) = \sum \varepsilon^{\beta\nu}\sigma_t
\]
\[
= (-1)^\beta (q^\varkappa - q^{\varkappa-1}) - (-1)^\beta q^{\varkappa-1} \sum_{1,\,q-1}^{\tau} \varrho^{\tau\beta}.
\]

Now the sum $\sum\limits_{1,\,q-1}^{\tau} \varrho^{\tau\beta}$ has, when $\beta$ is not divisible by $q$, the value $-1$, and, when $\beta$ is divisible by $q$, the value $q-1$, so that one obtains the following result:
\weq{12}{\begin{aligned}
(\varepsilon^{-\beta}, r)(\varepsilon^{\beta}, r) &= 0, & \beta &\equiv 0\ (\mathrm{mod}\ q) \\
&= (-1)^\beta q^\varkappa, & \beta\ \text{not} &\equiv 0\ (\mathrm{mod}\ q).
\end{aligned}}

If, then, $\beta$ is not divisible by $q$, neither of the factors $(\varepsilon^{-\beta}, r), (\varepsilon^{\beta}, r)$ can vanish; but if $\beta$ is divisible by $q$, at least one of the two factors vanishes. That they then both vanish is shown by the following direct consideration of the sum.

If $\beta$ is divisible by $q$, then for every integer $t$
\[
\varepsilon^{\beta[\nu + t\varphi(q^{\varkappa-1})]} = \varepsilon^{\beta\nu}
\]
\[
g^{t\varphi(q^{\varkappa-1})} \equiv 1 + \tau q^{\varkappa-1}\ (\mathrm{mod}\ q^\varkappa),
\]
in which $\tau$ now runs, together with $t$, through a complete system of residues with respect to the modulus $q$. If, then, under the assumption that $\beta$ is divisible by $q$, one replaces in (3) $\nu$ by $\nu + t\varphi(q^{\varkappa-1})$, i.\,e.\ $n$ by $n(1+\tau q^{\varkappa-1})$, it follows that
\[
(\varepsilon^\beta, r) = \sum^{n} \varepsilon^{\beta\nu} r^n r^{n q^{\varkappa-1}\tau}.
\]


% ===== scan leaf 87, printed 66 =====
% !PRINTED 66
\weberpage{66}
This sum is therefore independent of the arbitrarily chosen $\tau$. If we sum here from $\tau=0$ to $\tau=q-1$, then, since $\sum r^{n q^{\varkappa-1}\tau}=0$, we obtain:
\weq{13}{(\varepsilon^\beta, r) = 0.}
We therefore have the theorem:

1.~\gesperrt{If the degree of the root of unity $r$ is a power of an odd prime, then the resolvent $(\varepsilon^\beta, r)$ vanishes if and only if $\beta$ is divisible by $q$.}

We still have to consider the case in which the degree of the root of unity $r$ is a power of 2. We therefore set
\[
m = 2^\lambda
\]
and assume at first $\lambda \geqq 3$. Then for every odd number $n$ we can determine a pair of indices $\nu, \nu_1$ from the congruences
\[
n \equiv (-1)^{\nu_1} 5^\nu\ (\mathrm{mod}\ m)
\]
and indeed $\nu_1$ is determined to the modulus 2, and $\nu$ to the modulus $2^{\lambda-2}$. By the resolvents we understand in this case the sums
\weq{14}{\bigl((-1)^{\beta_1}, \Theta^\beta, r\bigr) = \sum^{n} (-1)^{\beta_1 \nu_1} \Theta^{\beta \nu} r^n,}
in which $\Theta$ denotes a primitive root of unity of degree $2^{\lambda-2}$. We now proceed quite similarly as before. If in (14) we replace $n$ by $nn'$ and denote by $\nu_1', \nu'$ the indices of $n'$, we obtain
\weq{15}{(-1)^{-\beta_1 \nu_1'} \Theta^{-\beta \nu'} \bigl((-1)^{\beta_1}, \Theta^\beta, r\bigr) = \sum^{n} (-1)^{\beta_1 \nu_1} \Theta^{\beta \nu} r^{n n'},}
and from this, by multiplication with $r^{n'}$ and summation over $n'$,
\weq{16}{\bigl((-1)^{-\beta_1}, \Theta^{-\beta}, r\bigr)\bigl((-1)^{\beta_1}, \Theta^\beta, r\bigr) = \sum^{n n'} (-1)^{\beta_1 \nu_1} \Theta^{\beta \nu} r^{n'(n+1)}.}
Now, for the reasons adduced above,
\[
\begin{aligned}
\sum^{n'} r^{n'(n+1)} &= 0, & &\text{if } n+1 \text{ is not divisible by } 2^{\lambda-1}\text{,} \\
&= 2^{\lambda-1} & &\text{for } n+1 \equiv 0\ (\mathrm{mod}\ 2^\lambda),\ \nu=0,\ \nu_1=1, \\
&= -2^{\lambda-1} & &\text{for } n+1 \equiv 2^{\lambda-1}\ (\mathrm{mod}\ 2^\lambda),\ \nu=2^{\lambda-3},\ \nu_1=1,
\end{aligned}
\]
and accordingly (16) gives
\weq{17}{\begin{aligned}
\bigl((-1)^{-\beta_1}, \Theta^{-\beta}, r\bigr)\bigl((-1)^{\beta_1}, \Theta^\beta, r\bigr) &= (-1)^{\beta_1} 2^\lambda, & \beta &\equiv 1\ (\mathrm{mod}\ 2) \\
&= 0, & \beta &\equiv 0\ (\mathrm{mod}\ 2).
\end{aligned}}


% ===== scan leaf 88, printed 67 =====
% !PRINTED 67
\weberpage{67}
\pdfbookmark[2]{§ 18. Cyclotomic Fields}{bm22}
That in the case of an even $\beta$ each of the two factors vanishes is seen by setting, in (15),
\[
n' = 1 + 2^{\lambda-1},
\]
and consequently
\[
r^{n'} = -r, \qquad \nu_1' = 0, \qquad \nu' = 2^{\lambda-3}
\]
there. Then this formula gives, since $n$ is odd:
\[
\bigl((-1)^{\beta_1}, \Theta^\beta, r\bigr) = -\Theta^{\beta 2^{\lambda-3}} \bigl((-1)^{\beta_1}, \Theta^\beta, r\bigr).
\]
For even $\beta$, however, $\Theta^{\beta 2^{\lambda-3}} = +1$, and consequently:
\weq{18}{\bigl((-1)^{\beta_1}, \Theta^\beta, r\bigr) = 0;}
hence:

2.~\gesperrt{If the degree of the root of unity $r$ is a power of 2 and greater than 4, then the resolvent $\bigl((-1)^{\beta_1}, \Theta^\beta, r\bigr)$ vanishes if and only if $\beta$ is even.}

In the case $m=4$, and so $r=i$, one has only the two resolvents $i+i^3, i-i^3$, which we can combine under the designation $((-1)^\beta, i)$, $\beta=0,1$, and $((-1)^\beta, i)$ is $=0$ if and only if $\beta=0$.

\begin{center}
\S.\,18.\\
\gesperrt{Cyclotomic Fields.}
\end{center}

The theory of \gesperrt{Abelian} groups opens to us a deeper insight into the theory of the roots of unity and of the algebraic numbers that spring from them.

Now let $m$ be any positive integer, decomposed into its prime factors:
\weq{1}{m = 2^\lambda q_1^{\varkappa_1} q_2^{\varkappa_2} \cdots,}
and let $r$ be a primitive $m$th root of unity. The case $\lambda=1$ we can exclude once and for all from our consideration; for if $m$ is odd and $r$ runs through the primitive $m$th roots of unity, no two opposite ones occur among them, and $-r$ runs through the primitive $2m$th roots of unity.

$r$ is the root of an irreducible \gesperrt{Abelian} equation with integral coefficients, of degree
\weq{2}{\nu = \varphi(m),}
as we have proved in \S.\,134 of the first volume.


% ===== scan leaf 89, printed 68 =====
% !PRINTED 68
\weberpage{68}
The totality of all rational functions of $r$ with rational numbers as coefficients is thus a number field $\Omega(r)$ of degree $\nu$, which we call a \gesperrt{cyclotomic field}. We wish, however, to take the concept of the cyclotomic field somewhat more generally, and to understand by it every field whose numbers are all rational functions of any roots of unity whatever.

The field $\Omega(r)$ of degree $\nu$, which consists of \gesperrt{all} rational functions of an $m$th root of unity $r$, we call, for more precise distinction, the \gesperrt{full cyclotomic field} of order $m$, and we denote it by $\Omega_m$.

Arbitrary roots of unity $r, r', r'' \ldots$ of arbitrary degrees $m, m', m'' \ldots$ can always be regarded as powers of one and the same root of unity $\varrho$, whose degree is the least common multiple of $m, m', m'' \ldots$. Accordingly a cyclotomic field which contains only rational functions of $r, r', r'' \ldots$ is a divisor of the full cyclotomic field $\Omega(\varrho)$, and we therefore obtain all cyclotomic fields that exist at all when we seek out all the divisors of all full cyclotomic fields.

The \gesperrt{Galois} group of the field $\Omega_m$ consists of all the substitutions
\weq{3}{(r, r^n),}
where $n$ denotes every relative prime to $m$, taken according to the modulus $m$. For the cyclotomic equation of degree $\nu$, whose roots are the $r^n$, is a normal equation and hence its own \gesperrt{Galois} resolvent. Since, if $a, b$ are two of these numbers $n$,
\[
(r, r^a)(r, r^b) = (r, r^{ab})
\]
this group is isomorphic with the group $\mathfrak{N}$ of all number classes $N$ of the numbers coprime to $m$, which we have considered in the preceding section.

If $\mathfrak{A}$ is a divisor of $\mathfrak{N}$, and $\varrho$ a function from $\Omega_m$ belonging to $\mathfrak{A}$, then the totality of the rational functions of $\varrho$, $\Omega(\varrho)$, is a field contained in $\Omega_m$.

If conversely $\Omega$ is a divisor of $\Omega_m$ and $\varrho$ a primitive number of the field $\Omega$, hence also a number in $\Omega_m$, then $\Omega$ can be represented as the totality of the rational functions of $\varrho$ and denoted by $\Omega(\varrho)$.


% ===== scan leaf 90, printed 69 =====
% !PRINTED 69
\weberpage{69}
This function $\varrho$ then belongs to a certain group $\mathfrak{A}$,
which is a divisor of $\mathfrak{N}$. We also call the group $\mathfrak{A}$ and
the field $\Omega(\varrho)$ \gesperrt{belonging together}, and from this we draw the
theorem:

\gesperrt{1. To every divisor $\mathfrak{A}$ of $\mathfrak{N}$ there belongs a certain
cyclotomic field $\Omega(\varrho)$ contained in $\Omega_m$, and
conversely to every divisor $\Omega(\varrho)$ of $\Omega_m$
there belongs a certain divisor $\mathfrak{A}$ of $\mathfrak{N}$ as group, in the
sense that, if $a$ is a number from $\mathfrak{A}$, all
numbers of the field $\Omega(\varrho)$ admit the permutations $(r,r^a)$,
and that conversely every number in $\Omega_m$
that admits these permutations is contained in $\Omega(\varrho)$.}

The Galois group of such a field $\Omega(\varrho)$ we obtain, according to Bd. I, §. 156, by decomposing $\mathfrak{N}$ into the system of
subsidiary groups:
\[
\mathfrak{N} = \mathfrak{A} + \mathfrak{A}_1 + \mathfrak{A}_2 + \cdots,
\]
and seeking the permutations produced among the subsidiary groups $\mathfrak{A}, \mathfrak{A}_1, \mathfrak{A}_2 \ldots$ by composition
with the elements of $\mathfrak{N}$. But this group is, according to §. 4, isomorphic with the
group $\mathfrak{N}/\mathfrak{A}$, and hence also isomorphic with the group reciprocal to $\mathfrak{A}$
(§. 12), which we shall denote by $\mathfrak{B}$. If $a$ is the degree of $\mathfrak{A}$
and $b$ that of $\mathfrak{B}$, then $ab = \nu$, and $\varrho$ satisfies an irreducible
Abelian equation of degree $b$.

\gesperrt{2. We thus obtain all cyclotomic fields
if for every modulus $m$ we form all the
divisors $\mathfrak{A}$ of the group $\mathfrak{N}$, seek for each of these
groups $\mathfrak{A}$ an associated function $\varrho$,
and from these derive the fields $\Omega(\varrho)$.}

But there remains the question whether in this process one and
the same field $\Omega$ can occur more than once, whereby we are led to
the investigation of the common divisors of two cyclotomic fields.

According to the definition given above, one must carefully distinguish between the group to which a field $\Omega$ \gesperrt{belongs}
and the Galois group of the field; the two groups are
reciprocal to each other. Thus the field $\Omega_m$ itself belongs to the
unit group, while its Galois group is $\mathfrak{N}$,


% ===== scan leaf 91, printed 70 =====
% !PRINTED 70
\weberpage{70}
The following theorem now holds:

\gesperrt{3. If $\Omega'$ and $\Omega''$ are two divisors of $\Omega_m$ which belong to the
groups $\mathfrak{A}'$ and $\mathfrak{A}''$, then $\Omega''$ is a divisor
of $\Omega'$ if and only if $\mathfrak{A}'$ is a divisor
of $\mathfrak{A}''$.}

For if $\Omega''$ is a divisor of $\Omega'$, then all numbers
of $\Omega''$ must admit the substitutions of the group $\mathfrak{A}'$ to which $\Omega'$ belongs; hence $\mathfrak{A}'$ must be contained in $\mathfrak{A}''$. And if conversely $\mathfrak{A}'$ is contained in $\mathfrak{A}''$,
then all numbers of $\Omega''$ admit
all the substitutions of $\mathfrak{A}'$ and are therefore rationally
representable by a primitive number of the field $\Omega'$; that is, $\Omega''$
is contained in $\Omega'$ (§. 144 of the first volume).

The field $\Omega_m$ contains as divisors all the fields $\Omega_{m_1}$ for
which $m_1$ is a divisor of $m$; for $\Omega_{m_1}$ consists of all
rational functions of $r^{\frac{m}{m_1}}$. Now we have
\[
r^{a\frac{m}{m_1}} = r^{\frac{m}{m_1}},
\]
if and only if
\weq{4}{a \equiv 1\,(\mathrm{mod}\ m_1).}

The numbers $a$ in $\mathfrak{N}$ that satisfy the congruence (4) thus
form the group to which the field $\Omega_{m_1}$ belongs. We will
denote this group by $\mathfrak{A}_{m_1}$ and symbolically put
\weq{5}{\mathfrak{A}_{m_1} \equiv 1\,(\mathrm{mod}\ m_1).}

Since the degree of the field $\Omega_{m_1}$ is equal to $\varphi(m_1)$, which is
at the same time the index of the divisor $\mathfrak{A}_{m_1}$ of $\mathfrak{N}$, the degree
of $\mathfrak{A}_{m_1}$ is equal to $\varphi(m):\varphi(m_1)$.

If $m = m_1 m_2$ and $m_1$ is relatively prime to $m_2$, one obtains
all the numbers of $\mathfrak{A}_{m_1}$ from $a = 1 + x m_1$, if $x$ is determined
so that $a$ becomes congruent modulo $m_2$ with all the numbers relatively prime to
$m_2$. From this there results, in agreement
with the above general determination, the degree of $\mathfrak{A}_{m_1}$ for
this case $=\varphi(m_2)$.

But if now $m_1$ and $m_2$ are any two divisors of $m$, and
$\mu$ is the greatest common divisor of $m_1$ and $m_2$, then
the greatest common divisor of the two fields $\Omega_{m_1}, \Omega_{m_2}$
will belong to the group which one may regard as the least common
multiple of the two groups
\weq{6}{\mathfrak{A}_{m_1} \equiv 1\,(\mathrm{mod}\ m_1),\quad \mathfrak{A}_{m_2} \equiv 1\,(\mathrm{mod}\ m_2)}
%CONT>>


% ===== scan leaf 92, printed 71 =====
% !PRINTED 71
\weberpage{71}
%<<CONT
obtains. But this is the group
\weq{7}{\mathfrak{A}_\mu \equiv 1\,(\mathrm{mod}\ \mu),}

that is, the group consisting of all those numbers of $\mathfrak{N}$ which satisfy
the congruence $a \equiv 1\,(\mathrm{mod}\ \mu)$. For first of all it is clear
that $\mathfrak{A}_{m_1}$ and $\mathfrak{A}_{m_2}$ are divisors of $\mathfrak{A}_\mu$. Hence the
least common multiple of $\mathfrak{A}_{m_1}$ and $\mathfrak{A}_{m_2}$ is also contained in $\mathfrak{A}_\mu$.
But if, on the other hand, $\alpha = 1 + \xi\mu$ is any number in $\mathfrak{A}_\mu$, then
one can determine the numbers $a_1 = 1 + x_1 m_1$ in $\mathfrak{A}_{m_1}$ and $a_2 = 1 + x_2 m_2$
in $\mathfrak{A}_{m_2}$ in such a way that $\alpha \equiv a_1 a_2\,(\mathrm{mod}\ m)$. One has
only to set $\xi\mu = x_1 m_1 + x_2 m_2$, which by Bd. I, §. 118 is always
possible. Hence $\mathfrak{A}_\mu$ is also contained in the least common multiple of $\mathfrak{A}_{m_1}$ and $\mathfrak{A}_{m_2}$, which is accordingly
$\mathfrak{A}_\mu$ itself. From this, then, the theorem:

\gesperrt{4. If $m_1$ and $m_2$ are any two natural numbers
and $\mu$ is the greatest common divisor of
$m_1$ and $m_2$, then the field $\Omega_\mu$ is the greatest
common divisor of $\Omega_{m_1}$ and $\Omega_{m_2}$, and
$\Omega_{\frac{m_1 m_2}{\mu}}$ their least common multiple.}

From this it now follows: if two full cyclotomic fields $\Omega_m$
and $\Omega_{m'}$ have a common divisor $\Omega$, and $m'$ is smaller
than $m$, then there must be a proper divisor $d$ of $m$ such that
$\Omega$ is also a divisor of $\Omega_d$.

We shall call a divisor of $\Omega_m$ \gesperrt{primary} if it is
not at the same time contained in a full cyclotomic field $\Omega_{m'}$ of
lower $m'$. It then follows that one obtains all
the non-primary divisors of $\Omega_m$ by letting the
index $d$ in $\Omega_d$ run through all the proper divisors of $m$ and
seeking out the divisors of $\Omega_d$.

If now we denote by $q_1, q_2, \ldots q_\tau$ all the distinct
primes dividing $m$ (2 included) and
set
\weq{8}{m = q_1 m_1 = q_2 m_2 = \cdots = q_\tau m_\tau,}

then every $d$ is a divisor of one of $m_1, m_2, \ldots m_\tau$, and we obtain the non-primary divisors of $\Omega_m$ if, for the fields
\weq{9}{\Omega_{m_1},\ \Omega_{m_2},\ \ldots\ \Omega_{m_\tau}}
we seek out all their divisors.

These fields, however, belong to the group
\weq{10}{Q_1 \equiv 1\,(\mathrm{mod}\ m_1),\ \ Q_2 \equiv 1\,(\mathrm{mod}\ m_2),\ \ldots\ Q_\tau \equiv 1\,(\mathrm{mod}\ m_\tau),}
%CONT>>


% ===== scan leaf 93, printed 72 =====
% !PRINTED 72
\weberpage{72}
\pdfbookmark[2]{§ 19. Primary and Non-primary Divisors of the Group N}{bm23}
%<<CONT
and thus by 3. a divisor $\Omega$ of $\Omega_m$ will fail to be primary
if and only if one of the groups $Q_1, Q_2, \ldots Q_\tau$ is contained
in the group $\mathfrak{A}$ belonging to $\Omega$. If, then, we call
those divisors of the group $\mathfrak{N}$ which contain none of the groups $Q_1, Q_2 \ldots Q_\tau$
as divisors primary divisors, we can state the theorem:

\gesperrt{5. To obtain all primary divisors of $\Omega_m$,
one has to seek out all primary divisors of the group $\mathfrak{N}$
and to form the associated fields.}

\gesperrt{6. But if one sets up all primary divisors of all
full cyclotomic fields, one obtains
every cyclotomic field, and each one
only once, and indeed each represented by
roots of unity of the lowest possible degree.}

The field $\Omega_{2m}$, when $m$ is odd, has no primary divisor at all and is identical with $\Omega_m$. In all other cases
$\Omega_m$ is at least its own primary divisor.

For what follows it is further necessary that we become clear
about the extent and the constitution of the groups $Q$.

\begin{center}
\S.\,19.

\gesperrt{Primary and non-primary divisors of the group $\mathfrak{N}$.}
\end{center}

We denote by $q$ any one of the primes dividing $m$
and set $m = q m'$. We then consider the
group
\[
Q \equiv 1\,(\mathrm{mod}\ m').
\]

To determine the conditions for a number in $Q$, we apply
the representation of the group $\mathfrak{N}$ by a basis and the
notation of the indices, as we introduced and explained them
in §. 16, and accordingly denote the group reciprocal to $Q$
by $P$.

The indices of a number in $Q$ we denote by
\[
\gamma_{-1}, \gamma_0, \gamma_1, \ldots \gamma_\mu,
\]
and those of a number in $P$ by
\[
\delta_{-1}, \delta_0, \delta_1, \ldots \delta_\mu.
\]


% ===== scan leaf 94, printed 73 =====
% !PRINTED 73
\weberpage{73}

Then the $\gamma$ must satisfy the condition:

\weq{1}{C_{-1}^{\gamma_{-1}} C_0^{\gamma_0} C_1^{\gamma_1} \dots C_\mu^{\gamma_\mu} \equiv 1 \, (\mathrm{mod}\ m').}

This condition demands that all indices, with the exception of the one corresponding
to the prime $q$, which we denote by $\gamma$, must be zero. With respect to
$\gamma$, however, a distinction is to be made according as $q$ still divides $m'$
or not, i.\,e. whether $q$ is a multiple or only a simple factor of $m$. If $q$ is
no longer contained in $m'$, then the congruence (1) contains no restriction at all
for $\gamma$, and $\gamma$ can take every value with respect to the modulus $c$,
i.\,e. every value $0, 1, \dots q - 2$.

But if $q$ is still contained in $m'$, and hence $m$ is divisible by $q^\varkappa$, and $\varkappa > 1$,
then the condition (1) demands:
\[
C^\gamma \equiv 1 \, (\mathrm{mod}\ q^{\varkappa-1}),
\]
hence
\weq{2}{\gamma \equiv 0 \, \left(\mathrm{mod}\ \dfrac{c}{q}\right),}
and $\gamma$ can therefore receive each of the values
\[
0, \dfrac{c}{q}, \dfrac{2c}{q}, \dots \dfrac{(q-1)c}{q}
\]
and no others. In the first case the degree of the group $Q$ is equal to $q - 1$, in the
second equal to $q$.

The indices $\delta$ of a number from the group $P$ are obtained, by \S.\,12, 7.
and \S.\,16, (7), from the condition that for all admissible $\gamma$
\weq{3}{\varepsilon_{-1}^{\gamma_{-1}\delta_{-1}} \varepsilon_0^{\gamma_0\delta_0} \varepsilon_1^{\gamma_1\delta_1} \dots \varepsilon_\mu^{\gamma_\mu\delta_\mu} = 1}
shall hold. By what has just been proved about the indices $\gamma$, however,
(3) demands only the one thing, that the index $\delta$ corresponding to the
prime $q$ shall be divisible by $q$ or by $q - 1$, according as $q$ divides $m$
several times or only once. We therefore emphasise the theorem:

7. \gesperrt{The group $P$ reciprocal to $Q$ is characterised by the fact that
the index $\delta$ corresponding to $q$ of all its numbers is divisible by $q$
or by $q - 1$, according as $q$ divides $m$ several times or only once.}

And from this:

8. \gesperrt{A non-primary divisor $\mathfrak{A}$ of $\mathfrak{N}$ is
characterised by the fact that in the reciprocal group $\mathfrak{B}$}
%CONT>>


% ===== scan leaf 95, printed 74 =====
% !PRINTED 74
\weberpage{74}
\pdfbookmark[2]{§ 20. The Cyclotomic Periods}{bm24}
%<<CONT
\gesperrt{the index $\beta$ corresponding to a prime $q$, of all the numbers $b$, is divisible by $q$ or by $q - 1$, according as $q$ occurs several times or only once among the prime factors of $m$.}

\begin{center}
\S.\,20.

\gesperrt{The cyclotomic periods.}
\end{center}

As the simplest functions by which one may attempt to represent the cyclotomic fields, the cyclotomic periods present themselves, which are an immediate generalisation of the Gauss periods considered in \S.\,167 of the first volume. By them we understand the following.

Let $\mathfrak{A}$ denote a divisor of the group $\mathfrak{N}$ of index $e$, and let $a$ run through the numbers of $\mathfrak{A}$. If $r$ is a primitive $m$th root of unity, then the sum
\weq{1}{\eta = \overset{a}{\Sigma}\, r^a}
is called a \gesperrt{cyclotomic period belonging to the group $\mathfrak{A}$} of index $e$.

If we make in (1) one of the substitutions $(r, r^a)$, then $\eta$ remains unchanged, as follows immediately from the group property of the $a$.

In order to decompose $\mathfrak{N}$ into the subsidiary groups belonging to $\mathfrak{A}$, we must choose the numbers $1, n_1, n_2, \ldots n_{e-1}$ from $\mathfrak{N}$ suitably, so that the decomposition
\weq{2}{\mathfrak{N} = \mathfrak{A} + \mathfrak{A} n_1 + \mathfrak{A} n_2 + \cdots + \mathfrak{A} n_{e-1}}
is obtained.
If we then make in $\eta$ the substitutions
\weq{3}{(r, r), (r, r^{n_1}), \ldots (r, r^{n_{e-1}}),}
then $\eta$ passes into the conjugate periods
\weq{4}{\eta, \eta_1, \ldots \eta_{e-1}}
and if these are all different from one another, then $\eta$ belongs to the group $\mathfrak{A}$. The field $\Omega(\eta)$ belonging to $\mathfrak{A}$ consists of all rational functions of $\eta$, and the numbers $\eta_1 \ldots \eta_{e-1}$ all belong to the same group $\mathfrak{A}$.

If, however, the quantities (4) are not all different from one another, then the number $\eta$ does not belong to the group $\mathfrak{A}$, but to a more comprehensive group $\mathfrak{A}'$, of which $\mathfrak{A}$ is a divisor
%CONT>>


% ===== scan leaf 96, printed 75 =====
% !PRINTED 75
\weberpage{75}
%<<CONT
is. To determine the conditions for this case, we extend somewhat
further the concept of the \gesperrt{resolvents}, as we have already
considered it in \S.\,17.

We let $n$ run through the sequence of the numbers of the group $\mathfrak{N}$,
and denote by $\chi(n)$ one of the characters of this
group. The extended resolvents are then, if $r$ is an
$m$th root of unity:
\weq{5}{(\chi, r) = \sum^{n} \chi(n) r^n.}
If we understand by $\mathfrak{B}$ the group reciprocal to $\mathfrak{A}$ and let
$b$ run through the numbers of $\mathfrak{B}$, then by \S.\,12, 7. there is
a number of characters $\chi_b$, equal to the degree of $\mathfrak{B}$, which are
distinguished by the fact that
\[
\chi_b(a) = 1
\]
for every number $a$ from the group $\mathfrak{A}$. From this it then follows,
by the definition of the characters \S.\,11, (4), for every $n$:
\weq{6}{\chi_b(a n) = \chi_b(n),}
that is, the character $\chi_b$ has one and the same value for all
numbers of each subsidiary group $\mathfrak{A} n_1, \mathfrak{A} n_2 \ldots$ Accordingly
the resolvent $(\chi_b, r)$ will depend only on the periods $\eta$, and
will take the expression:
\weq{7}{(\chi_b, r) = \eta + \chi_b(n_1) \eta_1 + \cdots + \chi_b(n_{e-1}) \eta_{e-1}.}
If $\eta$ belongs to a group $\mathfrak{A}'$, then the conjugate
numbers $\eta_1, \eta_2, \ldots \eta_{e-1}$ belong to the same group (because $\mathfrak{A}'$ is a
normal divisor of $\mathfrak{N}$, Bd.\ I, \S.\,154). If therefore $a'$ is any
number from $\mathfrak{A}'$, the quantities $\eta, \eta_1, \eta_2, \ldots \eta_{e-1}$ remain
unchanged by the substitution $(r, r^{a'})$, and by (7) we also have
\weq{8}{(\chi_b, r) = (\chi_b, r^{a'}).}
On the other hand, one obtains from (5), if one bears in mind that
$a' n$ runs through the whole group $\mathfrak{N}$ together with $n$,
\[
\begin{aligned}
(\chi_b, r) &= \sum^{n} \chi_b(n a') r^{n a'} = \chi_b(a') \sum^{n} \chi_b(n) r^{n a'}\\
&= \chi_b(a') (\chi_b, r^{a'}),
\end{aligned}
\]
hence by (8):
\weq{9}{(\chi_b, r) = \chi_b(a') (\chi_b, r).}
If $\mathfrak{A}$ is not identical with $\mathfrak{A}'$, but is a proper divisor
of $\mathfrak{A}'$, then the group $\mathfrak{B}'$ reciprocal to $\mathfrak{A}'$ is a proper divisor
of $\mathfrak{B}$ (by \S.\,12, 8.); if therefore $b$ is a number in $\mathfrak{B}$ which does not
%CONT>>


% ===== scan leaf 97, printed 76 =====
% !PRINTED 76
\weberpage{76}
%<<CONT
is at the same time contained in $\mathfrak{B}'$, then one can choose $a'$ so that $\chi_b(a')$ is not $= 1$. But then it follows from (9)
\weq{10}{(\chi_b, r) = 0;}
hence the theorem:

\gesperrt{1. If the period $\eta$ belongs not to $\mathfrak{A}$, but to a more comprehensive group $\mathfrak{A}'$, then every resolvent $(\chi_b, r)$ vanishes when $b$ is a number from $\mathfrak{B}$ that does not occur in $\mathfrak{B}'$.}

In order now to pursue further the conditions for this theorem, we must take into account the mode of formation of the characters.

We choose the notation as we introduced it at the close of \S.\,16. Then, if $\beta_{-1}, \beta_0, \beta_1, \ldots \beta_\mu$ are the indices of $b$ and $\nu_{-1}, \nu_0, \nu_1, \ldots \nu_\mu$ those of $n$,
\weq{11}{\chi_b(n) = \varepsilon_{-1}^{\beta_{-1}\,\nu_{-1}}\, \varepsilon_0^{\beta_0\nu_0}\, \varepsilon_1^{\beta_1\nu_1} \ldots \varepsilon_\mu^{\beta_\mu\nu_\mu}.}

If we denote by $r_0, r_1, \ldots r_\mu$ primitive roots of unity of the degrees $2^\lambda, q_1^{\varkappa_1}, q_2^{\varkappa_2}, \ldots q_\mu^{\varkappa_\mu}$, then we can represent every primitive $m$th root of unity in the form (Bd.\,I, \S.\,132)
\weq{12}{r = r_0 r_1 \ldots r_\mu,}
and if we determine $n_0, n_1, \ldots n_\mu$ from the congruences
\weq{13}{n \equiv n_0 \ (\mathrm{mod}\ 2^\lambda), \quad n \equiv n_1 \ (\mathrm{mod}\ q_1^{\varkappa_1}), \ldots n \equiv n_\mu \ (\mathrm{mod}\ q_\mu^{\varkappa_\mu}),}
then $(\chi_b, r)$ splits by (5) into the product of the following sums:
\weq{14}{\overset{n_0}{\Sigma}\, \varepsilon_{-1}^{\beta_{-1}\,\nu_{-1}}\, \varepsilon_0^{\beta_0\nu_0}\, r_0^{n_0}, \quad \overset{n_1}{\Sigma}\, \varepsilon_1^{\beta_1\nu_1}\, r_1^{n_1}, \ldots \overset{n_\mu}{\Sigma}\, \varepsilon_\mu^{\beta_\mu\nu_\mu}\, r_\mu^{n_\mu}.}

Now if $b$ is a number that occurs in $\mathfrak{B}$ but not in $\mathfrak{B}'$, then by theorem 1. one of these sums must vanish.

But for this, by \S.\,17, the necessary and sufficient condition is that one of the congruences
\weq{15}{\beta_0 \equiv 0\ (\mathrm{mod}\ 2), \quad \beta_1 \equiv 0\ (\mathrm{mod}\ q_1), \ldots \beta_\mu \equiv 0\ (\mathrm{mod}\ q_\mu)}
is satisfied, and indeed one whose modulus divides $m$ more than once.

This condition must first of all, when the hypothesis of theorem 1. holds, be fulfilled for every number $b$ that occurs in $\mathfrak{B}$ but not in $\mathfrak{B}'$. But if $b'$ is a number in $\mathfrak{B}'$, then for every integer exponent $x$ the product $b'^{x} b$ is in $\mathfrak{B}$, but
%CONT>>


% ===== scan leaf 98, printed 77 =====
% !PRINTED 77
\weberpage{77}

%<<CONT
not contained in $\mathfrak{B}'$, and hence, if $\beta_0', \beta_1', \dots \beta_\mu'$ are the
indices of $b'$, then for every $x$ one of the congruences

\weq{16}{
\begin{aligned}
x\beta_0' + \beta_0 &\equiv 0 \pmod{2} \\
x\beta_1' + \beta_1 &\equiv 0 \pmod{q_1} \\
&\cdots\cdots\cdots\cdots \\
x\beta_\mu' + \beta_\mu &\equiv 0 \pmod{q_\mu}
\end{aligned}
}

must be satisfied.

If now two numbers $\beta', \beta$ are not both divisible by a prime $q$, then one can always dispose of $x$ in such a way that $x\beta' + \beta$ does not become divisible by $q$; one need only assume $x$ not divisible by $q$ when $\beta$ is divisible by $q$, and $x$ divisible by $q$ when $\beta$ is not divisible by $q$. One can therefore also determine $x$ so that it satisfies several such requirements simultaneously. It therefore follows from (16) that either $\beta_0', \beta_0$ must be divisible by $2$, or $\beta_1', \beta_1$ by $q_1 \dots$ or $\beta_\mu', \beta_\mu$ by $q_\mu$; i.\,e. one of the congruences (15) must also be fulfilled when $b$ is any number from $\mathfrak{B}$, no matter whether it occurs in $\mathfrak{B}'$ or not. Finally it follows again from (16), if we understand by $b, b'$ two arbitrary numbers from $\mathfrak{B}$, that one and the same of the congruences (15) must hold for all numbers from $\mathfrak{B}$. We can now summarise this as follows:

\gesperrt{2. If the period $\eta$ does not belong to $\mathfrak{A}$, then among the primes (including 2) that divide $m$ more than once there must be one, $q$, such that for every number $b$ from $\mathfrak{B}$ the index $\beta$ corresponding to $q$ is divisible by $q$.}

We now consider the group $Q$ which, according to the notation (5) of \S.\,18, is defined by the congruence

\weq{17}{Q \equiv 1 \pmod{\dfrac{m}{q}}}

and denote by $P$ the group reciprocal to it.

By Theorem 7., \S.\,19, every group which, like $\mathfrak{B}$, has the property that the index $\beta$ corresponding to $q$ is divisible by $q$, is a divisor of the group $P$, and consequently, by \S.\,12, 8., $Q$ is a divisor of $\mathfrak{A}$; hence, by \S.\,19, 8., $\mathfrak{A}$ is a non-primary divisor of $\mathfrak{N}$, and with this we have arrived at the theorem:

\gesperrt{3. If $\mathfrak{A}$ is a primary divisor of $\mathfrak{N}$, then the period $\eta$ belongs to the group $\mathfrak{A}$.}


% ===== scan leaf 99, printed 78 =====
% !PRINTED 78
\weberpage{78}

All primary divisors of the full cyclotomic field $\Omega_m$ are accordingly representable in the form $\Omega(\eta)$, i.\,e. all numbers of such a divisor are rationally representable by the cyclotomic periods $\eta$, and since one finds the non-primary divisors of $\Omega_m$ again as primary divisors of lower cyclotomic fields, it follows:

\gesperrt{4. All cyclotomic fields are representable in the form $\Omega(\eta)$.}

Or also:

\gesperrt{5. Every rational function of roots of unity can be represented as a rational function of a cyclotomic period.}

And this theorem can also be expressed in the form:

\gesperrt{6. If $\mathfrak{A}$ is an arbitrary primary or non-primary divisor of $\mathfrak{N}$ of index $e$, then there is a divisor $m_1$ of $m$ and a cyclotomic period $\eta$ of index $e$ lying in $\Omega_{m_1}$, such that $\eta$, regarded as a function of $r$, is a function belonging to the group $\mathfrak{A}$, and thus, if $a$ belongs to $\mathfrak{A}$, remains unchanged under the substitution $(r, r^a)$ and receives $e$ different values for all the substitutions of $\mathfrak{N}$.}

If $a$ runs through a group $\mathfrak{A}$ which is a primary divisor of $\mathfrak{N}$, and we set

\[
\eta_h = \sum^a r^{ah},
\]

even when $h$ is not relatively prime to $m$, then we can represent these quantities $\eta_h$, which all admit the permutations of the group $\mathfrak{A}$, rationally by $\eta = \sum r^a$. On the other hand, however, every rational function of $\eta$ can also be represented linearly by the $\eta_h$. This will be proved once it is shown that the product of two $\eta_h$ can be expressed linearly by the $\eta_h$. But we have:

\[
\eta_h \eta_k = \sum^{a, a'} r^{ha + ka'};
\]

if one first takes the sum over $a'$ with $a$ held fixed, one can replace $a'$ by $aa'$, since $aa'$ runs through the group $\mathfrak{A}$ at the same time as $a'$. One thus obtains

\[
\eta_h \eta_k = \sum^{a, a'} r^{a(h + ka')} = \sum^{a'} \eta_{h + ka'}.
\]


% ===== scan leaf 100, printed 79 =====
% !PRINTED 79
\weberpage{79}
\pdfbookmark[2]{§ 21. Cyclotomic Fields with a Given Group}{bm25}

If the prime $q$ divides $m$ only once and $\mathfrak{A}$ is divisible by $Q$, then $\mathfrak{A}$ is indeed not primary; but the period $\eta$ nevertheless belongs to the group $\mathfrak{A}$.

In this case, however, $\eta$ can be represented by roots of unity of lower degree. This representation may be found in the following way. All numbers of $Q$ are of the form $1 + hm'$, where $h$ runs through the sequence of numbers $0, 1, \dots q - 1$, with the exception of the single value for which

\weq{18}{1 + hm' = kq}

becomes divisible by $q$. If, then, under the assumption that $\mathfrak{A}$ is divisible by $Q$, we set

\[
\mathfrak{A} = Q + Qa_1 + Qa_2 \dots,
\]

it follows that every number $a$ in $\mathfrak{A}$ is of the form:

\weq{19}{a = a'(1 + hm'),}

where $a'$ runs through the sequence of numbers $1, a_1, a_2 \dots$. From this it is further to be concluded that the residues of the numbers $1, a', a'' \dots$ with respect to the modulus $m'$ form a group. If we now take the sum over all values $h = 0, 1, \dots q - 1$ and then remove again the one value determined by (18), it follows:

\[
\eta = \sum^a r^a = \sum^{a'} \sum^h r^{a'(1 + hm')} - \sum^{a'} r^{a'kq},
\]

or, since $\sum^h r^{a'm'h} = 0$:

\[
\eta = -\sum^{a'} r^{a'kq}.
\]

Now $r^q$ is an $m'$th root of unity, and $\eta$ is therefore, apart from the sign, equal to a period formed from $m'$th roots of unity\footnote{Here the memoir of \gesperrt{Fuchs} is to be compared: Ueber die Perioden, welche aus den Wurzeln der Gleichung $\omega^n = 1$ gebildet sind, wenn $n$ eine zusammengesetzte Zahl ist. Crelle's Journal, Bd. 61 (1863).}.

\begin{center}
\S.\,21.

\gesperrt{Cyclotomic Fields with Given Group}.
\end{center}

In the foregoing, then, it has emerged that every cyclotomic period

\[
\eta = \sum^a r^a,
\]

in which $r$ is a primitive $m$th root of unity and $a$ runs through the numbers of a group $\mathfrak{A}$, when $\mathfrak{A}$ is a primary divisor of $\mathfrak{N}$
%CONT>>


% ===== scan leaf 101, printed 80 =====
% !PRINTED 80
\weberpage{80}

%<<CONT
is the root of an Abelian equation whose group is isomorphic with the group $\mathfrak{B}$ reciprocal to $\mathfrak{A}$.

The problem that we now wish to pose and solve is the following:

\gesperrt{I. The modulus $m$ and the group $\mathfrak{A}$ are to be determined in all possible ways so that $\mathfrak{B}$ has an arbitrarily given system of invariants.}

Or, what is the same thing:

\gesperrt{All cyclotomic fields of given group are to be determined.}

In this we always make the assumption that $\mathfrak{A}$ is to be a primary divisor of $\mathfrak{N}$, because otherwise we would obtain one and the same field several times. The problem splits into two parts, namely:

\gesperrt{II. Which moduli $m$ are suited to generate a group $\mathfrak{B}$ of given invariants?}

\gesperrt{III. How does one find this group $\mathfrak{B}$ and the associated group $\mathfrak{A}$?}

In this investigation we must alter the notation somewhat as against the earlier, so that clarity is not lost. The exceptional position of the prime 2, which always had to be taken into account in the considerations hitherto, made a somewhat cumbersome notation necessary. In what follows this distinction is no longer necessary to the same degree, and therefore the notation can now be simplified.

We denote the indices of a number $a$ from the group $\mathfrak{A}$ by

\[
\alpha_1, \alpha_2, \dots \alpha_\mu \quad \text{(indices of $a$)}
\]

and the corresponding index moduli by

\[
c_1, c_2, \dots c_\mu \quad \text{(index moduli)},
\]

so that $\mu$ is equal to the number of primes dividing $m$, or, if $m$ is divisible by 8, greater by 1.

Let further

\[
\omega_1, \omega_2, \dots \omega_\mu
\]

be primitive roots of unity of the degrees $c_1, c_2, \dots c_\mu$. Now if the indices of a number $b$ in the group $\mathfrak{B}$ reciprocal to $\mathfrak{A}$ are

\[
\beta_1, \beta_2, \dots \beta_\mu \quad \text{(indices of $b$),}
\]
%CONT>>


% ===== scan leaf 102, printed 81 =====
% !PRINTED 81
\weberpage{81}
%<<CONT
then the relation between the $\alpha$ and the $\beta$ is expressed by the
equation

\weq{1}{\omega_1^{\alpha_1\beta_1}\,\omega_2^{\alpha_2\beta_2}\,\dots\,\omega_\mu^{\alpha_\mu\beta_\mu}=1}

The invariants of the group $\mathfrak{B}$, which by \S.\,10 are all prime powers, we shall denote by
\[
i_1, i_2, \dots i_\nu \quad (\text{invariants of } \mathfrak{B})
\]
We take these invariants to be arbitrarily given
prime powers, and denote by $J$ their least common multiple, so that
\weq{2}{J=i_1i_1'=i_2i_2'=\dots=i_\nu i_\nu'}
can be set.

Moreover, $\mathfrak{A}$ shall be assumed to be a primary divisor of $\mathfrak{N}$, which by \S.\,19, 8. is equivalent to the condition:

\gesperrt{IV. None of the indices $\beta$ shall, for all numbers $b$, be divisible by $q$ when the corresponding prime $q$ divides $m$ more than once, or by $q-1$ when $q$ divides $m$ only once.}

By the definition of the invariants, the group $\mathfrak{B}$ must have a basis whose elements
\[
\begin{aligned}
&\phantom{\text{of the degrees}} \quad g_1, g_2, \dots g_\nu \quad (\text{basis of } \mathfrak{B})\\
&\text{of the degrees} \quad i_1, i_2, \dots i_\nu
\end{aligned}
\]
are, so that every number $b$ is contained once and only once in the form
\weq{3}{b\equiv g_1^{x_1}g_2^{x_2}\dots g_\nu^{x_\nu}\ (\mathrm{mod}\ m)}
when the exponents $x_1, x_2, \dots x_\nu$ each run through a complete system of residues with respect to the moduli $i_1, i_2, \dots i_\nu$.

All quantities of the form (3), if $g_1, g_2, \dots g_\nu$ are arbitrary numbers, certainly form a group under unrestricted variation of the integer exponents $x$. But if the $g_h$ are to be the elements of a basis of this group, and $i_h$ their degrees, then $1\equiv g_1^{x_1}g_2^{x_2}\dots g_\nu^{x_\nu}\ (\mathrm{mod}\ m)$ may be satisfied only when $x_h$ is divisible by $i_h$; hence:

\gesperrt{1. The necessary and sufficient condition that $g_1, g_2, \dots g_\nu$ be the elements of a basis of a group $\mathfrak{B}$ of the degrees $i_1, i_2, \dots i_\nu$ is that the congruence}
\weq{4}{g_1^{x_1}g_2^{x_2}\dots g_\nu^{x_\nu}\equiv 1\ (\mathrm{mod}\ m)}
%CONT>>


% ===== scan leaf 103, printed 82 =====
% !PRINTED 82
\weberpage{82}
%<<CONT
\gesperrt{holds then and only then, if at the same time the
congruences}
\weq{5}{x_1\equiv 0\ (\mathrm{mod}\ i_1),\quad x_2\equiv 0\ (\mathrm{mod}\ i_2),\quad \dots\ x_\nu\equiv 0\ (\mathrm{mod}\ i_\nu)}
\gesperrt{are satisfied.}

According to this, then, $J$ is the least positive number which, for all $b$, satisfies the congruence
\[
b^J\equiv 1\ (\mathrm{mod}\ m)
\]
or the least positive number for which, for all systems of the $\beta$,
\weq{6}{J\beta_1\equiv 0\ (\mathrm{mod}\ c_1),\quad J\beta_2\equiv 0\ (\mathrm{mod}\ c_2),\quad \dots\ J\beta_\mu\equiv 0\ (\mathrm{mod}\ c_\mu).}

From this there result the first conclusions about the composition of the number $m$.

a) \gesperrt{An odd prime $q$ which is not contained in $J$ can divide $m$ only simply.}

For if $m$ is divisible by $q^\varkappa$, then one of the quantities $c$ is, say, $c_1=\varphi(q^\varkappa)=q^{\varkappa-1}(q-1)$; hence if $\varkappa>1$, then $c_1$ is still divisible by $q$, and from (6) it then follows that all the $\beta_1$ would have to be divisible by $q$, which contradicts assumption IV.

b) \gesperrt{An odd prime $q$ which divides $J$ can be contained in $m$ at most once more than in $J$.}

For one can conclude in the same way that if $c_1=q^{\varkappa-1}(q-1)$, and $J$ is not divisible by $q^{\varkappa-1}$, all the $\beta_1$ would have to be divisible by $q$.

c) \gesperrt{If $J$ is odd, then $m$ must also be odd.}

For if $m$ is divisible by a power of 2, and hence at least by 4, then the index modulus corresponding to the number 2 is even, and the corresponding $\beta$ would by (6) all have to be even, contrary to requirement IV.

d) \gesperrt{If $J$ is divisible by a power of 2, then $m$ can contain the factor 2 at most twice more often than $J$.}

For if $m$ is divisible by $2^\lambda$ and $\lambda>2$, then two of the index moduli are
\[
c_1=2,\quad c_2=2^{\lambda-2}.
\]
%CONT>>


% ===== scan leaf 104, printed 83 =====
% !PRINTED 83
\weberpage{83}
%<<CONT
If, then, $J$ were not divisible by $2^{\lambda-2}$, all the $\beta_2$ would have to be divisible by 2, which again is contrary to requirement IV.

e) \gesperrt{If $q$ is a prime that divides $m$ simply, then $q-1$ must be divisible by at least one of the primes dividing $J$.}

For if $c_1=q-1$ were relatively prime to $J$, all the $\beta_1$ would have to be divisible by $q-1$, in contradiction with IV. The further conditions for $m$ will emerge later.

\gesperrt{Determination of the Group $\mathfrak{B}$:}

The group $\mathfrak{B}$ is determined when its basis $g_1, g_2, \dots g_\nu$ is given. We shall determine the elements of the basis by their indices, and accordingly introduce the following notation. Let
\weq{7}{
\begin{aligned}
&\gamma_{1,1}, \gamma_{1,2}, \dots \gamma_{1,\mu} \quad \text{the indices of } g_1\\
&\gamma_{2,1}, \gamma_{2,2}, \dots \gamma_{2,\mu} \quad \text{\textquotedbl{}} \quad \text{\textquotedbl{}} \quad \text{\textquotedbl{}} \quad g_2\\
&\dots\dots\dots\dots\dots\dots\dots\dots\dots\dots\dots\dots\dots\dots\\
&\gamma_{\nu,1}, \gamma_{\nu,2}, \dots \gamma_{\nu,\mu} \quad \text{\textquotedbl{}} \quad \text{\textquotedbl{}} \quad \text{\textquotedbl{}} \quad g_\nu.
\end{aligned}
}
The indices $\beta_1, \beta_2, \dots \beta_\mu$ of an arbitrary element $b$ are then obtained, by (3), in the form:
\weq{8}{
\begin{aligned}
\beta_1 &\equiv \gamma_{1,1}x_1+\gamma_{2,1}x_2+\dots+\gamma_{\nu,1}x_\nu\ (\mathrm{mod}\ c_1)\\
\beta_2 &\equiv \gamma_{1,2}x_1+\gamma_{2,2}x_2+\dots+\gamma_{\nu,2}x_\nu\ (\mathrm{mod}\ c_2)\\
&\dots\dots\dots\dots\dots\dots\dots\dots\dots\dots\dots\dots\dots\dots\\
\beta_\mu &\equiv \gamma_{1,\mu}x_1+\gamma_{2,\mu}x_2+\dots+\gamma_{\nu,\mu}x_\nu\ (\mathrm{mod}\ c_\mu),
\end{aligned}
}
and we now ask for the necessary and sufficient conditions on the numbers $\gamma_{h,k}$ in order that $g_1, g_2 \dots g_\nu$ be the basis of a group $\mathfrak{B}$ with the required properties. By 1., it is first of all necessary for this that:

\gesperrt{2. The numbers $\gamma_{h,k}$ must be so constituted that the congruences}
\weq{9}{
\begin{aligned}
\gamma_{1,1}x_1+\gamma_{2,1}x_2+\dots+\gamma_{\nu,1}x_\nu &\equiv 0\ (\mathrm{mod}\ c_1)\\
\gamma_{1,2}x_1+\gamma_{2,2}x_2+\dots+\gamma_{\nu,2}x_\nu &\equiv 0\ (\mathrm{mod}\ c_2)\\
&\dots\dots\dots\dots\dots\dots\dots\dots\dots\dots\dots\dots\\
\gamma_{1,\mu}x_1+\gamma_{2,\mu}x_2+\dots+\gamma_{\nu,\mu}x_\nu &\equiv 0\ (\mathrm{mod}\ c_\mu)
\end{aligned}
}
\gesperrt{are satisfied then and only then. when at the same time the congruences}
% [sic: Weber prints a full stop before the lower-case ``wenn''; a comma is grammatically required.]
\weq{10}{x_1\equiv 0\ (\mathrm{mod}\ i_1),\quad x_2\equiv 0\ (\mathrm{mod}\ i_2),\quad \dots\ x_\nu\equiv 0\ (\mathrm{mod}\ i_\nu)}
\gesperrt{hold.}


% ===== scan leaf 105, printed 84 =====
% !PRINTED 84
\weberpage{84}
If this condition is satisfied, then $g_1, g_2, \dots g_\nu$ is certainly the basis of a group with the invariants $i_1, i_2, \dots i_\nu$. But if the requirement set up in IV. is also to be satisfied, then it follows further:

\gesperrt{3. If $q_h$ is a prime dividing $m$, to which the index $\beta_h$ corresponds, then the $\nu$ quantities}
\[
\gamma_{1,h}, \gamma_{2,h}, \dots \gamma_{\nu,h},
\]
\gesperrt{must not, if $q_h$ divides $m$ more than once, all be divisible by $q_h$, and, if $q_h$ divides $m$ only once, must not all be divisible by $q_h-1$.}

The conditions 2., 3. are then accordingly the necessary and sufficient ones, and it now only remains to bring them into a form in which the possibility of their fulfilment can be judged.

We denote by $\delta_{k,h}$ the \gesperrt{greatest common divisor of $i_k$ and $c_h$ and set}
\weq{11}{i_k=\delta_{k,h}\,i_{k,h}, \quad c_h=\delta_{k,h}\,c_{k,h}, \quad \begin{aligned}k&=1,2,\dots\nu,\\ h&=1,2,\dots\mu,\end{aligned}}
so that $i_{k,h}$ and $c_{k,h}$ are \gesperrt{relatively prime}. It is then to be remarked that the quantities
\[
\delta_{k,h}, \quad i_{k,h}, \quad c_{k,h}
\]
are completely determined by $m$ and the invariants $i_k$. By the definition of the invariants, the $\delta_{k,h}$ and $i_{k,h}$ must be \gesperrt{prime powers}.

From the assumptions which we made in a) to d) about the number $m$, there results a consequence for the $c_{k,h}$ which we must emphasise.

\gesperrt{4. If $q_h^\varkappa$ is one of the factors of $m$ to which the index $\beta_h$ corresponds, so that $c_h=q_h^{\varkappa-1}(q_h-1)$, or, if $q=2$ and $\varkappa>2$, $c_h=2^{\varkappa-2}$, and if $\varkappa=2$, $c_h=2$, then, if $\varkappa>1$, the numbers}
\weq{12}{c_{1,h}, c_{2,h}, \dots c_{\nu,h}}
\gesperrt{cannot all be divisible by $q_h$, and, if $\varkappa=1$, cannot all be divisible by $q-1$.}
% [sic: Weber switches from q_h to q in the exceptional 2-adic clause and in the final q-1 clause; q_h is intended in both places.]

For if the quantities (12) are all divisible by $q_h$, then by (11) the numbers
\weq{13}{\delta_{1,h}.\ \delta_{2,h}, \dots \delta_{\nu,h}}
% [sic: Weber prints a full stop, rather than a comma, after the first term.]
%CONT>>


% ===== scan leaf 106, printed 85 =====
% !PRINTED 85
\weberpage{85}
%<<CONT
all not divisible by $q_h^{\varkappa-1}$, and the numbers
\weq{14}{i_{1,h},\ i_{2,h},\ \dots\ i_{\nu,h}}
are all not divisible by $q_h$; consequently, by (11), none of the
numbers $i_1, i_2, \dots i_\nu$, and hence also $J$, is divisible by $q_h^{\varkappa-1}$,
which contradicts the assumptions a) and b). For $q_h = 2$ one has
in this reasoning only to put $\varkappa - 1$ or $2$ (when $\varkappa = 2$) in place of $\varkappa$,
in order to obtain the same contradiction with c) and d).

    But if $\varkappa = 1$ and the numbers (12) are all divisible by $q_h - 1$,
then they must be equal to $q_h - 1$; the numbers (13) are
all $= 1$, and it would therefore follow that the numbers $i_k$, and
hence also $J$, were relatively prime to $q_h - 1$, which contradicts the hypothesis e).

    If we now return to the requirement stated in 2. for
the quantities $\gamma_{k,h}$, it results first of all that
the congruences (9) must be satisfied when $x_k = i_k$ and the
remaining $x = 0$ are put, thus:
\[
\gamma_{k,h}\, i_k \equiv 0\ (\mathrm{mod}\ c_h),
\]
and from this, taking account of (11):
\[
i_{k,h}\, \gamma_{k,h} \equiv 0\ (\mathrm{mod}\ c_{k,h}).
\]

    But because $i_{k,h}$ and $c_{k,h}$ are relatively prime, it follows that $\gamma_{k,h}$
must be divisible by $c_{k,h}$. We therefore introduce a new system
of integers $e_{k,h}$ by the equations
\weq{15}{\gamma_{k,h} = c_{k,h}\, e_{k,h}}
and now seek for these numbers the conditions following from 2. Since the numbers $\gamma_{k,h}$, as appears from their definition, are determined only with respect to the modulus $c_h$, $e_{k,h}$ is to be determined only with respect to
the modulus $\delta_{k,h}$. If we take one of the congruences (9):
\[
\gamma_{1,h}x_1 + \gamma_{2,h}x_2 + \cdots + \gamma_{\nu,h}x_\nu \equiv 0\ (\mathrm{mod}\ c_h),
\]
and introduce (15) into it, it takes the form
\weq{16}{c_{1,h}e_{1,h}x_1 + c_{2,h}e_{2,h}x_2 + \cdots + c_{\nu,h}e_{\nu,h}x_\nu \equiv 0\ (\mathrm{mod}\ c_h).}
Here we now put, according to (11):
\[
c_{k,h} = \frac{c_h}{\delta_{k,h}} = \frac{c_h\, i_{k,h}}{i_k}\,.
\]
Accordingly the congruences (16) are equivalent to
the requirement that
%CONT>>


% ===== scan leaf 107, printed 86 =====
% !PRINTED 86
\weberpage{86}
%<<CONT
\[
\frac{e_{1,h}i_{1,h}x_1}{i_1} + \frac{e_{2,h}i_{2,h}x_2}{i_2} + \cdots \frac{e_{\nu,h}i_{\nu,h}x_\nu}{i_\nu}
\]
must be integers, or, having regard to the notation (2), with the congruences:
\weq{17}{e_{1,h}i_{1,h}i'_1x_1 + e_{2,h}i_{2,h}i'_2x_2 + \cdots e_{\nu,h}i_{\nu,h}i'_\nu x_\nu \equiv 0\ (\mathrm{mod}\ J).}
These congruences are, as one sees, always satisfied when
$x_1$ is divisible by $i_1$, $x_2$ by $i_2 \dots$, $x_\nu$ by $i_\nu$, and
requirement 2. therefore now reduces to this:

    5. \gesperrt{that the $e_{k,h}$ are to be determined so that the congruences (17) shall be satisfied for no values of the $x$ other than those satisfying the congruences $x_k \equiv 0\ (\mathrm{mod}\ i_k)$.}

    To this is added, as a transformation of requirement 3.:

    6. The numbers
\[
c_{1,h}e_{1,h},\ c_{2,h}e_{2,h},\ \dots\ c_{\nu,h}e_{\nu,h}
\]
\gesperrt{must not all be divisible by $q_h$ when $\varkappa > 1$, and
by $(q_h - 1)$ when $\varkappa = 1$.}

    Requirement 6. contains, on account of 4., no impossibility.

    \gesperrt{By 5. and 6. the necessary and sufficient conditions for the numbers $e_{h,k}$ are expressed.}

    The question that now remains, therefore, is whether, and under what hypotheses, these requirements can be satisfied by a suitable choice of the $e_{h,k}$.
% [sic: Weber switches in these two sentences from the defined e_{k,h} to e_{h,k}; e_{k,h} is intended.]

    This question is greatly simplified by the fact that the
congruence (17) can be split into several. The invariants $i_k$
are, as we have assumed, prime powers. Let one of these
primes be denoted by $p$, and let
\weq{18}{i_1 = p^{\pi_1},\ i_2 = p^{\pi_2},\ \dots\ i_\varrho = p^{\pi_\varrho},}
while the remaining $i_k$, if there are still any, shall no longer
be divisible by $p$. The greatest of the exponents $\pi_1, \pi_2, \dots \pi_\varrho$ we shall denote by $\pi$.

    Then
\[
J = p^\pi J',
\]
and $J'$ is no longer divisible by $p$. Further,
\[
i'_1 = p^{\pi-\pi_1}J',\ i'_2 = p^{\pi-\pi_2}J',\ \dots\ i'_\varrho = p^{\pi-\pi_\varrho}J',
\]
while $i'_{\varrho+1}, \dots i'_\nu$ are divisible by $p^\pi$.


% ===== scan leaf 108, printed 87 =====
% !PRINTED 87
\weberpage{87}
Accordingly the congruence (17) has as a consequence:
\begin{gather}
e_{1,h}i_{1,h}p^{\pi-\pi_1}x_1 + e_{2,h}i_{2,h}p^{\pi-\pi_2}x_2 + \cdots e_{\varrho,h}i_{\varrho,h}p^{\pi-\pi_\varrho}x_\varrho \notag\\
\equiv 0\ (\mathrm{mod}\ p^\pi) \qquad h = 1, 2, \dots \mu, \tag{19}
\end{gather}
and the remaining primes dividing $J$ each yield just such a
system of congruences, which in their totality are completely
equivalent in meaning to the congruence (17). These
congruences (19) must not admit of being satisfied in any other way
than by the numbers
\[
p^{\pi-\pi_1}x_1,\ p^{\pi-\pi_2}x_2,\ \dots\ p^{\pi-\pi_\varrho}x_\varrho
\]
being divisible by $p^\pi$.

    If, then, we set for abbreviation:
\[
p^{\pi-\pi_1}x_1 = y_1,\ p^{\pi-\pi_2}x_2 = y_2,\ \dots\ p^{\pi-\pi_\varrho}x_\varrho = y_\varrho,
\]
then (19) takes the form:
\weq{20}{e_{1,h}i_{1,h}y_1 + e_{2,h}i_{2,h}y_2 + \cdots e_{\varrho,h}i_{\varrho,h}y_\varrho \equiv 0\ (\mathrm{mod}\ p^\pi),}
and these congruences are to have
\[
y_1 \equiv 0,\ y_2 \equiv 0,\ \dots\ y_\varrho \equiv 0\ (\mathrm{mod}\ p^\pi)
\]
as a necessary consequence.

    The system of linear congruences (20) can now be discussed
by the same means as a system of linear equations, namely by
means of determinants (cf. Chapter II of the first volume).

    For a better overview we shall write the congruences (20)
for the moment thus:
\begin{gather}
a_{1,1}y_1 + a_{2,1}y_2 + \cdots + a_{\varrho,1}y_\varrho \equiv 0 \notag\\
a_{1,2}y_1 + a_{2,2}y_2 + \cdots + a_{\varrho,2}y_\varrho \equiv 0\ (\mathrm{mod}\ p^\pi). \tag{21}\\
\dots\dots\dots\dots\dots\dots\dots\dots\dots\dots \notag\\
a_{1,\mu}y_1 + a_{2,\mu}y_2 + \cdots + a_{\varrho,\mu}y_\varrho \equiv 0 \notag
\end{gather}
If the coefficients $a_{k,h}$ are all divisible by $p$, then
these congruences are already satisfied when the $y$ are all only
divisible by $p^{\pi-1}$, and this case must therefore not occur.

    We shall now assume that from the matrix
\weq{22}{
\begin{matrix}
a_{1,1}, & a_{2,1}, & \dots & a_{\varrho,1} \\
a_{1,2}, & a_{2,2}, & \dots & a_{\varrho,2} \\
\dots & \dots & \dots & \dots \\
a_{1,\mu}, & a_{2,\mu}, & \dots & a_{\varrho,\mu}
\end{matrix}
}
%CONT>>


% ===== scan leaf 109, printed 88 =====
% !PRINTED 88
\weberpage{88}
%<<CONT
form a determinant of $\tau$ rows that is not divisible by $p$, while all determinants of more than $\tau$ rows,
if any such exist, are divisible by $p$. We may
assume that the indices are so arranged that
\[
\Delta = \sum \pm a_{1,1}a_{2,2} \dots a_{\tau,\tau}
\]
is not divisible by $p$. If we denote by $\Delta_{k,h}$ the subdeterminants of $\Delta$ and set
\[
\Delta_{k,1}a_{s,1} + \Delta_{k,2}a_{s,2} + \cdots \Delta_{k,\tau}a_{s,\tau} = D_{k,s},
\]
then $D_{k,s} = \Delta$ when $k = s$, and $= 0$ when $s \leqq \tau$ and is different from
$k$, and, when $s > \tau$, it is a $\tau$-rowed determinant of the matrix (22). With the aid of this notation, the following can
now be derived from the first $\tau$ congruences (21):
\begin{gather}
\Delta y_1 + D_{1,\tau+1}y_{\tau+1} + \cdots D_{1,\varrho}y_\varrho \equiv 0 \notag\\
\Delta y_2 + D_{2,\tau+1}y_{\tau+1} + \cdots D_{2,\varrho}y_\varrho \equiv 0\ (\mathrm{mod}\ p^\pi) \tag{23}\\
\dots\dots\dots\dots\dots\dots\dots\dots\dots \notag\\
\Delta y_\tau + D_{\tau,\tau+1}y_{\tau+1} + \cdots D_{\tau,\varrho}y_\varrho \equiv 0. \notag
\end{gather}
If for brevity we set
\[
A_{s,r} = \Delta a_{s,r} - \sum_{1,\tau}^{k} a_{k,r}D_{k,s} = \Delta a_{s,r} - \sum_{1,\tau}^{h}\sum_{1,\tau}^{k} \Delta_{k,h}a_{s,h}a_{k,r},
\]
then by Bd. I, \S.\,23, $A_{s,r}$ is a $(\tau+1)$-rowed determinant of the
matrix (22), and hence by hypothesis divisible by $p$.

    According to this notation we obtain from (23):
\[
\Delta (a_{1,r}y_1 + a_{2,r}y_2 + \cdots + a_{\varrho,r}y_\varrho)
\equiv A_{\tau+1,r}y_{\tau+1} + \cdots + A_{\varrho,r}y_\varrho\ (\mathrm{mod}\ p^\pi),
\]
and it thus appears that, when $\tau < \varrho$, the congruences (20)
are already satisfied when $y_{\tau+1}, \dots y_\varrho$ are assumed divisible by $p^{\pi-1}$, because all the $A_{s,r}$ are divisible by $p$. This,
then, is not admissible. Hence we must have $\tau = \varrho$, and this also
suffices, because (23) then reduces to
\[
\Delta y_1 \equiv 0,\ \Delta y_2 \equiv 0 \dots \Delta y_\varrho \equiv 0\ (\mathrm{mod}\ p^\pi)
\]
and consequently, since $\Delta$ is not divisible by $p$, the quantities $y_1, y_2, \dots y_\varrho$
must be divisible by $p^\pi$.

    Hereby the requirement 6. has been brought into another form, which, returning to our original notation
in (20), we can state as follows:
% [sic: Weber prints requirement 6; the determinant argument reformulates requirement 5, while requirement 6 is the separate primarity condition.]
%CONT>>


% ===== scan leaf 110, printed 89 =====
% !PRINTED 89
\weberpage{89}
%<<CONT

7. \gesperrt{The number} $\varrho$ \gesperrt{of the invariants divisible by a prime} $p$
\gesperrt{must be equal to or less than the number} $\mu$
\gesperrt{of the index moduli, and it must be possible to form from the matrix}

\weq{24}{
\begin{gathered}
e_{1,1}\,i_{1,1}, \quad e_{2,1}\,i_{2,1}, \; \dots \; e_{\varrho,1}\,i_{\varrho,1} \\
e_{1,2}\,i_{1,2}, \quad e_{2,2}\,i_{2,2}, \; \dots \; e_{\varrho,2}\,i_{\varrho,2} \\
\cdot\ \cdot\ \cdot\ \cdot\ \cdot\ \cdot\ \cdot\ \cdot\ \cdot\ \cdot \\
e_{1,\mu}\,i_{1,\mu}, \quad e_{2,\mu}\,i_{2,\mu}, \; \dots \; e_{\varrho,\mu}\,i_{\varrho,\mu}
\end{gathered}
}

\gesperrt{at least one determinant of} $\varrho$ \gesperrt{rows that is not
divisible by} $p$\gesperrt{.}

This involves first of all a requirement on the $i_{k,h}$, that is,
in the last instance a condition on the number $m$.

All terms of any $\varrho$-rowed determinant of the matrix (24)
contain a factor of the form

\[
i_{1,h_1}\,i_{2,h_2}, \; \dots \; i_{\varrho,h_\varrho},
\]

in which $h_1, h_2, \dots h_\varrho$ denote any combination of $\varrho$
distinct numbers of the sequence $1, 2, \dots \mu$. Now if all
these products were divisible by $p$, then certainly all determinants of the
matrix (24) of $\varrho$ rows would also be divisible by $p$, and so there
must be at least one such combination that is not divisible by $p$.
But since the $i_{k,h}$ can only be powers of $p$, the equalities

\[
i_{1,h_1} = 1, \quad i_{2,h_2} = 1, \; \dots \; i_{\varrho,h_\varrho} = 1
\]

must hold. If we now go back to the meaning of the $i_{k,h}$ in (11), we
can express this last requirement on the degree $m$ as
follows:

f) \gesperrt{If a prime} $p$ \gesperrt{divides} $\varrho$
\gesperrt{invariants} $i_1, i_2, \dots i_\varrho$\gesperrt{, then the number}
$\mu$ \gesperrt{of the index moduli must be equal to or greater than}
$\varrho$\gesperrt{, and it must be possible to select} $\varrho$ \gesperrt{of these
index moduli} $c_{h_1}, c_{h_2}, \dots c_{h_\varrho}$ \gesperrt{in such a way
that} $c_{h_1}$ \gesperrt{is divisible by} $i_1$, $c_{h_2}$
\gesperrt{by} $i_2 \dots$, $c_{h_\varrho}$ \gesperrt{by} $i_\varrho$ $^1$).

\noindent\rule{4cm}{0.4pt}

\footnotesize
$^1$) The proof of the fact that this requirement can always be
satisfied in infinitely many ways rests on the theorem
of number theory that among the numbers of an arithmetic progression
whose initial term and difference are without common divisor there occur
infinitely many primes -- a theorem which up to now has been established only by
\normalsize
%CONT>>


% ===== scan leaf 111, printed 90 =====
% !PRINTED 90
\weberpage{90}
%<<CONT

Herewith, however, everything is exhausted that we must require of $m$. For
when this condition f) is satisfied, we need only assume, e.\,g., the
$e_{1,h_1}, e_{2,h_2}, \dots e_{\varrho,h_\varrho}$ not divisible by $p$,
and the remaining $e$ divisible by $p$; then the determinant

\[
\begin{vmatrix}
e_{1,h_1}\,i_{1,h_1}, & \dots & e_{1,h_\varrho}\,i_{1,h_\varrho} \\
\cdot\ \cdot\ \cdot & \cdots & \cdot\ \cdot\ \cdot \\
e_{\varrho,h_1}\,i_{\varrho,h_1}, & \dots & e_{\varrho,h_\varrho}\,i_{\varrho,h_\varrho}
\end{vmatrix}
\equiv e_{1,h_1}\,e_{2,h_2} \dots e_{\varrho,h_\varrho} \pmod{p}
\]

is therefore certainly not divisible by $p$.

And under this assumption the condition 6. can also still be satisfied.
To see this, let $p^\pi$ denote any one of the invariants, say $i_1$, and
let $c_h$ denote the index modulus which by f) is divisible by $p^\pi$,
so that

\weq{25}{c_h = p^\pi c_{1,h}}

holds. If then $q_h$ is the prime divisor of $m$ belonging to the index
modulus $c_h$, then $p$ is either equal to $q_h$ or a divisor of $q_h - 1$.
By the assumption last made, $e_{1,h}$ is not divisible by $p$, while
$e_{2,h}, e_{3,h}, \dots e_{\nu,h}$ are divisible by $p$. The products

\[
c_{1,h}\,e_{1,h}, \quad c_{2,h}\,e_{2,h}, \; \dots \; c_{\nu,h}\,e_{\nu,h}
\]

are then not all to be divisible by $q_h$ or by $q_h - 1$.

If, first, $q_h$ is a simple divisor of $m$, then $c_h = q_h -
1$, and by (25) $c_{1,h}$ is not divisible by $q_h - 1$. One can
therefore still assume $e_{1,h}$ in such a way that $c_{1,h}\,e_{1,h}$ is not divisible by
$q_h - 1$. But if $q_h$ is a $\varkappa$-fold divisor of
$m$ and $\varkappa > 1$, and if $q_h = p$, then, when $p$ is odd,
$c_h = p^{\varkappa - 1}(p-1)$, and by the hypothesis b)
$\pi \leqq \varkappa - 1$. If then $c_h$ is to be divisible by $p^\pi$,
we must have $\pi = \varkappa - 1$, and (25) gives $c_{1,h} =
p - 1$. But if $p = 2$, then $c_h = 2$ or $= 2^{\varkappa - 1}$,
and hence by the hypothesis d) we must have $\pi = 1$ or $= \varkappa -
2$, and $c_{1,h} = 1$. One can therefore
% [sic: Weber prints pi <= kappa-1; hypothesis b) gives pi >= kappa-1, which is what makes divisibility force equality.]
% [sic: Weber prints c_h=2^{kappa-1} in the second 2-adic case; the index modulus is 2^{kappa-2}, as stated in Theorem 4, and that value is required for c_{1,h}=1.]

\noindent\rule{4cm}{0.4pt}

\footnotesize
\gesperrt{Dirichlet} with application of the integral calculus. (Abhandlungen
der Berliner Akademie 1837; \gesperrt{Dirichlet's Werke}, Bd.\,I, Nr.\,21;
\gesperrt{Dirichlet-Dedekind}, Vorlesungen \"uber Zahlentheorie, Supplement VI;
\gesperrt{Bachmann}, Analytische Zahlentheorie, Abschn.\,IV). If we retain the
notation used above, then, in order to satisfy the requirement f) for a
given $p$, one may take up $p^{\pi+1}$ as a factor in $m$. But then,
if $\varrho > 1$, a further $\varrho - 1$ prime factors
$q$ must occur in $m$, which are bound to congruences of the form $q \equiv 1
\pmod{i_1} \dots$.
\normalsize
%CONT>>


% ===== scan leaf 112, printed 91 =====
% !PRINTED 91
\weberpage{91}

%<<CONT
we can also in these cases take $e_{1,h}$ so that $c_{1,h}\,e_{1,h}$
does not become divisible by $q_h$.

The requirements herewith imposed on the $e_{k,h}$ thus consist only in
this, that they are to be divisible by certain primes and not divisible
by others, and so they are still compatible with one another in
infinitely many ways. But this last assumption about the $e_{k,h}$ had
only the purpose of showing that the earlier, more general requirements
can always be satisfied. There may still be many other ways of meeting
them.

\noindent\rule{4cm}{0.4pt}

\begin{center}
\gesperrt{Survey of the Results.}
\end{center}

If it is a question of representing all cyclotomic fields, and each one
only once, and indeed in the simplest possible way, by cyclotomic
periods, one proceeds as follows:

Take an arbitrary system of prime powers as invariants

\[
i_1, i_2, \dots i_v
\]

and then choose a number $m$ according to the following conditions.

If $p^\pi$ is the highest power of a prime $p$ that occurs among the
invariants, then take $p$ into $m$ at most $(\pi+1)$ times, or, if
$p = 2$, $(\pi+2)$ times.

A prime $q$ that does not divide the invariants at all may be taken
into $m$ only singly; but $m$ may have only such single prime factors
$q$ for which $q - 1$ is divisible by one of the prime factors $p$ of
the invariants. Hence $p$ itself, too, can occur singly in $m$ only if
$p - 1$ is divisible by one of the other $p$.

Further, $m$ must also satisfy the condition that among the index
moduli

\[
c_1, c_2, \dots c_\mu
\]

$\varrho$ distinct ones, say $c_1, c_2, \dots c_\varrho$, are divisible
by $i_1, i_2, \dots i_\varrho$, whenever $i_1, i_2, \dots i_\varrho$
are powers of one and the same prime. Then determine the quantities
$i_{k,h}$ and $c_{k,h}$ by the formulae (11) and divide the given
invariants into systems, so that all powers of one and the same prime
are united in one system.

For each of these systems form the matrix (24) and choose the numbers
$c_{k,h}$ so that from every such matrix a
% [sic: Weber prints c_{k,h}; these quantities were fixed by (11), whereas the matrix entries and the ensuing construction require choosing e_{k,h}.]
%CONT>>


% ===== scan leaf 113, printed 92 =====
% !PRINTED 92
\weberpage{92}
\pdfbookmark[2]{§ 22. Determination of the Group A}{bm26}

%<<CONT
determinant of $\varrho$ rows can be formed which is not divisible by the
prime in question, and that not all the products

\[
c_{1,h}\,e_{1,h}, \quad c_{2,h}\,e_{2,h}, \; \dots \; c_{\nu,h}\,e_{\nu,h}
\]

become divisible by $q_h$ or by $q_h - 1$. Then one sets

\[
\gamma_{k,h} = c_{k,h}\,e_{k,h},
\]

and thus has, by (7), the indices of a basis of an Abelian group
$\mathfrak{B}$, whose reciprocal group $\mathfrak{A}$ is a primary
divisor of the whole group $\mathfrak{N}$.

\noindent\rule{4cm}{0.4pt}

\begin{center}
\S.\,22.

\gesperrt{Determination of the group $\mathfrak{A}$.}
\end{center}

The final goal of this investigation is not so much the determination of
the group $\mathfrak{B}$ as that of the group $\mathfrak{A}$, from which
one can directly form the cyclotomic periods $\eta = \Sigma r^a$. This
problem can be solved in the following way: by (1), the indices $\alpha$
of the numbers in $\mathfrak{A}$ depend on the $\beta$ through the
equation

\weq{1}{\omega_1^{\alpha_1\beta_1}\,\omega_2^{\alpha_2\beta_2} \dots
\omega_\mu^{\alpha_\mu\beta_\mu} = 1,}

in which $\omega_1, \omega_2, \dots \omega_\mu$ denote fixed primitive
roots of unity of the degrees $c_1, c_2, \dots c_\mu$. But if one
understands by $C$ the least common multiple of $c_1, c_2, \dots c_\mu$
and sets

\weq{2}{C = c_1 c_1' = c_2 c_2' = \dots = c_\mu c_\mu',}

and denotes by $\omega$ a primitive $C$th root of unity, then instead of
(1) one can also set

\[
\omega^{c_1'\alpha_1\beta_1 + c_2'\alpha_2\beta_2 + \dots +
c_\mu'\alpha_\mu\beta_\mu} = 1,
\]

and therefore obtains for the $\alpha$ the congruence

\[
c_1'\alpha_1\beta_1 + c_2'\alpha_2\beta_2 + \dots +
c_\mu'\alpha_\mu\beta_\mu \equiv 0 \pmod{C},
\]

which is equivalent to the condition that the sums

\[
\frac{\alpha_1\beta_1}{c_1} + \frac{\alpha_2\beta_2}{c_2} + \dots +
\frac{\alpha_\mu\beta_\mu}{c_\mu} = \sum_{1,\mu}^{h} \frac{\alpha_h
\beta_h}{c_h}
\]

must be integers. If one substitutes in these for $\beta_k$ the
expressions \S.\,21, (8), it follows that also

\[
\sum^{k} x_k \sum^{h} \frac{\alpha_h\,\gamma_{k,h}}{c_h},
\]
%CONT>>


% ===== scan leaf 114, printed 93 =====
% !PRINTED 93
\weberpage{93}
%<<CONT
or, since the integers $x_k$ are arbitrary, the sums
\[
\sum_{1,\mu}^{h}\frac{\alpha_h\,\gamma_{k,h}}{c_h}\qquad k=1,2,\ldots \nu
\]
must be integers. If one then sets, according to (15) and (11) \S.\,21,
\[
\gamma_{k,h}=c_{k,h}\,e_{k,h},\quad c_h=\delta_{k,h}\,c_{k,h},\quad i_k=\delta_{k,h}\,i_{k,h},
\]
it follows that
\[
\sum^{h}\frac{\alpha_h\,e_{k,h}}{\delta_{k,h}}=\sum^{h}\frac{\alpha_h\,i_{k,h}\,e_{k,h}}{i_k}
\]
must be integers, and this requirement is equivalent to the system of congruences
\[
\sum^h \alpha_h\,i_{k,h}\,e_{k,h}\equiv 0\ (\mathrm{mod}\ i_k),
\]
or, written more fully,
\weq{3}{\begin{aligned}
&e_{1,1}i_{1,1}\alpha_1+e_{1,2}i_{1,2}\alpha_2+\cdots+e_{1,\mu}i_{1,\mu}\alpha_\mu\equiv 0\ (\mathrm{mod}\ i_1)\\
&e_{2,1}i_{2,1}\alpha_1+e_{2,2}i_{2,2}\alpha_2+\cdots+e_{2,\mu}i_{2,\mu}\alpha_\mu\equiv 0\ (\mathrm{mod}\ i_2)\\
&\cdots\cdots\cdots\cdots\cdots\cdots\cdots\cdots\cdots\cdots\cdots\cdots\cdots\cdots\cdots\cdots\cdots\cdots\cdots\cdots\\
&e_{\nu,1}i_{\nu,1}\alpha_1+e_{\nu,2}i_{\nu,2}\alpha_2+\cdots+e_{\nu,\mu}i_{\nu,\mu}\alpha_\mu\equiv 0\ (\mathrm{mod}\ i_\mu). % [sic: Weber prints i_\mu; the row pattern requires i_\nu.]
\end{aligned}}
By these congruences the numbers $\alpha_1, \alpha_2, \ldots \alpha_\mu$ are, as must be the case, determined only with respect to the moduli $c_1, c_2, \ldots c_\mu$, because $i_{k,h}c_h=i_k c_{k,h}$ is divisible by $i_k$; and \gesperrt{to find the indices of all numbers of the group $\mathfrak{A}$, one has simply to seek out all solutions of the congruences (3).}

\begin{center}
\rule{4cm}{0.4pt}
\end{center}


% ===== scan leaf 115, printed 94 =====
\weberpage{94}% page number not printed on this page
\pdfbookmark[1]{Fourth Section. Cubic and Biquadratic Abelian Fields.}{bm27}
\pdfbookmark[2]{§ 23. Cubic cyclotomic fields}{bm28}
% !PRINTED none
\begin{center}
\gesperrt{Fourth Chapter.}
\end{center}

\begin{center}
\gesperrt{Cubic and biquadratic Abelian fields.}
\end{center}

\begin{center}
\rule{3cm}{0.4pt}
\end{center}

\begin{center}
\S.\,23.
\end{center}

\begin{center}
\gesperrt{Cubic cyclotomic fields.}
\end{center}

The general investigations presented in the last sections admit manifold applications to special cases, of which the simplest shall be considered here. It is noteworthy in this connection that even the simplest cases, if one attacks them directly, lead almost into the same difficulties which we have overcome by the general investigation.

We shall first treat the problem of finding all \gesperrt{cubic cyclotomic fields, that is, all cyclotomic periods which satisfy a rational equation of degree $3$,} and indeed in such a way that, of the several expressions of which the same quantity is capable, the simplest is set down.

We have here only one invariant $i_1=3$, and according to the prescriptions in the summary of \S.\,21 there may be taken up into $m$ the factor 9, but not 3, and further primes in arbitrary number of the form $6N+1$, that is the primes 7, 13, 19, 31, 37, 43, 61, 67, 73, 79, 97 . . ., and so we thus obtain as suitable values of $m$ under 200:
\[
\begin{aligned}
m &= 7,\ 9,\ 13,\ 19,\ 31,\ 37,\ 43,\ 61,\ \mathbf{63},\ 67,\ 73,\ 79,\ \mathbf{91},\ 97,\ 103,\ 109,\\
&\quad \mathbf{117},\ 127,\ \mathbf{133},\ 139,\ 151,\ 157,\ 163,\ \mathbf{171},\ 181,\ 193,\ 199,
\end{aligned}
\]
among which the numbers with several prime divisors are distinguished by bold type. The smallest number $m$ which more than two primes divide is $7.9.13=819$.


% ===== scan leaf 116, printed 95 =====
% !PRINTED 95
\weberpage{95}
If, then, we understand by $q_1, q_2, q_3 \ldots$ numbers taken from the sequence consisting of the number 9 and the primes 7, 13, 19, 31 . . . , then the general expression for $m$ is:
\weq{1}{m = q_1\,q_2\,q_3\cdots q_\mu.}
The index moduli are, when $q_1=9$, $c_1=6$, $c_2=q_2-1$, $c_3=q_3-1 \ldots$ If 9 and 7 do not occur among the factors of $m$, the index modulus 6 drops out.

The quantities $\delta_{1,1}, \delta_{1,2}, \ldots \delta_{1,\mu}$ are the greatest common divisors of $i_1=3$ with the $c_1, c_2, \ldots c_\mu$; they are therefore all $=3$, and by \S.\,21, (11) it follows that:
\[
i_{1,h}=1,\quad c_{1,h}=\tfrac{1}{3}\,c_1.
\]
% [sic: Weber prints c_1 on the right; because h is arbitrary and c_{1,h}=c_h/delta_{1,h}, the intended subscript is h.]
The matrix \S.\,21, (24) here consists of only the single vertical row $e_{1,1}, e_{1,2}, \ldots e_{1,\mu}$, and these quantities may be chosen arbitrarily, provided only that at least one is among them that is indivisible by 3. These numbers $e_{1,1}, e_{1,2}, \ldots e_{1,\mu}$ are, moreover, determined only to the modulus 3, and may therefore be taken equal to 0, 1 or 2. But condition 6. of \S.\,21 also requires that, when $q_1=9$, $c_{1,1}e_{1,1}=2e_{1,1}$ shall not be divisible by 3, and, when $q_2$ is a prime, that
\[
c_{1,2}\,e_{1,2}=\tfrac{1}{3}(q_2-1)\,e_{1,2}
\]
shall not be divisible by $q_2-1$; that is to say, \gesperrt{none of the numbers} $e_{1,1}, e_{1,2}, \ldots e_{1,\mu}$ may be divisible by 3, and we can choose $\pm 1$ for each of them.

In order, then, to find the indices $\alpha_1, \alpha_2, \ldots \alpha_\mu$ of all the numbers $a$ that can occur in the period
\weq{2}{\eta = \sum r^a}
one has, by \S.\,22, (3), to take all the numbers $\alpha$ (to the modulus 3) that satisfy a congruence
\weq{3}{\pm\alpha_1\pm\alpha_2\pm\alpha_3\cdots\pm\alpha_\mu\equiv 0\ (\mathrm{mod}\ 3)}
where for \gesperrt{a definite group one combination of signs must be held fixed}.

The number of the possible cubic fields depends, for given $m$, only on the number $\mu$ of the index moduli, and, since one of the signs in (3) can be assumed arbitrarily, amounts to $2^{\mu-1}$.

The result would take an entirely similar shape if one were to put another prime $p$ in place of 3 for $i_1$. There would then merely enter into the congruence (3), in place of the double
%CONT>>


% ===== scan leaf 117, printed 96 =====
% !PRINTED 96
\weberpage{96}
%<<CONT
signs there enters in each case a complete system of residues of $p$ (with the exclusion of 0).

But we wish to remain with the assumption $i_1=3$, and to understand by $\varepsilon_1, \varepsilon_2, \ldots \varepsilon_\mu$ any combination of the numbers $\pm 1$. Then we can also write the congruence (3) thus
\weq{4}{\varepsilon_1\alpha_1+\varepsilon_2\alpha_2+\cdots+\varepsilon_\mu\alpha_\mu\equiv 0\ (\mathrm{mod}\ 3).}

The combination of signs $\varepsilon_1, \varepsilon_2, \ldots \varepsilon_\mu$, in which one sign can be assumed arbitrarily, since the simultaneous change of all the signs in (4) alters nothing, determines the individual group $\mathfrak{A}$, and evidently these groups $\mathfrak{A}$ are also all different from one another. For if, e.g., in one group $\varepsilon_1=\varepsilon_2=+1$, and in another $\varepsilon_1=+1$, $\varepsilon_2=-1$, then the first contains the combination of numbers:
\[
\alpha_1\equiv 1,\quad \alpha_2\equiv -1,\quad \alpha_3\equiv 0,\quad \ldots\ \alpha_\mu\equiv 0\ (\mathrm{mod}\ 3),
\]
which does not occur in the second.

In order to find the groups $\mathfrak{A}$ definitively, one takes as a basis a system of primitive roots $g_1, g_2, \ldots g_\mu$ of $q_1, q_2, \ldots q_\mu$ and determines $a$ to the modulus $m$ from the congruences:
\weq{5}{\begin{aligned}
a &\equiv g_1^{\alpha_1}\ (\mathrm{mod}\ q_1)\\
&\equiv g_2^{\alpha_2}\ (\mathrm{mod}\ q_2)\\
&\cdots\cdots\cdots\cdots\cdots\\
&\equiv g_\mu^{\alpha_m}\ (\mathrm{mod}\ q_\mu). % [sic: Weber prints \alpha_m; the pattern requires \alpha_\mu.]
\end{aligned}}
If we set, understanding by $\sigma$ one of the residues 0, $\pm 1$,
\weq{6}{\sigma \equiv \varepsilon_1\nu_1+\varepsilon_2\nu_2+\cdots+\varepsilon_\mu\nu_\mu\ (\mathrm{mod}\ 3),}
then $\mathfrak{A}$ consists of all numbers $a$ whose indices $\nu_1=\alpha_1, \nu_2=\alpha_2, \ldots \nu_\mu=\alpha_\mu$ satisfy the condition $\sigma=0$. The two other values $\sigma=1$, $\sigma=-1$ determine the two subsidiary groups
$\mathfrak{A}'$, $\mathfrak{A}''$ of $\mathfrak{A}$, so that
\[
\mathfrak{N} = \mathfrak{A} + \mathfrak{A}' + \mathfrak{A}''
\]
holds. If $a'$ and $a''$ are each a number from $\mathfrak{A}'$ and $\mathfrak{A}''$, then $\mathfrak{A}'=\mathfrak{A}a'$, $\mathfrak{A}''=\mathfrak{A}a''$.

We now denote by $\eta$ the cyclotomic period that belongs to $\mathfrak{A}$, and by $\eta'$, $\eta''$ the conjugate periods, thus
\weq{7}{\eta = \sum r^a,\quad \eta' = \sum r^{a'},\quad \eta'' = \sum r^{a''}.}

If then $\varrho$ is an imaginary third root of unity, we obtain the resolvent
\[
(\varrho, \eta) = \eta + \varrho\,\eta' + \varrho^2\,\eta'',
\]
%CONT>>


% ===== scan leaf 118, printed 97 =====
% !PRINTED 97
\weberpage{97}
%<<CONT
and for this, making use of the notation (6), we can set

\weq{8}{(\varrho, \eta) = \overset{n}{\Sigma}\, \varrho^{\sigma} r^n,}

wherein $n$ runs through all the numbers of the group $\mathfrak{N}$, and in $\sigma$ the
indices of $n$ are to be substituted each time for $\nu_1, \nu_2 \ldots \nu_\mu$.

Now let $r_1, r_2, \ldots r_\mu$ denote primitive roots of unity of the degrees
$q_1, q_2, \ldots q_\mu$, and let

\[ r = r_1 r_2, \ldots r_\mu \]

be set. Then the expression (8) can be decomposed into a product,
namely, if we set

\[ n \equiv n_h \ (\mathrm{mod}\ q_h) \]

into:

\weq{9}{(\varrho, \eta) = \Sigma\, \varrho^{\varepsilon_1 \nu_1} r_1^{n_1}\, \Sigma\, \varrho^{\varepsilon_2 \nu_2} r_2^{n_2} \ldots \Sigma\, r_\mu^{\varepsilon_\mu \nu_\mu} r_\mu^{n_\mu}.} % [sic: Weber prints r_\mu as the final first base; the parallel factors require \varrho.]

The factors of this product are now precisely the resolvents which we
investigated in \S.\,172 of the first volume, and we obtain, if we

\weq{10}{\overset{n_1}{\Sigma}\, \varrho^{\nu_h} r^{n_h} = (\varrho, \eta_h), \quad h = 1, 2, \ldots \mu} % [sic: Weber prints the upper bound n_1 and the base r^{n_h}; the general h factor requires n_h and r_h^{n_h}.]

set:

\weq{11}{(\varrho, \eta) = (\varrho^{\varepsilon_1}, \eta_1)(\varrho^{\varepsilon_2}, \eta_2) \ldots (\varrho^{\varepsilon_\mu}, \eta_\mu).}

If we now apply the formulae (5), (6), (25), (26) of the section just
cited, we obtain:

\weq{12}{(\varrho, \eta)(\varrho^{-1}, \eta) = m,}

\weq{13}{(\varrho, \eta)^3 = m(a + b\varrho), \quad (\varrho^2, \eta)^3 = m(a + b\varrho^2),}

wherein $a$ and $b$ are certain integers, which can be computed from the
formulae mentioned, but which turn out differently according to the choice
of the $\varepsilon$.

From (12) and (13) there then follows further:

\weq{14}{m = a^2 - ab + b^2.}

According to this we can set up the cubic equation whose roots are the
quantities $\eta, \eta', \eta''$. Let it read:

\weq{15}{\eta^3 - \alpha \eta^2 + \beta \eta - \gamma = 0,}

wherein $\alpha, \beta, \gamma$ are integers. As regards $\alpha$ first of
all, it is equal to the sum of all primitive $m$th
roots of unity, hence

\[ \alpha = \Sigma\, r_1\, \Sigma\, r_2 \ldots \Sigma\, r_\mu, \]

and this value is $= 0$ or $= (-1)^{\mu}$,
according as 9 occurs among the factors of $m$ or not. The two
other coefficients can be computed in just the same way as at the
%CONT>>


% ===== scan leaf 119, printed 98 =====
% !PRINTED 98
\weberpage{98}
%<<CONT
mentioned place. We wish here to arrange the calculation in a somewhat
altered form. We set:

\[ \xi = \eta - \frac{\alpha}{3}, \]

which has the consequence that $(\varrho, \eta) = (\varrho, \xi)$, on account of
$1 + \varrho + \varrho^2 = 0$. Then for $\xi$ there results the cubic
equation:

\weq{16}{\xi^3 + P\xi - Q = 0,}

wherein

\weq{17}{P = \beta - \frac{\alpha^2}{3}, \quad Q = \frac{2\alpha^3}{27} - \frac{\alpha\beta}{3} + \gamma.}

Now from (12) there results:

\[ \begin{aligned}
(\varrho, \xi)(\varrho^{-1}, \xi) = m &= \xi^2 + \xi'^2 + \xi''^2 - \xi\xi' - \xi\xi'' - \xi'\xi'' \\
&= -3P
\end{aligned} \]

or

\weq{18}{P = -\frac{m}{3},}

and from (13), after some simple calculations, if

\[ S = \xi^2\xi' + \xi'^2\xi'' + \xi''^2\xi, \quad S' = \xi\xi'^2 + \xi'\xi''^2 + \xi''\xi^2 \]

is set:

\[ \begin{aligned}
m(a + b\varrho) &= 9Q + 3\varrho S + 3\varrho^2 S' \\
m(a + b\varrho^2) &= 9Q + 3\varrho^2 S + 3\varrho S' \\
0 &= 9Q + 3S + 3S',
\end{aligned} \]

whence by addition:

\weq{19}{27Q = m(2a - b),}

and by subtraction of the first two:

\weq{20}{mb = 3(S - S') = 3(\xi - \xi')(\xi'' - \xi)(\xi'' - \xi').}

If $m$ is divisible by 9, then $\alpha = 0$, and $P$ and $Q$ are
\gesperrt{integers} and $P$ is divisible by 3. But if $m$ is not divisible by 9,
then only $3P$ and $27Q$ are integers, not divisible by 3.

If $D$ is the discriminant of the equation (15), and hence an integer, then

\weq{21}{\begin{aligned}
\sqrt{D} &= (\eta - \eta')(\eta'' - \eta)(\eta'' - \eta') = (\xi - \xi')(\xi'' - \xi)(\xi'' - \xi') \\
&= \frac{mb}{3}
\end{aligned}}

is also an integer, and hence, if $m$ is not divisible by 9, $b$ is
divisible by 3. Now by (14)
%CONT>>


% ===== scan leaf 120, printed 99 =====
% !PRINTED 99
\weberpage{99}
%<<CONT
\weq{22}{4m = (2a - b)^2 + 3b^2,}

and if $m$ is divisible by 9, then by (19) $2a - b$ is divisible by 3,
and from (22) it then follows that in this case too $b$ is divisible by 3.
We can therefore in all cases set

\[ 2a - b = A, \quad b = 3B, \]

so that $A, B$ are integers which satisfy the condition

\weq{23}{4m = A^2 + 27B^2}

and the equation for $\xi$ then becomes

\weq{24}{\xi^3 - \frac{m}{3}\xi - \frac{mA}{27} = 0,}

and for the discriminant of this equation there follows

\weq{25}{\sqrt{D} = mB.}

In order to carry out the computation for the simplest examples, we take for
$q_1 = 9, q_2 = 7, q_3 = 13$ from \S.\,172 of the first volume the values:

\[ \begin{aligned}
q_1 &= 9, & a + b\varrho &= 3\varrho, \\
q_2 &= 7, & a + b\varrho &= -(1 + 3\varrho), \\
q_3 &= 13, & a + b\varrho &= -(4 + 3\varrho).
\end{aligned} \]

From these there result for $m = 63$ the two values:

\[ \begin{aligned}
m = 63:\quad a + b\varrho &= -3\varrho(1 + 3\varrho) = 9 + 6\varrho, \\
a + b\varrho &= -3\varrho(1 + 3\varrho^2) = -9 - 3\varrho,
\end{aligned} \]

for

\[ \begin{aligned}
m = 91:\quad a + b\varrho &= (1 + 3\varrho)(4 + 3\varrho) = -5 + 6\varrho, \\
a + b\varrho &= (1 + 3\varrho)(4 + 3\varrho^2) = 10 + 9\varrho,
\end{aligned} \]

and for $m = 819$ the values:

\[ \begin{aligned}
a + b\varrho &= 3\varrho(1 + 3\varrho)(4 + 3\varrho) = -3(6 + 11\varrho), \\
&= 3\varrho(1 + 3\varrho)(4 + 3\varrho^2) = -3(9 - \varrho), \\
&= 3\varrho(1 + 3\varrho^2)(4 + 3\varrho) = 3(9 + 10\varrho), \\
&= 3\varrho(1 + 3\varrho^2)(4 + 3\varrho^2) = 3(6 - 5\varrho).
\end{aligned} \]

From these the cubic equations for $\eta$ in these three cases are found:

\[ \begin{aligned}
m = 63:\quad &\eta^3 - 21\eta - 28 = 0, & (\eta = \xi) \\
&\eta^3 - 21\eta + 35 = 0, \\
m = 91:\quad &\eta^3 - \eta^2 - 30\eta + 64 = 0, & (\eta = \xi + \tfrac{1}{3}) \\
&\eta^3 - \eta^2 - 30\eta - 27 = 0,
\end{aligned} \]
%CONT>>


% ===== scan leaf 121, printed 100 =====
% !PRINTED 100
\weberpage{100}
%<<CONT
\[ \begin{aligned}
m = 819:\quad &\eta^3 - 273\eta + 91 = 0, & (\eta = \xi) \\
&\eta^3 - 273\eta + 91 . 19 = 0, \\
&\eta^3 - 273\eta + 91 . 8 = 0, \\
&\eta^3 - 273\eta + 91 . 17 = 0.
\end{aligned} \]

The roots of these individual equations are best found from equation (9)
or (11), by arranging them according to powers of $\varrho$, making use,
however, only of $\varrho^3 = 1$, not of $\varrho^2 + \varrho + 1 = 0$.

Thus one obtains, e.\,g., for the case $m = 63$:

\[ \begin{aligned}
(\varrho, \eta) &= (\eta_1 + \varrho\eta_1' + \varrho^2\eta_1'')(\eta_2 + \varrho\eta_2' + \varrho^2\eta_2'') \\
\text{or} \\
&= (\eta_1 + \varrho\eta_1' + \varrho^2\eta_1'')(\eta_2 + \varrho\eta_2'' + \varrho^2\eta_2'),
\end{aligned} \]

in which

\[ \begin{aligned}
\eta_1 &= r_1 + r_1^{-1}, & \eta_1' &= r_1^2 + r_1^{-2}, & \eta_1'' &= r_1^4 + r_1^{-4} \\
\eta_2 &= r_2 + r_2^{-1}, & \eta_2' &= r_2^3 + r_2^{-3}, & \eta_2'' &= r_2^2 + r_2^{-2}
\end{aligned} \]

are to be set (Vol.\,I, \S.\,172), and one thus finds:

\[ \begin{aligned}
\eta &= \eta_1\eta_2 + \eta_1'\eta_2'' + \eta_1''\eta_2' \\
\text{or} \\
\eta &= \eta_1\eta_2 + \eta_1'\eta_2' + \eta_1''\eta_2''.
\end{aligned} \]

In order to find the groups $\mathfrak{A}$ from these, one has to determine
two numbers $x, y$ in such a way that $7x + 9y = 1$, and then

\[ r = r_1 r_2, \quad \text{thus} \quad r_1 = r^{7x}, \ r_2 = r^{9y} \]

is to be set. If one takes $x = 4, y = -3$, one obtains for $x$ the two values % [sic: Weber prints ``für x'', although x has just been fixed and values of \eta follow.]
of $\eta$:

\[ \begin{aligned}
\eta &= (r^{28} + r^{-28})(r^{27} + r^{-27}) + (r^7 + r^{-7})(r^9 + r^{-9}) \\
&\qquad + (r^{14} + r^{-14})(r^{18} + r^{-18}) \\
\eta &= (r^{28} + r^{-28})(r^{27} + r^{-27}) + (r^7 + r^{-7})(r^{18} + r^{-18}) \\
&\qquad + (r^{14} + r^{-14})(r^9 + r^{-9}),
\end{aligned} \]

hence

\[ \mathfrak{A} = \pm 1, \ \pm 8, \ \pm 2, \ \pm 16, \ \pm 4, \ \pm 32, \]

or

\[ \mathfrak{A} = \pm 1, \ \pm 8, \ \pm 11, \ \pm 25, \ \pm 5, \ \pm 23. \]

The two groups can be represented by the powers of 2 and of 11.

The cubic equations whose laws of formation we have now come to know
contain, together with their \gesperrt{Tschirnhausen transformations, all
cubic cyclotomic equations.} But these equations are also all
%CONT>>


% ===== scan leaf 122, printed 101 =====
% !PRINTED 101
\weberpage{101}
\pdfbookmark[2]{§ 24. Biquadratic cyclotomic fields}{bm29}
%<<CONT
different, and cannot be carried over into one another by \gesperrt{Tschirnhausen transformation}, because they belong to primary divisors of different full cyclotomic fields, or, if they are contained in the same full cyclotomic field, have different groups.

One can, to be sure, also form cyclotomic periods from other roots of unity, which lead to cubic equations. These periods, however, can be expressed by lower roots of unity, and are not primary. Thus one obtains, e.g., when $r$ is a 35th root of unity, a period of 8 terms whose exponents are formed from the powers of $-8$:

\[
\eta = r + r^{-1} + r^6 + r^{-6} + r^8 + r^{-8} + r^{13} + r^{-13},
\]
which therefore also satisfies a cubic equation.

But we have
\[
1 + r^7 + r^{-7} + r^{14} + r^{-14} = 0,
\]
and if one multiplies this equation by $(r^{15} + r^{-15})$:
\[
r^{15} + r^{-15} + r + r^{-1} + r^6 + r^{-6} + r^8 + r^{-8} + r^{13} + r^{-13} = 0,
\]
hence
\[
\eta = -(r^{15} + r^{-15}),
\]
which already occurs among the periods of the 7th roots of unity. This phenomenon is explained by the fact that $m = 35$ does not have the property required in \S.\,21, that $q - 1$ is divisible by 3 for all primes $q$ dividing $m$.

\begin{center}
\S.\,24.

\gesperrt{Biquadratic Cyclotomic Fields.}
\end{center}

With the biquadratic cyclotomic equations one has two kinds to distinguish. We can assume two invariants $i_1 = i_2 = 2$ or only one invariant $i_1 = 4$. We shall put the first case first.

For $m$ we have in this case, according to the summary in \S.\,21, the following conditions:

$m$ may be odd, or divisible by 4 or by 8, but by no higher power of 2. No odd prime can divide $m$ more than once, and there must be at least two index moduli present; hence, besides in
%CONT>>


% ===== scan leaf 123, printed 102 =====
% !PRINTED 102
\weberpage{102}
%<<CONT
the case $m = 8$, at least two different primes must be contained in $m$.

Let us set:
\[
m = q_1\, q_2 \cdots q_\mu,
\]
if $m$ is odd or divisible only by 4,
\[
m = q_2\, q_3 \cdots q_\mu,
\]
if $m$ is divisible by 8, so that $q_1 = 4$ may occur and in the second formula $q_2 = 8$, while the remaining $q$ are odd primes; then the index moduli in the three cases are:
\[
\begin{aligned}
& q_1 - 1,\quad q_2 - 1,\ \ldots,\ q_\mu - 1\\
& 2\phantom{{}_1 -{}1,}\ ,\quad q_2 - 1,\ \ldots,\ q_\mu - 1\\
& 2,\quad 2\phantom{{}_1 -{}1}\ ,\quad q_3 - 1,\ \ldots,\ q_\mu - 1.
\end{aligned}
\]
Since by this all the index moduli are divisible by 2, the conditions for $m$ are exhausted.

We have [\S.\,21, (11)]:
\weq{1}{\delta_{k,h} = 2,\ i_{k,h} = 1,\ c_{k,h} = \tfrac{1}{2} c_h,\ k = 1, 2,\ h = 1, 2, \ldots \mu.}
Now the quantities $e_{k,h}$ are to be determined in such a way that from the matrix
\weq{2}{\begin{array}{cccc}
e_{1,1}, & e_{1,2}, & \ldots & e_{1,\mu}\\
e_{2,1}, & e_{2,2}, & \ldots & e_{2,\mu}
\end{array}}
at least one odd two-rowed determinant can be formed.

In order that $\mathfrak{A}$ be a primary divisor of $\mathfrak{N}$, the condition \S.\,21, 6. is further added, where it is now to be distinguished whether $m$ is divisible by 8 or not. If $m$ is divisible by 8, the number of index moduli is greater by one than the number of prime factors of $m$, and there is an index modulus to which no prime factor of $m$ corresponds; in the other case the two numbers agree. According to \S.\,21, 6. the $e_{k,h}$ are to be chosen in such a way that in none of the pairs
\[
e_{1,1}, e_{2,1};\ e_{1,2}, e_{2,2};\ \ldots\ e_{1,\mu}, e_{2,\mu},
\]
to which a prime factor of $m$ corresponds are both numbers even. If $m$ is divisible by 8, there is among them a pair corresponding to the index modulus $c_k = 2$, on which no such restriction is imposed. If the numbers $e_{k,h}$ are (with respect to the modulus 2) so determined, then
%CONT>>


% ===== scan leaf 124, printed 103 =====
% !PRINTED 103
\weberpage{103}
%<<CONT
by \S.\,22 one obtains for the indices of the numbers in $\mathfrak{A}$ the two congruences:
\weq{3}{\begin{aligned}
e_{1,1}\alpha_1 + e_{1,2}\alpha_2 + \cdots + e_{1,\mu}\alpha_\mu &\equiv 0\\
e_{2,1}\alpha_1 + e_{2,2}\alpha_2 + \cdots + e_{2,\mu}\alpha_\mu &\equiv 0
\end{aligned}\ \ (\mathrm{mod}\ 2).}
To determine the number of possible groups $\mathfrak{A}$, one should bear in mind that the relations (3) can be replaced, without changing their content, by two mutually independent linear combinations. Accordingly one may take one of the pairs of numbers in the matrix (2), e.\,g. $e_{1,1}, e_{2,1} = 1, 0$, and thus put in place of (2):
\[
\begin{array}{cccc}
1, & e_{1,2}, & \ldots & e_{1,\mu}\\
0, & e_{2,2}, & \ldots & e_{2,\mu}.
\end{array}
\]
Now for the $e_{2,2} \ldots e_{2,\mu}$ one can put any combination of the numbers 0, 1, with the exception of the one in which all the $e_{2,h} = 0$. To every $e_{2,h} = 0$ the corresponding $e_{1,h} = 1$ must hold, while to an $e_{2,h} = 1$ there may correspond an $e_{1,h} = 0$ or $= 1$. The number of combinations arising in this way is:
\[
1 + (\mu - 1)2 + \frac{(\mu - 1)(\mu - 2)}{1.2} 2^2 + \cdots + 2^{\mu - 1} - 1 = 3^{\mu - 1} - 1.
\]
But now the second of the equations (3) can, without changing its meaning, also be added to the first, so that every two of the combinations obtained yield the same group. There thus remain
\weq{4}{\frac{3^{\mu - 1} - 1}{2}}
distinct combinations of the $e_{h,k}$, which also, as is easily seen, lead to distinct groups $\mathfrak{A}$. For suppose that the one time one has taken $e_{1,2}, e_{2,2} = 1, 0$, the other time $= 0, 1$; then $\alpha_1 = 0, \alpha_2 = 1, \alpha_3 = 0, \ldots \alpha_\mu = 0$ satisfies the first assumption, but does not satisfy the second. % [sic: Weber prints e_{h,k}; the defining matrix and congruences consistently use e_{k,h}.]

The number that we have just determined is exhaustive only when $m$ is not divisible by 8, because it is thereby assumed that we never have $e_{1,h}, e_{2,h} = 0, 0$. But if $m$ is divisible by 8, then this combination of values can occur in one pair, and we must therefore still add the cases which, through
\[
\begin{array}{ccccc}
1 & 0 & e_{1,3} & \ldots & e_{1,\mu}\\
0 & 0 & e_{2,3} & \ldots & e_{2,\mu}
\end{array}
\]
%CONT>>


% ===== scan leaf 125, printed 104 =====
% !PRINTED 104
\weberpage{104}
%<<CONT
are indicated, and whose number one finds, just as above, to be
\[
= \frac{3^{\mu - 2} - 1}{2}
\]
If, then, $m$ is divisible by 8, the total number of the groups $\mathfrak{A}$ is
\weq{5}{\frac{3^{\mu - 1} - 1}{2} + \frac{3^{\mu - 2} - 1}{2} = 2 . 3^{\mu - 2} - 1.}
The smallest values that $m$ can assume in this case are
\[
m = 8, 12, 15, 20, 21, 24, 28, 33, 35, 39, 40, 44 \ldots
\]
The smallest number that gives rise to more than one such group is $m = 24$, for which one has $\mu = 3$, and which therefore, by (5), gives five different groups. One obtains, namely, the following admissible and mutually distinct combinations of the $e_{h,k}$: % [sic: Weber again prints e_{h,k}; the displayed combinations and defining congruences use e_{k,h}.]
\[
\begin{array}{ccccc}
1\ 0\ 1 & 1\ 1\ 0 & 1\ 1\ 0 & 1\ 0\ 0 & 1\ 0\ 0\\
0\ 1\ 0, & 0\ 0\ 1, & 0\ 1\ 1, & 0\ 1\ 1, & 0\ 0\ 1,
\end{array}
\]
which leads to the following determinations concerning the indices:
\[
\begin{aligned}
& \alpha_1 \equiv \alpha_3,\quad \alpha_2 \equiv 0,\\
& \alpha_1 \equiv \alpha_2,\quad \alpha_3 \equiv 0,\\
& \alpha_1 \equiv \alpha_2 \equiv \alpha_3,\qquad (\mathrm{mod}\ 2)\\
& \alpha_1 \equiv 0,\quad \alpha_2 \equiv \alpha_3\\
& \alpha_1 \equiv 0,\quad \alpha_3 \equiv 0,
\end{aligned}
\]
and if one therefore determines the numbers $a$ by the congruences
\[
a \equiv (-1)^{\alpha_2} 5^{\alpha_1} (\mathrm{mod}\ 8), \quad a \equiv 2^{\alpha_3} (\mathrm{mod}\ 3)
\]
the following five two-term groups result:
\[
\begin{aligned}
\mathfrak{A}_1 &= 1,\ 5\\
\mathfrak{A}_2 &= 1,\ 19\\
\mathfrak{A}_3 &= 1,\ 11\\
\mathfrak{A}_4 &= 1,\ 23\\
\mathfrak{A}_5 &= 1,\ 7.
\end{aligned}
\]
To compute the periods $\eta$, we can proceed similarly as in the preceding section. We denote by $\nu_1, \nu_2, \ldots \nu_\mu$ the indices of a number $n$ and set
\[
\begin{aligned}
e_{1,1} \nu_1 + e_{1,2} \nu_2 + \cdots + e_{1,\mu} \nu_\mu &= \sigma_1,\\
e_{2,1} \nu_1 + e_{2,2} \nu_2 + \cdots + e_{2,\mu} \nu_\mu &= \sigma_2,
\end{aligned}
\]
%CONT>>


% ===== scan leaf 126, printed 105 =====
% !PRINTED 105
\weberpage{105}
%<<CONT
and obtain, if $\varepsilon_1$, $\varepsilon_2$ have the value 0 or 1, the resolvents

\weq{6}{\begin{aligned}
\Omega_{\varepsilon_1,\varepsilon_2} &= \sum^n (-1)^{\varepsilon_1\sigma_1+\varepsilon_2\sigma_2}\, r^n\\
&= \eta + (-1)^{\varepsilon_1}\eta' + (-1)^{\varepsilon_2}\eta'' + (-1)^{\varepsilon_1+\varepsilon_2}\eta'''.
\end{aligned}}

If now $m = q_1\, q_2 \ldots q_\mu$ is not divisible by 8, then, if $r_1, r_2, \ldots r_\mu$
denote primitive roots of the degrees $q_1, q_2, \ldots q_\mu$, we can
set
\[ r = r_1\, r_2 \ldots r_\mu \]
and further
\[ \begin{aligned}
n &\equiv n_1\ (\mathrm{mod}\ q_1)\\
&\equiv n_2\ (\mathrm{mod}\ q_2)\\
&\cdot\ \cdot\ \cdot\ \cdot\ \cdot\ \cdot\ \cdot\\
&\equiv n_\mu\ (\mathrm{mod}\ q_\mu),
\end{aligned} \]
from which it results:
\weq{7}{\Omega_{\varepsilon_1,\varepsilon_2} = \prod_{1,\mu}^{h} \sum^{n_h} (-1)^{\nu_h(\varepsilon_1 e_{1,h}+\varepsilon_2 e_{2,h})}\, r_h^{n_h}.}

The individual factors of this product are the \gesperrt{Gauss sums}
determined in \S.\,171 (Bd.\ I), which are there expressed by
square roots. Only when $q_1 = 4$ and
$\varepsilon_1\, e_{1,1} + \varepsilon_2\, e_{1,2} = 0$ is $\Omega_{\varepsilon_1,\varepsilon_2} = 0$. % [sic: Weber prints e_{1,2}; specializing the preceding general factor to h=1 requires e_{2,1}.]

But if $m$ is divisible by 8 and
\[ m = q_2\, q_3 \ldots q_\mu, \quad q_2 = 8, \]
then, setting $r = r_2\, r_3 \ldots r_\mu$, one obtains:
\[ \begin{aligned}
\Omega_{\varepsilon_1,\varepsilon_2} &= \sum^{n_2} (-1)^{\varepsilon_1(e_{1,1}\nu_1+e_{1,2}\nu_2)}(-1)^{\varepsilon_2(e_{2,1}\nu_1+e_{2,2}\nu_2)}\, r_2^{n_2}\\
&\times \prod_{3,\mu}^{h} \sum^{n_h} (-1)^{\nu_h(\varepsilon_1 e_{1,h}+\varepsilon_2 e_{2,h})}\, r_h^{n_h}.
\end{aligned} \]

In the five groups set up above, which belong to $m = 24$,
the periods can be computed from this formula, or also
easily directly, by setting
\[ r = e^{\frac{\pi i}{12}} = e^{\frac{\pi i}{3}}\, e^{-\frac{\pi i}{4}} \]
and one then obtains for the five quantities:
\[ \mathfrak{A}_1)\quad \eta = \frac{(1+i)\sqrt{3}}{\sqrt{2}} \]
\[ \mathfrak{A}_2)\quad \eta = \frac{\sqrt{3}-i}{\sqrt{2}} \]
%CONT>>


% ===== scan leaf 127, printed 106 =====
% !PRINTED 106
\weberpage{106}
%<<CONT
\[ \mathfrak{A}_3)\quad \eta = i\,\frac{\sqrt{3}-1}{\sqrt{2}} \]
\[ \mathfrak{A}_4)\quad \eta = \frac{1+\sqrt{3}}{\sqrt{2}} \]
\[ \mathfrak{A}_5)\quad \eta = \frac{1+i\sqrt{3}}{\sqrt{2}}. \]

Finally, the case of a single invariant $i_1 = 4$ shall now
be considered. The conditions for $m$ then consist simply in
this, that $m$ can be divisible by 4, 8, 16, but by no higher
power of 2, that odd primes can divide $m$ only once, and that
at least one of the index moduli must be divisible by 4. The
first values are thus
\[ m = 5,\ 13,\ 15,\ 16,\ 17,\ 20,\ 29,\ 33,\ 37,\ 39,\ 40\ \ldots \]

We shall assume that of the index moduli,
$c_1, c_2, \ldots c_\varrho$ are divisible by 4, and $c_{\varrho+1}, c_{\varrho+2}, \ldots c_\mu$ by 2, but
not by 4. Then $\varrho$ must be at least equal to 1,
and consequently
\[ \begin{aligned}
&\delta_{1,1} = \delta_{1,2} = \cdots = \delta_{1,\varrho} = 4, \quad \delta_{1,\varrho+1} = \cdots = \delta_{1,\mu} = 2\\
&i_{1,1} = i_{1,2} = \cdots = i_{1,\varrho} = 1, \quad i_{1,\varrho+1} = \cdots = i_{1,\mu} = 2\\
&c_{1,1} = \frac{c_1}{4}, \cdots\cdots, c_{1,\varrho} = \frac{c_\varrho}{4}, \quad c_{1,\varrho+1} = \frac{c_{\varrho+1}}{2}, \cdots, c_{1,\mu} = \frac{c_\mu}{2}.
\end{aligned} \]

Now the $e_{1,h}$ are to be determined so that of the numbers
\[ e_{1,1}, e_{1,2}, \ldots e_{1,\varrho} \]
at least one is odd, and that (on account of \S.\,21, 6.) of the
numbers
\[ e_{1,1}, e_{1,2}, \ldots e_{1,\varrho} \]
none is divisible by 4, and of the
\[ e_{1,\varrho+1}, e_{1,\varrho+2}, \ldots e_{1,\mu} \]
none is divisible by 2. An exception is again formed here,
in the case where $m$ is divisible by 8 or by 16, by the number $e_{1,h}$
which corresponds to none of the prime factors of $m$, and which can also
be even.

Then for the indices $\alpha$ of the numbers $a$ we obtain the congruence:
\[ e_{1,1}\alpha_1 + \cdots + e_{1,\varrho}\alpha_\varrho + 2\,(e_{1,\varrho+1}\alpha_{\varrho+1} + \cdots + e_{1,\mu}\alpha_\mu) \equiv 0\ (\mathrm{mod}\ 4). \]


% ===== scan leaf 128, printed 107 =====
% !PRINTED 107
\weberpage{107}
\pdfbookmark[2]{§ 25. Cubic Abelian equations}{bm30}
If we take as an example $m = 48$, we have two possibilities, namely:
\[ \begin{aligned}
e_{1,1} &= 1, & e_{1,2} &= 0, & e_{1,3} &= 1\\
e_{1,1} &= 1, & e_{1,2} &= 1, & e_{1,3} &= 1,
\end{aligned} \]
and from these one finds the two groups:
\[ \begin{aligned}
\mathfrak{A}_1 &= 1,\ 17,\ 31,\ 47\\
\mathfrak{A}_2 &= 1,\ \ 7,\ 41,\ 47.
\end{aligned} \]

\begin{center}
\S.\,25.
\end{center}

\begin{center}
\gesperrt{Cubic Abelian Equations.}
\end{center}

We shall conclude these considerations of the cubic and
biquadratic cyclotomic equations with the proof of a remarkable
theorem, which is suited to lend these equations
 a far higher and more general interest,
 and which is a special case of an entirely general theorem
that we shall come to know only later.

The theorem runs thus:

\gesperrt{All cubic and biquadratic Abelian
equations in the field of the rational numbers are
cyclotomic equations.}

The results of the two preceding sections thereby gain
 a higher value. For it then follows that
by the cyclotomic periods $\eta$ formed there \gesperrt{all Abelian
number fields} of the third and fourth degree are represented,
and that fields of this kind other than those discussed there do not
exist.

The possibility of proving this theorem for the special cases of the
cubic and biquadratic equations simply and without
further preparation rests on the fact that in the
fields of the third and fourth roots of unity $R(\varrho)$, $R(i)$ the same
 laws of the decomposition of numbers into their prime factors
hold as in the fields of the rational numbers, as is demonstrated in
\S\S.\,173, 174 of the first volume. It is therefore
also of interest to consider these special cases more closely,
because the course and the goal of the general proof
 to be supplied later are thereby more clearly recognised.


% ===== scan leaf 129, printed 108 =====
% !PRINTED 108
\weberpage{108}
Let us begin, then, with the cubic equations, and understand
by $x_0, x_1, x_2$ the roots of any \gesperrt{Abelian}
equation of the third degree. Since 3 is a prime, the
equation must be cyclic, i.\,e.\ \gesperrt{the cyclic functions
of the quantities} $x_0, x_1, x_2$ \gesperrt{are rational numbers.}

If we denote by $\varrho$ an imaginary third root of unity,
we have the two resolvents
\weq{1}{\begin{aligned}
(\varrho, x_0) &= x_0 + \varrho\, x_1 + \varrho^2 x_2\\
(\varrho^2, x_0) &= x_0 + \varrho^2 x_1 + \varrho\, x_2,
\end{aligned}}
whose cubes, then. must be numbers of the field $R(\varrho)$, and whose product % [sic: Weber prints a full stop after ``also'', although the sentence continues.]
must be a rational number.

Accordingly, denoting by $a, b, c$ rational (integer
or fractional) numbers, we set
\weq{2}{\begin{aligned}
(\varrho, x_0)^3 &= a + b\varrho\\
(\varrho^2, x_0)^3 &= a + b\varrho^2\\
(\varrho, x_0)(\varrho^2, x_0) &= c,
\end{aligned}}
and hence also
\weq{3}{(a+b\varrho)(a+b\varrho^2) = c^3.}

Now in the field $R(\varrho)$ we have, besides the six units
\[ \pm 1,\ \pm\varrho,\ \pm\varrho^2, \]
which can all be represented as powers of $-\varrho$, the
primes $\sqrt{-3} = \varrho - \varrho^2$, the real primes $q$ of the
form $3N+2$, and the two complex factors of the real
primes $p$ of the form $3N+1$. The decomposition of these
primes into the two complex factors $\pi, \pi'$:
\weq{4}{p = \pi\pi'}
we have presented in \S.\,172 of the first volume in the form
\weq{5}{p = \psi_1(\varrho)\,\psi_1(\varrho^2).}
We can therefore set
\weq{6}{\pi = \psi_1(\varrho), \quad \pi' = \psi_1(\varrho^2),}
and have thereby selected, among the various numbers associated
with $\pi$, one determinate number; and these numbers
$\pi, \pi'$ stand in a determinate relation to the cyclotomic
periods of the $p$th roots of unity $\eta$:
\weq{7}{(\varrho, \eta)^3 = p\pi, \quad (\varrho^2, \eta)^3 = p\pi'}
[Bd.\ I, \S.\,172, (5)].


% ===== scan leaf 130, printed 109 =====
% !PRINTED 109
\weberpage{109}

Now $a + b\varrho$ is an integral or fractional number of the field $R(\varrho)$.
If we decompose the numerator and denominator of this number into their prime
factors, cancel common factors, and if we understand by $q_1, q_2, \ldots$ real
primes of the form $3N+2$, and by $\pi_1, \pi_1', \pi_2, \pi_2', \ldots$
conjugate pairs of imaginary primes of the form (6), then we obtain

\weq{8}{\begin{aligned}
a + b\varrho &= (-\varrho)^\lambda (\sqrt{-3})^n q_1^{s_1} q_2^{s_2} \ldots \pi_1^{t_1} \pi_1'^{t_1'} \pi_2^{t_2} \pi_2'^{t_2'} \ldots \\
a + b\varrho^2 &= (-\varrho^2)^\lambda (-\sqrt{-3})^n q_1^{s_1} q_2^{s_2} \ldots \pi_1^{t_1'} \pi_1'^{t_1} \pi_2^{t_2'} \pi_2'^{t_2} \ldots
\end{aligned}}

wherein

\[
\lambda, \ n, \ s_1, \ s_2 \ldots t_1, \ t_1', \ t_2, \ t_2' \ldots
\]

are positive or negative integers.

But now by (3) the product of the two numbers (8) is the cube of a rational
number, and this gives, since a rational number is decomposable into prime
factors in only one way, the following conditions:

\weq{8}{\begin{aligned}
n &\equiv 0, \quad s_1 \equiv 0, \quad s_2 \equiv 0 \ldots \\
t_1+t_1' &\equiv 0, \quad t_2+t_2' \equiv 0 \ldots \quad \text{(mod 3),}
\end{aligned}} % [sic: Weber repeats equation number (8); (9) follows.]

and we therefore set

\weq{9}{\begin{aligned}
n &= 3\mu, \quad s_1 = 3\sigma_1, \quad s_2 = 3\sigma_2 \ldots \\
t_1+t_1' &= 3\tau_1, \quad t_2+t_2' = 3\tau_2 \ldots,
\end{aligned}}

so that $n, \sigma_1, \sigma_2 \ldots \tau_1, \tau_2 \ldots$ are integers. % [sic: Weber prints n, although n=3\mu has just been set and \mu is evidently intended.]
Now by the formulae (7):

\weq{10}{\begin{aligned}
\pi^t \pi'^{t'} &= p^{-(t+t')} (\varrho,\eta)^{3t} (\varrho^2,\eta)^{3t'} \\
\pi^{t'} \pi'^{t} &= p^{-(t+t')} (\varrho,\eta)^{3t'} (\varrho^2,\eta)^{3t}.
\end{aligned}}

If finally we further denote by $\varepsilon$ a $9$th root of unity, then we can
set

\[
-\varrho = (-\varepsilon)^3
\]

and thereby represent $a+b\varrho$ and $a+b\varrho^2$ as cubes. We thus
obtain by (2):

\weq{11}{(\varrho, x_0)^3 = (-\varepsilon)^{3\lambda} (\sqrt{-3})^{3\mu} q_1^{3\sigma_1} q_2^{3\sigma_2} \ldots p_1^{-3\tau_1} p_2^{-3\tau_2} \ldots H_1^3,}

wherein

\weq{12}{H_1 = (\varrho,\eta_1)^{t_1} (\varrho^2,\eta_1)^{t_1'} (\varrho,\eta_2)^{t_2} (\varrho^2,\eta_2)^{t_2'} \ldots}

is a cyclotomic number. If we denote by $H_2$ the number that arises from
$H_1$ by the interchange of $\varrho$ with $\varrho^2$, then $H_1 H_2$ is a
rational number.

If we extract the third root from (11), then it follows, since a third root of
unity $\varrho_1$ may here appear as a factor besides,

\weq{13}{(\varrho, x_0) = (-\varepsilon)^\lambda \varrho_1 (\sqrt{-3})^\mu q_1^{\sigma_1} q_2^{\sigma_2} \ldots p_1^{-\tau_1} p_2^{-\tau_2} \ldots H_1,}
%CONT>>


% ===== scan leaf 131, printed 110 =====
% !PRINTED 110
\weberpage{110}
\pdfbookmark[2]{§ 26. Biquadratic Abelian equations}{bm31}
%<<CONT
and in the same way we obtain

\weq{14}{(\varrho^2, x_0) = (-\varepsilon^2)^\lambda \varrho_2 (-\sqrt{-3})^u q_1^{\sigma_1} q_2^{\sigma_2} \ldots p_1^{-\tau_1} p_2^{-\tau_2} \ldots H_2.} % [sic: Weber prints exponent u; the conjugate of (13) requires \mu.]

Because the product of these two expressions must be rational, there results
further, between the roots of unity $\varrho, \varrho_1, \varrho_2$, the relation

\[
\varrho^\lambda \varrho_1 \varrho_2 = 1.
\]

If we set

\weq{15}{x_0 + x_1 + x_2 = A,}

then $A$ too is a rational number, and from (13), (14), (15) it follows that

\weq{16}{\begin{aligned}
x_0 &= 1/3 \left[A + (\varrho, x_0) + (\varrho^2, x_0)\right] \\
x_1 &= 1/3 \left[A + \varrho^2 (\varrho, x_0) + \varrho (\varrho^2, x_0)\right] \\
x_2 &= 1/3 \left[A + \varrho (\varrho, x_0) + \varrho^2 (\varrho^2, x_0)\right]
\end{aligned}}

are cyclotomic numbers. % [sic] printed "Kreistheilungzahlen", elsewhere on the page "Kreistheilungszahlen"
From the formulae (13), (14) we can see the composition of these numbers out of
the periods $\eta_1, \eta_2, \ldots$ But a closer examination of this
subject is no longer required, because already in \S.\,23 we have completely
constructed all cyclotomic numbers that satisfy cubic equations.

\noindent\hrulefill

\begin{center}
\S.\,26.

\gesperrt{Biquadratic Abelian Equations.}
\end{center}

The proof of the corresponding theorem for the biquadratic \gesperrt{Abelian} equations
can be carried out quite similarly. Only here two cases are to be
distinguished, namely equations with non-cyclic group, and
equations with cyclic group.

For the non-cyclic \gesperrt{Abelian} equations with the roots $x_0, x_1, x_2,
x_3$ the group is:

\weq{1}{1, \ \ (0,1)(2,3), \ \ (0,2)(1,3), \ \ (0,3)(1,2),}

and thus the three squares are

\weq{2}{\begin{aligned}
(x_0+x_1-x_2-x_3)^2 &= a \\
(x_0-x_1+x_2-x_3)^2 &= b \\
(x_0-x_1-x_2+x_3)^2 &= c
\end{aligned}}

and the product

\weq{3}{(x_0+x_1-x_2-x_3)(x_0-x_1+x_2-x_3)(x_0-x_1-x_2+x_3) = c,} % [sic: Weber prints RHS c, reusing c from the third square in (2); this rational product requires a distinct symbol.]
%CONT>>


% ===== scan leaf 132, printed 111 =====
% !PRINTED 111
\weberpage{111}
%<<CONT
and further the sum

\weq{4}{x_0 + x_1 + x_2 + x_3 = 4A}

are rational numbers.

If we extract the square roots from (2) and take into account that, by (3),
$\sqrt{a}$ differs from $\sqrt{bc}$ only by a rational factor, then by
addition it follows that:

\weq{5}{x_0 = A + B\sqrt{b} + C\sqrt{c} + D\sqrt{bc},}

in which $A, B, C, D$ are rational numbers; and from this one obtains $x_1,
x_2, x_3$ by changing the signs of $\sqrt{b}, \sqrt{c}$.

But by Band I, \S.\,171, all square roots can, if necessary with the aid of
$i$, which is itself a cyclotomic number, be expressed rationally through
the cyclotomic periods (the \gesperrt{Gauss} sums), so that the proof of
our theorem is already contained in (5).

\noindent\hrulefill

If the group of the biquadratic equation is cyclic, we can arrange the
roots so that the group consists of the permutations

\weq{6}{1, \ \ (0,1,2,3), \ \ (0,2)(1,3), \ \ (0,3,2,1)}

and then we are to set:

\weq{7}{\begin{aligned}
4A &= x_0+x_1+x_2+x_3 \\
(i,x_0) &= x_0+ix_1-x_2-ix_3 \\
(-1,x_0) &= x_0-x_1+x_2-x_3 \\
(-i,x_0) &= x_0-ix_1-x_2+ix_3;
\end{aligned}}

in these, $A$ and $(-1,x_0)^2=m$ are rational, and hence $(-1,x_0)$ can be
expressed through cyclotomic numbers. The fourth powers

\weq{8}{(i,x_0)^4 = a+bi, \quad (-i,x_0)^4 = a-bi}

are numbers of the field $R(i)$, and in order to be able to draw the same
conclusions as with the cubic equations, it thus only remains to represent
$a+bi$ and $a-bi$ as fourth powers of cyclotomic numbers. For this there
serves further the theorem that

\weq{9}{(i,x_0)(-i,x_0) = c}

is a rational number.


% ===== scan leaf 133, printed 112 =====
% !PRINTED 112
\weberpage{112}
In the field $R(i)$ we have (Bd. I, \S.\,173) the units $\pm 1, \pm i$, further,
as primes, the associated factors $1\pm i$ of 2, the real primes
$q$ of the form $4N+3$, and the two complex factors $\pi, \pi'$ of the real
primes $p$ of the form $4N+1$.

In Band I, \S.\,172 we found the decomposition

\[
p = \psi_1(i)\,\psi_1(-i)
\]

which permits us to set

\weq{10}{\pi=\psi_1(i), \quad \pi'=\psi_1(-i)}

and, if $\eta$ is an $N$-term period of $p$th
roots of unity,

\[
(i,\eta)^4 = p\,\psi_1(i)^2, \quad (-i,\eta)^4 = p\,\psi_1(-i)^2,
\]

or

\weq{11}{(i,\eta)^4=\pi^3\pi', \quad (-i,\eta)^4=\pi\pi'^3, \quad (i,\eta)(-i,\eta)=p.}

If we now decompose the numbers $a+bi$, $a-bi$ into their prime factors
in $R(i)$, as with the cubic equations, then, employing a corresponding
notation, we obtain

\weq{12}{\begin{aligned}
a+bi &= i^\lambda (1+i)^n q_1^{s_1} q_2^{s_2} \ldots \pi_1^{t_1} \pi_1'^{t_1'} \pi_2^{t_2} \pi_2'^{t_2'} \ldots \\
a-bi &= i^{-\lambda}(1-i)^n q_1^{s_1} q_2^{s_2} \ldots \pi_1^{t_1'} \pi_1'^{t_1} \pi_2^{t_2'} \pi_2'^{t_2} \ldots
\end{aligned}}

But the product of these two numbers must be the fourth power of a rational
number [by (8) and (9)], and from this it follows:

\weq{13}{\begin{aligned}
n &\equiv 0 & &\text{(mod 4)} \\
s_1 &\equiv 0, \quad s_2 \equiv 0 \ldots & &\text{(mod 2)} \\
t_1+t_1' &\equiv 0, \quad t_2+t_2' \equiv 0 \ldots & &\text{(mod 4)} \\
t_1-t_1' &\equiv 0, \quad t_2-t_2' \equiv 0 \ldots & &\text{(mod 2).}
\end{aligned}}

But now it follows from (11):

\weq{14}{\pi^t\pi'^{t'} = (\pi^3\pi')^{\frac{t-t'}{2}}(\pi\pi')^{\frac{-t+3t'}{2}} = (i,\eta)^{2(t-t')}\,p^{\frac{-t+3t'}{2}},}

and further

\weq{15}{(1+i)^n = (-1)^{\frac{n}{4}}\,2^{\frac{n}{2}},}

and accordingly, having regard to (13), we can therefore, if we denote by $H_1, H_2$
two cyclotomic numbers, namely the product of all the numbers occurring in the
decomposition of $a+bi$

\[
(i,\eta)^{\frac{t-t'}{2}},
\]
%CONT>>


% ===== scan leaf 134, printed 113 =====
% !PRINTED 113
\weberpage{113}
%<<CONT
and the number arising from it by the interchange of $i$ with $-i$,
and further by $c$ a rational number, we set
\[
\begin{aligned}
a + b i &= i^{\lambda}\, c^2 H_1^4\\
a - b i &= i^{-\lambda} c^2 H_2^4.
\end{aligned}
\]
From (8) it follows, by extraction of the $4$th root, whereby
a power of $i$ may enter as a factor,
\weq{15}{\begin{aligned}
(i, x_0) &= i^h \sqrt[4]{i}^{\lambda} \sqrt{c}\, H_1\\
(-i, x_0) &= i^{-h} \sqrt[4]{-i}^{\lambda} \sqrt{c}\, H_2\\
(-1, x_0) &= \sqrt{m},
\end{aligned}}
and from this, since $\sqrt[4]{i}$, $\sqrt{c}$ and $\sqrt{m}$ are cyclotomic
numbers, there results the correctness of our theorem in this case also.

Here too the further consideration of the structure of the expressions
occurring in (15) can serve to investigate more closely the composition
of the $x_0, x_1, x_2, x_3$ by the periods $\eta$; but this again can
lead to no other results than those already derived in \S.\,24.

We have thus herewith proved the theorem in general, \gesperrt{that
all Abelian fields of the third and fourth degrees
are cyclotomic fields, and that all these fields
can be represented rationally by the cyclotomic periods.}

\rule{3cm}{0.4pt}


% ===== scan leaf 135, printed 114 =====
\weberpage{114}% page number not printed on this page
\pdfbookmark[1]{Fifth Section. Constitution of General Groups.}{bm32}
\pdfbookmark[2]{§ 27. Construction of groups according to Cayley}{bm33}
% !PRINTED none
\weberpage{none}
\begin{center}
\gesperrt{Fifth Section.}

\gesperrt{Constitution of the General Groups.}
\end{center}

\rule{4cm}{0.4pt}

\begin{center}
\S.\,27.

\gesperrt{Formation of Groups according to Cayley.}
\end{center}

The general definition of the group, which was given in \S.\,1, still leaves much about the nature of this concept in the dark, and the various individual groups too, which we have come to know in the course of our considerations, give only hints of general laws, and show that the concept of a group contains nothing contradictory in itself. More is contained in the definition of the group than appears at first sight, and the number of possible groups that can be composed from a given number of elements is a very limited one. The general laws which prevail here have as yet been recognised only to the smallest part, so that every new special group, especially one of smaller number of terms, offers a new interest and invites detailed study.

\gesperrt{What groups are at all possible among a given number of elements, that is, for a given degree?} That is the general question at issue, from whose complete solution, however, we are still far removed. \gesperrt{Cayley} first took up this problem for the lowest degrees\footnote{On the theory of groups, as depending on the symbolic equation $\Theta^n = 1$. Philosophical Magazine, Vol.\ II, 1854. (\gesperrt{Cayley's} mathematical papers, Vol.\ II, 125.) American Journal of mathematics, Vol.\ I.}.


% ===== scan leaf 136, printed 115 =====
% !PRINTED 115
\weberpage{115}
For any arbitrary degree number $n$ we always have one, namely the cyclic group, which we obtain if we take $a^n = 1$ and assume the $n$ elements
\[
1, a, a^2 \ldots a^{n-1}
\]
to be distinct.

If $n$ is a prime, then this is the only group of degree $n$ (if isomorphic groups are regarded as identical). For the degree of every element of a group is a divisor of the degree of the group, and hence in a group of prime degree every element, with the exception of the unit element, has degree $n$.

To represent a group of degree $n$ completely, one would have to construct a square table with $n^2$ cells; arranged in $n$ rows and $n$ columns. One designates each row and each column by one of the given elements, and places in the cell in which these two intersect the composed element, where, say, the row indicator denotes the first component and the column indicator the second:
\[
\begin{array}{c!{\vrule width 0.8pt}c|c|c|c}
 & 1 & \alpha & \beta & \gamma \ldots \\
\noalign{\hrule height 0.8pt}
1 & 1 & \alpha & \beta & \gamma \ldots \\ \hline
\alpha & \alpha & \alpha^2 & \alpha\beta & \alpha\gamma \ldots \\ \hline
\beta & \beta & \beta\alpha & \beta^2 & \beta\gamma \ldots \\ \hline
\gamma & \gamma & \gamma\alpha & \gamma\beta & \gamma^2 \ldots \\
\multicolumn{5}{c}{\cdots\cdots\cdots\cdots\cdots\cdots\cdots\cdots\cdots\cdots}
\end{array}
\]

But the cells of such a table cannot be filled in with the elements quite arbitrarily; rather, one must keep to the associative law, so that, when one looks up $\alpha\beta\gamma$ by first finding in the row $\alpha$ and the column $\beta$ the element $(\alpha\beta)$, and then in the row $(\alpha\beta)$ and the column $\gamma$ the element $(\alpha\beta)\gamma$, one finds the same result as when one first looks up $(\beta\gamma)$ in the row $\beta$ and the column $\gamma$, and then $\alpha(\beta\gamma)$ in the row $\alpha$ and the column $(\beta\gamma)$.

For $n = 4$ we have, if an element of degree $4$ exists, the cyclic group
\[
1, \alpha, \alpha^2, \alpha^3,
\]
%CONT>>


% ===== scan leaf 137, printed 116 =====
% !PRINTED 116
\weberpage{116}
%<<CONT
and if all elements are of degree $1$ or $2$, the groups
\[
1, \alpha, \beta, \alpha\beta
\]
with the condition $\alpha\beta = \beta\alpha$. There are therefore only these two groups of degree $4$. If the elements are denoted by $1, \alpha, \beta, \gamma$, we have the two tables:
\[
\begin{array}{c|cccc}
 & 1 & \alpha & \beta & \gamma \\
\hline
1 & 1 & \alpha & \beta & \gamma \\
\alpha & \alpha & \beta & \gamma & 1 \\
\beta & \beta & \gamma & 1 & \alpha \\
\gamma & \gamma & 1 & \alpha & \beta
\end{array}
\qquad
\begin{array}{c|cccc}
 & 1 & \alpha & \beta & \gamma \\
\hline
1 & 1 & \alpha & \beta & \gamma \\
\alpha & \alpha & 1 & \gamma & \beta \\
\beta & \beta & \gamma & 1 & \alpha \\
\gamma & \gamma & \beta & \alpha & 1
\end{array}
\]

The case $n = 6$ can be completely settled in the following way: We denote the six elements by $1, \alpha, \beta, \gamma, \delta, \varepsilon$, so that 1 is the unit of the group. If among them there occurs an element of degree 6, the group is cyclic.

If we therefore set this case aside, the degrees of the elements $\alpha, \beta, \gamma, \delta, \varepsilon$ can only be 2 or 3. If $\alpha$ and $\beta$ are of degree $2$, then $\alpha\beta$ cannot be of degree $2$; for otherwise we should have
\[
\alpha = \alpha^{-1}, \quad \beta = \beta^{-1}, \quad \alpha\beta = \beta^{-1}\alpha^{-1} = \beta\alpha,
\]
and $1, \alpha, \beta, \alpha\beta$ would therefore be a group of degree $4$; but such a group cannot be a divisor of a group of degree $6$.

There must therefore occur at least one element of degree $3$, and if we denote such an element by $\alpha$, and $\gamma$ is not contained in $1, \alpha, \alpha^2$, the group arranges itself thus:
\weq{1}{S = 1, \alpha, \alpha^2, \gamma, \gamma\alpha, \gamma\alpha^2.}
To determine the condition that this be a group, we form
\[
\gamma S = \gamma, \gamma\alpha, \gamma\alpha^2, \gamma^2, \gamma^2\alpha, \gamma^2\alpha^2,
\]
which must be identical with $S$; and therefore
\[
\gamma^2 = 1 \text{ or } = \alpha \text{ or } = \alpha^2
\]
must hold. But if $\gamma^2 = \alpha$ or $= \alpha^2$, then $\gamma^3 = \alpha\gamma$ or $= \alpha^2\gamma$, and $\gamma^6$, but no lower power of $\gamma = 1$, i.\,e. $S$ is
%CONT>>


% ===== scan leaf 138, printed 117 =====
% !PRINTED 117
\weberpage{117}
\pdfbookmark[2]{§ 28. Relation of general groups to permutation groups}{bm34}
%<<CONT
cyclic ($S = 1, \gamma, \gamma^2, \gamma^3, \gamma^4, \gamma^5$). There remains, then, only the assumption that $\gamma$ is of degree $2$, so that
\weq{2}{\gamma^2 = 1}
holds.
Since, however, we can replace $\gamma$ by $\gamma\alpha$ or $\gamma\alpha^2$, these too must be of degree $2$, and it follows that:
\weq{3}{\gamma\alpha = \alpha^2\gamma, \quad \gamma\alpha^2 = \alpha\gamma,}
with whose help one can always reduce every compositum of arbitrarily many powers of $\alpha$ and $\gamma$ to one and only one of the elements $S$. And if one now denotes
\[
1, \alpha, \alpha^2, \gamma, \gamma\alpha, \gamma\alpha^2
\]
by
\[
1, \alpha, \beta, \gamma, \delta, \varepsilon
\]
one obtains the following table:
\[
\begin{array}{c!{\vrule width 0.8pt}c|c|c|c|c|c}
 & 1 & \alpha & \beta & \gamma & \delta & \varepsilon \\
\noalign{\hrule height 0.8pt}
1 & 1 & \alpha & \beta & \gamma & \delta & \varepsilon \\ \hline
\alpha & \alpha & \beta & 1 & \varepsilon & \gamma & \delta \\ \hline
\beta & \beta & 1 & \alpha & \delta & \varepsilon & \gamma \\ \hline
\gamma & \gamma & \delta & \varepsilon & 1 & \alpha & \beta \\ \hline
\delta & \delta & \varepsilon & \gamma & \beta & 1 & \alpha \\ \hline
\varepsilon & \varepsilon & \gamma & \delta & \alpha & \beta & 1
\end{array}
\]

In every row and in every column, as must be the case, all the elements occur.

\begin{center}
\S.\,28.

\gesperrt{Relation of the general groups to the permutation groups.}
\end{center}

The groups that we have hitherto employed most are the permutation groups; this kind of groups acquires a heightened significance for general group theory as well through the following considerations.

Let
\weq{1}{P = a_0, a_1, a_2, \ldots, a_{n-1}}
be an arbitrary group of degree $n$. If we take from $P$ any
%CONT>>


% ===== scan leaf 139, printed 118 =====
% !PRINTED 118
\weberpage{118}
%<<CONT
an element $b$, then the complex $Pb$ is completely identical with $P$. The two sequences
\[
A = a_0,\ a_1,\ a_2,\ \ldots,\ a_{n-1}
\]
\[
Ab = a_0b,\ a_1b,\ a_2b,\ \ldots,\ a_{n-1}b
\]
can therefore differ from one another only in their arrangement. The
passage from $A$ to $Ab$ is thus a permutation of $n$ figures
$0, 1, 2 \ldots n - 1$, and to every element $b$ of $P$ there corresponds such
a permutation, which we shall denote by $\pi_b$. To two different
elements $b, c$ there always correspond two different
permutations $\pi_b, \pi_c$, and to the unit element there corresponds the
identical permutation. If we set
\[
\pi_b = (A,\ Ab), \quad \pi_c = (A,\ Ac),
\]
then we can also denote $\pi_c$ by $(Ab, Abc)$, and from this
it follows that
\weq{2}{\pi_b \pi_c = \pi_{bc},}
i.\,e. the permutations $\pi$ form a permutation group isomorphic with $P$. This permutation group is transitive, since, for example,
the element $a_0$ can pass into any other element of $P$ whatever. We thus have the theorem:

1.\ \gesperrt{Every group of degree $n$ is isomorphic with a
transitive permutation group of $n$ figures.}

There are of course still other permutation groups isomorphic
with a given group, and it would be of particular
interest to form such a group with the smallest possible number of figures.
This problem cannot at present be solved in general.
We must content ourselves here with a few theorems.

We assume that $R$ is any divisor of $P$ of
index $j$, and denote the elements of $R$ by $c$; further, decomposing
$P$ into a system of subsidiary groups, we set
\weq{3}{P = R + Rb_1 + Rb_2 + \cdots + Rb_{j-1}.}
If then $a$ is any element of $P$, the two
systems
\weq{4}{\begin{aligned}
B &= R,\ Rb_1,\ Rb_2,\ \ldots,\ Rb_{j-1}\\
Ba &= Ra,\ Rb_1a,\ Rb_2a,\ \ldots,\ Rb_{j-1}a,
\end{aligned}}
agree, apart from their order, and so to every element $a$ there corresponds a definite permutation $\pi_a$ of the
$j$ figures $0, 1 \ldots j - 1$, which we denote by
\weq{5}{\pi_a = (B,\ Ba)}
%CONT>>


% ===== scan leaf 140, printed 119 =====
% !PRINTED 119
\weberpage{119}
%<<CONT
can denote. Here too the law
\weq{6}{\pi_a \pi_{a'} = \pi_{aa'},}
holds, if $a'$ is likewise any element from $P$. There remains
the question whether different elements $a$ can give rise to the
same permutation $\pi$. If $\pi_a = \pi_{a'}$, then it follows from (6):
\[
\pi_{a'a^{-1}} = 1,
\]
and if therefore we denote by $a_0$ all the elements satisfying
the condition
\weq{7}{\pi_{a_0} = 1}
then $a' = a_0 a$, and it comes down to determining the totality
of the elements $a_0$. But if
\weq{8}{Rba_0 = Rb}
is to hold, then
\[
ba_0 = cb
\]
or
\[
a_0 = b^{-1}cb
\]
must hold, i.\,e. $a_0$ must belong to the group $b^{-1}Rb$, and this is
also sufficient for (8) to hold.

From this it follows that the elements $a_0$ satisfying the
condition (7) fill the whole group $R_0$, which is the intersection
of all the divisors of $P$ conjugate to $R$:
\[
R,\ b_1^{-1}Rb_1,\ b_2^{-1}Rb_2,\ \ldots,\ b_{j-1}^{-1}Rb_{j-1},
\]
and, as we have proved earlier, is a normal divisor
of $P$.

From this there results, for the case that $R_0 = 1$, the theorem:

2.\ \gesperrt{If a group $P$ has a divisor of index $j$
which has no element other than the unit
element in common with its conjugate
divisors, then $P$ is (simply) isomorphic with a
transitive permutation group of $j$ figures.}

That the permutation group is transitive is seen from (4),
where each of the elements $b_1, b_2, \ldots, b_{j-1}$ can be put for $a$.

Theorem 2. always takes place when $P$ is a \gesperrt{simple group}
and $R$ a proper divisor of $P$, because then $R_0$, as a normal divisor of $P$, must reduce to the unit group\footnotemark.

Now let $R$ and $Q$ denote any two divisors of the group $P$,
$R$ of index $j$, $Q$ of degree $q$.
\footnotetext{\gesperrt{Hölder}, Mathematische Annalen, Bd.\ 40, S.\ 57.}


% ===== scan leaf 141, printed 120 =====
% !PRINTED 120
\weberpage{120}
We decompose $P$ as in (4) into the $j$ subsidiary groups $B$, denote
by $e$ all the elements of $Q$, and now investigate the permutation
\[
\pi_e = (B,\ Be),
\]
which is called forth by the elements $e$ among the $j$ elements (4).
By (5) and (6), these permutations certainly form a group $\Pi$
simply or multiply isomorphic with $Q$.

If $a_1$ is any element of $P$, then in the complex
\[
P_1 = Ra_1Q
\]
formed according to the composition of the parts, a certain number of the
subsidiary groups $B$ will be contained. We find this
number by seeking the condition that
$Ra_1 = Ra_1e$, or that $a_1e$ be contained in $Ra_1$, or $e$ in $a_1^{-1}Ra_1$.
If, then,
\[
Q_1 \text{ the intersection of } Q \text{ and } a_1^{-1}Ra_1
\]
is of index $h_1$ with respect to $Q$, then $P_1$ contains exactly $h_1$ and not
more of the complexes $B$.

If $a_2$ is a second element of $P$, then the complex
\[
P_2 = Ra_2Q
\]
either coincides entirely with $P_1$, or the two contain
no common element at all. For if $P_1$ and $P_2$ contain a
common element, then they have one of the complexes $B$ in common. But that is possible only if $a_2e$ is contained in $Ra_1e'$,
and hence $a_2$ in $P_1$; and then $P_1$ and $P_2$ are identical.
We can therefore continue forming the complexes $P_1, P_2 \ldots$ until the group $P$ is exhausted, and obtain the decomposition of the group $P$ with respect to the two groups $Q$ and $R$:
\weq{9}{\begin{aligned}
P &= Ra_1Q + Ra_2Q + \cdots\\
&= P_1 \phantom{{}={}} + P_2 \phantom{{}={}} + \cdots
\end{aligned}}

From the remark that, if $a_2$ does not occur in $Ra_1Q$,
then $a_2^{-1}$ is not contained in $Qa_1^{-1}R$, there results the second
decomposition:
\weq{10}{P = Qa_1^{-1}R + Qa_2^{-1}R + \cdots\footnotemark.}

The parts $P_1, P_2, \ldots$ of (9) contain respectively $h_1, h_2, \ldots$
elements of the sequence $B$, if $h_1, h_2, \ldots$ denote the indices of the intersections of $Q$ with $a_1^{-1}Ra_1, a_2^{-1}Ra_1, \ldots$ with respect to $Q$, and hence % [sic: Weber prints a_2^{-1}Ra_1; the conjugate pattern requires a_2^{-1}Ra_2.]
divisors of the degree $q$ of $Q$.
\footnotetext{I owe these decompositions to a communication from \gesperrt{Dedekind}.}


% ===== scan leaf 142, printed 121 =====
% !PRINTED 121
\weberpage{121}
\pdfbookmark[2]{§ 29. The first Sylow theorem}{bm35}
If now $e$ is an arbitrary element of $Q$, then $Qe = Q$,
and consequently also $Ra_iQe = Ra_iQ$. From this it follows that by
composition with $e$ the totality of the $h_i$ elements of which
$P_i$ is composed is not altered, and that they are therefore
only interchanged among one another by the permutation group $\Pi$.
But one also sees immediately that these $h_i$ elements $B$ are \gesperrt{transitively} connected by $\Pi$; for $P_1$ consists precisely
of all those complexes $B$ which can be derived from one among them by
application of elements of $Q$, and
hence each of these complexes $B$ of $P_i$ can be carried over into
every other by permutations of $\Pi$.

We wish to formulate the result obtained as a theorem
in the following way:

3.\ \gesperrt{If $P$ is a group of degree $n$, $R$ a divisor
of $P$ of index $j$, $Q$ a divisor of $P$ of
degree $q$, then with $Q$ a permutation group $\Pi$
of $j$ figures is simply or multiply isomorphic. The
$j$ figures split into rows of $h_1, h_2, h_3, \ldots$
figures each, transitively connected by $\Pi$, so that}
\weq{11}{j = h_1 + h_2 + h_3 + \cdots}
\gesperrt{and $h_1, h_2, h_3, \ldots$ are divisors of $q$. The
numbers $h_1, h_2, h_3, \ldots$ are the indices of divisors
of $Q$ which result as intersections of $Q$ with
the conjugate groups $a^{-1}Ra$.}

If the various groups $Q_1, Q_2, \ldots$ have, besides the
unit element, no common divisor, which with simple groups $P$ always occurs, the isomorphism is simple.

\begin{center}
\S.\,29.

\gesperrt{The first \gesperrt{Sylow} theorem.}
\end{center}

It is not easy, on the direct path taken by \gesperrt{Cayley},
to go beyond the lowest numbers of degree in the
formation of groups. For further investigations
in this direction, a theorem is of great use, which in a restricted
scope originates from \gesperrt{Cauchy}, but in general
was proved by \gesperrt{Sylow}. This important theorem we shall
%CONT>>


% ===== scan leaf 143, printed 122 =====
% !PRINTED 122
\weberpage{122}
%<<CONT
prove here by a method of \gesperrt{Frobenius} which rests only on the general concept of a group\footnote{\gesperrt{Cauchy}, Exerc. d'analyse, Tom. III. \gesperrt{Sylow}, Mathem. Annalen, Bd. 5. \gesperrt{Frobenius}, Journ. für Mathematik, Bd. 100. \gesperrt{Netto}, Mathem. Annalen, Bd. 13; Substitutionentheorie. \S.\,48.}.

The theorem can be stated simply thus:

\gesperrt{I. If $P$ is a group of any elements whatever of degree $n$, and $p^{\varkappa}$ is a prime power dividing $n$, then the group $P$ has a divisor of degree $p^{\varkappa}$.}

\gesperrt{Cauchy} has proved this theorem for $\varkappa = 1$. For the case that $n$ is a prime, it is evident; for $n = 4$ and $n = 6$ one can read it off easily from the compilations in \S.\,27, so that it may therefore be regarded as proved for the cases $n = 2, 3, 4, 5, 6, 7$. On the basis of this observation, complete induction can be applied, and we therefore assume \gesperrt{that the theorem has been proved for groups whose degree is lower than} $n$.

We shall denote the elements of the group $P$ by the letters $a, b, c, \ldots$ and understand by $x$ a symbol for a variable element in $P$ which runs through the whole group $P$.

We first gather together all those elements $a$ which satisfy the condition
\weq{1}{x^{-1}ax = a \quad \text{or} \quad ax = xa}
so that $a$ denotes every element that is \gesperrt{permutable} with all elements of $P$. To these belongs in any case the unit element, and the totality of these elements $a$ forms a group $A$ contained in $P$. For from
\[ x^{-1}ax = a,\quad x^{-1}a'x = a' \]
it follows that
\[ x^{-1}aa'x = aa'. \]

\gesperrt{$A$ is an Abelian group and a normal divisor of $P$}; for from the definition (1) it follows that $x^{-1}Ax = A$, and for any two elements $a, a'$ of $A$:
\[ aa' = a'a. \]
We denote by $\nu$ and $j$ the degree and index of this group, so that
\weq{2}{n = \nu j}
%CONT>>


% ===== scan leaf 144, printed 123 =====
% !PRINTED 123
\weberpage{123}
%<<CONT
is. Now if $c$ is an element of $P$ not belonging to $A$, then $x^{-1}cx$ will not be equal to $c$ for all elements $x$. But certain elements $b, b', \ldots$ will be found among the $x$, among them certainly the unit element, which satisfy the condition
\[ x^{-1}cx = c \]
and these elements form a group, because from
\weq{3}{b^{-1}cb = c,\quad b'^{-1}cb' = c}
it follows that also
\[ (b'b)^{-1}c(b'b) = b^{-1}b'^{-1}cb'b = c \]
must hold. We denote this group by $B$ and its degree and index by $\mu, \varepsilon$, so that
\weq{4}{n = \mu\varepsilon}
holds; $B$ is a proper divisor of $P$, and hence $\varepsilon$ is in any case greater than 1, since otherwise $c$ would be contained in $A$, contrary to the assumption.

If $g$ is an element of $P$ that does not occur in $B$, and $b$ an arbitrary element of $B$, then by (3):
\weq{5}{g^{-1}b^{-1}cbg = g^{-1}cg,}
which is different from $c$, because $g$ does not occur in $B$. If therefore one decomposes $P$ into the subsidiary groups:
\weq{6}{P = B + Bg_1 + Bg_2 + \cdots + Bg_{\varepsilon-1},}
then, on putting for $x$ in $x^{-1}cx$ all the elements of one of these subsidiary groups, one always obtains the same element $g^{-1}cg$, and one thus obtains $\varepsilon$ and no more elements, and each equally often:
\weq{7}{c,\ g_1^{-1}cg_1,\ g_2^{-1}cg_2,\ \ldots,\ g_{\varepsilon-1}^{-1}cg_{\varepsilon-1};}
all these elements, whose totality we shall denote by $C$, are different from one another and from the elements $a$. But $C$ is not a group, since the unit element does not occur among them.

If now the totality of the elements of $P$ is not yet exhausted by $A$ and $C$, one takes an element $c'$ occurring neither in $A$ nor in $C$ and forms by the same rule the group $B'$ of degree $\mu'$ of all elements which satisfy the condition
\[ b'^{-1}c'b' = c' \]
and the system $C'$:
\weq{8}{c',\ g_1'^{-1}c'g_1',\ g_2'^{-1}c'g_2',\ \ldots,\ g_{\varepsilon'-1}'^{-1}c'g_{\varepsilon'-1}',}
%CONT>>


% ===== scan leaf 145, printed 124 =====
% !PRINTED 124
\weberpage{124}
%<<CONT
and one then sees that these elements are all different from the elements $C$, since, if $g'^{-1}c'g' = g^{-1}cg$ were the case, we should also have $c' = g'g^{-1}cgg'^{-1}$, and $c'$ would thus, contrary to the assumption, already occur in $C$. In this way one continues with the formation of the groups $B, B', B'', \ldots$ and the systems $C, C', C'', \ldots$ until all the elements of $P$ are accommodated in the systems $A, C, C', C'', \ldots$. One then has the numerical equations:
\weq{9}{n = \nu j = \mu\varepsilon = \mu'\varepsilon' = \mu''\varepsilon''\ldots}
\weq{10}{n = \nu + \varepsilon + \varepsilon' + \varepsilon'' + \cdots}

On the basis of these formulae \gesperrt{Sylow}'s theorem can now be proved by complete induction, in which two cases are to be distinguished:

1) \gesperrt{The degree} $\nu$ \gesperrt{of} $A$ \gesperrt{is divisible by} $p$\gesperrt{.} In this case, since $A$ is an Abelian group, there is by \S.\,9, 2. an element $a$ in $A$ of degree $p$, and the cyclic group of degree $p$
\[ 1,\ a,\ a^2,\ \ldots,\ a^{p-1}, \]
which we denote by $Q$, is a normal divisor of $P$, because for an element $a$ of $A$ we should always have $x^{-1}ax = a$. The complementary group to $Q$, $P/Q$ is of degree $n:p$, and its degree is therefore smaller than $n$ and divisible by $p^{\varkappa-1}$. By our assumption there is therefore a divisor of $P/Q$ of degree $p^{\varkappa-1}$, and consequently by \S.\,6, 2. there is a divisor of $P$ of degree $p^{\varkappa}$.

2) \gesperrt{The degree} $\nu$ \gesperrt{of} $A$ \gesperrt{is not divisible by} $p$\gesperrt{.} Since $n$ is divisible by $p$, in this case by (10) the numbers $\varepsilon, \varepsilon', \varepsilon'', \ldots$, which are all greater than 1, cannot all be divisible by $p$. Now if $\varepsilon$ is not divisible by $p$, then $\mu$ is divisible by $p^{\varkappa}$. The degree of the group $B$ is smaller than $n$ and divisible by $p^{\varkappa}$, and so by assumption there is a divisor of $B$, which is therefore also a divisor of $P$, of degree $p^{\varkappa}$.

With this Theorem I is proved.

To this theorem, however, there are still essential supplements to be added.
%CONT>>


% ===== scan leaf 146, printed 125 =====
% !PRINTED 125
\weberpage{125}
\pdfbookmark[2]{§ 30. The second Sylow theorem}{bm36}
%<<CONT
\begin{center}
\S.\,30.\\
\gesperrt{The second Sylow theorem.}
\end{center}

Now let, as above, $P$ be a group of degree $n$ and $p^{\varkappa}$ the highest power of the prime $p$ that divides $n$; let $Q$ be a divisor of $P$ of degree $p^{\varkappa}$, which exists by Theorem I, \S.\,29.

The totality of all elements $c$ of $P$ that satisfy the condition
\weq{1}{c^{-1}Qc = Q,}
among which all the elements of $Q$ are certainly contained, forms a group; for from
\[ Qc = cQ,\quad Qc' = c'Q \]
it follows that $Qcc' = cc'Q$ also holds. These elements $c$ we may call \gesperrt{the elements of $P$ permutable with the group $Q$}. If we denote the group of the $c$ by $R$, then $R$ is a divisor of $P$, and $Q$ a divisor of $R$, and indeed a normal divisor. In particular cases $R$ can be identical with $Q$ as well as with $P$. But if $b$ is any element of $P$ not contained in $R$, then among the elements of the group $b^{-1}Qb$ there is certainly at least one that is not contained in $Q$.

If we denote by $r$ the index of $Q$ with respect to $R$, and by $j$ the index of $R$ with respect to $P$, then
\weq{2}{n = p^{\varkappa}rj,}
and $r$ and $j$ are not divisible by $p$.

The theorem that we have to prove reads:

\gesperrt{II. The index $j$ of $R$ is of the form $pk + 1$, where $k$ is any integer $[j \equiv 1\ (\mathrm{mod}\ p)]$; there are $j$ and no more distinct divisors of $P$ of degree $p^{\varkappa}$, which are all conjugate to one another. Every divisor of $P$ whose degree is a power of $p$ is contained in one of these conjugate groups.}

If $c$ is an element from $R$ of degree $\gamma$, and $s$ is the least positive exponent for which $c^s$ is contained in $Q$, then every other exponent $t$ for which $c^t$ belongs to $Q$ must be divisible by $s$
%CONT>>


% ===== scan leaf 147, printed 126 =====
% !PRINTED 126
\weberpage{126}

%<<CONT
be. For if we set $t = m s + s_1$, where $s_1 < s$, then $c^{s_1}$
is contained in $Q$, and $s_1$ must therefore be $= 0$. Accordingly $s$ is a
divisor of $\gamma$. But now the elements

\[
Q + c\,Q + c^2\,Q + \cdots c^{s-1}\,Q
\]

form, by (1), a group contained in $P$, and indeed of degree $p^{\varkappa} s$,
and since $p^{\varkappa} s$ must then be a divisor of $n$, it follows that $s$
cannot be divisible by $p$. Now if the degree $\gamma$ of $c$
were a power of $p$, then $s$ would also have to be a power of $p$, which
by the foregoing is not possible.

The degree of $c$ can therefore certainly not be a power of
$p$; i.e., \gesperrt{apart from the elements $Q$ there are in $R$ no
elements whose degree is a power of $p$.}

If now $b$ denotes an element of $P$ that does not belong to $R$,
then by the definition of $R$ there is in the group

\[
Q' = b^{-1}\,Q\,b
\]

at least one element that does not occur in $Q$, and the
degree of this element must be a power of $p$, because the
degree of $Q'$ is equal to $p^{\varkappa}$, hence a power of $p$. This element
can therefore not occur in $R$, \gesperrt{and $Q'$ is accordingly not a
divisor of $R$.}

If we choose the elements $b_1, b_2, \ldots, b_{j-1}$ from $P$ in such a way
that we can set

\weq{10}{P = R + R\,b_1 + R\,b_2 + \cdots + R\,b_{j-1}} % [sic: Weber jumps from equation (2) to (10) in this section and subsequently cites (10).]

then the $j - 1$ groups:

\[
b_1^{-1}\,Q\,b_1,\ \ b_2^{-1}\,Q\,b_2,\ \ \ldots,\ \ b_{j-1}^{-1}\,Q\,b_{j-1},
\]

which are all of degree $p^{\varkappa}$, are distinct from $Q$. But they are
also distinct from one another; for if we had, e.g.,

\[
b_1^{-1}\,Q\,b_1 = b_2^{-1}\,Q\,b_2,
\]

it would follow that:

\[
b_2\,b_1^{-1}\,Q\,b_1\,b_2^{-1} = Q,
\]

i.e., $b_2\,b_1^{-1} = c$ would be contained in $R$, hence $b_2 = c\,b_1$ in $R\,b_1$, which
by (10) is not possible. We therefore certainly have $j$ distinct
conjugate groups of degree $p^{\varkappa}$:

\weq{11}{Q,\ \ b_1^{-1}\,Q\,b_1,\ \ b_2^{-1}\,Q\,b_2,\ \ \ldots,\ \ b_{j-1}^{-1}\,Q\,b_{j-1}.}

If $Q' = b^{-1}\,Q\,b$ is one of them, then $R' = b^{-1}\,R\,b$ contains
the divisor $Q'$, but apart from this no element whose degree
%CONT>>


% ===== scan leaf 148, printed 127 =====
% !PRINTED 127
\weberpage{127}
\pdfbookmark[2]{§ 31. Groups of degree pα}{bm37}
%<<CONT
is a power of $p$. $R'$ is the totality of all elements of $P$ permutable with $Q'$.

We now further apply the theorem 3., \S.\,28, understanding by $P$, $R$, $Q$ the groups that are here denoted by the same letters. Since here $Q$ is a divisor of $R$, we have $h_1 = 1$. But of the conjugate groups $b^{-1}Rb$ none except $R$ is divisible by $Q$, and so the remaining $h_2, h_3, \ldots$ are all greater than $1$. But since they must be divisors of the degree of $Q$, they are powers of $p$, and from the equation

\[
j = h_1 + h_2 + h_3 + \cdots
\]

it therefore follows that $j \equiv 1\ (\mathrm{mod}\ p)$.

But if $Q_1$ is any divisor of $P$ whose degree is a power of $p$, we again apply the decomposition of $j$ according to the Theorem 3., \S.\,28, putting for the groups $P$, $R$, $Q$ occurring there the groups $P$, $R$, $Q_1$. Then again

\[
j = h_1 + h_2 + h_3 + \cdots,
\]

in which now $h_1, h_2, h_3, \ldots$ denote the indices (with respect to $Q$) of the intersections of $Q_1$ with the group $R$ and its conjugates, which are therefore likewise powers of $p$. % [sic: Weber prints "in Bezug auf Q"; after substituting Q_1 for theorem 3's Q, these indices are with respect to Q_1.]
But since $j \equiv 1\ (\mathrm{mod}\ p)$, for at least one of the numbers we must have $h = 1$, i.e. one of the groups $b_\lambda^{-1}Rb_\lambda$, and consequently also one of the groups (11), is divisible by $Q_1$. If $Q_1$ is of degree $p^\varkappa$, it is identical with one of the groups (11).

With this, then, the Theorem II is completely proved.

\begin{center}
\S.\,31.

\gesperrt{Groups of degree $p^\alpha$.}
\end{center}

\gesperrt{Sylow} has derived from his theorems a remarkable property of the groups whose degree is a power of a prime $p$, namely $p^\alpha$. Let $P$, then, be such a group and $\beta$ a positive number smaller than $\alpha$. Then by the theorem I, \S.\,29 the group $P$ contains a divisor $Q$ of degree $p^\beta$, whose index is $p^{\alpha-\beta}$. We apply the theorem 3., \S.\,28, putting for $R$ and $Q$ this group $Q$ of degree $p^\beta$; then $h_1, h_2, h_3, \ldots$ are
%CONT>>


% ===== scan leaf 149, printed 128 =====
% !PRINTED 128
\weberpage{128}
%<<CONT
the indices of $Q$ and of the divisors of $P$ conjugate to $Q$, taken with respect to $Q$. Hence $h_1 = 1$ and

\[
p^{\alpha-\beta} = h_1 + h_2 + h_3 \ldots % [sic: Weber prints no plus sign before the ellipsis; the preceding decomposition requires one.]
\]

The summands $h_1, h_2, h_3, \ldots$, being divisors of the degree of $Q$, can here only be powers of $p$; and since $h_1 = 1$; it follows that at least $p$ of these summands must be $= 1$, since otherwise their sum could not be divisible by $p$. But this means nothing else than that at least $p - 1$ of the conjugate groups $b^{-1}Qb$, where $b$ is not contained in $Q$, must be divisible by $Q$ and therefore equal to $Q$. But if

\[
b^{-1}Qb = Q
\]

then for every exponent $s$ we also have

\[
b^{-s}Qb^s = Q.
\]

If $b^r$ is the lowest power of $b$ contained in $Q$, then, since the degree of $b$ is a power of $p$, $r$ must also be a power of $p$, and

\[
b^{\frac{r}{p}} = c
\]

is an element of $P$ which does not itself occur in $Q$, but whose $p$th power is contained in $Q$, and for which at the same time the condition

\[
cQc^{-1} = Q
\]
holds.

From this it follows that the elements

\[
R = Q + Qc + Qc^2 + \cdots + Qc^{p-1}
\]

form a group of degree $p^{\beta+1}$, of which $Q$ is a normal divisor. We therefore have the theorem:

III. \gesperrt{If $P$ is a group of degree $p^\alpha$ and $Q$ a divisor of $P$ of degree $p^\beta$ ($\beta < \alpha$), then there is a divisor $R$ of $P$ of degree $p^{\beta+1}$, of which $Q$ is a normal divisor.}

If in this theorem one takes $\beta = \alpha - 1$, then $R$ coincides with $P$, and it results that $Q$ is a normal divisor of $P$. Applying the same theorem to the group $Q$, and so on, one obtains:

IV. \gesperrt{Every group whose degree is a power of a prime is metacyclic (\S.\,8).}


% ===== scan leaf 150, printed 129 =====
% !PRINTED 129
\weberpage{129}
\pdfbookmark[2]{§ 32. Theorem of Frobenius}{bm38}
\begin{center}
\S.\,32.

\gesperrt{Theorem of Frobenius.}
\end{center}

\gesperrt{Frobenius} has established a theorem which may be regarded as a counterpart to the fourth Sylow theorem, and which may be stated thus\footnote{\gesperrt{Frobenius}, Sitzungsberichte der Berl. Akademie, 4. Mai 1893.}:

V. \gesperrt{Every group whose degree is a product of nothing but distinct primes is metacyclic.}

We therefore assume that the degree $n$ of a group $P$ is \gesperrt{divisible by no square number}. If $t$ is the greatest prime divisor of $n$, then by the \gesperrt{Cauchy-Sylow} theorem (\S.\,29) there is a divisor $T$ of degree $t$ of $P$, which, because $t$ is a prime, is a cyclic group, and thus consists of the elements

\[
T = 1, \; a, \; a^2, \; \ldots \; a^{t-1}
\]

If it could now be shown that under the assumptions made $T$ is a normal divisor of $P$, then one could apply the same mode of reasoning to the group $P/T$, whose degree is $n:t$, and, if $t_1$ is the second-greatest prime dividing $n$, would conclude that there is a normal divisor $N$ of the group $P/T$ of degree $t_1$. But from this it would then follow (\S.\,6, 2.) that $P$ has a normal divisor $T_1$ of degree $t\,t_1$, of which $T$ itself is in turn a normal divisor. By considering the group $P/T_1$ in the same way, and so on, one obtains a composition series of $P$:

\[
P, \; T_k, \; T_{k-1}, \; \ldots \; T_1, \; T, \; 1,
\]

in which the index series consists of the prime factors of $n$ arranged in ascending order of magnitude, and $P$ proves to be metacyclic. Everything therefore comes down to showing that $T$ is a \gesperrt{normal divisor} of $P$.

This theorem is easier to prove if one regards it as a special case of a more general one than if one attacks it directly. Under the assumptions made about $P$ and $n$, this more general theorem reads:

$\alpha$) \gesperrt{If $n = \mu\nu$ is decomposed into two factors and every prime divisor of $\nu$ is greater than every prime}
%CONT>>


% ===== scan leaf 151, printed 130 =====
% !PRINTED 130
\weberpage{130}
%<<CONT
\gesperrt{divisor of $\mu$, then there are $\nu$ and no more elements in $P$ whose degree divides $\nu$.}

If we apply this theorem to $\nu = t$, it follows that apart from $T$ there can be no elements in $P$ whose degree $= t$. But if $c$ is any element of $P$, then the group $c^{-1}Tc$ conjugate to $T$ is likewise of degree $t$, and must therefore be identical with $T$. It is thereby proved that $T$ is a normal divisor of $P$.

For the proof of theorem $\alpha$) one applies complete induction. It is evidently correct when $\mu = 1$, $\nu = n$, and we therefore assume as proved:

$\alpha'$) \gesperrt{If $\mu'\nu'$ is divisible by no square, every prime divisor of $\nu'$ greater than every prime divisor of $\mu'$, and the number of prime divisors of $\mu'$ smaller than the number of prime divisors of $\mu$, then in every group of degree $\mu'\nu'$ there are exactly $\nu'$ elements whose degree is a divisor of $\nu'$.}

To derive from this the proof of theorem $\alpha$), we denote by $p$ the greatest prime divisor of $\mu$ and by $Q$ a group of degree $p$ contained in $P$, whose existence is established by the \gesperrt{Cauchy-Sylow} theorem. Then all prime factors of $\nu$ are greater than $p$. As in \S.\,30 we denote by $R$ the group which consists of all elements $c$ of $P$ satisfying the condition
\[
c^{-1}Qc = Q
\]
whose degree we denote, as above, by $pr$; $j$ signifies, as before, the index of the divisor $R$ of $P$. We set $r = r_1 r_2$, and understand by $r_1$ the greatest common divisor of $\mu$ and $r$, so that $r_2$ is the greatest common divisor of $\nu$ and $r$. Then
\[
n = \mu\nu = r_1 p r_2 j.
\]
By the hypothesis $\alpha'$), $P$ now contains exactly $p\nu$ elements whose degree is a divisor of $p\nu$; let the totality of these elements be denoted by $U$. But the group $R$ contains, likewise by $\alpha'$), exactly $pr_2$ elements whose degree is a divisor of $pr_2$, and hence a divisor of $p\nu$. This system we shall denote by $V$. Evidently all elements of $V$ are at the same time contained in $U$.


% ===== scan leaf 152, printed 131 =====
% !PRINTED 131
\weberpage{131}
If we can now also prove that there are in $U$ \gesperrt{exactly $(p-1)\nu$ elements whose degree is divisible by $p$,} then we have reached our goal; for it then follows that among the elements $U$, and hence also among the elements of $P$, there are exactly
\[
p\nu - (p-1)\nu = \nu
\]
whose degree is a divisor of $\nu$.

But this results from the following:

$\beta$) \gesperrt{If $v$ is an element of $V$, then all the elements $vQ$ (whose number is $p$) belong to $V$. Among them there are $p-1$ whose degree is divisible by $p$, and one whose degree is not divisible by $p$. In $V$ there are $(p-1)r_2$ elements whose degree is divisible by $p$, and $r_2$ elements whose degree is not divisible by $p$.}

For if $a$ is an element of $Q$ different from $1$, then, since $v$ belongs to $R$, so that $v^{-1}av$ is contained in $Q$, and since $Q$ consists of the powers of $a$,
\weq{1}{v^{-1}av = a^s.}
By repeated application there results from this, for every exponent $h$,
\weq{2}{v^{-h}av^h = a^{s^h}.}
If we take for $h$ the degree of the element $v$, which by the hypothesis is a divisor of $p\nu$, and thus set $v^h = 1$, it follows that
\weq{3}{a = a^{s^h}, \quad s^h \equiv 1 \ (\mathrm{mod}\ p).}
But no prime factors divide $h$ other than such as are equal to or greater than $p$, and hence $p-1$ is relatively prime to $h$, and one can satisfy the congruence
\weq{4}{xh \equiv 1 \ (\mathrm{mod}\ p-1)}
accordingly, from the \gesperrt{Fermat's} theorem there results:
\weq{5}{s^{xh} \equiv s \equiv 1 \ (\mathrm{mod}\ p),}
hence by (1):
\weq{6}{av = va,}
i.e.\ $a$ and $v$ are permutable. From this it follows, for every exponent $\lambda$,
\weq{7}{(va)^\lambda = v^\lambda a^\lambda.}


% ===== scan leaf 153, printed 132 =====
% !PRINTED 132
\weberpage{132}
If, as before, $h$ is the degree of $v$, and if we set $\lambda = h$ when $h$ is divisible by $p$, and $\lambda = hp$ when $h$ is not divisible by $p$, then it follows from (7) that $(va)^\lambda = 1$, i.e. the degree of $va$ is a divisor of $h$ or of $hp$, hence in any case a divisor of $p\nu$; i.e. $va$ \gesperrt{is contained in} $V$.

If we now replace $a$ in (7) by $a^x$ and let $x$ run through the sequence of numbers $0, 1, \ldots p-1$, then $a^x$ runs through the group $Q$ and $va^x$ through the system $vQ$, and for every $\lambda$ we have
\weq{8}{(va^x)^\lambda = v^\lambda a^{x\lambda}.}
Now if the degree $h$ of $v$ is not divisible by $p$, then
\[
(va^x)^h = a^{xh},
\]
and hence $= 1$ only when $x = 0$; on the other hand
\[
(va^x)^{ph} = 1,
\]
i.e. the degree of $va^x$, when $x$ is not $= 0$, is a divisor of $ph$, but not of $h$, hence divisible by $p$; on the other hand it is $= h$ when $x = 0$.

But if the degree $h$ of $v$ is divisible by $p$, then it is of the form $pg$, and $g$ is not divisible by $p$, because $p$ divides $n$, and hence also $h$, only once. Thus $v^g$ is an element of $R$ of degree $p$, and must therefore, by \S.\,30, be contained in $Q$. If accordingly we set
\[
v^g = a^y,
\]
then
\[
(va^x)^g = a^{y+gx},
\]
and this is $= 1$ if and only if $y + gx \equiv 0 \ (\mathrm{mod}\ p)$, which occurs for only \gesperrt{one} value of $x$. For this value of $x$ the degree of $va^x$ is a divisor of $g$; for the other values of $x$ it is a divisor of $pg$, but not of $g$, hence divisible by $p$. With this the first two assertions of theorem $\beta$) are proved.

In order to see the last part as well, one may observe that if $v_1, v_2$ denote two elements of $V$, the systems $v_1 Q, v_2 Q$ are either wholly identical, or have not a single element in common. From this it follows that $V$ can be represented in the following manner:
\[
V = v_1 Q + v_2 Q + \cdots + v_{r_2} Q
\]
(although $V$ is in general not a group), and in each of these $r_2$ partial systems $vQ$ there are $p-1$ elements whose degree
%CONT>>


% ===== scan leaf 154, printed 133 =====
% !PRINTED 133
\weberpage{133}
%<<CONT
divisible by $p$, and an element whose degree is not divisible by $p$.

By Theorem II, \S.\,30 there are $j$ and no more mutually distinct conjugate groups $Q, Q', Q'', \ldots Q^{(j-1)}$ in $P$, and $j$ mutually distinct conjugate groups $R, R', R'', \ldots R^{(j-1)}$.

According to this, the following theorem can be proved:

$\gamma$) \gesperrt{Every element of $P$ whose degree is divisible by $p$ occurs in one and only one of the groups $R, R', R'' \ldots R^{(j-1)}$.}

The groups $Q, Q', Q'' \ldots$ are cyclic, since their degree is a prime, and any two of them, because they are distinct from one another, contain no common element besides the unit element. Let
\[
Q = 1, a, a^2, \ldots a^{p-1}.
\]
The group $R$, which consists of all elements satisfying the condition
\[
c^{-1}Qc = Q
\]
contains, besides the powers of $a$, no element of degree $p$; and the corresponding statement holds for the other groups $Q', R', Q'', R'' \ldots$

Now if $b$ is an element of $P$ whose degree is divisible by $p$ and may be equal to $ps$, so that $s$ is not divisible by $p$, then $b^s$ is an element whose degree is $p$, and which therefore occurs in one and only one of the cyclic groups $Q, Q', Q'' \ldots$ But if $b^s = a$ is an element of $Q$, then
\[
b^{-1}ab = a,
\]
that is, $b$ occurs in the group $R$, and can occur in none of the other groups, e.g. in $R'$, because otherwise $b^s$ would also occur in $Q'$, which is not possible, since $Q$ and $Q'$ have only the unit element in common. Thus $\gamma$) is proved.

$\delta$) \gesperrt{In $U$ there are exactly $jr_2(p-1)$ elements whose degree is divisible by $p$.}

By $U$ we have understood the totality of the elements of $P$ whose degree is a divisor of $pr$. Now by $\beta$) there are in $R$ exactly $r_2(p-1)$ elements from $U$ whose degree is divisible by $p$, and since each of the $j$ groups $R, R', R'', \ldots$ can take the place of $R$, $\delta$) follows from $\gamma$). In order, then, still to show, which by the above is all that matters, that in $U$
%CONT>>


% ===== scan leaf 155, printed 134 =====
% !PRINTED 134
\weberpage{134}
\pdfbookmark[2]{§ 33. Groups of degree pα q}{bm39}
%<<CONT
exactly $(p-1)\nu$ elements occur whose degree is divisible by $p$, only the formula

\[ \varepsilon)\quad \nu = j\,r_2 \]

remains to be proved. This results from the following consideration. By $\alpha')$ the number of the elements $U$ is equal to $p\,\nu$. But among these elements $U$ there is certainly one, the unit element, whose degree is not divisible by $p$, and consequently by $\delta)$

\[ j\,r_2\,(p-1) < p\,\nu. \]

On the other hand $n = \mu\,\nu = p\,r_1 r_2 j$, and since $\nu$ is relatively prime to $p\,r_1$, $j\,r_2$ is divisible by $\nu$. Hence $j\,r_2 : \nu$ is a positive integer which is smaller than $p : p-1$ and consequently also smaller than 2, and which can therefore only be equal to 1. Thereby $\varepsilon)$ is proved, and with it the whole theorem.

\begin{center}
\S.\,33.

\gesperrt{Groups of degree $p^{\alpha} q$.}
\end{center}

\gesperrt{Frobenius} has stated the following theorem in the memoir mentioned:

VI. \gesperrt{Every group of degree $p^{\alpha} q$, if $p$ and $q$ are primes different from one another and $\alpha$ denotes an arbitrary positive exponent, is metacyclic.}

The following proof of this theorem comes from a communication by letter from \gesperrt{Frobenius} to the author.

First of all it is clear that it suffices to prove that no group $P$ of degree $p^{\alpha} q$ is simple. For our theorem is proved for $\alpha = 1$. If we therefore assume it already proved for every smaller $\alpha$, and suppose that $P$ has a proper normal divisor $Q$ which comprises more than the unit element, then the two groups $Q$, $P/Q$ are metacyclic by the assumption or by Theorem IV, and hence $P$ is also metacyclic.

We now assume, then, that the group $P$ of degree $p^{\alpha} q$ is simple, and we have to derive a contradiction from this assumption in order to establish its impossibility.

$\alpha)$ First, by Theorem I., \S.\,29, we have a divisor of $P$ of degree $q$. We will assume that there is a
%CONT>>


% ===== scan leaf 156, printed 135 =====
% !PRINTED 135
\weberpage{135}
%<<CONT
divisor of $P$ of degree $p^{\beta} q$, where $0 \leqq \beta < \alpha$ is assumed, and let $\beta$ be taken as large as possible. Then, because of the assumed simplicity of the group $P$, the theorem of \S.\,28, 2. is applicable, according to which $P$ is isomorphic with a permutation group of $p^{\alpha-\beta}$ figures. Among the permutations of this group there are also some whose degree is divisible by $q$, and which therefore, decomposed into their cycles, must contain a cycle of $q$ figures, and from this it follows:

\weq{1}{p^{\alpha-\beta} > q.}

$\beta)$ Secondly, now let $Q$ be a divisor of $P$ of degree $p^{\alpha}$ and index $q$. Since by hypothesis $P$ is simple, the group denoted by $R$ in \S.\,30, which consists of the elements of $P$ permutable with $Q$, cannot be identical with $P$, since otherwise $Q$ would be a normal divisor of $P$, which, since $P$ is assumed to be a simple group, is after all impossible. But since $Q$ is a divisor of $R$, and the index of $Q$ with respect to $P$ is a prime, we must have $R = Q$. By II, \S.\,30 there are therefore $q$ and no more mutually distinct conjugate groups

\[ Q,\ Q_1,\ Q_2,\ \dots\ Q_{q-1}, \]

which are all of degree $p^{\alpha}$. Among these groups we take two, say $Q$, $Q_1$, which have a greatest common divisor $Q_0$ of degree $p^{\gamma}$, and we assume these two groups so chosen that $\gamma$ becomes \gesperrt{as large as possible}. Then $0 \leqq \gamma < \alpha$. If we now apply to the group $P$ the theorem of \S.\,28, 3., putting the group $Q$ for the two groups $Q$, $R$ of that theorem, then

\[ q = h_1 + h_2 + h_3 + \cdots, \]

and in it $h_1 = 1$, and the remaining summands $h_2, h_3, \dots$ are the indices of greatest common divisors of $Q$ with one each of the groups $Q_1, Q_2, \dots$. Hence $h_1, h_2, \dots$ are divisible by $p^{\alpha-\gamma}$, % [sic: Weber includes h_1 among the divisible terms although h_1 = 1.]
and it follows that

\[ q \equiv 1 \pmod{p^{\alpha-\gamma}}, \]

hence

\weq{2}{q > p^{\alpha-\gamma},}

and consequently, on account of (1):

\weq{3}{\gamma > \beta,\quad \gamma > 0.}

$\gamma)$ If now, then, $Q_0$ is the intersection of $Q$ and $Q_1$ of degree $\gamma$, % [sic: Weber prints gamma although Q_0 was defined to have degree p^gamma.]
we can by \S.\,31, III. determine a divisor $R$ of $Q$ of degree $p^{\gamma+1}$, of which $Q_0$ is a normal divisor, and $R$
%CONT>>


% ===== scan leaf 157, printed 136 =====
% !PRINTED 136
\weberpage{136}
\pdfbookmark[2]{§ 34. Simple groups}{bm40}
%<<CONT
is then contained in none of the groups $Q_1, Q_2, \dots, Q_{q-1}$, because it was assumed that $Q$ has no divisor in common with any of these groups whose degree is greater than $p^{\gamma}$. In the same way we can find a divisor $R_1$ of $Q_1$ of degree $p^{\gamma+1}$, of which $Q_0$ is a normal divisor, and which is contained in none of the groups $Q, Q_2, \dots, Q_{q-1}$. Now $R$ as well as $R_1$ is contained in $P$. We consider the smallest group $S$ that contains $R$ and $R_1$ at the same time, which is in any case contained in $P$ (the least common multiple of $R$ and $R_1$). The degree of this group $S$ cannot be a power of $p$, since otherwise, by Theorem II., \S.\,30, $S$, and hence $R$ and $R_1$, would have to be contained in one of the groups $Q, Q_1, \dots, Q_{q-1}$. But $R$ is contained only in $Q$, $R_1$ only in $Q_1$; hence this assumption is inadmissible. The degree of $S$ must therefore be of the form $p^{\lambda} q$.

But $\lambda$ cannot be smaller than $\alpha$; for, since a group of degree $p^{\gamma+1}$ is contained in $S$, namely $R$, we have

\weq{4}{\lambda \geqq \gamma+1 > \beta.}

By the hypothesis made in $\alpha)$, however, no divisor of degree $p^{\lambda} q$ is contained in $P$ in which $\lambda$ is at the same time greater than $\beta$ and smaller than $\alpha$. By (4), therefore, necessarily $\lambda = \alpha$, i.\,e. $S$ is identical with $P$.

Now $Q_0$ is a normal divisor of $R$ and $R_1$. If, then, we collect all elements $c$ of $P$ that satisfy the condition

\[ c^{-1}\,Q_0\,c = Q_0 \]

into a group, this group contains $R$ as well as $R_1$, and is therefore the whole group $P$. Hence $Q_0$ is also a normal divisor of $P$. Since by (3) $\gamma$ is greater than zero, the degree of $Q_0$ is greater than 1, and we come upon a contradiction with the assumption that $P$ is simple.

\begin{center}
\S.\,34.

\gesperrt{Simple Groups.}
\end{center}

The theorems that we have come to know in the preceding sections admit fairly far-reaching conclusions about the nature of groups. They show that, at least for the lower degrees, the metacyclic (and cyclic) groups decidedly predominate. The
%CONT>>


% ===== scan leaf 158, printed 137 =====
% !PRINTED 137
\weberpage{137}
%<<CONT
few non-metacyclic groups contained among them, and especially the simple non-cyclic groups, therefore command all the greater interest.

It has not so far been possible to comprehend the degree numbers of the simple groups in a definite law, and one has been content to determine all simple groups up to certain limits of the degree numbers, on the one hand with the help of the theorems developed above, on the other hand by special methods. \gesperrt{Hölder} has carried out these investigations for the degree numbers up to 200, and \gesperrt{Cole} for the degree numbers up to 500\footnote{\gesperrt{Hölder}, Mathematische Annalen, Bd.\,40. \gesperrt{Cole}, American Journal, Vol.\,14.}. Of simple non-cyclic groups, only those of degrees 60, 168, 360, already known from other investigations, have emerged in this. Further, a simple group of degree 660 is also known, and recently \gesperrt{Cole} and \gesperrt{Moore} have called attention to a simple group of degree 504\footnote{Bulletin of the New York mathem. society, Oct.\,1893.}. These investigations become very laborious as the degree numbers grow. We shall confine ourselves here to the consideration of the degree numbers of the first hundred.

Here, by the three general theorems of \gesperrt{Sylow} and \gesperrt{Frobenius}, the majority of the groups prove to be metacyclic, and there remain only the degree numbers

\[ 36,\ 60,\ 72,\ 84,\ 90,\ 100 \]

But that a group of degree $36$ cannot be simple follows from \gesperrt{Sylow}'s theorem. For such a group would have to have a divisor of degree $9$, and would therefore have to be representable by the permutations of four figures. But that is impossible, because there are only 24 permutations of four figures. A group of degree $72 = 8 . 9$ must, by \S.\,30, II., since 8 is not $\equiv 1 \pmod{3}$, contain a divisor of degree 18, and would therefore, if it were simple, likewise be representable by the permutations of four figures, which is still less possible. That simple groups of degrees 84 and 100 cannot exist follows likewise at once from \gesperrt{Sylow}'s theorems, since, if we separate out the groups of degree $7$ or of degree $25$, no divisor remains that is congruent to 1 with respect to the modulus 7 or 5. That all these degree numbers can belong only to metacyclic groups then follows without further ado from \S.\,6, 2., from the fact that the groups whose degrees
%CONT>>


% ===== scan leaf 159, printed 138 =====
% !PRINTED 138
\weberpage{138}
%<<CONT
are proper divisors of these numbers, have already been proved to be metacyclic.

There remains only the investigation of the numbers 60 and 90. That a simple group of degree 60 exists we know already; namely the alternating group of the permutations of five letters (Bd.\,I, \S.\,177). But it is still questionable whether this is the only one.

A simple group $P$ of degree 60 contains, by \S.\,29, I., as divisor a group of degree 5 and a group of degree 3. If in the second \gesperrt{Sylow's} theorem (\S.\,30) one takes for $Q$ the group of degree 5, then $R$ must be of degree 10, because the index of $R$, being congruent to 1 with respect to the modulus 5, can only be $= 6$. Hence by \S.\,28, 2. the group $P$ of degree 60 can be represented as a transitive permutation group of six figures; but it cannot contain any cyclic permutation of only three figures, because as a simple group it must (by Bd.\,I, \S.\,158, 2.) be primitive, and therefore, if it contained a three-term cycle, it would by Bd.\,I, \S.\,153, 10. have to contain the whole alternating group of degree 360. We may therefore assume one of the permutations of degree 3 in the form:
\[
a = (1, 2, 3) (4, 5, 6).
\]
To this we will add a permutation of degree 5, $b$, which can only be a cyclic one. If we let the figure 6 be absent from it, then in place of $b$ we can take such a power of $b$ (which must of course also occur in $P$) that, say, 1, 2 become the first two figures, and then six possibilities remain for $b$:
\[
(1, 2, 3, 4, 5), \quad (1, 2, 4, 3, 5), \quad (1, 2, 4, 5, 3),
\]
\[
(1, 2, 3, 5, 4), \quad (1, 2, 5, 3, 4), \quad (1, 2, 5, 4, 3);
\]
of these six assumptions, however, the three standing in one row are again not essentially different (i.\,e. different only in the designation). For if, for example, in the second assumption $b = (1, 2, 4, 3, 5)$ one replaces the element $a$ by $a^2$ and $b$ by $b^3$, one obtains
\[
(1, 3, 2) (4, 6, 5); \quad (1, 3, 2, 5, 4),
\]
which, by the permutation of the figures 2 with 3 and 4 with 5, passes into the first assumption. The third assumption $b = (1, 2, 4, 5, 3)$ passes, if one replaces $b$ by $b^4$ and $a$ by $a^2$ and then
%CONT>>


% ===== scan leaf 160, printed 139 =====
% !PRINTED 139
\weberpage{139}
%<<CONT
1 with 2 and 4 with 5 are interchanged, into the first, and in the same way the three other assumptions about $b$ can be reduced to one another. There remain, then, only two possibilities to investigate:
\[
\begin{gathered}
a = (1, 2, 3) (4, 5, 6), \quad b = (1, 2, 3, 4, 5) \\
a = (1, 2, 3) (4, 5, 6), \quad b = (1, 2, 3, 5, 4).
\end{gathered}
\]
Of these, however, the first is to be rejected, because in
\[
a^2 b = (1, 4, 6)
\]
it would yield a three-term cycle, which cannot occur. There remains, then, only the single possibility:
\[
a = (1, 2, 3) (4, 5, 6), \quad b = (1, 2, 3, 5, 4),
\]
and there is therefore, if one regards equivalent groups as not different, only \gesperrt{one} simple group of the 60th degree.

These two elements $a, b$ may be regarded as generating elements of the whole group, and by their repeated composition one can derive the whole group of the 60th degree in a great variety of ways. Besides the unit element we obtain the permutations of the 5th degree:
\[
\begin{gathered}
(2, 3, 6, 5, 4), \quad (1, 3, 5, 6, 4), \quad (1, 2, 5, 4, 6), \\
(1, 2, 6, 3, 5), \quad (1, 2, 4, 6, 3), \quad (1, 2, 3, 5, 4),
\end{gathered}
\]
which with their powers yield 24; further the permutations of the 3rd degree:
\[
\begin{array}{ll}
(1, 2, 3) (4, 5, 6), & (1, 3, 5) (2, 4, 6), \\
(1, 2, 4) (3, 6, 5), & (1, 3, 6) (2, 4, 5), \\
(1, 2, 5) (3, 6, 4), & (1, 4, 5) (2, 3, 6), \\
(1, 2, 6) (3, 4, 5), & (1, 4, 6) (2, 3, 5), \\
(1, 3, 4) (2, 6, 5), & (1, 5, 6) (2, 3, 4),
\end{array}
\]
which together with their squares yield 20. Finally one obtains 15 more permutations of the 2nd degree:
\[
\begin{array}{lll}
(1, 2) (3, 4), & (1, 2) (5, 6), & (3, 4) (5, 6), \\
(1, 3) (4, 5), & (1, 3) (2, 6), & (2, 6) (4, 5), \\
(1, 4) (2, 5), & (1, 4) (3, 6), & (2, 5) (3, 6), \\
(1, 5) (2, 3), & (1, 5) (4, 6), & (2, 3) (4, 6), \\
(1, 6) (2, 4), & (1, 6) (3, 5), & (2, 4) (3, 5),
\end{array}
\]
of which each three standing in one row form, together with the unit element, a group of the 4th degree.

That this group can also be represented by the permutations of five figures now follows already from the fact
%CONT>>


% ===== scan leaf 161, printed 140 =====
% !PRINTED 140
\weberpage{140}
\pdfbookmark[2]{§ 35. Groups of degree p q}{bm41}
%<<CONT
that a simple permutation group of degree 60 on five figures actually exists, namely the alternating group.

But one can also verify the fact by exhibiting a divisor of degree 12 of $P$ (\S.\,28, 3.). One obtains this divisor of degree 12, e.g., from the fact that the group of degree 4:
\[
Q = 1, (1, 2) (3, 4), (1, 2) (5, 6), (3, 4) (5, 6)
\]
is permutable with the permutations of degree 3:
\[
c = (1, 3, 5) (2, 4, 6)
\]
that is, satisfies the condition $c^{-1} Q c = Q$, and one thus then obtains a group of degree 12:
\[
Q, \; Q c, \; Q c^2.
\]
The simple group of degree 60 is also called the \gesperrt{icosahedral group}, for a reason that we shall come to know later.

The same consideration also shows that a simple group of degree 90 does not exist. For this group would have to contain an element of degree 5 and an element of degree 3. Since $90 = 5 . 18$ and 18 has no divisor other than 6 that is congruent to 1 modulo 5, it follows that, if we choose for $Q$ a group of degree 5, $R$ must be of index 6, and the group $P$ would therefore be representable as a transitive permutation group of six figures. Just as above, one concludes that it would have to contain the two elements $a, b$, which, as we have seen, leads to the icosahedral group, which cannot be a divisor of a group of degree 90.

\begin{center}
\S.\,35.

\gesperrt{Groups of degree $pq$.}
\end{center}

It is now of further interest to investigate how, for a given degree number, the groups are constituted. For this the theorems derived in the preceding sections are the foundations. To be sure, we are as yet in a position to survey the existing groups completely only for the simpler degree numbers. A work by \gesperrt{Hölder}\footnote{\gesperrt{Hölder}, Die Gruppen der Ordnungen $p^3, pq^2, pqr, p^4$. Mathematische Annalen, Bd. 43.} goes furthest in this direction. We consider here only the simple case,
%CONT>>


% ===== scan leaf 162, printed 141 =====
% !PRINTED 141
\weberpage{141}
%<<CONT
that the degree $pq$ is the product of two primes, which may also be equal to one another.

Such a group $P$ is metacyclic and has a normal divisor $Q$ of degree $p$, which is therefore a cyclic group, which we represent thus:
\weq{1}{Q = 1, a, a^2, \dots a^{p-1} \quad (a^p = 1).}
If now $b$ is an element of $P$ not contained in $Q$, then
\[
b^{-1} Q b = Q.
\]
If $b^h$ is the lowest power of $b$ that occurs in $Q$, then $Q + Qb + \cdots Qb^{h-1}$ is a group, which must be identical with $P$, and consequently $h = q$ must hold. We can therefore represent $P$ thus:
\weq{2}{P = Q + Qb + Qb^2 + \cdots + Qb^{q-1}.}
If now $g$ is a primitive root of the prime $p$, and $\nu$ is an exponent that is for the present still undetermined, and is to be taken with respect to the modulus $p - 1$, then, since $b^{-1} a b$ is contained in $Q$, it follows that
\weq{3}{b^{-1} a b = a^{g^\nu},}
from which, for every exponent $\alpha$, there results:
\weq{4}{b^{-1} a^\alpha b = a^{\alpha g^\nu},}
and by repeated composition with $b$:
\weq{5}{b^{-\beta} a^\alpha b^\beta = a^{\alpha g^{\nu\beta}}.}
Now in any case $b^{pq} = 1$, and if we therefore set $\beta = pq$ in (5), it follows that for every $\alpha$
\[
a^{\alpha g^{\nu pq}} = a^\alpha,
\]
hence $g^{\nu pq} \equiv 1 \pmod{p}$, or
\weq{6}{\nu pq \equiv 0 \pmod{p - 1}}
must hold. From this one concludes further that, if $p - 1$ is not divisible by $q$, which always occurs when $p = q$, then $\nu$ is divisible by $p - 1$, and hence by (3)
\[
ab = ba,
\]
and in consequence also, for every $\alpha, \beta, \alpha', \beta'$:
\[
a^\alpha b^\beta = b^\beta a^\alpha, \quad a^\alpha b^\beta a^{\alpha'} b^{\beta'} = a^{\alpha'} b^{\beta'} a^\alpha b^\beta,
\]
i.e. that the group must be an \gesperrt{Abel}ian one. But we have already investigated this case in the second chapter of this volume and found that, when $p = q$, there are two groups (with
%CONT>>


% ===== scan leaf 163, printed 142 =====
% !PRINTED 142
\weberpage{142}
%<<CONT
the two invariants $p, p$ or with one invariant $p^2$), and,
if $p$ is different from $q$, one group (with the two invariants
$p, q$).

There thus remains for us the case
\[
p \equiv 1\ (\mathrm{mod}\ q)
\]
still to investigate, in which $\nu$ can have any value satisfying the condition
\[
\nu \equiv 0\ \left(\mathrm{mod}\ \frac{p-1}{q}\right)
\]
of which there are $q$ different ones with respect to the modulus
$p - 1$, namely:
\weq{7}{\nu = 0,\quad \frac{p-1}{q},\quad 2\,\frac{p-1}{q},\quad \dots\ \frac{(q-1)(p-1)}{q}.}
The case $\nu = 0$ leads here too to an \gesperrt{Abelian} group.
But the other values of $\nu$ are also admissible, and they lead
to just as many non-commutative groups. These non-commutative
groups are isomorphic to one another, and are reduced
to one another when one replaces $b$ by a $b^\beta$,
as the formula (5) shows.

We may assume $b^q = 1$, for $b^q$ is certainly contained in $Q$,
and $b^{pq} = 1$. If, then, $b^q$ is not $= 1$, we replace
$b$ by $b^p$.

One obtains one of the non-commutative groups by letting,
in
\weq{8}{\Theta = a^\alpha b^\beta}
$\alpha$ run through the sequence of numbers $0, 1, \dots p - 1$ and $\beta$ through the numbers
$0, 1, \dots q - 1$. The composition of two
elements of this group results, by (5), from
\weq{9}{b^\beta a^\alpha = a^{\alpha g^{-\nu\beta}} b^\beta, \quad a^p = 1,\ b^q = 1,}
by which every compositum of elements $\Theta$ can be reduced to the
form $\Theta$.

To show that the elements (8) really form a group under
the rules of composition (9), one has to prove the properties
of the group, namely that from $\Theta\Theta' = \Theta\Theta''$ and from $\Theta'\Theta = \Theta''\Theta$
it follows that $\Theta' = \Theta''$, and further the associative law
\[
\Theta(\Theta'\Theta'') = (\Theta\Theta')\Theta''
\]
Both, however, follow very easily from the composition
that follows from (9):
\weq{10}{a^\alpha b^\beta a^{\alpha'} b^{\beta'} = a^{\alpha+\alpha' g^{-\nu\beta}} b^{\beta+\beta'}.}


% ===== scan leaf 164, printed 143 =====
% !PRINTED 143
\weberpage{143}
\pdfbookmark[2]{§ 36. Bounds for the index of a divisor of the symmetric permutation group}{bm42}
\begin{center}
\S.\,36.

\gesperrt{Limits of the index of a divisor of the symmetric
permutation group.}
\end{center}

We conclude these considerations with the proof of a
theorem on permutation groups which, through the difficulty
that its proof at first presented, has attained a certain celebrity,
and which is of importance for the judgement of algebraic
questions\,$^{1}$).

It concerns the symmetric permutation group
 $P$ of $n$ figures and its divisors of smallest
possible index. We know that the group $P$ always has a
divisor of index 2, namely the alternating group.
Besides this, a divisor of index $n$ is also known, which comprises all
permutations of $P$ that leave one figure unchanged,
and which is therefore intransitive.

First of all the following theorem holds:

I. \gesperrt{The index of an imprimitive divisor of $P$ is
always greater than $n$, and the index of an intransitive
 divisor is equal to or greater
than $n$, and equal to $n$ only when the divisor
leaves one figure at rest and permutes the remaining
$n - 1$ figures in all possible ways.}

Let us assume that $Q$ is an imprimitive divisor of $P$
of index $j$, and that there are $r$ systems of imprimitivity of
$s$ figures each, so that $n = rs$. A number which the degree of the
group $Q$ certainly cannot exceed is obtained if we
count all the permutations in each single one of the $r$ systems and then
also all the permutations of the systems. The degree
of $Q$ is therefore less than or equal to the product $[\Pi(s)]^r\,\Pi(r)$,
if for every integer $n$
\[
\Pi(n) = 1.2.3\dots n
\]
%CONT>>

$^{1}$) \gesperrt{Bertrand}, Journal de Mathématiques, Tome XV (1845). \gesperrt{Serret},
Algèbre supér., Section IV, Chapitre III. C. \gesperrt{Jordan}, Traité des substitutions,
p.\,67. \gesperrt{Netto}, Substitutionentheorie, Capitel VI. Crelle's Journ.,
Bd. 100.


% ===== scan leaf 165, printed 144 =====
% !PRINTED 144
\weberpage{144}
%<<CONT
and consequently
\weq{1}{j \geqq \frac{\Pi(n)}{[\Pi(s)]^r\,\Pi(r)}.}
It is easy to see that this number, when neither of the
factors $r, s$ is equal to 1, is greater than $n$. This follows
if we write the quotient thus:
\[
\frac{\Pi(n)}{[\Pi(s)]^r\,\Pi(r)} = \frac{(r+1)(r+2)\dots n}{(2.3\dots s)^r} =
\]
\[
\frac{n}{2}\cdot\left(\frac{r+1}{2}\cdot\frac{r+2}{2}\dots\frac{2r-1}{2}\right)\left(\frac{2r}{3}\cdot\frac{2r+1}{3}\dots\frac{3r-1}{3}\right)\dots
\]
\[
\left(\frac{(s-1)r}{s}\cdot\frac{(s-1)r+1}{s}\dots\frac{n-1}{s}\right).
\]
For according to this
\weq{2}{j > \frac{n}{2}\left(\frac{r+1}{2}\right)^{r-1}\left(\frac{2r}{3}\right)^r \dots \left(\frac{(s-1)r}{s}\right)^r.}
The factors $\dfrac{r+1}{2}, \dfrac{2r}{3}, \dots \dfrac{(s-1)r}{s}$ are all equal to or
greater than 1, for we have
\[
\frac{(h-1)r}{h} - 1 = \frac{(h-1)(r-1)-1}{h},
\]
which is therefore always positive, and the first factor
\[
\left(\frac{r+1}{2}\right)^{r-1}
\]
is greater than 2 when $r > 2$. But if $r = 2$, then the
product of the two factors is:
\[
\left(\frac{r+1}{2}\right)^{r-1}\left(\frac{2r}{3}\right)^r = \frac{3}{2}\left(\frac{4}{3}\right)^2 = \frac{8}{3} > 2,
\]
and consequently in all circumstances $j > n$.

If, secondly, the group $Q$ is intransitive, and the
system of the $n$ figures splits into two systems of $a$ and $b$ figures respectively, so
that the figures of these two systems are permuted by $Q$ only among
 themselves, and $n = a + b$, then all the permutations
 of $Q$ are contained in the group that consists of all the permutations
 of the $a$ and of the $b$ figures; that is, the degree
of $Q$ is equal to or less than $\Pi(a)\Pi(b)$, and consequently the
index $j$ of $Q$
\weq{3}{j \geqq \frac{\Pi(n)}{\Pi(a)\,\Pi(b)} = \frac{n}{1}\,\frac{n-1}{2}\dots\frac{b+1}{a}.}


% ===== scan leaf 166, printed 145 =====
% !PRINTED 145
\weberpage{145}
If we assume, as we are free to do, that $b \geqq a$, then the
expression on the right side of this inequality is greater
than $n$, and equal to $n$ only when $b = n - 1, a = 1$. In
this special case $j = n$ can occur, but only when
the $n - 1$ elements are permuted by $Q$ in all possible ways,
which is expressed in Theorem I.

The expressions (3) for the lower limit of $j$ are nothing
other than the binomial coefficients $B_a^{(n)}$, whose formation
shows at once that up to the middle, i.e. as long as $2a \leqq n$,
they form an \gesperrt{increasing sequence of numbers}.

For further use we shall draw from this the
conclusion:

$\alpha$) \gesperrt{If $a = 1$, that is, if the group $Q$ leaves one figure
unchanged, then its degree is a divisor of
$\Pi(n-1)$; but if $a > 1$, then the degree of $Q$ is
$\leqq \Pi(n-2)\,\Pi(2)$.}

The theorem which we further wish to prove now reads:

\gesperrt{II. Besides the alternating group there is no
transitive and primitive divisor $Q$ of $P$
whose index is $\leqq n$, except in the
two cases $n = 4, n = 6$.}

In the proof we make use of the theorems (\S.\,153,
9., 10. of the first volume) that a transitive and primitive
divisor of $P$ which does not contain the whole alternating group,
and hence is not identical with the alternating or the symmetric
group itself, can contain no transposition and no
cyclic permutation of only three figures.

Let $P$ then be the symmetric permutation group of the
$n$ figures $0, 1, 2 \dots n - 1$, and $Q$ a primitive and transitive
divisor of $P$ of index $j$ which does not contain the alternating
group. Further, let $P$ be decomposed into the subsidiary groups:
\weq{4}{P = Q + Q\pi_1 + Q\pi_2 + \dots + Q\pi_{j-1}.}
We consider the system of transpositions:
\weq{5}{(0, 1), (0, 2), \dots (0, n-1),}
none of which can occur in $Q$. Now if $j < n$, then at least
two of these permutations, say $(0,1), (0,2)$, must occur in the
%CONT>>


% ===== scan leaf 167, printed 146 =====
% !PRINTED 146
\weberpage{146}
%<<CONT
the same subsidiary group, say in $Q\pi_1$, and consequently there are two
permutations $\varkappa_1, \varkappa_2$ in $Q$ which satisfy the condition

\[ \varkappa_1 \pi_1 = (0, 1), \quad \varkappa_2 \pi_1 = (0, 2) \]

It is therefore

\[ \varkappa_1 \pi_1 \pi_1^{-1} \varkappa_2^{-1} = \varkappa_1 \varkappa_2^{-1} = (0,1)(0,2) = (0,1,2) \]

contained in $Q$. But this contradicts our assumption that $Q$ is to contain no
three-term cycle. Accordingly $j$ cannot be $< n$. But if $j = n$, and hence the
degree of $Q$ equal to $\Pi(n-1)$, then the $n-1$ permutations (5) must occur in
$n-1$ different subsidiary groups, and $P$ can be represented thus:

\weq{6}{P = Q + Q(0,1) + Q(0,2) + \cdots + Q(0, n-1).}

If we consider any other transposition, e.g. $(2,3)$, this cannot occur
in $Q$, nor in $Q(0,2)$ or in $Q(0,3)$, since otherwise
$(2,3)(0,2) = (0,2,3)$ or $(2,3)(0,3) = (0,3,2)$ would occur in $Q$. We may therefore,
without impairing the generality, assume that $(2,3)$ occurs in $Q(0,1)$,
and from this it follows:

\gesperrt{$\beta$) The group $Q$ contains the transposition pair}

\weq{7}{(0,1)(2,3).}

The group $Q$ has a divisor $Q_0$, which consists of all those permutations
that leave the figure 0 in its place. Since $Q$ is transitive, it follows that if
$\varkappa_1, \varkappa_2, \ldots, \varkappa_{n-1}$ are elements of $Q$ which carry the 0
into 1, into 2, \ldots, into $n-1$, then

\weq{8}{Q = Q_0 + Q_0 \varkappa_1 + \cdots + Q_0 \varkappa_{n-1},}

and since $Q$ is of degree $\Pi(n-1)$, it follows from this that $Q_0$ is of degree

\weq{9}{g = \frac{\Pi(n-1)}{n}}

Since $g$ must be an integer, and hence $\Pi(n-1)$ divisible by $n$, we
conclude first of all that the case $j = n$ can never occur when $n$
is a prime, and for this case our Theorem II. is therefore proved.

In general, however, we can conclude:

\gesperrt{$\gamma$) If $n$ is not $= 4$, then the $n-1$ figures
$1, 2, \ldots, n-1$ are still transitively connected by $Q_0$.}


% ===== scan leaf 168, printed 147 =====
% !PRINTED 147
\weberpage{147}
For if $Q_0$ were intransitive in these $n-1$ figures, then by $\alpha$)

\[
\begin{aligned}
&\text{either } \Pi(n-2) \text{ divisible by } g,\\
&\text{or } \Pi(n-3)\,\Pi(2) \geqq g
\end{aligned}
\]

would have to hold; hence by (9):

\[
\begin{aligned}
&\text{either } \frac{n}{n-1} \text{ an integer, hence } n \geqq 2(n-1)\\
&\text{or } n^2 - 5n + 2 \leqq 0.
\end{aligned}
\]

The one is impossible when $n > 2$, the other when $n > 4$.

We set aside the case $n = 4$ and now consider the divisor $Q_{0,1}$ of
$Q$ which leaves the two figures 0, 1 unchanged.

Since $Q_{0,1}$, regarded as a permutation group of the $n-1$ figures, is a
transitive divisor of $Q_0$, the degree $g_1$ of $Q_{0,1}$ is, as one concludes
by applying the decomposition (8) once more, equal to $g : n-1$; hence

\weq{10}{g_1 = \frac{\Pi(n-1)}{n(n-1)} = \frac{\Pi(n-2)}{n}.}

From this we conclude, similarly as above:

\gesperrt{$\delta$) If $n$ is not $= 6$, then the $n-2$ figures
$2, 3, \ldots, n-1$ are transitively connected by $Q_{0,1}$.}

For if they were not, then by $\alpha$):

\[
\begin{aligned}
&\text{either } \Pi(n-3) \text{ divisible by } g_1,\\
&\text{or } \Pi(n-4)\,\Pi(2) \geqq g_1
\end{aligned}
\]

would have to hold; hence by (10):

\[
\begin{aligned}
&\text{either } \frac{n}{n-2} \text{ an integer, hence } n \geqq 2n - 4,\\
&\text{or } n^2 - 7n + 6 \leqq 0.
\end{aligned}
\]

The first is not possible when $n > 4$, the second when $n > 6$.

Now if $n > 6$, the following also holds:

\gesperrt{$\varepsilon$) The group $Q_{0,1,2}$, which leaves the three figures 0, 1, 2
unchanged, has the degree}

\[ g_2 = \frac{g_1}{n-2} = \frac{\Pi(n-3)}{n}, \]

\gesperrt{and it is not possible that yet a fourth figure 3 should remain unchanged
under the whole group $Q_{0,1,2}$.}


% ===== scan leaf 169, printed 148 =====
% !PRINTED 148
\weberpage{148}
The degree $g_2$ is obtained exactly like that of $Q_0$ and $Q_{0,1}$. But if
under $Q_{0,1,2}$ a fourth figure 3 were also to remain unchanged, then $g_2$
would be a divisor of the degree of the symmetric group of $n-4$ figures, hence
a divisor of $\Pi(n-4)$. Thus

\[ \frac{n\,\Pi(n-4)}{\Pi(n-3)} = \frac{n}{n-3} \]

would have to be an integer, hence $n \geqq 2n - 6$ or $n \leqq 6$.

From $\varepsilon$) we conclude that there is in $Q$ a permutation $\varkappa$
by which $0, 1, 2, 3$ passes into $0, 1, 2, 4$, where 4 is a figure different
from 3. By $\beta$), then, $Q$ also contains the permutation:

\[ \varkappa^{-1}(0,1)(2,3)\varkappa = (0,1)(2,4), \]

and consequently also:

\[ (0,1)(2,3)(0,1)(2,4) = (2,3,4), \]

that is, a three-term cycle, which contradicts the hypothesis. Hereby
Theorem II. is completely proved.

That the cases $n = 4$ and $n = 6$ really form exceptions follows from Bd. I,
\S.\,160 and \S.\,182, where we have seen \gesperrt{that the symmetric permutation group
of four figures possesses a transitive divisor of index 3 and that of six
figures a transitive divisor of index 6.}

\begin{center}
\rule{3cm}{0.4pt}
\end{center}


% ===== scan leaf 170, printed 149 =====
\weberpage{149}% page number not printed on this page
% !PRINTED none
% Part-title (Buch-Titel) page: mostly blank, with two centred rules framing
% the book/part heading. No running head, no printed page number.
\begin{center}
\gesperrt{SECOND BOOK.}

\rule{3cm}{0.4pt}

\vspace{2cm}

\gesperrt{LINEAR GROUPS.}

\vspace{2cm}

\rule{5cm}{0.4pt}
\end{center}


% ===== scan leaf 171, printed 150 =====
\weberpage{150}% page number not printed on this page
% !PRINTED none
% Blank page (verso). All three bands show only faint bleed-through ghosting
% from the facing/adjacent printed page; no legible running head, page
% number, or body text is present anywhere on this page image.


% ===== scan leaf 172, printed 151 =====
\weberpage{151}% page number not printed on this page
\pdfbookmark[0]{Second Book. Linear Groups.}{bm43}
\pdfbookmark[1]{Sixth Section. Groups of Linear Substitutions.}{bm44}
\pdfbookmark[2]{§ 37. Linear substitutions and their composition}{bm45}
% !PRINTED none
% Chapter-opening page: no running head/folio visible in any of the three
% bands (consistent with an Abschnitt-opening page).

\begin{center}
\gesperrt{Sixth Section.}

\gesperrt{Groups of Linear Substitutions.}
\end{center}

\begin{center}
\rule{3cm}{0.4pt}
\end{center}

\begin{center}
§. 37.

\gesperrt{Linear Substitutions and their Composition.}
\end{center}

One of the most effective means for the formation of groups, to which at the same time many applications lead, is furnished by the linear substitutions and their composition. We have already met with such linear substitutions repeatedly, and have considered them, for example, in the second chapter of the first volume on the occasion of the law of multiplication of determinants, and then in the eleventh chapter in connection with equivalent numbers.

By a \gesperrt{linear substitution} of $n$ variables we understand a system of equations by which a system of $n$ variables $y_1, y_2, \dots y_n$ is expressed linearly in terms of another system $x_1, x_2, \dots x_n$. According to the number of the variables we distinguish \gesperrt{unary, binary, ternary, quaternary} substitutions, and in general we shall call the number of the variables the \gesperrt{dimension} of the substitution. We confine ourselves for the present to \gesperrt{homogeneous substitutions}, so that, if the coefficients are denoted by $a_i^{(\varkappa)}$, the substitution is represented by the system of equations

\weq{1}{y_\varkappa = \sum_{1,n}^{i} a_i^{(\varkappa)} x_i \qquad \varkappa = 1, 2, \dots n}

Often the variables themselves do not matter, so that such a substitution is sufficiently characterised by its coefficients $a_i^{(\varkappa)}$. One sometimes also uses
%CONT>>


% ===== scan leaf 173, printed 152 =====
% !PRINTED 152
\weberpage{152}
%<<CONT
a simple letter, say $A$, to denote a substitution, and then sets

\weq{2}{A = \begin{pmatrix} a_1^{(1)}, & a_2^{(1)}, & \dots & a_n^{(1)} \\ a_1^{(2)}, & a_2^{(2)}, & \dots & a_n^{(2)} \\ \multicolumn{4}{c}{\cdots\cdots\cdots\cdots} \\ a_1^{(n)}, & a_2^{(n)}, & \dots & a_n^{(n)} \end{pmatrix}.}

We shall also write, in abbreviation, $A = (a_i^{(\varkappa)})$, and denote the determinant of the substitution coefficients by $|A|$. This determinant is always assumed to be different from zero. We also represent the system of equations (1) symbolically thus:

\weq{3}{(y_1, y_2, \dots y_n) = A (x_1, x_2, \dots x_n),}

or still more briefly:

\weq{4}{(y) = A (x).}

The substitution:

\weq{5}{y_1 = x_1, \quad y_2 = x_2, \quad \dots \quad y_n = x_n}

or

\weq{6}{J = \begin{pmatrix} 1, & 0, & \dots & 0 \\ 0, & 1, & \dots & 0 \\ \multicolumn{3}{c}{\cdots\cdots\cdots} \\ 0, & 0, & \dots & 1 \end{pmatrix}}

is called the \gesperrt{identical substitution}. A substitution of the form

\weq{7}{y_1 = \mu_1 x_1, \quad y_2 = \mu_2 x_2, \quad \dots, \quad y_n = \mu_n x_n}

or

\weq{8}{M = \begin{pmatrix} \mu_1, & 0, & \dots & 0 \\ 0, & \mu_2, & \dots & 0 \\ \multicolumn{3}{c}{\cdots\cdots\cdots} \\ 0, & 0, & \dots & \mu_n \end{pmatrix}}

shall be called a \gesperrt{multiplicative substitution}, or briefly a \gesperrt{multiplication}, and the elements $\mu_1, \mu_2, \dots, \mu_n$ the \gesperrt{multipliers}, an expression which is justified by the equations (7).

If we assume a second linear substitution of dimension $n$, by which a new system of variables $z$ is introduced:

\weq{9}{z_i = \sum_{1,n}^{h} b_h^{(i)} y_h, \quad (z) = B (y),}

and introduce for $y$ the values from (1), we obtain the $z$
%CONT>>


% ===== scan leaf 174, printed 153 =====
% !PRINTED 153
\weberpage{153}
%<<CONT
expressed by the $x$ by means of a new linear substitution $E$, which we may denote thus:

\weq{10}{z = E (x) = B A (x),}

and the substitution $E = BA$ is said to be \gesperrt{composed} or \gesperrt{compounded} of $B$ and $A$. But a distinction must here be made between $AB$ and $BA$. Two substitutions $A, B$ with the special property that $AB = BA$ are said to be \gesperrt{permutable} or \gesperrt{commutative} with one another. The substitution coefficients of $E$ are obtained by inserting the expressions (1) in (9):

% [sic: Weber prints i below the first sum although h is its summation index; 1,n is expected.]
\weq{11}{\begin{gathered} z_k = \sum_{i}^{h} e_h^{(k)} x_h, \\ e_h^{(k)} = \sum_{1,n}^{i} b_i^{(k)} a_h^{(i)}. \end{gathered}}

These formulae are exactly the same as those we used in \S. 27 of the first volume\footnote{Only there, for an easily discernible reason, the notation was somewhat different.}, and we can accordingly express the law of the composition of substitutions expressed in (11) as follows:

\gesperrt{1. To form the substitution $BA$ composed of two substitutions $B, A$, one proceeds exactly as if the two determinants $|B|, |A|$ were to be multiplied together according to the rule of multiplication. Here, in order to form the elements of a row, the elements of a row of the first component are to be multiplied by the corresponding elements of the columns of the second component and then added. This is also expressed briefly by saying that in the first component the summation is taken by rows, in the second by columns.}

From this rule follows the consequence:

\gesperrt{2. The determinant of a composed substitution is equal to the product of the determinants of the components.}


% ===== scan leaf 175, printed 154 =====
% !PRINTED 154
\weberpage{154}
To compose three substitutions of the same dimension, $C$, $B$, $A$, one must introduce the expressions (10) for $z$ into a new linear substitution

\weq{12}{(u) = C(z)}

and express the variables $u$ in terms of the $x$:

\weq{13}{(u) = CBA(x).}

Since it is evidently immaterial whether one introduces the expressions (10) into (12), or whether one first expresses $z$ in terms of $y$ by (9) and then $y$ in terms of $x$ by (4), the \gesperrt{associative law} holds for this composition:

\weq{14}{C(BA) = (CB)A = CBA,}

which can also easily be confirmed by calculation, if one forms the elements of $C(BA)$ and $(CB)A$ according to (11). For both one finds the expression:

\[
\sum^i \sum^h c_i^{(k)}\, b_h^{(i)}\, a_l^{(h)}.
\]

We therefore state the theorem:

\par
3.\enspace\gesperrt{In the composition of the linear substitutions the associative law holds, but not always the commutative law.}

A multiplication in which all the multipliers are equal to one another, that is

\[
N = \begin{pmatrix} \nu, & 0, & \ldots & 0 \\ 0, & \nu, & \ldots & 0 \\ \cdot & \cdot & \cdot & \cdot \\ 0, & 0, & \ldots & \nu \end{pmatrix}
\]

shall be called a \gesperrt{similarity substitution}, and two substitutions which, like $A$ and $AN$ or $A$ and $NA$, can be derived from one another by composition with a similarity substitution, are called \gesperrt{similar substitutions}. From the law of composition the following theorems result at once:

\par
4.\enspace\gesperrt{A similarity substitution is permutable with every substitution $A$ of the same dimension; two similarity substitutions, compounded with one another, again give a similarity substitution, and two substitutions similar to a third are also similar to one another.}


% ===== scan leaf 176, printed 155 =====
% !PRINTED 155
\weberpage{155}
Just as easily we obtain:

\par
5.\enspace\gesperrt{By composition with the identical substitution $J$ no substitution is changed.}

The substitution $J$ is therefore regarded in composition as a \gesperrt{unit}, and, wherever it occurs compounded with other substitutions, it may be omitted.

Now the following theorem holds further:

\par
6.\enspace\gesperrt{To every substitution $A$ there is one and only one inverse substitution $A^{-1}$, by which the condition}

\weq{15}{AA^{-1} = A^{-1}A = J}

\gesperrt{is satisfied.}

This last theorem follows from the fundamental formulae of the theory of determinants, if one determines the elements $\alpha_h^{(k)}$ of $A^{-1}$ from the linear equations

\weq{16}{\begin{aligned} \sum^i a_i^{(h)}\alpha_k^{(i)} &{}= 0, & h \gtrless k\\ &{}= 1, & h = k \end{aligned}}

from which, if we denote by $A_h^{(k)}$ the subdeterminants of $|A|$, there results

\weq{17}{|A|\,\alpha_h^{(k)} = A_k^{(h)}}

and which entail the other system

\weq{18}{\begin{aligned} \sum^i \alpha_i^{(h)} a_k^{(i)} &{}= 0, & h \gtrless k\\ &{}= 1, & h = k \end{aligned}}

The inverse substitution to $A$ is nothing other than the solution of the system of equations (1) of the direct substitution $A$, so that from $(y) = A(x)$ it follows that:

\weq{19}{(x) = A^{-1}(y).}

For the composition of the inverse substitutions there results from $ABB^{-1}A^{-1} = J$ the theorem:

\weq{20}{(AB)^{-1} = B^{-1}A^{-1}.}

If $A$, $B$, $C$ are three substitutions, then by composition with $A^{-1}$ it follows from each of the two equations

\[
AB = AC, \quad BA = CA,
\]

that $B = C$ must hold. Accordingly, for the composition of the substitutions the characteristic marks of a group are fulfilled, and thus the totality of all substitutions is
%CONT>>


% ===== scan leaf 177, printed 156 =====
% !PRINTED 156
\weberpage{156}
%<<CONT
of definite dimension form an (infinite) \gesperrt{group} (\S.\,1). If we pick out of this totality any set which is closed in itself, in such a way that any two of its elements yield by composition an element of the same set, then this set likewise forms a \gesperrt{group}, which may be finite or infinite.

If $L$ denotes a fixed substitution, then from every substitution $A$ of the same dimension one can derive a substitution

\weq{21}{A' = L^{-1}AL}

which is called the \gesperrt{transform of $A$ by $L$}.

If we set

\weq{22}{(x) = L(x'), \quad (y) = L(y'),}

then it follows from (4):

\weq{23}{(y') = A'(x'),}

so that the passage to the transformed substitution is equivalent to the simultaneous transformation of both sequences of variables by $L^{-1}$.

If

\[
A' = L^{-1}AL, \quad B' = L^{-1}BL,
\]

then from the laws of composition it follows that

\[
A'B' = L^{-1}ABL,
\]

and from this therefore the theorem:

\par
7.\enspace\gesperrt{If $A$ runs through the substitutions of a group, then for fixed $L$ the transform $L^{-1}AL$ runs through the substitutions of an isomorphic group.}

The \gesperrt{transposed substitutions} must be carefully distinguished from the inverse substitutions.

Namely, from every substitution $A$ one obtains another definite substitution, called the transposed substitution to $A$, which we shall for the moment denote by $A_1$, by making the rows of $A$ into columns and conversely; that is, by interchanging in (2) the upper and lower indices of the $a$.

If in the sum (11), $\sum^i b_i^{(k)} a_h^{(i)}$, one interchanges the lower with the upper indices and at the same time $a$ with $b$, one obtains an expression which by (11) is equal to $c_k^{(h)}$.


% ===== scan leaf 178, printed 157 =====
% !PRINTED 157
\weberpage{157}
This remark gives the rule according to which the transposed substitutions are composed, which is expressed in the theorem:

\par
8.\enspace\gesperrt{If $A_1$, $B_1$ are the transposes of $A$, $B$, then $A_1 B_1$ is the transpose of $BA$. In symbols:}

\weq{24}{(BA)_1 = A_1 B_1.}

If we multiply the substitution formulae (1) by an undetermined factor $\eta_k$ and then take the sum with respect to $k$, it follows that:

\weq{25}{\sum y_k \eta_k = \sum^i \sum^k x_i\, a_i^{(k)}\, \eta_k,}

or, if we

\weq{26}{\sum^k a_i^{(k)} \eta_k = \xi_i}

set,

\weq{27}{\sum y_k \eta_k = \sum x_i \xi_i.}

If, then, we represent the substitutions (1) by (4), the formulae (26) are the expression for the substitution

\[
(\xi) = A_1(\eta),
\]

and we can therefore also state the theorem:

\par
9.\enspace\gesperrt{If $A$ and $A_1$ are transposes of one another of dimension $n$, and if $x$, $y$, $\xi$, $\eta$ are four systems of variables which are related to one another by the substitutions}

\weq{28}{(y) = A(x), \quad (\xi) = A_1(\eta)}

\gesperrt{then the identity holds}

\weq{29}{y_1\eta_1 + y_2\eta_2 + \cdots y_n\eta_n = x_1\xi_1 + x_2\xi_2 + \cdots x_n\xi_n.}

If two sequences of variables

\[
\begin{gathered} x_1, x_2, \ldots x_n \\ \xi_1, \xi_2, \ldots \xi_n \end{gathered}
\]

are transformed simultaneously by the substitutions (28) into two new sequences

\[
\begin{gathered} y_1, y_2, \ldots y_n \\ \eta_1, \eta_2, \ldots \eta_n \end{gathered}
\]

then the two sequences of variables are called \gesperrt{contragredient} to one another, and the formulae (28) represent two mutually \gesperrt{contragredient transformations}.


% ===== scan leaf 179, printed 158 =====
% !PRINTED 158
\weberpage{158}
\pdfbookmark[2]{§ 38. Substitution of the ratios}{bm46}

\begin{center}
\S.\,38.

\gesperrt{Substitution of the Ratios.}
\end{center}

We have already seen above that two substitutions of the same dimension which are similar to a third are similar to one another. Accordingly we can unite all mutually similar substitutions of dimension $n$ into one class, and every substitution can be placed in one and only one such class. The individual substitutions of a class are called the \gesperrt{representatives} of the class, and every class is completely determined by any one of its representatives.

If $A$ is similar to $A'$ and $B$ similar to $B'$, then $AB$ is also similar to $A'B'$.

If, then, we denote by $\mathfrak{A}, \mathfrak{B}$ the classes to which $A, A'$ and $B, B'$ belong, one always arrives at the same class, whatever representatives from $\mathfrak{A}$ and from $\mathfrak{B}$ one may compose. This class, which is represented by $AB$ or $A'B'$, we therefore call composed of $\mathfrak{A}$ and $\mathfrak{B}$, and we denote it by $\mathfrak{A}\mathfrak{B}$. In this composition the same rules hold as in the composition of the substitutions themselves, and the totality of the classes therefore likewise forms a group.

This conception of the substitutions and their composition is always expedient when, as for example in projective geometry, only the ratios of the variables matter, and we therefore designate the substitution classes as \gesperrt{substitutions of the ratios.}

We shall, however, denote these substitutions of the ratios in the same way as the other substitutions, namely each class by a representative, whereby a misunderstanding will not easily arise. The representative can be chosen from very different points of view.

It is often advisable to take it so that the determinant is $= 1$. But in that case, for a substitution of dimension $n$, the multiplicative substitution $(\mu)$ is determined only up to an $n$th root of unity. If the representatives are so chosen, this property is preserved under composition, and these representatives of the classes therefore also form among themselves a substitution group.


% ===== scan leaf 180, printed 159 =====
% !PRINTED 159
\weberpage{159}
\pdfbookmark[2]{§ 39. Permutations as linear substitutions}{bm47}

In the region of binary substitutions, the substitution $A$:
\[
(y_1, y_2) = \begin{pmatrix} a, & b \\ c, & d \end{pmatrix} (x_1, x_2)
\]
regarded as a substitution of the ratios, is equivalent to the \gesperrt{linear fractional substitution}:
\[
\eta = \frac{a\,\xi + b}{c\,\xi + d},
\]
if we set $\xi = x_1 : x_2$ and $\eta = y_1 : y_2$; if a second substitution $A'$:
\[
(z_1, z_2) = \begin{pmatrix} a', & b' \\ c', & d' \end{pmatrix} (y_1, y_2),
\]
or
\[
\zeta = \frac{a'\eta + b'}{c'\eta + d'}
\]
is added, one obtains the composite substitution
\[
A'' = A'A,
\]
by which $\zeta$ is expressed in terms of $\xi$:
\[
\zeta = \frac{a''\xi + b''}{c''\xi + d''},
\]
according to the rules of composition, in the form
\[
\begin{pmatrix} a'', & b'' \\ c'', & d'' \end{pmatrix} = \begin{pmatrix} a', & b' \\ c', & d' \end{pmatrix} \begin{pmatrix} a, & b \\ c, & d \end{pmatrix} = \begin{pmatrix} a'a + b'c, & a'b + b'd \\ c'a + d'c, & c'b + d'd \end{pmatrix}.
\]
If such a substitution is multiplicative, that is, if it has the form
\[
\begin{pmatrix} a, & 0 \\ 0, & d \end{pmatrix},
\]
then we shall also call the ratio $a : d$ the \gesperrt{multiplier}.

\begin{center}
\S.\,39.

\gesperrt{Permutations as linear substitutions.}
\end{center}

The importance of the linear substitutions, and especially of the \gesperrt{finite} groups formed from them, for algebra results from the fact that the permutation groups of $n$ elements can be regarded as special cases of such substitution groups.


% ===== scan leaf 181, printed 160 =====
% !PRINTED 160
\weberpage{160}

Namely, if we denote by $x_1, x_2, \dots, x_n$ a system of $n$ variables, and by $\alpha_1, \alpha_2, \dots \alpha_n$ any arrangement of the $n$ figures $1, 2, \dots n$, then the equations:
\weq{1}{x_1 = x'_{\alpha_1}, \quad x_2 = x'_{\alpha_2}, \quad \dots \quad x_n = x'_{\alpha_n}}
determine a linear substitution of dimension $n$:
\weq{2}{A = \begin{pmatrix} a_1^{(1)}, & \dots & a_n^{(1)} \\ \cdot & \cdot & \cdot \\ a_1^{(n)}, & \dots & a_n^{(n)} \end{pmatrix}, \quad (x) = A(x'),}
in which in every row and in every column only one coefficient is different from zero, and this one has the value 1.

The determinant $|A|$ of the substitution is therefore $= \pm 1$. But if, after carrying out the substitution, one puts $x_i$ again for $x'_i$, the result is nothing other than the permutation
\[
\begin{pmatrix} 1, & 2, & \dots & n \\ \alpha_1, & \alpha_2, & \dots & \alpha_n \end{pmatrix}
\]
of the indices of $x$.

If $B$ is a second substitution formed in the same way
\weq{3}{(x') = B(x''),}
or more explicitly:
\weq{4}{x'_1 = x''_{\beta_1}, \quad x'_2 = x''_{\beta_2}, \quad \dots, \quad x'_n = x''_{\beta_n},}
then the composition, by the rules of \S.\,37, yields:
\weq{5}{x_1 = x''_{\beta_{\alpha_1}}, \quad x_2 = x''_{\beta_{\alpha_2}}, \quad \dots \quad x_n = x''_{\beta_{\alpha_n}},}
for which the abbreviated notation
\weq{6}{(x) = AB(x'')}
is to be used.

But by Bd.\,I, \S.\,148 (by the composition of the permutations):
\weq{7}{\begin{pmatrix} 1, & 2, & \dots & n \\ \alpha_1, & \alpha_2 & \dots & \alpha_n \end{pmatrix} \begin{pmatrix} 1, & 2, & \dots & n \\ \beta_1, & \beta_2 & \dots & \beta_n \end{pmatrix} = \begin{pmatrix} 1, & 2, & \dots & n \\ \beta_{\alpha_1}, & \beta_{\alpha_2} & \dots & \beta_{\alpha_n} \end{pmatrix}.}
If, then, we regard the substitutions $A, B$ as permutations of the indices, the substitution $AB$ has the same meaning as the composed permutation $AB$.

The permutation groups of $n$ figures are accordingly nothing other than a special case of \gesperrt{finite groups} of linear substitutions of dimension $n$.

The permutations of the first kind correspond to substitutions with the determinant $+1$, and the permutations of the second kind to substitutions with the determinant $-1$.


% ===== scan leaf 182, printed 161 =====
% !PRINTED 161
\weberpage{161}
\pdfbookmark[2]{§ 40. The invariants of finite groups of linear substitutions}{bm48}

This follows simply from the fact that a transposition, e.g. $(1, 2)$, corresponds to the substitution
\[
\begin{pmatrix} 0, & 1, & 0, & \dots & 0 \\ 1, & 0, & 0, & \dots & 0 \\ 0, & 0, & 1, & \dots & 0 \\ \cdot & \cdot & \cdot & \cdot & \cdot \\ 0, & 0, & 0, & \dots & 1 \end{pmatrix},
\]
whose determinant is $-1$, and from the fact that all permutations of the first kind can be composed from an even number of transpositions, and all permutations of the second kind from an odd number of transpositions.

\begin{center}
\S.\,40.

\gesperrt{The invariants of finite groups of linear substitutions.}
\end{center}

We now restrict our considerations to \gesperrt{finite groups} of linear homogeneous substitutions. The example of the permutation groups shows that such groups exist. Others we shall come to know later.

We denote such a group by $S$, and its substitutions by $A, B, C, \dots$ If in any function of the variables $x_1, x_2, \dots, x_n$ the variables $(x)$ are replaced by $A(x)$, we say that \gesperrt{the substitution $A$ is applied to the variables $(x)$}. We now first define the following concept:

1. \gesperrt{A form $\Phi(x_1, x_2, \dots, x_n)$ of $n$ variables is called an invariant or invariant form of the group $S$ if it remains unchanged when all the substitutions $A, B, C, \dots$ of the group $S$ are applied to the variables $(x)$.}

As before (Bd.\,I, \S.\,15), a form is here to be understood as an integral homogeneous function of the variables $x$.

We can also express this definition by the formulae:
\weq{1}{\Phi[A(x)] = \Phi[B(x)] = \Phi[C(x)] \dots}

That for every finite group $S$ there are invariant forms in arbitrary abundance is easy to see. One need only
%CONT>>


% ===== scan leaf 183, printed 162 =====
% !PRINTED 162
\weberpage{162}
%<<CONT
to take an arbitrary form $\varphi(x)$ of the $n$ variables $x$,
and to form the functions:
\weq{2}{\varphi[A(x)],\ \varphi[B(x)],\ \varphi[C(x)],\ \ldots}
for all substitutions of the group, and to take any
symmetric function of the forms (2) for $\Phi$.

For if one applies to $(x)$ a substitution of the group $S$,
the totality of the functions (2) does not change; only
their order becomes a different one.

One can conceive the notion of invariants still more generally,
as follows:

2. \gesperrt{A form $\mathit\Psi(x_1,x_2,\ldots,x_n)$ is called an
invariant of the group $S$ also when it takes on
constant factors when the substitutions
$A,B,C \ldots$ of the group $S$ are applied to the
variables $(x)$.}

This property is expressed by formulae as follows:
\weq{3}{\mathit\Psi[A(x)]=\alpha\,\mathit\Psi(x),\ \mathit\Psi[B(x)]=\beta\,\mathit\Psi(x),\ \mathit\Psi[C(x)]=\gamma\,\mathit\Psi(x)\ldots}
in which the coefficients $\alpha,\beta,\gamma\ldots$ are independent of the $x$.

When a distinction is necessary, we shall call the forms
defined in 1. \gesperrt{absolute invariants}, and those defined in 2. \gesperrt{relative
invariants} of the group $S$.

From the formulae (3) there results:
\weq{4}{\mathit\Psi[AB(x)] = \alpha\beta\,\mathit\Psi(x),}
and from this it follows that the factors $\alpha,\beta,\gamma,\ldots$ in their totality,
under composition by multiplication, must form a
group.

Between this (commutative) group and the group $S$
there is in general a multiple isomorphism, since
different substitutions of $S$ can lead to the same factor $\alpha$.

If $\mu$ is the degree of the group $S$, then the degree of each
element $A, B \ldots$ of $S$ is a divisor of $\mu$, and consequently
$A^\mu = B^\mu \cdots$ is equal to the identical substitution.

From this it follows that the factors $\alpha,\beta,\gamma\ldots$ are \gesperrt{$\mu$th roots of unity}.

If $e$ is the least positive number satisfying the condition
\weq{5}{\alpha^e = \beta^e = \gamma^e = \cdots = 1}
then the $\alpha, \beta, \gamma, \ldots$ are at the same time $e$th roots
%CONT>>


% ===== scan leaf 184, printed 163 =====
% !PRINTED 163
\weberpage{163}
%<<CONT
of unity, and $e$ shall be called the \gesperrt{index of the invariant} $\mathit\Psi(x)$.
If $\varepsilon$ is a primitive $e$th root of unity, we can set
\[
\alpha=\varepsilon^a,\quad \beta=\varepsilon^b,\quad \gamma=\varepsilon^c \ldots
\]
and the exponents $a,b,c\ldots$ can have no common
divisor with $e$. From this it follows that the
integers $x,y,z,\ldots$ can be determined so that
\[
ax+by+cz+\cdots \equiv 1\ (\mathrm{mod}\ e)
\]
holds (Bd. I, \S.\,118). Now since, on account of the group nature, among
the factors $\alpha,\beta,\gamma,\ldots$ the number
\[
\alpha^x\,\beta^y\,\gamma^z \cdots = \varepsilon
\]
also occurs, it follows that the totality of the factors $\alpha,\beta,\gamma\ldots$
must coincide with the powers of $\varepsilon$:
\[
1,\ \varepsilon,\ \varepsilon^2,\ \ldots\ \varepsilon^{e-1}
\]

If we seek in $S$ all substitutions $A$ that satisfy the condition
\weq{4}{\mathit\Psi[A(x)] = \mathit\Psi(x)}
to which the identical substitution certainly belongs, we
obtain a new group $T$, which is a divisor of $S$.
If $E$ then denotes a substitution of $S$ satisfying the condition
\weq{5}{\mathit\Psi[E(x)] = \varepsilon\,\mathit\Psi(x)}
then all substitutions $AE$
and $EA$ have the same property, and we obtain the decomposition
of $S$ into the subsidiary groups:
\weq{6}{S = T + TE + TE^2 + \cdots + TE^{e-1};}
the index of the divisor $T$ is therefore $=\varepsilon$. % [sic]
At the same time it follows further, since every substitution $E^{-1}AE$ satisfies condition (4), that
\weq{7}{E^{-1}TE = T,}
from which it appears that $T$ is a normal divisor of $S$. We
have thus proved the theorem:

3. \gesperrt{If $\mathit\Psi(x)$ is an invariant of the group $S$ of
index $e$, then all substitutions of $S$
by which $\mathit\Psi(x)$ remains unchanged form a normal divisor
 $T$ of $S$ of index $e$.}

For the group $T$, $\mathit\Psi(x)$ is an absolute invariant. The invariants
 of index 1 are the absolute invariants of $S$.
The group $T$ is called the \gesperrt{group belonging to the invariant
$\mathit\Psi(x)$.}


% ===== scan leaf 185, printed 164 =====
% !PRINTED 164
\weberpage{164}
If the group $S$ of degree $\mu$ has an invariant of index
 $\mu$, then $T$ becomes the unit group and $S$ is a cyclic
group, which consists of the elements $1, E, E^2 \ldots E^{\mu-1}$.

If we consider as an example the symmetric permutation group
 $P$ of $n$ elements $x_1,x_2,\ldots,x_n$, which by \S.\,39 is indeed contained
among the groups of linear substitutions, then we have
as absolute invariants the \gesperrt{symmetric functions} of the
$n$ variables $x$. The difference product
\[
\sqrt{\varDelta} = (x_1-x_2)(x_1-x_3)\ldots(x_{n-1}-x_n)
\]
is a relative invariant of index 2, to which the alternating
group belongs. Since the group $P$ has no normal divisor
apart from the alternating group and is also not cyclic,
there are no relative invariants of index higher than 2.

The absolute invariants, i.e. the symmetric functions,
are here rationally representable by a finite number
of such forms, namely by the symmetric fundamental functions,
and the invariants of index 2 are the product
of $\sqrt{\varDelta}$ with absolute invariants. We shall prove below that analogous theorems also
hold in the general case.

We close here with the proof of a general
theorem which in all applications is of great use for the formation
of the invariants of a group $S$.

In \S.\,60 of the first volume we defined, for an arbitrary form
of the $n$ variables $F(x_1,x_2,\ldots x_n) = F(x)$, certain forms of the same
 variables $C(x_1,x_2,\ldots x_n) = C(x)$ as covariants,
which were characterised by the property that, when $F(x)$ is transformed
by an arbitrary linear substitution into a new form $F'(y)$,
 the form $C$ satisfies the condition
\[
C'(y) = r^\lambda\,C(x),
\]
where $C'$ depends on the coefficients of $F'$ just as $C$ does on the
coefficients of $F$. Here $r$ denotes the substitution determinant,
and is thus a constant. The coefficients of the
form $F$ occur in $C$ only homogeneously. As an example of a covariant
 we cite the \gesperrt{Hessian determinant}.

If now $F(x)$ is an invariant of the group $S$, and the
$y$ are likewise connected with the $x$ by a substitution from $S$,
then the coefficients of $F$ differ only by
a common constant factor from the coefficients
of $F'$, and the same therefore also holds of the two forms $C$
%CONT>>


% ===== scan leaf 186, printed 165 =====
% !PRINTED 165
\weberpage{165}
\pdfbookmark[2]{§ 41. The theorem of Hilbert}{bm49}
%<<CONT
and $C'$. From this we conclude that $C$ also belongs to the invariants
of the group $S$, and state this as a theorem:

4. \gesperrt{If from an invariant form of the
group $S$ one forms arbitrary covariants, one obtains
new invariant forms of the group.}

\begin{center}
\S.\ 41.

\gesperrt{Hilbert's Theorem.}
\end{center}

The proof of the theorem on the finiteness of the system of invariants
 of a linear substitution group rests on a very
general theorem about systems of forms of any kind whatever, which
\gesperrt{Hilbert} discovered and applied with distinguished
success in manifold investigations on the finiteness
of systems of invariants, and to whose derivation we now wish
to proceed $^{1}$).

I. \gesperrt{If $\mathfrak{S}$ denotes any system of forms in the
$n$ variables $x_1,x_2,\ldots x_n$, in finite or
infinite number, then from $\mathfrak{S}$ a
finite number of forms $F_1,F_2,\ldots,F_\mu$ can be so
selected that every form $F$ of $\mathfrak{S}$ can be represented
by an expression}
\weq{1}{F = A_1F_1 + A_2F_2 + \cdots + A_\mu F_\mu}
\gesperrt{in which $A_1,A_2,\ldots A_\mu$ are
forms of the variables $x_1,x_2,\ldots x_n$.}

The definition of the system of forms $\mathfrak{S}$ must be so complete
that of every single form of the variables $x$ it is decided
whether it belongs to $\mathfrak{S}$ or not, but it is otherwise bound to no
assumption.

If $\mathfrak{S}$ consists of only a finite number of forms, then
our theorem is self-evident, for in this case one can
take all the forms of $\mathfrak{S}$ for $F_1,F_2,\ldots F_\mu$. The
proof will therefore only have to concern itself with the case of an infinite
system $\mathfrak{S}$.

For simplicity of expression we shall call the system of forms
 $F_1,F_2,\ldots F_\mu$ a \gesperrt{basis of the system} $\mathfrak{S}$.
The functions $A_1,A_2,\ldots A_\mu$ must be so constituted that

\bigskip
\noindent\rule{4cm}{0.4pt}

\noindent $^{1}$) \gesperrt{Hilbert}, ``Ueber die Theorie der algebraischen Formen''. Mathematische
 Annalen, Bd. 36 (1890).
%CONT>>


% ===== scan leaf 187, printed 166 =====
% !PRINTED 166
\weberpage{166}
%<<CONT
the $\mu$ products $A_1F_1, A_2F_2, \ldots, A_\mu F_\mu$ are all of the same degree, the degree of $F$. Of course, however, it is in general not required that, conversely, all functions of the form $A_1F_1+A_2F_2+\cdots+A_\mu F_\mu$ for arbitrary $A_i$ belong to the system $\mathfrak S$\footnote{This occurs only for special systems $\mathfrak S$, which are called modules.}.

The theorem that we have to prove is evident when we are dealing with functions of a single variable $x_1$. For then all forms of a system $\mathfrak S$ are powers of the variable $x_1$ with non-negative exponents, multiplied by some constant coefficients. Functions that vanish identically we need not take into account. If we then take for $F_1$ one of these functions of the \gesperrt{lowest possible degree}, then every other function of $\mathfrak S$ can be represented in the form of a product $A_1F_1$, in which $A_1$ is likewise a power of $x_1$ with non-negative exponent and constant coefficient.

In order, then, to arrive at the general proof by the application of complete induction, we first assume that Theorem I. has been proved correct for every system $\mathfrak S_0$ of forms of $n$ variables $x$, and consider first a system $\mathfrak S_r$ of forms $F$ which, besides the $x$, contain also an $(n+1)$st variable $y$, but not in a higher than the $r$th power, where $r$ is any positive integer. Every function $F$ can then be put in one way only into the form:
\weq{2}{F = y^r\varphi + \psi,}
in which the variable $y$ no longer occurs at all in $\varphi$, and occurs in $\psi$ at most up to the $(r-1)$st power.

When $F$ runs through the system $\mathfrak S_r$, then $\varphi$ runs through a certain system $\mathfrak S_0$, which by our assumption can be represented by a basis, let us assume in the form:
\weq{3}{\varphi = a_1\varphi_1 + a_2\varphi_2 + \cdots + a_\mu\varphi_\mu.}
Since now $\varphi_1,\varphi_2,\ldots,\varphi_\mu$ belong to the system $\mathfrak S_0$, there are functions $F_1, F_2,\ldots,F_\mu$ in $\mathfrak S_r$, such that
\weq{4}{F_1=y^r\varphi_1+\psi_1,\ F_2=y^r\varphi_2+\psi_2,\ \ldots,\ F_\mu=y^r\varphi_\mu+\psi_\mu,}
and such that $y$ occurs in $\psi_1,\psi_2,\ldots,\psi_\mu$ at most up to the $(r-1)$st power.


% ===== scan leaf 188, printed 167 =====
% !PRINTED 167
\weberpage{167}
From (2), (3), (4), however, it follows, if we further put
\[ \Psi = \psi - a_1\psi_1 - a_2\psi_2 \cdots - a_\mu\psi_\mu \]
that:
\weq{5}{F = a_1F_1+a_2F_2+\cdots+a_\mu F_\mu+\Psi,}
in which $\Psi$ contains the variable $y$ at most up to the degree $r-1$. All functions $\Psi$ that satisfy the condition (5) form a system $\mathfrak S_{r-1}$.

We now make the further assumption that Theorem I. is proved for functions of every system $\mathfrak S_{r-1}$. Then we can accordingly represent $\Psi$ by a basis in the form:
\weq{6}{\Psi = b_1\Psi_1+b_2\Psi_2+\cdots+b_\nu\Psi_\nu,}
in which $\Psi_1,\Psi_2,\ldots,\Psi_\nu$ are special functions of our system $\mathfrak S_{r-1}$. One can therefore dispose of the coefficients $a_{i,k}$ in such a way that the functions
\weq{7}{\begin{aligned}
F_1' &= a_{1,1}F_1+a_{2,1}F_2+\cdots+a_{\mu,1}F_\mu+\Psi_1\\
F_2' &= a_{1,2}F_1+a_{2,2}F_2+\cdots+a_{\mu,2}F_\mu+\Psi_2\\
&\cdots\cdots\cdots\cdots\cdots\cdots\cdots\cdots\cdots\cdots\cdots\cdots\cdots\cdots\\
F_\nu' &= a_{1,\nu}F_1+a_{2,\nu}F_2+\cdots+a_{\mu,\nu}F_\mu+\Psi_\nu
\end{aligned}}
are contained in $\mathfrak S_r$.

If one multiplies by $b_1,b_2,\ldots b_\nu$ and puts in the sum, according to (5) and (6):
\[ \begin{aligned}
b_1\Psi_1+b_2\Psi_2+\cdots+b_\nu\Psi_\nu &= \Psi\\
&= F - a_1F_1-a_2F_2-\cdots-a_\mu F_\mu,
\end{aligned} \]
then it follows that
\weq{8}{\begin{aligned}
F &= b_1F_1'+b_2F_2'+\cdots+b_\nu F_\nu'\\
&\quad+A_1F_1+A_2F_2+\cdots+A_\mu F_\mu,
\end{aligned}}
in which
\weq{9}{\begin{aligned}
A_1 &= a_1 - a_{1,1}b_1 - \cdots - a_{1,\nu}b_\nu\\
&\cdots\cdots\cdots\cdots\cdots\cdots\cdots\cdots\cdots\cdots\cdots\\
A_\mu &= a_\mu - a_{\mu,1}b_1 - \cdots - a_{\mu,\nu}b_\nu.
\end{aligned}}
The forms $F_1',\ldots F_\nu', F_1,\ldots F_\mu$ belong to the system $\mathfrak S_r$ and form a basis of this system, and Theorem I. is thus proved for every system $\mathfrak S_r$, under the presupposition that it holds for $\mathfrak S_0$.

Now let $\mathfrak S$ be an arbitrary system of forms in $n+1$ variables. We pick out any form $F_0$ of degree $r$ from this system, and put, for the variables on which the system depends,
\weq{10}{y,\ x_1+\lambda_1y,\ \ldots,\ x_n+\lambda_ny,}
%CONT>>


% ===== scan leaf 189, printed 168 =====
% !PRINTED 168
\weberpage{168}
\pdfbookmark[2]{§ 42. Finiteness of the system of invariants of a finite linear substitution group}{bm50}
%<<CONT
wherein $\lambda_1,\ldots\lambda_n$ are constants, over which we dispose in such a way that the coefficient of $y^r$ in $F_0$ does not vanish, i.e.\ that $F_0(1,\lambda_1,\ldots\lambda_n)$ becomes different from zero. Then the system $\mathfrak S$, like every form of the variables (10) generally, passes into a system of forms of the variables $y,x_1,\ldots x_n$, and conversely every form that depends on these variables can also be represented as a form of the linear combinations (10).

Any form $F$ of the system $\mathfrak S$ is now ordered according to powers of $y$, and then the division by $F_0$ is carried out with respect to $y$, whereby
\weq{11}{F = a_0F_0 + \Phi}
may result, so that $a_0$ is the quotient and $\Phi$ the remainder of the division. $a_0$ and $\Phi$ are \gesperrt{integral functions} of the variables $x, y$, because the coefficient of the highest power of $y$ in the divisor $F_0$ is constant. $\Phi$ does not exceed the degree $r-1$ with respect to $y$. Now if $F$ runs through the system $\mathfrak S$, the totality of the functions $\Phi$ defined by (11) forms a system $\mathfrak S_{r-1}$, of which we assume the representability by a basis as already proved. We can therefore set:
\weq{12}{\Phi = A_1\Phi_1 + \cdots + A_\mu\Phi_\mu,}
so that $\Phi_1,\ldots\Phi_\mu$ belong to the system $\mathfrak S_{r-1}$, i.e.\ in such a way that in the system $\mathfrak S$ the forms
\weq{13}{F_1 = a_1F_0+\Phi_1,\ \ldots,\ F_\mu = a_\mu F_0+\Phi_\mu}
can be determined. If we therefore set
\weq{14}{A_0 = a_0 - a_1A_1 - \cdots - a_\mu A_\mu,}
then from (11), (12) and (13) it follows:
\weq{15}{F = A_0F_0+A_1F_1+\cdots+A_\mu F_\mu,}
whereby Theorem I.\ is proved in general.

\begin{center}
\S. 42.\\
\gesperrt{Finiteness of the system of invariants of a finite linear substitution group.}
\end{center}

The proof of Theorem I.\ given in the preceding section is in itself subject to no exceptions whatever. But we lose nothing essential of its generality if we exclude once and for all identically vanishing forms. If
%CONT>>


% ===== scan leaf 190, printed 169 =====
% !PRINTED 169
\weberpage{169}
%<<CONT
further, the system $\mathfrak{S}$ contains forms of degree 0, that is, constants different from zero, then our theorem is self-evident, since we may take such a constant for $F_1$, and, if $A_1 = F : F_1$ is set, $F = A_1F_1$. In this form, however, the theorem is without content. But if we exclude from $\mathfrak{S}$ all constant forms, there remains a system $\mathfrak{S}'$ which no longer contains any forms of degree 0, and for which our theorem likewise holds. The basis $F_1, F_2, \ldots F_\mu$ then likewise contains no forms of degree 0, and we can therefore also express Theorem I. thus:

\gesperrt{II. All forms of the system $\mathfrak{S}$ of positive degree can be expressed by means of a basis $F_1, F_2, \ldots F_\mu$, whose elements are of positive degree, in the form:}
\weq{1}{F = \Phi_1F_1 + \Phi_2F_2 + \cdots + \Phi_\mu F_\mu.}
\gesperrt{The degrees of $\Phi_1, \Phi_2, \ldots \Phi_\mu$ are then lower than the degree of $F$.}

This theorem finds its most important applications in investigations of the possibility of representing all forms of a certain system $\mathfrak{S}$ as \gesperrt{integral rational functions of a finite number} among them. Such a system of forms is called a \gesperrt{finite} one (not in the sense that it consists of only a finite number of forms). The following theorem holds:

\gesperrt{III. If the coefficients $\Phi_1, \Phi_2, \ldots \Phi_\mu$ in the representation (1) of Theorem II. can be so chosen for every form $F$ in $\mathfrak{S}$ that, when they are not constant, they themselves belong to the system $\mathfrak{S}$, then the system $\mathfrak{S}$ is finite.}

This follows immediately from the fact that the degrees of the forms $\Phi_1, \Phi_2, \ldots \Phi_\mu$ are lower than the degree of $F$. If, then, one applies the representation (1) to the non-constant ones among the functions $\Phi$, one arrives at coefficients of still lower degree, and so, on repeated application of this procedure, one must finally come to constants.

From this there now follows, by a very simple mode of reasoning, which I owe to an oral communication from \gesperrt{Hurwitz}, the finiteness of the system of invariants $\mathfrak{J}$ of a finite group of linear substitutions $S$.


% ===== scan leaf 191, printed 170 =====
% !PRINTED 170
\weberpage{170}
Let $F_1, F_2, \ldots, F_\mu$ be a basis of the system $\mathfrak{J}$ determined according to II., and let $F$ be any other non-constant invariant of $S$. Then the forms $A_1, A_2, \ldots, A_\mu$ can be so determined that
\weq{2}{F = A_1F_1 + A_2F_2 + \cdots + A_\mu F_\mu}
holds. If one applies to this identical equation all the substitutions of the group $S$, then by hypothesis $F, F_1, F_2, \ldots F_\mu$ remain unchanged, while $A_\varkappa$ may pass over into $A_\varkappa, A_\varkappa', A_\varkappa'' \ldots$ If one forms the sum of the equations thus derived from (2) and sets, where $m$ denotes the degree of the group $S$,
\weq{3}{m\Phi_\varkappa = A_\varkappa + A_\varkappa' + A_\varkappa'' + \cdots,}
then it follows that:
\weq{4}{F = \Phi_1F_1 + \Phi_2F_2 + \cdots + \Phi_\mu F_\mu.}
But the functions $\Phi_1, \Phi_2, \ldots \Phi_\mu$ are by \S.\,40 invariants of $S$, and with this, according to III., the finiteness of the system $\mathfrak{J}$ is proved.

The same conclusion can also be applied to the relative invariants, as follows:

We denote by $F(x)$ the whole system of the functions which satisfy the conditions \S.\,40, (3):
\[
F[A(x)] = \alpha F(x), \quad F[B(x)] = \beta F(x) \ldots
\]
with fixed factors $\alpha, \beta, \ldots$, which, as we have seen, are roots of unity.

By \gesperrt{Hilbert's} theorem a special system of such functions $F_1, F_2, \ldots, F_m$ can be selected in such a way that one can set
\weq{5}{F = A_1F_1 + A_2F_2 + \cdots + A_mF_m}
where $A_1, A_2, \ldots, A_m$ are forms of the variables $(x)$. If one treats this formula just as the formula (2), by applying to it the substitutions of the group $S$ and then forming the sum, one obtains, corresponding to the formula (4):
\[
F = \Phi_1F_1 + \Phi_2F_2 + \cdots + \Phi_mF_m,
\]
in which the coefficients $\Phi_1, \Phi_2, \ldots, \Phi_m$ are \gesperrt{absolute invariants}.

By way of completion it must still be added that inhomogeneous functions of the variables can be invariants only when their individual homogeneous parts are invariants.


% ===== scan leaf 192, printed 171 =====
% !PRINTED 171
\weberpage{171}
\pdfbookmark[2]{§ 43. The form problem}{bm51}
Finally we may also inquire after fractional invariants. If
\[
\frac{F(x)}{F_1(x)}
\]
is such a fractional invariant, we first assume that the two integral rational functions $F(x), F_1(x)$ of the $n$ variables $x$ are freed of common divisors (Bd.\,I, \S.\,51).

If we restrict ourselves for the present to absolute invariants, and if accordingly
\[
\frac{F(x)}{F_1(x)} = \frac{F[A(x)]}{F_1[A(x)]},
\]
then, since the $x$ are here regarded as independent variables, numerator and denominator have no common divisor either on the right or on the left, and it follows that:
\[
F[A(x)] = \alpha F(x), \quad F_1[A(x)] = \alpha F_1(x),
\]
in which $\alpha$ is a constant factor which, as we have seen earlier, is a root of unity. Numerator and denominator of a fractional invariant must therefore themselves be invariants, if only relative ones.

But if we are not seeking precisely the simplest representation, we can also represent absolute fractional invariants as quotients of absolute integral invariants. We need only extend the fraction by a suitable power of the denominator, that is, if the $\alpha$ are $e$th roots of unity, to put
\[
\frac{F(x)}{F_1(x)} = \frac{F(x) F_1(x)^{e-1}}{[F_1(x)]^e}
\]
Accordingly we can also represent all relative invariants with a definite system of factors $\alpha, \beta, \ldots$ as the product of one of them with absolute invariants, which however may be fractional.

\begin{center}
\S.\,43.

\gesperrt{The Form Problem.}
\end{center}

In the preceding section it was proved that for a finite group of linear substitutions there is a finite number of independent invariants. If $n$ is the dimension of the linear
%CONT>>


% ===== scan leaf 193, printed 172 =====
% !PRINTED 172
\weberpage{172}
%<<CONT
substitution, these forms are homogeneous functions of
$n$ variables, and hence no more than $n$ mutually
independent ones can exist. This does not mean that all
these forms can be expressed \gesperrt{rationally} in terms of $n$ among them,
but between $n+1$ invariants there must always subsist a rational
equation, which results by elimination of the $n$ variables.

But there remains the converse question to investigate, whether
for a linear substitution group of dimension
$n$ there really always exist $n$ independent invariants.

If we put fixed values for the variables, then all the
invariants of the group thereby likewise receive
determinate values. The question that we still have to answer
is now this: whether to a system of values of the invariants
there exist determinate systems of values of the variables, and
indeed in finite number. Once this is proved, we can regard the
variables as (many-valued) algebraic functions of the invariants,
and the number of mutually independent
invariants cannot be smaller than the number
of the variables, because otherwise a part of the variables, when the
values of the invariants are given, would still remain
arbitrary.

If $(x)$ is a system of values of the variables corresponding
to a determinate system of values of the invariants, and $A$ a substitution
of the group $S$, then by the nature of the invariants
the system $A(x)$ corresponds to the same system of values of the invariants,
and our problem therefore has at least as many
solutions as the degree of the group amounts to. That this, however, is the
exact number of solutions appears from the following
considerations.

In special cases, i.e. for special values of the invariants,
several of these systems of values of the variables
may coincide. But this can happen only for such systems of values
of the $(x)$ for which a relation of the form
$(x) = A(x)$ subsists, for from $A(x) = B(x)$ there would follow $(x) = A^{-1}B(x)$,
which is contained in the form $(x) = A(x)$.

We assume a linear homogeneous function $\Theta$ of the variables
$x_1, x_2, \dots x_n$, which we denote thus:
\weq{1}{\Theta = \Theta(x) = c_1 x_1 + c_2 x_2 + \cdots + c_n x_n.}


% ===== scan leaf 194, printed 173 =====
% !PRINTED 173
\weberpage{173}
If $A$ denotes a substitution of the finite group $S$, we
can derive from $\Theta(x)$ a new function:
\weq{2}{\Theta[A(x)] = c_1' x_1 + c_2' x_2 + \cdots + c_n x_n} % [sic]
whose coefficients $c_i'$ are connected, by \S.\,37, 9., with the original
coefficients $c_i$ through the \gesperrt{transposed substitution corresponding to $A$},
\[ (c') = A_1(c) \]

If now $A, B, C, \dots$ are all the substitutions of the
group $S$, one can form in the same way the functions
\weq{3}{\Theta[A(x)],\ \Theta[B(x)],\ \Theta[C(x)],\ \dots}
which we also denote more briefly by
\weq{4}{\Theta,\ \Theta_1,\ \Theta_2,\ \dots,\ \Theta_{\mu-1}}
where $\mu$ is the degree of the group $S$.

Now one can dispose of the coefficients $c_i$, which up to now
are still entirely arbitrary, in such a way that no two of the functions
(4) become identical with one another (Bd.\ I, \S.\,143, 1.).
But from (3) it follows:

\gesperrt{1. If one applies simultaneously in the $\mu$ functions
(4) any substitution from $S$, the
totality of these functions does not change,
but they suffer only a permutation.}

From this it follows further:

\gesperrt{2. Every symmetric function of the quantities (4) is
an absolute invariant of the group $S$.}

We remark further that in place of the function $\Theta$
one could put any other function, even a non-linear, indeed a fractional
or inhomogeneous one, provided only that the $\mu$ functions
(4) are different from one another.

If we now form the product:
\weq{5}{\begin{aligned}
\Phi(t) &= (t-\Theta)(t-\Theta_1)\dots(t-\Theta_{\mu-1}) \\
&= t^\mu + A_1 t^{\mu-1} + A_2 t^{\mu-2} + \cdots + A_\mu,
\end{aligned}}
which is an integral rational function of degree $\mu$ in $t$, then
the coefficients $A_1, A_2, \dots, A_\mu$ are, by 2., invariants of the
group $S$, and the quantities $\Theta, \Theta_1, \dots, \Theta_{\mu-1}$ are the roots
of the equation
\weq{6}{\Phi(t) = 0.}


% ===== scan leaf 195, printed 174 =====
% !PRINTED 174
\weberpage{174}
We now take as domain of rationality $\Omega$ the field which
consists of the absolute invariants of the group $S$ and all
numbers\footnote{One may also restrict oneself to a particular number field, provided only that the substitution coefficients of the group $S$ are contained in it.}. To this domain of rationality belong the coefficients
of $\Phi(t)$, and we first prove the theorem:

\gesperrt{3. The function $\Phi(t)$ is irreducible in $\Omega$.}

For if $\Psi(t)$ is any function in $\Omega$ which vanishes for $t = \Theta$,
then in the equation $\Psi(\Theta) = 0$, since
the $x_1, x_2, \dots, x_n$ are independent variables, we may replace the system $(x)$
of these variables by $A(x)$, where $A$ is any substitution
of $S$. In this way $\Theta$ can be carried over into each of the functions
$\Theta_1, \Theta_2, \dots, \Theta_{\mu-1}$, while the coefficients
of $\Psi$ remain unchanged, and consequently $\Psi(\Theta_1) = 0$,
$\Psi(\Theta_2) = 0, \dots \Psi(\Theta_{\mu-1}) = 0$. Hence $\Psi(t)$ must be divisible by $\Phi(t)$,
whereby the irreducibility is proved.

Now let $\omega$ denote any function of the variables $x$, and let
\[ \omega,\ \omega_1,\ \omega_2,\ \dots,\ \omega_{\mu-1} \]
be the functions which arise from $\omega$ by application of the
substitutions of $S$; these functions need not all be different from one another.
If $t$ is a variable, then
\weq{7}{\Phi(t)\left(\frac{\omega}{t-\Theta} + \frac{\omega_1}{t-\Theta_1} + \cdots + \frac{\omega_{\mu-1}}{t-\Theta_{\mu-1}}\right) = \Psi(t)}
is an integral rational function of degree $(\mu-1)$ in $t$, and at the same time
it is an invariant of $S$, hence a function in $\Omega$.
If we set $t = \Theta$ in it, then by an inference already often
employed before (Bd.\ I, \S.\,143, 155) it follows that:
\weq{8}{\omega = \frac{\Psi(\Theta)}{\Phi'(\Theta)},}
in which the following theorem is contained:

\gesperrt{4. Every rational function of the variables $(x)$ can
be expressed rationally in terms of $\Theta$, and hence
belongs to the field $\Omega(\Theta)$.}

From this theorem we can derive a second proof
that all invariants of the group $S$ can be represented rationally by a
finite number of them, namely the coefficients of the function
$\Phi(t)$. For if we take an absolute in-
%CONT>>


% ===== scan leaf 196, printed 175 =====
% !PRINTED 175
\weberpage{175}
%<<CONT
variant $J$, by Theorem 4., as a rational function of $\Theta$,
then this function cannot change when one of the substitutions
$(\Theta, \Theta_1), (\Theta, \Theta_2), \dots$ is carried out, and this function
is therefore rationally expressible in terms of the coefficients of $\Phi(t)$.
Whether this representation is indeed possible by integral functions
would remain undecided in this proof.

Among the functions $\omega$ of the variables $x$ the
variables $x$ themselves are also contained, and the question from which we
set out, whether the variables $x$ can be regarded as algebraic functions
of the invariants, is thereby decided in the
affirmative.

The problem of representing the variables $x$ as algebraic functions
of the invariants of the group $S$, that is, the determination
of the field $\Omega(\Theta)$, is called, after \gesperrt{F. Klein}, the \gesperrt{form problem
of the group $S$}\footnote{Vorlesungen über das Ikosaëder, S.\,123. (Leipzig 1884.)}. If the number of variables is $n$, we call
it the form problem of the $n$th dimension.

The field $\Omega(\Theta)$ is an algebraic field over $\Omega$ completely determined
by the group $S$. It is a normal field,
for by 4. the conjugate quantities $\Theta, \Theta_1, \Theta_2, \dots \Theta_{\mu-1}$
are all contained in the field $\Omega(\Theta)$ itself. The equation $\Phi(t) = 0$
is a normal equation and is the \gesperrt{Galois resolvent
of the form problem} (Bd.\ I, \S.\,145).

\gesperrt{5. The Galois group of the form problem,
i.\,e.\ the Galois group of the equation
$\Phi(t) = 0$, is isomorphic with the group $S$.}

To prove this, we denote by $A, B$ two substitutions
from $S$ and by $AB = C$ the substitution composed
from them. Now if
\[ \Theta_1 = \Theta[A(x)],\ \Theta_2 = \Theta[B(x)],\ \Theta_3 = \Theta[C(x)], \]
then the substitution
\[ (\Theta, \Theta_2) = (\Theta_1, \Theta_3), \]
and consequently
\[ (\Theta, \Theta_1)(\Theta, \Theta_2) = (\Theta, \Theta_1)(\Theta_1, \Theta_3) = (\Theta, \Theta_3), \]
i.\,e.\ the group of the substitutions $(\Theta, \Theta_1), (\Theta, \Theta_2), \dots$ is
isomorphic with the group of the $A, B, \dots$.

If instead of the function $\Theta$ we take a function $\eta$ which does not
have all its values distinct, but admits the substitutions of a
%CONT>>


% ===== scan leaf 197, printed 176 =====
% !PRINTED 176
\weberpage{176}
\pdfbookmark[2]{§ 44. Klein's extension of the fundamental algebraic problem}{bm52}
%<<CONT
divisor $S'$ of $S$ of index $j$, but no others, then
$\eta$ satisfies an equation of degree $j$, which is to be regarded as a resolvent
of the form problem. Every function which likewise admits the
permutations of $S'$ is then a rational
function of $\eta$ and of the invariants of the group $S$. The
resolvent of the $\eta$ is a partial or total resolvent, according
as the group $S'$ has a common divisor with the divisors
of $S$ conjugate to it, or is relatively prime
to them (Bd. I, \S.\,156).

\begin{center}
\S.\ 44.

\gesperrt{Klein's Extension of the Fundamental Algebraic Problem.}
\end{center}

The considerations carried out in the preceding section
form a direct generalisation of the \gesperrt{Galois} theory
for a general equation of degree $n$; and this
theory is contained as a special case in the theory of linear substitution groups.

For according to \S.\ 39 the symmetric functions
of $n$ independent variables $x_1, x_2, \ldots, x_n$ are the invariants of the
group $S$ which consists of the permutations of these $n$ variables,
and the equation $\Phi(t) = 0$ is thus for this case, by Bd. I,
\S.\ 145, the \gesperrt{Galois} resolvent of the equation of degree $n$,
whose roots are the quantities $x_i$. We can therefore conceive the general problem of algebra, to solve an equation of degree $n$, as a form problem of a linear substitution group
of dimension $n$. But now there are special equations which can be solved by form problems of lower dimension,
in particular the \gesperrt{pure equations}
which are solvable by a form problem of the first dimension.

For linear homogeneous substitutions of one dimension are
only of the form
\weq{1}{x' = \alpha x,}
and if these are to form a finite group of degree $\mu$,
then the coefficients $\alpha$ must be roots of unity of degree $\mu$.
Conversely, if we let $\alpha$ in (1) run through all the $\mu$th roots of unity,
we have a group of degree $\mu$. This
group has an absolute invariant $x^{\mu}$, and when this is given,
one obtains $x$ as the $\mu$th root of it.


% ===== scan leaf 198, printed 177 =====
% !PRINTED 177
\weberpage{177}
\pdfbookmark[2]{§ 45. Influence of relative invariants}{bm53}
The solution of the general equations of degrees $2$, $3$ and
$4$ is thus reducible to form problems of only one dimension. We shall see in a later chapter
that the general equation of degree $5$ is reducible to a binary
form problem, and one may now, as an
immediate extension of the problem of reducing the solution of an
equation to pure equations, pose the question:
what is the least number of dimensions of a form problem by which a given equation can be solved?
The problem would then have to be formulated thus:

\gesperrt{From the roots of a given equation, rational functions are to be formed in the smallest possible number, in such a way that they pass into homogeneous linear functions of themselves when the roots are subjected to the permutations of the Galois group of the given equation.}

In this way F. \gesperrt{Klein} has extended the problem of the algebraic
solution of an equation. To open the way to an answer to the question, he has proved for the general equation
of degrees $6$ and $7$ that they can be reduced to quaternary form problems.

The general equation of degree $n$ is, as we have seen,
immediately equivalent to a form problem of $n$ dimensions.
But the question whether the general equation of degree $n$,
when $n$ is greater than 7, corresponds to a form problem of
fewer than $n$ dimensions has not yet been answered; it must probably be answered in the negative $^{1})$.

\begingroup
\renewcommand{\thefootnote}{1)}
\footnotetext[1]{F. \gesperrt{Klein}, Zur Theorie der allgemeinen Gleichungen $6^{\mathrm{ten}}$ und
$7^{\mathrm{ten}}$ Grades; Mathematische Annalen, Bd. XXVIII, S. 18. Cf. also ``The
Evanston Colloquium, Lectures on Mathematics by \gesperrt{Felix Klein}'', Lecture IX, London und New York, Macmillan and Co. (1894).}
\endgroup

\begin{center}
\S.\ 45.

\gesperrt{Influence of relative invariants.}
\end{center}

In the definition of the form problem in \S.\ 43 we have
used only the absolute invariants of the group $S$. Now there
are, as we have seen, also cases in which, besides the ab-
%CONT>>


% ===== scan leaf 199, printed 178 =====
% !PRINTED 178
\weberpage{178}
%<<CONT
solute ones there also exist relative invariants, and these can be
used for the simplification of the form problem.

If $r$ is such a relative invariant, then a certain power
of $r$, whose degree $e$ is a divisor of the degree $\mu$ of $S$,
will be an absolute invariant, and $r$ is thus determined by a \gesperrt{pure
equation} in $\Omega$.

All the substitutions of $S$ by which $r$ remains unchanged
form by themselves a group $T_r$, and the group $T_r$ is, as we
have seen in \S.\,40, a normal divisor of $S$.

Further, however, we see that every function of the variables $x$
that remains unchanged under the substitutions of the group $T_r$
can be expressed rationally in terms of $r$, that is, is contained in the
field $\Omega_r$ derived from $\Omega$ by adjunction of $r$.

For according to \S.\,40 there is a substitution $E$ in $S$, and
a primitive $e$th root of unity $\varepsilon$, such that $r$ passes under $E$ into $\varepsilon r$,
and all the values that $r$ can assume:
\[
r, \varepsilon r, \varepsilon^2 r, \ldots \varepsilon^{e-1} r,
\]
are obtained by repetition of the substitution $E$. If, further,
$\tau$ denotes a function of the $x$ that admits the substitutions of the group
$T_r$, and passes under $E$ and its powers into
\[
\tau, \tau_1, \tau_2, \ldots \tau_{e-1}
\]
then all these functions likewise admit the
substitutions of $T_r$, because $T_r$ is a normal divisor of $S$. The
function
\weq{1}{(\varepsilon^{-h}, \tau) = \tau + \varepsilon^{-h}\tau_1 + \cdots + \varepsilon^{-h(e-1)}\tau_{e-1}, \qquad (h = 0, 1, 2, \ldots e-1),}
which is formed on the analogy of the \gesperrt{Lagrange} resolvents (Bd.\,I,
\S.\,164), takes on the factor $\varepsilon^h$ upon application of the substitution $E$,
and if we therefore set
\weq{2}{(\varepsilon^{-h}, \tau) = r^h \Psi_h
}
then $\Psi_h$ is an absolute invariant, and hence contained in the
field $\Omega$.

From (1) and (2), however, it follows that
\weq{3}{e\tau = \Psi + r\Psi_1 + r^2\Psi_2 + \cdots + r^{e-1}\Psi_{e-1},}
whereby the assertion is proved.

If we set
\weq{4}{\mu = e\nu,}
%CONT>>


% ===== scan leaf 200, printed 179 =====
% !PRINTED 179
\weberpage{179}
\pdfbookmark[2]{§ 46. The extended concept of invariant}{bm54}
%<<CONT
then $\nu$ is the degree of the group $T_r$, and if the function $\Theta$, which
we applied in \S.\ 43 to the solution of the form problem,
passes, under the substitutions of $T_r$, into
\[ \Theta,\ \Theta_1,\ \ldots \Theta_{\nu-1} \]
then
\[ \Phi_r(t) = (t - \Theta)(t - \Theta_1) \ldots (t - \Theta_{\nu-1}) \]
is a function of $t$ in the extended domain of rationality
$\Omega_r$; the field $\Omega(\Theta)$ is identical with $\Omega_r(\Theta)$. It is therefore an algebraic field of degree $\mu$ over $\Omega$ and of degree $\nu$ over $\Omega_r$.
The form problem of degree $\mu$ is accordingly reduced, by the adjunction
of a radical, to a form problem of degree $\nu$.

\begin{center}
\S.\ 46.

\gesperrt{The extended concept of invariant.}
\end{center}

The relative invariants are at bottom a special case
of a more general concept, which we can apply to the reduction
of the form problem in a similar way to the relative
invariants.

We seek systems of homogeneous forms of equal
degree in the variables $x_1, x_2 \ldots x_n$:
\weq{1}{X_1, X_2, \ldots X_m,}
of such a nature that, under the applications of the substitutions of the group $S$, the functions $X_1, X_2, \ldots X_m$ undergo a
system $\Sigma$ of \gesperrt{linear substitutions}; the substitutions $\Sigma$ then likewise form a group, and indeed
a group which is simply or multiply isomorphic with $S$. If
$m = 1$, we obtain, as one sees, the concept of the relative
invariants. We now seek all the substitutions of $S$ which
leave each single one of the functions (1) unchanged. These substitutions form a group, which we denote by $T$, and
of which we wish to show that it is a normal divisor of
$S$. If we denote by $s$ a substitution of $S$ under which
the $X$ undergo the substitution $\sigma$, then under $s^{-1}$ the $X$ undergo
the substitution $\sigma^{-1}$. If further $\tau$ is a substitution of $T$,
then under $\tau$ the $X$ remain unchanged. Accordingly, under
$s^{-1}\tau s$ the $X$ undergo the substitution $\sigma^{-1}\sigma = 1$, i.\,e. they remain un-
%CONT>>


% ===== scan leaf 201, printed 180 =====
% !PRINTED 180
\weberpage{180}

%% [running head, not body text -- commented out 2026-08-07] Page furniture (running
%% head, folio, signature) is excluded from this edition by standing policy.
%% \noindent 180 \hfill \textit{Sechster Abschnitt.} \hfill \S.\,46.

%<<CONT
changed. $s^{-1}\tau s$ is therefore also contained in $T$, and $T$ is consequently a normal divisor of $S$.

We denote by $\mu, \nu$ the degrees of $S$ and $T$ and set
\[
\mu = e\nu;
\]
then $e$ is the index of $T$ and at the same time the degree of the group $\Sigma$.

Every absolute invariant of the group $T$, i.\,e. every function of $x$ that admits the substitutions of $T$, can be expressed rationally in $\Omega$ in terms of the $X$.

This follows by the method of inference already often applied: if we denote by $\varrho$ a function of the $X$ that takes $e$ distinct values $\varrho, \varrho_1, \ldots \varrho_{e-1}$, and set
\weq{2}{\Phi(t) = (t-\varrho)(t-\varrho_1)\ldots(t-\varrho_{e-1})}
then $\Phi(t)$ is a rational function of the variable $t$ in $\Omega$.

An absolute invariant $r$ of $T$ takes at most $e$ distinct values, which correspond to the values $\varrho, \varrho_1, \ldots \varrho_{e-1}$, namely $r, r_1, \ldots r_{e-1}$, where in certain circumstances equal values may also occur. The function of $t$:
\[
\frac{\Phi(t)\,r}{t-\varrho} + \frac{\Phi(t)\,r_1}{t-\varrho} + \cdots + \frac{\Phi(t)\,r_{e-1}}{t-\varrho_{e-1}} = \Psi(t) % [sic] second denominator printed t-\varrho, not t-\varrho_1
\]
is then contained in $\Omega$, and for $t=\varrho$ we obtain:
\weq{3}{r = \frac{\Psi(\varrho)}{\Phi'(\varrho)},}
whereby, since $\Phi'(\varrho)$ is different from zero, $r$ is rationally expressed by $\varrho$. Thus $r$ is contained in the field $\Omega(X_1, X_2, \ldots X_m)$.

The invariants of the group $T$ are at the same time invariants of the group $S$; by the solution of the form problem for the group $\Sigma$ the $X_1, \ldots X_m$ are then known, and hence also the invariants of the group $T$, and the form problem of the group $S$ is thus reduced to the successive solution of the two form problems of the groups $\Sigma$ and $T$, which are of lower degree than the form problem for $S$. A form problem which is decomposable in this way into two form problems of lower degree we may call an \gesperrt{imprimitive form problem}. Since for such a reduction it is necessary that $T$ be a normal divisor of $S$ different from $S$ and from the unit, the corresponding form problem, when $S$ is a simple group, is always primitive.


% ===== scan leaf 202, printed 181 =====
% !PRINTED 181
\weberpage{181}
\pdfbookmark[2]{§ 47. Normal forms}{bm55}

%% [running head, not body text -- commented out 2026-08-07] Page furniture (running
%% head, folio, signature) is excluded from this edition by standing policy.
%% \noindent \S.\,47. \hfill \textit{Normalformen.} \hfill 181

But we can also conclude conversely that, if $S$ possesses a normal divisor of index $e$, the form problem is imprimitive. For in order to obtain a system of the functions $X_1, \ldots X_m$, we need only take an invariant $\varrho$ of the group $T$ which contains $e$ different values in $S$, and we can take for $X_1, \ldots X_m$ simply the values $\varrho, \varrho_1, \ldots \varrho_{e-1}$. The group $\Sigma$ is then the permutation group called forth by $S$ among the $\varrho$.

In the sense of \S.\,44, however, it will always be a matter of making $m$ as small as possible, and a reduction of the problem in this sense will be achieved only when $m < n$.

\begin{center}
\S.\,47.

\gesperrt{Normal Forms.}
\end{center}

We must undertake one more general consideration before we pass on to special applications. We have already seen in \S.\,37 that from every group $S$ of linear substitutions we can derive infinitely many isomorphic groups, by transforming with an arbitrary substitution $L$ of the same dimension, that is, by forming the transformed group
\weq{1}{L^{-1} S L}
and this is equivalent to the introduction of other variables $y$ in place of $x$ by the substitution
\weq{2}{(x) = L(y).}

We can use this transformation to represent the group $S$ in a simple normal form, and then the invariants also receive certain fixed normal forms, which in many cases can assume very simple shapes. If we form for the normal form the resolvent of the form problem [\S.\,43, (8)], which we now wish to write in the form:
\weq{3}{\Phi(\Theta, A_1, A_2, \ldots) = 0,}
wherein $A_1, A_2, \ldots$ denote invariants of the group $S$, then this too, if we use the normal form, will receive a simple shape. The variables $x_i$ can, as we have seen
%CONT>>


% ===== scan leaf 203, printed 182 =====
% !PRINTED 182
\weberpage{182}

%% [running head, not body text -- commented out 2026-08-07] Page furniture (running
%% head, folio, signature) is excluded from this edition by standing policy.
%% \noindent 182 \hfill \textit{Sechster Abschnitt.} \hfill \S.\,47.

%<<CONT
rationally in terms of $\Theta$ and the $A_k$, and we set
\weq{4}{x_i = \varphi_i(\Theta, A_1, A_2, \ldots).}

These expressions too, once the function $\Theta$ has been decided upon, will assume fixed shapes for the normal form.

But now the case may also occur that the invariants are given not in the normal form, but in some arbitrary other form, which we shall call the general form. Then the task will be to find the substitution $L$ by which the general form is transformed into a normal form recognised from the outset as possible. This problem, as we shall now see, is none other than the form problem posed for the normal form itself.

Let us assume that the invariants in the general form, as functions of $y$, are $B_1, B_2, \ldots$, so that by the substitution (2) the identities
\weq{5}{A_1 = B_1, \ A_2 = B_2, \ldots}
are established. The task is, when the functions $A_i, B_i$ are given as to their form, to determine the substitution $L$ which makes the equations (5) into identities. This problem is solved if we assume equation (3) to be solved for a suitably chosen special system of values of the $A_i$. To prove this, we form the complete differentials of equations (3) and (4):
\[
0 = \Phi'(\Theta)\,d\Theta + \sum^s \Phi'(A_s)\,dA_s
\]
\[
dx_i = \varphi_i'(\Theta)\,d\Theta + \sum^s \varphi_i'(A_s)\,dA_s,
\]
in which $\Phi'(\Theta)$, $\Phi'(A_s)$, $\varphi_i'(\Theta)$, $\varphi_i'(A_s)$ denote the partial derivatives. If we eliminate $d\Theta$, it follows that:
\weq{6}{dx_i = \sum^s \frac{\varphi_i'(A_s)\,\Phi'(\Theta) - \Phi'(A_s)\,\varphi_i'(\Theta)}{\Phi'(\Theta)}\, dA_s.}

Now, in consequence of the equations (5):
\weq{7}{dA_s = \sum^k B_s'(y_k)\, dy_k,}
and if we therefore
\weq{8}{x_i = \alpha_{1,i} y_1 + \cdots + \alpha_{n,i} y_n}
\weq{9}{dx_i = \alpha_{1,i}\, dy_1 + \cdots + \alpha_{n,i}\, dy_n}
%CONT>>


% ===== scan leaf 204, printed 183 =====
% !PRINTED 183
\weberpage{183}

%% [running head, not body text -- commented out 2026-08-07] Page furniture (running
%% head, folio, signature) is excluded from this edition by standing policy.
%% \noindent \S.\,47. \hfill \textit{Normalformen.} \hfill 183

%<<CONT
then the comparison of (6) with (9) yields:
\weq{10}{\alpha_{k,i} = \sum^s \frac{\varphi_i'(A_s)\,\Phi'(\Theta) - \Phi'(A_s)\,\varphi_i'(\Theta)}{\Phi'(\Theta)}\, B_s'(y_k).}

The right-hand side of these equations must therefore reduce to a constant, and we can find its value if we set for the $y$ any special system of values, which is bound only by the one condition that $\Phi'(\Theta)$ does not vanish. For this special system of values the values of the $A_i$ are determined by the equations (5), and $\Theta$ is known, if for this special system of values of the $A_i$ we assume the form problem (3) to be solved. Then by (10) the coefficients $\alpha_{k,i}$, and with them the substitution $L$, are completely determined.

\noindent\hrulefill


% ===== scan leaf 205, printed 184 =====
\weberpage{184}% page number not printed on this page
\pdfbookmark[1]{Seventh Section. Groups of Binary Linear Substitutions.}{bm56}
\pdfbookmark[2]{§ 48. Ternary orthogonal substitutions}{bm57}
% !PRINTED none

\begin{center}
\gesperrt{Seventh Section.}

Groups of binary linear substitutions.

\rule{4cm}{0.4pt}

\S.\,48.

\gesperrt{Ternary orthogonal substitutions.}
\end{center}

It is now our task to separate out narrower groups from the
totality of the linear substitutions, and finally to arrive at
finite groups. Such narrower groups, which may still be infinite,
are obtained if one makes the requirement that given homogeneous
functions of the variables shall remain invariant. We shall not,
however, pursue the problem further in this generality, but pass
at once to the consideration of the most important special case.
We shall restrict ourselves to ternary substitutions and require
that a quadratic form of non-vanishing determinant shall remain
invariant. Since, as we have seen earlier (Bd. I, \S.\,57), every
such quadratic form can be transformed by linear transformation
into a sum of squares, we do not restrict the problem further if
we take for this quadratic form the sum of the squares. Such
substitutions are called \gesperrt{orthogonal}.

The substitution

\weq{1}{(y_1, y_2, y_3) = A (x_1, x_2, x_3)}

is therefore orthogonal if the substitution coefficients are so determined that the identity holds:

\weq{2}{y_1^2 + y_2^2 + y_3^2 = x_1^2 + x_2^2 + x_3^2.}

If one substitutes the expressions (1) in (2) and sets the
coefficients of corresponding terms equal to one another, then
%CONT>>


% ===== scan leaf 206, printed 185 =====
% !PRINTED 185
\weberpage{185}
%<<CONT
one obtains six relations among the nine coefficients
of $A$.

These relations read, if

\weq{3}{A = \begin{pmatrix} a_1, & a_2, & a_3\\ b_1, & b_2, & b_3\\ c_1, & c_2, & c_3 \end{pmatrix}}

is assumed:

\weq{4}{\begin{array}{r@{}c@{}l@{\qquad}r@{}c@{}l}
a_1^2+b_1^2+c_1^2 &{}={}& 1, & a_2a_3+b_2b_3+c_2c_3 &{}={}& 0\\
a_2^2+b_2^2+c_2^2 &{}={}& 1, & a_3a_1+b_3b_1+c_3c_1 &{}={}& 0\\
a_3^2+b_3^2+c_3^2 &{}={}& 1, & a_1a_2+b_1b_2+c_1c_2 &{}={}& 0,
\end{array}}

and are well known from analytic geometry. For the square
of the substitution determinant $|A|$ these relations yield, by
the multiplication rule of determinants, the value 1, and consequently
$|A|$ has the value $\pm 1$.

From the formulae (4) it follows that the \gesperrt{inverse substitution} of $A$

\[
A^{-1} = \begin{pmatrix} a_1, & b_1, & c_1\\ a_2, & b_2, & c_2\\ a_3, & b_3, & c_3 \end{pmatrix}
\]

is identical with the transposed one, and this property could
also serve as a definition of the orthogonal substitutions.

The totality of the orthogonal substitutions forms
a group. Among them there is contained a narrower group,
distinguished by the value

\weq{5}{|A| = + 1}

which we shall designate as the group of the \gesperrt{proper
orthogonal substitutions}.

One may interpret the linear substitution (1) as the transition
from one rectangular coordinate system to a second
with the same origin, if one regards $x_1, x_2, x_3$ and
$y_1, y_2, y_3$ as rectangular coordinates of one and the same
(variable) point in the first and in the second coordinate system. By $A$ the mutual position of the
two coordinate systems is determined, and conversely. If
the condition (5) holds, as we shall now assume, then
the first coordinate system can be brought into coincidence with the second by rotation about a fixed axis through a
definite angle of rotation.


% ===== scan leaf 207, printed 186 =====
% !PRINTED 186
\weberpage{186}
For if we set

\[
D = \begin{vmatrix} a_1 - 1, & a_2, & a_3\\ b_1, & b_2 - 1, & b_3\\ c_1, & c_2, & c_3 - 1 \end{vmatrix},
\]

then from the relations (4) we obtain:

\[
|A|\, D = - D,
\]

and thus, if (5) holds, $D = 0$. Accordingly, three quantities
$\lambda, \mu, \nu$ can be determined from the equations:

\weq{6}{\begin{array}{l}
\lambda = a_1 \lambda + a_2 \mu + a_3 \nu\\
\mu = b_1 \lambda + b_2 \mu + b_3 \nu\\
\nu = c_1 \lambda + c_2 \mu + c_3 \nu\\
\lambda^2 + \mu^2 + \nu^2 = 1
\end{array}}

and these three quantities $\lambda, \mu, \nu$ determine the direction of a straight line which makes the same
angles with the axes $y_1, y_2, y_3$ as with the axes $x_1, x_2, x_3$. A line laid in this
direction through the origin of coordinates is
the axis of rotation.

In the case $|A| = - 1$ this conclusion no longer applies.

Besides this there is yet another interpretation of the orthogonal
substitution $(y) = A (x)$, according to which $(x)$ and $(y)$ are the coordinates
of two different points $x$ and $y$, referred to one and
the same coordinate system. When $x$ sweeps over a part of space (a line,
surface or body), the corresponding
point $y$ fills a congruent part of space (if $|A| = - 1$, one
that is its mirror image). This is shown by the following consideration.

The group of the \gesperrt{proper} orthogonal substitutions
is equivalent to a group that can be formed from the different
positions of a body rotatable about a fixed point.
This group is obtained by taking an arbitrary position $E$ as the unit, from which one passes to any other
position $A$ by a rotation about a definite axis through
a definite angle. If $B$ is a third position, then by
the composed position $AB$ is to be understood the position
one obtains by carrying out the rotation that led to $A$
not from the unit position, but from the
position $B$. For if one takes a system of axes standing in rigid
connection with the body, e.g. the system
of the principal axes of inertia, and denotes by $(x)$ the coordinates
%CONT>>


% ===== scan leaf 208, printed 187 =====
% !PRINTED 187
\weberpage{187}
%<<CONT
of an arbitrary point fixed in space, referred to this
system of axes in the unit position $E$, one obtains the coordinates of the same point, referred to the system of axes in the
position $A$ or $B$, by two orthogonal substitutions $A(x)$
and $B(x)$.

Accordingly $AB(x)$ are the coordinates, referred to the
system $A$, of a point $y$ with the coordinates $(y) = B(x)$
in the system $E$. The point $y$ thus has in the system $E$ the same
coordinates as the point $x$ in the system $B$; that is, $y$ lies with respect to the
unit position as $x$ does to the position $B$. If we now carry out the
rotation that leads from $E$ to $B$ in such a way that we hold the point $x$
fast, but let the point $y$ and the system of axes $A$ accompany
the rotation, then the point $y$ arrives at $x$, and the coordinates of $x$, referred to the new situation of the system $A$, are
the same as those of the point $y$ with respect to the original
situation of $A$, that is, $AB(x)$. But this new situation of the system $A$
can evidently also be reached by performing the rotation
that leads to the position $A$ not from the unit position,
but from the position $B$.

From the group of the proper ones one obtains the totality
of all orthogonal substitutions by composition with
\gesperrt{one} improper one, say with $(x_1, x_2, x_3) = (y_1, y_2, - y_3)$,
which, according to the second geometric conception, is a reflection
in the plane $x_1, x_2$, in which the point $y$ is the mirror image
of the point $x$.

Of special interest now are the \gesperrt{finite groups} contained in the totality
of the proper orthogonal substitutions, into which we shall enter more closely later.
Here we wish only to remark the following about the geometric side of this
question. To such a finite group $G$
of degree $g$ of orthogonal substitutions there corresponds a
finite group of positions of a body. If one imagines
the body fixed in these various positions simultaneously
and the whole united into a new rigid body, one
obtains a figure which can cover itself in a finite number $g$ of
distinct positions. For every
group one can find bodies of infinitely many different shapes.
The most intuitive and best-known relations result if one assumes the body bounded by plane faces.


% ===== scan leaf 209, printed 188 =====
% !PRINTED 188
\weberpage{188}

Such figures are the regular pyramids, the double pyramids
and the regular solids (tetrahedron, octahedron, cube,
dodecahedron and icosahedron).

The regular pyramid of $g$ sides is brought into coincidence with
itself by rotation about the principal axis through an angle
$2\pi:g$ and, by repetition of this rotation, in $g$ different ways.

If $g$ is even, a regular double pyramid of $g$ side faces can be
brought into coincidence with itself not only by rotation about its
principal axis through the angle $4\pi:g$, but also (in $\tfrac{1}{2}g$
different ways) by rotation about an axis lying in the equatorial
plane, with an angle of rotation of 180 degrees.

The tetrahedron admits 12 different positions in which it
occupies the same space.

To obtain them, denote the place of a vertex in the unit position
by 1. One then obtains three positions of the tetrahedron in which
the vertex 1 remains fixed. But one can bring any vertex to the
place of 1, whereby the number is quadrupled.

The same number is also found if one holds fast the plane of a
side face, in which case one again finds three positions of the
tetrahedron, and brings each side face into the initial plane.

Finally one can also count in such a way that one brings each
edge into coincidence with a fixed segment in two ways. One proceeds
in the same manner with the remaining regular solids. One thus finds
the degree of the group

\begin{quote}
equal to the product of the number of vertices by the number of the
edges or side faces meeting in a vertex, or

equal to the product of the number of side faces by the number of
sides or vertices of a bounding face, or

equal to twice the number of edges.
\end{quote}

Each of these counts yields

\begin{quote}
for the tetrahedron 12 positions,\\
for the octahedron and the cube 24 positions,\\
for the dodecahedron and icosahedron 60 positions.
\end{quote}


% ===== scan leaf 210, printed 189 =====
% !PRINTED 189
\weberpage{189}
\pdfbookmark[2]{§ 49. Linear fractional substitutions}{bm58}

Let it finally be remarked that, instead of the \gesperrt{positions} of the
body, the rotations themselves can also be conceived as elements
of the group.

\begin{center}
\S.\,49.

\gesperrt{Linear fractional substitutions.}
\end{center}

The group of orthogonal ternary substitutions is, as shall
now be proved, isomorphic with the group of linear fractional
substitutions of one variable, or of the binary substitutions of
the ratios (\S.\,38).

We denote a ternary orthogonal substitution
by
\weq{1}{(y_1,\ y_2,\ y_3) = A\,(x_1,\ x_2,\ x_3),}
\weq{2}{y_1^2 + y_2^2 + y_3^2 = x_1^2 + x_2^2 + x_3^2.}
To this we first apply, following \S.\,37, (21), the
transformation by the fixed (not orthogonal) substitution
\weq{3}{L = \left\{\begin{array}{ccc}
i, & 0, & 0 \\[4pt]
0, & \dfrac{1}{\sqrt{2}}, & \dfrac{1}{\sqrt{2}} \\[6pt]
0, & \dfrac{i}{\sqrt{2}}, & \dfrac{-i}{\sqrt{2}}
\end{array}\right\},\quad
L^{-1} = \left\{\begin{array}{ccc}
-i, & 0, & 0 \\[4pt]
0, & \dfrac{1}{\sqrt{2}}, & \dfrac{-i}{\sqrt{2}} \\[6pt]
0, & \dfrac{1}{\sqrt{2}}, & \dfrac{i}{\sqrt{2}}
\end{array}\right\}}
in which $i = \sqrt{-1}$, and whose determinant has the value 1,
and set
\weq{4}{\begin{aligned}
(y_1,\ y_2,\ y_3) &= L\,(y_1',\ y_2',\ y_3')\\
(x_1,\ x_2,\ x_3) &= L\,(x_1',\ x_2',\ x_3')' % [sic] source prints a stray prime after the closing parenthesis
\end{aligned}}
\weq{5}{(y_1',\ y_2',\ y_3') = A'\,(x_1',\ x_2',\ x_3'),}
in which
\weq{6}{A' = L^{-1} A L}
likewise denotes a linear substitution, whose determinant
$|A'| = |A|$, and hence is equal to $\pm 1$. The relation
(2) passes, by the substitutions (4), into
\weq{7}{-y_1'^2 + 2\,y_2'\,y_3' = -x_1'^2 + 2\,x_2'\,x_3'}
from which six relations between the coefficients of $A'$
result, which however we shall not set up here, since we can
find the substitution $A'$ more simply in another way.


% ===== scan leaf 211, printed 190 =====
% !PRINTED 190
\weberpage{190}

If, namely, we understand by $\xi_1, \xi_2$ new variables and set
\weq{8}{x_1' = \sqrt{2}\,\xi_1\,\xi_2,\quad x_2' = \xi_1^2,\quad x_3' = \xi_2^2}
then $x_1'^2 - 2\,x_2'\,x_3'$ vanishes identically, and
$y_1', y_2', y_3'$ pass, by (5) and (8), into binary quadratic
forms of the variables $\xi_1, \xi_2$, which by (7) must identically
satisfy the equation
\weq{9}{y_1'^2 - 2\,y_2'\,y_3' = 0}
that is, $y_2' y_3'$ must become a
perfect square. The quadratic forms $y_1', y_2',
y_3'$ we now decompose each into two linear factors. Here
$y_2'$ and $y_3'$ can have no common factor, since otherwise the
other factor too would agree in both functions, and hence
the three coefficients of $y_2'$ and $y_3'$ would have to be
proportional to one another. But then $|A'|$ would vanish,
which is not possible. Accordingly it follows from (9) that $y_2'$
and $y_3'$ must be squares of linear functions. If we set
\weq{10}{\eta_1 = \alpha\,\xi_1 + \beta\,\xi_2,\quad \eta_2 = \gamma\,\xi_1 + \delta\,\xi_2,}
then accordingly, having regard to (9), we may set
\weq{11}{y_1' = \sqrt{2}\,\eta_1\,\eta_2,\quad y_2' = \eta_1^2,\quad y_3' = \eta_2^2}
If we carry out the multiplications and then carry out the
substitution (8) in the reverse sense, we obtain
from (10) and (11):
\[
\begin{aligned}
y_1' &= (\alpha\delta + \beta\gamma)\,x_1' + \sqrt{2}\,\alpha\gamma\,x_2' + \sqrt{2}\,\beta\delta\,x_3', \\
y_2' &= \sqrt{2}\,\alpha\beta\,x_1' + \alpha^2\,x_2' + \beta^2\,x_3', \\
y_3' &= \sqrt{2}\,\gamma\delta\,x_1' + \gamma^2\,x_2' + \delta^2\,x_3'.
\end{aligned}
\]
Thus the substitution $A'$ reads:
\weq{12}{A' = \begin{pmatrix}
\alpha\delta + \beta\gamma, & \sqrt{2}\,\alpha\gamma, & \sqrt{2}\,\beta\delta \\
\sqrt{2}\,\alpha\beta, & \alpha^2, & \beta^2 \\
\sqrt{2}\,\gamma\delta, & \gamma^2, & \delta^2
\end{pmatrix}.}
By this substitution the condition (7) is not yet completely
satisfied; it follows only that the right and left sides
differ by a constant factor, which depends on the
coefficients $\alpha, \beta, \gamma, \delta$. Now, however, the coefficient
of $x_1'^2$ in the combination $y_1'^2 - 2\,y_2'\,y_3'$ is:
\[
(\alpha\delta + \beta\gamma)^2 - 4\,\alpha\delta\beta\gamma = (\alpha\delta - \beta\gamma)^2,
\]
%CONT>>


% ===== scan leaf 212, printed 191 =====
% !PRINTED 191
\weberpage{191}
%<<CONT
and we must therefore set the determinant $\alpha\delta - \beta\gamma = \pm 1$.
If one computes the determinant $|A'|$, one obtains for it
$(\alpha\delta - \beta\gamma)^3$, so that, if we take account only
of the proper orthogonal substitutions, we must set
\weq{13}{\alpha\delta - \beta\gamma = 1}
while
\weq{14}{\alpha\delta - \beta\gamma = -1}
corresponds to the improper orthogonal substitutions.
By this the substitution $A'$ is completely determined, and is
reduced to the binary substitution
\weq{15}{(\eta_1, \eta_2) = \begin{pmatrix} \alpha, & \beta \\ \gamma, & \delta \end{pmatrix} (\xi_1, \xi_2),\quad \alpha\delta - \beta\gamma = \pm 1,}
or, if $\eta_1 : \eta_2 = \eta$, $\xi_1 : \xi_2 = \xi$ is set,
to the linear fractional substitution
\weq{16}{\eta = \dfrac{\alpha\xi + \beta}{\gamma\xi + \delta}.}
The following point remains to be noted. The two substitutions
$A, A'$ are transformations of one another, and thus correspond
to each other one-to-one.

The numbers $\alpha, \beta, \gamma, \delta$ are determined by $A'$
only up to a common sign. If, then, we consider the
linear substitution
\weq{17}{\begin{pmatrix} \alpha, & \beta \\ \gamma, & \delta \end{pmatrix}}
with one of the two conditions (13) and (14), then the substitution
$A'$, and hence also $A$, is indeed uniquely determined by this;
but conversely to each $A'$ there correspond two
substitutions of the form (17):
\[
\begin{pmatrix} \alpha, & \beta \\ \gamma, & \delta \end{pmatrix}
\quad\text{and}\quad
\begin{pmatrix} -\alpha, & -\beta \\ -\gamma, & -\delta \end{pmatrix}.
\]
But we again obtain a one-to-one correspondence if we
combine these two substitutions into a common concept, which
we denote by $A''$.

Finally, as regards the substitution (16), which we denote by
\weq{18}{A''' = \begin{pmatrix} \alpha, & \beta \\ \gamma, & \delta \end{pmatrix}}
and, following \S.\,38, conceive as a substitution of the ratios,
we can remove a common factor of
$\alpha, \beta, \gamma, \delta$
%CONT>>


% ===== scan leaf 213, printed 192 =====
% !PRINTED 192
\weberpage{192}
%<<CONT
be chosen in two ways, satisfying condition (13) or condition (14), respectively.
Accordingly we obtain from every substitution $A'''$ two
orthogonal substitutions $A$, of which the one is proper, the
other improper.

The three substitutions
\weq{19}{A,\ A',\ A''}
are thus in one-to-one correspondence; the substitutions

\weq{20}{A,\ A',\ A'''}
however only when we further impose the requirement that $A$
is to be a proper orthogonal substitution.

If now
\weq{21}{B,\ B',\ B''}
is a second system of substitutions of the form (19), then
$AB$, $A'B'$ are likewise two substitutions corresponding to each other, as
follows immediately from the meaning of the $A'$ as the transformed substitution
 of the $A$. But it remains to be shown
that
\weq{22}{AB,\ A'B',\ A''B''}
is also an associated system of the form (19), and hence that
the groups of the $A, A', A''$ are isomorphic. From this it then follows
of itself, from the relation in which $A''$ and $A'''$ stand to one
 another, that (on restriction to proper orthogonal
substitutions $A$) the group of the $A'''$ is also isomorphic with them.

To prove this, let us set:
\weq{23}{(\eta_1,\eta_2) = A''(\xi_1,\xi_2),\quad (\xi_1,\xi_2) = B''(\zeta_1,\zeta_2),}
whence it follows that
\weq{24}{(\eta_1,\eta_2) = A''B''(\zeta_1,\zeta_2).}

From $A''$, $B''$, $A''B''$ we now derive, according to (12), three ternary
substitutions $A'$, $B'$, $C'$. Thus we obtain, by (8):
\weq{25}{\begin{aligned}
(\sqrt{2}\,\eta_1\eta_2,\ \eta_1^2,\ \eta_2^2) &= A'(\sqrt{2}\,\xi_1\xi_2,\ \xi_1^2,\ \xi_2^2)\\
(\sqrt{2}\,\xi_1\xi_2,\ \xi_1^2,\ \xi_2^2) &= B'(\sqrt{2}\,\zeta_1\zeta_2,\ \zeta_1^2,\ \zeta_2^2)\\
(\sqrt{2}\,\eta_1\eta_2,\ \eta_1^2,\ \eta_2^2) &= C'(\sqrt{2}\,\zeta_1\zeta_2,\ \zeta_1^2,\ \zeta_2^2).
\end{aligned}}

These formulae are identities in respect of $\zeta_1, \zeta_2$, and they
must therefore remain correct when $\sqrt{2}\,\zeta_1\zeta_2, \zeta_1^2, \zeta_2^2$ are replaced by three
%CONT>>


% ===== scan leaf 214, printed 193 =====
% !PRINTED 193
\weberpage{193}
\pdfbookmark[2]{§ 50. Conditions of reality}{bm59}
%<<CONT
independent variables $z_1', z_2', z_3'$. If we denote
the expressions which thereby result for
\[ \sqrt{2}\,\xi_1\xi_2,\ \xi_1^2,\ \xi_2^2;\quad \sqrt{2}\,\eta_1\eta_2,\ \eta_1^2,\ \eta_2^2 \]
by
\[ x_1', x_2', x_3';\quad y_1', y_2', y_3', \]
then the formulae (25) yield:
\weq{26}{\begin{aligned}
(y_1', y_2', y_3') &= A'(x_1', x_2', x_3')\\
(x_1', x_2', x_3') &= B'(z_1', z_2', z_3')\\
(y_1', y_2', y_3') &= C'(z_1', z_2', z_3'),
\end{aligned}}
that is,
\[ C' = A'B', \]
q. e. d.

\begin{center}
\S.\,50.

\gesperrt{Reality conditions.}
\end{center}

There still remains one question for us to answer. So far
we have nowhere taken account of the reality of the coefficients.
But if any geometric application is to be made, it is
necessary that the orthogonal substitution $A$ be real. It therefore
remains to investigate to what conditions the substitution coefficients
$\alpha, \beta, \gamma, \delta$ in $A''$ are to
be subjected, in order that the coefficients of $A$ become real.

To decide this question, we form, according to \S.\,49, (3), (12),
the composition
\[ A = LA'L^{-1}, \]
and obtain:
\weq{1}{
A = \left\{
\begin{array}{ccc}
\alpha\delta+\beta\gamma, & i(\alpha\gamma+\beta\delta), & \alpha\gamma-\beta\delta\\[4pt]
-i(\alpha\beta+\gamma\delta), & \dfrac{\alpha^2+\beta^2+\gamma^2+\delta^2}{2}, & i\,\dfrac{-\alpha^2+\beta^2-\gamma^2+\delta^2}{2}\\[8pt]
\alpha\beta-\gamma\delta, & i\,\dfrac{\alpha^2+\beta^2-\gamma^2-\delta^2}{2}, & \dfrac{\alpha^2-\beta^2-\gamma^2+\delta^2}{2}
\end{array}
\right\}.
}

The coefficients of this substitution are thus to be real,
and moreover
\weq{2}{\alpha\delta - \beta\gamma = \pm 1.}

If one of the four coefficients $\alpha, \beta, \gamma, \delta$ vanishes,
then a second must also vanish. For if, e. g.,
$\beta = 0$, then $i\alpha\gamma$ and $\alpha\gamma$ must be real. But this, together with (2)
%CONT>>


% ===== scan leaf 215, printed 194 =====
% !PRINTED 194
\weberpage{194}
%<<CONT
is compatible only if $\gamma = 0$. Likewise $\alpha$ and $\delta$ can
only vanish together.

If now $\beta$ is zero and $\gamma = 0$, then $\alpha^2 + \delta^2$ and $i(\alpha^2 - \delta^2)$
must be real, hence $\alpha^2$ and $\delta^2$ conjugate imaginary, and by (2) their
product $= 1$. According as the upper or the lower
sign holds in (2), $\alpha, \pm\delta$ must be conjugate imaginary. In the same way
it follows that, if $\alpha$ and $\delta = 0$, $\beta$ and $\mp\gamma$ must be conjugate
imaginary.

Let us now assume, then, that none of the four numbers $\alpha, \beta, \gamma, \delta$
vanishes. Then, in order that the quantities
\[
\begin{aligned}
&\alpha\delta+\beta\gamma,\quad i(\alpha\gamma+\beta\delta),\quad -i(\alpha\beta+\gamma\delta)\\
&\alpha\delta-\beta\gamma,\quad \alpha\gamma-\beta\delta,\quad \alpha\beta-\gamma\delta
\end{aligned}
\]
may be real, we must have
\[
\begin{aligned}
&\alpha\delta \text{ and } \beta\gamma \text{ real,}\\
&\alpha\gamma \text{ conjugate imaginary with } -\beta\delta,\\
&\alpha\beta \quad \text{\textquotedbl{}}\quad \text{\textquotedbl{}}\quad \text{\textquotedbl{}}\quad -\gamma\delta
\end{aligned}
\]
and accordingly also
\[
\alpha^2 = \frac{\alpha\gamma \cdot \alpha\beta}{\beta\gamma} \text{ conjugate imaginary with } \delta^2 = \frac{\beta\delta \cdot \gamma\delta}{\beta\gamma}
\]
and
\[
\beta^2 = \frac{\alpha\beta \cdot \delta\beta}{\alpha\delta} \quad \text{\textquotedbl{}}\quad \text{\textquotedbl{}}\quad \text{\textquotedbl{}}\quad \gamma^2 = \frac{\gamma\delta \cdot \alpha\gamma}{\alpha\delta};
\]
thus $\alpha$ must be conjugate with $\pm\delta$, and $\beta$ conjugate with $\mp\gamma$. That
the same sign should stand in both cases contradicts the requirement
that $\alpha\gamma$ and $-\beta\delta$ be conjugate. Since the product
of two conjugate imaginary quantities is always positive, the upper
or the lower signs hold according as the upper or the lower
sign stands in (2), that is, according as $A$ is properly
or improperly orthogonal. We thus arrive at
the following general result, in which the special
cases of vanishing $\alpha, \beta, \gamma, \delta$ are also included:

\gesperrt{The necessary and sufficient condition
that the orthogonal substitution (1), (2) be real
consists in this, that $\alpha, \pm\delta$ and $\beta, \mp\gamma$ form two conjugate
imaginary pairs, in which the upper signs hold
for the properly, the lower for the improperly
orthogonal substitutions.}


% ===== scan leaf 216, printed 195 =====
% !PRINTED 195
\weberpage{195}
\pdfbookmark[2]{§ 51. Finite groups of linear fractional substitutions. Poles of the groups}{bm60}
\begin{center}
\S.\,51.

\gesperrt{Finite groups of linear fractional substitutions.}

\gesperrt{Poles of the groups.}
\end{center}

From the developments so far it follows that, once all finite
groups of linear fractional substitutions have been found, all
finite groups of properly orthogonal ternary substitutions, and
all finite groups of binary linear substitutions with determinant
1, can be found from them without further ado.

Besides these there are also finite groups which contain,
along with the properly orthogonal substitutions, improperly
orthogonal ones as well.

These groups, to which we shall return later, cannot be
derived from the groups of linear fractional substitutions alone,
because with the fractional substitutions the distinction between
properly and improperly orthogonal substitutions is effaced.

We therefore now seek first of all to determine \gesperrt{all finite
groups of linear fractional substitutions} $^{1)}$.

We denote by $G$ such a group of degree $n$,
and by
\weq{1}{x,\ \Theta_1(x),\ \Theta_2(x),\ \ldots,\ \Theta_{n-1}(x)}
its elements, in which the $\Theta$ are symbols for linear functions
\weq{2}{\Theta(x) = \frac{\alpha x + \beta}{\gamma x + \delta}}
for which the notation
\weq{3}{\Theta = \begin{pmatrix}\alpha, & \beta\\ \gamma, & \delta\end{pmatrix}}
is also used.

From every group of the form (1) we can, by \S.\,37, derive
a whole family of isomorphic groups, if we denote by
$L$ an arbitrary linear substitution:
\weq{4}{1,\ L\Theta_1L^{-1},\ L\Theta_2L^{-1},\ \ldots,\ L\Theta_{n-1}L^{-1},}
%CONT>>

\par\noindent\rule{4cm}{0.4pt}\par
$^{1)}$ On the theory of these groups compare especially:
\gesperrt{Gordan}, Ueber endliche Gruppen linearer Transformationen einer
Veränderlichen. Mathematische Annalen, Bd. XII, S. 23 (1877);
\gesperrt{Klein}, Vorlesungen über das Ikosaëder (Leipzig 1884).


% ===== scan leaf 217, printed 196 =====
% !PRINTED 196
\weberpage{196}
%<<CONT
and we regard our problem as solved when a representative of each
of these families has been determined.

We shall later make use of this freedom in order to represent
the groups found as simply as possible.

The determinant
\[
\alpha\delta - \beta\gamma = \Delta
\]
must be different from zero, and if simplicity demands it,
we may indeed assume it $= 1$, which we shall occasionally
do.

In the sense of group theory we denote the repetition
of a substitution by exponents: $\Theta, \Theta^2, \Theta^3, \ldots$,
by which, then, are to be understood not powers, but always
linear substitutions of the group (1) again. The identical
substitution $x = x$, which is to be regarded as the unit of the
group, is also denoted by $\Theta^0$ or by 1.

Two inverse substitutions $\Theta, \Theta^{-1}$ can be represented
in the form:
\[
\Theta = \begin{pmatrix} \alpha, & \beta \\ \gamma, & \delta \end{pmatrix}, \quad
\Theta^{-1} = \begin{pmatrix} \delta, & -\beta \\ -\gamma, & \alpha \end{pmatrix}.
\]

Since the group is by hypothesis finite, each of its
elements has a definite \gesperrt{degree}, i.e.\ for each $\Theta$ there
is a definite least positive number $t$ for which $\Theta^t = 1$.
This number $t$ must be a divisor of $n$.

For what follows it is of importance to consider, for each
of the substitutions of the group (excepting the identical
substitution), the values of the variable that become equal to
their transforms, that is, the roots of the equations
\weq{5}{x = \Theta(x).}
These values we shall call the \gesperrt{poles} of the substitution $\Theta$.

The equation (5) is quadratic and, if one substitutes for
$\Theta$ the expression (2), takes the form:
\weq{6}{\gamma x^2 + (\delta - \alpha) x - \beta = 0.}
This equation has two roots, which are equal to one another
only when
\weq{7}{(\delta - \alpha)^2 + 4\beta\gamma = 0}
or
\weq{8}{(\alpha + \delta)^2 = 4\Delta}
%CONT>>


% ===== scan leaf 218, printed 197 =====
% !PRINTED 197
\weberpage{197}
%<<CONT
is. It is easy to see that this case, if $\Theta$ has a finite
degree, cannot occur.

First suppose that $\beta$ is zero or $\gamma = 0$. It follows from (7) that
$\alpha = \delta$ must hold, and both can be assumed $= 1$.

Now if $\Theta$ is one of the two substitutions
\[
\begin{pmatrix} 1, & 0 \\ \gamma, & 1 \end{pmatrix}, \quad
\begin{pmatrix} 1, & \beta \\ 0, & 1 \end{pmatrix}
\]
then for every $\lambda$:
\[
\Theta^\lambda = \begin{pmatrix} 1, & 0 \\ \lambda\gamma, & 1 \end{pmatrix} \text{ or } \begin{pmatrix} 1, & \lambda\beta \\ 0, & 1 \end{pmatrix},
\]
as one easily finds by complete induction. Therefore,
since $\beta$ and $\gamma$ cannot both be zero if $\Theta$ is not
the identical substitution, $\Theta^\lambda$ cannot equal 1 for any positive $\lambda$,
as it would have to if $\lambda$ were equal to the degree of
$\Theta$.

But if $\beta$ is different from zero, then, assuming now
$\Delta = 1$ and, by (8), $\alpha + \delta = 2$, we set
\[
L = \begin{pmatrix} \beta, & 0 \\ 1-\alpha, & \beta^{-1} \end{pmatrix},
\]
and obtain
\[
\begin{aligned}
L^{-1}\Theta L &= \begin{pmatrix} \beta^{-1}, & 0 \\ \alpha-1, & \beta \end{pmatrix} \begin{pmatrix} \alpha, & \beta \\ \gamma, & \delta \end{pmatrix} \begin{pmatrix} \beta, & 0 \\ 1-\alpha, & \beta^{-1} \end{pmatrix} \\
&= \begin{pmatrix} 1, & \beta^{-1} \\ 0, & 1 \end{pmatrix},
\end{aligned}
\]
hence
\[
\Theta = L \begin{pmatrix} 1, & \beta^{-1} \\ 0, & 1 \end{pmatrix} L^{-1},
\]
and from this
\[
\Theta^\lambda = L \begin{pmatrix} 1, & \lambda\beta^{-1} \\ 0, & 1 \end{pmatrix} L^{-1},
\]
which again cannot equal 1 for any positive $\lambda$. We
therefore have the theorem:

\par\noindent 1.\ \gesperrt{Each of the $n-1$ non-identical substitutions
of a group of degree $n$ has two poles different
from one another.}

Every non-identical substitution of the group $G$ thus gives us
two poles. But one and the same value can occur as a pole
in several different substitutions. If we count
one of these values $h$ times when it occurs as a pole in $h$ substitutions of the
group, then the number of poles results
%CONT>>


% ===== scan leaf 219, printed 198 =====
% !PRINTED 198
\weberpage{198}
%<<CONT
equal to $2n-2$. These values shall be called the \gesperrt{poles of the
group}. The exact enumeration of these poles gives us the
most important information about the number and nature of the
possible finite groups.

If the substitution coefficients $\beta, \gamma$ are equal to zero, then
$\Theta$ is a multiplicative substitution
\[
\begin{pmatrix} \alpha, & 0 \\ 0, & \delta \end{pmatrix}.
\]

In order that this substitution can be of finite degree,
the quotient $\alpha : \delta = \varepsilon$ must be a \gesperrt{root of unity
whose degree is a divisor of $n$}. As poles of this substitution
we regard $x = 0$ and $x = \infty$, both of which satisfy the following condition:
\[
x = \varepsilon x
\]

By (4) one can derive from $G$ an isomorphic group
\[
L G L^{-1} = G'
\]
if for $L$ an arbitrary linear substitution
\[
L = \begin{pmatrix} A, & B \\ C, & D \end{pmatrix}
\]
is taken. The poles of this group are obtained by applying
the substitution $L$ to the poles of $G$. If, then,
$a$ is one of the poles of $G$, then
\[
a' = \frac{Aa + B}{Ca + D}
\]
is the corresponding pole of $G'$.

One can determine the substitution coefficients $A, B, C, D$ in such a way
that three of the poles of $G'$ receive prescribed values.
In order, for example, to make the poles $0, \infty, 1$ of $G'$ correspond
to the poles $a, b, c$ of $G$, one sets
\[
L(x) = \frac{c-b}{c-a} \, \frac{x-a}{x-b}.
\]
If $a$ and $b$ are the poles of one and the same substitution $\Theta$
of $G$, then the corresponding substitution of $G'$ is accordingly
multiplicative.

From this we obtain the theorem:

\par\noindent 2.\ \gesperrt{For every group of linear substitutions one
can find a transformed group in which any one
of the given substitutions except the identity}
%CONT>>


% ===== scan leaf 220, printed 199 =====
% !PRINTED 199
\weberpage{199}
\pdfbookmark[2]{§ 52. The different kinds of possible groups}{bm61}
%<<CONT
\gesperrt{corresponds to a multiplication. The transforming substitution $L$ is thereby
itself determined up to a multiplication.}

\begin{center}
\S.\,52.

\gesperrt{The different kinds of possible groups.}
\end{center}

Let $a$ be one of the poles of the group $G$, and let us
assume that there are $\nu - 1$ and no more elements in $G$,
$\Theta_1, \Theta_2, \ldots, \Theta_{\nu-1}$, such that
\weq{1}{a = \Theta_1(a) = \Theta_2(a) \cdots = \Theta_{\nu-1}(a)}
holds. Such a pole shall be called a $\nu$-fold pole. It is
now first of all clear that the elements
\weq{2}{1, \Theta_1, \Theta_2, \ldots, \Theta_{\nu-1}}
form a group by themselves, and indeed a divisor of $G$;
for from
\[
a = \Theta_1(a), \quad a = \Theta_2(a)
\]
it follows, if on the right-hand side of the second equation
one replaces $a$ by $\Theta_1(a)$, which is equal to it:
\[
a = \Theta_2 \Theta_1(a).
\]
\gesperrt{Hence $\nu$ must be a divisor of $n$:}
\weq{3}{n = \nu\mu.}
We denote the group (2) by $Q$.

It can be proved that this group must be cyclic, that is,
consists of the powers of one of its elements.

For if by transformation according to Theorem 2.
we send $a$ to infinity, the substitutions (2) take the
form:
\[
x, \; \varepsilon_1 x + c_1, \; \varepsilon_2 x + c_2, \ldots, \varepsilon_{\nu-1} x + c_{\nu-1},
\]
and the composition of two among them gives:
\[
\begin{aligned}
\Theta_1 \Theta_2 &= \varepsilon_1 \varepsilon_2 x + (c_2 \varepsilon_1 + c_1) \\
\Theta^\lambda &= \varepsilon^\lambda x + c(\varepsilon^{\lambda-1} + \varepsilon^{\lambda-2} + \cdots + 1) \\
\Theta^{-1} &= \varepsilon^{-1} x - c \varepsilon^{-1}.
\end{aligned}
\]
From this one concludes that all $\varepsilon_1, \varepsilon_2, \ldots, \varepsilon_{\nu-1}$ must be roots
of unity of degree $\nu$. Moreover, no two of the $\varepsilon$ can be
equal to one another. For if, e.g., $\varepsilon_1 = \varepsilon_2$, then we should have
\[
\Theta_1 \Theta_2^{-1} = x + (c_1 - c_2).
\]


% ===== scan leaf 221, printed 200 =====
% !PRINTED 200
\weberpage{200}
This substitution must occur in $G$, and by Theorem 1 we must have $c_1 = c_2$, and hence $\Theta_1$ identical with $\Theta_2$. The numbers $1, \varepsilon_1, \varepsilon_2, \ldots, \varepsilon_{\nu-1}$ must therefore together contain all the $\nu$th roots of unity, and hence coincide with the powers of a primitive one among them. Since, on the other hand, all the powers of one of the quantities $\Theta$ must occur in group (1), it follows that, if the $\varepsilon$ occurring in $\Theta$ is a primitive $\nu$th root of unity, the whole group (1) is exhausted by the powers of $\Theta$. % [sic/reference: both occurrences of ``group (1)'' are printed; the stabilizer group Q in (2), named in the next sentence, is required by the argument.]
Hence this group $Q$ consists of the elements

\weq{4}{1, \Theta, \Theta^2, \ldots, \Theta^{\nu-1}.}

If we take (following Theorem 2) $\Theta$ as a multiplicative substitution, then 0 is the second pole of $\Theta$, and at the same time the second pole of all the substitutions (4). This pole is therefore likewise at least \gesperrt{$\nu$-fold}. But since the two poles of $\Theta$ can be interchanged in this consideration, we obtain the theorems:

3. \gesperrt{A $\nu$-fold pole determines a cyclic group $Q$ of degree $\nu$ contained in $G$.}

4. \gesperrt{The two poles of any substitution $\Theta$ of the group $G$ are poles of equal multiplicity, and are the only two poles of the group $Q$.}

By \S.\,2, $\mu - 1$ elements

\weq{5}{\psi_1, \psi_2, \ldots, \psi_{\mu-1}}

can be so selected in $G$ that $\psi_1 Q, \psi_2 Q, \ldots, \psi_{\mu-1} Q$ are the cosets of $Q$, and hence

\weq{6}{G = Q + \psi_1 Q + \psi_2 Q + \cdots + \psi_{\mu-1} Q}

holds. We now set

\weq{7}{a_1 = \psi_1(a), a_2 = \psi_2(a), \ldots, a_{\mu-1} = \psi_{\mu-1}(a).}

The quantities $a_1, a_2, \ldots, a_{\mu-1}$ are not only all different from $a$, but also different from one another. For if, say,

\[ \psi_1(a) = \psi_2(a), \]

it would follow that

\[ a = \psi_1^{-1}\psi_1(a) = \psi_1^{-1}\psi_2(a), \]

and $\psi_1^{-1}\psi_2$ would be contained in $Q$, hence $\psi_2$ in $\psi_1 Q$, which is contrary to the assumption. It is now easy to show that these values $a_1, a_2, \ldots, a_{\mu-1}$ are all $\nu$-fold poles of the group. It suffices to show this for $a_1$.


% ===== scan leaf 222, printed 201 =====
% !PRINTED 201
\weberpage{201}
We take any one of the equations (1)

\weq{8}{a = \Theta(a)}

and in the equation

\weq{9}{a_1 = \psi_1(a)}

replace $a$ by the equal value $\Theta(a)$. We thus obtain

\weq{10}{a_1 = \psi_1 \Theta(a).}

But by (9) we have $a = \psi_1^{-1}(a_1)$, and consequently from (10) there results:

\weq{11}{a_1 = \psi_1 \Theta \psi_1^{-1}(a_1).}

From this one sees, since we have $\nu - 1$ distinct substitutions $\Theta$, that $a_1$ is a $\nu$-fold pole, and indeed one belonging to the group $\psi_1 Q \psi_1^{-1}$ conjugate to $Q$. For this reason we call

\weq{12}{a, a_1, a_2, \ldots, a_{\mu-1}}

\gesperrt{a system of conjugate $\nu$-fold poles of $G$.}

An arbitrary substitution $\chi$ of $G$ can, by (6), always be brought into the form $\psi_i \Theta$, so that $\Theta$ belongs to $Q$, and from this it follows that $\chi(a) = \psi_i \Theta(a) = a_i$; and likewise, if one sets $\psi_i \Theta \psi_k = \psi_h \Theta'$, so that $\Theta'$ also belongs to $Q$:

\[ \chi(a_k) = \psi_i \Theta \psi_k(a) = \psi_h \Theta'(a) = a_h. \]

Since, further, two quantities $\chi(a_h), \chi(a_k)$ can be equal to one another only when $a_h = a_k$, the theorem follows:

5. \gesperrt{The two sequences}

\[ a, \quad a_1, \quad a_2, \quad \ldots, \quad a_{\mu-1} \]
\[ \chi(a), \quad \chi(a_1), \quad \chi(a_2), \quad \ldots, \quad \chi(a_{\mu-1}) \]

\gesperrt{agree with one another, apart from the order, whatever substitution from $G$ may be taken for $\chi$.}

If now $b$ is a pole not contained in the system (12), then $\chi(b)$ also cannot occur in (12), and from this it results that two systems of conjugate poles, if they are not entirely identical, have no pole in common.

Accordingly we can arrange all the poles of the group $G$ into systems of conjugate poles, and by \S.\,51 we obtain their total number $2n - 2$ if we count each $\nu$-fold pole $(\nu - 1)$ times.


% ===== scan leaf 223, printed 202 =====
% !PRINTED 202
\weberpage{202}
This remark gives us an essential limitation of the numbers $\nu$. For according to it:

\weq{13}{2n - 2 = \mu(\nu - 1) + \mu'(\nu' - 1) + \mu''(\nu'' - 1) + \cdots}

or, if we denote by $h$ the number of systems of conjugate poles, and set $n = \mu \nu = \mu' \nu' = \cdots$:

\weq{14}{2n - 2 = nh - \mu - \mu' - \mu'' - \cdots}

The numbers $\nu, \nu', \nu'' \ldots$ are at least equal to 2, hence

\[ 1 \leqq \mu \leqq \frac{n}{2}, \quad 1 \leqq \mu' \leqq \frac{n}{2}, \cdots \]

and therefore by (14)

\[ \frac{nh}{2} \leqq 2n - 2 \leqq (n-1)h, \]

or

\[ 2 \leqq h \leqq 4 - \frac{4}{n}. \]

Thus $h$ can have only one of the two values 2 or 3, and we find five ways of satisfying equation (13).

If first $h = 2$, then it follows from (14) that

\[ \mu + \mu' = 2, \quad \mu = \mu' = 1, \quad \nu = \nu' = n. \]

We therefore have:

I. \gesperrt{Cyclotomic group} or \gesperrt{cyclic group:}

\[ \nu = \nu' = n, \quad n \text{ arbitrary.} \]
\[ \mu = \mu' = 1, \]

If further $h = 3$, then it follows from (14):

\weq{15}{\mu + \mu' + \mu'' = n + 2.}

From this it is to be concluded that at least one of the numbers $\nu, \nu', \nu''$ must be equal to 2. For if they were all $\geqq 3$, we should have

\[ \mu \leqq \frac{n}{3}, \quad \mu' \leqq \frac{n}{3}, \quad \mu'' \leqq \frac{n}{3}, \]

hence $\mu + \mu' + \mu'' \leqq n$, which stands in contradiction with (15).

If then $\nu = 2$, $\mu = \dfrac{n}{2}$, it follows from (15):

\weq{16}{\mu' + \mu'' = \frac{n}{2} + 2.}

We assume first that $\nu' = 2$ as well. If we then set $\nu'' = m$, it follows from (16):

\[ \mu'' = 2, \quad n = 2m, \]
%CONT>>


% ===== scan leaf 224, printed 203 =====
% !PRINTED 203
\weberpage{203}
\pdfbookmark[2]{§ 53. Transformation of the substitutions of G to simple forms}{bm62}
%<<CONT
and we obtain a second possibility:

\gesperrt{II. Dihedral group,} $\nu = \nu' = 2, \nu'' = m, n = 2m, m \geqq 2$

\[ \mu = \mu' = m, \quad \mu'' = 2. \]

If, further, neither of the numbers $\nu', \nu''$ is equal to 2, then one of them must be equal to 3. For if they are both $\geqq 4$, then

\[ \mu' + \mu'' \leqq \frac{n}{2}, \]

which stands in contradiction with (16). If then $\nu' = 3$, $\mu' = \dfrac{n}{3}$, it follows from (16):

\weq{17}{\mu'' = \frac{n}{6} + 2, \quad \nu'' = \frac{6n}{n+12},}

from which it follows that we must have $\nu'' < 6$, and hence can have only one of the values 3, 4, 5.

Accordingly we obtain three more possible cases:

\gesperrt{III. Tetrahedral group:} $\nu = 2, \nu' = 3, \nu'' = 3, n = 12$

\[ \mu = 6, \quad \mu' = 4, \quad \mu'' = 4. \]

\gesperrt{IV. Octahedral group:} $\nu = 2, \nu' = 3, \nu'' = 4, n = 24$

\[ \mu = 12, \quad \mu' = 8, \quad \mu'' = 6. \]

\gesperrt{V. Icosahedral group:} $\nu = 2, \nu' = 3, \nu'' = 5, n = 60$

\[ \mu = 30, \quad \mu' = 20, \quad \mu'' = 12. \]

With this all possibilities are exhausted.

It is, to be sure, not yet proved that these groups, which we have provisionally listed under their customary names, and which we gather under the common name of the \gesperrt{polyhedral groups}, actually exist, nor how great their manifoldness is. Only the following considerations lead to this proof.

\begin{center}
\S.\,53.

\gesperrt{Transformation of the substitutions of $G$ to simple forms.}
\end{center}

Before we proceed to the definitive establishment of these groups, we derive a theorem which restricts still further the possibility of the substitutions that occur.


% ===== scan leaf 225, printed 204 =====
% !PRINTED 204
\weberpage{204}

If $a$ is a $\nu$-fold pole of the group $G$ and

\weq{1}{a = \Theta(a) = \Theta^2(a)\cdots = \Theta^{\nu-1}(a),}

then by \S.\,52, 4. the second pole $a'$ of $\Theta$ is likewise $\nu$-fold,
and we have

\weq{2}{a' = \Theta(a') = \Theta^2(a')\cdots = \Theta^{\nu-1}(a').}

Now if the group $G$ has more than one system of conjugate
$\nu$-fold poles, then $a, a'$ may occur either in the same
or else in different ones of these systems.

But if there is only \gesperrt{one} system of $\nu$-fold poles, then $a$
and $a'$ must occur in the same system, and there must therefore
exist in $G$ a substitution $\psi$ not contained among the powers of $\Theta$,
such that

\weq{3}{a' = \psi(a)}

holds. From this it follows, by applying (1) and (2):
\[
\psi(a) = \Theta\psi(a), \quad a = \psi^{-1}\Theta\psi(a),
\]
from which we conclude that $\psi^{-1}\Theta\psi$ occurs among the powers of
$\Theta$, and hence, by (2), the following relations must hold:
\[
a' = \psi^{-1}\Theta\psi(a'), \quad \psi(a') = \Theta\psi(a')
\]
Thus $\psi(a')$ is also a pole of $\Theta$, and since $\psi(a')$
cannot be $= a'$, because otherwise $a'$ would be a pole of $\psi$, which
cannot be, since $\psi$ does not occur among the powers of $\Theta$,
we have

\weq{4}{\psi(a') = a.}

If now, by a transformation of the group, one brings the poles
$a$ and $a'$ to $0$ and $\infty$, then $\Theta$ takes the form
\[
\Theta(x) = \varepsilon x,
\]
in which $\varepsilon$ is a primitive $\nu$th root of unity, and $\psi$ must, on account of
(3) and (4), have the form:
\[
\psi(x) = \frac{c}{x},
\]
in which $c$ is a constant. By a further transformation
of the group one can give the constant $c$ any value whatever,
e.g. the value $1$, by transformation with the substitution $x\sqrt{c}$ for $x$.
We thus obtain the following theorem:

1.~\gesperrt{If in the group $G$ only one system of conjugate
 $\nu$-fold poles occurs, then one can}
%CONT>>


% ===== scan leaf 226, printed 205 =====
% !PRINTED 205
\weberpage{205}
\pdfbookmark[2]{§ 54. The fundamental forms}{bm63}
%<<CONT
\gesperrt{transform the group so that it contains the two
substitutions}
\[
\Theta = \varepsilon x, \quad \psi = \frac{c}{x}
\]
\gesperrt{where $\varepsilon$ denotes a primitive $\nu$th root of unity,
and $c$ can be an arbitrarily prescribed
value, e.g. also $1$.}

The hypothesis of this theorem is, in the cases enumerated
in \S.\,52, always satisfied for one of the various values $\nu$,
except in the case of the cyclic group and of the
dihedral group with $m = 2$.

Finally we prove also the following theorem:

2.~\gesperrt{If $a, a'$ are the poles of a substitution $\Theta$ of $G$,
then the poles conjugate to $a$ and $a'$,}
\[
b = \chi(a), \quad b' = \chi(a'),
\]
\gesperrt{where $\chi$ is an arbitrary substitution of $G$, are the
two poles of one and the same substitution,
namely of the substitution $\chi\Theta\chi^{-1}$.}

The correctness follows immediately from inspection
of the equations:
\[
\begin{aligned}
a &= \Theta(a), & \chi\Theta\chi^{-1}\chi(a) &= \chi\Theta(a) = \chi(a) \\
a' &= \Theta(a'), & \chi\Theta\chi^{-1}\chi(a') &= \chi\Theta(a') = \chi(a').
\end{aligned}
\]

\begin{center}
\S.\,54.

\gesperrt{The Fundamental Forms.}
\end{center}

In order to arrive at the definitive enumeration of all finite groups
$G$, it is necessary to enter more closely into the invariants of the
corresponding binary substitution groups.

We must therefore, besides the linear fractional substitutions


\weq{1}{y = \frac{\alpha x + \beta}{\gamma x + \delta},}

which we have denoted by $A'''$ in \S.\,49, also the binary
substitutions

\weq{2}{(y_1, y_2) = \begin{pmatrix}\alpha, & \beta \\ \gamma, & \delta\end{pmatrix}(x_1, x_2)}
%CONT>>


% ===== scan leaf 227, printed 206 =====
% !PRINTED 206
\weberpage{206}
%<<CONT
consider, which were there denoted by $A''$, and which, if one sets
$y_1 : y_2 = y$, $x_1 : x_2 = x$, lead again to the substitution (1).
In (2) it is always assumed that

\weq{3}{\alpha\delta - \beta\gamma = 1}

but the two substitutions
\[
\begin{pmatrix}\alpha, & \beta \\ \gamma, & \delta\end{pmatrix}, \quad
\begin{pmatrix}-\alpha, & -\beta \\ -\gamma, & -\delta\end{pmatrix}
\]
are regarded as identical.

It can lead to no misunderstanding if we denote both kinds
of substitutions alike by a symbol such as
\[
\Theta = \begin{pmatrix}\alpha, & \beta \\ \gamma, & \delta\end{pmatrix}
\]
since then in both the rule of composition is exactly
the same.

Now if

\weq{4}{a, a_1, a_2, \ldots, a_{\mu-1}}

is a system of conjugate poles of the group $G$, then

\weq{5}{f(x) = (x-a)(x-a_1)\cdots(x-a_{\mu-1})}

is an integral function of degree $\mu$ whose roots are those conjugate
poles. This function $f(x)$ has the following property:

If we take any substitution of the group $G$
\[
\psi(x) = \frac{\alpha x + \beta}{\gamma x + \delta},
\]
then the roots of the function

\weq{6}{(\gamma x + \delta)^{\mu} f[\psi(x)]}

are the quantities $\psi^{-1}(a), \psi^{-1}(a_1), \ldots, \psi^{-1}(a_{\mu-1})$. But these,
apart from the order, agree by \S.\,52, 5. with the
quantities $a, a_1, a_2, \ldots, a_{\mu-1}$, and consequently the
equations (5) and (6) can differ only by a constant factor.
If therefore we set

\weq{7}{x_2^{\mu} f\left(\frac{x_1}{x_2}\right) = f(x_1, x_2)}

then this form $f$ remains unchanged, apart from a constant
factor, when any substitution of the group $G$ is applied
to $(x_1, x_2)$.

By the definition \S.\,40, 2. $f(x_1, x_2)$ is therefore an
\gesperrt{invariant of the group} $G$, and we have the theorem:


% ===== scan leaf 228, printed 207 =====
% !PRINTED 207
\weberpage{207}

1.~\gesperrt{To every system of conjugate poles of the group $G$
there corresponds an invariant form, whose degree
is equal to the number of poles of the system, and
whose roots are precisely these poles.}

We denote by $f_1, f_2, \ldots$ the invariant forms belonging
to the various systems of conjugate poles, which are of the
degrees $\mu, \mu', \ldots$ and which we shall call the \gesperrt{fundamental
forms} of the group.

If now $F(x_1, x_2)$ is any invariant form of the
group $G$ and $f(x, y)$ a fundamental form, and if $F(x, 1) = 0$
has a common root with $f(x, 1) = 0$, then all the roots
of $f = 0$ must at the same time be roots of $F = 0$. For
by hypothesis the two functions $F(x, 1)$ and
$F[\psi(x), 1]$ have the same roots. Hence if $F(a, 1) = 0$,
then also $F[\psi(a), 1] = F(a_1, 1) = 0$. We therefore have the
following theorem:

2.~\gesperrt{If an invariant form $F(x_1, x_2)$ belonging to $G$
has a divisor in common with one of the fundamental forms $f(x_1, x_2)$,
then $F(x_1, x_2)$ is divisible by $f(x_1, x_2)$.}

If $F(x_1, x_2)$ is an invariant form of the group $G$, and $\xi$
a root of the equation $F(x, 1) = 0$, then, as
$\psi$ runs through the elements of the group $G$, all the
$n$ quantities $\psi(\xi)$ are also roots of the same equation. If therefore the
degree of $F$ is lower than the degree of the group, these
quantities cannot all be different from one another, and it
follows for any two different substitutions
$\chi, \psi$ of $G$ that
\[
\chi(\xi) = \psi(\xi)
\]
or
\[
\xi = \chi^{-1}\psi(\xi),
\]
that is, $\xi$ must occur among the poles of the group $G$, and
$F$ is therefore, by Theorem 2., divisible by one of the fundamental
forms. Since we can apply the same mode of inference to the
quotient of the division, it follows:

3.~\gesperrt{An invariant form $F$ of the group $G$, whose
degree is lower than the degree of the group,
is, apart from a constant factor, a
product of powers of the fundamental forms.}
%CONT>>


% ===== scan leaf 229, printed 208 =====
% !PRINTED 208
\weberpage{208}
%<<CONT
It is here always assumed that the coefficients
of $F(x_1, x_2)$ are not all equal to zero. If we drop this
assumption, we can therefore state Theorem 3.\ also
as follows:

4.\ \gesperrt{If $F(x_1, x_2)$ is an invariant form of the group $G$
of lower degree than $G$, which cannot be
represented as a product of the fundamental forms,
then $F(x_1, x_2)$ must vanish identically.}

\rule{4cm}{0.4pt}


% ===== scan leaf 230, printed 209 =====
\weberpage{209}% page number not printed on this page
\pdfbookmark[1]{Eighth Section. The Polyhedral Groups.}{bm64}
\pdfbookmark[2]{§ 55. The cyclic groups and the dihedral groups}{bm65}
% !PRINTED none

\begin{center}
\gesperrt{Eighth Section.}

\textbf{\gesperrt{The Polyhedral Groups.}}

\rule{3cm}{0.4pt}

\S.\ 55.
\end{center}

\begin{center}
\gesperrt{The cyclic groups and the dihedral groups.}
\end{center}

We now proceed to apply the general principles to the actual formation of the
various polyhedral groups, and thus first to show that the kinds of these groups
recognised as possible in \S.\,52 all exist.

For the cyclic groups of degree $n$ we have only two
systems of conjugate poles, each of which contains only one $(n-1)$-fold
pole. These two poles must therefore be the common
poles of all substitutions of the group, and by
\S.\,51, 2. may be taken as $0$ and $\infty$. We then have
only the two linear fundamental forms
\weq{1}{f_1 = x_1, \quad f_2 = x_2.}
All substitutions of the group are multiplicative, and the
multiplier must be an $n$th root of unity. They must therefore
have the form:
\weq{2}{x, \varepsilon x, \varepsilon^2 x, \ldots, \varepsilon^{n-1} x,}
if $\varepsilon$ is a primitive $n$th root of unity.

We thus obtain in fact for every $n$ a cyclic
group, which we shall denote by $C_n$:
\weq{3}{C_n = \begin{pmatrix} \varepsilon^r, & 0 \\ 0, & 1 \end{pmatrix}, \quad r = 0, 1, \ldots, n-1,}
or written with the determinant $1$:
\weq{4}{C_n = \begin{pmatrix} \varepsilon^{1/2\,r}, & 0 \\ 0, & \varepsilon^{-1/2\,r} \end{pmatrix}.}


% ===== scan leaf 231, printed 210 =====
% !PRINTED 210
\weberpage{210}
For the dihedral group of degree $n = 2m$, which we shall denote by $D_m$, we have, by the preceding section, two systems of $m$ conjugate two-fold poles each, and one system of two conjugate $m$-fold poles. The latter must therefore, by the theorem of \S.\,53, be the poles of a substitution. If we let them coincide with $0$ and $\infty$, we obtain the fundamental form of the second degree:
\[
f_1 = x_1 x_2.
\]
But this can be an invariant form of the group only if all the substitutions are contained in one of the two forms
\[
\begin{pmatrix} \alpha, & 0 \\ 0, & \delta \end{pmatrix}, \quad \begin{pmatrix} 0, & \beta \\ \gamma, & 0 \end{pmatrix}
\]
and here $\alpha:\delta$ and $-\beta:\gamma$ must be $m$th roots of unity. We therefore obtain the substitutions of the group, if $\varepsilon$ is a primitive $m$th root of unity:
\[
x, \varepsilon x, \varepsilon^2 x, \ldots, \varepsilon^{m-1} x
\]
\[
\frac{1}{x}, \frac{\varepsilon}{x}, \frac{\varepsilon^2}{x}, \cdots, \frac{\varepsilon^{m-1}}{x},
\]
i.\,e.
\weq{5}{D_m = \begin{pmatrix} \varepsilon^r, & 0 \\ 0, & 1 \end{pmatrix}, \begin{pmatrix} 0, & \varepsilon^r \\ 1, & 0 \end{pmatrix},}
or, represented with the determinant $+1$:
\weq{6}{D_m = \begin{pmatrix} \varepsilon^{1/2\,r}, & 0 \\ 0, & \varepsilon^{-1/2\,r} \end{pmatrix}, \begin{pmatrix} 0, & i\varepsilon^{1/2\,r} \\ i\varepsilon^{-1/2\,r}, & 0 \end{pmatrix}, \quad r = 0, 1, \ldots, m-1.}
One sees at once that these substitutions really form a group.

Besides $0$ and $\infty$ we have here also the poles $\pm \varepsilon^{1/2\,h}$ arising from the equations
\[
x = \frac{\varepsilon^h}{x}
\]
from which one further obtains the two fundamental forms of degree $m$:
\weq{7}{f_2 = x_1^m + x_2^m, \quad f_3 = x_1^m - x_2^m}

From this one sees that the dihedral groups and the cyclic groups are entirely different from one another, since their fundamental forms have different degrees.


% ===== scan leaf 232, printed 211 =====
% !PRINTED 211
\weberpage{211}
By transformation by means of the substitution

\[
\begin{pmatrix} e^{\frac{\pi i}{4}}, & 0 \\ 0, & e^{-\frac{\pi i}{4}} \end{pmatrix}
\]

the dihedral group $D_m$ passes over into

\[
\begin{pmatrix} \varepsilon^{1/2\,\lambda}, & 0 \\ 0, & \varepsilon^{-1/2\,\lambda} \end{pmatrix}, \quad
\begin{pmatrix} 0, & \varepsilon^{1/2\,\lambda} \\ -\varepsilon^{-1/2\,\lambda}, & 0 \end{pmatrix}
\]

or into

\[
\begin{gathered}
x, \; \varepsilon x, \; \varepsilon^2 x, \; \ldots, \; \varepsilon^{m-1} x\\[4pt]
\dfrac{-1}{x}, \; \dfrac{-\varepsilon}{x}, \; \dfrac{-\varepsilon^2}{x}, \; \cdots, \; \dfrac{-\varepsilon^{m-1}}{x}
\end{gathered}
\]

and the fundamental forms for this representation are:

\[
f_1 = x_1 x_2, \quad f_2 = x_1^m + i^m x_2^m, \quad f_3 = x_1^m - i^m x_2^m.
\]

The dihedral group $D_m$ contains as divisor the cyclic
group $C_m$, and indeed $C_m$ is a \gesperrt{normal divisor} of $D_m$.

By transformation with any multiplicative
substitution

\[
\begin{pmatrix} \alpha, & 0 \\ 0, & \delta \end{pmatrix}
\]

$C_m$ does not change, while the substitutions of $D_m$ not contained
in $C_m$ change their form; for we have:

\[
\begin{gathered}
\begin{pmatrix} \delta, & 0 \\ 0, & \alpha \end{pmatrix}
\begin{pmatrix} \varepsilon, & 0 \\ 0, & 1 \end{pmatrix}
\begin{pmatrix} \alpha, & 0 \\ 0, & \delta \end{pmatrix}
=
\begin{pmatrix} \varepsilon, & 0 \\ 0, & 1 \end{pmatrix},\\[6pt]
\begin{pmatrix} \delta, & 0 \\ 0, & \alpha \end{pmatrix}
\begin{pmatrix} 0, & \varepsilon \\ 1, & 0 \end{pmatrix}
\begin{pmatrix} \alpha, & 0 \\ 0, & \delta \end{pmatrix}
=
\begin{pmatrix} 0, & \varepsilon \delta^2 \\ \alpha^2, & 0 \end{pmatrix}.
\end{gathered}
\]

Apart from such a multiplicative substitution, however,
the group $D_m$ is completely determined by the group $C_m$
contained in it. For if we set

\[
D_m = C_m + \varphi C_m,
\]
in which

\[
\varphi = \begin{pmatrix} \alpha, & \beta \\ \gamma, & \delta \end{pmatrix}
\]

is put, then, since $C_m$ must be a normal divisor of $D_m$, we must
have $\varphi C_m = C_m \varphi$, i.\,e.\ for every exponent $r$ there must
be an exponent $s$, and conversely, such that

\[
\begin{pmatrix} \alpha, & \beta \\ \gamma, & \delta \end{pmatrix}
\begin{pmatrix} \varepsilon^r, & 0 \\ 0, & 1 \end{pmatrix}
=
\begin{pmatrix} \varepsilon^s, & 0 \\ 0, & 1 \end{pmatrix}
\begin{pmatrix} \alpha, & \beta \\ \gamma, & \delta \end{pmatrix}
\]
%CONT>>


% ===== scan leaf 233, printed 212 =====
% !PRINTED 212
\weberpage{212}
\pdfbookmark[2]{§ 56. The tetrahedral group}{bm66}
%<<CONT
or
\[
\begin{pmatrix} \alpha\,\varepsilon^r, & \beta \\ \gamma\,\varepsilon^r, & \delta \end{pmatrix} = \begin{pmatrix} \varepsilon^s\alpha, & \varepsilon^s\beta \\ \gamma, & \delta \end{pmatrix}.
\]

If $\beta$ is zero and $\gamma = 0$, then $\varphi$ itself would be of the form $\varepsilon x$,
and $\varepsilon$ a $2m$th root of unity, and $D_m$ would therefore be a cyclic
group and not a dihedral group. But if $\beta$ or $\gamma$ is different
from zero, it follows that $\varepsilon^s = \varepsilon^{-r}$ and $\alpha = \delta = 0$, hence
\[
\varphi = \begin{pmatrix} 0, & \beta \\ \gamma, & 0 \end{pmatrix},
\]
which by a multiplicative transformation can be brought to the form
$\begin{pmatrix} 0, & 1 \\ 1, & 0 \end{pmatrix}$.

Among the dihedral groups the group $D_2$ is especially
to be singled out, in which there are three systems of two two-fold
poles each. It consists of the four substitutions
\[
\begin{pmatrix} 1, & 0 \\ 0, & 1 \end{pmatrix},\quad
\begin{pmatrix} -1, & 0 \\ 0, & 1 \end{pmatrix},\quad
\begin{pmatrix} 0, & 1 \\ 1, & 0 \end{pmatrix},\quad
\begin{pmatrix} 0, & 1 \\ -1, & 0 \end{pmatrix}
\]
and is called the \gesperrt{four-group}.

\begin{center}
\S.\,56. \\[6pt]
\gesperrt{The tetrahedral group.}
\end{center}

For the tetrahedral group we have, by \S.\,52, III., a
system of six conjugate two-fold poles. We can
therefore assume, by \S.\,53, that in the group the two substitutions

\[
\Theta(x) = -x, \quad \psi(x) = \frac{1}{x}
\]
occur, and consequently also
\[
\psi_1(x) = \psi\,\Theta(x) = -\frac{1}{x}.
\]

The two substitutions $\psi$ and $\psi_1$ give the poles $\pm 1$
and $\pm i$, and these must, since $\psi$ and $\psi_1$ are of the $2$nd degree
and in $G$ only two-fold or three-fold poles occur at all,
belong to the two-fold ones (\S.\,52, 3.). The substitution $\Theta(x)$
gives the two-fold poles $0, \infty$. Accordingly the equation
on which the two-fold poles depend reads $x(x^4-1) = 0$ and we
obtain the first fundamental form of the $6$th degree:
\weq{1}{f(x_1, x_2) = x_1 x_2 (x_1^4 - x_2^4).}


% ===== scan leaf 234, printed 213 =====
% !PRINTED 213
\weberpage{213}
To find the remaining substitutions of the group, we observe that by the theorem \S.\,53, 2., if we take $\pm 1$ for $a, a'$ and $0, \infty$ for $b, b'$, we can infer the existence of a substitution $\chi$ in $G$ which satisfies the conditions
\[
\chi(-1) = 0, \quad \chi(+1) = \infty,
\]
and that $\chi$ must therefore have the form:
\weq{2}{\chi = \lambda\,\frac{x+1}{x-1},}
where $\lambda$ is a constant factor.

If we substitute the six conjugate poles $0, \infty, +1, -1,$
$+i, -i$ into $\chi$, we must obtain the same values in another order (\S.\,52, 5.). But these values are
\[
-\lambda,\ \lambda,\ \infty,\ 0,\ -\lambda i,\ \lambda i.
\]

Hence $\lambda = \pm 1$ or $= \pm i$ must hold. The values $\pm 1$ are not admissible, since otherwise
\[
\chi(i) = \mp i, \quad \Theta\chi(i) = \pm i,
\]
and so $\pm i$ would be the poles of $\chi$ or of $\Theta\chi$, whereas they are in fact only two-fold and are the poles of $\Theta\psi$. Hence $\lambda = \pm i$ must hold, and we may take the upper sign without restriction of generality, for the lower sign then corresponds to the substitution $\Theta\chi$. Therefore
\[
\chi(x) = i\,\frac{x+1}{x-1}.
\]

From this one obtains
\[
\chi^2(x) = \frac{x+i}{x-i}, \quad \chi^3(x) = x,
\]
and there result the 12 substitutions:
\weq{3}{\begin{array}{cccc}
1, & \Theta, & \psi, & \Theta\psi \\
\chi, & \Theta\chi, & \psi\chi, & \Theta\psi\chi \\
\chi^2, & \Theta\chi^2, & \psi\chi^2, & \Theta\psi\chi^2,
\end{array}}
or
\weq{4}{\pm x,\ \pm\frac{1}{x},\ \pm i\,\frac{x+1}{x-1},\ \pm i\,\frac{x-1}{x+1},\ \pm\frac{x+i}{x-i},\ \pm\frac{x-i}{x+i}.}

That these substitutions really form a group, the \gesperrt{tetrahedral group}, of which $\Theta, \psi, \chi$ are the generating elements, follows from
\weq{5}{\chi\Theta = \Theta\psi\chi,\ \chi^2\Theta = \psi\chi^2,\ \psi\Theta = \Theta\psi,\ \chi\psi = \Theta\chi,\ \chi^2\psi = \Theta\psi\chi^2,}
%CONT>>


% ===== scan leaf 235, printed 214 =====
% !PRINTED 214
\weberpage{214}
%<<CONT
which, in view of $\Theta^2 = \psi^2 = \chi^3 = 1$, follow from the three
relations
\[
\psi\Theta = \Theta\psi, \quad \chi\Theta = \Theta\psi\chi, \quad \chi\psi = \Theta\chi
\]

From (5) it further follows that the four-group
\weq{6}{1,\ \Theta,\ \psi,\ \Theta\psi}
is a \gesperrt{normal divisor} of the tetrahedral group.

The substitutions of the tetrahedral group can, in each of the
two forms
\[
\chi^\lambda \psi^\mu \Theta^\nu, \quad \Theta^\nu \psi^\mu \chi^\lambda
\]
\[
\lambda = 0, 1, 2, \quad \mu = 0, 1, \quad \nu = 0, 1
\]
be represented.

To obtain the two fundamental forms of the fourth degree $\Phi_1, \Phi_2$
belonging to the three-fold poles, one can either compute
from (5) the poles still missing, or one can proceed as
follows:

Since the fundamental forms $\Phi_1, \Phi_2$ must admit the substitutions $\Theta$ and $\psi$,
and cannot be divisible by $x_1$ and $x_2$,
they can only be of the following form:
\[
\begin{aligned}
\Phi_1 &= x_1^4 + m_1 x_1^2 x_2^2 + x_2^4 \\
\Phi_2 &= x_1^4 + m_2 x_1^2 x_2^2 + x_2^4,
\end{aligned}
\]
in which $m_1, m_2$ are constants still to be determined. But if one forms
the \gesperrt{Hessian} covariant of $\Phi_1$, one obtains by
\S.\,40, 4. an invariant of the group:
\[
\frac{1}{12}\left[\frac{\partial^2 \Phi_1}{\partial x_1^2}\,\frac{\partial^2 \Phi_1}{\partial x_2^2} - \left(\frac{\partial^2 \Phi_1}{\partial x_1\,\partial x_2}\right)^2\right] = 2m_1(x_1^4+x_2^4) + (12-m_1^2)\,x_1^2 x_2^2.
\]

This form, since no other fundamental forms of the fourth
degree exist, can agree only with $\Phi_1$ or with $\Phi_2$
(apart from a constant factor). Hence we must have
\[
\frac{12-m_1^2}{2m_1} = m_1 \ \text{or}\ m_2
\]
The first assumption, however, gives $m_1 = \pm 2$. Then $\Phi_1$
would be a square, which cannot be. There remains, then, only:
\[
m_1^2 + 2m_1 m_2 = 12,
\]
and likewise from $\Phi_2$:
\[
m_2^2 + 2m_1 m_2 = 12,
\]
hence
\[
m_1 = -m_2 = \pm 2\sqrt{-3},
\]
%CONT>>


% ===== scan leaf 236, printed 215 =====
% !PRINTED 215
\weberpage{215}
%<<CONT
and the two fundamental forms become:
\weq{7}{\begin{aligned}
\Phi_1 &= x_1^4 + 2\sqrt{-3}\,x_1^2 x_2^2 + x_2^4 \\
\Phi_2 &= x_1^4 - 2\sqrt{-3}\,x_1^2 x_2^2 + x_2^4.
\end{aligned}}

Between the three fundamental forms $f, \Phi_1, \Phi_2$ a relation can be derived by eliminating $x_1, x_2$. From (7) one finds
\[
2(x_1^4+x_2^4) = \Phi_1 + \Phi_2
\]
\[
4\sqrt{-3}\,x_1^2 x_2^2 = \Phi_1 - \Phi_2,
\]
and from (1)
\[
f^2 = x_1^2 x_2^2 (x_1^4+x_2^4)^2 - 4x_1^6 x_2^6,
\]
from which one easily computes
\weq{8}{12\sqrt{-3}\,f^2 = \Phi_1^3 - \Phi_2^3.}

The substitutions $\Theta, \psi, \chi$, when represented with determinant
1, take the form:
\weq{9}{\begin{pmatrix} i, & 0 \\ 0, & -i \end{pmatrix},\quad
\begin{pmatrix} 0, & i \\ i, & 0 \end{pmatrix},\quad
\begin{pmatrix} \dfrac{1-i}{2}, & \dfrac{1-i}{2} \\[6pt] -\dfrac{1+i}{2}, & \dfrac{1+i}{2} \end{pmatrix},}
from which it follows, by the conditions established in \S.\,50, that
the corresponding group of orthogonal ternary substitutions
is real.

If one applies the substitutions (9) with the determinant 1
to the fundamental form $f(x_1, x_2)$, a very simple
\gesperrt{calculation shows that the fundamental form $f$ is an absolute invariant}
of the binary group $G$.

The \gesperrt{Hessian} determinant of $f$ also has this property:
\[
\begin{aligned}
H &= -\frac{1}{25}\left[f''(x_1,x_1)\,f''(x_2,x_2) - f''(x_1,x_2)^2\right] \\
&= (x_1^4+x_2^4)^2 + 12\,x_1^4 x_2^4,
\end{aligned}
\]
which is nothing other than the product of the two functions
$\Phi_1, \Phi_2$; while the functions $\Phi_1$ and $\Phi_2$ themselves, under the third
of the substitutions (9), acquire a third root of unity as a
factor.


% ===== scan leaf 237, printed 216 =====
% !PRINTED 216
\weberpage{216}
\pdfbookmark[2]{§ 57. The octahedral group}{bm67}

\begin{center}
\S.\,57.

\gesperrt{The Octahedral Group.}
\end{center}

In the octahedral group there occurs a system of six conjugate
four-fold poles. We accordingly have a substitution of degree 4
whose period
\[
1, \ \Theta, \ \Theta^2, \ \Theta^3, \ \Theta^4 = 1
\]
gives a pair of these conjugate four-fold poles. By \S.\,53, 1.
we may assume in the group the substitutions:
\weq{1}{\Theta(x) = ix, \quad \psi(x) = \frac{1}{x},}
and we have
\weq{2}{\psi^2 = 1, \quad \Theta\psi = \psi\Theta^3, \quad \Theta^2\psi = \psi\Theta^2, \quad \Theta^3\psi = \psi\Theta.}
The fundamental form of degree 6 belonging to the system of the
four-fold poles must remain unchanged, up to a constant factor,
under the substitutions $\Theta$ and $\psi$, and it must moreover
contain the factor $x_1x_2$, and hence be of one of the two forms
\[
x_1x_2(x_1^4 - x_2^4), \quad x_1x_2(x_1^4 + x_2^4).
\]
It is immaterial which of the two assumptions we pursue, since
the one passes into the other by the substitution of $ix_2$ for $x_2$,
which corresponds to the passage to a transformed group. If we take
\weq{3}{f(x_1, x_2) = x_1x_2(x_1^4 - x_2^4)}
as the fundamental form of degree 6, then the six four-fold poles are
\weq{4}{0, \ \infty, \ 1, \ -1, \ i, \ -i.}
Now by \S.\,53, 2. there must be in $G$ a substitution $\chi$
such that
\[
\chi(-1) = 0, \quad \chi(+1) = \infty
\]
and if one determines the constant factor as with the tetrahedral
group [\S.\,56, (2)], it is found that
\weq{5}{\chi(x) = i\,\frac{x+1}{x-1}.}
From this we obtain
\[
\chi^2(x) = \frac{x+i}{x-i}, \quad \chi^3(x) = x,
\]
%CONT>>


% ===== scan leaf 238, printed 217 =====
% !PRINTED 217
\weberpage{217}

%<<CONT
and the relations
\weq{6}{\Theta\chi = \chi^2\psi\Theta^3, \quad \Theta^2\chi = \chi\psi, \quad \Theta^3\chi = \chi^2\Theta \\
\psi\chi = \chi\psi\Theta^2, \quad \psi\chi^2 = \chi^2\Theta^2.}
These relations, in combination with (2), show that the
substitutions
\weq{7}{\chi^{\lambda}\psi^{\mu}\Theta^{\nu}, \quad \lambda = 0, 1, 2; \quad \mu = 0, 1; \quad \nu = 0, 1, 2, 3}
do in fact form a group, because with their aid every
substitution of the form
\[
\chi^{\lambda}\psi^{\mu}\Theta^{\nu}\chi, \quad \chi^{\lambda}\psi^{\mu}\Theta^{\nu}\psi,
\]
and consequently, by repetition, every substitution
\[
\chi^{\lambda}\psi^{\mu}\Theta^{\nu}\,\chi^{\lambda'}\psi^{\mu'}\Theta^{\nu'}
\]
can be brought into the form (7). The group can be represented
explicitly in the form
\weq{8}{i^{\nu}x, \quad \frac{i^{\nu}}{x}, \quad i^{\nu}\,\frac{x - i^{\nu'}}{x + i^{\nu'}}, \quad \nu, \nu' = 0, 1, 2, 3}
and is of the 24th degree. It is the
\gesperrt{octahedral group}.

The group has a normal divisor of the 12th degree, which can
be represented in the form $\chi^{\lambda}\psi^{\mu}\Theta^{2\nu}$ and which is a
tetrahedral group.

The representation becomes in a certain respect simpler if,
in place of $\psi$, one introduces a new element
\weq{9}{\omega = \psi\Theta = \frac{-i}{x}}
which, just like $\psi$, is of the second order. Then
the group can be represented by
\[
\chi^{\lambda}\omega^{\mu}\Theta^{\nu},
\]
and in place of the relations (2) and (6) the following appear:
\[
\omega\chi = \chi^2\omega, \quad \omega\chi^2 = \chi\omega
\]
\weq{10}{\Theta\chi = \chi^2\omega\Theta^2, \; \Theta^2\chi = \chi\omega\Theta^3, \; \Theta^3\chi = \chi^2\Theta, \; \Theta^2\chi^2 = \chi^2\omega\Theta \\
\Theta\omega = \omega\Theta^3, \quad \Theta^2\omega = \omega\Theta^2, \quad \Theta^3\omega = \omega\Theta.}
But all these relations follow as consequences
from the four:
\weq{11}{\omega\chi = \chi^2\omega, \quad \Theta\omega = \omega\Theta^3, \quad \Theta\chi = \chi^2\omega\Theta^2, \quad \Theta^2\chi = \chi\omega\Theta^3,}
in combination with the formulae expressing the degrees:
\weq{12}{\chi^3 = 1, \quad \omega^2 = 1, \quad \Theta^4 = 1.}


% ===== scan leaf 239, printed 218 =====
% !PRINTED 218
\weberpage{218}

For it follows first from the second of the relations (11):
\[
\Theta^2\omega = \Theta\omega\Theta^3 = \omega\Theta^2, \quad \Theta^3\omega = \Theta\omega\Theta^2 = \omega\Theta,
\]
further from the last of (11):
\[
\Theta^3\chi = \Theta\chi\omega\Theta^3 = \chi^2\omega\Theta^2\omega\Theta^3 = \chi^2\Theta,
\]
and further:
\[
\omega\chi^2 = \chi^2\omega\chi = \chi\omega,
\]
\[
\Theta^2\chi^2 = \chi\omega\Theta^3\chi = \chi\omega\chi^2\Theta = \chi^2\omega\Theta.
\]
If now we have any three elements $\chi$, $\omega$, $\Theta$ that can be
composed according to some rule, and if in this $\chi$ is of the
third, $\omega$ of the second, $\Theta$ of the fourth degree, then it follows
from the validity of the relations (11) that the elements
$\chi^{\lambda}\omega^{\mu}\Theta^{\nu}$ form a group of degree 24, which is isomorphic
to the octahedral group. For this it only remains to show
that from these assumptions it follows that the 24 elements
$\chi^{\lambda}\omega^{\mu}\Theta^{\nu}$ are all different from one another. Let us assume
that
\[
\chi^{\lambda}\omega^{\mu}\Theta^{\nu} = \chi^{\lambda'}\omega^{\mu'}\Theta^{\nu'},
\]
then it would follow that
\[
\chi^{\lambda - \lambda'} = \omega^{\mu'}\Theta^{\nu' - \nu}\omega^{-\mu},
\]
and by (11):
\[
\chi^{\lambda - \lambda'} = \omega^{\mu' - \mu}\Theta^{\pm(\nu - \nu')}.
\]
But now it follows from (11) that $\omega\Theta$, $\omega\Theta^2$, $\omega\Theta^3$ are of the
second degree, and since $\chi$ is of the third degree, this relation
can hold only if $\lambda - \lambda'$ is divisible by 3, $\mu - \mu'$ by 2 and
$\nu - \nu'$ by 4, and hence $\chi^{\lambda} = \chi^{\lambda'}$, $\omega^{\mu} = \omega^{\mu'}$, $\Theta^{\nu} = \Theta^{\nu'}$.

By the relation following from (10),
\[
\omega = \chi\Theta\chi\Theta^2
\]
we can compose $\omega$ from $\chi$ and $\Theta$, and may therefore regard $\chi$ and
$\Theta$ as \gesperrt{generating substitutions} of the group.

The two still missing fundamental forms of the eighth and twelfth
degrees are found very easily if one first forms the Hessian
covariant of $f$.

One thus obtains the fundamental form of the eighth degree:
\weq{13}{W = x_1^8 + 14\,x_1^4x_2^4 + x_2^8,}
and if one forms from $f$ and $W$ the functional determinant,
which is again a covariant of $f$ (Bd. I, \S.\,59, 60), there
results the fundamental form of the twelfth degree:
\weq{14}{K = x_1^{12} - 33\,x_1^8x_2^4 - 33\,x_1^4x_2^8 + x_2^{12}.}


% ===== scan leaf 240, printed 219 =====
% !PRINTED 219
\weberpage{219}

The roots of $f$ correspond to the six vertices of the octahedron or
the six faces of the cube; the roots of $W$ to the eight faces of the octahedron
or vertices of the cube, and the roots of $K$, both for the octahedron
and for the cube, to the 12 edges (cf. \S.\,48).

If one represents the forms $W$, $K$ thus:
\[
W = (x_1^4 + x_2^4)^2 + 12\,x_1^4x_2^4, \quad K = (x_1^4 + x_2^4)^3 - 36\,x_1^4x_2^4(x_1^4 + x_2^4),
\]
then the identical relation subsisting between the three functions $f$, $W$, $K$
is easily derived:
\weq{15}{W^3 - K^2 = 108\,f^4.}
The octahedral group contains, as one sees, the tetrahedral group,
and arises from it by the addition of the one substitution
 $\Theta(x) = ix$. The fundamental forms of the octahedron are therefore
at the same time invariant forms of the tetrahedron, and in fact
the forms $f$ of the tetrahedron and octahedron agree with one another
 exactly, and we have $W = \Phi_1\Phi_2$,
\[
K = \tfrac{1}{2}(\Phi_1^3 + \Phi_2^3), \quad (\S.\,56).
\]
Written with the determinant 1, the two generating
 substitutions of the octahedral group read as follows:
\[
\Theta = \begin{pmatrix} \sqrt{i}, & 0 \\ 0, & \dfrac{1}{\sqrt{i}} \end{pmatrix}, \quad
\chi = \begin{pmatrix} \dfrac{1}{\sqrt{2i}}, & \dfrac{1}{\sqrt{2i}} \\[6pt] \dfrac{1}{i\sqrt{2i}}, & \dfrac{-1}{i\sqrt{2i}} \end{pmatrix},
\]
in which $\sqrt{2i} = 1 + i$ is to be set.

Under the substitution $\Theta$ the form $f(x_1, x_2)$ now \gesperrt{changes}
its \gesperrt{sign}, as formula (3) immediately shows. By the
application of $\chi$ the linear factors of $f$ are altered as
 follows:
\[
\begin{gathered}
x_1, \quad x_2, \quad x_1 + x_2, \quad x_1 - x_2, \quad x_1 + ix_2, \quad x_1 - ix_2 \\
\frac{x_1 + x_2}{\sqrt{2i}}, \quad \frac{x_1 - x_2}{i\sqrt{2i}}, \quad -i(x_1 + ix_2), \quad (x_1 - ix_2), \quad \frac{2x_1}{\sqrt{2i}}, \quad \frac{2x_2}{\sqrt{2i}},
\end{gathered}
\]
and from this it follows that $f(x_1, x_2)$ \gesperrt{remains unchanged} under
application of the substitution $\chi$, and thereby the
changes of the form $f$ under all the other octahedral substitutions
 are determined at the same time.

The \gesperrt{Hessian} covariant $W$ of $f$ remains unchanged
when $f$ is changed into $-f$, and consequently $W$ remains under the
%CONT>>


% ===== scan leaf 241, printed 220 =====
% !PRINTED 220
\weberpage{220}
\pdfbookmark[2]{§ 58. The icosahedral group}{bm68}
%<<CONT
octahedral substitutions absolutely unchanged, while $K$ again
suffers the same changes of sign as $f$.

\begin{center}
\S.\,58.

\gesperrt{The icosahedral group.}
\end{center}

Since in the icosahedral group we have only one system of conjugate
five-fold poles, we can (\S.\,53, 1.) assume in this group
the two substitutions

\weq{1}{\Theta(x) = \varepsilon x, \quad \psi(x) = \frac{-1}{x}}

where $\varepsilon$ denotes a primitive fifth root of unity.
We here take, as we are free to do, the substitution $\psi$ in the
form $-1:x$, because the formulae thereby become simpler.
The fundamental form of degree 12 belonging to the system of the
five-fold poles must, since it admits the substitutions (1), be of
the form:

\weq{2}{f(x_1,x_2) = x_1x_2(x_1^{10} + m\,x_1^5x_2^5 - x_2^{10}),}

and it remains to determine the constant
factor $m$. The two other fundamental forms are of degree 20
and degree 30. We can now determine $m$ by the theorem \S.\,54, 4.
if we can form a covariant of $f(x)$
whose degree is lower than 60, and which cannot be represented
as a product of powers of three functions of degree 12, 20, 30,
and which by the theorem mentioned must be identically
zero.

Now a covariant of degree 16 of the
form $f$ can easily be formed, if, following Bd.\ I, \S.\,60, we take the fourth polar of the
form $f(x_1,x_2)$:

\[
12\,.\,11\,.\,10\,.\,9\;P_4(x,\xi) = u_0\xi_1^4 + u_1\xi_1^3\xi_2 + u_2\xi_1^2\xi_2^2 + u_3\xi_1\xi_2^3 + u_4\xi_2^4,
\]

where

\[
u_0 = \frac{\partial^4 f}{\partial x_1^4}, \quad u_1 = 4\,\frac{\partial^4 f}{\partial x_1^3\,\partial x_2}, \quad u_2 = 6\,\frac{\partial^4 f}{\partial x_1^2\,\partial x_2^2},
\]
\[
u_3 = 4\,\frac{\partial^4 f}{\partial x_1\,\partial x_2^3}, \quad u_4 = \frac{\partial^4 f}{\partial x_2^4}.
\]

From this we obtain a covariant of degree 16 of $f$ as
%CONT>>


% ===== scan leaf 242, printed 221 =====
% !PRINTED 221
\weberpage{221}
%<<CONT
first invariant of the form biquadratic with respect to the variables $\xi_1, \xi_2$, namely (Bd.\ I, \S.\,64):

\weq{3}{u_2^2 - 3\,u_1u_3 + 12\,u_0u_4.}

But since a form of degree 16 cannot be represented as a product of forms of degrees 12, 20 and 30, this covariant must vanish identically. Now here

\[
u_0 = 11\,.\,10\,.\,9\,.\,8\,x_1^7x_2 + 6\,.\,5\,.\,4\,.\,3\,m\,x_1^2x_2^6,
\]
\[
u_1 = 4\,.\,11\,.\,10\,.\,9\,x_1^8 + 4\,.\,6\,.\,5\,.\,4\,.\,6\,m\,x_1^3x_2^5,
\]
\[
u_2 = 6^3\,.\,5^2\,m\,x_1^4x_2^4,
\]
\[
u_3 = -4\,.\,11\,.\,10\,.\,9\,x_2^8 + 4\,.\,6\,.\,5\,.\,4\,.\,6\,m\,x_2^3x_1^5,
\]
\[
u_4 = -11\,.\,10\,.\,9\,.\,8\,x_2^7x_1 + 6\,.\,5\,.\,4\,.\,3\,m\,x_2^2x_1^6,
\]

and if we form from these the covariant (3) and set the coefficient of $x_1^8x_2^8$ equal to 0, we obtain $m = \pm 11$. Both numbers are here admissible. One choice is transformed into the other by replacing $x_1$ with $-x_1$, and hence by a transformation of the group. We therefore set:

\weq{4}{f(x_1,x_2) = x_1x_2(x_1^{10} + 11\,x_1^5x_2^5 - x_2^{10}).}

The ten five-fold poles still missing are therefore obtained by solution of the equation

\[
x^{10} + 11\,x^5 - 1 = 0,
\]

from which one finds the following values:

\[
x^5 = \frac{-11 \pm 5\sqrt5}{2} = \left(\frac{-1 \pm \sqrt5}{2}\right)^5
\]

If we therefore set

\weq{5}{\begin{aligned}
\omega &= \varepsilon + \varepsilon^4 = 2\cos\frac{2\pi}{5} = \frac{-1+\sqrt5}{2} \\
\omega' &= \varepsilon^2 + \varepsilon^3 = 2\cos\frac{4\pi}{5} = \frac{-1-\sqrt5}{2},
\end{aligned}}

so that

\weq{6}{\omega^2 + \omega = 1, \quad \omega + \omega' = -1, \quad \omega\omega' = -1}

holds, then the five-fold poles other than 0 and $\infty$ are:

\weq{7}{\xi = \varepsilon^\lambda\omega, \quad \varepsilon^\lambda\omega' = \frac{-\varepsilon^\lambda}{\omega}, \quad \lambda = 0, 1, 2, 3, 4.}

Now we apply the theorem \S.\,53, 2., according to which, if $a, a'$ are two poles of a substitution $\Theta$, and $b$ is a pole conjugate to $a$, there must exist a substitution $\chi$ in the group such
%CONT>>


% ===== scan leaf 243, printed 222 =====
% !PRINTED 222
\weberpage{222}
%<<CONT
that $b = \chi(a)$, $b' = \chi(a')$ are the two poles of the substitution
$\chi^{-1}\Theta\chi$. % [sic/mathematics: with b=chi(a), theorem 53.2 requires chi Theta chi^{-1}; the source prints the inverse conjugation order and the ensuing calculation follows that printed order.]
In this we can take $\infty$ and 0 for $a$ and $a'$. and $\omega$ for $b$. % [sic: the source prints a full stop before the lower-case ``und''; preserved here rather than silently regularized.]
Then $b'$ becomes a pole $\xi$ from the sequence (7),
still to be determined more closely, and for $\chi$ we obtain the form

\[
\chi(x) = \frac{\alpha x + \beta}{\gamma x + \delta} = \begin{pmatrix} \alpha, & \beta \\ \gamma, & \delta \end{pmatrix},
\]

in which

\weq{8}{\omega = \frac{\alpha}{\gamma}, \quad \xi = \frac{\beta}{\delta}}

is to be set.
Now we can take for $\Theta$ any of the substitutions $\varepsilon^h x$,
provided only that $h$ is not divisible by 5, and thus obtain the substitution


\[
\begin{aligned}
\chi^{-1}\Theta\chi &= \begin{pmatrix} \delta, & -\beta \\ -\gamma, & \alpha \end{pmatrix}\begin{pmatrix} \varepsilon^h, & 0 \\ 0, & 1 \end{pmatrix}\begin{pmatrix} \alpha, & \beta \\ \gamma, & \delta \end{pmatrix} \\
&= \begin{pmatrix} \varepsilon^h\alpha\delta - \beta\gamma, & (\varepsilon^h-1)\beta\delta \\ -(\varepsilon^h-1)\alpha\gamma, & \alpha\delta - \varepsilon^h\beta\gamma \end{pmatrix},
\end{aligned}
\]

whose poles must be $\omega$ and $\xi$. By \S.\,51, (6), then, $\omega$ and $\xi$
must be the roots of the quadratic equation

\[
\alpha\gamma x^2 + (\alpha\delta + \beta\gamma)x + \beta\delta = 0
\]

from which

\[
\frac{\beta\delta}{\alpha\gamma} = \omega\xi, \quad \frac{\delta}{\gamma} + \frac{\beta}{\alpha} = -\omega - \xi
\]

follows. From the second of these equations there results by means of (8):

\[
(\alpha + \delta)(\alpha\delta + \beta\gamma) = 0,
\]

and since $\alpha\delta + \beta\gamma$ cannot vanish, because otherwise $\omega + \xi = 0$
would have to hold, which by (7) is not possible, we have $\alpha + \delta = 0$.
Accordingly we obtain for $\chi$, if for the sake of simplicity we assume
$\gamma = 1$:

\weq{9}{\chi = \begin{pmatrix} \omega, & \beta \\ 1, & -\omega \end{pmatrix}.}

Here $\beta$, which by (8) has the value $-\omega\xi$, is still to be
determined.

If we assume the substitution $\chi$ as known, then
we can form the whole icosahedral group.

Namely, if one lets $r$ and $s$ run through the numbers 0, 1, 2, 3, 4,
one has in this group the substitutions

\weq{10}{\Theta^r, \quad \psi\Theta^r, \quad \Theta^r\chi\Theta^s, \quad \Theta^r\chi\psi\Theta^s,}
%CONT>>


% ===== scan leaf 244, printed 223 =====
% !PRINTED 223
\weberpage{223}
%<<CONT
There are exactly 60 of them. That they are all different
from one another is seen when one represents them in the form:

\weq{11}{\Theta^r = \begin{pmatrix} \varepsilon^r, & 0 \\ 0, & 1 \end{pmatrix},}

\weq{12}{\psi\Theta^r = \begin{pmatrix} 0, & 1 \\ -\varepsilon^r, & 0 \end{pmatrix},}

\weq{13}{\Theta^r\chi\Theta^s = \begin{pmatrix} \varepsilon^r\omega, & \varepsilon^{r-s}\beta \\ 1, & -\varepsilon^{-s}\omega \end{pmatrix},}

\weq{14}{\begin{aligned}
\Theta^r\chi\psi\Theta^s &= \begin{pmatrix} -\varepsilon^{r+s}\beta, & \varepsilon^r\omega \\ \varepsilon^s\omega, & 1 \end{pmatrix} \\
&= \begin{pmatrix} -\varepsilon^r\beta\omega^{-1}, & \varepsilon^{r-s} \\ 1, & \varepsilon^{-s}\omega^{-1} \end{pmatrix},
\end{aligned}}

of which, since $-\omega^2$ is no power of $\varepsilon$, evidently no
two are equal to one another.

The icosahedral group has, as we have seen, 12 five-fold
poles. Each pair of these poles is the pair of poles of four substitutions
of degree 5, which together with the identity form a
five-term cycle. Consequently there are in the group
24 elements of degree 5. Likewise there are 20 three-fold poles,
which lead to 20 substitutions of the third order, and 30 two-fold
poles, which two by two are the poles of one substitution
of the second order each, so that there are 15 substitutions of the second
order. Together with the identity this gives

\[
24 + 20 + 15 + 1 = 60.
\]

The determination of $\beta$ in the substitution $\chi$ now results
from a consideration of the degrees of (13) and (14). To this end
we first form, for an arbitrary linear substitution $S = \begin{pmatrix} a, & b \\ c, & d \end{pmatrix}$,
the second and third repetition:

\[
\begin{aligned}
S^2 &= \begin{pmatrix} a^2 + bc, & b(a+d) \\ c(a+d), & bc + d^2 \end{pmatrix} \\
S^3 &= \begin{pmatrix} a^3 + 2abc + bcd, & b(a^2+bc+ad+d^2) \\ c(a^2+bc+ad+d^2), & d^3+abc+2bcd \end{pmatrix},
\end{aligned}
\]

and if we leave aside the case that $b$ or $c$ vanishes,
which does not come into consideration here, we obtain the
necessary and sufficient condition:
%CONT>>


% ===== scan leaf 245, printed 224 =====
% !PRINTED 224
\weberpage{224}

%<<CONT
for a substitution of the second order
\weq{15}{a + d = 0,}
and for a substitution of the third order
\weq{16}{a^2 + bc + ad + d^2 = 0.}

Now the substitutions (11), apart from the identical one contained among them, are of the fifth order. The five substitutions (12) are of the second order. Of the substitutions (13), by (15), the five contained in the form $\Theta^r\chi\Theta^{-r}$ (and only these) are of the second order, and consequently five more of the substitutions (14) must also be of the second order. But by (15) this is only possible \gesperrt{if $\beta$ is a power of $\varepsilon$}.

The substitutions of the third order must now distribute themselves over the forms (13) and (14). For these the conditions resulting from (16) are:

that (13) be of the third order
\weq{17}{\omega^2(\varepsilon^{r+s} + \varepsilon^{-(r+s)} - 1) = -\beta,}
and that (14) be of the third order
\weq{18}{\omega^2 = \beta - \varepsilon^{-(r+s)} - \beta^2\varepsilon^{r+s}.}

But if either of the two conditions (17) or (18) is to be satisfiable, we must have:
\[
\beta = 1
\]
For $\omega^2$ is real, and the whole left side of (17) is therefore likewise real; consequently $\beta$ must be real, and, since it is a power of $\varepsilon$, must be equal to 1.

But if (18) is to be satisfied, the left side must not change when one passes to the conjugate imaginary quantities; that is, we must have
\[
\beta - \varepsilon^{-(r+s)} - \beta^2\varepsilon^{r+s} = \beta^{-1} - \varepsilon^{r+s} - \beta^{-2}\varepsilon^{-(r+s)}
\]
or
\[
(\beta - \beta^{-1})(1 - \beta^{-1}\varepsilon^{-(r+s)} - \beta\varepsilon^{r+s}) = 0
\]
and this, because $\beta$ is a power of $\varepsilon$, is possible only if $\beta = \beta^{-1}$, and hence $\beta = 1$.

Accordingly five of the substitutions (14) are of the second order, namely $\Theta^r\chi\psi\Theta^{-r}$, and it follows that (13) is of the third order if
\[
\omega^2(\varepsilon^{r+s} + \varepsilon^{-(r+s)} - 1) = -1,
\]
%CONT>>


% ===== scan leaf 246, printed 225 =====
% !PRINTED 225
\weberpage{225}
%<<CONT
or, if one multiplies by $\omega'^2$, by (5) and (6)
\[
\begin{aligned}
\varepsilon^{r+s} + \varepsilon^{-(r+s)} - 1 &= -\varepsilon - \varepsilon^{-1} - 2 \\
\varepsilon + \varepsilon^{-1} + \varepsilon^{r+s} + \varepsilon^{-(r+s)} + 1 &= 0,
\end{aligned}
\]
a condition which is fulfilled if and only if
$r + s \equiv \pm 2 \pmod{5}$.

If (14) is to be of the third order, then by
(18) and (5) we must have
\[
\begin{aligned}
\varepsilon^2 + \varepsilon^{-2} + 2 &= 1 - \varepsilon^{-(r+s)} - \varepsilon^{r+s} \\
\varepsilon^2 + \varepsilon^{-2} + \varepsilon^{r+s} + \varepsilon^{-(r+s)} + 1 &= 0
\end{aligned}
\]
and this condition is satisfied if and only if
$r + s \equiv \pm 1 \pmod{5}$.

Accordingly we obtain, as must be the case, in (13) and (14)
exactly 10 substitutions of the second and 20 substitutions of the third
order, and the remaining cases are therefore of the fifth order.

If we summarise the result of this consideration, we
have:

The substitution $\chi$ must have the expression:
\weq{19}{\chi = \begin{pmatrix} \omega, & 1 \\ 1, & -\omega \end{pmatrix},}
and among the 60 substitutions (10):
\weq{20}{\Theta^r, \; \psi\Theta^r, \; \Theta^r\chi\Theta^s, \; \Theta^r\chi\psi\Theta^s}
there occur, besides the identity:

15 \gesperrt{substitutions} of \gesperrt{degree} $2$:

$\psi\Theta^r, \; \Theta^r\chi\Theta^{-r}, \; \Theta^r\chi\psi\Theta^{-r}, \; r = 0, 1, 2, 3, 4,$
\weq{21}{
\begin{aligned}
&\text{20 }\text{\gesperrt{substitutions}}\text{ of }\text{\gesperrt{degree} }3\text{:} \\
&\Theta^r\chi\Theta^s, \quad r+s \equiv \pm 2 \pmod{5} \\
&\Theta^r\chi\psi\Theta^s, \quad r+s \equiv \pm 1 \pmod{5}, \\
&\text{24 }\text{\gesperrt{substitutions}}\text{ of }\text{\gesperrt{degree} }5\text{:} \\
&\Theta^r, \; \Theta^r\chi\Theta^s, \quad r+s \equiv \pm 1 \pmod{5} \\
&\Theta^r\chi\psi\Theta^s, \quad r+s \equiv \pm 2 \pmod{5}.
\end{aligned}
}
In explicit form one obtains for the substitutions (20)
the expression:
\weq{22}{
\begin{pmatrix} \varepsilon^r, & 0 \\ 0, & 1 \end{pmatrix}, \;
\begin{pmatrix} 0, & \varepsilon^r \\ -1, & 0 \end{pmatrix}, \;
\begin{pmatrix} \varepsilon^r\omega, & \varepsilon^{r-s} \\ 1, & -\varepsilon^{-s}\omega \end{pmatrix}, \;
\begin{pmatrix} -\varepsilon^{r+s}, & \varepsilon^r\omega \\ \varepsilon^s\omega, & 1 \end{pmatrix}
}
or
\[
\varepsilon^r x, \quad \frac{-\varepsilon^r}{x}, \quad \frac{\varepsilon^r\omega x + \varepsilon^{r-s}}{x - \varepsilon^{-s}\omega}, \quad \frac{-\varepsilon^{r+s} x + \varepsilon^r\omega}{\varepsilon^s\omega x + 1}.
\]


% ===== scan leaf 247, printed 226 =====
% !PRINTED 226
\weberpage{226}
In order, then, finally to establish the existence of the icosahedral group, it remains to show that the totality of the substitutions (20) forms a group. But this follows from the laws of composition now to be derived.

First of all, one obtains very simply from the meaning of $\Theta$, $\psi$, $\chi$:
\[
\Theta = \begin{pmatrix} \varepsilon, & 0 \\ 0, & 1 \end{pmatrix}, \quad
\psi = \begin{pmatrix} 0, & 1 \\ -1, & 0 \end{pmatrix}, \quad
\chi = \begin{pmatrix} \omega, & 1 \\ 1, & -\omega \end{pmatrix}
\]
\weq{23}{\psi\Theta = \Theta^{-1}\psi, \quad \chi\psi = \psi\chi.}

Further, for an arbitrary exponent $\lambda$ one obtains:
\[
\chi\Theta^\lambda\chi = \begin{pmatrix} \varepsilon^\lambda\omega^2+1, & (\varepsilon^\lambda-1)\omega \\ (\varepsilon^\lambda-1)\omega, & \varepsilon^\lambda+\omega^2 \end{pmatrix},
\]
or, if according to (5) and (6) one sets
\weq{24}{\omega = \varepsilon+\varepsilon^{-1}, \quad \omega' = \varepsilon^2+\varepsilon^{-2}, \quad \omega\omega' = -1}
then:
\[
\chi\Theta^\lambda\chi = \begin{pmatrix} \varepsilon^\lambda\omega-\omega', & \varepsilon^\lambda-1 \\ \varepsilon^\lambda-1, & -\varepsilon^\lambda\omega'+\omega \end{pmatrix}.
\]

If we first set $\lambda=1$ in this, it follows by (24), (13), (14):
\[
\chi\Theta\chi = \begin{pmatrix} 1-\varepsilon^3, & \varepsilon-1 \\ \varepsilon-1, & \varepsilon-\varepsilon^3 \end{pmatrix}
= \begin{pmatrix} \varepsilon^3(\varepsilon+1), & 1 \\ 1, & -\varepsilon(\varepsilon+1) \end{pmatrix}
\]
\[
= \begin{pmatrix} \varepsilon\omega', & 1 \\ 1, & -\varepsilon^{-1}\omega' \end{pmatrix}
= \begin{pmatrix} -\varepsilon^2, & \varepsilon\omega \\ \varepsilon\omega, & 1 \end{pmatrix}
= \Theta\chi\psi\Theta,
\]
from which, on taking inverses on both sides,
\[
\chi\Theta^{-1}\chi = \Theta^{-1}\chi\psi\Theta^{-1}
\]
follows. If, on the other hand, one sets $\lambda=2$, it follows:
\[
\chi\Theta^2\chi = \begin{pmatrix} \varepsilon-\varepsilon^2, & \varepsilon^2-1 \\ \varepsilon^2-1, & \varepsilon-1 \end{pmatrix}
= \begin{pmatrix} \varepsilon^2(\varepsilon^2+1), & 1 \\ 1, & -\varepsilon(\varepsilon^2+1) \end{pmatrix}
\]
\[
= \begin{pmatrix} \varepsilon^3\omega, & 1 \\ 1, & -\varepsilon^2\omega \end{pmatrix}
= \Theta^{-2}\chi\Theta^{-2},
\]
in which again $\Theta$ may be replaced by $\Theta^{-1}$. One therefore has the composition formulae:

\weq{25}{\chi\Theta^{\pm1}\chi = \Theta^{\pm1}\chi\psi\Theta^{\pm1}, \quad \chi\Theta^{\pm2}\chi = \Theta^{\mp2}\chi\Theta^{\mp2}.}

By application of the formulae (23) and (25) one can now compose any two of the substitutions (20), and one always arrives again at a substitution of the form (20), whereby
%CONT>>


% ===== scan leaf 248, printed 227 =====
% !PRINTED 227
\weberpage{227}
\pdfbookmark[2]{§ 59. The divisors of the icosahedral group}{bm69}
%<<CONT
the group nature has been proved and the icosahedral group has been formed.

From the first formula (25) there results further
\[
\psi = \chi\Theta^{\mp1}\chi\Theta^{\pm1}\chi\Theta^{\mp1},
\]
from which one concludes that $\psi$ can be derived from $\chi$ and $\Theta$, and hence that $\chi$ and $\Theta$ alone can already be regarded as \gesperrt{generating substitutions of the icosahedral group}.

If we wish to represent the substitution $\chi$ with the determinant 1, we observe the relation
\[
-\omega^2-1 = \omega-2 = (\varepsilon^2-\varepsilon^{-2})^2,
\]
and obtain, if we further substitute $\varepsilon+\varepsilon^{-1}$ for $\omega$:
\weq{26}{\chi = \begin{pmatrix} \dfrac{1}{\varepsilon-\varepsilon^{-1}}, & \dfrac{1}{\varepsilon^2-\varepsilon^{-2}} \\[2mm] \dfrac{1}{\varepsilon^2-\varepsilon^{-2}}, & \dfrac{-1}{\varepsilon-\varepsilon^{-1}} \end{pmatrix}.}

\S.\,59.

\gesperrt{The Divisors of the Icosahedral Group.}

The divisors of the icosahedral group must be sought among the lower polyhedral groups. Among these are found first the cyclic groups $C_n$, each of which consists of the powers of one of the icosahedral substitutions. The enumeration of the number of these groups follows at once from the arrangement of the substitutions according to their degrees, as we have given it in the preceding section; namely:
\[
\begin{aligned}
&15 \text{ groups } C_2, \\
&10 \quad \text{„} \quad C_3, \\
&\phantom{0}6 \quad \text{„} \quad C_5.
\end{aligned}
\]
There are further contained in the icosahedral group dihedral groups $D_2$, $D_3$, $D_5$, as whose representatives we set up the following:
\[
\begin{aligned}
D_2 &: 1, \psi, \chi, \psi\chi, \\
D_3 &: 1, \chi\Theta^2, \Theta^{-2}\chi, \psi\chi, \psi\Theta^2, \Theta^{-2}\psi\chi\Theta^2, \\
D_5 &: \Theta^r, \psi\Theta^r, \quad r=0,1,2,3,4,5.
\end{aligned}
\]
% [sic: Weber prints r=0,1,2,3,4,5; since \Theta^5=1, the last value duplicates r=0.]


% ===== scan leaf 249, printed 228 =====
% !PRINTED 228
\weberpage{228}
We obtain the number of dihedral groups from the fact that a dihedral group is completely determined by the cyclic group contained in it (\S.\,55). For the four-group, however, one must further observe that the same group is obtained whether one starts from $\psi$, from $\chi$, or from $\psi\chi$, so the number obtained directly must be divided by 3. Thus there are 5 groups $D_2$, 10 groups $D_3$, and 6 groups $D_5$.

Of particular importance, however, are the tetrahedral groups contained in the icosahedral group. We shall determine one of them, which we denote by $Q$.

The tetrahedral group must (by \S.\,56) contain a four-group $D_2$. We start from such a group $D_2$ and choose for it
\[
1, \psi, \chi, \psi\chi.
\]

Now further, a substitution of the third order must occur in $Q$. We may assume it in one of the two forms $\Theta^r\chi\Theta^s$ or $\Theta^r\chi\psi\Theta^s$. Since both assumptions lead to the same result, we choose as the substitution of the third order
\[
\varphi = \Theta^r\chi\Theta^s, \quad r+s\equiv\pm2 \ (\mathrm{mod}\ 5) \quad [\S.\,58, (21)].
\]

For neither of the numbers $r, s$, however, is congruence to zero ($\equiv0$) possible, since otherwise either $\chi\varphi$ or $\varphi\chi$ would be a power of $\Theta$, and hence of degree $5$, whereas no element of degree $5$ can occur in $Q$.

We further form, by \S.\,58, (23), (25):
\[
\begin{aligned}
\chi\varphi &= \Theta^r\chi\psi\Theta^{r+s}, \quad r\equiv\pm1 \\
\chi\varphi &= \Theta^{-r}\chi\Theta^{-r+s}, \quad r\equiv\pm2,
\end{aligned}
\]
whereby, for the reason cited above, $r\equiv-s$ is excluded in the first case and $r\equiv s$ in the second, and in the first case $2r+s\equiv\pm1$, in the second $-2r+s\equiv\pm2$ is required. Accordingly the four possible cases
\[
r\equiv\pm1, \quad s\equiv\pm2, \quad r\equiv\pm2, \quad s\equiv\pm1 \ (\mathrm{mod}\ 5)
\]
remain.

If one of the four substitutions of the third order so determined occurs in $Q$, then the three others occur as well. For if we set
\weq{1}{\varphi = \Theta\chi\Theta^2,}
then by \S.\,58, (23), (25), it follows that:
\weq{2}{\varphi^{-1} = \Theta^{-2}\chi\Theta^{-1}, \quad \varphi\chi = \Theta^{-1}\chi\Theta^{-2}, \quad \chi\varphi^{-1} = \Theta^2\chi\Theta.}


% ===== scan leaf 250, printed 229 =====
% !PRINTED 229
\weberpage{229}
If, in addition, one also forms $\varphi\psi$, $\varphi^{-1}\psi$, $\varphi\chi\psi$, $\chi\varphi^{-1}\psi$, then it follows that:
\[
\begin{aligned}
\varphi\psi &= \Theta\chi\psi\Theta^{-2}, \quad \varphi^{-1}\psi = \Theta^{-2}\chi\psi\Theta, \quad \varphi\chi\psi = \Theta^{-1}\chi\psi\Theta^2, \\
&\chi\varphi^{-1}\psi = \Theta^2\chi\psi\Theta^{-1},
\end{aligned}
\]
from which one sees that one would come to no other result if one wished to assume the substitution of the third order from which one starts in the second form $\Theta^r\chi\psi\Theta^s$. The whole group $Q$ is thus completely determined by the assumed four-group $D_2$, and one obtains it in the form
\weq{3}{
\begin{array}{ll}
1 & \psi \\
\varphi=\Theta\chi\Theta^2, & \varphi\psi=\Theta\chi\psi\Theta^{-2}, \\
\varphi^{-1}=\Theta^{-2}\chi\Theta^{-1}, & \varphi^{-1}\psi=\Theta^{-2}\chi\psi\Theta, \\
\chi & \chi\psi \\
\varphi\chi=\Theta^{-1}\chi\Theta^{-2}, & \varphi\chi\psi=\Theta^{-1}\chi\psi\Theta^2, \\
\varphi^{-1}\chi=\Theta^2\chi\psi\Theta^{-1}, & \varphi^{-1}\chi\psi=\Theta^2\chi\Theta.
\end{array}
}

One can derive this group from the substitutions $\varphi, \psi, \chi$ as the generators and put it into the form:
\weq{4}{Q = \varphi^r\chi^s\psi^t, \qquad \begin{gathered}r=0,1,2,\\ s=0,1; \quad t=0,1.\end{gathered}}

That a tetrahedral group is really represented in this way follows from the compositions, which one easily finds from (1) by means of \S.\,58, (23), (25):
\weq{5}{
\begin{aligned}
\psi\varphi &= \varphi\chi\psi, \quad \chi\varphi = \varphi\psi, \\
\psi\varphi^2 &= \chi\varphi^2\psi, \quad \chi\varphi^2 = \varphi^2\chi\psi. \quad \psi\chi = \chi\psi. % [sic: Weber prints a full stop between the two relations, unlike the commas separating the parallel relations; the printed punctuation is preserved.]
\end{aligned}
}

Since five four-groups $D_2$ are contained in the tetrahedral group, the icosahedral group has five tetrahedral groups as divisors. % [sic: Weber prints Tetraëdergruppe; mathematically Ikosaëdergruppe is intended.]
These we can derive from $Q$ by transformation by means of the powers of $\Theta$, and we obtain them in the form
\[
\Theta^{-r}Q\Theta^r \quad r=0,1,2,3,4.
\]

Apart from the identity, these groups have no substitution common to all of them. For the two groups $Q$ and $\Theta^{-1}Q\Theta$ have in common with one another, apart from the identity, only the two elements
\[
\Theta^{-2}\chi\Theta^{-1}, \quad \Theta\chi\Theta^2
\]
and of these neither occurs in $\Theta Q\Theta^{-1}$.

From this it now follows, by Theorem 2. in \S.\,28, that the icosahedral group is isomorphic to a permutation group
%CONT>>


% ===== scan leaf 251, printed 230 =====
% !PRINTED 230
\weberpage{230}
\pdfbookmark[2]{§ 60. The fundamental forms of the icosahedral group}{bm70}
%<<CONT
of order 60 on five letters. This permutation group is obtained by right-multiplying the right cosets
\[
Q, \quad Q\Theta, \quad Q\Theta^2, \quad Q\Theta^3, \quad Q\Theta^4
\]
by all the elements $\sigma$ of the icosahedral group, that is, by forming
\[
Q\sigma, \quad Q\Theta\sigma, \quad Q\Theta^2\sigma, \quad Q\Theta^3\sigma, \quad Q\Theta^4\sigma
\]
and examining the resulting permutation of these right cosets. For the generating substitutions of the icosahedral group $\sigma = \Theta, \psi, \chi$ we thus obtain the permutations
\[
(1, \ \Theta, \ \Theta^2, \ \Theta^3, \ \Theta^4) \qquad \sigma = \Theta
\]
\[
\begin{pmatrix} Q, & Q\Theta, & Q\Theta^2, & Q\Theta^3, & Q\Theta^4 \\ Q, & Q\Theta^4, & Q\Theta^3, & Q\Theta^2, & Q\Theta \end{pmatrix} \qquad \sigma = \psi
\]
\[
\begin{pmatrix} Q, & Q\Theta, & Q\Theta^2, & Q\Theta^3, & Q\Theta^4 \\ Q, & Q\Theta^3, & Q\Theta^4, & Q\Theta, & Q\Theta^2 \end{pmatrix} \qquad \sigma = \chi.
\]

The last is found from the formulae (1) and (2), according to which
\[
\Theta\chi = \varphi\Theta^3, \quad \Theta^2\chi = \chi\varphi^2\Theta^4, \quad \Theta^3\chi = \varphi^2\Theta, \quad \Theta^4\chi = \varphi\chi\Theta^2.
\]

One sees that these permutations all belong to the first kind, and that therefore the permutation group in question can be no other than the alternating group on five letters (Bd. I, \S.\,177). From this it also follows that the icosahedral group is simple, and consequently coincides with the group which we already designated provisionally as the icosahedral group in \S.\,34.

\begin{center}
\S.\,60.\\[2pt]
\gesperrt{The Fundamental Forms of the Icosahedral Group.}
\end{center}

We have found above one of the three fundamental forms of the icosahedral group, $f$:
\weq{1}{f = x_1 x_2 (x_1^{10} + 11 x_1^5 x_2^5 - x_2^{10}).}
We therefore still lack two of these forms, one of the 20th and one of the 30th degree.

The fundamental form of degree 20 is obtained as the \gesperrt{Hessian} covariant of $f$. We set it
\weq{2}{H = \frac{1}{121}\,[f''(x_1,x_1)\, f''(x_2,x_2) - f''(x_1,x_2)^2],}
%CONT>>


% ===== scan leaf 252, printed 231 =====
% !PRINTED 231
\weberpage{231}
%<<CONT
and find by a simple calculation:
\weq{3}{H = -(x_1^{20} + x_2^{20}) + 228\,(x_1^{15}x_2^5 - x_2^{15}x_1^5) - 494\,x_1^{10}x_2^{10}.}
The fundamental form of degree 30 we can define as the functional determinant of $H$ and $f$. If we set
\weq{4}{T = \frac{1}{20}\,[f'(x_1)\,H'(x_2) - f'(x_2)\,H'(x_1)],}
then one finds
\weq{5}{T = (x_1^{30} + x_2^{30}) + 522\,(x_1^{25}x_2^5 - x_1^5x_2^{25}) - 10005\,(x_1^{20}x_2^{10} + x_2^{20}x_1^{10}).}
There is also an identity among these three forms. To find it, we set
\[
\lambda = x_1^{10} - x_2^{10}, \quad \mu = x_1^5 x_2^5,
\]
and obtain, if we take out the factor $x_1^{10} + x_2^{10}$ from $T$ and then form $T^2$:
\[
f^5 = \mu\,(\lambda + 11\mu)^5,
\]
\[
H = -\lambda^2 + 228\,\lambda\mu - 496\,\mu^2,
\]
\[
T^2 = (\lambda^2 + 4\mu^2)\,(\lambda^2 + 522\,\lambda\mu - 10004\,\mu^2)^2.
\]
From this one obtains by numerical computation
\weq{6}{T^2 + H^3 = 1728\,f^5,}
which is the relation sought.

Since the icosahedral group, as we have seen, is simple, it can by \S.\,40 have no relative invariants. From this it follows that the icosahedral forms $f$, $H$, $T$ remain absolutely unchanged when one subjects them to the binary icosahedral substitutions represented with determinant $1$. But the same can also be demonstrated directly in the following way, without making use of those general theorems.

The relation (6) can be presented in the form:
\weq{7}{\frac{T^2}{f^5} + \frac{H^3}{f^5} = 1728,}
and the two quotients $T^2 : f^5$, $H^3 : f^5$ we can regard as functions of the single variable $x = x_1 : x_2$. If we apply to this variable a substitution of the icosahedral group, these two quotients change only by constant factors; that is, if we
\weq{8}{\frac{T^2}{f^5} = \Phi(x), \quad \frac{H^3}{f^5} = \Psi(x), \quad y = \frac{\alpha x + \beta}{\gamma x + \delta}}
%CONT>>


% ===== scan leaf 253, printed 232 =====
% !PRINTED 232
\weberpage{232}
\pdfbookmark[2]{§ 61. The invariants of the icosahedron}{bm71}
%<<CONT
and denote by $h$, $k$ two constants, then
\[
\Phi(y) = h\,\Phi(x), \quad \Psi(y) = k\,\Psi(x),
\]
provided that $\begin{pmatrix} \alpha, & \beta \\ \gamma, & \delta \end{pmatrix}$ is one of the icosahedral substitutions. But now, since $y$ as well as $x$ is an independent variable, by (7) the identity
\[
\Phi(x) + \Psi(x) = h\,\Phi(x) + k\,\Psi(x) = 1728,
\]
holds, and since $\Phi$ and $\Psi$ do not stand in a constant ratio, we must have $h = k = 1$; that is, the quotients $\Phi$ and $\Psi$ remain absolutely unchanged under the icosahedral substitutions.

From this it follows further \gesperrt{that $f$, $H$, $T$ remain absolutely unchanged under the homogeneous icosahedral substitutions with determinant $1$.}

For if $f$ passes under such a substitution into $cf$, then $H$ and $T$ pass into $c^2H$ and $c^3T$ [by (2) and (4)], and the quotients $\Phi$ and $\Psi$ into $c\Phi$ and $c\Psi$. Since, on the other hand, $\Phi$ and $\Psi$ remain unchanged, $c = 1$.

If we set $\Psi(x) = z$, then $\Phi(x) = 1728 - z$, and from (8) we obtain the two equations:
\weq{9}{H^3 - z f^5 = 0,}
\weq{10}{T^2 - (1728 - z)\,f^5 = 0,}
of which, on account of the identity (6), the one follows from the other, so that they are really only two different forms of one and the same equation. If we now regard $z$ in it as given, we have an equation of degree 60 for $x$, which is called the \gesperrt{icosahedral equation}.

We shall return to the properties of this equation and to its relation to the general equation of degree 5 in a later section.

\begin{center}
\S.\,61.\\[2pt]
\gesperrt{The Invariants of the Icosahedron.}
\end{center}

By the fundamental forms of the polyhedral groups with which we have so far become acquainted, the region of the invariants of the groups in question is, as we can now prove, exhausted. We
%CONT>>


% ===== scan leaf 254, printed 233 =====
% !PRINTED 233
\weberpage{233}
%<<CONT
restrict ourselves in this proof to the consideration of the icosahedral group, since for the other groups entirely similar inferences are to be made, which we may leave to the reader. The circumstance that in the tetrahedral and octahedral groups relative invariants occur alongside the absolute ones, while the icosahedral group, as a simple group, has only absolute invariants, is of no essential influence here.

We prove, then, that all invariants of the icosahedral group can be represented as integral rational functions of the three forms $f$, $H$, $T$.

To this we are led by the following chain of inferences:

1. \gesperrt{No two of the icosahedral forms $f$, $H$, $T$ have a common divisor.}

This follows immediately from the definition of these functions, according to which the equations $f = 0$, $H = 0$, $T = 0$ furnish the five-fold, three-fold and two-fold poles, which are all different from one another.

2. \gesperrt{A double root of the icosahedral equation}
\weq{1}{H^3 - z f^5 = 0}
\gesperrt{is necessarily a root of $f$, $H$ or $T$, and can therefore occur only for one of the values $z = 0, \infty, 1728$.}

For if (1) has a double root, then together with the function the first derivative, or, if one employs the homogeneous form, by \gesperrt{Euler}'s theory [Bd. I, \S.\,17, (5)] the two derivatives with respect to $x_1$ and $x_2$, must vanish at the same time, hence
\[
3H^2 H'(x_1) - 5z f^4 f'(x_1) = 0,
\]
\[
3H^2 H'(x_2) - 5z f^4 f'(x_2) = 0;
\]
consequently, since by 1. $H$ and $f$ do not vanish at the same time, only three possibilities remain: either $H = 0$, $z = 0$, or $f = 0$, $z = \infty$, or finally
\[
H'(x_1)\,f'(x_2) - H'(x_2)\,f'(x_1) = 0,
\]
i.e. $T = 0$, $z = 1728$ [\S.\,60, (4)]. This also yields the evident theorem:
%CONT>>


% ===== scan leaf 255, printed 234 =====
% !PRINTED 234
\weberpage{234}
\pdfbookmark[2]{§ 62. Polyhedral groups of the second kind. Crystallographic groups}{bm72}
%<<CONT

\gesperrt{3. If} $J(x_1, x_2)$ \gesperrt{is any invariant form of the icosahedral group, and} $\xi$ \gesperrt{one of its roots, so that} $J(\xi, 1) = 0$, \gesperrt{then all those quantities are also roots of} $J$ \gesperrt{which arise from} $\xi$ \gesperrt{through the fractional icosahedral substitutions.}

If therefore $J$ has a divisor in common with one of the functions $f$, $H$, $T$, then $J$ is divisible by the form in question.

We begin by factoring the highest possible powers of the three fundamental forms $f$, $H$, $T$ out of $J$, so that $J$ is coprime to these functions. If then $\xi$ is a root of $J$, we can determine $\eta$ so that $\xi$ is a root of the form
\[
\Theta = H^3 - \eta f^5
\]
and $\xi$ is then, by 2., a simple root of $\Theta$. Then, by 3., all the remaining roots of $\Theta$ must also at the same time be roots of $J$, i.e. $J$ is divisible by $\Theta$. The quotient is again an invariant form of the icosahedron, and the argument can be repeated. We thus arrive at the theorem:

\gesperrt{4. Every invariant of the icosahedron can be represented in the form:}
\weq{2}{J(x_1,x_2) = C f^\alpha H^\beta T^\gamma F(H^3, f^5),}
\gesperrt{wherein $\alpha$, $\beta$, $\gamma$ denote non-negative integers, $C$ a constant, and $F$ an integral homogeneous function.}

\begin{center}
\S.\,62.\\
\gesperrt{Polyhedral groups of the second kind. Crystallographic groups.}
\end{center}

We have already, in \S.\,51, called attention to groups of linear ternary substitutions which contain, besides the properly orthogonal substitutions, also improperly orthogonal ones, and which we call \gesperrt{groups of the second kind}.

We now wish to designate the finite ones among them as \gesperrt{polyhedral groups} of the second kind, or also as \gesperrt{extended polyhedral groups}, and the substitutions contained in them by the
%CONT>>


% ===== scan leaf 256, printed 235 =====
% !PRINTED 235
\weberpage{235}

%<<CONT
\gesperrt{determinant} $-1$ \gesperrt{as substitutions of the second kind}.

According to the results of \S.\,50, these groups are isomorphic with the groups of binary linear substitutions with determinant $+1$ and $-1$, whereas for the fractional substitutions, which led us to the polyhedral groups, the two kinds are not distinguishable.

The finite groups of the second kind are, apart from their general group-theoretic interest, important on account of their application in crystallography. Their complete determination no longer presents any essential difficulty. By substitutions we understand simply binary linear substitutions with determinant $\pm 1$, and we recall that two such substitutions whose coefficients differ only in their common sign, such as
\[
\begin{pmatrix} \alpha, & \beta \\ \gamma, & \delta \end{pmatrix}, \quad
\begin{pmatrix} -\alpha, & -\beta \\ -\gamma, & -\delta \end{pmatrix},
\]
are regarded as identical (\S.\,49).

Since a group of linear substitutions contains the identical substitution, which is of determinant $+1$, in every group $Q$ of the second kind there must occur, besides the improper substitutions, proper ones as well, and these proper substitutions form by themselves a group $P$. If $\varphi$ is any substitution of the second kind occurring in $Q$, then every composition of $\varphi$ with a substitution of the second kind from $Q$ occurs in $P$, and hence we can always decompose $Q$ thus:
\weq{1}{Q = P + P\varphi.}
The group $P$ must be one of the polyhedral groups considered earlier. We must have $\varphi^{-1}P\varphi = P$, and $P$ is therefore a \gesperrt{normal divisor} of $Q$. Conversely, if $P$ is a polyhedral group and $\varphi$ a substitution of the second kind that satisfies the condition $\varphi^{-1}P\varphi = P$, then $Q = P + P\varphi$ is a \gesperrt{group of the second kind}.

Here too we regard two groups that arise from one another by transformation as not essentially different. In order to form all $Q$, we have to go through the various cases.

I. If $P$ is the unit group, which we denote as the cyclic group of the first degree by $C_1$, so that $P$ consists only of the
%CONT>>


% ===== scan leaf 257, printed 236 =====
% !PRINTED 236
\weberpage{236}

%<<CONT
identical substitution, then we must have $\varphi^2 = 1$, and this is also sufficient. Here we can, by transformation, bring $\varphi$ to the form $\begin{pmatrix} \alpha, & 0 \\ 0, & -\alpha^{-1} \end{pmatrix}$, and obtain as condition $\alpha^2 = \pm 1$, and accordingly two forms of the substitution $\varphi$:
\[
\varphi' = \begin{pmatrix} i, & 0 \\ 0, & i \end{pmatrix}, \quad
\varphi'' = \begin{pmatrix} 1, & 0 \\ 0, & -1 \end{pmatrix},
\]
of which the first is called the \gesperrt{inversion}, the second the \gesperrt{reflection}. If one carries them over, according to \S.\,50, (1), to a rectangular coordinate transformation, then $\varphi'$ signifies the simultaneous change of sign of all three coordinates, $\varphi''$ the replacement of $x$ by $-x$. From this there result the two groups of the second kind:
\weq{2}{C_1' = \begin{pmatrix} i^h, & 0 \\ 0, & i^h \end{pmatrix}, \quad C_1'' = \begin{pmatrix} 1, & 0 \\ 0, & (-1)^h \end{pmatrix}, \quad h = 0, 1.}

II. Let $P$ be a cyclic group $C_n$, $n > 1$.
If we set
\weq{3}{c = \begin{pmatrix} \varepsilon, & 0 \\ 0, & \varepsilon^{-1} \end{pmatrix}, \quad \varepsilon = e^{\frac{\pi i}{n}},}
then by \S.\,55:
\weq{4}{C_n = 1, c, c^2, \ldots, c^{n-1},}
in which
\[
c^s = \begin{pmatrix} \varepsilon^s, & 0 \\ 0, & \varepsilon^{-s} \end{pmatrix};
\]
if now
\weq{5}{\varphi = \begin{pmatrix} \alpha, & \beta \\ \gamma, & \delta \end{pmatrix}, \quad \alpha\delta - \beta\gamma = -1}
holds, then we first form
\weq{6}{\varphi^{-1} c \varphi = \begin{pmatrix} \alpha\delta\varepsilon - \beta\gamma\varepsilon^{-1}, & \beta\delta(\varepsilon - \varepsilon^{-1}) \\ -\alpha\gamma(\varepsilon - \varepsilon^{-1}), & \alpha\delta\varepsilon^{-1} - \beta\gamma\varepsilon \end{pmatrix},}
\weq{7}{\varphi^2 = \begin{pmatrix} \alpha^2 + \beta\gamma, & (\alpha + \delta)\beta \\ (\alpha + \delta)\gamma, & \beta\gamma + \delta^2 \end{pmatrix}.}
These two substitutions must occur in the sequence (4), and hence we must have $\beta\delta = 0$, $\alpha\gamma = 0$; thus either $\beta$ and $\gamma$ are both zero, or $\alpha$ is zero together with $\delta = 0$. If $\beta = \gamma = 0$, it follows from (7) that $\alpha$ is of the form $\varepsilon^{1/2\,r}$, and hence
\[
\varphi = \begin{pmatrix} e^{\frac{\pi i}{2n}\,r}, & 0 \\ 0, & -e^{-\frac{\pi i}{2n}\,r} \end{pmatrix}
\]
%CONT>>


% ===== scan leaf 258, printed 237 =====
% !PRINTED 237
\weberpage{237}

%<<CONT
must hold, in which $r$ denotes an integer. According as $r$ is taken odd or even, we obtain two different groups, which cannot be reduced to one another by transformation:
\weq{1}{C_n' = \begin{pmatrix} e^{\frac{\pi i h}{2n}}, & 0 \\ 0, & (-1)^h e^{-\frac{\pi i h}{2n}} \end{pmatrix} \quad h = 0, 1, \ldots, 2n-1,}
\weq{2}{C_n'' = \begin{pmatrix} e^{\frac{\pi i h}{n}}, & 0 \\ 0, & \pm e^{-\frac{\pi i h}{n}} \end{pmatrix} \quad h = 0, 1, \ldots, n-1. \quad \text{(both signs)}}
If $n$ is even, then $C_1'$ and $C_1''$ are both contained in $C_n''$; but if $n$ is odd, then $C_1'$ is contained in $C_n'$, and $C_1''$ in $C_n''$.

If next $\alpha = \delta = 0$, $\beta\gamma = 1$, then we can transform $\varphi$ by the transformation
\[
\begin{pmatrix} \beta^{-1/2}, & 0 \\ 0, & \beta^{1/2} \end{pmatrix}
\begin{pmatrix} 0, & \beta \\ \gamma, & 0 \end{pmatrix}
\begin{pmatrix} \beta^{1/2}, & 0 \\ 0, & \beta^{-1/2} \end{pmatrix}
= \begin{pmatrix} 0, & 1 \\ 1, & 0 \end{pmatrix},
\]
by which the group $P$ remains unchanged, and we obtain yet a third group:
\weq{3}{C_n''' = \begin{pmatrix} e^{\frac{\pi i h}{n}}, & 0 \\ 0, & e^{-\frac{\pi i h}{n}} \end{pmatrix}, \quad \begin{pmatrix} 0, & e^{\frac{\pi i h}{n}} \\ e^{-\frac{\pi i h}{n}}, & 0 \end{pmatrix} \qquad h = 0, 1, \ldots, n-1.}

III. Let $P$ be the dihedral group $D_n$, which consists of the substitutions
\weq{8}{\begin{aligned} D_n = {}& 1, c, c^2, \ldots, c^{n-1} \\ & d, cd, c^2d, \ldots, c^{n-1}d, \\ = {}& C_n + C_n d \end{aligned} \qquad d = \begin{pmatrix} 0, & i \\ i, & 0 \end{pmatrix}}
The substitution $c$ is of degree $n$, and the $c^\lambda d$ are all of degree $2$. But $\varphi^{-1}c\varphi$ is of degree $n$, and must therefore, if $n > 2$, be contained among the powers of $c$. Consequently the preceding considerations are applicable to the determination of $\varphi$, and one of the groups $C_n'$, $C_n''$, $C_n'''$ must be contained in the group $Q$. Since $d$ does not occur in these groups, there result for the dihedral groups of the second kind the forms
\weq{9}{C_n' + C_n'd, \quad C_n'' + C_n''d, \quad C_n''' + C_n'''d,}
%CONT>>


% ===== scan leaf 259, printed 238 =====
% !PRINTED 238
\weberpage{238}
%<<CONT
of which the first two are
\[
1)\quad D_n' = \begin{pmatrix} e^{\frac{\pi i h}{2n}}, & 0 \\ 0, & (-1)^h e^{-\frac{\pi i h}{2n}} \end{pmatrix},\quad
\begin{pmatrix} 0, & -e^{\frac{\pi i h}{2n}} \\ (-1)^{h+n} e^{-\frac{\pi i h}{2n}}, & 0 \end{pmatrix}
\]
\[
h = 0, 1, \ldots, 2n-1,
\]
\[
2)\quad D_n'' = \begin{pmatrix} e^{\frac{\pi i h}{n}}, & 0 \\ 0, & \pm\, e^{-\frac{\pi i h}{n}} \end{pmatrix},\quad
\begin{pmatrix} 0, & e^{\frac{\pi i (2h+n)}{2n}} \\ \pm\, e^{-\frac{\pi i (2h+n)}{2n}}, & 0 \end{pmatrix}
\]
\[
h = 0, 1, \ldots, n-1\text{,}
\]
while the third group in (9) is
\[
\begin{pmatrix} e^{\frac{\pi i h}{n}}, & 0 \\ 0, & e^{-\frac{\pi i h}{n}} \end{pmatrix},\quad
\begin{pmatrix} 0, & e^{\frac{\pi i h}{n}} \\ e^{-\frac{\pi i h}{n}}, & 0 \end{pmatrix},
\]
\[
\begin{pmatrix} 0, & e^{\frac{\pi i}{2n}(2h+n)} \\ -e^{-\frac{\pi i}{2n}(2h+n)}, & 0 \end{pmatrix},\quad
\begin{pmatrix} e^{\frac{\pi i}{2n}(2h+n)}, & 0 \\ 0, & -e^{-\frac{\pi i}{2n}(2h+n)} \end{pmatrix}
\]
\[
h = 0, 1, \ldots, n-1\text{.}
\]
For even $n$ this coincides with $D_n''$, and for odd $n$ with $D_n'$.

The case $n=2$, where $c = \begin{pmatrix} i, & 0 \\ 0, & -i \end{pmatrix}$, forms only an apparent exception, because in this case
\weq{10}{\varphi^{-1} c \varphi = d \quad \text{or} \quad \varphi^{-1} c \varphi = c d}
can hold. If one pursues these assumptions by a simple calculation, two further groups result:
\weq{11}{
\begin{aligned}
&1,\ c,\ d,\ cd,\ \varphi_1,\ \varphi_1 c,\ \varphi_1 d,\ \varphi_1 c d\\
&1,\ c,\ d,\ cd,\ \varphi_2,\ \varphi_2 c,\ \varphi_2 d,\ \varphi_2 c d,
\end{aligned}
}
in which
\[
\varphi_1 = \begin{pmatrix} \dfrac{1}{\sqrt2}, & \dfrac{1}{\sqrt2} \\ \dfrac{1}{\sqrt2}, & \dfrac{-1}{\sqrt2} \end{pmatrix},\quad
\varphi_2 = \begin{pmatrix} \dfrac{1}{\sqrt2}, & -\dfrac{i}{\sqrt2} \\ \dfrac{i}{\sqrt2}, & -\dfrac{1}{\sqrt2} \end{pmatrix}.
\]


% ===== scan leaf 260, printed 239 =====
% !PRINTED 239
\weberpage{239}
But we have
\[
\varphi_2^{-1} \varphi_1 \varphi_2 = \begin{pmatrix} 0, & e^{\frac{\pi i}{4}} \\ e^{-\frac{\pi i}{4}}, & 0 \end{pmatrix}
\]
\[
\varphi_1^{-1} \varphi_2 \varphi_1 = \begin{pmatrix} 0, & e^{\frac{\pi i}{4}} \\ e^{-\frac{\pi i}{4}}, & 0 \end{pmatrix},
\]
\[
\begin{aligned}
&c\varphi_1 = \varphi_1 d,\quad d\varphi_1 = \varphi_1 c,\quad cd\varphi_1 = \varphi_1 dc\\
&c\varphi_2 = \varphi_2 dc,\quad d\varphi_2 = \varphi_2 d,\quad cd\varphi_2 = \varphi_2 c,
\end{aligned}
\]
and from these formulae it follows that the groups (11) are transformed into $D_2'$ by transformation with $\varphi_1$ and $\varphi_2$.

In the determination of the remaining polyhedral groups of the second kind the following general remarks are of use.
The inversion $j = \begin{pmatrix} i, & 0 \\ 0, & i \end{pmatrix}$ is a similarity substitution (\S.\,37) and therefore commutes with every other substitution, and consequently from every polyhedral group of the first kind we can derive at least one group of the second kind:
\weq{12}{P + Pj.}
In order to find the other groups of this kind that may still exist, we recall that by \gesperrt{Sylow}'s theorem (\S.\,29, I.) every group $G$ must contain a group whose degree is the highest power of 2 that divides the degree of $G$. But the extended tetrahedral, octahedral and icosahedral groups have the degrees 24, 48, 120, and must therefore contain a divisor of degree 8, 16, 8.

Now let $T$ be the tetrahedral, $O$ the octahedral and $J$ the icosahedral group. The groups $T$ and $J$ contain a four-group $D_2$ but no substitution of the fourth order, while $O$ contains a dihedral group $D_4$ but no substitution of the eighth order. Each extended polyhedral group must therefore contain one of the extended dihedral groups considered under III.
% [sic] the German prints \"o here (Di\"odergruppe) vs \"e elsewhere on this page (Di\"edergruppen) — reproduced as read in the German column.

Of the extended dihedral groups, however, $D_2''$ arises from $D_2$ by inversion, while
\[
\begin{aligned}
D_2' = &\begin{pmatrix}1, & 0\\ 0, & 1\end{pmatrix},\ \begin{pmatrix}\sqrt{i}, & 0\\ 0, & +i\sqrt{i}\end{pmatrix},\ \begin{pmatrix}i, & 0\\ 0, & -i\end{pmatrix},\ \begin{pmatrix}i\sqrt{i}, & 0\\ 0, & +\sqrt{i}\end{pmatrix},\\
&\begin{pmatrix}0, & -1\\ 1, & 0\end{pmatrix},\ \begin{pmatrix}0, & -i\sqrt{i}\\ \sqrt{i}, & 0\end{pmatrix},\ \begin{pmatrix}0, & i\\ i, & 0\end{pmatrix},\ \begin{pmatrix}0, & -\sqrt{i}\\ i\sqrt{i}, & 0\end{pmatrix}
\end{aligned}
\]
%CONT>>


% ===== scan leaf 261, printed 240 =====
% !PRINTED 240
\weberpage{240}
%<<CONT
arises through composition of $D_2$ with
\[
j_1 = \begin{pmatrix} 0, & -\sqrt{i} \\ i\sqrt{i}, & 0 \end{pmatrix}
\]
and thus yields
\[
D_2' = D_2 + D_2 j_1\text{.}
\]
Now in general:
\[
j_1^{-1} \begin{pmatrix} \alpha, & \beta \\ \gamma, & \delta \end{pmatrix} j_1 = \begin{pmatrix} \delta, & \gamma i \\ -\beta i, & \alpha \end{pmatrix},
\]
and from this one concludes by \S.\,56, (4), that $j_1^{-1} T j_1 = T$, and hence that from the tetrahedral group two extended groups $T' = T+Tj$ and $T'' = T+Tj_1$ can be derived; but also only two, because by \S.\,56 the tetrahedral group is completely determined by a four-group contained in it. For if, instead of the extension $D_2'$, we choose the extension of the four-group to one of the groups (11), we obtain an extended tetrahedral group $T'''$, which passes into $T''$ by transformation with $\varphi_1$ or $\varphi_2$.

Since the substitution
\[
\begin{pmatrix} 0, & e^{\frac{\pi i}{8}} \\ e^{-\frac{\pi i}{8}}, & 0 \end{pmatrix}
\]
does not commute with the octahedral group \S.\,57, (8), no further group of the second kind can be derived from the octahedral group except by inversion, and the same holds of the icosahedral group. To apply this mode of reasoning, one must transform the icosahedral group in such a way that one of the four-groups contained in it takes the simplest form, which is somewhat lengthy, but by no means difficult, and shall not be carried out further here.

We thus obtain the following further groups of the second kind:
\[
\begin{aligned}
&\text{IV.}\quad 1)\ T' = T+Tj, \qquad 2)\ T'' = T+Tj_1.\\
&\text{V.}\qquad\qquad O' = O+Oj.\\
&\text{IV.}\qquad\qquad J' = J+Jj.
\end{aligned}
\]
% [sic] the third line is printed "IV." again (as the first); expected "VI." given the sequence — reproduced as printed.


% ===== scan leaf 262, printed 241 =====
% !PRINTED 241
\weberpage{241}
Herein are contained the 32 symmetry systems of the crystallographers. They are the groups:
\[
\begin{array}{ll}
C_1,\ C_1',\ C_1'' & D_2,\ D_2',\ D_2''\\
C_2,\ C_2',\ C_2'',\ C_3''' & D_3,\ D_3',\ D_3''\\ % [sic: Weber prints C_3''' in the C_2 row; C_2''' is mathematically expected.]
C_3,\ C_3',\ C_3'',\ C_3''' & D_4,\ D_4'\\
C_4,\ C_4',\ C_4'' & D_6,\ D_6'\\
C_6,\ C_6',\ C_6'' & T,\ T',\ T''\\
 & O,\ O'.
\end{array}
\]
The remaining polyhedral groups are excluded in crystallography by the crystallographic law of rational indices\footnote{Cf. \gesperrt{Sch\"onfliess}, Krystallsysteme u. Krystallstructur. Leipzig 1891.}.


% ===== scan leaf 263, printed 242 =====
\weberpage{242}% page number not printed on this page
\pdfbookmark[1]{Ninth Section. Congruence Groups.}{bm73}
\pdfbookmark[2]{§ 63. Function congruences}{bm74}
% !PRINTED none
% Section-opening page: no page number is visible in the running head of any of
% the three bands. Page carries the Abschnitt title, the \S.\,63 sub-heading, the
% opening paragraphs of the section, and a footnote citing the literature.

\gesperrt{Ninth Section.}

\gesperrt{Congruence Groups.}

\S.\,63.

\gesperrt{Function Congruences.}

From the linear substitutions, whose formation and composition we have come to know in the earlier sections, yet another, wholly different kind of finite group can be derived: the \gesperrt{congruence groups}.
% [sic] Weber prints the plural verb "lassen" with the singular head "eine ... Art"; the English gives the intended grammatical sense.

These congruence groups have a manifold interest. They stem in the first place from the theory of elliptic functions, where they present themselves as \gesperrt{Galois} groups of the modular equations. But they are also of importance for general group theory, because they give us a means of forming whole sequences of \gesperrt{simple groups}. This is all the more remarkable since, as we have seen in the fourth section, simple groups, at least among the lower degrees, are very rare.

This theory of congruence groups is not only greatly generalized; its governing laws also emerge much more clearly when one relies on an extension of the concept of congruence originating with \gesperrt{Gauss}, which has since been developed and applied several times to problems of group theory, especially by \gesperrt{Galois}\footnote{\gesperrt{Galois}, ``Sur la théorie des nombres''. Bulletin des sciences math. de Ferussac. 1830. (Cf. the note to \S.\,151 of the first volume.) --- \gesperrt{Schönemann}, Grundzüge einer allgemeinen Theorie der höheren Congruenzen, deren Modulus eine reelle Primzahl ist. (Crelle's Journ. f. Mathematik, Bd. 31, 1846.) --- \gesperrt{Dedekind}, Abriss einer Theorie der höheren Congruenzen in Bezug auf einen reellen Primzahlmodulus. (Crelle's Journ. f. Mathematik, Bd. 54, 1856.)}.


% ===== scan leaf 264, printed 243 =====
% !PRINTED 243
\weberpage{243}
In order to gain a foundation for this theory, we consider an integral function of a variable quantity $t$

\weq{1}{f(t) = a_0 t^n + a_1 t^{n-1} + \cdots + a_{n-1} t + a_n,}

whose coefficients are to be integers. We choose in addition a prime $p$ as modulus, and assume that $a_0$ is not divisible by $p$.

Two integral functions of $t$ in which corresponding coefficients are congruent with respect to the modulus $p$ shall themselves be called \gesperrt{congruent with respect to the modulus $p$}.

The function $f(t)$ is called \gesperrt{reducible with respect to the modulus $p$} if there exist two integral functions $\varphi(t)$, $\psi(t)$ with integral coefficients, in each of which at least one term depending on $t$ has a coefficient not divisible by $p$, such that

\weq{2}{f(t) \equiv \varphi(t)\,\psi(t)\ (\mathrm{mod}\ p)}

holds. In $\varphi(t)$ and $\psi(t)$ one may omit all terms whose coefficients are divisible by $p$, and then the degree of $\varphi(t)\,\psi(t)$ must agree with the degree of $f(t)$. Hence the degrees of $\varphi(t)$ and $\psi(t)$ must be lower than $n$. If there are no such functions $\varphi(t)$, $\psi(t)$, then $f(t)$ is called \gesperrt{irreducible with respect to the modulus $p$}.

A function irreducible with respect to the modulus $p$ is certainly always absolutely irreducible in the field of rational numbers. But the converse need not hold.

Thus, for example, the function of degree $2$, $t^2 - R$, is reducible with respect to $p$ if $R$ is a quadratic residue of $p$; for if $a^2 \equiv R\ (\mathrm{mod}\ p)$, then

\[
t^2 - R \equiv (t-a)(t+a)\ (\mathrm{mod}\ p).
\]

On the other hand $t^2 - N$, if $N$ is a non-residue of $p$, is irreducible, because otherwise $t^2 - N$ would have to become divisible by $p$ for a rational $t$.

Another example is offered by the functions

\[
t^2 + t + 1, \quad t^3 + t + 1,
\]

which are irreducible for the modulus $2$. For were they reducible, one of the factors would have to be linear, and there would have to be an integer which, substituted for $t$, makes these functions even numbers. But that is evidently impossible, since both functions are odd numbers for $t = 0$ and $t = 1$.


% ===== scan leaf 265, printed 244 =====
% !PRINTED 244
\weberpage{244}
If

\weq{3}{P = t^n + p_1 t^{n-1} + p_2 t^{n-2} + \cdots + p_{n-1} t + p_n}

denotes a function of degree $n$ irreducible with respect to the modulus $p$, in which for the sake of simplicity we assume the coefficient of the highest power $t^n$ to be equal to $1$, then from every other integral function $F(t)$ with integral coefficients a residue $\Phi(t)$ can be derived whose degree is lower than $n$, by setting (according to \S.\,3 of the first volume)

\weq{4}{F(t) = QP + \Phi(t)}

Here $\Phi(t)$ has integral coefficients, and we designate its relation to the function $F(t)$ as a congruence with respect to the modulus $P$.

\gesperrt{Accordingly, two functions $F(t)$, $F_1(t)$ that have the same residue $\Phi(t)$ are called congruent with respect to the modulus $P$}, which is expressed by the formula

\[
F(t) \equiv F_1(t)\ (\mathrm{mod}\ P)
\]

This is equivalent to saying that $F(t) - F_1(t)$ is divisible by $P$.

If two functions $F(t)$ and $F_1(t)$ do not give the same residue, but two residues that are congruent with respect to the modulus $p$, then the functions $F$, $F_1$ are called congruent with respect to the double modulus $P$, $p$, and one expresses this by a formula as follows:

\weq{5}{F(t) \equiv F_1(t)\ (\mathrm{mod}\ P, p).}

For this kind of congruence the following theorem holds:

\gesperrt{1. The product of two functions $F(t)$, $F_1(t)$ cannot be congruent to zero unless one of the factors is congruent to zero.}

The proof follows from the algorithm of the greatest common divisor. For if $P_1$ is an integral function of lower degree than $P$ which is not congruent to zero with respect to the modulus $p$, then one can determine a sequence of functions of the same kind $P_2$, $P_3$, $\ldots$ of decreasing degree, and the quotients $Q_1$, $Q_2$, $\ldots$ likewise as integral functions, in such a way that

\weq{6}{P \equiv Q_1 P_1 + P_2,\quad P_1 \equiv Q_2 P_2 + P_3, \ldots,\quad P_{\nu-2} \equiv Q_{\nu-1} P_{\nu-1} + P_\nu\ (\mathrm{mod}\ p)}

holds, and the sequence of these equations breaks off when $P_\nu$ has become a constant (an integer). This constant $P_\nu$
%CONT>>


% ===== scan leaf 266, printed 245 =====
% !PRINTED 245
\weberpage{245}
\pdfbookmark[2]{§ 64. Congruence fields}{bm75}
%<<CONT
cannot, however, be $\equiv 0\ (\mathrm{mod}\ p)$; for otherwise a congruence could be derived from (6):

\[
P \equiv T P_{\nu-1}\ (\mathrm{mod}\ p),
\]

in which $T$ is a non-constant function of $t$, and $P$ would not be irreducible with respect to the modulus $p$.

Now if $Q$ is any integral function of $t$ that satisfies the condition

\[
Q P_1 \equiv 0\ (\mathrm{mod}\ P, p)
\]

then, on multiplying (6) by $Q$, it follows that:

\[
Q P_2 \equiv 0,\ Q P_3 \equiv 0,\ \ldots,\ Q P_\nu \equiv 0\ (\mathrm{mod}\ P, p),
\]

hence $Q \equiv 0$. And if we now take for $P_1$, $Q$ the residues of the functions $F(t)$, $F_1(t)$, then Theorem 1. is hereby proved.

The residues of all functions $F(t)$ with respect to the modulus $P$ are of the form

\weq{7}{X = X(t) = x_0 + x_1 t + \cdots + x_{n-1} t^{n-1},}

and if one wishes to have only those among them that are incongruent with respect to the modulus $p$, it suffices to let the coefficients $x_0, x_1, \ldots x_{n-1}$ run through the sequence of numbers $0, 1, 2, \ldots, p-1$.

\gesperrt{There are therefore $p^n$ and no more functions incongruent with respect to the moduli $P$, $p$.}

\S.\,64.

\gesperrt{Congruence fields.}

It is only another mode of expression for the considerations carried out in the preceding section if, following \gesperrt{Galois}, one introduces an imaginary number $\varepsilon$ which satisfies the congruence

\weq{1}{P(\varepsilon) \equiv 0\ (\mathrm{mod}\ p)}

There is no integral rational number of this kind as soon as $n > 1$, because otherwise $P(t) - P(\varepsilon) = (t-\varepsilon)Q$ would be divisible by $t-\varepsilon$, hence $P(t) \equiv (t-\varepsilon)Q$, and $P(t)$ would not be irreducible. Nevertheless one can, just as in ordinary number theory one introduces the imaginary unit $i = \sqrt{-1}$, make use of such a symbol $\varepsilon$ in calculation. One calculates with it, so far as addition, subtraction and multiplication are concerned, according to the rules of literal calculation, and in doing so may use the equation $P(\varepsilon) = 0$ at pleasure. One can even take $\varepsilon$ to be simply a root of the equation $P(x) = 0$
%CONT>>


% ===== scan leaf 267, printed 246 =====
% !PRINTED 246
\weberpage{246}
%<<CONT
in the ordinary sense. All numbers at which one arrives in this way can be brought to the form

\weq{2}{\alpha = a_0 + a_1\varepsilon + \cdots + a_{n-1}\varepsilon^{n-1}}

Since in calculating with such numbers the modulus $p$ is always held fixed, we shall designate two numbers that are congruent with respect to the modulus $p$ simply as \gesperrt{equal}. The objects of the calculation are then properly not the individual numbers themselves, but the \gesperrt{classes of numbers} consisting of all numbers congruent to one another. These numbers $\alpha$ are called the \gesperrt{Galois imaginaries}. The system of such imaginary numbers belonging to a definite $\varepsilon$ we denote for brevity by $\mathfrak{C}$.

We then have first the theorem:

\gesperrt{2. There are $p^n$ and no more distinct numbers in $\mathfrak{C}$.}

But from Theorem 1. it follows further:

\gesperrt{3. The product of two or more numbers from $\mathfrak{C}$ is equal to zero if and only if at least one of the factors is equal to zero.}

We obtain the complete system of numbers $\mathfrak{C}$ if the rational numbers $a_0, a_1, \ldots, a_{n-1}$ in (2) are each made to run through a complete system of residues with respect to the modulus $p$. Every system $\mathfrak{C}$ contains the residues of the natural numbers $0, 1, \ldots, p-1$, and for $n = 1$, $\mathfrak{C}$ is identical with this system.

If $\alpha$ is a fixed number in $\mathfrak{C}$ different from zero, then $\alpha\xi$ runs through the full system $\mathfrak{C}$ together with $\xi$. For it follows from 3. that $\alpha\xi$ can be equal to $\alpha\xi'$ only when $\xi = \xi'$. From this it follows:

\gesperrt{4. If $\alpha, \beta$ are two given numbers in $\mathfrak{C}$ and $\alpha$ is different from zero, then there is one and only one number $\gamma$ in $\mathfrak{C}$ for which the condition}

\[
\alpha\gamma = \beta
\]

\gesperrt{is satisfied.}

With this the operation of division in the system $\mathfrak{C}$ is also shown to be permissible. We denote the number $\gamma$ whose existence is asserted in 4. by

\[
\beta : \alpha \quad \text{or} \quad \frac{\beta}{\alpha}.
\]


% ===== scan leaf 268, printed 247 =====
% !PRINTED 247
\weberpage{247}
One may call the system $\mathfrak{C}$ a \gesperrt{finite field}, since it contains only finitely many numbers and has the characteristic property of a field: the arithmetic operations can be performed without restriction, except for division by zero (Bd.\,I, \S.\,139)\textsuperscript{1)}. We shall therefore call $\mathfrak{C}$ a \gesperrt{congruence field}.

Let $n$ be called the \gesperrt{degree} and $p$ the \gesperrt{modulus} of the field $\mathfrak{C}$.

If $\alpha$ is a number in $\mathfrak{C}$ different from zero, then the product $\alpha\xi = \eta$ runs through the full number system $\mathfrak{C}$ together with $\xi$. If we exclude $\xi = 0$, then $\eta = 0$ also does not occur, and the product $\Pi$ of all the $\xi$ agrees with the product of all the $\eta$ and is different from zero. But by multiplication of all the equations $\alpha\xi = \eta$ it follows that

\[
\alpha^{p^n-1}\Pi = \Pi,
\]

and after cancelling the common factor $\Pi$ there results

\gesperrt{5. Fermat's theorem:}

\weq{3}{\alpha^{p^n-1} = 1.}

If one multiplies by $\alpha$, one obtains this theorem in the form

\weq{4}{\alpha^{p^n} = \alpha,}

in which it still holds also for $\alpha = 0$. If $\alpha$ is different from zero, then on account of (3) there is a \gesperrt{least} positive exponent $e$ for which

\weq{5}{\alpha^e = 1}

holds, and for which consequently the powers $1, \alpha, \alpha^2, \ldots, \alpha^{e-1}$ are different from one another. One then says that $\alpha$ \gesperrt{belongs to the exponent} $e$. This exponent $e$ must be a divisor of $p^n - 1$. For if $\alpha^m = 1$ for any exponent $m$, we set $m = qe + e'$, where $q$ is an integer and $e' < e$. Then $\alpha^{e'} = 1$ as well, and consequently $e' = 0$ must hold, i.\,e. $m$ must be divisible by $e$. Among these exponents $m$ there is found also $p^n - 1$. If we therefore set

\[
p^n - 1 = ef,
\]

then $\alpha^h$ will always belong to the exponent $e$ whenever $h$ is relatively prime to $e$. For then and only then is $he$ the least positive multiple of $h$ divisible by $e$.

\textsuperscript{1)} See the author's memoir: Die allgemeinen Grundlagen der \gesperrt{Galois'schen} Gleichungstheorie. Mathematische Annalen, Bd.\,43.


% ===== scan leaf 269, printed 248 =====
% !PRINTED 248
\weberpage{248}
If, then, any element $\alpha$ belongs to the exponent $e$, the number of such elements equals the number of positive integers less than $e$ and relatively prime to $e$. This number, which we have already denoted by $\varphi(e)$, satisfies, as $e$ runs through all divisors of $p^n - 1$, the relation

\weq{6}{\textstyle\sum \varphi(e) = p^n - 1}

(Bd.\,I, \S.\,132, 133).

If $\psi(e)$ is the number of elements $\alpha$ that belong to the exponent $e$, then $\psi(e) = \varphi(e)$ or $= 0$, and since the total number of elements $\alpha$ is $p^n - 1$ and every element $\alpha$ belongs to one and only one exponent, we have also

\weq{7}{\textstyle\sum \psi(e) = p^n - 1,}

and from (6) and (7) it follows that $\psi(e)$ is \gesperrt{always equal to} $\varphi(e)$.

There are therefore $\varphi(p^n - 1)$ numbers $\gamma$ in $\mathfrak{C}$ that belong to the exponent $p^n - 1$, and which consequently have the property that the powers

\weq{8}{1, \gamma, \gamma^2, \ldots, \gamma^{p^n-2}}

are all different from one another. By this sequence, therefore, the totality of the numbers in $\mathfrak{C}$ different from zero is exhausted.

Such numbers $\gamma$ are called \gesperrt{primitive roots of $\mathfrak{C}$}.

Among the numbers in $\mathfrak{C}$ there are some that can be represented as the square of another number in $\mathfrak{C}$, which we call \gesperrt{squares}, and others for which this is not the case, which we call \gesperrt{non-squares}.

When $p = 2$, every number in $\mathfrak{C}$ is a square, as formula (4) shows.

But when $p$ is odd, one half of the numbers in $\mathfrak{C}$ different from zero consists of squares, the other of non-squares.

For two opposite numbers, such as $+\alpha$ and $-\alpha$, are different from one another when $p$ is odd, and yet give the same square. If, then, one raises all the numbers of $\mathfrak{C}$ to the square, one obtains at most $\frac{1}{2}(p^n - 1)$ different squares. Now, on the other hand, $\beta^2$ can be equal to $\alpha^2$ only when $\beta = \pm\alpha$; for from $\beta^2 - \alpha^2 = (\beta - \alpha)(\beta + \alpha) = 0$ it follows that $\beta - \alpha$ or $\beta + \alpha$ must vanish. There are therefore really $\frac{1}{2}(p^n - 1)$ squares and just as many non-squares.


% ===== scan leaf 270, printed 249 =====
% !PRINTED 249
\weberpage{249}
If one counts zero among the squares, one obtains the theorem:

\gesperrt{6. If $p = 2$, all numbers in $\mathfrak{C}$ are squares. If $p$ is odd, there are $\frac{1}{2}(p^n + 1)$ squares, $\frac{1}{2}(p^n - 1)$ non-squares in $\mathfrak{C}$.}

For odd $p$ a primitive root must always be a non-square, and in the sequence (8) the numbers with even exponents are the squares, those with odd exponents the non-squares.

The product and the quotient of two squares are each evidently again squares.

From this it follows that the product of a non-square and a square different from zero is a non-square. For if the product $\alpha\beta$ is a square and the one factor $\alpha$ is a square, then $\beta = \alpha\beta : \alpha$ must also be a square.

From this one concludes further that the product of two non-squares is a square. For if $\beta$ denotes any non-square, let the factor $\xi$ in the product $\beta\xi$ run through all the numbers of $\mathfrak{C}$ different from zero. The product then runs through the same sequence of numbers, and since $\beta\xi$ is a non-square for every square $\xi$, it must be a square for every non-square $\xi$.

All this also follows very simply from the representation (8) of the numbers of $\mathfrak{C}$ by a primitive root.

We close these general considerations with the theorem, which we shall need later:

\gesperrt{7. Every non-square number is the sum of two squares.}

Here $p$ is to be assumed odd, since only then are there non-square numbers.

First of all, every number $\beta$ can then be represented as the difference of two squares, as the identity

\[
\left(\beta + \tfrac{1}{4}\right)^2 - \left(\beta - \tfrac{1}{4}\right)^2 = \beta
\]

shows. Now if $-1$ is a square in $\mathfrak{C}$, one sets

\[
\left(\beta + \tfrac{1}{4}\right)^2 = \xi^2, \quad -\left(\beta - \tfrac{1}{4}\right)^2 = \eta^2,
\]

and obtains

\[
\beta = \xi^2 + \eta^2.
\]

But if $-1$, and therefore also $p - 1$, is a non-square, then consider the sequence of the natural numbers $1, 2, \ldots, p - 1$, whose
%CONT>>


% ===== scan leaf 271, printed 250 =====
% !PRINTED 250
\weberpage{250}
\pdfbookmark[2]{§ 65. Congruence groups in the field E}{bm76}
%<<CONT
first term is a square and whose last is a non-square, and choose in it a square term followed by a non-square.

If then $a = \xi^2$ is a square and $\xi^2 + 1$ a non-square, then, if $\beta$ is a non-square, $\beta : \xi^2 + 1 = \eta^2$ is a square, and from this it follows:
\[ \beta = \xi^2 \eta^2 + \eta^2 \]
q. e. d.

\begin{center}
\S.\,65.

\gesperrt{Congruence groups in the field $\mathfrak{C}$.}
\end{center}

In every congruence field $\mathfrak{C}$ one can form a finite group of linear substitutions, if in
\weq{1}{A = \begin{pmatrix} \alpha, & \beta \\ \gamma, & \delta \end{pmatrix}}
one lets the elements $\alpha, \beta, \gamma, \delta$ run through all the numbers of $\mathfrak{C}$, and if one composes any two of these substitutions according to the rule
\weq{2}{\begin{pmatrix} \alpha, & \beta \\ \gamma, & \delta \end{pmatrix} \begin{pmatrix} \alpha', & \beta' \\ \gamma', & \delta' \end{pmatrix} = \begin{pmatrix} \alpha\alpha' + \beta\gamma', & \alpha\beta' + \beta\delta' \\ \gamma\alpha' + \delta\gamma', & \gamma\beta' + \delta\delta' \end{pmatrix}}
The degree of this group is $p^{4n}$, if $p$ is the modulus and $n$ the degree of $\mathfrak{C}$.

But in it there is contained a group of lower degree. If we call $\alpha\delta - \beta\gamma$ the determinant of the substitution (1), it follows from (2) that the determinant of the substitution $AA'$ composed from $A$ and $A'$ is equal to the product of the determinants of $A$ and $A'$. If therefore we impose on the numbers $\alpha, \beta, \gamma, \delta$ the condition that
\weq{3}{\alpha\delta - \beta\gamma = 1}
shall hold, then these elements $A$ likewise form a group. The degree of this group is equal to the number of solutions of (3). To find this number, we first observe that for $\alpha, \beta$ we may take any two numbers from $\mathfrak{C}$ which are not both equal to zero. For if, say, $\alpha$ is different from zero, one can set $\gamma = 0, \delta = 1 : \alpha$ and thereby satisfy the condition (3). The number of usable pairs of numbers $\alpha, \beta$ is therefore $p^{2n} - 1$. But if for a pair of numbers $\alpha, \beta$ a solution $\gamma_0, \delta_0$ of (3) has been found, then all the others must satisfy the condition
\[ \alpha (\delta - \delta_0) = \beta (\gamma - \gamma_0) \]
%CONT>>


% ===== scan leaf 272, printed 251 =====
% !PRINTED 251
\weberpage{251}
%<<CONT
If one therefore sets $\gamma - \gamma_0 = \alpha\xi$ (in case $\alpha$ is different from zero), it follows that $\delta - \delta_0 = \beta\xi$, hence
\[ \gamma = \gamma_0 + \alpha\xi, \quad \delta = \delta_0 + \beta\xi, \]
and conversely every pair of numbers $\gamma, \delta$ contained in this form also satisfies it. This evidently holds also in the excepted case $\alpha = 0$. Since we can now put $p^n$ different numbers for $\xi$, the degree of our group is therefore
\weq{4}{p^n (p^{2n} - 1).}

If $p$ is odd, we can derive a still smaller group:

For if $p = 2$, a number $\alpha$ is not different from $-\alpha$, and hence the elements
\weq{5}{\begin{pmatrix} \alpha, & \beta \\ \gamma, & \delta \end{pmatrix}, \quad \begin{pmatrix} -\alpha, & -\beta \\ -\gamma, & -\delta \end{pmatrix}}
are also identical with one another. But if $p$ is odd, the two elements (5) are different from one another, and the whole system of these substitutions splits into pairs of the form (5).

If we take two such pairs and compose each element of the one pair with each element of the other according to the rule (2), we obtain only two different elements, which again form such a pair. These pairs can therefore themselves be regarded as elements of a new group of degree
\weq{6}{\dfrac{p^n (p^{2n} - 1)}{2}}
We can also express this by saying that two elements $A$ whose corresponding numbers differ only in sign are to be regarded as identical.

The group so defined, whose degree is expressed for $p = 2$ by (4) and for an odd $p$ by (6), we shall call \gesperrt{the congruence group belonging to the field $\mathfrak{C}$, and denote it by $E$}.

For the closer investigation of this group it is of importance to know a system of generating substitutions. If, following \S.\,64, (2), we represent the numbers of $\mathfrak{C}$ in the form:
\weq{7}{\xi = x_0 + x_1 \varepsilon_1 + \cdots + x_{n-1} \varepsilon^{n-1},} % [sic] first power printed as subscript \varepsilon_1, not \varepsilon or \varepsilon^1
in which the $x_i$ are rational numbers taken with respect to the modulus $p$
%CONT>>


% ===== scan leaf 273, printed 252 =====
% !PRINTED 252
\weberpage{252}
%<<CONT
denote, then we obtain a system of generating substitutions in the form:
\weq{8}{A_0 = \begin{pmatrix} 0, & 1 \\ -1, & 0 \end{pmatrix}, \quad B_h = \begin{pmatrix} 1, & 0 \\ \varepsilon^h, & 1 \end{pmatrix}, \quad h = 0, 1, \ldots, n-1.}

To prove this, we observe that
\[ \begin{pmatrix} 1, & 0 \\ \gamma, & 1 \end{pmatrix} \begin{pmatrix} 1, & 0 \\ \gamma', & 1 \end{pmatrix} = \begin{pmatrix} 1, & 0 \\ \gamma + \gamma', & 1 \end{pmatrix} \]
holds; repeated application therefore shows that, if
\[ \gamma = c_0 + c_1 \varepsilon + \cdots + c_{n-1} \varepsilon^{n-1} \]
is an arbitrary number in $\mathfrak{C}$ whose coefficients $c$ are integers, then
\weq{9}{\begin{pmatrix} 1, & 0 \\ \gamma, & 1 \end{pmatrix} = B_0^{c_0} B_1^{c_1} \cdots B_{n-1}^{c_{n-1}}.}

Further, by composition with $A_0$ it follows that:
\weq{10}{\begin{gathered} \begin{pmatrix} 0, & 1 \\ -1, & 0 \end{pmatrix}\begin{pmatrix} 1, & 0 \\ \alpha, & 1 \end{pmatrix} = \begin{pmatrix} \alpha, & 1 \\ -1, & 0 \end{pmatrix}, \quad \begin{pmatrix} 1, & 0 \\ \delta, & 1 \end{pmatrix}\begin{pmatrix} 0, & 1 \\ -1, & 0 \end{pmatrix} = \begin{pmatrix} 0, & 1 \\ -1, & \delta \end{pmatrix} \\ \begin{pmatrix} 0, & 1 \\ -1, & 0 \end{pmatrix}\begin{pmatrix} 1, & 0 \\ -\beta, & 1 \end{pmatrix}\begin{pmatrix} 0, & -1 \\ 1, & 0 \end{pmatrix} = \begin{pmatrix} 1, & \beta \\ 0, & 1 \end{pmatrix}. \end{gathered}}

From this it appears that from the elements (8) one can compose all substitutions of the form
\[ \begin{pmatrix} \alpha, & 1 \\ -1, & 0 \end{pmatrix}, \quad \begin{pmatrix} 1, & \beta \\ 0, & 1 \end{pmatrix}, \quad \begin{pmatrix} 1, & 0 \\ \gamma', & 1 \end{pmatrix}, \quad \begin{pmatrix} 0, & 1 \\ -1, & \delta \end{pmatrix} \]
in which $\alpha, \beta, \gamma, \delta$ are arbitrary numbers in $\mathfrak{C}$. Further,
\weq{11}{\begin{pmatrix} \alpha, & 1 \\ -1, & 0 \end{pmatrix}\begin{pmatrix} 1, & \beta \\ 0, & 1 \end{pmatrix}\begin{pmatrix} 1, & 0 \\ \gamma, & 1 \end{pmatrix} = \begin{pmatrix} \alpha + \gamma + \alpha\beta\gamma, & \alpha\beta + 1 \\ -\beta\gamma - 1, & -\beta \end{pmatrix}.}

If one assumes $\beta$ different from zero, one can choose $\alpha$ and $\gamma$ so that $\alpha\beta + 1$, $-\beta\gamma - 1$ become arbitrary numbers, and by (11) one can therefore compose from $A_0, B_h$ all substitutions $\begin{pmatrix} \alpha, & \beta \\ \gamma, & \delta \end{pmatrix}$ in which $\delta$ is different from zero. The restriction that $\delta$ be different from zero can, however, by the formula
\[ \begin{pmatrix} \alpha, & \beta \\ \gamma, & 0 \end{pmatrix} = \begin{pmatrix} -\beta, & \alpha \\ 0, & \gamma \end{pmatrix}\begin{pmatrix} 0, & -1 \\ 1, & 0 \end{pmatrix} \]
be dropped.


% ===== scan leaf 274, printed 253 =====
% !PRINTED 253
\weberpage{253}
The group $E$ contains a divisor of index $p^n + 1$, which consists of all substitutions of the form
\[ \begin{pmatrix} \alpha, & \beta \\ 0, & \delta \end{pmatrix} \]
in which $\beta$ denotes an arbitrary number in $\mathfrak{C}$, and $\alpha$ any number different from zero, and $\delta = \alpha^{-1}$.

To this divisor there corresponds a permutation group on $p^n + 1$ letters isomorphic to the group $E$, which one can form as follows:

The linear expression
\weq{12}{\eta = \dfrac{\alpha\xi + \beta}{\gamma\xi + \delta} \quad \text{or} \quad \xi = \dfrac{\delta\eta - \beta}{-\gamma\eta + \alpha}}
gives for every number $\xi$ in $\mathfrak{C}$ a corresponding number $\eta$, except when $\gamma\xi + \delta = 0$. Likewise, to every $\eta$ there is a definite $\xi$, except for $\gamma\eta - \alpha = 0$. To remove this exception, it suffices to extend the system $\mathfrak{C}$ by adjoining \gesperrt{one} element, which we call ``infinity'' and denote by $\infty$. With this symbol one reckons according to the following rules, in which $\alpha$ denotes any element of the extended system $E$:
% [sic] Weber prints ordinary E both here and in "system E" below, although both references mean the extended field \mathfrak{C}; preserved as printed. Ordinary E otherwise denotes the congruence group.
\[
\begin{array}{l@{\quad}l}
\alpha + \infty = \alpha - \infty = \infty, & \text{except for } \alpha = \infty \\
\alpha \cdot \infty = \infty & \text{\textquotedbl}\ \ \text{\textquotedbl}\ \ \alpha = 0 \\
\alpha : 0 = \infty & \text{\textquotedbl}\ \ \text{\textquotedbl}\ \ \alpha = 0 \\
\alpha : \infty = 0 & \text{\textquotedbl}\ \ \text{\textquotedbl}\ \ \alpha = \infty.
\end{array}
\]
Then the result of every calculation with the four arithmetic operations always leads to a definite number of the system $E$, and only the combinations $\infty \pm \infty, 0 . \infty, 0 : 0, \infty : \infty$ remain undefined.

By (12), to the number $\xi = -\delta : \gamma$ there then corresponds the number $\eta = \infty$, and to the number $\xi = \infty$ the number $\eta = \alpha : \gamma$, and to every substitution $A$ in $E$ there corresponds a definite permutation of the $p^n + 1$ elements of the extended system $\mathfrak{C}$. One sees easily that only the identical substitution leaves all these elements unchanged, and that consequently two different substitutions in $E$ always give rise to different permutations. The isomorphism is therefore one-to-one.


% ===== scan leaf 275, printed 254 =====
% !PRINTED 254
\weberpage{254}
\pdfbookmark[2]{§ 66. Simplicity of the Group E}{bm77}

\begin{center}
\S.\,66.\\
\gesperrt{Simplicity of the Group $E$.}
\end{center}

The first question in a more thorough investigation of the group $E$ is that of a possibly existing normal divisor. It can be proved that, apart from two quite simple exceptions, no such normal divisor exists.

Let us assume that $G$ is a normal divisor of $E$ which contains at least one substitution different from the identity,
\weq{1}{A = \begin{pmatrix} \alpha, & \beta \\ \gamma, & \delta \end{pmatrix}}
By the definition of a normal divisor, if $T$ is an arbitrary element of $E$, then
\weq{2}{TAT^{-1}}
is also contained in $G$, and it is to be shown that in this way, and by composition of such substitutions, one can derive all the elements of $E$, from which it then follows that $G$ must be identical with $E$. For this, however, it suffices to establish the occurrence of the generating substitutions $A_0$, $B_h$ [\S.\,65, (8)] in $G$.

We proceed to this step by step.

1) In $G$ there occurs a substitution of the form
\weq{3}{\begin{pmatrix} 0, & \beta \\ \gamma, & \delta \end{pmatrix}}
To show this, we form according to (2):
\weq{4}{\begin{pmatrix} \xi, & 1 \\ -1, & 0 \end{pmatrix}\begin{pmatrix} \alpha, & \beta \\ \gamma, & \delta \end{pmatrix}\begin{pmatrix} 0, & -1 \\ 1, & \xi \end{pmatrix} = \begin{pmatrix} \beta\xi+\delta, & \beta\xi^2+(\delta-\alpha)\xi-\gamma \\ -\beta, & -\beta\xi+\alpha \end{pmatrix}.}
Now if $\beta$ is different from zero, one can determine $\xi$ from $\beta\xi+\delta=0$ and the goal is attained. But if $\beta=0$, then $\gamma=0$, $\delta-\alpha=0$ do not hold simultaneously, since otherwise $A$ would be the identical substitution; one can therefore take $\xi$ in such a way that $(\delta-\alpha)\xi-\gamma$ is different from zero. Then by (4) one has formed a substitution $A$ in $G$ in which $\beta$ is not zero, and to which the formula (4) can now be applied once more, in order to make $\alpha$ zero.

2) In $G$ there occurs a substitution of the form
\[
\begin{pmatrix} 0, & 1 \\ -1, & \xi \end{pmatrix}
\]
%CONT>>


% ===== scan leaf 276, printed 255 =====
% !PRINTED 255
\weberpage{255}
%<<CONT
To prove this, we form according to (2):
\weq{5}{\begin{pmatrix} \lambda, & \mu \\ \nu, & \varrho \end{pmatrix}\begin{pmatrix} 0, & \beta \\ \gamma, & \delta \end{pmatrix}\begin{pmatrix} \varrho, & -\mu \\ -\nu, & \lambda \end{pmatrix} = \begin{pmatrix} \mu\varrho\gamma-\lambda\nu\beta-\mu\nu\delta, & -\mu^2\gamma+\lambda^2\beta+\lambda\mu\delta \\ \varrho^2\gamma-\nu^2\beta-\varrho\nu\delta, & -\varrho\mu\gamma+\lambda\nu\beta+\lambda\varrho\delta \end{pmatrix},}
and it therefore only remains to show that $\lambda,\mu,\nu,\varrho$ can be determined so that the three equations:
\weq{6}{
\begin{aligned}
\lambda\varrho-\mu\nu &= 1\\
\mu\varrho\gamma-\lambda\nu\beta-\mu\nu\delta &= 0\\
-\mu^2\gamma+\lambda^2\beta+\lambda\mu\delta &= 1
\end{aligned}
}
are satisfied. But from the last two equations it follows, if one multiplies them first by $\mu,\varrho$, then by $\lambda,\nu$ and adds each time, using the first equation:
\weq{7}{\lambda\beta+\mu\delta = \varrho, \quad \mu\gamma = \nu,}
and if one substitutes these values of $\varrho,\nu$ in the first equation (6):
\weq{8}{\lambda^2\beta+\lambda\mu\delta-\mu^2\gamma = 1.}
But this equation can always be satisfied. For on account of $\beta\gamma=-1$ neither of the two numbers $\beta,\gamma$ can vanish. Now if $\beta$ is a square, we can set $\mu=0$ and determine $\lambda$ from $\lambda^2\beta=1$. This is \gesperrt{always} possible for $p=2$. But if $p$ is odd, then, if $\beta$ is a non-square, we obtain from (8):
\[
(2\lambda\beta+\delta\mu)^2+\mu^2(4-\delta^2) = 4\beta.
\]
Now if, in the first place, $4-\delta^2$ is a non-square, one can determine $\mu$ from $\mu^2(4-\delta^2)=4\beta$ and then $\lambda$ from $2\lambda\beta+\delta\mu=0$. But if $4-\delta^2$ is a square, one decomposes $4\beta$, according to \S.\,64, 7., into the sum of two squares $4\beta=\sigma^2+\tau^2$ and determines $\mu$ and $\lambda$ from
\[
\mu^2(4-\delta^2) = \sigma^2, \quad 2\lambda\beta+\delta\mu = \tau.
\]
Having thus found $\lambda$ and $\mu$, one obtains $\varrho$ and $\nu$ from (7), and thereby the equations (6) are in fact satisfied. With this the goal is attained.
% [sic] The printed proof omits the case 4-\delta^2=0. Over F_5, \beta=\gamma=\delta=2 has \beta\gamma=-1, while (8), 2\lambda^2+2\lambda\mu-2\mu^2=1, has no solution. Thus the printed argument does not prove the asserted universal solvability; the source text is preserved faithfully.

3) The group $G$ contains the substitution
\[
A_0 = \begin{pmatrix} 0, & 1 \\ -1, & 0 \end{pmatrix},
\]
%CONT>>


% ===== scan leaf 277, printed 256 =====
% !PRINTED 256
\weberpage{256}
%<<CONT
with the sole exception of the case $n=1$, $p=2$. First we have in $G$ the substitutions
\[
\begin{pmatrix} 0, & -1 \\ 1, & 0 \end{pmatrix}\begin{pmatrix} 0, & 1 \\ -1, & \xi \end{pmatrix}\begin{pmatrix} 0, & 1 \\ -1, & 0 \end{pmatrix} = \begin{pmatrix} \xi, & 1 \\ -1, & 0 \end{pmatrix},
\]
hence also
\[
\begin{pmatrix} 0, & 1 \\ -1, & \xi \end{pmatrix}\begin{pmatrix} \xi, & 1 \\ -1, & 0 \end{pmatrix} = \begin{pmatrix} 1, & 0 \\ 2\xi, & 1 \end{pmatrix},
\]
and consequently
\[
\begin{pmatrix} 1, & 0 \\ 2\xi, & 1 \end{pmatrix}\begin{pmatrix} 0, & 1 \\ -1, & \xi \end{pmatrix} = \begin{pmatrix} 0, & 1 \\ -1, & 3\xi \end{pmatrix},
\]
\[
\begin{pmatrix} 1, & 0 \\ 2\xi, & 1 \end{pmatrix}\begin{pmatrix} 0, & 1 \\ -1, & 3\xi \end{pmatrix} = \begin{pmatrix} 0, & 1 \\ -1, & 5\xi \end{pmatrix} \ \text{and so on,}
\]
\[
\begin{pmatrix} 0, & 1 \\ -1, & m\xi \end{pmatrix}
\]
for every odd integer $m$; if therefore $p$ is odd, we can take $m=p$ and obtain $A_0$. But if $p=2$, then, denoting by $\varrho$ a number in $\mathfrak{C}$ still to be determined, we have the following substitutions in $G$:
\[
\begin{pmatrix} 0, & \varrho \\ -\varrho^{-1}, & 0 \end{pmatrix}\begin{pmatrix} 0, & 1 \\ -1, & \xi \end{pmatrix}\begin{pmatrix} 0, & -\varrho \\ \varrho^{-1}, & 0 \end{pmatrix} = \begin{pmatrix} \xi, & \varrho^2 \\ -\varrho^{-2}, & 0 \end{pmatrix}.
\]
Since here every number is a square, one can also replace $\varrho^2$ by $\varrho$, and thus finds that in $G$ the substitution
\[
\begin{pmatrix} \xi, & \varrho \\ -\varrho^{-1}, & 0 \end{pmatrix}
\]
for every nonzero $\varrho$, and consequently also its inverse
\[
\begin{pmatrix} 0, & -\varrho \\ \varrho^{-1}, & \xi \end{pmatrix}
\]
must occur. If, then, $\varrho,\sigma$ denote two numbers different from zero, we have in $G$ also
\weq{9}{\begin{pmatrix} 0, & \varrho^{-1}\sigma^{-1} \\ -\varrho\sigma, & \xi \end{pmatrix}\begin{pmatrix} \xi, & \sigma^{-1} \\ -\sigma, & 0 \end{pmatrix}\begin{pmatrix} 0, & \varrho \\ -\varrho^{-1}, & \xi \end{pmatrix} = \begin{pmatrix} 0, & 1 \\ -1, & \xi\varrho(1+\sigma+\sigma\varrho) \end{pmatrix}.}


% ===== scan leaf 278, printed 257 =====
% !PRINTED 257
\weberpage{257}

If now $n$ is greater than 1, then one can take $\varrho$ and $1+\varrho$ to be different from zero, and then set $\sigma=-1:(1+\varrho)$, whereby the substitution (9) goes over into $A_0$.

But if $n=1$ and $p=2$, then the \gesperrt{first exceptional case} occurs; for then one of the two numbers $\varrho$ and $1+\varrho$ is always equal to zero. This case leads to a group of degree $6$
\[
E = \begin{pmatrix}1,&0\\0,&1\end{pmatrix}, \begin{pmatrix}1,&0\\1,&1\end{pmatrix}, \begin{pmatrix}0,&1\\1,&1\end{pmatrix}, \begin{pmatrix}1,&1\\1,&0\end{pmatrix}, \begin{pmatrix}1,&1\\0,&1\end{pmatrix}, \begin{pmatrix}0,&1\\1,&0\end{pmatrix},
\]
in which a normal divisor of degree $3$ is contained:
\[
G = \begin{pmatrix}1,&0\\0,&1\end{pmatrix}, \begin{pmatrix}0,&1\\1,&1\end{pmatrix}, \begin{pmatrix}1,&1\\1,&0\end{pmatrix}.
\]

We now set this exceptional case aside. We then have further:

4) The group $G$ contains every substitution of the form
\weq{10}{\begin{pmatrix} 1, & 0 \\ \xi, & 1 \end{pmatrix},}
where $\xi$ is an arbitrary number in $\mathfrak{E}$, except in the case $n=1$, $p=3$.

For by 3), $G$ contains the substitution
\[
\begin{pmatrix} 0, & \varrho \\ -\varrho^{-1}, & 0 \end{pmatrix}\begin{pmatrix} 0, & 1 \\ -1, & 0 \end{pmatrix}\begin{pmatrix} 0, & -\varrho \\ \varrho^{-1}, & 0 \end{pmatrix}\begin{pmatrix} 0, & 1 \\ -1, & 0 \end{pmatrix} = \begin{pmatrix} \varrho^2, & 0 \\ 0, & \varrho^{-2} \end{pmatrix},
\]
if $\varrho$ is an arbitrary number in $\mathfrak{E}$ different from zero. From this it follows that
\weq{11}{\begin{pmatrix} 1, & 0 \\ -\lambda, & 1 \end{pmatrix}\begin{pmatrix} \varrho^{-2}, & 0 \\ 0, & \varrho^2 \end{pmatrix}\begin{pmatrix} 1, & 0 \\ \lambda, & 1 \end{pmatrix}\begin{pmatrix} \varrho^2, & 0 \\ 0, & \varrho^{-2} \end{pmatrix} = \begin{pmatrix} 1, & 0 \\ \lambda(\varrho^4-1), & 1 \end{pmatrix}}
also occurs in $G$ for every arbitrary $\lambda$. If now one can take $\varrho$ different from zero in such a way that $\varrho^4-1$ does not vanish, then for every $\xi$ one can determine the number $\lambda$ from $\lambda(\varrho^4-1)=\xi$, and thus obtains from (11) every substitution of the form (10). In order to judge the possibility of this, we take a primitive root $\gamma$ of $\mathfrak{E}$ and set $\varrho=\gamma^h$.

If one can take the exponent $h$ so that $4h$ is not divisible by $p^n-1$, then $\varrho$ satisfies our requirement. This is always possible when $p=2$, $n>1$, and when $p^n-1>4$, and so only the two cases $n=1$, $p=3$ and $n=1$, $p=5$ still remain doubtful.
% [signature line "Weber, Algebra. II.  17" omitted -- gathering signature, not body text]


% ===== scan leaf 279, printed 258 =====
% !PRINTED 258
\weberpage{258}

In the case $n=1$, $p_{,}=5$, however, we have in $G$ by (2) % [sic] odd subscript comma printed after p
and 3) the substitution

\[
\begin{pmatrix} 1, & 1 \\ 2, & -2 \end{pmatrix}
\begin{pmatrix} 0, & 1 \\ -1, & 0 \end{pmatrix}
\begin{pmatrix} -2, & -1 \\ -2, & 1 \end{pmatrix}
\begin{pmatrix} 0, & 1 \\ -1, & 0 \end{pmatrix}
=
\begin{pmatrix} -2, & 0 \\ 0, & 2 \end{pmatrix},
\]
and consequently also
\weq{12}{
\begin{pmatrix} 1, & 0 \\ -\xi, & 1 \end{pmatrix}
\begin{pmatrix} -2, & 0 \\ 0, & 2 \end{pmatrix}
\begin{pmatrix} 1, & 0 \\ \xi, & 1 \end{pmatrix}
\begin{pmatrix} 2, & 0 \\ 0, & -2 \end{pmatrix}
=
\begin{pmatrix} 1, & 0 \\ -2\xi, & 1 \end{pmatrix},
}
hence, since $\xi$ and consequently also $-2\xi$ is arbitrary, again the substitutions (10) and with them therefore all the generators of the group $E$,
\gesperrt{and so in all these cases $G$ is identical with $E$ and
consequently the group $E$ is simple}.

The case $n=1$, $p=3$, however, presents the second genuine
exception. In this case the group $E$ is of degree $12$
and it has the normal divisor of degree $4$

\[
G =
\begin{pmatrix} 1, & 0 \\ 0, & 1 \end{pmatrix},
\begin{pmatrix} 0, & 1 \\ -1, & 0 \end{pmatrix},
\begin{pmatrix} -1, & 1 \\ 1, & 1 \end{pmatrix},
\begin{pmatrix} 1, & 1 \\ 1, & -1 \end{pmatrix}.
\]

To convince oneself of this, one need only compare the two
subsidiary groups:
\[
\begin{aligned}
\begin{pmatrix} 1, & 0 \\ 1, & 1 \end{pmatrix} G &=
\begin{pmatrix} 1, & 0 \\ 1, & 1 \end{pmatrix},
\begin{pmatrix} 0, & 1 \\ -1, & 1 \end{pmatrix},
\begin{pmatrix} -1, & 1 \\ 0, & -1 \end{pmatrix},
\begin{pmatrix} 1, & 1 \\ -1, & 0 \end{pmatrix} \\
\begin{pmatrix} 1, & 0 \\ -1, & 1 \end{pmatrix} G &=
\begin{pmatrix} 1, & 0 \\ -1, & 1 \end{pmatrix},
\begin{pmatrix} 0, & 1 \\ -1, & -1 \end{pmatrix},
\begin{pmatrix} -1, & 1 \\ -1, & 0 \end{pmatrix},
\begin{pmatrix} 1, & 1 \\ 0, & 1 \end{pmatrix}
\end{aligned}
\]
with $G\begin{pmatrix} 1, & 0 \\ 1, & 1 \end{pmatrix}$, $G\begin{pmatrix} 1, & 0 \\ -1, & 1 \end{pmatrix}$.

With this it is therefore proved in general that, apart from the two
cases $n=1$, $p=2$ and $n=1$, $p=3$, the group $E$
is simple. For the first cases one obtains for the degrees $g$
of these simple groups:

\[
\begin{array}{lll@{\qquad}lll}
n=1, & p=\ 5, & g=\ \ 60 & n=2, & p=2, & g=\ \ 60 \\
     & p=\ 7, & g=\ 168  &      & p=3, & g=\ 360 \\
     & p=11,  & g=\ 660  &      & p=5, & g=7800 \\
     & p=13,  & g=1092   & n=3, & p=2, & g=\ 504 \\
     &        &          &      & p=3, & g=9828 \\
     &        &          & n=4, & p=2, & g=4080^{1)}.
\end{array}
\]

\rule{4cm}{0.4pt}

$^{1)}$ See Bulletin of the New York math. Society. October 1893. In
the report on the mathematical congress in Chicago there are
two notes on these groups by \gesperrt{Cole} and \gesperrt{Moore}.


% ===== scan leaf 280, printed 259 =====
% !PRINTED 259
\weberpage{259}
\pdfbookmark[2]{§ 67. Congruence fields of the second degree}{bm78}

\begin{center}
\S.\,67.

\gesperrt{Congruence Fields of the Second Degree.}
\end{center}

In every congruence field $\mathfrak{E}_{n,p}$ for the modulus $p$ of degree
$n$, the congruence field of the first degree $\mathfrak{E}_{1,p}$, which
consists of the rational numbers taken with respect to the modulus $p$, is
contained as a divisor. Hence in the group of linear
substitutions $E_{n,p}$ there is also contained as a divisor a group which
consists of all substitutions of the form

\weq{1}{A = \begin{pmatrix} a, & b \\ c, & d \end{pmatrix}, \quad ab-cd=1 % [sic] printed ab-cd, not ad-bc
}
where $a, b, c, d$ are integral rational numbers taken with respect to the
modulus $p$.

This group, which appears as the group of the modular equations
in the theory of elliptic functions$^{1)}$, is the
true goal of our considerations. But for the
investigation of this group, and especially for the search for
its divisors, it offers the greatest advantage not to consider this group
independently by itself, but as a divisor of a
group $E_{2,p}$. Namely, if we denote by $N$ any fixed
quadratic non-residue of $p$, then, as already remarked
above, $t^2-N$ is irreducible with respect to the modulus $p$, and we obtain
a congruence field of degree $2$ when, by the equation

\weq{2}{\varepsilon^2 = N}
we introduce a \gesperrt{Galois} imaginary number.

For the sake of clarity it shall be stipulated for the next
considerations that the small Latin letters denote
integral rational numbers (taken with respect to the modulus $p$), which we
here also call \gesperrt{real numbers}, and the Greek letters
numbers of the field $\mathfrak{E}_{2,p}$.

This field $\mathfrak{E}_{2,p}$ has the property, very important for us,
that \gesperrt{in it all real numbers are squares}.

For if $a$ is a quadratic residue of $p$, we can
set $a=x^2$; and if $b$ is a quadratic non-residue, then
$x^2=bN$ can be satisfied, hence $b=(x\varepsilon^{-1})^2$.

\rule{4cm}{0.4pt}

$^{1)}$ Compare the author's book: ``Elliptische Functionen und
algebraische Zahlen''. Braunschweig 1891.


% ===== scan leaf 281, printed 260 =====
% !PRINTED 260
\weberpage{260}

A number of the form $a\varepsilon$ shall be called a \gesperrt{purely imaginary}
\gesperrt{number}, and two numbers of the form $a+b\varepsilon$, $a-b\varepsilon$ \gesperrt{con-}
\gesperrt{jugate imaginary numbers}. The transfer of these designations
from the theory of the ordinary complex numbers
to the numbers of the field $\mathfrak{E}_{2,p}$ is justified by the
agreement of the rules of calculation, and can lead to no misunderstanding,
since in these considerations there is no
question of the ordinary complex numbers.

Since always, when $N$ is a quadratic non-residue,
$N^{1/2(p-1)}=-1$, it follows from (2):

\weq{3}{\varepsilon^p = -\varepsilon.}

If one applies the binomial theorem to the $p$th power of the
binomial $\alpha=a+b\varepsilon$, and observes that all the binomial coefficients,
with the exception of the first and the last, are divisible
by $p$, and hence here are $=0$, and further that by \gesperrt{Fermat's}
theorem $a^p=a$, $b^p=b$, then by (3) we obtain:

\weq{4}{\alpha = a+b\varepsilon, \quad \alpha^p = a-b\varepsilon,}
and consequently, by multiplication:

\weq{5}{\alpha^{p+1} = a^2 - Nb^2,}
from which it follows that $\alpha^{p+1}$ is always a real number.

\gesperrt{If $\gamma$ is a primitive root of the field $\mathfrak{E}_{2,p}$ [\S.\,64, (8)],
then $\gamma^r$ is real when and only when $r$ is divisible by $p+1$,
and $\gamma^{p+1}$ is a primitive root of the
prime $p$ in the ordinary sense}.

\gesperrt{The necessary and sufficient condition
for a number $\alpha$ in $\mathfrak{E}_{2,p}$ to be real consists in the
equation $\alpha^p=\alpha$}.

\gesperrt{Every real number is representable as the $(p+1)$th power of a number
in $\mathfrak{E}_{2,p}$}.

For that $\gamma^r$ is real when $r$ is divisible by $p+1$
follows from (5). But that it can be real in no other way results
from the fact that by Fermat's theorem every real number
different from zero satisfies the condition $a^{p-1}=1$, so that,
if $\gamma^r$ is real, $\gamma^{r(p-1)}=1$ must hold, and therefore
$r(p-1)$ must be divisible by $p^2-1$, and consequently $r$ by $p+1$.
Since, further, $\gamma^{r(p+1)}$ is $=1$ only when $r$ is divisible by $p-1$,
$\gamma^{p+1}$ is a primitive root of $p$.

\gesperrt{The concept of the roots of unity can also be transferred to
the numbers of the field $\mathfrak{E}_{2,p}$}.


% ===== scan leaf 282, printed 261 =====
% !PRINTED 261
\weberpage{261}
\pdfbookmark[2]{§ 68. The real linear congruence group Lp}{bm79}

If $m$ is a divisor of $p^2-1$, and

\[
\varrho = \gamma^{\frac{p^2-1}{m}},
\]
then $\varrho^m=1$ and there is no lower power of $\varrho$ than the $m$th
that becomes equal to 1. We therefore call $\varrho$ a
primitive $m$th root of unity (in $\mathfrak{E}_{2,p}$). All other $m$th
roots of unity are then powers $\varrho^s$ of $\varrho$, and among these
those and only those are primitive for which $s$ is relatively prime to $m$.

This follows from the fact that every number $\xi$ different from zero
is representable in the form $\xi=\gamma^r$, and that $\xi^m=\gamma^{rm}$ is
$=1$ if and only if $rm$ is divisible by $p^2-1$.

Every number different from zero in $\mathfrak{E}_{2,p}$ is a root of unity
 of degree $p^2-1$.

\begin{center}
\S.\,68.

\gesperrt{The real linear congruence group $L_p$.}
\end{center}

We now proceed, equipped with these means, to the
investigation of the real congruence group $L_p$, which consists of all
substitutions of the form

\[
A = \begin{pmatrix} a, & b \\ c, & d \end{pmatrix}
\]
in which $a, b, c, d$ are real integers taken with respect to the modulus $p$,
which satisfy the condition $ad-bc\equiv 1$.
The group $L_p$ is, according to the earlier notation, to be denoted
by $E_{1,p}$, and its degree, if we exclude the case $p=2$, is:

\[
\frac{p(p^2-1)}{2}.
\]

If $A, U$ are any two elements of a group, we have
already earlier called

\[
U^{-1}AU = A'
\]
the element transformed from $A$ by $U$. $A'^2, A'^3, \ldots$
are transformed from $A^2, A^3, \ldots$, from which it follows that all
elements obtained from one another by transformation have the same
degree. If we take the elements $A$ from any
group $G$, then the transformed elements form a group

\[
U^{-1}GU = G',
\]
isomorphic with $G$, which we call the transformed group of $G$.


% ===== scan leaf 283, printed 262 =====
% !PRINTED 262
\weberpage{262}
Now let $A$ be any element of the group $L_p$,
and $U = \begin{pmatrix} \lambda, & \mu \\ \nu, & \varrho \end{pmatrix}$ an element of $E_{2,p}$. We wish to derive from $A$, by
transformation with $U$, a certain normal form.

It shall first be attempted to transform $A$ into the normal form
\weq{1}{S = \begin{pmatrix} \sigma, & 0 \\ 0, & \sigma^{-1} \end{pmatrix}}
From the equation
\[
\begin{pmatrix} a, & b \\ c, & d \end{pmatrix} \begin{pmatrix} \lambda, & \mu \\ \nu, & \varrho \end{pmatrix} = \begin{pmatrix} \lambda, & \mu \\ \nu, & \varrho \end{pmatrix} \begin{pmatrix} \sigma, & 0 \\ 0, & \sigma^{-1} \end{pmatrix}
\]
one obtains
\weq{2}{\begin{aligned}
&(a-\sigma)\,\lambda + b\,\nu = 0, \quad (a-\sigma^{-1})\,\mu + b\,\varrho = 0,\\
&c\,\lambda + (d-\sigma)\,\nu = 0, \quad\ \ c\,\mu + (d-\sigma^{-1})\,\varrho = 0,
\end{aligned}}
and from this it follows that $\sigma$ and $\sigma^{-1}$ must be the roots of the quadratic
equation
\[
(a-\sigma)\,(d-\sigma) - b\,c = 0
\]
or
\weq{3}{\sigma^2 - \sigma\,(a+d) + 1 = 0}
Once $\sigma$ and $\sigma^{-1}$ have been determined from this, one finds
from (2) the ratios $\lambda : \nu$ and $\mu : \varrho$. The quantities $\lambda, \mu, \nu, \varrho$
themselves are then still to be determined so that $\lambda\varrho - \mu\nu = 1$.
This is always possible if the values determined from (2)
do not satisfy the equation $\lambda\varrho = \mu\nu$. If $b$ and $c$ are both $= 0$,
then $A$ already has the form (1). But if one of the two numbers
$b, c$, say $b$, is different from zero, then from (2) we obtain:
\weq{4}{\frac{\nu}{\lambda} = -\frac{(a-\sigma)}{b}, \quad \frac{\varrho}{\mu} = -\frac{(a-\sigma^{-1})}{b},}
and we can have $\lambda\varrho = \mu\nu$, or $\nu : \lambda = \varrho : \mu$, only if
$\sigma = \sigma^{-1}$, that is, $\sigma = \pm 1$, i.\,e. by (3), if $a + d = \pm 2$.

But equation (3) is always solvable in the field $\mathfrak{E}_{2,p}$, because
every \gesperrt{real} number in this field is a square, and it yields
\weq{5}{\begin{aligned}
\sigma &= \frac{a+d}{2} + \sqrt{\left(\frac{a+d}{2}\right)^2 - 1}\\
\sigma^{-1} &= \frac{a+d}{2} - \sqrt{\left(\frac{a+d}{2}\right)^2 - 1}
\end{aligned}}


% ===== scan leaf 284, printed 263 =====
% !PRINTED 263
\weberpage{263}
We thus arrive at the first result:

1.~\gesperrt{A substitution $A$ can always be transformed to the normal form $S$ when $a+d$ is not $=\pm 2$.
$\sigma$ is real or imaginary, according as $(a+d)^2 - 4$ is a quadratic residue or non-residue
of $p$.}

In the case that still remains, where $a+d = \pm 2$, the
normal form $S$ cannot be produced; on the other hand, in
this case we can bring $A$ to another normal form, namely
\weq{6}{T = \begin{pmatrix} 1, & t \\ 0, & 1 \end{pmatrix}}
in which $t$ is real.

To prove this, we may first assume
\weq{7}{a + d = +2}
because the other case $a+d = -2$ is reduced to this one by
changing all the signs of $a, b, c, d$.

But from (7) it follows further that
\weq{8}{a - 1 = -(d-1), \quad (a-1)(d-1) = b\,c.}
If then, first, $b$ or $c = 0$, we have $a = d = 1$
and either
\[
A = \begin{pmatrix} 1, & b \\ 0, & 1 \end{pmatrix},
\]
which already has the normal form $T$, or
\[
A = \begin{pmatrix} 1, & 0 \\ c, & 1 \end{pmatrix} = \begin{pmatrix} 0, & 1 \\ -1, & 0 \end{pmatrix} \begin{pmatrix} 1, & -c \\ 0, & 1 \end{pmatrix} \begin{pmatrix} 0, & -1 \\ 1, & 0 \end{pmatrix},
\]
so that $\begin{pmatrix} 0, & -1 \\ 1, & 0 \end{pmatrix} A \begin{pmatrix} 0, & 1 \\ -1, & 0 \end{pmatrix}$ has the normal form $T$. But if $b$ and $c$ are different from zero,
we obtain from (8) the transformation
\[
\begin{pmatrix} (d-1)\,b^{-1}, & 0, \\ -1, & (a-1)\,c^{-1} \end{pmatrix} \begin{pmatrix} a, & b \\ c, & d \end{pmatrix} \begin{pmatrix} (a-1)\,c^{-1}, & 0, \\ 1, & (d-1)\,b^{-1} \end{pmatrix} = \begin{pmatrix} 1, & -c \\ 0, & 1 \end{pmatrix},
\]
which has the form $T$. With this, then, it is proved:

2.~\gesperrt{A substitution $A$ in which $a+d = \pm 2$ can always be brought
to the normal form $T$ by transformation.}


% ===== scan leaf 285, printed 264 =====
% !PRINTED 264
\weberpage{264}
From this the degrees of the elements of our group can easily be
determined. We have three cases to distinguish here.

I.~\gesperrt{$a + d = \pm 2$.} \gesperrt{Every substitution of this kind can be
transformed into the normal form $T$. But for every
exponent $r$}
\[
T^r = \begin{pmatrix} 1, & r\,t \\ 0, & 1 \end{pmatrix},
\]
\gesperrt{and consequently the degree of this substitution is $p$, except
when $t = 0$, that is, when $T$ is the identical substitution.}

The number of elements of this kind is easy to determine.
If we first take $c = 0$, we can choose $b$ arbitrarily,
and for $a, d$ we obtain from $a + d = \pm 2$, $a\,d = 1$ the two
determinations $a = d = \pm 1$. This gives $2p$ forms, the identical
substitution included. If we next take $a$ and $c$
arbitrary, but $c$ different from zero, which gives $p\,(p-1)$ possibilities,
then from $a\,d - b\,c = 1$ it follows that
\[
b = -c^{-1} + a\,d\,c^{-1},
\]
and for each $a$, $d$ can be determined in two ways. The
total number $2\,p^2$ that we obtain in this way must still be halved, because
here the substitutions occur as two different ones in which
all signs are opposite; and we thus obtain $p^2$ substitutions
of this kind (the identical one included).

II.~\gesperrt{$(a+d)^2 - 4$ a quadratic residue of $p$. Such a
substitution can be transformed into the normal form $S$
with real $\sigma$.}

If we understand by $\gamma$ a primitive root of the field
$\mathfrak{E}_{2,p}$, then
\[
\gamma^{\frac{p^2-1}{2}} = -1,
\]
and we can determine $r$ so that
\[
\sigma = \gamma^{r(p+1)}, \quad \sigma^{\frac{p-1}{2}} = (-1)^r
\]
holds (\S.\,67). By (5) we then obtain
\weq{9}{a + d = \gamma^{r(p+1)} + \gamma^{-r(p+1)},}
and since by (5) the two values $\sigma, \sigma^{-1}$ are, apart from
their order, determined by $a+d$, two values
of the expression (9) become equal or opposite to one another
only when the corresponding values, with positive or
%CONT>>


% ===== scan leaf 286, printed 265 =====
% !PRINTED 265
\weberpage{265}
%<<CONT
taken with negative sign, are congruent to $r$ modulo $\frac{1}{2}(p-1)$.
One therefore obtains all values of this expression that are distinct
from one another and from $\pm 2$ by setting
\weq{10}{r = 1, 2, \cdots \frac{p-3}{2}}
Since here for every exponent $k$
\[
S^k = \begin{pmatrix} \sigma^k, & 0 \\ 0, & \sigma^{-k} \end{pmatrix}
\]
the degree of $S$, and hence also of every substitution of the
II\textsuperscript{nd} kind, is either $\frac{1}{2}(p-1)$ or a divisor of it, according
as $r$ is relatively prime to $\frac{1}{2}(p-1)$ or not.

To determine the number of these substitutions, we proceed
as above. If first $c = 0$, it follows that
\[
a = \gamma^{\pm r(p+1)}, \quad d = \gamma^{\mp r(p+1)},
\]
and $b$ can have $p$ values. The number $r$ takes the values (10),
and accordingly this yields $p\,(p-3)$ substitutions. If next
$c$ is different from zero, then $p\,(p-1)$ is the number of the
distinct assumptions for $a, c$. If $a$ is assumed, we can
determine $d$ from (9) and set
\[
b = -c^{-1} + a\,d\,c^{-1},
\]
from which we obtain $\frac{1}{2}\,p\,(p-1)(p-3)$ determinations,
and so, together with those counted for $c = 0$, $\frac{1}{2}\,p\,(p-3)(p+1)$.
This number too is still to be halved, so that we obtain
\[
\frac{1}{4}\,p\,(p-3)(p+1)
\]
substitutions of the II\textsuperscript{nd} kind.

III.~$(a+d)^2 - 4$ a quadratic non-residue of $p$. In this
case $\sigma$ is not real, but conjugate to its reciprocal value.
Hence by \S.\,67, (4) $\sigma^p = \sigma^{-1}$, and consequently
\[
\sigma^{\frac{p+1}{2}} = \pm 1.
\]
From this it follows that the degree of such a substitution
is $\frac{1}{2}(p+1)$ or a divisor of this number.
In this third case we set
\weq{11}{\begin{aligned}
&\sigma = \gamma^{r(p-1)}, \quad a+d = \gamma^{r(p-1)} + \gamma^{-r(p-1)}\\
&r = 1, 2, \ldots, \tfrac{1}{2}(p-1).
\end{aligned}}
Here the case $c = 0$ or $b = 0$ cannot occur, because
from $a\,d = 1$ it would follow that $a = \sigma^{\pm 1}$, $d = \sigma^{\mp 1}$
%CONT>>


% ===== scan leaf 287, printed 266 =====
% !PRINTED 266
\weberpage{266}
\pdfbookmark[2]{§ 69. Imaginary form of the group Lp}{bm80}
%<<CONT
would have to be, and $a$ and $d$ would not come out real. Accordingly
we take for $a$ and $c$, as above, the $p\,(p-1)$ different
systems of values, and obtain from (11) for each of them
$1/2\,(p-1)$ determinations of $b$ and $d$; this number too is to be
halved, and there result
\[
1/4\;p\,(p-1)\,(p-1)
\]
substitutions of the III\textsuperscript{rd} kind. The total number of all substitutions
 of the group $L_p$ is accordingly, as it must be,
\[
p^2 + 1/4\,p\,(p-3)\,(p+1) + 1/4\,p\,(p-1)^2 = 1/2\,p\,(p^2-1).
\]

\begin{center}
§. 69.

\gesperrt{Imaginary form of the group $L_p$.}
\end{center}

The substitutions of the first two kinds of the group $L_p$
have the property that the normal form $T$ or $S$ corresponding
 to them is likewise contained in $L_p$, and that one can
transform them into these normal forms by \gesperrt{real} substitutions
, that is, by substitutions which themselves occur in $L_p$
. Neither of these holds for the substitutions of the third
kind.

One can, however, as we shall now prove, transform the whole
group $L_p$ into another group $\Lambda_p$ in such a way that $\Lambda_p$ is
\gesperrt{a divisor of the group $E_{2,p}$ isomorphic with $L_p$}, that the
normal forms of the transformations of the third kind are contained
in $\Lambda_p$, and that every substitution of the third kind in $\Lambda_p$ can
be transformed into the normal form by substitutions of
$\Lambda_p$ itself.

In order to find such a transformation of $L_p$, we consider
the substitution
\weq{1}{S = \begin{pmatrix} \sigma, & 0 \\ 0, & \sigma^{-1} \end{pmatrix}}
under the assumption that $\sigma$ and $\sigma^{-1}$ are conjugate imaginary
, and seek to determine the substitution
\weq{2}{R = \begin{pmatrix} \lambda, & \mu \\ \nu, & \varrho \end{pmatrix}}
in such a way that
\weq{3}{R\,S\,R^{-1}}
%CONT>>


% ===== scan leaf 288, printed 267 =====
% !PRINTED 267
\weberpage{267}
%<<CONT
becomes real. But by (1) and (2) we obtain
\[
R\,S\,R^{-1} = \begin{pmatrix} \lambda\,\varrho\,\sigma - \mu\,\nu\,\sigma^{-1}, & -\lambda\,\mu\,(\sigma-\sigma^{-1}) \\ \varrho\,\nu\,(\sigma-\sigma^{-1}), & -\mu\,\nu\,\sigma+\varrho\,\lambda\,\sigma^{-1} \end{pmatrix},
\]
and this becomes real when $\lambda\mu$ and $\varrho\nu$ are purely imaginary, and $\lambda\varrho$ and
$-\mu\nu$ are conjugate imaginary. These conditions can be satisfied
in several ways: most simply, no doubt, and at the same time
in such a way that the determinant of $R = 1$, if one sets
\weq{4}{R = \begin{pmatrix} \frac{1}{2}, & -\frac{1}{2}\,\varepsilon^{-1} \\ \varepsilon, & 1 \end{pmatrix}}
If, on the other hand,
\[
A = \begin{pmatrix} a, & b \\ c, & d \end{pmatrix}
\]
is an arbitrary real substitution from $L_p$, then by (4),
since $\varepsilon^2 = N$, we obtain:
\[
R^{-1}\,A\,R =
\]
\weq{5}{\begin{pmatrix} \dfrac{a+d}{2}+b\,\varepsilon+\dfrac{c}{4}\,\varepsilon^{-1}, & -\dfrac{a-d}{2}\,\varepsilon^{-1}+b-\dfrac{c}{4}\,N^{-1} \\[1.2ex] -N\left(\dfrac{a-d}{2}\,\varepsilon^{-1}+b-\dfrac{c}{4}\,N^{-1}\right), & \dfrac{a+d}{2}-b\,\varepsilon-\dfrac{c}{4}\,\varepsilon^{-1} \end{pmatrix},}
for which we may also set
\weq{6}{A = \begin{pmatrix} \alpha, & \beta \\ -N\beta', & \alpha' \end{pmatrix} = \begin{pmatrix} \alpha, & \beta \\ -N\beta^p, & \alpha^p \end{pmatrix}, \quad \alpha\alpha'+N\beta\beta' = 1}
in which $\alpha,\alpha'$ and $\beta,\beta'$ are two conjugate imaginary pairs,
so that by §.\,65, (4) one may set $\alpha' = \alpha^p$, $\beta' = \beta^p$.

If, conversely, we take any substitution $A$ of the
form (6), then
\[
A = R\,A\,R^{-1} =
\]
\weq{7}{\begin{pmatrix} \dfrac{\alpha+\alpha'}{2}+\varepsilon\,\dfrac{\beta'-\beta}{2}, & \dfrac{\beta+\beta'}{4}+\varepsilon^{-1}\,\dfrac{\alpha-\alpha'}{4} \\[1.2ex] -N(\beta+\beta')+\varepsilon(\alpha-\alpha'), & \dfrac{\alpha+\alpha'}{2}-\varepsilon\,\dfrac{\beta'-\beta}{2} \end{pmatrix}}
results as a \gesperrt{real} substitution from $L_p$.

\gesperrt{From this it follows that, when $A$ runs through the whole group $L_p$,
the substitution $A$ determined by (5) and (6) runs through a group
isomorphic with $L_p$, which we denote by $\Lambda_p$}.


% ===== scan leaf 289, printed 268 =====
% !PRINTED 268
\weberpage{268}
In the group $\Lambda_p$ the normal form $S$ for the substitutions
of the third kind is contained.

If we now assume a substitution $A$ of the third kind in
the group $\Lambda_p$, then, as shall now be proved, we can transform
it into the normal form by transformation with a substitution
$U$ which itself belongs to the group $\Lambda_p$. For if we set
\weq{8}{A = \begin{pmatrix} \alpha, & \beta \\ -N\beta', & \alpha' \end{pmatrix}, \quad U = \begin{pmatrix} \lambda, & \mu \\ -N\mu', & \lambda' \end{pmatrix},}
\weq{9}{\alpha\alpha'+N\beta\beta' = 1, \qquad \lambda\lambda'+N\mu\mu' = 1,}
in which $\lambda,\lambda'$ and $\mu,\mu'$ are conjugate pairs, the
equations (3), (4), §.\,68 take the form:
\weq{10}{(\sigma-\alpha)(\sigma-\alpha')+N\beta\beta' = \sigma^2-\sigma(\alpha+\alpha')+1 = 0,}
\weq{11}{N\,\dfrac{\mu'}{\lambda} = \dfrac{\alpha-\sigma}{\beta}, \qquad \dfrac{\lambda'}{\mu} = -\dfrac{\alpha-\sigma^{-1}}{\beta},}
and the equation (10), since $A$ is to be of the third kind,
must have two roots reciprocal to one another and conjugate
imaginary.

But from the second equation (11) it follows, by passing
to the conjugate imaginary numbers, that
\[
\dfrac{\lambda}{\mu'} = -\dfrac{\alpha'-\sigma}{\beta'},
\]
and by (10) this is a consequence of the first equation (11).
If we therefore set
\[
N\mu' = \varkappa\,(\alpha-\sigma), \qquad N\mu = \varkappa'\,(\alpha'-\sigma^{-1})
\]
\[
\lambda = \varkappa\beta, \qquad \lambda' = \varkappa'\beta',
\]
in which $\varkappa,\varkappa'$ are two conjugate imaginary numbers, the
equations (11) are satisfied, and one can then determine $\varkappa\varkappa'$
so that $\lambda\lambda'+N\mu\mu' = 1$.

From this there follows yet another important result. If we
denote by $\gamma$ a primitive root of the field $\mathfrak{E}_{2,p}$, then for
the substitutions of the normal form for the second and
third kind we may assume the expressions (§.\,68):
\[
S_2^r = \begin{pmatrix} \gamma^{(p+1)r}, & 0 \\ 0, & \gamma^{-(p+1)r} \end{pmatrix}, \qquad S_3^r = \begin{pmatrix} \gamma^{(p-1)r}, & 0 \\ 0, & \gamma^{-(p-1)r} \end{pmatrix}.
\]

If now $A$ is a substitution of the second kind from $L_p$ and $A$ a
%CONT>>


% ===== scan leaf 290, printed 269 =====
% !PRINTED 269
\weberpage{269}
%<<CONT
substitution of the third kind from $\Lambda_p$, then we can determine
the substitutions $U$, $V$ from $L_p$ or from $\Lambda_p$ in such a way that
\[
A = U\,S_2^r\,U^{-1}, \qquad A = V\,S_3^r\,V^{-1}.
\]
If, then, corresponding to the value $r=1$, we set
\[
A_1 = U\,S_2\,U^{-1}, \qquad A_1 = V\,S_3\,V^{-1},
\]
it follows that
\[
A = A_1^r, \qquad A = A_1^r,
\]
and $A_1$ occurs in $L_p$, $A_1$ in $\Lambda_p$.

If, further, $B$ is a substitution of the first kind from $L_p$, then,
setting
\[
T_1 = \begin{pmatrix} 1, & 1 \\ 0, & 1 \end{pmatrix}
\]
one can again determine $U$ from $L_p$ in such a way that
\[
B = U\,T_1^r\,U^{-1} = B_1^r
\]
holds. From this results the theorem:

\gesperrt{In the group $L_p$ (or $\Lambda_p$) the substitutions
of the first, second and third kinds can, with the inclusion
of the identity, be arranged in cycles of $p$, $\dfrac{p-1}{2}$, $\dfrac{p+1}{2}$
terms.}

Two such cycles can have no term in common besides the
unit element.

This is at once evident for the substitutions of the first
kind, whose degree is equal to $p$; for with these the whole
cycle can be derived by raising to powers from any element
different from unity.

But if a common element occurs in two cycles of the
second or third kind, then the group can be so transformed
that the two cycles take the form:
\weq{12}{\begin{aligned} &1, \quad U\,S\,U^{-1}, \quad U\,S^2\,U^{-1}, \ \ldots \\ &1, \quad S, \quad S^2, \ \ldots, \end{aligned}}
in which $S = \begin{pmatrix} \sigma, & 0 \\ 0, & \sigma^{-1} \end{pmatrix}$ has the normal form.

If now, for any two exponents
$r,t$ different from zero,
\[
S^r = U\,S^t\,U^{-1},
\]
%CONT>>


% ===== scan leaf 291, printed 270 =====
% !PRINTED 270
\weberpage{270}
%<<CONT
then, if $U = \begin{pmatrix}\lambda, & \mu\\ \nu, & \varrho\end{pmatrix}$ is set, there results as condition
\[
\begin{aligned}
\sigma^r\lambda &= \pm\sigma^t\lambda, & \sigma^r\mu &= \pm\sigma^{-t}\mu\\
\sigma^{-r}\nu &= \pm\sigma^t\nu, & \sigma^{-r}\varrho &= \pm\sigma^{-t}\varrho,
\end{aligned}
\]
where the same sign must hold throughout. From this it follows that either
\[
\sigma^r = \pm\sigma^t, \quad \nu = 0, \quad \mu = 0, \quad \lambda\varrho = 1
\]
or
\[
\sigma^r = \pm\sigma^{-t}, \quad \lambda = 0, \quad \varrho = 0, \quad \mu\nu = -1
\]
must be the case, and then we obtain:
\[
USU^{-1} = S \quad \text{or} \quad = S^{-1},
\]
and in both cases, therefore, the two cycles (12) agree completely with one another.

According to this, from the numbers at the close of \S.\,68 there results for the number of cycles of the first, second and third kind:
\[
p+1, \quad \frac{p(p+1)}{2}, \quad \frac{p(p-1)}{2}.
\]

It is expedient to consider, beside the groups $L_p$ and $\Lambda_p$, yet a third group in $E_{2,p}$, which is isomorphic with these, even though it cannot be reduced to them by transformation.

Since $N$ is real, we can by \S.\,67 determine a number $\nu$ in $\mathfrak{E}_{2,p}$ so that $N = \nu^{p+1}$. We now let there correspond to the substitution
\[
A = \begin{pmatrix}\alpha, & \beta\\ -N\beta^p, & \alpha^p\end{pmatrix}
\]
from $\Lambda_p$ a substitution $B$, which is formed thus:
\weq{13}{B = \begin{pmatrix}\alpha, & \nu\beta\\ -\nu^p\beta^p, & \alpha^p\end{pmatrix}, \quad \alpha^{p+1} + N\beta^{p+1} = 1.}
If we compose two substitutions $B$, $B_1$ of the form (13), we obtain
\[
\begin{aligned}
BB_1 &= \begin{pmatrix}\alpha, & \nu\beta\\ -\nu^p\beta^p, & \alpha^p\end{pmatrix}\begin{pmatrix}\alpha_1, & \nu\beta_1\\ -\nu^p\beta_1^p, & \alpha_1^p\end{pmatrix}\\
&= \begin{pmatrix}\alpha\alpha_1 - N\beta\beta_1^p, & \nu(\alpha\beta_1+\beta\alpha_1^p)\\ -\nu^p(\beta^p\alpha_1+\alpha^p\beta_1^p), & \alpha^p\alpha_1^p - N\beta^p\beta_1\end{pmatrix};
\end{aligned}
\]
in these,

$\alpha\alpha_1 - N\beta\beta_1^p, \ \alpha^p\alpha_1^p - N\beta^p\beta_1$ and $\alpha\beta_1+\beta\alpha_1^p, \ \beta^p\alpha_1+\alpha^p\beta_1^p$
%CONT>>


% ===== scan leaf 292, printed 271 =====
% !PRINTED 271
\weberpage{271}
\pdfbookmark[2]{§ 70. Divisors of the group Lp whose degree is divisible by p}{bm81}
%<<CONT
two conjugate pairs, and therefore $BB_1$ also has the form (13) and stands in the same relation to $AA_1$ as $B$ to $A$ and $B_1$ to $A_1$.

\gesperrt{From this it follows that the totality of the substitutions $B$ forms a group isomorphic with $\Lambda_p$, and this group we shall denote by $\Gamma_p$.}

If we put $\beta$ in place of $\nu\beta$ and denote by $\alpha', \beta'$ the quantities conjugate to $\alpha, \beta$, then we may also take the substitutions of the group $\Gamma_p$ in the simpler form:
\weq{14}{B = \begin{pmatrix}\alpha, & \beta\\ -\beta', & \alpha'\end{pmatrix}, \quad \alpha\alpha' + \beta\beta' = 1.}
In order to pass from one of these substitutions $B$ to the corresponding real substitution $A$, one first derives from (14) the corresponding substitution $A$:
\weq{15}{A = \begin{pmatrix}\alpha, & \nu^{-1}\beta\\ -\nu\beta', & \alpha'\end{pmatrix},}
and from it forms, according to (3) and (7):
\weq{16}{A = RAR^{-1}.}
In spite of the isomorphism, the group $\Gamma_p$ cannot be transformed into $L_p$ or $\Lambda_p$, at least not by substitutions whose numbers belong to the field $\mathfrak{E}_{2,p}$. To carry out the transformation, one would have to go into a higher field, in which \gesperrt{all} numbers of $\mathfrak{E}_{2,p}$ are squares.

\begin{center}
\S.\,70.\\
\gesperrt{Divisors of the group $L_p$ whose degree is}\\
\gesperrt{divisible by $p$.}
\end{center}

It is a problem of the greatest interest to determine all divisors of the group $L_p$. On the solution of this problem depends in what way the group $L_p$ can be represented by permutations of figures, and hence in what algebraic equations these groups can occur (\S.\,28, 2.).

The cyclic groups contained in $L_p$ we have already considered in the preceding sections. We have seen that the degree of a cyclic group is either equal to $p$ or a divisor of $1/2\,(p-1)$ or of $1/2\,(p+1)$
%CONT>>


% ===== scan leaf 293, printed 272 =====
% !PRINTED 272
\weberpage{272}
%<<CONT
is, and every divisor of these numbers also occurs among the degrees of the cyclic groups. The index of a cyclic divisor can never be smaller than $1/2\,(p^2-1)$, $p(p+1)$, $p(p-1)$.

If the degree of a divisor $G$ of $L_p$ is divisible by $p$, then it contains a cyclic group of degree $p$, and by \S.\,68 we can transform the group $G$ in such a way that this cyclic group consists of the substitutions
\[
T = \begin{pmatrix}1, & t\\ 0, & 1\end{pmatrix}, \quad t = 0, 1, 2, \ldots, p-1
\]
If now $G$ contains a further substitution
\[
A = \begin{pmatrix}a, & b\\ c, & d\end{pmatrix}
\]
in which $c$ is different from zero, then
\weq{1}{\begin{pmatrix}1, & -ac^{-1}\\ 0, & 1\end{pmatrix}\begin{pmatrix}a, & b\\ c, & d\end{pmatrix}\begin{pmatrix}1, & -dc^{-1}\\ 0, & 1\end{pmatrix} = \begin{pmatrix}0, & -c^{-1}\\ c, & 0\end{pmatrix}}
also occurs in it. From this it follows further that
\[
\begin{pmatrix}0, & -c^{-1}\\ c, & 0\end{pmatrix}\begin{pmatrix}1, & t\\ 0, & 1\end{pmatrix}\begin{pmatrix}0, & -c^{-1}\\ c, & 0\end{pmatrix} = \begin{pmatrix}1, & 0\\ -c^2 t, & 1\end{pmatrix},
\]
and consequently every substitution
\[
U = \begin{pmatrix}1, & 0\\ u, & 1\end{pmatrix}, \quad u = 0, 1, 2, \ldots, p-1
\]
also occurs. If one puts $U$ for $u = -1$ in place of $A$ in the formula (1), it follows that $G$ also contains the substitution $\begin{pmatrix}0, & 1\\ -1, & 0\end{pmatrix}$, which together with $U$ gives a system of generating elements of $L_p$ (\S.\,65), so that $G$ is therefore identical with $L_p$. From this it follows that $G$, if it is not to be identical with $L_p$, can contain only substitutions of the form
\[
A = \begin{pmatrix}a, & b\\ 0, & a^{-1}\end{pmatrix}
\]
Conversely, all these substitutions form a group of degree $1/2\,p(p-1)$, and hence a divisor of index $p+1$.

Among these there are also contained certain divisors of greater index, which one obtains by taking for $a$ not all values, but all values of the form $a^s$, where $s$ is a divisor of $1/2\,(p-1)$. The degree of such a group is $1/2\,p(p-1) : s$.


% ===== scan leaf 294, printed 273 =====
% !PRINTED 273
\weberpage{273}
\pdfbookmark[2]{§ 71. Divisors of the group Lp whose degree is not divisible by p}{bm82}
If we conceive the group $L_p$ as the group of the linear substitutions
\[
\eta = \frac{a\xi+b}{c\xi+d}
\]
and let $\xi, \eta$ in it run through the numbers $\infty, 0, 1, \ldots, p-1$ according to \S.\,65, then $G$ gives us the group of the integral linear substitutions
\[
\eta = a^2\xi + ab.
\]
These are the substitutions that leave the figure $\infty$ unchanged, and hence yield a permutation group of only $p$ figures. These are the same special metacyclic groups that we considered in the seventeenth chapter of the first volume.

The various groups transformed from $G$ leave some one of the remaining $p+1$ indices unchanged, and the total number of these groups is $p+1$.

\begin{center}
\S.\,71.\\
\gesperrt{Divisors of the group $L_p$ whose degree is not divisible}\\
\gesperrt{by $p$.}
\end{center}

To find the remaining divisors of $L_p$ whose degree is not divisible by $p$, we can go, step by step, the same way that led us in the seventh and eighth chapters to the determination of the polyhedral groups. We may content ourselves here with bringing out the principal points of the derivation, and refer for the proofs to the place named, where entirely the same conclusions were to be drawn.

For this it is required to conceive the group $L_p$, as set forth in \S.\,65, as a group of linear fractional substitutions
\weq{1}{\Theta(\xi) = \frac{a\xi+b}{c\xi+d}}
and to assign to the variable $\xi$ in it all $p^2+1$ numerical values of the congruence field $\mathfrak{E}_{2,p}$, including $\infty$. We thereby gain the advantage that the equation
\weq{2}{\xi = \frac{a\xi+b}{c\xi+d}}
%CONT>>


% ===== scan leaf 295, printed 274 =====
% !PRINTED 274
\weberpage{274}
%<<CONT
always has two roots. These two roots are distinct from one another
if we exclude substitutions of degree $p$, which indeed cannot occur
in a group whose degree is not divisible by $p$, and which by
\S.\,68, 1.\ are the only ones for which the two roots of (2)
coincide.

These two roots we call the poles of $\Theta$.

We investigate a group $G$ whose degree $n$ is not divisible by $p$,
and which is to be a divisor of $L_p$. If in $G$ the substitutions
\[
\Theta_1,\ \Theta_2,\ \ldots,\ \Theta_{\nu-1},
\]
occur, but no others, which have the same pole $\alpha$, then we
call $\alpha$ \gesperrt{a $\nu$-fold pole} (\S.\,52).

Now if $S$ is any substitution of the group $\mathfrak{E}_{2,p}$, then,
if we transform $L_p$ by $S^{-1}$, we obtain a group $S L_p S^{-1}$
isomorphic with $L_p$, and $G$ passes over, by the same
transformation, into an isomorphic group $S G S^{-1}$.

Although the substitutions of this latter group have no real
coefficients, still each of them has two poles, which one obtains
by applying the substitution $S$ to the poles of $\Theta$. For, if
$\alpha$ is a pole of $\Theta$, we have:
\[
S\,\Theta\,S^{-1}\,S(\alpha) = S\,\Theta(\alpha) = S(\alpha).
\]
We therefore have, as in \S.\,51, 2.:

\gesperrt{1. The group $G$ can be so transformed that an arbitrary
one of its substitutions receives the poles 0 and $\infty$, so that
this substitution becomes multiplicative.}

We now enumerate in order the theorems of \S.\,52, so far as they
come into question here, and shall add an elaboration only there
where the altered hypotheses require it:

\gesperrt{2. If $\alpha$ is a $\nu$-fold pole, the substitutions 1,
$\Theta_1$, $\Theta_2$, . . ., $\Theta_{\nu-1}$ form a cyclic
group. Their substitutions can all (by transformation) be put into
the form:}
\weq{3}{\Theta(\xi) = \gamma^{h\frac{p^2-1}{\nu}}\,\xi, \quad h = 0, 1, \ldots, \nu-1 \quad (\S.\,52, 3.).}


% ===== scan leaf 296, printed 275 =====
% !PRINTED 275
\weberpage{275}
Here $\gamma$ denotes a primitive root, hence $\gamma^{\frac{p^2-1}{\nu}}$
a primitive $\nu$th root of unity of the congruence field
$\mathfrak{E}_{2,p}$ (\S.\,67).

3. \gesperrt{Both poles of a substitution $\Theta$ are of equal multiplicity}
(\S.\,52, 4.).

If $Q$ is the cyclic group of degree $\nu$ of the powers
of $\Theta$, and $n = \nu\mu$, then one obtains
\[
G = Q + \psi_1 Q + \cdots + \psi_{\mu-1} Q,
\]
where $\psi_1, \ldots, \psi_{\mu-1}$ are substitutions from $G$, and

4. \gesperrt{the numbers}
\weq{4}{\alpha,\ \psi_1(\alpha) = \alpha_1,\ \psi_2(\alpha) = \alpha_2,\ \ldots,\ \psi_{\mu-1}(\alpha) = \alpha_{\mu-1}}
\gesperrt{form a system of conjugate poles of equal multiplicity of the group $G$.}

All the poles of the group $G$ can thus be gathered into systems
of conjugate poles, and in this way we obtain the indeterminate
equation
\weq{5}{\begin{aligned}
2n - 2 &= \mu(\nu-1) + \mu'(\nu'-1) + \mu''(\nu''-1) \ldots \\
&= nh - \mu - \mu' - \mu'' \ldots,
\end{aligned}}
where $h$ is the number of systems of conjugate poles; and of this
equation we have seen in \S.\,52 that it admits only a limited
number of solutions.

In order to find the divisors of the congruence groups corresponding
to these solutions, we may simply put, in the formulae of the
polyhedral groups set up in the eighth chapter, for the roots of
unity occurring there the roots of unity of the field
$\mathfrak{E}_{2,p}$ defined in \S.\,67, with which indeed
calculation proceeds by the same rules. In this, of course, only such
roots of unity can occur whose degree is a divisor of $p^2 - 1$.

Finally, for every such group it remains to be investigated whether
its substitutions are contained in $L_p$ or in a group isomorphic
with it. We then have the following cases of solutions of
equation (5), where $\gamma$ always denotes a primitive root of the field
$\mathfrak{E}_{2,p}$ (\S.\,55).

\gesperrt{I. Cyclic groups of degree $n$,}\quad
$\begin{aligned} h &= 2,\ \nu = \nu' = n \\ \mu &= \mu' = 1 \end{aligned}$
\[
C_n = \begin{pmatrix} \gamma^{\frac{p^2-1}{2n}\lambda}, & 0 \\[2pt] 0, & \gamma^{-\frac{(p^2-1)}{2n}\lambda} \end{pmatrix},
\quad \lambda = 0, 1, \ldots, n-1.
\]


% ===== scan leaf 297, printed 276 =====
% !PRINTED 276
\weberpage{276}
This group is real when $n$ is a divisor of $1/2\,(p-1)$, because
then and only then the exponents of $\gamma$ are divisible by $p+1$.
The group is contained in $\Gamma_p$ and at the same time in $\Lambda_p$ (\S.\,69)
when $\gamma^{\frac{p^2-1}{2n}}$, $\gamma^{-\frac{p^2-1}{2n}}$
are conjugate, hence when
\[
\gamma^{-\frac{p^2-1}{2n}} = \gamma^{p\frac{p^2-1}{2n}}
\quad \text{or} \quad
\gamma^{(p+1)\frac{p^2-1}{2n}} = 1
\]
that is, when $n$ is a divisor of $1/2\,(p+1)$. This agrees with
\S.\,68, according to which no cyclic groups occur other than those
whose degree is equal to $p$ or a divisor of $1/2\,(p-1)$ or $1/2\,(p+1)$.
In the remaining polyhedral groups $h = 3$, and
we have:

II. \gesperrt{Dihedral groups} of degree $2m$,\quad
$\begin{aligned} \nu &= \nu' = 2,\ \nu'' = m \\ \mu &= \mu' = m,\ \mu'' = 2 \end{aligned}$
\[
D_m = \begin{pmatrix} \gamma^{\frac{p^2-1}{2m}\lambda}, & 0 \\[2pt] 0, & \gamma^{-\frac{p^2-1}{2m}\lambda} \end{pmatrix},
\begin{pmatrix} 0, & \gamma^{\frac{p^2-1}{2m}\lambda} \\[2pt] -\gamma^{-\frac{p^2-1}{2m}\lambda}, & 0 \end{pmatrix},
\quad \lambda = 0, 1, \ldots, m-1.
\]
This group is contained in $L_p$ when $p \equiv 1 \pmod{2m}$, and
in $\Gamma_p$ when $p \equiv -1 \pmod{2m}$, as in Case I.

Here the second representation of the dihedral group given in \S.\,55
has been chosen, in which the distinction of the cases becomes somewhat
simpler than with the first.

III. \gesperrt{Tetrahedral group,}\quad
$\begin{aligned} \nu &= 2,\ \nu' = 3,\ \nu'' = 3,\ n = 12 \\ \mu &= 6,\ \mu' = 4,\ \mu'' = 4. \end{aligned}$

To represent this group, we observe that in the field
$\mathfrak{E}_{2,p}$ an $8$th root of unity
$\gamma^{\frac{p^2-1}{8}}$ always exists. We thus obtain, by \S.\,56, the
three generating substitutions
\[
\Theta = \begin{pmatrix} \gamma^{\frac{p^2-1}{4}}, & 0 \\[2pt] 0, & \gamma^{-\frac{p^2-1}{4}} \end{pmatrix},
\quad
\psi = \begin{pmatrix} 0, & \gamma^{\frac{p^2-1}{4}} \\[2pt] -\gamma^{-\frac{p^2-1}{4}}, & 0 \end{pmatrix},
\quad
\psi_1 = \psi\,\Theta = \begin{pmatrix} 0, & 1 \\ -1, & 0 \end{pmatrix}
\]
\[
\chi = \begin{pmatrix} \dfrac{1}{\sqrt{-2}}\,\gamma^{\frac{p^2-1}{8}}, & \dfrac{1}{\sqrt{-2}}\,\gamma^{\frac{p^2-1}{8}} \\[8pt] \dfrac{1}{\sqrt{-2}}\,\gamma^{-\frac{p^2-1}{8}}, & -\dfrac{1}{\sqrt{-2}}\,\gamma^{-\frac{p^2-1}{8}} \end{pmatrix}.
\]
1) If $p \equiv 1 \pmod 4$, then $\gamma^{\frac{p^2-1}{4}}$ is real,
and $\gamma^{\frac{p^2-1}{8}}$ is real or purely imaginary according
as $p \equiv 1$ or $\equiv 5 \pmod 8$.


% ===== scan leaf 298, printed 277 =====
% !PRINTED 277
\weberpage{277}
In agreement with this, $\sqrt{-2}$ is likewise real or purely imaginary
(Bd.\,I, \S.\,138, 4., 6.), and consequently in these cases $\Theta$,
$\psi$, $\chi$, and hence the whole tetrahedral group, are real.

2) But if $p \equiv 3 \pmod 4$, then $\gamma^{\frac{p^2-1}{4}}$,
$\gamma^{-\frac{p^2-1}{4}}$ are conjugate imaginary, because
$\gamma^{(p+1)\frac{p^2-1}{4}} = 1$.

If $p \equiv 3 \pmod 8$, then $\sqrt{-2}$ is real and
$\gamma^{\frac{p^2-1}{8}(p+1)} = -1$, hence $\gamma^{\frac{p^2-1}{8}}$
and $-\gamma^{\frac{p^2-1}{8}}$ are conjugate imaginary; and if $p \equiv 7
\pmod 8$, then $\sqrt{-2}$ is purely imaginary, $\gamma^{\frac{p^2-1}{8}}$
and $\gamma^{-\frac{p^2-1}{8}}$ are conjugate imaginary, and consequently
$\Theta$, $\psi_1$, $\chi$ belong in these cases to the group $\Gamma_p$.
There is thus in \gesperrt{all cases in the group} $L_p$ a tetrahedral group.

For $p = 5$, for example, since $-2$ is a quadratic non-residue of 5,
one may take $\varepsilon = \sqrt{-2}$, and, since $\pm 2$ are the two
primitive roots of 5, say
\[
\gamma^6 = -2, \quad \gamma^3 = \sqrt{-2}, \quad \gamma = 1 + \sqrt{-2}
\]
then it follows that
\weq{6}{\begin{aligned}
&\Theta = \begin{pmatrix} 2, & 0 \\ 0, & 3 \end{pmatrix},\quad
\psi = \begin{pmatrix} 0, & 2 \\ 2, & 0 \end{pmatrix},\quad
\psi_1 = \begin{pmatrix} 0, & 1 \\ -1, & 0 \end{pmatrix} \\
&\chi = \begin{pmatrix} 1, & 1 \\ 2, & -2 \end{pmatrix},\quad
\chi^2 = \begin{pmatrix} 2, & 1 \\ 2, & -1 \end{pmatrix},
\end{aligned}}
and from this results the tetrahedral group sought for $p = 5$
[\S.\,56, (3)]:
\weq{7}{\begin{aligned}
&\begin{pmatrix} 1, & 0 \\ 0, & 1 \end{pmatrix},
\begin{pmatrix} 2, & 0 \\ 0, & -2 \end{pmatrix},
\begin{pmatrix} 0, & 2 \\ 2, & 0 \end{pmatrix},
\begin{pmatrix} 0, & 1 \\ -1, & 0 \end{pmatrix} \\
&\begin{pmatrix} 1, & 1 \\ 2, & -2 \end{pmatrix},
\begin{pmatrix} 2, & 2 \\ 1, & -1 \end{pmatrix},
\begin{pmatrix} -1, & 1 \\ 2, & 2 \end{pmatrix},
\begin{pmatrix} -2, & 2 \\ 1, & 1 \end{pmatrix} \\
&\begin{pmatrix} 2, & 1 \\ 2, & -1 \end{pmatrix},
\begin{pmatrix} -1, & 2 \\ 1, & 2 \end{pmatrix},
\begin{pmatrix} 1, & 2 \\ 1, & -2 \end{pmatrix},
\begin{pmatrix} -2, & 1 \\ 2, & 1 \end{pmatrix}.
\end{aligned}}
This group is a divisor of $L_5$ of index 5.

IV. \gesperrt{Octahedral group.}

In the octahedral group a cyclic group of degree 4 is contained, and
such a group can therefore exist only when $1/2\,(p-1)$ or
$1/2\,(p+1)$ is divisible by 4, that is, when $p \equiv \pm 1 \pmod 8$.


% ===== scan leaf 299, printed 278 =====
% !PRINTED 278
\weberpage{278}
Under this assumption the formulae of \S.\,57 lead to the setting up
of an octahedral group.

The generating substitutions are:
\[
\Theta = \begin{pmatrix} \gamma^{\frac{p^2-1}{8}}, & 0 \\ 0, & \gamma^{-\frac{p^2-1}{8}} \end{pmatrix}, \quad
\psi = \begin{pmatrix} 0, & \gamma^{\frac{p^2-1}{4}} \\ \gamma^{\frac{p^2-1}{4}}, & 0 \end{pmatrix}
\]
\weq{8}{
\chi = \begin{pmatrix} \dfrac{1}{\sqrt{2}}\,\gamma^{-\frac{p^2-1}{8}}, & \dfrac{1}{\sqrt{2}}\,\gamma^{-\frac{p^2-1}{8}} \\[2mm] \dfrac{-1}{\sqrt{2}}\,\gamma^{\frac{p^2-1}{8}}, & \dfrac{1}{\sqrt{2}}\,\gamma^{\frac{p^2-1}{8}} \end{pmatrix}.
}
This group is real when $p\equiv 1\ (\mathrm{mod}\ 8)$, and imaginary, but contained in $\Gamma_p$, when $p\equiv -1\ (\mathrm{mod}\ 8)$.

The first example is $p=7$. Here $-1$ is a quadratic non-residue, and one can therefore set $\varepsilon=\sqrt{-1}=i$. If we take $\gamma=2+i$, $\gamma^7=2-i$, $\gamma^6=2\,(1+i)$, $\gamma^8=5$, then $\gamma$ is a primitive root, since $\gamma,\gamma^2,\dots,\gamma^7$ are different from one another and $5$ is a real primitive root of $7$, so that every non-vanishing number of the field $\mathfrak{E}_{2,7}$ can be represented in the form
\[
5^{\mu}\gamma^{\nu} = \gamma^{8\mu+\nu}
\]
if
\[
\begin{gathered}
\mu = 0,\ 1,\ 2,\ 3,\ 4,\ 5,\\
\nu = 0,\ 1,\ 2,\ 3,\ 4,\ 5,\ 6,\ 7
\end{gathered}
\]
is set.

Further, we then have $\sqrt{2}=3$, and consequently the substitutions $\Theta,\psi,\chi$ are:
\weq{9}{
\begin{pmatrix} 2\,(1+i), & 0 \\ 0, & 2\,(1-i) \end{pmatrix}, \
\begin{pmatrix} 0, & i \\ i, & 0 \end{pmatrix}, \
\begin{pmatrix} 4\,(1-i), & 4\,(1-i) \\ -4\,(1+i), & 4\,(1+i) \end{pmatrix}.
}
If one wishes to pass to the real form, one must proceed according to \S.\,69. One passes from the substitutions $B$ to the $A$ by setting $\nu=\gamma^3=2-3i$, and obtains from (9) for $\Theta,\psi,\chi$ in the group $\Lambda_p$:
\weq{10}{
\begin{pmatrix} 2+2i, & 0 \\ 0, & 2-2i \end{pmatrix}, \
\begin{pmatrix} 0, & 3-2i \\ 3+2i, & 0 \end{pmatrix}, \
\begin{pmatrix} 4-4i, & 1-4i \\ 1+4i, & 4+4i \end{pmatrix}.
}
Here, further,
\[
R = \begin{pmatrix} 4, & 4i \\ i, & 1 \end{pmatrix}
\]
%CONT>>


% ===== scan leaf 300, printed 279 =====
% !PRINTED 279
\weberpage{279}
%<<CONT
to be set, from which, by forming $R\,A\,R^{-1}$, the generating substitutions of the octahedral group result in real form:
\[
\Theta = \begin{pmatrix} 2, & 1 \\ 3, & 2 \end{pmatrix}, \
\psi = \begin{pmatrix} 2, & 2 \\ 1, & -2 \end{pmatrix}, \
\omega = \psi\,\Theta = \begin{pmatrix} 3, & -1 \\ 3, & -3 \end{pmatrix}
\]
\weq{11}{
\chi = \begin{pmatrix} 0, & 2 \\ 3, & 1 \end{pmatrix}, \
\chi^2 = \begin{pmatrix} -1, & 2 \\ 3, & 0 \end{pmatrix}.
}
It is now easy, following \S.\,57, to form the complete group, which it is not necessary to carry out further.

\gesperrt{The octahedral group for $p=7$ is a divisor of $L_7$ of index 7.}

V. \gesperrt{Icosahedral group.}

In the case $p=5$, $L_p$ is itself an icosahedral group. For other values of $p$, an icosahedral group can be contained in $L_p$ only if $p^2-1$ is divisible by 5, that is, $p\equiv\pm 1\ (\mathrm{mod}\ 5)$. Then from \S.\,58 we obtain the generating substitutions of such a group.

If, for brevity, we set
\[
\gamma^{\frac{p^2-1}{10}} = \varrho,
\]
then we have to substitute $\varrho^2$ for $\varepsilon$ in the formulae of \S.\,58, and obtain
\weq{12}{
\Theta = \begin{pmatrix} \varrho, & 0 \\ 0, & \varrho^{-1} \end{pmatrix}, \
\psi = \begin{pmatrix} 0, & 1 \\ -1, & 0 \end{pmatrix}, \quad
\chi = \begin{pmatrix} \dfrac{-1}{\varrho^2-\varrho^{-2}}, & \dfrac{1}{\varrho-\varrho^{-1}} \\[2mm] \dfrac{1}{\varrho-\varrho^{-1}}, & \dfrac{1}{\varrho^2-\varrho^{-2}} \end{pmatrix}\ [\S.\,58,\,(26)].
}
If $p\equiv 1\ (\mathrm{mod}\ 5)$, then $\varrho$, and with it the whole group, is real. But if $p\equiv -1\ (\mathrm{mod}\ 5)$, then $\varrho$ and $\varrho^{-1}$ are conjugate, and consequently the group found is contained in $\Gamma_p$, and in order to obtain the icosahedral group in $L_p$ we must first transform according to \S.\,69.

\gesperrt{For $p=11$ we obtain immediately a real icosahedral group of index 11.}


% ===== scan leaf 301, printed 280 =====
% !PRINTED 280
\weberpage{280}
For $\varrho=\gamma^{12}$ we may choose an arbitrary primitive root of 11, e.g.\ 8. Then $\varrho^{-1}=7$ and
\[
\varrho-\varrho^{-1}=1,\quad \varrho^2-\varrho^{-2}=4,
\]
hence
\weq{13}{
\Theta = \begin{pmatrix} 3, & 0 \\ 0, & 4 \end{pmatrix}, \
\psi = \begin{pmatrix} 0, & 1 \\ -1, & 0 \end{pmatrix}, \
\chi = \begin{pmatrix} -3, & 1 \\ 1, & 3 \end{pmatrix}.
}
With this it is thus established which kinds of divisors can occur in the group $L_p$, and it has also been shown that all these divisors are actually present. The question of how one obtains, for each type, all the divisors that are possible at all, we have left aside here. We wish to make only the following remarks about it.

If one transforms a group $G$ that has been found by any substitution of the group $L_p$, one obtains a conjugate divisor. But if one transforms by an imaginary substitution of the group $E_p$, it can happen that the transformed group of $G$ nevertheless becomes real and is not conjugate to $G$ within $L_p$ (although both are conjugate divisors of $E_p$). Thus from the octahedral and the icosahedral group one obtains in general a second, non-conjugate one, if one transforms
with $\begin{pmatrix} \varepsilon, & 0 \\ 0, & \varepsilon^{-1} \end{pmatrix}$. For the tetrahedral group this yields a new group only when $p\equiv\pm 1\ (\mathrm{mod}\ 8)$. The second group results from the formula:
\weq{14}{
\begin{pmatrix} \varepsilon, & 0 \\ 0, & \varepsilon^{-1} \end{pmatrix}
\begin{pmatrix} a, & b \\ c, & d \end{pmatrix}
\begin{pmatrix} \varepsilon^{-1}, & 0 \\ 0, & \varepsilon \end{pmatrix}
=
\begin{pmatrix} a, & bN \\ cN^{-1}, & d \end{pmatrix}
\quad (N=\varepsilon^2).
}
In order to exhibit these circumstances in the examples $p=5,7,11$, it is expedient to transform the group (6) found above for $p=7$ in such a way that the substitution $\chi$, which is of the third order, receives the normal form $S$. One finds this transformation easily, and obtains for this case:
\[
\begin{pmatrix} 1, & 1 \\ -2, & -1 \end{pmatrix}
\begin{pmatrix} 0, & 2 \\ 3, & 1 \end{pmatrix}
\begin{pmatrix} -1, & -1 \\ 2, & 1 \end{pmatrix}
=
\begin{pmatrix} 3, & 0 \\ 0, & 5 \end{pmatrix}
= \chi'
\]
\weq{15}{
\begin{pmatrix} 1, & 1 \\ -2, & -1 \end{pmatrix}
\begin{pmatrix} 2, & 1 \\ 3, & 2 \end{pmatrix}
\begin{pmatrix} -1, & -1 \\ 2, & 1 \end{pmatrix}
=
\begin{pmatrix} 1, & -2 \\ -1, & 3 \end{pmatrix}
= \Theta'
}
\[
\begin{pmatrix} 1, & 1 \\ -2, & -1 \end{pmatrix}
\begin{pmatrix} 2, & 2 \\ 1, & -2 \end{pmatrix}
\begin{pmatrix} -1, & -1 \\ 2, & 1 \end{pmatrix}
=
\begin{pmatrix} -3, & -3 \\ 1, & 3 \end{pmatrix}
= \psi',
\]
%CONT>>


% ===== scan leaf 302, printed 281 =====
% !PRINTED 281
\weberpage{281}

%<<CONT
from which one obtains, by composition,

\weq{16}{\psi' \Theta' \chi' = \begin{pmatrix} 0, & 1 \\ -1, & 0 \end{pmatrix}}

We may accordingly, in the three cases, take the following generating substitutions of the tetrahedral, octahedral and icosahedral group [(6), (13), (15)]:

\[
p = 5. \quad \begin{pmatrix} 2, & 0 \\ 0, & 3 \end{pmatrix}, \ \begin{pmatrix} 0, & 1 \\ -1, & 0 \end{pmatrix}, \ \begin{pmatrix} 2, & 2 \\ 1, & -1 \end{pmatrix}
\]
\[
p = 7. \quad \begin{pmatrix} 3, & 0 \\ 0, & 5 \end{pmatrix}, \ \begin{pmatrix} 0, & 1 \\ -1, & 0 \end{pmatrix}, \ \begin{pmatrix} 1, & -2 \\ -1, & 3 \end{pmatrix}
\]
\[
p = 11. \quad \begin{pmatrix} 3, & 0 \\ 0, & 4 \end{pmatrix}, \ \begin{pmatrix} 0, & 1 \\ -1, & 0 \end{pmatrix}, \ \begin{pmatrix} -3, & 1 \\ 1, & 3 \end{pmatrix}.
\]

By repeated composition, which can be arranged in very many different ways, the complete groups are easily obtained from these:

\[
\begin{gathered}
p = 5.\\[2pt]
\begin{pmatrix} x, & 0 \\ 0, & x^{-1} \end{pmatrix}, \ \begin{pmatrix} 0, & x \\ -x^{-1}, & 0 \end{pmatrix}, \ \begin{pmatrix} x, & y \\ 2y^{-1}, & 3x^{-1} \end{pmatrix} \quad \begin{aligned} x &= 1, 2\\ y &= \pm 1, \pm 2 \end{aligned}
\end{gathered}
\]
\[
\begin{gathered}
p = 7.\\[2pt]
\begin{pmatrix} x, & 0 \\ 0, & x^{-1} \end{pmatrix}, \ \begin{pmatrix} 0, & x \\ -x^{-1}, & 0 \end{pmatrix}, \ \begin{pmatrix} x, & -xz \\ -2x^{-1}z^{-1}, & 3x^{-1} \end{pmatrix}, \ \begin{pmatrix} xz, & x \\ -3x^{-1}, & -2x^{-1}z^{-1} \end{pmatrix}\\[2pt]
x = 1, 2, 3; \quad z = 1, 2, 4\\
(z \text{ quadratic residue of } 7).
\end{gathered}
\]
\[
\begin{gathered}
p = 11.\\[2pt]
\begin{pmatrix} x, & 0 \\ 0, & x^{-1} \end{pmatrix}, \ \begin{pmatrix} 0, & x \\ -x^{-1}, & 0 \end{pmatrix}, \ \begin{pmatrix} x, & -xz \\ -x^{-1}z^{-1}, & 2x^{-1} \end{pmatrix}, \ \begin{pmatrix} xz, & x \\ -2x^{-1}, & -x^{-1}z^{-1} \end{pmatrix}\\[2pt]
x = 1, 2, 3, 4, 5; \quad z = 1, 3, 4, 5, 9\\
(z \text{ quadratic residue of } 11).
\end{gathered}
\]

The application of the transformation (14) leads, for $p = 5$, to no new group; for $p = 7, 11$ one obtains in each case one other group, which arises from these when one puts for $z$ the sequence of the quadratic non-residues instead of the residues\textsuperscript{1}).

\bigskip
\hrule
\bigskip

\textsuperscript{1}) The property of the congruence groups of being simple was already known to \gesperrt{Galois}. Likewise the divisors of index $p$ for $p = 5, 7, 11$ were known to him. These groups have been investigated more thoroughly by
%CONT>>


% ===== scan leaf 303, printed 282 =====
% !PRINTED 282
\weberpage{282}
\pdfbookmark[2]{§ 72. Constitution of the group L7 of degree 168}{bm83}

\begin{center}
\S.\,72.

\gesperrt{Constitution of the Group $L_7$ of Degree 168.}
\end{center}

Among the simple groups found here, the next after the icosahedral group, which recurs here as $L_5$, is the group $L_7$ of degree 168, which, as we have seen, contains an octahedral group as a divisor. Since this remarkable group will meet us several times more later on, we wish here to enter somewhat more closely into its structure.

In \S.\,57 the octahedral group is represented by three elements $\chi, \omega, \Theta$ in the form:

\weq{1}{\chi^\lambda \omega^u \Theta^v, \quad \lambda = 0, 1, 2; \ \mu = 0, 1; \ \nu = 0, 1, 2, 3.} % [sic] formula prints Latin u, v; the parameter list prints Greek mu, nu

Between these elements there subsist the relations

\weq{2}{\omega\chi = \chi^2\omega, \quad \Theta\omega = \omega\Theta^3, \quad \Theta\chi = \chi^2\omega\Theta^2, \quad \Theta^2\chi = \chi\omega\Theta^3,}

and these conditions, when the degrees 3, 2, 4 of the elements $\chi, \omega, \Theta$ are added to them, we have shown to be sufficient to characterise the system (1) as an octahedral group.

From (2) we have derived in \S.\,57, as consequences, the formulae:

\weq{3}{\begin{gathered}
\omega\chi = \chi^2\omega, \quad \omega\chi^2 = \chi\omega, \quad \Theta^2\chi^2 = \chi^2\omega\Theta,\\
\Theta\chi = \chi^2\omega\Theta^2, \quad \Theta^2\chi = \chi\omega\Theta^3, \quad \Theta^3\chi = \chi^2\Theta,\\
\Theta\omega = \omega\Theta^3, \quad \Theta^2\omega = \omega\Theta^2, \quad \Theta^3\omega = \omega\Theta, \quad \Theta\chi^2 = \chi\Theta^3,
\end{gathered}}

which we shall make use of several times later.

In \S.\,71 we have also represented the octahedral group as a congruence group with respect to the modulus 7, and in the formulae (11) there we have found

\weq{4}{\chi = \begin{pmatrix} 0, & 2 \\ 3, & 1 \end{pmatrix}, \ \omega = \begin{pmatrix} 3, & -1 \\ 3, & -3 \end{pmatrix}, \ \Theta = \begin{pmatrix} 2, & 1 \\ 3, & 2 \end{pmatrix}}

For what follows, however, it is expedient, in the interest of a simpler calculation, to further trans-
%CONT>>

\bigskip
\hrule
\bigskip

%<<CONT
\gesperrt{Serret} (Cours d'alg\`ebre supérieure) and by C. \gesperrt{Jordan} (Traité des Substitutions). Cf. \gesperrt{Weber}, Elliptische Functionen etc., \S.\,84 f. The complete enumeration of all divisors is due to \gesperrt{Gierster} (Math. Ann., Bd. XVIII). A detailed treatment of the congruence groups in ``\gesperrt{Klein-Fricke}, Vorlesungen über die Theorie der elliptischen Modulfunctionen'', Bd. I, Leipzig 1890.


% ===== scan leaf 304, printed 283 =====
% !PRINTED 283
\weberpage{283}

%<<CONT
form. For it has already been shown earlier that the group $L_7$ contains divisors of degree 21 whose index is $= 8$, and which are of particular importance precisely for our purpose. Among the conjugate divisors of index 8, however, the simplest and the most convenient for calculation is the one consisting of all substitutions

\weq{5}{\begin{pmatrix} a, & b \\ 0, & a^{-1} \end{pmatrix} \quad (\mathrm{mod}\ 7)}

and we wish to arrange the transformation so that $\chi$ appears in this form. This we achieve by transformation of the substitutions (4) by means of $\begin{pmatrix} 1, & 1 \\ -1, & 0 \end{pmatrix}$, and thereby there results from (4)

\weq{6}{\chi = \begin{pmatrix} 2, & 3 \\ 0, & -3 \end{pmatrix}, \ \omega = \begin{pmatrix} 1, & -3 \\ 3, & -1 \end{pmatrix}, \ \Theta = \begin{pmatrix} 1, & 3 \\ -2, & 2 \end{pmatrix}.}

We now further assemble, for a better overview, the whole octahedral group into a table, whose six rows contain the elements $\Theta^\lambda, \omega\Theta^\lambda, \chi\Theta^\lambda, \chi\omega\Theta^\lambda, \chi^2\Theta^\lambda, \chi^2\omega\Theta^\lambda$ for $\lambda = 0, 1, 2, 3$:

\weq{7}{\begin{gathered}
\begin{pmatrix} 1, & 0 \\ 0, & 1 \end{pmatrix}, \ \begin{pmatrix} 1, & 3 \\ -2, & 2 \end{pmatrix}, \ \begin{pmatrix} 2, & 2 \\ 1, & -2 \end{pmatrix}, \ \begin{pmatrix} 2, & -3 \\ 2, & 1 \end{pmatrix};\\
\begin{pmatrix} 1, & -3 \\ 3, & -1 \end{pmatrix}, \ \begin{pmatrix} 0, & 3 \\ 2, & 0 \end{pmatrix}, \ \begin{pmatrix} -1, & 1 \\ -2, & 1 \end{pmatrix}, \ \begin{pmatrix} 3, & 1 \\ -3, & -3 \end{pmatrix};\\
\begin{pmatrix} 2, & 3 \\ 0, & -3 \end{pmatrix}, \ \begin{pmatrix} 3, & -2 \\ -1, & 1 \end{pmatrix}, \ \begin{pmatrix} 0, & 2 \\ 3, & 1 \end{pmatrix}, \ \begin{pmatrix} 3, & -3 \\ 1, & -3 \end{pmatrix};\\
\begin{pmatrix} 3, & 2 \\ 2, & -3 \end{pmatrix}, \ \begin{pmatrix} 1, & 1 \\ -1, & 0 \end{pmatrix}, \ \begin{pmatrix} 1, & 2 \\ 1, & 3 \end{pmatrix}, \ \begin{pmatrix} -3, & 0 \\ 2, & 2 \end{pmatrix};\\
\begin{pmatrix} 3, & 3 \\ 0, & -2 \end{pmatrix}, \ \begin{pmatrix} -3, & 1 \\ -3, & 3 \end{pmatrix}, \ \begin{pmatrix} 2, & 0 \\ -2, & -3 \end{pmatrix}, \ \begin{pmatrix} -2, & 1 \\ 3, & -2 \end{pmatrix};\\
\begin{pmatrix} -2, & 2 \\ 1, & 2 \end{pmatrix}, \ \begin{pmatrix} -1, & 2 \\ 3, & 0 \end{pmatrix}, \ \begin{pmatrix} 2, & 1 \\ 3, & 2 \end{pmatrix}, \ \begin{pmatrix} 0, & 1 \\ -1, & -1 \end{pmatrix}.
\end{gathered}}
To obtain the whole group $L_7$ we must still add an element of degree $7$, and for this we choose

\weq{8}{\tau = \begin{pmatrix} 1, & 1 \\ 0, & 1 \end{pmatrix}.}


% ===== scan leaf 305, printed 284 =====
% !PRINTED 284
\weberpage{284}

Then the entire group $L_7$ is represented thus:

\weq{9}{\tau^\varrho \chi^\lambda \omega^\mu \Theta^\nu, \quad \begin{aligned} \varrho &= 0, 1, 2, 3, 4, 5. 6\\ \lambda &= 0, 1, 2; \ \mu = 0, 1\\ \nu &= 0, 1, 2, 3. \end{aligned}} % [sic] full stop printed between 5 and 6

That all these elements (9) are different from one another follows simply from the fact that no power of $\tau$ whose exponent is not divisible by 7 can be contained in the octahedral group, because the octahedral group has no element of degree $7$.

The group (5) of degree $21$ contained in the group $L_7$ is then represented in the form $\tau^\varrho \chi^\lambda$.

In order to characterise the composition of the elements (9), it evidently suffices if for every $\varrho$ the elements

\weq{10}{\Theta\tau^\varrho, \ \omega\tau^\varrho, \ \chi\tau^\varrho}

are represented in the form (9), because thereby, together with the compositions in the octahedral group, the symbolic product of any two elements of the form (9) can again be represented in the form (9).

Now since

\weq{11}{\tau^\varrho = \begin{pmatrix} 1, & \varrho \\ 0, & 1 \end{pmatrix}.}

one first obtains very simply from (6)

\weq{12}{\chi\tau^\varrho = \tau^{4\varrho}\chi,}

and from this

\weq{13}{\chi^2\tau^\varrho = \tau^{2\varrho}\chi^2,}

and one sees easily that these six relations are a consequence of the single one:

\weq{14}{\chi\tau = \tau^4\chi.}

The remaining compositions (10), however, are obtained only by actual computation in the individual cases, in which the table (7) renders good service. One applies it in such a way that for each value $\varrho$ one forms the elements

\[
\tau^{\varrho'} \omega \tau^\varrho, \quad \tau^{\varrho'} \Theta \tau^\varrho
\]

and determines $\varrho'$ so that these elements are found in the table, which is always possible in only one way.


% ===== scan leaf 306, printed 285 =====
% !PRINTED 285
\weberpage{285}
One thus finds, by an easy, if somewhat laborious, calculation
\weq{15}{\begin{array}{rl@{\qquad}rl}
\omega\tau &{}= \tau^2\chi^2\omega\Theta^2, & \Theta\tau &{}= \tau^3\omega\Theta\\
\omega\tau^2 &{}= \tau\chi^2\Theta^3, & \Theta\tau^2 &{}= \tau\chi\omega\Theta^3\\
\omega\tau^3 &{}= \tau^5\chi\Theta^2, & \Theta\tau^3 &{}= \tau^5\chi\omega\\
\omega\tau^4 &{}= \tau^4\chi^2\Theta, & \Theta\tau^4 &{}= \tau^6\omega\Theta^2\\
\omega\tau^5 &{}= \tau^3\chi^2\omega\Theta, & \Theta\tau^5 &{}= \tau^2\Theta^3\\
\omega\tau^6 &{}= \tau^6\omega\Theta^3, & \Theta\tau^6 &{}= \tau^4\chi^2\Theta^2.
\end{array}}
But all these twelve relations can be derived as consequences
from four of them, which may be selected in manifold ways.
One can, for example, choose for these four fundamental
relations the following:
\weq{16}{\begin{array}{c}
\omega\tau = \tau^2\chi^2\omega\Theta^2, \quad \omega\tau^3 = \tau^5\chi\Theta^2, \quad \omega\tau^4 = \tau^4\chi^2\Theta,\\
\Theta\tau = \tau^3\omega\Theta,
\end{array}}
from which, with the use of the formulae (3), (12), (13), all the others
follow easily, e.g.:
\[
\omega\tau^5 = \tau^4\chi^2\Theta\tau = \tau^4\chi^2\tau^3\omega\Theta = \tau^3\chi^2\omega\Theta,
\]
and so the rest.

Just as we have seen earlier that we can regard $\chi$ and $\Theta$
as generating elements of the octahedral group, so we can
now regard the two elements $\omega, \tau$ as generators of the group
$L_7$, for it follows from (15):
\weq{17}{\Theta^3 = \omega\tau\omega\tau^6, \quad \chi = \tau^2\Theta\tau^3\omega.}
Summing up, we obtain the following result:

\gesperrt{I. If four elements $\tau, \chi, \omega, \Theta$ of the degrees 7, 3, 2, 4
are given, in whose composition the relations
 hold:}
\[
\begin{array}{ccc}
\omega\chi = \chi^2\omega, & \Theta\omega = \omega\Theta^3, & \Theta\chi = \chi^2\omega\Theta^2,\\
\Theta^2\chi = \chi\omega\Theta^3, & \omega\tau = \tau^2\chi^2\omega\Theta^2, & \omega\tau^3 = \tau^5\chi\Theta^2,\\
\omega\tau^4 = \tau^4\chi^2\Theta, & \Theta\tau = \tau^3\omega\Theta, & \chi\tau = \tau^4\chi,
\end{array}
\]
\gesperrt{then the 168 elements}
\[
\sigma = \tau^\varrho \chi^\lambda \omega^u \Theta^v,
\]
\gesperrt{when $\varrho, \lambda, \mu, \nu$ run through complete systems of residues
 with respect to the moduli 7, 3, 2, 4, form a simple group
of degree 168, and $\omega$ and $\tau$ can be regarded as generating
 elements of this group.} % [sic] exponents in the enunciation text (mu, nu) do not match the exponents u, v used in the preceding formula for sigma


% ===== scan leaf 307, printed 286 =====
% !PRINTED 286
\weberpage{286}
Among the divisors of this group are to be singled out the
octahedral group
\[
\chi^\lambda \omega^u \Theta^v
\]
of index 7, then the elements
\[
\tau^\varrho \chi^\lambda,
\]
which by the relations (12) and (13) form a group of degree 21,
and thus a divisor of the total group of index 8;
further the group of degree 8, $\omega^u \Theta^v$. The entire group
can also be represented in another order as follows:
\[
\chi^\lambda \omega^u \Theta^v \tau^\varrho, \quad \omega^u \Theta^v \chi^\lambda \tau^\varrho, \quad \tau^\varrho \omega^u \Theta^v \chi^\lambda.
\]

\bigskip
\begin{center}
\rule{4cm}{0.4pt}
\end{center}


% ===== scan leaf 308, printed 287 =====
\weberpage{287}% page number not printed on this page
% !PRINTED none
% Part-title page (no running head, no page number).
\begin{center}
\gesperrt{THIRD BOOK.}

\rule{2.5cm}{0.4pt}

\vspace{1.5cm}

\gesperrt{APPLICATIONS}

\vspace{0.5cm}

\gesperrt{OF}

\vspace{0.5cm}

\gesperrt{GROUP THEORY.}

\vspace{1.5cm}

\rule{5cm}{0.4pt}
\end{center}


% ===== scan leaf 309, printed 288 =====
\weberpage{288}% page number not printed on this page
% !PRINTED none
% Blank page (verso of the part-title leaf). No text or page number is printed.


% ===== scan leaf 310, printed 289 =====
\weberpage{289}% page number not printed on this page
\pdfbookmark[0]{Third Book. Applications of Group Theory.}{bm84}
\pdfbookmark[1]{Tenth Section. General Theory of Metacyclic Equations.}{bm85}
\pdfbookmark[2]{§ 73. The resolvents of the composition series}{bm86}
% !PRINTED none
\weberpage{none}

\begin{center}
\gesperrt{Tenth Chapter.}

\bigskip

\gesperrt{General Theory of Metacyclic Equations.}
\end{center}

\bigskip
\begin{center}
\rule{4cm}{0.4pt}
\end{center}
\bigskip

\begin{center}
\S.\,73.

\medskip

\gesperrt{The Resolvents of the Composition Series.}
\end{center}

\bigskip

From the theorems of general group theory that are treated in the first Book of this
volume, important algebraic consequences result when they are applied to the
\gesperrt{Galois} group of an equation.

An arbitrary field $\Omega$ is now assumed as domain of rationality; let
$f(x) = 0$ be an equation in this field without multiple roots, and $P$ its
\gesperrt{Galois} group. Let this group $P$ have a normal divisor $Q$. We denote by
$n$ the degree of $P$ and by $j$ the index of the divisor $Q$, so that $n : j$ is the
degree of $Q$.

We have seen in the first volume (\S.\,156) that the group of $f(x)$ is reduced to
$Q$ by the adjunction of a root of an irreducible normal equation of degree $j$,
and this auxiliary equation of degree $j$ we have designated
the \gesperrt{partial resolvent}.

The \gesperrt{Galois} group of this partial resolvent we obtained as follows: if
$\psi$ denotes a function of the roots of $f(x)$ belonging to the group $Q$, and if $P$
is decomposed into the subsidiary groups
\[
P = Q + Qa + Qb + Qc + \cdots
\]
where, then, $a, b, c, \ldots$ are certain permutations from $P$, then $\psi$ passes,
under the whole subsidiary group $Qa$, into one and the same function $\psi_a$, and the
\gesperrt{Galois} group of the resolvent consists of the substitutions
\[
(\psi, \psi), (\psi, \psi_a), (\psi, \psi_b), (\psi, \psi_c), \ldots
\]


% ===== scan leaf 311, printed 290 =====
% !PRINTED 290
\weberpage{290}

If one composes two of these substitutions, one must observe that
\[
(\psi, \psi_b) = (\psi_a, \psi_{ab})
\]
holds, so that one has
\[
(\psi, \psi_a)(\psi, \psi_b) = (\psi, \psi_{ab})
\]
These substitutions therefore compose in entirely the same way as the subsidiary groups do according to \S.\,4 of this volume, and from this it follows:

\noindent 1. \gesperrt{The Galois group of the partial resolvent belonging to $Q$
is isomorphic with the group $P/Q$ complementary to $Q$.}

Let us now take any composition series of $P$ with the associated
index series (\S.\,6):
\weq{1}{
\begin{gathered}
P,\ P_1,\ P_2,\ \ldots,\ P_{\mu-1},\ 1 \\
j_1,\ j_2,\ \ldots,\ j_{\mu-1},\ j_\mu,
\end{gathered}
}
then, by what has just been said, the group of the equation $f = 0$ is reduced from $P$ to $P_1$
by the adjunction of \gesperrt{one} root of a normal equation of degree $j_1$. Then
it is reduced to $P_2$ by a root of a normal equation of degree $j_2$,
and so on, and finally the equation is completely solved by a root
of a normal equation of degree $j_\mu$.

On this procedure of solution a new light now falls from the theorem on the
invariability of the index series (\S.\,6, I.).

\noindent 2. \gesperrt{To solve the equation $f = 0$, one has to adjoin, one after
another, one root each of a normal equation of the degrees $j_1, j_2, \ldots, j_\mu$.
The degrees of these resolvents can indeed be altered in their order,
but not in their totality.}

The numbers $j_1, j_2, \ldots, j_\mu$ are thus far more deeply connected with the nature of an
equation, or more generally with the fields defined by the equation,
than, say, the degree of the equation; for while the degree can be
changed in manifold ways by transformation, when one introduces functions of the roots as new
unknowns, the numbers $j_1, j_2, \ldots, j_\mu$ always remain preserved and
are to be regarded as true \gesperrt{invariants} of the normal field defined by
the equation.


% ===== scan leaf 312, printed 291 =====
% !PRINTED 291
\weberpage{291}

To these invariants belongs also the degree $n$ of the group $P$ itself, which is
determined by the $j$ as follows:
\weq{2}{n = j_1 j_2 \cdots j_{\mu-1} j_\mu.}
For, by the meaning of the indices, $n : j_1$ is the degree of $P_1$, $n : j_1 j_2$
the degree of $P_2$, and so on, and since the last of the degrees is equal to $1$,
formula (2) results.

In general, in a composition series of $P$ each term is a normal divisor only of
the one immediately preceding it. But there may also occur individual terms which
are normal divisors of terms preceding them further back as well. Especially
important are those terms which are normal divisors of $P$ itself, and hence also
of all the terms of the composition series that precede them. We shall see what
algebraic consequences are to be drawn from this circumstance.

Let $P_\nu$ be a term of the series (1) which is at the same time a normal divisor
of $P$. The index of $P_\nu$ with respect to $P$ is $j_1 j_2 \ldots j_\nu$, and
this product is the degree of the partial resolvent by which the group $P$ is
reduced to $P_\nu$, which we shall denote by $\chi(y) = 0$. The group of this
resolvent, which is a normal equation, we obtain by Theorem 1. in the form
$P/P_\nu$.

We can now easily find a composition series for this group, together with the
corresponding index series, namely:
\weq{3}{
\begin{gathered}
P/P_\nu,\ P_1/P_\nu,\ P_2/P_\nu,\ \ldots,\ P_{\nu-1}/P_\nu,\ 1 \\
j_1,\quad j_2,\quad \ldots,\quad j_{\nu-1},\quad j_\nu.
\end{gathered}
}
For, by Theorem 1., \S.\,6, $P_1/P_\nu$ is a normal divisor of $P/P_\nu$ of
index $j_1$, and it is a greatest normal divisor,
because by 2., \S.\,6 no normal divisor of $P/P_\nu$ can stand above $P_1/P_\nu$
if no normal divisor of $P$ stands above $P_1$. And one can conclude in the same
way with respect to the following terms of (3).

From this we draw a further important consequence. We have already shown in
\S.\,7 that, if $Q$ is any normal divisor of $P$, a composition series of $P$ can
be found in which $Q$ occurs. If now $R$ is another normal divisor of $P$ and at
the same time a divisor of $Q$, then the composition series of $P$ can be so
arranged that $Q$ and $R$ occur in it.

We suppose such a composition series to be determined:
\weq{4}{P,\ Q',\ Q'',\ \ldots,\ Q,\ \ldots,\ R.}


% ===== scan leaf 313, printed 292 =====
% !PRINTED 292
\weberpage{292}
\pdfbookmark[2]{§ 74. Metacyclic equations}{bm87}

According to (3) we can derive from this the composition series of the two groups $P/Q$
and $P/R$, which we shall indicate thus:
\weq{5}{P/Q,\ Q'/Q,\ Q''/Q,\ \ldots,\ 1}
\weq{6}{P/R,\ Q'/R,\ Q''/R,\ \ldots,\ Q/R \ldots}
Now the index of the divisor $Q'/Q$ of $P/Q$ is equal to the index of the divisor $Q'$
of $P$, and hence also equal to the index of the divisor $Q'/R$ of $P/R$ (\S.\,6, 1.), and
the same holds of the following terms. We therefore state the theorem:

\noindent 3. \gesperrt{If $Q$ and $R$ are normal divisors of $P$, and $R$ is a
divisor of $Q$, then the index series of $P/Q$ is a part of the index series of
$P/R$..} % [sic] double period as printed

If the indices $j_1, j_2, \ldots, j_\nu$ are all primes, then the solution
of the resolvent $\chi(y) = 0$ is reducible to the solution of a chain of cyclic
equations of degrees $j_1, j_2, \ldots, j_\nu$; this resolvent is
therefore metacyclic in the sense in which we settled this concept in the seventeenth chapter
of the first volume, according to which the metacyclic equations are
identical with those otherwise called algebraically solvable.

An irreducible equation $f(x) = 0$ is metacyclic if the index series of its
group consists entirely of primes. At the place mentioned we investigated the
conditions for metacyclic equations of prime degree, and
must now carry these considerations through for the general case, in which
the degree of the equation is arbitrarily composite.

\begin{center}
\S.\,74.

\medskip

\gesperrt{Metacyclic Equations.}
\end{center}

\bigskip

If $f(x) = 0$ is an irreducible equation of degree $m$ and $P$ is its
\gesperrt{Galois} group, then $P$ is transitive. Let a composition and index series
for $P$ be
\weq{1}{
\begin{gathered}
P,\ P_1,\ P_2,\ \ldots,\ P_{\mu-1},\ 1 \\
j_1,\ j_2,\ \ldots,\ j_{\mu-1},\ j_\mu.
\end{gathered}
}
Since the last group $1$ is intransitive, an intransitive group must at some point
occur in the composition series; let $P_\lambda$ therefore be the first
intransitive group of the sequence (1). Then
%CONT>>


% ===== scan leaf 314, printed 293 =====
% !PRINTED 293
\weberpage{293}
%<<CONT
all the following groups $P_{\lambda+1}, P_{\lambda+2}, \ldots$ as well, which are indeed all
divisors of $P_\lambda$, are intransitive.

We must now recall the theorems 1., 2., 3. in \S.\,158 of the first volume. Since
$P_\lambda$ is a normal divisor of $P_{\lambda-1}$, it follows from those theorems
that $P_{\lambda-1}$ is imprimitive, and if by the necessary adjunctions the
group of $f(x) = 0$ has been reduced to $P_{\lambda-1}$, then by the further
adjunction of a root of a normal equation of degree $j_\lambda$ the function
$f(x)$ will split into several irreducible factors
\weq{2}{f(x) = f_1(x)\, f_2(x) \ldots f_h(x)}
which are all of equal degree and whose number $h$ is a divisor of
$j_\lambda$.

If we denote by $m$ the degree of $f(x)$, and by $m_1$ the degree of $f_1(x)$, then
\weq{3}{m = h\, m_1.}
By the cited theorem 1. in \S.\,158, Bd. I, the equations
$f_1 = 0, f_2 = 0, \ldots, f_h = 0$ all have the same group, which we shall denote
by $Q$. If $\alpha, \alpha_1, \alpha_2, \ldots, \alpha_{m_1-1}$ are the roots of
$f_1 = 0$, then we obtain the group $Q$ by collecting the permutations of the $\alpha$
that are brought about by $P_\lambda$. The first question concerns
the relation of the degrees $p_\lambda$ and $q$ of the two groups $P_\lambda$ and $Q$.

There may be several distinct permutations in $P_\lambda$ that generate the same
element in $Q$, that is, which contain the same permutation of the $\alpha$.
If $\pi_1, \pi_2$ are two distinct permutations from $P_\lambda$ that bring about
the same permutation among the $\alpha$, then $\pi_2 \pi_1^{-1} = \pi_0$ will be a
permutation that leaves the $\alpha$ at rest, and thus
$\pi_2 = \pi_0 \pi_1$. If conversely $\pi_0$ is any permutation from
$P_\lambda$ that does not permute the roots $\alpha$, then the permutation
$\pi_2 = \pi_0 \pi_1$ will produce the same change among the $\alpha$ as $\pi_1$.

From this it follows that every permutation of the $\alpha$ is generated equally often by the
permutations of $P_\lambda$, namely just as often as the $\alpha$
remain unchanged by permutations from $P_\lambda$. If we denote this number
by $p_0$, then
\weq{4}{p_\lambda = p_0\, q,}
\gesperrt{or the degree $p_\lambda$ of the group $P_\lambda$ is a multiple of the degree $q$
of $Q$.}


% ===== scan leaf 315, printed 294 =====
% !PRINTED 294
\weberpage{294}
We have now proved, in §.\,154, 7. of the first volume, the theorem that the degree
of a transitive permutation group is always divisible by the number of permuted
figures. Here, however, $f_1(x)$ is irreducible, and hence the group
$Q$ is transitive. Thus $q$ is divisible by the degree of $f_1(x)$, i.e. by $m_1$,
and consequently, by (4), $p_\lambda$ is also divisible by $m_1$.

The degree $p_{\lambda-1}$ of $P_{\lambda-1}$ we obtain, from the meaning of
$j_\lambda$, in the form:
\weq{5}{p_{\lambda-1} = j_\lambda\, p_\lambda.}
Now $p_\lambda$ is a multiple of $m_1$, $j_\lambda$ a multiple of $h$,
and $h\, m_1 = m$ is equal to the degree of $f(x)$. Accordingly
\weq{6}{p_{\lambda-1} = k\, m}
is a multiple of $m$.

But from the meaning of the $j$ it follows [§.\,73, (2)] that
$p_{\lambda-1} = j_\lambda j_{\lambda-1} \ldots j_\mu$, % [sic] second factor $j_{\lambda-1}$ as printed; equation (7) gives $j_{\lambda+1}$.
and consequently we have the
relation
\weq{7}{j_\lambda j_{\lambda+1} j_{\lambda+2} \ldots j_\mu = k\, m,}
from which a very remarkable consequence can be drawn.

Suppose that the irreducible function $f(x)$ becomes reducible by successive adjunction
of roots of cyclic equations; then, by §.\,176, Bd. I, the composition
series (1) can be assumed so ordered that the indices
$j_1, j_2, \ldots, j_\lambda$ are primes. Thus, since $h$ is a divisor of
$j_\lambda$, we have $j_\lambda = h$, and $j_\lambda$ is, by (3), a prime
dividing $m$.

If now, besides $j_\lambda$, a second prime $p$ divides $m$, then
by (7) one of the factors $j_\lambda, j_{\lambda+1}, j_{\lambda+2}, \ldots, j_\mu$
must be divisible by $p$, and so they certainly cannot all be equal to $j_\lambda$.

From this, however, it follows by Theorem III, §.\,7, that one can find a composition series of
$P_\lambda$:
\weq{8}{P_\lambda,\ P_{\lambda+1},\ \ldots,\ P_{\mu-1},\ 1}
such that among the groups $P_\lambda, P_{\lambda+1}, \ldots,
P_{\mu-1}$ one, say $P_\nu$, is a \gesperrt{normal divisor of} $P$.

But this $P_\nu$, as a divisor of the intransitive group $P_\lambda$, is itself
intransitive, and hence $P$ itself must be \gesperrt{imprimitive} (Bd. I, §.\,158, 2).

We state this in the following form as a theorem:

\gesperrt{1. If several different primes divide the degree of an
irreducible equation, then}
%CONT>>


% ===== scan leaf 316, printed 295 =====
% !PRINTED 295
\weberpage{295}
%<<CONT
\gesperrt{this equation can become reducible by successive adjunction of roots
of cyclic equations only if it is imprimitive.}

We can state something further about the manner of the reduction, its possibility
being presupposed. If $P_\nu$ is the first group of the sequence (8) that is a
normal divisor of $P$, then by Theorem III, §.\,7, just cited, with the correct
arrangement of the composition series:
\weq{9}{j_\lambda = j_{\lambda+1} \cdots = j_{\nu+1},}
hence all equal to the same prime, and the resolvent $\chi = 0$, by which the
group $P$ is reduced to $P_\nu$, consequently has a metacyclic group.

If we denote by $A, B, \ldots, S$ the systems of intransitivity of the group
$P_\nu$, whose number shall be $s$, and set
\weq{10}{m = r\, s,}
then by adjunction of a root of $\chi = 0$ the function $f(x)$ will split into
$s$ factors of degree $r$.

The $A, B, \ldots, S$ are systems of imprimitivity of $P$ (Bd. I, §.\,158, 2.).

If we denote by $Q$ the totality of all permutations of $P$ which leave the
individual systems $A, B, \ldots, S$ in their places and interchange only the
elements of the systems among themselves, then $Q$, as we saw in §.\,158 of the
first volume, is a normal divisor of $P$, and by adjunction of the roots of an
auxiliary equation $\varphi(y) = 0$ of degree $s$, $f(x)$ splits into $s$
factors of degree $r$, to which the individual systems $A, B, \ldots, S$ belong
as roots:
\[
f(x) = f(x, y_\alpha)\, f(x, y_\beta)\, f(x, y_\gamma) \ldots
\]
Since by the auxiliary equation $\varphi(y) = 0$ the group $P$ is reduced to
$Q$, the group of $\varphi(y) = 0$ is isomorphic with $P/Q$. On the other hand,
however, the intransitive group $P_\nu$ is also a divisor of $Q$, since the
systems $A, B, \ldots, S$ are not displaced by $P_\nu$, and we can therefore
apply the theorem §.\,73, 3. to the three groups $P, Q, P_\nu$. Since the index
series of $P/P_\nu$ by hypothesis consists entirely of primes, by that
theorem the same holds of $P/Q$. This group is metacyclic, and hence the
auxiliary equation $\varphi(y) = 0$ is also metacyclic. There follows, then,
the theorem:
%CONT>>


% ===== scan leaf 317, printed 296 =====
% !PRINTED 296
\weberpage{296}
%<<CONT
\gesperrt{2. If an irreducible equation, in whose degree more than one prime
divides, splits into factors by successive adjunction of radicals, then a
splitting into $s$ factors of degree $r$ is brought about by adjunction of the
roots of a metacyclic equation of degree $r$.}

This theorem contains as a special case the theorem due to \gesperrt{Abel},
that an irreducible equation whose degree $m$ is not the power of a prime can
be solvable by radicals only if it splits into factors by the adjunction of
the roots of a solvable equation of lower degree, whose degree is a divisor of
$m$\textsuperscript{1)}. With this the question of the solution of an equation
by radicals, or even of the mere reduction of an equation by radicals, is
extraordinarily simplified. One can in fact restrict the further investigation
to equations whose degree is a prime or a prime power, because all other cases
are reduced to these by repeated application of theorem 2.

Thus, for example, the question of all irreducible equations of degree $6$
solvable by radicals admits a downright trivial answer:

\gesperrt{In order to obtain all metacyclic equations of degree $6$ in any field
$\Omega$, adjoin to the field $\Omega$ a square root and form in the extended
field all cubic equations, or adjoin the root of a cubic equation and form in
the extended field all quadratic equations.}

As an example of such an equation we cite the one treated by \gesperrt{Hesse},
on which depend the circular sections of a surface of degree $2$ not
referred to its principal axes. This problem is solved by square roots, if the
principal axes have first been determined by the solution of a cubic
equation\textsuperscript{2)}.

\noindent\rule{4cm}{0.4pt}\par
\textsuperscript{1)} \gesperrt{Abel} gives the theorem without proof in the
memoir ``Sur la rés. algébr. des équations''. Oeuvres complètes
1881, Bd. II, S. 217. Cf. also the memoir by \gesperrt{Galois} in
\gesperrt{Liouville's} Journal, Bd. 11.\par
\textsuperscript{2)} \gesperrt{Hesse}, ``Ueber die Auflösung derjenigen
Gleichungen $6^{\mathrm{ten}}$ Grades etc.'' Crelle's Journal, Bd. 41
(1851).


% ===== scan leaf 318, printed 297 =====
% !PRINTED 297
\weberpage{297}
\pdfbookmark[2]{§ 75. Metacyclic equations whose degree is a power of a prime}{bm88}
\begin{center}
\S.\,75.\\
\gesperrt{Metacyclic equations whose degree is a power of a prime.}
\end{center}

After the last theorems, the interest of further investigations on metacyclic equations concentrates chiefly on the case in which the degree of the equation is a power $p^k$ of a prime $p$. We assume the equation to be irreducible, and restrict ourselves to the consideration of \gesperrt{primitive equations}; for the solution of the imprimitive ones reduces to the successive solution of two (or more) primitive equations, whose degrees are likewise powers of $p$, and which, if the original equation is metacyclic, must also be metacyclic.

We now assume, then, that $P$ is the group of an irreducible primitive metacyclic equation $f(x) = 0$ of degree $p^k$. By Bd. I, \S.\,158, 2., not only $P$ itself, but all its normal divisors (with the exception of the unit group) must still be \gesperrt{transitive}. We now apply Theorem IV., proved in general in \S.\,8, according to which $P$ possesses a normal divisor $Q$ consisting entirely of commutative elements; and if several such divisors are present, we understand by $Q$ one of them whose degree is as low as possible, but still greater than 1. This group $Q$ is therefore likewise still transitive. But $Q$ can have no further proper divisor different from unity that is at the same time a normal divisor of $P$, because otherwise this divisor would take the place of $Q$.

When, by the appropriate adjunctions, the group of our equation has been reduced from $P$ to $Q$, then $f(x) = 0$ has become an \gesperrt{Abelian equation} in the extended domain of rationality, and since it has still remained irreducible, by Bd. I, \S.\,162 the degree of the equation is equal to the degree of the group; i.\,e. $Q$ is an \gesperrt{Abelian group} of degree $p^k$. The degrees of all the elements of $Q$ are therefore powers of $p$. We wish to prove that, apart from the unit element, \gesperrt{only elements of degree $p$ itself occur in $Q$.}

Let us assume that $p^\lambda$ is the highest degree that occurs among the elements of $Q$, and that $\lambda > 1$. Then all the elements of $Q$ whose degree is a divisor of $p^{\lambda-1}$ form a
%CONT>>


% ===== scan leaf 319, printed 298 =====
% !PRINTED 298
\weberpage{298}
\pdfbookmark[2]{§ 76. Representation of the Abelian group Q}{bm89}
%<<CONT
divisor $Q'$ of $Q$ which is different both from $Q$ itself and from the unit group. But this divisor $Q'$ must be a normal divisor of $P$, because, if $\pi$ is contained in $Q$ and $\gamma$ in $P$, then $\gamma^{-1}\pi\gamma$, which is of the same degree as $\pi$, must likewise be contained in $Q$, and hence belongs to $Q'$ if $\pi$ belongs to $Q'$. But since, as remarked above, $Q$ can have no proper divisor that is greater than the unit group and at the same time a normal divisor of $P$, we must have $\lambda = 1$.

\begin{center}
\S.\,76.\\
\gesperrt{Representation of the Abelian group $Q$.}
\end{center}

The considerations of the preceding section have shown that a metacyclic group $P$, which can occur as the \gesperrt{Galois} group of an irreducible primitive equation $f(x) = 0$ of degree $n = p^k$, must contain an \gesperrt{Abelian group} $Q$ as a normal divisor, whose degree is $p^k$, and which, apart from the unit element, contains only elements of degree $p$.

By \S.\,9 this group $Q$ can be represented by a basis consisting of $k$ elements $A_1, A_2, \ldots, A_k$ of degree $p$, so that every element of $Q$ takes the form:
\weq{1}{A_1^{z_1} A_2^{z_2} \ldots A_k^{z_k},}
and here $z_1, z_2, \ldots, z_k$ run, independently of one another, each through a complete system of residues to the modulus $p$. Now, in order to find the corresponding permutation group, we start from an arbitrary root $x$ of $f(x)$, denote the root into which $x$ passes under the permutation (1) by
\weq{2}{[z_1, z_2, \ldots, z_k],}
and stipulate that this symbol shall not change its meaning when the numbers $z_1, z_2, \ldots, z_k$ are altered arbitrarily by multiples of $p$. To the root $x$ from which we started there then belongs the symbol $[0, 0, \ldots, 0]$, and by the symbol (2) every root of $f(x)$ is represented once and only once.

If $x$ passes into $x'$ under a permutation $A'$, and $x'$ into $x''$ under $A''$, then $x$ passes into $x''$ under the permutation $A' A''$. From this it follows, however, that $[z_1, z_2, \ldots, z_k]$, under the permutation
\weq{3}{A = A_1^{\alpha_1} A_2^{\alpha_2} \ldots A_k^{\alpha_k}}
%CONT>>


% ===== scan leaf 320, printed 299 =====
% !PRINTED 299
\weberpage{299}
\pdfbookmark[2]{§ 77. Analytic Representation of the Permutations}{bm90}
%<<CONT
into $[z_1+\alpha_1, z_2+\alpha_2, \ldots, z_k+\alpha_k]$. Accordingly we can also denote the permutation $A$ thus:
\weq{4}{\begin{pmatrix} z_1, & z_2, & \ldots, & z_k \\ z_1+\alpha_1, & z_2+\alpha_2, & \ldots, & z_k+\alpha_k \end{pmatrix},}
and one obtains the whole group $Q$, if one lets $\alpha_1, \alpha_2, \ldots, \alpha_k$ each run through a complete system of residues to the modulus $p$.

The group $Q$, whose mode of formation is now established, must be a normal divisor of $P$, and from this a general form can be derived in which all permutations of $P$ must be contained.

For this we require a lemma on the analytic representation of the permutations, which we first derive.

\begin{center}
\S.\,77.\\
\gesperrt{Analytic Representation of the Permutations.}
\end{center}

\gesperrt{Lemma. If under any permutation of the quantities (2), \S.\,76, $z_1$ passes into $z_1'$, $z_2$ into $z_2'$, \ldots, $z_k$ into $z_k'$, then one can always set
\weq{1}{
\begin{aligned}
z_1' &\equiv \varphi_1(z_1, z_2, \ldots, z_k)\\
z_2' &\equiv \varphi_2(z_1, z_2, \ldots, z_k)\\
&\;\cdots\cdots\cdots\cdots\\
z_k' &\equiv \varphi_k(z_1, z_2, \ldots; z_k)
\end{aligned}
\quad (\mathrm{mod}\ p).
}
where $\varphi_1, \varphi_2, \ldots, \varphi_k$ denote integral rational functions of the variables $z_1, z_2, \ldots, z_k$ with integral coefficients, which do not exceed the degree $p - 1$ with respect to any of the variables.}

The theorem is a generalisation of the corresponding theorem in \S.\,180 of the first volume. There the theorem is proved for $k = 1$. We preface the general proof, which rests on complete induction, with the following considerations.

Since in (1) we have to put for the arguments $z_i$ all combinations of $k$ sequences to the modulus $p$, there result $p^k$ such systems. Each of the functions $\varphi$ has $p^k$ undetermined coefficients, and if therefore the $z_k'$ are given, we obtain $kp^k$ linear congruences for just as many unknown numbers, namely the coefficients of the $\varphi_k$. It remains only to show
%CONT>>


% ===== scan leaf 321, printed 300 =====
% !PRINTED 300
\weberpage{300}
%<<CONT
that these congruences are independent of one another. Here we can consider each of the functions $\varphi$ by itself. If we therefore set
\weq{2}{z' \equiv \varphi(z_1, z_2, \ldots, z_k) \quad (\mathrm{mod}\ p),}
then, if we know for each combination of the $z_i$ the corresponding value of $z'$, we have $p^k$ linear congruences for the determination of the coefficients of $\varphi$. It is to be proved that these congruences are always solvable, where we may assume that this possibility has already been established for functions of fewer variables.

If we now order $\varphi$ according to powers of the variable $z_1$, we obtain
\weq{3}{z' \equiv \chi_0 z_1^{p-1} + \chi_1 z_1^{p-2} + \cdots + \chi_{p-1} \quad (\mathrm{mod}\ p),}
in which the coefficients $\chi_0, \chi_1, \ldots, \chi_{p-1}$ are functions of the same kind as $\varphi$, except that they depend on one variable fewer.

If we hold fixed any one of the combinations of the $z_2, \ldots, z_k$, and set $z_1 = 0, 1, \ldots, p - 1$, we obtain from (3) a system of $p$ linear congruences whose determinant is not divisible by $p$, and from it we can, as in \S.\,180 of the first volume, determine the values of $\chi_0, \chi_1, \ldots, \chi_{p-1}$ for this combination of the $z_2, \ldots, z_k$; and hence for every combination of the variables the values of the function $\chi_i$ (with respect to the modulus $p$) are known. By hypothesis, however, we can determine from these the coefficients of the functions $\chi_i$, which indeed depend on only $k - 1$ variables, and with this the lemma is proved.

Accordingly we can represent every permutation of the roots $[z_1, \ldots, z_k]$ thus:
\weq{4}{\begin{pmatrix} z_1, & z_2, & \ldots \\ \varphi_1(z_1, z_2, \ldots), & \varphi_2(z_1, z_2, \ldots), & \ldots \end{pmatrix},}
or in a still more abbreviated notation:
\weq{5}{\begin{pmatrix} z \\ \varphi(z) \end{pmatrix}.}

In the symbol (4) we can, without changing its meaning, replace $z_1, z_2, \ldots$ by any other combination $z_1', z_2', \ldots$, and if therefore, by the lemma,
\[
z_i' = \psi_i(z_1, z_2, \ldots)
\]
%CONT>>


% ===== scan leaf 322, printed 301 =====
% !PRINTED 301
\weberpage{301}
\pdfbookmark[2]{§ 78. Representation of the metacyclic group P}{bm91}
%<<CONT
then we can also represent (4) or (5) as follows:

\weq{6}{\begin{pmatrix} \psi_1(z_1, z_2, \dots), & \psi_2(z_1, z_2, \dots), \dots \\ \varphi_1(\psi_1, \psi_2, \dots), & \varphi_2(\psi_1, \psi_2, \dots), \dots \end{pmatrix} = \begin{pmatrix} \psi(z) \\ \varphi[\psi(z)] \end{pmatrix}.}

This leads to the composition of the permutations:

\weq{7}{\begin{pmatrix} z \\ \psi(z) \end{pmatrix}\begin{pmatrix} z \\ \varphi(z) \end{pmatrix} = \begin{pmatrix} z \\ \varphi[\psi(z)] \end{pmatrix}.}

\begin{center}
§.\,78.

\gesperrt{Representation of the metacyclic group $P$.}
\end{center}

Now the group $Q$ is to be a normal divisor of the group $P$; that is, if $A$ is a permutation from $Q$ and $B$ a permutation from $P$, then a second permutation $A'$ from $Q$ is to be determinable so that

\weq{1}{AB = BA'}

holds. If, in the abbreviated notation of the preceding section, we set

\[
A = \begin{pmatrix} z \\ z+\alpha \end{pmatrix}, \quad A' = \begin{pmatrix} z \\ z+\alpha' \end{pmatrix}, \quad B = \begin{pmatrix} z \\ \varphi(z) \end{pmatrix},
\]

then we obtain

\[
AB = \begin{pmatrix} z \\ \varphi(z+\alpha) \end{pmatrix}, \quad BA' = \begin{pmatrix} z \\ \varphi(z)+\alpha' \end{pmatrix},
\]

or

\[
\varphi(z+\alpha) \equiv \varphi(z) + \alpha'.
\]

This congruence, however, is only an abbreviated symbol for the system of congruences to the modulus $p$:

\weq{2}{\begin{aligned}
\varphi_1(z_1+\alpha_1, z_2+\alpha_2, \dots) &\equiv \varphi_1(z_1, z_2, \dots) + \alpha_1' \\
\varphi_2(z_1+\alpha_1, z_2+\alpha_2, \dots) &\equiv \varphi_2(z_1, z_2, \dots) + \alpha_2' \\
&\cdots\cdots\cdots\cdots\cdots\cdots\cdots\cdots\cdots\cdots
\end{aligned}}

Herein the system of numbers $\alpha_1, \alpha_2, \dots$ may be given arbitrarily to the modulus $p$; $\alpha_1', \alpha_2', \dots$ are thereby determined.

If in the first of the formulae (2) we set one of the numbers $\alpha_1, \alpha_2, \dots$ in turn equal to 1 and the rest equal to 0, the result may be:

\weq{3}{\begin{aligned}
\varphi_1(z_1+1, z_2 \dots) &\equiv \varphi_1(z_1, z_2, \dots) + \alpha_{1,1} \\
\varphi_1(z_1, z_2+1, \dots) &\equiv \varphi_1(z_1, z_2, \dots) + \alpha_{1,2} \\
&\cdots\cdots\cdots\cdots\cdots\cdots\cdots\cdots\cdots\cdots
\end{aligned}}


% ===== scan leaf 323, printed 302 =====
% !PRINTED 302
\weberpage{302}
If one applies the first of these formulae $\alpha_1$ times, the second $\alpha_2$ times
in succession, and so on, it follows that:

\[
\begin{aligned}
\varphi_1(z_1+\alpha_1, z_2, \dots) &\equiv \varphi_1(z_1, z_2, \dots) + \alpha_1\alpha_{1,1} \\
\varphi_1(z_1, z_2+\alpha_2, \dots) &\equiv \varphi_1(z_1, z_2, \dots) + \alpha_2\alpha_{1,2} \\
&\cdots\cdots\cdots\cdots\cdots\cdots\cdots\cdots\cdots,
\end{aligned}
\]

and if in the first one changes $z_2$ into $z_2+\alpha_2$, and applies the
second, and continues in this way, one finally obtains:

\weq{4}{\begin{aligned}
\varphi_1(z_1+\alpha_1, z_2+\alpha_2, \dots) &\equiv \\
\varphi_1(z_1, z_2, \dots) + \alpha_1\alpha_{1,1} + \alpha_2\alpha_{1,2} + \cdots
\end{aligned}}

In this we now set $z_1, z_2, \cdots = 0$ and then write in place of
$\alpha_1, \alpha_2, \dots$ again $z_1, z_2, \dots$. Further, we denote
$\varphi_1(0, 0, \dots)$ by $\alpha_1$ and thus obtain from (4):

\weq{5}{\varphi_1(z_1, z_2, \dots) \equiv \alpha_1 + z_1\alpha_{1,1} + z_2\alpha_{1,2} + \cdots,}

that is, $\varphi_1$ must be a \gesperrt{linear function}.

In exactly the same way we proceed with all the congruences (2) and thus arrive at the following final result:

\gesperrt{I. All permutations of the Galois group $P$ of a primitive irreducible metacyclic equation of degree $p^k$}

\[
\begin{pmatrix} z_1, z_2, \dots, z_k \\ z_1', z_2', \dots, z_k' \end{pmatrix}
\]

\gesperrt{are of the form}

\weq{6}{\begin{aligned}
z_1' &\equiv \alpha_{1,1}z_1 + \alpha_{1,2}z_2 + \cdots + \alpha_{1,k}z_k + \alpha_1 \\
z_2' &\equiv \alpha_{2,1}z_1 + \alpha_{2,2}z_2 + \cdots + \alpha_{2,k}z_k + \alpha_2 \quad (\mathrm{mod}\ p). \\
&\cdots\cdots\cdots\cdots\cdots\cdots\cdots\cdots\cdots\cdots
\end{aligned}}

\[
z_k' \equiv \alpha_{k,1}z_1 + \alpha_{k,2}z_2 + \cdots + \alpha_{k,k}z_k + \alpha_k
\]

The system of congruences (6) represents a permutation always and only when the determinant

\[
\sum \pm \alpha_{1,1}\alpha_{2,2}\cdots\alpha_{k,k}
\]

is not divisible by $p$, because then and only then does there also conversely belong to every system of the $z'$ a determinate system of the $z$ (with respect to the modulus $p$).

The totality of all permutations (6) is in fact a group, which is called the \gesperrt{general linear congruence group}, and in it the metacyclic groups $P$ must be con-
%CONT>>


% ===== scan leaf 324, printed 303 =====
% !PRINTED 303
\weberpage{303}
%<<CONT
tained. But not all linear congruence groups are, conversely, also metacyclic, and to separate these off from them is yet a further problem, which up to now has been solved only in special cases $^{1)}$.

But in order at least to formulate the problem more precisely, we add the following considerations.

We denote the complete linear congruence group (6) by $R$, and separate off from it the group, contained within it, of the linear homogeneous substitutions

\weq{7}{\begin{aligned}
z_1' &\equiv \alpha_{1,1}z_1 + \alpha_{1,2}z_2 + \cdots + \alpha_{1,k}z_k \\
z_2' &\equiv \alpha_{2,1}z_1 + \alpha_{2,2}z_2 + \cdots + \alpha_{2,k}z_k \quad (\mathrm{mod}\ p) \\
&\cdots\cdots\cdots\cdots\cdots\cdots\cdots\cdots\cdots\cdots \\
z_k' &\equiv \alpha_{k,1}z_1 + \alpha_{k,2}z_2 + \cdots + \alpha_{k,k}z_k
\end{aligned}}

with non-vanishing determinant, which is called the \gesperrt{linear homogeneous congruence group}, and which we shall denote by $S$. Further, in (6) there is contained the $k$-fold cyclic group $Q$, which the substitutions

\weq{8}{z_1' = z_1 + \alpha_1, \quad z_2' = z_2 + \alpha_2, \dots, \quad z_k' = z_k + \alpha_k}

constitute.

If now $\sigma$ denotes an element of $S$ and $\gamma$ an element of $Q$, it follows without difficulty from the rule of composition \S.\,77, (7) that every substitution $\pi$ of the group $R$ can be represented \gesperrt{in one and only one way} in each of the two forms

\weq{9}{\pi = \sigma\gamma = \gamma'\sigma}

where, if $\sigma, \gamma$ are determined by (7), (8), $\gamma'$ denotes the substitution

\weq{10}{z_1' = z_1 + \alpha_1', \quad z_2' = z_2 + \alpha_2', \dots, \quad z_k' = z_k + \alpha_k'}

where the $\alpha_i'$ are derived from the $\alpha_i$ by the substitution $\sigma$.

$Q$ is a normal divisor of $R$; for if $\gamma_1$ is an arbitrary element of $Q$, then by repeated application of (9):

\[
\sigma\gamma\gamma_1 = \gamma'\gamma_1'\sigma = \gamma'\gamma_1'\sigma\gamma^{-1}\gamma = \gamma'\gamma_1'\gamma'^{-1}\sigma\gamma,
\]

hence, if

\[
\sigma\gamma = \pi, \quad \gamma'\gamma_1'\gamma'^{-1} = \gamma_1''
\]

is put,

\[
\pi\gamma_1 = \gamma_1''\pi,
\]
%CONT>>

\begin{center}
\rule{4cm}{0.4pt}
\end{center}

$^{1)}$ \gesperrt{C. Jordan}, Liouville's Journ. de Math\'em. S\'er. (2), Bd. XIII.


% ===== scan leaf 325, printed 304 =====
% !PRINTED 304
\weberpage{304}
%<<CONT
and consequently

\weq{11}{\pi^{-1}Q\pi = Q.}

Now if $\pi$ runs through any divisor $R_1$ of $R$ which for its part contains the group $Q$, then $Q$ is a normal divisor of $R_1$, and the substitution $\sigma$ occurring in (9) must run through a divisor $S_1$ of $S$. One may set:

\[
R_1 = S_1 Q = Q S_1,
\]

and since the substitutions contained in the form (9) are all different from one another, the degree of $R_1$ is equal to the product of the degrees of $S_1$ and $Q$. If $R_2$ is a normal divisor of $R_1$ which likewise still contains $Q$ (as a normal divisor), then in the same way

\[
R_2 = S_2 Q,
\]

and $S_2$ is a normal divisor of $S_1$.

For by hypothesis, for every element $\pi_1$ of $R_1$,

\[
\pi_1^{-1}S_2 Q \pi_1 = S_2 Q,
\]

hence by (11):

\[
\pi_1^{-1}S_2 Q \pi_1 = \pi_1^{-1}S_2\pi_1 Q = S_2 Q,
\]

and consequently

\[
\pi_1^{-1}S_2\pi_1 = S_2,
\]

and since this holds also for $\pi_1 = \sigma_1$, where $\sigma_1$ is contained in $S_1$, $S_2$ is a normal divisor of $S_1$. Conversely, if $S_2$ is a normal divisor of $S_1$, then $S_2 Q$ is a normal divisor of $S_1 Q$. For if $S_2\sigma_1 = \sigma_1 S_2$, it follows that

\[
\sigma_1 S_2 Q = S_2\sigma_1 Q = S_2 Q\sigma_1,
\]

and from this

\[
\gamma\sigma_1 S_2 Q = \gamma S_2 Q\sigma_1 = S_2\gamma' Q\sigma_1 = S_2 Q\sigma_1 = S_2 Q\gamma^{-1}\gamma\sigma_1 = S_2 Q\gamma\sigma_1
\]

(because $\gamma' Q = Q\gamma^{-1} = Q$); i.e.

\[
\pi_1 S_2 Q = S_2 Q\pi_1,
\]

q. e. d. But from this the following results:

\gesperrt{II. If $P$ is the Galois group of a primitive irreducible metacyclic equation of degree $p^k$, then $P$ is representable in the form:}

\[
P = TQ,
\]

\gesperrt{wherein $T$ is a metacyclic group contained in the congruence group $S$.}


% ===== scan leaf 326, printed 305 =====
% !PRINTED 305
\weberpage{305}

The problem of finding all these metacyclic equations is
thus reduced to finding all metacyclic divisors of the homogeneous
congruence group $S$.

If we take as an example the case $p = 3$, $k = 2$, and thus
ask for the metacyclic equations of degree 9, then by this
theorem we are concerned with the group $S$ of the linear
substitutions taken modulo 3

\weq{12}{\sigma = \begin{pmatrix} a, & b \\ c, & d \end{pmatrix},}

with which we occupied ourselves, from a somewhat different
point of view, in \S.\,66.

There we first considered a divisor $E$ of $S$, which consists
of all substitutions $\sigma$ with the condition $a\,d - b\,c$
$\equiv 1$ (mod 3), and moreover combined two substitutions

\weq{13}{\begin{pmatrix} a, & b \\ c, & d \end{pmatrix}, \quad \begin{pmatrix} -a, & -b \\ -c, & -d \end{pmatrix}}

into a single element. The group so specialised was of
degree 12 and had a normal divisor of index 3:

\[ G = \begin{pmatrix} 1, & 0 \\ 0, & 1 \end{pmatrix}, \begin{pmatrix} 0, & 1 \\ -1, & 0 \end{pmatrix}, \begin{pmatrix} -1, & 1 \\ 1, & 1 \end{pmatrix}, \begin{pmatrix} 1, & 1 \\ 1, & -1 \end{pmatrix}, \]

from which $E$ was composed thus:

\[ E = G + \begin{pmatrix} 1, & 0 \\ 1, & 1 \end{pmatrix} G + \begin{pmatrix} 1, & 0 \\ -1, & 1 \end{pmatrix} G. \]

In the group $S$, however, the two substitutions (13)
count as distinct, and moreover the substitutions
$\sigma$ whose determinant is $a\,d - b\,c \equiv -1$ (mod 3) must be
added, whereby the degree of the group rises to 48.
$G$ is extended, by the change of the signs of
all its elements, to a group of degree 8, and $E$ to a
group of degree 24, of which the extended $G$ is normal divisor.
Here the whole group $S$ is now \gesperrt{metacyclic}, as one
sees from the following composition, in which the extended
group $G$ is denoted by $S_2$ and the extended group $E$ by $S_1$:

\[ S_4 = \begin{pmatrix} 1, & 0 \\ 0, & 1 \end{pmatrix}, \begin{pmatrix} -1, & 0 \\ 0, & -1 \end{pmatrix} \]
%CONT>>


% ===== scan leaf 327, printed 306 =====
% !PRINTED 306
\weberpage{306}
\pdfbookmark[2]{§ 79. Ternary linear congruence group for the modulus 2}{bm92}
%<<CONT
\[
\begin{aligned}
S_3 &= S_4 + \begin{pmatrix} -1, & 1 \\ 1, & 1 \end{pmatrix} S_4 \\
S_2 &= S_3 + \begin{pmatrix} 0, & 1 \\ -1, & 0 \end{pmatrix} S_3 \\
S_1 &= S_2 + \begin{pmatrix} 1, & 0 \\ 1, & 1 \end{pmatrix} S_2 + \begin{pmatrix} 1, & 0 \\ -1, & 1 \end{pmatrix} S_2 \\
S &= S_1 + \begin{pmatrix} -1, & 0 \\ 0, & 1 \end{pmatrix} S_1.
\end{aligned}
\]

We accordingly have the composition and index series of $S$

\[
\begin{gathered}
S,\ S_1,\ S_2,\ S_3,\ S_4, \\
2 \quad 3 \quad 2 \quad 2
\end{gathered}
\]

\gesperrt{and we thus arrive here at the result that all equations
of degree 9 with linear congruence group
are metacyclic.}

\begin{center}
\S.\,79.

\gesperrt{Ternary linear congruence group for the modulus 2.}
\end{center}

Manifold interesting relations arise when
we consider the ternary linear congruence group $R$ for the modulus 2,
which can be conceived as a transitive permutation group of eight
elements, and which must contain the group of the metacyclic
equations of degree 8. By what was
proved above, it is above all a matter of considering the homogeneous
congruence group $S$, which consists of the elements

\weq{1}{\sigma = \begin{pmatrix} a, & b, & c \\ a_1, & b_1, & c_1 \\ a_2, & b_2, & c_2 \end{pmatrix}}

where the figures $a, b, c, \ldots$ are to be taken according to the modulus 2, hence $= 0$
or $= 1$, and which are compounded according to the rule given
in \S.\,37. Only such systems
of the numbers $a, b, c, \ldots$ come into consideration here whose determinant
is $= 1$, i.\,e. odd.

To determine the degree of the group $S$, observe that
the first row of $\sigma$ can have seven different forms:

\[
(1,0,0),\ (0,1,0),\ (0,0,1),\ (0,1,1),\ (1,0,1),\ (1,1,0),\ (1,1,1).
\]


% ===== scan leaf 328, printed 307 =====
% !PRINTED 307
\weberpage{307}

For the first assumption $(a, b, c) = (1, 0, 0)$ the determinant
of $\sigma$ becomes equal to $b_1 c_2 - c_1 b_2$, which can become $= 1$ in six
different ways, namely:

\[
\begin{pmatrix} b_1, & c_1 \\ b_2, & c_2 \end{pmatrix} = \begin{pmatrix} 1, & 0 \\ 0, & 1 \end{pmatrix}, \begin{pmatrix} 1, & 0 \\ 1, & 1 \end{pmatrix}, \begin{pmatrix} 0, & 1 \\ 1, & 0 \end{pmatrix}, \begin{pmatrix} 0, & 1 \\ 1, & 1 \end{pmatrix}, \begin{pmatrix} 1, & 1 \\ 1, & 0 \end{pmatrix}, \begin{pmatrix} 1, & 1 \\ 0, & 1 \end{pmatrix}.
\]

Then $a_1, a_2$ can further be assumed in four different
ways. Hence, when the first row is $(1, 0, 0)$,
one has 24 different forms of $\sigma$. The same number results
for the assumption of the first row in the forms $(0, 1, 0), (0, 0, 1)$.
But one obtains the same for the other assumptions
about the first row. For if this is, e.\,g., $(0, 1, 1)$, then

\[
a_1 (b_2 - c_2) - a_2 (b_1 - c_1) \equiv 1 \text{ (mod 2)}
\]

must hold, which again yields six possibilities for $a_1, b_1 - c_1, a_2, b_2 - c_2$,
and since one can assume $c_1, c_2$ in four ways, one
again obtains 24 possibilities. Finally, when the
first row is $(1, 1, 1)$,

\[
(a_1 - c_1)(b_2 - c_2) - (a_2 - c_2)(b_1 - c_1) \equiv 1 \text{ (mod 2)}
\]

must hold, which leads to the same number. The total number of all different
substitutions $\sigma$, i.\,e. the degree of $S$, is therefore 168.
The number 168 as the degree number of a group has already
met us once, namely with the group $L_7$ (\S.\,66, 72), and hence
the conjecture suggests itself that the groups $S$ and $L_7$
might be isomorphic. If this conjecture is confirmed,
the investigation of the group $S$ would thereby be essentially
facilitated, since we already know the divisors of the
group $L_7$.

But it is Theorem I. in \S.\,72 that serves to test this
conjecture.

To this end we seek to select generating elements $\chi, \omega, \Theta, \tau$
of the group $S$ in such a way that they satisfy the characteristic
conditions of Theorem I., \S.\,72. This can be done in
several ways. A general procedure for it is, to be sure, lacking,
but it succeeds easily by a few trials.

One first chooses an element of degree 7 arbitrarily
and forms from it the period, say
%CONT>>


% ===== scan leaf 329, printed 308 =====
% !PRINTED 308
\weberpage{308}
\pdfbookmark[2]{§ 80. Resolvents of the equation of the 8th degree}{bm93}
%<<CONT
\weq{2}{
\begin{aligned}
\tau &= \begin{pmatrix} 1, & 0, & 1 \\ 1, & 0, & 0 \\ 0, & 1, & 0 \end{pmatrix}, &
\tau^2 &= \begin{pmatrix} 1, & 1, & 1 \\ 1, & 0, & 1 \\ 1, & 0, & 0 \end{pmatrix}, &
\tau^3 &= \begin{pmatrix} 0, & 1, & 1 \\ 1, & 1, & 1 \\ 1, & 0, & 1 \end{pmatrix}. \\
\tau^4 &= \begin{pmatrix} 1, & 1, & 0 \\ 0, & 1, & 1 \\ 1, & 1, & 1 \end{pmatrix}, &
\tau^5 &= \begin{pmatrix} 0, & 0, & 1 \\ 1, & 1, & 0 \\ 0, & 1, & 1 \end{pmatrix}, &
\tau^6 &= \begin{pmatrix} 0, & 1, & 0 \\ 0, & 0, & 1 \\ 1, & 1, & 0 \end{pmatrix}
\end{aligned}
}

Next one selects, among the elements of degree 2, of which
there are 21, a suitable one for $\omega$; this one, for example, is suitable

\weq{3}{\omega = \begin{pmatrix} 1, & 0, & 0 \\ 0, & 1, & 0 \\ 0, & 1, & 1 \end{pmatrix};}

and then, according to the formulae \S.\,72, (17), the elements $\Theta$
and $\chi$ are determined:

\weq{4}{\Theta = \begin{pmatrix} 1, & 1, & 1 \\ 0, & 1, & 0 \\ 0, & 1, & 1 \end{pmatrix}, \quad \Theta^2 = \begin{pmatrix} 1, & 1, & 0 \\ 0, & 1, & 0 \\ 0, & 0, & 1 \end{pmatrix}, \quad \Theta^3 = \begin{pmatrix} 1, & 0, & 1 \\ 0, & 1, & 0 \\ 0, & 1, & 1 \end{pmatrix}}

\weq{5}{\chi = \begin{pmatrix} 1, & 0, & 0 \\ 0, & 0, & 1 \\ 0, & 1, & 1 \end{pmatrix}, \quad \chi^2 = \begin{pmatrix} 1, & 0, & 0 \\ 0, & 1, & 1 \\ 0, & 1, & 0 \end{pmatrix}.}

It is then a matter of a simple calculation to verify the formulae
of Theorem I., \S.\,72, whereby the agreement
of the groups $S$ and $L_7$ is established.

From this it follows that the group $S$, like $L_7$, is \gesperrt{simple}.

\begin{center}
\S.\,80.

\gesperrt{Resolvents of the Equation of Degree 8.}
\end{center}

In order now to make the application to the theory of equations
of degree 8, we denote as before by $Q$ the
cyclic group of degree 8: % [sic] the displayed translation group is elementary abelian of order 8, not cyclic.

\[
\begin{pmatrix} z_1, & z_2, & z_3 \\ z_1 + h_1, & z_2 + h_2, & z_3 + h_3 \end{pmatrix} \text{ (mod 2)},
\]

or more briefly

\[
(z_i,\ z_i + h_i),
\]

and by $S$ the ternary congruence group considered in \S.\,79,
and we set

\[
R = S\,Q,
\]
%CONT>>


% ===== scan leaf 330, printed 309 =====
% !PRINTED 309
\weberpage{309}
%% [page furniture, not body text -- commented out 2026-08-10] Page furniture (running
%% head, folio, signature) is excluded from this edition by standing policy.
%% \noindent \S.\,80. Equations of the 8th Degree.

%<<CONT
so that $R$ is the entire ternary linear congruence group for the modulus 2.

We consider a system of eight independent variables $X_{z_1,z_2,z_3}$, in which the indices $z_1, z_2, z_3$ are to be taken with respect to the modulus 2.

From these we form the seven functions (\gesperrt{Lagrange} resolvents, Bd.\,I, \S.\,164).
\weq{1}{\overset{z}{\Sigma} (-1)^{\Sigma z\xi} X_{z_1,z_2,z_3} = \Psi_{\xi_1,\xi_2,\xi_3},}
in which, for abbreviation,
\[
\Sigma z\xi = z_1\xi_1 + z_2\xi_2 + z_3\xi_3
\]
is set, and $\xi_1, \xi_2, \xi_3$ each run through a complete system of residues with respect to the modulus 2, but with the exclusion of the combination $0, 0, 0$.

If we apply the substitutions of the group $Q$ to the indices $z_i$ of the quantities $X$, the $X$ undergo the permutations of a group, which we likewise denote by $Q$. The changes of the $\Psi$ that correspond to these permutations follow from (1), if we apply to the indices of the $X$ the substitution $(z_i, z_i+h)$:
\weq{2}{\Sigma (-1)^{\Sigma z\xi} X_{z_1+h_1, z_2+h_2, z_3+h_3} = \overset{z}{\Sigma} (-1)^{\Sigma(z+h)\xi} X_{z_1,z_2,z_3} = (-1)^{\Sigma h\xi} \Psi_{\xi_1,\xi_2,\xi_3}.}

\gesperrt{1. The functions $\Psi$ therefore change only their sign under the permutations of $Q$, and the squares of the $\Psi$ remain unchanged by $Q$.}

We have further to determine the influence of a substitution $\sigma$ from $S$ upon the functions $\Psi$.

For this purpose let us denote by $\varrho$ the \gesperrt{transposed substitution} to $\sigma$, and set
\weq{3}{z' = \sigma(z), \quad \xi = \varrho(\xi'), \quad \xi' = \varrho^{-1}(\xi),}
then (\S.\,37, 9.):
\weq{4}{\Sigma \xi z = \Sigma \xi' z'.}

If we make the substitutions $\sigma$ of the group $S$ in the indices of $X$, we obtain a permutation group of the $X$, to be denoted by $S$, and from (1) it follows that
\[
\overset{z}{\Sigma} (-1)^{\Sigma z\xi} X_{z_1', z_2', z_3'} = \overset{z}{\Sigma} (-1)^{\Sigma \xi' z'} X_{z_1', z_2', z_3'},
\]
%CONT>>


% ===== scan leaf 331, printed 310 =====
% !PRINTED 310
\weberpage{310}
\pdfbookmark[2]{§ 81. The triple systems of the resolvents}{bm94}
%% [page furniture, not body text -- commented out 2026-08-10] Page furniture (running
%% head, folio, signature) is excluded from this edition by standing policy.
%% \noindent Zehnter Abschnitt. \hfill \S.\,81.

%<<CONT
and since the system of the $z_i'$ runs through the same sequence of values as the system of the $z_i$, this sum is equal to $\Psi_{\xi_1',\xi_2',\xi_3'}$, and we obtain:

\gesperrt{2. The application of the permutation $\sigma$ to the $X$ in the functions $\Psi$ is equivalent to the application of $\varrho^{-1}$ to the indices $\xi_1, \xi_2, \xi_3$.}

\begin{center}
\S.\,81.\\
\gesperrt{The triple systems of the resolvents.}
\end{center}

We now denote the resolvent $\Psi_{\xi_1,\xi_2,\xi_3}$ more briefly by $\Psi^\xi$ and remark that, by formula (2), \S.\,80, besides the squares of these functions certain products also admit the permutations of the group $Q$. To these belongs first of all the product of all the quantities $\Psi$, which we shall denote by $A$:
\weq{1}{A = \overset{\xi}{\prod} \Psi_\xi,}
but then also the products of three, $\Psi_\xi \Psi_\eta \Psi_\zeta$, and consequently also their reciprocal values
\weq{2}{v = \frac{1}{\Psi_\xi \Psi_\eta \Psi_\zeta},}
when the indices satisfy the conditions:
\weq{3}{\begin{aligned}
\xi_1 + \eta_1 + \zeta_1 &\equiv 0\\
\xi_2 + \eta_2 + \zeta_2 &\equiv 0 \quad (\mathrm{mod}\ 2).\\
\xi_3 + \eta_3 + \zeta_3 &\equiv 0
\end{aligned}}
Such a system of three functions $\Psi$ we call a \gesperrt{triple}. Likewise we shall call three index systems $(\xi), (\eta), (\zeta)$ that satisfy the conditions (3) a triple.

There are in all seven and no more such triples; for among the three systems $(\xi), (\eta), (\zeta)$ no two can be equal to one another, and by two of them the third is uniquely determined. One can therefore take for $(\xi), (\eta)$ any pair from the seven possible systems $(\xi)$, but then obtains each triple six times. The total number is therefore $7 . 6 : 6 = 7$.

In order to characterise the individual triple systems, we introduce three numbers $\alpha_1, \alpha_2, \alpha_3$, to be taken modulo 2, which do not all three vanish, and remark that the congruence
\weq{4}{\alpha_1 \xi_1 + \alpha_2 \xi_2 + \alpha_3 \xi_3 \equiv 0 \ (\mathrm{mod}\ 2)}
%CONT>>


% ===== scan leaf 332, printed 311 =====
% !PRINTED 311
\weberpage{311}
%% [page furniture, not body text -- commented out 2026-08-10] Page furniture (running
%% head, folio, signature) is excluded from this edition by standing policy.
%% \noindent \S.\,81. \hfill Triple systems.

%<<CONT
are satisfied for three and only three systems of the unknowns $\xi_1, \xi_2, \xi_3$, and that these three form a triple, so that the triple is determined by the three congruences
\weq{5}{\Sigma \alpha\xi \equiv 0, \quad \Sigma \alpha\eta \equiv 0, \quad \Sigma \alpha\zeta \equiv 0 \quad (\mathrm{mod}\ 2)}
One recognises this without lengthy general considerations if one considers individually the cases in which only one $\alpha$, or two, or all three $\alpha \equiv 1$ (mod 2).

Since there are seven distinct systems of numbers $\alpha_1, \alpha_2, \alpha_3$, a triple is completely determined by the relations (5) from each such system of the $\alpha$, and we can denote the function $v$ unambiguously by $v_{\alpha_1,\alpha_2,\alpha_3}$ or more briefly by $v_\alpha$:
\weq{6}{v_\alpha = \frac{1}{\Psi_\xi \Psi_\eta \Psi_\zeta}.}

\gesperrt{3. The functions $v_\alpha$, regarded as functions of the $X$, admit the permutations of the group $Q$.}

If we apply to the $z$ a substitution $\sigma$ from $S$, then by 2.\ (\S.\,80) the $\xi, \eta, \zeta$ simultaneously undergo the substitution $\varrho^{-1}$; from this, however, it follows by (5) that the $\alpha$ undergo the substitution $\sigma$, and we therefore have:

\gesperrt{4. The application of the substitution $\sigma$ from $S$ to the indices $z$ of $X$ has as its consequence the application of the same substitution to the indices $\alpha$ of $v_\alpha$.}

If, in the sums (5), in which $(\xi), (\eta), (\zeta)$ is the triple belonging to $(\alpha)$, we put for $(\alpha)$ another system $(\beta)$, then on account of (3)
\[
\Sigma \beta\xi + \Sigma \beta\eta + \Sigma \beta\zeta \equiv 0.
\]

Now if $\beta$ is different from $\alpha$, these three sums cannot all three be congruent to 0, because otherwise two distinct systems $(\alpha), (\beta)$ would lead to the same triple; \gesperrt{consequently, of these sums, two must be odd and one even.} We therefore have the theorem:

\gesperrt{5. If $(\xi)$ runs through the triple belonging to $(\alpha)$, and $(\beta)$ is a system different from $(\alpha)$, then among the sums $\Sigma \beta\xi$ two are odd and one is even.}


% ===== scan leaf 333, printed 312 =====
% !PRINTED 312
\weberpage{312}
%% [page furniture, not body text -- commented out 2026-08-10] Page furniture (running
%% head, folio, signature) is excluded from this edition by standing policy.
%% \noindent Zehnter Abschnitt. \hfill \S.\,81.

If in the sum $\Sigma \alpha\xi$ we make all seven possible assumptions for the system $(\xi)$, then one obtains the value 0 three times and consequently the value 1 four times.

The congruence
\weq{7}{\Sigma \alpha\xi \equiv 1 \ (\mathrm{mod}\ 2)}
therefore has four and only four different solutions for $\xi$. The totality of these four solutions we shall call a \gesperrt{quadruple}. Every triple determines a quadruple uniquely and conversely, and so there are also seven different quadruples, which by (7) are completely determined by the number system $(\alpha)$. We now prove the theorem:

\gesperrt{6. If $(\xi)$ runs through the quadruple belonging to $(\alpha)$, and $(\beta)$ is a system different from $(\alpha)$, then among the four sums $\Sigma \beta\xi$ two are even and two are odd.}

For if $(\xi)$ runs through the totality of all seven systems, then among the sums $\Sigma \beta\xi$ there are found three even and four odd; of these, the triple belonging to $(\alpha)$ furnishes two odd and one even; there thus remain for the quadruple two even and two odd.

We now change the notation for the quadruples and take for the system $(\alpha)$ another symbol $(h_1, h_2, h_3) = (h)$. The quadruple belonging to $(h)$ we denote by $(\alpha), (\beta), (\gamma), (\delta)$, so that the four congruences:
\weq{8}{\Sigma h\alpha \equiv 1, \quad \Sigma h\beta \equiv 1, \quad \Sigma h\gamma \equiv 1, \quad \Sigma h\delta \equiv 1 \quad (\mathrm{mod}\ 2)}
hold.

In the denominator of the product of the four functions
\[
v_\alpha, \quad v_\beta, \quad v_\gamma, \quad v_\delta
\]
there occur by (6) in all 12 factors $\Psi_\xi$, and indeed each factor $\Psi_\xi$ occurs among them as often as the number of even numbers among the sums
\[
\Sigma \xi\alpha, \quad \Sigma \xi\beta, \quad \Sigma \xi\gamma, \quad \Sigma \xi\delta
\]
amounts to, i.\,e.\ $\Psi_h$ does not occur at all, while every other $\Psi_\xi$ occurs twice (by 6.).

If we now also make use of the relation (1), it follows that:
\weq{9}{\Psi_h^2 = A^2 v_\alpha v_\beta v_\gamma v_\delta.}


% ===== scan leaf 334, printed 313 =====
% !PRINTED 313
\weberpage{313}
\pdfbookmark[2]{§ 82. Application to equations of the 8th degree}{bm95}

For the application we assemble the seven triples, from which one can read off the quadruples as the complement to the full system. In the first column stands the symbol ($\alpha$) to which the triple belongs:

\[
\begin{array}{c||ccc|}
1\,0\,0 & 0\,1\,0 & 0\,0\,1 & 0\,1\,1 \\
0\,1\,0 & 1\,0\,0 & 0\,0\,1 & 1\,0\,1 \\
0\,0\,1 & 1\,0\,0 & 0\,1\,0 & 1\,1\,0 \\
0\,1\,1 & 1\,0\,0 & 0\,1\,1 & 1\,1\,1 \\
1\,0\,1 & 0\,1\,0 & 1\,0\,1 & 1\,1\,1 \\
1\,1\,0 & 0\,0\,1 & 1\,1\,0 & 1\,1\,1 \\
1\,1\,1 & 0\,1\,1 & 1\,0\,1 & 1\,1\,0 \\
\end{array}
\]

If we denote the symbols ($\alpha$), in the order in which they stand in the first column of this table, by $1, 2, 3, 4, 5, 6, 7$, then we can represent the table for the triples and quadruples more simply thus:

\[
\begin{array}{c|ccc|cccc}
1 & 2 & 3 & 4 & 1 & 5 & 6 & 7 \\
2 & 1 & 3 & 5 & 2 & 4 & 6 & 7 \\
3 & 1 & 2 & 6 & 3 & 4 & 5 & 7 \\
4 & 1 & 4 & 7 & 2 & 3 & 5 & 6 \\
5 & 2 & 5 & 7 & 1 & 3 & 4 & 6 \\
6 & 3 & 6 & 7 & 1 & 2 & 4 & 5 \\
7 & 4 & 5 & 6 & 1 & 2 & 3 & 7 \\
\end{array}
\]

\begin{center}
\S.\,82.

\gesperrt{Application to equations of degree $8$.}
\end{center}

In the formulae derived above we now put for the variables $X_{z_1,z_2,z_3}$ the roots of an equation of degree $8$, of which we assume that its \gesperrt{Galois} group $P = TQ$, in the fixed domain of rationality $\Omega$, is contained in the linear congruence group $R = SQ$, so that $T$ is a divisor of $S$. All functions that admit the permutations of this group are rational.

To these belongs first of all the sum of the $X$:

\weq{1}{\sum^{z} X_{z_1, z_2, z_3} = B,}
%CONT>>


% ===== scan leaf 335, printed 314 =====
% !PRINTED 314
\weberpage{314}
\pdfbookmark[2]{§ 83. Metacyclic equations of the 8th degree}{bm96}
%<<CONT
further the quantity $A$ defined by (1) of the previous section; but then also the symmetric functions of the seven quantities $\Psi^2$, and further, if we now add for the sake of simplicity the assumption, to which we shall return, that none of the quantities $\Psi$ vanishes, the symmetric functions of the quantities $v$; and not only the symmetric functions of the $v$, but all functions of the $v$ that admit the permutations of the group $T$.

The seven quantities

\weq{2}{\begin{aligned}
v_{1,0,0} &= v_1, & v_{0,1,0} &= v_2, & v_{0,0,1} &= v_3, \\
v_{0,1,1} &= v_4, & v_{1,0,1} &= v_5, & v_{1,1,0} &= v_6, & v_{1,1,1} &= v_7
\end{aligned}}

are then the roots of a rational equation of degree $7$ with the group $T$, and by addition of equation (1) and the equations [\S.\,80, (1)]:

\[
\Sigma (-1)^{\Sigma z\xi} X_{z_1, z_2, z_3} = \Psi_{\xi_1, \xi_2, \xi_3}
\]

there results, having regard to (9) of the previous section:

\weq{3}{8\,X_{0,0,0} = A\Bigl\{\begin{aligned}[t]
&\sqrt{v_1}\sqrt{v_5}\sqrt{v_6}\sqrt{v_7} + \sqrt{v_2}\sqrt{v_4}\sqrt{v_6}\sqrt{v_7} \\
&+ \sqrt{v_3}\sqrt{v_4}\sqrt{v_5}\sqrt{v_7} + \sqrt{v_2}\sqrt{v_3}\sqrt{v_5}\sqrt{v_6} \\
&+ \sqrt{v_1}\sqrt{v_3}\sqrt{v_4}\sqrt{v_6} + \sqrt{v_1}\sqrt{v_2}\sqrt{v_4}\sqrt{v_5} \\
&+ \sqrt{v_1}\sqrt{v_2}\sqrt{v_3}\sqrt{v_7}
\end{aligned}\Bigr\} + B.}

There are 15 changes of sign of the seven square roots $\sqrt{v}$ under which this expression remains unchanged, namely when all seven signs are changed simultaneously, or when the signs of a triple or of a quadruple are changed. \gesperrt{Consequently the expression (3) takes only eight different values, however the seven square roots occurring in it may be determined.}

\begin{center}
\S.\,83.

\gesperrt{Metacyclic Equations of Degree $8$.}
\end{center}

Now in order, after these preparations, to find the primitive metacyclic equations of degree $8$, we have first to determine the metacyclic divisors of the group $S$. The group $S$ itself is, as already mentioned, not metacyclic. Its
%CONT>>


% ===== scan leaf 336, printed 315 =====
% !PRINTED 315
\weberpage{315}
%<<CONT
divisors we have already considered (\S.\,70 to 72), and all proper divisors of $S$ are metacyclic. Among them are divisors of degree $21$ and $7$, which we shall denote by $T_{21}$, $T_7$, further divisors of degree $24$, and still other divisors whose degree numbers divide 24, all of which we shall here comprise under the symbol $T_{24}$.

These groups $T_{24}$ cannot connect the seven quantities $(\xi) = (\xi_1, \xi_2, \xi_3)$ transitively with one another, because the degree of a transitive permutation group of seven elements must be divisible by 7 (Bd.\,I, \S.\,154, 7.). We shall show that the permutation groups of the $X$ corresponding to the groups $T_{24}\,Q$ are imprimitive. % [sic] Full stop before the syntactic continuation with ``and that.''
and that these groups therefore drop out of our consideration\footnote{From the literature on the equations of degree $8$ the following are to be mentioned: \gesperrt{Nöther}, Mathem. Annalen, Bd. XV. \gesperrt{Wilshaus}, Ueber die algebraische Auflösbarkeit der Gleichungen $8^{\mathrm{ten}}$ Grades. Dissertation. Marburg 1888.}.

Since all permutations $\sigma$ which leave two of the seven elements $(\xi)$ unchanged, or merely permute them among themselves, leave a third element unchanged, namely the complement to the triple, two cases are possible:

\begin{itemize}
\item[1)] the group $T_{24}$ leaves one element at rest, or
\item[2)] the seven elements are decomposed into two parts of three and four elements, of which each part is permuted only among itself.
\end{itemize}

In the latter case we may regard the three elements as belonging to a triple, and accordingly the four elements as belonging to a quadruple. For if three elements which do not belong to a triple are permuted among themselves, then the three complements of each pair of elements to the triple are also permuted among themselves, and the one remaining element stays unchanged. We thus come back to case 1).

We now add to the seven systems $(\xi)$ the eighth $(0,0,0) = (0)$ and denote the eight quantities $X_{\xi_1, \xi_2, \xi_3}$ by

\weq{1}{0, 1, 2, 3, 4, 5, 6, 7.}

In case 1) two of these eight elements, say 0, 1, then remain unchanged under the substitutions of $T_{24}$. There is, further, in $Q$ one and only one substitution $\gamma_0$ by which 0 is interchanged with 1.


% ===== scan leaf 337, printed 316 =====
% !PRINTED 316
\weberpage{316}
The two substitutions $1, \gamma_0$ form a group $Q_0$. Any substitution $\gamma_1$ of $Q$ that does not occur in $Q_0$ carries 0, 1 into two elements different from them, say 2, 3, and by the two substitutions $Q_0\gamma_1$ the pair 0, 1 passes into 2, 3 or into 3, 2. If, further, by $\gamma_2$ the pair 0, 1 passes into a third pair 4, 5, then 4, 5 is different both from 0, 1 and from 2, 3, and so we can decompose $Q$ into the subsidiary groups:

\[
Q = Q_0 + Q_0\gamma_1 + Q_0\gamma_2 + Q_0\gamma_3,
\]

whereby the eight elements (1) are decomposed into four pairs

\weq{2}{0, 1;\quad 2, 3;\quad 4, 5;\quad 6, 7}

so that by the two substitutions of one of the subsidiary groups $Q_0\gamma_i$ the pair 0, 1 passes into one of these four pairs.

Now if $k$ is any substitution from the group

\[
P = T_{24}\,Q,
\]

then the elements $\tau_0, \tau_1, \tau_2, \tau_3$ from $T_{24}$ and $\gamma_0', \gamma_1', \gamma_2', \gamma_3'$ from $Q$ can be so determined that

\[
k = \tau_0\gamma_0', \quad \gamma_1 k = \tau_1\gamma_1', \quad \gamma_2 k = \tau_2\gamma_2', \quad \gamma_3 k = \tau_3\gamma_3'
\]

or

\[
k = \tau_0\gamma_0' = \gamma_1^{-1}\tau_1\gamma_1' = \gamma_2^{-1}\tau_2\gamma_2' = \gamma_3^{-1}\tau_3\gamma_3',
\]

and this representation shows that by $k$ any one of the pairs (2) is carried into a pair of the same system; for, e.\,g., by $\gamma_1^{-1}$ the pair 2, 3 is carried into 0, 1, by $\tau_1$ the pair 0, 1 remains unchanged, and by $\gamma_1'$, which belongs to $Q$, the pair 0, 1 is carried into one of the pairs (2).

The group $T_{24}\,Q$ is therefore imprimitive, and indeed in such a way that we have four systems of imprimitivity. The roots 0, 1 satisfy a quadratic equation whose coefficients depend on the roots of a biquadratic equation.

Let us now pass to the case 2), in which the elements of a triple 1, 2, 3 and of a quadruple 4, 5, 6, 7 are each transitively connected among themselves by $T_{24}$. The eight quantities (1) here split into two quadruples:

\weq{3}{0, 1, 2, 3;\quad 4, 5, 6, 7,}

and by the substitutions of $Q$ the quantities of one of these quadruples are either only interchanged among themselves, or they are carried into the quantities of the other quadruple.


% ===== scan leaf 338, printed 317 =====
% !PRINTED 317
\weberpage{317}

To see this, we denote, in the earlier notation, 1, 2, 3 by $(\xi)$, $(\eta)$, $(\zeta)$; since these three quantities form a triple, there is [by \S.\,81, (5)] a system $(\alpha)$ such that
\weq{4}{\Sigma\alpha\xi \equiv 0, \quad \Sigma\alpha\eta \equiv 0, \quad \Sigma\alpha\zeta \equiv 0 \ (\mathrm{mod}\ 2)}
holds. If we now apply a substitution from $Q$:
\[
\begin{pmatrix} \xi_1, & \xi_2, & \xi_3 \\ \xi_1', & \xi_2', & \xi_3' \end{pmatrix},
\]
in which we may take $\xi_i' = \xi_i + h_i$, then it follows from (4):
\[
\Sigma\alpha\xi' \equiv \Sigma\alpha\eta' \equiv \Sigma\alpha\zeta' \equiv \Sigma\alpha h,
\]
that is, the $(\xi')$, $(\eta')$, $(\zeta')$ form the triple belonging to $(\alpha)$ if $(h)$ itself belongs to this triple, or, in the other case, $(h)$, $(\xi')$, $(\eta')$, $(\zeta')$, is the quadruple belonging to $(\alpha)$.

If now again $k$ denotes any substitution from $P$ and $\gamma$ a substitution from $Q$ by which the quadruple 0, 1, 2, 3 passes into 4, 5, 6, 7, then we can determine the elements $\tau'$, $\tau''$ from $T_{24}$ and $\gamma'$, $\gamma''$ from $Q$ so that
\[
k = \tau'\gamma' = \gamma^{-1}\tau''\gamma'',
\]
whereby, as above, it is proved that the two systems (3) are imprimitively connected by $k$.

After adjunction of a square root to the domain of rationality, the quantities 0, 1, 2, 3 are then the roots of a biquadratic equation.

By these considerations, for primitive equations of degree 8 the possibility is also excluded that one of the seven quantities $\Psi_{\xi_1,\xi_2,\xi_3}$ [\S.\,80, (1)] vanishes. For these functions are rational functions of the roots $X$; hence if one of them vanishes, every permutation of the group of the equation can be applied to the rational equation $\Psi = 0$. From this it follows, by the theorem \S.\,80, 2., that either all seven functions $\Psi$ must vanish, in which case the roots would be rational, or that the substitutions $\varrho^{-1}$, applied to the $(\xi)$, must form an intransitive group. The degree of the group could then not be divisible by 7. Now the group of the $\varrho^{-1}$ is isomorphic with the group $T$ of the $\sigma$ (\S.\,37, 8.), and therefore the degree of $T$ is also not divisible by 7, and consequently
%CONT>>


% ===== scan leaf 339, printed 318 =====
% !PRINTED 318
\weberpage{318}
%<<CONT
$T$ is also intransitive. From this it follows, as shown above, that $TQ$
is imprimitive.

As the group for primitive metacyclic equations of degree 8
there remain, therefore, only the two types
\[
T_7 Q, \quad T_{21} Q
\]
In the first case the quantities $v$ occurring in the formula [\S.\,82, (3)]
are the roots of a cyclic equation of degree 7
with the group $(z, z+b)$, in the second of a metacyclic one with the
group $(z, az+b)$ (Bd. I, \S.\,180), in which $z, a, b$ are taken according to the modulus 7
and $a$ runs through the values 1, 2, 4.

In order to obtain an assignment of the $v_i$, as they occur in the formula (3), \S.\,82,
to the $v_z$, in which $z$ is the index occurring in the groups $(z, z+b)$ or $(z, az+b)$,
we can start from an arbitrary
substitution of the seventh order, say from
\[
\tau = \begin{pmatrix} 1, & 0, & 1 \\ 1, & 0, & 0 \\ 0, & 1, & 0 \end{pmatrix}
\quad [\S.\,79, (2)].
\]

If we then take $v_{1,0,0}$ for $v_0$, one finds by repeated
application of $\tau$ to $1, 0, 0$ the following arrangement:
\[
(1,0,0)=0, \quad (1,1,0)=1, \quad (1,1,1)=2, \quad (0,1,1)=3, \quad (1,0,1)=4,
\]
\[
(0,1,0)=5, \quad (0,0,1)=6,
\]
and by \S.\,82, (2) one has therefore to replace
\[
v_1, v_2, v_3, v_4, v_5, v_6, v_7
\]
by
\[
v_0, v_5, v_6, v_3, v_4, v_1, v_2
\]
By this there results from \S.\,82, (3):
\[
8X = A \sum_{0,\,6}^{z} \sqrt{v_z}\,\sqrt{v_{z+1}}\,\sqrt{v_{z+2}}\,\sqrt{v_{z+4}} + B,
\]
in which $A, B$ are rational quantities, and $v_z$ are the roots of a cyclic or
of a metacyclic equation of the group $(z, az+b)$ of degree 7.

Conversely, it is also not difficult to prove that the eight quantities
derived from this by change of the signs are always the roots
of a primitive metacyclic equation
%CONT>>


% ===== scan leaf 340, printed 319 =====
% !PRINTED 319
\weberpage{319}
\pdfbookmark[2]{§ 84. Biquadratic equations}{bm97}
%<<CONT
are of the 8th degree, which shall not be entered into more closely
here (cf. Bd. I, \S.\,185).

\begin{center}
\S.\ 84.

\gesperrt{Biquadratic Equations.}
\end{center}

We have restricted ourselves to the consideration of the primitive
metacyclic equations of the 8th degree, because the imprimitive ones
come back to equations of the 4th degree. In order, then, to be able to
represent the roots of the latter also in the manner of the preceding
section, we have still to seek an analogous representation of the roots
of biquadratic equations. Here the same considerations as in
\S.\,82 are applicable, only in an essentially simpler form.

We denote the roots of the biquadratic equation by
\weq{1}{X_{0,0}, X_{1,0}, X_{0,1}, X_{1,1},}

and can then represent the whole permutation group of the 24th
degree as a binary linear congruence group for the modulus 2:
\weq{2}{P = SQ.}

The homogeneous group $S$ contained therein is of the 6th
degree and consists of the substitutions
\weq{3}{\sigma = \begin{pmatrix}1, & 0\\0, & 1\end{pmatrix},
\begin{pmatrix}1, & 0\\1, & 1\end{pmatrix},
\begin{pmatrix}0, & 1\\1, & 0\end{pmatrix},
\begin{pmatrix}0, & 1\\1, & 1\end{pmatrix},
\begin{pmatrix}1, & 1\\1, & 0\end{pmatrix},
\begin{pmatrix}1, & 1\\0, & 1\end{pmatrix},}

and one may regard
\weq{4}{\sigma_1 = \begin{pmatrix}0, & 1\\1, & 0\end{pmatrix}, \quad
\sigma_2 = \begin{pmatrix}1, & 1\\0, & 1\end{pmatrix}}

as the generating substitutions of the group $S$.

We now define, as in \S.\,80, the three quantities:
\weq{5}{\begin{aligned}
\Psi_{1,0} &= X_{0,0} - X_{1,0} + X_{0,1} - X_{1,1}\\
\Psi_{0,1} &= X_{0,0} + X_{1,0} - X_{0,1} - X_{1,1}\\
\Psi_{1,1} &= X_{0,0} - X_{1,0} - X_{0.1} + X_{1,1}, % [sic] printed with a period after 0 instead of a comma
\end{aligned}}

so that $\Psi^2_{1,0}, \Psi^2_{0,1}, \Psi^2_{1,1}$, but also the product
$\Psi_{1,0}\Psi_{0,1}\Psi_{1,1}$, remain unchanged by the substitutions
of the cyclic group $Q$.


% ===== scan leaf 341, printed 320 =====
% !PRINTED 320
\weberpage{320}
If one applies the substitution $\sigma$ to the indices $\xi, \eta$ of the $X$,
the indices of the $\Psi$ undergo [as in \S.\,80, (2)] the substitution
$\varrho^{-1}$, where $\varrho$ is the transposed substitution of $\sigma$.

But, corresponding to the substitutions (4), we have
\weq{6}{\varrho_1^{-1} = \begin{pmatrix}0, & 1\\1, & 0\end{pmatrix}, \quad
\varrho_2^{-1} = \begin{pmatrix}1, & 0\\1, & 1\end{pmatrix}.}

Assuming now, for the sake of simplicity, that the three quantities (5) are
different from zero, we set
\weq{7}{\begin{aligned}
v_{1,0} &= \frac{\Psi_{0,1}\Psi_{1,1}}{\Psi_{1,0}}
= \frac{\Psi_{0,1}\Psi_{1,0}\Psi_{1,1}}{\Psi_{1,0}^2};\\
v_{0,1} &= \frac{\Psi_{1,0}\Psi_{1,1}}{\Psi_{0,1}}
= \frac{\Psi_{0,1}\Psi_{1,0}\Psi_{1,1}}{\Psi_{0,1}^2};\\
v_{1,1} &= \frac{\Psi_{1,0}\Psi_{0,1}}{\Psi_{1,1}}
= \frac{\Psi_{1,0}\Psi_{0,1}\Psi_{1,1}}{\Psi_{1,1}^2}.
\end{aligned}}

The quantities $v_{\xi,\eta}$ remain unchanged under the substitutions of $Q$,
and undergo, through $\sigma_1, \sigma_2$, the permutations determined by (6)
\[
\begin{array}{l l l l}
 & v_{1,0}, & v_{0,1}, & v_{1,1}\\
\sigma_1 & v_{0,1}, & v_{1,0}, & v_{1,1}\\
\sigma_2 & v_{1,1}, & v_{0,1}, & v_{1,0}.
\end{array}
\]

From this it follows that the symmetric functions of the three quantities
$v_{1,0}, v_{0,1}, v_{1,1}$ remain unchanged under the permutations $P$, and
are consequently contained in the domain of rationality of the symmetric
functions of the $X$. The three quantities $v$ are therefore the roots of a
cubic equation, and there would be no difficulty in forming this cubic
resolvent of the general biquadratic equation.

But from (7) it follows further that
\[
v_{1,0}\,v_{0,1} = \psi_{1,1}^2, \quad v_{1,0}\,v_{1,1} = \psi_{0,1}^2,
\quad v_{0,1}\,v_{1,1} = \psi_{1,0}^2, % [sic] printed with lowercase psi, elsewhere capital Psi is used for these quantities
\]
and accordingly, if we further set the rational quantity
\[
X_{0,0} + X_{1,0} + X_{0,1} + X_{1,1} = a
\]
it follows from (5) that
\weq{8}{4X_{0,0} = a + \sqrt{v_{1,0}}\,\sqrt{v_{0,1}}
+ \sqrt{v_{1,0}}\,\sqrt{v_{1,1}} + \sqrt{v_{0,1}}\,\sqrt{v_{1,1}}.}

The expression (8) gives only four different values when we determine the
signs of the three square roots occurring in it in all possible ways, and one
can also con-
%CONT>>


% ===== scan leaf 342, printed 321 =====
% !PRINTED 321
\weberpage{321}
%<<CONT
versely show without difficulty that the four quantities
$X$ contained in the form (8), when the $v_{1,0}, v_{0,1}, v_{1,1}$ are the roots of any
cubic equation, satisfy a biquadratic equation with
rational coefficients\textsuperscript{1)}.

\noindent\rule{0.35\linewidth}{0.4pt}

\footnotesize
\noindent\textsuperscript{1)} This form of the solution of the biquadratic equation is the
analogue of the \gesperrt{Cayley} form of the \gesperrt{Cardan formula} (Bd. I,
\S.\,36). I do not know whether this form of the solution of biquadratic
equations is otherwise already known in elementary algebra. I
first found it in an English collection of problems ``Mathematical
Tripos, Part I, June 1894'', where the cubic resolvent
of the $v$ is also formed.
\normalsize


% ===== scan leaf 343, printed 322 =====
\weberpage{322}% page number not printed on this page
\pdfbookmark[1]{Eleventh Section. The Points of Inflection of a Curve of the Third Order.}{bm98}
\pdfbookmark[2]{§ 85. Ternary forms and algebraic curves}{bm99}
% !PRINTED none
% Part-title (Abschnitt-opening) page: no running head / page number printed on this leaf.

\begin{center}
\gesperrt{Eleventh Chapter.}

\textbf{\gesperrt{The Points of Inflection of a Curve of the Third Order.}}

\noindent\rule{0.25\linewidth}{0.4pt}

\S.\,85.

\gesperrt{Ternary Forms and Algebraic Curves.}
\end{center}

In order to be able, in the further course of our exposition,
to present to the reader some of the manifold and interesting
applications of algebra to geometric problems, without referring
him to other means, the necessary geometric foundations for
these applications shall here be derived in brief. In doing so we
restrict ourselves to the geometry of the plane, and even there
we set aside everything that is not directly required for our
applications\textsuperscript{1)}.

We employ the homogeneous triangular coordinates, which are
defined by this, that the position of a point $x$ is determined by
its distances from the three sides of the coordinate triangle,
measured in three arbitrary units of measure (or multiplied by
three arbitrary constant factors). These three quantities are called
the \gesperrt{coordinates} of the point $x$, and are usually denoted
by $x_1, x_2, x_3$ or by $y_1, y_2, y_3$ etc.

Between these three coordinates of the point $x$ there subsists
a linear, non-homogeneous relation, which one obtains by setting
the area of the coordinate triangle $(1, 2, 3)$ equal to the
sum of the areas of the three triangles $(x, 2, 3)$, $(x, 3, 1)$, $(x, 1, 2)$.
A coordinate, and hence also the expression for the
area of the corresponding one of the three last-named tri-
%CONT>>

\noindent\rule{0.25\linewidth}{0.4pt}

\footnotesize
\noindent\textsuperscript{1)} Fuller details may be found in the textbooks of analytic
geometry, among which the work of \gesperrt{George Salmon}, ``Higher plane
curves'', German by \gesperrt{Fiedler}, deserves special mention.
\normalsize


% ===== scan leaf 344, printed 323 =====
% !PRINTED 323
\weberpage{323}
%<<CONT
angles change their sign as often as the point $x$ passes through a side
of the coordinate triangle.

With the help of this non-homogeneous relation, unity can be
represented linearly and homogeneously by $x_1, x_2, x_3$, and thereby
every non-homogeneous function can be transformed into a homogeneous
one.

According to these determinations, the coordinates $x_1, x_2, x_3$
are at first not independent variables, but are bound to the relation
mentioned. But we now extend their meaning so that we can regard
them as independent variables, if we stipulate that we shall use them
only in homogeneous forms or equations;
and that not $x_1, x_2, x_3$ themselves,
but only the ratios $x_1 : x_2 : x_3$ shall determine the position of the point $x$.
Then we can disregard entirely the relation mentioned
between the $x_1, x_2, x_3$, and regard the $x_1, x_2, x_3$ as independent variables, for which only the single combination of values $x_1 = 0, x_2 = 0, x_3 = 0$ is excluded. The
values $h x_1, h x_2, h x_3$ then determine, for every $h$ different from zero, one and the same point $x$. Strictly speaking, the
coordinates $x_1, x_2, x_3$ of the point $x$ are then not the distances
from the coordinate axes described above themselves, but any quantities whatever that are proportional to these distances.

Every linear substitution

\weq{1}{\begin{aligned}
x_1 &= \alpha_1^{(1)} y_1 + \alpha_1^{(2)} y_2 + \alpha_1^{(3)} y_3 \\
x_2 &= \alpha_2^{(1)} y_1 + \alpha_2^{(2)} y_2 + \alpha_2^{(3)} y_3 \\
x_3 &= \alpha_3^{(1)} y_1 + \alpha_3^{(2)} y_2 + \alpha_3^{(3)} y_3,
\end{aligned}}

\noindent whose determinant

\weq{2}{r = \Sigma \pm \alpha_1^{(1)} \alpha_2^{(2)} \alpha_3^{(3)}}

\noindent is different from zero, can be interpreted geometrically in two
ways: first as a \gesperrt{mapping}, in that we regard $x_1, x_2, x_3$
and $y_1, y_2, y_3$ as coordinates of two points $x, y$ in the same
coordinate system, so that to every point $y$ there corresponds a point $x$
and conversely; and then also as a \gesperrt{coordinate transformation}, if we regard $x_1, x_2, x_3$ and $y_1, y_2, y_3$ as coordinates of one and the same point, referred to two different
coordinate systems. For the considerations that come next
we shall hold to the latter interpretation.


% ===== scan leaf 345, printed 324 =====
% !PRINTED 324
\weberpage{324}
\pdfbookmark[2]{§ 86. Singular points. Points of inflection. Double tangents}{bm100}
Any ternary form of degree $n$, $f(x_1, x_2, x_3)$, set equal to zero,
represents a curve of the $n$th order, which has $n$ points of intersection
with a straight line. We denote this curve briefly as the curve $f$.

By the substitution (1), $f$ passes into another form of the same degree in the variables $y_1, y_2, y_3$:

\weq{3}{f(x_1, x_2, x_3) = \varphi(y_1, y_2, y_3),}

\noindent and $\varphi = 0$ represents the same curve as $f = 0$, referred to
the coordinate system $y_1, y_2, y_3$.

If we form the derivatives of the identity (3) with respect to the
variables $y_1, y_2, y_3$, it follows, having regard to (1):

\weq{4}{\begin{aligned}
\varphi'(y_1) &= \alpha_1^{(1)} f'(x_1) + \alpha_2^{(1)} f'(x_2) + \alpha_3^{(1)} f'(x_3) \\
\varphi'(y_2) &= \alpha_1^{(2)} f'(x_1) + \alpha_2^{(2)} f'(x_2) + \alpha_3^{(2)} f'(x_3) \\
\varphi'(y_3) &= \alpha_1^{(3)} f'(x_1) + \alpha_2^{(3)} f'(x_2) + \alpha_3^{(3)} f'(x_3).
\end{aligned}}

\begin{center}
\S.\,86.

\gesperrt{Singular Points. Points of Inflection. Double Tangents.}
\end{center}

If for a point $x$ the three equations

\weq{1}{f'(x_1) = 0, \quad f'(x_2) = 0, \quad f'(x_3) = 0}

\noindent are satisfied simultaneously, then by \gesperrt{Euler's theorem}

\weq{2}{n f = x_1 f'(x_1) + x_2 f'(x_2) + x_3 f'(x_3)}

\noindent (Bd. I, \S.\,17) this point lies on the curve $f$. The equations \S.\,85,(4)
show that at the same point $\varphi'(y_1), \varphi'(y_2), \varphi'(y_3)$
must also vanish, hence that the property of a point characterised by the
equations (1) is preserved under linear transformation, and therefore, as we
also express it, belongs to the \gesperrt{invariant} properties.

Not on every curve of the $n$th order do points occur which satisfy
the conditions (1), as is shown, e.g., by the assumption $f(x_1, x_2, x_3)$
$= x_1^n + x_2^n + x_3^n$. Rather, if even a single such point is to exist,
an equation between the coefficients of $f$ must be satisfied, which one
obtains by elimination of $x_1, x_2, x_3$ from the three equations (1) (Bd. I, \S.\,50).
One can represent this relation in the form that a certain integral
rational and homogeneous function of all the coefficients of $f$ must be
equal to zero, and this function
%CONT>>


% ===== scan leaf 346, printed 325 =====
% !PRINTED 325
\weberpage{325}
%<<CONT
is called the discriminant of $f$. By this definition it is determined only up to a numerical factor, and up to now it can be computed only in the simplest cases.

The points of the curve $f$ whose coordinates satisfy the conditions (1) are called \gesperrt{singular points}.

One may also raise the question under what condition it is possible that infinitely many singular points occur on a curve. For this to happen, first of all the result of eliminating one of the unknowns, say $x_3$, from the two equations $f'(x_1) = 0$, $f'(x_2) = 0$ must vanish identically, i.\,e. the two functions $f'(x_1)$ and $f'(x_2)$, regarded as functions of $x_3$, must have a common divisor, and by Bd. I, \S.\,51 they must therefore also have a (non-constant) common divisor in the region of the forms $x_1, x_2, x_3$. This divisor must then in turn, as follows from the same considerations, have a common divisor with $f$, or, what is the same thing, with $f'(x_3)$, and it thus results that infinitely many singular points can be present only when the four forms $f, f'(x_1), f'(x_2), f'(x_3)$ have a common divisor.

Now let $u$ be such a common divisor, which we assume to be \gesperrt{irreducible}.

We set
\[
f(x_1, x_2, x_3) = uv
\]
\[
f'(x_1) = u\,\frac{\partial v}{\partial x_1} + v\,\frac{\partial u}{\partial x_1}
\]
\[
f'(x_2) = u\,\frac{\partial v}{\partial x_2} + v\,\frac{\partial u}{\partial x_2}
\]
\[
f'(x_3) = u\,\frac{\partial v}{\partial x_3} + v\,\frac{\partial u}{\partial x_3},
\]
and these equations show, since $u$ is irreducible and therefore
\[
u, \quad \frac{\partial u}{\partial x_1}, \quad \frac{\partial u}{\partial x_2}, \quad \frac{\partial u}{\partial x_3}
\]
can have no common divisor, that $v$ must be divisible by $u$; i.\,e. infinitely many singular points can occur only when $f(x_1, x_2, x_3)$ has a \gesperrt{quadratic divisor}. In this case we say that the curve $f$ contains a doubly counting component.


% ===== scan leaf 347, printed 326 =====
% !PRINTED 326
\weberpage{326}
In order to recognise the geometric meaning of the singular points, we imagine the function $f$ ordered by powers of the one variable $x_3$:

\weq{2}{f(x_1, x_2, x_3) = x_3^n f_0 + x_3^{n-1} f_1 + x_3^{n-2} f_2 + \cdots,}
in which $f_0, f_1, f_2, \dots$ are binary forms of the variables $x_1, x_2$, whose degree is indicated by the index, so that $f_0$ is a constant, $f_1$ a linear, $f_2$ a quadratic form, and so on.

Now if a singular point is present, then on account of the invariance of this property we may assume that this point falls at the corner $x_1 = 0$, $x_2 = 0$ of the coordinate triangle, which we shall denote by $\xi_3$. Then $f_0$ and $f_1$ must vanish, and if we then set $x_1 = 0$, the equation $f = 0$ acquires the factor $x_2^2$. This says that two of the points of intersection of the line $x_1 = 0$ coincide in the singular point. But the line $x_2 = 0$ can be any straight line passing through the singular point, and so it follows that the singular point has the property that \gesperrt{every} straight line passing through it cuts the curve there in two coincident points. % [sic: the source prints $x_2=0$ here; the preceding derivation uses the cutting line $x_1=0$]

For this reason the singular points are also called \gesperrt{double points}; by which, however, it is not excluded that more than two points of intersection may also coincide, whereby the \gesperrt{higher singularities} arise.

To the form of the equation (2) we attach a few further simple considerations that will be important in what follows.

If the point $\xi_3$ lies on the curve $f$ but is not a singular point, then the constant $f_0$ must vanish, and $f_1 = 0$ is the equation of the \gesperrt{tangent} of the curve at this point. If we take this tangent for the line $x_1 = 0$, then the equation of the curve, if we assume a constant factor equal to 1, reads:

\weq{3}{f = x_3^{n-1} x_1 + x_3^{n-2} f_2 + \cdots = 0.}
If the quadratic form $f_2$ is divisible by $x_1$, then the equation (3) acquires for $x_1 = 0$ the factor $x_2^3$, and the point $\xi_3$ counts as a threefold point of intersection of $x_1 = 0$ with the curve. Such a point is called a \gesperrt{point of inflection} or \gesperrt{inflexion point} of the curve $f$, and $x_1 = 0$ is called the \gesperrt{inflectional tangent}. The points of inflection do not belong to the singular points.


% ===== scan leaf 348, printed 327 =====
% !PRINTED 327
\weberpage{327}
\pdfbookmark[2]{§ 87. Fundamental covariants of a ternary form}{bm101}
If the point $\xi_3$ is a point of inflection and $x_1 = 0$ is the inflectional tangent, then $f$ has the form

\weq{4}{f = x_1 \varphi + x_2^3 \psi,}
in which $\varphi$ is a form of degree $(n-1)$, $\psi$ a form of degree $(n-3)$.

Conversely, if $f$ has this form, the point $\xi_3$ is a point of inflection, and $x_1 = 0$ is the inflectional tangent.

A straight line that touches the curve $f$ in two separate points is called a \gesperrt{double tangent}. If we take such a double tangent as the line $x_3 = 0$ and the points of contact as the vertices of the coordinate triangle lying on $x_3 = 0$, then, when $x_3 = 0$ is substituted into $f = 0$, an equation must arise that has both $x_1$ and $x_2$ as quadratic factors. The function $f$ must therefore have the form:

\weq{4}{f = x_3 \varphi + x_1^2 x_2^2 \psi,} % [sic] number (4) reprinted; page also numbers preceding display (4)
in which $\varphi$ is a form of degree $(n-1)$, $\psi$ a form of degree $(n-4)$.

Somewhat more generally we may also say: if $x_3 = 0$ is a double tangent, then $f$ can be represented thus:

\weq{5}{f = x_3 \varphi + \chi^2 \psi,}
in which $\chi = 0$ is the equation of a conic section. The points of contact of the double tangent are the points of intersection of $x_3 = 0$ and $\chi = 0$.

\begin{center}
\S.\,87.

\gesperrt{Fundamental Covariants of a Ternary Form.}
\end{center}

In \S.\,60 of the first volume we became acquainted with the concept of the invariants and covariants of a form. In geometry too these forms play a great role.

If by the linear substitution \S.\,85, (1) the ternary form $f(x_1, x_2, x_3)$, which we also denote in abbreviated fashion by $f(x)$, passes into $\varphi(y)$, and if we denote by $a$ the coefficients of $f$, by $b$ the coefficients of $\varphi$, and by $r$ the substitution determinant, then a form $\Phi(x, a)$ that is homogeneous both with respect to the $x$ and with respect to the $a$ is called a \gesperrt{covariant} of $f(x)$ if it satisfies the condition

\weq{1}{\Phi(y, b) = r^\lambda \Phi(x, a)}
%CONT>>


% ===== scan leaf 349, printed 328 =====
% !PRINTED 328
\weberpage{328}
%<<CONT
is satisfied. If in particular the form is independent of the variables $x$ and depends only on the coefficients $a$, it becomes an \gesperrt{invariant}. The exponent $\lambda$, which is called the \gesperrt{weight} of the covariant, is an integer.

In Vol.\ I, \S.\,59 we became acquainted with a covariant of weight 2 that is present in all forms whose degree is greater than 1, namely the \gesperrt{Hesse determinant}:

\weq{2}{H(x_1, x_2, x_3) = \begin{vmatrix} f''(x_1, x_1), & f''(x_1, x_2), & f''(x_1, x_3) \\ f''(x_2, x_1), & f''(x_2, x_2), & f''(x_2, x_3) \\ f''(x_3, x_1), & f''(x_3, x_2), & f''(x_3, x_3) \end{vmatrix}.}

Besides this covariant $H$ we shall consider two other general covariants, which exist for all ternary forms of degree higher than $2$, and which can also easily be generalised to forms of arbitrarily many variables. The first of these forms, if $H$ denotes the \gesperrt{Hesse} covariant as above, is:

\weq{3}{C_1(x, a) = \begin{vmatrix} f''(x_1, x_1), & f''(x_1, x_2), & f''(x_1, x_3), & H'(x_1) \\ f''(x_2, x_1), & f''(x_2, x_2), & f''(x_2, x_3), & H'(x_2) \\ f''(x_3, x_1), & f''(x_3, x_2), & f''(x_3, x_3), & H'(x_3) \\ H'(x_1), & H'(x_2), & H'(x_3) & 0 \end{vmatrix}.}

To demonstrate the invariant property for this form, we make in it the substitution \S.\,85, (1). Thereby $f(x, a)$ passes over into $\varphi(y, b)$ and $r^2 H(x, a)$ into $H(y, b)$. If we then understand by $H'(y_1), H'(y_2), H'(y_3)$ the derivatives of $H(y, b)$, we have the formulae

\weq{4}{
\begin{aligned}
\varphi''(y_\mu, y_\nu) &= \sum\limits^{i,k}_{i} f''(x_i, x_k)\, \alpha_i^{(\mu)} \alpha_k^{(\nu)} \\
H'(y_\nu) &= r^2 \sum H'(x_i)\, \alpha_i^{(\nu)},
\end{aligned}
}
in which $\mu, \nu, i, k$ run from 1 to 3.

If one now multiplies the determinant $C_1(x, a)$ twice in succession by the substitution determinant, written in the form

\[
r = \begin{vmatrix} \alpha_1^{(1)}, & \alpha_1^{(2)}, & \alpha_1^{(3)}, & 0 \\ \alpha_2^{(1)}, & \alpha_2^{(2)}, & \alpha_2^{(3)}, & 0 \\ \alpha_3^{(1)}, & \alpha_3^{(2)}, & \alpha_3^{(3)}, & 0 \\ 0, & 0, & 0, & 1 \end{vmatrix}
\]
by taking the one
%CONT>>


% ===== scan leaf 350, printed 329 =====
% !PRINTED 329
\weberpage{329}
\pdfbookmark[2]{§ 88. The Hessian curve}{bm102}
%<<CONT
time by rows, the other time by columns, it follows by the
formulas (4):

\weq{5}{C_1\,(y,b) = r^6\,C_1\,(x,a),}

whereby the invariant property is expressed and the weight
determined to be $= 6$.

The third covariant $C_2$ that we consider is the functional determinant
 of the three forms $f$, $H$, $C_1$, thus

\weq{6}{C_2\,(x, a) = \begin{vmatrix} f'\,(x_1), & H'\,(x_1), & C_1'\,(x_1) \\ f'\,(x_2), & H'\,(x_2), & C_1'\,(x_2) \\ f'\,(x_3), & H'\,(x_3), & C_1'\,(x_3) \end{vmatrix},}

which, as already follows from \S.\,59 of the first volume, has the
invariant property:

\weq{7}{C_2\,(y,b) = r^9\,C_2\,(x,a),}

and which therefore has the weight 9.

If we denote by $\lambda$ the weight, by $\nu$ the degree in the
variables, by $\mu$ the degree in the coefficients, the following
compilation results:

\begin{center}
\begin{tabular}{lccc}
 & $\nu$, & $\mu$, & $\lambda$ \\
\hline
$f$ & $n$, & 1, & 0 \\
$H$ & $3n-6$, & 3, & 2 \\
$C_1$ & $8n-18$, & 8, & 6 \\
$C_2$ & $12n-27$, & 12, & 9.
\end{tabular}
\end{center}

\begin{center}
\S.\,88.

\gesperrt{The Hesse Curve.}
\end{center}

The first of the covariants introduced, $H$, has a simple
geometric meaning, with which we now wish to become acquainted.
This investigation is essentially facilitated by the fact that,
on account of the invariant property of the form $H$, we lose
nothing in generality if we give the coordinate system any
special position with respect to the curve $f$.

If we place the vertex $\xi_3$ of the coordinate triangle at a
point of the curve $f$ that does not belong to the singular ones, and
choose the tangent at this point as the $x_1$-axis, then the form $f$
takes, by \S.\,86, (3), the expression:

\weq{1}{f = h\,x_3^{n-1}x_1 + x_3^{n-2}(a\,x_1^2+2\,b\,x_1x_2+c\,x_2^2)+x_3^{n-3}f_3+\cdots,}
%CONT>>


% ===== scan leaf 351, printed 330 =====
% !PRINTED 330
\weberpage{330}
%<<CONT
in which $h$, $a$, $b$, $c$ are constants, of which $h$ is different from zero,
and $f_3$ is a binary cubic form in $x_1, x_2$.

If, in accordance with this, we compute the function $H$ arranged
according to descending powers of $x_3$:

\weq{2}{H = x_3^{3n-6}\,H_0 + x_3^{3n-7}\,H_1 + \cdots,}

then we find

\weq{3}{\begin{aligned} H_0 &= -\,2\,(n-1)^2\,h^2\,c \\ H_1 &= -(n-1)^2\,h^2\,f_3''\,(x_2,x_2) + 4\,(n-1)(n-2)\,h\,(b^2-a\,c)\,x_1. \end{aligned}}

The curve represented by the equation $H=0$ is called
the \gesperrt{Hesse curve} of the curve $f$. The formulae (2) and (3)
show that the Hesse curve passes through the point $\xi_3$
when and only when $H_0=0$, that is $c=0$, i.\,e.\ by \S.\,86,
when this point is a point of inflection. From the expression
for $H_1$ one can conclude further that the inflectional tangent
$x_1=0$ is at the same time a tangent of the curve $H=0$
when and only when $f_3''\,(x_2,x_2)$, and consequently $f_3$ itself, is divisible
by $x_1$. In this case, as appears from (3), $x_1=0$ would be a
\gesperrt{tangent having four-point contact} at the point $\xi_3$. This
case can occur among curves of the third order only
when $f$ is divisible by $x_1$, hence never for an irreducible
curve of the third order.

Besides the points of inflection, the \gesperrt{Hesse} curve also passes
through the singular points of the curve $f$, which shall not
be investigated further here.

If, then, we restrict our considerations to curves $f$ without
a singular point, we can say that every point of
intersection of the curve $f$ with its \gesperrt{Hesse} curve gives a
point of inflection of the curve $f$, and that this curve $f$ has no
other points of inflection. If points with a tangent of four-point
contact occur, the curves $f$
and $H$ touch one another, and we can regard such points as
having arisen through the coincidence of two points of inflection. If
we further apply the theorem of Bezout (Bd.\,I, \S.\,49), we obtain
the result:

\begin{quote}
\gesperrt{A curve of the $n$th order without a singular point
has $3\,n\,(n-2)$ inflexion points;}
\end{quote}

and in particular:

\begin{quote}
\gesperrt{A curve of the third order without a singular point
has nine separately situated inflexion points.}
\end{quote}


% ===== scan leaf 352, printed 331 =====
% !PRINTED 331
\weberpage{331}
\pdfbookmark[2]{§ 89. Inflexion points of a curve of the third order}{bm103}
\begin{center}
\S.\,89.

\gesperrt{Inflexion points of a curve of the third order.}
\end{center}

Let $f=0$ now be the equation of a curve of the third order
without a singular point. We place two of the vertices of the
coordinate triangle $\xi_1, \xi_2$ at two of the nine inflexion points,
and take the corresponding inflexion tangents for the
sides $x_2, x_1$. Then, if $x_1=0$ or $x_2=0$ is set, $f$ must
reduce to the term $h\,x_3^3$, in which $h$ is a constant,
and consequently all terms that have a coefficient different
from zero in the cubic form, excepting
$h\,x_3^3$, must be divisible by $x_1\,x_2$. Accordingly $f$ takes the form

\weq{1}{f\,(x_1,x_2,x_3) = x_1\,x_2\,(a_1\,x_1+a_2\,x_2+a_3\,x_3) + h\,x_3^3,}

in which $a_1, a_2, a_3$ are constants. From this, however, it follows that
the straight line

\[
a_1\,x_1 + a_1\,x_2 + a_3\,x_3 = 0
\] % [sic] printed coefficient of x_2 is a_1; expected a_2

is likewise an inflexion tangent of the curve $f$, and that its point of
intersection with the line $x_3=0$ is a third inflexion point.
This point is different from the first two $\xi_1, \xi_2$, because
neither $a_1$ nor $a_2 = 0$ can hold if the curve has no singular
point. We shall express this important result in the form
of a theorem:

1. \gesperrt{If one joins any two inflexion points of a
curve of the third order without a singular point
by a straight line, then this straight line cuts
the curve in a third inflexion point.}

Since there are nine inflexion points, four straight lines
issue from each of these points, each of which contains two
further inflexion points. There are nine such pencils
of four lines, and because each of these lines occurs in three pencils,
the number of the lines is $9\,.\,4 : 3 = 12$. We
can accordingly supplement Theorem 1. as follows:

2. \gesperrt{The nine inflexion points of a curve of the third
order lie by threes on twelve straight
lines, and through each inflexion point pass
four of these lines.}


% ===== scan leaf 353, printed 332 =====
% !PRINTED 332
\weberpage{332}
From the expression (1) we can derive a \gesperrt{canonical representation}
 for the cubic form $f$:
We denote by $\varepsilon$ an imaginary third root of unity
and introduce two variables $z_1, z_2$, which are defined by the equations

\weq{2}{\begin{aligned} a_1\,x_1 &= \varepsilon\,z_1 + \varepsilon^2\,z_2 - \tfrac{1}{3}\,a_3\,x_3 \\ a_2\,x_2 &= \varepsilon^2\,z_1 + \varepsilon\,z_2 - \tfrac{1}{3}\,a_3\,x_3 \end{aligned}}

from which it further follows:

\weq{3}{-\,(a_1\,x_1 + a_2\,x_2 + a_3\,x_3) = z_1 + z_2 - \tfrac{1}{3}\,a_3\,x_3.}

The variables $z_1, z_2$ are determined by (2) as linear functions
of $x_1, x_2, x_3$, and the linear functions $z_1, z_2, x_3$
are linearly independent as functions of $x_1, x_2, x_3$. If we
introduce (2) and (3) into (1), we obtain

\[
-\,a_1\,a_2\,f = z_1^3 + z_2^3 - \left(a_1\,a_2\,h + \frac{a_3^3}{27}\right)x_3^3 + a_3\,z_1\,z_2\,x_3.
\]

One gives this expression a more symmetric shape
 by introducing for $z_1, z_2, x_3$ three new linear functions
$y_1, y_2, y_3$, which differ from these only by constant factors
:

\[
f = \varphi\,(y_1,y_2,y_3) = \alpha_1\,y_1^3 + \alpha_2\,y_2^3 + \alpha_3\,y_3^3 + 6\,m\,y_1\,y_2\,y_3,
\]

where $\alpha_1, \alpha_2, \alpha_3, m$ denote constants.

If no singular point is present, the
coefficients $\alpha_1, \alpha_2, \alpha_3$ cannot vanish, and one can,
without restriction of generality, take them equal to one another, say $=h$.
One could even set them $=1$, which however, in order to
maintain homogeneity, is for the present better not done.
We thus have the \gesperrt{canonical form} for the ternary
cubic form:

\weq{4}{\varphi\,(y_1,y_2,y_3) = h\,(y_1^3+y_2^3+y_3^3) + 6\,m\,y_1\,y_2\,y_3.}

According to our point of departure, the line $x_3$, or what
is the same thing, $y_3$, was the line joining three inflexion points.
But there are twelve such lines, and we can therefore choose the
line $y_3$ in twelve ways. If, however, we put $y_1$ or $y_2$ in
place of $y_3$, we obtain the same representation in the
canonical form, and there are thus only four essentially different
 kinds of this representation (apart from the arbitrary
 $h$), i.\,e. four ways of determining the lines $y_1, y_2, y_3$. We
therefore have the theorem:
%CONT>>


% ===== scan leaf 354, printed 333 =====
% !PRINTED 333
\weberpage{333}
\pdfbookmark[2]{§ 90. Transformation of the cubic form to the canonical form}{bm104}
%<<CONT

\noindent 3.\ \gesperrt{A ternary cubic form $f(x_1,x_2,x_3)$ whose discriminant is different from zero can be brought, in four different ways, into the canonical form}
\[
\varphi(y_1,y_2,y_3) = h(y_1^3+y_2^3+y_3^3) + 6\,m\,y_1y_2y_3
\]
\gesperrt{by transformation.}

Since on each of the lines $y_1=0$, $y_2=0$, $y_3=0$ there lie three inflexion points, the problem of the inflexion points is essentially identical with that of the transformation of the cubic form $f$ to the canonical form, with which we shall occupy ourselves first $^1$).

\begin{center}
\S.\,90.

\gesperrt{Transformation of the cubic form to the canonical form.}
\end{center}

In order to carry out the transformation of the ternary cubic form to the canonical form, that is, to reduce it to the solution of certain algebraic equations, we must first form the covariants for the canonical form.

Accordingly, let
\[
\varphi(y) = h(y_1^3+y_2^3+y_3^3) + 6\,m\,y_1y_2y_3,
\]
for which, as an abbreviation, we also set:
\weq{1}{\varphi(y) = h\,\Sigma y_1^3 + 6\,m\,y_1y_2y_3,}
where by the symbol $\Sigma$ we understand a sum over three terms formed from the first by cyclic permutation of the variables $y_1,y_2,y_3$.

To simplify the formulae we shall multiply the \gesperrt{Hesse} covariant by a suitable numerical factor, and accordingly
\weq{2}{\Delta(x) = \frac{1}{6^3}\left|\begin{array}{ccc} f_{11}, & f_{12}, & f_{13} \\ f_{21}, & f_{22}, & f_{23} \\ f_{31}, & f_{32}, & f_{33} \end{array}\right|}

\medskip
\noindent\rule{4cm}{0.4pt}

\smallskip
\noindent $^1$) Of the literature on the points of inflection of curves of the third order and on the theory of ternary cubic forms we mention, besides the older works of \gesperrt{Newton} and \gesperrt{Mac Laurin}: \gesperrt{Hesse}, ``Ueber die Elimination etc.'' and ``Ueber die Wendepunkte der Curven dritter Ordnung'', \gesperrt{Crelle's} Journal, Bd. 28 (1844). \gesperrt{Aronhold}, Theorie der homogenen Functionen 3\textsuperscript{ten} Grades von drei Veränderlichen, \gesperrt{Crelle's} Journal, Bd. 55 (1858). \gesperrt{Salmon}, Higher plane curves (deutsch von \gesperrt{Fiedler}).
%CONT>>


% ===== scan leaf 355, printed 334 =====
% !PRINTED 334
\weberpage{334}
%<<CONT
set, where $f_{ik}$ stands for $f''(x_i,x_k)$. For the canonical form one obtains from this by a simple calculation
\weq{3}{\Delta(y) = -h\,m^2\,\Sigma y_1^3 + (h^3+2m^3)\,y_1y_2y_3.}
From this we derive the second covariant $C_1$ (\S.\,87), into which we likewise introduce a numerical factor, and accordingly
\weq{4}{J(x) = \frac{1}{36}\left|\begin{array}{cccc} f_{11}, & f_{12}, & f_{13}, & \Delta_1 \\ f_{21}, & f_{22}, & f_{23}, & \Delta_2 \\ f_{31}, & f_{32}, & f_{33}, & \Delta_3 \\ \Delta_1, & \Delta_2, & \Delta_3, & 0 \end{array}\right|}
set, where $\Delta_i$ stands for $\Delta'(x_i)$. This covariant too must be represented in the canonical form, which presents no difficulty if one forms the determinant according to the formula Bd. I, \S.\,23, (12). One finds thus:
\weq{5}{\begin{aligned}
J(y) &= -h^2(h^3+8m^3)^2\,\Sigma y_2^3y_3^3\\
&\quad + 3m^2(5h^6+26h^3m^3-4m^6)\,y_1^2y_2^2y_3^2\\
&\quad + hm(2h^6+5h^3m^3+20m^6)\,y_1y_2y_3\,\Sigma y_1^3\\
&\quad + 9h^2m^6(\Sigma y_1^3)^2.
\end{aligned}}
Accordingly the three symmetric functions $y_1y_2y_3$, $\Sigma y_1^3$, $\Sigma y_2^3y_3^3$ can be expressed through the three forms $\varphi(y)$, $\Delta(y)$, $J(y)$, and hence also through the corresponding forms in the original variables $x$, which we denote briefly by $f,\Delta,J$, according to the invariant relations
\weq{6}{\varphi(y) = f,\quad \Delta(y) = r^2\Delta,\quad J(y) = r^6 J}
One obtains first from (1) and (3):
\weq{7}{(h^3+8m^3)\,y_1y_2y_3 = m^2 f + r^2\Delta,}
\weq{8}{h(h^3+8m^3)\,\Sigma y_1^3 = (h^3+2m^3)\,f - 6\,m\,r^2\Delta,}
and then from (5)
\weq{9}{\begin{aligned}
h^2(h^3+8m^3)^2\,\Sigma y_2^3y_3^3 &= (2h^3+m^3)\,m^3 f^2 + m(2h^3-5m^3)\,r^2 f\Delta\\
&\quad + 3m^2 r^4 \Delta^2 - r^6 J.
\end{aligned}}
Very simple now is the calculation of the third covariant $C_2$ (\S.\,87) for the canonical form.

We introduce here too a simplifying numerical factor, and set
\weq{10}{K = K(x) = \frac{1}{9}\left|\begin{array}{ccc} f_1, & \Delta_1, & J_1 \\ f_2, & \Delta_2, & J_2 \\ f_3, & \Delta_3, & J_3 \end{array}\right|,}
%CONT>>


% ===== scan leaf 356, printed 335 =====
% !PRINTED 335
\weberpage{335}
%<<CONT
in which $f_k,\Delta_k,J_k$ denote the derivatives with respect to $x_k$. For the canonical form we then have
\weq{11}{K(y) = r^9 K.}
To find $K(y)$, we need only form the functional determinant from the left-hand sides of the formulae (7), (8), (9), from which, by application of simple determinant theorems, we obtain (Bd. I, \S.\,22, \S.\,27):
\[
K(y) = -h^3(h^3+8m^3)^3 \left|\begin{array}{ccc} 1, & y_1^3, & y_1^6 \\ 1, & y_2^3, & y_2^6 \\ 1, & y_3^3, & y_3^6 \end{array}\right|,
\]
and consequently
\weq{12}{r^9 K = h^3(h^3+8m^3)^3\,(y_1^3-y_2^3)(y_1^3-y_3^3)(y_2^3-y_3^3).}
It is still to be remarked concerning these formulae that the factor $h^3+8m^3$ occurring in them cannot vanish if the curve $f$ has no singular point. For a singular point is determined by the three equations:
\[
h\,y_1^2 + 2m\,y_2y_3 = 0,\quad h\,y_2^2 + 2m\,y_3y_1 = 0,\quad h\,y_3^2 + 2m\,y_1y_2 = 0,
\]
and these three equations are compatible with one another when $h^3+8m^3=0$. (If, for example, $h=-2m$, they are satisfied when $y_1=y_2=y_3$.)

The coefficient $h$ can have any prescribed value different from zero for a given form $f$, and only when this value is given is the problem determined. For the sake of simplicity we will now set $h=1$, and obtain from (7) to (9) the result:

If one sets
\weq{13}{\begin{aligned}
P &= \frac{(2+m^3)m^3f^2+(2-5m^3)mr^2f\Delta+3m^2r^4\Delta^2-r^6J}{(1+8m^3)^2}, \\
Q &= \frac{(1+2m^3)f-6mr^2\Delta}{1+8m^3}, \\
R &= \frac{m^2f+r^2\Delta}{1+8m^3},
\end{aligned}}
then $y_1^3,y_2^3,y_3^3$ are the roots of the cubic equation
\weq{14}{u^3 - Qu^2 + Pu - R^3 = 0,}
and the discriminant of this cubic equation is:
\weq{15}{D_1 = \frac{r^{18}K^2}{(1+8m^3)^6}.}


% ===== scan leaf 357, printed 336 =====
% !PRINTED 336
\weberpage{336}
\pdfbookmark[2]{§ 91. The invariants of the curve of the third order and the biquadratic equation}{bm105}

The coefficients of the cubic equation are known as soon as $m^3, mr^2, r^6$ are, or also when $m$ and $r^2$ are known.

If the equation (14) is then solved, $y_1,y_2,y_3$ are determined up to third roots of unity, which however must satisfy the relation $y_1y_2y_3=R$; and further than this these quantities cannot, by the nature of the matter, be determined.

If the curve $f$ is a real curve, then from the formulae (13) and (15) there follows further the theorem that, \gesperrt{when $m$ and $r$ are real, $D_1$ is positive, and consequently the roots of the equation (14) are all three real.}

If one expresses the discriminant $D_1$ in terms of the coefficients of the equation $P,Q,R$, then from the formula (15) there flows also the result that the square of the covariant $K$ can be expressed as an integral rational function of the covariants $f,\Delta,J$. The expression becomes rather complicated and shall not be pursued further here.

But in order to form the biquadratic equation on which $m$ and $r$ depend, it is necessary to enter more closely into the invariants of the third order.

\begin{center}
\S.\,91.

\gesperrt{The invariants of the curve of the third order and the biquadratic equation.}
\end{center}

All curves of the third order whose equation in the canonical form can be represented linearly by the same quantities $\Sigma y_1^3$, $y_1y_2y_3$ have the same points as points of inflection — for example, therefore, the curves $f=0$ and $\Delta=0$, and more generally all curves of the family
\weq{1}{F = \lambda f + 6\mu\Delta = 0,}
where $\lambda,\mu$ are arbitrary parameters. All curves of this family intersect in nine fixed points, and such a family is called a \gesperrt{pencil of curves}. The nine fixed points, which are the points of inflection for all curves of the pencil, are called the \gesperrt{fundamental points}. If we form from the form $F$ the \gesperrt{Hesse} determinant, this form too is expressed linearly by $\Sigma y_1^3$ and $y_1y_2y_3$, and can therefore by \S.\,90, (7) and (8) also be expressed linearly by $f$ and $\Delta$, i.e. we obtain a form of the same pencil.


% ===== scan leaf 358, printed 337 =====
% !PRINTED 337
\weberpage{337}
We therefore set this determinant
\weq{2}{\Delta_{\lambda,\mu}(x) = Lf + M\Delta,}

wherein $L$ and $M$ are binary forms of the third degree in
$\lambda, \mu$.

If we apply an arbitrary linear transformation to the variables $x$,
then, on account of the invariant property of the function $\Delta$,
\weq{3}{\Delta_{\lambda,\mu}(y,b) = r^2\Delta_{\lambda,\mu\,r^2}(x,a),}

and from this it follows that the coefficients of the individual powers
and products of $\lambda, \mu$ are covariants of the original form, and
the coefficients $L, M$ of $f$ and $\Delta$ in these expressions, which
depend rationally on the coefficients $a$ of the form $f$, are therefore
\gesperrt{invariants} of this form. We have in this a means of actually computing
these invariants.

For if we denote by $a_i, A_i, B_i$ the coefficients of the
individual powers and products of the variables $x$ in $f, \Delta$ and
$\Delta_{\lambda,\mu}$, then from (2) there results a system of ten
equations:
\[
L a_i + M A_i = B_i,
\]
from which $L$ and $M$ can be determined. One obtains, if $i, k$ are two
distinct indices,
\weq{4}{\begin{aligned}
L\,(a_i A_k - a_k A_i) &= B_i A_k - B_k A_i\\
M\,(a_i A_k - a_k A_i) &= -B_i a_k + B_k a_i.
\end{aligned}}

From this one concludes further that $L, M$ are \gesperrt{integral}
functions of the coefficients $a_i$. For in (4) the right-hand sides are
integral functions, and if therefore $L$ or $M$, represented in the
simplest form, still had a denominator, then all the determinants
$a_iA_k - a_kA_i$, which are integral functions of the fourth degree in
the $a$, would have to have a common divisor. But that this is not
possible can be shown in any special case, e.g. for the form
\[
f = 3x_3(\alpha x_1^2 + \beta x_2^2) + a x_1^3 + b x_2^3,
\]
whose Hesse form is:
\[
-\beta^2(\alpha x_3 + a x_1)x_2^2 - \alpha^2(\beta x_3 + b x_2)x_1^2
\]
Here among the $a_iA_k - a_kA_i$ there occur, e.g., the two
functions $a^2\beta^2, b^2\alpha^2$, which have no common
divisor.


% ===== scan leaf 359, printed 338 =====
% !PRINTED 338
\weberpage{338}
% running head: 338 | Elfter Abschnitt. | §. 91.

In order to compute $\varDelta_{\lambda,\mu}$ for the canonical form, one has,
in \S.\,90, (3), to replace
\[
h, \qquad m
\]
by
\[
h(\lambda-6m^2\mu), \qquad m\lambda+(h^3+2m^3)\mu
\]
and one thus finds:
\[
\varDelta_{\lambda,\mu} = -\,h(\lambda-6m^2\mu)\,[m\lambda+(h^3+2m^3)\mu]^2\,\Sigma y_1^3
+ \{h^3(\lambda-6m^2\mu)^3+2[m\lambda+(h^3+2m^3)\mu]^3\}\,y_1y_2y_3,
\]
from which, by \S.\,90, (7), (8), with regard to (2), it follows after some
simple calculations that:
\weq{5}{\begin{gathered}
L = -2\lambda^2\mu(h^3-m^3)m - \lambda\mu^2(h^6-20h^3m^3-8m^6) \\
\qquad{}+ 8\mu^3m^2(h^3-m^3)^2, \\
M = \lambda^3+12\lambda\mu^2m(h^3-m^3)+2\mu^3(h^6-20h^3m^3-8m^6).
\end{gathered}}

In this way, then, one obtains only two invariants, of degree
$4$ and degree $6$, namely, in the canonical form:
\weq{6}{m(h^3-m^3), \quad h^6-20h^3m^3-8m^6.}

The weight of these invariants is [by Bd. I, \S.\,60, (3)]
likewise 4 and 6, and if we therefore denote them in the general
form by $S$ and $T$ and set $h=1$, we find
\weq{7}{\begin{gathered}
m(1-m^3) = r^4S, \\
1-20m^3-8m^6 = r^6T.
\end{gathered}}

For $L$ and $M$ one obtains from (5) the representation:
\weq{8}{\begin{gathered}
L = -2S\lambda^2\mu - T\lambda\mu^2 + 8S^2\mu^3 \\
M = \lambda^3+12S\lambda\mu^2+2T\mu^3.
\end{gathered}}

The functions $S$ and $T$ of degree $4$ and degree $6$ were
first computed by \gesperrt{Aronhold}. We shall not set down the long
expressions here, but from now on regard these two
quantities as known and given.

From (7) a biquadratic equation can be derived
for the ratio $m^2:r^2$. One finds from it, namely:
\weq{9}{\begin{gathered}
m^5-m^8 = r^4m^4S \\
m^2-2m^5+m^8 = r^8S^2 \\
m^2-20m^5-8m^8 = r^6m^2T.
\end{gathered}}

If one eliminates $m^2$ and $m^5$ from these, one finds
\weq{10}{27m^8+18Sm^4r^4+Tm^2r^6-S^2r^8=0.}


% ===== scan leaf 360, printed 339 =====
% !PRINTED 339
\weberpage{339}
% running head: §. 91. | Invariants of the Curve of the Third Order. | 339

This equation takes a simpler form through the
substitution
\weq{11}{z = \frac{3m^2}{r^2}.}
It then becomes
\weq{12}{z^4+6Sz^2+Tz-3S^2=0.}
This biquadratic equation has the property that
its first invariant $A=0$, and that its discriminant
(Bd. I, \S.\,47):
\weq{13}{D_2 = -27(T^2+64S^3)^2 = -27D^2,}
is thus always negative. The quantity $D=T^2+64S^3$ takes for
the canonical form (5) the expression:
\weq{14}{h^3(h^3+8m^3)^3,}
and hence cannot vanish so long as the curve $f$ has no
singular point. (It is in fact nothing other than
the discriminant of the form $f$.)

From (13) it now follows that the biquadratic equation
 (12) has, for real curves $f$, a negative discriminant, and
consequently (Bd. I, \S.\,78) two real and two conjugate imaginary
roots, and since the product of all four roots, $-3S^2$,
is negative, \gesperrt{of the real roots one is positive},
one negative.

From $z$ one can calculate, with the help of (11) and the first equation
 (7), the two quantities $m$ and $r$. One obtains in this way
\weq{15}{m^3 = \frac{z^2}{9S+z^2}, \quad mr^2 = \frac{3z}{9S+z^2}, \quad r^6 = \frac{27z}{(9S+z^2)^2},}
and here, as one concludes from (12) by setting
$z^2=-9S$, the quantity $9S+z^2$ could vanish only if either
 $S$ or $D$ is zero. With this the coefficients $P$, $Q$, $R^3$
of the cubic equation \S.\,90, (14) are uniquely determined by $z$ and by
known quantities, and for the real positive root $z$
the values $m$ and $r$ become real, and consequently $y_1$, $y_2$, $y_3$ real as well.

A slight exception is formed by the case $S=0$, in which one
root of equation (12) vanishes, and the second real
root becomes $-\sqrt[3]{T}$, and for the corresponding values of $m$
and $r^6$ one obtains from (7) and (11)
\[
m = 0, 1; \quad r^6 = \frac{1}{T}, \ \frac{-27}{T};
\]
%CONT>>


% ===== scan leaf 361, printed 340 =====
% !PRINTED 340
\weberpage{340}
% running head: 340 | Elfter Abschnitt. | §. 91.
%<<CONT
according, then, as $T$ is positive or negative, the first or the
second value of $r$ turns out real, and in essentials everything
remains as before.

From this results the theorem:

\gesperrt{Every real ternary cubic form with non-vanishing
discriminant can be brought by a real
transformation to the real canonical form.}

One can now easily give the coordinates of the points of inflection
in the canonical coordinate system. One obtains them
by setting, in the equation
\[
y_1^3+y_2^3+y_3^3+6m\,y_1y_2y_3=0
\]
$y_1$, $y_2$, $y_3=0$ in turn. Since one of the non-vanishing
coordinates can be set $=1$, one finds, if $\varepsilon$
is an imaginary third root of unity, the following values
of the coordinates $(y_1, y_2, y_3)$:
\weq{16}{\begin{array}{ccc}
(0, 1, -1), & (0, 1, -\varepsilon), & (0, 1, -\varepsilon^2), \\
(-1, 0, 1), & (-\varepsilon, 0, 1), & (-\varepsilon^2, 0, 1), \\
(1, -1, 0), & (1, -\varepsilon, 0), & (1, -\varepsilon^2, 0),
\end{array}}
and it follows from this that \gesperrt{always three of the inflexion points
are real and six imaginary.}

In the table (16), the three points standing in any one row
lie on a straight line, namely on the lines
\[
y_1=0, \quad y_2=0, \quad y_3=0.
\]

It is easy, after this, to assemble also the remaining triples of points
that lie on a straight line, and thereby
to determine these twelve lines. One has only always to
assemble such points whose coordinates give a vanishing
determinant. One can arrange these twelve lines
in triples in four ways, so that in each triple all
nine inflexion points occur. The first of these triples is
formed by the rows of (16). The three others read thus:
\weq{17}{\begin{array}{ccc}
(0, 1, -1), & (-1, 0, 1), & (1, -1, 0), \\
(0, 1, -\varepsilon), & (-\varepsilon, 0, 1), & (1, -\varepsilon, 0), \\
(0, 1, -\varepsilon^2), & (-\varepsilon^2, 0, 1), & (1, -\varepsilon^2, 0),
\end{array}}
%CONT>>


% ===== scan leaf 362, printed 341 =====
% !PRINTED 341
\weberpage{341}
% running head: §. 91. | Reality. | 341
%<<CONT

\weq{18}{\begin{array}{ccc}
(0, 1, -1), & (-\varepsilon, 0, 1), & (1, -\varepsilon^2, 0) \\
(-1, 0, 1), & (1, -\varepsilon, 0), & (0, 1, -\varepsilon^2) \\
(1, -1, 0), & (0, 1, -\varepsilon), & (-\varepsilon^2, 0, 1),
\end{array}}
\weq{19}{\begin{array}{ccc}
(0, 1, -1), & (-\varepsilon^2, 0, 1), & (1, -\varepsilon, 0) \\
(-1, 0, 1), & (1, -\varepsilon^2, 0), & (0, 1, -\varepsilon) \\
(1, -1, 0), & (0, 1, -\varepsilon^2), & (-\varepsilon, 0, 1).
\end{array}}

Of the straight lines (16), each contains one real and two
conjugate imaginary points, and all three lines are real. Of the
lines (17), the first contains the three real points and is real,
while the other two are conjugate imaginary. In (18) and (19)
all three lines are imaginary, and indeed the lines of (18) are
conjugate to the lines of (19).

A still clearer notation is the following. If
$\xi$ and $\eta$ each run through a complete system of residues with respect to the modulus 3,
say $0, 1, 2$, and two numbers congruent with respect to the modulus 3
are regarded as not distinct, then there are nine
combinations $(\xi, \eta)$, which we can use as a designation for the nine
points of inflection. If we stipulate that $(\xi, \eta)$ and
$(\eta, \xi)$ denote conjugate imaginary points of inflection, and hence $(0, 0)$, $(1, 1)$, $(2, 2)$
the three real ones, then from (16) to (19) we obtain the following
 table, in which again the points of inflection standing in one row
 lie on a straight line:
\weq{20}{\begin{array}{ccc@{\qquad}ccc}
(0,0) & (1,2) & (2,1) & (0,0) & (1,1) & (2,2) \\
(1,1) & (2,0) & (0,2) & (1,2) & (2,0) & (0,1) \\
(2,2) & (0,1) & (1,0) & (2,1) & (0,2) & (1,0) \\[4pt]
(0,0) & (2,0) & (1,0) & (0,0) & (0,2) & (0,1) \\
(1,1) & (0,1) & (2,1) & (1,1) & (1,0) & (1,2) \\
(2,2) & (1,2) & (0,2) & (2,2) & (2,1) & (2,0).
\end{array}}

One sees from this that three points of inflection $(\xi_1, \eta_1)$, $(\xi_2, \eta_2)$,
$(\xi_3, \eta_3)$ lie on a straight line then and only then, when
the two congruences
\weq{21}{\begin{aligned}
\xi_1+\xi_2+\xi_3 &\equiv 0 \\
\eta_1+\eta_2+\eta_3 &\equiv 0
\end{aligned} \quad (\mathrm{mod}\ 3)}
are satisfied.


% ===== scan leaf 363, printed 342 =====
% !PRINTED 342
\weberpage{342}
\pdfbookmark[2]{§ 92. Triple equations}{bm106}

One can also derive the designation of the four systems of straight
lines from the single scheme:
\[
\begin{array}{ccc}
(0,\,0) & (1,\,2) & (2,\,1)\\
(1,\,1) & (2,\,0) & (0,\,2)\\
(2,\,2) & (0,\,1) & (1,\,0)
\end{array}
\]
by setting up the rows, the columns, and the combinations of three
that correspond to the positive and negative terms in the formation
of the determinant.

\S.\,92.

\gesperrt{Triple equations.}

The system of the nine points of inflection has the property
that from any two of them one can derive a wholly determinate
third in a rational way. This expresses itself as a property of the
equation of degree $9$ on which the determination of the points of
inflection depends, which has, e.g., the abscissae of these points
as its roots. We set up the following definition:

\gesperrt{An equation without equal roots is called a triple
equation if every two of its roots determine a third,
so that each root of such a triple can be expressed
rationally in terms of the other two, and indeed in such
a way that in an equation}
\weq{1}{x_3 = \Theta\,(x_1, x_2),}
\gesperrt{in which $\Theta$ denotes a fixed rational function, one may put for
$x_1, x_2, x_3$ the roots of any triple whatever
in any order}\footnote{The triple equations of degree $9$ were first treated by
\gesperrt{Hesse} in the memoir: ``Algebraische Auflösung derjenigen Gleichung
$9^{\text{ten}}$ Grades etc.'' in Crelle's Journal, Bd.\,34 (1847).}.

The equation on which the points of inflection depend is a
triple equation of degree $9$. But there are also triple equations
of other degrees, e.g. of degree $7$, to which belong those considered
in the preceding chapter\footnote{Cf. \gesperrt{N\"other}, Mathem. Annalen, Bd.\,15 (1879): ``Ueber die Gleichungen
$8^{\text{ten}}$ Grades und ihr Auftreten in der Theorie der Curven vierter
Ordnung.''}.


% ===== scan leaf 364, printed 343 =====
% !PRINTED 343
\weberpage{343}

We shall therefore consider here only the triple equations of degree $9$,
and first show how the designation and arrangement of the roots
that we became acquainted with in the previous section for the
points of inflection can be derived from the triple property.

Each of the nine roots occurs in four triples, and in all there
therefore exist $4.9:3=12$ distinct triples. Let us take an
arbitrary one of these triples and denote its
roots by
\[
(0,\,0)\ (1,\,2)\ (2,\,1).
\]
Then let $(1,\,1)$ be an arbitrary fourth root. With the first
three this determines, to begin with, three triples, which we denote by
\[
(1,\,1)\ (0,\,0)\ (2,\,2), \quad (1,\,1)\ (1,\,2)\ (1,\,0), \quad (1,\,1)\ (2,\,1)\ (0,\,1)
\]
and there remain two more roots, which form the
fourth triple with $(1,\,1)$, and which we now denote by
\[
(1,\,1)\ (2,\,0)\ (0,\,2)
\]

We now seek the triple determined by $(0,\,0), (2,\,0)$.
In this triple no roots can occur which occur with either of the
two roots $(0,\,0), (2,\,0)$ in one of the triples already determined, % [sic] German source prints feminine "einer" with neuter "Tripel"; expected "einem"
hence not
\[
(1,\,2)\ (2,\,1)\ (1,\,1)\ (2,\,2)\ (0,\,2),
\]
and for the third root of this triple only $(0,\,1)$
or $(1,\,0)$ remains, and the same two roots are also all that remain
for the triple determined by $(0,\,0)\,(0,\,2)$. But since we can
still interchange $(0,\,2)$ with $(2,\,0)$, we do not restrict
generality if the two further triples
\[
(0,\,0)\ (2,\,0)\ (1,\,0), \quad (0,\,0)\ (0,\,2)\ (0,\,1)
\]
are assumed. We now determine the triples that contain $(1,\,0)$,
of which two are already known. The two others
must contain the roots $(2,\,2), (0,\,1), (0,\,2), (2,\,1)$. But since
$(0,\,1)\,(1,\,0)\,(0,\,2)$ and $(0,\,1)\,(1,\,0)\,(2,\,1)$ are not triples [because
$(0,\,1)\,(0,\,2)$ determines the triple $(0,\,0)\,(0,\,2)\,(0,\,1)$ and $(0,\,1)\,(2,\,1)$ the
triple $(1,\,1)\,(2,\,1)\,(0,\,1)$], we have the further
triples
\[
(2,\,2)\ (0,\,1)\ (1,\,0), \quad (1,\,0)\ (0,\,2)\ (2,\,1).
\]
In the same way we obtain
\[
(0,\,1)\ (2,\,0)\ (1,\,2).
\]


% ===== scan leaf 365, printed 344 =====
% !PRINTED 344
\weberpage{344}
\pdfbookmark[2]{§ 93. The group of the triple equations}{bm107}

Finally we form one more triple containing $(0,\,2)$ and one
containing $(2,\,0)$,
\[
(0,\,2)\ (1,\,2)\ (2,\,2), \quad (2,\,0)\ (2,\,1)\ (2,\,2),
\]
with which all twelve triples, and with them the arrangement \S.\,91, (20),
are determined.

\S.\,93.

\gesperrt{The group of the triple equations.}

We shall now determine the \gesperrt{Galois} group $P$ of the triple
equations. First of all it is clear that in this group only
such permutations occur under which every triple passes again
into a triple. For if we assume that the roots of a
triple $x_1, x_2, x_3$ pass by a permutation into $y_1, y_2, y_3$,
then, since this permutation is applicable to the equation (1) (\S.\,92),
it follows that
\[
y_1 = \Theta\,(y_2, y_3),
\]
and consequently $y_1$ is the third root of the triple determined
by $y_2, y_3$.

If we now pass over to the notation for the roots:
\weq{1}{x = (\xi, \eta)}
then it results from the compilation of the triples in
§§.\,91 and 92 that three roots
\weq{2}{x_1 = (\xi_1, \eta_1), \quad x_2 = (\xi_2, \eta_2), \quad x_3 = (\xi_3, \eta_3)}
always form a triple then, and only then, when
\weq{3}{\begin{array}{l}
\xi_1 + \xi_2 + \xi_3 \equiv 0\\
\eta_1 + \eta_2 + \eta_3 \equiv 0
\end{array} \quad (\mathrm{mod}\ 3).}
If we call a permutation group of the $(\xi, \eta)$ whose substitutions
all belong to the complete linear congruence group for
the modulus 3 (\S.\,78) a \gesperrt{linear group}, then the following theorem
holds:

\gesperrt{1. The group $P$ of a triple equation is always
linear.}

According to \S.\,77, every permutation of the quantities $(\xi, \eta)$
whatever can be represented in the form:
\weq{4}{\xi' \equiv \varphi\,(\xi, \eta), \quad \eta' \equiv \psi\,(\xi, \eta) \quad (\mathrm{mod}\ 3),}
%CONT>>


% ===== scan leaf 366, printed 345 =====
% !PRINTED 345
\weberpage{345}

%<<CONT
in which $\varphi, \psi$ are integral functions of the variables $\xi, \eta$
with integral coefficients, whose degree with respect to neither of the variables
exceeds the value 2. If we order these functions according to $\eta$, we may
set:
\weq{5}{\xi' = A\eta^2 + B\eta + C, \quad \eta' = A'\eta^2 + B'\eta + C',}
in which $A, B, C, A', B', C'$ are integral functions of $\xi$, at most
of degree $2$. In the equations (4) we shall hold the
variable $\xi$ fixed and set $\eta = 0, 1, 2$; that is, we apply
(5) to the triple $(\xi, 0)\,(\xi, 1)\,(\xi, 2)$. The corresponding
$(\xi_1', \eta_1'), (\xi_2', \eta_2'), (\xi_3', \eta_3')$ must then satisfy the relations (3),
hence:
\[
\begin{array}{ll}
\xi_1' \equiv C, & \eta_1' \equiv C',\\
\xi_2' \equiv A + B + C, & \eta_2' \equiv A' + B' + C',\\
\xi_3' \equiv A - B + C, & \eta_3' \equiv A' - B' + C',
\end{array}
\]
and from this it follows that $A$ and $A'$ must be congruent to
zero for every $\xi$.

In the same way one may conclude that no term with $\xi^2$ can
occur in the substitutions (4), and that these substitutions
 must therefore have the form:
\[
\begin{array}{l}
\xi' \equiv m\,\xi\,\eta + a\,\xi + b\,\eta + c\\
\eta' \equiv m'\,\xi\,\eta + a'\,\xi + b'\,\eta + c'
\end{array} \quad (\mathrm{mod}\ 3).
\]
But if one applies these substitutions to the triple $(0,\,0)\,(1,\,1)\,(2,\,2)$,
and sets the sum of the corresponding $\xi'$ and $\eta'$ equal to
zero, it follows that we must have $m = m' = 0$, and hence that
every permutation of the group $P$ must be contained in the form
\weq{6}{\begin{array}{l}
\xi' \equiv a\,\xi + b\,\eta + c\\
\eta' \equiv a'\,\xi + b'\,\eta + c'
\end{array} \quad (\mathrm{mod}\ 3)}
as was asserted.

\gesperrt{2. Conversely, by every linear substitution
of the form (6) three quantities $(\xi, \eta)$ of a triple
pass into three $(\xi', \eta')$, which likewise form a
triple.}

For from (6) it follows, for any three pairs of numbers $(\xi_1, \eta_1)$,
$(\xi_2, \eta_2), (\xi_3, \eta_3)$:
\weq{7}{\begin{array}{l}
\xi_1' + \xi_2' + \xi_3' \equiv a\,(\xi_1 + \xi_2 + \xi_3) + b\,(\eta_1 + \eta_2 + \eta_3)\\
\eta_1' + \eta_2' + \eta_3' \equiv a'\,(\xi_1 + \xi_2 + \xi_3) + b'\,(\eta_1 + \eta_2 + \eta_3)
\end{array} \quad (\mathrm{mod}\ 3);}
%CONT>>


% ===== scan leaf 367, printed 346 =====
% !PRINTED 346
\weberpage{346}
%<<CONT
thus if $\xi_1 + \xi_2 + \xi_3$ and $\eta_1 + \eta_2 + \eta_3$ are congruent to zero,
then the same holds of the left sides of (7).

Now conversely, under every linear substitution of the form (6), three quantities
($\xi, \eta$) of a triple always pass again into a triple, and from this it can be
proved that every equation of degree 9 whose Galois group $P$ is contained in the
complete linear congruence group is a triple equation, provided the further condition
is added that by the group $P$ any two given roots pass into two other, likewise
arbitrarily given roots — or, what amounts to the same thing, if $P$ is
\gesperrt{doubly transitive}.

For if $P$ consists entirely of linear substitutions, then by the adjunction of two
arbitrary roots $x_1 = (\xi_1, \eta_1)$, $x_2 = (\xi_2, \eta_2)$, $P$ reduces to a
group $P_1$ which leaves unchanged not only $x_1, x_2$, but also the third root
$x_3 = (\xi_3, \eta_3)$ of the triple determined by $x_1, x_2$; consequently $x_3$
belongs to the new domain of rationality, and hence $x_3$ is rationally expressible
in terms of $x_1, x_2$ in the form

\weq{8}{x_3 = \Theta (x_1, x_2).}

Now if the group $P$ is doubly transitive, there is a permutation in $P$ by which
$x_1, x_2$ pass into two arbitrary other roots $y_1, y_2$, and by which consequently
$x_3$ also passes into the third root of the triple $y_1, y_2, y_3$; and this
permutation is applicable to (8). Hence also

\[
y_3 = \Theta (y_1, y_2).
\]

We state this as a theorem:

3. \gesperrt{Every equation of degree 9 whose group is linear
and doubly transitive is a triple equation.}

By the permutations of a doubly transitive linear group every triple can be carried
into every other, and indeed in arbitrary arrangement.

The hypothesis of double transitivity can also be expressed thus: that all the
permutations of $P$ which leave one root unchanged still connect the remaining roots
transitively with one another. If we denote by $P_0$ the group contained in $P$ by
whose permutations one of the roots $x_0$ remains unchanged, then $P_0$ must still be
transitive with respect to the other roots. But if the group $P$ is only simply
transitive,
%CONT>>


% ===== scan leaf 368, printed 347 =====
% !PRINTED 347
\weberpage{347}
%<<CONT
so $P_0$ is no longer transitive, that is, after adjunction of $x_0$
the equation for the remaining eight roots splits into several factors. If the
group $P$ is transitive at all, then $x_0$ can here be any one of the roots;
for if $\pi$ is a permutation of $P$ by which $x_0$ passes into $x_1$,
then by $P_1 = \pi^{-1} P_0 \pi$ the root $x_1$ remains unchanged, and if $P_0$
is intransitive for eight roots, then $P_1$ must also be intransitive for eight roots,
since the permutations of $P_1$ differ from those of $P_0$ only in the designation of the
roots (Bd. I, §.\,154, 6.).

The simply transitive ones among the linear groups have a more special character
than the group of the triple equations.

In §.\,78 we became acquainted with a decomposition of the general linear congruence
group for the modulus 3, from which we can survey these relations.

The \gesperrt{Abelian} group $Q$, which consists of the substitutions

\weq{9}{\xi' \equiv \xi + \alpha, \quad \eta' \equiv \eta + \beta \quad (\mathrm{mod}\ 3)}

is certainly only simply transitive; for if a root is to remain unchanged
by (9), then we must have $\alpha \equiv 0$, $\beta \equiv 0$, and all roots remain
unchanged. Every root is rationally expressible in terms of every other.

If we now seek out in the group $P$ the group $P_0$ of the substitutions by which
$(0,0)$ remains unchanged, then $P_0$ contains only homogeneous substitutions and is
therefore the greatest common divisor of $P$ and the homogeneous congruence group $S$,
and we have $P = P_0 Q$.

If we form for $S$ the composition series (§.\,78):

\[
S,\ S_1,\ S_2,\ S_3,\ S_4,
\]

then $S_4 = \begin{pmatrix} 1, & 0 \\ 0, & 1 \end{pmatrix}, \begin{pmatrix} -1, & 0 \\ 0, & -1 \end{pmatrix}$,
and by $S_4$ the roots $(1,1)$, $(2,2)$ merely pass into one another. By $S_4 Q$
the triple $(0,0)$ $(1,1)$ $(2,2)$ can indeed undergo every arrangement of its roots.
But it can pass over only into the two other triples $(1,0)$ $(2,1)$ $(0,2)$;
$(0,1)$ $(1,2)$ $(2,0)$. $S_4 Q$ is therefore only \gesperrt{simply transitive} and
is moreover \gesperrt{imprimitive}.

By the group $S_3$, which arises from $S_4$ by the addition of the
substitution

\[
\begin{pmatrix} -1, & 1 \\ 1, & 1 \end{pmatrix}
\]

%CONT>>


% ===== scan leaf 369, printed 348 =====
% !PRINTED 348
\weberpage{348}
\pdfbookmark[2]{§ 94. Reality conditions of the triple equations}{bm108}
%<<CONT
arises, $(1,1)$, $(2,2)$ can still pass over into $(0,2)$, $(0,1)$, but not into
$(2,0)$, $(1,0)$ or into $(1,2)$, $(2,1)$. Hence the group $S_3 Q$ is likewise not yet
doubly transitive, but it is primitive. If, however, in order to form $S_2$,
we adjoin the substitution

\[
\begin{pmatrix} 0, & 1 \\ -1, & 0 \end{pmatrix}
\]

then $(1,1)$, $(2,2)$ passes over into $(0,2)$, $(0,1)$ and into $(1,2)$, $(2,1)$ and
$(2,0)$, $(1,0)$. Hence the group $S_2 Q$ is doubly transitive, and consequently so are
$S_1 Q$ and $S Q$. The triple equations therefore first begin with $S_2 Q$.
The equations with narrower groups are not yet proper triple equations.

We can form still other linear groups contained in $S Q$, e.\,g. if we
denote by $S'$ the cyclic group of degree 3:

\[
S' = \begin{pmatrix} 1, & 0 \\ 0, & 1 \end{pmatrix}, \begin{pmatrix} 1, & 0 \\ 1, & 1 \end{pmatrix}, \begin{pmatrix} 1, & 0 \\ -1, & 1 \end{pmatrix}
\]

the group of degree 27, $S' Q = Q S'$. Through the group $S'$,
$(1,1)$ $(2,2)$ passes over into $(1,1)$ $(2,2)$, $(1,2)$ $(2,1)$, $(1,0)$ $(2,0)$, while
$(0,1)$ $(0,2)$ does not arise from it. Hence this too is not the group of a
triple equation.

Here there would still be a region for further investigations into these special
metacyclic equations of degree 9 and the representation of their
roots according to the principles of \gesperrt{Kronecker} (Bd. I, eighteenth chapter).

\begin{center}
\S.\,94.\\[2pt]
\gesperrt{Reality Relations of the Triple Equations.}
\end{center}

If the domain of rationality is real, we can draw some general conclusions
about the reality of the roots of a triple equation of degree 9.
For if in this case two real roots occur in one triple, then,
as follows from the equation $x_3 = \Theta (x_1, x_2)$, the third root must also be
real. If, therefore, there is any imaginary root at all, then in each of the
four triples to which this root belongs there must occur at least one further imaginary
root, and consequently the number of imaginary roots is at least equal to 5,
or, since the number of imaginary roots
%CONT>>


% ===== scan leaf 370, printed 349 =====
% !PRINTED 349
\weberpage{349}
%<<CONT
must be even — equal to 6. Two real or two conjugate imaginary roots
$x_1, x_2$ always belong to a triple whose third root is real; this
follows from the equation:

\[
x_3 = \Theta (x_1, x_2) = \Theta (x_2, x_1).
\]

We shall accordingly call a triple which contains two real or two conjugate
imaginary roots a \gesperrt{real triple}. If any imaginary roots occur in a
triple, then the conjugate imaginary roots also form a triple, and two such
triples shall be called \gesperrt{conjugate imaginary triples}.

There are then three possibilities:

1. \gesperrt{All real roots.}

2. \gesperrt{One real and eight imaginary roots.}

In this case the four pairs of conjugate imaginary roots must determine four
triples, all of which contain the same real root as third. If we denote
this real root by $(0,0)$, then

\[
(1,2)\ (2,1),\quad (1,1)\ (2,2),\quad (1,0)\ (2,0),\quad (0,1)\ (0,2)
\]

must be the four pairs of conjugate roots.

3. \gesperrt{Three real and six imaginary roots.}

The three real roots must here form a triple. If we take for the real
roots

\weq{1}{(0,0)\ (1,1)\ (2,2),}

then in the two rows

\weq{2}{\begin{aligned} &(1,2)\ (0,2)\ (1,0) \\ &(2,1)\ (2,0)\ (0,1) \end{aligned}}

conjugate imaginary roots cannot occur in the same row, because otherwise
the third root of the row in question would be real.

It is to be shown that the three real triples determined by the conjugate
pairs must together contain all three real points.

For suppose that, among the three triples

\[
(0,0)\ (1,2)\ (2,1),\quad (0,0)\ (2,0)\ (1,0),\quad (0,0)\ (0,1)\ (0,2)
\]

two real ones occur; then the third must also be real, i.\,e.

\[
(1,2)\ (2,1),\quad (2,0)\ (1,0),\quad (0,1)\ (0,2)
\]
%CONT>>


% ===== scan leaf 371, printed 350 =====
% !PRINTED 350
\weberpage{350}
%<<CONT
be conjugate pairs. But then the three roots $(1,1)$ $(2,0)$ $(0,2)$ form a triple, whose conjugates $(1,1)$ $(1,0)$ $(0,1)$ form no triple, which is impossible. Each of the three real roots $(0,0)$, $(1,1)$, $(2,2)$ therefore occurs, besides in (1), in one and only one real triple, whose two other roots are conjugate imaginary; and the third case thus leads to the arrangement of the real and imaginary roots that we have come to know in the points of inflection of the curve of the third order.

But that the cases 1. and 2. can also occur is shown by the following consideration: From a general equation of degree $9$ one obtains a triple equation by adjunction of a function belonging to the linear group $SQ$. Such a function is, e.g.:

\[
\begin{aligned}
v &= (0,0)(1,1)(2,2) + (0,0)(1,2)(2,1) + (0,0)(1,0)(2,0) \\
&\quad + (0,0)(0,1)(0,2) + (1,1)(1,2)(1,0) + (1,1)(2,1)(0,1) \\
&\quad + (1,1)(0,2)(2,0) + (1,2)(0,1)(2,0) + (2,2)(2,1)(2,0) \\
&\quad + (2,2)(1,2)(0,2) + (2,2)(0,1)(1,0) + (2,1)(0,2)(1,0),
\end{aligned}
\]

and this function is real in the cases 1., 2., 3. Hence the extended domain of rationality, in which our equation is a triple equation, is also real.

\begin{center}
\rule{0.22\textwidth}{0.4pt}
\end{center}


% ===== scan leaf 372, printed 351 =====
\weberpage{351}% page number not printed on this page
\pdfbookmark[1]{Twelfth Section. The Double Tangents of a Curve of the Fourth Order.}{bm109}
\pdfbookmark[2]{§ 95. Number of the double tangents of a curve of the fourth order}{bm110}
% !PRINTED none
% Abschnitt-opening page: no running head/page number printed on this recto.
\begin{center}
\gesperrt{Twelfth Chapter.}

\textbf{The Double Tangents of a Curve of the Fourth Order.}

\rule{0.3\linewidth}{0.4pt}

§. 95.

\gesperrt{Number of the Double Tangents of a Curve\\
of the Fourth Order.}
\end{center}

A geometrical problem which, on account of its manifold relations to other regions, is of particular importance, and which at the same time, in a similar way to the problem of the points of inflection of the curves of the third order, offers remarkable algebraic relations, is that of the double tangents of the curves of the fourth order.

By a \gesperrt{double tangent} of a curve one understands, as has already been remarked in \S.\,86, a straight line which touches the curve in two distinct points. In curves of lower than the fourth order, double tangents cannot occur.

In the curves of the fourth order a double tangent has, besides the points of contact, no point in common with the curve. In special cases the two points of contact may also coincide. Then we have lines with four-point contact.

The determination of the double tangents will, as an algebraic problem, depend on a certain algebraic equation whose degree is equal to the number of the double tangents, and the first question is that of the degree of this equation, that is, of the number of the double tangents. We restrict ourselves here to the consideration of curves of the fourth order, although the way that we take is applicable also to curves of higher order
%CONT>>


% ===== scan leaf 373, printed 352 =====
% !PRINTED 352
\weberpage{352}
%<<CONT
is. We also assume that the curve of the fourth order is free of singular points\footnote{\gesperrt{Jacobi}, \glqq Ueber die Anzahl der Doppeltangenten ebener algebraischer Curven\grqq{}, \gesperrt{Crelle's} Journal, Bd.\,40 (1850). \gesperrt{Clebsch}, \glqq Bemerkung zu \gesperrt{Jacobi's} Beweis für die Anzahl der Doppeltangenten\grqq{}, ebend., Bd.\,63 (1864). From the rather extensive literature on the theory of the double tangents of a curve of the fourth order the following writings are to be singled out. First, several papers by \gesperrt{Hesse} in \gesperrt{Crelle's} Journal, Bd.\,41, 49, 52. Further \gesperrt{Steiner}, Eigenschaften der Curven vierter Ordnung rücksichtlich ihrer Doppeltangenten, \gesperrt{Crelle's} Journal, Bd.\,49. \gesperrt{Aronhold}, Monatsber. d. Berl. Akademie 1864. \gesperrt{Geiser}, \glqq Ueber die Doppeltangenten einer ebenen Curve 4\,ten Grades\grqq{}, Math. Annalen, Bd.\,1 (1868). \gesperrt{Cayley}, \gesperrt{Crelle's} Journal, Bd.\,68 (1868). The work of \gesperrt{Salmon} cited above is also to be singled out here. In more recent times the theory of the double tangents has been developed further in connection with the theory of the \gesperrt{Abelian} functions. \gesperrt{Riemann}, Gesammelte mathematische Werke (2.\,Aufl., 1892) XXXI, aus dem Nachlass. \gesperrt{Weber}, Theorie der \gesperrt{Abel'schen} Functionen vom Geschlecht 3, Berlin 1876. \gesperrt{Frobenius}, \gesperrt{Crelle's} Journal, Bd.\,99 u.\,103. The algebraic side of the question is treated in two memoirs: \gesperrt{Nöther}, Math. Annalen, Bd.\,15 and \gesperrt{Weber}, ebend., Bd.\,23.}. Now let $f(x_1, x_2, x_3)$, or more briefly $f(x)$, be a ternary form of degree 4 which, set equal to zero, represents a curve of the fourth order without singular point, which we call the curve $f$.

If in this equation
\weq{1}{f(x_1, x_2, x_3) = 0}
we put in place of the variables $x_1, x_2, x_3$ expressions of the form

\[
x_1+t y_1, \quad x_2+t y_2, \quad x_3+t y_3,
\]

then there results an equation of degree 4 with respect to $t$:
\weq{2}{f(x+ty) = 0,}
whose roots, if $(x)$ and $(y)$ are two fixed points, when substituted in $x+ty$, give the coordinates of the points of intersection of the connecting line of the points $(x), (y)$, which we shall call the line $(x,y)$, with the curve $f$.

Now let the function $f(x+ty)$ be ordered by powers of $t$. This gives (by Bd.\ I, \S.\,60):
\weq{3}{f(x+yt) = P_0(x,y)+4tP_1(x,y)+6t^2P_2(x,y)+4t^3P_3(x,y)+t^4P_4(x,y),}
in which the $P_\nu(x,y)$ are the polars of $f(x)$, and hence homogeneous
%CONT>>


% ===== scan leaf 374, printed 353 =====
% !PRINTED 353
\weberpage{353}
%<<CONT
functions of order $\nu$ with respect to $y$, and of order $(4-\nu)$ with respect to $x$. Namely, we have

\weq{4}{
\begin{aligned}
P_0(x,y) &= f(x) \\
P_1(x,y) &= \tfrac{1}{4}\sum_i y_i f'(x_i) \\
P_2(x,y) &= \tfrac{1}{12}\sum_i y_i y_k f''(x_i x_k) \\
P_3(x,y) &= \tfrac{1}{4}\sum x_i f'(y_i) \\
P_4(x,y) &= f(y).
\end{aligned}
}

If $x$ lies on the curve $f$, then $P_0=0$, and one root of equation (2) vanishes. If we further take $(y)$ on the tangent at the point $(x)$, then $P_1(x,y)=0$ as well, and two roots of (2) vanish. The two remaining ones are determined, by (3), through the quadratic equation
\weq{5}{6P_2+4tP_3+t^2P_4=0}
If this equation has two equal roots, the line $(x,y)$ will have yet a second point of contact with the curve $f$, i.\,e. it will be a double tangent. But the condition for this is that the discriminant of equation (5) vanish, or that
\weq{6}{R(x,y) = 2P_3^2 - 3P_4P_2}
be equal to zero. Here $R(x,y)$ is a function homogeneous with respect to each of the two systems $(x)$ and $(y)$, and indeed of degree $2$ for the $x$, of degree $6$ for the $y$. But it is still to be remarked that the equation $R=0$ is also satisfied when $(x)$ is an arbitrary point of the curve, and $(y)$ coincides with $(x)$, because then

\[
P_2(x,x)=P_3(x,x)=P_4(x,x)=f(x)=0
\]

holds. Otherwise, however, $R$ will vanish only when the line $(x,y)$ is a double tangent. We therefore set up the theorem:

1. \gesperrt{The equation $R(x,y)=0$ is, when $(x)$ is a point of the curve $f$ and $(y)$ is a point of the tangent at $(x)$ different from $(x)$, the necessary and sufficient condition for $(x)$ to be a point of contact of a double tangent.}

It is now a matter of deriving from this condition another which contains only the $x$ alone, and which for no
%CONT>>


% ===== scan leaf 375, printed 354 =====
% !PRINTED 354
\weberpage{354}
%<<CONT
points other than the points of contact of the double tangents
is satisfied.

The variables $y$ satisfy the tangent equation

\weq{7}{y_1 f'(x_1) + y_2 f'(x_2) + y_3 f'(x_3) = 0.}

If we denote by $b_1, b_2, b_3$ any three arbitrary constants, and set

\weq{8}{b_x = b_1 x_1 + b_2 x_2 + b_3 x_3,}

so that $b_x = 0$ is the equation of an arbitrary straight line, then we may take, in addition to (7), the equation

\weq{9}{b_1 y_1 + b_2 y_2 + b_3 y_3 = 0}

as well, and thus regard $(y)$ as the point of intersection of the tangent at $(x)$ with the arbitrary straight line $b_x$. We then obtain

\weq{10}{\begin{aligned}
y_1 &= b_2 f'(x_3) - b_3 f'(x_2),\\
y_2 &= b_3 f'(x_1) - b_1 f'(x_3),\\
y_3 &= b_1 f'(x_2) - b_2 f'(x_1),
\end{aligned}}

and hereby the condition (7) is identically satisfied. By this substitution $R(x,y)$ passes into a homogeneous function of degree $20$ in the $x$ and of degree $6$ in the $b$, which we shall denote by $D(x,b)$. But this function $D(x,b)$ vanishes not only at the points of contact of the double tangents, but also at the four points of intersection of the line $b_x$ with the curve $f$. This condition must now be transformed further, so that it becomes independent of the $b$.

If we pass from the point $(y)$ to any other point of the tangent, we can achieve this by replacing $y$ by $\mu x + \lambda y$, where $\lambda, \mu$ denote two arbitrary parameters. Thereby $f(x+ty)$ passes into

\[
f[(1+\mu t)x + \lambda t y] = (1+\mu t)^4 f\left(x + \frac{\lambda t}{1+\mu t}\, y\right)
\]

and the development (3) yields, if we set

\[
P_0(x,y),\ P_1(x,y),\ P_1(x,\mu x + \lambda y)
\]

equal to zero:

\[
\begin{aligned}
&6t^2(1+\mu t)^2\lambda^2 P_2(x,y) + 4t^3(1+\mu t)\lambda^3 P_3(x,y) + t^4\lambda^4 P_4(x,y)\\
&= 6t^2 P_2(x,\mu x + \lambda y) + 4t^3 P_3(x,\mu x + \lambda y) + t^4 P_4(x,\mu x + \lambda y),
\end{aligned}
\]
%CONT>>


% ===== scan leaf 376, printed 355 =====
% !PRINTED 355
\weberpage{355}
%<<CONT
thus, if we arrange both sides according to powers of $t$:

\[
\begin{aligned}
P_2(x,\mu x + \lambda y) &= \lambda^2 P_2(x,y)\\
P_3(x,\mu x + \lambda y) &= 3\lambda^2\mu\, P_2(x,y) + \lambda^3 P_3(x,y)\\
P_4(x,\mu x + \lambda y) &= 6\lambda^2\mu^2 P_2(x,y) + 4\lambda^3\mu\, P_3(x,y) + \lambda^4 P_4(x,y),
\end{aligned}
\]

and from this one obtains, by (6),

\weq{11}{R(x,\mu x+\lambda y) = \lambda^6 R(x,y).}

But this formula does not hold identically; it holds only under the assumption that $(x)$ is a point of the curve, that is, that $f(x) = 0$.

The point $(\mu x+\lambda y)$ can be an entirely arbitrary point of the tangent at $(x)$, and hence, if $a_1, a_2, a_3$ are three arbitrary quantities, we can so dispose of $\mu, \lambda$ that

\weq{12}{\begin{aligned}
\mu x_1 + \lambda y_1 &= a_2 f'(x_3) - a_3 f'(x_2)\\
\mu x_2 + \lambda y_2 &= a_3 f'(x_1) - a_1 f'(x_3)\\
\mu x_3 + \lambda y_3 &= a_1 f'(x_2) - a_2 f'(x_1)
\end{aligned}}

holds. Then $(\mu x+\lambda y)$ is the point of intersection of the tangent with the line $a_x = 0$, if

\weq{13}{a_x = a_1 x_1 + a_2 x_2 + a_3 x_3}

is set. By this, $R(x,\mu x+\lambda y)$ passes over into $D(x,a)$.

Now, in order to determine $\lambda$ and $\mu$, we multiply the equations (10) and (12) by $a_1, a_2, a_3$, then (12) by $b_1, b_2, b_3$, and add in each case. In view of $b_y = 0$ we thereby obtain:

\[
\begin{aligned}
a_y &= -\mu b_x = \Sigma \pm a_1 b_2 f'(x_3)\\
\mu a_x + \lambda a_y &= 0,
\end{aligned}
\]

and from this

\[
\lambda b_x = a_x.
\]

According to this, equation (11) can be represented as follows:

\weq{14}{\frac{D(x,a)}{a_x^6} = \frac{D(x,b)}{b_x^6},}

and in this form it shows that the quotient $D(x,a) : a_x^6$ is independent of the arbitrary quantities $a$. But equation (14) is not an identity; it is satisfied only so long as $f(x) = 0$. We can, however, derive an identity from it if we assume that $f(x)$ is irreducible (or at least contains no multiply counting factor). For the form of degree $26$

\[
b_x^6\, D(x,a) - a_x^6\, D(x,b)
\]
%CONT>>


% ===== scan leaf 377, printed 356 =====
% !PRINTED 356
\weberpage{356}
%<<CONT
equal to zero for all values of $(x)$ satisfying the equation $f(x)=0$, and consequently it must be divisible by $f(x)$. We can therefore, if we denote by $\Phi(x,a,b)$ a form of degree $22$, set:

\weq{15}{b_x^6\, D(x,a) - a_x^6\, D(x,b) = f(x)\, \Phi(x,a,b).}

Let us imagine for a moment that a new coordinate system is introduced in (15), in which the lines $a_x, b_x$ are two axes, of which we assume that they do not intersect on the curve $f$; then it follows from the comparison of the right and left sides of the identical equation (15) that no term can occur in $\Phi$ that is not multiplied either by $a_x^6$ or by $b_x^6$, and consequently $\Phi(x)$ has the form

\[
\Phi(x,a,b) = a_x^6\, \Phi_2(x) - b_x^6\, \Phi_1(x),
\]

in which $\Phi_1, \Phi_2$ are forms of order $16$. Accordingly the identical equation (15) takes the form

\[
b_x^6\,[D(x,a) + f(x)\,\Phi_1(x)] = a_x^6\,[D(x,b) + f(x)\,\Phi_2(x)],
\]

and it shows that $D(x,a) + f(x)\,\Phi_1(x)$ is divisible by $a_x^6$. There exists, therefore, a form of degree $14$, $\chi(x,a,b)$, such that

\weq{16}{\frac{D(x,a)}{a_x^6} = \chi(x,a,b) - f(x)\,\frac{\Phi_1(x,a,b)}{a_x^6}.}

If we here substitute for the $a$ and $b$ any special, e.g. rational, values $c, d$, and denote $\chi(x,c,d)$, which then depends only on $x$ and on the coefficients of the form $f(x)$, by $\chi(x)$, then it follows:

\[
\frac{D(x,c)}{c_x^6} = \chi(x) - f(x)\,\frac{\Phi_1(x,c,d)}{c_x^6},
\]

and if we subtract this from (16) and use the relation (15) for $b=c$:

\weq{17}{\chi(x,a,b) - \chi(x) = f(x)\,\Psi(x),}

in which $\Psi(x)$ is a function that in any case can contain no other denominator than a product of powers of $a_x$ and $c_x$. But since $f(x)$ is not divisible by $a_x$ and $c_x$, $\Psi(x)$ must be an integral function, and (16) yields, if we denote by $\Theta(x)$ a new form of $x$ (of degree 16):

\weq{18}{\frac{D(x,a) - f(x)\,\Theta(x)}{a_x^6} = \chi(x).}


% ===== scan leaf 378, printed 357 =====
% !PRINTED 357
\weberpage{357}
\pdfbookmark[2]{§ 96. The Steiner complexes}{bm111}
This equation shows, however, that the curve of degree $14$:

\weq{19}{\chi(x) = 0,}

which is entirely independent of the $a$, passes through the points of contact of the double tangents, but through no other point of the curve $f(x)$.

The function $\chi(x)$ depends rationally on the coefficients of the equation $f(x)$. There are infinitely many different such curves, among which covariants of $f(x)$ are also found. As \gesperrt{Hesse} has shown, they can be expressed in a simple way by determinants.

Here we draw from this the conclusion:

2.\ \gesperrt{A curve of the fourth order without a singular point has 28 double tangents.}

But one objection to this conclusion must still be met. It would be conceivable that the curve $\chi$ touches the curve $f$, or that the curve $f$ passes through singular points of the curve $\chi$. Then the number of points of intersection, and possibly also the number of double tangents, would be diminished, so that our mode of inference really teaches only that a curve of the fourth order cannot have more than 28 double tangents. But this objection will dispose of itself through the following considerations, in that we shall become acquainted with forms of the equation of the curve in which the existence of 28 distinct double tangents is evident.

\begin{center}
\S.\,96.

\gesperrt{The Steiner Complexes.}
\end{center}

If $x_1 = 0$ is the equation of any double tangent of the curve of the fourth order $f = 0$, which we call for short the double tangent $x_1$, then, when one sets $x_1 = 0$, $f$ must pass into a square. From this it results that $f$ must be of the form:

\weq{1}{f = x_1 V - u^2,}

in which $V$ is a cubic, $u$ a quadratic form. But the function $f$ can be given the shape (1) in a triply infinite number of ways. For if $p$ denotes an arbitrary
%CONT>>


% ===== scan leaf 379, printed 358 =====
% !PRINTED 358
\weberpage{358}
%<<CONT
linear function containing three arbitrary constants, it
follows from (1):
\[
f = x_1 (V + 2pu + x_1p^2) - (u + px_1)^2,
\]
which is again of the form (1).

If now $y_1 = 0$ is the equation of a second double tangent distinct from $x_1$, then we can choose the constants in $p$ (and indeed in infinitely many different ways) so that the conic section $u + px_1 = 0$ passes through the two points of contact of $y_1$, and we may therefore assume that the function $u$ has already been so chosen in the form (1). Then the cubic form $V$ must also vanish at the same points.

But still more: since the line $y_1 = 0$ is to be a double tangent, it must be the case that, if we introduce $x_1, y_1$ and any third linear function $z$ independent of them as variables, at the points of contact
\[
f'(x_1) = 0, \quad f'(z) = 0
\]
and from this it follows by (1):
\[
V'(x_1) = 0, \quad V'(z) = 0,
\]
that is, the curve of the third order $V = 0$ is touched by the line $y_1 = 0$ at the points of contact with $f$, and this is possible only if $V$ splits and contains the function $y_1$ as a divisor. Accordingly we obtain a new shape of the function $f$:
\weq{2}{f = x_1 y_1 v - u^2,}
in which $u, v$ are functions of the second degree.

If $x_1, y_1$ are any arbitrary linear forms, and $u, v$ quadratic forms of three variables, then the function $f$ determined by (2), set equal to zero, represents a curve of the fourth order, of which $x_1$ and $y_1$ are two double tangents.

If conversely $x_1$ and $y_1$ are two arbitrary double tangents of a curve of the fourth order $f = 0$, then $f$ can be brought to the form (2).

Let us imagine the coefficients of $x_1$ and $y_1$ as variable, and let these quantities vary in such a way that $x_1$ and $y_1$ intersect on the curve $f$; then a double point will occur.

The double tangents degenerate into two tangents issuing from the double point.


% ===== scan leaf 380, printed 359 =====
% !PRINTED 359
\weberpage{359}
If the double tangents fall together into one line, two double points appear.

On the other hand, without singular points arising, the two points of contact of a double tangent may very well fall together, and the double tangent thus pass over into a tangent with four-point contact. This occurs when the line $x_1$ touches the conic $u$.

The form (2), too, can still be produced in infinitely many ways while retaining $x_1, y_1$.

To see this, let us assume that
\[
f = x_1 y_1 v - u^2 = x_1 y_1 v_1 - u_1^2,
\]
from which the identical equation follows:
\weq{3}{x_1 y_1 (v - v_1) = (u - u_1)(u + u_1).}
Now if $u - u_1$ were divisible by $x_1$, and $u + u_1$ by $y_1$, then at the point of intersection of $x_1$ and $y_1$, $u$, and consequently $f$, would also vanish. This point of intersection would lie on the curve $f$, which, so long as the curve has no singular point, is not possible. Hence one of the two factors $u - u_1$, $u + u_1$ in (3) must be divisible by $x_1 y_1$. Since the sign of $u_1$ is not yet determined, we may assume that $u - u_1$ is divisible by $x_1 y_1$, and therefore identical with $x_1 y_1$ up to a constant factor. Then, if $\lambda$ is such a factor, we have
\[
u_1 = u + \lambda x_1 y_1,
\]
and the function $f$ takes the form:
\weq{4}{f = x_1 y_1 (v + 2\lambda u + \lambda^2 x_1 y_1) - (u + \lambda x_1 y_1)^2.}
Conversely, for every arbitrary $\lambda$ the formulae (2) and (4) are identical with one another.

Now if we can determine $\lambda$ so that the quadratic function $v + 2\lambda u + \lambda^2 x_1 y_1$ splits into two linear factors, then, setting $u + \lambda x_1 y_1 = u_1$, $f$ takes by (4) the shape
\weq{5}{f = x_1 y_1 x_2 y_2 - u_1^2,}
and this form shows that $x_2, y_2$ are \gesperrt{two further double tangents}, and that the eight points of contact of the double tangents $x_1, y_1, x_2, y_2$ lie on a conic $u_1 = 0$.


% ===== scan leaf 381, printed 360 =====
% !PRINTED 360
\weberpage{360}
Now the necessary and sufficient condition for a ternary quadratic form to split into two linear factors is that the determinant of the quadratic form vanish (Bd. I, \S.\,56).

If we express the forms $u, v$ in terms of $x_1, y_1$ and some third variable $y$, and denote the coefficients in $v$ and $u$ by $a_{ik}, b_{ik}$, so that $2a_{23}, 2b_{23}$ become the coefficients of $x_1 y_1$, we obtain the condition for the splitting of $v + 2\lambda u + \lambda^2 x_1 y_1$ in the form:
\weq{6}{\begin{vmatrix} a_{11}+2\lambda b_{11}, & a_{12}+2\lambda b_{12}, & a_{13}+2\lambda b_{13} \\ a_{21}+2\lambda b_{21}, & a_{22}+2\lambda b_{22}, & a_{23}+2\lambda b_{23}+\tfrac{1}{2}\lambda^2 \\ a_{31}+2\lambda b_{31}, & a_{32}+2\lambda b_{32}+\tfrac{1}{2}\lambda^2, & a_{33}+2\lambda b_{33} \end{vmatrix} = 0,}
and this is an equation of degree $5$ with respect to $\lambda$, which therefore gives us five decompositions:
\[
x_2 y_2,\ x_3 y_3,\ x_4 y_4,\ x_5 y_5,\ x_6 y_6.
\]
The coefficients in $f$ can, as (6) shows, be so disposed that all the functions $x_i y_i$ thus determined are different from one another, and from this it follows, as shown above, that they remain different so long as the curve $f$ has no singular points. Hence the following theorem holds:

1. \gesperrt{To every pair of double tangents $x_1 y_1$ there belong five further pairs $x_i y_i$ of such a kind that the eight points of contact of $x_1 y_1 x_i y_i$ lie on a conic section.}

The six pairs of double tangents that are determined in this way:
\[
x_1 y_1,\ x_2 y_2,\ x_3 y_3,\ x_4 y_4,\ x_5 y_5,\ x_6 y_6,
\]
we shall call a \gesperrt{Steiner complex}\footnote{Hitherto the expression \gesperrt{Steiner group} has commonly been used for this. But since the word ``group'' has an entirely definite other meaning in algebra, we prefer to avoid this expression here.}.

Let us consider the pairs $x_1 y_1, x_2 y_2, x_3 y_3$ of such a complex, and set
\[
f = x_1 y_1 x_2 y_2 - u_1^2 = x_1 y_1 x_3 y_3 - u_2^2,
\]
%CONT>>


% ===== scan leaf 382, printed 361 =====
% !PRINTED 361
\weberpage{361}
%<<CONT
then there follows from this the identity
\[
x_1 y_1 (x_2 y_2 - x_3 y_3) = (u_1 - u_2)(u_1 + u_2),
\]
and from this, as above,
\[
\begin{aligned}
u_1 - u_2 &= h x_1 y_1\\
u_1 + u_2 &= \frac{x_2 y_2 - x_3 y_3}{h},
\end{aligned}
\]
in which $h$ denotes a constant; hence
\[
2u_1 = h x_1 y_1 + \frac{x_2 y_2 - x_3 y_3}{h},
\]
and accordingly
\weq{7}{4f = 4 x_1 y_1 x_2 y_2 - \left(h x_1 y_1 + \frac{x_2 y_2 - x_3 y_3}{h}\right)^2.}
We can give this form of the equation a more elegant shape if we replace
\[
h x_1,\quad \frac{x_2}{h},\quad \frac{x_3}{h}
\]
by
\[
x_1,\quad x_2,\quad x_3
\]
respectively. Then the new $x_1, x_2, x_3$, set equal to zero, denote the same lines as the original ones, since the two indeed differ only by a constant factor, and from (7) there results:
\weq{8}{\begin{aligned}
-4f &= (x_1 y_1 + x_2 y_2 - x_3 y_3)^2 - 4 x_1 y_1 x_2 y_2\\
&= x_1^2 y_1^2 + x_2^2 y_2^2 + x_3^2 y_3^2\\
&\quad - 2 x_2 y_2 x_3 y_3 - 2 x_1 y_1 x_2 y_2 - 2 x_1 y_1 x_3 y_3,
\end{aligned}}
and the equation $f = 0$ can also be represented in the elegant irrational form
\[
x_1 y_1 + x_2 y_2 - x_3 y_3 = -2\sqrt{x_1 y_1 x_2 y_2}
\]
or
\weq{9}{\sqrt{x_1 y_1} + \sqrt{x_2 y_2} + \sqrt{x_3 y_3} = 0}
From (9) we can again derive the form (8) backwards. But since (9) is entirely symmetric, we can interchange the three pairs, and obtain, e.g., also
\[
4 x_2 y_2 x_3 y_3 = (-x_1 y_1 + x_2 y_2 + x_3 y_3)^2,
\]
from which it is to be seen that the points of contact of $x_2 y_2 x_3 y_3$ also lie on a conic section, and that in the complex which one obtains from $x_2 y_2$, not only the pair $x_1 y_1$, but
%CONT>>


% ===== scan leaf 383, printed 362 =====
% !PRINTED 362
\weberpage{362}
%<<CONT
but also the pair $x_3 y_3$, and consequently all the pairs $x_i y_i$, must occur, so that this complex is not different at all from the one derived from $x_1 y_1$. We therefore have the theorem:

\gesperrt{2. The pairs of a \gesperrt{Steiner} complex have the property that the eight points of contact of any two of these pairs lie on a conic, and that one always obtains the same complex, from whichever of the six pairs one may start.}

From this it follows that three double tangents of a \gesperrt{Steiner} complex, such as $x_1, y_1, x_2$, of which two form a pair of the complex, have their six points of contact on a conic. On the other hand, there are again systems of three double tangents (triples) whose six points of contact do not lie on a conic.

According to a designation introduced by \gesperrt{Frobenius}, three double tangents, such as $x_1, y_1, x_2$, whose points of contact lie on a conic, form a \gesperrt{syzygetic triple}. Three double tangents whose six points of contact do not lie on a conic form an \gesperrt{azygetic triple}. Correspondingly, we shall call four double tangents whose eight points of contact lie on a conic a \gesperrt{syzygetic quadruple}, and any system of double tangents of which every three are azygetic an \gesperrt{azygetic system}.

Here now the following important theorem holds:

\gesperrt{3. Three double tangents which occur in a \gesperrt{Steiner} complex in such a way that no two of them form a pair, such as $x_1, x_2, x_3$, are always azygetic.}

The proof of this theorem follows simply from the form of the equation (9). From it one sees first of all that the three lines $x_1, x_2, x_3$ do not intersect in one point; for such a point of intersection would lie on the curve $f$ and would therefore be a singular point. We can therefore introduce $x_1, x_2, x_3$ as coordinates and accordingly
\[
\begin{gathered}
y_1 = \alpha_1 x_1 + \alpha_2 x_2 + \alpha_3 x_3\\
y_2 = \beta_1 x_1 + \beta_2 x_2 + \beta_3 x_3\\
y_3 = \gamma_1 x_1 + \gamma_2 x_2 + \gamma_3 x_3
\end{gathered}
\]
%CONT>>


% ===== scan leaf 384, printed 363 =====
% !PRINTED 363
\weberpage{363}
\pdfbookmark[2]{§ 97. Pairs and triples of complexes}{bm112}
%<<CONT
put. In this none of the three constants $\alpha_1, \beta_2, \gamma_3$ can vanish. For if, e.g., $\alpha_1 = 0$, then by (9) the lines $x_2, x_3, y_1$ intersect on the curve $f$, and this point of intersection would have to be a singular point.

By equation (9), however, the coordinates of the points of contact of $x_1, x_2, x_3$ result from the following three pairs of equations, of which each pair gives two points of contact:
\[
\begin{gathered}
x_1 = 0,\quad x_2(\beta_2 x_2 + \beta_3 x_3) - x_3(\gamma_2 x_2 + \gamma_3 x_3) = 0\\
x_2 = 0,\quad x_3(\gamma_3 x_3 + \gamma_1 x_1) - x_1(\alpha_3 x_3 + \alpha_1 x_1) = 0\\
x_3 = 0,\quad x_1(\alpha_1 x_1 + \alpha_2 x_2) - x_2(\beta_1 x_1 + \beta_2 x_2) = 0.
\end{gathered}
\]
If now these six points are to lie on a conic $\varphi = 0$, then $\varphi$, apart from a constant factor $h$, must pass over, for $x_1 = 0$, into the left part of the second equation of the first pair $x_2(\beta_2 x_2 + \beta_3 x_3) - x_3(\gamma_2 x_2 + \gamma_3 x_3)$. If we therefore denote by $a_1, a_2, a_3$ the coefficients of $x_1^2, x_2^2, x_3^2$ in $\varphi$, then we must have
\weq{10}{
\begin{gathered}
a_2 = h_1 \beta_2,\quad a_3 = -h_1 \gamma_3\\
a_3 = h_2 \gamma_3,\quad a_1 = -h_2 \alpha_1\\
a_1 = h_3 \alpha_1,\quad a_2 = -h_3 \beta_2
\end{gathered}
}
and here $h_1, h_2, h_3$ are three constants.

But from (10) it follows that
\[
h_2 = -h_3,\quad h_3 = -h_1,\quad h_1 = -h_2,
\]
and this would only be possible if $h_1 = h_2 = h_3 = 0$.
But this is the case only when $a_1, a_2, a_3$ vanish, that is, when the conic $\varphi$ passes through the points of intersection of the three straight lines $x_1, x_2, x_3$. But it is to pass through the points of contact of these three double tangents, which are certainly different from their three points of intersection. Thus our assumption is shown to be impossible, and Theorem 3. is proved.

\begin{center}
§. 97.

\gesperrt{Complex pairs and complex triples.}
\end{center}

The theorems of the preceding section give an excellent means of exploring the manifold geometric and algebraic relations of the double tangents and of
%CONT>>


% ===== scan leaf 385, printed 364 =====
% !PRINTED 364
\weberpage{364}
%<<CONT
to present. We confine ourselves here to what is required for the attainment of our chief aim, namely the determination of the \gesperrt{Galois} group of the problem.

We start from an arbitrarily selected \gesperrt{Steiner} complex:
\weq{1}{x_1 y_1,\ x_2 y_2,\ x_3 y_3,\ x_4 y_4,\ x_5 y_5,\ x_6 y_6.}
The form of the equation \S.\,96, (5) then shows that in the complex determined by the pair $x_1 x_2$ the pair $y_1 y_2$, and in the complex determined by $x_1 y_2$ the pair $x_2 y_1$, must occur.

We therefore consider, besides (1), the two further complexes:
\weq{2}{x_1 x_2,\quad y_1 y_2,\quad \ldots}
\weq{3}{x_1 y_2,\quad x_2 y_1,\quad \ldots}
Each of the two complexes (2), (3) contains, besides those already known, eight further double tangents, and these must all be different from the $x_i, y_i$, because, if e.g. $x_3$ occurred in (2), then $x_1 x_2 x_3$ would be syzygetic, which contradicts theorem 3., (\S.\,96). Likewise the two complexes (2) and (3) can contain no common double tangent besides $x_1, x_2, y_1, y_2$. For if, say, $z$ occurred in both, then $x_1 x_2 z$ would be syzygetic by (2) and azygetic by (3), which is a contradiction. From this it follows:

\gesperrt{4. In the three complexes (1), (2), (3) taken together all 28 double tangents occur.}

From this it follows that every double tangent which forms a syzygetic triple with any pair $x_1 y_1$ must occur in the complex $x_1 y_1$; that therefore every syzygetic triple is completed by a determinate further double tangent into a syzygetic quadruple; and that consequently a conic which passes through the points of contact of three double tangents cuts the curve $f$ in the points of contact of a fourth double tangent.

Two pairs of a complex always form a syzygetic quadruple, and however one may divide such a quadruple into two pairs, both pairs always belong to the same complex.

Since from 28 things one can form 14.27 pairs, since every pair of double tangents determines a complex, and in every
%CONT>>


% ===== scan leaf 386, printed 365 =====
% !PRINTED 365
\weberpage{365}
%<<CONT
each complex, the total number of the complexes is $14.27 : 6 = 63$.

\gesperrt{5. There are 63 \gesperrt{Steiner} complexes.}

If from the pairs of the complex (1) we pick out any other pair of pairs instead of $x_1 y_1, x_2 y_2$, then from it, following the scheme of (2) and (3), we can each time form two new complexes. Since there are 15 such pairs of pairs, we obtain 30 new complexes of the type (2), (3), which are all derived in the same way from the complex (1), and which are all different from one another.

If we now take any one of the double tangents $z_1$ newly entering in (2) and (3), we obtain two new complexes by starting from the two pairs $x_1 z_1, y_1 z_1$, and since we can choose $z_1$ in 16 different ways, there result in this way 32 new complexes, with which the total number of all the complexes is exhausted. But the distribution of the double tangents among the complexes $x_1 z_1, y_1 z_1$ must still be investigated more closely.

Without restricting generality, since if necessary we can interchange $x_1$ with $y_1$, we can make the assumption that $z_1$ occurs in the complex (2), and there in the pair $z_1 z_2$.

Then we obtain the two complexes:
\weq{4}{x_1 z_1,\quad x_2 z_2,\quad \ldots}
\weq{5}{y_1 z_1,\quad y_2 z_2,\quad \ldots}
We now consider the three complexes:
\weq{4}{x_1 z_1,\quad x_2 z_2,\quad \ldots}
\weq{2}{x_1 x_2,\quad z_1 z_2,\quad \ldots}
\weq{4a}{x_1 z_2,\quad x_2 z_1,\quad \ldots,}
which by Theorem 4. must contain all the double tangents, hence among them also $x_3, y_3$. But these do not occur in (2), and likewise they cannot both occur in (4) or both in (4a), since otherwise $x_1, x_3, y_3$ would have to be azygetic, whereas [by (1)] they are syzygetic. Since if need be we may interchange $x_3$ and $y_3$ in the notation, we can assume that $x_3$ occurs in (4), and indeed in a pair $x_3 z_3$.

Then, since $z_1 z_2 z_3$ are azygetic, $z_3$ cannot occur in the complex (2) and must consequently be contained in (3).


% ===== scan leaf 387, printed 366 =====
% !PRINTED 366
\weberpage{366}
Now we have the two complexes:

\weq{4}{x_2 z_2,\ x_3 z_3,\ \dots}
\weq{4\text{b}}{x_2 x_3,\ z_2 z_3,\ \dots{},}

and since $x_2 x_3 y_2 y_3$ are a syzygetic quadruple, the group (4 b) also contains the pair $y_2 y_3$, and consequently $y_2 y_3 z_2 z_3$ are also a syzygetic quadruple. From this it follows further that $y_2 z_2$ and $y_3 z_3$ belong to the same complex. Since the same consideration as for $x_3 y_3$ can also be carried out for the pairs $x_4 y_4, x_5 y_5, x_6 y_6$, the complexes (4), (5) can accordingly be formed completely, and they take the expression:

\weq{4}{x_1 z_1,\ x_2 z_2,\ x_3 z_3,\ x_4 z_4,\ x_5 z_5,\ x_6 z_6}
\weq{5}{y_1 z_1,\ y_2 z_2,\ y_3 z_3,\ y_4 z_4,\ y_5 z_5,\ y_6 z_6.}

Here $z_1 z_2$ form a pair of the complex (2), and $z_3, z_4, z_5, z_6$, which form an azygetic system, all occur in the complex (3), in which no two are paired.

Since, as shown above, we can derive from the arbitrarily assumed complex (1), after the type of (2), (3), (4), (5), all complexes that exist at all, the theorem results:

6. \gesperrt{Any two complexes have either a syzygetic quadruple or an azygetic system of six double tangents in common.}

Two complexes that have a syzygetic quadruple in common we shall call a \gesperrt{syzygetic complex pair}. To every such pair there is a third complex which contains the same quadruple. Three such complexes we call a \gesperrt{syzygetic complex triple} [e.\,g. (1), (2), (3)].

Likewise we call two complexes of the second kind, that is, such as have six azygetic elements in common, an \gesperrt{azygetic complex pair}.

Every azygetic complex pair is likewise completed to a triple by a determinate complex which forms an azygetic pair with each of the two complexes [as (1), (4), (5)]. Such a one we call an \gesperrt{azygetic complex triple}.

In a syzygetic triple all 28 double tangents occur, in an azygetic one only 18.


% ===== scan leaf 388, printed 367 =====
% !PRINTED 367
\weberpage{367}
\pdfbookmark[2]{§ 98. The Aronhold systems of seven}{bm113}
\S.\,98.

\gesperrt{The Aronhold Seven-Systems.}

Of particular importance are the azygetic systems of seven double tangents, which were first considered by \gesperrt{Aronhold}, and which are therefore called \gesperrt{Aronhold seven-systems}. We also call them \gesperrt{complete seven-systems}, or briefly \gesperrt{complete systems}.

That such systems exist is shown by the compilations of the preceding section. For if

\weq{1}{\begin{aligned}
&x_1 y_1,\ x_2 y_2,\ x_3 y_3,\ x_4 y_4,\ x_5 y_5,\ x_6 y_6\\
&x_1 z_1,\ x_2 z_2,\ x_3 z_3,\ x_4 z_4,\ x_5 z_5,\ x_6 z_6
\end{aligned}}

is an azygetic complex pair, then

\[ x_1,\ x_2,\ x_3,\ x_4,\ x_5,\ y_6,\ z_6 \]

is a complete seven-system (because in the complex $y_6 z_6$ none of the $x$ occurs). It will be shown that no azygetic systems of more than seven elements exist. First of all we have the theorem:

7. \gesperrt{Any six elements of a complete system occur in one and only one complex.}

We first prove the second part of the assertion, that is, we prove that, if

\weq{2}{x_1,\ x_2,\ x_3,\ x_4,\ x_5,\ x_6,\ x_7}

is a complete system, then $x_1, x_2, x_3, x_4, x_5, x_6$ cannot occur in two different complexes. If we suppose this to be possible, then the two complexes must form an azygetic pair like (1). In the syzygetic triple

\[\begin{aligned}
&y_1 z_1,\ y_2 z_2,\ y_3 z_3,\ y_4 z_4,\ y_5 z_5,\ y_6 z_6\\
&y_1 y_2,\ z_1 z_2,\ x_1 x_2\\
&y_1 z_2,\ y_2 z_1
\end{aligned}\]

the double tangents $x_3, x_4, x_5, x_6, x_7$ must occur, and since they occur neither in the first nor in the second of these complexes, they must occur in the third. But since only four pairs still remain in this complex, at
%CONT>>


% ===== scan leaf 389, printed 368 =====
% !PRINTED 368
\weberpage{368}
%<<CONT
least two of the $x_3, \dots x_7$ would form a pair in it. But that is not possible, since this pair would otherwise form a syzygetic triple with one of the remaining $x$.

To prove that the double tangents $x_1, x_2, x_3, x_4, x_5, x_6$ always occur in one complex, we take out two arbitrary ones of them, $x_1 x_2$, and choose another pair $y_1 y_2$ occurring in the complex $x_1 x_2$, which is possible in five ways. From these we form the syzygetic complex triple

\weq{2}{\begin{aligned}
\alpha)\ & x_1 x_2,\ y_1 y_2\\
\beta)\ & x_1 y_1,\ x_2 y_2\\
\gamma)\ & x_1 y_2,\ x_2 y_1,
\end{aligned}}

which represents five different such triples. Now in $\beta$) and $\gamma$) the $x_3, x_4, x_5, x_6, x_7$ must occur, and indeed no two of them paired. We now show first that these five $x$ cannot distribute themselves two and three over the two complexes $\beta$), $\gamma$), but only one and four. For suppose the complexes $\beta$), $\gamma$) were composed thus:

\weq{2}{\begin{aligned}
\beta)\ & x_1 y_1,\ x_2 y_2,\ x_3 y_3,\ x_4 y_4,\ x_5 y_5,\ \cdot\\
\gamma)\ & x_1 y_2,\ x_2 y_1,\ x_6 y_6,\ x_7 y_7,\ \cdot\quad\cdot\text{,}
\end{aligned}}

then we can further consider the complex

\weq{2}{\delta)\quad x_6 x_7,\ y_6 y_7\ \cdots{},}

which together with $\beta$) forms a syzygetic pair, because neither $x_1$ nor $y_1$ occurs in $\delta$), so the pair cannot be azygetic. Hence $\beta$) and $\delta$) must have a syzygetic quadruple in common, and since in two pairs of $\beta$) at least one $x$ different from $x_6, x_7$ occurs, then this $x$ must also occur in $\delta$), which is impossible, because no $x$ is syzygetic with $x_6, x_7$. There thus remains for $\beta$), $\gamma$) only one composition, like the following:

\weq{3}{\begin{aligned}
\beta)\ & x_1 y_1,\ x_2 y_2,\ x_3 y_3,\ x_4 y_4,\ x_5 y_5,\ x_6 y_6\\
\gamma)\ & x_1 y_2,\ x_2 y_1,\ x_7 y_7,\ \cdot\quad\cdot\quad\cdot
\end{aligned}}

That we make precisely this assumption, and do not take the six $x$ into $\gamma$), is no restriction, since we can, if necessary, interchange $y_1$ with $y_2$.

If now, keeping $x_1 x_2$ fixed, we let the other pairs of (2) $\alpha$) take the place of $y_1 y_2$, we obtain
%CONT>>


% ===== scan leaf 390, printed 369 =====
% !PRINTED 369
\weberpage{369}
\pdfbookmark[2]{§ 99. The Hesse-Cayley algorithm for the designation of the double tangents}{bm114}
%<<CONT
five complex formations of the type (3). But two such complex formations can differ only in this, that in place of $x_7$ there occurs each time another of the elements $x_3, x_4, x_5, x_6, x_7$, and all these possibilities must also occur, because otherwise we should obtain two different complexes which contain the same six elements of the complete system of the $x$, which is not possible, as we have proved. Accordingly we may assume that the $x_1, x_2, x_3, x_4, x_5, x_6$ occurring in the complex (3) $\beta$) are \gesperrt{any six elements whatever} of the complete system, which was to be proved.

\S.\,99.

\gesperrt{The Hesse-Cayley algorithm for the notation of the double tangents.}

Theorem 7. of the preceding section leads to a mode of notation for the double tangents by which a very perspicuous representation of all these relations is possible, and which has been developed by \gesperrt{Cayley} (following \gesperrt{Hesse}).

We take as a basis a complete seven-system

\[ x_1,\ x_2,\ x_3,\ x_4,\ x_5,\ x_6,\ x_7 \]

whose elements we denote simply by the figures 1, 2, 3, 4, 5, 6, 7. We separate off an arbitrary double tangent, say $x_1$, of this system, and form the complex which contains the remaining six:

\weq{1}{x_2 y_2,\ x_3 y_3,\ x_4 y_4,\ x_5 y_5,\ x_6 y_6,\ x_7 y_7.}

The double tangent $y_2$ is then completely determined by the chosen $x_1$ and by $x_2$, and can therefore be denoted by $[1\,2]$. In the same way we denote $y_3$ by $[1\,3]$, and so on. The double tangents $[1\,2], [1\,3], \dots, [1\,7]$ are hereby completely determined and different from one another. From this determination it is also evident what is to be understood in general by $[\mu\,\nu]$, when $\mu, \nu$ denote two different figures from the sequence 1 to 7.

It is now first to be shown that $[\mu\,\nu] = [\nu\,\mu]$. It suffices if we prove that $[1\,2] = [2\,1]$. For this we need only form the complex which contains $x_1, x_3, x_4, x_5, x_6, x_7$
%CONT>>


% ===== scan leaf 391, printed 370 =====
% !PRINTED 370
\weberpage{370}
% Source-faithful three-stroke Cayley configuration sign.
\providecommand{\weberthreeup}{\mathord{\tikz[baseline=-0.2ex,x=0.23em,y=0.23em,line width=0.45pt]{\draw (0,0)--(-1,1) (0,0)--(0,1.08) (0,0)--(1,1);}}}
%<<CONT
contains, and which forms an azygetic pair with (1). It must
therefore be of the form:
\weq{2}{x_1 y_2,\ x_3 z_3,\ x_4 z_4,\ x_5 z_5,\ x_6 z_6,\ x_7 z_7,}
and accordingly $y_2$ is also to be denoted by $[2\,1]$, q. e. d.

The $z$ in (2) are different from the $y$ in (1), and since, e.g.,
$z_3 = [2\,3]$, it follows that in general two $[\mu\,\nu]$ which do not
agree in both figures $\mu, \nu$ are different from one another.
Since twenty-one pairs can be formed from seven figures, one
obtains in this way all the double tangents, and each only once.

From this mode of representation it also follows that no
azygetic systems of more than seven double tangents exist;
for if we add to the seven $x$ any further double tangent, for
which, the figures 1 to 7 being equally entitled, we may choose
say $y_2 = [1\,2]$, then $y_2 x_1 x_2$ is a syzygetic triple.

It is expedient to introduce an eighth figure 8, and to
denote the elements of the original system of seven not by the
single figures, but by the pairs
\[
[1\,8],\ [2\,8],\ [3\,8],\ [4\,8],\ [5\,8],\ [6\,8],\ [7\,8]
\]
where it shall then also hold that $[1\,8] = [8\,1]$, and so forth.
Then all 28 double tangents are denoted uniformly by the pairs
$[\mu\,\nu]$, in which $\mu$ and $\nu$ denote two different figures of the
sequence 1 to 8.

Here it is very helpful for clarity to employ an intuitive
notation\footnote{After \gesperrt{Cayley}; cf. \gesperrt{Salmon}, ``Higher plane curves''.}. One indicates
a double tangent $[\mu\,\nu]$ by a single stroke $\mid$, at
whose ends one imagines the two figures $\mu, \nu$ placed.
Then two strokes without a common point, $\mid\mid$, denote two
double tangents in whose designation $[\mu\,\nu]$ no common
figure occurs, and two strokes issuing from one point, $\vee$, two
double tangents which have a common figure in their symbols $[\mu\,\nu]$.
After this the more complicated signs that we employ later are
self-explanatory. For a triple of double tangents we have, e.g., the
following five signs $\mid\mid\mid$, $\sqcap$, $\triangle$, $\weberthreeup$, $\vee\mid$.


% ===== scan leaf 392, printed 371 =====
% !PRINTED 371
\weberpage{371}
% Source-faithful three-stroke Cayley configuration sign.
\providecommand{\weberthreeup}{\mathord{\tikz[baseline=-0.2ex,x=0.23em,y=0.23em,line width=0.45pt]{\draw (0,0)--(-1,1) (0,0)--(0,1.08) (0,0)--(1,1);}}}
It shall now first be investigated how the azygetic and the syzygetic triples
are distinguished in this notation.

Here it is to be observed that, according to the explanation
given up to now, the figure 8 occupies a special place, while
the figures 1 to 7 appear entirely alike and can be permuted
arbitrarily.

From the two complexes (1), (2) we further derive the completion
 to an azygetic complex triple, namely:
\weq{3}{
\begin{aligned}
&x_2 y_2,\ x_3 y_3,\ x_4 y_4,\ x_5 y_5,\ x_6 y_6,\ x_7 y_7,\\
&x_1 y_2,\ x_3 z_3,\ x_4 z_4,\ x_5 z_5,\ x_6 z_6,\ x_7 z_7,\\
&x_1 x_2,\ y_3 z_3,\ y_4 z_4,\ y_5 z_5,\ y_6 z_6,\ y_7 z_7,
\end{aligned}
}
and now go through the individual symbols for the double-tangent
 triples.

We begin with the symbol $\weberthreeup$, for which, in view of
the exceptional position of the figure 8, we have three types to
consider:
\weq{4}{
\begin{aligned}
&[1\,8]\,[2\,8]\,[3\,8] = x_1 x_2 x_3\\
&[1\,5]\,[1\,6]\,[1\,7] = y_5 y_6 y_7\\
&[1\,6]\,[1\,7]\,[1\,8] = y_6 y_7 x_1,
\end{aligned}
}
and a glance at the three complexes (3) shows (by the theorem
\S.\,96, 3.) that all these triples are azygetic.

Secondly we consider the symbol $\vee\mid$, for which four
types are to be taken into account:
\weq{5}{
\begin{aligned}
&[8\,3]\,[8\,4]\,[1\,2] = x_3 x_4 y_2\\
&[8\,4]\,[1\,4]\,[2\,3] = x_4 y_4 z_3\\
&[1\,2]\,[1\,3]\,[4\,8] = y_2 y_3 x_4\\
&[1\,4]\,[1\,5]\,[2\,3] = y_4 y_5 z_3,
\end{aligned}
}
and these triples too are azygetic.

For the symbol $\triangle$ there is to be considered:
\weq{6}{
\begin{aligned}
&[1\,8]\,[2\,8]\,[1\,2] = x_1 x_2 y_2\\
&[1\,2]\,[1\,3]\,[2\,3] = y_2 y_3 z_3,
\end{aligned}
}
which likewise prove to be azygetic, because $y_2$ does not occur in the complex
 (3), which contains the pairs $x_1 x_2$, $y_3 z_3$.

For the symbol $\mid\mid\mid$ there are two types:
\weq{7}{
\begin{aligned}
&[1\,3]\,[2\,4]\,[5\,8] = y_3 z_4 x_5\\
&[1\,3]\,[2\,4]\,[5\,6] = y_3 z_4 [5\,6].
\end{aligned}
}


% ===== scan leaf 393, printed 372 =====
% !PRINTED 372
\weberpage{372}
% Source-faithful three-stroke Cayley configuration sign.
\providecommand{\weberthreeup}{\mathord{\tikz[baseline=-0.2ex,x=0.23em,y=0.23em,line width=0.45pt]{\draw (0,0)--(-1,1) (0,0)--(0,1.08) (0,0)--(1,1);}}}
If we consider the syzygetic complex triple
\weq{8}{
\begin{aligned}
&y_3 z_4,\ y_4 z_3,\ \cdot\ \cdot\ \cdot\ \cdot\\
&y_3 y_4,\ z_3 z_4,\ x_3 x_4\ \cdot\ \cdot\\
&y_3 z_3,\ y_4 z_4,\ \cdot\ \cdot\ \cdot\ \cdot,
\end{aligned}
}
in which all double tangents must occur, we find $[5\,8]$ and
$[5\,6]$ not in the last two complexes of (8),
because
\weq{9}{
\begin{aligned}
&y_3 z_3 [5\,8] = [1\,3]\,[2\,3]\,[5\,8],\quad y_3 y_4 [5\,8] = [1\,3]\,[1\,4]\,[5\,8]\\
&y_3 z_3 [5\,6] = [1\,3]\,[2\,3]\,[5\,6],\quad y_3 y_4 [5\,6] = [1\,3]\,[1\,4]\,[5\,6]
\end{aligned}
}
have the symbol $\vee\mid$ and are therefore azygetic. Consequently
$[5\,8]$ and $[5\,6]$ occur in the first of the complexes (8), and the
two triples (7) are syzygetic.

Finally we have also to consider the symbol $\sqcap$, which
again gives three types:
\weq{10}{
\begin{aligned}
&[1\,2]\,[2\,3]\,[3\,4] = y_2 z_3 [3\,4]\\
&[1\,2]\,[2\,3]\,[3\,8] = y_2 z_3 x_3\\
&[1\,8]\,[8\,2]\,[2\,3] = x_1 x_2 z_3.
\end{aligned}
}
These triples too are syzygetic, which for the last two is to be
seen immediately from the complex triple (3), and for
the first follows from the syzygetic complex triple
\[
\begin{aligned}
&y_2 z_3,\ x_1 x_3,\ \cdot\ \cdot\ \cdot\ \cdot\\
&y_2 x_1,\ z_3 x_3,\ \cdot\ \cdot\ \cdot\ \cdot\\
&y_2 x_3,\ z_3 x_1,\ \cdot\ \cdot\ \cdot\ \cdot,
\end{aligned}
\]
of which the last two do not contain $[3\,4]$, because
$y_2 x_1 [3\,4]$ and $y_2 x_3 [3\,4]$ both have the symbol $\vee\mid$.

Here, however, the exceptional position of the figure 8 has entirely
disappeared, and we arrive at the result:

\gesperrt{Among the triples of double tangents, those with the
symbols}
\[
\mid\mid\mid,\ \sqcap
\]
\gesperrt{are syzygetic, and those with the symbols}
\[
\weberthreeup,\ \triangle,\ \vee\mid
\]
\gesperrt{azygetic.}

From this one very easily obtains the symbols for all the
syzygetic and azygetic quadruples:
%CONT>>


% ===== scan leaf 394, printed 373 =====
% !PRINTED 373
\weberpage{373}
%<<CONT
% Source-faithful Cayley configuration signs: three/four concurrent strokes and a seven-rayed star.
\providecommand{\weberthreeup}{\mathord{\tikz[baseline=-0.2ex,x=0.23em,y=0.23em,line width=0.45pt]{\draw (0,0)--(-1,1) (0,0)--(0,1.08) (0,0)--(1,1);}}}
\providecommand{\weberfourup}{\mathord{\tikz[baseline=-0.2ex,x=0.23em,y=0.23em,line width=0.45pt]{\draw (0,0)--(-1.05,1) (0,0)--(-0.35,1.08) (0,0)--(0.35,1.08) (0,0)--(1.05,1);}}}
\providecommand{\webersevenstar}{\mathord{\tikz[baseline=-0.35ex,x=0.21em,y=0.21em,line width=0.45pt]{\draw (0,0)--(90:1) (0,0)--(38.57:1) (0,0)--(-12.86:1) (0,0)--(-64.29:1) (0,0)--(-115.71:1) (0,0)--(-167.14:1) (0,0)--(141.43:1);}}}
\gesperrt{The quadruples of double tangents with the
signs}
\[
\mid\mid\mid\mid,\ \square
\]
\gesperrt{are syzygetic, and those with the signs}
\[
\weberfourup,\ \triangle\mid,\ \weberthreeup\mid,\ \vee\vee
\]
\gesperrt{azygetic.}

All remaining quadruples are neither syzygetic nor azygetic.

After this it is easy to form the signs for all the complete
seven-systems.

Such a sign must consist of seven strokes, which cannot have
more than eight vertices and which form a figure
from which neither of the two figures $\mid\mid\mid$, $\sqcap$ can be
detached. From this it follows first that this figure can consist of not more than
two separate parts, because otherwise the figure $\mid\mid\mid$
would be contained in it, and that no part can have more than one centre
from which several strokes issue, except when
this part is the triangle $\triangle$, because otherwise the figure $\sqcap$ would
occur.

Now if the figure is in one part, it must be a seven-rayed
star $\webersevenstar$, and since one can choose any of the eight figures as
middle point, these are eight possibilities, of
which one is the assumption considered above:
\[
[1\,8]\,[2\,8]\,[3\,8]\,[4\,8]\,[5\,8]\,[6\,8]\,[7\,8].
\]
But if the figure is in two parts, it is first to be excluded
that the one part consists of one or of two strokes or of
a three-rayed star; for in these cases the other part
would have to be a star with six, five or four rays.
But for that not enough of the eight figures remain.
It therefore only remains that the one part is a
triangle, the other a four-rayed star, $\triangle\,\weberfourup$, and this
assumption is also in fact always admissible. A representative
of such a system is:
\[
[1\,2]\,[1\,3]\,[2\,3]\,[4\,5]\,[4\,6]\,[4\,7]\,[4\,8].
\]
The triangle can be chosen in $8.7.6 : 2.3 = 56$ different
ways, and for every triangle there are still five possibilities
of taking the four-rayed star, since one can make any of the
remaining five points the centre. The number
%CONT>>


% ===== scan leaf 395, printed 374 =====
% !PRINTED 374
\weberpage{374}
\pdfbookmark[2]{§ 100. Rational determination of the curve from a complete system of seven}{bm115}
% Source-faithful Cayley configuration signs: a four-stroke star and a seven-rayed star.
\providecommand{\weberfourup}{\mathord{\tikz[baseline=-0.2ex,x=0.23em,y=0.23em,line width=0.45pt]{\draw (0,0)--(-1.05,1) (0,0)--(-0.35,1.08) (0,0)--(0.35,1.08) (0,0)--(1.05,1);}}}
\providecommand{\webersevenstar}{\mathord{\tikz[baseline=-0.35ex,x=0.21em,y=0.21em,line width=0.45pt]{\draw (0,0)--(90:1) (0,0)--(38.57:1) (0,0)--(-12.86:1) (0,0)--(-64.29:1) (0,0)--(-115.71:1) (0,0)--(-167.14:1) (0,0)--(141.43:1);}}}
%<<CONT
of these determinations is therefore 280, and we arrive at the following result:

\gesperrt{There are in all 288 Aronhold systems, whose symbols have one of the two forms}

\[
\webersevenstar;\ \Delta\,\weberfourup.
\]

From these symbols, however, one must not conclude that these systems of seven split into two different kinds. The difference between the two figures lies only in the notation. In fact, we may start from any one whatever of the complete systems of seven in order to derive the notation.

From this one easily finds the notation for the \gesperrt{Steiner} complexes, which are to be formed after one of the two following types:

\[
[1\,2][3\,4],\ [1\,3][2\,4],\ [1\,4][2\,3],\ [5\,6][7\,8],\ [5\,7][6\,8],\ [5\,8][6\,7],
\]
\[
[1\,7][1\,8],\ [2\,7][2\,8],\ [3\,7][3\,8],\ [4\,7][4\,8],\ [5\,7][5\,8],\ [6\,7][6\,8].
\]

\begin{center}
\S.\,100.

\gesperrt{Rational determination of the curve from a complete system of seven.}
\end{center}

The great significance of the \gesperrt{Aronhold} systems for the problem of the double tangents is expressed in the following two theorems:

I. \gesperrt{If for a curve of the fourth order a complete system of seven is given, then from it all the remaining double tangents can be rationally determined.}

II. \gesperrt{If seven arbitrary straight lines in a plane are given, then in general, i.\,e. when certain rational functions of the coefficients in the equations of these lines do not vanish, one can determine in a rational way the equation of a curve of the fourth order without singular point, for which the given seven lines form an Aronhold system.}

In order to prove the first theorem, it would, in view of the complete equal entitlement of the figures 1 to 7, suffice to
%CONT>>


% ===== scan leaf 396, printed 375 =====
% !PRINTED 375
\weberpage{375}
%<<CONT
determination of an eighth double tangent from a given complete system. But we do better to determine three simultaneously. So let now

\weq{1}{x_1,\ x_2,\ x_3,\ x_4,\ x_5,\ x_6,\ x_7}

be a complete seven-system assumed as known. We imagine the three \gesperrt{Steiner} complexes formed, in which six at a time of the double tangents (1) are contained, with the exclusion first of $x_1$, then of $x_2$, lastly of $x_3$.

The 15 double tangents newly entering there we denote by the letter $\xi$, and obtain these three complexes [by \S.\,99, (1), (2)] in the form:

\weq{2}{\begin{aligned}
x_2\xi_3 &,\ x_3\xi_2,\ x_4\xi_{41},\ x_5\xi_{51},\ x_6\xi_{61},\ x_7\xi_{71}\\
x_3\xi_1 &,\ x_1\xi_3,\ x_4\xi_{42},\ x_5\xi_{52},\ x_6\xi_{62},\ x_7\xi_{72}\\
x_1\xi_2 &,\ x_2\xi_1,\ x_4\xi_{43},\ x_5\xi_{53},\ x_6\xi_{63},\ x_7\xi_{73},
\end{aligned}}

and by the notation of \S.\,99 we have

\[
\xi_1 = [2\,3],\ \xi_2 = [3\,1],\ \xi_3 = [1\,2],\ \xi_{41} = [1\,4],\ \ldots
\]

From (2) there results yet another \gesperrt{Steiner} complex, which forms a syzygetic pair with each of the complexes (2):

\weq{3}{x_1\xi_1,\ x_2\xi_2,\ x_3\xi_3,\ \ldots}

and this complex contains none of the double tangents $x_4, x_5, x_6, x_7$.

If now we take the functions sought, $\xi_1, \xi_2, \xi_3$, multiplied by the suitable constant factors, we can put the equation of the curve of the fourth order, by \S.\,96, (9), into the form

\weq{4}{\sqrt{x_1\xi_1} + \sqrt{x_2\xi_2} + \sqrt{x_3\xi_3} = 0}

and thus take the rational function $f$, which, set equal to zero, gives the equation of the curve, in each of the three mutually identical forms

\weq{5}{f = 4x_2\xi_2x_3\xi_3 - u_1^2 = 4x_3\xi_3x_1\xi_1 - u_2^2 = 4x_1\xi_1x_2\xi_2 - u_3^2,}

where

\weq{6}{\begin{aligned}
u_1 &= -x_1\xi_1 + x_2\xi_2 + x_3\xi_3\\
u_2 &= x_1\xi_1 - x_2\xi_2 + x_3\xi_3\\
u_3 &= x_1\xi_1 + x_2\xi_2 - x_3\xi_3
\end{aligned}}

is set. Now [by (2)] $x_2\xi_3\, x_4\xi_{41}$ also form a syzygetic quadruple, and consequently we can, if over a
%CONT>>


% ===== scan leaf 397, printed 376 =====
% !PRINTED 376
\weberpage{376}
%<<CONT
the constant factor in $\xi_{41}$ is suitably disposed of, also represent $f$ in the form:

\weq{7}{f = 4x_2\xi_3\,x_4\xi_{41} - v^2,}

in which $v$ is a quadratic form. From this and from the first representation in (5), however, there results the identity

\[
4x_2\xi_3(x_3\xi_2 - x_4\xi_{41}) = (u_1 - v)(u_1 + v),
\]

and from this we conclude that one of the two factors on the right must be divisible by $x_2\xi_3$ [\S.\,96, (3)]. If we assume that this is $u_1 - v$, which can be achieved by disposing of the sign of $v$, it follows that a constant factor $\lambda$ different from zero must exist, such that
% [sic] The source prints \lambda here; the following display and later prose use \lambda_1.

\[
\begin{aligned}
u_1 - v &= 2\lambda_1 x_2\xi_3,\\
u_1 + v &= \frac{2(x_3\xi_2 - x_4\xi_{41})}{\lambda_1},
\end{aligned}
\]

from which

\[
u_1 = \lambda_1 x_2\xi_3 + \frac{x_3\xi_2 - x_4\xi_{41}}{\lambda_1}
\]

follows. If we replace $u_1$ here by its expression from (6), we obtain

\[
x_4\xi_{41} = x_3\xi_2 - \lambda_1(-x_1\xi_1 + x_2\xi_2 + x_3\xi_3) + \lambda_1^2 x_2\xi_3.
\]

Of such equations we obtain, to begin with, three, if we permute the indices 1, 2, 3 cyclically and put in place of the unknown constant $\lambda_1$ three constants $\lambda_1, \lambda_2, \lambda_3$:

\weq{8}{\begin{aligned}
x_4\xi_{41} &= x_3\xi_2 - \lambda_1(-x_1\xi_1 + x_2\xi_2 + x_3\xi_3) + \lambda_1^2 x_2\xi_3\\
x_4\xi_{42} &= x_1\xi_3 - \lambda_2(x_1\xi_1 - x_2\xi_2 + x_3\xi_3) + \lambda_2^2 x_3\xi_1\\
x_4\xi_{43} &= x_2\xi_1 - \lambda_3(x_1\xi_1 + x_2\xi_2 - x_3\xi_3) + \lambda_3^2 x_1\xi_2.
\end{aligned}}

In order to determine the constants $\lambda_1, \lambda_2, \lambda_3$, we divide the last two of these equations by $\lambda_2$ and $\lambda_3$ and add. In this way follows the identical equation

\weq{9}{\begin{aligned}
x_4\left(\frac{\xi_{42}}{\lambda_2} + \frac{\xi_{43}}{\lambda_3}\right) &=\\
x_1\left(-2\xi_1 + \lambda_3\xi_2 + \frac{\xi_3}{\lambda_2}\right) &+ \xi_1\left(\lambda_2 x_3 + \frac{x_2}{\lambda_3}\right).
\end{aligned}}

Since now $x_4, x_1, \xi_1$ are azygetic, and consequently cannot pass through one point, the line

\[
\lambda_2 x_3 + \frac{x_2}{\lambda_3} = 0
\]
%CONT>>


% ===== scan leaf 398, printed 377 =====
% !PRINTED 377
\weberpage{377}
%<<CONT
pass through the intersection of $x_1$ and $x_4$, and consequently $x_4$ is linearly composed of $x_1$ and $\lambda_2 x_3 + \dfrac{x_2}{\lambda_3}$. Now, however, we can express $x_4$ linearly in terms of $x_1, x_2, x_3$, say in the form

\weq{10}{x_4 = a_1 x_1 + a_2 x_2 + a_3 x_3,}

and in it $a_1, a_2, a_3$ are to be regarded as known, and none of these constants can vanish. Hence $\lambda_2\lambda_3 = a_3 : a_2$, and if we introduce a new constant $h_1$, then

\weq{11}{\lambda_2 = h_1 a_3,\quad \frac{1}{\lambda_3} = h_1 a_2,}

and equation (9) yields the identity:

\[
h_1 x_4\left(\frac{\xi_{42}}{\lambda_2} + \frac{\xi_{43}}{\lambda_3} - h_1\xi_1\right) = x_1\left(-h_1(2+a_1h_1)\xi_1 + \frac{\xi_2}{a_2} + \frac{\xi_3}{a_3}\right).
\]

From this one concludes further, if $k_1$ is a new constant,

\[
k_1 x_4 = -h_1(2+a_1h_1)\xi_1 + \frac{\xi_2}{a_2} + \frac{\xi_3}{a_3},
\]

\weq{12}{\frac{\xi_{42}}{\lambda_2} + \frac{\xi_{43}}{\lambda_3} = h_1\xi_1 + \frac{k_1}{h_1} x_1.}

If one interchanges the figures 1, 2, 3 cyclically in this, then three such systems result from (8); first:

\weq{13}{\begin{aligned}
k_1 x_4 &= -h_1(2+a_1h_1)\xi_1 + \frac{\xi_2}{a_2} + \frac{\xi_3}{a_3}\\
k_2 x_4 &= \frac{\xi_1}{a_1} - h_2(2+a_2h_2)\xi_2 + \frac{\xi_3}{a_3}\\
k_3 x_4 &= \frac{\xi_1}{a_1} + \frac{\xi_2}{a_2} - h_3(2+a_3h_3)\xi_3.
\end{aligned}}

Since these three expressions for $x_4$ must be identical with one another, it follows that:

\[
-\frac{h_1(2+a_1h_1)}{k_1} = \frac{1}{k_2 a_1} = \frac{1}{k_3 a_1},
\]

hence $k_2 = k_3$ and likewise $k_2 = k_1$. Thus $k_1, k_2, k_3$ are equal to one another and we write $k$ for them. Then it follows further that

\[
a_1 h_1(2+a_1h_1) + 1 = (a_1h_1+1)^2 = 0,
\]

hence $h_1 a_1 = -1$, and likewise

\[
h_1 = \frac{-1}{a_1},\quad h_2 = \frac{-1}{a_2},\quad h_3 = \frac{-1}{a_3}.
\]


% ===== scan leaf 399, printed 378 =====
% !PRINTED 378
\weberpage{378}
Accordingly the formulae (13) yield in agreement

\[
k\,x_4 = \frac{\xi_1}{a_1} + \frac{\xi_2}{a_2} + \frac{\xi_3}{a_3}
\]
or
\weq{14}{k\,(a_1\,x_1 + a_2\,x_2 + a_3\,x_3) = \frac{\xi_1}{a_1} + \frac{\xi_2}{a_2} + \frac{\xi_3}{a_3}.}
From (11), however, there follows further
\weq{15}{\lambda_1 = -\frac{a_2}{a_3}, \quad \lambda_2 = -\frac{a_3}{a_1}, \quad \lambda_3 = -\frac{a_1}{a_2}.}
From (12) there then result, moreover, the three relations
\[
\begin{aligned}
-\frac{\xi_{42}}{\lambda_2} - \frac{\xi_{43}}{\lambda_3} &= \frac{\xi_1}{a_1} + k\,a_1\,x_1 \\
-\frac{\xi_{43}}{\lambda_3} - \frac{\xi_{41}}{\lambda_1} &= \frac{\xi_2}{a_2} + k\,a_2\,x_2 \\
-\frac{\xi_{41}}{\lambda_1} - \frac{\xi_{42}}{\lambda_2} &= \frac{\xi_3}{a_3} + k\,a_3\,x_3.
\end{aligned}
\]
From these, by addition, with regard to (14),
\[
\frac{\xi_{41}}{\lambda_1} + \frac{\xi_{42}}{\lambda_2} + \frac{\xi_{43}}{\lambda_3} = -k\,(a_1\,x_1 + a_2\,x_2 + a_3\,x_3),
\]
and consequently
\weq{16}{
\begin{aligned}
\frac{\xi_{41}}{\lambda_1} &= \frac{\xi_1}{a_1} - k\,(a_2\,x_2 + a_3\,x_3) \\
\frac{\xi_{42}}{\lambda_2} &= \frac{\xi_2}{a_2} - k\,(a_3\,x_3 + a_1\,x_1) \\
\frac{\xi_{43}}{\lambda_3} &= \frac{\xi_3}{a_3} - k\,(a_1\,x_1 + a_2\,x_2).
\end{aligned}
}
In these formulae the constant $k$ is still unknown. We determine it from the remark that we can quadruple the whole system of formulae considered so far, by letting $x_4$, $x_5$, $x_6$, $x_7$ take the place of $x_4$.

Corresponding to formula (10), we wish to express these four functions linearly in terms of $x_1$, $x_2$, $x_3$ in the manner:
\weq{17}{
\begin{aligned}
x_4 &= a_{4,1}\,x_1 + a_{4,2}\,x_2 + a_{4,3}\,x_3 \\
x_5 &= a_{5,1}\,x_1 + a_{5,2}\,x_2 + a_{5,3}\,x_3 \\
x_6 &= a_{6,1}\,x_1 + a_{6,2}\,x_2 + a_{6,3}\,x_3 \\
x_7 &= a_{7,1}\,x_1 + a_{7,2}\,x_2 + a_{7,3}\,x_3,
\end{aligned}
}
%CONT>>


% ===== scan leaf 400, printed 379 =====
% !PRINTED 379
\weberpage{379}
%<<CONT
in which the coefficients $a$ are to be regarded as known. Then from (14) we obtain four relations, if in place of $k$ we let four different constants $k_4$, $k_5$, $k_6$, $k_7$ enter:
\weq{18}{
\begin{aligned}
k_4\,(a_{4,1}\,x_1+a_{4,2}\,x_2+a_{4,3}\,x_3) &= \frac{\xi_1}{a_{4,1}}+\frac{\xi_2}{a_{4,2}}+\frac{\xi_3}{a_{4,3}} \\
k_5\,(a_{5,1}\,x_1+a_{5,2}\,x_2+a_{5,3}\,x_3) &= \frac{\xi_1}{a_{5,1}}+\frac{\xi_2}{a_{5,2}}+\frac{\xi_3}{a_{5,3}} \\
k_6\,(a_{6,1}\,x_1+a_{6,2}\,x_2+a_{6,3}\,x_3) &= \frac{\xi_1}{a_{6,1}}+\frac{\xi_2}{a_{6,2}}+\frac{\xi_3}{a_{6,3}} \\
k_7\,(a_{7,1}\,x_1+a_{7,2}\,x_2+a_{7,3}\,x_3) &= \frac{\xi_1}{a_{7,1}}+\frac{\xi_2}{a_{7,2}}+\frac{\xi_3}{a_{7,3}},
\end{aligned}
}
and now the constants $k$ are to be determined so that of these four equations one follows from the other three. The conditions for this can be formed in a symmetric manner by introducing four factors $l_4$, $l_5$, $l_6$, $l_7$, whose ratios are determined from the equations:
\weq{19}{\frac{l_4}{a_{4,i}} + \frac{l_5}{a_{5,i}} + \frac{l_6}{a_{6,i}} + \frac{l_7}{a_{7,i}} = 0, \quad i = 1, 2, 3,}
and which must then also satisfy the three equations
\weq{20}{k_4\,l_4\,a_{4,i} + k_5\,l_5\,a_{5,i} + k_6\,l_6\,a_{6,i} + k_7\,l_7\,a_{7,i} = 0, \quad i = 1, 2, 3}
from which the ratios of the $k$ are determined. A common factor of the four quantities $k$ remains, by the nature of the matter, undetermined and can be assumed arbitrarily. According to this, the functions $\xi_1$, $\xi_2$, $\xi_3$ can be rationally determined from the equations (18), and by (16) the functions $\xi_{4i}$, $\xi_{5i}$, $\xi_{6i}$, $\xi_{7i}$ are then also determined.

There remain six double tangents still wanting, which can be obtained by suitable permutations among the functions of the system (1). With this the Theorem I. placed at the head is proved.

\par\noindent\rule{0.3\textwidth}{0.4pt}\par

To see the correctness of Theorem II. as well, one need only follow our analysis backwards, regarding the coefficients $a_{k,i}$ as independent variables. Then by equations (18), (19), (20) the functions $\xi_1$, $\xi_2$, $\xi_3$ are
%CONT>>


% ===== scan leaf 401, printed 380 =====
% !PRINTED 380
\weberpage{380}
\pdfbookmark[2]{§ 101. The Galois group of the problem of the double tangents}{bm116}
%<<CONT
rationally determined by these $a_{k,i}$, and from (16) and (17) one then obtains $\xi_{4i}$, $\xi_{5i}$, $\xi_{6i}$, $\xi_{7i}$, $x_4$, $x_5$, $x_6$, $x_7$.

By substituting the $\xi_1$, $\xi_2$, $\xi_3$ into equation (4) one obtains the equation of a curve of the fourth order, whose coefficients are rational functions of the $a_{k,i}$, and the discriminant of this equation cannot vanish identically, because conversely, as we have seen, one can derive a system of equations (18), (19), (20) from the equation of a curve of the fourth order without a singular point.

From the equations (18), (16) there then also follow the formulae (7), and those which arise from them by permutation of 4 with 5, 6, 7; from which it is to be concluded that $x_4$, $x_5$, $x_6$, $x_7$ together with $x_1$, $x_2$, $x_3$ form a complete seven-system.

\S.\,101.

\gesperrt{The Galois group of the double tangent problem.}

The theorems that we have derived are sufficient to determine the \gesperrt{Galois} group of the algebraic equation on which the double tangents depend. We here take as \gesperrt{domain of rationality the field which consists of all rational functions of the 14 ratios of the coefficients of a general ternary form of degree $4$}, in which these coefficients count as \gesperrt{independent variables}. The equation of degree $28$ we may then think of as formed, say, in the following way: we take as unknowns the abscissae of the points of intersection of the double tangents with an arbitrary fixed straight line $L$. By means of the roots $\xi$ of this equation, which we call the \gesperrt{double tangent equation}, the double tangents are then rationally representable.

If one uses a \gesperrt{Cartesian} coordinate system $x$, $y$, whose $x$-axis is the line $L$, the equation of a double tangent takes the form
\weq{1}{y = \Theta\,(x - \xi),}
and the double tangent equation can be formed, if $F(x,y) = 0$ is the equation of the curve of the fourth order, by seeking the conditions that the function of $x$
%CONT>>


% ===== scan leaf 402, printed 381 =====
% !PRINTED 381
\weberpage{381}
%<<CONT
of degree $4$, $F[x, \Theta(x-\xi)]$, be a complete square. This gives two equations between $\xi$ and $\Theta$, from which, by elimination of $\Theta$, one obtains the double-tangent equation. Since to each $\xi$ there belongs only one value of $\Theta$ (so long as no two double tangents intersect on the line $L$), $\Theta$ can be expressed rationally in terms of $\xi$, and indeed in a form
\weq{2}{\Theta = \varphi\,(\xi),}
which holds for every pair $\xi$, $\Theta$ that belong together.

The roots of the double-tangent equation arrange themselves, just as the corresponding double tangents do, into complexes, seven-systems, and so on. We denote these roots, just like the double tangents in \S.\,99, by the symbol $[ik]$, in which $i$, $k$ run through the pairs of the figures 1 to 8 and $[ik] = [ki]$.

If we regard any two of the double tangents, $p$, $q$, as known, then we can derive from them by a rational path the equation $u = 0$ of a conic which passes through the four points of contact of these double tangents, and we can therefore set up the equation \S.\,96, (2)
\[
f = p\,q\,v - u^2
\]
in rational form. Then, by \S.\,96, (4), there are among the pencil of conics
\[
v + 2\,\lambda\,u + \lambda^2\,p\,q = 0
\]
five which split into a pair of lines, and if we form the product
\[
\prod_{\lambda} (v + 2\,\lambda\,u + \lambda^2\,p\,q) = \Phi
\]
over the five roots of the equation of degree $5$ on which $\lambda$ depends [\S.\,96, (6)], then this product is likewise expressible rationally in terms of the coefficients of $f$ and of $p$, $q$. But this product is a form of degree $10$, which splits into ten linear factors, which together with $p\,q$ form a \gesperrt{Steiner} complex.

The coefficients in $\Phi$ can now also be expressed rationally in terms of the two roots $\xi_1$, $\xi_2$ of the double-tangent equation corresponding to $p$, $q$, and if we then seek the abscissae of the points of intersection of the line $L$ with $\Phi = 0$, we obtain an equation of degree $10$ for $\xi$
\[
X(\xi_1, \xi_2, \xi) = 0,
\]
%CONT>>


% ===== scan leaf 403, printed 382 =====
% !PRINTED 382
\weberpage{382}
%<<CONT
whose roots are the ten roots syzygetic with $\xi_1\xi_2$. If $\xi_3$ denotes one of these, then

\weq{3}{X(\xi_1,\xi_2,\xi_3) = 0,}

and from this follows the theorem

\gesperrt{1. There is a rational function of three variables $X(\xi_1,\xi_2,\xi_3)$ which vanishes when $\xi_1\xi_2\xi_3$ is any syzygetic triple of roots of the double tangent equation, and which does not vanish when $\xi_1\xi_2\xi_3$ is an azygetic triple.}

By \S.\,100, by means of a complete seven-system of these roots

\weq{4}{\xi_1,\ \xi_2,\ \xi_3,\ \xi_4,\ \xi_5,\ \xi_6,\ \xi_7}

all the roots can be expressed rationally, and indeed in such a way that, e.g.,

\weq{5}{\xi_{12} = \Psi(\xi_1,\xi_2 \mid \xi_3,\xi_4,\xi_5,\xi_6,\xi_7)}

becomes a root, where $\Psi$ denotes a rational function which does not change when $\xi_1$ and $\xi_2$ are interchanged, or when $\xi_3,\xi_4,\xi_5,\xi_6,\xi_7$ are permuted arbitrarily. But if in (5), the function $\Psi$ being held fixed, another pair $\xi_i,\xi_k$ is put in place of $\xi_1,\xi_2$, one obtains another root $\xi_{ik}$. The roots (4) may also be denoted by $[1\,8],\dots[7\,8]$ and $\xi_{ik}$ by $[i\,k]$.

By every permutation of the roots (1), according to (2), the whole system of the roots $[i\,k]$ undergoes a certain permutation.

But if one replaces the complete seven-system (4) by another, the formula (5) yields a definite other root, and the whole system of the roots $[i\,k]$ is subjected to a second permutation. It is then first of all easy to prove:

\gesperrt{2. The permutation group $P$ of the roots of the double tangent equation, which one obtains when in (4) and (5) one puts in place of $\xi_1,\xi_2,\dots\xi_7$ all complete systems, each in every arbitrary order, is the Galois group of the double tangent equation.}

To demonstrate this, we have two things to show:
%CONT>>


% ===== scan leaf 404, printed 383 =====
% !PRINTED 383
\weberpage{383}
%<<CONT
\gesperrt{a) Every permutation $\pi$ of the roots of the double tangent equation which is applicable to all rational equations between these roots belongs to $P$.}

This follows thus: If $\pi$ is applicable to all rational equations between the roots, then the same holds of the powers of $\pi$. By 1., a syzygetic triple can never be carried by $\pi$ into an azygetic one, or conversely; for if $\xi_1\xi_2\xi_3$ is a syzygetic triple, $\pi$ is applicable to the equation (3); hence $\xi_1\xi_2\xi_3$ cannot pass into an azygetic triple. And the converse also is not possible, because otherwise a syzygetic triple would be carried into an azygetic one by $\pi^{-1}$. Hence any complete seven-system also passes, under $\pi$, again into such a system in some arrangement or other, and if one then applies $\pi$ to all equations of the form (5), there results a permutation of the $\xi_i, \xi_{ik}$ which belongs to $P$.

\gesperrt{b) Every rational equation between the roots of the double tangent equation admits all permutations of the group $P$.}

A rational relation between the roots has the form

\weq{6}{\Phi(\xi_1,\dots,\xi_7,\xi_{12},\dots,a\dots) = 0,}

in which $\Phi$ is a rational function, and $a\dots$ denote the ratios of the coefficients of $f$. Herein one can, by (5), express the $\xi_{12},\xi_{13},\dots$ rationally in terms of $\xi_1,\xi_2,\dots,\xi_7$, and by \S.\,100, II., since the curve $f$ is rationally determined by a complete seven-system of its double tangents, the $a$ can then be expressed rationally (with only numerical coefficients) in terms of the seven pairs of quantities (2)

\[
\xi_1, \Theta_1;\ \xi_2, \Theta_2;\ \dots;\ \xi_7, \Theta_7
\]

But since these can be given entirely arbitrarily, the equation (6) must pass into an identity under these substitutions. The equation (6) must therefore also remain correct if one replaces in it the $\xi_1,\xi_2,\dots,\xi_7$ by another seven-system, or also by the same one in another order, and at the same time carries out among the $\xi_{12},\dots$ the permutation prescribed by the formula (5), that is, if one performs among the roots of the double tangent equation a permutation of the group $P$.


% ===== scan leaf 405, printed 384 =====
% !PRINTED 384
\weberpage{384}
\pdfbookmark[2]{§ 102. Representation of the group}{bm117}
\providecommand{\weberfourup}{\mathord{\tikz[baseline=-0.2ex,x=0.23em,y=0.23em,line width=0.45pt]{\draw (0,0)--(-1.05,1) (0,0)--(-0.35,1.08) (0,0)--(0.35,1.08) (0,0)--(1.05,1);}}}
\begin{center}
\S.\,102.

\gesperrt{Representation of the Group.}
\end{center}

The degree of the group $P$ can be stated at once. Since there are 288 complete seven-systems, and since the elements of such a system can be permuted in $\Pi(7)$ ways, the degree of the group is:

\[
288\,\Pi(7) = 36\,\Pi(8) = 1\,451\,520.
\]

In forming the permutations of the roots we make use of two kinds of notation. First

\weq{1}{\xi_1,\ \xi_2,\ \dots,\ \xi_7,\ \xi_{ik},}

in which $i,k$ runs from 1 to 7, and the uniform notation $[i\,k]$, where $i,k$ runs from 1 to 8, and where $\xi_1,\xi_2\dots$ are to be denoted by $[1\,8],[2\,8],\dots$.

If we now permute two figures from the sequence $1,2,\dots,7$, e.\,g. 1 with 2, then $\xi_1$ passes into $\xi_2$, and by the definition of \S.\,99 $\xi_{12}$ remains unchanged, $\xi_{13}$ passes into $\xi_{23}$, and so on.

If, then, we permute the figures 1 to 7 arbitrarily in the sequence of the roots $[i\,k]$, we obtain nothing but permutations of the group $P$.

Now we have seen in \S.\,100 that, in the arrangement into syzygetic and azygetic triples, and in consequence also into complexes and complete systems, the figure 8 appears as completely equally entitled with the remaining figures 1 to 7, and since the assignment of the root $\xi_{ik}$ to the pair $\xi_i\xi_k$ [by the formula \S.\,101, (5)] depends by \S.\,99 only on this arrangement, we can also permute the figures $1,2,\dots,7,8$ without leaving the group $P$. There is therefore contained in $P$ a divisor which is isomorphic with the symmetric group of the permutations of eight elements, which therefore has the index 36, and which we shall denote by $S$.

To determine the still missing permutations of $P$, we let a new seven-system of the type $\Delta\,\weberfourup$ take the place of the seven-system $\xi_1,\dots,\xi_7$, say as follows:

% [sic] The source prints a full stop after [6 8] in the first row.
\weq{2}{\begin{pmatrix} [1\,8], & [2\,8], & [3\,8], & [4\,8], & [5\,8], & [6\,8]. & [7\,8] \\ {}[2\,3], & [3\,1], & [1\,2], & [4\,8], & [5\,8], & [6\,8], & [7\,8] \end{pmatrix},}
%CONT>>


% ===== scan leaf 406, printed 385 =====
% !PRINTED 385
\weberpage{385}
%<<CONT
and denote the permutation thereby determined in a twofold way by

\weq{3}{\Pi_{1,2,3,8} = \Pi_{4,5,6,7},}

in that we stipulate that $\Pi_{\alpha_1\alpha_2\alpha_3\alpha_4}$ and $\Pi_{\beta_1\beta_2\beta_3\beta_4}$ shall mean the same thing when $\beta_1,\beta_2,\beta_3,\beta_4$ is any permutation of $\alpha_1,\alpha_2,\alpha_3,\alpha_4$, or when $\alpha_1,\alpha_2,\alpha_3,\alpha_4,\beta_1,\beta_2,\beta_3,\beta_4$ together comprise all eight figures.

By this symbol the permutation (2) is unambiguously designated, and the number of the distinct permutations of this kind amounts to exactly 35, so that we can represent the whole group $P$ by the subsidiary groups thus:

\weq{4}{P = S + \Sigma\, S\, \Pi_{\alpha_1\alpha_2\alpha_3\alpha_4},}

wherein $\alpha_1,\alpha_2,\alpha_3,\alpha_4$ run through all systems of four figures from the sequence $1,2,\dots,8$, in which an arbitrary figure, e.\,g. $\alpha_4 = 8$, can be held fixed.

To recognise the influence of $\Pi_{1,2,3,8}$ on any root $[i\,k]$, it suffices to consider the three roots $[1\,2], [1\,4], [4\,5]$, because $1,2,3$ on the one hand, $4,5,6,7$ on the other, occur entirely alike in $\Pi_{1,2,3,8}$.

These permutations are obtained simply from the remark that $[1\,2], [1\,4]$ are, by \S.\,99, the roots joined with $[2\,8]$ and $[4\,8]$ in the complex

\weq{5}{[2\,8][1\,2],\ [3\,8][1\,3],\ [4\,8][1\,4],\ [5\,8][1\,5],\ [6\,8][1\,6],\ [7\,8][1\,7]}

and that likewise $[4\,5]$ is the root joined with $[5\,8]$ in the complex

\weq{6}{[1\,8][1\,4],\ [2\,8][2\,4],\ [3\,8][3\,4],\ [5\,8][5\,4],\ [6\,8][6\,4],\ [7\,8][7\,4]}

But by the permutation (2) the complexes (5) and (6) pass over into the following, as one easily finds from the representation of the complexes in \S.\,99 [the complex (5) remains unchanged as a whole]:

\[
[3\,1][3\,8],\ [2\,1][2\,8],\ [4\,8][1\,4],\ [5\,8][5\,1],\ [6\,8][6\,1],\ [7\,8][7\,1]
\]

\[
[2\,3][1\,4],\ [3\,1][2\,4],\ [1\,2][3\,4],\ [5\,8][6\,7],\ [6\,8][5\,7],\ [7\,8][5\,6],
\]

from which one obtains the following permutations effected by $\Pi_{1,2,3,8}$:

\weq{7}{\begin{array}{l} [1\,2],\ [1\,4],\ [4\,5] \\ {}[3\,8],\ [1\,4],\ [6\,7]. \end{array}}

From (2) and (7) the following general and simple rule can now be derived.


% ===== scan leaf 407, printed 386 =====
% !PRINTED 386
\weberpage{386}

3.\ \gesperrt{If $\alpha_1, \alpha_2, \alpha_3, \alpha_4, \beta_1, \beta_2, \beta_3, \beta_4$ together denote all eight figures, then, in order to determine the influence of}
\[
\Pi_{\alpha_1\alpha_2\alpha_3\alpha_4} = \Pi_{\beta_1\beta_2\beta_3\beta_4}
\]
\gesperrt{on any root $[\mu\,\nu]$, one has to distinguish whether $\mu, \nu$ both occur among the $\alpha$ or both among the $\beta$, or whether the figure $\mu$ occurs among the $\alpha$ and $\nu$ among the $\beta$. In the first cases $[\mu\,\nu]$ passes into $[\mu'\,\nu']$, where $\mu, \nu, \mu', \nu'$ denote either all $\alpha$ or all $\beta$; in the third case $[\mu\,\nu]$ remains unchanged.}

In all cases the permutations $\Pi_{\alpha_1\alpha_2\alpha_3\alpha_4}$ are only of the second degree, and hence on a single repetition lead back to the identity.

The group $P$ is thereby completely determined and represented, and to establish the laws of composition in $P$ only a few more formulae are needed, which follow easily from the rule set up above. Here it is to be remarked that one can always take two distinct permutations $\Pi_{\alpha_1\alpha_2\alpha_3\alpha_4}$ in such a way that they have two or three figures in common in the index, since, if they have only one figure in common, one may take for the index of one of them its complement.

Now let
\[
\sigma = \begin{pmatrix} 1, & 2, & 3, & 4, & 5, & 6, & 7, & 8 \\ \alpha_1, & \alpha_2, & \alpha_3, & \alpha_4, & \alpha_5, & \alpha_6, & \alpha_7, & \alpha_8 \end{pmatrix}
\]
be any permutation of the eight figures, and let $(\alpha, \beta)$ denote the transposition of the two figures $\alpha$ and $\beta$, hence an element of $S$. Then
\weq{8}{\Pi_{1,2,3,4}\,\sigma = \sigma\,\Pi_{\alpha_1\alpha_2\alpha_3\alpha_4},}
\weq{9}{\Pi_{1,2,3,4}\,\Pi_{1,2,3,4} = 1,}
\weq{10}{\Pi_{1,2,3,4}\,\Pi_{1,2,3,5} = (4,5)\,\Pi_{1,2,3,4},}
\weq{11}{\Pi_{1,2,3,4}\,\Pi_{1,2,5,6} = (1,2)(3,4)(5,6)(7,8)\,\Pi_{1,2,7,8}.}

One proves these formulae easily by rule 3., by going through the individual cases, in which of course only quite a small number of types have to be considered; thus, for example, $[1\,4]$ passes by $\Pi_{1,2,3,4}$ into $[2\,3]$, and this by $\Pi_{1,2,3,5}$ into $[1\,5]$, and $[1\,5]$ is no longer altered by $\Pi_{1,2,3,4}$; consequently there is effected by
\weq{12}{\Pi_{1,2,3,4}\,\Pi_{1,2,3,5}\,\Pi_{1,2,3,4}}
%CONT>>


% ===== scan leaf 408, printed 387 =====
% !PRINTED 387
\weberpage{387}
\pdfbookmark[2]{§ 103. Simplicity of the group of the problem of the double tangents}{bm118}

%<<CONT
the interchange of $[1\,4]$ with $[1\,5]$ in agreement with formula (10). Just as easily one recognises that, e.g., $[1\,2]$ is not altered by (12).

\begin{center}
\S.\,103.

\gesperrt{Simplicity of the group of the double tangent problem.}
\end{center}

The representation of the group $P$ which we developed in the preceding section now furnishes us a quite simple proof that this group has \gesperrt{no normal divisor}, and hence, in the manner of expression we used earlier, is \gesperrt{simple}; from which it then follows that the group cannot be lowered by the adjunction of irrationalities with smaller group, and thus, for example, not by cyclic equations.

We represent $P$ in the form (4), \S.\,102:
\weq{1}{P = S + \Sigma\, S\, \Pi_{\alpha_1, \alpha_2, \alpha_3, \alpha_4},}
in which $S$ is the whole group of all permutations of eight figures. The group $S$ has a normal divisor $S'$ of index 2, namely the alternating group of the eight figures, which for its part is simple, and every normal divisor of $S$ that does not consist of the identical permutation alone must contain the whole group $S'$. We set
\weq{2}{S = S' + S'',}
in which $S'' = S'\sigma$, where $\sigma$ denotes any permutation of the second kind, e.g. a transposition.

We now assume that $Q$ is a normal divisor of $P$ which does not consist of the identical substitution alone.

The greatest common divisor of $Q$ and $S$ is then a normal divisor of $S$, and must therefore, if it is not the identical group, contain the group $S'$. From this it follows:

1.\ \gesperrt{If $Q$ contains a non-identical permutation from $S$, then $Q$ contains the whole group $S'$.}

We then prove that $Q$, if it contains the group $S'$, must be identical with $P$.

For if $Q$ contains the whole group $S'$, then it contains
%CONT>>


% ===== scan leaf 409, printed 388 =====
% !PRINTED 388
\weberpage{388}

%<<CONT
it, as a normal divisor of $P$, also contains, when $(4,5,6)$ is a three-term cycle from $S'$ [by \S.\,102, (8), (10)]:
\[
\Pi_{1,2,3,5}\,(5,4,6)\,\Pi_{1,2,3,5} = (5,4,6)\,\Pi_{1,2,3,4}\,\Pi_{1,2,3,5}
\]
\[
= (5,4,6)(4,5)\,\Pi_{1,2,3,4},
\]
and consequently also
\[
S'(4,5)\,\Pi_{1,2,3,4} = S''\,\Pi_{1,2,3,4}.
\]
Now if
\[
\sigma = \begin{pmatrix} 1, & 2, & 3, & 4, & 5, & 6, & 7, & 8 \\ \alpha_1, & \alpha_2, & \alpha_3, & \alpha_4, & \alpha_5, & \alpha_6, & \alpha_7, & \alpha_8 \end{pmatrix}
\]
is a permutation from $S$ in which $\alpha_1, \alpha_2, \alpha_3, \alpha_4$ are arbitrarily given, then one can still choose the arrangement of the $\alpha_5, \alpha_6, \alpha_7, \alpha_8$ so that $\sigma$ belongs at will to $S'$ or to $S''$. If one chooses $\sigma$ in $S'$ and observes that then $S'' = S''\sigma^{-1}$, it follows that
\[
S''\,\sigma^{-1}\,\Pi_{1,2,3,4}\,\sigma = S''\,\Pi_{\alpha_1,\alpha_2,\alpha_3,\alpha_4}
\]
must also be contained in $Q$. But now we have also
\[
\Pi_{1,2,3,4}\,(3,5)(4,6)\,\Pi_{1,2,3,4} = (3,5)(4,6)\,\Pi_{1,2,5,6}\,\Pi_{1,2,3,4}
\]
\[
= (3,5)(4,6)(1,2)(3,4)(5,6)(7,8)\,\Pi_{1,2,7,8}
\]
\[
= (1,2)(3,6)(4,5)(7,8)\,\Pi_{1,2,7,8},
\]
and since $(1,2)(3,6)(4,5)(7,8)$ belongs to $S'$, $Q$ also contains $S'\,\Pi_{1,2,7,8}$, and consequently, as above, all $S'\,\Pi_{\alpha_1,\alpha_2,\alpha_3,\alpha_4}$, hence also all $S\,\Pi_{\alpha_1\alpha_2\alpha_3\alpha_4}$, and therefore also
\[
S\,\Pi_{\alpha_1,\alpha_2,\alpha_3,\alpha_4}\,\Pi_{\alpha_1,\alpha_2,\alpha_3,\alpha_4} = S,
\]
that is, $Q$ comprises the whole group $P$. From this it follows, in combination with 1.:

2.\ \gesperrt{If a normal divisor $Q$ of $P$ has, besides the identical permutation, any permutation whatever in common with $S$, then $Q$ is identical with $P$.}

It may now further be asked whether $Q$ can contain a permutation from one of the subsidiary groups of (1), that is, an element of the form $\sigma\,\Pi_{1,2,3,4}$, without being identical with $P$. If first $\sigma = 1$, so that $Q$ contains the element $\Pi_{1,2,3,4}$, then as a normal divisor of $P$ it also contains all the other $\Pi_{\alpha_1\alpha_2\alpha_3\alpha_4}$, as may be seen from the formula \S.\,102, (8), and with them also $\Pi_{1,2,3,4}\,\Pi_{1,2,3,5}\,\Pi_{1,2,3,4} = (4,5)$, and is therefore by 2. identical with $P$.


% ===== scan leaf 410, printed 389 =====
% !PRINTED 389
\weberpage{389}

But if $\sigma$ is different from 1, then one can choose a pair of figures $\alpha, \beta$, both among the $1,2,3,4$ or both among the $5,6,7,8$, in such a way that the transposition
\[
\sigma\,(\alpha,\beta)\,\sigma^{-1} = (\alpha', \beta')
\]
is different from $(\alpha, \beta)$. For if one first takes for $\alpha$ a figure changed by $\sigma^{-1}$, then $\alpha'$ is different from $\alpha$ (Bd.\,I, \S.\,154, 6.); if one then chooses for $\beta$ a figure which does not pass into $\alpha$ under $\sigma^{-1}$, and which occurs together with $\alpha$ in the first or in the second half of the eight figures — such a figure certainly always exists, because only one figure passes into $\alpha$ under $\sigma^{-1}$ — then $(\alpha', \beta')$ is different from $(\alpha, \beta)$. But it then follows that in the group $Q$ the element $(\alpha,\beta)\,\sigma\,\Pi_{1,2,3,4}\,(\alpha,\beta)$, and consequently also [by \S.\,102, (8), (9)]
\[
(\alpha,\beta)\,\sigma\,\Pi_{1,2,3,4}\,(\alpha,\beta)\,\Pi_{1,2,3,4}\,\sigma^{-1} = (\alpha,\beta)\,\sigma\,(\alpha,\beta)\,\sigma^{-1} = (\alpha,\beta)(\alpha',\beta')
\]
occurs, and this is a permutation from $S$ different from the identity. With this, then, it is proved:

3.\ \gesperrt{The group $P$ of the double-tangent equation is simple.}

As a special application we may point out that the group $P$ can effect among the 28 roots of the double-tangent equation only permutations of the first kind, and that \gesperrt{consequently the discriminant of this equation is a square}.

For if we denote for the moment by $G$ the group of all permutations of the 28 roots, and by $A$ the alternating group contained in it as a normal divisor, then the greatest common divisor of $A$ and $P$ is a normal divisor of $P$ and must therefore be identical with $P$; that is, $P$ is contained in $A$.

If one seeks out the group of the permutations which leave one of the roots of the double-tangent equation, say $[1\,2]$, unchanged, one finds a group which is still transitive on the remaining 27 roots\footnote{By a geometrical theorem of \gesperrt{Geiser}, this group is isomorphic with the group of the equation of degree 27, on which the solution of the problem of the 27 straight lines on a surface of the third order depends. (Mathem. Annalen, Bd.\,I.).}.


% ===== scan leaf 411, printed 390 =====
% !PRINTED 390
\weberpage{390}

To prove this, it suffices to show that by the permutations
of this group any root, say [1\ 3], can pass into
every other (with the exception of [1\ 2]). Now
[1,\,3] % [sic] comma printed here, unlike the bracket pairs elsewhere on this page which are space-separated without a comma
passes, by permutations of the group $S$ by which
1, 2 remain unchanged, into [1\ 4], \dots, [1\ 8]; likewise [2\ 3] into
[2\ 4], \dots, [2\ 8]. It therefore remains to show that [1\ 3]
can also be carried into [2\ 3] and into [4\ 5]. But this is shown
by inspection of the three following complete systems of seven:

\[
\begin{aligned}
&[1\ 2]\ [1\ 3]\ [1\ 4]\ [1\ 5]\ [1\ 6]\ [1\ 7]\ [1\ 8]\\
&[1\ 2]\ [2\ 3]\ [2\ 4]\ [2\ 5]\ [2\ 6]\ [2\ 7]\ [2\ 8]\\
&[1\ 2]\ [4\ 5]\ [3\ 4]\ [3\ 5]\ [1\ 6]\ [1\ 7]\ [1\ 8].
\end{aligned}
\]

According to this the group $P$ is doubly transitive. But it
cannot be more than doubly transitive. For if one leaves
two roots unchanged, a root syzygetic with these two
cannot be carried into an azygetic one.
By the adjunction of two roots the double tangent equation
becomes reducible. A factor of degree $10$ separates off,
whose roots, together with the two given ones, form a \gesperrt{Steiner}
complex (the function $X$ in Theorem 1., \S.\,101).

Among the divisors of the group $P$ the
group $S$ of index 36 is especially remarkable. This group is still
transitive, because by it [1\ 2] can be carried into any arbitrary root $[\varkappa\ \lambda]$.
But it is only simply transitive,
because when [1\ 2] is held fixed the root [1\ 3] can no longer
pass into [4\ 5].

By the adjunction of \gesperrt{one} root of an equation of degree $36$,
the group of the double tangent equation is therefore reduced to the
group of a general equation of degree $8$.

If one sets, e.\,g.,

\[
v_1 = [1\ 2]\ [1\ 3]\ [1\ 4]\ [1\ 5]\ [1\ 6]\ [1\ 7]\ [1\ 8],
\]

and defines $v_2, v_3, \dots, v_8$ correspondingly, then every symmetric
function of these eight quantities, e.\,g.

\[
u = v_1 + v_2 + v_3 + v_4 + v_5 + v_6 + v_7 + v_8,
\]

is a root of an equation of degree $36$. If one adjoins to the
problem the quantity $u$, the $v_1, v_2, \dots, v_8$ become the roots
of an equation of degree $8$, which has no affect.


% ===== scan leaf 412, printed 391 =====
% !PRINTED 391
\weberpage{391}
\pdfbookmark[2]{§ 104. Reality of the double tangents}{bm119}

\S.\,104.

\gesperrt{Reality of the double tangents.}

We wish further to discuss the question of how, in the case
of a real curve of the fourth order without singular point, the double
tangents can behave with regard to their reality. We
now assume, then, the coefficients of the equation of the curve
of the fourth order to be real, i.\,e. we presuppose a real domain
of rationality.

If then a double tangent $\xi$ is imaginary, the conjugate
line $\xi'$ must also be a double tangent.

If in a rational equation between the roots of the
double-tangent equation we replace every imaginary root by
the conjugate, there arises again a correct equation.
It follows from this that to syzygetic or azygetic
systems of % [sic] the German prints "vom" here where "von" would be expected
double tangents there belong conjugate systems of the same
character, which one obtains by replacing $i$ everywhere
by $-i$.

A \gesperrt{Steiner} complex $C$ therefore passes, by the transition
from $i$ to $-i$, into a \gesperrt{Steiner} complex $C'$, and we
distinguish two cases:

If $C$ is identical with $C'$, we call $C = C'$ a \gesperrt{real
complex}.

But if $C$ is different from $C'$, they form a \gesperrt{conjugate
complex pair}.

The pairs of a real complex consist either of
two real double tangents, or they are conjugate pairs $\xi, \xi'$,
or they contain two non-conjugate imaginary double tangents
$\xi, \eta$. In the latter case the pair formed from
the conjugate double tangents $\xi', \eta'$ must then also occur
in the complex. Two such pairs $(\xi, \eta)$, $(\xi', \eta')$ we shall call a
\gesperrt{conjugate double pair}.

Never does a real root $x$ occur paired with an imaginary
$\xi$ in a real complex, because otherwise, beside
the pair $(x, \xi)$, the conjugate pair $(x, \xi')$ would also have to occur in the same
complex, which is impossible, since these two
pairs would contain a common element $x$.

An \gesperrt{Aronhold} seven-system $S$ passes, by the interchange
of all its elements with the conjugates, into just such a
%CONT>>


% ===== scan leaf 413, printed 392 =====
% !PRINTED 392
\weberpage{392}

%<<CONT
into such a system $S'$, and we call such a system
\gesperrt{real} when $S$ is identical with $S'$. A real seven-system
therefore contains, for every root $\xi$ occurring in it, the conjugate
$\xi'$, and consequently must contain at least one, but always an
odd number, of real double tangents.

We now derive the principal theorems in order:

\gesperrt{1. Not all the double tangents can be imaginary.}

For let $\xi, \xi'$ be a conjugate pair and $\eta$ a third
imaginary double tangent, such that the three $\xi, \xi', \eta$ are \gesperrt{azygetic}.
(This would certainly be possible if all the double tangents
were imaginary.) If $\eta'$ is conjugate to $\eta$, we have the two
azygetic triples:

\[
\xi,\ \xi',\ \eta;\quad \xi,\ \xi',\ \eta'.
\]

Since, accordingly, $\xi'$ does not occur in the \gesperrt{Steiner}
complex determined by the pair $\xi\eta$, this complex is
imaginary.

We set it beside its conjugate complex:

\weq{1}{
\begin{aligned}
&1)\quad \xi\eta,\ \xi_1\eta_1,\ \xi_2\eta_2,\ \xi_3\eta_3,\ \xi_4\eta_4,\ \xi_5\eta_5\\
&2)\quad \xi'\eta',\ \xi_1'\eta_1',\ \xi_2'\eta_2',\ \xi_3'\eta_3',\ \xi_4'\eta_4',\ \xi_5'\eta_5',
\end{aligned}
}

where by $\xi_i', \eta_i'$ we understand the elements conjugate to $\xi_i, \eta_i$,
so that, if $\xi_i$ is real, we are to set $\xi_i = \xi_i'$, and conversely
from $\xi_i = \xi_i'$ it also follows that $\xi_i$ is real.

Two possibilities are now to be distinguished. If,
\gesperrt{first}, the complexes 1), 2) form a syzygetic pair (\S.\,97),
then they have four syzygetic elements in common. Among these
$\xi, \eta, \xi', \eta'$ cannot occur, and we therefore do not restrict
generality if we assume that $\xi_1, \eta_1, \xi_2, \eta_2$ are
the common elements. But these cannot in their totality
be different from $\xi_1', \eta_1', \xi_2', \eta_2'$, since the
passage to the accented letters, i.\,e. to the conjugate
imaginary quantities, by which 1) goes over into 2), is everywhere
permitted.

Now the assumption $\xi_1 = \eta_1'$ is excluded, because from it
$\xi_1' = \eta_1$ would follow and 1) would not be different from 2).

If $\xi_1 = \xi_1'$, then $\xi_1$ is \gesperrt{real}.

But if $\xi_1 = \xi_2'$, $\xi_2 = \xi_1'$,
then $\eta_1 = \eta_1'$, hence $\eta_1$ is \gesperrt{real},
and from this the assumption $\xi_1 = \eta_2'$, $\xi_1' = \eta_2$, $\xi_2 = \xi_2'$ is not
essentially different.


% ===== scan leaf 414, printed 393 =====
% !PRINTED 393
\weberpage{393}
This case therefore always leads to a real double tangent.

But if, \gesperrt{secondly}, the pair of complexes 1), 2) is azygetic, then
every pair of the one complex contains an element that also occurs in the
other complex, and an element that does not occur in the other. If, then,
say $\eta$ occurs in the complex 2), we may, without impairing generality,
assume $\eta = \eta_1'$, $\eta' = \eta_1$,
and we obtain the following pairs in 1) and 2)
\weq{2}{
\begin{array}{rl}
1) & \xi\eta,\ \xi_1\eta' \\
2) & \xi_1'\eta,\ \xi'\eta'.
\end{array}
}
If, further, $\eta_2$ occurs in 1) and 2), then $\eta_2$ cannot be equal to
$\xi_2'$, since otherwise 1) would contain the pair $\xi_2\xi_2'$ and would
consequently be real.

If $\eta_2 = \eta_2'$, then $\eta_2$ is \gesperrt{real}. But if $\eta_2$ is equal
to an element of the three last pairs of 2), then we may without
restriction assume it $= \eta_3'$, hence $\eta_2 = \eta_3'$, $\eta_3 = \eta_2'$.

But then, by the theorem of \S.\,96, 3., both $\eta, \eta_2, \eta_2'$ and
$\eta, \eta', \eta_2$ are azygetic. The two pairs $\eta\eta_2$, $\eta'\eta_2'$
determine two complexes, of which the second contains neither $\eta$ nor $\eta_2$,
and which therefore form a \gesperrt{syzygetic pair}. With this we are led back
to the first assumption, and our theorem 1. is proved.

There are therefore always at least \gesperrt{two real} double tangents.

\gesperrt{2. There is always at least one system of four real
syzygetic double tangents.}

In the proof of this theorem we start from a real pair $x\,y$, which
by 1. always exists, and consider the real complex that is determined
by this pair. Since the conjugate double pairs are two pairs each,
there must, among the six pairs of this complex,
either occur a second real pair, and then theorem
2. holds, or it must contain a conjugate pair $\xi\xi'$.

In the latter case we consider the syzygetic triple of complexes
determined by these two pairs, in which all 28 double tangents
occur (\S.\,97, 4.):
\[
\begin{array}{rl}
1) & x\,y,\ \xi\xi' \\
2) & x\xi,\ y\xi',\ \xi_1\eta_1 \\
3) & x\xi',\ y\xi,\ \xi_1'\eta_1',
\end{array}
\]
%CONT>>


% ===== scan leaf 415, printed 394 =====
% !PRINTED 394
\weberpage{394}
%<<CONT
of which the last two are conjugate imaginary, and which consequently,
since apart from $x, y, \xi, \xi'$ they have no common element, contain
only imaginary pairs besides, of which we have listed one, $\xi_1\eta_1$,
together with its conjugate $\xi_1'\eta_1'$.

It therefore suffices if we can prove the existence of one further real
double tangent. For this must then occur in 1), and must in this complex
belong to a real pair.

We now form in addition the two conjugate complexes:
\[
\begin{array}{rl}
4) & x\xi_1,\ \xi\eta_1 \\
5) & x\xi_1',\ \xi'\eta_1',
\end{array}
\]
in which $y$ certainly does not occur, since otherwise $x, \xi, y$ would be
azygetic (\S.\,96, 3.). For the same reason $\xi$ does not occur in 5),
nor $\xi'$ in 4).

Now if, \gesperrt{in the first case}, the complexes 4) and 5) are syzygetic, then
$\xi_1$ must occur in 5) and $\xi_1'$ in 4), and if 4) contains the pair $p\xi_1'$,
then the pair $p\xi_1$ must occur in 5). Thus $p$ is identical with its
conjugate element, that is, \gesperrt{real}.

If, \gesperrt{secondly}, the complexes 4) and 5) are azygetic, then $\xi_1, \xi$ do
not occur in 5), and $\eta_1$ must be contained in 5).
If $\eta_1$ is paired in 5) with an imaginary element $\xi_2'$, and
if in 4) yet another imaginary pair $\xi_3\eta_3$ occurs, then these
complexes continue in the following way:
\[
\begin{array}{rl}
4) & x\xi_1,\ \xi\eta_1,\ \xi_2\eta_1',\ \xi_3\eta_3,\ \xi_3'\eta_4 \\
5) & x\xi_1',\ \xi'\eta_1',\ \xi_2'\eta_1,\ \xi_3'\eta_3',\ \xi_3\eta_4',
\end{array}
\]
in which $\eta_4$ is also imaginary, and the accented letters are always
understood to mean the elements conjugate to the unaccented ones.
The two complexes then contain only one further pair each, $\xi_5\eta_5$,
$\xi_5'\eta_5'$, which must contain a common element. But it cannot
be that $\xi_5 = \eta_5'$, since otherwise also $\xi_5' = \eta_5$, and
hence the two pairs would be identical. Therefore it must be that $\xi_5 = \xi_5'$ (or
$\eta_5 = \eta_5'$); that is, one of these two double tangents
is \gesperrt{real}, and with this our theorem 2. is proved.

3. \gesperrt{If, besides the four real double tangents $x, y, p, q$,
whose existence theorem 2. asserts, yet another one is present,
then there exist eight real double tangents, which form four pairs
in a Steiner complex.}


% ===== scan leaf 416, printed 395 =====
% !PRINTED 395
\weberpage{395}
We form the real syzygetic complex triple in which all
double tangents must occur:
\[
\begin{array}{rl}
1) & x\,y,\ p\,q \\
2) & x\,p,\ y\,q \\
3) & x\,q,\ y\,p.
\end{array}
\]
If a fifth real double tangent $r$ is present, we may assume
that it occurs in complex 1). But then this complex must
contain a third real pair $r\,s$, and either a fourth real pair,
in which case Theorem 3. already holds, or a conjugate
pair $\varrho\varrho'$.

If complex 2) does not contain solely real double tangents —
a case in which Theorem 3. would likewise hold — then an
imaginary pair $\xi\eta$ must occur in 2), and accordingly we thus consider the
following complexes:
\[
\begin{array}{rl}
1) & x\,y,\ p\,q,\ r\,s,\ \varrho\varrho' \\
2) & x\,p,\ y\,q,\ \xi\eta \\
4) & x\,r,\ y\,s.
\end{array}
\]
The two complexes 2), 4), however, are \gesperrt{azygetic}, since e.g. $r$ does not
occur in 2) [because it occurs in 1)]. Consequently 4) can also contain
only one of the two $\xi, \eta$, and consequently, since the complex
4) is real, $\xi, \eta$ cannot be conjugate. From this there results a
conjugate double pair for each of 2) and 4), and we have:
\[
\begin{array}{rl}
2) & x\,p,\ y\,q,\ \xi\eta,\ \xi'\eta' \\
4) & x\,r,\ y\,s,\ \xi\zeta,\ \xi'\zeta'.
\end{array}
\]
wherein $\zeta, \zeta'$ again denotes a pair of conjugate double tangents.

Now we form the real complex, syzygetic with 2) as well
as with 4),
\[
5)\quad \xi\xi',\ \eta\eta',\ \zeta\zeta',
\]
which can contain none of the elements $x, y, p, q, r, s$, and further the
two conjugate imaginary complexes
\[
\begin{array}{rl}
6) & x\xi,\ p\eta,\ r\zeta \\
7) & x\xi',\ p\eta',\ r\zeta',
\end{array}
\]
and 5), 6), 7) form an \gesperrt{azygetic triple}. $y, q, s$ do not occur in 6)
and 7), for otherwise they would have to be azygetic with two each
of the $x, p, r$, which by 2) and 4) is impossible.


% ===== scan leaf 417, printed 396 =====
% !PRINTED 396
\weberpage{396}
Thus each of the complexes 6), 7) is azygetic with the complex 1),
and hence one and only one of the two double tangents
$\varrho, \varrho'$ must occur in 6); if this is $\varrho$, then $\varrho'$ is contained in
7). If $\varrho\sigma$ is a pair of 6), then $\sigma$ is different from its
conjugate $\sigma'$, because otherwise $\sigma\varrho'$ would have to occur in 7) and consequently
$\varrho\varrho'$ in 5), which is not the case. Thus, if $\tau, \tau'$ are two further conjugate imaginary double tangents,
we have the complexes:
\[
\begin{array}{rl}
6) & x\xi,\ p\eta,\ r\zeta,\ \varrho\sigma,\ \sigma'\tau \\
7) & x\xi',\ p\eta',\ r\zeta',\ \varrho'\sigma',\ \sigma\tau'.
\end{array}
\]
If $t\lambda$ is the last pair of the complex 6), then the pair
$t'\lambda'$ occurs in 7), and these two pairs must contain a common
element. But since $\lambda$ cannot be equal to $t'$, because otherwise
both pairs would be identical, we must have $t = t'$ (or $\lambda = \lambda'$, which
is not essentially different); thus $t$ is real, and we
obtain:
\[
\begin{array}{rl}
6) & x\xi,\ p\eta,\ r\zeta,\ \varrho\sigma,\ \sigma'\tau,\ t\lambda \\
7) & x\xi',\ p\eta',\ r\zeta',\ \varrho'\sigma',\ \sigma\tau',\ t\lambda'.
\end{array}
\]
We now return to the syzygetic complex triple 1), 2), 3).
Since from 6) we conclude that $x\,p\,t$ and $x\,r\,t$ are azygetic,
$t$ can occur neither in 2) nor in 3), and must therefore be
contained in 1). But since 1) is a real complex, $t$ must appear in it
paired with a real double tangent $u$, and the
complex 1) becomes
\[
x\,y,\ p\,q,\ r\,s,\ t\,u,
\]
whereby our Theorem 3. is proved.

\gesperrt{4. If five real pairs occur in a complex, then
all the double tangents are real.}

Let us start from a complex with five real pairs
\[
1)\quad x_1y_1,\ x_2y_2,\ x_3y_3,\ x_4y_4,\ x_5y_5,
\]
and assume first that besides the last pair of this
complex there exists an imaginary double tangent $\xi$; then
a conjugate double tangent $\xi'$ is also present, and the real complex 2) determined by the
pair $\xi\xi'$ is syzygetic with 1) [because
$\xi$ and $\xi'$ do not occur in 1)]. But the four common
elements of 1) and 2) must be real
%CONT>>


% ===== scan leaf 418, printed 397 =====
% !PRINTED 397
\weberpage{397}
%<<CONT
be, because 2), as a real complex, cannot contain a real double tangent
paired with an imaginary one.

Let then
\[
2)\quad \xi\xi',\ x_1x_2,\ y_1y_2,
\]
and from this we derive the two complexes:
\[
\begin{aligned}
3)&\quad x_1\xi,\ x_2\xi'\\
4)&\quad x_1\xi',\ x_2\xi,
\end{aligned}
\]
which are syzygetic with one another, and hence have no element in
common besides the four indicated. But they are at the same time
conjugate, and therefore no real double tangent besides $x_1, x_2$ can
occur in 3). Now, however, 1) and 3) are azygetic, because
$\xi$ does not occur in 1); hence one element from each pair of 1)
must occur in 3). Thus 3) contains real elements besides $x_1, x_2$,
whereby the contradiction results.

We conclude from this that the two complexes which are determined
by the pairs $x_1x_2$, $y_1y_2$ and which form with 1) a syzygetic
triple of complexes contain \gesperrt{only real elements}.

If we let one of these complexes, say $x_1x_2$, take the place
of 1) and then repeat our argument, it follows
that all the double tangents whatsoever, and so also those of the
last pair of 1), must be real.

5. \gesperrt{If more than eight real double tangents
are present, then there are sixteen which form
four real pairs each in a syzygetic triple of
complexes.}

If we assume more than eight real double tangents, then by
the theorems 3., 4. we can assemble the following syzygetic triple
of complexes:
\[
\begin{aligned}
1)&\quad x_1y_1,\ x_2y_2,\ x_3y_3,\ x_4y_4\\
2)&\quad x_1x_2,\ y_1y_2,\ z_1z_2\\
3)&\quad x_1y_2,\ x_2y_1,
\end{aligned}
\]
in which the $x_i, y_i, z_i$ are real.

With this we combine the complex
\[
4)\quad x_1z_1,\ x_2z_2,
\]
which is azygetic with 1), and can therefore contain one and only
one element from each of the pairs $x_3y_3, x_4y_4$, say $x_3$ and $x_4$.
But since 4) is a real complex, there must be two further
%CONT>>


% ===== scan leaf 419, printed 398 =====
% !PRINTED 398
\weberpage{398}
\pdfbookmark[2]{§ 105. Proof of the existence of the four cases}{bm120}
%<<CONT
real double tangents $z_3, z_4$, so that the complex 4) continues
as follows:
\[
4)\quad x_1z_1,\ x_2z_2,\ x_3z_3,\ x_4z_4;
\]
$z_3$ and $z_4$ do not occur in 2), nor in 1) either, and
must therefore be contained in 3).

Neither can they occur paired in 3), because
$x_1, z_3, z_4$ are azygetic. From this, if $t_3, t_4$ are two
further real double tangents, we conclude that the complex 3)
is composed as follows:
\[
3)\quad x_1y_2,\ x_2y_1,\ z_3t_3,\ z_4t_4.
\]

If one now considers, in place of the complex 4), the complex
\[
5)\quad x_1z_3,\ y_2t_3,
\]
then one can conclude in exactly the same way that there is a
fourth real pair $u_1u_2$ in the complex 2), and with this 5. is proved.

If we gather the result together, we recognise that,
with regard to the reality of the double tangents of a real
curve of the fourth order, only four cases are possible:

1) \gesperrt{Four real syzygetic double tangents.}

2) \gesperrt{Eight real double tangents, namely four pairs
of a Steiner complex.}

3) \gesperrt{Sixteen real double tangents, which form four
real pairs each in a syzygetic complex triple.}

4) \gesperrt{Twenty-eight real double tangents.}

\begin{center}
\S.\,105.\\
\gesperrt{Proof of the existence of the four cases.}
\end{center}

In the preceding section we have at first proved only
that there can be no cases other than 1), 2), 3), 4).
But that these four cases are indeed all possible is now
also easy to show, on the basis of the theorem of \S.\,100, according
to which from seven arbitrarily given straight lines one can
derive, by a \gesperrt{rational way}, a curve of the fourth order for
which the given lines form an \gesperrt{Aronhold} system,
from which all the double tangents can be derived \gesperrt{rationally}.

Among the seven given straight lines imaginary ones may
also occur, and as a condition for the reality of the curve
%CONT>>


% ===== scan leaf 420, printed 399 =====
% !PRINTED 399
\weberpage{399}
%<<CONT
(i.e. of the coefficients in its equation) there follows this one, that
the system which one obtains, when every imaginary line is
replaced by its conjugate, is again an Aronhold system
of the same curve. For then the coefficients do not change
when $i$ is replaced by $-i$.

The curve therefore certainly becomes real if, in the given
system of seven, the conjugate of every imaginary line
occurs; then they form a \gesperrt{real system of seven}.

Such a real system can now contain:

1) one real and three pairs of conjugate imaginary lines;

2) three real and two pairs of conjugate imaginary lines;

3) five real and one pair of conjugate imaginary lines;

4) seven real lines.

We shall see that these four assumptions lead, in the same
order, to the four cases enumerated in the previous
section.

This demonstration becomes very simple if one makes use of the
mode of designating the double tangents which we set forth in \S.\,99,
according to which we can form the Steiner complexes
immediately from the designation.

1) In the first case we designate the given seven
lines by
\[
0,\ 1,\ 1',\ 2,\ 2',\ 3,\ 3',
\]
assume $0$ to be real, and $1$ conjugate imaginary with $1'$, $2$ with $2'$,
$3$ with $3'$. The remaining 21 double tangents are then designated
by the symbols $[01], [01'], [11'], \dots$.

According to the rule of \S.\,99 we obtain the two
following Steiner complexes:
\[
\begin{aligned}
(1)&\quad 0[01],\ 1'[11'],\ 2[12],\ 2'[12'],\ 3[13],\ 3'[13']\\
(1')&\quad 0[01'],\ 1[11'],\ 2[1'2],\ 2'[1'2'],\ 3[1'3],\ 3'[1'3'].
\end{aligned}
\]

The complex conjugate to (1) must contain the elements $0,1,2,2',3,3'$,
and is therefore, by \S.\,98, 7., identical with (1'), which
likewise contains these six elements.

From this it follows that $[11']$ is real, and that
\[
\begin{aligned}
&{}[01]\quad [01']\\
&{}[12]\quad [1'2']\\
&{}[12']\quad [1'2]
\end{aligned}
\]
%CONT>>


% ===== scan leaf 421, printed 400 =====
% !PRINTED 400
\weberpage{400}
%<<CONT
are conjugate pairs. Since in this consideration one may interchange 1 with 2
and 3, it follows that $0, [11'], [22'], [33']$
are real and all the remaining double tangents imaginary.

By \S.\,99 we can easily form a Steiner complex
that contains the four real double tangents:
\[
0[11'],\ [22'][33'],\ 1[01'],\ 1'[01],\ [23][2'3'],\ [23'][2'3],
\]
from which it is further to be seen that this complex contains, besides the real
pairs, two conjugate pairs and one conjugate double pair.

This is therefore the case 1) of the preceding section.

2) Let the seven given lines be:
\[
0,\ 1,\ 2,\ 3,\ 3',\ 4,\ 4',
\]
and among them $0, 1, 2$ real, $3$ conjugate to $3'$ and $4$ to $4'$. We
form three complexes:
\[
\begin{aligned}
(0)&\quad 1[01],\ 2[02],\ 3[03],\ 3'[03'],\ 4[04],\ 4'[04']\\
(3)&\quad 0[03],\ 1[13],\ 2[23],\ 3'[33'],\ 4[34],\ 4'[34']\\
(3')&\quad 0[03'],\ 1[13'],\ 2[23'],\ 3[33'],\ 4[3'4],\ 4'[3'4'].
\end{aligned}
\]

The complex conjugate to (0) contains $1, 2, 3, 3', 4, 4'$ and
is consequently identical with (0), i.e. the complex (0) is real.
From this it follows that $[01], [02]$ are real, and $[03]$ is conjugate to $[03']$ and $[04]$
to $[04']$. Further it results in the same
way that (3), (3') are conjugate complexes, and that therefore
$[33']$ is real, and $[34]$ with $[3'4']$, $[34']$ with $[3'4]$ are conjugate imaginary.
Since one may interchange 0, 1, 2 arbitrarily, and likewise 3
with 4, we obtain the following compilation:

\begin{center}
\gesperrt{Real double tangents:}
\end{center}
\[
0,\ 1,\ 2,\ [12],\ [20],\ [01],\ [33'],\ [44'].
\]

\begin{center}
\gesperrt{Conjugate pairs:}
\end{center}
\[
\begin{aligned}
3\ \ 4 &\quad [03][04][13][14][23][24][34][34']\\
3'\ 4' &\quad [03'][04'][13'][14'][23'][24'][3'4'][3'4].
\end{aligned}
\]

The Steiner complex that contains four real pairs
is here
\[
0[12],\ 1[20],\ 2[01],\ [33'][44'],\ [34][3'4'],\ [34'][3'4],
\]
and it contains in addition two conjugate pairs. This is therefore the
case 2) of \S.\,104.


% ===== scan leaf 422, printed 401 =====
% !PRINTED 401
\weberpage{401}

3) Let the given lines be:
\[
0,\ 1,\ 2,\ 3,\ 4,\ 5,\ 5',
\]
and of these let 0, 1, 2, 3, 4 be real, 5 and 5' conjugate imaginary. The
consideration of the three complexes:
\[
\begin{aligned}
(0)\quad & 1\,[1\,0],\ 2\,[2\,0],\ 3\,[3\,0],\ 4\,[4\,0],\ 5\,[5\,0],\ 5'\,[5'\,0] \\
(5)\quad & 0\,[0\,5],\ 1\,[1\,5],\ 2\,[2\,5],\ 3\,[3\,5],\ 4\,[4\,5],\ 5'\,[5'\,5] \\
(5')\quad & 0\,[0\,5'],\ 1\,[1\,5'],\ 2\,[2\,5'],\ 3\,[3\,5'],\ 4\,[4\,5'],\ 5\ [5\,5'],
\end{aligned}
\]
of which the first is real and the two others conjugate imaginary,
and of those derived from (0) by permutation of 0, 1, 2, 3, 4,
gives, exactly as above, the following result:

\begin{center}
\gesperrt{Real double tangents:}
\end{center}
\[
\begin{aligned}
&0,\quad 1,\quad 2,\quad 3,\quad 4,\quad [0\,1],\ [0\,2],\ [0\,3] \\
&[0\,4]\ [1\,2],\ [1\,3]\ [1\,4],\ [2\,3],\ [2\,4],\ [3\,4],\ [5\,5'].
\end{aligned}
\]

\begin{center}
\gesperrt{Conjugate pairs:}
\end{center}
\[
\begin{aligned}
&5,\ [0\,5],\ [1\,5],\ [2\,5],\ [3\,5],\ [4\,5] \\
&5',\ [0\,5'],\ [1\,5'],\ [2\,5'];\ [3\,5'],\ [4\,5'], % [sic] semicolon printed where a comma is expected
\end{aligned}
\]
and we find a syzygetic triple of complexes:
\[
\begin{aligned}
&0,\quad 1,\quad [0\,2][1\,2],\ [0\,3][1\,3],\ [0\,4][1\,4\,],\ [0\,5][1\,5\,],\ [0\,5'][1\,5'] \\
&2,\quad 3,\quad [0\,2][0\,3],\ [1\,2][1\,3],\ [4\,2][4\,3\,],\ [2\,5][3\,5\,],\ [2\,5'][3\,5'] \\
&[0\,1],\ [2\,3],\ [0\,2][1\,3],\ [0\,3][1\,2],\quad 4\ \ [5\,5'],\quad 5\ \ [4\,5'],\quad 5'\ \ [4\,5],
\end{aligned}
\]
of which each complex contains four further real pairs and one
imaginary double pair. But several such triples can be formed,
since 0, 1, 2, 3, 4 may be permuted arbitrarily.
This is the third case of \S.\,104.

4) That, finally, when we suppose all elements of the given
system of seven to be real, all the remaining double tangents
also turn out real, is an immediate consequence of the rational
representation (\S.\,100).

With this it is shown that the four cases of the preceding
section can indeed all occur. Whether the mode of generation
discussed here is the only possible one — in other
words, whether for every real curve of the fourth order a real
\gesperrt{Aronhold} system exists — this question we must leave
undecided.


% ===== scan leaf 423, printed 402 =====
% !PRINTED 402
\weberpage{402}

It must also be pointed out here that from the reality
of a double tangent the reality of the points of contact
by no means follows, because these points of contact depend
in turn on a quadratic equation. The real double tangents
thus split again into two kinds, those with real and
those with imaginary points of contact. What possible
cases are to be distinguished here — this question we do not
pursue further \textsuperscript{1}).

\begin{center}
\rule{0.3\linewidth}{0.4pt}
\end{center}

\textsuperscript{1}) Cf. on the whole question of the reality of the double tangents
from the geometrical point of view: \gesperrt{Zeuthen}, ``Sur les diff\'erentes formes
des courbes planes du quatri\`eme ordre''. Mathem. Ann., Bd. VII (1873).


% ===== scan leaf 424, printed 403 =====
\weberpage{403}% page number not printed on this page
\pdfbookmark[1]{Thirteenth Section. General Theory of the Equation of the Fifth Degree.}{bm121}
\pdfbookmark[2]{§ 106. Statement of the problem}{bm122}
% !PRINTED none

\begin{center}
\gesperrt{Thirteenth Chapter.}

\medskip
\gesperrt{General Theory of the Equation of the Fifth Degree.}

\medskip
\rule{0.3\linewidth}{0.4pt}

\medskip
\S.\,106.

\gesperrt{Statement of the Question.}
\end{center}

We have seen in \S.\,44 that beyond the question of the
\gesperrt{Galois} group of an algebraic equation there is a further-reaching
problem, namely the question of the linear substitution group
of the smallest possible number of dimensions to whose form
problem the given equation can be reduced. For the metacyclic
equations, and so in particular for the general equations of the
3rd and 4th degree, this number of dimensions is equal to 1,
which is just what is expressed by saying that these equations
are solvable by radicals. The equations next to be investigated
are then those of the 5th degree, and it will appear that the
solution of the general equation of the 5th degree leads to a
binary linear substitution group, namely to the icosahedron
problem.

But connected with this there is yet another question.
The general equation of the 5th degree contains five coefficients,
which can be regarded as independent variables, and
consequently the root of such an equation is an algebraic
function of five variables. Now one can, however, already
diminish the number of these variables by the simple linear
substitutions. By the \gesperrt{Tschirnhausen} transformation, e.\,g.
to the \gesperrt{Jerrard} form, one can even make the equation depend
on only one variable coefficient (one \gesperrt{parameter}), and one can
thus, through the mediation of equations which do not attain
the 5th degree, make the root of a general equation of the
5th degree depend on an algebraic function of one variable.


% ===== scan leaf 425, printed 404 =====
% !PRINTED 404
\weberpage{404}
\pdfbookmark[2]{§ 107. Lüroth's theorem}{bm123}

The question to which we are led here is thus how one can
reduce the solution of an algebraic equation, whose coefficients
depend on a certain number of variables, to an equation
with the smallest possible number of parameters, and in
particular, under what circumstances and by what means
one can reduce it to equations with only one
parameter.

We accordingly now pose the question, under what
assumptions, for a \gesperrt{general equation of degree $n$,
resolvents with only one parameter exist.}

Let $x_0, x_1, \ldots, x_{n-1}$ be a system of $n$ independent variables,
\weq{1}{u = \Phi(x_0, x_1, \ldots, x_{n-1})}
a rational function of these variables which, under the permutations
of the \gesperrt{alternating group} of the $n$ letters $x$, passes over into the
mutually distinct functions
\weq{2}{u,\ u_1, u_2, \ldots u_{\nu-1}}
where $\nu$ may be equal to or less than the degree of the alternating
group, and $z$ is a function of the variables $x$
which remains unchanged under the alternating group. The
question is when a rational equation of degree $\nu$ with respect
to $u$ exists,
\weq{3}{F(u, z) = 0,}
which is identically satisfied by the $\nu$ functions (2)? The
coefficients in (3) must be independent of the $x$ and hence pure
numbers. For the present we let no restriction of the domain of rationality
enter in the region of numbers.

Equation (3) is, when $\nu > 1$, a resolvent of the
equation of degree $n$ whose roots are the $x$, which depends only on the
\gesperrt{one parameter} $z$. We exclude the case $\nu = 1$.

\begin{center}
\S.\,107.

\gesperrt{Theorem of L\"uroth.}
\end{center}

First the following lemma is to be proved \textsuperscript{1}):

1. \gesperrt{If $u, u_1, u_2, \ldots, u_{\nu-1}$ is a system of rational functions
 of one variable $t$, then there is a rational
function of the $u, u_1, \ldots, u_{\nu-1}$:}
\[
\vartheta = \chi(u, u_1, \ldots, u_{\nu-1})
\]

\begin{center}
\rule{0.3\linewidth}{0.4pt}
\end{center}

\textsuperscript{1}) \gesperrt{L\"uroth}, Mathematische Annalen, Bd. IX.
%CONT>>


% ===== scan leaf 426, printed 405 =====
% !PRINTED 405
\weberpage{405}
%<<CONT
\gesperrt{of such a kind that $u, u_1, \ldots, u_{\nu-1}$ can be expressed
rationally through $\vartheta$ alone.}

To prove it, we set

\weq{1}{u = \dfrac{\varphi(t)}{\psi(t)}, \quad u_1 = \dfrac{\varphi_1(t)}{\psi_1(t)}, \quad \cdots, \quad u_{\nu-1} = \dfrac{\varphi_{\nu-1}(t)}{\psi_{\nu-1}(t)},}

and understand by $\varphi(t)$, $\psi(t)$ integral functions of $t$
without common divisor; likewise by $\varphi_1(t)$,
$\psi_1(t)$, and so on. Moreover we may further assume that
none of the functions $u, u_1, \ldots, u_{\nu-1}$ is a constant.
Alongside the variable $t$ we introduce a second variable $\tau$ and
set:

\weq{2}{v = \dfrac{\varphi(\tau)}{\psi(\tau)}, \quad v_1 = \dfrac{\varphi_1(\tau)}{\psi_1(\tau)}, \quad \cdots, \quad v_{\nu-1} = \dfrac{\varphi_{\nu-1}(\tau)}{\psi_{\nu-1}(\tau)}.}

Now we form the integral functions of the two variables $t, \tau$,
of which none vanishes identically:

\weq{3}{\begin{aligned}
\varphi(t)\,\psi(\tau) - \varphi(\tau)\,\psi(t) &= \Phi(t,\tau) = -\Phi(\tau,t)\\
\varphi_1(t)\,\psi_1(\tau) - \varphi_1(\tau)\,\psi_1(t) &= \Phi_1(t,\tau) = -\Phi_1(\tau,t)\\
&\cdots\cdots\cdots\cdots\cdots\cdots\cdots\cdots\cdots\cdots\cdots\cdots\cdots\cdots\\
\varphi_{\nu-1}(t)\,\psi_{\nu-1}(\tau) - \varphi_{\nu-1}(\tau)\,\psi_{\nu-1}(t) &= \Phi_{\nu-1}(t,\tau) = -\Phi_{\nu-1}(\tau,t).
\end{aligned}}

If we regard them as functions of $t$, they all vanish
when $t = \tau$, and hence they have a greatest
common divisor, which with respect to $t$ is at least of the
first degree, and which shall be of degree $\mu$:

\weq{4}{R = R(t, \tau).}

This function $R$ is rational also with respect to $\tau$, and if
one arranges it so that $\tau$ occurs neither in the denominator nor in a
superfluous factor, then $R(t, \tau)$ can at most change its sign
under the interchange of $t$ and $\tau$.
For the functions (3), regarded also as functions of the two
variables $t, \tau$, are divisible by $R(t, \tau)$, and, since they
are alternating, by $R(\tau, t)$ (Bd.\,I, \S.\,51). Hence
$R(t, \tau)$ is divisible by $R(\tau, t)$, and conversely, and
consequently the two differ only by a constant factor,
which must be $= \pm 1$, because interchanging $t$ with
$\tau$ once more restores the original function.

It can further be proved that none of the functions (3) (for
indeterminate $\tau$), regarded as a function of $t$, contains a factor
%CONT>>


% ===== scan leaf 427, printed 406 =====
% !PRINTED 406
\weberpage{406}
%<<CONT

more than once, and that in consequence of this $R\,(t, \tau)$ also is
divisible by no square. For suppose that
\[
\begin{aligned}
\varPhi\,(t,\tau) &= \varphi\,(t)\,\psi\,(\tau) - \varphi\,(\tau)\,\psi\,(t) \\
\varPhi'(t,\tau) &= \varphi'(t)\,\psi\,(\tau) - \varphi\,(\tau)\,\psi'(t)
\end{aligned}
\]
had, as functions of $t$, a common divisor $P$; then $P$ would also
have to be a divisor of
\[
\psi'(t)\,\varPhi - \psi(t)\,\varPhi' = [\varphi\,(t)\,\psi'(t) - \psi\,(t)\,\varphi'(t)]\,\psi\,(\tau)
\]
Hence $P$ would be a divisor of $\varphi\,(t)\,\psi'(t) - \psi\,(t)\,\varphi'(t)$, and
could therefore be assumed to be independent of $\tau$.

If we now take two values $\tau_1, \tau_2$ of $\tau$ such that
$\varPhi\,(\tau_2, \tau_1)$ becomes different from zero, then by (3)
\[
\varPhi\,(t,\tau_1)\,\varphi\,(\tau_2) - \varPhi\,(t,\tau_2)\,\varphi\,(\tau_1) = \varphi\,(t)\,\varPhi\,(\tau_2,\tau_1).
\]
In this the left side is divisible by $P$, and hence
$\varphi\,(t)$, and consequently $\psi\,(t)$, is also divisible by $P$, which contradicts
 the assumption that these two functions are to be without
 common divisor.

From this it follows incidentally that $R\,(t,\tau) = -R\,(\tau,t)$ must hold,
since $R$ is divisible by $t - \tau$, but not by $(t-\tau)^2$.

To form the function $R\,(t,\tau)$, one may seek out the greatest
common divisor of the functions $\varphi_h(t) - v_h\,\psi_h(t)$
, from which it follows that $R\,(t,\tau)$, apart from a factor
independent of $t$, can be represented rationally in terms of $v, v_1, \ldots, v_{\nu-1}$.
If, then, $a$ and $b$ are any two fixed numerical
values,
\weq{5}{\frac{R\,(a,\tau)}{R\,(b,\tau)} = \vartheta = \chi\,(v, v_1, \ldots, v_{\nu-1})}
is a rational function of $v, v_1, \ldots, v_{\nu-1}$. Moreover, if $\xi$
denotes a new variable,
\weq{6}{R\,(a,\xi) - \vartheta\,R\,(b,\xi) = X}
is of degree $\mu$ with respect to $\xi$.

If $\tau$ is an arbitrary value, let the roots of the
equation $R\,(t,\tau) = 0$ be
\weq{7}{t = \tau, \tau', \tau'', \ldots}
so that for every index $h$
\[
\varPhi_h\,(\tau, \tau), \quad \varPhi_h\,(\tau', \tau), \quad \varPhi_h\,(\tau'', \tau), \ldots
\]
vanish.


% ===== scan leaf 428, printed 407 =====
% !PRINTED 407
\weberpage{407}
\pdfbookmark[2]{§ 108. Resolvents with one parameter}{bm124}
If $\tau', \tau''$ are any two of these roots, then from (3) it follows that:
\[
\begin{aligned}
\varphi_h\,(\tau')\,\psi_h\,(\tau) - \varphi_h\,(\tau)\,\psi_h\,(\tau') &= 0\\
\varphi_h\,(\tau'')\,\psi_h\,(\tau) - \varphi_h\,(\tau)\,\psi_h\,(\tau'') &= 0,
\end{aligned}
\]
and from this, since $\psi_h\,(\tau)$ and $\varphi_h\,(\tau)$ cannot vanish
at the same time,
\[
\varphi_h\,(\tau')\,\psi_h\,(\tau'') - \varphi_h\,(\tau'')\,\psi_h\,(\tau') = 0,
\]
i.\,e., if $\tau', \tau''$ are any two of the quantities (7), we have
\[
\Phi_h\,(\tau', \tau'') = 0.
\]

From this it follows that the $\nu$ functions $\Phi_h\,(t,\tau')$, apart from a
factor independent of $t$, have the same greatest common divisor as the
functions $\Phi_h\,(t,\tau)$. The same holds for the functions
$\Phi_h\,(t,\tau'')$ and so on, or the equations
\[
R\,(t,\tau) = 0, \quad R\,(t,\tau') = 0, \quad R\,(t,\tau'') = 0, \ldots
\]
all have the same roots. From this it follows, by (5):
\[
\vartheta = \dfrac{R\,(a,\tau)}{R\,(b,\tau)} = \dfrac{R\,(a,\tau')}{R\,(b,\tau')} = \dfrac{R\,(a,\tau'')}{R\,(b,\tau'')} = \cdots,
\]
and consequently, by (6), the $\mu$ values
\[
\xi = \tau, \tau', \tau'', \ldots
\]
are the roots of the equation $X = 0$. On the other hand, from
$\Phi_h\,(\tau',\tau) = 0$, $\Phi_h\,(\tau'',\tau) = 0, \ldots$ it follows that
\[
v_h = \dfrac{\varphi_h\,(\tau)}{\psi_h\,(\tau)} = \dfrac{\varphi_h\,(\tau')}{\psi_h\,(\tau')} = \dfrac{\varphi_h\,(\tau'')}{\psi_h\,(\tau'')} \cdots = \dfrac{1}{\mu}\,\Sigma\,\dfrac{\varphi_h\,(\tau)}{\psi_h\,(\tau)},
\]
and hence $v_h$, as a symmetric function of the roots of (6), can be
represented rationally in terms of the coefficients of that equation,
i.\,e. rationally in terms of $\vartheta$. If we again interchange $t$ with
$\tau$, we obtain by (1) and (2) expressions of the form

\weq{10}{\begin{aligned}
u &= f\,(\vartheta), \quad u_1 = f_1\,(\vartheta), \ldots, u_{\nu-1} = f_{\nu-1}(\vartheta)\\
\vartheta &= \chi\,(u, u_1, \ldots, u_{\nu-1}),
\end{aligned}}

in which $f, f_1, \ldots, f_{\nu-1}, \chi$ denote rational functions.

\begin{center}
\S.\ 108.\\
\gesperrt{Resolvents with one parameter.}
\end{center}

We now return to our initial assumptions (\S.\,106), and denote by
$u, u_1, \ldots, u_{\nu-1}$ a system
%CONT>>


% ===== scan leaf 429, printed 408 =====
% !PRINTED 408
\weberpage{408}
%<<CONT
rational functions of $x_0, x_1, \ldots, x_{n-1}$ which arise from one of
them by the permutations of the alternating group, and are roots of the
equation

\weq{1}{F\,(u, z) = 0}

where $z$ is a rational function of the $x$ which remains unchanged by the
permutations of the alternating group.

We prove the following second lemma:

\gesperrt{2. If the rational function $\Psi\,(u, u_1, \ldots, u_{\nu-1})$
vanishes identically when for the
\[
x_0, x_1, \ldots, x_{n-1}
\]
certain rational functions of a variable $t$}

\weq{2}{x_h = \varphi_h\,(t)}

\gesperrt{are put, and if by this substitution $z$ does not become
independent of $t$, then $\Psi$ vanishes identically also as a function of
the independent variables $x_0, x_1, \ldots, x_{\nu-1}$.}
% [sic] The source prints x_{\nu-1} here, although the independent variables were introduced through x_{n-1}.

The Galois group of equation (1) in the field of the rational functions of
$z$ will consist of certain permutations of the roots
$u, u_1, \ldots, u_{\nu-1}$. If we carry out these permutations in the
function $\Psi$, and form the product

\weq{3}{\Pi\,\Psi\,(u, u_1, \ldots, u_{\nu-1}) = f\,(z)}

of all the functions so obtained, there results a rational function
$f\,(z)$ of $z$. Now if one of the factors of the product (3) vanishes
identically after the substitution (2), and $z$ is not independent of $t$,
then $f\,(z)$ must be identically equal to zero, and consequently one of
the factors of the product (3) must also vanish identically as a function
of the $x$. But if one of these factors is identically equal to zero, then
all the other factors vanish as well, because in every rational equation
between the $u$ one can carry out all the permutations of the Galois
group. With this the theorem 2. is proved.

This theorem, taken together with the theorem 1. proved in the preceding
section, now gives the following result:

\gesperrt{3. If $u, u_1, \ldots, u_{\nu-1}$ are the roots of an equation
(1) depending on a parameter $z$, and if $u, u_1, \ldots,
u_{\nu-1}, z$ are rational functions of the}
%CONT>>


% ===== scan leaf 430, printed 409 =====
% !PRINTED 409
\weberpage{409}
%<<CONT
\gesperrt{independent variables $x_0, x_1, \ldots, x_{n-1}$, one can
set}
\weq{4}{u = f(\vartheta),\ u_1 = f_1(\vartheta), \ldots, u_{\nu-1} = f_{\nu-1}(\vartheta)}
\gesperrt{where $f, f_1, \ldots, f_{\nu-1}$ are rational functions
of $\vartheta$ and}
\weq{5}{\vartheta = \chi(u, u_1, \ldots, u_{\nu-1}) = \psi(x_0, x_1, \ldots, x_{n-1})}
\gesperrt{is a rational function of the $u$, or also of the
$x$.}

For, by Theorem 1., if we replace
$x_0, x_1, \ldots, x_{n-1}$ by rational functions of a variable $t$,
in such a way that $z$ does not become independent of $t$, we can first represent the functions
$u_h$ by the formulae (4), (5). The relations
\[
u_h = f_h(\vartheta) = f_h[\chi(u, u_1, \ldots, u_{\nu-1})]
\]
are then identically satisfied with respect to $t$, and must therefore,
by Theorem 2., also be identical in the variables $x$,
q. e. d.

4.~\gesperrt{If of two variables $\vartheta, \vartheta_1$ each is a ratio\-
nal function of the other, then they are linear
functions of one another.}

For the proof we assume that
\weq{6}{\vartheta = \frac{\varphi(\vartheta_1)}{\psi(\vartheta_1)},\ \vartheta_1 = \frac{\varphi_1(\vartheta)}{\psi_1(\vartheta)},}
and understand by $\varphi, \psi$ two integral functions without common
 divisor, and likewise by $\varphi_1, \psi_1$. If we arrange the second of the
equations (6) in the form $\varphi_1(\vartheta) - \vartheta_1\psi_1(\vartheta) = 0$ according to powers
 of $\vartheta$, it may take the form:
\[
a_0\vartheta^m + a_1\vartheta^{m-1} + \cdots + a_{m-1}\vartheta + a_m = 0,
\]
in which the coefficients $a_0, a_1, \ldots, a_m$ are \gesperrt{integral linear functions}
 of $\vartheta_1$. If we substitute in it for $\vartheta$ the value
from the first of the equations (6), there follows the equation, identical
 with respect to $\vartheta_1$,
\[
a_0\varphi^m + a_1\varphi^{m-1}\psi + \cdots + a_{m-1}\varphi\psi^{m-1} + a_m\psi^m = 0.
\]
Accordingly $a_0\varphi^m$ must be divisible by $\psi$, and since $\varphi$
and $\psi$ are assumed to be relatively prime, $a_0$ must be divisible by $\psi$,
and hence $\psi$ constant or linear. In the same way we conclude
that $\varphi$ must be constant or linear, and since not both
%CONT>>


% ===== scan leaf 431, printed 410 =====
% !PRINTED 410
\weberpage{410}
\pdfbookmark[2]{§ 109. Group of the resolvents with one parameter}{bm125}
%<<CONT
functions can be constant, then by (6) Theorem 4.
is proved.

We apply this theorem to the function $\vartheta$ occurring in 3.
If we perform with the variables $x$ any permutation
of the alternating group, the functions
$u, u_1, \ldots, u_{\nu-1}$ likewise undergo a permutation. By
(5), $\vartheta$ passes under this permutation into another function $\vartheta_1$,
which by (4) is rationally expressible in terms of $\vartheta$. But
$\vartheta$ is likewise rationally expressible in terms of $\vartheta_1$, since in
Theorem 3. one may replace $\vartheta$ by $\vartheta_1$, and $\vartheta$ also arises from $\vartheta_1$
by a permutation of the alternating group. Hence
by 4., $\vartheta$ and $\vartheta_1$ are \gesperrt{linear functions} of one another.
From this the following result is obtained:

5.~\gesperrt{If for a general equation of degree $n$
a resolvent with one parameter exists,
then there is a rational function $\vartheta$
of $n$ variables $x$ which, under the permutations
of the alternating group, passes into a linear function
of itself}
\weq{7}{\chi(\vartheta) = \frac{a\vartheta + b}{c\vartheta + d}}
\gesperrt{where $a, b, c, d$ are constants, i.\,e.\ numbers.}

\begin{center}
\S.\,109.

\gesperrt{Group of the resolvents with one parameter.}
\end{center}

The function $\vartheta$ that occurs in Theorem 5. of the preceding section
is, by Theorem 3., itself the root of a resolvent
with \gesperrt{one parameter} $\Phi(\vartheta, z) = 0$. This equation,
since every root is rationally expressible in terms of every other,
is a normal equation, and its \gesperrt{Galois} group is isomorphic
with the group of the linear substitutions $\chi(\vartheta)$ of
Theorem 5. Since one can now carry out in the identical equation $\Phi(\vartheta, z) = 0$
all the permutations of the alternating group $A$ of the variables $x$,
whereby $z$ does not change, while $\vartheta$ passes into every
other root of $\Phi$, the group of the linear
substitutions $\chi(\vartheta)$, i.\,e.\ the \gesperrt{Galois} group of the equa-
%CONT>>


% ===== scan leaf 432, printed 411 =====
% !PRINTED 411
\weberpage{411}
%<<CONT
tion $\Phi$, isomorphic (simply or multiply) with the alternating permutation
group of $n$ figures. Let the group of the linear substitutions
 $\chi(\vartheta)$ be denoted by $L$. Since it is
finite, it must be identical with one of the polyhedral groups
considered in the seventh chapter.

For simplification we now make use of the theorem proved
in \S.\,51, 2., according to which the group $L$ can be
so transformed that an arbitrary one of the non-identical
substitutions of $L$ becomes a multiplication. But a transformation
 of the group $L$ is equivalent to this, that for $\vartheta$
a linear function of $\vartheta$ is put, to which the same property
 belongs as to the original function $\vartheta$.

If, then, we denote by $\vartheta_\pi$ the function of the variables
$x_0, x_1, \ldots, x_{n-1}$ that arises from $\vartheta$ by application of the permutation $\pi$,
we can choose $\vartheta$ so that for a definite,
but arbitrary permutation $\pi$
\[
\vartheta_\pi = \varepsilon\vartheta
\]
holds. If one applies $\pi$ repeatedly, it follows that
\[
\vartheta_{\pi^2} = \varepsilon^2\vartheta,\ \vartheta_{\pi^3} = \varepsilon^3\vartheta, \ldots,
\]
and if therefore $p$ is the degree of the permutation $\pi$, then $\varepsilon$ is a
$p$th root of unity.

If first $n = 3$, the alternating group consists of
the powers of the cyclic permutation $\gamma = (0, 1, 2)$. We
may assume that the substitution corresponding to the permutation,
 $\chi(\vartheta)$, is multiplicative, that is, that $\vartheta_\gamma = \varepsilon\vartheta$, where $\varepsilon$
is a third root of unity. Hence $\vartheta^3$ is an alternating
(or symmetric) function. As resolvent with one parameter
 we thus obtain the pure equation
\[
\vartheta^3 = z,
\]
where $z$ is an alternating function. We may take for $\vartheta$,
say, the \gesperrt{Lagrange} resolvent $x_0 + \varepsilon x_1 + \varepsilon^2 x_2$, and
obtain the well-known reduction of the cubic equation to
a pure equation (Bd. I, \S.\,159).

We pass to the case $n = 4$, in which the alternating
 group $A$ contains the permutations
\[
1,\ \alpha_1 = (0, 1)(2, 3),\ \alpha_2 = (0, 2)(1, 3),\ \alpha_3 = (0, 3)(1, 2)
\]
which form a four-group $B$; besides these there occur
%CONT>>


% ===== scan leaf 433, printed 412 =====
% !PRINTED 412
\weberpage{412}
%<<CONT
in $A$ there are also eight cyclic permutations of three figures each,
$\gamma = (0, 1, 2) \ldots$, and $A$ can be represented thus:
\[
A = B + B\gamma + B\gamma^2.
\]
Any one permutation $\alpha$ together with a permutation $\gamma$ may be
regarded as generators of the group. For if, say,
\[
\gamma = (0, 1, 2),\ \alpha_1 = (0, 1)(2, 3),
\]
then
\weq{1}{\alpha_1\gamma\alpha_1 = \gamma' = (0, 3, 1),}
and from $(0, 1, 2), (0, 3, 1)$ the whole group $A$ can be derived
(Bd. I, \S.\,153, 7.).

Likewise one may take $\gamma$ and $\gamma'$ as generating permutations
of $A$. In particular,
\weq{2}{\gamma'\gamma\gamma' = \alpha_1.}
We choose $\vartheta$ so that to the permutation $\gamma$ there corresponds a multiplication
$\varepsilon\vartheta$, in which, since $\gamma$ is of degree 3, $\varepsilon$ denotes a third root of unity.
 We set
\weq{3}{\vartheta = \frac{\varphi(x_0, x_1, x_2, x_3)}{\psi(x_0, x_1, x_2, x_3)},}
and understand by $\varphi, \psi$ two integral functions of the four
variables $x$ without common divisor.

From the identical equation
\[
\frac{\varphi_\gamma}{\psi_\gamma} = \varepsilon\,\frac{\varphi}{\psi}
\]
it then follows that
\weq{4}{\varphi_\gamma = \varepsilon_1\varphi,\ \psi_\gamma = \varepsilon_2\varphi} % [sic]
must hold, in which $\varepsilon_1, \varepsilon_2$ are likewise third roots of unity.

We apply one of the permutations $\alpha$ to $\vartheta$, and obtain
 from $\vartheta_\alpha = \chi(\vartheta)$:
\weq{5}{\frac{\varphi_\alpha}{\psi_\alpha} = \frac{a\varphi + b\psi}{c\varphi + d\psi}.}
Now since $\varphi, \psi$, hence also $a\varphi + b\psi$ and $c\varphi + d\psi$, and
likewise $\varphi_\alpha, \psi_\alpha$, are without common divisor, in the
identity (5) the numerator must equal the numerator and the denominator the
denominator, at least up to a constant factor.
This constant factor we can absorb into the constants $a, b, c, d$
and obtain
\[
\text{6)}\quad \varphi_\alpha = a\varphi + b\psi,\ \psi_\alpha = c\varphi + d\psi.
\]


% ===== scan leaf 434, printed 413 =====
% !PRINTED 413
\weberpage{413}
If we here take $\alpha=\alpha_1$ and $\alpha=\alpha_2$, and set for
brevity $\varphi_{\alpha_1}=\varphi_1$, $\varphi_{\alpha_2}=\varphi_2$, then
from the two equations
\[
\varphi_1 = a_1\varphi + b_1\psi, \quad \varphi_2 = a_2\varphi + b_2\psi
\]
we can eliminate $\psi$ and obtain an equation of the form:
\weq{7}{h\varphi + h_1\varphi_1 + h_2\varphi_2 = 0,}
in which $h$, $h_1$, $h_2$ are constants, of which at least two are
different from zero. To the identical equation (7) we can now apply the
permutations $\alpha_1$, $\alpha_2$, and obtain, since
$\alpha_1\alpha_2 = \alpha_2\alpha_1 = \alpha_3$, $\alpha_1\alpha_1 = 1$,
$\alpha_2\alpha_2 = 1$,
\[
\begin{array}{l|l}
h\varphi + h_1\varphi_1 + h_2\varphi_2 = 0 & h,\\
h\varphi_1 + h_1\varphi + h_2\varphi_3 = 0 & h_1,\\
h\varphi_2 + h_1\varphi_3 + h_2\varphi = 0 & -h_2,
\end{array}
\]
from which, by multiplication with the factors standing beside them and
addition:
\weq{8}{(h^2 + h_1^2 - h_2^2)\varphi + 2hh_1\varphi_1 = 0.}
If, then, $h$ and $h_1$ are different from zero, $\varphi$ and $\varphi_1$
differ only by a constant factor (which must be $=\pm 1$, since $\alpha_1$
is of degree 2). But if $h$ or $h_1=0$, then from (7) it follows that
$\varphi_1=\pm\varphi_2$ or $\varphi=\pm\varphi_2$, and the application of
$\alpha_1$ to the first equation gives $\varphi=\pm\varphi_3$. Hence in
every case \gesperrt{one of the relations} holds:
\weq{9}{\varphi = \pm\varphi_1, \quad \varphi = \pm\varphi_2, \quad \varphi = \pm\varphi_3.}
Since now one can equally well regard $\gamma,\alpha_1$ as $\gamma,\alpha_2$,
as also $\gamma,\alpha_3$, as generating elements of the group $A$, this
result, taken together with (4), yields the theorem:

1. \gesperrt{The function $\varphi$ has the property of changing, under all
permutations of the alternating group, only by a constant factor.}

But the very same conclusion is applicable to $\psi$ as well, and hence
also to the quotient $\vartheta$ of the two functions. Accordingly, for
every permutation $\pi$ of the alternating group,
\weq{10}{\vartheta_\pi = \varepsilon\vartheta.}
Herein, when $\pi$ belongs to the cyclic permutations $\gamma$,
$\varepsilon$ is a third root of unity. But from this it is further to be
concluded, because all permutations $\pi$ are composed of $\gamma$, $\gamma'$
%CONT>>


% ===== scan leaf 435, printed 414 =====
% !PRINTED 414
\weberpage{414}
%<<CONT
can be so chosen that the constant $\varepsilon$ occurring in (10) is \gesperrt{always}
a third root of unity. If $\pi$ belongs to the $\alpha$, then $\varepsilon$
is at the same time a second root of unity and must therefore be $=1$.
From this we infer the theorem:

2. \gesperrt{The function $\vartheta$ remains unchanged under the
permutations of the four-group $A$.}

The third power of $\vartheta$ admits the permutations of the alternating
group, and $\vartheta$ is therefore the root of a pure cubic equation
$\vartheta^3 = z$. This, then, is the resolvent $\Phi(\vartheta,z) = 0$,
which in this case is a \gesperrt{partial resolvent} (Bd. I, \S.\,161).

From these considerations there now results without difficulty the following theorem,
whose proof is the goal of this whole
investigation:

3. \gesperrt{A general equation of degree higher than
4 has no resolvents with one parameter.}

For if $n$ is greater than 4, then, if we regard $\vartheta$ as a function
of any four of the variables $x$, we can apply Theorem 2.
It then follows that $\vartheta$ remains unchanged when any
pair of transpositions is performed among the variables. Since one
can now compose the whole alternating group out of transposition pairs
(Bd. I, \S.\,153, 8.), it follows that $\vartheta$ remains unchanged by
the alternating group altogether, and is therefore an alternating or
a symmetric function. $\Phi(\vartheta,z) = 0$ reduces to
the identity $\vartheta = z$ and is not a resolvent.

Let it further be remarked that, if the symmetric group is put in place
of the alternating group, this theorem holds a fortiori.

The theorem whose proof we have developed here was first
stated by \gesperrt{Kronecker}, without his publishing a proof
of it. The first proof was given by \gesperrt{F. Klein}, which
was then further simplified by \gesperrt{Gordan}\footnote{\gesperrt{F. Klein}, Vorlesungen über das Ikosaëder. The
proof by \gesperrt{Gordan}, which our presentation follows in its main
features, is found in Bd. 29 der mathematischen Annalen (1887).}.


% ===== scan leaf 436, printed 415 =====
% !PRINTED 415
\weberpage{415}
\pdfbookmark[2]{§ 110. The icosahedral equation}{bm126}
\begin{center}
\S.\,110.

\gesperrt{The Icosahedral Equation.}
\end{center}

The result of these considerations is, in the first place, so far as it
concerns the general equations of higher than the 4th degree, a negative
one. We do not succeed, by the rational formation of resolvents, in
reducing the general equation of the 5th degree to an equation with \gesperrt{one}
parameter. If we nevertheless do not yet wish to give up this goal,
nothing remains but that we no longer regard the variables
$x_0, x_1, \ldots, x_{n-1}$ as completely independent variables, and that
we let certain equations subsist between them. We are then, to be sure,
no longer dealing with the general equation of the degree in question,
but with a more special one, in which the coefficients are not
independent of one another, and the second question presents itself, how
and by what irrational functions the general equation can be reduced to
this special one. If this succeeds by means of equations of lower degree,
e.g. by square roots, a reduction of the problem will at any rate have
been achieved.

We therefore now assume that the variables $x_0, x_1, \ldots, x_{n-1}$
are connected with one another by an arbitrary number of relations
\weq{1}{A(x_0, x_1, \ldots, x_{n-1}) = 0, \quad B(x_0, x_1, \ldots, x_{n-1}) = 0, \ldots}
in which the $A, B, \ldots$ are rational functions of the $x$, which are
so constituted in number and mutual dependence that the values of the $x$
are not numerically fixed, but are still variable. Moreover, we take as
a basis a permutation group $\mathfrak{A}$ of the
$x_0, x_1, \ldots, x_{n-1}$, and regard the functions of the $x$ that
admit the permutations of $\mathfrak{A}$, and only these, as rational,
so that $\mathfrak{A}$ is the Galois group of the equation whose roots
are the $x_0, x_1, \ldots, x_{n-1}$. \gesperrt{The relations (1) must then be so
constituted that they too admit the permutations $\mathfrak{A}$.}

Under these presuppositions we ask after the possibility of determining
two rational functions $u, z$ of $x_0, x_1, \ldots, x_{n-1}$ in such a
way that between them an equation
\weq{2}{F(u, z) = 0}
%CONT>>


% ===== scan leaf 437, printed 416 =====
% !PRINTED 416
\weberpage{416}
%<<CONT
consists, which is satisfied for all values of the $x$ that satisfy
the relations (1), and in which $z$, always with regard to the
relations (1), remains unchanged when the $x$ are subjected to the
permutations of the group $\mathfrak{A}$. The function $u$ is thereby
to pass into $\nu$ distinct functions
\weq{3}{u, u_1, \ldots, u_{\nu-1}}
so that the quantities (3) are the roots of the equation (2), and,
when $\nu > 1$, (2) is a resolvent of the equation for the $x$.

In this generality we have as yet no point of attack for the solution
of the question. We therefore add a further restrictive assumption,
namely: \gesperrt{It is to be possible to set, for the $x_h$, rational
functions of a variable $t$}
\weq{4}{x_h = \varphi_h(t)}
\gesperrt{so that the relations (1) are identically satisfied, and that at the
same time $z$ does not become independent of
$t$}\footnote{If one regards the $x$ and the $u$ as the coordinates
each of a point $X$ and $U$ in a space of $n$ and of $\nu$
dimensions, then the relations (1) determine for $x$ a manifold of
fewer than $n$ dimensions. To each point $X$ there corresponds a
point $U$, and to the totality of the $X$, by (2), a manifold of the
$U$ of one dimension (a curve). Our assumption is that this is a
so-called rational or unicursal curve, or (in the language of the
theory of functions) a curve of genus zero.}.

To give an example here at once, one with which we shall later occupy
ourselves in detail, we mention the case of five variables
$x_0, x_1, x_2, x_3, x_4$ which satisfy the conditions
\weq{5}{\Sigma x = 0, \quad \Sigma x^2 = 0}
and which are therefore the roots of a \gesperrt{principal equation} of degree 5
(Bd. I, \S.\,54). That the assumption we have made holds in this
case is shown by the representation in Bd. I, \S.\,74, (4):
\[
y = t_3 F_0 + t_2(\alpha_2 F_1 + \alpha_1 F_2 + \alpha_0 F_3),
\]
which, when the determinations given there are made, put for the $x$,
satisfies the conditions (5) for all values of the ratio
$t = t_2 : t_3$.


% ===== scan leaf 438, printed 417 =====
% !PRINTED 417
\weberpage{417}
It is also shown there how the general equation of degree $5$ can be transformed into one corresponding to these assumptions.

Under the assumption (4), the subsistence of the relations (1) has no influence on the conclusions of \S.\,108, and we then have the following theorems.

Let $x_0, x_1, \ldots, x_{n-1}$ be a system of variables which is subject to the relations $A=0$, $B=0$, \ldots\ but otherwise independent, and let $\mathfrak{A}$ be a permutation group of the $x$ whose permutations leave the functions $A, B, \ldots$ unchanged.

\gesperrt{If the equation whose roots are the quantities $x$ is to possess a resolvent}
\weq{6}{F(u,z) = 0}
\gesperrt{which, apart from numerical coefficients, contains only one parameter that is a function of the $x$ unchanged by the group $\mathfrak{A}$, then there must exist a function $\vartheta$ of the variables $x_0, \ldots, x_{n-1}$ which, under the permutations $\alpha$ of $\mathfrak{A}$ (having regard to the relations $A=0$, $B=0$, \ldots), undergoes linear substitutions}
\weq{7}{\vartheta_\alpha = \dfrac{a\vartheta+b}{c\vartheta+d},}
\gesperrt{in which $a, b, c, d$ are numerical coefficients. The roots of the equation (6) are rationally expressible in terms of $\vartheta$, and consequently $\vartheta$ is also the root of an equation with one parameter}
\weq{8}{\Phi(\vartheta, z) = 0.}

If (6) is a total resolvent, which always occurs when the group $\mathfrak{A}$ is simple, then the $x$ too can be expressed rationally in terms of $\vartheta$ and the functions belonging to the group $\mathfrak{A}$.

Up to this point, then, there is complete agreement with the results of \S.\,108.

But the restriction upon which we came in \S.\,109 will here no longer necessarily arise, because now it is no longer possible, as in the case of completely independent variables, to conclude from the equality of two fractional functions the agreement of numerator and denominator.

The linear substitutions (7) form a group $P$, (simply or multiply) isomorphic with $\mathfrak{A}$, which with one of our
%CONT>>


% ===== scan leaf 439, printed 418 =====
% !PRINTED 418
\weberpage{418}
%<<CONT
\gesperrt{polyhedral groups} must be identical. We wish here to treat in more detail only the most interesting case, that $P$ is the \gesperrt{icosahedral group}, which moreover we may assume, without restriction of generality, in the normal form (\S.\,58).

For this group we have derived the three fundamental forms in \S.\,60. If we denote the variables of these fundamental forms by $y_1, y_2$, they are the three homogeneous functions of the $12$th, $20$th and $30$th degrees:
\weq{9}{\begin{aligned}
f(y) &= y_1 y_2 (y_1^{10} + 11 y_1^5 y_2^5 - y_2^{10}) \\
H(y) &= -y_1^{20} - y_2^{20} + 228 (y_1^{15} y_2^5 - y_2^{15} y_1^5) - 494 y_1^{10} y_2^{10} \\
T(y) &= y_1^{30} + y_2^{30} + 522 (y_1^{25} y_2^5 - y_1^5 y_2^{25}) - 10005 (y_1^{20} y_2^{10} + y_2^{20} y_1^{10}),
\end{aligned}}
between which the relation
\weq{10}{T^2 + H^3 = 1728 f^5}
further subsists.

If in the quotient $T^2 : f^5$ we put for the ratio $y_1 : y_2$ one of the 60 values $\vartheta_\alpha$ that arise from $\vartheta$ by the icosahedral substitutions, this quotient always receives the same value, which can therefore be expressed rationally in terms of the coefficients of equation (8), that is, rationally in terms of $z$. We may set it directly equal to $z$, and obtain for the resolvent (8) the form
\weq{11}{T^2 - zf^5 = 0,}
which on account of (10) is equivalent to
\[
H^3 - (1728 - z) f^5 = 0
\]
This equation is homogeneous in $y_1, y_2$ and of the $60$th degree, and if we set $y_2 = 1$, its 60 roots are $y_1 = \vartheta_\alpha$. It is the \gesperrt{icosahedral equation}, which was already defined in \S.\,60, and which therefore, written out in full in the form (11), runs thus:
\[
(\vartheta^{30} + 522 \vartheta^{25} - 10005 \vartheta^{20} - 10005 \vartheta^{10} - 522 \vartheta^5 + 1)^2 = z \vartheta^5 (\vartheta^{10} + 11 \vartheta^5 - 1)^5.
\]

The problem, as we have now framed it, comes down to the question:

\gesperrt{Which equations can be solved by means of the icosahedral equation?}


% ===== scan leaf 440, printed 419 =====
% !PRINTED 419
\weberpage{419}
\pdfbookmark[2]{§ 111. The resolvents of the icosahedral equation}{bm127}
\begin{center}
\S.\,111.

\gesperrt{The Resolvents of the Icosahedral Equation.}
\end{center}

The formulation of the problem which we gave at the close of the preceding section leads us to the question of the various resolvents of the icosahedral equation. To every divisor of the icosahedral group there corresponds such a resolvent, whose degree is equal to the index of the divisor, and is therefore always a divisor of 60. Now the icosahedral group contains (\S.\,59):

\begin{tabular}{@{}ll c l@{}}
1) Tetrahedral groups (\S.\,56): & Resolvents & of the $5$th & degree \\
2) Dihedral groups $D_5$ (\S.\,55): & \textquotedbl & of the $6$th & \textquotedbl \\
3) Dihedral groups $D_3$: & \textquotedbl & of the $10$th & \textquotedbl \\
4) Cyclic groups $C_5$: & \textquotedbl & of the $12$th & \textquotedbl \\
5) Four-groups $D_2$: & \textquotedbl & of the $15$th & \textquotedbl \\
6) Cyclic groups $C_3$: & \textquotedbl & of the $20$th & \textquotedbl \\
7) Cyclic groups $C_2$: & \textquotedbl & of the $30$th & \textquotedbl \\
\end{tabular}

Since the icosahedral group is simple, all these equations are \gesperrt{total resolvents} of one another; we confine ourselves here to the consideration of the most important among them.

The resolvents of lowest, namely of degree $5$, are those belonging to the tetrahedral group.

To find them, we must seek a function of $\vartheta$ which remains unchanged by the substitutions of the tetrahedral group, while it is five-valued in the icosahedral group.

We take the icosahedral group in the normal form set up in \S.\,58, (22), and take out the tetrahedral group $Q$ (\S.\,59) contained in it, which we have constructed over the four-group
\weq{1}{1, \psi, \chi, \psi\chi}
in which, if $\varepsilon$ is an imaginary fifth root of unity, $\psi$ and $\chi$ have the meaning:
\weq{2}{\psi(x) = \dfrac{-1}{x}, \quad \chi(x) = \dfrac{\omega x+1}{x-\omega}, \quad \omega = \varepsilon + \varepsilon^{-1}.}

If we set $\Theta(x) = \varepsilon x$, the whole icosahedral group becomes
\weq{3}{P = Q + Q\Theta + Q\Theta^2 + Q\Theta^3 + Q\Theta^4,}
and
\weq{4}{Q_h = \Theta^{-h} Q \Theta^h}
are the tetrahedral groups conjugate to $Q$.


% ===== scan leaf 441, printed 420 =====
% !PRINTED 420
\weberpage{420}
Now in \S.\,56 we have become acquainted with two forms $f$ and $H$ of degrees $6$ and $8$, which remain absolutely unchanged by the tetrahedral substitutions when these are taken homogeneous and with determinant 1. To be sure, we cannot apply these forms directly in the form in which they are given there, because there the tetrahedral group was assumed in a different shape. But we find the first of these forms very simply if we recall that its roots were the two-fold poles of the tetrahedral group, and hence the roots of the three equations
\[
x = \psi(x), \quad x = \chi(x), \quad x = \psi\chi(x)
\]
or
\[
x^2 + 1 = 0, \quad x^2 - 2\omega x - 1 = 0, \quad \omega x^2 + 2x - \omega = 0,
\]
and thus, with the use of the relation
\[
\omega^2 + \omega = 1
\]
one obtains the equation for the two-fold poles
\[
(x^2+1)(x^4+2x^3-6x^2-2x+1) = 0.
\]

Hence the sought tetrahedral form of degree $6$ in the homogeneous variables $x_1, x_2$ becomes
\weq{5}{t(x_1,x_2) = x_1^6 + 2x_1^5x_2 - 5x_1^4x_2^2 - 5x_1^2x_2^4 - 2x_1x_2^5 + x_2^6.}

The second tetrahedral form that we seek is the \gesperrt{Hesse} determinant of this.

If we set
\[
\tilde\omega(x_1,x_2) = \dfrac{1}{5^2 \cdot 4^2} \left\{ t''(x_1,x_1)\, t''(x_2,x_2) - [t''(x_1,x_2)]^2 \right\},
\]
then a simple calculation yields
\weq{6}{\tilde\omega(x_1,x_2) = -(x_1^8+x_2^8) + (x_1^7x_2 - x_1x_2^7) - 7(x_1^6x_2^2+x_1^2x_2^6) - 7(x_1^5x_2^3-x_1^3x_2^5).}

We further adduce the three icosahedral forms of degrees $12$, $20$ and $30$, which indeed also remain unchanged by the tetrahedral substitutions (\S.\,60):
\weq{7}{f(x_1,x_2) = x_1x_2(x_1^{10}+11x_1^5x_2^5-x_2^{10})}
\weq{8}{H(x_1,x_2) = -(x_1^{20}+x_2^{20}) + 228(x_1^{15}x_2^5-x_1^5x_2^{15}) - 494x_1^{10}x_2^{10}}
\weq{9}{T(x_1,x_2) = (x_1^{30}+x_2^{30}) + 522(x_1^{25}x_2^5-x_1^5x_2^{25}) - 10005(x_1^{20}x_2^{10}+x_2^{20}x_1^{10}).}


% ===== scan leaf 442, printed 421 =====
% !PRINTED 421
\weberpage{421}

If we form from these invariant forms of the tetrahedron functions of $x_1:x_2$ and set $\vartheta = x_1:x_2$, we obtain functions which remain unchanged by the linear fractional tetrahedral substitutions, and which are therefore roots of resolvents of degree 5 of the icosahedral equation. The icosahedral equation itself is obtained in two forms, if one sets
\weq{10}{H^3 = zf^5,\quad T^2 = z_1f^5,\quad z + z_1 = 1728}
and the resolvents depend rationally on $z$.

If for the moment we denote by $\varphi(x_1, x_2)$ any absolutely invariant form of the tetrahedron, then this form passes, by any one of the substitutions $Q\Theta$, into
\weq{11}{\varphi_\varepsilon = \varphi(\varepsilon^{-2}x_1,\ \varepsilon^2x_2)}
and this function remains unchanged by the substitutions of the group $\Theta^{-1}Q\Theta$.

\gesperrt{From this it follows that every symmetric function of the five forms $\varphi_\varepsilon$ is an invariant of the icosahedral group, and can therefore, by the theorem of \S.\,61, be expressed by the fundamental forms $f, H, T$.}

This theorem permits a comparatively easy computation of the resolvents.

The simplest functions that we can use as roots of resolvents are
\weq{12}{r = \frac{t^2}{f},\quad u = \frac{f^2t}{T},\quad v = \frac{f\tilde\omega}{H}.}

If we first observe that the symmetric fundamental functions of the five functions $t_\varepsilon$ are icosahedral invariants of the degrees $6, 12, 18, 30$, then, since the degree numbers $6, 18$ do not occur among these invariants, there follows from the theorems in \S.\,61 an equation of the form
\weq{13}{t^5 + aft^3 + bf^2t + cT = 0,}
in which $a, b, c$ are numerical coefficients. To determine them, we denote by $S(\varphi)$ the sum of the five functions $\varphi_\varepsilon$ in (11) and by $\Pi(\varphi)$ their product, and obtain from the \gesperrt{Newton} formulae (Bd.\,I, \S.\,42):
\weq{14}{af = -\tfrac{1}{2}S(t^2),\quad bf^2 = \tfrac{1}{2}a^2f^2 - \tfrac{1}{4}S(t^4),\quad cT = -\Pi(t).}

In the sums $S(\varphi)$ only those terms $x_1^h x_2^k$ have a coefficient different from zero in which $h - k$ is divisible by 5. One therefore obtains $a, b, c$ by equating the
%CONT>>


% ===== scan leaf 443, printed 422 =====
% !PRINTED 422
\weberpage{422}
\pdfbookmark[2]{§ 112. The principal resolvent of the fifth degree}{bm128}

%<<CONT
coefficients of $x_1^{11}x_2,\ x_1^{22}x_2^2,\ x_1^{30}$ on both sides of the equations (14), and thus finds the relation
\weq{15}{t^5 - 10ft^3 + 45f^2t - T = 0,}
or also
\weq{16}{t^2(t^4 - 10ft^2 + 45f^2)^2 - T^2 = 0.}

From (16) and (15) there now result, by (10) and (12), the equations of the fifth degree for $r$ and for $u$
\weq{17}{r(r^2 - 10r + 45)^2 = z_1,}
\weq{18}{u^5 - \frac{10u^3}{z_1} + \frac{45u}{z_1^2} - \frac{1}{z_1^2} = 0.}

The last of these equations passes, by the substitution
\weq{19}{y = -\frac{1}{3u},\quad z_1 = 27\gamma}
into the form
\weq{20}{y^5 + 15y^4 - 10\gamma y^2 + 3\gamma^2 = 0}
which we have already come to know in \S.\,75 of the first volume as a normal form of the equation of the fifth degree.

\begin{center}
\S.\,112.\\[4pt]
\gesperrt{The Principal Resolvent of the Fifth Degree.}
\end{center}

If one makes use of the functions $u, v$ simultaneously, one can derive a family of resolvents which have the form of the principal equation of the fifth degree, and which can be brought directly into agreement with an arbitrarily given principal equation of the fifth degree\footnote{\gesperrt{Kiepert}, G\"ottinger Nachrichten 1878. Crelle's Journal, Bd.\,87. \gesperrt{Klein}, Ikosa\"eder, S.\,106.}.

Since for the icosahedron there are no invariants of degree 8, 14, 16, 22 or 28, the functions $S(\tilde\omega)$, $S(t\tilde\omega)$, $S(\tilde\omega^2)$, $S(t\tilde\omega^2)$, $S(t^2\tilde\omega^2)$ must vanish, and if we therefore set
\weq{1}{Y = \alpha\tilde\omega + \beta t\tilde\omega}
then the functions
\[ S(Y),\quad S(Y^2) \]
will also vanish for arbitrary values of the parameters $\alpha, \beta$, and from this there results an identical equation of the form
%CONT>>


% ===== scan leaf 444, printed 423 =====
% !PRINTED 423
\weberpage{423}

%<<CONT
\weq{2}{Y^5 + 5aY^2 + 5bY + c = 0,}
in which $a, b, c$ denote homogeneous functions of degrees 3, 4, 5 in $\alpha, \beta$, whose coefficients are icosahedral invariants.

We can easily determine the function $c$ from the formula
\[ c = -\Pi(\tilde\omega)\,\Pi(\alpha + \beta t). \]

For, by (15) of the preceding section, we have for an indeterminate $\lambda$
\[ \lambda^5 - 10f\lambda^3 + 45f^2\lambda - T = \Pi(\lambda - t), \]
hence, if one sets $\lambda = -\alpha : \beta$,
\[ \Pi(\alpha + \beta t) = \alpha^5 - 10f\alpha^3\beta^2 + 45f^2\alpha\beta^4 + T\beta^5, \]
and further, by comparison of a single term, there follows at once [\S.\,111, (6), (8)]
\[ \Pi(\tilde\omega) = -H^2, \]
hence
\weq{3}{c = H^2(\alpha^5 - 10f\alpha^3\beta^2 + 45f^2\alpha\beta^4 + T\beta^5).}

The coefficients $a, b$ are perhaps most simply computed with the help of the \gesperrt{Newton} formulae (Bd.\,I, \S.\,42):
\[
\begin{aligned}
-15a &= S(\alpha\tilde\omega + \beta t\tilde\omega)^3\\
-20b &= S(\alpha\tilde\omega + \beta t\tilde\omega)^4,
\end{aligned}
\]
by applying the formulae which result, after a short calculation, from the comparison of one term each of the right and left sides:
\[
\begin{array}{ll}
S(\tilde\omega^3) = -3.5.8f^2, & S(t\tilde\omega^3) = -5T,\\
S(t^2\tilde\omega^3) = -5.72f^3, & S(t^3\tilde\omega^3) = -3.5fT,\\
S(\tilde\omega^4) = 4.5fH, & S(t^2\tilde\omega^4) = -5.12f^2H,\\
S(t^3\tilde\omega^4) = -5HT, & S(t^4\tilde\omega^4) = -4.5.27f^3H.
\end{array}
\]

Thus one finds
\weq{4}{a = 8f^2\alpha^3 + T\alpha^2\beta + 72f^3\alpha\beta^2 + fT\beta^3}
\weq{5}{b = -fH\alpha^4 + 18f^2H\alpha^2\beta^2 + HT\alpha\beta^3 + 27f^3H\beta^4.}

In order to obtain from this the resolvent of the icosahedral equation, we must introduce in place of $\tilde\omega$ and $t$ the functions $u, v$ [\S.\,111, (12)].

If we set, denoting by $\lambda, \mu$ any two indeterminate quantities,
\weq{6}{\alpha = \frac{\lambda f}{H},\quad \beta = \frac{\mu f^3}{HT},}
%CONT>>


% ===== scan leaf 445, printed 424 =====
% !PRINTED 424
\weberpage{424}

%<<CONT
it follows by (1) and \S.\,111, (12):
\weq{7}{Y = \lambda v + \mu vu,}
and if, as in \S.\,111, we set
\weq{8}{\frac{H^3}{f^5} = z,\quad \frac{T^2}{f^5} = z_1,\quad z + z_1 = 1728}
then the expressions (3), (4), (5) pass into the following:
\weq{9}{az = 8\lambda^3 + \lambda^2\mu + \frac{72\lambda\mu^2 + \mu^3}{z_1}}
\[
bz = -\lambda^4 + \frac{18\lambda^2\mu^2 + \lambda\mu^3}{z_1} + 27\frac{\mu^4}{z_1^2}
\]
\[
cz = \lambda^5 - 10\frac{\lambda^3\mu^2}{z_1} + \frac{45\lambda\mu^4 + \mu^5}{z_1^2},
\]
if $Y$ is the root of the equation
\weq{10}{Y^5 + 5aY^2 + 5bY + c = 0}
This equation has the form of a principal equation of degree 5, and it can represent every principal equation if we may take $a, b, c$ arbitrarily. If, then, $a, b, c$ are assumed to be arbitrarily given, we have to determine from (9) the three quantities $\lambda, \mu, z_1$ (and $z = 1728 - z_1$), and it is a very remarkable peculiarity of this system of equations that these unknowns can be expressed \gesperrt{rationally} in terms of the given quantities $a, b, c$ and the discriminant of the equation (10). To prove this, we first derive simpler equations from (9).

If we multiply the second of them by $\lambda$ and add it to the third, it follows that
\weq{11}{z_1(\lambda b + c) = \mu^2a,}
and if we multiply the third by $\lambda$ and the second by $\mu^2 : z_1$ and subtract,
\weq{12}{z\left(\lambda c - \frac{\mu^2}{z_1}b\right) = \left(\lambda^2 - 3\frac{\mu^2}{z_1}\right)^3.}

Likewise, from the first and second we obtain
\[
z\,\frac{a\lambda + 8b}{\mu} = \lambda^3 + \frac{216\lambda^2\mu}{z_1} + 9\frac{\lambda\mu^2}{z_1} + \frac{216\mu^3}{z_1^2}.
\]


% ===== scan leaf 446, printed 425 =====
% !PRINTED 425
\weberpage{425}
This equation yields, if one squares both sides:
\[
\begin{aligned}
z_1 z^2\left(\frac{a\lambda+8b}{\mu}\right)^2 &= \lambda^6 z_1 + 2.216\,\lambda^5\mu + \left(18+\frac{216^2}{z_1}\right)\lambda^4\mu^2 \\
&\quad + \frac{2.2160}{z_1}\,\lambda^3\mu^3 + \left(\frac{81}{z_1}+\frac{2.216^2}{z_1^2}\right)\lambda^2\mu^4 \\
&\quad + \frac{18.216}{z_1^2}\,\lambda\,\mu^5 + \frac{216^2\mu^6}{z_1^3},
\end{aligned}
\]
and if one subtracts this from the formula derived from the first equation (9):
\[
\begin{aligned}
27a^2z^2 &= 1728\,\lambda^6 + 2.216\,\lambda^5\mu + 27\left(1+\frac{16.72}{z_1}\right)\lambda^4\mu^2 \\
&\quad + \frac{2.2160}{z_1}\,\lambda^3\mu^3 + 27\left(\frac{2}{z_1}+\frac{72^2}{z_1^2}\right)\lambda^2\mu^4 \\
&\quad + \frac{18.216}{z_1^2}\,\lambda\,\mu^5 + \frac{27\mu^6}{z_1^2},
\end{aligned}
\]
then, having regard to (8), it follows that:
\weq{13}{z\left[27a^2 - z_1\left(\frac{a\lambda+8b}{\mu}\right)^2\right] = \left(\lambda^2-\frac{3\mu^2}{z_1}\right)^3.}
From this there results further, by comparison with (12),
\weq{14}{27a^2 - \frac{z_1}{\mu^2}(a\lambda+8b)^2 = \lambda c - \frac{\mu^2}{z_1}b,}
and if one eliminates from this, by means of (11), the ratio $\mu^2:z_1$, there results the quadratic equation for $\lambda$:
\weq{15}{\begin{split} \lambda^2(a^4+abc-b^3) - \lambda(11a^3b-ac^2+2b^2c) \\ - 27a^3c + 64a^2b^2 - bc^2 = 0. \end{split}}
The discriminant of this quadratic equation is
\[
\begin{aligned}
\Delta &= (11a^3b - ac^2 + 2b^2c)^2 \\
&\quad + 4(a^4+abc-b^3)(27a^3c - 64a^2b^2 + bc^2) \\
&= a^2(108a^5c - 135a^4b^2 + 90a^2bc^2 \\
&\quad - 320ab^3c + 256b^5 + c^4).
\end{aligned}
\]

In \S.\,74 of the first volume the discriminant $D$ of the principal equation of degree $5$ is derived. If in the formula (3) found there one puts
\[
a_0=1,\quad a_3=5a,\quad a_4=5b,\quad a_5=c
\]
%CONT>>


% ===== scan leaf 447, printed 426 =====
% !PRINTED 426
\weberpage{426}
\pdfbookmark[2]{§ 113. Resolvents of the sixth degree}{bm129}
%<<CONT
is substituted, one finds
\[
D = 5^5(108a^5c-135a^4b^2+90a^2bc^2 - 320ab^3c+256b^5+c^4),
\]
so that
\weq{16}{5^5\Delta = a^2D}
results, and from (15) one obtains for $\lambda$ the expression
\[
\lambda = \frac{11a^3b - ac^2 + 2b^2c + \frac{1}{25}a\sqrt{\frac{1}{5}D}}{2(a^4+abc-b^3)}
\]
in which the square root may take both signs.

One sees that, besides the square root of the discriminant, which is rationally expressible in terms of the roots, $\sqrt5$ also occurs in it.

Once $\lambda$ has been computed, one obtains from (11) the ratio $z_1:\mu^2$, and from (12) the last unknown $z$, represented \gesperrt{rationally} in terms of $\lambda$.

\begin{center}
\S.\,113.

\gesperrt{Resolvents of the Sixth Degree.}
\end{center}

To find the resolvents of degree $6$ of the icosahedral equation, we start from one of the dihedral groups $D_5$ contained in the icosahedral group (\S.\,111, 2.). The representation of the icosahedral group [\S.\,58, (20)] gives us at once the decomposition into the subsidiary groups; if
\weq{1}{Q = \Theta^r,\ \Theta^r\psi,\ r=0,1,2,3,4}
is the dihedral group from which we start, then we obtain the full icosahedral group
\weq{2}{P = Q + Q\chi + Q\chi\Theta + Q\chi\Theta^2 + Q\chi\Theta^3 + Q\chi\Theta^4.}
If now $U_\infty$ is a function which remains unchanged by the substitutions of $Q$ and which is carried by
\[
\chi,\ \chi\Theta,\ \chi\Theta^2,\ \chi\Theta^3,\ \chi\Theta^4
\]
into
\[
U_0,\ U_1,\ U_2,\ U_3,\ U_4
\]
then the $U_\infty, U_0, U_1, U_2, U_3, U_4$ are the roots of a resolvent of degree $6$. (On the notation $U_\infty$ cf. \S.\,65.)


% ===== scan leaf 448, printed 427 =====
% !PRINTED 427
\weberpage{427}
To obtain the \gesperrt{Galois} group of this equation of degree $6$, we have to examine the influence which the application of the icosahedral substitutions has on the system of subsidiary groups (2). For this it suffices to consider the substitutions $\Theta$, $\psi$, $\chi$. But, if we attend to the order of the subsidiary groups in (2), we have by \S.\,58, (23), (25):
\[
\begin{aligned}
P\Theta &= Q + Q\chi\Theta + Q\chi\Theta^2 + Q\chi\Theta^3 + Q\chi\Theta^4 + Q\chi \\
P\psi &= Q + Q\chi + Q\chi\Theta^4 + Q\chi\Theta^3 + Q\chi\Theta^2 + Q\chi\Theta \\
P\chi &= Q\chi + Q + Q\chi\Theta + Q\chi\Theta^3 + Q\chi\Theta^2 + Q\chi\Theta^4,
\end{aligned}
\]
and from this it follows that the indices of $U$ are permuted as follows:
\[
\begin{array}{c@{\quad}cccccc}
 & \infty, & 0, & 1, & 2, & 3, & 4 \\
\cline{2-7}
\Theta) & \infty & 1 & 2 & 3 & 4 & 0 \\
\psi) & \infty & 0 & 4 & 3 & 2 & 1 \\
\chi) & 0 & \infty & 1 & 3 & 2 & 4.
\end{array}
\]
If we denote the index generally by $\xi$, and take $\xi$, as in the seventh chapter, with respect to the modulus 5, so that two numbers congruent with respect to the modulus 5 count as not distinct, then we can represent these permutations thus:
\[
\Theta = (\xi,\ \xi+1),\quad \psi = (\xi,\ -\xi),\quad \chi = \left(\xi,\ \frac{1}{\xi}\right)\!.
\]
\gesperrt{But from this it follows that the group of our equation of degree} $6$ \gesperrt{coincides with the congruence group} $L_5$ (\S.\,68)\gesperrt{.}

To arrive at the formation of resolvents of degree $6$, we must proceed in a way similar to that used for the resolvents of degree $5$.

We have already become acquainted in \S.\,55 with the fundamental forms of the dihedral groups. We now represent them in the form:
\[
\varphi_1 = x_1x_2,\quad \varphi_2 = x_1^5+ix_2^5,\quad \varphi_3 = x_1^5-ix_2^5.
\]
If we take $\Theta$ and $\psi$ with the determinant 1, that is
\[
\Theta = \begin{pmatrix} e^{\frac{\pi i}{5}}, & 0 \\ 0, & e^{-\frac{\pi i}{5}} \end{pmatrix},\quad \psi = \begin{pmatrix} 0, & 1 \\ -1, & 0 \end{pmatrix};
\]
then the fundamental forms $\varphi_1, \varphi_2, \varphi_3$ change in the following way:
\[
\begin{array}{c@{\quad}ccc}
 & \varphi_1 & \varphi_2 & \varphi_3 \\
\cline{2-4}
\Theta) & \varphi_1 & -\varphi_2 & -\varphi_3 \\
\psi) & -\varphi_1 & -i\varphi_2 & i\varphi_3.
\end{array}
\]


% ===== scan leaf 449, printed 428 =====
% !PRINTED 428
\weberpage{428}
Thus $\varphi_1^2, \varphi_2^4, \varphi_3^4, \varphi_2\varphi_3$ are absolutely invariant, and if therefore, as in \S.\,111, we understand by $f, H, T$ the icosahedral invariants, we can introduce the following functions as roots of resolvents of degree $6$:
\[
\frac{\varphi_1^2H}{f^2},\quad \frac{\varphi_2^4}{H},\quad \frac{\varphi_3^4}{H},\quad \frac{\varphi_2\varphi_3H}{T}.
\]
In order to obtain the expressions for the remaining roots, the influence of the substitutions $\Theta^r\chi$, taken with determinant 1, on the functions $\varphi_1, \varphi_2, \varphi_3$ must be investigated. We shall carry this out only for the function $\varphi_1$, which leads to the simplest resolvent of degree $6$.

If we apply in $\varphi_1$ the substitution $\chi$ in the form \S.\,58, (26), that is, replace $x_1, x_2$ by
\[
\frac{x_1}{\varepsilon-\varepsilon^{-1}}+\frac{x_2}{\varepsilon^2-\varepsilon^{-2}},\quad \frac{x_1}{\varepsilon^2-\varepsilon^{-2}}-\frac{x_2}{\varepsilon-\varepsilon^{-1}}
\]
then $\varphi_1$ goes over into
\[
\frac{-x_1^2+x_1x_2+x_2^2}{\sqrt5}
\]
We accordingly set
\weq{3}{w_\infty = 5\varphi_1^2 = 5x_1^2x_2^2,}
and obtain from this, by applying the substitutions $\chi\Theta^r$, the five further forms
\weq{4}{w_r = (\varepsilon^rx_1^2+x_1x_2-\varepsilon^{-r}x_2^2)^2.}
Now the symmetric functions of the six forms $w$ are invariant icosahedral forms, and from this one can easily form the coefficients of the equation whose roots are the $w$. By comparison of the degrees and of one coefficient in each case, we obtain for the power sums of the $w$
\[
S(w)=0,\ S(w^2)=0,\ S(w^3)=30f,\ S(w^4)=0,\ S(w^5)=-5H,
\]
and thus the product of all six $w$
\[
\Pi(w)=5f^2,
\]
and accordingly, by the \gesperrt{Newton} formulae (Bd.\,I, \S.\,42), there results for $w$ the equation
\weq{5}{w^6-10fw^3+Hw+5f^2=0.}
If we therefore set
\weq{6}{U=\frac{wH}{f^2},\ z=\frac{H^3}{f^2},}
% [sic] The source prints $z=H^3/f^2$; the homogeneous relation and the continuation indicate that $f^5$ was intended.
%CONT>>


% ===== scan leaf 450, printed 429 =====
% !PRINTED 429
\weberpage{429}
\pdfbookmark[2]{§ 114. On the solution of the icosahedral equation by transcendental functions}{bm130}
%<<CONT
there follows from this the equation of degree 6 for $U$:
\weq{7}{U^6 - 10\,z\,U^3 + z^2\,U + 5\,z^2 = 0}
as the \gesperrt{resolvent of degree} 6 of the icosahedral equation.

If, following (3) and (4), we set
\weq{8}{\begin{aligned}
\sqrt{w_\infty} &= \sqrt{5}\,x_1\,x_2 \\
\sqrt{w_r} &= \varepsilon^r x_1^2 + x_1 x_2 - \varepsilon^{-r} x_2^2,
\end{aligned}}
then from this the following relations result:
\weq{9}{\begin{aligned}
\sqrt{5\,w_\infty} &= \sqrt{w_0} + \sqrt{w_1} + \sqrt{w_2} + \sqrt{w_3} + \sqrt{w_4} \\
0 &= \sqrt{w_0} + \varepsilon^2\sqrt{w_1} + \varepsilon^4\sqrt{w_2} + \varepsilon\sqrt{w_3} + \varepsilon^3\sqrt{w_4} \\
0 &= \sqrt{w_0} + \varepsilon^3\sqrt{w_1} + \varepsilon\sqrt{w_2} + \varepsilon^4\sqrt{w_3} + \varepsilon^2\sqrt{w_4}
\end{aligned}}
\weq{10}{\frac{x_1\sqrt{5}}{x_2} = \frac{\sqrt{w_0} + \varepsilon^4\sqrt{w_1} + \varepsilon^3\sqrt{w_2} + \varepsilon^2\sqrt{w_3} + \varepsilon\sqrt{w_4}}{\sqrt{w_\infty}},}
and in the equations (9), (10), since they are homogeneous, the $U$ may also be put in place of the $w$.

The equations of degree 6 whose roots satisfy the relations (9) were first pointed out by \gesperrt{Jacobi}; they are therefore called \gesperrt{Jacobi} equations of degree 6\footnote{Suite des notices sur les fonctions elliptiques. \gesperrt{Crelle's} Journal, Bd. III (1828); Werke, Bd. I, S. 261.}.

\begin{center}
\S.\,114.\\[4pt]
\gesperrt{On the solution of the icosahedral equation by transcendental functions.}
\end{center}

We wish here, without entering into proofs or more detailed developments, to point briefly to the relations that exist between the investigations we have carried out and the theory of certain transcendental functions, in particular the theory of elliptic functions and the theory of the hypergeometric series, from which results the solution of the icosahedral equation, and with it of all its resolvents, by transcendental functions. This leads out of the region of algebra, but stands to it in about the same relation as when one employs logarithms in the computation of root quantities, or the division of the angle in the case of the cubic equations.


% ===== scan leaf 451, printed 430 =====
% !PRINTED 430
\weberpage{430}
We must confine ourselves here to a brief statement of the results, and shall make no applications of them in what follows, so that this exposition, which can be fully intelligible only to a reader somewhat familiar with the theory of transcendental functions, may also be passed over without disturbing the continuity.

Let $\omega$ denote a variable whose imaginary part is positive, so that
\weq{1}{q = e^{\pi i \omega}}
is a quantity whose absolute value is a proper fraction. We shall consider the following three functions of this variable $\omega$:
\weq{2}{\eta(\omega) = q^{1/12}\prod_{1,\infty}^{\nu}(1 - q^{2\nu}) = q^{1/12}\sum_{-\infty,\infty}^{\nu}(-1)^\nu q^{3\nu^2+\nu} = \sum_{-\infty,\infty}^{\nu}(-1)^\nu q^{\frac{(6\nu+1)^2}{12}},}
\weq{3}{f(\omega) = q^{-\frac{1}{24}}\prod(1 + q^{2\nu-1}),}
\weq{4}{\gamma_2(\omega) = \frac{f(\omega)^{24} - 16}{f(\omega)^8}\footnote{Cf.: \gesperrt{Weber}, ``Elliptische Functionen und algebraische Zahlen''. Braunschweig 1891. S. 63, 86, 149.}.}
In the transformation theory of the elliptic functions\footnote{Ibid., S. 263, 316.} an equation of degree 6 is derived, whose roots are the six quantities
\weq{5}{v_\infty = 5\left(\frac{\eta(5\omega)}{\eta(\omega)}\right)^2, \qquad v_\xi = \left(\frac{\eta\left(\dfrac{-24\xi+\omega}{5}\right)}{\eta(\omega)}\right)^2}
when $\xi$ runs through a complete system of residues to the modulus 5. This equation reads
\weq{6}{v^6 + 10v^3 - \gamma_2(\omega)v + 5 = 0,}
and, as one sees, has great similarity to the resolvent of degree 6 considered in the preceding section
\weq{7}{U^6 - 10zU^3 + z^2U + 5z^2 = 0.}
These two equations pass directly into one another, if one puts
\weq{8}{z = \gamma_2(\omega)^3}
\weq{9}{U = -\gamma_2(\omega)v}
%CONT>>


% ===== scan leaf 452, printed 431 =====
% !PRINTED 431
\weberpage{431}
%<<CONT
sets. Equation (6) has the same \gesperrt{Galois} group as (7). We can now put any one of the roots of (7) for $U_\infty$, and consequently we can let $U_\infty$ and $v_\infty$ in (9) correspond to one another. Then we can take any one of the remaining roots of (7) for $U_0$, because we can replace $x_1, x_2$ in \S.\,113, (4) by $\varepsilon^r x_1$, $\varepsilon^{-r} x_2$. We therefore let $U_0$ and $v_0$ correspond to one another. Finally, we can still dispose of the imaginary fifth root of unity $\varepsilon$ in such a way that $U_1$ and $v_1$ correspond, and then it follows from the agreement of the group that $U_\xi$ and $v_\xi$ correspond throughout, and hence that
\weq{10}{U_\xi = -\gamma_2(\omega)v_\xi, \qquad \xi = \infty, 0, 1, 2, 3, 4}
\gesperrt{If, then, we determine} $\omega$ \gesperrt{for a given} $z$ \gesperrt{from the transcendental equation (8), then (10) and (5) give us the complete solution of the icosahedral resolvent of the sixth degree.}

By (2) and (5) we obtain, for $\xi = 0, 1, 2, 3, 4$,
\weq{11}{\eta(\omega)\sqrt{v_\xi} = \sum_{-\infty,\infty}^{\nu}(-1)^\nu e^{\frac{\pi i}{60}(\omega - 24\xi)(6\nu+1)^2}.}
The exponent $(6\nu+1)^2$ can give only the residue $0, 1, -1$ with respect to the modulus 5. If, then, we let $\nu_1, \nu_2$ run through all the numbers which satisfy the conditions
\[
\nu_1 \equiv 1, 2, \qquad \nu_2 \equiv 0, 3 \pmod{5}
\]
then we can decompose the expression (11) as follows:
\weq{12}{\eta(\omega)\sqrt{v_\xi} = \eta(5\omega) + \varepsilon^\xi \sum^{\nu_1}(-1)^{\nu_1} e^{\frac{\pi i}{60}\omega(6\nu_1+1)^2} + \varepsilon^{-\xi}\sum^{\nu_2}(-1)^{\nu_2} e^{\frac{\pi i}{60}\omega(6\nu_2+1)^2},}
if $\varepsilon = e^{\frac{2\pi i}{5}}$ is set. From this one sees that the quantities $\sqrt{v_\xi}$ satisfy the relations (9) of the preceding section, and the formula (10), \S.\,113 gives for the root $x = x_1 : x_2$ of the icosahedral equation the expression:
\weq{13}{\eta(5\omega)x = \sum^{\nu_1}(-1)^{\nu_1} e^{\frac{\pi i}{60}\omega(6\nu_1+1)^2}.}
The remaining roots are obtained by applying the icosahedral substitutions to $x$. One can, however, also
%CONT>>


% ===== scan leaf 453, printed 432 =====
% !PRINTED 432
\weberpage{432}
%<<CONT
by replacing the variable $\omega$ by an equivalent value (Bd. I, \S.\,120).

\begin{center}
\rule{0.25\linewidth}{0.4pt}
\end{center}

Finally, we wish here also to mention briefly the solution of the icosahedral equation by hypergeometric series. Following \gesperrt{Gauss}, we understand by
\[
F(\alpha, \beta, \gamma, \xi)
\]
the infinite series
\[
1 + \frac{\alpha\beta}{1.\gamma}\xi + \frac{\alpha(\alpha+1)\beta(\beta+1)}{1.2\,\gamma(\gamma+1)}\xi^2 + \frac{\alpha(\alpha+1)(\alpha+2)\beta(\beta+1)(\beta+2)}{1.2.3\,\gamma(\gamma+1)(\gamma+2)}\xi^2 + \cdots
\]
% [sic] The source clearly repeats $\xi^2$ in the third term; the hypergeometric series requires $\xi^3$.
If we set in it
\[
z = 1728\,\xi,
\]
then the root of the icosahedral equation
\[
H^3 - zf^5 = 0
\]
obtains the expression
\[
x = \xi^{\frac{1}{3}}\,\frac{F\left(\dfrac{11}{60}, \dfrac{19}{60}, \dfrac{4}{3}, \xi\right)}{F\left(\dfrac{11}{60}, \dfrac{-1}{60}, \dfrac{1}{3}, \xi\right)}.
\]
These series converge only so long as the absolute value of $\xi$ is a proper fraction, but they can be replaced, for other values of $\xi$, by similar series.\footnote{\gesperrt{H. A. Schwarz}, ``Ueber diejenigen Fälle, in denen die \gesperrt{Gauss}'sche Reihe $F(\alpha,\beta,\gamma,x)$ eine algebraische Function ihres vierten Elementes ist''. Crelle's Journal, Bd. 75, 1872 (gesammelte mathematische Werke, Bd. II). \gesperrt{F. Klein}, Vorlesungen über das Ikosaëder, Abschnitt I, Cap. III.}

\begin{center}
\rule{0.25\linewidth}{0.4pt}
\end{center}


% ===== scan leaf 454, printed 433 =====
\weberpage{433}% page number not printed on this page
\pdfbookmark[1]{Fourteenth Section. Groups of Linear Ternary Substitutions.}{bm131}
\pdfbookmark[2]{§ 115. Ternary linear substitution group of the 168th degree}{bm132}
% !PRINTED none
% Chapter-opening page: no running head/page number is printed at the top of this
% page; instead the page carries the Abschnitt title block, then the §.115 heading
% and the start of the section text. (Weber suppresses the folio number on
% Abschnitt-opening pages.)

\begin{center}
\gesperrt{Fourteenth Chapter.}

\gesperrt{Groups of Linear Ternary Substitutions.}

\rule{0.3\linewidth}{0.4pt}

\S.\,115.

\gesperrt{Ternary linear substitution group of degree 168.}
\end{center}

The question of the totality of all possible finite groups of
linear substitutions in several dimensions, in particular those of
three and four dimensions, has been treated by \gesperrt{C. Jordan},
who, much as was done for the binary substitutions, has set up a
finite number of types of such groups\footnote{\gesperrt{C. Jordan}, Mémoire sur les équations différentielles linéaires à
intégrale algébrique. Crelle's Journal für Mathematik, Bd. 84 (1878). Sur
la détermination des groupes d'ordre fini etc. Atti della R. Accademia di
Napoli, Vol. VIII (1879). \gesperrt{Valentiner}, Kj\"ob. Skrift (6) V (1889).}. These general investigations
we cannot pursue here; we content ourselves with considering one
example of a ternary group in more detail.

In the chapter on the congruence groups (\S.\,66) we became
acquainted with a sequence of \gesperrt{simple} groups, of which the first,
of degree 60, proved to be isomorphic with the icosahedral group.
We now wish to consider the next of these simple groups, of
degree 168, which we have denoted by $L_7$, and attempt to
construct a ternary linear substitution group isomorphic with it.

To this we are led by the Theorem \S.\,72, I., if we seek out
four elements $\tau, \chi, \omega, \Theta$ which, under composition, satisfy the
conditions of this theorem.

We will first assume a substitution $\tau$ in the normal form
\weq{1}{\tau = \begin{pmatrix} \varepsilon_1, & 0, & 0 \\ 0, & \varepsilon_2, & 0 \\ 0, & 0, & \varepsilon_3 \end{pmatrix},}
%CONT>>


% ===== scan leaf 455, printed 434 =====
% !PRINTED 434
\weberpage{434}
%<<CONT
and in it, since $\tau$ is to be of degree 7, the $\varepsilon_1, \varepsilon_2, \varepsilon_3$
must be seventh roots of unity which satisfy the condition
\weq{2}{\varepsilon_1 \varepsilon_2 \varepsilon_3 = 1}

We shall now seek to determine the substitution $\chi$ so
that the condition $\chi \tau = \tau^4 \chi$, or
\weq{3}{\tau \chi = \chi \tau^2}
is fulfilled. If we take $\chi$ in the form
\[
\chi = \begin{pmatrix} \alpha_1, & \alpha_2, & \alpha_3 \\ \beta_1, & \beta_2, & \beta_3 \\ \gamma_1, & \gamma_2, & \gamma_3 \end{pmatrix}
\]
then the condition (3) yields
\[
\begin{pmatrix} \alpha_1 \varepsilon_1, & \alpha_2 \varepsilon_1, & \alpha_3 \varepsilon_1 \\ \beta_1 \varepsilon_2, & \beta_2 \varepsilon_2, & \beta_3 \varepsilon_2 \\ \gamma_1 \varepsilon_3, & \gamma_2 \varepsilon_3, & \gamma_3 \varepsilon_3 \end{pmatrix}
=
\begin{pmatrix} \alpha_1 \varepsilon_1^2, & \alpha_2 \varepsilon_2^2, & \alpha_3 \varepsilon_3^2 \\ \beta_1 \varepsilon_1^2, & \beta_2 \varepsilon_2^2, & \beta_3 \varepsilon_3^2 \\ \gamma_1 \varepsilon_1^2, & \gamma_2 \varepsilon_2^2, & \gamma_3 \varepsilon_3^2 \end{pmatrix}.
\]

Now since $\alpha_1, \beta_1, \gamma_1$ cannot all three vanish, $\varepsilon_1^2$
must agree with one of the three quantities $\varepsilon_1, \varepsilon_2, \varepsilon_3$.
If we had $\varepsilon_1 = \varepsilon_1^2$, hence $\varepsilon_1 = 1$, and consequently $\varepsilon_3 = \varepsilon_2^{-1}$, then
we could have neither $\varepsilon_2$ nor $\varepsilon_3 = 1$, since otherwise $\tau$ would be the identical
substitution. We should therefore have to have $\beta_1 = \gamma_1 = \alpha_2 = \alpha_3 = 0$ and $\varepsilon_2^2 = \varepsilon_3$,
$\varepsilon_2^3 = 1$, which is impossible. There remain, therefore, by (2),
two possible assumptions:
\[
\varepsilon_1 = \varepsilon, \quad \varepsilon_2 = \varepsilon^2, \quad \varepsilon_3 = \varepsilon^4,
\]
\[
\varepsilon_1 = \varepsilon, \quad \varepsilon_2 = \varepsilon^4, \quad \varepsilon_3 = \varepsilon^2,
\]
in which $\varepsilon$ is an imaginary seventh root of unity, and these
two assumptions lead to two \gesperrt{different} groups, which are
built quite alike, and moreover are also isomorphic. We pursue the first assumption further here, and accordingly set
\weq{4}{\tau = \begin{pmatrix} \varepsilon, & 0, & 0 \\ 0, & \varepsilon^2, & 0 \\ 0, & 0, & \varepsilon^4 \end{pmatrix}, \quad \chi = \begin{pmatrix} 0, & 0, & 1 \\ 1, & 0, & 0 \\ 0, & 1, & 0 \end{pmatrix}, \quad \chi^2 = \begin{pmatrix} 0, & 1, & 0 \\ 0, & 0, & 1 \\ 1, & 0, & 0 \end{pmatrix}, \quad \chi^3 = 1.}

That we have set the non-vanishing quantities $\alpha_3, \beta_1, \gamma_2$ equal to 1
is no restriction, since we can reduce every other assumption to this one by a transformation of the whole
group.

The meaning of $\chi$ in this form is that of a cyclic
permutation of the variables. The substitutions $\tau, \chi$ together
generate a group $\tau^\varrho \chi^\lambda$ of degree 21.


% ===== scan leaf 456, printed 435 =====
% !PRINTED 435
\weberpage{435}
Now, in order to determine $\omega$, we make use of the relation
\[
\omega \chi = \chi^2 \omega,
\]
and obtain, if
\[
\omega = \begin{pmatrix} \alpha_1, & \alpha_2, & \alpha_3 \\ \beta_1, & \beta_2, & \beta_3 \\ \gamma_1, & \gamma_2, & \gamma_3 \end{pmatrix}
\]
is set,
\[
\begin{pmatrix} \alpha_2, & \alpha_3, & \alpha_1 \\ \beta_2, & \beta_3, & \beta_1 \\ \gamma_2, & \gamma_3, & \gamma_1 \end{pmatrix}
=
\begin{pmatrix} \beta_1, & \beta_2, & \beta_3 \\ \gamma_1, & \gamma_2, & \gamma_3 \\ \alpha_1, & \alpha_2, & \alpha_3 \end{pmatrix},
\]
hence
\[
\begin{aligned}
\alpha_1 &= \beta_3 = \gamma_2 = \alpha, \\
\alpha_2 &= \beta_1 = \gamma_3 = \beta, \\
\alpha_3 &= \beta_2 = \gamma_1 = \gamma,
\end{aligned}
\]
and $\omega$ thus takes the form
\weq{5}{\omega = \begin{pmatrix} \alpha, & \beta, & \gamma \\ \beta, & \gamma, & \alpha \\ \gamma, & \alpha, & \beta \end{pmatrix}.}

If we express the condition that $\omega$ is to be of degree 2,
the relations
\weq{6}{\begin{aligned} \alpha^2 + \beta^2 + \gamma^2 &= 1 \\ \beta\gamma + \gamma\alpha + \alpha\beta &= 0. \end{aligned}}
result.

From (6) it follows that $(\alpha + \beta + \gamma)^2 = 1$, and if one sets the
determinant of $\omega$ equal to 1, then with the use of (6) it follows that
$(\alpha + \beta + \gamma)^3 = -1$, consequently
\weq{7}{\alpha + \beta + \gamma = -1.}

From (6) and (7) we further infer the relations:
\weq{8}{\alpha = \beta\gamma - \alpha^2, \quad \beta = \gamma\alpha - \beta^2, \quad \gamma = \alpha\beta - \gamma^2}
\weq{9}{\alpha^3 + \beta^3 + \gamma^3 - 3\alpha\beta\gamma = -1,}
of which the latter, which one obtains from
\[
(\alpha + \beta + \gamma)^3 - 3(\beta\gamma + \gamma\alpha + \alpha\beta)(\alpha + \beta + \gamma) = -1
\]
represents the negative value of the determinant of $\omega$.

Finally we further form, according to the relations [\S.\,72, (15)]
\[
\Theta = \chi \tau^3 \omega \tau^4, \quad \Theta^3 = \chi \tau^6 \omega \tau^2 \text{:}
\]
\weq{10}{\Theta = \begin{pmatrix} \gamma \varepsilon^2, & \alpha \varepsilon^6, & \beta \\ \alpha, & \beta \varepsilon^4, & \gamma \varepsilon^5 \\ \beta \varepsilon^3, & \gamma, & \alpha \varepsilon \end{pmatrix}, \quad \Theta^{-1} = \begin{pmatrix} \gamma \varepsilon^5, & \alpha, & \beta \varepsilon^4 \\ \alpha \varepsilon, & \beta \varepsilon^3, & \gamma \\ \beta, & \gamma \varepsilon^2, & \alpha \varepsilon^6 \end{pmatrix}.}


% ===== scan leaf 457, printed 436 =====
% !PRINTED 436
\weberpage{436}
If we now seek the condition that $\Theta^4 = 1$,
we have to set $\Theta^2 = \Theta^{-2}$. If, then, we form $\Theta^2$ and $\Theta^{-2}$
according to (10), one may first try to bring about the agreement
by setting one of the three numbers $\alpha, \beta, \gamma$, say $\gamma = 0$;
but then by (6) and (7) a second, say
$\beta = 0$, must also vanish, and the third must be $\alpha = -1$. But then $\Theta^2$ does not
agree with $\Theta^{-2}$. Hence all three quantities $\alpha, \beta, \gamma$
are different from zero. If we now set the second and third
terms of the first row in the substitutions
$\Theta^2$ and $\Theta^{-2}$ derived from (10) equal to one another, we obtain, when the factors
$\gamma$ and $\alpha$ are discarded:
\[
\alpha(\varepsilon - \varepsilon^{-2}) = \beta(\varepsilon^{-1} - 1), \quad \gamma(\varepsilon^4 - 1) = \beta(\varepsilon^3 - \varepsilon)
\]
or
\[
\alpha(\varepsilon^2 - \varepsilon^{-2}) = \beta(\varepsilon^4 - \varepsilon^{-4}), \quad \gamma(\varepsilon^2 - \varepsilon^{-2}) = \beta(\varepsilon - \varepsilon^{-1}),
\]
and from this, if $h$ denotes an undetermined factor:
\weq{11}{\begin{aligned} \alpha &= h(\varepsilon^4 - \varepsilon^{-4}) \\ \beta &= h(\varepsilon^2 - \varepsilon^{-2}) \\ \gamma &= h(\varepsilon - \varepsilon^{-1}). \end{aligned}}

The factor $h$ is determined according to (7) from the equation
\weq{12}{-1 = h(\varepsilon + \varepsilon^2 + \varepsilon^4 - \varepsilon^{-1} - \varepsilon^{-2} - \varepsilon^{-4})}
and by Bd. I, \S.\,171 we obtain
\weq{13}{h = \frac{i}{\sqrt{7}} = \frac{-1}{\sqrt{-7}}.}

The sign of $\sqrt{7}$ depends on the choice of $\varepsilon$ and
is positive if, e.\,g.,
\[
\varepsilon = e^{\frac{2\pi i}{7}}
\]
is taken. One then obtains for $\alpha, \beta, \gamma$
\[
\alpha = \frac{-2\sin\dfrac{8\pi}{7}}{\sqrt{7}}, \quad \beta = \frac{-2\sin\dfrac{4\pi}{7}}{\sqrt{7}}, \quad \gamma = \frac{-2\sin\dfrac{2\pi}{7}}{\sqrt{7}}.
\]

From (11) and (12) there follow further the formulae
\weq{14}{\begin{aligned}
\alpha\varepsilon + \beta\varepsilon^4 + \gamma\varepsilon^2 &= 1 \\
\alpha\varepsilon^{-1} + \beta\varepsilon^{-4} + \gamma\varepsilon^{-2} &= 1, \\
\alpha + \beta\varepsilon^{-1} + \gamma\varepsilon^2 &= 0 \\
\alpha + \beta\varepsilon + \gamma\varepsilon^{-2} &= 0; \\
\alpha\beta\gamma &= \tfrac{1}{7},
\end{aligned}}
%CONT>>


% ===== scan leaf 458, printed 437 =====
% !PRINTED 437
\weberpage{437}
%<<CONT
and $\alpha$, $\beta$, $\gamma$ are the roots of the cubic equation
\[ \alpha^3 + \alpha^2 - \frac{1}{7} = 0. \]
By \S.\,72, (15), (12), (13) one may now set
\weq{15}{\Theta^2 = \tau^4\,\omega\,\tau^5\,\chi}
and from this one easily obtains, by (4) and (5),
\weq{16}{\Theta^2 = \begin{pmatrix} \beta, & \gamma\varepsilon^3, & \alpha\varepsilon^2 \\ \gamma\varepsilon^{-3}, & \alpha, & \beta\varepsilon^{-1} \\ \alpha\varepsilon^{-2}, & \beta\varepsilon, & \gamma \end{pmatrix}.}

If one compares this with the result that one obtains when $\Theta^2$ is formed directly from (10), one finds:
\weq{17}{\begin{aligned}
\alpha &= \alpha^2\varepsilon^{-1} + \beta^2\varepsilon + \gamma^2\varepsilon^{-2}, & \alpha &= \beta\gamma + \gamma\alpha\varepsilon^2 + \alpha\beta\varepsilon^{-1} \\
\beta &= \alpha^2\varepsilon^{-1} + \beta^2\varepsilon^3 + \gamma^2\varepsilon^{-3}, & \beta &= \beta\gamma\varepsilon^3 + \gamma\alpha + \alpha\beta\varepsilon \\
\gamma &= \alpha^2\varepsilon^2 + \beta^2\varepsilon^3 + \gamma^2\varepsilon^{-2}, & \gamma &= \beta\gamma\varepsilon^{-3} + \gamma\alpha\varepsilon^{-2} + \alpha\beta,
\end{aligned}}

and the comparison with the $\Theta^{-2}$ formed from (10) yields an entirely similar system of formulae, which arises from (17) when $\varepsilon$ is interchanged with $\varepsilon^{-1}$, namely:
\[ \begin{aligned}
\alpha &= \alpha^2\varepsilon + \beta^2\varepsilon^{-1} + \gamma^2\varepsilon^2, & \alpha &= \beta\gamma + \gamma\alpha\varepsilon^{-2} + \alpha\beta\varepsilon \\
\beta &= \alpha^2\varepsilon + \beta^2\varepsilon^{-3} + \gamma^2\varepsilon^3, & \beta &= \beta\gamma\varepsilon^{-3} + \gamma\alpha + \alpha\beta\varepsilon^{-1} \\
\gamma &= \alpha^2\varepsilon^{-2} + \beta^2\varepsilon^{-3} + \gamma^2\varepsilon^2, & \gamma &= \beta\gamma\varepsilon^3 + \gamma\alpha\varepsilon^2 + \alpha\beta.
\end{aligned} \]

But in order that no circular reasoning lie herein, it must be shown that the relations (17) are actually satisfied in consequence of the determinations (11), (12); for only then are the validity of the relation (15) and the equality of $\Theta^2$ with $\Theta^{-2}$ established. For this it suffices to derive two of these relations, say
\[ \begin{aligned}
\alpha &= \alpha^2\varepsilon^{-1} + \beta^2\varepsilon + \gamma^2\varepsilon^{-2} \\
\alpha &= \beta\gamma + \gamma\alpha\varepsilon^2 + \alpha\beta\varepsilon^{-1},
\end{aligned} \]
for once these have been established, one may interchange $\varepsilon$ with $\varepsilon^{-1}$ in them, whereby $\alpha$, $\beta$, $\gamma$ are not changed, and then $\varepsilon$ with $\varepsilon^2$ and with $\varepsilon^4$, whereby $\alpha$, $\beta$, $\gamma$ pass into $\gamma$, $\alpha$, $\beta$ and into $\beta$, $\gamma$, $\alpha$; and in this way one obtains the whole system of formulae (17).

These two relations, however, result from a simple calculation by (11), (12) in the form:
\[ \begin{aligned}
&-(\varepsilon + \varepsilon^2 + \varepsilon^4 - \varepsilon^{-1} - \varepsilon^{-2} - \varepsilon^{-4})(\varepsilon^4 - \varepsilon^{-4}) \\
&= \varepsilon^{-1}(\varepsilon^4 - \varepsilon^{-4})^2 + \varepsilon(\varepsilon^2 - \varepsilon^{-2})^2 + \varepsilon^{-2}(\varepsilon - \varepsilon^{-1})^2 \\
&= (\varepsilon - \varepsilon^{-1})(\varepsilon^2 - \varepsilon^{-2}) + \varepsilon^2(\varepsilon - \varepsilon^{-1})(\varepsilon^4 - \varepsilon^{-4}) \\
&\phantom{=} + \varepsilon^{-1}(\varepsilon^2 - \varepsilon^{-2})(\varepsilon^4 - \varepsilon^{-4}).
\end{aligned} \]


% ===== scan leaf 459, printed 438 =====
% !PRINTED 438
\weberpage{438}
\pdfbookmark[2]{§ 116. Poles and axes of the ternary groups}{bm133}
From these formulae one easily obtains, by (5), (8), (14):
\weq{18}{\begin{aligned}
\omega\Theta &= \begin{pmatrix} \gamma, & \alpha\varepsilon^{-2}, & \beta\varepsilon \\ \alpha\varepsilon^2, & \beta, & \gamma\varepsilon^3 \\ \beta\varepsilon^{-1}, & \gamma\varepsilon^{-3}, & \alpha \end{pmatrix}, \quad
\omega\Theta^2 = \begin{pmatrix} \gamma, & \alpha\varepsilon, & \beta\varepsilon^3 \\ \alpha\varepsilon^{-1}, & \beta, & \gamma\varepsilon^2 \\ \beta\varepsilon^{-3}, & \gamma\varepsilon^{-2}, & \alpha \end{pmatrix}, \\
\Theta\omega &= \omega\Theta^3 = \begin{pmatrix} \alpha, & \beta\varepsilon^{-1}, & \gamma\varepsilon^{-3} \\ \beta\varepsilon, & \gamma, & \alpha\varepsilon^{-2} \\ \gamma\varepsilon^3, & \alpha\varepsilon^2, & \beta \end{pmatrix}.
\end{aligned}}

After this it is now very simple to confirm the characteristic relations of Theorem I., \S.\,72 for our group by actual computation, and with this the ternary group of degree 168 is therefore established\footnote{Cf. F. \gesperrt{Klein}, ``Ueber die Transformation siebenter Ordnung der elliptischen Functionen''. ``Ueber die Auflösung gewisser Gleichungen vom 7\textsuperscript{ten} und 8\textsuperscript{ten} Grade''. Mathem. Annalen, Bd. 14 u. 15. \gesperrt{Klein-Fricke}, Vorlesungen über Modul-Functionen, Bd. I, dritter Abschnitt, Capitel VII. \gesperrt{Gordan}, ``Ueber Gleichungen 7\textsuperscript{ten} Grades mit einer Gruppe von 168 Substitutionen''. Mathem. Annalen, Bd. 20, 25.}).

\begin{center}
\S.\,116.\\[2pt]
\gesperrt{Poles and Axes of the Ternary Groups.}
\end{center}

We shall now, for the sake of brevity, make use of a geometrical mode of expression, by regarding the variables $x_1, x_2, x_3$ as triangular coordinates of a point $x$ in a plane, although imaginary values of $x_1, x_2, x_3$ are also to be admissible.

If $A$ denotes a ternary linear substitution with determinant 1, and if
\weq{1}{(y) = A(x),}
then $y_1, y_2, y_3$ are the coordinates of a second point $(y)$, referred to the same coordinate system, which is derived from $(x)$ by the substitution $A$.

When $x$ traverses a straight line, then $y$ also traverses a straight line, and if the equations of these two lines are
\weq{2}{u_1x_1 + u_2x_2 + u_3x_3 = 0, \quad v_1y_1 + v_2y_2 + v_3y_3 = 0}
then it follows, by substituting (1) in (2), that $u_1, u_2, u_3$ are connected with $v_1, v_2, v_3$ by the \gesperrt{transposed substitution} of $A$:
\[ (u) = A_1(v) \]
%CONT>>


% ===== scan leaf 460, printed 439 =====
% !PRINTED 439
\weberpage{439}
%<<CONT
are connected (\S.\,37, 9.). Thus by $A$ not only is a point derived from every point, but also a straight line from every straight line.

The relation of the points $y$ to the points $x$ expressed by (1) can be designated as a \gesperrt{mapping of the plane into itself}, in which to every point there corresponds a definite other point, and to every straight line a straight line. Such a mapping is also called in geometry a collineation. The point $y$ is called the \gesperrt{image} of the point $x$, and the straight line $v$ the \gesperrt{image} of the straight line $u$. If one has the images of two points, then the connecting line of these images is the image of the connecting line of the two original points.

If $S$ denotes a second ternary linear substitution and
\[ A' = S^{-1}AS \]
the substitution transformed from $A$ by $S$, then
\[ (y') = A'(x') \]
is equivalent to
\[ (y) = A(x), \]
if
\[ (y) = S(y'), \quad (x) = S(x') \]
is set. Now instead of defining a mapping by $S$, one can also understand by it a \gesperrt{coordinate transformation}, by understanding by $x_1', x_2', x_3'$ the coordinates of the point $(x)$ in a new coordinate system determined by $S$. Then the connection between the two points $x$, $y$ is expressed by the transformed substitution $A'$ just as by $A$, and the transformation of the substitutions is thus equivalent to the transformation of the coordinate system. The composition of the transformed substitutions takes place, as we have already seen earlier, just like that of the original ones, and thus by transformation there arises from every group an isomorphic group.

Of great importance is now the investigation of the points that remain unchanged under the mapping by a substitution $A$.

Let us take $A$ in the form:
\[ A = \begin{pmatrix} \alpha, & \beta, & \gamma \\ \alpha', & \beta', & \gamma' \\ \alpha'', & \beta'', & \gamma'' \end{pmatrix}, \]
%CONT>>


% ===== scan leaf 461, printed 440 =====
% !PRINTED 440
\weberpage{440}
%<<CONT
we thus obtain, as the condition that the point $y$ coincide with the point $x$ — since a point is uniquely determined by the \gesperrt{ratios} of its coordinates — the condition that, for a suitably determined coefficient $\lambda$:
\weq{3}{\begin{aligned}
\lambda x_1 &= \alpha\,x_1 + \beta\,x_2 + \gamma\,x_3 \\
\lambda x_2 &= \alpha'x_1 + \beta'x_2 + \gamma'x_3 \\
\lambda x_3 &= \alpha''x_1 + \beta''x_2 + \gamma''x_3,
\end{aligned}}
from which, by elimination of $x_1, x_2, x_3$, one obtains the \gesperrt{cubic} equation for $\lambda$:
\weq{4}{\Lambda = \begin{vmatrix} \alpha - \lambda, & \beta, & \gamma \\ \alpha', & \beta' - \lambda, & \gamma' \\ \alpha'', & \beta'', & \gamma'' - \lambda \end{vmatrix} = 0.}

There are therefore in general \gesperrt{three points} which are their own images, and these we shall call (as with the binary substitutions) the \gesperrt{poles of the substitution $A$}. The three lines joining these poles are then likewise their own images, but in such a way that on each of these three lines there lie only two points which are mapped onto themselves. The images of the remaining points are displaced along the line. These lines we shall call the \gesperrt{axes of the substitution $A$}.

But special cases can now occur which deserve to be noted. Two or even all three poles may coincide in one point, as the example
\[ \begin{pmatrix} \alpha & 0 & 0 \\ \alpha' & \alpha & 0 \\ \alpha'' & \beta'' & \alpha^{-2} \end{pmatrix} \]
shows. Here, if $\alpha'$, $\alpha''$, $\beta''$ are different from zero, two poles coincide in one point, or, if $\alpha = 1$, all three.

More important still is another special case, namely that in which infinitely many points are their own images. If we set aside the case in which all points are their own images, which occurs only in multiplication, then the points which are their own images, if there are infinitely many of them, must lie on a straight line; for this case can occur only if, for a root of (4), the three equations (3) follow from one of them, or, what is the same thing, if together with $\Lambda$ all the first subdeterminants
%CONT>>


% ===== scan leaf 462, printed 441 =====
% !PRINTED 441
\weberpage{441}
%<<CONT
vanish. Such a straight line we shall call a \gesperrt{principal axis of the substitution} $A$. But only certain special substitutions possess principal axes. A substitution with a principal axis has, besides this, also \gesperrt{a pole}, which in special cases may also fall into the principal axis.

In order to survey these relations, let us place the principal axis in the line $x_1=0$; then, if $x_1=0$, we must have $y_1=0$, $y_2:y_3=x_2:x_3$. From this it follows that
\[
\begin{aligned}
y_1 &= \alpha\, x_1\\
y_2 &= \alpha'\, x_1 + \beta\, x_2\\
y_3 &= \alpha''\, x_1 + \beta\, x_3.
\end{aligned}
\]
The pole existing outside the principal axis we obtain by setting $y_1=\alpha\,x_1$, $y_2=\alpha\,x_2$, $y_3=\alpha\,x_3$; then we obtain
\[
(\beta-\alpha)\,x_2+\alpha'\,x_1=0,\quad (\beta-\alpha)\,x_3+\alpha''\,x_1=0,
\]
and this point will fall into the line $x_1=0$ when, and only when, $\beta=\alpha$. This case cannot occur with finite groups (except in the case of multiplication).

If a pole distinct from the principal axis exists, we can place it at the vertex $x_2=0$, $x_3=0$, and the substitution takes the form
\[
y_1=\alpha\,x_1,\quad y_2=\beta\,x_2,\quad y_3=\beta\,x_3,
\]
where $\alpha$ is different from $\beta$.

A straight line $u_1x_1+u_2x_2+u_3x_3=0$ is its own image when, and only when, either $u_2=u_3=0$ or $u_1=0$.

\gesperrt{If, then, a substitution has a pole and a principal axis, there remain unchanged, besides the latter, all the straight lines, and only those, which pass through the pole.}

It should further be remarked that, if three separate poles are present, they cannot fall on a straight line without the line becoming a principal axis. For if a straight line is its own image, then, unless the line corresponds to itself point for point, only two points can lie on it which are their own images.

Summing all this up, we find the following kinds of ternary linear substitutions, where coinciding poles are counted only once:
%CONT>>


% ===== scan leaf 463, printed 442 =====
% !PRINTED 442
\weberpage{442}
\pdfbookmark[2]{§ 117. Application to the group G168. Sevenfold poles}{bm134}
%<<CONT

\begin{tabular}{@{}l l l l@{}}
1. & Substitutions with three poles, & & \\
2. & \textquotedbl & \textquotedbl & two poles,\\
3. & \textquotedbl & \textquotedbl & one pole,\\
4. & \textquotedbl & \textquotedbl & one principal axis and one pole,\\
5. & \textquotedbl & \textquotedbl & one principal axis.\\
\end{tabular}

As the last case, the multiplication would still have to be enumerated, for which every arbitrary point is its own image.

The kind of the substitution is preserved under every transformation.

If a point remains unchanged by a substitution $A$, then it also remains unchanged under every repetition of $A$, hence under $A^2, A^3,\dots$. The poles of $A$ are therefore always found among the points that are their own images under the mapping by $A^2$. But the case may well occur that, for example, $A$ has three poles, while a power, $A^k$, possesses a principal axis. Then two of the poles of $A$ must lie on the principal axis of $A^k$, and the third pole of $A$ is the single pole of $A^k$. For if $A$ is a substitution with three poles, then, by placing the poles at the vertices of the coordinate triangle, we can put it in the normal form
\[
A=\begin{pmatrix}\alpha, & 0, & 0\\ 0, & \beta, & 0\\ 0, & 0, & \gamma\end{pmatrix}
\]
in which $\alpha,\beta,\gamma$ are different from one another. If now $\beta^k=\gamma^k$, so that $\beta$ differs from $\gamma$ by a $k$th root of unity as a factor, then $A^k$ has the line $x_1=0$ for its axis, and the opposite vertex of the coordinate triangle for its pole.

If $A$ is of finite degree $\mu$, then $\alpha,\beta,\gamma$ are $\mu$th roots of unity. If $A^k$ is the lowest power of $A$ that has a principal axis, then $\beta:\gamma$ must be a primitive $k$th and at the same time $\mu$th root of unity, and consequently $k$ must be a proper divisor of $\mu$; if $k$ is relatively prime to $\mu$, then $A^k$ has the same three poles as $A$.

\begin{center}
\S.\,117.

\gesperrt{Application to the group $G_{168}$. Sevenfold poles.}
\end{center}

We shall now seek out the poles and axes of the substitutions of our group $G_{168}$.


% ===== scan leaf 464, printed 443 =====
% !PRINTED 443
\weberpage{443}
This work is considerably lightened by the fact that, from the consideration of the congruence group $L_7$, the degrees and cycles of the group are already known to us (\S.\,68). According to this there are $p^2-1=48$ elements of degree $7$, $A_7$, $\frac14 p(p-3)(p+1)=56$ elements of degree $3$, $A_3$, and $\frac14 p(p-1)^2=63$ elements of degree $4$ or $2$, $A_4, A_2$, in $G_{168}$.

The elements $A_7$ arrange themselves (with the inclusion of the identical substitution) into eight cycles of seven terms, the $A_3$ into 28 cycles of three elements, and the $A_4$ and $A_2$ into 21 cycles of four terms. Each of these last cycles contains one element $A_2$ and two elements $A_4$, and so, collecting these together, we obtain:
\[
\begin{aligned}
&48 \text{ elements of the } 7 \text{th degree}\\
&56 \text{ „} \quad 3 \text{rd „}\\
&42 \text{ „} \quad 4 \text{th „}\\
&21 \text{ „} \quad 2 \text{nd „}.
\end{aligned}
\]
If a point remains unchanged by $\nu$ substitutions of the group (the identical one included), then we shall call such a point a \gesperrt{$\nu$-fold pole} (\S.\,52).

For establishing the poles and axes it suffices to consider one cycle of each kind, since from it the others are derived by transformation (\S.\,69).

If we take the substitution $\tau$ of degree $7$ [\S.\,115, (4)], we find the vertices of the coordinate triangle as three separate poles. All powers of $\tau$ whose exponents are not divisible by 7 have the same three poles.

By applying the substitutions $\chi,\chi^2$ to the coordinates of these three poles, they pass cyclically into one another; but if one applies the substitutions $\Theta,\Theta^2,\Theta^3,\omega,\omega\Theta,\omega\Theta^2,\omega\Theta^3$, which are completely formed in \S.\,115, (5), (10), (16), (18), then the whole coordinate triangle passes into one entirely different from it (e.g.\ by $\omega$ into the triangle with the vertex coordinates $\alpha,\beta,\gamma$; $\beta,\gamma,\alpha$; $\gamma,\alpha,\beta$), and from this it is to be concluded that these poles are not more than seven-fold. There are eight triples of such seven-fold poles, all of which are different from one another, and each of them can be transformed into every other by a substitution of the group $G_{168}$.

We thus have the theorem:
%CONT>>


% ===== scan leaf 465, printed 444 =====
% !PRINTED 444
\weberpage{444}
\pdfbookmark[2]{§ 118. The principal axes}{bm135}
%<<CONT
\gesperrt{1. There are 24 sevenfold poles, which, corresponding to the eight seven-term cycles, divide into eight triples. Each of these 24 points can be transformed into every other by substitutions of the group $G_{168}$.}

\begin{center}
\S.\,118.

\gesperrt{The Principal Axes.}
\end{center}

We pass on to the consideration of the substitutions of degrees $4$ and $2$, as whose representatives we choose $\Theta$ and $\Theta^2$. According to \S.\,115, (10) and \S.\,116 we have to solve for $\Theta$ the cubic equation:
\weq{1}{
\begin{vmatrix}
\gamma\varepsilon^2-\lambda, & \alpha\varepsilon^6, & \beta\\
\alpha, & \beta\varepsilon^4-\lambda, & \gamma\varepsilon^5\\
\beta\varepsilon^3, & \gamma, & \alpha\varepsilon-\lambda
\end{vmatrix} = 0,
}
which, developed, reads:
\weq{2}{
\begin{aligned}
&\lambda^3-\lambda^2(\alpha\varepsilon+\beta\varepsilon^4+\gamma\varepsilon^2)\\
&\quad -\lambda\left[\varepsilon^6(\alpha^2-\beta\gamma)+\varepsilon^3(\beta^2-\gamma\alpha)+\varepsilon^5(\gamma^2-\alpha\beta)\right]\\
&\quad +\alpha^3+\beta^3+\gamma^3-3\alpha\beta\gamma=0.
\end{aligned}
}
But this equation reduces, with the use of the relations (8), (9), (14) of \S.\,115, to
\weq{3}{\lambda^3-\lambda^2+\lambda-1=0,}
and has the roots
\[
\lambda=1,\quad \lambda=\pm i,
\]
and from these the coordinates of the three poles can be computed, which result as rational functions of $\varepsilon$ and $i$.

For the poles of $\Theta^2$ we obtain, according to \S.\,115, (16), first the cubic equation
\[
\begin{vmatrix}
\beta-\lambda, & \gamma\varepsilon^3, & \alpha\varepsilon^2\\
\gamma\varepsilon^{-3}, & \alpha-\lambda, & \beta\varepsilon^{-1}\\
\alpha\varepsilon^{-2}, & \beta\varepsilon, & \gamma-\lambda
\end{vmatrix} = 0,
\]
which, by the relations between $\alpha,\beta,\gamma$, is easily brought into the form
\weq{4}{(\lambda+1)^2(\lambda-1)=0}
and therefore has a double root $\lambda=-1$ and a simple root $\lambda=1$. If we set the value $\lambda=-1$
%CONT>>


% ===== scan leaf 466, printed 445 =====
% !PRINTED 445
\weberpage{445}
%<<CONT
into the equations by which the unchanged point is determined [\S.\,116, (3)], we obtain:

\[
\begin{array}{l}
(\beta+1)\,x_1 + \gamma\,\varepsilon^3 x_2 + \alpha\,\varepsilon^2 x_3 = 0,\\
\gamma\,\varepsilon^{-3} x_1 + (\alpha+1)\,x_2 + \beta\,\varepsilon^{-1} x_3 = 0,\\
\alpha\,\varepsilon^{-2} x_1 + \beta\,\varepsilon\, x_2 + (\gamma+1)\,x_3 = 0.
\end{array}
\]
These three equations all reduce to the single one which one derives, e.g., from the first by multiplication by $\beta$ and application of \S.\,115, (8):

\weq{5}{\alpha\,\gamma\,x_1 + \beta\,\gamma\,\varepsilon^3 x_2 + \alpha\,\beta\,\varepsilon^2 x_3 = 0,}
and which represents a straight line. This line is therefore a \gesperrt{principal axis}.

We obtain 21 such principal axes, since we have 21 substitutions of degree 2.

The functions which, set equal to zero, represent the 21 principal axes can easily be formed from (5), if one transforms the function

\weq{6}{A = \alpha\,\gamma\,x_1 + \beta\,\gamma\,\varepsilon^3 x_2 + \alpha\,\beta\,\varepsilon^2 x_3}
by the substitutions $\tau^r$, $\tau^r\chi$, $\tau^r\chi^2$ ($r = 0, 1, 2, 3, 4, 5, 6$). One thus obtains the 21 functions:

\weq{7}{\begin{array}{l}
A_{1,r} = \alpha\,\gamma\,\varepsilon^r x_1 + \gamma\,\beta\,\varepsilon^{2r+3} x_2 + \alpha\,\beta\,\varepsilon^{4r+2} x_3\\
A_{2,r} = \beta\,\gamma\,\varepsilon^{r+3} x_1 + \alpha\,\beta\,\varepsilon^{2r+2} x_2 + \alpha\,\gamma\,\varepsilon^{4r} x_3\\
A_{3,r} = \alpha\,\beta\,\varepsilon^{r+2} x_1 + \alpha\,\gamma\,\varepsilon^{2r} x_2 + \beta\,\gamma\,\varepsilon^{4r+3} x_3,
\end{array}}
from which it is easy to see that the 21 principal axes are in fact all different from one another.

We summarise this as follows:

\noindent 2.\ \gesperrt{There belong 21 different principal axes to the group $G_{168}$, of which each can be transformed into every other by substitutions of the group.}

Since there are only 21 principal axes, each of them must remain unchanged under eight substitutions, and these eight substitutions form a group. To determine this group, we apply to the expression $A$ represented in (6) the substitution $\omega$ [\S.\,115, (5)], and thereby it passes over into

\weq{8}{\begin{array}{l}
\gamma\,(\alpha^2 + \beta^2\varepsilon^3 + \alpha\,\beta\,\varepsilon^2)\,x_1 + \beta\,(\alpha\,\gamma + \gamma^2\varepsilon^3 + \alpha^2\varepsilon^2)\,x_2\\
{}+ \alpha\,(\gamma^2 + \beta\,\gamma\,\varepsilon^3 + \beta^2\varepsilon^2)\,x_3.
\end{array}}


% ===== scan leaf 467, printed 446 =====
% !PRINTED 446
\weberpage{446}
But by \S.\,115, (14) and (8) we have

\[
\alpha^2 + \beta^2\varepsilon^3 + \alpha\,\beta\,\varepsilon^2 = \alpha^2 + \beta\,(\alpha\,\varepsilon^2 + \beta\,\varepsilon^3) = \alpha^2 - \beta\,\gamma = -\,\alpha,
\]
and if one transforms the three coefficients of the expression (8) in this way, there results

\[
-\,\alpha\,\gamma\,x_1 - \beta\,\gamma\,\varepsilon^3 x_2 - \alpha\,\beta\,\varepsilon^2 x_3,
\]
that is, $A$ changes its sign under application of the substitution $\omega$, and the principal axis $A$ therefore remains unchanged by the substitution $\omega$.

Since $A$ moreover remains unchanged by the substitution $\Theta$, because two poles of $\Theta$ lie on $A$, we have the whole group of the 8th degree by which $A$ remains unchanged, which we call the \gesperrt{group of} $A$ and denote by $G_a$, in the form

\weq{9}{\omega^u\,\Theta^v.}
Among these eight substitutions there are only two, namely the identical one and $\Theta^2$, which leave all points of the line $A$ unchanged. Under the others only two points each remain at rest. For if one calculates, by the path followed here, from \S.\,115, (5), (16) the principal axes of $\omega$, $\omega\Theta$, $\omega\Theta^2$, $\omega\Theta^3$, one obtains, in the notation (7),

\[
A_{2,1},\ A_{3,0},\ A_{3,3},\ A_{2,0},
\]
which are all different from the principal axis $A_{1,0}$ of $\Theta^2$.

Besides the principal axis $A$, the substitution $\Theta^2$ has yet an isolated pole, which we obtain when we choose the root $\lambda = 1$ of the cubic equation (4). Then for the coordinates of this pole there result the equations

\[
\begin{array}{l}
(\beta-1)\,x_1 + \gamma\,\varepsilon^3 x_2 + \alpha\,\varepsilon^2 x_3 = 0,\\
\gamma\,\varepsilon^{-3} x_2 + (\alpha-1)\,x_2 + \beta\,\varepsilon^{-1} x_3 = 0,\\ % [sic] printed x_2 twice, expected x_1 in first term
\alpha\,\varepsilon^{-2} x_1 + \beta\,\varepsilon\, x_2 + (\gamma-1)\,x_3 = 0,
\end{array}
\]
and from these one obtains

\weq{10}{x_1 : x_2 : x_3 = \alpha\,\gamma\,\varepsilon^4 : \beta\,\gamma\,\varepsilon : \alpha\,\beta\,\varepsilon^2.}
Of such points there are 21. The point (10) we denote for the moment by $a$. Now the point $a$ remains unchanged not only by $\Theta$, but also, as a calculation on the basis of the formulae \S.\,115, (17), (14) shows, by $\omega$, and consequently by the whole group $G_a$, and is therefore a pole of order at least eight. But since it can be brought by 21 substitutions into 21 different positions, it is also not \gesperrt{more} than
%CONT>>


% ===== scan leaf 468, printed 447 =====
% !PRINTED 447
\weberpage{447}
%<<CONT
eight-fold. Since it remains unchanged by the substitutions of degree 2 $\omega$, $\Theta\omega$, $\Theta^2\omega$, $\Theta^3\omega$, and since the poles of these substitutions lie on $A$ (because $A$ remains unchanged by them), while $a$ does not lie on $A$, the principal axes of these four substitutions $\omega$, $\Theta\omega$, $\Theta^2\omega$, $\Theta^3\omega$ must pass through the point $a$.

The poles of the four substitutions $\omega$, $\Theta\omega$, $\Theta^2\omega$, $\Theta^3\omega$, which we shall call $a_1$, $a_2$, $a_3$, $a_4$, are, as images of the point $a$, likewise eight-fold poles, and must, like $a$, be poles of substitutions of the fourth order, all of which are in any case different from $\Theta$ and $\Theta^3$, because otherwise $\Theta^2$ would have to be equal to one of those four substitutions, which is not the case. Under $\omega$ only two points of the axis $A$, namely the poles of $\omega$ and of $\Theta^2\omega$, remain unchanged.

Let us now consider the two other poles $c_1$, $c_2$ of the substitution $\Theta$. These points likewise lie on the principal axis $A$, but are different from the points $a_1$, $a_2$, $a_3$, $a_4$, because these, as already remarked, do not occur among the poles of $\Theta$. Now $\omega\Theta\omega = \Theta$, hence

\[
\omega\,\Theta(x) = \Theta\,\omega(x),
\]
and if now $(x) = \Theta(x)$, then also

\[
\omega(x) = \Theta\,\omega(x);
\]
that is, if $(x)$ are the coordinates of a pole of $\Theta$, then the $\omega(x)$ have the same meaning. Under the transformation $\omega$, therefore, the poles of $\Theta$ merely pass into one another. The pole $a$ remains unchanged under $\omega$; $c_1$ and $c_2$, on the other hand, cannot remain unchanged, because they do not occur among the poles of $\omega$, and must therefore pass into one another.

We thus have the following overview:

\begin{center}
\begin{tabular}{@{}lllllll@{}}
Unchanged under & $\omega$ and $\Theta^2\omega$ & on & $A$ & only the points & $a_1, a_3$ & \\
\textquotedbl\ \textquotedbl & $\Theta\omega$\ \textquotedbl\ $\Theta^3\omega$ & \textquotedbl & $A$ & \textquotedbl\ \textquotedbl\ \textquotedbl & $a_2, a_4$ & \\
\textquotedbl\ \textquotedbl & $\Theta$\ \textquotedbl\ $\Theta^3$ & \textquotedbl & $A$ & \textquotedbl\ \textquotedbl\ \textquotedbl & $c_1, c_2.$ &
\end{tabular}
\end{center}

A point $x$ of $A$ different from these six passes into itself only under the substitutions 1 and $\Theta^2$, and may therefore be called a \gesperrt{two-fold pole}. Then the following theorem holds:

\noindent 3.\ \gesperrt{A two-fold pole passes, under the group of degree 8 derived from $\omega$ and $\Theta$, into four different positions.}


% ===== scan leaf 469, printed 448 =====
% !PRINTED 448
\weberpage{448}
\pdfbookmark[2]{§ 119. The threefold and sixfold poles}{bm136}
If, namely, $x$ is a two-fold pole, and $x$ passes by $\omega$ into $x'$ and by $\Theta$ into $x''$, then not only are $x'$ and $x''$ different from $x$, but they are also different from one another, because, if they were identical, $x$ would remain unchanged by $\omega\Theta^3$. By $\omega\Theta$, $x$ then passes into a fourth position $x'''$, and $x'''$ is different both from $x'$ and from $x''$, because otherwise $x$ would remain unchanged by $\omega\Theta\omega = \Theta$ or by $\omega\Theta\Theta^{-1} = \omega$. Since $x$ remains unchanged by $\Theta^2$, $x$ passes by $\Theta^3$ into $x''$, by $\Theta^2\omega = \omega\Theta^2$ into $x'$, and by $\omega\Theta^3$ into $x'''$.

Since we have two points $c_1$, $c_2$ on each principal axis, there are 42 in all, and each of them can be carried into every other by substitutions of the group. From this it follows that these poles are not more than four-fold. Hence the theorem:

\noindent 4.\ \gesperrt{There are 21 eight-fold and 42 four-fold poles. On each principal axis lie four eight-fold and two four-fold poles. Through each eight-fold pole pass four principal axes. Every eight-fold pole can be carried into every other eight-fold pole, and every four-fold pole into every other four-fold pole.}

\begin{center}
\S.\,119.

\gesperrt{The three- and six-fold poles.}
\end{center}

We now turn to the consideration of the substitutions of the third order and of their poles, and choose as representative of such a substitution

\[
\chi = \begin{pmatrix} 0 & 0 & 1\\ 1 & 0 & 0\\ 0 & 1 & 0 \end{pmatrix},
\]
whose poles are obtained from

\[
\lambda\,x_1 = x_3, \quad \lambda\,x_2 = x_1, \quad \lambda\,x_3 = x_2
\]
From this it follows that $\lambda^3 = 1$; hence $\lambda$ is a third root of unity, and if therefore $\varrho$ denotes an imaginary third root of unity, the three poles result:

\weq{1}{\begin{array}{l}
x_1 = x_2 = x_3\\
x_1 = \varrho\,x_2 = \varrho^2 x_3\\
x_1 = \varrho^2 x_2 = \varrho\,x_3.
\end{array}}


% ===== scan leaf 470, printed 449 =====
% !PRINTED 449
\weberpage{449}
Between these three poles, however, there is an essential difference. For if we apply the substitution $\omega$ to them, the first of them remains unchanged. It is therefore at least six-fold and belongs to the substitutions
\[
1, \chi, \chi^2, \omega, \omega\chi, \omega\chi^2.
\]
The two other poles, namely
\weq{2}{x_1 = \varrho x_2 = \varrho^2 x_3, \quad x_1 = \varrho^2 x_2 = \varrho x_3,}
pass into one another under the substitution $\omega$, and we must show that they are altered by every other substitution of the group except $\chi, \chi^2$, and hence that they are only three-fold. If for the moment we denote the two points (2) by $\pi$ and $\pi'$, we can first say that all substitutions of $G_{168}$ which leave $\pi$ unchanged must form a group $G_\pi$,
and indeed a divisor of $G_{168}$; the group $G_\pi$ in turn contains the cyclic group 1, $\chi$, $\chi^2$.

If $\sigma$ is any substitution of $G_\pi$, then $\pi$ occurs among the poles of $\sigma$. From this it follows that $\sigma$ can be neither of degree $7$ nor of degree $4$. For the coordinates of the poles of these two degrees are rational functions of $\varepsilon, i$, while the coordinates of $\pi$ contain the irrationality $\varrho$, which is independent of these (Bd. I, \S.\,134, cf. the supplement at the end of this volume). But that $\sigma$ also cannot be of degree $2$ follows from the fact that $\pi$ lies on none of the principal axes. For if it lay on a principal axis, then, because $\varrho$ cannot be expressed rationally in terms of $\varepsilon$, as is to be derived from (2) by means of the equation
\[
1 + \varrho + \varrho^2 = 0
\]
one of the 21 relations of \S.\,118, (7) would have to hold, namely
\[
\begin{aligned}
\alpha &= \beta\varepsilon^{r+3} &&= \gamma\varepsilon^{-2r+1}, \\
\alpha &= \beta\varepsilon^{-3r+3} &&= \gamma\varepsilon^{-r+1}, \\
\alpha &= \beta\varepsilon^{2r+3} &&= \gamma\varepsilon^{3r+1},
\end{aligned}
\]
of which evidently none is possible. The group $G_\pi$ therefore contains, besides the identical substitution, only substitutions of degree $3$, and the degree of $G_\pi$ must be a power of 3. But since 168 is not divisible by 9, the degree of $G_\pi$ must be equal to 3. From this it follows:

5. \gesperrt{There are 56 three-fold poles, which belong two by two to the same substitution, and accordingly arrange themselves in 28 pairs. Each of these poles}
%CONT>>


% ===== scan leaf 471, printed 450 =====
% !PRINTED 450
\weberpage{450}
%<<CONT
\gesperrt{can be transformed into every other by substitutions of the group} $G_{168}$\gesperrt{.}

If we set
\weq{3}{T = x_1 + x_2 + x_3,}
then $T = 0$ is the equation of the line joining the points $\pi, \pi'$. The function $T$ remains unchanged by $\chi$ and changes its sign by $\omega$. But by the 28 substitutions $\Theta^s\tau^r$ it passes over into 28 \gesperrt{different} functions $T_{s,r}$, which we can easily represent, with the aid of the formulae of \S.\,115, (11), (14), as follows:
\weq{4}{
\begin{aligned}
T_{0,r} &= \varepsilon^r x_1 + \varepsilon^{2r} x_2 + \varepsilon^{4r} x_3 \\
h\,T_{1,r} &= \beta^2\varepsilon^{1+r} x_1 + \alpha^2\varepsilon^{2+2r} x_2 + \gamma^2\varepsilon^{4+4r} x_3 \\
-h\,T_{2,r} &= \alpha^2\varepsilon^{2+r} x_1 + \gamma^2\varepsilon^{4+2r} x_2 + \beta^2\varepsilon^{1+4r} x_3 \\
h\,T_{3,r} &= \gamma^2\varepsilon^{4+r} x_1 + \beta^2\varepsilon^{1+2r} x_2 + \alpha^2\varepsilon^{2+4r} x_3,
\end{aligned}
}
and these, set equal to zero, represent 28 different straight lines, which can all be derived from one of them by substitutions of the group.

On the line $T$ lie the three points with the coordinates
\[
\begin{aligned}
\alpha\gamma &: \beta\gamma : \beta\alpha \\
\beta\alpha &: \alpha\gamma : \beta\gamma \\
\beta\gamma &: \beta\alpha : \alpha\gamma,
\end{aligned}
\]
which arise from (10), \S.\,118 by the substitutions $\tau^{-4}, \tau^{-4}\chi, \tau^{-4}\chi^2$, and thus belong to the eight-fold poles.

6. \gesperrt{On every line $T$ lie three eight-fold and two three-fold poles.}

It still remains to investigate the third pole $\pi_0$ of the substitution $\chi$, with the coordinates
\[
x_1 = x_2 = x_3
\]
of which we have already proved that it remains unchanged by the six substitutions
\weq{5}{1, \chi, \chi^2; \omega, \omega\chi, \omega\chi^2}
The question now arises whether this pole is not more than six-fold.

If besides (5) there are still other substitutions which leave the pole $\pi_0$ unchanged, then these must necessarily be of degree $2$, since $\pi_0$ is neither a four-fold nor a seven-fold pole, and because, further, the degree of the group belonging to $\pi_0$ cannot be divisible by 9.


% ===== scan leaf 472, printed 451 =====
% !PRINTED 451
\weberpage{451}
\pdfbookmark[2]{§ 120. The configuration of the group G168}{bm137}
Now the point $\pi_0$ lies on the three principal axes of the substitutions of degree $2$, $\omega, \omega\chi, \omega\chi^2$. If there is yet a further substitution of degree $2$ by which $\pi_0$ is left unchanged, then $\pi_0$ must lie on the principal axis of this substitution.

But if we set $x_1 = x_2 = x_3$ in the equations of the principal axes [\S.\,118, (7)], each of these three equations is satisfied for only one value of $r$; hence $\pi_0$ cannot lie on more than three of the principal axes, and $\pi_0$ is six-fold.

From this the theorem:

7. \gesperrt{There are 42 six-fold poles, which can all be derived from one of them by substitutions of the group. Through each of these poles pass three principal axes of the group.}

The pole $\pi_0$ lies on none of the lines $T_{s,r}$. This is easily shown by setting $x_1 = x_2 = x_3$ in the expressions (4), and then putting in for $\alpha^2, \beta^2, \gamma^2$ their expressions in terms of $\varepsilon$. One then sees easily (with an application of the irreducibility of the cyclotomic equation) that none of these expressions vanishes.

\S.\,120.

\gesperrt{The Configuration of the Group $G_{168}$.}

The totality of the straight lines and points that we have considered in the preceding sections we shall call the \gesperrt{configuration of the group} $G_{168}$. The word here has the same meaning in which it has been used in recent times in geometry\footnote{The expression was introduced by \gesperrt{Reye}: \glqq{}Geometrie der Lage.\grqq{} Bd. I, 2. Aufl. (1876).}.

In what follows we describe this configuration somewhat more closely, without adding anything new, only in order to let the theorems of the last sections stand out more vividly and more clearly.

We have in this configuration:
\[
\begin{aligned}
&\text{21 lines } A \text{ (the principal axes)}, \\
&\text{21 points } a \text{ (the eight-fold poles)}, \\
&\text{28 lines } T, \\
&\text{28 points } t \text{ (the six-fold poles)}.
\end{aligned}
\]


% ===== scan leaf 473, printed 452 =====
% !PRINTED 452
\weberpage{452}
The whole system is carried over into itself by 168 substitutions of the group.

\[
\begin{aligned}
&\text{On every line } A \text{ lie four points } a. \\
&\text{Through every point } a \text{ pass four lines } A. \\
&\text{On every line } T \text{ lie three points } a. \\
&\text{Through every point } t \text{ pass three lines } A.
\end{aligned}
\]

\gesperrt{The points} $a$ \gesperrt{and} $t$ \gesperrt{together form the complete system of all points of intersection of the lines} $A$\gesperrt{.}

For 21 straight lines intersect in 210 points, and of these, three or six coincide when three or four of these lines pass through one point. But we have
\[
210 = 21.6 + 28.3.
\]

\gesperrt{On every line} $A$ \gesperrt{lie four points} $t$\gesperrt{.}

For since through every point $t$ there pass three lines $A$, and on every line $A$ equally many points $t$ must lie (because every line $A$ is transformable into every other), it follows, if $x$ is the number of these points, that $x . 21 : 3 = 28$, hence $x = 4$. In the same way one concludes:

\gesperrt{Through every point} $a$ \gesperrt{pass four lines} $T$\gesperrt{.}

None of the points $t$ lies on a line $T$.
If we denote the $\nu$-fold poles by $P_\nu$, we thus have the following systems of poles:
\[
\begin{aligned}
&21 \text{ eightfold poles } P_8, \\
&28 \text{ sixfold} \;\text{\textquotedbl}\; P_6, \\
&42 \text{ fourfold} \;\text{\textquotedbl}\; P_4, \\
&56 \text{ threefold} \;\text{\textquotedbl}\; P_3, \\
&24 \text{ sevenfold} \;\text{\textquotedbl}\; P_7, \\
&\text{infinitely many twofold poles } P_2.
\end{aligned}
\]

The $P_2$ are all the points of the principal axes distinct from $P_4, P_6, P_8$; of the points $P_4$, two at a time lie on a line $A$; of the points $P_3$, two at a time lie on a line $T$; the points $P_7$ do not lie on the lines $A$ (that they also do not lie on $T$ will appear later).

The system of the axes corresponds exactly to the system of the poles. Every straight line that passes through an eightfold pole is an axis (among them the principal axes, four lines $T$ and two connecting lines of the $P_8$ with $P_4$). Accordingly
%CONT>>


% ===== scan leaf 474, printed 453 =====
% !PRINTED 453
\weberpage{453}
\pdfbookmark[2]{§ 121. Invariant curves of the group G168}{bm138}

%<<CONT
the $P_8$ may fittingly be called the \gesperrt{principal poles}, and the axes passing through them may be denoted by $A_8, A_6, A_4, A_2$.

To be singled out further are the 28 pairs of connecting lines of a $P_6$ with the two associated $P_3$, which we may denote by $A_3$, and finally the 24 connecting lines of each two $P_7$ belonging together, as whose representatives the sides of the coordinate triangle are to be regarded, and which may be denoted by $A_7$.

\begin{center}
\S.\,121.\\
\gesperrt{Invariant curves of the group $G_{168}$.}
\end{center}

A form of degree $\mu$, $\varphi\,(x_1, x_2, x_3)$, will in general pass over, under the substitutions of the group, into 168 different forms. For special forms of $\varphi$, however, this number can diminish, and then the substitutions under which $\varphi$ remains unchanged form a group $G'$ contained in $G_{168}$, whose degree must be a divisor of 168. The number of forms into which $\varphi$ passes over is the index of the divisor $G'$, and each of these different forms of $\varphi$ remains unchanged under a group conjugate to $G'$. The form $\varphi$ is an absolute invariant of the group $G'$.

The equation $\varphi = 0$ represents a curve referred to the coordinate triangle $x_1, x_2, x_3$, and this curve is mapped by the substitutions of $G_{168}$ onto other curves. If the 168 image curves are not all different from one another, then the curve $\varphi$ remains unchanged under the substitutions of a group $G'$, or changes only by a constant factor, and is thus a relative invariant of $G'$. A straight line remains unchanged under substitutions other than the identical one only when it belongs to the axes described in the preceding section.

All properties and relations between points, lines and curves that are indestructible under linear transformation (the so-called projective properties of geometry) are preserved in the images. If therefore, e.\,g., a straight line is a tangent or inflectional tangent or double tangent of a curve, then all images of the straight line stand in the same relation to the images of the curve. Likewise, if a point
%CONT>>


% ===== scan leaf 475, printed 454 =====
% !PRINTED 454
\weberpage{454}
\pdfbookmark[2]{§ 122. The first invariant of the group G168 and the fundamental curve}{bm139}

%<<CONT
double point, point of inflection, cusp, point of contact of a double tangent, etc.

Among the forms $\varphi$, those are now above all of importance to us which remain unchanged under the whole group, the invariants, of which there are here, since the group is simple, only absolute ones (\S.\,40). If $\Phi\,(x_1, x_2, x_3)$ is such a form, then the curve represented by the equation $\Phi = 0$ shall be called an invariant curve of the group. There are 168 mappings of the plane onto itself under which all points of the curve pass into points of the same curve. If any point lies on this curve, then all its image points also lie on it.

In general the number of points of the curve thus connected with one another is 168, in any case not greater. If it is smaller, then the points must be \gesperrt{poles}, and the number of image points is a divisor of 168 (namely 24 for the $P_7$, 21 for the $P_8$, 28 for the $P_6$, 42 for the $P_4$, 56 for the $P_3$, and 84 for a system of $P_2$ belonging together). If a pole of one of these kinds lies on the curve, then all poles of the same kind lie on it.

\begin{center}
\S.\,122.\\
\gesperrt{The first invariant of the group $G_{168}$ and the fundamental curve.}
\end{center}

In order now to form the invariants of our group, we seek out forms of the three variables $x_1, x_2, x_3$ that remain unchanged by the application of the three substitutions $\chi, \tau, \omega$ (\S.\,115). This suffices, since the whole group can be composed from these three substitutions (\S.\,72). Now $\chi$ is a cyclic permutation of the three variables $x_1, x_2, x_3$, and $\tau$ denotes the substitution

\[
\begin{pmatrix} x_1, & x_2, & x_3 \\ \varepsilon x_1, & \varepsilon^2 x_2, & \varepsilon^4 x_3 \end{pmatrix}.
\]

If, then, a term $x_1^{h_1} x_2^{h_2} x_3^{h_3}$ occurs in an invariant, in which $h_1, h_2, h_3$ are non-negative integers, then invariability under $\tau$ requires that

\weq{1}{h_1 + 2h_2 + 4h_3 \equiv 0\ (\mathrm{mod}\ 7)}
%CONT>>


% ===== scan leaf 476, printed 455 =====
% !PRINTED 455
\weberpage{455}

%<<CONT
and the invariability under $\chi$, that besides this one term the two corresponding ones

\[
x_2^{h_1} x_3^{h_2} x_1^{h_3}, \quad x_3^{h_1} x_1^{h_2} x_2^{h_3}
\]

also occur in the function, unless $h_1 = h_2 = h_3$. To this is added the condition of invariability under $\omega$. We now wish to see how we can satisfy these requirements, and first of all in such a way that the degree $m$ of the invariant, that is, the sum $h_1 + h_2 + h_3$, becomes as small as possible.

The condition (1) evidently cannot be satisfied if $m < 3$. If this sum is $= 3$, then $h_1 = h_2 = h_3 = 1$ must hold; but the product $x_1 x_2 x_3$ is evidently not unchanged under $\omega$. The smallest value that comes into consideration is therefore $m = 4$, and we therefore have to seek out all non-negative solutions of

\[
\begin{aligned}
h_1 + h_2 + h_3 &= 4 \\
h_1 + 2h_2 + 4h_3 &\equiv 0\ (\mathrm{mod}\ 7)
\end{aligned}
\]
%CONT>>

If we eliminate $h_1$, it follows that

\[
h_2 + 3h_3 \equiv 3\ (\mathrm{mod}\ 7),
\]

and from this result the only possible solutions:

\[
\begin{aligned}
h_3 &= 0, & h_2 &= 3, & h_1 &= 1 \\
h_3 &= 1, & h_2 &= 0, & h_1 &= 3 \\
h_3 &= 3, & h_2 &= 1, & h_1 &= 0,
\end{aligned}
\]

and consequently the \gesperrt{only} form of degree $4$ unchanged by $\chi$ and $\tau$

\weq{2}{f\,(x_1, x_2, x_3) = x_1^3 x_3 + x_2^3 x_1 + x_3^3 x_2.}

In order to test the influence of the substitution $\omega$ on the function $f$, we set

\weq{3}{
\begin{aligned}
x_1 &= \alpha y_1 + \beta y_2 + \gamma y_3 \\
x_2 &= \beta y_1 + \gamma y_2 + \alpha y_3 \\
x_3 &= \gamma y_1 + \alpha y_2 + \beta y_3,
\end{aligned}
}

in which the coefficients $\alpha, \beta, \gamma$ have the meaning fixed in \S.\,115, and assume that by (3) the transformation

\weq{4}{F\,(y_1, y_2, y_3) = f\,(x_1, x_2, x_3)}

is accomplished. We form the second derivatives of $F$ with respect to $y_1, y_2, y_3$, with regard to (3) and to the formulae

\[
\begin{aligned}
&\tfrac{1}{6} f''(x_1, x_1) = x_1 x_3, \quad \tfrac{1}{6} f''(x_2, x_2) = x_2 x_1, \quad \tfrac{1}{6} f''(x_3, x_3) = x_3 x_2 \\
&\tfrac{1}{3} f''(x_2, x_3) = x_3^2, \quad \tfrac{1}{3} f''(x_3, x_1) = x_1^2, \quad \tfrac{1}{3} f''(x_1, x_2) = x_2^2.
\end{aligned}
\]


% ===== scan leaf 477, printed 456 =====
% !PRINTED 456
\weberpage{456}

%<<CONT
From this we obtain:

\[
\begin{aligned}
\tfrac{1}{6} F''(y_1, y_1) &= x_1 x_3 \alpha^2 + x_2 x_1 \beta^2 + x_3 x_2 \gamma^2 \\
&\quad + x_3^2 \beta\gamma + x_1^2 \alpha\gamma + x_2^2 \alpha\beta, \\
\tfrac{1}{3} F'''(y_2, y_3) &= 2x_1 x_3 \beta\gamma + 2x_2 x_1 \alpha\gamma + 2x_3 x_2 \alpha\beta \\
&\quad + x_3^2 (\beta\gamma + \alpha^2) + x_1^2 (\alpha\gamma + \beta^2) + x_2^2 (\alpha\beta + \gamma^2).
\end{aligned}
\]
% [sic] mixed term printed with triple prime F''' here against F'' just below for the same pairing

Now since the solutions of the system (3) are of the same form ($y_1 = \alpha x_1 + \beta x_2 + \gamma x_3, \dots$), one obtains from this:

\[
\begin{aligned}
\tfrac{1}{6} F''(y_1, y_1) - y_1 y_3 &= \\
x_1 x_3 (\alpha^2 - \alpha\beta - \gamma^2) + x_2 x_1 (\beta^2 - \beta\gamma - \alpha^2) &+ x_2 x_3 (\gamma^2 - \beta^2 - \alpha\gamma), \\
\tfrac{1}{3} F''(y_2, y_3) - y_3^2 &= \\
x_1^2 (\beta^2 + \alpha\gamma - \gamma^2) + x_2^2 (\gamma^2 + \alpha\beta - \alpha^2) &+ x_3^2 (\alpha^2 + \beta\gamma - \beta^2).
\end{aligned}
\]

But the coefficients on the right-hand side of these equations $(\alpha^2 - \alpha\beta - \gamma^2) \dots$ turn out, from the values of $\alpha, \beta, \gamma$ [\S.\,115, (11)], to vanish, and so, if we further apply a cyclic permutation of the $x$, which entails a cyclic permutation of the $y$, we obtain:

\[
\begin{aligned}
&\tfrac{1}{6} F'''(y_1, y_1) = y_1 y_3, \quad \tfrac{1}{6} F'''(y_2, y_2) = y_2 y_1, \quad \tfrac{1}{6} F'''(y_3, y_3) = y_3 y_2, \\
&\tfrac{1}{3} F''(y_2, y_3) = y_3^2, \quad \tfrac{1}{3} F''(y_3, y_1) = y_1^2, \quad \tfrac{1}{3} F''(y_1, y_2) = y_2^2.
\end{aligned}
\]
% [sic] prime-count on F wavers between the two lines above, reproduced as printed

Accordingly, by \gesperrt{Euler}'s theorem [Bd.\ I, \S.\,17, (6)], we obtain:

\[
F\,(y_1, y_2, y_3) = y_1^3 y_3 + y_2^3 y_1 + y_3^3 y_2,
\]

and with this it is proved that the function $f\,(x_1, x_2, x_3)$ is in fact an invariant of our group $G_{168}$. Apart from a constant factor that may be arbitrarily attached, it is the only invariant of degree $4$.

The equation

\weq{5}{f\,(x_1, x_2, x_3) = 0}

signifies, if $x_1, x_2, x_3$ are coordinates in the plane, a \gesperrt{curve of the fourth order}, which we shall call the \gesperrt{fundamental curve} of the group, or the curve $f$; and there are 168 mappings of the plane onto itself under which to every point of this curve there corresponds a point of the curve. If a point lies on the fundamental curve, then all its image points lie on it as well. The number of distinct image points of a point can sink to a smaller number than 168 — and indeed always to a divisor of 168 — only if the point is a pole.


% ===== scan leaf 478, printed 457 =====
% !PRINTED 457
\weberpage{457}
The straight line $x_1 = 0$ cuts the curve in three coincident
points at $x_3 = 0$ and in a fourth point $x_2 = 0$ separated from
them. The vertices of the coordinate triangle are therefore points of
inflection of the curve, and the sides are the inflectional tangents.
If we denote the sides of the coordinate triangle by 1, 2, 3 and the
opposite vertices by the same figures, then side 1 is the inflectional
tangent at point 2, side 2 the inflectional tangent at point 3, and
side 3 the inflectional tangent at point 1.

We now investigate the position of the poles and axes with respect to
the fundamental curve. Since the vertices of the coordinate triangle
belong to the sevenfold poles, we conclude first of all
\gesperrt{that all sevenfold poles $P_7$ lie on the fundamental
curve.}

They are, like the vertices themselves, all \gesperrt{points of inflection} of
the curve, and these arrange themselves into eight inflection-point
triangles.

\gesperrt{The threefold poles} $P_3$ \gesperrt{likewise lie on the fundamental curve;
for if one sets}
\[
x_2 = \varrho\,x_1, \quad x_3 = \varrho^2 x_1
\]
or
\[
x_2 = \varrho^2 x_1, \quad x_3 = \varrho\,x_1,
\]
in which $\varrho$ is a cubic root of unity, then
$f(x_1, x_2, x_3)$ reduces to
\[
1 + \varrho + \varrho^2 = 0.
\]
To find the points of intersection of their connecting line
\[
T = x_1 + x_2 + x_3 = 0
\]
with the fundamental curve, one sets $x_3 = -x_1 - x_2$, whereby
$f(x_1, x_2, x_3) = 0$ passes into
\[
x_1^3(x_1+x_2) - x_2^3 x_1 + (x_1+x_2)^3 x_2 = (x_1^2+x_1x_2+x_2^2)^2 = 0
\]
Since the left side is a square, the four points of intersection
coincide twice in pairs, and it follows that the line $T$ is a
double tangent of the fundamental curve.

\gesperrt{The 28 lines $T$ are double tangents of the fundamental
curve; their points of contact are the 56 poles $P_3$.}

From this it follows incidentally that the points $P_7$ do not lie
on the lines $T$, since a double tangent can have no further point
of intersection with the curve besides the points of contact.

\gesperrt{The six- and eightfold poles do not lie on the fundamental
curve.}


% ===== scan leaf 479, printed 458 =====
% !PRINTED 458
\weberpage{458}
For the points $P_6$ this is immediately evident, if one sets
[according to \S.\,119, (1)] $x_1 = x_2 = x_3 = 1$, whereby
$f(x_1, x_2, x_3) = 3$, and hence does not vanish. To demonstrate the same for
the $P_8$, one sets, according to \S.\,118 (10):
\[
x_1 = \alpha\gamma\varepsilon^4, \quad x_2 = \beta\gamma\varepsilon,
\quad x_3 = \alpha\beta\varepsilon^2,
\]
whereby, since $\alpha\beta\gamma = \frac{1}{7}$ [\S.\,115, (14)],
\[
f(x_1, x_2, x_3) = \frac{1}{7}(\alpha^3\gamma^2 + \beta^3\alpha^2 + \gamma^3\beta^2)
\]
If one sets in this [according to \S.\,115, (8)]:
\[
\alpha^3 = \alpha\beta\gamma - \alpha^2, \quad \beta^3 = \alpha\beta\gamma - \beta^2,
\quad \gamma^3 = \alpha\beta\gamma - \gamma^2,
\]
one obtains
\weq{6}{f(x_1, x_2, x_3) = \tfrac{1}{7}\alpha\beta\gamma(\alpha^2+\beta^2+\gamma^2) - \tfrac{1}{7}(\alpha^2\gamma^2+\beta^2\alpha^2+\gamma^2\beta^2).}
But from \S.\,115, (6), (7) it follows that
\[
\alpha^2 + \beta^2 + \gamma^2 = 1,
\]
\[
\alpha^2\gamma^2 + \beta^2\alpha^2 + \gamma^2\beta^2 = -\alpha\beta\gamma(\alpha+\beta+\gamma) = \frac{1}{7},
\]
and therefore
\[
f(x_1, x_2, x_3) = \frac{1}{42}
\]
is different from zero.

The intersection points of the 21 principal axes with one another are, in their entirety,
the poles $P_6$ and $P_8$, and consequently two principal axes never
intersect on the fundamental curve. The intersection points of the principal axes with the
fundamental curve can therefore only be four-fold or two-fold poles. To
decide this question, we make the following consideration. If among
the intersection points of the principal axis $A$ with the fundamental curve a $P_2$
occurs, then all four intersection points are distinct from one another and
are all two-fold poles. But if a $P_4$ lies on $A$ and $f$,
then the second $P_4$ lying on $A$ also lies on $f$, and there can be
no other intersection point of $A$ and $f$, because such a point would be
a $P_2$ and would entail four further $P_2$. The four
intersection points of $A$ and $f$ must therefore coincide in pairs, and
$A$ is a double tangent. But now we have seen that the 28
lines $T$ are double tangents, and a curve of the fourth order can
have no more than 28 double tangents. Hence the assumption is
inadmissible, and it follows that the intersection points of the fundamental curve with
the principal axes are two-fold poles.

If, then, we denote the special two-fold poles that lie on the
curve $f$ by $\bar P_2$, we have the following distinguished
systems of points:
%CONT>>


% ===== scan leaf 480, printed 459 =====
% !PRINTED 459
\weberpage{459}
\pdfbookmark[2]{§ 123. The higher invariants}{bm140}
%<<CONT
\[
24\,P_7,\ 56\,P_3,\ 84\,\bar P_2 \ \text{ on the curve } f.
\]
\[
21\,P_8,\ 42\,P_4,\ 28\,P_6 \ \text{ not on the curve } f.
\]
All other points of the curve $f$ are carried by the substitutions of
the group into 168 distinct points.

From this it follows, moreover, immediately that the curve $f$ can
have no double point, and hence still less can it split into curves of
lower degree. For suppose it had a double point; then all its image
points would also have to be double points, and since a curve of the
fourth order, even when it splits, cannot have more than six double
points unless it has infinitely many double points, i.\,e. doubly
counted parts of the curve, $f$ would have to be the square of a
quadratic form. The root of $f$ would then also have to be an
invariant, whereas there are no quadratic invariants.

One can, incidentally, also easily show directly that the fundamental
curve has no double point; for the conditions for a double point
are:
\[
f'(x_1) = 3x_1^2x_3 + x_2^3 = 0,
\]
\[
f'(x_2) = 3x_2^2x_1 + x_3^3 = 0,
\]
\[
f'(x_3) = 3x_3^2x_2 + x_1^3 = 0,
\]
and these equations cannot be satisfied otherwise than when
$x_1, x_2, x_3$ vanish (\S.\,86).

For the sake of brevity, we shall call a system of points, each of
which can be transformed into every other by substitutions of the
group, a \gesperrt{system of connected points}.

\S.\,123.

\gesperrt{The higher invariants.}

Further invariants of our group can be derived, according to Theorem 4.,
\S.\,40, by forming covariants of the form $f$. We fix our attention
on the covariants discussed in general terms in \S.\,87, and first form
the \gesperrt{Hesse} covariant
\weq{1}{\Delta = \frac{1}{54}
\begin{vmatrix}
f''(x_1,x_1), & f''(x_1,x_2), & f''(x_1,x_3) \\
f''(x_2,x_1), & f''(x_2,x_2), & f''(x_2,x_3) \\
f''(x_3,x_1), & f''(x_3,x_2), & f''(x_3,x_3)
\end{vmatrix},}
%CONT>>


% ===== scan leaf 481, printed 460 =====
% !PRINTED 460
\weberpage{460}
%<<CONT
which, developed, takes the form
\weq{2}{\Delta = 5x_1^2x_2^2x_3^2 - x_1x_3^5 - x_2x_1^5 - x_3x_2^5.}
On the curve $\Delta$, then, lie the vertices of the coordinate triangle,
and consequently all the sevenfold poles $P_7$ lie on $\Delta$ and
form the complete system of the points of intersection of $f$ and $\Delta$.

A further covariant, namely one of degree 14,
we obtain from \S.\,87, (3), namely the determinant:
\weq{3}{C = \frac{1}{9}
\begin{vmatrix}
f''(x_1,x_1), & f''(x_1,x_2), & f''(x_1,x_3), & \Delta'(x_1) \\
f''(x_2,x_1), & f''(x_2,x_2), & f''(x_2,x_3), & \Delta'(x_2) \\
f''(x_3,x_1), & f''(x_3,x_2), & f''(x_3,x_3), & \Delta'(x_3) \\
\Delta'(x_1), & \Delta'(x_2), & \Delta'(x_3), & 0
\end{vmatrix}.}
The determinant $C$ can, though somewhat laboriously, be computed
from (3), and takes the expression:
\weq{4}{C = \Sigma x_1^{14} - 34x_1x_2x_3\,\Sigma x_1^{10}x_3 - 250x_1x_2x_3\,\Sigma x_1^8x_2^3 + 375x_1^2x_2^2x_3^2\,\Sigma x_1^6x_2^2 + 18\,\Sigma x_1^7x_2^7 - 126x_1^3x_2^3x_3^3\,\Sigma x_1^3x_2^2,}
in which the symbol $\Sigma$ means that the sum is to be taken of the
three terms which are obtained from the first when $x_1, x_2, x_3$
are cyclically interchanged. The computation of the first term
$x_1^{14}$ already suffices to show that the curve $C = 0$
does \gesperrt{not} pass through the vertices of the coordinate triangle.

A fourth invariant, namely one of degree 21,
we obtain according to \S.\,87, (6), by forming the functional determinant
of the three forms $f$, $\Delta$, $C$:
\weq{5}{K = \frac{1}{14}
\begin{vmatrix}
f'(x_1), & \Delta'(x_1), & C'(x_1) \\
f'(x_2), & \Delta'(x_2), & C'(x_2) \\
f'(x_3), & \Delta'(x_3), & C'(x_3)
\end{vmatrix},}
which likewise can be computed from this. We give here only the
first three terms, from which one sees that this curve too
does not pass through the vertices of the coordinate triangle
\weq{6}{K = x_1^{21} + x_2^{21} + x_3^{21} + \cdots{}^{1)}.}

$^{1)}$ The fully computed expressions are to be found in the
memoir of \gesperrt{Gordan}: ``Ueber die typische Darstellung
der ternären biquadratischen Form $f = x_1^3x_2 + x_2^3x_3 + x_3^3x_1$''
[Mathem. Annalen, Bd. XVII, S. 366 (1880)]. Also in \gesperrt{Klein-Fricke},
Modulfunctionen, Bd. I, S. 734.


% ===== scan leaf 482, printed 461 =====
% !PRINTED 461
\weberpage{461}
\pdfbookmark[2]{§ 124. The complete system of invariants}{bm141}
\begin{center}
§.\,124.

\gesperrt{The full system of invariants.}
\end{center}

We can now prove that the forms $f, \Delta, C, K$ are a
full system of invariants of the group, i.\,e. that all invariants
of the group can be represented as integral rational functions of these
four. We rely in this on the theorem of
\gesperrt{Bezout} (Bd.\,I, §.\,49), that two curves of the $m$th
and $n$th order which have more than $mn$ points in common
must have a common curve part. Points of contact are
in this to be counted as double or multiple. All intersection points of two
invariant curves either form one connected system, or they
split into systems of connected points, and all points of one
such system are to be counted equally often.

Let $\Phi$ be an invariant of the $m$th order which does not contain the function
$f$ as a factor, and among the intersection points of the curves
$\Phi$ and $f$ let the poles $P_7$ occur $h_1$ times, the poles $P_3$ $h_2$ times,
the poles $\overline{P}_2$ $h_3$ times. Besides these there shall also occur $h_4$
systems of 168 connected points each, of which moreover
(in the case of contact) several systems may coincide into one counted multiply.
Then the number of all intersection points of the two curves is
\[
4m = 24h_1 + 56h_2 + 84h_3 + 168h_4,
\]
and hence
\weq{1}{m = 6h_1 + 14h_2 + 21h_3 + 42h_4,}
in which $h_1, h_2, h_3, h_4$ denote non-negative integers, which
also cannot all vanish. The smallest value that $m$
can have is therefore 6, and an invariant curve of the $6$th
order must pass through the 24 points $P_7$, as we have
already proved for $\Delta$. If $\Delta'$ is a second invariant
of the $6$th order, which must therefore also pass through the points $P_7$,
then we can determine the constant $a$ in $\Delta' - a\Delta$ so
that the invariant curve $\Delta' - a\Delta = 0$ passes through
some $25$th point of $f$. But then, if
$\Delta' - a\Delta$ does not vanish identically,
$\Delta' - a\Delta$ must by \gesperrt{Bezout}'s theorem be divisible
by $f$. The quotient would be an invariant of the second order, which does not
exist; consequently $\Delta' - a\Delta$ must be identically zero.


% ===== scan leaf 483, printed 462 =====
% !PRINTED 462
\weberpage{462}
If $h_1 = 2$, $h_2 = h_3 = h_4 = 0$, then $m = 12$. An
invariant curve $\Phi$ of the $12$th order must touch the curve $f$ at
the 24 points $P_7$, and if in $\Phi - a\Delta^2$ we determine the
constant $a$ suitably, it follows that this function
must be divisible by $f$. The quotient, as an invariant
of the $8$th order, can only be of the form $bf^2$, and accordingly
$\Phi$ is of the form $a\Delta^2 + bf^3$. In the same way we can conclude
that an invariant of the $18$th order must have the form
$a\Delta^3 + b\Delta f^3$.

\begin{center}
\gesperrt{All invariants of the $6$th, $12$th, $18$th
order are therefore rational functions of $\Delta$ and $f$.}
\end{center}

The next value that $m$ can have according to (1) is $m = 14$. In
this case $h_2 = 1$, while $h_1, h_3, h_4$ are equal to zero.
All invariants of the $14$th order therefore pass through the 56
points $P_3$, and these form the complete system of intersection points
of such an invariant curve with the fundamental curve. This holds also of
the invariant $C$. If we have a second invariant $C'$ of the $14$th
order, we can again, as above, determine the constant $a$ so
that $C' - aC$ is divisible by $f$. The quotient is
an invariant of the $10$th order, and from this it follows, since apart
from $f\Delta$ there is no invariant of the $10$th order:

\begin{center}
\gesperrt{Every invariant of the $14$th order is representable in the
form:}
\end{center}
\[
aC + bf^2\Delta,
\]
\begin{center}
\gesperrt{wherein $a, b$ are constants. Conversely, every expression of
this form is an invariant of the $14$th order.}
\end{center}

Next, $m$ can according to (1) have the value 20, namely for
$h_1 = h_2 = 1$. An invariant of the $20$th order must therefore pass
through the points $P_7$ and $P_3$, that is, through the 80
points of intersection of $f$ with $\Delta C$. From this we can conclude, just as
before, that an invariant of the $20$th degree is representable in
the form
\[
f(a\Delta^2 + bf^3) + cC\Delta.
\]
wherein $a, b, c$ are arbitrary constants. There is no
independent invariant of the $20$th order.

Let us now assume that $K'$ and $K$ are two invariants
of the $21$st order; both must by (1) pass through the 84 points
$\overline{P}_2$, and consequently one can determine $a$ so that
$K' - aK$ becomes divisible by $f$. The quotient would be an invariant
%CONT>>


% ===== scan leaf 484, printed 463 =====
% !PRINTED 463
\weberpage{463}
%<<CONT
of degree $17$, which does not exist, and consequently $K'$ is
identical with $aK$.

\begin{center}
\gesperrt{There is therefore, apart from a constant factor, only
one invariant of degree $21$.}
\end{center}

But now the system of the 21 principal axes is also an invariant
of degree $21$, and from this it is to be concluded:

\begin{center}
\gesperrt{The invariant $K$ splits into 21 linear factors which,
set equal to zero, represent the principal axes of the group.}
\end{center}

We can then form an invariant of degree $42$,
namely
\weq{2}{\Delta^7 - kC^3 = Q,}
in which $k$ is an arbitrary constant, and this constant can
be determined so that the curve $Q$ passes through an arbitrarily given
point on $f$, and it must then also pass through all the points
connected with this one. Thus we can determine $k$ in $Q$ so
that the curve $Q$ cuts out of the curve $f$ an arbitrarily
given system of connected points.

If now $\Phi$ is an arbitrary invariant not divisible by $f$,
which passes $h_1, h_2, h_3$ times through the poles $P_7, P_3, \overline{P}_2$,
and moreover through arbitrary systems $S_1, S_2, \ldots$
of connected points on $f$, which may also partly coincide,
then we first form, according to (2), the forms
$Q_1, Q_2, \ldots$, which vanish in the systems $S_1, S_2, \ldots$,
and then the form
\[
\Psi = \Phi - a\Delta^{h_1} C^{h_2} K^{h_3} Q_1 Q_2 \ldots
\]

The curve $\Psi$ passes, for every value of the constant $a$, through
all the points of intersection of $\Phi$ with $f$, and if we therefore
determine $a$ so that $\Psi$ passes through some point of $f$
different from these, then $\Psi$ is divisible by $f$, hence
\[
\Phi = a\Delta^{h_1} C^{h_2} K^{h_3} Q_1 Q_2 \cdots + f\Phi_1.
\]
In this, $\Phi_1$ is now again an invariant, but of lower
degree than $\Phi$, and by complete induction the
theorem is hereby proved:

\begin{center}
\gesperrt{Every invariant of the group can be represented as an
integral rational function of the four fundamental invariants
$f, \Delta, C, K$.}
\end{center}

If one determines the constant $k$ in (2) so that the curve $Q$ passes
through one of the poles $\overline{P}_2$, then it can pass through no point not
%CONT>>


% ===== scan leaf 485, printed 464 =====
% !PRINTED 464
\weberpage{464}
%<<CONT
connected with $\overline{P}_2$ on the curve $f$; for it
cannot pass through the points $P_7, P_3$, because in these either
$\Delta$ or $C$ vanishes, but neither can it pass through a
point of the fundamental curve which is not a pole, because it would then
have to pass through the 168 connected points, and could not pass
through the pole $\overline{P}_2$. But then the constant $h$
can be determined so that $K^2 - hQ$ becomes divisible by $f$.

The quotient is an invariant of degree lower than
$42$, but in any case of even degree, and if we therefore
represent it by the fundamental invariants, $K$ cannot
occur in it. From this it follows:

\begin{center}
\gesperrt{The invariant $K^2$ can be expressed rationally in terms of
$f, \Delta, C$.}
\end{center}

If, following this theorem, we represent $K^2$ in the form:
\weq{3}{K^2 = \Sigma\, af^{\nu}\Delta^{\nu_1}C^{\nu_2},}
then in this sum, in which the $a$ are numerical
coefficients, only such terms can occur in which
\[
2\nu + 3\nu_1 + 7\nu_2 = 21,
\]
and by now counting off which values of $\nu, \nu_1, \nu_2$
can occur, we obtain the result that between the
forms
\[
\begin{gathered}
K^2,\quad \Delta^7,\quad C^3,\quad f\Delta^4 C,\quad f^2\Delta C^2,\quad f^3\Delta^5,\\
f^4\Delta^2 C,\quad f^6\Delta^3,\quad f^7 C,\quad f^9\Delta
\end{gathered}
\]
there exists a linear relation with numerical coefficients.
If we represent this relation in the form:
\weq{4}{K^2 = \Phi\Delta + \Psi C,}
in which $\Phi, \Psi$ are likewise invariants, we can draw
from it a further geometric conclusion.

The curves $\Delta$ and $C$ certainly do not intersect on the curve
$f$, because the intersection points of $\Delta$ with $f$ are, all of them, the
$P_7$, those of $C$ with $f$ the $P_3$. But if $\Delta$ and
$C = 0$, then $K = 0$ also, and consequently
\gesperrt{all intersection points of $\Delta$ and $C$ lie on the 21 principal axes of the
group}\footnote{The relation (3) has been completely computed by \gesperrt{Gordan} by the
methods of the theory of invariants (Mathem.\
Annalen, Bd.\ XVII, S.\,371).

We give here the formula there obtained in
the symbols used by us:
\[
K^2 = C^3 - 88f^2\Delta C^2 + 16.63f\Delta^4 C + 17.64f^4\Delta^2 C
- 256f^7 C + 27.64\Delta^7 - 128.469f^3\Delta^5 + 43.512f^6\Delta^3
- 2048f^9\Delta.
\]}.


% ===== scan leaf 486, printed 465 =====
% !PRINTED 465
\weberpage{465}

The relation (3) can be used to eliminate from an expression in the invariants all powers of $K$ except the first, and from this there follows further:

\begin{quote}
\gesperrt{An invariant of even degree can be represented rationally in terms of $f, \Delta, C$; an invariant of odd degree, as a product of $K$ with an invariant of even degree.}
\end{quote}

\begin{center}
\rule{4cm}{0.4pt}
\end{center}


% ===== scan leaf 487, printed 466 =====
\weberpage{466}% page number not printed on this page
\pdfbookmark[1]{Fifteenth Section. The Form Problem of the Group G168 and the Theory of Equations of the Seventh Degree.}{bm142}
\pdfbookmark[2]{§ 125. The resolvents of the form problem}{bm143}
% !PRINTED none

\begin{center}
\gesperrt{Fifteenth Chapter.}

\textbf{The Form Problem of the Group $G_{168}$ and the Theory of Equations of the Seventh Degree.}

\rule{4cm}{0.4pt}

\S.\,125.

\gesperrt{The resolvents of the form problem.}
\end{center}

The form problem for the group $G_{168}$ (\S.\,43) yields, according to the general theory, first of all an equation of degree 168, whose coefficients are invariants. But every divisor of the group leads to a resolvent of lower degree, and indeed of degree equal to the index of the divisor.

We shall here consider only the two most interesting cases of such divisors, namely the octahedral group

\weq{1}{G_{24} = \chi^\lambda \omega^\mu \Theta^\nu,}

and the group of degree 21

\weq{2}{G_{21} = \tau^\varrho \chi^\lambda \ (\S.\,72)}

which lead us to resolvents of the seventh and eighth degrees.

As regards first the resolvents of the seventh degree, in forming them we may start from the principal axes of the group. A principal axis, namely, remains unchanged by the group $\omega^\mu \Theta^\nu$, and is carried by the whole group $G_{24}$ into three different lines. There must therefore be seven triples of principal axes, which are determined by an equation of the seventh degree.

If, following \S.\,118, (7), we set:

\weq{3}{\begin{aligned}
A_1 &= \alpha\gamma x_1 + \gamma\beta\varepsilon^3 x_2 + \alpha\beta\varepsilon^2 x_3 \\
A_2 &= \beta\gamma\varepsilon^3 x_1 + \alpha\beta\varepsilon^2 x_2 + \alpha\gamma x_3 \\
A_3 &= \alpha\beta\varepsilon^2 x_1 + \alpha\gamma x_2 + \beta\gamma\varepsilon^3 x_3,
\end{aligned}}
%CONT>>


% ===== scan leaf 488, printed 467 =====
% !PRINTED 467
\weberpage{467}

%<<CONT
then $A_1=0$, $A_2=0$, $A_3=0$ are the equations of three of these axes. Under the substitution $\chi$ the functions $A_1, A_2, A_3$ pass cyclically into one another. Under the substitution $\Theta$ the functions $A_1, A_2, A_3$ undergo the following permutation:

\[
\begin{pmatrix} A_1, & A_2, & A_3 \\ A_1, & -A_3, & A_2 \end{pmatrix},
\]

as a simple calculation shows, if one actually carries out the substitution $\Theta$ and makes use of the formulae of \S.\,115.

To indicate at least the course of the calculation, let it be remarked that $A_1$ passes, under the substitution $\Theta$ [\S.\,115, (10)], into

\[
\alpha(\gamma^2\varepsilon^2 + \beta\gamma\varepsilon^3 + \beta^2\varepsilon^5)x_1 + \gamma(\alpha^2\varepsilon^6 + \beta^2 + \alpha\beta\varepsilon^2)x_2 + \beta(\alpha\gamma + \gamma^2\varepsilon + \alpha^2\varepsilon^3)x_3
\]

But by \S.\,115, (8), (17), (11):

\[
\begin{aligned}
\gamma^2\varepsilon^2 + \beta\gamma\varepsilon^3 + \beta^2\varepsilon^5 &= \gamma^2\varepsilon^2 + \alpha^2\varepsilon^3 + \beta^2\varepsilon^5 + \alpha\varepsilon^3 \\
&= \alpha(\varepsilon^4+\varepsilon^{-4}) = h(\varepsilon-\varepsilon^{-1}) = \gamma,
\end{aligned}
\]

and hence the coefficient of $x_1$ becomes equal to $\alpha\gamma$, as in $A_1$, and in the same way one transforms the remaining expressions. Now since $\omega$ can be composed from $\Theta$ and $\chi$ (\S.\,72), it follows that the quantities $A_1^2, A_2^2, A_3^2$ are only permuted among themselves by the group $G_{24}$, and that accordingly their symmetric functions are roots of resolvents of degree 7.

The simplest symmetric function of these quantities is the sum

\weq{4}{A_1^2 + A_2^2 + A_3^2,}

and this, multiplied by a suitable numerical factor, we shall introduce as the unknown of the resolvent.

If we order the sum (4) according to $x_1, x_2, x_3$, we obtain for it an expression of the form

\[
\lambda(x_1^2+x_2^2+x_3^2) + \mu(x_2x_3+x_3x_1+x_1x_2),
\]

in which, by (3):

\[
\begin{aligned}
\lambda &= (\alpha^2\gamma^2\varepsilon^{-1} + \beta^2\gamma^2\varepsilon^{-2} + \alpha^2\beta^2\varepsilon^3)\varepsilon \\
\mu &= 2\alpha\beta\gamma(\alpha\varepsilon+\beta\varepsilon^{-3}+\gamma\varepsilon^2)\varepsilon.
\end{aligned}
\]

Here the factor $\varepsilon$ has been taken out of the bracket, in order that the other factor may remain unchanged under the interchange of $\varepsilon$ and $\varepsilon^2$, and may therefore be expressed rationally in terms of $\sqrt{-7}$ (Bd. I, \S.\,171).


% ===== scan leaf 489, printed 468 =====
% !PRINTED 468
\weberpage{468}

According to the formulae \S.\,115, (14) one obtains first of all very simply:

\[
\mu = \frac{2\varepsilon}{7},
\]

and for $\lambda$ one obtains, if one substitutes the values \S.\,115, (11) for $\alpha, \beta, \gamma$, and uses the formulae \S.\,115, (12), (13):

\weq{5}{\begin{aligned}
\lambda &= -\varepsilon\,\frac{1+\sqrt{-7}}{14} \\
\frac{\mu}{\lambda} &= -\frac{1-\sqrt{-7}}{2}.
\end{aligned}}

If therefore one sets

\[
A_1^2+A_2^2+A_3^2 = \lambda z,
\]

then it follows that

\weq{6}{z = x_1^2+x_2^2+x_3^2 - \frac{1-\sqrt{-7}}{2}(x_2x_3+x_3x_1+x_1x_2),}

and this function we shall introduce as a root of the resolvent of degree 7. The remaining roots are obtained from it by the substitutions $\tau^r$, so that they all, in the common form

\weq{7}{\begin{aligned}
z_r &= \varepsilon^{2r}x_1^2 + \varepsilon^{4r}x_2^2 + \varepsilon^{r}x_3^2 \\
&\quad - \frac{1-\sqrt{-7}}{2}(\varepsilon^{-r}x_2x_3+\varepsilon^{-2r}x_3x_1+\varepsilon^{-4r}x_1x_2) \\
&\qquad r=0, 1, 2, 3, 4, 5, 6
\end{aligned}}

are contained.

The coefficients of the equation of degree 7 whose roots are the seven quantities $z_r$ are invariants of our group, whose degree can easily be specified.

For if $a_\nu$ is the coefficient of $z^{7-\nu}$ in this equation, after the coefficient of the seventh power has been brought to 1, then $a_\nu$ is an invariant of degree $2\nu$. Hence first of all we must have $a_1=0$, because there is no quadratic invariant. The remaining coefficients are expressed linearly and homogeneously in terms of the following combinations of invariants:

\weq{8}{\begin{aligned}
a_2 &\text{ by } f, \\
a_3 &\text{ \textquotedbl} \ \varDelta, \\
a_4 &\text{ \textquotedbl} \ f^2, \\
a_5 &\text{ \textquotedbl} \ \varDelta f, \\
a_6 &\text{ \textquotedbl} \ \varDelta^2, f^3, \\
a_7 &\text{ \textquotedbl} \ C, \varDelta f^2
\end{aligned}}


% ===== scan leaf 490, printed 469 =====
% !PRINTED 469
\weberpage{469}

and there are thus eight numerical coefficients to compute, which
can be found by comparing a few terms in the expressions of the
$a_\nu$ in terms of the roots on the one hand, and in terms of the invariants (\S.\,123)
on the other, and which, besides rational numbers, can contain only
$\sqrt{-7}$.

We shall at first compute these coefficients only under the
assumption that $f=0$ ${}^1$).

One then needs only to compute the first terms (with $x_1^5x_2$, $x_1^{10}x_2^2$)
in the power sums $\Sigma z_r^3$ and $\Sigma z_r^6$ and to apply the
\gesperrt{Newton} formulae. The last coefficient results
from the one term $x_1^{14}$ in the product of the $z_r$. One
finds first
\[
\Sigma z_r^3 = -3\cdot7\cdot\frac{1-\sqrt{-7}}{2}\,x_1^5x_2\dots
\]
\[
\Sigma z_r^6 = \left[6\cdot7 + 15\cdot7\left(\frac{1-\sqrt{-7}}{2}\right)^2\right]x_1^{10}x_2^2\dots
\]
and from these the sought resolvent
\weq{9}{z^7 - 7\cdot\frac{1-\sqrt{-7}}{2}\,\Delta z^4 - 7\cdot\frac{5+\sqrt{-7}}{2}\,\Delta^2 z - C = 0.}
If one wishes to reduce this equation to another which
depends only on the ratios of the $x_1, x_2, x_3$, one sets
\weq{10}{z = \frac{Cu}{\Delta^2}, \quad g = \frac{\Delta^7}{C^3},}
whereby one obtains from (9) an equation in which only the
one parameter $g$ still occurs, namely
\weq{11}{u^7 - 7\,\frac{1-\sqrt{-7}}{2}\,g\,u^4 - 7\,\frac{5+\sqrt{-7}}{2}\,g^2\,u - g^2 = 0.}
If one makes the same substitution in the general resolvent,
in which $f$ is not set $=0$, one has to introduce a
further parameter
\weq{12}{h = \frac{f\Delta^4}{C^2}}
as well, and the coefficients of the resolvent become, apart from
%CONT>>

\par\noindent\rule{0.3\linewidth}{0.4pt}\par
${}^1$) This case is of particular interest because, like
the icosahedral equation, it leads to the transformation equations from the
theory of elliptic functions.


% ===== scan leaf 491, printed 470 =====
% !PRINTED 470
\weberpage{470}

%<<CONT
numerical factors, as follows easily from (8), are in order
\[
h,\ g,\ h^2,\ hg,\ (g^2,h^3),\ (g^2,h^2g),
\]
in which $(g^2,h^3)$, $(g^2,h^2g)$ denote linear homogeneous expressions
with numerical coefficients.

The calculation can be carried out just as in the special case
above. For simplification one may set $x_3=0$ and still obtains
equations enough for the determination of all the coefficients.
One finds in this way the coefficients in order:
\[
7\,\frac{1-\sqrt{-7}}{2}\,h,
\]
\[
-7\,\frac{1-\sqrt{-7}}{2}\,g,
\]
\[
-7\,(4+\sqrt{-7})\,h^2,
\]
\[
14\,(2+\sqrt{-7})\,hg,
\]
\[
-7\,\frac{5+\sqrt{-7}}{2}\,g^2 - 7\cdot\frac{7+3\sqrt{-7}}{2}\,h^3,
\]
\[
-g^2 + \frac{167-7\sqrt{-7}}{2}\,gh^2.
\]
The functions $g, h$ are fractional invariants, which depend only
on the ratios $x_1:x_2:x_3$, and we can express every other invariant
of the same property \gesperrt{rationally} in terms of $g, h$. For if
we represent such an invariant as a quotient of two forms of equal
degree without common divisor, then numerator and denominator
must be integral invariants of equal degree, because two fractional
functions of the variables $x$ represented in simplest form can be
equal to one another only when numerator and denominator are
separately equal to one another up to constant factors. Now here
numerator and denominator cannot be of odd degree, because
otherwise they would have the common factor $K$ (\S.\,126). Hence
numerator and denominator are representable rationally in terms of $f, C, \Delta$;
and if a term occurring in the numerator or in the denominator is
\weq{13}{f^a\,\Delta^b\,C^c,}
and $m$ is the degree of numerator and denominator, then
\weq{14}{4a + 6b + 14c = m.}


% ===== scan leaf 492, printed 471 =====
% !PRINTED 471
\weberpage{471}
\pdfbookmark[2]{§ 126. Reduction of the general resolvent of the seventh degree to the special one}{bm144}
If we now set in (13)

\[
\varDelta = (g\,C^3)^{\frac{1}{7}}, \quad f = h\,C^2\,(g\,C^3)^{-\frac{4}{7}},
\]

then, with the use of (14), the expression (13) becomes

\[
g^{b+2c}\,h^a\,C^{\frac{m}{14}}\,g^{-\frac{m}{7}},
\]

and in the numerator and denominator the factor $C^{\frac{m}{14}}
g^{-\frac{m}{7}}$ can be cancelled, so that everything is expressed
rationally in terms of $g$ and $h$.

Just like $g, h$, the function $u$ also depends only on the ratio of
the variables $x_1 : x_2 : x_3$, and it follows easily that every
other function with the same property which, like $u$, admits the
substitutions of the group $G_{24}$, is representable rationally in
terms of $u, g, h$. For such a function can first of all, by the
general theorems of \S.\,43, be represented as a rational function of
$u$ and the invariants, and indeed in only one way as an integral
function of $u$ which does not exceed the $6$th degree (Bd.\,I,
\S.\,142). The coefficients in this representation are invariants,
and since $u$ depends only on the ratios, the coefficients too can
depend only on the ratios, and are therefore representable rationally
in terms of $g, h$.

\begin{center}
\S.\ 126.

\gesperrt{Reduction of the general resolvent of the seventh degree to
the special one.}
\end{center}

In the preceding section we became acquainted with two forms of the
resolvent of degree $7$ of the form problem, of which the one, the
special, holds for the case that $f = 0$, that is, for the case that
the point sought lies on the fundamental curve, and the general for
the case that it has an arbitrary position.

The quantities $u, g, h$, when they refer to the point $(x)$, we
shall denote by

\weq{1}{u_x = \dfrac{\varDelta^2 z}{C}, \quad g_x = \dfrac{\varDelta^7}{C^3}, \quad h_x = \dfrac{f\varDelta^4}{C^2}}

The general resolvent shall then be denoted by

\weq{2}{R(u_x, g_x, h_x) = R_x = 0}

and the special one arises from it when
%CONT>>


% ===== scan leaf 493, printed 472 =====
% !PRINTED 472
\weberpage{472}

%<<CONT
one sets $h_x = 0$. The quantities $u_x, g_x, h_x$ depend only on the
variables $x_1, x_2, x_3$.

Now, as \gesperrt{Klein} has remarked loc.\,cit.\ ${}^1$), the solution
of the general resolvent can be reduced to that of the special one,
if one adjoins the root of a \gesperrt{biquadratic equation}.

To prove this, we introduce beside the point $(x)$
a second point $(y)$ and form the polar of $f$:
\weq{3}{f_1(x,y) = y_1f'(x_1) + y_2f'(x_2) + y_3f'(x_3).}
If we then subject $(x)$ and $(y)$ simultaneously to the same substitution
of the group $G_{168}$, the function $f_1(x,y)$ remains
unchanged (Bd.\ I, \S.\,60). We now take the point $(x)$ arbitrarily,
but require of the point $(y)$ that it shall lie simultaneously
on the fundamental curve and on the polar of the point $(x)$,
so that simultaneously
\weq{4}{f(y_1,y_2,y_3) = 0, \quad f_1(x,y) = 0}
shall hold. Then to each point $x$ there correspond four points $y$, and
if we put for $(x)$ a point connected with it, then
each of these four points $(y)$ likewise passes into a connected
point. The function $h_y$ is now $=0$ and $g_y$ is a four-valued
function of the point $(x)$. Symmetric functions
of these four values remain unchanged under the substitutions
of $G_{168}$, and consequently $g_y$ is the \gesperrt{root of a biquadratic
equation whose coefficients depend rationally on the $g_x, h_x$}.
The root of this biquadratic equation must be
adjoined.

The function $u_y$ is a root of the special resolvent
\weq{5}{R(u_y, g_y, 0) = R_y = 0.}
We now denote the four points $y$ belonging to the same $x$
by $y, y', y'', y'''$ and form the symmetric function of these
four points
\weq{6}{(t-u_y)(t-u_{y'})(t-u_{y''})(t-u_{y'''}) = \Phi(t)}
for an indeterminate $t$. This function remains unchanged under
all substitutions of the group $G_{24}$, and is therefore rationally
expressible in terms of $u_x, g_x, h_x$.

\par\noindent\rule{0.3\linewidth}{0.4pt}\par
${}^1$) Mathematische Annalen, Bd.\ XV, S.\ 280.


% ===== scan leaf 494, printed 473 =====
% !PRINTED 473
\weberpage{473}
\pdfbookmark[2]{§ 127. Permutation group of seven digits of degree 168}{bm145}
Since, if we reserve to ourselves the determination of $x$, equation (3) can represent any arbitrary straight line, we can choose $x$ in such a way that among the four points $y, y', y'', y'''$ no two connected points occur. If we then set $u_y$ for the indeterminate quantity $t$ in (6), we obtain a rational equation

\weq{7}{\Psi (u_y, u_x, g, h) = 0,}

and this equation is no longer satisfied if for $(x)$ and $(y)$ we make a substitution from $G_{168}$ by which $u_x$ is changed and consequently $u_y$ is carried over into a value different from $u_y, u_{y'}, u_{y''}, u_{y'''}$.

Equation (7), regarded as an equation for $u_x$, therefore has only one root in common with the general resolvent $R_x$, and if one seeks the greatest common divisor of $R$ and $\Psi$, one obtains $u_x$ expressed rationally in terms of $u_y, g, h$. With this the general resolvent is reduced to the special one.

\begin{center}\rule{0.3\linewidth}{0.4pt}\end{center}

The group $G_{168}$ contains, further, a divisor of degree 21, $G_{21}$, which is generated by the substitutions $\chi, \tau$, and which gives rise to a resolvent of degree 8. As a root of this resolvent one can simply take the product $x_1 x_2 x_3$. The coefficients of this resolvent are invariants which rise up to degree 24. We shall not enter further into this resolvent here.

\begin{center}
\S.\,127.

\gesperrt{Permutation group of seven figures of degree 168.}
\end{center}

Since, as we have seen, the form problem of the ternary substitution group of degree 168 has a resolvent of degree 7, and the Galois resolvent of this equation of degree 7 is consequently of degree 168, it follows from this that in the general permutation group of seven figures, whose degree is $1.2.3.4.5.6.7 = 5040 = 168.30$, a divisor of degree 168 must be contained. The existence of this divisor was first
%CONT>>


% ===== scan leaf 495, printed 474 =====
% !PRINTED 474
\weberpage{474}
%<<CONT
\gesperrt{Kronecker} recognised\footnote{\gesperrt{Kronecker}, \glqq Ueber Gleichungen 7\textsuperscript{ten} Grades\grqq{}. Monatsberichte der Berliner Akademie 1858.}, and its closer investigation is of great importance for the general theory of the equation of degree 7.

We can obtain this permutation group by seeking out the permutations which are produced by the substitutions of the group $G_{168}$ among the quantities $z_0, z_1, z_2, z_3, z_4, z_5, z_6$ (\S.\,125). Here it is then to be observed, with regard to composition, that if the two linear substitutions $\xi_1, \xi_2$ effect the permutations $\pi_1, \pi_2$, the composed permutation $\pi_1 \pi_2$ is produced by $\xi_2 \xi_1$. For according to the definition one obtains $\xi_2 \xi_1 (x)$ by applying to the variables $(x)$ first the substitution $\xi_1$ and to the result the substitution $\xi_2$. Likewise $\pi_1 \pi_2$ denotes the permutation which results when one applies to the seven figures first $\pi_1$ and then $\pi_2$. We thus obtain, corresponding to the group $G_{168}$, a permutation group of degree 168, which we denote by $P_{168}$, and these two groups are isomorphic.

Since we have recognised $\tau$ and $\omega$ as generating elements of the group $G_{168}$ (\S.\,72, I.), it suffices to determine the permutations corresponding to these two substitutions in order to derive from them the whole group $P_{168}$.

Now from the substitution $\tau$ there arises the cyclic permutation of the seven figures

\[
(\tau) \qquad (0, 1, 2, 3, 4, 5, 6)
\]

and the influence of $\omega$ follows from the remark that $z_0$ remains unaltered by the substitutions of the octahedral group $\chi^\lambda \omega^u \Theta^v$. Hence, to obtain the alteration of $z_r$ by $\omega$, we have only to ascertain the influence of $\omega \tau^r$ on $z_0$. But this influence follows immediately from the formulae \S.\,72, (15), according to which, e.g., $\omega \tau = \tau^2 \chi^2 \omega \Theta^2$, so that $z_1$ passes into $z_2$, and so on. Accordingly, to the substitution $\omega$ corresponds the permutation

\[
(\omega) \qquad \begin{pmatrix} 0, & 1, & 2, & 3, & 4, & 5, & 6 \\ 0, & 2, & 1, & 5, & 4, & 3, & 6 \end{pmatrix} = (1, 2)\,(3, 5).
\]

One can represent the group $P_{168}$ by the congruence group of ter-
%CONT>>


% ===== scan leaf 496, printed 475 =====
% !PRINTED 475
\weberpage{475}
%<<CONT
nary linear substitutions for the modulus 2, which we investigated in \S.\,79\footnote{According to an oral communication, this is the way in which \gesperrt{Kronecker} formed them.}.

For this purpose we have to denote the seven quantities by three indices $(x_1, x_2, x_3)$, which are taken with respect to the modulus 2 and where the combination $(0, 0, 0)$ is excluded. One thus obtains the seven quantities

\weq{1}{(1,0,0),\ (0,1,0),\ (0,0,1);\ (0,1,1),\ (1,0,1),\ (1,1,0),\ (1,1,1),}

and if one replaces the $x_1, x_2, x_3$ by the linear combinations

\weq{2}{a x_1 + b x_2 + c x_3,\ a_1 x_1 + b_1 x_2 + c_1 x_3,\ a_2 x_1 + b_2 x_2 + c_2 x_3}

in which the substitution

\weq{3}{\begin{pmatrix} a, & b, & c \\ a_1, & b_1, & c_1 \\ a_2, & b_2, & c_2 \end{pmatrix}}

is to be taken with respect to the modulus 2, one obtains the permutation group $P_{168}$.

If, as in \S.\,79, we choose for $\tau$ the substitution of the congruence group

\weq{4}{\tau = \begin{pmatrix} 1, & 0, & 1 \\ 1, & 0, & 0 \\ 0, & 1, & 0 \end{pmatrix}}

then we can denote the quantities $z$ in such a way that they pass cyclically into one another by the application of $\tau$ and its repetitions, where we may choose an arbitrary one of the quantities (1) for $z_0$, say thus:

\weq{5}{z_0 = (1,0,0),\ z_1 = (1,1,0),\ z_2 = (1,1,1),\ z_3 = (0,1,1),\ z_4 = (1,0,1),\ z_5 = (0,1,0),\ z_6 = (0,0,1).}

Then $\omega$ is to be chosen so that $z_0, z_4, z_6$ remain unchanged by it and $z_1$ is interchanged with $z_2$, and $z_3$ with $z_5$. This gives

\weq{6}{\omega = \begin{pmatrix} 1, & 0, & 0 \\ 0, & 1, & 0 \\ 0, & 1, & 1 \end{pmatrix},}

and from this one obtains, by \S.\,72, (17):

\weq{7}{\chi = \begin{pmatrix} 1, & 0, & 0 \\ 0, & 0, & 1 \\ 0, & 1, & 1 \end{pmatrix}, \quad \Theta = \begin{pmatrix} 1, & 1, & 1 \\ 0, & 1, & 0 \\ 0, & 1, & 1 \end{pmatrix}.}


% ===== scan leaf 497, printed 476 =====
% !PRINTED 476
\weberpage{476}
\pdfbookmark[2]{§ 128. Equations of the seventh degree with a group of the 168th degree}{bm146}
The two generating permutations $\tau, \omega$ belong to the first kind (Bd. I, \S.\,153), and consequently $P_{168}$ is a divisor of the alternating \gesperrt{permutation group} of seven figures.

\begin{center}
\S.\,128.

\gesperrt{Equations of the seventh degree with a group of degree 168.}
\end{center}

We now turn to the general theory of the special kind of equations of degree 7 whose \gesperrt{Galois} group reduces to $P_{168}$.

Let $\lambda_0, \lambda_1, \lambda_2, \lambda_3, \lambda_4, \lambda_5, \lambda_6$ first be arbitrary quantities, and let

\weq{1}{\tau = (\lambda_0, \lambda_1, \lambda_2, \lambda_3, \lambda_4, \lambda_5, \lambda_6)}

be a cyclic permutation of degree 7. If we further introduce the transposition pair

\weq{2}{\omega = (\lambda_1, \lambda_2)\,(\lambda_3, \lambda_5)}

then these two permutations generate, through their compositions and repetitions, the whole group $P_{168}$.

If, following \S.\,72, (17), one composes $\chi$ and $\Theta$ from them, one obtains (since, as already remarked above, the composition of the permutations must take place in the reverse order to that of the substitutions) the following permutations of the seven indices:

\[
\begin{aligned}
\Theta^3 &= \tau^6 \omega \tau \omega = (1, 2, 5, 3)\,(4, 6),
\qquad \Theta = (1, 3, 5, 2)\,(4, 6) \\
\chi &= \omega \tau^3 \Theta \tau^2 = (1, 4, 2)\,(3, 5, 6),
\end{aligned}
\]

expressions which can also be derived very easily from \S.\,127, (7). One sees that $\lambda_0$ remains unchanged by $\omega, \Theta, \chi$, and consequently by the whole octahedral group $P_{24}$ contained in $P_{168}$. It is easy to form a function of the seven quantities $\lambda$ that belongs to the group $P_{168}$.

Take, for example, the product $\lambda_0 \lambda_4 \lambda_6$, which consists of the $\lambda$ left untouched by $\omega$, and form the sum of the products that arise from it by application of the cyclic permutation $\tau$ and its repetitions:

\weq{3}{v = \lambda_0 \lambda_4 \lambda_6 + \lambda_1 \lambda_5 \lambda_0 + \lambda_2 \lambda_6 \lambda_1 + \lambda_3 \lambda_0 \lambda_2 + \lambda_4 \lambda_1 \lambda_3 + \lambda_5 \lambda_2 \lambda_4 + \lambda_6 \lambda_3 \lambda_5.}


% ===== scan leaf 498, printed 477 =====
% !PRINTED 477
\weberpage{477}
If one applies to this the substitutions $\omega$ and $\tau$, one sees that $v$ remains unchanged and hence admits all permutations of the group $P_{168}$.

Now $P_{168}$ is a divisor of the symmetric permutation group of seven elements of index 30, and a divisor of the alternating group of index 15. Consequently $v$ is a root of an equation of the 30th degree, whose coefficients are symmetric functions of the $\lambda$, and a root of an equation of the 15th degree, whose coefficients also contain the difference product of the $\lambda$. If the $\lambda$ are the roots of an equation of the 7th degree without affect, then $v$ is the root of a resolvent of the 30th degree, which by adjunction of the square root of the discriminant splits into two factors of the 15th degree.

If, besides the symmetric functions of the $\lambda$, the quantity $v$ belongs to the domain of rationality --- whether it is rational from the outset, or the domain of rationality is extended by the adjunction of $v$ --- then the $\lambda$ are the roots of a special equation of the 7th degree whose \gesperrt{Galois} group is of the 168th degree. What we now still have to prove is that this special kind of equations of the 7th degree can be reduced to the form problem of the group $G_{168}$.

To achieve this, we must seek to form three rational functions $X_1$, $X_2$, $X_3$ of the roots $\lambda$ which, when the permutations of the group $P_{168}$ are carried out, undergo the corresponding linear substitutions of the group $G_{168}$, i.\,e. which, if $\pi$ is a permutation from $P_{168}$ and $A$ is the corresponding substitution from $G_{168}$, pass into $A(X_1, X_2, X_3)$ upon carrying out the permutation $\pi$.

If we substitute these functions $X_1, X_2, X_3$ for the variables in the invariants of the group $G_{168}$, these invariants pass into functions of the $\lambda_i$ which remain unchanged under the permutations of the group $P_{168}$, and which consequently belong to the domain of rationality. The computation of the values of the functions $X_1, X_2, X_3$ from these values of the invariants is then the form problem for the $G_{168}$. Any function of the $X_1, X_2, X_3$ that is changed by the substitutions of $G_{168}$ is a root of a resolvent of the given equation of the 7th degree, and indeed, since the groups $G_{168}$ and $P_{168}$ are simple, a total resolvent.


% ===== scan leaf 499, printed 478 =====
% !PRINTED 478
\weberpage{478}
\pdfbookmark[2]{§ 129. Contragredient groups}{bm147}
If one has any system of non-vanishing functions $X_1, X_2, X_3$, one can therefore always form such resolvents. By the form problem of the group $G_{168}$, then, the equation of degree 7 with the group $P_{168}$ is solved at the same time.

In order to reduce the solution of these equations of degree 7 to the special form problem as it occurs in the elliptic functions (with $f = 0$), an accessory equation of degree 4 must then still be solved, just as, in order to reduce the general equation of degree 5 to the icosahedral equation, we found an accessory square root necessary.

Everything now depends, then, on determining the functions $X_1, X_2, X_3$ of the roots $\lambda$ in accordance with these requirements. To make this possible, we must make use of some simple theorems from the general theory of invariants, which we shall derive in the following section, so far as they come into consideration for our task.

\begin{center}
§. 129.

\gesperrt{Contragredient Groups.}
\end{center}

We have already explained the concept of the contragredient transformation in §.\,37. Namely, if

\weq{1}{A = \begin{pmatrix} a_1, & b_1, & c_1 \\ a_2, & b_2, & c_2 \\ a_3, & b_3, & c_3 \end{pmatrix}, \quad A_1 = \begin{pmatrix} a_1, & a_2, & a_3 \\ b_1, & b_2, & b_3 \\ c_1, & c_2, & c_3 \end{pmatrix}}

are two substitutions transposed to each other, and $x_1, x_2, x_3$ and $\xi_1, \xi_2, \xi_3$ are two sequences of variables which, by the substitutions

\weq{2}{(y) = A(x), \quad \xi = A_1(\eta)}

are transformed into two new sequences of variables $y_1, y_2, y_3$ and $\eta_1, \eta_2, \eta_3$, then we have called these two sequences of variables, and likewise their transformations, \gesperrt{contragredient}.

By the substitution $y = A(x)$ every function $\Phi(x_1, x_2, x_3)$ of the $(x)$ is transformed into a function of the $(y)$, and the formation of the derivatives yields:

\weq{3}{\frac{\partial \Phi}{\partial x_1} = a_1 \frac{\partial \Phi}{\partial y_1} + a_2 \frac{\partial \Phi}{\partial y_2} + a_3 \frac{\partial \Phi}{\partial y_3}, \cdots}
%CONT>>


% ===== scan leaf 500, printed 479 =====
% !PRINTED 479
\weberpage{479}
%<<CONT
or, in our abbreviated notation:

\[
\left(\frac{\partial \Phi}{\partial x_1}, \frac{\partial \Phi}{\partial x_2}, \frac{\partial \Phi}{\partial x_3}\right) = A_1 \left(\frac{\partial \Phi}{\partial y_1}, \frac{\partial \Phi}{\partial y_2}, \frac{\partial \Phi}{\partial y_3}\right).
\]

This proves the theorem:

1. \gesperrt{The sequences of variables}

\[
(x_1, x_2, x_3) \quad \text{and} \quad \left(\frac{\partial \Phi}{\partial x_1}, \frac{\partial \Phi}{\partial x_2}, \frac{\partial \Phi}{\partial x_3}\right)
\]

\gesperrt{are contragredient.}

By repeated application of this theorem, the higher differential quotients with respect to the $x$ can also be formed by means of those with respect to the $y$, for which one obtains the following rule:

2. \gesperrt{To express the $m$th derivatives}

\[
\frac{\partial^m \Phi}{\partial x_1^\alpha \partial x_2^\beta \partial x_3^\gamma}, \quad \alpha + \beta + \gamma = m
\]

\gesperrt{by means of the derivatives with respect to $y$, replace, in the developed expression}

\weq{4}{\xi_1^\alpha \xi_2^\beta \xi_3^\gamma = \sum^{\varkappa, \lambda, \mu} C^{\alpha, \beta, \gamma}_{\varkappa, \lambda, \mu} \eta_1^\varkappa \eta_2^\lambda \eta_3^\mu, \quad \varkappa + \lambda + \mu = m}

\gesperrt{the products}

\[
\begin{aligned}
&\xi_1^\alpha \xi_2^\beta \xi_3^\gamma \quad \text{by} \quad \frac{\partial^m \Phi}{\partial x_1^\alpha \partial x_2^\beta \partial x_3^\gamma} \\
&\eta_1^\varkappa \eta_2^\lambda \eta_3^\mu \quad \text{by} \quad \frac{\partial^m \Phi}{\partial y_1^\varkappa \partial y_2^\lambda \partial y_3^\mu},
\end{aligned}
\]

\gesperrt{hence}

\weq{5}{\frac{\partial^m \Phi}{\partial x_1^\alpha \partial x_2^\beta \partial x_3^\gamma} = \sum C^{\alpha, \beta, \gamma}_{\varkappa, \lambda, \mu} \frac{\partial^m \Phi}{\partial y_1^\varkappa \partial y_2^\lambda \partial y_3^\mu},}

\gesperrt{where under the summation sign $\varkappa, \lambda, \mu$ run through all non-negative values satisfying the condition $\varkappa + \lambda + \mu = m$.}

The coefficients $C^{\alpha, \beta, \gamma}_{\varkappa, \lambda, \mu}$ are integral rational functions of the substitution coefficients $a_1, a_2, \ldots$

To prove this rule in general, one need only regard formula (5) as proved for $m - 1$ in place of $m$, and, with application of formula (3), form the derivative with respect to one of the variables $x$, and in doing so take the relation following from the definition (4)

\[
C^{\alpha, \beta, \gamma}_{\varkappa, \lambda, \mu} = C^{\alpha - 1, \beta, \gamma}_{\varkappa - 1, \lambda, \mu}\, a_1 + C^{\alpha - 1, \beta, \gamma}_{\varkappa, \lambda - 1, \mu}\, a_2 + C^{\alpha - 1, \beta, \gamma}_{\varkappa, \lambda, \mu - 1}\, a_3
\]

into account.


% ===== scan leaf 501, printed 480 =====
% !PRINTED 480
\weberpage{480}
We can now also generalise this theorem in the following way:

3. \gesperrt{If by the substitution $y = A(x)$ any form $\varphi(x)$ goes over into $\Phi(y)$, and if any second form $\psi(\xi)$ goes over by the transposed substitution $\xi = A_1(\eta)$ into $\Psi(\eta)$, then one obtains a new transformation by $A$ if one makes in $\psi(\xi)$ and $\Psi(\eta)$ the interchanges}

\[
\begin{aligned}
&\xi_1^\alpha \xi_2^\beta \xi_3^\gamma \quad \text{with} \quad \frac{\partial^{\alpha + \beta + \gamma} \varphi}{\partial x_1^\alpha \partial x_2^\beta \partial x_3^\gamma} \\
&\eta_1^\alpha \eta_2^\beta \eta_3^\gamma \quad \text{with} \quad \frac{\partial^{\alpha + \beta + \gamma} \Phi}{\partial y_1^\alpha \partial y_2^\beta \partial y_3^\gamma},
\end{aligned}
\]

\gesperrt{that is, if, as one may also express it, one puts in $\psi$ and $\Psi$ for the powers and products of the variables $\xi, \eta$ the corresponding derivatives of $\varphi, \Phi$ with respect to the variables $x, y$.}

From the rule of composition of linear substitutions there now follows immediately the following theorem:

4. \gesperrt{If $A$ runs through a group $G$, then the transposed substitution $A_1$ runs through a group $G_1$. If $A, B$ are two elements of $G$ and $A_1, B_1$ the corresponding elements of $G_1$, then $AB$ and $B_1 A_1$ are corresponding elements. The groups $G$ and $G_1$ are said to be contragredient to each other.}

The two groups $G$ and $G_1$ are, however, related to each other isomorphically only if one lets not the element $A$ but the element $A^{-1}$ correspond to $A_1$; for then $A_1 B_1$ corresponds to the element $A^{-1} B^{-1} = (BA)^{-1}$.

The invariants of the group $G_1$ are called contravariants of the group $G$. Accordingly the invariants of $G$ are also the contravariants of $G_1$.

From (3) there then follows the following theorem:

5. \gesperrt{If in a contravariant of $G$ one replaces the powers and products of the variables by the corresponding derivatives of an invariant, one again obtains an invariant of $G$}


% ===== scan leaf 502, printed 481 =====
% !PRINTED 481
\weberpage{481}
\pdfbookmark[2]{§ 130. Solution of the equation of the seventh degree with the group P168 by means of the form problem of the group G168}{bm148}

And likewise:

6. \gesperrt{If in an invariant of $G$ one replaces the powers and products of the variables by the corresponding derivatives of a contravariant, the result is again a contravariant.}

\begin{center}
\S. 130.

\gesperrt{Solution of the equation of the seventh degree with the group $P_{168}$ by the form problem of the group $G_{168}$.}
\end{center}

The linear substitution group $G_{168}$ has the remarkable property of being contragredient to itself. For the generating substitutions $\tau$, $\omega$ of $G_{168}$ (\S.\,115) remain unchanged under transposition, and hence if one transposes any substitution of the group, one obtains a substitution which likewise occurs in the group.

The contravariants of $G_{168}$ are therefore (apart from the designation of the variables) identical with its invariants.

But the octahedral group $G_{24}$ contained in the group $G_{168}$, which consists of the substitutions $\chi^\lambda\omega^\mu\Theta^\nu$, is distinct from its contragredient group; for by \S.\,72, (17) we have, e.\,g.,
\[
\Theta^3 = \omega\tau\omega\tau^6,
\]
and the substitution transposed to it,
\[
\Theta_1^3 = \tau^6\omega\tau\omega
\]
can, by \S.\,72, (15), be brought into the form $\tau\chi^2\Theta^2$ and is therefore not contained in $G_{24}$.

Now let $\eta_0$ be any function of the quantities $\lambda_0, \lambda_1, \ldots, \lambda_6$ belonging to the group $P_{24}$ (\S.\,128), e.\,g. the root $\lambda_0$ itself. Under the cyclic permutations $\tau^\nu$ let $\eta_0$ pass into $\eta_0, \eta_1, \ldots, \eta_6$. In the resolvents
\weq{1}{\tilde\omega_r = \sum_{0,6}^{\nu} \varepsilon^{\nu r}\eta_\nu \qquad r = 1, 2, \ldots, 6}
we then have a system of functions which under the permutations of $P_{168}$ substitute linearly, though not ternarily; since one can indeed express the $\eta_\nu$ themselves linearly in terms of the $\tilde\omega_r$.


% ===== scan leaf 503, printed 482 =====
% !PRINTED 482
\weberpage{482}

A system of three functions which undergo ternary substitution can be formed in the following way\footnote{F.\ \gesperrt{Klein}, \glqq Ueber die Auflösung gewisser Gleichungen vom 7\textsuperscript{ten} und 8\textsuperscript{ten} Grade\grqq. Mathem. Annalen, Bd. XV (1879).}.

We first introduce a system of auxiliary variables $x_1, x_2, x_3$, which we subject to the substitutions of the group $G_{168}$, and from these we form the function $z_0$ belonging to the group $G_{24}$, together with its conjugates $z_r$ [\S.\,125, (7)]:
\weq{2}{z_r = \varepsilon^{2r}x_1^2 + \varepsilon^{4r}x_2^2 + \varepsilon^r x_3^2 - \frac{1-\sqrt{-7}}{2}(\varepsilon^{-r}x_2x_3 + \varepsilon^{-2r}x_3x_1 + \varepsilon^{-4r}x_1x_2).}

To this we now take a function $\eta_0$ of the roots $\lambda_r$ of our equation of the 7th degree which belongs to the group $P_{24}$, e.\,g.\ a rational function of $\lambda_0$, and the conjugate values $\eta_0, \eta_1, \ldots, \eta_6$, and form the sum
\weq{3}{\psi = \eta_0 z_0 + \eta_1 z_1 + \eta_2 z_2 + \eta_3 z_3 + \eta_4 z_4 + \eta_5 z_5 + \eta_6 z_6,}
which is a quadratic function of the $x$, whose coefficients depend on the roots $\lambda_r$.

\gesperrt{This function $\psi$ does not change when the variables $(x)$ are subjected to a substitution of the group $G_{168}$ and the roots $\lambda_r$ at the same time to the corresponding permutation from $P_{168}$.} For by this simultaneous operation only the summands in the sum (2) are interchanged among one another, and hence the sum itself is not changed.

We may therefore designate $\psi$ as a \gesperrt{simultaneous invariant} of the groups $G_{168}$ and $P_{168}$.

If we arrange the function $\psi$ according to the variables $x_1, x_2, x_3$, we obtain
\weq{4}{\psi = p_1x_1^2 + p_2x_2^2 + p_3x_3^2 + 2q_1x_2x_3 + 2q_2x_3x_1 + 2q_3x_1x_2,}
where for abbreviation
\weq{5}{\begin{aligned}
p_1 &= \sum^r \varepsilon^{2r}\eta_r, & q_1 &= -\frac{1-\sqrt{-7}}{4}\sum^r \varepsilon^{-r}\eta_r\\
p_2 &= \sum^r \varepsilon^{4r}\eta_r, & q_2 &= -\frac{1-\sqrt{-7}}{4}\sum^r \varepsilon^{-2r}\eta_r\\
p_3 &= \sum^r \varepsilon^{r}\eta_r, & q_3 &= -\frac{1-\sqrt{-7}}{4}\sum^r \varepsilon^{-4r}\eta_r
\end{aligned}}


% ===== scan leaf 504, printed 483 =====
% !PRINTED 483
\weberpage{483}

is set, so that the $p, q$ are thus functions of the roots $\lambda_r$. Now we choose three different functions $\eta_0$ which satisfy the conditions stated so far, and denote them by $\eta_0, \eta_0', \eta_0''$.

Let the forms $\psi$ derived from these three functions be $\psi, \psi', \psi''$, and their coefficients (5) $p_i, q_i$; $p_i', q_i'$; $p_i'', q_i''$. Then not only is each of the three functions $\psi, \psi', \psi''$ a simultaneous invariant of the groups $P_{168}, G_{168}$, but so also is their functional determinant (Bd.\,I, \S.\,59):
\weq{6}{\Psi = \frac{1}{S}\begin{vmatrix} \dfrac{\partial\psi}{\partial x_1}, & \dfrac{\partial\psi}{\partial x_2}, & \dfrac{\partial\psi}{\partial x_3}\\[2mm] \dfrac{\partial\psi'}{\partial x_1}, & \dfrac{\partial\psi'}{\partial x_2}, & \dfrac{\partial\psi'}{\partial x_3}\\[2mm] \dfrac{\partial\psi''}{\partial x_1}, & \dfrac{\partial\psi''}{\partial x_2}, & \dfrac{\partial\psi''}{\partial x_3} \end{vmatrix}.}

The determinant $\Psi$ is a form of degree 3 in the variables $x$, which, according to the notation of Bd.\,I, \S.\,15, (4), may have the expression:
\weq{7}{\Psi = \sum A_{h,i,k}\, x_h x_i x_k.}

The coefficients $A_{h,i,k}$ are then functions of the roots $\lambda_r$ which depend linearly and homogeneously on each of the three systems $\eta_r, \eta_r', \eta_r''$. Such functions are called \gesperrt{trilinear}.

Now let $\xi_1, \xi_2, \xi_3$ denote a system of variables contragredient to $(x)$, so that the biquadratic form
\weq{8}{f(\xi_1, \xi_2, \xi_3) = \xi_1^3\xi_3 + \xi_2^3\xi_1 + \xi_3^3\xi_2}
is a contravariant of $G_{168}$. If we now in (7) replace
\[
x_h x_i x_k \quad \text{by} \quad \frac{1}{6}\frac{\partial^3 f}{\partial\xi_h\partial\xi_i\partial\xi_k}
\]
we obtain a linear form
\weq{9}{L = \frac{1}{6}\sum A_{h,i,k}\,\frac{\partial^3 f}{\partial\xi_h\partial\xi_i\partial\xi_k} = X_1\xi_1 + X_2\xi_2 + X_3\xi_3,}
in which the $X_1, X_2, X_3$ are functions of $\lambda_r$, and by the theorem of \S.\,129, 3.\ $L$ remains unchanged when the roots $\lambda_r$ are transformed by any permutation $\pi$ of the group $P_{168}$ and at the same time the variables $(\xi)$ together with the variables $(x)$ are transformed contragrediently by the corresponding substitution. % [sic] "irgend einer Permutation" as printed (dative for expected accusative)


% ===== scan leaf 505, printed 484 =====
% !PRINTED 484
\weberpage{484}
\pdfbookmark[2]{§ 131. Possibility of determining the functions X1, X2, X3}{bm149}

For if under $\pi$ the coefficient $A_{h,i,k}$ passes into $A_{h,i,k}'$, then by means of the transformation $(y) = A(x)$:
\[
\sum A_{h,i,k}' \, y_h y_i y_k = \sum A_{h,i,k} \, x_h x_i x_k,
\]
and by means of the substitution $\xi = A_1(\eta)$ the identity $f(\xi_1, \xi_2, \xi_3) = f(\eta_1, \eta_2, \eta_3)$ holds. Consequently, by the theorem of \S.\,129, 3., there results:
\[
\sum A_{h,i,k}' \, \frac{\partial^3 f}{\partial\eta_h\partial\eta_i\partial\eta_k} = \sum A_{h,i,k} \, \frac{\partial^3 f}{\partial\xi_h\partial\xi_i\partial\xi_k}.
\]

If, then, $\pi$ denotes any permutation of the group $P_{168}$ by which $X_1, X_2, X_3$ pass into $Y_1, Y_2, Y_3$, and
\[
A = \begin{pmatrix} a_1, & b_1, & c_1\\ a_2, & b_2, & c_2\\ a_3, & b_3, & c_3 \end{pmatrix}
\]
is the corresponding substitution from the group $G_{168}$, then we have to set:
\[
\begin{aligned}
\xi_1 &= a_1\eta_1 + a_2\eta_2 + a_3\eta_3\\
\xi_2 &= b_1\eta_1 + b_2\eta_2 + b_3\eta_3\\
\xi_3 &= c_1\eta_1 + c_2\eta_2 + c_3\eta_3,
\end{aligned}
\]
and the invariant property of $L$ gives the relation
\weq{10}{Y_1\eta_1 + Y_2\eta_2 + Y_3\eta_3 = X_1\xi_1 + X_2\xi_2 + X_3\xi_3,}
or, developed,
\weq{11}{\begin{aligned} Y_1 &= a_1X_1 + b_1X_2 + c_1X_3\\ Y_2 &= a_2X_1 + b_2X_2 + c_2X_3\\ Y_3 &= a_3X_1 + b_3X_2 + c_3X_3, \end{aligned}}
i.\,e. the permutation $\pi$, applied to the functions $X_1, X_2, X_3$, has the same effect as the linear substitution $A(X_1, X_2, X_3)$, and accordingly the $X_1, X_2, X_3$ are functions of the kind our problem demands (end of \S.\,128).

\begin{center}
\S. 131.

\gesperrt{Possibility of the determination of the functions $X_1, X_2, X_3$.}
\end{center}

In order to secure completely the reduction of the equation of degree 7 with the group $P_{168}$ to the form problem of the group $G_{168}$, one thing still remains:
%CONT>>


% ===== scan leaf 506, printed 485 =====
% !PRINTED 485
\weberpage{485}
%<<CONT
For what still remains is the proof that one can dispose of $\eta_0, \eta_0', \eta_0''$ (§.\,130) in such a way that the functions $X_1, X_2, X_3$ do not vanish identically. For this we must consider the mode of formation of the quantities $X$ somewhat more closely.

If, in accordance with (8) (§.\,130), we set in (9):
\[
\frac{\partial^3 f}{\partial \xi_1^3} = 6\,\xi_3, \quad \frac{\partial^3 f}{\partial \xi_1^2 \xi_2} = 0, \quad \frac{\partial^3 f}{\partial \xi_1^2 \partial \xi_3} = 6\,\xi_1, \quad \frac{\partial^3 f}{\partial \xi_1 \partial \xi_2 \partial \xi_3} = 0, \dots, % [sic] source omits the second \partial before \xi_2
\]
then it results that

\weq{1}{\begin{aligned}
X_1 &= A_{2,2,2} + 3\,A_{1,1,3}, \quad X_2 = A_{3,3,3} + 3\,A_{2,2,1}, \\
X_3 &= A_{1,1,1} + 3\,A_{3,3,2}.
\end{aligned}}

Now, however, by (4) and (6) (§.\,130) we further have:

\weq{2}{\Psi =
\begin{vmatrix}
p_1 x_1 + q_3 x_2 + q_2 x_3, & q_3 x_1 + p_2 x_2 + q_1 x_3, & q_2 x_1 + q_1 x_2 + p_3 x_3 \\
p_1' x_1 + q_3' x_2 + q_2' x_3, & q_3' x_1 + p_2' x_2 + q_1' x_3, & q_2' x_1 + q_1' x_2 + p_3' x_3 \\
p_1'' x_1 + q_3'' x_2 + q_2'' x_3, & q_3'' x_1 + p_2'' x_2 + q_1'' x_3, & q_2'' x_1 + q_1'' x_2 + p_3'' x_3
\end{vmatrix}.}

This can easily be ordered according to powers and products of the $x$, and if we accordingly use the notation
\[
(p_1, q_3, q_2) =
\begin{vmatrix}
p_1, & q_3, & q_2 \\
p_1', & q_3', & q_2' \\
p_1'', & q_3'', & q_2''
\end{vmatrix} \cdots,
\]
one obtains
\[
\begin{aligned}
A_{1,1,1} &= (p_1, q_3, q_2), & 3\,A_{3,3,2} &= (q_3, q_1, p_3) + (q_2, p_2, p_3) \\
A_{2,2,2} &= (q_3, p_2, q_1), & 3\,A_{1,1,3} &= (p_1, q_1, q_2) + (p_1, q_3, p_3) \\
A_{3,3,3} &= (q_2, q_1, p_3), & 3\,A_{2,2,1} &= (q_3, p_2, q_2) + (p_1, p_2, q_1),
\end{aligned}
\]
and from this, by (1),

\weq{3}{\begin{aligned}
X_1 &= (q_3, p_2, q_1) + (p_1, q_1, q_2) + (p_1, q_3, p_3) \\
X_2 &= (q_2, q_1, p_3) + (q_3, p_2, q_2) + (p_1, p_2, q_1) \\
X_3 &= (p_1, q_3, q_2) + (q_3, q_1, p_3) + (q_2, p_2, p_3).
\end{aligned}}

These functions $X_1, X_2, X_3$ are, as follows from what was remarked above [§.\,130, (5)], trilinear forms of the three sequences of variables $\eta_r, \eta_r', \eta_r''$, whose coefficients can be expressed rationally in terms of the seventh root of unity $\varepsilon$. But the coefficients of these forms are certainly not all equal to zero. For by (3) the quantities $p, q$ can be so chosen that $X_1, X_2, X_3$ do not become equal to zero, and the 21 quantities $\eta_r, \eta_r', \eta_r''$ can be determined from §.\,130, (5) in such a way that the 18 quantities $p, q$ (and moreover the three sums $\Sigma \eta_r$) obtain arbitrarily prescribed values
%CONT>>


% ===== scan leaf 507, printed 486 =====
% !PRINTED 486
\weberpage{486}
%<<CONT
If we then make, for the seven variables $\eta_r$, the substitution

\weq{4}{\eta_r = a_0 + a_1 \lambda_r + a_2 \lambda_r^2 + a_3 \lambda_r^3 + a_4 \lambda_r^4 + a_5 \lambda_r^5 + a_6 \lambda_r^6,}

whose determinant, being the product of all the differences $\lambda_i - \lambda_k$, is different from zero, and substitute correspondingly

\weq{5}{\eta_r = \sum_{0,\,6}^{s} a_s \lambda_r^s, \quad \eta_r' = \sum_{0,\,6}^{s} a_s' \lambda_r^s, \quad \eta_r'' = \sum_{0,\,6}^{s} a_s'' \lambda_r^s,}

then $X_1$ thereby goes over into a trilinear form of the three sequences of variables $a_s, a_s', a_s''$, which cannot vanish identically, since conversely the $a_s, a_s', a_s''$ are in turn linearly expressible in terms of the $\eta_r, \eta_r', \eta_r''$. Now one can assume for the variables $a_s, a_s', a_s''$ such rational numerical values (Bd. I, §.\,143) that the functions $X_1, X_2, X_3$ take values different from zero, and then (5) represents a suitable assumption for the functions $\eta_r, \eta_r', \eta_r''$\footnote{Cf., on the question hereby settled: \gesperrt{Burckhardt}, ``Ueber einen fundamentalen Satz der Lehre von den endlichen Gruppen linearer Substitutionen''. Mathem. Annalen, Bd.\ 42 (1892).}.


% ===== scan leaf 508, printed 487 =====
\weberpage{487}% page number not printed on this page
% !PRINTED none
% Part-title page (no running head, no page number).
\begin{center}
\gesperrt{BOOK FOUR.}

\rule{0.3\textwidth}{0.4pt}

\gesperrt{ALGEBRAIC NUMBERS.}

\rule{0.3\textwidth}{0.4pt}
\end{center}


% ===== scan leaf 509, printed 488 =====
\weberpage{488}% page number not printed on this page
% !PRINTED none
% Blank verso (no running head, page number, or body text).


% ===== scan leaf 510, printed 489 =====
\weberpage{489}% page number not printed on this page
\pdfbookmark[0]{Fourth Book. Algebraic Numbers.}{bm150}
\pdfbookmark[1]{Sixteenth Section. Numbers and Functionals of an Algebraic Field.}{bm151}
\pdfbookmark[2]{§ 132. Definition of algebraic numbers}{bm152}
% !PRINTED none
\begin{center}
\gesperrt{Sixteenth Chapter.}

\textbf{Numbers and Functionals of an Algebraic Field.}

\rule{4cm}{0.4pt}
\end{center}

\begin{center}
\S.\,132.

\gesperrt{Definition of the algebraic numbers.}
\end{center}

An algebraic equation
\[
F(x) = x^n + a_1 x^{n-1} + \cdots + a_{n-1} x + a_n = 0,
\]
whose coefficients $a_1, a_2, \ldots, a_n$ are rational numbers, we call for brevity a rational equation. It has, as we have proved in earlier chapters, always $n$ or fewer roots, but always at least one. How one can come as near as one pleases to each of these roots by rational numbers, say by decimal fractions or by continued fractions, if need be with the aid of the imaginary unit $i = \sqrt{-1}$, that is, how one can compute the values of the roots approximately, has been shown in the second Book of the first volume.

In the following considerations we shall not be concerned with these numerical values, but with the arithmetical laws to which these numbers are subject, and which can be derived from the definition itself and not from the numerical values. We therefore now place the following definition at the head:

\gesperrt{A number $\Theta$ which satisfies a rational equation}
\weq{1}{F(\Theta) = 0}
\gesperrt{is called an algebraic number.}

Every algebraic equation with rational coefficients furnishes us such algebraic numbers, which can therefore be given in arbitrary quantity. The question whether there are also numbers that are not algebraic will occupy us later.


% ===== scan leaf 511, printed 490 =====
% !PRINTED 490
\weberpage{490}
\pdfbookmark[2]{§ 133. Algebraic integers}{bm153}

An algebraic number satisfies not merely one, but infinitely many rational equations; for if we multiply two arbitrary functions of the form $F(x)$ with one another, we obtain a function of the same form, which vanishes for $x = \Theta$ if one of the factors has this property.

Among all the rational equations that an algebraic number satisfies there is one of lowest possible degree, $f(\Theta) = 0$, where $f(x)$ has the form:
\weq{2}{f(x) = x^n + a_1 x^{n-1} + a_2 x^{n-2} + \cdots + a_n,}
and there can also be only one such equation, if, as always hitherto, we take the coefficient of the highest power of $x$ equal to 1.

For if $f(x), f_1(x)$ are two functions of the form (2) of the same degree $n$, then $f(x) - f_1(x)$ is of degree lower than $n$, and if both $f(\Theta)$ and $f_1(\Theta)$ vanish, then $f(\Theta) - f_1(\Theta)$ vanishes also; if, then, this difference does not vanish identically, $\Theta$ satisfies an equation of degree lower than $n$, which is contrary to the assumption.

\gesperrt{The function $f(x)$ is irreducible in the field of rational numbers.}

For if $f(x)$ splits into two rational factors $f_1(x)$ and $f_2(x)$, each of which is of lower degree than $f(x)$, then $\Theta$ satisfies one of the two equations $f_1(\Theta) = 0$, $f_2(\Theta) = 0$, which contradicts our assumption.

\gesperrt{If $n$ is the degree of the rational equation of lowest degree that the number $\Theta$ satisfies, we call $\Theta$ an algebraic number of degree $n$.}

\begin{center}
\S.\,133.

\gesperrt{Integral algebraic numbers.}
\end{center}

\gesperrt{An algebraic number $\Theta$ is called an integral algebraic number if it satisfies a rational equation}
\weq{1}{\Theta^m + A_1 \Theta^{m-1} + \cdots + A_{m-1} \Theta + A_m = 0}
%CONT>>


% ===== scan leaf 512, printed 491 =====
% !PRINTED 491
\weberpage{491}
%<<CONT
\gesperrt{; its coefficients $A_1, A_2, \ldots, A_m$ are integers.}

We remark that, according to this definition, in order to characterise an algebraic number $\Theta$ as integral, it suffices if among the infinitely many equations of the form (1) which $\Theta$ satisfies there is \gesperrt{one} whose coefficients are integers.

The integral algebraic numbers include as a special case the ordinary integers, which for distinction we call \gesperrt{integral rational numbers}. The positive integral rational numbers we also call, following a widespread usage, \gesperrt{natural numbers}.

By integers without qualification we then understand integral algebraic numbers, rational as well as irrational.

1. \gesperrt{An integral algebraic number which is at the same time rational is necessarily an integral rational number.}

For suppose that $\Theta = P:Q$ is a rational fraction, and that $P, Q$ are integral rational numbers without common divisor, say $Q$ positive; then from (1) there results
\[
P^m + A_1 P^{m-1} Q + A_2 P^{m-2} Q^2 + \cdots + A_m Q^m = 0,
\]
and from this it is to be seen that every prime divisor of $Q$ would have to be contained in $P$. Hence we must have $Q = 1$, and $\Theta = P$ is an integral rational number.

2. \gesperrt{The sum, difference and product of two integers are again integers.}

To prove this principal theorem, we assume that $\alpha, \beta$ are two integers which satisfy the equations
\weq{2}{\begin{array}{l}
\alpha^\mu + a_1 \alpha^{\mu-1} + \cdots + a_{\mu-1} \alpha + a_\mu = 0 \\
\beta^\nu + b_1 \beta^{\nu-1} + \cdots + b_{\nu-1} \beta + b_\nu = 0
\end{array}}
and make one of the three assumptions
\[
\omega = \alpha + \beta, \ \alpha - \beta, \ \alpha \beta.
\]
Then we set $\mu \nu = m$ and denote the $m$ quantities
\[
\alpha^r \beta^s \quad \begin{array}{l} r = 0, 1, \ldots, \mu - 1 \\ s = 0, 1, \ldots, \nu - 1 \end{array}
\]
in any order by $\omega_1, \omega_2, \ldots, \omega_m$.


% ===== scan leaf 513, printed 492 =====
% !PRINTED 492
\weberpage{492}

Then the products $\omega \omega_1, \omega \omega_2, \ldots, \omega \omega_m$ can, with the help of the equations (2), be put into the form
\[
\omega \omega_r = c_{r,1} \omega_1 + c_{r,2} \omega_2 + \cdots + c_{r,m} \omega_m \quad r = 1, 2, \ldots, m,
\]
in which the coefficients $c_{s,r}$ are integral rational numbers. But if one eliminates the $\omega_1, \omega_2, \ldots, \omega_m$ from these equations, it follows that
\[
\begin{vmatrix}
c_{1,1} - \omega, & c_{1,2}, & \ldots, & c_{1,m} \\
c_{2,1}, & c_{2,2} - \omega, & \ldots, & c_{2,m} \\
\cdots & \cdots & \cdots & \cdots \\
c_{m,1}, & c_{m,2}, & \ldots, & c_{m,m} - \omega
\end{vmatrix}
= 0,
\]
which, developed, contains the form
\[
\omega^m + C_1 \omega^{m-1} + \cdots + C_m = 0,
\]
in which the $C_1, C_2, \ldots, C_m$ are likewise integral rational numbers. But this shows that $\omega$ is an integer, as was to be proved.

3. \gesperrt{If $f(x)$ is a function irreducible in the field of rational numbers, and if a root $\Theta$ of $f(x) = 0$ is an integer, then all roots of $f(x)$ are integers.}

For if a rational function $F(x)$ vanishes for $x = \alpha$, then $F(x)$ is divisible by $f(x)$, and all roots of $f(x)$ are at the same time roots of $F(x)$ (Bd.\,I, \S.\,141). Now if $\alpha$ is an integer, there is a function
\[
F(x) = x^m + A_1 x^{m-1} + \cdots + A_m
\]
with integral coefficients $A_1, \ldots, A_m$ which vanishes for $x = \alpha$, and $F(x)$ therefore also vanishes for all the other roots of $f(x)$, which accordingly are all integers.

If
\[
f(x) = x^n + a_1 x^{n-1} + \cdots + a_n,
\]
then the $a_1, a_2, \ldots, a_n$ are composed by multiplication and addition from the roots of $f(x)$, and are therefore, by 2., integral rational numbers. From this it follows:

4. \gesperrt{If $\Theta$ is an integral algebraic number, then the equation of lowest degree $f(\Theta) = 0$ [\S.\,132, (2)] has integral coefficients.}


% ===== scan leaf 514, printed 493 =====
% !PRINTED 493
\weberpage{493}
\pdfbookmark[2]{§ 134. Algebraic fields}{bm154}

We prove further the theorem:

5. \gesperrt{Every algebraic number $\Theta$ can be transformed by multiplication
with a natural number into an integral algebraic number.}

For if
\[
\Theta^m + A_1 \Theta^{m-1} + \cdots + A_m = 0,
\]
and $A_1, \ldots, A_m$ are rational numbers with the common denominator $a$, then
by multiplication with $a^m$ one obtains:
\[
(a\Theta)^m + A_1 a (a\Theta)^{m-1} + A_2 a^2 (a\Theta)^{m-2} + \cdots + A_m a^m = 0,
\]
from which it follows that $a\Theta$ is an \gesperrt{integer}.

\begin{center}
\S.\,134.

\gesperrt{Algebraic Fields.}
\end{center}

In the thirteenth chapter of the first volume we have seen how, from every
root $\Theta$ of an equation $f(\Theta) = 0$ of degree $n$ irreducible in any
field $\Omega$, one derives an algebraic field $\Omega(\Theta)$ over $\Omega$.
The $n$ fields derived from the $n$ roots of this equation, which may also in
part or all be identical, we have called \gesperrt{conjugate fields}.

If we denote by $R$ the field of rational numbers, then by our definition every
algebraic number of degree $n$, $\Theta$, gives rise to an algebraic field
$R(\Theta)$ over $R$, which from now on we shall call briefly an
algebraic number field of degree $n$. We have also already proved
earlier (Bd.\,I, \S.\,143) that one can always determine an \gesperrt{algebraic number
field which contains a finite number of arbitrarily given algebraic numbers.}

This theorem is emphasised at this point in order to point out that the
generality of a consideration concerning any finite number of algebraic numbers
is not impaired by supposing these numbers all to lie in one algebraic number
field.

Every number $\omega$ of such a field can be represented as an integral
function $\varphi(\Theta)$ of $\Theta$ with rational coefficients, and to every
number $\omega$ there corresponds in each of the $n$ conjugate fields a
determinate number. These \gesperrt{conjugate numbers} can in part be equal to one
another, and we have accordingly primitive and
%CONT>>


% ===== scan leaf 515, printed 494 =====
% !PRINTED 494
\weberpage{494}

%<<CONT
imprimitive numbers of the field. A symmetric function of the conjugate
numbers is a rational number. Among these symmetric functions two are of
especial importance, the sum and the product, of which the first is called
the \gesperrt{trace}, the second the \gesperrt{norm} of $\omega$. These two
numbers are denoted by $S(\omega)$ and $N(\omega)$.

Here, if among the conjugate numbers the same numbers occur several times,
these equal numbers are taken into the sum or the product as many times as
the degree of their frequency indicates.

Since in every such number field the four fundamental operations of
arithmetic can be carried out just as in the field of rational numbers, one
may also raise the question how far the fundamental arithmetical laws known
from the theory of rational numbers hold good in an arbitrary algebraic
number field. The first concern here is the decomposition of the integers
into their \gesperrt{prime factors}.

Since this decomposition does in general \gesperrt{not} succeed with the
numbers of the algebraic field themselves, an extension of the material of
computation is necessary in order to restore the simple laws, and such an
extension is possible in different senses. But these different extensions
must admit of being reduced to one another, or, more precisely stated, of
being placed in a unique relation to one another.

\gesperrt{Kummer} first solved the great problem, for the algebraic numbers
formed from roots of unity (the cyclotomic numbers), by the creation of the
\gesperrt{ideal numbers}\textsuperscript{1)}. Another entirely general
solution, subject to no exception, the problem has found through
\gesperrt{Dedekind}, who regarded the \gesperrt{ideals}\textsuperscript{2)} so named by him

\vspace{1em}\hrule\vspace{0.5em}

\noindent\textsuperscript{1)} \gesperrt{Kummer}, Theorie der idealen
Primfactoren der complexen Zahlen etc. \gesperrt{Crelle}'s Journal, Bd.\,35, 1846,
Bd.\,40, 1850. Abhandlungen der Berliner Akademie 1856. Sur la th\'eorie des
nombres complexes compos\'es de racines de l'unit\'e et de nombres entiers.
\gesperrt{Liouville}'s Journal, Bd.\,16, 1851.

\noindent\textsuperscript{2)} \gesperrt{Dedekind}, in the last Supplement of
the 2nd, 3rd and 4th editions of \gesperrt{Dirichlet's} Vorlesungen \"uber
Zahlentheorie (Braunschweig 1871, 1879, 1894). To be compared also: Sur la
th\'eorie des nombres entiers alg\'ebriques in the Bulletin of
\gesperrt{Darboux} and \gesperrt{Ho\"uel} (1\textsuperscript{re} s\'er. XI,
1877). Ueber den Zusammenhang zwischen der Theorie der Ideale und der
h\"oheren Congruenzen (Abhandlungen der Ges. der Wissensch. in G\"ottingen,
Bd.\,23,
%CONT>>


% ===== scan leaf 516, printed 495 =====
% !PRINTED 495
\weberpage{495}
\pdfbookmark[2]{§ 135. Integral functions in an algebraic field}{bm155}

%<<CONT
as the simplest elements of computation. \gesperrt{Kronecker}\textsuperscript{1)} has struck out on a
path different from these.

The theory of \gesperrt{Dedekind} has been presented by its founder
comprehensively, and with ever-increasing simplicity and perfection, in the
last supplement of the three most recent editions of \gesperrt{Dirichlet's}
Vorlesungen \"uber Zahlentheorie.

The theory of \gesperrt{Kronecker} first became known to the wider circle
of mathematicians in the year 1882, through the Festschrift for
\gesperrt{Kummer's} jubilee, and even now it is still not quite easily
accessible\textsuperscript{2)}.

As regards the \gesperrt{Dedekind} theory, then, we may refer the reader
who wishes to penetrate more deeply into this interesting subject to the
works named. Here, on the other hand, the attempt shall be made to
establish the connection between the various theories, and, on the basis
of the \gesperrt{Kronecker} assumptions and means, to develop from them as
much as is necessary for the algebraic applications that follow.

\begin{center}
\S.\,135.

\gesperrt{Integral Functions in an Algebraic Field.}
\end{center}

Already in the first volume we have had occasion several times to
introduce integral functions of an arbitrary number of independent
variables, and have also (in \S.\,141 f.)
%CONT>>

\vspace{1em}\hrule\vspace{0.5em}

\noindent 1878). Ueber die Discriminanten endlicher K\"orper (ebend.,
Bd.\,29, 1882). Ueber die Anzahl der Ideal-Classen in den verschiedenen
Ordnungen eines endlichen K\"orpers (Braunschweig 1877). Festschrift zur
S\"acularfeier des Geburtstages von \gesperrt{Gauss}. ``Zur Theorie der
Ideale'' and ``Ueber die Begr\"undung der Idealtheorie''. Nachrichten d.
Ges. d. Wissensch. in G\"ottingen 1894, 1895. To these belong also the
treatises of \gesperrt{Hilbert}, ``Ueber die Zerlegung der Ideale
etc.'', Mathem. Annalen, Bd.\,44, 1893. ``Grundz\"uge einer Theorie der
\gesperrt{Galois}'schen Zahlk\"orper'', G\"ottinger Nachrichten 1894.
\gesperrt{Hurwitz}, ``Zur Theorie der Ideale''. ``Ueber einen
Fundamentalsatz etc.'', G\"ottinger Nachrichten 1894, 1895.

\noindent\textsuperscript{1)} \gesperrt{Kronecker}, Grundz\"uge einer
arithmetischen Theorie der algebraischen Gr\"ossen. Festschrift zu
\gesperrt{Kummer's} 50j\"ahrigem Doctor-Jubil\"aum. Berlin 1882. (Also in
Bd.\,92 of \gesperrt{Crelle}'s Journal.) To be mentioned here are also the works of
\gesperrt{Hensel} in volumes 101, 103, 105, 111, 113 of \gesperrt{Crelle}'s Journal.

\noindent\textsuperscript{2)} Cf. the preface to the 4th edition of
\gesperrt{Dirichlet-Dedekind's} Vorlesungen \"uber Zahlentheorie.


% ===== scan leaf 517, printed 496 =====
% !PRINTED 496
\weberpage{496}

%<<CONT
the case in which the coefficients belong to a definite field
$\Omega$. We now make the assumption that this field $\Omega$ is an
\gesperrt{algebraic number field}, and accordingly consider expressions
$\varphi(x, y, z, \ldots)$ which are represented as a sum of terms of the
form
\[
\alpha x^r y^s z^t \ldots
\]
in which the exponents $r, s, t, \ldots$ are positive, or at least
non-negative, integers, while the coefficients $\alpha$ denote numbers in
$\Omega$. Such an expression
\weq{1}{\varphi(x, y, z, \ldots) = \sum \alpha x^r y^s z^t \ldots}
we call an \gesperrt{integral function in $\Omega$}. We always take the
expression to be so ordered and collected that the same combination of the
exponents $r, s, t, \ldots$ does not occur twice in it, and we call two such
expressions equal to one another only when they contain the same products
$x^r y^s z^t \ldots$, affected with the same coefficients. \gesperrt{An integral
function is set equal to zero when and only when all its coefficients are
zero.}

Several integral functions always yield integral functions again by
addition, subtraction and multiplication. By Bd.\,I, \S.\,143, I. one can
substitute for the variables $x, y, z, \ldots$ such rational numerical
values that one, or any number, of given integral functions in $\Omega$
different from zero receive non-vanishing numerical values (in $\Omega$).
From this it follows \gesperrt{that a product of several integral functions vanishes
only when one of its factors vanishes.}

The sum $r + s + t + \cdots$ of the exponents in a term of the expression
(1) is called the \gesperrt{degree of this term}, and the greatest value
which the degree of a term with non-vanishing coefficient in $\varphi$
assumes is called the \gesperrt{degree of the function $\varphi$}.

The degree of a product of two or more integral functions is equal to the
sum of the degrees of the individual factors.

For if in each of the factors one collects the sum of the terms of highest
degree into a homogeneous function, one obtains the terms of highest degree
of the product by multiplying all these homogeneous functions together.


% ===== scan leaf 518, printed 497 =====
% !PRINTED 497
\weberpage{497}

The product of these homogeneous functions cannot, by what was proved above, vanish if none of the factors vanishes, and its degree is equal to the sum of the degrees of the individual factors.

The numbers of the field $\Omega$ are included among the functions. One obtains them by letting either the degree or the number of variables sink down to zero.

Concerning the integral functions in $\Omega$ we have derived several important theorems in the first volume, of which we here single out the theorem \S.\,141, IV:

1. \gesperrt{Among the integral functions in $\Omega$ a distinction must be made between reducible and irreducible ones, of which the former can be represented as a product of several integral functions in $\Omega$, the others not. The reducible functions $\varphi$ can be decomposed into a finite number of irreducible factors, which are themselves integral functions in $\Omega$, and the irreducible factors of $\varphi$ are, apart from constant factors, i.\,e. numbers in $\Omega$, uniquely determined by $\varphi$ itself.}

As special cases, among the integral functions in $\Omega$ are also contained the integral functions in the field of the rational numbers $R$:
\weq{2}{\Phi(x,y,z,\ldots) = \sum a\,x^r y^s z^t \ldots,}
in which the coefficients $a$ are rational numbers.

If these coefficients are \gesperrt{integers without a common divisor}, then this function is called an \gesperrt{original or primitive} one, of which we have proved the following theorem in Bd. I, \S.\,2:

2. \gesperrt{The product of two primitive functions is again a primitive function.}

If the coefficients of the function (2), which we shall for the moment denote by
\[ a_0, a_1, a_2, \ldots \]
are integers with the greatest common divisor $m$, then, if $m>1$, $\Phi$ is an imprimitive \gesperrt{integral integer-coefficient function of divisor $m$, % [sic] "ganze ganzzahlige" so printed in the German
and the divisor of a primitive function is $= 1$.}


% ===== scan leaf 519, printed 498 =====
% !PRINTED 498
\weberpage{498}

If we set
\weq{3}{a_0 = m\,e_0,\quad a_1 = m\,e_1,\quad a_2 = m\,e_2 \ldots,}
then the $e_0, e_1, e_2, \ldots$ are integers without a common divisor;
\weq{4}{E(x,y,z,\ldots) = \sum e\,x^r y^s z^t \ldots}
is a primitive function, and we have
\[ \Phi = m\,E. \]

We have considered these functions $\Phi$ in \S.\,2 of the first volume, and have there proved of them the theorem:

3. \gesperrt{The divisor of a product of two or more integral functions $\Phi$ is equal to the product of the divisors of the individual factors.}

The integral functions in $\Omega$ depend, besides on the variables, on an algebraic number $\Theta$, but otherwise contain only rational numerical coefficients. If we denote such a function by $\varphi(\Theta, x, y, z, \ldots)$, we obtain the \gesperrt{conjugate functions} $\varphi, \varphi_1, \varphi_2, \ldots$ or
\weq{5}{\varphi(\Theta,x,y,z,\ldots),\quad \varphi(\Theta_1,x,y,z,\ldots),\quad \varphi(\Theta_2,x,y,z,\ldots),\ \ldots}
when we substitute for $\Theta$ all the roots of the irreducible equation $f(x)=0$ [\S.\,132, (2)]. These conjugate functions may also in part be equal to one another.

They are all equal to one another when $\varphi$ is a function in $R$; and conversely $\varphi$ is a function in $R$ when the conjugate functions are all equal to one another; for in that case, if we understand by $S(\varphi)$ the sum of the conjugate functions (the trace),
\[ n\varphi = S(\varphi), \]
and $S(\varphi)$ is an integral function whose coefficients are symmetric functions of the $n$ roots $\Theta$, that is, rational numbers.

Among the integral functions in $R$ there belongs also the \gesperrt{norm of} $\varphi$, that is, the product
\weq{6}{N(\varphi) = \varphi\,\varphi_1\,\varphi_2 \ldots,}
for all the coefficients of this function are symmetric functions of the $\Theta$. This function is divisible by $\varphi$, and if we set
\[ N(\varphi) = \varphi\,\varphi' \]
then both $\varphi$ and $\varphi'$ are integral functions in $\Omega$. For $\varphi'$, as the product $\varphi_1\varphi_2\ldots$, is an integral function, and the coeffi-
%CONT>>


% ===== scan leaf 520, printed 499 =====
% !PRINTED 499
\weberpage{499}
\pdfbookmark[2]{§ 136. The functionals of an algebraic field Ω and the extended field Ω}{bm156}
%<<CONT
cients of $\varphi'$ are symmetric functions of the roots of the equation
\[ \frac{f(x)}{x - \Theta} = 0, \]
which are themselves contained in $\Omega$.

From the definition it follows that the norm of a product is equal to the product of the norms of the factors, and thus that, if $\varphi, \psi$ are functions in $\Omega$,
\weq{7}{N(\varphi\psi) = N(\varphi)\,N(\psi)}
holds.

\begin{center}
\S.\,136.

\gesperrt{The functionals of an algebraic field $\Omega$ and the extended field $\overline{\Omega}$.}
\end{center}

The variables that are used in the theory of algebraic numbers do not have the significance of symbols for variable sequences of numbers, as one is accustomed to from the theory of functions; rather, they are merely symbols of calculation without any independent meaning. In the functions of these variables what really matters is only the systems of coefficients, and the variables are used only in order to apply the well-known and familiar rules of literal calculation to these systems of coefficients. This does not, of course, exclude that occasionally the numbers are also considered which one obtains when the variables are replaced by certain numbers, e.g. by rational numbers.

\gesperrt{Accordingly we introduce into our considerations both integral and fractional functions of arbitrarily many variables, whose coefficients are numbers of an algebraic field $\Omega$, and stipulate that these functions are computed with as literal calculation prescribes.}

Every such function $\omega$ can be represented as a quotient of two integral functions in $\Omega$,
\weq{1}{\omega = \frac{\varphi}{\psi},}
where $\psi$ must always be assumed different from zero. Two such functions are only equal to one another:
\[ \frac{\varphi}{\psi} = \frac{\varphi_1}{\psi_1}, \]
%CONT>>


% ===== scan leaf 521, printed 500 =====
% !PRINTED 500
\weberpage{500}
%<<CONT
if $\varphi\psi_1 = \varphi_1\psi$. If the two functions $\varphi, \psi$ have no common divisor, then $\varphi:\psi$ is called an \gesperrt{irreducible fraction}. Among the various representations of a function $\omega$ there is \gesperrt{one} by an irreducible fraction, and this we shall call the \gesperrt{simplest representation}. In it, numerator and denominator are completely determined by $\omega$ itself, up to a common factor, which may be an arbitrary number in $\Omega$ (Bd. I, \S.\,141).

Such a function $\omega$ we shall call a \gesperrt{functional of the field} $\Omega$. The integral functions and the numbers themselves are contained among these as special cases. For the numbers, \gesperrt{every} representation as a quotient of two numbers in $\Omega$ is to be regarded as a simplest representation.

The four fundamental operations of arithmetic can be applied to the functionals to the same extent as to the numbers, and the totality of all functionals of the field $\Omega$ is therefore likewise a field, which we denote by $\overline{\Omega}$ and shall call the \gesperrt{functional field} of $\Omega$.\footnote{It is natural to designate the functionals of the field $\Omega$, with which one calculates as with the numbers in $\Omega$, outright as \gesperrt{ideal numbers} of the field $\Omega$. This mode of expression would attach itself on the one hand to the ideal factors of \gesperrt{Kummer}, on the other to the ideals introduced by \gesperrt{Dedekind}, which stand to the functionals in a very close relation, to be discussed later. However tempting such a terminology would be in some respects, it proves nevertheless not to be expedient in others, because it would apparently attribute to the functionals --- which are after all only a means to an end, not in themselves an object of our interest --- a greater importance than in reality belongs to them. In introducing the word ``Functional'' we follow a suggestion of \gesperrt{Dedekind}.}

The functional field $\overline{\Omega}$ contains the number field $\Omega$ as a divisor.

A functional whose coefficients are rational numbers is called a \gesperrt{rational functional}. The totality of all rational functionals is the functional field $\overline{R}$ of the field $R$ of the rational numbers, and the field $\overline{R}$ is contained in every algebraic functional field $\overline{\Omega}$.

Every rational functional $A$ can be represented as a quotient of two integral functions in $R$, to which one can also give integral coefficients, by multiplying numerator and denominator by the common denominator of all the coefficients.


% ===== scan leaf 522, printed 501 =====
% !PRINTED 501
\weberpage{501}

If in this representation numerator and denominator are imprimitive, one can take out the divisor
and obtains a representation of the form
\weq{2}{A = a\,\frac{E_1(x,y,z,\ldots)}{E_2(x,y,z,\ldots)} = a\,\frac{E_1}{E_2} = aE,}
in which $a$ is a \gesperrt{positive}, integral or fractional rational number,
while $E_1, E_2$ are primitive functions in $R$. From this
it follows:

1. \gesperrt{Every rational functional can be transformed,
by multiplication with a primitive function in $R$, into an
integral function in $R$, and several rational functionals
can be represented as fractions whose common denominator
is a primitive function.}

The positive number $a$ and the quotient $E_1:E_2$ are completely
determined by $A$. For if we put $a$ into the form of a
rational fraction $q_1:q_2$ and assume that
\[
\frac{q_1}{q_2}\,\frac{E_1}{E_2} = \frac{q_1'}{q_2'}\,\frac{E_1'}{E_2'},
\]
then it follows that
\[
q_1 q_2' E_1 E_2' = q_2 q_1' E_2 E_1'.
\]
% [sic: Weber prints ``Functionen $R$'' without ``in''.]
Here, then, we have two integral functions $R$ with integral
coefficients which are equal to one another, and whose divisor,
since $E_1 E_2'$ and $E_2 E_1'$ are primitive functions, is $q_1 q_2'$
or $q_2 q_1'$. Consequently $q_1 q_2' = q_2 q_1'$, and hence
also
\[
\frac{q_1}{q_2} = \frac{q_1'}{q_2'}, \quad \frac{E_1}{E_2} = \frac{E_1'}{E_2'}.
\]

2. \gesperrt{We call the positive rational number $a$ the
absolute value of the functional $A$.}

In the case in which $A$ is itself a number, $a$ is the absolute
value of $A$ in the ordinary sense of this word, and $E$
is to be set $= +1$ or $= -1$, according as $A$ is positive or
negative. Thus $E$ is nothing more than the sign of
$A$. In the far-reaching generalisation of these elementary
concepts lies the strangeness which our definition offers at first
glance. It will, however, prove in what follows to be thoroughly
appropriate and useful.

From \S.\,135, 3. there results the theorem:
%CONT>>


% ===== scan leaf 523, printed 502 =====
% !PRINTED 502
\weberpage{502}
%<<CONT

3. \gesperrt{The absolute value of a product of two or
more rational functionals is equal to the product of the
absolute values of the factors.}

If $\omega$ is a functional of the field $\Omega$, then $N(\omega)$ is a rational
functional, and the \gesperrt{norm of a product or of a quotient
of two functionals is equal to the product or the quotient
of the norms of the parts.}

4. \gesperrt{We call the absolute value of the rational
functional $N(\omega)$ the absolute norm $N_a(\omega)$ of
$\omega$ and set}
\weq{3}{N(\omega) = N_a(\omega)\,E(\omega).}
$N_a(\omega)$ is always a positive rational number, and $E(\omega)$
is the quotient of two primitive integral functions. From 3.,
if $\alpha, \beta$ are two functionals in $\Omega$, there follows the formula
\weq{4}{N_a(\alpha\beta) = N_a(\alpha)\,N_a(\beta),}
or the theorem:

5. \gesperrt{The absolute norm of a product of functionals
in $\Omega$ is equal to the product of the absolute norms of the
factors.}

If we replace all the coefficients of a functional in $\Omega$ by
the corresponding numbers of a field $\Omega_1$ conjugate to $\Omega$,
we obtain a \gesperrt{conjugate functional}. The entire
functional field $\overline{\Omega}$ thereby passes over into a
\gesperrt{conjugate functional field} $\overline{\Omega}_1$. Since every equation
between functionals of the field $\Omega$ is at bottom only the
combination of a series of equations between numbers of the
field $\Omega$, we have the theorem (Bd. I, \S.\,141):

6. \gesperrt{Every equation between functionals of the field
$\overline{\Omega}$ remains correct if in place of all the functionals the
corresponding elements of a conjugate field $\overline{\Omega}_1$
are put.}

If $t$ denotes a variable not occurring in $\omega$, then
the function $N(t-\omega)$ has the form
\weq{5}{\Phi(t) = t^n + A_1 t^{n-1} + A_2 t^{n-2} + \cdots + A_n,}
in which the coefficients $A_i$ are rational functionals, and this
function $\Phi(t)$ vanishes when $\omega$ is put for $t$. There is,
however, not merely one such function $\Phi(t)$ which for $t = \omega$
%CONT>>


% ===== scan leaf 524, printed 503 =====
% !PRINTED 503
\weberpage{503}
%<<CONT
vanishes, but arbitrarily many, since one can multiply each
$\Phi(t)$ by an arbitrary function of the same form. It can
also happen that already a product of fewer than $n$ factors
of $N(t-\omega)$ acquires rational coefficients, from which
functions $\Phi(t)$ of degree lower than $n$ arise, which
vanish for $t = \omega$. From this it follows:

7. \gesperrt{For every functional $\omega$ of the field $\Omega$
there are infinitely many functions $\Phi(t)$ which
vanish for $t = \omega$, in which the highest power of $t$
has the coefficient 1, and whose remaining
coefficients are contained in $\overline{R}$.}

We then say that $\omega$ is a root of the equation
\[
\Phi(t) = 0.
\]

Among the various functions $\Phi(t)$ whose existence is
asserted in Theorem 7. there is one and only one $F(t)$ of
lowest possible degree. For if there existed two such
functions $F(t)$ and $F_1(t)$ of the same degree $m$, the
difference $F(t) - F_1(t)$ would be at most of degree
$m - 1$ with respect to $t$ and would vanish for $t = \omega$.
If one divides by the coefficient of the highest power of $t$,
one obtains a function $\Phi(t)$ of lower degree than $F(t)$,
which by the assumption about $F(t)$ does not exist.

The function $F(t)$ cannot be decomposed in the field
$\overline{R}$ into factors that are integral functions of $t$.
For if it split into several factors of the same form
$F_1(t)$, $F_2(t)$, then one of these factors, which are after
all of lower degree than $F(t)$ itself, would have to vanish
for $t = \omega$, contrary to our assumption.

Every function $\Phi(t)$ of Theorem 7. is then divisible by
this irreducible function $F(t)$, so that the quotient
$\Phi(t):F(t)$ is an integral function of $t$ in $\overline{R}$,
in which the highest power of $t$ has the coefficient 1
(cf. Bd. I, \S.\,141, I.).

Among the functions $\Phi(t)$ there is also, as already
remarked, the norm $N(t-\omega)$, and consequently
$N(t-\omega)$ is divisible by $F(t)$. If the two functions are
of the same degree, then $N(t-\omega) = F(t)$. But if the
degree of $F(t)$ is lower than $n$, the quotient
$N(t-\omega):F(t)$ vanishes at least for one more of the
functionals conjugate to $\omega$, and consequently for all of
them; consequently also for $t = \omega$, and hence this
quotient
%CONT>>


% ===== scan leaf 525, printed 504 =====
% !PRINTED 504
\weberpage{504}
\pdfbookmark[2]{§ 137. Integral functionals}{bm157}
%<<CONT
once more by $F(t)$, that is, $N(t-\omega)$ is divisible by $F(t)^2$. By continuing this method of inference one recognises that $N(t-\omega)$ must be a power of $F(t)$, and that consequently the degree of $F(t)$ is a divisor of $n$ (Bd. I, \S.\,144).

We sum this up once more in the theorem:

8. \gesperrt{Every functional in $\Omega$ is the root of one and only one irreducible equation $F(t) = 0$ in $\overline{R}$ and $N(t-\omega)$ is a power of $F(t)$.}

\begin{center}
\S.\ 137.

\gesperrt{Integral Functionals.}
\end{center}

1. \textbf{Definition:} \gesperrt{A rational functional shall be called integral if its absolute value is an integer.}

From this definition it follows first of all that the sum, the difference and the product of two, and consequently also of arbitrarily many, integral rational functionals are again integral. For the product this is an immediate consequence of the theorem \S.\,136, 3. But in order to carry out the proof for the sum and the difference, we represent two integral rational functionals $A_1, A_2$ by fractional functions in such a way that they receive a primitive form as common denominator:
\[
A_1 = a_1\,\frac{E_1}{E}, \quad A_2 = a_2\,\frac{E_2}{E}.
\]
Then
\[
A_1 \pm A_2 = \frac{a_1 E_1 \pm a_2 E_2}{E},
\]
since now $a_1, a_2$ are integers, the divisor of the integral function $a_1 E_1 \pm a_2 E_2$ is the absolute value of $A_1 \pm A_2$. This is an integer, and hence $A_1 \pm A_2$ is an integral functional.

For constant rational functionals, that is, for rational numbers, Definition 1. gives the \gesperrt{integral rational numbers.}

We further set up, just as in \S.\,133 for the numbers, the following definition of the integral functionals of the field $\Omega$.

2. \textbf{Definition:} \gesperrt{A functional $\omega$ from $\Omega$ is called integral if it is the root of an equation}
\[
\Phi(t) = t^m + A_1 t^{m-1} + A_2 t^{m-2} + \cdots + A_m = 0
\]
%CONT>>


% ===== scan leaf 526, printed 505 =====
% !PRINTED 505
\weberpage{505}
%<<CONT
\gesperrt{in which the coefficients $A_1, A_2, \ldots, A_m$ are integral rational functionals.}

Functionals which do not belong among the integral ones we shall occasionally also, for the sake of brevity, designate as fractional functionals.

In justification of this definition we first prove the theorem:

3. \gesperrt{If $\omega$ is an integral functional according to Definition 2. and at the same time rational, then it is an integral rational functional (according to Definition 1.).}

Suppose that $\omega$ is a fractional rational functional, and hence that the absolute value of $\omega$ is a rational fraction $p:q$, in which $p$ and $q$ are integral rational numbers without a common divisor; then $q\omega$ is an integral rational functional whose absolute value is $=p$, and therefore relatively prime to $q$. But if, on the other hand, $\omega$ is at the same time integral in the sense of Definition 2., then we can set
\[
\omega^m = -(A_1\omega^{m-1} + A_2\omega^{m-2} \ldots)
\]
from which
\[
(q\omega)^m = -q[A_1(q\omega)^{m-1} + A_2\,q(q\omega)^{m-2} + \cdots) % [sic] opening square bracket is closed by a round parenthesis as printed
\]
follows.

But from this it results that the absolute value $p^m$ of $(q\omega)^m$ must be divisible by $q$, which is possible only if $q=1$. Definition 1. is therefore contained in Definition 2. as a special case.

If
\[
\Phi(t) = t^m + A_1t^{m-1} + A_2t^{m-2} + \cdots + A_m
\]
is a function of $t$ in which the coefficients $A_1, A_2, \ldots, A_m$ are fractional rational functionals with the variables $x, y, \ldots$, but free of the variable $t$, then one can determine a function $aE(x,y,\ldots)$ in which $a$ denotes the common denominator of the absolute values of $A_1, A_2, \ldots$ and $E(x,y,\ldots)$ a primitive integral function, so that
\weq{1}{aE(x\ y,\ldots)\,\Phi(t) = P(t,x,y,\ldots)} % [sic] printed without comma between x and y
is itself a primitive integral function [§.\,136, 2.].

If $\Phi(t)$ splits into two factors $\Phi_1(t)$, $\Phi_2(t)$ in $R$, that is, if
\[
\Phi(t) = \Phi_1(t)\,\Phi_2(t),
\]
then let one determine accordingly for the two functions $\Phi_1, \Phi_2$ the factors $a_1E_1, a_2E_2$, so that
%CONT>>


% ===== scan leaf 527, printed 506 =====
% !PRINTED 506
\weberpage{506}
%<<CONT
\[
a_1E_1\Phi_1 = P_1, \quad a_2E_2\Phi_2 = P_2
\]
become primitive functions, and then
\weq{2}{a_1a_2E_1E_2\Phi = P_1P_2}
is also a primitive function. From (1) and (2) it follows that
\weq{3}{a_1a_2E_1E_2P = aEP_1P_2}
and consequently
\weq{4}{a = a_1a_2.}
If therefore $a=1$, the natural numbers $a_1, a_2$ must also be $=1$, and we have the theorem:

4. \gesperrt{If $\omega$ is an integral functional in $\Omega$ and $\Phi(t)$ a function of $t$ with integral coefficients in $\overline{R}$ which vanishes for $t=\omega$, then every rational divisor of $\Phi(t)$ also has integral coefficients in $\overline{R}$; in particular, the irreducible function $F(t)$, of which $\omega$ is a root according to §.\,136, 8., has integral rational functionals as coefficients.}

This theorem is of fundamental importance because it teaches that a functional $\omega$ which in the field $\Omega$ belongs among the fractional ones cannot, as might be supposed, appear in a higher field as an integral functional.

If a functional $\omega$ is not integral, it satisfies an equation
\[
\Phi(\omega) = \omega^m + A_1\omega^{m-1} + A_2\omega^{m-2} + \cdots + A_m = 0,
\]
in which the coefficients $A_1, A_2, \ldots, A_m$ are rational, but not integral, functionals. If $a$ is the common denominator of the absolute values of $A_1, A_2, \ldots, A_m$, then $aA_1, aA_2, \ldots, aA_m$ are integral functionals. Now
\[
a^m\Phi(\omega) = (a\omega)^m + aA_1(a\omega)^{m-1} + a^2A_2(a\omega)^{m-2} + \cdots
+ a^mA_m = 0,
\]
and $a\omega$ is therefore an integral functional. From this it follows:

5. \gesperrt{Every functional $\omega$ of the field $\Omega$ can be transformed into an integral functional by multiplication with an integral rational number, and therefore every functional $\omega$ can be represented as a quotient of two integral functionals. This representation is possible in infinitely many different ways, among others in such a way that the denominator is a natural number.}


% ===== scan leaf 528, printed 507 =====
% !PRINTED 507
\weberpage{507}
6. \gesperrt{If $\omega$ is a functional in $\Omega$, and if there are $m$ quantities $\omega_1, \omega_2, \ldots, \omega_m$, which do not all vanish, of such a kind that the $m$ products $\omega\omega_i$ for $i=1,2,\ldots,m$ can be put into the form:}
\weq{5}{\omega\,\omega_i = A_{1,i}\omega_1 + A_{2,i}\omega_2 + \cdots + A_{m,i}\omega_m,}
\gesperrt{in which the $m^2$ symbols $A_{k,i}$ are integral rational functionals, then $\omega$ is an integral functional.}

Here $m$ is any natural integer. In the application the $\omega_i$ are always numbers or functionals in $\Omega$; for the validity of the theorem, however, only the practicability of the multiplication indicated in (5) is essential.

The theorem is a simple consequence of the definition of the integral functionals; for since the $\omega_i$ do not all vanish, by the determinant theorem (Bd. I, §.\,24, II.) we must have
\weq{6}{\begin{vmatrix}
A_{1,1}-\omega, & A_{2,1}, & \ldots, & A_{m,1} \\
A_{1,2}, & A_{2,2}-\omega, & \ldots, & A_{m,2} \\
\hdotsfor{4} \\
A_{1,m}, & A_{2,m}, & \ldots, & A_{m,m}-\omega
\end{vmatrix} = 0}
By arranging according to powers of $\omega$ there results from this an equation
\weq{7}{\omega^m + A_1\omega^{m-1} + \cdots + A_m = 0,}
in which the coefficients $A_1, A_2, \ldots$ are composed from the $A_{i,k}$ by addition and multiplication, and are therefore, by 1., integral rational functionals; accordingly $\omega$ is also an integral functional.

For the last coefficient $A_m$ in the equation (7) we obtain the expression
\weq{8}{(-1)^mA_m = \Sigma \pm A_{1,1}A_{2,2}\ldots A_{m,m},}
\gesperrt{which is therefore equal to the product of all the roots of the equation (7).}

Theorem 4. is important as a criterion for integral functionals; it serves us here for the proof of the following theorem: % [sic: Weber prints Theorem 4.; the determinant criterion just proved is Theorem 6.]

7. \gesperrt{If $\Psi(x,y,\ldots)$ is an integral function whose coefficients are integral rational numbers or functionals, but free of the variables $x,y,\ldots$, and if further $\alpha,\beta,\ldots$ are integral functionals in $\Omega$, then}
\weq{9}{\omega = \Psi(\alpha,\beta,\ldots)}
\gesperrt{is also an integral functional.}


% ===== scan leaf 529, printed 508 =====
% !PRINTED 508
\weberpage{508}
Namely, let $\mu, \nu, \ldots$ be the degrees of the equations with integral coefficients which (by 2.) the functionals $\alpha, \beta, \ldots$ satisfy, and $m=\mu\nu\ldots$ By $\omega_1, \omega_2, \ldots, \omega_m$ we understand the $m$ quantities
\[
\alpha^r\beta^s\ldots \qquad r=0,1,\ldots,\mu-1, \quad s=0,1,\ldots,\nu-1, \ldots
\]
Then, by means of the equations of degrees $\mu, \nu, \ldots$ which the numbers $\alpha, \beta, \ldots$ satisfy, all products $\alpha^r\beta^s\ldots$ in which one of the exponents $r, s, \ldots$ is greater than $\mu-1, \nu-1, \ldots$ can be expressed linearly in the form (5) in terms of $\omega_1, \omega_2, \ldots, \omega_m$, and the same therefore also holds of every functional $\omega$, and consequently also of the products $\omega\omega_1, \omega\omega_2, \ldots, \omega\omega_m$. From this it follows by 6. that $\omega$ is integral, as we wished to prove.

As a special case of the theorem 7., from which, incidentally, the general theorem can easily be inferred again, we emphasise:

8. \gesperrt{By addition, subtraction and multiplication of integral functionals there always arise integral functionals again.}

We have already remarked earlier (§.\,134) that one can always determine an algebraic field which contains arbitrarily given algebraic numbers. From this it follows that in the theorem 7. the $\alpha, \beta, \ldots$ may be arbitrary integral algebraic numbers or functionals, and the corresponding statement holds of the theorem 8.

From the definition 2. it follows further, by the theorem §.\,136, 8.:

9. \gesperrt{If $\omega$ is an integral functional of the field $\Omega$, then all functionals conjugate to $\omega$ are also integral.}

As a particular consequence of this theorem, the following may further be mentioned:

10. \gesperrt{The absolute norm of an integral functional is a natural integer.}

Finally we prove also the following theorem:

11. \gesperrt{If a functional $\omega$ satisfies an equation of the form}
\weq{10}{\omega^m + \alpha_1\omega^{m-1} + \cdots + \alpha_m = 0}
\gesperrt{in which $\alpha_1, \alpha_2, \ldots, \alpha_m$ are integral functionals in $\Omega$, then $\omega$ is also an integral functional.}


% ===== scan leaf 530, printed 509 =====
% !PRINTED 509
\weberpage{509}
\pdfbookmark[2]{§ 138. Divisibility. Associated functionals. Units}{bm158}

To prove it, we denote by $t$ a variable which does not occur in the $\alpha$ and in $\omega$, and set

\weq{11}{\varphi(t) = t^m + \alpha_1 t^{m-1} + \cdots + \alpha_m.}

Then, if $n$ is the degree of the field $\Omega$,
\[
\Phi(t) = N[\varphi(t)]
\]
is an integral function of $t$ of degree $mn$, whose coefficients are integral rational functionals. At the same time $\Phi(\omega) = 0$, and consequently $\omega$ is an integer.

\gesperrt{From this it follows further that Theorem 6. remains correct if the coefficients $A_{k,i}$ are integral functionals not in $R$, but in $\Omega$.}

If an integral functional $\omega$ is at the same time a number, then it is an integer, because in this case $N(t-\omega)$ has integral rational numbers as coefficients. The integers which we considered in \S.\,133 are accordingly contained as special cases among the integral functionals.

\S.\,138.

\gesperrt{Divisibility. Associated functionals. Units.}

The integral functionals are subject to laws of divisibility similar to those of the integral rational numbers. To recognise them, we place the following definition at the head:

1. \gesperrt{An integral functional $\alpha$ is called divisible by another integral functional $\beta$, different from zero, if the quotient $\alpha:\beta=\gamma$ is an integral functional.}

Then $\alpha=\beta\gamma$, and $\beta$ and $\gamma$ are called \gesperrt{divisors} of $\alpha$, and one also says that \gesperrt{$\beta$ and $\gamma$ divide $\alpha$}. The number 0 is divisible by \gesperrt{every integral} functional.

From this definition the following fundamental theorems on divisibility are obtained without difficulty:

2. \gesperrt{If $\alpha$ and $\alpha_1$ are divisible by $\beta$, then $\alpha\pm\alpha_1$ is also divisible by $\beta$.}

For if $\alpha=\beta\gamma$, $\alpha_1=\beta\gamma_1$, then $\alpha\pm\alpha_1=\beta(\gamma\pm\gamma_1)$, and if $\gamma$ and $\gamma_1$ are integral, then $\gamma\pm\gamma_1$ is also integral. % [sic] source prints ``nnd'' for ``und''


% ===== scan leaf 531, printed 510 =====
% !PRINTED 510
\weberpage{510}

3. \gesperrt{If $\alpha$ is divisible by $\beta$, and $\beta$ is divisible by
$\gamma$, then $\alpha$ is also divisible by $\gamma$.}

For according to the assumption there are two integral functionals $\varkappa$,
$\lambda$ which satisfy the conditions $\alpha = \varkappa\beta$, $\beta = \lambda\gamma$.
Accordingly we have also $\alpha = \varkappa\lambda\gamma$, and since $\varkappa\lambda$
is integral, $\alpha$ is divisible by $\gamma$.

A product $\beta\gamma$ of two integral functionals is divisible both by $\beta$ and by
$\gamma$, and consequently, if $\beta$ is divisible by $\alpha$, then $\beta\gamma$ is
also divisible by $\alpha$. From these theorems it follows:

4. \gesperrt{If $\alpha$, $\beta$, \ldots\ are divisible by $\delta$, and if $\xi$,
$\eta$, \ldots\ are arbitrary integral functionals, then $\xi\alpha + \eta\beta + \ldots$
is divisible by $\delta$.}

These definitions and theorems extend, of course, also to the case in which numbers
take the place of functionals, and thus we obtain the divisibility of
integers.

5. \gesperrt{Two integral functionals $\alpha$, $\beta$ which are mutually divisible
by one another are called associated.}

6. \gesperrt{An integral functional $\varepsilon$ associated with the natural number 1,
i.\,e. every divisor of the number 1, is called a unit.}

According as $\varepsilon$ is a functional or a number, $\varepsilon$ is a
\gesperrt{functional} or a \gesperrt{numerical} unit.

In the field $R$ of the rational numbers, the primitive integral functions and the
quotients of two among them are to be regarded as the functional units. Of numerical
units there are in $R$ only two, namely $+\,1$ and $-\,1$.

Concerning the concepts hereby introduced, which stand in close mutual relation,
we derive a series of theorems.

7. \gesperrt{If $\alpha$, $\beta$ are associated functionals, then the quotients
$\beta : \alpha$ and $\alpha : \beta$ are units.}

For if we set $\beta = \alpha\varepsilon$, then $\varepsilon$ is an integral
functional and $1 : \varepsilon = \alpha : \beta$ is likewise integral; hence
$\varepsilon$ is a divisor of the number 1, i.\,e. $\varepsilon$ is a unit.

8. \gesperrt{If $\alpha$ is an integral functional and $\varepsilon$ a unit, then
$\alpha$ and $\alpha\varepsilon$ are associated.}

This follows immediately from the definition; for $\alpha : \alpha\varepsilon =
1 : \varepsilon$ and $\alpha\varepsilon : \alpha = \varepsilon$ are both integral
functionals.


% ===== scan leaf 532, printed 511 =====
% !PRINTED 511
\weberpage{511}

A unit is an integral functional whose reciprocal is likewise integral, and this reciprocal functional is itself a unit. Any integral functional whatever is divisible by a unit. In general the theorem holds:

9. \gesperrt{The product and the quotient of two units are again units.}

For if $\varepsilon_1, \varepsilon_2$ are two units, then $\varepsilon_1\varepsilon_2$ and $1:\varepsilon_1\varepsilon_2$ are integral functionals, hence $\varepsilon_1\varepsilon_2$ is a unit, and likewise $\varepsilon_1:\varepsilon_2$ and $\varepsilon_2:\varepsilon_1$ are integral.

10. \gesperrt{If an integral functional $\mu$ is divisible by another, $\alpha$, then every functional $\mu'$ associated with $\mu$ is also divisible by every functional $\alpha'$ associated with $\alpha$.}

For if $\mu:\alpha$ is integral and $\varepsilon, \varepsilon_1$ are units, then $\mu\varepsilon:\alpha\varepsilon_1$ is also integral.

11. \gesperrt{If $\alpha$ is associated with $\beta$ and with $\gamma$, then $\beta$ and $\gamma$ are also associated with one another.}

For if $\beta=\alpha\varepsilon$, $\gamma=\alpha\varepsilon_1$, then $\beta=\gamma\varepsilon:\varepsilon_1$, and $\varepsilon:\varepsilon_1$ is, by 8., a unit.

If $\alpha$ is divisible by $\beta$, then the absolute norm of $\alpha$ is divisible by the absolute norm of $\beta$; for from $\alpha=\beta\gamma$ it follows, by \S.\,136, 5.:

\weq{1}{N_a(\alpha) = N_a(\beta)\, N_a(\gamma).}

From this it results further that the absolute norm of a unit must be $=1$. But the converse also holds:

12. \gesperrt{An integral functional whose absolute norm is $=1$ is a unit.}

For if $\varepsilon$ is an integral functional with the absolute norm 1, then $N(\varepsilon)$ is an integral rational functional with the absolute value 1, and therefore $1:N(\varepsilon)$ is also an integral functional. If we then set
\[
N(\varepsilon) = \varepsilon\varepsilon',
\]
then $\varepsilon'$, as a product of integral functionals (the conjugates of $\varepsilon$), is itself integral, and consequently
\[
\frac{1}{\varepsilon} = \frac{\varepsilon'}{N(\varepsilon)}
\]
is also an integral functional, hence $\varepsilon$ is a unit.


% ===== scan leaf 533, printed 512 =====
% !PRINTED 512
\weberpage{512}
\pdfbookmark[2]{§ 139. Greatest common divisor}{bm159}

From this we further conclude, by formula (1), that \gesperrt{associated functionals have the same absolute norm}.

According to this, an integral rational functional is always associated with its absolute value. If then $\alpha$ is any integral functional of the field $\overline{\Omega}$, then $N(\alpha)$ and $N_a(\alpha)$ are also associated. Now since in $N(\alpha)=\alpha\alpha'$ the factor $\alpha'$, being a product of integral functionals, is itself integral, $N(\alpha)$, and consequently also the natural number $N_a(\alpha)$, is divisible by $\alpha$. We therefore formulate the further theorem:

13. \gesperrt{There exist natural integers (in infinite number), among them the absolute norm of $\alpha$, which are divisible by an arbitrary integral functional $\alpha$. If $a$ is the least among the natural numbers divisible by $\alpha$, then every integral rational number divisible by $\alpha$ is divisible by $a$.}

For if $m$ is an integral rational number divisible by $\alpha$, then the residue of the division of $m$ by $a$ is also an integral rational number divisible by $\alpha$. Since this is smaller than $a$, it must be $=0$.

\S.\,139.

\gesperrt{Greatest common divisor.}

Let $\alpha, \beta, \ldots$ denote integral functionals in $\Omega$ different from zero, in arbitrary but finite number, and $x, y, \ldots$ variables which do not occur in $\alpha, \beta, \ldots$. Then

\weq{1}{\delta = \alpha x + \beta y + \cdots}

is likewise an integral functional in $\Omega$. The norm of $\delta$ is an integral homogeneous function of degree $n$ of the variables $x, y, \ldots$, whose coefficients are integral rational functionals. If we denote the absolute norm of $\delta$ by $D$, then

\weq{2}{N(\delta) = D\,E(x, y, \ldots) = D\,E,}

and in this $E(x, y, \ldots) = E$ is an integral rational function of the variables $x, y, \ldots$ and in general a fractional function of the variables occurring in $\alpha, \beta, \ldots$, but in any case a \gesperrt{rational unit} [\S.\,136, (2)].


% ===== scan leaf 534, printed 513 =====
% !PRINTED 513
\weberpage{513}

We shall now understand by $x_0, y_0, \dots$ any integral rational numbers whatever. Then
\[ \delta_0 = \alpha x_0 + \beta y_0 + \cdots \]
is also an integral functional, and we shall prove \gesperrt{that $\delta_0$ is divisible by $\delta$.}

For this purpose we need only consider the integral functional
\[ \delta t - \delta_0 = \alpha (x t - x_0) + \beta (y t - y_0) \dots \]
in which $t$ denotes a new variable. Then [by (2)]:

\weq{3}{N (\delta t - \delta_0) = D E (x t - x_0, y t - y_0, \dots),}

or, if we arrange according to descending powers of $t$:

\weq{4}{N (\delta t - \delta_0) = D (t^n E + t^{n-1} E_1 + t^{n-2} E_2 + \cdots),}

in which $E_1, E_2, \dots$ are, by \S.\,137, 7., 8., integral rational functionals. If we now set
\[ C_1 = \frac{E_1}{E}, \quad C_2 = \frac{E_2}{E}, \cdots, \]
then, since $E$ is a unit, $C_1, C_2, \dots$ are also integral rational functionals; and if we further set
\[ \frac{\delta_0}{\delta} = \eta \]
it follows from (4) that
\[ N (\delta) N (t - \eta) = D E (t^n + C_1 t^{n-1} + C_2 t^{n-2} + \cdots), \]
or, on account of (2),

\weq{5}{N (t - \eta) = t^n + C_1 t^{n-1} + C_2 t^{n-2} + \cdots}

Now since this function vanishes when we set $t = \eta$, it follows that $\eta$ is an integral functional (\S.\,137, 2.), and it is thus proved that $\delta_0$ is divisible by $\delta$.

Since $x_0, y_0, \dots$ may denote arbitrary integral rational numbers, we conclude from this that the functionals $\alpha, \beta, \dots$ are themselves divisible by $\delta$, and hence that $\delta$ is a common divisor of the functionals $\alpha, \beta, \dots$.

On the other hand, however, $\delta$ is also (by \S.\,138, 4.) divisible by every common divisor of $\alpha, \beta, \dots$, and $\delta$ therefore has the \gesperrt{characteristic properties of the greatest common divisor of $\alpha, \beta, \dots$} The same properties, however, belong by \S.\,138, 7. to every functional associated with $\delta$; and likewise two functionals $\delta, \delta'$ that have this double prop-
%CONT>>


% ===== scan leaf 535, printed 514 =====
% !PRINTED 514
\weberpage{514}
%<<CONT
erty are divisible by one another, hence associated, and we accordingly set up the definition:

1. \gesperrt{The functional $\delta = \alpha x + \beta y + \cdots$ and every functional associated with it is called a greatest common divisor of $\alpha, \beta, \dots$}

If $\delta$ is a unit, we also say that $\alpha, \beta, \dots$ are without common divisor, for then there is, apart from the units, no integral functional which divides all of $\alpha, \beta, \dots$

2. \gesperrt{Two functionals $\alpha, \beta$ which have no common divisor, for which therefore $\alpha x + \beta y$ is a unit, are called relative prime functionals or coprime.}

From these definitions the following theorems result very simply:

3. \gesperrt{If integral functionals $\xi, \eta, \dots$ exist in $\Omega$ such that
\[ \alpha \xi + \beta \eta + \cdots = \varepsilon \]
is a unit, then the integral functionals $\alpha, \beta, \dots$ are without common divisor.}

For if $\alpha, \beta, \dots$ have a common divisor $\delta$, then (by \S.\,138) $\delta$ is also a divisor of $\alpha \xi + \beta \eta + \cdots$, and $\delta$ must therefore also be a unit.

4. \gesperrt{If $\alpha, \beta, \gamma$ are three integral functionals and $\alpha$ is relatively prime to $\beta$ and to $\gamma$, then $\alpha$ is also relatively prime to $\beta \gamma$.}

For by assumption, if $x, y, u, v$ are four variables,
\[ \varepsilon = \alpha x + \beta y, \quad \varepsilon_1 = \alpha u + \gamma v \]
are units. Accordingly
\[ \alpha (\alpha u x + \gamma v x + \beta u y) + \beta \gamma v y = \varepsilon \varepsilon_1 \]
is also a unit. But since $\alpha u x + \gamma v x + \beta u y$ and $v y$ are integral, it follows from Theorem 3. that $\alpha$ and $\beta \gamma$ are relatively prime.

To this is joined the proof of the following very important theorem:

5. \gesperrt{If $\alpha, \beta, \mu$ are integral functionals, $\alpha$ relatively prime to $\beta$ and $\alpha \mu$ divisible by $\beta$, then $\mu$ is divisible by $\beta$.}


% ===== scan leaf 536, printed 515 =====
% !PRINTED 515
\weberpage{515}
\pdfbookmark[2]{§ 140. Prime Functionals in the Field Ω}{bm160}

For by the assumption concerning $\alpha, \beta$,
\[ \alpha x + \beta y = \varepsilon \]
is a unit. By multiplication with $\mu$ it follows from this that
\[ \alpha \mu x + \beta \mu y = \varepsilon \mu, \]
and since $\alpha \mu$ and $\beta \mu$ are by assumption divisible by $\beta$, it follows that $\varepsilon \mu$, and consequently also the $\mu$ associated with $\varepsilon \mu$, is divisible by $\beta$.

\begin{center}
\S.\,140.

\gesperrt{Prime functionals in the field $\Omega$.}
\end{center}

By the theorems of the preceding section we have gained the means to derive the laws of divisibility of the integral functionals in the field $\Omega$ in exactly the same way as is applied in the elements of arithmetic to the natural integers. We define as follows:

1. \gesperrt{An integral functional $\pi$ of the field $\Omega$ which is not a unit is called a prime functional if, apart from the units, it is divisible only by the functionals associated with itself. Every integral functional in $\Omega$ which has other divisors besides these is called composite.}

The concept of the prime functional is accordingly essentially dependent on the field $\Omega$. The prime functionals of one field may very well be composite in another, extended field.

If $\omega$ is an arbitrary integral functional and $\pi$ a prime functional of the field $\Omega$, only two cases are possible: either $\omega$ is relatively prime to $\pi$, or $\omega$ is divisible by $\pi$; for a common divisor of $\omega$ and $\pi$ can only be either a unit or associated with $\pi$, and in the latter case $\omega$ is divisible by $\pi$. From this follows the theorem:

2. \gesperrt{If the product $\alpha \beta$ of two integral functionals $\alpha, \beta$ in $\Omega$ is divisible by a prime functional $\pi$, then one of the two factors must be divisible by $\pi$.}


% ===== scan leaf 537, printed 516 =====
% !PRINTED 516
\weberpage{516}
\pdfbookmark[2]{§ 141. Decomposition of the Integral and Fractional Functionals into Prime Factors}{bm161}

For if $\alpha$ and $\beta$ are both not divisible by $\pi$, and hence both relatively prime to $\pi$, then by \S.\,139, 4. $\alpha \beta$ is also relatively prime to $\pi$.

From this it follows by repeated application that, if a product of \gesperrt{several} factors is divisible by $\pi$, at least one of the factors must be divisible by $\pi$.

3. \gesperrt{The least natural integer $p$ that is divisible by a prime functional $\pi$ in $\Omega$ is a natural prime, and the absolute norm of $\pi$ is a power of $p$. Every integral rational number divisible by $\pi$ is also divisible by $p$.}

By \S.\,138, 13. there are natural numbers that are divisible by $\pi$. The least among them, $p$, cannot be decomposable into two natural factors that are greater than 1; for if $p = p_1 p_2$, then either $p_1$ or $p_2$ must be divisible by $\pi$ (by 2.). But if neither of the factors $p_1, p_2 = 1$, then they are both less than $p$, which contradicts the assumption about $p$. Consequently $p$ is a natural prime. That every natural number $m$ divisible by $\pi$ is divisible by $p$ we have already proved in \S.\,138, 13.

Let us now set

\weq{1}{p = \pi \omega,}

then $\omega$ is an integral functional, and if we take the absolute norms on both sides, it follows, since the absolute norm of a natural number is the $n$th power of this number:

\weq{2}{p^n = N_a (\pi) N_a (\omega).}

From this follows the second part of our theorem, that the natural integer $N_a (\pi)$ is a power of $p$.

If accordingly we set

\weq{3}{N_a (\pi) = p^f,}

then $f$ is a number that can have only one of the values $1, 2, 3, \dots, n$, if $n$ denotes the degree of the field $\Omega$.

The number $f$ is called the \gesperrt{degree} of the prime functional $\pi$.

\begin{center}
\S.\,141.

\gesperrt{Decomposition of the integral and fractional functionals into prime factors.}
\end{center}

1. \gesperrt{Every integral functional $\omega$ of the field $\Omega$ different from zero that is not a unit is divisible by a prime functional.}


% ===== scan leaf 538, printed 517 =====
% !PRINTED 517
\weberpage{517}
If $\omega$ is prime, the theorem is evident, because $\omega$ is divisible by itself. But if $\omega$ is not prime, then it is divisible by an integral functional $\omega_1$ which is neither a unit nor associated with $\omega$. If, then,
\weq{1}{\omega = \omega_1 \alpha,}
neither $\omega_1$ nor $\alpha$ is a unit. But from (1) follows
\[
N_a(\omega) = N_a(\omega_1)\, N_a(\beta), % [sic] printed with $\beta$, though (1) introduces $\alpha$
\]
and since $N_a(\beta)$ is greater than 1, we have
\weq{2}{N_a(\omega_1) < N_a(\omega).}

If $\omega_1$ is not yet prime, one can apply the same inference to $\omega_1$ and finds a divisor $\omega_2$ of $\omega_1$ such that
\weq{3}{N_a(\omega_2) < N_a(\omega_1) < N_a(\omega)}
holds. But since there are only a finite number of natural numbers that are smaller than $N_a(\omega)$, the sequence of functionals $\omega, \omega_1, \omega_2, \ldots$ must break off, and this is possible only if the last of them is a prime functional, whereby Theorem 1. is proved.

2. \gesperrt{Every integral functional $\omega$ in the field $\Omega$ that is different from zero and from the units can be decomposed into a finite number of prime factors.}

For if $\pi_1$ is a prime factor of $\omega$ and
\weq{4}{\omega = \pi_1 \omega_1,}
then, since $N_a(\pi_1) > 1$, we have
\[
N_a(\omega) > N_a(\omega_1).
\]

If $\omega_1$ is not a unit, it is divisible by a prime functional $\pi_2$, and from
\[
\omega_1 = \pi_2 \omega_2
\]
follows
\[
N_a(\omega_1) > N_a(\omega_2).
\]

Continuing in this way, one obtains a sequence of integral functionals $\omega_1, \omega_2, \ldots$, whose absolute norms continually decrease, and this sequence therefore breaks off with a unit. If $\pi_\nu$ is the last of them that is not a unit, then $\pi_\nu$ is itself a prime functional, and it follows that
\weq{5}{\omega = \pi_1 \pi_2 \ldots \pi_\nu,}
q. e. d.


% ===== scan leaf 539, printed 518 =====
% !PRINTED 518
\weberpage{518}
3. \gesperrt{An integral functional $\omega$ in the field $\Omega$ is decomposable into prime factors in only one way, if associated prime factors are regarded as not different. Associated functionals contain the same prime factors.}

For let us assume that the two products of prime factors
\weq{6}{\pi_1 \pi_2 \ldots \pi_\nu, \qquad \varkappa_1 \varkappa_2 \ldots \varkappa_\mu}
are associated with one another; then the product $\pi_1 \pi_2 \ldots \pi_\nu$ is divisible by the prime factor $\varkappa_1$, and hence, by \S.\,140, 2., one of the factors, say $\pi_1$, must be divisible by $\varkappa_1$ and consequently be associated with $\varkappa_1$. Then the products
\[
\pi_2 \ldots \pi_\nu, \qquad \varkappa_2 \ldots \varkappa_\mu
\]
are also associated, and consequently one of the factors of the first product, say $\pi_2$, is divisible by $\varkappa_2$ and therefore associated with $\varkappa_2$. In this way one can continue to reason, and it results not only that the number of the $\varkappa$ must agree with the number of the $\pi$, but also that the $\varkappa$ are individually associated with the $\pi$.

Among the prime factors of an integral functional $\omega$ the same one can occur several times, and these mutually equal (or associated) factors can be combined into a power. If $\omega$ is divisible by $\pi^h$, we say that \gesperrt{the prime functional $\pi$ divides $\omega$ $h$ times.}

According to this it has a perfectly definite sense to speak of \gesperrt{the prime factors} of an integral functional. From what has been proved we further conclude:

4. \gesperrt{An integral functional $\alpha$ is divisible by another $\beta$ then and only then, when all prime factors of $\beta$ occur among the prime factors of $\alpha$, and each of them divides $\alpha$ at least as often as it divides $\beta$.}

For if $\alpha$ is divisible by $\beta$, and $\beta$ by $\pi^h$, then $\alpha$ is also divisible by $\pi^h$.

If $\alpha, \beta, \ldots$ are integral functionals, $\pi_1, \pi_2, \ldots$ different prime functionals, and $\varepsilon_1, \varepsilon_2, \ldots$ units, then we can determine exponents $a_1, a_2, \ldots, b_1, b_2, \ldots$ in such a way that
\weq{7}{\begin{aligned}
\varepsilon_1 \alpha &= \pi_1^{a_1} \pi_2^{a_2} \ldots \\
\varepsilon_2 \beta &= \pi_1^{b_1} \pi_2^{b_2} \ldots \\
&\ldots\ldots\ldots\ldots
\end{aligned}}
%CONT>>


% ===== scan leaf 540, printed 519 =====
% !PRINTED 519
\weberpage{519}
%<<CONT
holds, provided we set the exponent equal to zero when one of the prime factors does not divide the functional in question. A common divisor of the numbers $\alpha, \beta, \ldots$ can contain no prime factors other than $\pi_1, \pi_2, \ldots$, and every common divisor of $\alpha, \beta, \ldots$ has the form
\weq{8}{\varepsilon \delta = \pi_1^a \pi_2^b \ldots,}
in which $a$ may not be greater than the least of the numbers $a_1, b_1, \ldots$, and $b$ not greater than the least of the numbers $a_2, b_2, \ldots$, and so forth.

If $a$ is the least among the numbers $a_1, b_1, \ldots$, and $b$ the least among the numbers $a_2, b_2, \ldots$, then the number $\delta$ represented in (8) is the greatest common divisor of the numbers $\alpha, \beta, \ldots$ Expressed in words:

One obtains the greatest common divisor of several integral functionals $\alpha, \beta, \ldots$ by forming a product of prime factors, into which one takes up each prime factor as often as it divides each of the numbers $\alpha, \beta, \ldots$.

Two numbers are relatively prime when they contain no common prime factor.

Correspondingly we define as the \gesperrt{least common multiplum} $\mu$ of the functionals $\alpha, \beta, \ldots$ a product of prime factors, into which we take up a factor $\pi$ so often that it divides none of the functionals $\alpha, \beta, \ldots$ more often than it divides $\mu$. This functional $\mu$ then has, just like every functional associated with $\mu$, the twofold and, as we may add, characteristic property that it is divisible by each of the functionals $\alpha, \beta, \ldots$, and that every other functional which is divisible by $\alpha, \beta, \ldots$ is also divisible by $\mu$.

The least common multiplum of two relative prime functionals is their product.

These theorems can be applied in order to represent fractional functionals \gesperrt{in the simplest shape, or as reduced fractions}, by decomposing numerator and denominator into their prime factors and cancelling the greatest common divisor. Numerator and denominator of a reduced fraction are completely determined by the fraction itself, apart from a common unit factor, which remains undetermined.

In the same way one can bring any number of given fractions to a common denominator, the principal denominator, by choosing a common multiplum of all the given denominators as the common denominator.


% ===== scan leaf 541, printed 520 =====
% !PRINTED 520
\weberpage{520}
\pdfbookmark[2]{§ 142. Integral Functions in the Field Ω}{bm162}
If the given functionals are actually fractions, that is, if they do not all reduce to integral functionals, then the common denominator must contain at least one prime factor that does not occur in all the numerators.

All this is in perfect agreement with the rules of elementary arithmetic, and the methods of proof that we have employed here are also essentially the same as those used there. The central point of the deduction is, on the one hand, the far-reaching generalisation of the concept of the unit, and on the other, the definition, founded upon it, of the greatest common divisor in \S.\,139.

\begin{center}
\S.\,142.\\
\gesperrt{Integral Functions in the Field $\overline{\Omega}$.}
\end{center}

The means we now have at our disposal suffice to apply to the functionals the mode of reasoning of which we made use in \S.\,2 of the first Band for the proof of \gesperrt{Gauss}'s theorem.

In fact, the theorem on which that proof essentially rests --- that a product of two integers is divisible by a prime only when this prime divides one of the factors --- has been proved valid also for the prime functionals now introduced.

We may accordingly, by following the proof just mentioned step by step, regard the theorem as proved:

1. \gesperrt{If two integral functions of a variable $t$, of degrees $h$ and $k$:}
\[
\varphi = \alpha_0 t^h + \alpha_1 t^{h-1} + \cdots + \alpha_h
\]
\[
\psi = \beta_0 t^k + \beta_1 t^{k-1} + \cdots + \beta_k,
\]
\gesperrt{whose coefficients are integral functionals that do not contain the variable $t$, have a product}
\[
\chi = \gamma_0 t^{h+k} + \gamma_1 t^{h+k-1} + \cdots + \gamma_{h+k}
\]
\gesperrt{in which the coefficients $\gamma_0, \gamma_1, \ldots, \gamma_{h+k}$ have a common prime divisor $\pi$, then $\pi$ must divide either all the coefficients of $\varphi$ or all the coefficients of $\psi$.}


% ===== scan leaf 542, printed 521 =====
% !PRINTED 521
\weberpage{521}

From this theorem we here draw an important consequence: The functions

\weq{1}{\varphi = \varphi_0 t^h + \varphi_1 t^{h-1} + \cdots + \varphi_h,}

in which the coefficients $\varphi_0, \varphi_1, \ldots, \varphi_h$ are integral or fractional functionals in $\Omega$, but are assumed to be free of the variables $t$, themselves belong to the functionals in the field $\Omega$. Of them the following theorem holds:

2. \gesperrt{A functional $\varphi$ is integral only when the coefficients $\varphi_0, \varphi_1, \ldots, \varphi_h$ are integral functionals.}

To prove it, we assume that $\varphi_0, \varphi_1, \ldots, \varphi_h$ are not all integral at once. If we determine their common denominator $\mu$ and set

\[
\mu\varphi_0 = \alpha_0, \quad \mu\varphi_1 = \alpha_1, \quad \ldots, \quad \mu\varphi_h = \alpha_h,
\]

then $\alpha_0, \alpha_1, \ldots, \alpha_h$ are integral functionals, and $\mu$ contains at least one prime factor $\pi$ which does not at the same time divide all the numerators $\alpha_0,, \alpha_1, \ldots, \alpha_h$ % [sic] doubled comma in print
(cf. the close of the preceding section).

Then the function

\weq{2}{\chi = \mu\varphi = \alpha_0 t^h + \alpha_1 t^{h-1} + \cdots + \alpha_h,}

whose coefficients are now integral, is certainly an integral functional.

If we now assume that $\varphi$ itself is integral, then $\mu\varphi$ is divisible by $\pi$, while yet not all the coefficients $\alpha_0, \alpha_1, \ldots, \alpha_h$ are divisible by $\pi$. Now if $\varphi$ is an integral functional, it satisfies an equation of the form

\weq{3}{\varphi^m = C_1 \varphi^{m-1} + C_2 \varphi^{m-2} + \cdots + C_m,}

in which the coefficients $C_1, C_2, \ldots, C_m$ are integral rational functionals.

Following \S.\,136, (2), we put these functionals in the form

\[
C_1 = \frac{a_1 E_1}{E}, \quad C_2 = \frac{a_2 E_2}{E}, \quad \ldots, \quad C_m = \frac{a_m E_m}{E},
\]

in which $a_1, a_2, \ldots, a_m$ are the absolute values of $C_1, C_2, \ldots, C_m$, hence natural integers (or zero), and $E, E_1, \ldots, E_m$ are primitive integral functions, hence units.

Then equation (3) yields:

\[
E\varphi^m = a_1 E_1 \varphi^{m-1} + a_2 E_2 \varphi^{m-2} + \cdots + a_m E_m,
\]

and by multiplication with $\mu^m$:

\weq{4}{E\chi^m = \mu\left(a_1 E_1 \chi^{m-1} + a_2 E_2 \mu \chi^{m-2} + \cdots + a_m E_m \mu^{m-1}\right).}


% ===== scan leaf 543, printed 522 =====
% !PRINTED 522
\weberpage{522}

Here there now stand, on the right hand and on the left, integral functions of $t$ whose coefficients are integral functionals. On the right side all these coefficients have the factor $\mu$, hence also the factor $\pi$, while by assumption not all the coefficients of $\chi$ have this factor. Since $E$ is a unit, the coefficients of $E$ likewise do not contain the factor $\pi$, and consequently, by Theorem 1., the coefficients of $E\chi^m$ also cannot all be divisible by $\pi$, which the equation (4) would nevertheless require. From this it follows that our assumption, that $\varphi$ is integral while $\varphi_0, \varphi_1, \ldots, \varphi_h$ by contrast are not all integral, is inadmissible, and Theorem 2. is thus proved.

If we now assume that in $\varphi$ the coefficients $\varphi_0, \varphi_1, \ldots$ are themselves in turn integral functions of a variable, and apply Theorem 2. repeatedly to this, we arrive at the conclusion:

3. \gesperrt{An integral rational function of arbitrarily many variables, whose coefficients are numbers or functionals with other variables, is an integral functional only when the coefficients are integral.}

We can represent every functional $\omega$ as a quotient of two integral functions in such a way that the denominator becomes a primitive function in the field $R$.

For if $\varphi, \psi$ are integral functions in $\Omega$, and if

\[
\omega = \frac{\varphi}{\psi}, \quad N(\psi) = \psi\psi',
\]

then we can multiply both terms of the fraction $\omega$ by $\psi'$ and obtain in the denominator an integral function with rational coefficients. We can then reckon the divisor of this function to the numerator of $\omega$, and we obtain, if $E$ denotes a primitive function and $\chi$ an integral function in $\Omega$:

\[
E\omega = \chi,
\]

that is, one can transform every functional $\omega$ in $\Omega$, by multiplication with a primitive function $E$, into an integral function of the variables whose coefficients are numbers in $\Omega$.

With the aid of Theorem 3. we obtain from this:

4. \gesperrt{Every integral functional $\omega$ in the field $\Omega$ is associated with an integral function $\chi$, whose coeffi-}
%CONT>>


% ===== scan leaf 544, printed 523 =====
% !PRINTED 523
\weberpage{523}
\pdfbookmark[2]{§ 143. The Prime Factors of the Numbers of the Field Ω}{bm163}

%<<CONT
\gesperrt{cients are integers in $\Omega$, and passes over into $\chi$ through multiplication with a rational functional which is a unit.}

\begin{center}
\S.\,143.

\gesperrt{The prime factors of the numbers of the field $\Omega$.}
\end{center}

Since the numbers are also contained among the functionals of the field $\Omega$, it follows from the results of \S.\,141 that the integers of the field $\Omega$ can likewise be decomposed into prime factors, but into prime factors which in general are not numbers but functionals. It is, however, not excluded that in special cases numbers too may occur among the prime functionals, which are then called \gesperrt{primes of the field $\Omega$}.

All equations between functionals are in the last instance identical equations between integral rational functions, and resolve themselves into a sequence of equations between numbers. They therefore remain correct when other symbols are put in place of the variables.

From this it follows that an integral functional, a unit, a prime functional do not cease to be integral functionals, units, prime functionals when other symbols are put in place of the variables.

If, then, we decompose a number into its prime factors, and if functionals occur among these prime factors, the variables in the equation representing this decomposition can be replaced by arbitrary other variables, and from this follows the generalisation:

1. \gesperrt{The prime factors of an integral functional $\omega$ can be so represented that they do not contain the variables on which $\omega$ depends.}

For every integral functional $\omega$ is a divisor of numbers, even of rational numbers, for example of the absolute norm of $\omega$. The prime factors of $\omega$ are therefore to be sought among the prime factors of one of these numbers, and can accordingly be represented by variables which are different from those occurring in $\omega$.


% ===== scan leaf 545, printed 524 =====
% !PRINTED 524
\weberpage{524}

By \S.\,142, 4., if $\omega$ is a given integral functional, $E$ a unit, and $\varphi$ an integral function in $\Omega$, we can set:

\weq{1}{E\omega = \varphi(x, y, \ldots).}

If, on the other hand, we decompose $\omega$ into its prime factors, then by 1. we may assume these prime factors to be free of the variables $x, y, \ldots$, and if we again form the product of these prime factors, we obtain a number $\omega_1$ associated with $\omega$, which is free of the variables $x, y, \mathinner{.\,.}$.

If we order the quotient $\varphi : \omega_1$, which is an integer (indeed a unit), according to the variables $x, y, \ldots$, then we conclude from \S.\,142, 3. that all the coefficients of the function $\varphi$ are divisible by $\omega_1$, and consequently also by $\omega$ and $\varphi$.

If, on the other hand, all the coefficients of $\varphi$ are divisible by some factor $\delta$ in $\Omega$, then $\varphi$ too is divisible by $\delta$, and from this it follows:

2. \gesperrt{An integral function $\varphi$ with integers as coefficients is the greatest common divisor of all its coefficients.}

The integer $\varphi$ therefore passes, by \S.\,139, into an associated number when the individual powers and products of the variables are each replaced by a single variable. From this it also follows as a corollary that every functional $\omega$ passes into an associated one when the variables are in any way denoted differently.

We now have to prove the following important theorem:

3. \gesperrt{If $\omega$ is an arbitrary integral functional, one can choose a number $\alpha$ divisible by $\omega$ in such a way that the quotient $\alpha : \omega$ is relatively prime to an arbitrarily given functional $\mu$.}

We first prove that there is an integer $\alpha$ which is divisible by $\omega$, but not by $\omega\pi$, where $\pi$ is an arbitrary prime functional.

For if, following \S.\,142, 4., we form an integral function $\varphi$ associated with $\omega$, then the coefficients of this function are indeed all divisible by $\omega$, but not all divisible by $\omega\pi$, since otherwise $\varphi$, and hence $\omega$ itself, would be divisible by $\omega\pi$, which is not possible. Among the coefficients of $\varphi$ there is therefore at least one which has the required property.


% ===== scan leaf 546, printed 525 =====
% !PRINTED 525
\weberpage{525}
Now let $\pi_1, \pi_2, \pi_3, \ldots$ be an arbitrary number of given prime functionals, different from one another. We set

\[
\omega_1 = \omega\pi_2\pi_3\ldots, \quad \omega_2 = \omega\pi_1\pi_3\ldots, \quad \omega_3 = \omega\pi_1\pi_2\ldots,
\]
and determine, by what has just been proved, the integers $\alpha_1, \alpha_2, \alpha_3, \ldots$ in $\Omega$, so that
\[
\begin{array}{l}
\alpha_1 \text{ becomes divisible by } \omega_1 \text{, but not by } \omega_1\pi_1\\
\alpha_2 \quad \text{\textquotedbl} \quad \text{\textquotedbl} \quad \text{\textquotedbl} \quad \omega_2 \quad \text{\textquotedbl} \quad \text{\textquotedbl} \quad \text{\textquotedbl} \quad \omega_2\pi_2\\
\alpha_3 \quad \text{\textquotedbl} \quad \text{\textquotedbl} \quad \text{\textquotedbl} \quad \omega_3 \quad \text{\textquotedbl} \quad \text{\textquotedbl} \quad \text{\textquotedbl} \quad \omega_3\pi_3\\
\cdots\cdots\cdots\cdots\cdots\cdots\cdots\cdots\cdots\cdots\cdots\cdots\cdots\cdots\cdots\cdots,
\end{array}
\]
and from these derive the integer
\[
\alpha = \alpha_1 + \alpha_2 + \alpha_3 + \cdots
\]
This number is evidently divisible by $\omega$, since all the summands $\alpha_1, \alpha_2, \alpha_3, \ldots$ are divisible by $\omega$. But it is not divisible by $\omega\pi_1$, because $\alpha_2, \alpha_3, \ldots$, though not $\alpha_1$, are divisible by $\omega\pi_1$; and likewise it is not divisible by $\omega\pi_2, \omega\pi_3, \ldots$. If, therefore, we set
\[
\alpha = \omega\eta
\]
then $\eta$ is an integral functional which is not divisible by $\pi_1, \pi_2, \pi_3, \ldots$, and which therefore, when $\pi_1, \pi_2, \pi_3, \ldots$ are the prime factors of $\mu$ different from one another, is relatively prime to $\mu$.

Now let us assume, arbitrarily, an integer $\beta$ in $\Omega$ divisible by $\omega$, and let us set $\beta = \omega\mu$; then we can determine $\alpha = \omega\eta$ so that $\eta$ becomes relatively prime to $\mu$, and then $\omega$ is the greatest common divisor of $\alpha$ and $\beta$. From this it follows:

4. \gesperrt{Every integral functional $\omega$ of the field $\Omega$ is the greatest common divisor of two integers, and consequently $\omega$ is associated with a binary linear form $\alpha x + \beta y$, in which $\alpha$ and $\beta$ are integers.}

\begin{center}
\rule{4cm}{0.4pt}
\end{center}

A general remark, one that reaches over into another region, may still find its place here.

It is the principal result of this chapter that the integral algebraic numbers in a definite field can be decomposed into prime factors in a unique manner, exactly in the same way as is known for the integral rational numbers; though indeed only by the fact that the content of the field
%CONT>>


% ===== scan leaf 547, printed 526 =====
% !PRINTED 526
\weberpage{526}
%<<CONT
(by the adjunction of variables) is enlarged. There thus arises an extended field, in which the laws of decomposability hold purely.

The starting point of the definitions was formed by the field $R$ of the rational numbers, and if we wish to give ourselves an account of the question on which properties of the field $R$ the possibility of this extension rests, we find that it is on the one hand the existence of the integers in $R$, and on the other hand the unique decomposability of these integers into prime factors, that alone were used in the whole deduction. If, then, we let any other field to which these two properties belong take the place of the field $R$, we shall be able to draw the same conclusions. If we take for $R$ another algebraic number field, we obtain, to be sure, nothing new; but we certainly do if, for example, we put in place of the field $R$ the field of the rational functions of one variable, to which, indeed, the two fundamental properties also belong. In this way we gain a starting point for the theory of the algebraic functions of one variable\footnote{Cf. \gesperrt{Dedekind-Weber}, Theorie der algebraischen Functionen einer Veränderlichen. \gesperrt{Crelle's} Journal, Bd. 92.}.

\begin{center}
\rule{4cm}{0.4pt}
\end{center}


% ===== scan leaf 548, printed 527 =====
\weberpage{527}% page number not printed on this page
\pdfbookmark[1]{Seventeenth Section. Theory of Algebraic Fields.}{bm164}
\pdfbookmark[2]{§ 144. Basis of an Algebraic Number Field. Discriminants}{bm165}
% !PRINTED none

\begin{center}
\gesperrt{Seventeenth Chapter.}

\gesperrt{Theory of Algebraic Fields.}

\rule{4cm}{0.4pt}

\S.\,144.

\gesperrt{Basis of an algebraic number field.}

\gesperrt{Discriminants.}
\end{center}

Let
\weq{1}{f(\Theta) = \Theta^n + a_1\Theta^{n-1} + \cdots + a_n = 0}
be the irreducible equation of degree $n$ with rational coefficients $a_1, a_2, \ldots, a_n$, which defines for us an algebraic field $\Omega = R(\Theta)$. The field $\Omega$ is then the totality of all numbers of the form
\weq{2}{\omega = h_1 + h_2\Theta + h_3\Theta^2 + \cdots + h_n\Theta^{n-1},}
in which $h_1, h_2, \ldots, h_n$ are \gesperrt{rational numbers}. But if we put rational functionals for $h_1, h_2, \ldots, h_n$, we obtain from (2) all \gesperrt{functionals} of the field $\Omega$.

If we consider an arbitrary system of $n$ numbers
\weq{3}{\omega_r = h_{1,r} + h_{2,r}\Theta + \ldots + h_{n,r}\Theta^{n-1}, \quad r = 1,2,\ldots,n}
under the assumption that the determinant
\weq{4}{H = \Sigma \pm h_{1,1}h_{2,2}\ldots h_{n,n}}
is different from zero, then, on account of the irreducibility of $f$, an equation of the form
\weq{5}{k_1\omega_1 + k_2\omega_2 + \cdots + k_n\omega_n = 0,}
in which $k_1, k_2, \ldots, k_n$ are rational numbers, can subsist only when these coefficients are all zero, and this still holds even when \gesperrt{rational functionals} are admitted for the coefficients in equation (5).


% ===== scan leaf 549, printed 528 =====
% !PRINTED 528
\weberpage{528}

If, however, we eliminate from the equations (2) and (3) the powers of $\Theta$, there results a relation of the form
\weq{6}{\omega = k_1\omega_1 + k_2\omega_2 + \cdots + k_n\omega_n.}

Every number and every functional of the field $\Omega$ can therefore be represented in the form (6), if one puts rational numbers or functionals for $k_1, k_2, \ldots, k_n$. For a given $\omega$ this representation is possible [because of (5)] in only one way, and $\omega$ is a number when the $k_1, k_2, \ldots, k_n$ are numbers, but a functional when functionals also occur among the $k$.

A system of numbers such as
\[
\omega_1, \omega_2, \ldots, \omega_n,
\]
we call a \gesperrt{basis of the field $\Omega$}. The powers of $\Theta$ also form such a basis:
\[
1, \Theta, \Theta^2, \ldots, \Theta^{n-1}.
\]

If we denote by $\omega_{r,1}, \omega_{r,2}, \ldots, \omega_{r,n}$ the numbers conjugate to $\omega_r$, then the square of the determinant
\weq{7}{\Delta(\omega_1, \omega_2, \ldots, \omega_n) = (\Sigma \pm \omega_{1,1}\omega_{2,2}\ldots\omega_{n,n})^2}
is a symmetric function of the conjugate values $\Theta_1, \Theta_2, \ldots, \Theta_n$ and is consequently a rational number.

This number is called the \gesperrt{discriminant of the system} $\omega_1, \omega_2, \ldots, \omega_n$. In particular,
\weq{8}{\Delta(1, \Theta, \Theta^2, \ldots, \Theta^{n-1}) = \begin{vmatrix} 1, & \Theta_1, & \ldots, & \Theta_1^{\,n-1}\\ 1, & \Theta_2, & \ldots, & \Theta_2^{\,n-1}\\ \multicolumn{4}{c}{\cdots\cdots\cdots\cdots\cdots\cdots}\\ 1, & \Theta_n, & \ldots, & \Theta_n^{\,n-1} \end{vmatrix}^2}
is the discriminant of the equation (1) (Bd.\ I, \S.\,46), for which, if $N$ is the symbol for the norm, one can also put
\weq{9}{(-1)^{\frac{n(n-1)}{2}} N f'(\Theta)}
This number is therefore certainly different from zero. By the multiplication theorem of determinants (Bd.\ I, \S.\,27) there results from (3) and (7):
\weq{10}{\Delta(\omega_1, \omega_2, \ldots, \omega_n) = H^2\Delta(1, \Theta, \Theta^2, \ldots, \Theta^{n-1}),}
from which one concludes that the discriminant of a basis of $\Omega$ is always different from zero. Since $H$ is a rational number, it follows that the ratio of the discriminants of different bases is the square of a rational number, and that the
%CONT>>


% ===== scan leaf 550, printed 529 =====
% !PRINTED 529
\weberpage{529}
\pdfbookmark[2]{§ 145. The Minimal Basis and the Field Discriminant}{bm166}
%<<CONT
discriminants of all bases of $\Omega$ have the same sign. Formula (10) also shows that any system of $n$ numbers $\omega_r$ of the field $\Omega$ is a basis of $\Omega$ whenever the determinant square (7) does not vanish.

If $\omega_1, \omega_2, \ldots, \omega_n$ is a basis of $\Omega$, and if $c_{r,s}$ denote rational numbers, then the $n$ numbers
\[
\omega_r' = c_{r,1}\omega_1 + c_{r,2}\omega_2 + \cdots + c_{r,n}\omega_n \qquad r = 1, 2, \ldots, n
\]
also form a basis of $\Omega$, provided the determinant
\[
C = \Sigma \pm c_{1,1}c_{2,2} \ldots c_{n,n}
\]
does not vanish, and we have
\weq{11}{\varDelta(\omega_1', \omega_2', \ldots, \omega_n') = C^2 \varDelta(\omega_1, \omega_2, \ldots, \omega_n).}

If, for example, we multiply the elements $\omega_1, \omega_2, \ldots, \omega_n$ of a basis by rational coefficients $c_1, c_2, \ldots, c_n$, none of which vanishes, we obtain a new basis
\[
c_1\omega_1, c_2\omega_2, \ldots, c_n\omega_n.
\]

\begin{center}
\S.\,145.

\gesperrt{The Minimal Basis and the Field Discriminant.}
\end{center}

By the remark just made, a basis of $\Omega$ does not lose the property of being a basis when each of its numbers is multiplied by a rational number different from zero. Now, according to \S.\,133,5., every number can be converted into an integer by multiplication by an integral rational number, and from this it follows \gesperrt{that there are bases of $\Omega$ whose elements are all integers. The discriminant of such a basis is an integral rational number.} This integral rational number is different from zero. It changes when another integral basis is chosen, but for a definite field it retains an unchanged sign.

Among all these integers, which can occur as discriminants of an integral basis, and which all stand to one another in quadratic ratio, one must now be the smallest in absolute value. This smallest discriminant we denote by $\varDelta$, and we call it the \gesperrt{fundamental number}, or also the \gesperrt{discriminant of the field} $\Omega$.


% ===== scan leaf 551, printed 530 =====
% !PRINTED 530
\weberpage{530}
This $\varDelta$ is a positive or negative, but never vanishing, integral rational number completely determined by $\Omega$, and there is always a basis of $\Omega$ consisting of integers $\omega_1, \omega_2, \ldots, \omega_n$ whose discriminant is equal to $\varDelta$.

Such a basis we shall briefly call a \gesperrt{minimal basis of} $\Omega$.

If we understand by $k_1, k_2, \ldots, k_n$ any integral rational numbers whatever, then every number of the shape
\weq{1}{\omega = k_1\omega_1 + k_2\omega_2 + \cdots + k_n\omega_n}
is an integral algebraic number, and we now prove the fundamental theorem:

1. \gesperrt{If $\omega_1, \omega_2, \ldots, \omega_n$ is a minimal basis, then all the integers of the field $\Omega$ are contained in the form (1).}

Since $\omega_1, \omega_2, \ldots, \omega_n$ is a basis, every number in $\Omega$ can first of all be represented in the form (1) if rational fractions are admitted for $k_1, k_2, \ldots, k_n$. Let us assume, then, that an integer $\omega$ is representable in the form
\weq{2}{\omega = \frac{k_1\omega_1 + k_2\omega_2 + \cdots + k_n\omega_n}{k}}
where $k_1, k_2, \ldots, k_n, k$ are integral rational numbers such that not all of $k_1, k_2, \ldots, k_n$ have a common divisor with $k$. If $p$ is any natural prime dividing $k$, and $k = pk'$, then at least one of the coefficients $k_1, k_2, \ldots, k_n$ must be indivisible by $p$. Suppose, say, that $k_1$ is not divisible by $p$; then the integral rational number $l$ can be determined so that $l k_1 \equiv 1 \; (\mathrm{mod}\ p)$, or $(l k_1 - 1)$ becomes divisible by $p$. It then follows from (2):
\weq{3}{l k' \omega - \frac{l k_1 - 1}{p}\omega_1 = \frac{\omega_1 + l k_2\omega_2 + \cdots + l k_n\omega_n}{p} = \omega_1',}
and $\omega_1'$ is likewise an integral algebraic number. If we set
\weq{4}{\omega_2' = \omega_2, \ldots, \omega_n' = \omega_n,}
then the numbers $\omega_1', \omega_2', \ldots, \omega_n'$ form an \gesperrt{integral basis} of $\Omega$, because the numbers $\omega_1, \omega_2, \ldots, \omega_n$, and consequently all the numbers $\omega$, can be expressed linearly in terms of $\omega_1', \omega_2', \ldots, \omega_n'$, and formula (11) of the preceding section yields
\weq{5}{\varDelta(\omega_1', \omega_2', \ldots, \omega_n') = \frac{1}{p^2}\varDelta(\omega_1, \omega_2, \ldots, \omega_n).}


% ===== scan leaf 552, printed 531 =====
% !PRINTED 531
\weberpage{531}
The discriminant $\varDelta(\omega_1', \omega_2', \ldots, \omega_n')$ is therefore smaller than $\varDelta(\omega_1, \omega_2, \ldots, \omega_n)$, and this contradicts the assumption that $\omega_1, \omega_2, \ldots, \omega_n$ is a minimal basis. With this our theorem is proved.

If one denotes the totality of all \gesperrt{integers} of the field $\Omega$ by $\mathfrak{o}$, then one obtains all the numbers of $\mathfrak{o}$, and each only once, if in
\weq{6}{k_1\omega_1 + k_2\omega_2 + \cdots + k_n\omega_n}
one lets the coefficients $k_1, k_2, \ldots, k_n$ run through all the integral rational numbers.

For this reason a minimal basis of $\Omega$ is also called a \gesperrt{basis of $\mathfrak{o}$}.

\gesperrt{The discriminant of a basis of $\mathfrak{o}$ is equal to the fundamental number of the field $\Omega$.}

By formula (11) of the preceding section we can derive from a basis of $\mathfrak{o}$ arbitrarily many others by linear substitution:
\weq{7}{(\omega_1', \omega_2', \ldots, \omega_n') = C(\omega_1, \omega_2, \ldots, \omega_n),}
if $C$ (in the notation of \S.\,37) denotes a linear substitution with rational integral coefficients and the determinant $\pm 1$.

Since, by the meaning of the basis of $\mathfrak{o}$, all integers in $\Omega$, and hence also the elements $\omega_1', \omega_2', \ldots, \omega_n'$ of a second basis of $\mathfrak{o}$, are representable linearly with integral rational coefficients by $\omega_1, \omega_2, \ldots, \omega_n$, it follows conversely that one can derive all the others from one basis of $\mathfrak{o}$ by such linear substitutions with the determinant $\pm 1$.

For if
\[
(\omega_1, \omega_2, \ldots, \omega_n) = C'(\omega_1', \omega_2', \ldots, \omega_n'),
\]
then the composed substitution $CC'$ must be the identical one, and the product of the two determinants is therefore 1, and hence each of them $= \pm 1$.

Since among the numbers of $\mathfrak{o}$ the number 1 is always contained, one can, if $\omega_1, \omega_2, \ldots, \omega_n$ is a basis of $\mathfrak{o}$, determine the \gesperrt{integral rational numbers} $c_1, c_2, \ldots, c_n$ so that the relation
\weq{8}{c_1\omega_1 + c_2\omega_2 + \cdots + c_n\omega_n = 1}
is satisfied.


% ===== scan leaf 553, printed 532 =====
% !PRINTED 532
\weberpage{532}
\pdfbookmark[2]{§ 146. The Bases of the Functionals}{bm167}
But the following theorem also holds with respect to the \gesperrt{functionals}:

2. \gesperrt{If $\omega_1, \omega_2, \ldots, \omega_n$ is a basis of $\mathfrak{o}$, then in the form}
\weq{9}{\omega = u_1\omega_1 + u_2\omega_2 + \cdots + u_n\omega_n,}
\gesperrt{in which $u_1, u_2, \ldots, u_n$ are integral rational functionals, all integral functionals in $\Omega$ are contained.}

For we have already shown above that all functionals whatsoever are contained in the form (9) when the coefficients $u_1, u_2, \ldots, u_n$ are integral or fractional rational functionals. But we can always determine an integral primitive function $e$ in such a way that
\[
e u_1 = y_1, \; e u_2 = y_2, \; \ldots, \; e u_n = y_n
\]
are integral functions of the variables, and then
\weq{10}{e\omega = y_1\omega_1 + y_2\omega_2 + \cdots + y_n\omega_n,}
and since $e$ is a unit, $e\omega$ is integral together with $\omega$. Now since the coefficients of the powers and products of the variables in the function (10) must, by \S.\,142, 3., be integers, it follows by 1. that the coefficients in the functions $y_1, y_2, \ldots, y_n$ must be integral rational numbers, and that consequently $u_1, u_2, \ldots, u_n$ are integral rational functionals.

\begin{center}
\S.\,146.

\gesperrt{The Bases of the Functionals.}
\end{center}

\gesperrt{If $\mu$ is an integral functional of the field $\Omega$, then by a basis of $\mu$ we understand a system of $n$ integers in $\Omega$:}
\weq{1}{\alpha_1, \alpha_2, \ldots, \alpha_n,}
\gesperrt{which is a basis of the field $\Omega$, and which has the property that in the form}
\weq{2}{\alpha = x_1\alpha_1 + x_2\alpha_2 + \cdots + x_n\alpha_n}
\gesperrt{all the integers of the field $\Omega$ that are divisible by $\mu$, and no others, are contained, when integral rational numbers are put for $x_1, x_2, \ldots, x_n$.}

It shall now be proved that every integral functional has a basis.


% ===== scan leaf 554, printed 533 =====
% !PRINTED 533
\weberpage{533}

First of all it is clear that a basis of a functional $\mu$ is at the
same time a basis of all the functionals associated with $\mu$; further,
that every element $\alpha_1, \alpha_2, \ldots$ of such a basis must be
divisible by $\mu$; and finally, that from one basis of $\mu$ all the
others can be derived by applying a linear substitution with integral
rational coefficients and determinant $\pm 1$.

In order now to form a basis of $\mu$, we start from a minimal basis
(basis of $\mathfrak{o}$), which, as above, we denote by
\[
\omega_1, \omega_2, \ldots, \omega_n
\]

Since there are positive integral rational numbers that are divisible by
$\mu$, there are also integral rational numbers $a$ for which the
product $a\,\omega_1$ is divisible by $\mu$. The \gesperrt{least} positive among
these numbers we will denote by $a_{1,1}$, and we set
\[
\alpha_1 = a_{1,1}\,\omega_1
\]
Thereupon we denote by $a_{2,2}$ the \gesperrt{least positive} integral rational
number for which an integral rational $a_{1,2}$ can be determined in
such a way that
\[
\alpha_2 = a_{1,2}\,\omega_1 + a_{2,2}\,\omega_2
\]
becomes divisible by $\mu$, and so we continue. In this way we form the
system
\weq{3}{
\begin{aligned}
\alpha_1 &= a_{1,1}\,\omega_1 \\
\alpha_2 &= a_{1,2}\,\omega_1 + a_{2,2}\,\omega_2 \\
&\cdots\cdots\cdots\cdots\cdots\cdots\cdots\cdots\cdots\cdots \\
\alpha_n &= a_{1,n}\,\omega_1 + a_{2,n}\,\omega_2 + \cdots + a_{n,n}\,\omega_n,
\end{aligned}
}
\gesperrt{in which $a_{1,1}, a_{2,2}, \ldots, a_{n,n}$ are the least
positive integral rational numbers that make possible a determination of
the integral rational numbers $a_{1,2}, \ldots, a_{n-1,n}$ such that the
integers $\alpha_1, \alpha_2, \ldots, \alpha_n$ become divisible by $\mu$.}

Every number $a_{r,r}$ is, according to this, defined
independently of the others by the fact that it is the least natural
number for which the associated integral rational numbers
$a_{1,r}, a_{2,r}, \ldots, a_{r-1,r}$ can be so determined that
\weq{4}{
\alpha_r = a_{1,r}\,\omega_1 + a_{2,r}\,\omega_2 + \cdots +
a_{r-1,r}\,\omega_{r-1} + a_{r,r}\,\omega_r
}
becomes divisible by $\mu$. In the particular case where
$\omega_r$ is itself divisible by $\mu$, one has accordingly to set
$a_{r,r} = 1$, and
%CONT>>


% ===== scan leaf 555, printed 534 =====
% !PRINTED 534
\weberpage{534}

%<<CONT
the $a_{1,r}, a_{2,r}, \ldots, a_{r-1,r}$ can all be assumed equal to
zero.

If we denote by $\Delta$ the fundamental number of the field $\Omega$,
the discriminant of the system $\alpha_1, \alpha_2, \ldots, \alpha_n$
determined by (3) is obtained according to the formula \S.\,144, (11):
\weq{5}{
\Delta(\alpha_1, \alpha_2, \ldots, \alpha_n) = a_{1,1}^2\, a_{2,2}^2
\ldots a_{n,n}^2\, \Delta.
}
This is an integral rational number different from zero, and
consequently the quantities $\alpha$ are a basis of $\Omega$.

In order, then, to show that the system $\alpha_1, \alpha_2, \ldots,
\alpha_n$ so determined is a basis of $\mu$, it still remains to be
proved that every integer $\alpha$ divisible by $\mu$ can be
represented in the form (2). To carry out this proof, let us assume
any index $r \leqq n$, and seek the condition for an integer of the
form
\weq{6}{
\gamma_r = h_1\,\omega_1 + h_2\,\omega_2 + \cdots + h_r\,\omega_r
}
to be divisible by $\mu$, when $h_1, h_2, \ldots, h_r$ are integral
rational numbers.

First of all it follows that $h_r$ must be divisible by $a_{r,r}$.
For if we denote by $q_r$ the residue of the division of $h_r$ by
$a_{r,r}$, and set
\[
h_r = l_r\,a_{r,r} + q_r, \qquad 0 \leqq q_r < a_{r,r},
\]
then by (6) we obtain
\[
\gamma_r - l_r\,\alpha_r = (h_1 - l_r\,a_{1,r})\,\omega_1 +
(h_2 - l_r\,a_{2,r})\,\omega_2 + \cdots + q_r\,\omega_r,
\]
and this number would also have to be divisible by $\mu$. By the
definition of $a_{r,r}$, however, this is possible only if $q_r = 0$.
But then $\gamma_r - l_r\,\alpha_r$ receives the expression:
\weq{7}{
\gamma_r - l_r\,\alpha_r = h_1'\,\omega_1 + h_1'\,\omega_2 + \cdots +
h_{r-1}'\,\omega_{r-1}, % [sic]
}
and thus becomes of the same form as (6), only that $r - 1$ takes the
place of $r$; and in (7) the same mode of reasoning is to be repeated.
Accordingly, by complete induction we obtain the theorem:

\gesperrt{An integer of the field $\Omega$ representable in the form
$h_1\,\omega_1 + h_2\,\omega_2 + \cdots + h_r\,\omega_r$ is always
then and only then divisible by $\mu$, when it is representable in
the form}
\[
l_1\,\alpha_1 + l_2\,\alpha_2 + \cdots + l_r\,\alpha_r
\]
\gesperrt{in which the coefficients $l_1, l_2, \ldots, l_r$ are
integral rational numbers.}

If one sets $r = n$ in this theorem, one has the proof that the
system of numbers (3) is a basis of $\mu$.


% ===== scan leaf 556, printed 535 =====
% !PRINTED 535
\weberpage{535}

We have already seen above how, from one basis of $\mu$, all the
others can be derived.

We now understand by $\alpha_1, \alpha_2, \ldots, \alpha_n$ an
arbitrary basis of $\mu$, and represent the functional $\mu$,
according to \S.\,142, 4., as a quotient of two integral functions
whose denominator is a rational unit. The coefficients of the
numerator are then integers divisible by $\mu$ (\S.\,143, 2.), and can
therefore be represented in the form (2).

If we gather this representation together suitably, there thus results
for $\mu$ an expression
\weq{8}{
\mu = \alpha_1 u_1 + \alpha_2 u_2 + \cdots + \alpha_n u_n,
}
in which the $u_1, u_2, \ldots, u_n$ are \gesperrt{integral rational
functionals}.

With this we wish to compare the linear form
\weq{9}{
\lambda = \alpha_1 t_1 + \alpha_2 t_2 + \cdots + \alpha_n t_n
}
in which $t_1, t_2, \ldots, t_n$ are variables. This linear form is
(by \S.\,139) the greatest common divisor of the numbers
$\alpha_1, \alpha_2, \ldots, \alpha_n$, and is divisible by $\mu$,
because the numbers $\alpha_1, \alpha_2, \ldots, \alpha_n$ are
divisible by $\mu$. On the other hand, however, as the representation
(8) shows, $\mu$ is also divisible by $\lambda$, and consequently the
two functionals $\mu$ and $\lambda$ are associated with one another.

\gesperrt{The functional $\lambda$ we shall call a basis form of the
functional $\mu$; $\lambda$ is then at the same time a basis form of
all functionals associated with $\mu$.}

A basis form $\lambda$ of $\mu$ has the property that from
it one obtains all the integers in $\Omega$ divisible by $\mu$ (or by
$\lambda$) when one sets integral rational numbers for the
variables.

The basis $\alpha_1, \alpha_2, \ldots, \alpha_n$ stands to the
functional $\mu$ in a relation similar to that in which the basis
$\omega_1, \omega_2, \ldots, \omega_n$ of the system $\mathfrak{o}$ of all
integers in $\Omega$ stands to the units. In fact the linear form
\weq{10}{
\tau = \omega_1 t_1 + \omega_2 t_2 + \cdots + \omega_n t_n
}
is a unit; for it is the greatest common divisor of
$\omega_1, \omega_2, \ldots, \omega_n$, and must therefore, as the
formula \S.\,145, (8) shows, be a divisor of 1, and hence a unit.

\gesperrt{Accordingly we shall call the functional $\tau$ a basis form of
$\mathfrak{o}$.}


% ===== scan leaf 557, printed 536 =====
% !PRINTED 536
\weberpage{536}
\pdfbookmark[2]{§ 147. The Absolute Norms of the Functionals}{bm168}

From such a basis form one obtains all the integers of the field
$\Omega$ if one puts integral rational numbers for the variables.

The fundamental form $\tau$ is the root of an irreducible equation
of degree $n$
\[
F(t) = N(t - \tau) = 0,
\]
in which the coefficients of the powers of $t$ are integral rational (and
homogeneous) functions of the variables $t_1, t_2, \ldots, t_n$.

\begin{center}
\S. 147.

\gesperrt{The absolute norms of the functionals.}
\end{center}

The elements $\alpha_1, \alpha_2, \ldots, \alpha_n$ of a basis of the
functional $\mu$ can, as integers in $\Omega$, be expressed linearly
and with integral coefficients through a basis $\omega_1, \omega_2,
\ldots, \omega_n$ of $\rho$ in the form
\weq{1}{
(\alpha_1, \alpha_2, \ldots, \alpha_n) = A(\omega_1, \omega_2, \ldots,
\omega_n),
}
where $A$ denotes a linear substitution with integral rational
coefficients. A special case of this is offered by the formulae
(3) of the preceding section.

A basis form $\lambda$ of $\mu$ we shall represent as follows:
\weq{2}{
\lambda = \sum^{\nu} \alpha_\nu t_\nu,
}
and we write the substitution (1) more fully:
\weq{3}{
\alpha_s = \sum^{\nu} a_{s,\nu}\,\omega_\nu,
}
where the $t_\nu$ are variables, the $a_{s,\nu}$ are the substitution
coefficients, and the summation letter $\nu$ runs from 1 to $n$.

Since the products $\alpha_s\,\omega_r$ are all divisible by $\mu$, they
can, according to the meaning of the basis, be represented in the form
\weq{4}{
\alpha_\nu\,\omega_r = \sum^{s} g^{(s)}_{\nu,r}\,\alpha_s,
}
where the $g^{(s)}_{\nu,r}$ are integral rational numbers. From this one
then obtains, by (2),
\weq{5}{
\lambda\,\omega_r = \sum^{s} \alpha_s\,t_{s,r},
}
if
\weq{6}{
t_{s,r} = \sum^{\nu} g^{(s)}_{\nu,r}\,t_\nu
}
are integral rational linear forms.

If one again substitutes the expressions (3) in (5), it follows that
\weq{7}{
\lambda\,\omega_r = \sum^{\nu} \omega_\nu \sum^{s} a_{s,\nu}\,t_{s,r}.
}


% ===== scan leaf 558, printed 537 =====
% !PRINTED 537
\weberpage{537}

If we eliminate the $\omega_\nu$ from these linear equations, we
can represent the equation of degree $n$ for $\lambda$ in determinant
form, and we obtain the product of the roots of this equation,
that is, the norm of $\lambda$, as the determinant formed from
the $n^2$ quantities
\[
\sum^{s} a_{s,\nu}\,t_{s,r}
\]
(\S.\,137, 8.). This determinant, however, can be decomposed
according to the multiplication theorem of determinants
(see Bd. I, \S.\,27), and gives, when
\[
A = \Sigma \pm a_{1,1}a_{2,2}\ldots a_{n,n}, \quad T = \Sigma \pm t_{1,1}t_{2,2}\ldots t_{n,n}
\]
is set,
\weq{8}{N(\lambda) = A\,T.}

Here $T$ is an integral function of the variables $t_\nu$ with integral
rational numerical coefficients, which we shall now further
prove to be primitive.

Let us assume that some natural prime $p$ divides the
divisor of $T$. Then we can determine $n$ integral rational
forms $y_1, y_2, \ldots, y_n$ of the variables $t$ in such a way that the
$n$ sums
\weq{9}{u_s = t_{s,1}y_1 + t_{s,2}y_2 + \cdots + t_{s,n}y_n}
are all divisible by $p$, without all $y_1, y_2, \ldots, y_n$ being
divisible by $p$.

The correctness of this assertion follows from elementary
theorems on determinants (Bd. I, second chapter).

For if the $t_{s,r}$ are all divisible by $p$, we can take
the $y_r$ quite arbitrarily, e.g.\ equal to 1.

Otherwise we assume that, besides the determinant $T$,
all the $m$-rowed subdeterminants are also divisible by $p$
$(m \leqq n)$,
and that among the $(m-1)$-rowed subdeterminants
at least one is not divisible by $p$. If then,
say, among the $(m-1)$-rowed determinants of the matrix
\[
\begin{array}{cccc}
t_{1,1}, & t_{1,2}, & \ldots, & t_{1,m} \\
\multicolumn{4}{c}{\cdots\cdots\cdots\cdots\cdots\cdots\cdots\cdots\cdots\cdots} \\
t_{m-1,1}, & t_{m-1,2}, & \ldots, & t_{m-1,m}
\end{array}
\]
there is one not divisible by $p$, then for $y_1, y_2, \ldots, y_m$ we put
precisely these $(m-1)$-rowed determinants, and take $y_{m+1}, \ldots, y_n$
$= 0$. These $y$ then satisfy the requirement laid down.


% ===== scan leaf 559, printed 538 =====
% !PRINTED 538
\weberpage{538}

If the $y$ are then determined in this way, we set
\weq{10}{\omega = \omega_1 y_1 + \omega_2 y_2 + \cdots + \omega_n y_n,}
and derive from (5) the equation
\weq{11}{\lambda\,\omega = \sum^{s} \alpha_s u_s.}

Now since the $\alpha_s$ are divisible by $\lambda$ and the $u_s$ by $p$,
it follows that $\omega$ is divisible by $p$, and hence, by
\S.\,145, 2., that $y_1, y_2, \ldots, y_n$ must be divisible by $p$, which
is contrary to our assumption.

This proves that $T$ is a primitive function, \gesperrt{and
that therefore the absolute value of the determinant $A$
is equal to the absolute norm of $\lambda$, and consequently also
to that of $\mu$} (\S.\,138).

If we apply the result to the special basis (3), \S.\,146,
we obtain the formula:
\weq{12}{N_a(\mu) = a_{1,1}\,a_{2,2}\ldots a_{n,n}.}

To this we append the following remarks: If
in a basis form of a functional $\mu$
\[
\lambda = \alpha_1 t_1 + \alpha_2 t_2 + \cdots + \alpha_n t_n
\]
one applies to the variables $t_1, t_2, \ldots, t_n$ a linear
substitution with integral coefficients and determinant $\pm 1$, a
new basis form of $\mu$ arises. For to all integral rational values
of the variables $t_1, t_2, \ldots, t_n$ there correspond integral rational
values of the new variables, and conversely.

But the application of a linear substitution to the variables $t$
is equivalent to the application of the transposed
substitution to the coefficients $\alpha_1, \alpha_2, \ldots, \alpha_n$ (\S.\,37), that is, to
the transition to a new basis of $\mu$:
\weq{13}{(\beta_1, \beta_2, \ldots, \beta_n) = (B)\,(\alpha_1, \alpha_2, \ldots, \alpha_n),}
which then, through composition with (1), yields:
\weq{14}{(\beta_1, \beta_2, \ldots, \beta_n) = (B)\,(A)\,(\omega_1, \omega_2, \ldots, \omega_n).}

If the discriminant (\S.\,144) of the numbers $\beta$ determined
by a substitution (13) agrees with the discriminant of the $\alpha$,
then the substitution determinant must be $B = \pm 1$;
and we thus arrive at the theorems:
\begin{enumerate}
\item[1.] \gesperrt{The discriminant of a basis of $\mu$ is equal
to the square of the absolute norm of $\mu$,
multiplied by the fundamental number of the field;}
\end{enumerate}
%CONT>>


% ===== scan leaf 560, printed 539 =====
% !PRINTED 539
\weberpage{539}
\pdfbookmark[2]{§ 148. Complete Residue System with Respect to a Modulus}{bm169}
%<<CONT
\begin{enumerate}
\item[] and conversely:
\item[2.] \gesperrt{If $\beta_1, \beta_2, \ldots, \beta_n$ is a system of integers
divisible by $\mu$ whose discriminant is equal
to the square of the absolute norm of $\mu$,
multiplied by the fundamental number of the field,
then $\beta_1, \beta_2, \ldots, \beta_n$ is a basis of $\mu$.}
\end{enumerate}

\S.\,148.

\gesperrt{Complete system of residues with respect to a modulus.}

\begin{enumerate}
\item[1.] \textbf{Definition:} \gesperrt{Two integral algebraic numbers
$\xi, \eta$ whose difference $\xi - \eta$ is divisible
by an integral functional $\mu$ different from zero
are called congruent to one another with respect
to the modulus $\mu$.}
\end{enumerate}

The concept of congruence can also be extended to functionals,
which, however, we shall not do for the present. On the other hand,
it is essential to take into account, as moduli of the congruence,
not merely the numbers but also the functionals. Every congruence
remains valid if the modulus is replaced by an associated
functional. In the modulus, arbitrary unit factors may be added
or omitted.

For the congruence we use the \gesperrt{Gauss} symbol
(Bd. I, \S.\,115)
\weq{1}{\xi \equiv \eta \ (\mathrm{mod}\ \mu).}

Such a congruence is equivalent in meaning to the equation
\weq{2}{\xi = \eta + \omega\mu,}
in which $\omega$ can be any integral functional.

From this representation one then recognises at once that, just
as in congruences between rational numbers, if in an expression
composed of integers by \gesperrt{addition}, \gesperrt{subtraction} and \gesperrt{multiplication}
every number is replaced by a number congruent to it with respect
to the modulus $\mu$, the result is a number congruent with respect
to the same modulus; in symbols:

If $\xi_1, \eta_1, \xi_2, \eta_2, \ldots$ are integers which satisfy the congruences
\[
\xi_1 \equiv \eta_1, \ \xi_2 \equiv \eta_2, \ldots \ (\mathrm{mod}\ \mu)
\]
%CONT>>


% ===== scan leaf 561, printed 540 =====
% !PRINTED 540
\weberpage{540}
%<<CONT
satisfy, and if $\psi(x_1, x_2, \ldots)$ is an integral function with
integral rational coefficients, then also
\weq{3}{\psi(\xi_1, \xi_2, \ldots) \equiv \psi(\eta_1, \eta_2, \ldots) \ (\mathrm{mod}\ \mu).}

\begin{enumerate}
\item[2.] \gesperrt{If $m$ is the least positive integral rational
number that is divisible by $\mu$, then two integral
rational numbers congruent with respect to the
modulus $\mu$ are also congruent with respect to the
modulus $m$} (\S.\,138, 13.).
\end{enumerate}

A principal theorem on the congruences, to the proof of which we
now proceed, is the following:

\begin{enumerate}
\item[3.] \gesperrt{The number of the integers of a field $\Omega$
that are incongruent with respect to a modulus
$\mu$ is finite, and in fact equal to the absolute
norm of $\mu$.}
\end{enumerate}

To prove it, we take a basis of the functional $\mu$, and in fact
we choose precisely the special basis determined in \S.\,146, (3):
\weq{4}{
\begin{array}{l}
\alpha_1 = a_{1,1}\,\omega_1 \\
\alpha_2 = a_{1,2}\,\omega_1 + a_{2,2}\,\omega_2 \\
\cdots\cdots\cdots\cdots\cdots\cdots\cdots\cdots\cdots\cdots \\
\alpha_n = a_{1,n}\,\omega_1 + a_{2,n}\,\omega_2 + \cdots + a_{n,n}\,\omega_n,
\end{array}
}
where $\omega_1, \omega_2, \ldots, \omega_n$ is a basis of $\mathfrak{o}$.

Every integer $\eta$ in $\Omega$ can then be represented in the
form:
\weq{5}{\eta = y_1\omega_1 + y_2\omega_2 + \cdots + y_n\omega_n,}
where the $y_1, y_2, \ldots, y_n$ are integral rational numbers. With
this we combine the system of numbers
\weq{6}{\xi = x_1\omega_1 + x_2\omega_2 + \cdots + x_n\omega_n,}
in which we let $x_1$ run through a complete system of residues with
respect to the modulus $a_{1,1}$, $x_2$ with respect to the modulus
$a_{2,2} \ldots$, $x_n$ with respect to the modulus $a_{n,n}$. The
number of the numbers $\xi$ so formed amounts to $a_{1,1}a_{2,2}\ldots a_{n,n}$,
a number which by \S.\,147, (12) is equal to the absolute norm
of $\mu$.

We shall show that every number $\eta$ is congruent with respect
to the modulus $\mu$ to one and only one of these numbers $\xi$.

If $\eta - \xi$ is to be divisible by $\mu$, then according to the
meaning of the basis this difference must in the form
\weq{7}{\eta - \xi = h_1\alpha_1 + h_2\alpha_2 + \cdots + h_n\alpha_n}
%CONT>>


% ===== scan leaf 562, printed 541 =====
% !PRINTED 541
\weberpage{541}
\pdfbookmark[2]{§ 149. Congruences}{bm170}
%<<CONT
be representable, in which the $h_1, h_2, \ldots, h_n$ are integral
rational numbers; and when this condition is fulfilled, then conversely
$\eta$ is also congruent to $\xi$ with respect to the modulus $\mu$. If
we order the right-hand side by means of (4) according to $\omega_1,
\omega_2, \ldots, \omega_n$, then from (5) and (6) there results:
\weq{8}{\begin{array}{rclllll}
y_1 &=& x_1 &+ h_1a_{1,1} &+ h_2a_{1,2} &+ \cdots &+ h_na_{1,n} \\
y_2 &=& x_2 & &+ h_2a_{2,2} &+ \cdots &+ h_na_{2,n} \\
\multicolumn{7}{c}{.\quad.\quad.\quad.\quad.\quad.\quad.\quad.\quad.\quad.\quad.\quad.\quad.\quad.\quad.} \\
y_n &=& x_n & & & &+ h_na_{n,n}.
\end{array}}

But in this lies the proof of our assertion. For if $y_1, y_2,
\ldots, y_n$ are any given integral rational numbers, then first of
all, from
\[
y_n = x_n + h_na_{n,n}
\]
$x_n$ and $h_n$ are completely determined. Thereupon we obtain from
the last equation but one
\[
y_{n-1} - h_na_{n-1,n} = x_{n-1} + h_{n-1}a_{n-1,n-1},
\]
whereby $x_{n-1}$ and $h_{n-1}$ are fully determined. The next
preceding equation yields
\[
y_{n-2} - h_na_{n-2,n} - h_{n-1}a_{n-2,n-1} = x_{n-2} + h_{n-2}a_{n-2,n-2},
\]
whereby $x_{n-2}$ and $h_{n-2}$ are fully determined, and we continue
in this way until all the $x_i$ and $h_i$ are determined.

Accordingly we may suitably call the totality of the numbers $\xi$, or
any system of numbers congruent to these, a \gesperrt{complete
system of residues with respect to the modulus $\mu$.}

The number of individuals of a complete system of residues for the
modulus $\mu$ is equal to the absolute norm of $\mu$.

\begin{center}
\S.\,149.

\gesperrt{Congruences.}
\end{center}

From the existence of a complete system of residues with respect to a
modulus $\mu$, which we have proved in the preceding section, there
results a series of theorems that are analogous to well-known
elementary theorems of rational number theory.

If $\alpha$ denotes any integer in $\Omega$ and $\mu$ an integral
functional serving as modulus, and if $\alpha$ is relatively prime to
$\mu$, then two numbers $\alpha\xi$ and $\alpha\xi'$ are only then
congruent with respect to the
%CONT>>


% ===== scan leaf 563, printed 542 =====
% !PRINTED 542
\weberpage{542}
%<<CONT
modulus $\mu$ when the integers $\xi, \xi'$ are congruent. From this
it follows that if $\xi$ runs through a complete system of residues
with respect to the modulus $\mu$, the same also holds of the product
$\alpha\xi$, whereby the theorem is proved:

1. \gesperrt{If $\alpha$ is an integer in $\Omega$, relatively prime
to the modulus $\mu$, and further $\gamma$ is an arbitrary integer in
$\Omega$, then the congruence:}
\weq{1}{\alpha\xi \equiv \gamma \ (\mathrm{mod}\ \mu)}
\gesperrt{is always solvable by an integer $\xi$, and also by only
one, if a complete system of residues with respect to the modulus
$\mu$ is prescribed for $\xi$.}

If in this $\mu$ is itself a number, which we denote by $\beta$, then
the quotient
\[
\frac{\alpha\xi - \gamma}{\beta} = -\eta
\]
is also an integer, and the foregoing theorem then takes the form:

2. \gesperrt{If $\alpha, \beta, \gamma$ are three integers in $\Omega$
and $\alpha, \beta$ are relatively prime, then one can determine two
other integers $\xi, \eta$ in $\Omega$ so that}
\weq{2}{\alpha\xi + \beta\eta = \gamma}
\gesperrt{holds. In particular, one can therefore also, for two
arbitrary relative primes $\alpha, \beta$, satisfy the equation}
\weq{3}{\alpha\xi + \beta\eta = 1}
\gesperrt{by integers $\xi, \eta$.}

This theorem can be carried over to a system of several numbers.

If the numbers $\alpha, \beta, \gamma, \ldots$ in $\mathfrak{o}$ have
no common divisor, then we can first determine numbers $\eta_1, \zeta_1,
\ldots$ in $\mathfrak{o}$ so that
\[
\beta_1 = \beta\eta_1 + \gamma\zeta_1 + \cdots
\]
becomes relatively prime to $\alpha$. For if $\pi_1, \pi_2,
\ldots$ are the distinct prime factors of $\alpha$, no one of which
can divide all of $\beta, \gamma, \ldots$, and if, say, $\beta$ is not
divisible by $\pi_1$, we have to take $\eta_1$ indivisible by $\pi_1$,
the remaining numbers $\zeta_1, \ldots$ divisible by $\pi_1$, and so
forth, and we obtain for each of the numbers $\eta_1, \zeta_1, \ldots$
the condition that it be divisible by some of the prime factors
%CONT>>


% ===== scan leaf 564, printed 543 =====
% !PRINTED 543
\weberpage{543}
%<<CONT
$\pi_i$, but not divisible by others, and this requirement can always
be satisfied according to \S.\,143, 3. Then by 2. we can determine
the integer $\tau$ so that
\[
\alpha\xi + \beta_1\tau = 1
\]
holds, and if we set $\tau\eta_1 = \eta$, $\tau\zeta_1 = \zeta, \ldots$,
we obtain the theorem:

3. \gesperrt{If $\alpha, \beta, \gamma, \ldots$ are numbers in
$\mathfrak{o}$ without a common divisor, then other numbers
$\xi, \eta, \zeta, \ldots$ in $\mathfrak{o}$ can be determined so that}
\weq{4}{\alpha\xi + \beta\eta + \gamma\zeta + \cdots = 1}
\gesperrt{holds.}

From this we can further infer the following theorem:

4. \gesperrt{If $\alpha, \beta, \gamma, \ldots$ are arbitrary numbers
in $\mathfrak{o}$, and $\mu$ is a number in $\mathfrak{o}$ divisible
by the greatest common divisor of all these numbers, then the numbers
$\xi, \eta, \zeta, \ldots$ in $\mathfrak{o}$ can be determined so that}
\weq{5}{\mu = \alpha\xi + \beta\eta + \gamma\zeta + \cdots}
\gesperrt{holds.}

For if, according to \S.\,139, we take the greatest common divisor of
the numbers $\alpha, \beta, \gamma, \ldots$ in the form
\[
\delta = \alpha u + \beta v + \gamma w \ldots
\]
in which $u, v, w, \ldots$ are variables, then by assumption there is
an integral functional $\omega$ such that $\mu = \delta \omega$. We
represent the functional $\omega$ as a quotient of two integral
functions $\varphi : \varepsilon$, of which $\varepsilon$ is a unit,
and consequently $\varphi$ a function with integral coefficients, and
so obtain
\weq{6}{\varepsilon\mu = (\alpha u + \beta v + \gamma w + \cdots)\varphi.}
If one arranges both sides of this equation according to the variables
occurring in them, one obtains, when $\varepsilon_1, \varepsilon_2, \ldots$
are the coefficients in $\varepsilon$, a system of equations of the
following form:
\weq{7}{\begin{aligned}
\varepsilon_1\mu &= \alpha\xi_1 + \beta\eta_1 + \gamma\zeta_1 + \cdots \\
\varepsilon_2\mu &= \alpha\xi_2 + \beta\eta_2 + \gamma\zeta_2 + \cdots \\
& \text{.\,.\,.\,.\,.\,.\,.\,.\,.\,.\,.\,.\,.\,.\,.\,.\,.\,.\,.\,.\,.\,.\,.,}
\end{aligned}}
in which the $\xi_i, \eta_i, \zeta_i, \ldots$ are numbers in
$\mathfrak{o}$. But now the numbers $\varepsilon_1, \varepsilon_2, \ldots$,
as coefficients of a unit, have no common
%CONT>>


% ===== scan leaf 565, printed 544 =====
% !PRINTED 544
\weberpage{544}
\pdfbookmark[2]{§ 150. Fermat's Theorem}{bm171}
%<<CONT
mon divisor, and consequently, by 3., one can determine the numbers
$\tau_1, \tau_2, \ldots$ in $\mathfrak{o}$ so that
\weq{8}{\varepsilon_1\tau_1 + \varepsilon_2\tau_2 + \cdots = 1}
holds. But from (7), if one multiplies by $\tau_1, \tau_2, \ldots$
and adds, and then sets
\[
\xi = \tau_1\xi_1 + \tau_2\xi_2 + \cdots, \quad
\eta = \tau_1\eta_1 + \tau_2\eta_2 + \cdots
\]
there follows, by means of (8), the equation (5) which was to be proved.

\begin{center}
\S.\,150.

\gesperrt{Fermat's Theorem.}
\end{center}

From the theory of congruences, consequences can be drawn which
correspond exactly to the theorems of rational number theory derived
from \gesperrt{Fermat's theorem}; the most important of these must be
discussed here, for the sake of later application.

By $\pi$ we shall understand a prime functional of degree $f$; that
is, if $p$ is the natural prime divisible by $\pi$, we set
\weq{1}{N_a(\pi) = p^f}
(\S.\,140). If now $\alpha$ is any number in $\mathfrak{o}$ not
divisible by $\pi$, then, as we have already seen in the preceding
section, the product $\alpha\xi$ will run through a complete system of
residues with respect to the modulus $\pi$ at the same time as the
number $\xi$. If we leave out the number divisible by $\pi$, there
remain $N_a(\pi) - 1$ numbers, and if we form the product, it follows
that
\weq{2}{\alpha^{p^f-1}\Pi(\xi) \equiv \Pi(\xi) \ (\mathrm{mod}\ \pi),}
where $\Pi(\xi)$ denotes the product of all the numbers of a complete
system of residues (with the exclusion of zero), and is therefore not
divisible by $\pi$. Accordingly, from (2) follows
\weq{3}{\alpha^{p^f-1} \equiv 1 \ (\mathrm{mod}\ \pi),}
or, if one multiplies by $\alpha$,
\weq{4}{\alpha^{p^f} \equiv \alpha \ (\mathrm{mod}\ \pi),}
and in the last form the theorem holds even when $\alpha$ is divisible
by $\pi$.

The formulae (3), (4), which correspond exactly to
\gesperrt{Fermat's theorem} (Bd.\,I, \S.\,136), shall here too be
called \gesperrt{Fermat's theorem}.


% ===== scan leaf 566, printed 545 =====
% !PRINTED 545
\weberpage{545}
We state it thus:

1. \gesperrt{If $\pi$ is a prime divisor of the natural prime $p$, then for every integer $\omega$ in the field $\Omega$}
\[ \omega^{N_a(\pi)} \equiv \omega \pmod{\pi}. \]

Important consequences now attach to this:

2. \gesperrt{If $f(t)$ denotes an integral function of degree $m$ whose coefficients are integers in $\Omega$, and $\pi$ a prime functional, then the congruence}
\weq{5}{f(t) \equiv 0 \pmod{\pi}}
\gesperrt{has at most $m$ roots, that is, there are at most $m$ incongruent integers in $\Omega$ which, put for $t$, satisfy the congruence.}

For if $\alpha$ denotes any number in $\mathfrak{o}$, we can set
\weq{6}{f(t) = (t-\alpha)\,f_1(t) + f(\alpha)}
in which $f_1(t)$ is just such a function as $f(t)$, but only of degree $(m-1)$. But if $f(\alpha) \equiv 0 \pmod{\pi}$, then every root of (5) must satisfy the congruence
\[ (t-\alpha)\,f_1(t) \equiv 0 \pmod{\pi} \]
It must therefore either be congruent with $\alpha$, or be a root of the congruence of degree $(m-1)$
\[ f_1(t) \equiv 0 \pmod{\pi} \]
If we assume our theorem as proved for congruences of degree $(m-1)$, then it accordingly holds also for congruences of degree $m$; and since it holds for congruences of the first degree, it is generally true.

Every number in $\mathfrak{o}$ not divisible by $\pi$ satisfies, as we have seen, the congruence
\weq{7}{\omega^{p^f-1} \equiv 1 \pmod{\pi}.}
Now if $a$ is the least natural number for which the congruence
\weq{8}{\omega^a \equiv 1 \pmod{\pi}}
is satisfied, then it can be shown by the method of inference already often applied that every other exponent $l$ for which $\omega^l \equiv 1$ must be a multiple of $a$. For if $l$ were not divisible by $a$, then, if $a'$ is the residue of the division of $a$ by $l$, we should also have $\omega^{a'} \equiv 1$, which by the assumption about $a$ is not possible. % [sic] the source has ``division of $a$ by $l$''; the minimality argument requires the remainder of $l$ by $a$
Hence $a$ is a divisor of $p^f-1$, and we call $\omega$ a \gesperrt{number belonging to the exponent $a$}.


% ===== scan leaf 567, printed 546 =====
% !PRINTED 546
\weberpage{546}
If $\omega$ belongs to the exponent $a$, then the powers
\weq{9}{1, \omega, \omega^2, \ldots, \omega^{a-1}}
are all incongruent, and therefore, since they all satisfy the congruence (8) (by 2.), they form the totality of the roots of this congruence. Among the numbers (9), therefore, all other numbers $\omega$ belonging to the exponent $a$ must be sought. But $\omega^l$ will belong to the exponent $a$ only when $l$ is relatively prime to $a$, and it follows:

If there are any numbers $\omega$ at all that belong to the exponent $a$, then their number is as great as the number of relative primes to $a$ in the sequence of numbers $0, 1, 2, \ldots, a-1$. We denote this number, as we did earlier (Bd. I, \S.\,132), by $\varphi(a)$. But that to every divisor $a$ of $p^f-1$ there always belongs at least one number $\omega$, and consequently $\varphi(a)$ numbers, can be proved in just the same way as the corresponding theorem of rational number theory is proved in \S.\,136 of the first volume.

The numbers $\omega$ that belong to the exponent $p^f-1$, of which by the foregoing there are always $\varphi(p^f-1)$, are called \gesperrt{primitive roots} of $\pi$. If $\gamma$ is such a primitive root, then the powers
\[ 1, \gamma, \gamma^2, \ldots, \gamma^{p^f-2} \]
form a complete system of residues to the modulus $\pi$, with the exclusion of the number divisible by $\pi$.

By \gesperrt{Fermat}'s theorem for rational numbers, $t^{p-1}-1$ is divisible by $p$, and consequently also by $\pi$, for $t = 1, 2, \ldots, p-1$. The congruence
\weq{10}{t^{p-1} \equiv 1 \pmod{\pi}}
therefore has the roots
\[ 1, 2, \ldots, p-1, \]
and these are incongruent with one another, since by \S.\,140,3. a rational number is divisible by $\pi$ only when it is divisible by $p$. By Theorem 2., therefore, the congruence (10) has no roots other than these.

If we multiply the congruence (10) by $t$ as well, it follows that the congruence of degree $p$
\[ t^p - t \equiv 0 \pmod{\pi} \]
has the $p$ roots
\[ 0, 1, 2, \ldots, p-1, \]
%CONT>>


% ===== scan leaf 568, printed 547 =====
% !PRINTED 547
\weberpage{547}
\pdfbookmark[2]{§ 151. The Dedekind Ideals}{bm172}
%<<CONT
and has no others. In this lies the proof of the following theorem:

3. \gesperrt{A number $\omega$ in $\mathfrak{o}$ is then and only then congruent to a rational number with respect to the modulus $\pi$, when the condition}
\[ \omega^p \equiv \omega \pmod{\pi} \]
\gesperrt{is satisfied.}

If one observes further that the polynomial coefficients in the $p$th power of a polynomial are all divisible by $p$, with the exception of those which belong to the $p$th powers of the individual terms of the polynomial, there results the following further theorem, which relates to integral functions in $\Omega$ of arbitrary variables $x, y, \ldots$:

4. \gesperrt{If $\psi(x,y,\ldots)$ is an integral function of the variables $x, y, \ldots$ with integral coefficients from $\Omega$, then the holding of the congruence}
\[ [\psi(x,y,\ldots)]^p \equiv \psi(x^p, y^p, \ldots) \pmod{\pi} \]
\gesperrt{is the necessary and sufficient condition that all coefficients of $\psi$ are congruent to integral rational numbers with respect to the modulus $\pi$.}

\begin{center}
\S.\,151.\\
\gesperrt{The Dedekind Ideals.}
\end{center}

\gesperrt{Dedekind} founds the theory of algebraic numbers, in the papers already mentioned above, on the concept of the \gesperrt{ideal}.

We now wish to prove that the theory of ideals agrees in essence with the theory of functionals, by showing how the transition from the one to the other can be effected.

The system of all integers of an algebraic number field $\Omega$ shall, as above, be denoted by $\mathfrak{o}$\footnote{Cf. \gesperrt{Dirichlet-Dedekind}, Vorlesungen über Zahlentheorie, in \S.\,167 of the third, \S.\,177 of the fourth edition.}. A system of numbers $\mathfrak{a}$ contained in $\mathfrak{o}$ is called an \gesperrt{ideal} if it satisfies the two requirements:
%CONT>>


% ===== scan leaf 569, printed 548 =====
% !PRINTED 548
\weberpage{548}
%<<CONT
I. \gesperrt{The sum and difference of any two numbers in $\mathfrak{a}$ always yield numbers in $\mathfrak{a}$ again.}

II. \gesperrt{The product of any number in $\mathfrak{a}$ and a number in $\mathfrak{o}$ belongs to the system $\mathfrak{a}$.}

This requirement would be satisfied by the system consisting of the single number zero, which, however, for the sake of simplicity is not designated as an ideal.

The system $\mathfrak{o}$, on the other hand, is a proper ideal. Likewise the system $\mathfrak{o}\mu$ of all numbers in $\mathfrak{o}$ divisible by a definite number $\mu$ is an ideal, and such a one is called a \gesperrt{principal ideal}.

By the product $\mathfrak{a}\mathfrak{b}$ of two ideals $\mathfrak{a}$ and $\mathfrak{b}$ one understands the totality of all numbers which one obtains when one multiplies any number $\alpha$ from $\mathfrak{a}$ by a number $\beta$ from $\mathfrak{b}$ and adds an arbitrary number of such products of numbers, that is, the totality of all numbers of the form $\Sigma \alpha\beta$. That this product $\mathfrak{a}\mathfrak{b}$ is again an ideal is immediately evident. According to this definition, for example, $\mathfrak{o}\mathfrak{a} = \mathfrak{a}$, and the ideal $\mathfrak{o}$ plays the role of the unit in this multiplication.

One can now relate the ideals and the functionals to one another in such a way that the following laws prevail:

1. \gesperrt{To every integral functional there corresponds a definite ideal, and to associated functionals there corresponds the same ideal.}

2. \gesperrt{To every ideal there correspond infinitely many, but only associated, integral functionals.}

3. \gesperrt{To the product of two or more integral functionals there correspond the products of the ideals corresponding to the factors.}

4. \gesperrt{To an integer there corresponds a principal ideal.}

5. \gesperrt{To the units there corresponds the ideal $\mathfrak{o}$.}

To define this correspondence, we first assign to the system of all units the ideal $\mathfrak{o}$. If, further, $\varphi$ is any integral functional which is not a unit, then the totality of all integers $\alpha$ of the field $\Omega$ divisible by $\varphi$ satisfies, by \S.\,138, the requirements I, II, and is therefore an ideal, which we denote by $\mathfrak{a}$\footnote{For the notation of the ideals we use, with \gesperrt{Dedekind}, the small \gesperrt{German letters}.} and assign to the functional $\varphi$
%CONT>>


% ===== scan leaf 570, printed 549 =====
% !PRINTED 549
\weberpage{549}
% running head: \S.\,151. Ideale. 549
%<<CONT
assign. The same ideal $\mathfrak{a}$ is then also assigned to all the
functionals associated with $\varphi$. This assignment has the
properties 1., 4., 5.

If $\varphi$ and $\varphi_1$ are two non-associated functionals, then certainly
one of them, say $\varphi$, is not divisible by the other $\varphi_1$,
and consequently (by \S.\,143, 3.) there are integers that are
divisible by $\varphi$ but not by $\varphi_1$. Consequently, \gesperrt{different}
ideals $\mathfrak{a}$, $\mathfrak{a}_1$ are always assigned to non-associated functionals.

But it must now also be shown that in this way \gesperrt{all}
ideals of the field $\Omega$ can be obtained; in other words, that
every ideal $\mathfrak{a}$ different from $\mathfrak{o}$ consists of the totality
of the numbers divisible by a certain functional factor.

We therefore now start from any ideal $\mathfrak{a}$ and
choose from it an arbitrary finite set of numbers,
$\alpha_1, \alpha_2, \ldots, \alpha_r$, whose greatest common divisor may be
$\delta_r$. This functional $\delta_r$ has a finite number of prime
factors.

Now if there is a number $\alpha_{r+1}$ in $\mathfrak{a}$ that is not divisible
by $\delta_r$, then the greatest common divisor $\delta_{r+1}$ of
$\delta_r$ and $\alpha_{r+1}$ has fewer prime factors than $\delta_r$. If we
continue with this mode of reasoning, we come to the result that
a finite number of numbers $\alpha_1, \alpha_2, \ldots, \alpha_m$ can be chosen
from $\mathfrak{a}$ in such a way that the greatest common divisor $\delta$ of these
numbers \gesperrt{divides all the numbers of $\mathfrak{a}$.}

On the other hand, every number in $\mathfrak{o}$ divisible by $\delta$ belongs to $\mathfrak{a}$.
For by \S.\,149, 4. every number $\alpha$ divisible by $\delta$ can be
put in the form

\[
\alpha = \alpha_1 \xi_1 + \alpha_2 \xi_2 + \cdots + \alpha_m \xi_m,
\]

in which $\xi_1, \xi_2, \ldots, \xi_m$ are numbers in $\mathfrak{o}$, and consequently,
by I. and II., $\alpha$ belongs to the ideal $\mathfrak{a}$.

From this it results that, when $\delta$ is a unit, the
ideal $\mathfrak{a}$ is identical with $\mathfrak{o}$. Every ideal $\mathfrak{a}$ is thus
characterised by the fact that all its numbers have a certain greatest
common divisor.

It remains still to be shown that, when the functionals $\varphi, \psi$
correspond to the two ideals $\mathfrak{a}, \mathfrak{b}$, the product $\varphi\psi$
corresponds to the ideal $\mathfrak{a}\mathfrak{b}$.


% ===== scan leaf 571, printed 550 =====
% !PRINTED 550
\weberpage{550}
% running head: 550 Siebzehnter Abschnitt. \S.\,151.
Since all numbers from $\mathfrak{a}\mathfrak{b}$ are of the form $\Sigma\alpha\beta$, it is first of all
clear that all these numbers are divisible by $\varphi\psi$.

But if, according to \S.\,143, 4., we represent the functional $\varphi$ as the
greatest common divisor of two numbers $\alpha_1, \alpha_2$, then
these numbers, being divisible by $\varphi$, belong to the ideal $\mathfrak{a}$,
and, since a unit factor in $\varphi$ is of no
consequence, we can set

\[
\varphi = \alpha_1 x_1 + \alpha_2 x_2
\]

if $x_1, x_2$ are variables. In the same way, if
$\beta_1, \beta_2$ denote two numbers from $\mathfrak{b}$ and $y_1, y_2$ variables, we can set

\[
\psi = \beta_1 y_1 + \beta_2 y_2
\]

and from this we obtain

\[
\varphi\psi = \alpha_1 \beta_1 x_1 y_1 + \alpha_1 \beta_2 x_1 y_2 + \alpha_2 \beta_1 x_2 y_1 + \alpha_2 \beta_2 x_2 y_2.
\]

Thus by \S.\,143, 2., $\varphi\psi$ is the greatest common
divisor of the four numbers $\alpha_1\beta_1, \alpha_1\beta_2, \alpha_2\beta_1, \alpha_2\beta_2$, which belong to the ideal
$\mathfrak{a}\mathfrak{b}$, and consequently $\varphi\psi$ is the \gesperrt{greatest common
 divisor} of all numbers of the ideal $\mathfrak{a}\mathfrak{b}$.

With this the mutual assignment of the functionals and
ideals has been effected in accordance with requirements 1. to 5.

Under this assignment, \gesperrt{prime ideals} correspond to the
\gesperrt{prime functionals}, and the decomposition of the ideals into prime factors, and
in general the laws of divisibility of the ideals, result
in complete agreement with the corresponding theorems
from the theory of the functionals.

With this complete agreement it can lead to no
inconvenience if we combine the whole system of all functionals
associated with one another into one common concept
 and use for it the name \gesperrt{ideal}.

The system of all functionals associated with an integer
 is then a \gesperrt{principal ideal}.

If any number or functional is divisible by the
functionals of an ideal, we call it \gesperrt{divisible
by the ideal}; and if a functional $\varphi$ is decomposed into factors
to which the ideals $\mathfrak{a}, \mathfrak{b}, \ldots$ correspond, then,
the unit factors being omitted in the notation,
 we also set:

\weq{1}{\varphi = \mathfrak{a}\mathfrak{b} \ldots}

A basis of the functional is at the same time a basis of the ideal,
and the absolute norm of the representing functional agrees
%CONT>>


% ===== scan leaf 572, printed 551 =====
% !PRINTED 551
\weberpage{551}
% running head: \S.\,151. Ideals. 551
%<<CONT
with the number which in \gesperrt{Dedekind} is called the \gesperrt{norm of the ideal}.
It shall therefore be so named here as well, and accordingly,
when the functional $\varphi$ belongs to the ideal $\mathfrak{a}$, we set

\weq{2}{N_a(\varphi) = N(\mathfrak{a}).}

The norm of an ideal is thus always a natural number.

If $\mathfrak{p}$ is a prime ideal and $p$ is the natural prime
divisible by $\mathfrak{p}$, then

\weq{3}{N(\mathfrak{p}) = p^f,}

and $f$ is called the \gesperrt{degree} of the prime ideal $\mathfrak{p}$.

In the congruences (\S.\,149, 150), according to these stipulations,
the corresponding ideals may also be regarded
as moduli in place of the functionals.

Finally, let a word be added on \gesperrt{Kummer's}
creation of the ideal numbers\footnote{Cf. \gesperrt{Dirichlet-Dedekind}, Zahlentheorie, \S.\,176.}.

If $\alpha, \beta$ are any two numbers from $\mathfrak{o}$ which have the greatest
common divisor $\chi$, and if $\beta = \chi\varphi$, then
the congruence

\weq{4}{\alpha\omega \equiv 0 \pmod{\beta}}

will be satisfied only for those numbers $\omega$ from $\mathfrak{o}$ that are
divisible by $\varphi$. But the functional $\varphi$ will in general
not be associated with any number. Now instead of using the functionals,
one can create for oneself a new system of concepts, in which
one combines all the numbers of $\mathfrak{o}$ into pairs
$(\alpha, \beta)$, and can call these pairs ``\gesperrt{ideal numbers}'' in $\mathfrak{o}$;
in exactly the same way in which, by pairing the integral
rational numbers, one has formed the fractions, or, by pairing the real
numbers, the complex numbers.

One then calls the number $\omega$ divisible by the ideal number $(\alpha, \beta)$
if it satisfies the congruence (1). % [sic] printed reference reads \textquotedbl{}(1)\textquotedbl{}; the congruence introduced on this page is numbered (4)
This is the standpoint
of \gesperrt{Kummer's} theory. But in order that this definition
of divisibility be fruitful, it must still be established
under what assumptions one is to regard two ideal numbers with
different elements, $(\alpha, \beta), (\alpha', \beta')$, as equal,
how the actual numbers in $\mathfrak{o}$ are to be ranged among the ideal
numbers, and finally what is to be understood by the product of two ideal
numbers, so that a number divisible by an ideal number
%CONT>>


% ===== scan leaf 573, printed 552 =====
% !PRINTED 552
\weberpage{552}
\pdfbookmark[2]{§ 152. Equivalence}{bm173}
% running head: 552 Siebzehnter Abschnitt. \S.\,152.
%<<CONT
divisible number can be represented as a product of ideal
numbers. All these questions are answered through the mediation
of the theory of functionals (or the theory of ideals), if one
assigns to the ideal number $(\alpha, \beta)$ the functional $\varphi$, and they can
scarcely be decided in any simpler way. The introduction
of the system of ideal numbers then offers no further
essential advantages, and one can just as well regard the
functionals themselves, or the ideals, as elements of the
manifold that is to be newly introduced.

\begin{center}
\S.\,152.

\gesperrt{Equivalence.}
\end{center}

We now call two functionals that differ only by a
unit factor, even when they are fractional,
\gesperrt{associated}, and set up the following definition:

1. \gesperrt{Two integral or fractional functionals $\varphi, \psi$
in the field $\Omega$ are called equivalent if their
quotient $\varphi : \psi$ is associated with a number.}

Thus the two functionals $\varphi$ and $\psi$ are called equivalent
if there exist a unit $\varepsilon$ and a number $\alpha$ in $\Omega$ such that

\weq{1}{\frac{\varphi}{\psi} = \alpha\varepsilon}

holds. If $\varphi$ is equivalent with $\psi$ and with $\psi_1$, it follows from (1) that

\[
\frac{\varphi}{\psi} = \alpha\varepsilon, \qquad \frac{\varphi}{\psi_1} = \alpha_1\varepsilon_1,
\]

consequently

\[
\frac{\psi}{\psi_1} = \frac{\alpha_1}{\alpha} \frac{\varepsilon_1}{\varepsilon},
\]

and since $\alpha_1 : \alpha$ is a number and $\varepsilon_1 : \varepsilon$ a unit, the
first theorem follows:

2. \gesperrt{Two functionals that are equivalent to a third
 are also equivalent to one another
.}

If, in accordance with this, one divides all the functionals of
the field $\Omega$ into \gesperrt{classes}, by casting two functionals into the
same or into different classes according as they are equivalent
or not, it results that two of these classes which have a single
%CONT>>


% ===== scan leaf 574, printed 553 =====
% !PRINTED 553
\weberpage{553}
%<<CONT
contain a common element must be completely identical, and the division into classes is therefore entirely unambiguous. Every class is completely determined by an arbitrary functional contained in it, a \gesperrt{representative}.

3. \gesperrt{Two functionals associated with one another are also equivalent, and therefore occur in the same class.}

For if $\varphi$ and $\psi$ are associated, then their quotient $\varphi:\psi$ is a unit, and $\varphi$ and $\psi$ are therefore also equivalent.

But a class contains not merely the individual functionals $\varphi$, but all the \gesperrt{ideals} determined by these functionals, and when a more precise designation is necessary we therefore call these classes \gesperrt{functional classes} or, still more frequently, in accordance with the usual usage, \gesperrt{ideal classes}.

4. \gesperrt{The totality of all integral and fractional numbers of the field $\Omega$, combined with all the units and the products of numbers with units, form among themselves a class, which is called the principal class.}

As representative of the principal class one may consider, e.g., the number 1. The principal class, conceived as an ideal class, contains the ideal $\mathfrak{o}$ and will therefore in what follows be denoted by the letter $O$.

5. \gesperrt{In every ideal class there are integral functionals.}

For by the definition, if $\varphi$ is any functional and $\alpha$ a number, $\varphi$ is equivalent to $\alpha\varphi$. But by \S.\,137, 5. we can determine the number $\alpha$, even rational, so that $\alpha\varphi$ becomes an integral functional.

6. \gesperrt{From every ideal class $C$ we can select a representative $\varphi$ which is not only itself an integral functional, but is also relatively prime to an arbitrarily given integral functional $\omega$.}

If, in order to prove this theorem, we first take by 5. an arbitrary integral representative $\varphi$ of the class $C$ and an integer $\alpha$ divisible by $\varphi$, then

\weq{2}{\varphi\chi = \alpha,}
%CONT>>


% ===== scan leaf 575, printed 554 =====
% !PRINTED 554
\weberpage{554}
\pdfbookmark[2]{§ 153. The Class Number of the Field Ω}{bm174}
%<<CONT
and $\chi$ an integral functional. We now choose (according to \S.\,143, 3.) a number $\beta$ divisible by $\chi$ in such a way that $\beta:\chi$ becomes relatively prime to $\omega$, and set

\weq{3}{\psi\chi = \beta.}

Since $\varphi:\psi$ is now a number, $\psi$ is equivalent to $\varphi$, and $\psi$ is a representative, coprime to $\omega$, of the class $C^1$).

\S.\,153.

\gesperrt{The class number of the field $\Omega$.}

We come now to the proof of the important theorem:

1. \gesperrt{The number of ideal classes of a field $\Omega$ is finite.}

This theorem is equivalent to the following:

2. \gesperrt{In every class there are integral functionals whose absolute norm does not exceed a certain finite number depending only on the nature of the field $\Omega$.}

For since every integral functional is a factor of its absolute norm, and every integer has only a finite number of non-associated factors, there are, apart from the associated ones, only a finite number of integral functionals that have a given number as absolute norm. If now it can be proved that in every class there occurs an integral functional whose absolute norm under a finite number

\begin{footnotesize}
$^1$) We wish here to point out in passing an analogy of the equivalence of ideals. The whole theory of algebraic numbers can be carried over, with the necessary modifications, to the theory of algebraic functions, which in turn finds its geometric expression in the theory of algebraic curves. To the ideals there then correspond systems of points on a fixed fundamental curve, and to the principal ideals full systems of intersection points of the fundamental curve with another algebraic curve. When two systems of points complete each other to a full system of intersection points, which would find its expression in the above formula $\varphi\chi = \alpha$, each of the two systems of points $\varphi$, $\chi$ is called the \gesperrt{residue} of the other. Two systems of points which, like $\varphi$, $\psi$, have the same residue are called corresidual in geometry. This concept thus corresponds to equivalence. (Cf. \gesperrt{Brill} u. \gesperrt{Nöther}, Mathem. Annalen, Bd. 7. \gesperrt{Salmon}, ``higher plane Curves'', deutsch von \gesperrt{Fiedler}.)
\end{footnotesize}
%CONT>>


% ===== scan leaf 576, printed 555 =====
% !PRINTED 555
\weberpage{555}
%<<CONT
lies, the finiteness of the number of classes is thereby proved.

If, as hitherto, we let $\omega_1, \omega_2, \ldots, \omega_n$ denote a minimal basis of $\Omega$, then all the integers of the field are obtained from

\weq{1}{\omega = \omega_1 x_1 + \omega_2 x_2 + \cdots + \omega_n x_n}

when we assign integral rational numerical values to the $x_1, x_2, \ldots, x_n$. If we assume a positive integer $k$ and stipulate that none of the numbers $x_i$ shall step outside the interval $\pm k$, then the \gesperrt{absolute value} of $\omega$ will not lie above the limit

\[
(|\omega_1| + |\omega_2| + \cdots + |\omega_n|) k = rk
\]

provided that by $|\omega_1|, |\omega_2|, \ldots$ are understood the absolute values (Einleitung, Bd. I, S. 18) of the (real or complex) quantities $\omega_i$, and $r$ denotes the sum $|\omega_1| + |\omega_2| + \cdots$. If we form the expression (1) for the $n$ conjugate fields, and take the product, then, denoting by $R$ a positive real number that lies above the product of the $n$ values $r$, we obtain

\weq{2}{N_a(\omega) < R k^n.}

This number $R$ depends only on the nature of the field $\Omega$, and not on $k$.

Now let $\mu$ be any integral functional in $\Omega$, and $N_a(\mu)$ its absolute norm. If we determine the integer $k$ so that

\weq{3}{k^n \leqq N_a(\mu) < (k+1)^n,}

and if we further assign to the numbers $x_i$ in (1) the values $0, 1, 2, \ldots, k$, then the number of the distinct values that arise from (1) is $(k+1)^n$, and hence greater than $N_a(\mu)$. But by \S.\,148, 3. the number of the numbers incongruent with respect to the modulus $\mu$ is equal to $N_a(\mu)$, and consequently among the numbers $\omega$ so determined there must occur at least two distinct numbers congruent with respect to the modulus $\mu$. If then $\omega' \equiv \omega'' (\mathrm{mod}\ \mu)$, the difference

\weq{4}{\alpha = \omega' - \omega'' = (x_1' - x_1'') \omega_1 + \cdots + (x_n' - x_n'') \omega_n}

will be divisible by $\mu$, and at the same time the integers

\[
x_1' - x_1'' = a_1, \ldots, x_n' - x_n'' = a_n
\]
%CONT>>


% ===== scan leaf 577, printed 556 =====
% !PRINTED 556
\weberpage{556}
%<<CONT
taken absolutely, not greater than $k$. There is a number different from zero and divisible by $\mu$:

\weq{5}{\alpha = a_1 \omega_1 + a_2 \omega_2 + \cdots + a_n \omega_n,}

in which the integer coefficients $a_1, a_2 \ldots a_n$ do not exceed the limits $\pm k$, and consequently by (2) and (3)

\weq{6}{N_a(\alpha) < R k^n < R N_a(\mu).}

Now since $\alpha$ is divisible by $\mu$, we set

\weq{7}{\alpha = \mu\varphi,\ N_a(\alpha) = N_a(\mu) N_a(\varphi),}

and obtain from (6)

\weq{8}{N_a(\varphi) < R.}

If now $\psi$ is a representative of an arbitrarily given class $C$, we choose $\mu$ so that

\[
\beta = \mu\psi
\]

is a number, and if then, in accordance with (7),

\[
\alpha = \mu\varphi
\]

then

\[
\frac{\varphi}{\psi} = \frac{\alpha}{\beta},
\]

and thus $\varphi$ and $\psi$ are equivalent. Hence $\varphi$ is likewise a representative of the class $C$, and it satisfies the condition (8). Thus, as was to be proved, there occurs in every class a functional whose absolute norm lies below $R$.

The number of ideal classes, which we shall denote by $h$, is accordingly a natural number peculiar to the field $\Omega$, which is called the \gesperrt{class number}.

In the simplest case, where the class number is equal to 1, every functional is associated with a number; that is, every (integral or fractional) functional can be transformed into a unit by separating off a numerical factor. In this case every integer of the field $\Omega$ can be decomposed into \gesperrt{prime factors}, and these fields have a theory which agrees in essentials with rational number theory. For such fields the introduction of the functionals and ideals is not necessary.

Here belong, besides the field of the rational numbers, among others the field of the \gesperrt{Gauss} imaginary numbers and the field of the numbers formed from third roots of unity (Bd. I, \S.\,173, 174).


% ===== scan leaf 578, printed 557 =====
% !PRINTED 557
\weberpage{557}
\pdfbookmark[2]{§ 154. The Group of Ideal Classes}{bm175}

\S.\,154.

\begin{center}
\gesperrt{The Group of the Ideal Classes.}
\end{center}

The ideal classes can be compounded on the basis of the following theorem:

1.~\gesperrt{If $\varphi$, $\psi$, $\varphi_1$, $\psi_1$ are functionals, and $\varphi$ is equivalent to $\varphi_1$, $\psi$ equivalent to $\psi_1$, then $\varphi\psi$ is also equivalent to $\varphi_1 \psi_1$.}

The correctness of this follows immediately from the definition. For if $\varphi:\varphi_1$ and $\psi:\psi_1$ are associated with numbers, then $\varphi\psi:\varphi_1\psi_1$ is also associated with a number.

In the same way the converse theorem also follows:

2.~\gesperrt{If $\varphi$ is equivalent to $\varphi_1$ and $\varphi\psi$ equivalent to $\varphi_1\psi_1$, then $\psi$ is also equivalent to $\psi_1$.}

If, then, we consider two ideal classes $A$, $B$, which may also be identical, and form the product $\varphi\psi$ of any functional $\varphi$ from $A$ and a functional $\psi$ from $B$, the class $C$ in which the product $\varphi\psi$ occurs is independent of the choice of $\varphi$ and $\psi$, and the class $C$ is completely determined by the two classes $A$, $B$. \gesperrt{We call $C$ compounded from $A$ and $B$}, and write symbolically

\weq{1}{C = AB = BA.}

The class $C$ contains all products of an element of $A$ with an element of $B$, but it may also contain still other functionals.

Since this composition is derived from true multiplication, the laws of multiplication also hold for this composition, namely the commutative and the associative law.

It follows further from Theorem 2. that if $AB = AB_1$, then $B = B_1$ must also hold, and consequently the ideal classes, under this composition, generate a \gesperrt{finite Abelian group of degree $h$}, to which we can apply all the theorems that we have come to know about such groups in the second chapter of this Band.

The unit of this group is the principal class $O$ already defined in \S.\,152, which indeed, as we have seen, has the represen-
%CONT>>


% ===== scan leaf 579, printed 558 =====
% !PRINTED 558
\weberpage{558}
%<<CONT
tative 1. To define opposite classes, we take a representative $\varphi$ of a class $A$ and a number $\alpha$ divisible by $\varphi$. If now $\alpha = \varphi\chi$, then $\chi$ is a representative of the class $A^{-1}$, and we have $AA^{-1} = O$.

If $A$ is an arbitrary class, and $h$ the class number, then always

\[ A^h = O, \]

and if $k$ is the least positive number that satisfies the condition

\[ A^k = O \]

then $k$ is a divisor of $h$. From this results the theorem:

3.~\gesperrt{Every functional $\varphi$ in $\Omega$ belongs to a definite exponent $k$, which is a divisor of the class number $h$, so that $\varphi^k$ is associated with a number in $\Omega$.}

\begin{center}
\rule{4cm}{0.4pt}
\end{center}


% ===== scan leaf 580, printed 559 =====
% !PRINTED 559
\weberpage{559}
\pdfbookmark[1]{Eighteenth Section. Discriminants.}{bm176}
\pdfbookmark[2]{§ 155. Minimum of a Quadratic Form}{bm177}
\begin{center}
\gesperrt{Eighteenth Chapter.}

\gesperrt{Discriminants.}

\rule{0.3\textwidth}{0.4pt}

\S.\,155.

\gesperrt{Minimum of a Quadratic Form.}
\end{center}

On an entirely new path \gesperrt{Minkowski} has proved the finiteness of the number of classes\footnote{Ueber die positiven quadratischen Formen etc. \gesperrt{Crelle}, Bd.\ 107, 1891. More thoroughly still in the book by \gesperrt{Minkowski}, ``Geometrie der Zahlen'', Leipzig 1896.}, a path which starts from peculiar considerations, capable of numerous applications, and in doing so leads to the proof, long sought in vain, of a theorem on the field discriminant. We shall here communicate briefly the principal points of these investigations.

\gesperrt{Minkowski} starts from the consideration of a quadratic form of $n$ variables

\weq{1}{f(x_1, x_2, \ldots, x_n) = \Sigma^{i,k} a_{i,k}\, x_i x_k}

with real coefficients $a_{i,k}$, of which it is assumed that it can be decomposed into $n$ and no fewer positive squares of linear functions (a positive form, Bd.\ I, \S.\,81 to 83). Let one of these representations be

\weq{2}{f(x_1, x_2, \ldots, x_n) = \xi_1^2 + \xi_2^2 + \cdots + \xi_n^2,}
wherein
\weq{3}{\xi_i = \alpha_{1,i} x_1 + \alpha_{2,i} x_2 + \cdots + \alpha_{n,i} x_n}

is a system of linear functions of the $x_i$ with non-vanishing determinant. If we set

\weq{4}{A = \Sigma \pm \alpha_{1,1}\alpha_{2,2}\ldots\alpha_{n,n},}
%CONT>>


% ===== scan leaf 581, printed 560 =====
% !PRINTED 560
\weberpage{560}
%<<CONT
then the determinant $D$ of the quadratic form (1) (Bd.\,I, \S.\,59) is
\weq{5}{D = \Sigma \pm a_{1,1}a_{2,2}\ldots a_{n,n} = A^2.}

The $\alpha_{i,k}$ are then likewise real numbers, which can be determined in infinitely many ways. We assume an arbitrary, but henceforth unalterable, transformation (2), (3).

Let the system of equations (3), solved for the $x$, yield:
\weq{6}{x_i = \beta_{1,i}\xi_1 + \cdots + \beta_{n,i}\xi_n.}

We shall now assign to the variables $x_i$ only integral rational numerical values. It is then clear that the function $f$ can vanish only when all the $x_i$ vanish, and that its value, if we exclude this case, cannot sink below a certain finite limiting value $M$, which we call the \gesperrt{minimum of the form} $f$. For if the value of $f$ sinks below a value $\alpha^2$, then by (2) all the $\xi_i$ must sink below $\alpha$ in absolute value, and by (6), for a sufficiently small $\alpha$, all the $x_i$ would be pressed down below 1 and, since they are integers, would then all have to be equal to zero.

What matters above all is not so much to determine the minimum exactly, as to find for it as small an upper limit as possible, beyond which it cannot go.

We now avail ourselves, with \gesperrt{Minkowski}, of a mode of expression borrowed from a geometry in space of $n$ dimensions, which leads extraordinarily simply to the desired result. For the cases $n=2$, $n=3$ the expressions there employed are completely intuitive and of the very simplest kind. We recommend the reader to make the steps now to be carried out intuitively clear to himself in these two special cases. For greater values of $n$, to be sure, actual intuition is lost. These expressions are then to be regarded as a short notation for certain analytic relations, the analogy with ordinary space still easing their comprehension. There is not the slightest difficulty in replacing all these expressions by purely analytic formulae; but the brevity of expression and the ease of comprehension would suffer greatly by it.

We now designate the variables $\xi_1, \ldots, \xi_n$ as coor-
%CONT>>


% ===== scan leaf 582, printed 561 =====
% !PRINTED 561
\weberpage{561}
%<<CONT
dinates of a point $\xi$ in a space of $n$ dimensions
$R_\xi$, and call the expression

\weq{7}{\varrho = \sqrt{(\xi_1' - \xi_1'')^2 + (\xi_2' - \xi_2'')^2 + \cdots + (\xi_n' - \xi_n'')^2},}

the \gesperrt{distance} of the two points $\xi'$, $\xi''$.

If in (3) we let the variables $x_i$ run through all integer
values, the associated values of the $\xi_i$ give us a system of
points in the space $R_\xi$. The point that corresponds to the
values $x_i = 0$ is the initial point of the coordinates, or the
\gesperrt{origin}. The totality of these points remains unchanged if we
replace $x_i$ by $x_i + x_i^{(0)}$, where $x_1^{(0)}, x_2^{(0)},
\ldots, x_n^{(0)}$ denotes a fixed system of integers. We can also
express this fact by saying that the system of points of the
$\xi_i$ comes to cover itself when it undergoes a parallel
displacement by which the origin arrives at the place of an
arbitrarily given other point of the system; or again: The points
of the system are situated with respect to each of its points in
the same way as with respect to every other.

We shall call this system a \gesperrt{point lattice} and its
individual points the \gesperrt{lattice points}.

By (2) and (7) the value of $f$ is the square of the distance of
the point $\xi$ from the origin, and consequently $\sqrt{M}$ is
the \gesperrt{smallest distance between two lattice points} that occurs in
the whole system.

Besides the distances in the space $R_n$ we have also to consider
volumes, i.\,e. the values of the $n$-fold integral

\weq{8}{V_\xi = \int\!\!\int \ldots \int d\xi_1\, d\xi_2 \ldots d\xi_n,}

in which the variables $\xi_1, \xi_2, \ldots, \xi_n$ fill some
finite region $G_\xi$, which is determined, say, by one or more
inequalities. Thus the region bounded by the condition

\weq{9}{\xi_1^2 + \xi_2^2 + \cdots + \xi_n^2 \leqq 1 \text{ or also } \Sigma(\xi_i - \xi_i^0)^2 \leqq 1}

for arbitrary values $\xi_i^0$ is an \gesperrt{$n$-dimensional sphere} of
radius 1, whose volume we shall denote by $K_n$. From every
region $G_\xi$ we can derive a \gesperrt{similar region} $G_\xi^{(t)}$ by
enlarging the coordinates all in the ratio $t^{\frac{1}{n}} : 1$.
If $V^{(t)}$ is the volume of the new region, then by (8)

\weq{10}{V^{(t)} = t\,V.}


% ===== scan leaf 583, printed 562 =====
% !PRINTED 562
\weberpage{562}
If, besides the space $R_\xi$, we consider also a second space $R_x$,
in which the variables $x$ determined by (3) or (6) are the
coordinates, then to every point of the one space there corresponds
one and only one point of the other space. To a finite % [sic: the page prints ``endliche'' for ``endlichen''; the inflection error has no English counterpart]
region $G_\xi$ there corresponds a finite region $G_x$.

In order to ascertain the relation between the volumes $V_\xi$ and $V_x$,
we must treat the integral (8) according to the rules for the
transformation of multiple integrals, for which the general formula holds:

\weq{11}{\int\!\int\!\ldots\!\int d\xi_1\, d\xi_2 \ldots d\xi_n
= \int\!\int\!\cdots\!\int
\sqrt{\left(\Sigma\pm\frac{d\xi_1}{dx_1}\frac{d\xi_2}{dx_2}\cdots\frac{d\xi_n}{dx_n}\right)^{2}}\,
dx_1\, dx_2 \cdots dx_n\ \footnotemark.}

\footnotetext{The general formulation of this transformation is due to
\gesperrt{Jacobi} (De determinantibus functionalibus; \gesperrt{Crelle},
Bd. 22, 1891, % [sic: Weber prints 1891.]
gesammelte Werke, Bd. 3). The derivation is to be found
in most of the more detailed textbooks of the integral calculus, e.g.
\gesperrt{Serret-Harnack}, Bd. 2; \gesperrt{Lipschitz}, Lehrbuch der
Analysis Bd. 2; also \gesperrt{Baltzer}, Determinanten.}

If in this we insert for the functional determinant the value (4) or
(5), there results

\weq{12}{V_\xi = \sqrt{D}\ V_x,}

in which now the integral

\weq{13}{V_x = \int\!\int\!\cdots\!\int d x_1\, d x_2 \cdots d x_n}

refers to the region $G_x$.

The integral $V_x$ we may conceive as the limiting value of a sum of
elements $\Delta x_1 \Delta x_2 \ldots \Delta x_n$, and
if we give to all the increments $\Delta x_i$ the quantity $1 : t^{\frac1n}$,
then each of these elements will have the value $1 : t$.
If $Z_t$ denotes the number of these elements, then $V_x$ is the
limiting value which the ratio $Z_t : t$ approaches as $t$
increases indefinitely; thus

\weq{14}{V_x = \mathrm{Lim}_{t\,=\,\infty} \frac{Z_t}{t}.}

Now the number of elements $Z_t$ is just as great as the number of
points in the interior of $G_x$ whose coordinates have the
expression

\[
x_1 = h_1 t^{-\frac1n},\ x_2 = h_2 t^{-\frac1n},\ \ldots,\ x_n = h_n t^{-\frac1n}
\]

with integer $h_1, h_2, \ldots, h_n$, or again equal to the
%CONT>>


% ===== scan leaf 584, printed 563 =====
% !PRINTED 563
\weberpage{563}
%<<CONT
number of points with integral coordinates that lie in a region
$G_x^{(t)}$ enlarged in the ratio $t^{\frac1n} : 1$, and this
number is just as great as the number of lattice points that lie
in the region $G_\xi^{(t)}$. We thus have the following theorem:

\gesperrt{If $G_\xi$ is any finite region in the space $R_\xi$
and $G_\xi^{(t)}$ the region derived from $G_\xi$ by
enlargement in the linear ratio $t^{\frac1n} : 1$; if $Z_t$
is the number of lattice points lying in the interior of
$G_\xi^{(t)}$, then the volume determined by $G_\xi$ is}

\weq{15}{V_\xi = \sqrt{D}\ \mathrm{Lim}_{t\,=\,\infty} \frac{Z_t}{t}.}

We now place about every lattice point, as centre, an
$n$-dimensional sphere of radius $\frac12\sqrt{M}$. Since $\sqrt{M}$
is the least distance between any two lattice points, no two of
these spheres can have a region in common; they can at most touch
one another in a point, and none of these spheres contains,
besides its centre, a further lattice point. Such a sphere
has, by (10), the volume

\[
K_n \left(\frac{\sqrt{M}}{2}\right)^{n},
\]

and the total volume of all those spheres whose centre lies in the
interior of $G_\xi^{(t)}$ is

\[
V_1^{(t)} = Z_t\, K_n \left(\frac{\sqrt{M}}{2}\right)^{n}.
\]

But this volume $V_1^{(t)}$ is evidently smaller than $V^{(t)}$ —
at least if we disregard those spheres which project partly out of
the region $G^{(t)}$, and whose total volume, divided by $t$,
converges to zero as $t$ increases indefinitely — because its
region fills only a part of $G^{(t)}$; and consequently $V$ too is
greater than the limiting value of $V_1^{(t)} : t$, thus

\[
V > \mathrm{Lim}\, \frac{Z_t\, K_n}{t} \left(\frac{\sqrt{M}}{2}\right)^{n},
\]

and from this, by (15),

\weq{16}{\sqrt{D} > K_n \left(\frac{\sqrt{M}}{2}\right)^{n}.}

Since it is now easy to determine the volume $K_n$, the problem is
hereby solved, and an upper limit for $M$ found.


% ===== scan leaf 585, printed 564 =====
% !PRINTED 564
\weberpage{564}
We shall not insert here the exact value for $K_n$, but rather a
smaller value that can be expressed more simply, whereby the
inequality (11) becomes all the more correct.

By the definition we have

\weq{17}{K_n = \int\!\int\!\cdots\!\int d\xi_1\, d\xi_2 \ldots d\xi_n,}

with the boundary

\weq{18}{\xi_1^2 + \xi_2^2 + \cdots + \xi_n^2 \leqq 1.}

The boundary equation (18) is certainly satisfied if we restrict
the $\xi_i$ to the intervals

\weq{19}{-\frac{1}{\sqrt{n}} < \xi_i < \frac{1}{\sqrt{n}}}

which geometrically has the meaning that we replace the sphere by
the inscribed cube. Accordingly we have

\[
K_n > \int\limits_{-\frac{1}{\sqrt n}}^{\frac{1}{\sqrt n}}
\int\limits_{-\frac{1}{\sqrt n}}^{\frac{1}{\sqrt n}}
\cdots
\int\limits_{-\frac{1}{\sqrt n}}^{\frac{1}{\sqrt n}}
d\xi_1\, d\xi_2 \cdots d\xi_n,
\]

and when this integration is carried out,

\weq{20}{K_n > \left(\frac{2}{\sqrt{n}}\right)^{n}.}

The upper limit determined in this way is never attained by $K_n$,
if we disregard the case $n = 1$.

The comparison of (20) with (16) yields a very simple expression
for the upper limiting value of $M$, namely

\weq{21}{M < n \sqrt[n]{D},}

and $M$ is (apart from the case $n = 1$) actually smaller than, and
not equal to, the value $n \sqrt[n]{D}$.

This shows itself clearly if one seeks a still somewhat narrower
limit for $M$.

The condition (18) is certainly then fulfilled if we set

\weq{22}{-\alpha < \xi_n < \alpha,\qquad \xi_1^2 + \xi_2^2 + \cdots + \xi_{n-1}^2 < 1 - \alpha^2}

where $\alpha$ may be an arbitrary positive proper fraction.
(The geometric meaning of this is that one replaces the sphere
by an inscribed cylinder of height $2\alpha$). Thus we have

\[
\begin{gathered}
K_n > 2\alpha \int\!\int\!\cdots\!\int d\xi_1\, d\xi_2 \ldots d\xi_{n-1}\\
\xi_1^2 + \xi_2^2 + \cdots + \xi_{n-1}^2 < 1 - \alpha^2.
\end{gathered}
\]


% ===== scan leaf 586, printed 565 =====
% !PRINTED 565
\weberpage{565}
If in the $(n-1)$-fold integral one replaces $\xi_i$ by
$\sqrt{1-\alpha^2}\,\xi_i$, one finds for it the value
$\sqrt{(1-\alpha^2)}^{\,n-1} K_{n-1}$, and thus one finds the
inequality

\[
K_n > 2\alpha \sqrt{(1-\alpha^2)}^{\,n-1} K_{n-1}.
\]

Now we wish to choose $\alpha$ so that $2\alpha
\sqrt{(1-\alpha^2)}^{\,n-1}$ becomes as great as possible, and for
this one finds, by the ordinary rule, by setting the differential
quotient equal to zero,

\[
\alpha = \frac{1}{\sqrt{n}}.
\]

From this

\[
K_n > 2\, \frac{\sqrt{(n-1)}^{\,n-1}}{\sqrt{n}^{\,n}}\, K_{n-1}
\]

\[
\left(\frac{\sqrt{n}}{2}\right)^{n} K_n > \left(\frac{\sqrt{n-1}}{2}\right)^{n-1} K_{n-1}.
\]

If one therefore sets

\[
K_n = \left(\frac{2}{\sqrt{\Theta_n\, n}}\right)^{n},
\]

then

\weq{23}{\Theta_n^n < \Theta_{n-1}^{n-1},}

hence $\Theta_n^n$ is a number that decreases as $n$ increases. At the
same time there results from (16)

\weq{24}{M < \Theta_n\, n\, \sqrt[n]{D}.}

For $n = 2$, $K_n$ is equal to the area of the circle with radius 1,
that is, to the \gesperrt{Ludolphian} number $\pi$, and from this
there results

\[
\Theta_2 = \frac{2}{\pi} < 0{,}636,\qquad \Theta_2^2 < 0{,}404,
\]

% [sic: Weber prints references (22) and (23); the bound uses (24) and (23).]
and hence, by (22) and (23), the following determination of the limit
holds, valid for all values of $n$ greater than 1:

\weq{25}{M < n \sqrt[n]{0{,}404\, D}\ \footnotemark.}

\footnotetext{For special applications it is useful to know the limit
as exactly as possible. For this purpose serves an exact computation
of the integral $K_n$ according to its definition (17), (18), which
can easily be carried out by various methods of the integral
calculus, e.g. by means of the discontinuous factor of
\gesperrt{Dirichlet}:

\[
K_n = \frac{\Gamma\left(\frac12\right)^{n}}{\Gamma\left(1+\frac{n}{2}\right)},
\]
}


% ===== scan leaf 587, printed 566 =====
% !PRINTED 566
\weberpage{566}
\pdfbookmark[2]{§ 156. Application to Algebraic Fields}{bm178}

\begin{center}
\S.\,156.

\gesperrt{Application to algebraic fields.}
\end{center}

From the minimum for a positive quadratic form one can
derive a minimum of linear forms, by conceiving a positive
quadratic form as a sum of squares of linear forms. On this
rest the following applications to the theory of algebraic
fields, which we preface with a lemma:

1. \gesperrt{If $a_1, a_2, \ldots, a_n$ are real positive numerical values, then}
\weq{1}{a_1 + a_2 + \cdots + a_n \geqq n \sqrt[n]{a_1 a_2 \ldots a_n}.}

This theorem is evidently correct for $n=2$, for
$a_1 + a_2 - 2\sqrt{a_1 a_2} = (\sqrt{a_1} - \sqrt{a_2})^2$, and thus never negative. We
therefore apply complete induction, by assuming formula
(1) as already proved for $n-1$ terms. If we then assume
that $a_n$ is exceeded in magnitude by none of the remaining
$a$, then
\[ a_n \geqq \sqrt[n-1]{a_1 a_2 \ldots a_{n-1}} = b, \]
and
\[ a_1 + a_2 + \cdots + a_{n-1} \geqq (n-1)\,b. \]
Hence
\[ a_1 + a_2 + \cdots + a_n \geqq a_n + (n-1)\,b, \]
and consequently
\[ \frac{a_1 + a_2 + \cdots + a_n}{n} \geqq b + \frac{a_n - b}{n}, \]

\begin{footnotesize}
\noindent\rule{4cm}{0.4pt}

in which $\Gamma(x)$ denotes the gamma function, for which the equations
\[ \Gamma(x) = (x-1)\,\Gamma(x-1), \quad \Gamma\left(\frac{1}{2}\right) = \sqrt{\pi} \]
hold, so that $K_n$ is expressed by rational numbers and by
$\sqrt{\pi}$.

(Cf. e.\,g. Vorlesungen \"uber die Theorie der bestimmten Integrale
nach \gesperrt{Lejeune-Dirichlet} von G.\,F. \gesperrt{Meyer} herausgegeben, \S.\,175.
Leipzig 1871).

From this there results
\[ \Theta_n = \frac{4}{n\pi}\,\Gamma\left(1 + \frac{n}{2}\right)^{\frac{2}{n}}, \]
which, by the approximation formulae for the $\Gamma$-function, for large
values of $n$ agrees nearly with
\[ \frac{2}{\pi e} = 0{,}234\ldots \]
(\gesperrt{Minkowski}, Crelle, Bd. 107, S. 292).
\end{footnotesize}
%CONT>>


% ===== scan leaf 588, printed 567 =====
% !PRINTED 567
\weberpage{567}
%<<CONT

and by the binomial theorem, since $a_n - b$ is not negative,
\[ \left(\frac{a_1 + a_2 + \cdots + a_n}{n}\right)^n \geqq b^n + b^{n-1}(a_n - b) = a_n b^{n-1}, \]
but this coincides with the inequality (1) that is to be proved.

Now let $\omega_1, \omega_2, \ldots, \omega_n$ be a minimal basis of an algebraic
field $\Omega$. Further, let $\varphi$ be an arbitrary integral
functional.

We assume a basis of $\varphi$ (\S.\,146):
\weq{2}{(\beta_1, \beta_2, \ldots, \beta_n) = \begin{pmatrix} b_{1,1} & \ldots & b_{1,n} \\ \cdot & \cdots & \cdot \\ b_{n,1} & \ldots & b_{n,n} \end{pmatrix} (\omega_1, \omega_2, \ldots, \omega_n),}
so that
\weq{3}{N_a(\varphi) = \pm \Sigma \pm b_{1,1} \ldots b_{n,n}}
holds. From the linear form associated with $\varphi$
\weq{4}{\eta = \beta_1 x_1 + \beta_2 x_2 + \cdots + \beta_n x_n}
we then obtain all the integers of the field $\Omega$ that are
divisible by $\varphi$, when all possible combinations of integers
are put for the $x_1, x_2, \ldots, x_n$.

From (4) there arise $n$ different forms $\eta_1, \eta_2, \ldots, \eta_n$
when we pass over to the $n$ conjugate fields, $\Omega_1, \Omega_2, \ldots, \Omega_n$.
The square of the determinant of this system of
$n$ forms is, by \S.\,147, equal to the discriminant
\weq{5}{\Delta(\beta_1, \beta_2, \ldots, \beta_n) = \pm N_a(\varphi)^2 \Delta,}
if $\Delta$ denotes the fundamental number of the field.

The conjugate fields $\Omega_1, \Omega_2 \ldots$ may be partly real, partly
conjugate imaginary in pairs. As already several times before, we
denote by $|\omega|$ the \gesperrt{absolute value} of a real
or complex quantity, and, understanding by
$x_1, x_2, \ldots, x_n$ real variables, we consider the quadratic form
\weq{6}{f(x_1, x_2, \ldots, x_n) = |\eta_1|^2 + |\eta_2|^2 + \cdots + |\eta_n|^2 = S\,|\eta|^2.}
If $\Omega$ is real, then
\weq{7}{|\eta|^2 = \eta^2,}
hence the square of a real linear form. But if $\Omega, \Omega'$
are two conjugate imaginary ones among the conjugate fields, then
the corresponding linear forms $\eta, \eta'$ are also conjugate
imaginary, and we have
\weq{8}{|\eta|^2 = |\eta'|^2 = \eta \eta'}
%CONT>>


% ===== scan leaf 589, printed 568 =====
% !PRINTED 568
\weberpage{568}
%<<CONT
the sum of two squares of real linear forms. The product
$\eta \eta'$ then appears with the factor 2 in the function $f$.
Hence $f$ is in all circumstances a positive form.

But by (7), (8), $f$ can also be regarded as a quadratic form
of the variables $\eta_1, \eta_2, \ldots, \eta_n$,
\weq{9}{f = F(\eta_1, \eta_2, \ldots, \eta_n) = \sum^{i,k} a_{i,k}\,\eta_i \eta_k}
and in it, if $\eta_i$ is real, $a_{i,i} = 1$, and if
$\eta_i, \eta_k$ is a conjugate imaginary pair, $a_{i,k} = 1$. All the remaining
coefficients are $= 0$. The determinant of $F$ is therefore $= \pm 1$,
because in every row and in every column only one element
different from zero, with the value 1, occurs. By Bd. I,
\S.\,56 the determinant of the form $f$ is equal to the product of the
determinant of $F$ and the square of the determinant of the
linear forms $\eta$; hence by (5)
\weq{10}{= \pm [N_a(\varphi)]^2 \Delta,}
where the sign is to be chosen so that the value comes out positive.

By \S.\,155, (24) one can now put for the $x_i$ such integral
rational numbers that
\weq{11}{f(x_1, x_2, \ldots, x_n) < \Theta_n\, n \sqrt[n]{\pm \Delta N_a(\varphi)^2}}
holds, where $\Theta_n$, if $n$ is greater than 1, is a proper fraction and
always less than $\sqrt[n]{0{,}404}$. If one puts these values of $x_i$ into
(4), then $\eta$ passes over into an integer $\beta$ divisible by $\varphi$,
and with this we therefore have the theorem:

2. \gesperrt{If $\varphi$ is any integral functional in $\Omega$, then there}
\gesperrt{is an integer $\beta$ in $\Omega$ divisible by $\varphi$ with}
\gesperrt{the property that}
\weq{12}{S\,|\eta|^2 < \Theta_n\, n \sqrt[n]{\pm \Delta [N_a(\varphi)]^2}}
\gesperrt{holds, where $\Delta$ is the fundamental number of the field and}
\gesperrt{$\pm \Delta$ is positive.}

The product of the conjugate values $\eta$ is equal to the norm
of $\eta$, and hence by Lemma 1. [with regard to
(7) and (8)]
\weq{13}{n \sqrt[n]{[N(\eta)]^2} \leqq S\,|\eta|^2,}
and so by (12) there follows the inequality relating to the
absolute values:
\weq{14}{N(\eta)^2 < \Theta_n^n\, [N_a(\varphi)]^2\, \Delta.}
Since now $\eta$ is divisible by $\varphi$, we can
\[ \varepsilon \eta = \varphi \psi, \quad \pm N(\eta) = N_a(\varphi)\, N_a(\psi) \]
%CONT>>


% ===== scan leaf 590, printed 569 =====
% !PRINTED 569
\weberpage{569}
\pdfbookmark[2]{§ 157. Prime Factors of the Natural Prime Numbers}{bm179}
%<<CONT
set, wherein $\psi$ is an integral functional and $\varepsilon$ a unit, and
obtain from (14)
\weq{15}{N_a(\psi) < \sqrt{\pm 0{,}404\, \Delta}.}

Since we may set for $\varphi$ any integral functional, $\psi$ can be
a representative of an ideal of an arbitrary class, and
we have thus proved the important theorem:

3. \gesperrt{In every ideal class we can find a repre-}
\gesperrt{sentative whose norm is less than the}
\gesperrt{square root of the absolute value of the}
\gesperrt{fundamental number of the field.}

The upper limit for the norm of an ideal of a class that was
used in \S.\,153 for the proof of the finiteness of the class number
is hereby determined in a much sharper manner.

But the inequality (15) admits yet another very remarkable
conclusion. For since $N_a(\psi)$ is a positive integer, and thus
at least equal to 1, it follows from (15)
that, apart from the case $n=1$, $\Delta$ must be absolutely greater than
1, and \gesperrt{that therefore, besides the field of the rational
numbers, there is no field whose fundamental number would be $\pm 1$.} This
is the proof of this important theorem, earlier long sought
in vain, first supplied by \gesperrt{Minkowski}\footnote{In a letter to \gesperrt{Hermite} (Comptes rendus d. Pariser Akademie,
26. Januar 1891), and also in the book mentioned above, \gesperrt{Min-}
\gesperrt{kowski} finds a far narrower limit, in that he considers not only the minimum
of the quadratic forms, but that of still other functions as well,
such as $\eta_1^p + \eta_2^p + \cdots + \eta_n^p$, wherein $p$ is an arbitrary positive number, which
does not even need to be integral. For $p=1$ there results, if $n-\nu$
is the number of conjugate imaginary pairs:
\[ N_a(\psi) < \left(\frac{4}{\pi}\right)^{n-\nu} \frac{1.2.3\ldots n}{n^n}\, \sqrt{\pm \Delta}. \]
% [sic: Weber prints ``keiner'' for ``kleiner'' (less than).]
From this still further conclusions may be drawn, such as these: the fundamental
 number of a quadratic field must be greater than 4 or less than
$-2$; the fundamental number of a cubic field must be greater than 20
or less than $-12$, etc.}.

\begin{center}
\S.\,157.

\gesperrt{Prime factors of the natural primes.}
\end{center}

A more exact investigation of the basis form $\tau$ of $\mathfrak{o}$ of the
 field $\Omega$, explained in \S.\,146, shall now give us the means,
%CONT>>


% ===== scan leaf 591, printed 570 =====
% !PRINTED 570
\weberpage{570}
%<<CONT
actually to decompose any arbitrarily given natural prime $p$ into its
prime factors, and thereby actually to represent all prime ideals of
the field $\Omega$, or at least representatives of all the
prime ideals\footnote{We follow here a paper by \gesperrt{Hensel}: ``Untersuchung der
Fundamentalgleichung einer Gattung für eine reelle Primzahl als Modul
und Bestimmung der Theiler der Discriminante'' (\gesperrt{Crelle}'s Journ., Bd. 113).
Also to be mentioned are the other papers by \gesperrt{Hensel}, Ibid.,
Bd. 101, 104, 105, 113.}.

Let $\mathfrak p$ be an arbitrary prime ideal of degree $f$, so that
\weq{1}{N(\mathfrak p) = p^f}
where $p$ denotes a natural prime which is divisible by the
prime factor $\mathfrak p$. $f$ is a positive exponent
$\leqq n$. The least integral rational number that is divisible by
$\mathfrak p$ is $p$, and every other integral rational number that is
divisible by $\mathfrak p$ is therefore also divisible by $p$
(\S.\,138).

We must now also consider congruences with the modulus $\mathfrak p$
between integral functions of arbitrary variables with coefficients in
$\mathfrak o$, and we remark that by \S.\,143, 2. two such functions
are \gesperrt{congruent then and only then, when the corresponding
coefficients are congruent.}

We denote the field of rational numbers by $R$, and accordingly we
also call a function with rational coefficients a \gesperrt{function
in} $R$.

The basis form of $\mathfrak o$
\weq{2}{\tau = \omega_1 t_1 + \omega_2 t_2 + \cdots + \omega_n t_n}
satisfies, as we have seen in \S.\,146, an equation of degree
$n$, $F(\tau) = 0$, whose coefficients are integral
rational functions of $t_1, t_2, \ldots, t_n$. Thus $F(\tau)$ is in
any case divisible by $\mathfrak p$, and from this it follows that
there are integral functions $\Phi(t)$ in $R$, which besides $t$
may contain any other variables, and which by the substitution
$t = \tau$ pass over into functionals in $\mathfrak o$ divisible by
$\mathfrak p$, and which therefore satisfy the congruence
\weq{3}{\Phi(\tau) \equiv 0\ (\mathrm{mod}\ \mathfrak p)}
and so we shall be able to say that $\tau$ \gesperrt{is a root of the
congruence}
\weq{4}{\Phi(t) \equiv 0\ (\mathrm{mod}\ \mathfrak p).}

The function $\Phi$ will certainly have to contain the variables
$t_1, t_2, \ldots, t_n$; but it can also contain still other
%CONT>>


% ===== scan leaf 592, printed 571 =====
% !PRINTED 571
\weberpage{571}
%<<CONT
and if we thus denote the variables of $\Phi$ by $t, u_1, u_2,
\ldots$, we shall also set
\weq{5}{\Phi(t) = \Phi(t, u_1, u_2, \ldots),}
or more briefly $\Phi(t, u)$, and this function contains integral
rational numerical coefficients. If we raise the function $\Phi(t)$
to the $p$th power and apply the formula \S.\,150, 4., there follows
from the congruence (3) a new congruence
\weq{6}{\Phi(\tau^p, u_1^p, u_2^p, \ldots) \equiv 0\ (\mathrm{mod}\ \mathfrak p).}

But in the same way we also have
\weq{7}{\tau^p \equiv \omega_1^p t_1^p + \omega_2^p t_2^p + \cdots + \omega_n^p t_n^p.}

If we substitute this in the congruence (6), there arises a
function divisible by $\mathfrak p$, in which the variables
$u_1, u_2, \ldots$, among which indeed the $t_1, t_2, \ldots, t_n$
are included, occur only in the $p$th power. The congruence must
therefore remain correct if we replace the $u_1^p, u_2^p,
\ldots$ and hence also the $t_1^p, t_2^p, \ldots, t_n^p$ by
independent variables $u_1, u_2, \ldots t_1, t_2, \ldots, t_n$
(\S.\,143).

If we accordingly set
\weq{8}{\tau_1 = \omega_1^p t_1 + \omega_2^p t_2 + \cdots + \omega_n^p t_n,}
then from (6) there results
\[
\Phi(\tau_1, u_1, u_2, \ldots) \equiv 0\ (\mathrm{mod}\ \mathfrak p),
\]
and consequently $\tau_1$ too is a root of the congruence (4).

This same procedure can be applied repeatedly, and we find
that
\[
\tau_2 = \omega_1^{p^2} t_1 + \omega_2^{p^2} t_2 + \cdots + \omega_p^{p^2} t_n % [sic: Weber prints \omega_p here; (9) has \omega_n.]
\]
is also a root of the congruence (4), and so on.

If, then, we
define a system of forms $\tau_r$ by
\weq{9}{\tau_r = \omega_1^{p^r} t_1 + \omega_2^{p^r} t_2 + \cdots + \omega_n^{p^r} t_n}
for arbitrary positive exponents $r$, all these quantities
$\tau_r$ are at the same time roots of the congruence (4). It
remains to answer the question how many of these forms $\tau_r$
are different from one another.

By \S.\,150, (4) it is certain that
\[
\tau_r \equiv \tau_{r'}\ (\mathrm{mod}\ \mathfrak p),
\]
if
\[
r \equiv r'\ (\mathrm{mod}\ f)
\]
holds, and consequently among the $\tau_r$ there are certainly not
more than $f$ that are different with respect to the modulus $\mathfrak p$
\weq{10}{\tau, \tau_1, \tau_2, \ldots, \tau_{f-1}.}


% ===== scan leaf 593, printed 572 =====
% !PRINTED 572
\weberpage{572}
That these forms are, however, really different from one another
follows from the fact that (by \S.\,145, 150) we can put for the
variables $t_1, t_2, \ldots, t_n$ integral rational numbers such
that $\tau$ passes over into a primitive root $\gamma$ of the prime
ideal $\mathfrak p$. By the same substitution the quantities (10)
become
\weq{11}{\equiv \gamma, \gamma^p, \gamma^{p^2}, \ldots, \gamma^{p^{f-1}}\ (\mathrm{mod}\ \mathfrak p),}
which are all different from one another with respect to the modulus
$\mathfrak p$. Hence, too, no two of the forms (10) can be congruent
with respect to the modulus $\mathfrak p$, since otherwise the two
corresponding numbers (11) would also turn out congruent.

We sum this up as a theorem thus:

1. \gesperrt{Every congruence (4), of which one root is $t = \tau$,
has the $f$ different roots}
\[
\tau, \tau_1, \tau_2, \ldots, \tau_{f-1}.
\]

Accordingly, dividing $\Phi(t)$ algebraically by $t - \tau$, we can
set
\[
\Phi(t) \equiv (t - \tau)\, \Phi_1(t)\ (\mathrm{mod}\ \mathfrak p)
\]
and $\tau_1, \tau_2, \ldots, \tau_{f-1}$ are roots of
$\Phi_1(t) \equiv 0$, in which, however, $\Phi_1(t)$ does not yet
have rational coefficients, but coefficients in $\mathfrak o$. If we
divide $\Phi_1(t)$ again by $t - \tau_1$, and continue in this way,
it follows finally that, if we now set the function of degree $f$
\weq{12}{\Pi(t) = (t - \tau)(t - \tau_1) \ldots (t - \tau_{f-1})}
then
\weq{13}{\Phi(t) \equiv \Pi(t)\, \Phi_0(t)\ (\mathrm{mod}\ \mathfrak p),}
in which $\Phi_0(t)$ is an integral function with coefficients in
$\mathfrak o$.

The function $\Pi(t)$ depends on the variables $t, t_1, \ldots, t_n$,
and to express this we set
\[
\Pi(t) = \Pi(t, t_1, \ldots, t_n).
\]

The coefficients of this form are numbers in $\mathfrak o$, and it
can further be shown that they are congruent with integral rational
numbers with respect to the modulus $\mathfrak p$. This proof is
obtained by raising to the $p$th power:
\[
[\Pi(t)]^p \equiv (t^p - \tau^p)(t^p - \tau_1^p) \ldots (t^p - \tau_{f-1}^p).
\]

Now, however, by (9)
\[
\tau_r^p \equiv \omega_1^{p^{r+1}} t_1^p + \omega_2^{p^{r+1}} t_2^p + \cdots + \omega_n^{p^{r+1}} t_n^p\ (\mathrm{mod}\ \mathfrak p),
\]
and so we obtain $\tau_r^p$ from $\tau_{r+1}$ when we replace $t_i$
by $t_i^p$,
%CONT>>


% ===== scan leaf 594, printed 573 =====
% !PRINTED 573
\weberpage{573}
%<<CONT
and $\tau_f$ is congruent with $\tau$. Accordingly we obtain the
congruence:
\[
[\Pi(t, t_1, t_2, \ldots, t_n)]^p \equiv \Pi(t^p, t_1^p, t_2^p, \ldots, t_n^p)\ (\mathrm{mod}\ \mathfrak p).
\]
By \S.\,150, 4. the theorem is thus proved:

2. \gesperrt{The form}
\[
\Pi(t) = (t - \tau)(t - \tau_1) \ldots (t - \tau_{f-1})
\]
\gesperrt{is congruent, with respect to the modulus $\mathfrak p$, to
an integral and homogeneous form of degree $f$ in $R$ in the variables
$t, t_1, t_2, \ldots, t_n$.}

This integral rational form, whose coefficients are completely
determined up to multiples of the prime $p$, we shall denote by $P(t)$.

Then from (13) follows
\weq{14}{\Phi(t) \equiv P(t)\, \Phi_0(t)\ (\mathrm{mod}\ \mathfrak p),}
and since now $\Phi(t)$ is an integral function in $R$, raising to the
power $p$ yields:
\[
\begin{aligned}
\Phi(t^p, u^p) &\equiv P(t^p, u^p)\ \Phi_0(t^p, u^p) \\
[\Phi(t, u)]^p &\equiv [P(t, u)]^p\ [\Phi_0(t, u)]^p
\end{aligned}
\qquad (\mathrm{mod}\ \mathfrak p).
\]
But because the left-hand sides are congruent with respect to the
modulus $\mathfrak p$, it follows that
\[
[P(t)]^p\ [\Phi_0(t, u)^p - \Phi_0(t^p, u^p)] \equiv 0\ (\mathrm{mod}\ \mathfrak p),
\]
and from this, since $P(t)$ is not divisible by $\mathfrak p$ (because
the coefficient of $t^f$ has the value 1):
\weq{15}{[\Phi_0(t, u)]^p \equiv \Phi_0(t^p, u^p),}
from which it follows by \S.\,150, 4. that $\Phi_0(t)$ too is
congruent to an integral rational form with respect to the modulus
$\mathfrak p$.

Accordingly we may in (14) also take $\Phi_0(t)$ to be an integral
form in $R$, and then by \S.\,140, 3. the congruence (14) must hold
not only for the modulus $\mathfrak p$, but for the modulus $p$.
From this we obtain the theorem:

3. \gesperrt{Among the forms $\Phi(t)$ which for $t = \tau$ pass over
into a functional divisible by $\mathfrak p$, the form $P(t)$, of
degree $f$ with respect to $t$, is the form of lowest degree, and if
$\Phi(t)$ is any one among them, then an integral function $\Phi_0(t)$
in $R$ can be so determined that}
\weq{15}{\Phi(t) \equiv P(t)\, \Phi_0(t)\ (\mathrm{mod}\ p)} % [sic] number (15) reused for a second equation on this page, as printed
\gesperrt{holds.}


% ===== scan leaf 595, printed 574 =====
% !PRINTED 574
\weberpage{574}
If we omit from $\Phi(t)$ and $\Phi_0(t)$ the terms whose coefficients are divisible by $p$, then the degree of $\Phi(t)$ with respect to $t$ is greater by $f$ than the degree of $\Phi_0(t)$.

$P(t)$ is an integral function in $R$ of the variables $t, t_1, t_2, \ldots, t_n$, which by the substitution $t=\tau$ goes over into a functional divisible by $\mathfrak p$, which of course is no longer contained in $R$.

The significance of this form $P(t)$ now emerges still more clearly when we prove the theorem:

4. \gesperrt{The prime factor $\mathfrak p$ is the greatest common divisor of $p$ and $P(\tau)$.}

For this we have only to show that, if $p$ is divisible by $\mathfrak p \mathfrak p_1$, where $\mathfrak p, \mathfrak p_1$ are two equal or different prime factors, then $P(\tau)$ is divisible by $\mathfrak p$, but not by $\mathfrak p \mathfrak p_1$.

By \S.\,143, 3. there always exists a \gesperrt{number} $\xi$ in $\mathfrak{o}$ which is divisible by $\mathfrak p$, but not by $\mathfrak p \mathfrak p_1$. This number will have the form
\weq{16}{\xi = a_1 \omega_1 + a_2 \omega_2 + \cdots + a_n \omega_n,}
wherein the $a_1, a_2, \ldots, a_n$ are integral rational numbers, and it therefore arises from the form $\tau$ by the substitution
\weq{17}{(t_1, t_2, \ldots, t_n) = (a_1, a_2, \ldots, a_n).}
If we make the same substitution in the forms $\tau_r$ defined in (9), then $\tau_r$ goes over into a number $\xi_r$ which, by \gesperrt{Fermat}'s theorem, satisfies the congruence
\[ \xi^{p^r} \equiv \xi_r \pmod{\mathfrak p} \]
and is thus certainly also divisible by $\mathfrak p$. If, therefore, in the form
\[ \Pi(t) = (t-\tau)(t-\tau_1) \cdots (t - \tau_{f-1}) \]
we put $\tau_r + \xi_r$ in place of $\tau_r$, then this form remains congruent to itself with respect to the modulus $\mathfrak p$. But this substitution comes to the same thing as substituting $t_1+a_1, t_2+a_2, \ldots, t_n+a_n$ in place of $t_1, t_2, \ldots, t_n$, and we accordingly obtain
\[ \Pi(t, t_1+a_1, \ldots, t_n+a_n) \equiv \Pi(t, t_1, \ldots, t_n) \pmod{\mathfrak p}, \]
% [sic: Weber prints ``nnd'' for ``und'' (and).]
and if we replace $\Pi$ by the congruent form $P$
\[ P(t, t_1+a_1, \ldots, t_n+a_n) \equiv P(t, t_1, \ldots, t_n) \pmod{\mathfrak p}. \]


% ===== scan leaf 596, printed 575 =====
% !PRINTED 575
\weberpage{575}
But since the form $P$ has none but \gesperrt{rational} coefficients, the latter congruence must also hold to the modulus $p$, and thus
\weq{18}{P(t, t_1+a_1, \ldots, t_n+a_n) \equiv P(t, t_1, \ldots, t_n) \pmod{p}.}
This congruence holds for variable $t, t_1, \ldots, t_n$.

Let us now assume that, contrary to the theorem to be proved, $P(\tau)$ is divisible by $\mathfrak p \mathfrak p_1$; then the congruence
\weq{19}{P(\tau, t_1, \ldots, t_n) \equiv 0 \pmod{\mathfrak p \mathfrak p_1},}
holds, and this remains correct when for the independent variables $t_1, t_2, \ldots, t_n$ the substitution $t_1+a_1, t_2+a_2, \ldots, t_n+a_n$ is made; and since $\tau$ thereby passes into $\tau+\xi$, it follows that
\weq{20}{P(\tau+\xi, t_1+a_1, \ldots, t_n+a_n) \equiv 0 \pmod{\mathfrak p \mathfrak p_1}.}
If, on the other hand, we make the substitution $t=\tau+\xi$ in the congruence (18), it follows that
\weq{21}{P(\tau+\xi, t_1, \ldots, t_n) = P(\tau+\xi) \equiv 0 \pmod{\mathfrak p \mathfrak p_1}.}
Let us now first assume that $\mathfrak p_1$ is different from $\mathfrak p$; then we can reason as follows.

We set $t_1=t_2=\cdots=t_n=0$ in (21), and hence also $\tau=0$. But by this $P(t)$ passes into $t^f$ (apart from multiples of $p$), and from (21) we obtain
\[ \xi^f \equiv 0 \pmod{\mathfrak p \mathfrak p_1}, \]
which, however, contradicts the assumption that $\xi$ is not to be divisible by $\mathfrak p_1$.

But if $\mathfrak p_1 = \mathfrak p$, we arrange (21) according to powers of $\xi$ and obtain
\weq{22}{P(\tau+\xi) = P(\tau) + \xi P'(\tau) + \frac{\xi^2}{2} P''(\tau) \ldots,}
where $P'(t), P''(t), \ldots$ are the derivatives of $P(t)$, and where it is to be observed that the forms
% [sic: Weber prints $P'''(\tau)$ here although the sentence refers to derivatives of $P(t)$.]
\[ \frac{1}{2} P''(t), \ \frac{1}{2.3} P'''(\tau) \ldots \]
are, in spite of the apparent denominators, integral forms in $R$. Now since by (19) and (21) $P(\tau+\xi)$ and $P(\tau)$ are divisible by $\mathfrak p^2$, and likewise, by assumption, $\xi^2, \xi^3, \ldots$, while $\xi$ is not divisible by $\mathfrak p^2$, it follows from (22) that
\[ P'(\tau) \equiv 0 \pmod{\mathfrak p}, \]
from which, on account of
\[ P(t) \equiv \Pi(t) \pmod{\mathfrak p} \]
%CONT>>


% ===== scan leaf 597, printed 576 =====
% !PRINTED 576
\weberpage{576}
%<<CONT
according to the meaning (12) of $\Pi(t)$, it follows that:
\[ \Pi'(\tau) \equiv (\tau-\tau_1)(\tau-\tau_2) \cdots (\tau-\tau_{f-1}) \equiv 0 \pmod{\mathfrak p}. \]
But this is impossible, because the $\tau, \tau_1, \ldots \tau_{f-1}$, as we have seen, are incongruent with respect to the modulus $\mathfrak p$. With this, then, our Theorem 4. is completely proved. We can accordingly, if $x, y$ denote two new variables, form a functional $\pi$ of the ideal $\mathfrak p$:
\weq{23}{\pi = x p + y P(\tau).}
We now consider integral forms $\Phi(t)$ in $R$ which by the substitution $t=\tau$ become divisible not only by $\mathfrak p$, but by the natural prime $p$, and which therefore satisfy the congruence
\weq{24}{\Phi(\tau) \equiv 0 \pmod p}
Let the prime $p$ be decomposed into its prime factors in $\Omega$ in the following manner:
\weq{25}{p = \mathfrak p \mathfrak p_1 \mathfrak p_2 \cdots,}
where
\[ \mathfrak p, \mathfrak p_1, \mathfrak p_2, \ldots \]
are equal or different prime ideals of the degrees
\[ f, f_1, f_2, \ldots \]
To these prime factors there corresponds (by Theorem 3.) a sequence of integral rational forms
\[ P(t), P_1(t), P_2(t), \ldots \]
of the degrees $f, f_1, f_2, \ldots$, and if, say, $\mathfrak p$ is identical with $\mathfrak p_1$, then $P(t)$ is also identical with $P_1(t)$. If one takes the norm on the right and on the left in (25), and takes account of the formula $N(\mathfrak p) = p^f$ [\S.\,140, (3), \S.\,151, (3)], it follows that
\weq{26}{n = f + f_1 + f_2 + \cdots}
Now if $\Phi(t)$ is an integral rational form satisfying the condition (24), it follows from Theorem 3. that
\weq{27}{\Phi(t) \equiv P(t) \Phi_1(t) \pmod p,}
where $\Phi_1(t)$ is an integral function in $R$ which satisfies the condition
\[ P(\tau) \Phi_1(\tau) \equiv 0 \pmod{\mathfrak p \mathfrak p_1 \mathfrak p_2 \ldots} \]
By Theorem 4. it follows from this that
\weq{28}{\Phi_1(\tau) \equiv 0 \pmod{\mathfrak p_1 \mathfrak p_2 \ldots},}
%CONT>>


% ===== scan leaf 598, printed 577 =====
% !PRINTED 577
\weberpage{577}
%<<CONT
and from this one concludes again by Theorem 3. (applied to $\mathfrak p_1$)
\[ \Phi_1(t) \equiv P_1(t) \Phi_2(t) \pmod p. \]
In this the same consideration can be repeated which leads to the congruence
\[ \Phi_2(t) \equiv 0 \pmod{\mathfrak p_2 \ldots} \]
from which, again by 3.,
\[ \Phi_2(t) \equiv P_2(t) \Phi_3(t) \pmod p \]
is to be concluded. If one continues in this way until all the prime factors of $p$ have been taken into account, the following theorem results:

5. \gesperrt{If $\Phi(t)$ is an integral function in $R$ which becomes divisible by $p$ by the substitution $t=\tau$, then another integral function $\Phi_0(t)$ in $R$ can be determined so that}
\weq{29}{\Phi(t) \equiv \Phi_0(t) P(t) P_1(t) P_2(t) \ldots \pmod p}
\gesperrt{holds.}

\gesperrt{The product $P(t) P_1(t) P_2(t) \ldots$ is of degree $n = f+f_1+f_2+\cdots$ and is the integral rational form of lowest degree which becomes divisible by $p$ by the substitution $t=\tau$.}

Among the functions $\Phi(t)$ occurring in Theorem 5. there is also the integral function of degree $n$
\[ F(t) = N(t-\tau), \]
which vanishes for $t=\tau$. For this function $\Phi_0(t)$ becomes independent of $t$, and since both in $F(t)$ and in $P(t), P_1(t), P_2(t) \ldots$ the highest power of $t$ has the coefficient 1, it follows that $\Phi_0(t) = 1$. Hence we have the congruence
\weq{30}{N(t-\tau) \equiv P(t) P_1(t) P_2(t) \ldots \pmod p.}
If in the decomposition (25) of the prime $p$ we collect the equal prime factors into powers, then, if $\mathfrak p_1, \mathfrak p_2, \ldots$ are different prime ideals of the degrees $f_1, f_2, \ldots$, we can take the positive exponents $e_1, e_2, \ldots$ in such a way that
\weq{31}{p = \mathfrak p_1^{e_1} \mathfrak p_2^{e_2} \ldots,}
and
\weq{32}{n = e_1 f_1 + e_2 f_2 + \cdots}
hold. Formula (30) then assumes the shape
\weq{33}{N(t-\tau) \equiv [P_1(t)]^{e_1} [P_2(t)]^{e_2} \ldots \pmod p.}


% ===== scan leaf 599, printed 578 =====
% !PRINTED 578
\weberpage{578}
\pdfbookmark[2]{§ 158. Dedekind's Theorem on the Field Discriminant}{bm180}

\begin{center}
\S.\,158.

\gesperrt{Dedekind's theorem on the field discriminant.}
\end{center}

The natural primes which divide the discriminant $\Delta$ of the field $\Omega$ have, with regard to their decomposition into prime factors, a special character, concerning which a theorem of \gesperrt{Dedekind} gives the most concise information; to its derivation we now proceed\footnote{\gesperrt{Dedekind}, ``Ueber die Discriminanten endlicher K\"orper'' im XXIX. Bande der Abhandlungen der Gesellschaft der Wissenschaften in G\"ottingen (1882).}.

Following \S.\,145, 2., we form the powers of the basis form
\[
\tau = t_1\omega_1 + t_2\omega_2 + \cdots + t_n\omega_n,
\]
and obtain the expressions:
\weq{1}{\tau^k = u_{1,k}\,\omega_1 + u_{2,k}\,\omega_2 + \cdots + u_{n,k}\,\omega_n,}
in which the $u_{i,k}$ are integral rational functions of the variables $t_1, t_2, \ldots, t_n$. As above, we set
\weq{2}{F(t) = N(t-\tau),}
so that $F(\tau)=0$. Now, following \S.\,144, we form the discriminant
\weq{3}{\Delta(1,\tau,\tau^2,\ldots,\tau^{n-1}) = (-1)^{\frac{n(n-1)}{2}}\, N F'(\tau),}
for which, according to (1), the value $\Delta U^2$ results, where $\Delta$ is the field discriminant and
\weq{4}{U = \left|\begin{array}{cccc} u_{1,0}, & u_{2,0}, & \cdots, & u_{n,0} \\ u_{1,1}, & u_{2,1}, & \cdots, & u_{n,1} \\ \cdots & \cdots & \cdots & \cdots \\ u_{1,n-1}, & u_{2,n-1}, & \cdots, & u_{n,n-1} \end{array}\right|,}
and is thus an integral function in $R$.

We must show that the form $U$ is a unit. Suppose that some prime $p$ divides $U$; then, as was already proved in \S.\,147, we can determine a system of integral functions $y_0, y_1, \ldots, y_{n-1}$ in $R$, \gesperrt{not all divisible by $p$}, in such a way that for $i=1,2,\ldots,n$
\[
y_0\, u_{i,0} + y_1\, u_{i,1} + \cdots + y_{n-1}\, u_{i,n-1} \equiv 0\ (\mathrm{mod}\ p)
\]
holds. But then from (1) it follows that
\[
y_0 + y_1\tau + \cdots + y_{n-1}\tau^{n-1} \equiv 0\ (\mathrm{mod}\ p).
\]
%CONT>>


% ===== scan leaf 600, printed 579 =====
% !PRINTED 579
\weberpage{579}

%<<CONT
But this contradicts Theorem 5. of the preceding section, according to which $\tau$ can satisfy, with respect to a prime modulus $p$, no congruence of degree lower than $n$.

Accordingly it follows from (3) \gesperrt{that the fundamental number $\Delta$ of the field, apart from the sign, is the absolute norm of the function $F'(\tau)$}, and that therefore
\weq{5}{\pm \Delta = N_a[F'(\tau)]}
is to be set.

If now, instead of the $\omega_i$, another basis $\omega_i'$ of $\mathfrak{o}$ is taken as the foundation, then in place of $\tau$ there appears another form
\weq{6}{\tau' = \omega_1' t_1 + \omega_2' t_2 + \cdots + \omega_n' t_n,}
for which we may also set
\weq{7}{\tau' = \omega_1 t_1' + \omega_2 t_2' + \cdots + \omega_n t_n',}
and in these the $\omega_i'$ are connected with the $\omega_i$ by a linear substitution with the determinant $\pm 1$ (\S.\,145):
\weq{8}{(\omega_1', \omega_2', \ldots, \omega_n') = C(\omega_1, \omega_2, \ldots, \omega_n).}
If one carries out this substitution in (6), and arranges according to $\omega_1, \omega_2, \ldots, \omega_n$, there results
\weq{9}{(t_1', t_2', \ldots, t_n') = C_1(t_1, t_2, \ldots, t_n),}
where $C_1$ is the transposed substitution of $C$ (\S.\,37). If one now forms the function $F(t)$ for the function $\tau'$, there may result $F_1(t)$; let the derivative for $t=\tau'$ be $F_1'(\tau')$. If we now set
\[
F'(\tau) = \Psi(t_1, t_2, \ldots, t_n),
\]
then $\Psi$ is an integral function in $R$, and on account of (7) there results
\[
F_1'(\tau') = \Psi(t_1', t_2', \ldots, t_n').
\]

% [sic: The source prints $F'(\tau')$ as the first functional; the preceding definitions compare $F'(\tau)=\Psi(t_1,\ldots,t_n)$ with $F_1'(\tau')=\Psi(t_1',\ldots,t_n')$. The primed argument is retained as printed.]
Since the substitution (9) is reversible, it follows from this that the two functionals $F'(\tau')$ and $F_1'(\tau')$ are mutually divisible by one another (\S.\,143), and that they are consequently associated.

The function $F'(\tau)$, whose absolute norm is equal to the absolute value of the fundamental number, we therefore call the \gesperrt{fundamental functional}, and the ideal represented by $F'(\tau)$ the \gesperrt{fundamental ideal} of the field $\Omega$. Now let $\mathfrak{p}^e$ be the highest power of the prime ideal $\mathfrak{p}$ which divides $p$, and let $e \geqq 1$; then by 5. of the preceding section we can
\weq{10}{F(t) \equiv P(t)^e\, \Phi(t)\ (\mathrm{mod}\ p)}
set, where $\Phi(t)$ is an integral function in $R$ of such a nature that $\Phi(\tau)$ is not divisible by $\mathfrak{p}$. From (10), however,
%CONT>>


% ===== scan leaf 601, printed 580 =====
% !PRINTED 580
\weberpage{580}

%<<CONT
follows, when we form the derivative with respect to $t$, which is evidently permissible:
\[
F'(t) \equiv e\,[P(t)]^{e-1}\,P'(t)\,\Phi(t) + [P(t)]^e\,\Phi'(t)\ (\mathrm{mod}\ p).
\]
If we set $t=\tau$ in this, then $P'(t)$ passes over into a form $P'(\tau)$ not divisible by $\mathfrak{p}$ (which we have already shown and used in the proof of \S.\,157, 4.), and we obtain
% [sic: Weber prints $F_1(\tau)$ here; the preceding factorization and following sentence use $\Phi(\tau)$.]
\weq{11}{F'(\tau) \equiv e\, P(\tau)^{e-1}\, P'(\tau)\, F_1(\tau)\ (\mathrm{mod}\ \mathfrak{p}^e).}
Now (by \S.\,157, 4.) $P(\tau)$ is divisible by $\mathfrak{p}$ but not by $\mathfrak{p}^2$, while $P'(\tau)$ and $\Phi(\tau)$ are not divisible by $\mathfrak{p}$, and so the formula (11) thus gives us the proof of the following theorem:

1. \gesperrt{If $\mathfrak{p}$ is an arbitrary prime ideal, $p$ the natural prime divisible by $\mathfrak{p}$, and $\mathfrak{p}^e$ the highest power of $\mathfrak{p}$ that divides $p$, then the fundamental form $F'(\tau)$ is in every case divisible by $\mathfrak{p}^{e-1}$; if, further, the exponent $e$ is not divisible by $p$, then $F'(\tau)$ is not divisible by $\mathfrak{p}^e$; but if $e$ is divisible by $p$, then $F'(\tau)$ is divisible by $\mathfrak{p}^e$ and perhaps by still higher powers of $\mathfrak{p}$.}

If $e$ is greater than 1, then by the foregoing $F'(\tau)$ is divisible by $\mathfrak{p}$, and consequently $N_a[F'(\tau)]$, and hence also the fundamental number $\Delta$, is divisible by $p$.

But if all prime ideals $\mathfrak{p}$ divide $p$ only in the first power, then $F'(\tau)$ is relatively prime to $p$. If one then decomposes $F'(\tau)$ into its prime factors, there occurs among them none whose norm is a power of $p$, and consequently $N_a[F'(\tau)]$ and $\Delta$ are likewise not divisible by $p$. From this there then follows the theorem:

2. \gesperrt{A natural prime $p$ is then and only then divisible in the field $\Omega$ by the square of a prime factor, when $p$ divides the fundamental number of $\Omega$.}

From these considerations we can draw still other important conclusions.

If we solve the system of linear equations (1) with respect to $\omega_1, \omega_2, \ldots, \omega_n$, we obtain expressions with the denominator $U$, which, as we have seen, is a unit. If one substitutes these expressions in
\weq{12}{\omega = x_1\omega_1 + x_2\omega_2 + \cdots + x_n\omega_n,}
%CONT>>


% ===== scan leaf 602, printed 581 =====
% !PRINTED 581
\weberpage{581}

%<<CONT
in which the $x_1, x_2, \ldots, x_n$ are integral rational numbers or integral rational functionals, one obtains an expression of the form
\weq{13}{\omega = A_0 + A_1\tau + A_2\tau^2 + \cdots + A_{n-1}\tau^{n-1},}
in which the $A_0, A_1, \ldots, A_{n-1}$ likewise denote \gesperrt{integral rational functionals}. By \S.\,145, 2., however, every integer and every integral functional of the field $\Omega$ is represented in this form, and we can therefore state the theorem:

3. \gesperrt{The powers}
\[
1, \tau, \tau^2, \ldots, \tau^{n-1}
\]
\gesperrt{form a basis of the integral functionals of the field $\Omega$.}

Instead of the powers of $\tau$, the functions
\[
F_0, F_1, F_2, \ldots, F_{n-1}
\]
can also be introduced as elements of the basis. These are defined by the equation
\weq{14}{\frac{F(t)}{t-\tau} = t^{n-1}F_0(\tau) + t^{n-2}F_1(\tau) + \cdots + F_{n-1}(\tau)}
and, as was already shown in \S.\,4 of the first Band, if
\[
F(t) = t^n + a_1t^{n-1} + \cdots + a_n
\]
they have the expression:
\weq{15}{\begin{aligned}
F_0(t) &= 1 \\
F_1(t) &= t + a_1 \\
F_2(t) &= t^2 + a_1 t + a_2 \\
&\cdots\cdots\cdots\cdots\cdots\cdots\cdots\cdots\cdots\cdots \\
F_{n-1}(t) &= t^{n-1} + a_1 t^{n-2} + a_2 t^{n-3} + \cdots + a_{n-1},
\end{aligned}}
so that the powers $1, t, t^2, \ldots, t^{n-1}$ can be represented without denominator by $F_0, F_1, \ldots, F_{n-1}$, and so that every integer or every integral functional $\omega$ also receives the expression:
\weq{16}{\omega = B_0 F_0(\tau) + B_1 F_1(\tau) + \cdots + B_{n-1} F_{n-1}(\tau),}
whose coefficients $B_0, B_1, \ldots, B_{n-1}$ are \gesperrt{integral rational functionals}.

If we now consider the integral function of $t$:
\[
S\,\frac{F(t)}{(t-\tau)\,F'(\tau)},
\]
% [sic: The German source prints a lowered comma after "so" at the line break; the German transcription retains it, while the English uses natural sentence punctuation.]
in which $S$ is the symbol for the trace explained in \S.\,134, it follows that it receives the value 1 for $t=\tau_1,\tau_2,\ldots,\tau_n$
%CONT>>


% ===== scan leaf 603, printed 582 =====
% !PRINTED 582
\weberpage{582}
%<<CONT
and that, since it is only of degree $(n-1)$, it must therefore be identically $=1$.
From this it follows by (14):
\[
1 = t^{n-1} S\,\frac{F_0(\tau)}{F'(\tau)} + t^{n-2} S\,\frac{F_1(\tau)}{F'(\tau)} + \cdots + S\,\frac{F_{n-1}(\tau)}{F'(\tau)}
\]
or
\weq{17}{\begin{aligned} S\,\frac{F_\nu(\tau)}{F'(\tau)} &= 0, \quad S\,\frac{F_{n-1}(\tau)}{F'(\tau)} = 1,\\ \nu &= 0, 1, \ldots, n-2. \end{aligned}}
From this, with the use of (16), we obtain:
\weq{18}{S\,\frac{\omega}{F'(\tau)} = B_{n-1},}
or the theorem:

4. \gesperrt{If $\omega$ is an integer or an integral functional of the field $\Omega$, then the trace of $\dfrac{\omega}{F'(\tau)}$ is an integral rational functional.}

\begin{center}
\rule{0.2\linewidth}{0.4pt}
\end{center}


% ===== scan leaf 604, printed 583 =====
\weberpage{583}% page number not printed on this page
\pdfbookmark[1]{Nineteenth Section. Relation of a Field to Its Divisors.}{bm181}
\pdfbookmark[2]{§ 159. Relative Norms}{bm182}
% !PRINTED none
\begin{center}
\gesperrt{Nineteenth Chapter.}

\gesperrt{Relations of a Field to its Divisors.}

\rule{0.2\linewidth}{0.4pt}

\S.\,159.

\gesperrt{Relative Norms.}
\end{center}

In this chapter we shall consider the modifications that occur in the definitions and theorems of the theory of algebraic numbers when, in place of the \gesperrt{absolute domain of rationality}, that is, of the field of the rational numbers, an arbitrary algebraic number field is substituted. In this connection some important supplements will also result for the general theory of prime factors, which find their expression in a series of theorems established by \gesperrt{Dedekind} and \gesperrt{Hilbert}${}^{1}$).

Let $R$, then, be an algebraic number field of degree $m$, and let
\weq{1}{f(\Theta) = \Theta^n + \alpha_1\Theta^{n-1} + \cdots + \alpha_n = 0}
be an equation of degree $n$, rational and irreducible in $R$. The field $\Omega = R(\Theta)$ is an algebraic field \gesperrt{over} $R$, which with respect to $R$ is of degree $n$ (Bd. I, \S.\,142). In the absolute domain of rationality $\Omega$ is of degree $mn$, and if $\varrho$ is a primitive number of the field $R$, then
\weq{2}{\Theta^s \varrho^t, \qquad \begin{aligned} s &= 0, 1, \ldots, n-1\\ t &= 0, 1, \ldots, m-1 \end{aligned}}
is a basis of the field $\Omega$; for, owing to the irreducibility of the
%CONT>>

${}^{1}$) \gesperrt{Dedekind}, ``Ueber die Discriminanten endlicher K\"orper'', Abhandlungen der G\"ottinger Gesellschaft der Wissenschaften (1882). ``Zur Theorie der Ideale'', Nachrichten der Gesellschaft der Wissenschaften zu G\"ottingen, 1894, Nr. 4. \gesperrt{Hilbert}, ``Grundz\"uge einer Theorie des \gesperrt{Galois'schen} Zahlk\"orpers''. Ibid. 1894, Nr. 3.


% ===== scan leaf 605, printed 584 =====
% !PRINTED 584
\weberpage{584}
%<<CONT
equation (1), no linear relation with rational numerical coefficients can subsist between the $mn$ numbers (2).

Every number $\omega$ of the field $\Omega$ can be represented as an integral function of degree $(n-1)$ of $\Theta$ with coefficients in $R$, and every such number $\omega$ satisfies an equation irreducible in $R$ of degree at most $n$; $\omega$ is then and only then an integer, when this equation, after the coefficient of the highest power has been brought to 1, has integers in $R$ as coefficients.

The totality of all integers in $\Omega$ we denote, as before, by $\mathfrak{o}$.

If, without altering the numbers of the field $R$, we take for $\Theta$ all the roots $\Theta_1, \Theta_2, \ldots, \Theta_n$ of equation (1), we obtain $n$ fields
\weq{3}{\Omega_1 = R(\Theta_1), \ \Omega_2 = R(\Theta_2), \ \ldots, \ \Omega_n = R(\Theta_n),}
\gesperrt{which are conjugate with respect to the field $R$.}

If $\omega_1, \omega_2, \ldots, \omega_n$ are corresponding numbers of these fields, the product
\weq{4}{\mathfrak{N}_R(\omega) = \omega_1 \omega_2 \ldots \omega_n}
is called \gesperrt{the partial norm (relative norm) taken with respect to $R$} of the number $\omega$.

A number $\eta$ of the field $R$ has an ordinary norm $N_R(\eta)$, taken in this field, which is a rational number. Now the relative norm $\mathfrak{N}_R(\omega)$ is such a number $\eta$, and if we take its norm in the field $R$, we obtain the \gesperrt{total norm} of the number $\omega$ in the field of rational numbers:
\weq{5}{N_\Omega(\omega) = N_R\,\mathfrak{N}_R(\omega).}

This results if one expresses $\omega$ by the powers (2) of $\Theta$ and $\varrho$, and then puts for $\varrho$ the $m$ conjugate values, and for each of these $\varrho$ the values $\Theta$ conjugate with respect to $R$.

In the same sense the \gesperrt{functionals} of the field $\Omega$, and the \gesperrt{ideals} represented by these, now also have \gesperrt{their partial norms}. The partial norm of a functional in $\Omega$ is a functional in $R$, and accordingly the partial norm of an ideal in $\Omega$ is an ideal in $R$.

Of the conjugate ideals two or more can also be equal to one another, and indeed this can occur even when the conjugate functionals representing the ideals are not identical, provided only that they are associated. If one
%CONT>>


% ===== scan leaf 606, printed 585 =====
% !PRINTED 585
\weberpage{585}
%<<CONT
therefore takes only the product of all the mutually distinct ones among the conjugate ideals, this need by no means be an ideal in $R$.

What has been set forth here for the norms holds also for the other symmetric functions, in particular for the traces. We call the sum
\weq{6}{\mathfrak{S}_R(\omega) = \omega_1 + \omega_2 + \cdots + \omega_n}
the \gesperrt{partial trace (or relative trace) of $\omega$ with respect to} $R$. It is a number or a functional in $R$, and for the total trace we obtain
\weq{7}{S_\Omega(\omega) = S_R\,\mathfrak{S}_R(\omega),}
in which the meaning of the symbols is immediately intelligible.

A particularly important relation results, however, from the consideration of the fundamental ideals.

Let $\tau$ be a basis form of $\Omega$ (thus a linear form in $mn$ variables, \S.\,146), and let
\[
F(t) = N_\Omega(t - \tau)
\]
be the irreducible function of degree $mn$ (in the absolute domain of rationality) whose root is $\tau$. Then in \S.\,158 we designated the ideal defined by the functional $F'(\tau)$ as the \gesperrt{fundamental ideal} of the field $\Omega$.

If now $\sigma$ is a basis form of $R$ and
\[
\psi(t) = N_R(t - \sigma),
\]
then the fundamental ideal of the field $R$ is defined by $\psi'(\sigma)$.

From the function $F(t)$, however, there splits off in the field $R$ an irreducible factor of degree $n$,
\[
f(t) = \mathfrak{N}_R(t - \tau),
\]
which vanishes for $t = \tau$. By $f'(\tau)$ we define the \gesperrt{partial fundamental ideal of $\Omega$ with respect to} $R$, and now prove the theorem:
\weq{8}{F'(\tau) = \varepsilon\, f'(\tau)\, \psi'(\sigma),}
in which $\varepsilon$ is a unit.

To prove it, we set, following the notation employed in \S.\,158,
\weq{9}{\frac{F(t)}{t - \tau} = t^{mn-1}F_0(\tau) + t^{mn-2}F_1(\tau) + \cdots + F_{mn-1}(\tau),}
\weq{10}{\frac{f(t)}{t - \tau} = t^{n-1}f_0(\tau) + t^{n-2}f_1(\tau) + \cdots + f_{n-1}(\tau),}
\weq{11}{\frac{\psi(t)}{t - \sigma} = t^{m-1}\psi_0(\sigma) + t^{m-2}\psi_1(\sigma) + \cdots + \psi_{m-1}(\sigma).}


% ===== scan leaf 607, printed 586 =====
% !PRINTED 586
\weberpage{586}
Here $F_\nu$, $\psi_\nu$ are integral functionals of the absolute domain of rationality,
while $f_\nu$ are integral functionals in $R$. The expressions for $F_\nu, f_\nu, \psi_\nu$
result from the formulae \S.\,158, (15). To these functionals one can apply the same
inferences that led to the Theorem \S.\,158, 4., and it results that:

\weq{12}{
\begin{aligned}
\mathfrak{S}_R \frac{f_\nu(\tau)}{f'(\tau)} &= 0, & \nu &= 0, 1, \ldots, n-2.\\
\mathfrak{S}_R \frac{f_{n-1}(\tau)}{f'(\tau)} &= 1.
\end{aligned}
}
If now one divides the functions $F_\nu(t)$, when they are of degree higher than $(n-1)$,
by $f(t)$ and takes the remainder of the division; and then in this remainder expresses
the powers of $t$ [according to \S.\,158, (15), applied to $f_\nu$] by the functions
$f_\nu(t)$, there results the representation

\weq{13}{
F_\nu(\tau) = \alpha_0 f_0 + \alpha_1 f_1 + \cdots + \alpha_{n-1} f_{n-1},
}
in which the $\alpha_0, \alpha_1, \ldots, \alpha_{n-1}$ are \gesperrt{integral functionals} in $R$. This
representation holds for every root $\tau$ of $f(t) = 0$, that is, for $\tau = \tau_1, \tau_2,
\ldots, \tau_n$. From this, however, there results by (12):
\[
\mathfrak{S}_R \left( \frac{F_\nu(\tau)}{f'(\tau)} \right) = \alpha_{n-1}.
\]
If one here divides by $\psi'(\sigma)$ and forms the trace $S_R$, one obtains from (7),
with application of Theorem 4., \S.\,158:
\weq{14}{
S_\Omega \left( \frac{F_\nu(\tau)}{f'(\tau)\,\psi'(\sigma)} \right) = a_\nu,
}
in which $a_\nu$ denotes an integral rational functional.

By (9) we now have, for every root $\tau_i$ of $F(t) = 0$ with the exception of $\tau_i = \tau$,
\[
\tau^{n\,m-1} F_0(\tau_i) + \tau^{n\,m-2} F_1(\tau_i) + \cdots + F_{n\,m-1}(\tau_i) = 0,
\]
and
\[
\tau^{n\,m-1} F_0(\tau) + \tau^{n\,m-2} F_1(\tau) + \cdots + F_{n\,m-1}(\tau) = F'(\tau),
\]
so that, when one multiplies (14) by $\tau^{n\,m-\nu-1}$ and takes the sum with respect
to $\nu$,
\weq{15}{
\frac{F'(\tau)}{f'(\tau)\,\psi'(\sigma)} = a_0 \tau^{n\,m-1} + a_1 \tau^{n\,m-2} + \cdots
}
results. \gesperrt{Therefore $F'(\tau)$ is divisible by $f'(\tau)\,\psi'(\sigma)$.}


% ===== scan leaf 608, printed 587 =====
% !PRINTED 587
\weberpage{587}
On the other hand, $f_\nu(\tau)\,\psi_\mu(\sigma)$ is an integral functional of the field
$\Omega$, and consequently, by the theorem \S.\,158, 4. already applied above:

\weq{16}{
\begin{aligned}
S_\Omega \left( \frac{f_\nu(\tau)\,\psi_\mu(\sigma)}{F'(\tau)} \right) &= b_{\mu,\nu} &
\mu &= 0, 1, \ldots, m-1\\
& & \nu &= 0, 1, \ldots, n-1
\end{aligned}
}
is an integral rational functional. If we multiply this equation by
$\tau^{n-\nu-1}\sigma^{m-\mu-1}$ and take the sum over all $\mu$ and $\nu$, then,
as before — since for two roots $\sigma_k, \tau_i$ of $\psi=0, f=0$ different
from $\sigma, \tau$
\[
\begin{aligned}
\sigma^{m-1}\psi_0(\sigma_k) + \sigma^{m-2}\psi_1(\sigma_k) + \cdots + \psi_{m-1}(\sigma_k) &= 0,\\
\tau^{n-1} f_0(\tau_i) + \tau^{n-2} f_1(\tau_i) + \cdots + f_{n-1}(\tau_i) &= 0,
\end{aligned}
\]
and
\[
\begin{aligned}
\sigma^{m-1}\psi_0(\sigma) + \sigma^{m-2}\psi_1(\sigma) + \cdots + \psi_{m-1}(\sigma) &= \psi'(\sigma),\\
\tau^{n-1} f_0(\tau) + \tau^{n-2} f_1(\tau) + \cdots + f_{n-1}(\tau) &= f'(\tau)
\end{aligned}
\]
hold — we obtain the equation
\weq{17}{
\frac{f'(\tau)\,\psi'(\sigma)}{F'(\tau)} = \sum^{\mu,\nu} b_{\mu,\nu}\, \tau^{n-\nu-1}\sigma^{m-\mu-1}.
}

Hence $f'(\tau)\,\psi'(\sigma)$ too is divisible by $F'(\tau)$, and our Theorem (8)
is therefore proved.

If we denote by $G_\Omega, G_R$ the fundamental ideals of the fields $\Omega, R$, and by
$\mathfrak{G}_R$ the relative fundamental ideal of $\Omega$ with respect to $R$, then the
theorem just proved may also be given the following expression:
\weq{18}{
G_\Omega = \mathfrak{G}_R\, G_R.
}
It is clear how this theorem may be generalised: Let
\[
\Omega_1, \Omega_2, \Omega_3, \ldots
\]
be a sequence of fields, each of which contains the following ones as divisors. Further, let
$\mathfrak{G}_{i,k}$ be the fundamental ideal of $\Omega_i$ with respect to
$\Omega_k$ ($k > i$); then
\weq{19}{
\mathfrak{G}_{i,k} = \mathfrak{G}_{i,i+1}\,\mathfrak{G}_{i+1,i+2} \cdots \mathfrak{G}_{k-1,k}.
}
The deeper significance of these fundamental ideals will emerge still more clearly in one
of the following sections.

But we wish to call attention already here to a remarkable phenomenon: namely, that it can
happen that $f'(\tau)$ is a unit, which by \gesperrt{Minkowski's} theorem is not
possible when $R$ is the field of rational numbers. In these cases, of which the theory of com-
%CONT>>


% ===== scan leaf 609, printed 588 =====
% !PRINTED 588
\weberpage{588}
\pdfbookmark[2]{§ 160. Prime Ideals in the Relatively Normal Field}{bm183}
%<<CONT
plex multiplication of the elliptic functions gives examples, the
partial fundamental ideal $\mathfrak{G}_R$ is to be set equal to 1.

If the $n$ fields (3) are identical with one another, then $\Omega$ is called a
\gesperrt{normal field with respect to} $R$ (relatively normal, if $R$ is not the absolute
domain of rationality). The totality of the $n$ substitutions $(\Theta, \Theta_n)$
that effect the transition of the numbers of the field $\Omega$ to the conjugate
numbers forms a group of degree $n$ --- as we have already seen in \S.\,147 of the
first volume --- which is nothing other than the \gesperrt{Galois} group of the
equation (1) in the domain of rationality $R$, and which we here call the \gesperrt{group of
the field $\Omega$ (with respect to $R$)}.

If in particular this group is commutative, then $\Omega$ is an
\gesperrt{Abelian field with respect to} $R$ (a relatively \gesperrt{Abelian field}). In this
case (1) is an \gesperrt{Abelian} equation in the domain of rationality $R$.

\begin{center}
\S.\,160.\\[6pt]
\gesperrt{Prime ideals in the relatively normal field.}
\end{center}

We now assume that the field $\Omega$ is normal with respect to the arbitrary
field $R$, and under this assumption we investigate the prime ideals of the field
$\Omega$.

Let the group of $\Omega$ with respect to $R$ be $\Phi$, and let the substitutions
of $\Phi$ be denoted by $\varphi, \varphi_1, \ldots$ We further denote by
\weq{1}{
\omega \mid \varphi
}
\gesperrt{the number that arises from $\omega$ by the substitution $\varphi$}, so that $\omega$
and $\omega \mid \varphi$ are contained in $\Omega$.

Just as the numbers $\omega$ do, so also all functionals, and with them at the same
time all ideals $\mathfrak{a}$ of the field $\Omega$, pass by a substitution
$\varphi$ into definite functionals and ideals $\mathfrak{a} \mid \varphi$; and if
$\omega$ is an integer divisible by $\mathfrak{a}$, then $\omega \mid \varphi$ is
divisible by $\mathfrak{a} \mid \varphi$.

If $\mathfrak{a}$ is any ideal in $\Omega$, then there are certain substitutions
$\psi$ in $\Phi$ that satisfy the condition
\weq{2}{
\mathfrak{a} \mid \psi = \mathfrak{a}
}
among them always the identical substitution, and these substitutions $\psi$
form a group $\Psi$, which is a divisor of $\Phi$
%CONT>>


% ===== scan leaf 610, printed 589 =====
% !PRINTED 589
\weberpage{589}
%<<CONT
is. \gesperrt{We call $\Psi$ the group belonging to the ideal $\mathfrak{a}$, and we also say
that $\mathfrak{a}$ belongs to the group $\Psi$.}

Since in every equation between numbers and functionals in $\Omega$ one may carry out all
the substitutions of the group $\Phi$, whereby the quantities of the field $R$
remain unchanged, without the equation ceasing to hold, it follows that
units, integral functionals, associated functionals, and functionals divisible by one
another do not lose these properties when any one of the substitutions of
$\Phi$ is carried out. We single out the theorem:

1. \gesperrt{If $\mathfrak{p}$ is a prime ideal, then all the ideals $\mathfrak{p} \mid \varphi$
conjugate to $\mathfrak{p}$ with respect to $R$ are also prime ideals. If $\mathfrak{a}$
is divisible by any power of $\mathfrak{p}$ then $\mathfrak{a} \mid \varphi$ is divisible
by the same power of $\mathfrak{p} \mid \varphi$.} % [sic] German source omits the comma after "theilbar" and prints a stray comma after $\mathfrak p\mid\varphi$

If $\mathfrak{p}$ is a prime ideal in $\Omega$, there is one and only one natural prime
$p$ divisible by $\mathfrak{p}$.

If we decompose this into its prime factors in $R$, one of these prime factors, which
we denote by $\mathfrak{P}$, must be divisible by $\mathfrak{p}$.

If now $\mathfrak{A}$ is any ideal in $R$ which is at the same time divisible by $\mathfrak{p}$,
then the greatest common divisor of $\mathfrak{P}$ and
$\mathfrak{A}$, which by \S.\,139 is also contained in $R$, is divisible by $\mathfrak{p}$,
and is therefore certainly not a unit. $\mathfrak{P}$ and $\mathfrak{A}$ are thus
not relatively prime, and $\mathfrak{A}$ is consequently divisible by the prime ideal $\mathfrak{P}$.
If $\mathfrak{A}$ is also a prime ideal, then $\mathfrak{A}$ and
$\mathfrak{P}$ must be identical. With this, having regard to 1., it is proved:

2. \gesperrt{If $\mathfrak{p}$ is a prime ideal in $\Omega$, there is one and only one
prime ideal $\mathfrak{P}$ in $R$ that is divisible by $\mathfrak{p}$, and every ideal
in $R$ divisible by $\mathfrak{p}$ is at the same time divisible by $\mathfrak{P}$. The
prime ideal $\mathfrak{P}$ is at the same time divisible by all the
prime ideals $\mathfrak{p} \mid \varphi$ conjugate to $\mathfrak{p}$.}

The partial norm $\mathfrak{N}_R(\mathfrak{p})$ is divisible by $\mathfrak{p}$, and consequently, as
an ideal in $R$, by $\mathfrak{P}$. But by the theorem 1. it can, moreover, be
divisible by no ideal in $R$ different from $\mathfrak{P}$. For if $\mathfrak{P}'$
is any prime divisor of $\mathfrak{N}_R(\mathfrak{p})$, then $\mathfrak{P}'$ is divisible by
at least one of the prime factors $\mathfrak{p} \mid \varphi$ and is hence equal to $\mathfrak{P}$. Consequently
%CONT>>


% ===== scan leaf 611, printed 590 =====
% !PRINTED 590
\weberpage{590}
%<<CONT
the partial norm of $\mathfrak{p}$ is a power of $\mathfrak{P}$. We set
\weq{3}{\mathfrak{N}_R(\mathfrak{p}) = \mathfrak{P}^f,}
and call $f$ \gesperrt{the degree of $\mathfrak{p}$ with respect to the field $R$.}

The norm of $\mathfrak{P}$ in the field $R$ is likewise a power of $p$, which we shall denote by $P$; thus
\weq{4}{N_R(\mathfrak{P}) = P.}

For the total norm of $\mathfrak{P}$ in the field $\Omega$ we then obtain, by \S.\,159, (5) % [sic] source names $\mathfrak P$ here, while (5) and the norm identity require $\mathfrak p$
\weq{5}{N_\Omega(\mathfrak{p}) = P^f.}

If $\mathfrak{p}^g$ is the highest power of $\mathfrak{p}$ that divides $\mathfrak{P}$, then by 1. $\mathfrak{P}$ is also divisible by the $g$th power of all the prime factors conjugate to $\mathfrak{p}$, and by no higher power. Now if $\mathfrak{p}_1, \mathfrak{p}_2, \ldots, \mathfrak{p}_e$ are \gesperrt{the mutually distinct ones} among the prime ideals conjugate to $\mathfrak{p}$, then
\weq{6}{\mathfrak{P} = (\mathfrak{p}_1\mathfrak{p}_2\ldots\mathfrak{p}_e)^g,}
and if one takes the partial norm on both sides, and, as before, denotes by $n$ the degree of the field $\Omega$ with respect to $R$, so that the partial norm of $\mathfrak{P}$ becomes equal to $\mathfrak{P}^n$, then it follows from (3):
\weq{7}{n = efg.}
The natural numbers $e, f, g$ are therefore divisors of $n$.

If $\mathfrak{p}$ passes into $\mathfrak{p}_1$ under $\varphi_1$, that is, if
\[
\mathfrak{p}_1 = \mathfrak{p} \,|\, \varphi_1,
\]
then $\mathfrak{p} = \mathfrak{p}_1 \,|\, \varphi_1^{-1}$, and if $\Psi$ is the group associated with $\mathfrak{p}$, then for every substitution $\psi$ contained in $\Psi$ we also have
\[
\mathfrak{p}_1 = \mathfrak{p} \,|\, \psi \varphi_1, \quad \mathfrak{p}_1 = \mathfrak{p}_1 \,|\, \varphi_1^{-1} \psi \varphi_1.
\]

If, conversely, $\mathfrak{p}_1 \,|\, \varphi = \mathfrak{p}_1$, then it follows that $\mathfrak{p} \,|\, \varphi_1 \varphi \varphi_1^{-1} = \mathfrak{p}$, that is, $\varphi_1 \varphi \varphi_1^{-1} = \psi$ or $\varphi = \varphi_1^{-1} \psi \varphi_1$.

Hence $\mathfrak{p}_1$ is associated with the group $\varphi_1^{-1} \Psi \varphi_1$, and if
\[
\mathfrak{p}_1 = \mathfrak{p} \,|\, \varphi_1, \quad \mathfrak{p}_2 = \mathfrak{p} \,|\, \varphi_2, \ldots, \mathfrak{p}_e = \mathfrak{p} \,|\, \varphi_e
\]
then
\weq{8}{\Phi = \Psi\varphi_1 + \Psi\varphi_2 + \cdots + \Psi\varphi_e.}
$\Psi$ is thus a divisor of $\Phi$ of index $e$ and of degree $gf$.


% ===== scan leaf 612, printed 591 =====
% !PRINTED 591
\weberpage{591}
\pdfbookmark[2]{§ 161. Primitive Roots of the Prime Ideals}{bm184}
\begin{center}
\S.\,161.\\[4pt]
\gesperrt{Primitive Roots of the Prime Ideals.}
\end{center}

From the formulae (4), (5) of the preceding section it follows, having regard to \S.\,150, that for every integer $\eta$ of the field $R$
\[
\eta^P \equiv \eta \pmod{\mathfrak{P}},
\]
hence also
\weq{1}{\eta^P \equiv \eta \pmod{\mathfrak{p}},}
and for every number $\omega$ in $\mathfrak{o}$
\weq{2}{\omega^{P^f} \equiv \omega \pmod{\mathfrak{p}}.}
The number of integers in $R$ incongruent with respect to $\mathfrak{P}$ is $P$, and the number of numbers in $\Omega$ incongruent with respect to $\mathfrak{p}$ is $P^f$.

We have already shown earlier (\S.\,150) that for every prime ideal $\mathfrak{p}$ there are \gesperrt{primitive roots} $\gamma$, that is, numbers in $\Omega$ which satisfy the congruence
\weq{3}{\gamma^{P^f-1} \equiv 1 \pmod{\mathfrak{p}}}
and at the same time have the property that no lower power with positive exponent becomes congruent to unity. Then every integer in $\Omega$ not divisible by $\mathfrak{p}$ is congruent to one and only one of the powers of $\gamma$
\weq{4}{1, \gamma, \gamma^2, \ldots, \gamma^{P^f-2}}
with respect to the modulus $\mathfrak{p}$.

Every integer $\Theta$ in $\Omega$ satisfies an equation of degree at most $n$, whose coefficients are integers in $R$.

This equation is at the same time a congruence with respect to the modulus $\mathfrak{p}$, and among all such congruences there will be one of \gesperrt{lowest possible degree}
\[
F(\Theta) \equiv 0 \pmod{\mathfrak{p}}
\]
in which, moreover, we may assume the coefficient of the highest power of $\Theta$ to be equal to 1. For if the highest coefficient (indivisible by $\mathfrak{p}$, and consequently by $\mathfrak{P}$) is $\eta_0$, then we can always satisfy the congruence
\[
\eta_0 \eta \equiv 1 \pmod{\mathfrak{P}}
\]
by an integer $\eta$ in $R$, and have then only to multiply the function $F$ by this $\eta$. We can
%CONT>>


% ===== scan leaf 613, printed 592 =====
% !PRINTED 592
\weberpage{592}
%<<CONT
thus, when $t$ is a variable, assume the function $F$ in the form:
\[
F(t) = t^\nu + \alpha_1 t^{\nu-1} + \alpha_2 t^{\nu-2} + \cdots + \alpha_\nu,
\]
whose coefficients $\alpha_1, \alpha_2, \ldots, \alpha_\nu$ are integers in $R$. By division we can then, for every other integral function $F_1(t)$, determine the quotient $Q(t)$ and the residue $\varphi(t)$ in such a way that
\weq{5}{F_1(t) = Q(t) F(t) + \varphi(t),}
wherein
\weq{6}{\varphi(t) = \varrho_0 + \varrho_1 t + \cdots + \varrho_{\nu-1} t^{\nu-1}}
is a function of degree $(\nu-1)$ uniquely determined by $F_1(t)$, with integral coefficients $\varrho$.

If we set for $\Theta$ the primitive root $\gamma$, and for $F_1(t)$ a power of $t$, it results from (4), (5) and (6) that every number $\omega$ not divisible by $\mathfrak{p}$ satisfies a congruence
\weq{7}{\omega \equiv \varrho_0 + \varrho_1 \gamma + \cdots + \varrho_{\nu-1} \gamma^{\nu-1} \pmod{\mathfrak{p}}}
and since the $\varrho$ may also be $= 0$, this still holds for an $\omega$ divisible by $\mathfrak{p}$.

The coefficients $\varrho$ in the expression (7) are completely determined by $\omega$ itself with respect to the modulus $\mathfrak{P}$, since by our assumption $\gamma$ is to satisfy no congruence in $R$ of degree lower than $\nu$ whose coefficients are not all divisible by $\mathfrak{P}$. Consequently, since each of the coefficients in (7) can have only $P$ incongruent values, there are $P^\nu$ and no more incongruent numbers $\omega$, and from this it follows that we must have $\nu = f$.

But now from a congruence $F(\gamma) \equiv 0 \pmod{\mathfrak{p}}$ there follows, by repeated raising to powers, with regard to (1):
\[
F(\gamma^P) \equiv 0, \ F(\gamma^{P^2}) \equiv 0, \ldots \pmod{\mathfrak{p}},
\]
and we have thus proved the theorem (\S.\,150, 2.):

\gesperrt{A primitive root $\gamma$ of $\mathfrak{p}$ satisfies, with respect to the modulus $\mathfrak{p}$, a congruence in $R$ of degree $f$, of which}
\[
\gamma, \gamma^P, \gamma^{P^2}, \ldots, \gamma^{P^f-1}
\]
\gesperrt{are all the roots.}
% [sic: The source prints $\gamma^{P^f-1}$ as the last root, with -1 on the P-exponent level rather than inside the exponent of P. The P-power orbit requires $\gamma^{P^{f-1}}$; the printed expression is retained.]

This theorem can also be given the expression:

\gesperrt{The product}
% [sic: The last factor repeats the source's $\gamma^{P^f-1}$; the P-power-orbit factor would be $\gamma^{P^{f-1}}$.]
\weq{8}{(t-\gamma)(t-\gamma^P) \ldots (t-\gamma^{P^f-1})}
\gesperrt{is congruent, with respect to the modulus $\mathfrak{p}$, to an integral function $F(t)$ in $R$.}


% ===== scan leaf 614, printed 593 =====
% !PRINTED 593
\weberpage{593}
\pdfbookmark[2]{§ 162. The Partial Fundamental Ideal}{bm185}
\begin{center}
\S.\,162.\\[4pt]
\gesperrt{The Partial Fundamental Ideal.}
\end{center}

Now let
\weq{1}{\tau = t_1\omega_1 + t_2\omega_2 + \cdots + t_{mn}\omega_{mn}}
be a basis form of $\mathfrak{o}$. Then there will be certain substitutions $\chi$ in $\Phi$, and among them always the identical one, by which the congruence
\weq{2}{\tau \,|\, \chi \equiv \tau \pmod{\mathfrak{p}}}
is satisfied. All these substitutions form a group $X$ contained in $\Phi$, which is also a divisor of the group $\Psi$ to which $\mathfrak{p}$ belongs. For from $\tau$ one obtains (by \S.\,146) a basis form of $\mathfrak{p}$ if one puts for the variables $t$ certain linear functions of new variables with integral rational coefficients; if then $\pi$ is such a basis form of $\mathfrak{p}$, it follows from (2) that $\pi \,|\, \chi$ is divisible by $\pi$, and hence also that $\mathfrak{p} \,|\, \chi = \mathfrak{p}$.

Now $\tau$ satisfies an equation of degree $n$, $f(t) = 0$, whose coefficients are integral functions in $R$ with the variables $t_1, t_2 \ldots$; and if $t$ denotes a new variable, and $\tau_1, \tau_2, \ldots, \tau_n$ are the forms conjugate to $\tau$ with respect to $R$, we have
\weq{3}{f(t, t_1, t_2, \ldots) = f(t) = (t-\tau_1)(t-\tau_2) \ldots (t-\tau_n).}
From this we now obtain, if we repeatedly raise to the power $p$ and apply to the numerical coefficients of $f$ the formula \S.\,161, (1):
\weq{4}{f(t^P, t_1^P, t_2^P, \ldots) \equiv [f(t)]^P \equiv (t^P-\tau_1')(t^P-\tau_2') \ldots (t^P-\tau_n') \pmod{\mathfrak{p}},}
where $\tau_1', \tau_2', \ldots, \tau_n'$ are derived from $\tau_1, \tau_2, \ldots, \tau_n$ by replacing the variables $t_1, t_2, \ldots$ by $t_1^P, t_2^P, \ldots$.

If we now put $t = \tau$ in (4), then $f(t)$ vanishes, and it follows that one of the factors of the last product must be divisible by $\mathfrak{p}$. But this can be expressed by saying that there is in $\Phi$ a substitution $\psi_0$ by which the condition
\weq{5}{\tau^P \equiv \tau' \,|\, \psi_0 \pmod{\mathfrak{p}}}
is satisfied.

From $\tau$ all the integers in $\Omega$ arise, when one puts integral rational numbers for the variables. If one then applies
%CONT>>


% ===== scan leaf 615, printed 594 =====
% !PRINTED 594
\weberpage{594}
%<<CONT
\gesperrt{Fermat's} theorem for rational numbers as well, one recognises that the congruences (2) and (5) are \gesperrt{equivalent in meaning} to the formulae valid for \gesperrt{every number} $\omega$ in $\mathfrak{o}$

\weq{6}{\omega\mkern1mu|\mkern1mu\chi \equiv \omega,\quad \omega^P \equiv \omega\mkern1mu|\mkern1mu\psi_0\ (\mathrm{mod}\ \mathfrak{p}).}
From the second congruence (6) it follows that, if $\omega$ is divisible by $\mathfrak{p}$, then $\omega\mkern1mu|\mkern1mu\psi_0$ must always be divisible by $\mathfrak{p}$ as well. Consequently $\mathfrak{p}$ is divisible by $\mathfrak{p}\mkern1mu|\mkern1mu\psi_0$, and, being a prime ideal, $= \mathfrak{p}\mkern1mu|\mkern1mu\psi_0$. From this it follows that $\psi_0$ is contained in the group $\Psi$.

We now understand by $\gamma$ a primitive root of $\mathfrak{p}$, and apply the formula proved at the end of \S.\,161:
% [sic: The source repeats $\gamma^{P^f-1}$ in the last factor, as on printed page 592; the P-power orbit instead requires $\gamma^{P^{f-1}}$. The printed expression is retained.]
\weq{7}{(t-\gamma)(t-\gamma^P)\ldots(t-\gamma^{P^f-1}) \equiv F(t)\ (\mathrm{mod}\ \mathfrak{p}),}
in which $F(t)$ is an integral function of $t$ in $R$.

From this it follows that
\[ F(\gamma) \equiv 0\ (\mathrm{mod}\ \mathfrak{p}), \]
and if now $\psi$ is any substitution from $\Psi$:
\[ F(\gamma\mkern1mu|\mkern1mu\psi) \equiv 0\ (\mathrm{mod}\ \mathfrak{p}). \]

But from this it results that $\gamma\mkern1mu|\mkern1mu\psi$ must be congruent, with respect to $\mathfrak{p}$, to one of the powers of $\gamma$ occurring in (7), and so, say,
\weq{8}{\gamma^{P^\nu} \equiv \gamma\mkern1mu|\mkern1mu\psi\ (\mathrm{mod}\ \mathfrak{p}).}
Now every number $\omega$ in $\mathfrak{o}$ not divisible by $\mathfrak{p}$ is congruent to a power of $\gamma$, and if one therefore raises (8) to this power, it follows that
\weq{9}{\omega^{P^\nu} \equiv \omega\mkern1mu|\mkern1mu\psi\ (\mathrm{mod}\ \mathfrak{p}),}
and this congruence evidently holds even when $\omega$ is divisible by $\mathfrak{p}$, and hence for \gesperrt{all numbers} in $\mathfrak{o}$.

If we now apply the second of the congruences (6) $\nu$ times in succession, it follows that
\weq{10}{\omega^{P^\nu} \equiv \omega\mkern1mu|\mkern1mu\psi_0^\nu\ (\mathrm{mod}\ \mathfrak{p}),}
hence
\weq{11}{\omega\mkern1mu|\mkern1mu\psi \equiv \omega\mkern1mu|\mkern1mu\psi_0^\nu\ (\mathrm{mod}\ \mathfrak{p}),}
and if we replace $\omega$ by $\omega\mkern1mu|\mkern1mu\psi^{-1}$, or apply the substitution $\psi^{-1}$ to (11):
\weq{12}{\omega \equiv \omega\mkern1mu|\mkern1mu\psi^{-1}\psi_0^\nu \equiv \omega\mkern1mu|\mkern1mu\psi_0^\nu\psi^{-1}\ (\mathrm{mod}\ \mathfrak{p}).}
Thus both $\psi^{-1}\psi_0^\nu$ and $\psi_0^\nu\psi^{-1}$ are contained in $X$,
%CONT>>


% ===== scan leaf 616, printed 595 =====
% !PRINTED 595
\weberpage{595}
%<<CONT
from which it follows that \gesperrt{in one of the systems} (subsidiary groups)
\[ X,\ X\psi_0,\ X\psi_0^2,\ \ldots \]
\gesperrt{every substitution $\psi$ of $\Psi$ is contained.}

If $\mathfrak{p}$ is of degree $f$ (with respect to $R$), then
\[ \omega^{P^f} \equiv \omega\ (\mathrm{mod}\ \mathfrak{p}), \]
and $P^f$ is the lowest power of $P$ with positive exponent that satisfies this congruence for all $\omega$.

If we therefore set $\nu = f$ in (10), it follows that $\psi_0^f$ is contained in $X$, but that no power of $\psi_0$ with a lower positive exponent has this property. The subsidiary groups
\[ X,\ X\psi_0,\ X\psi_0^2,\ \ldots,\ X\psi_0^{f-1} \]
are therefore all different from one another, and we obtain:
\weq{13}{\Psi = X + X\psi_0 + X\psi_0^2 + \cdots + X\psi_0^{f-1}.}

By (12), however, the totality of the substitutions $\psi_0^\nu X$ is identical with $X\psi_0^\nu$ (for every $\nu$); hence
\weq{14}{\psi_0^\nu X = X\psi_0^\nu,}
from which it follows that \gesperrt{$X$ is a normal divisor of $\Psi$.}

\gesperrt{Since, as we have seen above, $\Psi$ is of degree $gf$, $X$ is of degree $g$.}

Now in the group $X$ there are substitutions $\chi_1$, among them always the identical one, for which the congruence
\weq{15}{\tau\mkern1mu|\mkern1mu\chi_1 \equiv \tau\ (\mathrm{mod}\ \mathfrak{p}^2)}
is satisfied, and these form a group $X^{(1)}$, whose degree may be denoted by $g_1$.

\gesperrt{It is to be proved that $g_1$ is a power of $p$.} For if we denote by $\pi$ a prime functional divisible by $\mathfrak{p}$, and by $\alpha$ an integral functional, then it follows from (15):
\weq{16}{\tau\mkern1mu|\mkern1mu\chi_1 = \tau + \pi^2\alpha,}
and since likewise, by (15),
\[ \alpha\mkern1mu|\mkern1mu\chi_1 \equiv \alpha\ (\mathrm{mod}\ \mathfrak{p}^2),\quad (\pi\mkern1mu|\mkern1mu\chi_1)^2 \equiv \pi^2\ (\mathrm{mod}\ \mathfrak{p}^3) \]
holds, it results from (16) that
\[ \tau\mkern1mu|\mkern1mu\chi_1^2 \equiv \tau + 2\pi^2\alpha\ (\mathrm{mod}\ \mathfrak{p}^3), \]
and in general, for every positive exponent $h$,
\[ \tau\mkern1mu|\mkern1mu\chi_1^h \equiv \tau + h\pi^2\alpha\ (\mathrm{mod}\ \mathfrak{p}^3). \]


% ===== scan leaf 617, printed 596 =====
% !PRINTED 596
\weberpage{596}
If one sets $h = p$, this yields
\[ \tau\mkern1mu|\mkern1mu\chi_1^p \equiv \tau\ (\mathrm{mod}\ \mathfrak{p}^3), \]
and in the same way one concludes from this, for every exponent $\nu$,
\weq{17}{\tau\mkern1mu|\mkern1mu\chi_1^{p^\nu} \equiv \tau\ (\mathrm{mod}\ \mathfrak{p}^{\nu+2}).}
To prove this formula in general, one applies complete induction. One assumes (17) as proved and sets
\[ \tau\mkern1mu|\mkern1mu\chi_1^{p^\nu} = \tau + \pi^{\nu+2}\alpha; \]
from this, by repeated application,
\[
\begin{aligned}
\tau\mkern1mu|\mkern1mu\chi_1^{2p^\nu} &\equiv \tau + 2\pi^{\nu+2}\alpha \\
\tau\mkern1mu|\mkern1mu\chi_1^{3p^\nu} &\equiv \tau + 3\pi^{\nu+2}\alpha\ (\mathrm{mod}\ \mathfrak{p}^{\nu+3}) \\
&\cdots\cdots\cdots\cdots \\
\tau\mkern1mu|\mkern1mu\chi_1^{p^{\nu+1}} &\equiv \tau + p\pi^{\nu+2}\alpha,
\end{aligned}
\]
and thus the formula (17) for $\nu + 1$.

Now one can always take $\nu$ so great that, of the differences
\[ \tau - \tau_1,\ \tau - \tau_2,\ \tau - \tau_3,\ \ldots, \]
if $\tau$ is different from $\tau_1, \tau_2, \tau_3, \ldots$, none is divisible by $\mathfrak{p}^{\nu-2}$, and then it follows from (17) that $\chi_1^{p^\nu}$ must be the identical substitution. \gesperrt{According to this, the degree of every element $\chi_1$ is a power of $p$, and consequently, by the Cauchy-Sylow theorem (\S.\,29, I.), no prime different from $p$ can divide the degree of $X^{(1)}$, as was to be proved.}

If, then, $g$ is not divisible by $p$, then $X^{(1)}$ is certainly the unit group.

One can continue in this way and form a group $X^{(2)}$ of degree $g_2$ from the congruence
\[ \tau\mkern1mu|\mkern1mu\chi_2 \equiv \tau\ (\mathrm{mod}\ \mathfrak{p}^3), \]
$X^{(2)}$ is a divisor of $X^{(1)}$, hence $g_2$ is likewise a power of $p$, and so one must finally arrive at the unit group. Among the successive groups $X^{(1)}, X^{(2)}, \ldots$ several may under certain circumstances also be equal to one another\footnote{\gesperrt{Hilbert} calls $\Psi$ the decomposition group, $X$ the inertia group, $X^{(1)}$ the ramification group of the ideal $\mathfrak{p}$. In the memoir mentioned above it is further proved that $g_1$ is the highest power of $p$ dividing $g$, that $X^{(1)}$ is a normal divisor of $X$, and that the whole group $\Psi$ is metacyclic.}.


% ===== scan leaf 618, printed 597 =====
% !PRINTED 597
\weberpage{597}
\pdfbookmark[2]{§ 163. The Ideals in the Divisors of the Field Ω}{bm186}
According to this, the highest power of the prime ideal $\mathfrak{p}$ that divides the partial fundamental ideal $\mathfrak{G}_R$ can be given. By \S.\,159, namely, this ideal is defined by the product
\[ f'(\tau) = (\tau - \tau_1)(\tau - \tau_2)\ldots(\tau - \tau_{n-1}), \]
in which $\tau_1, \tau_2, \ldots, \tau_{n-1}$ are those among the functionals conjugate to $\tau$ that are different from $\tau$. One of the factors of this product, $\tau - \tau\mkern1mu|\mkern1mu\varphi$, is then and only then divisible by $\mathfrak{p}$, when $\varphi$ belongs to $X$. Now since the identical substitution is excluded, it follows that $\mathfrak{G}$ contains the factor $\mathfrak{p}$ at least $g-1$ times.

But $\tau - \tau\mkern1mu|\mkern1mu\chi$ is then and only then divisible by $\mathfrak{p}^2$, when $\chi$ belongs to $X^{(1)}$, and consequently $\mathfrak{G}$ is divisible by $\mathfrak{p}^{g-1+g_1-1}$. We can continue in this way and find the exact power of $\mathfrak{p}$ that is contained in $\mathfrak{G}$. It has the exponent
\[ (g-1) + (g_1-1) + (g_2-1) + \cdots \]
Each of the numbers $g, g_1, g_2, \ldots$ is a divisor of the one preceding, and $g_1, g_2, \ldots$ are powers of $p$, and the sequence continues until one of the numbers has become $g = 1$. The composition of the fundamental ideal $\mathfrak{G}$ is thus completely determined by the decomposition of the group $X$.

The numbers $g, g_1, g_2, \ldots$ are moreover the same for all the conjugate prime ideals $\mathfrak{p}_1, \mathfrak{p}_2, \ldots, \mathfrak{p}_e$, and it therefore follows that
\[ \mathfrak{G} = \Pi(\mathfrak{p}_1\mathfrak{p}_2\ldots\mathfrak{p}_e)^{g-1+g_1-1+g_2-2\cdots}, \] % [sic] exponent pattern breaks from g_2-1 to g_2-2 as printed
where the product sign refers to the prime factors of all the primes $p$ of the field $\Omega$. In this, of course, only those primes $p$, present in finite number, for which $g > 1$ need be taken into account.

The partial norm of the ideal $\mathfrak{G}$
\[ \mathfrak{D} = \mathfrak{N}_R(\mathfrak{G}) = \Pi\, \mathfrak{P}^{ef(g-1+g_1-1+g_2-1+\cdots)} \]
is called the partial discriminant of the field $\Omega$ with respect to the field $R$. It is an ideal lying in the field $R$, which, in the case that $R$ is the field of rational numbers, passes over, apart from a unit factor, into the fundamental number of the field $\Omega$ (\S.\,158).

\begin{center}
\S.\,163.

\gesperrt{The Ideals in the Divisors of the Field $\Omega$.}
\end{center}

If the group $\Phi$ of the field $\Omega$ has a divisor $\Phi'$, then to $\Phi'$ there belongs a field $\Omega'$, which is a divisor of $\Omega$ and
%CONT>>


% ===== scan leaf 619, printed 598 =====
% !PRINTED 598
\weberpage{598}
%<<CONT
in turn contains $R$ as a divisor. The numbers of $\Omega'$ remain
unchanged under the substitutions $\Phi'$, and $\Phi'$ is the group of the
field $\Omega$ with respect to $\Omega'$ (cf. Bd. I, \S.\,155). We denote
the degree of $\Phi'$, which is a divisor of $n$, by $n'$, and set

\weq{1}{n = m'\,n'.}

Then $n'$ is the degree of $\Omega$ with respect to $\Omega'$, and $m'$
the degree of $\Omega'$ with respect to $R$.

It is to be observed here that while $\Omega$ is a normal field
also with respect to $\Omega'$, yet $\Omega'$ will in general not be a
normal field with respect to $R$.

The prime ideals in the field $\Omega'$ can now be derived completely
from those of $\Omega$.

If $\mathfrak{p}$ is any prime ideal in $\Omega$, then there is \gesperrt{one and
only one ideal $\mathfrak{p}'$ in $\Omega'$ divisible by $\mathfrak{p}$}, and $\mathfrak{p}'$, as a
divisor of $\mathfrak{P}$, can be divisible by no ideals other than those
conjugate to $\mathfrak{p}$. If $\varphi'$ runs through the substitutions of $\Phi'$, then $\mathfrak{p}'$,
since it remains unchanged under $\varphi'$, is also divisible by $\mathfrak{p}\,|\,\varphi'$.

If therefore $\mathfrak{p}_1'$ is divisible by $\mathfrak{p}\,|\,\varphi_1$, it is also divisible by
$\mathfrak{p}\,|\,\psi\,\varphi_1\,\varphi'$, if $\psi$ denotes an arbitrary element of the group
$\boldsymbol{\Psi}$ (\S.\,162).

If, therefore, $\varphi$ is any element of the system

\[
\Phi_1 = \boldsymbol{\Psi}\,\varphi_1\,\Phi'
\]

then $\mathfrak{p}_1'$ is divisible by $\mathfrak{p}\,|\,\varphi$. Accordingly the
group decomposition that we set forth in \S.\,28 of this Band now
comes into consideration.

If $\varphi_2$ is an element of $\Phi$ not contained in $\Phi_1$, the
system

\[
\Phi_2 = \boldsymbol{\Psi}\,\varphi_2\,\Phi'
\]

has no element at all in common with $\Phi_1$. If there is yet another
element $\varphi_3$, which occurs neither in $\Phi_1$ nor in $\Phi_2$, one forms
in the same way

\[
\Phi_3 = \boldsymbol{\Psi}\,\varphi_3\,\Phi',
\]

and continues in this manner until the whole group $\Phi$ is exhausted. One
thus obtains the decomposition

\weq{2}{\Phi = \Phi_1 + \Phi_2 + \Phi_3 + \cdots + \Phi_{e'}.}

We can determine the degree of one of these systems $\Phi_r$ according
to \S.\,28. It is equal to the product of the degree of $\Phi$ multiplied by % [sic] source prints $\Phi$; the equivalent expression below requires $\boldsymbol\Psi$
the index of the intersection $\Psi_r'$ of $\Phi'$ with $\Psi_r = \varphi_r^{-1}\boldsymbol{\Psi}\varphi_r$ with
respect to $\Phi'$, or, what is the same thing, equal to the product of the
degrees of $\boldsymbol{\Psi}$ and $\Phi'$, divided by the
%CONT>>


% ===== scan leaf 620, printed 599 =====
% !PRINTED 599
\weberpage{599}
%<<CONT
degree $h_r$ of the intersection of $\Phi'$ and $\Psi_r$. The degree of $\Phi_r$
is therefore equal to $f g n' : h_r$.

If we now denote by $\varphi_r$ all the elements of $\Phi_r$, then all
the prime ideals $\mathfrak{p}\,|\,\varphi_r$ in $\Omega$ lead to the same prime ideal
$\mathfrak{p}_r'$ in $\Omega'$. But it remains to be shown that two different ones
of these systems $\Phi_1, \Phi_2$ lead to \gesperrt{different} prime ideals $\mathfrak{p}_1', \mathfrak{p}_2'$.

First of all it is evident that all the $\mathfrak{p}\,|\,\varphi_2$ are different from
the $\mathfrak{p}\,|\,\varphi_1$. For if $\mathfrak{p}\,|\,\varphi_1 = \mathfrak{p}\,|\,\varphi_2$, then also
$\mathfrak{p} = \mathfrak{p}\,|\,\varphi_2\,\varphi_1^{-1}$, so that $\varphi_2\,\varphi_1^{-1}$ is contained in $\boldsymbol{\Psi}$, and consequently
$\varphi_2$ in $\boldsymbol{\Psi}\varphi_1$, that is, in $\Phi_1$, contrary to the assumption. We
may therefore assume an integer $\alpha$ in $\Omega$ which is divisible by
$\mathfrak{p}\,|\,\varphi_1$, but by none of the ideals $\mathfrak{p}\,|\,\varphi_2$. Then none of the
numbers $\alpha\,|\,\varphi'$ is divisible by $\mathfrak{p}\,|\,\varphi_2$ either, since otherwise
$\mathfrak{p}\,|\,\varphi_1\,\varphi' = \mathfrak{p}\,|\,\varphi_2$, and hence, contrary to the assumption, $\varphi_2$ would be
contained in $\Phi_1$. The product $\alpha'$ of all the mutually different
$\alpha\,|\,\varphi'$, which is a number contained in $\Omega'$, is divisible by
$\mathfrak{p}_1'$, but not by $\mathfrak{p}_2'$. Consequently $\mathfrak{p}_1'$ is different from $\mathfrak{p}_2'$,
and it thus results that $e'$ and
no more \gesperrt{different} prime ideals $\mathfrak{p}'$ divide $\mathfrak{P}$.
We therefore set

\weq{3}{\mathfrak{P} = \mathfrak{p}_1'^{\,a_1}\,\mathfrak{p}_2'^{\,a_2}\cdots\mathfrak{p}_{e'}'^{\,a_{e'}},}

in which the exponents $a_1, a_2, \ldots, a_{e'}$ are natural numbers still
to be determined more closely.

Now, in order to decompose the prime ideals $\mathfrak{p}_r'$ in the field $\Omega$,
we can use the results of \S.\,160, 162 by simply letting $\Omega'$
take the place of $R$. If then $X_r'$ denotes the intersection
of $\Phi'$ with $X_r = \varphi_r^{-1}\,X\,\varphi_r$, then $X_r$ is the totality of all % [sic] source omits the prime; the next sentence correctly refers to $X_r'$
% [sic: The source says "Körpers" although $\Phi'$ denotes the group, not the field.]
substitutions of the field $\Phi'$ which satisfy the condition

\[
\tau\,|\,\chi_r' \equiv \tau \ (\mathrm{mod}\ \mathfrak{p}_r)
\]

If we therefore denote the degree of $X_r'$ by $g_r$, then
$\mathfrak{p}_r'$ is divisible by $\mathfrak{p}_r^{g_r}$, but by no higher power of $\mathfrak{p}_r$.

The intersection $\Psi_r'$ of $\Psi_r$ with $\Phi'$, whose degree we have
already denoted above by $h_r$, is the totality of all substitutions
$\psi_r'$ in $\Phi'$ which satisfy the condition

\[
\mathfrak{p}_r\,|\,\psi_r' = \mathfrak{p}_r
\]

and if we therefore decompose $\Phi'$ into the subsidiary groups

\weq{4}{\Phi' = \Psi_r'\varphi_{r,1}' + \Psi_r'\varphi_{r,2}' + \cdots + \Psi_r'\varphi_{r,e_r}} % [sic] source omits the prime on the final $\varphi$, unlike the two preceding terms

then

\weq{5}{n' = e_r\,h_r.}


% ===== scan leaf 621, printed 600 =====
% !PRINTED 600
\weberpage{600}
If now we set

\[
\mathfrak{p}_{r,s} = \mathfrak{p}_r\,|\,\varphi_{r,s}
\]

then by \S.\,160, (6),

\weq{6}{\mathfrak{p}_r' = (\mathfrak{p}_{r,1}\,\mathfrak{p}_{r,2}\cdots\mathfrak{p}_{r,e_r})^{g_r},}

and the substitution of (6) in (3) yields, through comparison
with \S.\,160, (6),

\weq{7}{g = a_r\,g_r, \quad e_1 + e_2 + \cdots + e_{e'} = e,}

whereby the exponents $a_r$ are defined.

The partial norm of $\mathfrak{p}_{r,s}$, taken with respect to the field $\Omega'$,
is by \S.\,160, (3), (7):

\weq{8}{\mathfrak{N}_{\Omega'}(\mathfrak{p}_{r,s}) = \mathfrak{p}_r'^{\,f_r},}

if $f_r$ is defined by the equation

\weq{9}{n' = e_r\,f_r\,g_r, \quad (h_r = f_r\,g_r).}

We now wish to decompose the group $\Phi$ into subsidiary groups
with respect to $\Phi'$, and, if $\vartheta_1, \vartheta_2, \ldots, \vartheta_{m'}$ are suitably
chosen substitutions from $\Phi$, we set:

\weq{10}{\Phi = \Phi'\vartheta_1 + \Phi'\vartheta_2 + \cdots + \Phi'\vartheta_{m'},}

so that, by the substitutions $\vartheta_1, \vartheta_2, \ldots, \vartheta_{m'}$, the field
$\Omega'$ passes over into $m'$ conjugate (equal or different) fields

\[
\Omega_1', \ \Omega_2', \ \ldots, \ \Omega_{m'}'.
\]

The field $\Omega_t'$ then belongs to the group $\vartheta_t^{-1}\Phi'\vartheta_t$.

If now, by $\vartheta_t$, the ideals $\mathfrak{p}_r'$ and $\mathfrak{p}_{r,s}$ pass over into $\mathfrak{p}_{r,t}'$ and
$\mathfrak{p}_{r,s,t}$, then by \S.\,160, (6)

\weq{11}{\mathfrak{p}_{r,t}' = (\mathfrak{p}_{r,1,t}\,\mathfrak{p}_{r,2,t}\cdots\mathfrak{p}_{r,e_r,t})^{g_r},}

and the product of all these ideals for $t = 1, 2, \ldots, m'$ is the
partial norm of $\mathfrak{p}_r'$ taken in the field $\Omega'$ with respect to the
field $R$, which we denote by $\mathfrak{N}_R'(\mathfrak{p}_r')$. It must be a power
of $\mathfrak{P}$, and we set

\weq{12}{\mathfrak{N}_R'(\mathfrak{p}_r') = \mathfrak{P}^{\,f_r'}.}

To find $f_r'$, we need only decompose $\mathfrak{P}$ in (12), according to \S.\,160, (6),
into its prime factors $\mathfrak{p}$, whereby there result $e\,g\,f_r'$ prime factors
$\mathfrak{p}$ in $\mathfrak{P}^{\,f_r'}$. If, on the other hand, we form the product
of the $m'$ factors (11), we obtain $g_r\,e_r\,m'$ such prime factors.
Hence

\[
e\,g\,f_r' = m'\,g_r\,e_r,
\]

and therefore, by (1), (9) and \S.\,160, (7):

\weq{13}{f = f_r\,f_r'.}

\begin{center}
\rule{0.3\linewidth}{0.4pt}
\end{center}


% ===== scan leaf 622, printed 601 =====
\weberpage{601}% page number not printed on this page
\pdfbookmark[1]{Twentieth Section. Quadratic Fields.}{bm187}
\pdfbookmark[2]{§ 164. Basis of a Quadratic Field}{bm188}
% !PRINTED none
\begin{center}
\gesperrt{Twentieth Chapter.}

\gesperrt{Quadratic Fields.}
\end{center}

\begin{center}
\rule{0.2\linewidth}{0.4pt}
\end{center}

\begin{center}
\S.\,164.

\gesperrt{Basis of a quadratic field.}
\end{center}

We now wish to apply the results of the general theory of algebraic numbers to the particular case of the fields of degree $2$,
which are also called quadratic fields. A quadratic field $\Omega$
springs from a quadratic equation with integral coefficients, and is
thus derived from the field of the rational numbers by the adjunction
of a square root $\sqrt{d}$. Here $d$ may be assumed to be a positive
or negative integer without square divisor, and the single value
$d = 1$ is to be excluded. According as $d$ is positive or negative,
we obtain a real or an imaginary quadratic field.

We obtain the conjugate field if in all the numbers in $\Omega$ we
change $\sqrt{d}$ into $-\sqrt{d}$. By this, however, we obtain no new
field. The quadratic field is therefore an (absolute) normal field
(\S.\,160).

The numbers of the field $\Omega$ are none other than the quadratic
irrational numbers defined in \S.\,122 of the first Band, which are
all contained in the form $x + y\sqrt{d}$, where $x$ and $y$ denote
integral or fractional rational numbers. Every irrational number
$\omega$ of this field satisfies a quadratic equation

\weq{1}{c\,\omega^2 = a + b\,\omega,}

in which $a, b, c$ are integral rational numbers without common
divisor, of which $c$ may be assumed to be positive, and the
combination

\weq{2}{D = b^2 + 4\,a\,c}
%CONT>>


% ===== scan leaf 623, printed 602 =====
% !PRINTED 602
\weberpage{602}
%<<CONT
we have called the discriminant of the number $\omega$. The ratio $D:d$ is an integral rational square number. The two roots of (1) receive the expression

\weq{3}{\omega = \frac{b+\sqrt{D}}{2c}, \quad \omega' = \frac{b-\sqrt{D}}{2c}.}
If we set
\weq{4}{D = e^2 d,}
then
\weq{5}{\omega = \frac{b+e\sqrt{d}}{2c}, \quad \omega' = \frac{b-e\sqrt{d}}{2c},}
and on account of (1) these are then and only then integers, when $c=1$.

Now $D\equiv 0$ or $\equiv 1\ (\mathrm{mod}\ 4)$, and if, therefore, $D$ is even, it is divisible by 4; $b$ is in this case even, and since $d$ cannot be divisible by 4, $e$ must also be even.

But if $D$ is odd, then $b$ and $e$ are also odd. This case, however, as (4) shows, can occur only when $d\equiv 1\ (\mathrm{mod}\ 4)$.

From this we obtain the following result:

1. \gesperrt{If $d\equiv 2$ or $\equiv 3\ (\mathrm{mod}\ 4)$, the number}
\[ \omega = x + y\sqrt{d} \]
\gesperrt{is then and only then an integer, when the rational numbers $x, y$ are integers.} (Quadratic irrational numbers of the first kind, after \gesperrt{Dirichlet}.)

2. \gesperrt{If $d\equiv 1\ (\mathrm{mod}\ 4)$, then}
\[ \omega = \frac{x+y\sqrt{d}}{2} \]
\gesperrt{is then and only then an integer, when the rational numbers $x,y$ are either two even or two odd integers.} (Quadratic irrational numbers of the first or second kind, according as $e$ is even or odd.)

We can also express the latter thus:

3. \gesperrt{If $d\equiv 1\ (\mathrm{mod}\ 4)$, then}
\[ \omega = x + y\,\frac{1+\sqrt{d}}{2} \]
\gesperrt{is then and only then an integer, when the rational numbers $x$ and $y$ are integers.}


% ===== scan leaf 624, printed 603 =====
% !PRINTED 603
\weberpage{603}
\pdfbookmark[2]{§ 165. The Prime Ideals in Ω}{bm189}
From this we obtain a basis of $\mathfrak{o}$:

a) \quad $1, \sqrt{d},$ \hfill $d\equiv 2$ or $\equiv 3\ (\mathrm{mod}\ 4),$

b) \quad $1, \dfrac{1+\sqrt{d}}{2},$ \hfill $d\equiv 1\ (\mathrm{mod}\ 4).$

If we therefore set

\weq{6}{\begin{aligned}
\text{a)}\quad \Theta &= \sqrt{d}, & d &\equiv 2, 3\ (\mathrm{mod}\ 4)\\
\text{b)}\quad \Theta &= \frac{1+\sqrt{d}}{2}, & d &\equiv 1\ (\mathrm{mod}\ 4)
\end{aligned}}
then $1, \Theta$ is a basis of $\mathfrak{o}$.

For the fundamental number we accordingly obtain, if $\Theta'$ is conjugate to $\Theta$,
\[ \Delta = \left|\begin{array}{cc} 1, & \Theta \\ 1, & \Theta' \end{array}\right|^2 = (\Theta'-\Theta)^2, \]
hence

\weq{7}{\begin{aligned}
\text{a)}\quad \Delta &= 4d, & d &\equiv 2, 3\ (\mathrm{mod}\ 4)\\
\text{b)}\quad \Delta &= d, & d &\equiv 1\ (\mathrm{mod}\ 4),
\end{aligned}}
\weq{8}{\Theta - \Theta' = \sqrt{\Delta}.}

The fundamental number is therefore $\equiv 0$ or $\equiv 1\ (\mathrm{mod}\ 4)$, and, if it is divisible by 4, must be $\equiv 8, 12\ (\mathrm{mod}\ 16)$\footnote{An integral rational number $\Delta$ that is congruent to 0 or to 1 with respect to the modulus 4 has been called by \gesperrt{Kronecker} a number of discriminant form. If $\Delta$ has the property that no square factor can be separated off from it in such a way that a number of discriminant form remains, then $\Delta$ is called a \gesperrt{stem discriminant}. The fundamental number of a quadratic field is therefore always a \gesperrt{stem discriminant}, and conversely every stem discriminant is also the fundamental number of a quadratic field. Cf. \gesperrt{H. Weber}, Zahlentheoretische Untersuchungen aus dem Gebiete der elliptischen Functionen. I, II, III. Nachrichten der Gesellschaft der Wissenschaften zu G\"ottingen. 1893.}. It can therefore (in agreement with the theorem of \gesperrt{Minkowski}) be neither $=\pm 1$ nor $=\pm 2$. The fundamental number with the smallest absolute value is $\Delta = -3.$

\begin{center}
\S.\,165.\\
\gesperrt{The Prime Ideals in $\Omega$.}
\end{center}

Now let $\varphi$ denote an arbitrary integral functional in $\Omega$. Following \S.\,146, we wish to find a basis $\alpha_1, \alpha_2$ for $\varphi$, and with it the
%CONT>>


% ===== scan leaf 625, printed 604 =====
% !PRINTED 604
\weberpage{604}
%<<CONT
corresponding basis form $\alpha_1 t_1 + \alpha_2 t_2$ that is associated with $\varphi$. By \S.\,146, (3), if $\omega_1, \omega_2$ is a basis of $\mathfrak{o}$, we have to determine the integers $a_{1,1}, a_{1,2}, a_{2,2}$ in such a way that $a_{1,1}, a_{2,2}$ are positive and as small as possible, and that
\[ \alpha_1 = a_{1,1}\omega_1, \quad \alpha_2 = a_{1,2}\omega_1 + a_{2,2}\omega_2 \]
become divisible by $\varphi$. If we therefore set $(\omega_1, \omega_2) = (1, \Theta)$, then

\weq{1}{\alpha_1 = a_{1,1}, \quad \alpha_2 = a_{1,2} + a_{2,2}\Theta.}

Thus $a_{1,1}$ is the least positive rational number divisible by $\varphi$, and $a_{2,2}$ is the least positive number for which $a_{2,2}\Theta$ becomes congruent to a rational number according to the modulus $\varphi$.

We wish to apply this to the case where $\varphi$ is a prime functional $\pi$. Such a functional can here be only of the first or of the second degree, according as its absolute norm is equal to $p$ or $=p^2$. We therefore distinguish the two cases:

1) \quad $N_a(\pi) = p$

or

2) \quad $N_a(\pi) = p^2.$

Since here $N(p) = p^2$, in the first case $p$ is decomposable into two prime factors, while in the second case $p$ contains only one prime factor; that is, $p$ is still to be regarded as a prime in the field $\Omega$ itself.

If we now determine the basis (1) for $\varphi = \pi$, we obtain first of all $a_{1,1} = p$, and $a_{2,2}$ is the least positive integral rational number for which the product $a_{2,2}\Theta$ becomes congruent to a rational number modulo $\pi$. But for this, by \S.\,150, 3, the necessary and sufficient condition is
\[ (a_{2,2}\Theta)^p \equiv a_{2,2}\Theta \pmod{\pi}, \]
or, since by \gesperrt{Fermat}'s theorem $a_{2,2}^p \equiv a_{2,2}$,

\weq{2}{a_{2,2}(\Theta^p - \Theta) \equiv 0 \pmod{\pi}.}

Accordingly either $\Theta^p - \Theta \equiv 0$, and then $a_{2,2} = 1$, or $\Theta^p - \Theta$ is not $\equiv 0$, and then $a_{2,2} = p$.

In the first case $\Theta$, and with it every number $\omega$, is congruent to a rational number. There are therefore no more than $p$ numbers incongruent according to the modulus $\pi$, and consequently (\S.\,148, 3.) $p^f = p$, $f = 1$, that is, $p$ is decomposable into two prime factors of the first degree.


% ===== scan leaf 626, printed 605 =====
% !PRINTED 605
\weberpage{605}
Conversely, in this case 1) the $p$ rational numbers $0, 1, 2, \ldots, p-1$ form a complete system of residues with respect to the modulus $\pi$, and consequently every number in $\mathfrak{o}$, and hence also $\Theta$, is congruent with a rational number. Thus we have

1) \quad $N_a(\pi) = p$ \quad $a_{2,2} = 1$

2) \quad $N_a(\pi) = p^2$ \quad $a_{2,2} = p.$

In the first case $a_{1,2}$ is to be determined from the condition

\weq{3}{a_{1,2} + \Theta \equiv 0 \pmod{\pi}}
while in the second $a_{1,2} = 0$ may be taken. The basis form of $\pi$ thus results, if $u, v$ are the variables, as

1) \quad $\pi = pu + (a_{1,2}+\Theta)v$

2) \quad $\pi = p(u+\Theta v),$

from which one sees that $\pi$ in case 2), as was to be expected, is a \gesperrt{principal functional}.

In order to determine the number $a_{1,2}$ in case 1), we form the norm of $a_{1,2}+\Theta$, which must be divisible by $p$. If $\Theta$ passes over into $\Theta'$ through a change of the sign of $\sqrt{d}$, then by \S.\,164, (8) $\Theta - \Theta' = \sqrt{\Delta}$, and for the norm of $a_{1,2}+\Theta$ we obtain by \S.\,164 (6), a), b):

\weq{4}{N(a_{1,2}+\Theta) = (a_{1,2}+\Theta)(a_{1,2}+\Theta')}
\[ {} = a_{1,2}^2 - d, \quad d\equiv 2, 3\ (\mathrm{mod}\ 4) \]
\[ {} = a_{1,2}^2 + a_{1,2} + \frac{1-d}{4}, \quad d\equiv 1\ (\mathrm{mod}\ 4). \]

If this number is divisible by $p$, then one of the two factors $a_{1,2}+\Theta$, $a_{1,2}+\Theta'$ is divisible by $\pi$. But since $\Theta' = -\Theta$ or $=1-\Theta$, and therefore

a) \quad $a_{1,2} + \Theta' = -(-a_{1,2}+\Theta)$

b) \quad $a_{1,2} + \Theta' = -(-a_{1,2}-1+\Theta)$

one can always take $a_{1,2}$ in such a way that precisely $a_{1,2}+\Theta$ becomes divisible by $\pi$, and $\Theta$ is thus actually congruent with a rational number. The necessary and sufficient condition for $p$ to contain two prime factors is consequently that $a_{1,2}^2 - d$ or $a_{1,2}^2 + a_{1,2} + \frac14(1-d)$ becomes divisible by $p$ by a suitable assumption about $a_{1,2}$. By multiplication by 4 these two conditions pass over into
\[ 4a_{1,2}^2 - 4d \equiv 0, \quad (2a_{1,2}+1)^2 - d \equiv 0 \pmod{4p}, \]
%CONT>>


% ===== scan leaf 627, printed 606 =====
% !PRINTED 606
\weberpage{606}
%<<CONT
and when one considers that the fundamental number $\Delta$ is equal to $4d$ in case a), and equal to $d$ in case b), then, if we set
\weq{5}{x = 2a_{1,2} \text{ or } {}= 2a_{1,2}+1}
we obtain the following result:

1. \gesperrt{The necessary and sufficient condition for $p$ to contain two prime factors is that the congruence}
\weq{6}{x^2 \equiv \Delta\ (\mathrm{mod}\ 4p)}
\gesperrt{can be satisfied by an integral rational number $x$, or that the fundamental number $\Delta$ is a quadratic residue of $4p$.}

If we denote by $\pi, \pi'$ the two conjugate forms
\weq{7}{\pi = pu + (a_{1,2}+\Theta)v,\quad \pi' = pu + (a_{1,2}+\Theta')v,}
then multiplication yields
\weq{8}{\pi\pi' = p\left(pu^2 + xuv + \frac{x^2-\Delta}{4p}v^2\right).}

Now in this,
\[
pu^2 + xuv + \frac{x^2-\Delta}{4p}v^2
\]
is a unit, because, if $p$ does not divide $\Delta$, $x$ is not divisible by $p$, and if $p$ does divide $\Delta$, then $x$ is indeed divisible by $p$, but $(x^2-\Delta:4p)$ is not. For if $p$ is odd, $p^2$ cannot divide $\Delta$; and if $p=2$ and $\Delta$ is divisible by 2, then $\Delta = 4d$, $d \equiv 2, 3\ (\mathrm{mod}\ 4)$. But if we had
\[
x^2 - \Delta \equiv 0\ (\mathrm{mod}\ 16),
\]
it would follow that
\[
d \equiv \left(\frac{x}{2}\right)^2 (\mathrm{mod}\ 4),
\]
whereas $\left(\dfrac{x}{2}\right)^2$, being a square, cannot be congruent to 2 or 3. Accordingly $p$ is decomposed into the two prime factors $\pi$ and $\pi'$.

The two factors $\pi$ and $\pi'$ of $p$ can also be associated with one another. This occurs only when $\pi'$, and consequently also $\pi - \pi'$, is divisible by $\pi$. But now by (7)
\[
\pi - \pi' = (\Theta - \Theta')v,
\]
and consequently
\[
(\Theta - \Theta')^2 = \Delta
\]
%CONT>>


% ===== scan leaf 628, printed 607 =====
% !PRINTED 607
\weberpage{607}
\pdfbookmark[2]{§ 166. Functionals in the Quadratic Field and Quadratic Irrational Numbers}{bm190}
%<<CONT
be divisible by $p$, and this is also the sufficient condition.

2. \gesperrt{Thus $p$ is only then associated with the square of a prime functional, when $p$ divides the fundamental number $\Delta$, in agreement with the general theorem (\S.\,158).}

\begin{center}
\S.\,166.

\gesperrt{Functionals in the quadratic field and quadratic irrational numbers.}
\end{center}

We now consider in our quadratic field an arbitrary integral functional and form its basis form
\[
\varphi = \alpha_1 t_1 + \alpha_2 t_2,
\]
by determining $\alpha_1, \alpha_2$ according to \S.\,165:
\weq{1}{\alpha_1 = a_{1,1}, \quad \alpha_2 = a_{1,2} + a_{2,2}\Theta.}
According to \S.\,146 we then have
\weq{2}{N_a(\varphi) = a_{1,1}a_{2,2}.}
If we give the form $\varphi$ the expression
\weq{3}{\varphi = a_{1,1}t_1 + (a_{1,2}+a_{2,2}\Theta)t_2,}
and form the norm, there results
\weq{4}{N(\varphi) = a_{1,1}^2 t_1^2 + [2a_{1,1}a_{1,2} + a_{1,1}a_{2,2}(\Theta+\Theta')]t_1 t_2 + t_2^2 N(\alpha_2)}
and the divisor of this form is therefore, according to (2), equal to $a_{1,1}a_{2,2}$.
If therefore we set
\weq{5}{\begin{aligned}
a_{1,1}^2 &= c\,a_{1,1}a_{2,2}\\
2a_{1,1}a_{1,2}+a_{1,1}a_{2,2}(\Theta+\Theta') &= b\,a_{1,1}a_{2,2}\\
N(\alpha_2) &= -a\,a_{1,1}a_{2,2}
\end{aligned}}
then $a, b, c$ are integral rational numbers without common divisor, and the quotient
\weq{6}{\frac{\alpha_2}{\alpha_1} = \omega}
is a root of the quadratic equation
\weq{7}{c\omega^2 = a + b\omega.}
% [sic: The source prints the incomplete subscript $a_{1,}$ here; equations (1) and (5) identify the intended quantity as $a_{1,2}$.]
If we substitute in the expressions (1) the values of $a_{1,1}$ and $a_{1,}$ from (5), it follows that
\weq{8}{\begin{aligned}
2\alpha_2 &= a_{2,2}(b+\Theta-\Theta')\\
\alpha_1 &= c\,a_{2,2}.
\end{aligned}}


% ===== scan leaf 629, printed 608 =====
% !PRINTED 608
\weberpage{608}
Now by \S.\,164, (8) the difference $\Theta - \Theta'$ is equal to the square root of the fundamental number of $\Omega$, so that from (6) and (8) we obtain:
\weq{9}{\omega = \frac{b+\sqrt{\Delta}}{2c}, \quad \Delta = b^2 + 4ac,}
wherein the sign of $\sqrt{\Delta}$ is determined by the assumption made in \S.\,164, (6) about the sign of $\sqrt{d}$.

Thus $\omega$ is a \gesperrt{quadratic irrational number} belonging to the discriminant $\Delta$, such as we have considered in \S.\,122 of the first Band, and the form $\varphi$ can, if we set $a_{2,2}=e$, be expressed thus:
\weq{10}{\varphi = ec(t_1+\omega t_2).}

Herein $c$ is the least positive integral rational number for which the product $c\omega$ becomes an integer, and every other integral rational number $k$ by which the product $k\omega$ becomes an integer must be divisible by $c$.

One recognises this most simply if, by \S.\,164, one puts $\omega$ in one of the two forms, according as $b$ is even or odd:
\weq{11}{\omega = \frac{1}{c}\left(\frac{b}{2}+\Theta\right), \quad \frac{1}{c}\left(\frac{b-1}{2}+\Theta\right),}
from which, since $1, \Theta$ is a basis of $\mathfrak{o}$, what has been said follows.

On this is now founded the proof that, if $\omega$ denotes any quadratic irrational number belonging to the discriminant $\Delta$, $c$ the least positive integral rational number that makes $c\omega$ an integer, and $e$ an arbitrary positive integral rational number, then conversely, by
\weq{12}{\varphi = ec(t_1+\omega t_2)}
a basis form is always represented.

For this we have only to show that
\[
\alpha_1 = ec, \quad \alpha_2 = ec\omega
\]
is a basis of the functional $\varphi$ defined by (12).

If for this purpose we first form the norm of $\varphi$, we obtain by (7)
\weq{13}{N(\varphi) = e^2 c(ct_1^2 + bt_1t_2 - at_2^2),}
and consequently the absolute norm
\weq{14}{N_a(\varphi) = e^2 c.}


% ===== scan leaf 630, printed 609 =====
% !PRINTED 609
\weberpage{609}
Now if
\weq{15}{\alpha_1 = a_{1,1}, \quad \alpha_2 = a_{1,2}+a_{2,2}\Theta}
is a basis of $\varphi$, then on the other hand (by \S.\,146)
\weq{16}{N_a(\varphi) = a_{1,1}a_{2,2},}
and $a_{1,1}$ is the least positive integral rational number that is divisible by $\varphi$.

Now we can show, on the other hand, that this number is $= ec$, from which $ec = a_{1,1}$ follows, and, by (14) and (16), $e = a_{2,2}$.

For if an integral rational number $m$ is to be divisible by $\varphi$, then, if $\omega'$ is conjugate to $\omega$, since
\[
ct_1^2 + bt_1t_2 - at_2^2
\]
is a unit functional,
\[
\frac{m(ct_1^2+bt_1t_2-at_2^2)}{\varphi} = \frac{m}{e}(t_1+\omega' t_2)
\]
must be an integral functional. Consequently
\[
\frac{m}{e} \quad \text{and} \quad \frac{m}{e}\omega'
\]
must be integers, from which it follows that $m : e$ is an integral rational number divisible by $c$, and hence $m$ one divisible by $ec$. Since, on the other hand, $ec$ is divisible by $\varphi$, it follows that $ec$ is the least number satisfying these requirements. Moreover, it also results that the integer $ec\omega$ is divisible by $\varphi$. For in the representation (12), $ec\omega$ is a coefficient of the integral function $\varphi$, and is therefore divisible by $\varphi$ (\S.\,142).

If we express $\Theta$ in (15) in terms of $\omega$, then by (11) we obtain
\[
\begin{aligned}
\alpha_2 &= a_{1,2} - \frac{eb}{2} + ec\omega, \quad b \text{ even}\\
&= a_{1,2} - \frac{e(b-1)}{2} + ec\omega, \quad b \text{ odd,}
\end{aligned}
\]
and thus
\[
a_{1,2} \equiv \frac{eb}{2} \quad \text{or} \quad \equiv \frac{e(b-1)}{2}
\]
with respect to the modulus $\varphi$, and consequently, since both sides are rational, also with respect to the modulus $a_{1,1}=ec$. Accordingly $a_{1,2}$ may be set equal to the value $\frac{1}{2}eb$ or $\frac{1}{2}e(b-1)$, and we obtain $\alpha_2 = ec\omega$, which was to be proved.

In order to give a clear summary of what has hereby been proved, we may state the following theorems:
%CONT>>


% ===== scan leaf 631, printed 610 =====
% !PRINTED 610
\weberpage{610}
\pdfbookmark[2]{§ 167. Equivalent Functionals and Equivalent Numbers}{bm191}
%<<CONT

1. \gesperrt{If $\Omega$ is a quadratic field with the fundamental number $\varDelta$, and $\varphi$ an integral functional in this field, then a quadratic irrational number belonging to the discriminant $\varDelta$}

\weq{17}{\omega = \frac{b + \sqrt{\varDelta}}{2c}}

\gesperrt{and a positive integral rational number $e$ can be so determined that}

\weq{18}{ec(t_1 + \omega t_2)}

\gesperrt{becomes a basis form of $\varphi$. Conversely, if $\omega$ is any quadratic irrational number belonging to the discriminant $\varDelta$, the form (18) is always a basis form of a functional $\varphi$.}

\begin{center}
§. 167.

\gesperrt{Equivalent Functionals and Equivalent Numbers.}
\end{center}

Now if $\omega, \omega_1$ are any two quadratic irrational numbers belonging to the discriminant $\varDelta$, we obtain two basis forms

\weq{1}{\begin{aligned} \varphi &= ec(t_1+\omega t_2)\\ \varphi_1 &= e_1c_1(t_1+\omega_1 t_2). \end{aligned}}

If now the two numbers $\omega, \omega_1$ are equivalent in the sense in which we stipulated this concept in \S.\,120 of the first Band, then

\weq{2}{\omega_1 = \frac{p\omega+q}{r\omega+s},}

in which $p,q,r,s$ are integral rational numbers which satisfy the condition

\weq{3}{ps - qr = \pm 1.}

If we then make the substitution (2) in $\varphi_1$, it follows that

\[ (r\omega+s)\varphi_1 = e_1c_1[st_1+qt_2+\omega(rt_1+pt_2)], \]

or, if we set

\weq{4}{u_1 = st_1+qt_2,\quad u_2 = rt_1+pt_2,}

then

\[ (r\omega+s)\varphi_1 = e_1c_1(u_1+\omega u_2). \]

Now $ec(u_1+\omega u_2)$ is associated with $\varphi$, because (4) is an
% German breaks mid-word at the page end: "eine ganz-" (word completed on the next page)
%CONT>>


% ===== scan leaf 632, printed 611 =====
% !PRINTED 611
\weberpage{611}

%<<CONT
integral linear substitution with the determinant $\pm 1$ (\S.\,147), and consequently the two forms

\[ (r\omega+s)\,ec\,\varphi_1,\quad e_1c_1\varphi \]

are associated with one another, and this has the consequence that the two forms $\varphi,\varphi_1$ are equivalent in the sense defined in \S.\,152. With this we have proved:

2. \gesperrt{If the quadratic irrational numbers $\omega,\omega_1$ are equivalent, then the corresponding functionals $\varphi,\varphi_1$ are also equivalent.}

But the converse theorem can also be proved:

3. \gesperrt{If the functionals $\varphi,\varphi_1$ are equivalent, then the quadratic irrational numbers $\omega$ and $\omega_1$ are also equivalent.}

To carry out this proof, we denote by $\psi$ a functional of the class that is reciprocal to the class of $\varphi$, so that $\varphi\psi$ is equivalent to a number. Then we can set:

\weq{5}{\varphi\psi = \varepsilon\xi,}

in which $\varepsilon$ is a unit and $\xi$ a \gesperrt{number} in $\mathfrak{o}$. Herein we can assume $\varphi$ and $\psi$ both to be basis forms, and accordingly set

\weq{6}{\varphi = ec(t_1+t_2\omega),\quad \psi = e'c'(t_1+t_2\eta)}

in which $e', c', \eta$ have the same meaning for $\psi$ as $e,c,\omega$ have for $\varphi$. For abbreviation we set [\S.\,166, (13)]:

\weq{7}{N_a(\varphi) = e^2c = P,\quad N_a(\psi) = e'^2c' = Q.}

If we denote by $\psi'$ the form conjugate to $\psi$, and for the moment set the unit form

\[ c't_1^2 + b't_1t_2 - a't_2^2 = \varepsilon_1, \]

then it follows from (5) by multiplication with $\psi'$:

\[ Q\varepsilon_1\varphi = \varepsilon\xi\psi', \]

and consequently

\weq{8}{\frac{Q\varphi}{\xi} = \frac{\varepsilon}{\varepsilon_1}\psi' = \chi,}

in which $\chi$ and $\psi'$ are associated. If we set

\weq{9}{\frac{Qec}{\xi} = \beta_1,\quad \frac{Qec\omega}{\xi} = \beta_2,}
%CONT>>


% ===== scan leaf 633, printed 612 =====
% !PRINTED 612
\weberpage{612}

%<<CONT
then

\[ \chi = \frac{Q\varphi}{\xi} = \beta_1t_1+\beta_2t_2, \]

and consequently $\beta_1,\beta_2$ are integers divisible by $\psi'$.

If we form from (9) the discriminant $\varDelta(\beta_1,\beta_2)$, we obtain

\weq{10}{\varDelta(\beta_1,\beta_2) = \begin{vmatrix}\beta_1, & \beta_2\\ \beta_1', & \beta_2'\end{vmatrix}^2 = \frac{e^4c^4Q^4}{(\xi\xi')^2}(\omega-\omega')^2.}

But from (5) and (7) it follows that

\[ \xi\xi' = N(\xi) = PQ = e^2cQ, \]

and if we further set $\omega-\omega' = \dfrac{\sqrt{\varDelta}}{c}$, we find

\weq{11}{\varDelta(\beta_1,\beta_2) = Q^2\varDelta.}

But from this it follows, by the theorem of \S.\,147, 2., that $\beta_1,\beta_2$ is a basis of $\psi'$.

On the other hand, by (6), $e'c', e'c'\eta'$ is also a basis of $\psi'$, and consequently the integral rational numbers $p,q,r,s$ can be so determined that

\weq{12}{\begin{aligned}
\beta_1 &= \frac{Qec}{\xi} = e'c'(r\eta'+s)\\
\beta_2 &= \frac{Qec}{\xi}\,\omega = e'c'(p\eta'+q),
\end{aligned}}

and at the same time

\weq{13}{ps-qr = \pm 1.}

But from (12) it follows by division:

\weq{14}{\omega = \frac{p\eta'+q}{r\eta'+s},}

that is, $\omega$ is equivalent to $\eta'$.

If we now replace $\varphi$ by the equivalent form $\varphi_1$, there subsists an equation like (5):

\[ \varphi_1\psi = \varepsilon_1\xi_1, \]

in which $\psi$ has remained unchanged. Consequently $\omega_1$ is also equivalent to $\eta$, and hence $\omega$ and $\omega_1$ are equivalent to one another. % [sic] source prints $\eta$; repeating the derivation with unchanged $\psi$ requires $\eta'$
But this is the assertion expressed in our Theorem 3., which is therewith proved.

According to this, then, to every ideal class of the field $\Omega$ there corresponds a class of equivalent numbers $\omega$, and conversely; and we may therefore add the further theorem:
%CONT>>


% ===== scan leaf 634, printed 613 =====
% !PRINTED 613
\weberpage{613}
\pdfbookmark[2]{§ 168. Composition of the Quadratic Irrational Numbers}{bm192}
%<<CONT

4. \gesperrt{The number of non-equivalent irrational numbers of the discriminant $\varDelta$ is equal to the number of classes of forms of the field $\Omega$}\footnote{According to the remark made at the close of \S.\,128 of the first volume on the \gesperrt{Gauss} theory of quadratic forms, the class number of the quadratic forms according to the \gesperrt{Gauss} definition does not always agree with the number of ideal classes of the field $\Omega$. To establish the agreement, one must in some cases divide the ideal classes further into two classes each. In our theory of the equivalence of the numbers $\omega$, where proper equivalence is not distinguished from improper, this further division is not required. Cf. \gesperrt{Dirichlet-Dedekind}, Zahlentheorie. 4. Auflage, S. 578, 584.}\gesperrt{.}

\begin{center}
§. 168.

\gesperrt{Composition of the quadratic irrational numbers.}
\end{center}

We have seen in \S.\,166 that for every integral functional $\varphi$ in the field $\Omega$ a certain quadratic irrational number $\omega$ of the discriminant $\varDelta$ can be found.

But we must now still investigate how far this number $\omega$ is determined by the field $\Omega$ and by the functional $\varphi$.

By \S.\,166, (1), the basis of $\varphi$

\weq{1}{\alpha_1 = a_{1,1},\quad \alpha_2 = a_{1,2}+a_{2,2}\Theta}

is defined by the condition that $a_{1,1}$ is the least positive integral rational number that is divisible by $\varphi$, and $a_{2,2}$ is then determined by

\weq{2}{N_a(\varphi) = a_{1,1}a_{2,2}}

and thus the two numbers $a_{1,1},a_{2,2}$ are uniquely determined by $\varphi$. It remains to ask how far $a_{1,2}$ is also determined. But the number $a_{1,2}$ is given, with respect to the modulus $\varphi$, by the congruence $a_{1,2}+a_{2,2}\Theta\equiv 0$, and since $a_{1,2}$ must also be an integral rational number, $a_{1,2}$ is determined up to a multiple of $a_{1,1}$. This multiple of $a_{1,1}$ remains in fact arbitrary, since, if $\alpha_1,\alpha_2$ is a basis of $\varphi$, the same holds, for every integral rational $k$, of $\alpha_1,\alpha_2+k\alpha_1$.

Now since $\alpha_2:\alpha_1=\omega$ is the quadratic irrational number sought, we obtain the theorem:
%CONT>>


% ===== scan leaf 635, printed 614 =====
% !PRINTED 614
\weberpage{614}
%<<CONT

5. \gesperrt{To every integral functional of the field $\Omega$ there belongs one and only one sequence of quadratic irrational numbers of the discriminant $\Delta$, and the numbers of this sequence differ by arbitrary integral rational numbers and are consequently equivalent to one another. Associated forms lead to the same sequence.}

If $\varphi_1$, $\varphi_2$ denote two integral functionals in $\Omega$, and we put
\[
\varphi_1 \varphi_2 = \varphi_3,
\]
then we can determine for each of these three forms a quadratic irrational number $\omega_1$, $\omega_2$, $\omega_3$, and $\omega_3$ is determined by $\omega_1$, $\omega_2$ up to an additive integral rational number. One then calls $\omega_3$ \gesperrt{compounded} or composed from $\omega_1$ and $\omega_2$, and sets, in symbolic notation,
\weq{3}{\omega_1 \omega_2 = \omega_3.}

If $\psi_1$ is equivalent to $\varphi_1$, and $\psi_2$ to $\varphi_2$, then by \S.\,152 $\psi_1 \psi_2 = \psi_3$ is also equivalent to $\varphi_3$.

Now if $\eta_1$, $\eta_2$, $\eta_3$ are the irrational numbers associated by Theorem 5. with the forms $\psi_1$, $\psi_2$, $\psi_3$, then by 3., \S.\,167 the numbers $\eta_1$, $\omega_1$; $\eta_2$, $\omega_2$; $\eta_3$, $\omega_3$ form three pairs of equivalent numbers, and we can therefore state the following theorem:

6. \gesperrt{If in a composition of quadratic irrational numbers each element is replaced by an equivalent number, then the compounded number also passes into an equivalent one.}

In this way the classes of the quadratic irrational numbers of a discriminant $\Delta$ are combined into an \gesperrt{Abelian} group isomorphic with the group of ideal classes.

In order actually to find the number compounded from two given numbers, one has to proceed as follows.

We shall denote the bases, formed according to (1), of the forms $\varphi_1$, $\varphi_2$, $\varphi_3$ as follows:
\[
\begin{array}{lll}
\varphi_1) & \alpha_1 = a_{1,1} & \alpha_2 = a_{1,2} + a_{2,2}\,\Theta \\
\varphi_2) & \beta_1 = b_{1,1} & \beta_2 = b_{1,2} + b_{2,2}\,\Theta \\
\varphi_3) & \gamma_1 = c_{1,1} & \gamma_2 = c_{1,2} + c_{2,2}\,\Theta,
\end{array}
\]
so that
\[
\omega_1 = \frac{\alpha_2}{\alpha_1}, \quad \omega_2 = \frac{\beta_2}{\beta_1}, \quad \omega_3 = \frac{\gamma_2}{\gamma_1}
\]
%CONT>>


% ===== scan leaf 636, printed 615 =====
% !PRINTED 615
\weberpage{615}

%<<CONT
is formed. Then $c_{1,1}$ is the least positive integral rational number which is divisible by $\varphi_1 \varphi_2$, and this must in any case be a divisor of $a_{1,1} b_{1,1}$. If we have
\[
a_{1,1} b_{1,1} = k\, c_{1,1},
\]
where $k$ is a positive integral rational number, then, since $N_a(\varphi_1)\, N_a(\varphi_2) = N_a(\varphi_3)$ must hold,
\[
k\, a_{2,2} b_{2,2} = c_{2,2}.
\]
The last number $c_{1,2}$ is then to be determined from the congruence
\weq{4}{c_{1,2} + c_{2,2}\,\Theta \equiv 0 \ (\mathrm{mod}\ \varphi_1 \varphi_2).}

If it is only a matter of finding a representative of the composed class, one proceeds most simply as follows.

If $\varphi_1$ is assumed arbitrarily, one can (by \S.\,152, 6.) choose $\varphi_2$ in its class so that it becomes relatively prime to $a_{1,1}$. Then $\varphi_1$ is also relatively prime to $\varphi_2$, and $b_{1,1}$ relatively prime to $a_{1,1}$.

For if $a_{1,1}$ and $b_{1,1}$ have a greatest common divisor $d$, then this is also relatively prime to $\varphi_2$, and hence $b_{1,1} : d$ is divisible by $\varphi_2$; but if $d > 1$, this contradicts the definition of $b_{1,1}$. In the same way it also follows conversely that, if $a_{1,1}$ and $b_{1,1}$ are relatively prime, then $\varphi_2$ is relatively prime to $a_{1,1}$. For if none of the rational prime factors of $b_{1,1}$ divides $a_{1,1}$, then none of the prime factors dividing $\varphi_2$ can divide $a_{1,1}$ either.

In this case we now have:
\weq{5}{a_{1,1} b_{1,1} = c_{1,1}, \quad a_{2,2} b_{2,2} = c_{2,2}.}
The congruence (4) then becomes
\weq{6}{c_{1,2} + a_{2,2} b_{2,2}\,\Theta \equiv 0 \ (\mathrm{mod}\ \varphi_1 \varphi_2),}
and since
\[
a_{1,2} + a_{2,2}\,\Theta \equiv 0 \ (\mathrm{mod}\ \varphi_1), \quad b_{1,2} + b_{2,2}\,\Theta \equiv 0 \ (\mathrm{mod}\ \varphi_2)
\]
holds, (6) splits into the two congruences:
\[
\begin{aligned}
c_{1,2} &\equiv b_{2,2} a_{1,2} \ (\mathrm{mod}\ \varphi_1) \\
&\equiv a_{2,2} b_{1,2} \ (\mathrm{mod}\ \varphi_2).
\end{aligned}
\]
But by \S.\,138, 13. these congruences are equivalent to the ordinary numerical congruences
\weq{7}{
\begin{aligned}
c_{1,2} &\equiv b_{2,2} a_{1,2} \ (\mathrm{mod}\ a_{1,1}) \\
&\equiv a_{2,2} b_{1,2} \ (\mathrm{mod}\ b_{1,1}),
\end{aligned}}
by which $c_{1,2}$ is determined with respect to the modulus $c_{1,1}$.


% ===== scan leaf 637, printed 616 =====
% !PRINTED 616
\weberpage{616}

It only remains to show how one can determine $\varphi_2$ in accordance with the requirement laid down here. Since two forms which differ from one another only by a numerical factor are equivalent, $\varphi_2$ is equivalent to
\[
c_2\, (t_1 + \omega_2 t_2),
\]
[\S.\,166, (10)], and if we take this form for $\varphi_2$, then $b_{1,1} = c_2$. But the number $\omega_2$ can always be chosen among those equivalent to it in such a way that $c_2$ becomes relatively prime to $a_{1,1}$, and then $a_{1,1}$ is also relatively prime to $\varphi_2$.

For if one replaces $\omega_2$ by
\[
\frac{p\, \omega_2 + q}{r\, \omega_2 + s}, \quad p s - q r = \pm 1,
\]
then by Bd.\ I, \S.\,122, (13), $c_2$ passes over into
\[
c_2' = -a_2 r^2 + b_2 r s + c_2 s^2
\]
and one can now choose the integral rational numbers $r, s$ in such a way that they are coprime to one another and that $c_2'$ is not divisible by any given primes whatever.

For if $n$ is such a prime, then one takes, when $n$ divides $a_2$ and $c_2$, and hence does not divide $b_2$, $r$ and $s$ not divisible by $n$; when $n$ does not divide $a_2$, $s$ divisible by $n$ and $r$ not divisible by $n$; and when $n$ does not divide $c_2$, $r$ divisible by $n$, $s$ not divisible by $n$. (If $a_2$ and $c_2$ are both not divisible by $n$, one has the choice between the two latter assumptions.) These requirements are compatible with one another for different arbitrarily given primes $n$, and $p, q$ can then be determined from the condition $p s - q r = \pm 1$.

We have derived these laws of the composition of the quadratic irrational numbers here only under the assumption that the discriminant is at the same time the fundamental number of a quadratic field, i.e. that $\Delta$ is a stem discriminant. But the same laws also hold, with slight modifications, in general $^{1)}$.

\noindent\rule{4cm}{0.4pt}

$^{1)}$ \gesperrt{Gauss}, Disquisitiones arithm., Art.\,234 ff. \gesperrt{Dirichlet-Dedekind}, Zahlentheorie. 4. Auflage, \S.\,187.

\begin{center}
\rule{0.2\linewidth}{0.4pt}
\end{center}


% ===== scan leaf 638, printed 617 =====
\weberpage{617}% page number not printed on this page
\pdfbookmark[1]{Twenty-first Section. Cyclotomic Fields.}{bm193}
\pdfbookmark[2]{§ 169. Decomposition of the prime number q in the cyclotomic field Ωq}{bm194}
% !PRINTED none
% Abschnitt-title page (part divider): no running head/page number visible in
% the scan for this page; §. 169 opens on it.

\begin{center}
\gesperrt{Twenty-first Chapter.}

\gesperrt{Cyclotomic Fields.}
\end{center}

\begin{center}
\rule{4cm}{0.4pt}
\end{center}

\begin{center}
\S.\,169.

\gesperrt{Decomposition of the prime $q$ into prime factors in the cyclotomic field $\Omega_{q^{\varkappa}}$.}
\end{center}

We make a second application of the general theory of algebraic numbers to the cyclotomic fields, and thereby gain the means of bringing the theory of the \gesperrt{Abelian} number fields in general to a beautiful conclusion.

The considerations that we have first to undertake can be applied, with small modifications, to every full cyclotomic field (\S.\,18). For the sake of simplicity, however, we restrict ourselves here to the case in which the degree of the root of unity is a power of a prime, which suffices for the intended application.

Let $q$ be a natural prime (with the inclusion of 2) and
\weq{1}{m = q^{\varkappa}}
a power of $q$, whose positive exponent $\varkappa$ is assumed to be greater than 1 for $q=2$.

Let $r$ further be a primitive $m$th root of unity and $\Omega_m$ the full cyclotomic field for the exponent $m$, whose degree is
\weq{2}{\mu = \varphi(m) = q^{\varkappa-1}(q-1)}
(Bd. I, \S.\,134).

We denote throughout by $n$ the $\mu$ numbers of a complete system of residues for the modulus $m$, with the exclusion of the numbers divisible by $q$; and the $\mu$ numbers $r^n$ are then the roots of the irreducible equation $f(x)=0$ of degree $\mu$, in which
\weq{3}{f(x) = \dfrac{x^m-1}{x^{\frac{m}{q}}-1} = x^{q^{\varkappa-1}(q-1)} + x^{q^{\varkappa-1}(q-2)} + \cdots + 1}
%CONT>>


% ===== scan leaf 639, printed 618 =====
% !PRINTED 618
\weberpage{618}
%<<CONT
is to be set. $r$ is therefore an \gesperrt{integer} in $\Omega_m$, and its norm is $=+1$; consequently $r$ is a unit.

$\Omega_m$ is a normal field, which under the $\mu$ substitutions
\[ s_n = (r, r^n) \]
passes over into itself.

These $\mu$ substitutions $s_n$ form an \gesperrt{Abelian} group $\mathfrak{N}$, which is the \gesperrt{Galois} group of the field $\Omega_m$. It is isomorphic with the group, likewise to be denoted by $\mathfrak{N}$, of the number classes $n$ taken with respect to the modulus $m$ (\S.\,16).

The function $f(x)$ can be decomposed into the $\mu$ factors $x - r^n$, and we therefore set
\weq{4}{f(x) = \prod\limits^{n} (x - r^n).}

If we set $x=1$ in this, and observe that [by (3)] $f(1) = q$, it follows that
\weq{5}{q = \prod\limits^{n} (1 - r^n).}

The numbers
\weq{6}{\sigma_n = 1 - r^n, \quad \sigma = 1 - r}
are, however, integers of the field $\Omega_m$, and formula (5) shows first of all \gesperrt{that $q$ is decomposable in the field $\Omega_m$ into $\mu$ factors $\sigma_n$.}

We prove first that the $\mu$ numbers $\sigma_n$ are associated with one another. If $n, n'$ denote two numbers not divisible by $q$, we can determine a natural number $a$ in such a way that
\[ n' \equiv a n \ (\mathrm{mod}\ m) \]
holds. But then
\[ \frac{\sigma_{n'}}{\sigma_n} = \frac{1 - r^{an}}{1 - r^n} = 1 + r^n + r^{2n} + \cdots + r^{(a-1)n} \]
is an integer, and hence $\sigma_{n'}$ is divisible by $\sigma_n$. Since $n$ can here be interchanged with $n'$, $\sigma_n$ is also divisible by $\sigma_{n'}$, and hence the two numbers are associated (\S.\,138). If therefore $\varepsilon$ denotes a unit, then by (5):
\weq{7}{q = \varepsilon \sigma^{\mu}, \quad N(\sigma) = q.}

Thus the natural prime $q$ is associated with the $\mu$th power of an integer $\sigma$ in the field $\Omega_m$. But this number $\sigma$ is still a prime in the field $\Omega_m$ as well, and indeed of the first degree. For if $\sigma$ has any divisor $\sigma'$, the norm of $\sigma'$ must be a divisor of the norm of $\sigma$; hence, if $\sigma'$ is not a
%CONT>>


% ===== scan leaf 640, printed 619 =====
% !PRINTED 619
\weberpage{619}
%<<CONT
unit, then $N(\sigma') = q$. Consequently the norm of $\sigma : \sigma'$ is equal to 1, i.\,e. $\sigma : \sigma'$ is a unit and $\sigma$ is associated with $\sigma'$. We thus have the first theorem:

\begin{quote}
I. \gesperrt{The natural prime $q$ is associated, in the full cyclotomic field $\Omega_m$, with the $\mu$th power of a prime factor of the first degree.}
\end{quote}

If we denote by $m_1$ a divisor of $m$, smaller than $m$ and greater than 1, and set
\weq{8}{m = m_1 m_2}
so that $m_1$ and $m_2$ are themselves powers of $q$, then there results the decomposition:
\[
% [sic: The printed product needs the factor (-1)^{m_1-1} on its right-hand side. For even m_1 (the permitted q=2 case), the product as printed equals 1-x^{m_1}, not x^{m_1}-1.]
x^{m_1} - 1 = (x-1)(xr^{m_2}-1)(xr^{2m_2}-1)\ldots(xr^{(m_1-1)m_2}-1),
\]
and if we set $x = r$ in it, by (6):
% [sic: The prefactor in (9) should be -1 for every m_1. The printed (-1)^{m_1} repeats the preceding sign error when q=2; for m=4,m_1=2,r=i, the two printed sides are -2 and +2. The following norm identity (10) is unaffected because mu is even.]
\weq{9}{r^{m_1} - 1 = (-1)^{m_1}\prod_{0,\,m_1-1}^{s}\sigma_{1+sm_2},}
and from this, since $\mu$ is always even,
\weq{10}{N(r^{m_1}-1) = q^{m_1}.}

From these results the discriminant $\Delta$ of the equation (3) is easily formed, which by Bd.\,I, \S.\,47 has the expression:
\[
\Delta = (-1)^{\frac{\mu(\mu-1)}{2}} N f'(r).
\]

For by (3):
\[
f'(r) = \dfrac{m\,r^{m-1}}{r^{\frac{m}{q}} - 1},
\]
hence, since the norm of $r$ is equal to 1:
\[
N f'(r) = \dfrac{m^{\mu}}{q^{\frac{m}{q}}},
\]
and consequently (since $\mu$ is always even):
\weq{11}{\Delta = (-1)^{\frac{\mu}{2}} q^{q^{\varkappa-1}[\varkappa(q-1)-1]}.}

The discriminant $\Delta$ is thus, apart from the sign, a power of $q$.


% ===== scan leaf 641, printed 620 =====
% !PRINTED 620
\weberpage{620}
\pdfbookmark[2]{§ 170. Minimal basis of the field Ωm}{bm195}

\begin{center}
\S. 170.

\gesperrt{Minimal basis of the field $\Omega_m$.}
\end{center}

The theorems that were proved in the preceding section lead us to the following theorem:

II. \gesperrt{The numbers}
\weq{1}{1, r, r^2, \ldots, r^{\mu-1}}
\gesperrt{form a basis of the system $\mathfrak{o}$ of integers in $\Omega_m$.}

To show this it suffices, by \S.\,145, to prove that the integer in $\Omega_m$
\weq{2}{\omega = x_0 + x_1 r + x_2 r^2 + \cdots + x_{\mu-1} r^{\mu-1},}
in which $x_0, x_1, \ldots, x_{\mu-1}$ are integral rational numbers, can be divisible by a rational prime $p$ only when the numbers $x_0, x_1, \ldots, x_{\mu-1}$ are all divisible by $p$.

Let us assume, then, that $\omega$ is divisible by $p$; then all the numbers $\omega_n$ that arise from $\omega$ by one of the $\mu$ substitutions $(r, r^n)$ are also divisible by $p$, and thus from (2) we obtain a system of $\mu$ equations which are linear with respect to the $x_i$:
\weq{3}{\omega_n = x_0 + x_1 r^n + x_2 r^{2n} + \cdots + x_{\mu-1} r^{(\mu-1)n}.}

But the determinant of this system, that is, the determinant formed from the $\mu$ rows
\[ 1, r^n, r^{2n}, \ldots, r^{(\mu-1)n} \]
is, by Bd. I, \S.\,46, equal to the square root $\sqrt{\varDelta}$, and if we now solve the system of equations (3), it follows that all the rational numbers $\varDelta x_i$ must be divisible by $p$.

If now $p$ is not $=q$, then $\varDelta$ is not divisible by $p$, and all the $x_i$ must be divisible by $p$.

But if $p = q$, we set
\[ f(t) = x_0 + x_1 t + x_2 t^2 + \cdots + x_{\mu-1} t^{\mu-1}, \]
so that $\omega_n = f(r^n)$, and in this we set $r = 1 - \sigma$ [\S.\,169, (6)]. Then, according to the \gesperrt{Taylor} development:
\weq{4}{\begin{aligned}
\omega &= f(1-\sigma) = f(1) - \sigma f'(1) + \dfrac{\sigma^2}{1.2} f''(1) - \cdots\\
&\qquad - \dfrac{\sigma^{\mu-1}}{\Pi(\mu-1)} f^{(\mu-1)}(1);
\end{aligned}}
%CONT>>


% ===== scan leaf 642, printed 621 =====
% !PRINTED 621
\weberpage{621}
%<<CONT
and herein the coefficients
\[
f(1), f'(1), \tfrac12 f''(1), \cdots, \dfrac{1}{\Pi(\mu-1)} f^{(\mu-1)}(1)
\]
are integral rational numbers, namely
\begin{align*}
\dfrac{1}{\Pi(\mu-1)} f^{(\mu-1)}(1) &= x_{\mu-1} \\[6pt]
\dfrac{1}{\Pi(\mu-2)} f^{(\mu-2)}(1) &= x_{\mu-2}+(\mu-1)x_{\mu-1} \\[6pt]
\dfrac{1}{\Pi(\mu-3)} f^{(\mu-3)}(1) &= x_{\mu-3}+(\mu-2)x_{\mu-2}+\dfrac{(\mu-1)(\mu-2)}{1.2}x_{\mu-1} \tag{5}\\[6pt]
\cdots\cdots\cdots &\cdots\cdots\cdots\cdots\cdots\cdots\cdots \\[6pt]
f(1) &= x_0+x_1+\cdots\cdots\cdots\cdots+x_{\mu-1}.
\end{align*}
Now since $\omega$ is divisible by $q$, and consequently by $\sigma^\mu$,
$f(1)$ must by (4) be divisible by $q$. But since $\mu>1$, and
consequently $q:\sigma$ is still divisible by $\sigma$, it follows that
\[
f'(1) - \dfrac{\sigma}{1.2} f''(1) \ldots
\]
also is divisible by $\sigma$, and consequently $f'(1)$ by $q$; and reasoning
onward in this way, it results that all the rational
numbers (5) are divisible by $q$, and consequently the
$x_{\mu-1}, x_{\mu-2}, \ldots, x_0$ are also divisible by $q$. With this Theorem II. is completely proved.

Instead of the sequence (1) one may also take as a basis of $\mathfrak{o}$ an
arbitrary sequence of $\mu$ successive powers of $r$:
\weq{6}{r^a, r^{a+1}, r^{a+2}, \ldots, r^{a+\mu-1},}
which results when one multiplies all the terms of the sequence (1)
by $r^a$.

The correctness of this remark follows from the fact that
the $r^{\pm a}$ are integers, and that therefore $\omega$ and $r^a\omega$ are always simultaneously integral or fractional numbers.

For example, the numbers
\weq{7}{r^{-\frac12\mu+1}, r^{-\frac12\mu+2}, \ldots, r^{-1}, 1, r, r^2, \ldots, r^{\frac12\mu}}
form a basis of $\mathfrak{o}$.

With this the fundamental number $\Delta$ of the field $\Omega_m$ is also determined.
By the general definition (\S.\,145) it is equal to the discriminant % [sic] German source prints "Discirminante"; "Discriminante" is meant
\[
\Delta(1, r, r^2, \ldots, r^{\mu-1}),
\]
i.e. equal to the discriminant of the cyclotomic equation determined in \S.\,169.


% ===== scan leaf 643, printed 622 =====
% !PRINTED 622
\weberpage{622}
\pdfbookmark[2]{§ 171. The prime ideals in the field Ωm}{bm196}
\begin{center}
\S.\,171.

\gesperrt{The Prime Ideals in the Field $\Omega_m$.}
\end{center}

% [sic: The printed reference 164 should be 169, theorem I and (7), where the prime-power factorisation is established.]
% [Source shorthand: q is associated with sigma^mu, not generally equal to it as an element; the corresponding principal ideals are equal.]
We have seen in \S.\,164 that the natural prime $q$
is the $\mu$th power of a prime $\sigma$ in the field $\Omega_m$.

In order to determine all the prime ideals that exist in the
field $\Omega_m$, it therefore remains to decompose into prime factors
all the natural primes $p$ different from $q$, and to select from
among them those that are not associated.

The foundation for this investigation is formed by the theorems
of \S.\,162, when the field there denoted by $R$ is replaced by
the field of rational numbers.

Every number $\omega$ of the field $\Omega_m$ is carried by one of the $\mu$ substitutions
\[
s_n = (r, r^n)
\]
into a determinate other number $\omega_n$, which is likewise contained
in $\Omega_m$. If
\weq{1}{\omega = \omega_1 = a_0 + a_1 r + a_2 r^2 + \cdots + a_{\mu-1} r^{\mu-1},}
then
\weq{2}{\omega_n = a_0 + a_1 r^n + a_2 r^{2n} + \cdots + a_{\mu-1} r^{(\mu-1)n}.}
If one of these numbers $\omega_n$ is an integer, as we shall now
assume, then so are all the others, and this occurs always then
and only then, when the rational numbers
$a_0, a_1, \ldots, a_{\mu-1}$ are integers.

If $h, k$ are any two exponents, then
\[
r^h - r^k = r^k(r^{h-k} - 1),
\]
and by \S.\,169, (9) this is associated with a power of $\sigma$,
hence, if $\mathfrak{p}$ denotes a prime ideal dividing $p$, it is not
divisible by $\mathfrak{p}$, except when $h \equiv k \pmod{m}$. From this it results
% [sic: The following assertion needs congruence for every integral omega, not for one fixed omega (omega=1 is a counterexample). Choosing omega=r proves the required assertion about the inertia group X, so the unramifiedness conclusion remains valid.]
that two numbers $\omega_h, \omega_k$ can be congruent with respect to the modulus $\mathfrak{p}$ only when $h \equiv k \pmod{m}$, and that therefore the
group $X$ of \S.\,162, which satisfies the condition
\[
\omega\,|\,\chi \equiv \omega \pmod{\mathfrak{p}}
\]
contains only the \gesperrt{identical substitution}. But from this it
follows that $g = 1$, i.\,e. that $p$ \gesperrt{is not divisible by the square
of a prime ideal.}


% ===== scan leaf 644, printed 623 =====
% !PRINTED 623
\weberpage{623}
\pdfbookmark[2]{§ 172. The conjugate prime ideals}{bm197}
If we form from (2) the $p$th power of $\omega$, and observe
\gesperrt{Fermat's} theorem for rational numbers $[a_0^p \equiv a_0 \pmod{p}]$,
we obtain
\[
\omega^p \equiv a_0 + a_1 r^p + a_2 r^{2p} + \cdots + a_{\mu-1} r^{(\mu-1)p} \pmod{p},
\]
or by (2)
\weq{3}{\omega^p \equiv \omega_p \pmod{p}.}

This determines the group $\Psi$ to which the ideal $\mathfrak{p}$
belongs. For by (3) we have also
\[
\omega^p \equiv \omega_p \pmod{\mathfrak{p}},
\]
hence $\psi_0 = (r, r^p)$, and by \S.\,162, (13) the group $\Psi$ consists
of the powers of this substitution, so far as they are different
from one another.

If therefore $f$ is the least positive exponent for which
\weq{4}{p^f \equiv 1 \pmod{m}}
holds, that is, if $p$ belongs to the exponent $f$ for the modulus $m$, then
\weq{5}{\Psi = (r,r), (r,r^p), (r,r^{p^2}), \ldots (r,r^{p^{f-1}})}
and it is of degree $f$. Accordingly $\mathfrak{p}$ is also a prime ideal
of degree $f$, and
\weq{6}{N(\mathfrak{p}) = p^f.}

The number $e$ of the conjugate prime factors of $f$ that are different
from one another is $\mu : f$, and we therefore have the theorem: % [sic] source prints $f$ after "of"; mathematically $p$ is required

III. \gesperrt{If $p$ is a prime different from $q$ which belongs to the exponent $f$ for the modulus $m$, and if further $\mu = \varphi(m) = ef$, then $p$ splits in $\Omega_m$ into $e$ conjugate prime factors of degree $f$, different from one another.}

\begin{center}
\S.\,172.

\gesperrt{The conjugate prime ideals.}
\end{center}

We must enter somewhat more closely into the formation of the conjugate prime factors $\mathfrak{p}_n$ of a prime $p$ different from $q$,
which by the theorem last proved splits in $\Omega_m$ into $e$
prime factors different from one another, when
\weq{1}{\varphi(m) = ef}
holds, and $f$ is the least positive exponent for which
\weq{2}{p^f \equiv 1 \pmod{m}.}


% ===== scan leaf 645, printed 624 =====
% !PRINTED 624
\weberpage{624}
The group $\Psi$ is isomorphic with a group, contained in $\mathfrak{N}$ (\S.\,169), of integral
rational numbers taken with respect to the modulus $m$, which we denote by $\mathfrak{A}$, and
which consists of all numbers $a$ satisfying a
congruence
\weq{3}{a \equiv p^h \pmod{m}}
and which we can therefore represent symbolically by
\weq{4}{\mathfrak{A} \equiv p^h \pmod{m}}
if $h$ denotes an arbitrary integral exponent (cf. \S.\,18).

If, then, we select the numbers $\xi_1, \xi_2, \ldots, \xi_e$ from $\mathfrak{N}$
suitably, we can set:
\weq{5}{\mathfrak{N} = \mathfrak{A}\xi_1 + \mathfrak{A}\xi_2 + \mathfrak{A}\xi_3 + \cdots + \mathfrak{A}\xi}
% [sic] source prints the final \xi without the expected subscript e; (6)--(7) supply the context
and to each of these subsidiary groups there corresponds one of the conjugate functionals $\mathfrak{p}_{\xi_1}, \mathfrak{p}_{\xi_2}, \ldots, \mathfrak{p}_{\xi_e}$, which are all different from one another and whose product is associated with $p$.
% [sic] "Functionale" as printed; context elsewhere on the page uses "Primfactoren"/"Primideale" for this notion

We therefore set
\weq{6}{p = \mathfrak{p}_{\xi_1} \mathfrak{p}_{\xi_2} \ldots \mathfrak{p}_{\xi_e}.}
But in order to characterise more exactly the system of numbers
\weq{7}{\xi_1, \xi_2, \ldots, \xi_e,}
on which the matter chiefly depends,
we remark that for this purpose we can take any system of $e$ numbers of the
group $\mathfrak{N}$, provided only that no two among them,
$\xi, \xi'$, satisfy a congruence
\[
\xi'\xi^{-1} \equiv p^h \pmod{m}.
\]

In order to find such a system of numbers, it is well to separate the case
of an odd $m$ from the other, in which $m$ is a
power of 2.

1. If first $m = q^{\varkappa}$ is odd, we take a
primitive root $g$ of $m$ (\S.\,14) and set
\weq{8}{p \equiv g^{\gamma} \pmod{m}.}
For $f$ is [by (1), (2)] the smallest positive number that satisfies the
condition
\[
\gamma f \equiv 0 \pmod{\mu}
\]
and consequently $e$ is the greatest common divisor
of $\mu$ and $\gamma$, and any number $n \equiv g^{\lambda}$ is only then
%CONT>>


% ===== scan leaf 646, printed 625 =====
% !PRINTED 625
\weberpage{625}
%<<CONT
congruent to a power of $p$ when $\lambda$ is divisible by $e$. For only then can the congruence
\[ \lambda \equiv \gamma x \, (\mathrm{mod}\ \mu) \]
be satisfied. If we set
\[ \gamma = e \nu, \]
then $\nu$ is relatively prime to $f$. From this it follows that the sequence of numbers
\weq{9}{1, g, g^2, \ldots, g^{e-1}}
can be taken for the $\xi$. For if $h, h'$ are two different numbers of the sequence $0, 1, 2, \ldots, e-1$, then $h'-h$ is never divisible by $e$, and consequently
\[ \xi' \xi^{-1} = g^{h'-h} \]
is never congruent to a power of $p$.

We shall further distinguish the two cases in which $p-1$ is or is not divisible by $q$.

a) If $p \equiv 1 \, (\mathrm{mod}\ q)$, then $\gamma \equiv 0 \, (\mathrm{mod}\ q-1)$, and consequently $e$ is divisible by $q-1$. We can therefore set
\[ e = \varphi(q^{\varkappa_1}) = q^{\varkappa_1 - 1}(q-1) \]
where $0 < \varkappa_1 \leqq \varkappa$.

When $\varkappa_1 = \varkappa$, and hence $e = \varphi(m)$, $f=1$, and therefore $p-1$ is divisible by $m$, then $\xi$ runs through the whole group $\mathfrak{N}$, and all the prime functionals $\mathfrak{p}_n$ are different from one another.

If $\varkappa_1 < \varkappa$, then $\nu$ is not divisible by $q$, since otherwise $e$ would not be the greatest common divisor of $\mu$ and $\gamma$, and
\[ p - 1 \equiv g^{\nu\varphi(q^{\varkappa_1})} - 1 \, (\mathrm{mod}\ m) \]
is divisible by $q^{\varkappa_1}$, but by no higher power of $q$. If we therefore set $q^{\varkappa_1} = m_1$ and
\weq{10}{m = m_1 m_2,}
then $m_1$ is the \gesperrt{greatest common divisor of $p-1$ and $m$}. The prime $p$ splits in the field $\Omega_m$ into $\varphi(m_1)$ different prime factors. But $p$ splits into just as many prime factors in the field $\Omega_{m_1}$, which is a divisor of the field $\Omega_m$ and consists of all rational functions of the $m_1$th root of unity $r_1 = r^{m_2}$. From this it follows that the prime factors of $p$ in $\Omega_{m_1}$ are not further decomposable in the field $\Omega_m$ either, and that we can therefore represent the prime factors $p$ in $\Omega_m$ as functionals of the field $\Omega_{m_1}$.
% [sic] Source prints 'Primfactoren p' without 'von'; retain the printed wording.

Here $\xi$ runs through a complete system of residues not divisible by $q$ with respect to the modulus $m_1$, and the $\mathfrak{p}_\xi$ are ideals in $\Omega_{m_1}$.


% ===== scan leaf 647, printed 626 =====
% !PRINTED 626
\weberpage{626}
b) If, conversely, $\gamma$ is divisible by $q-1$, it follows from (8) that $p-1$ is divisible by $q$. If, then, we assume that $p-1$ is not divisible by $q$, then $e$ is also not divisible by $q-1$, and $\xi$ runs through the sequence of numbers (9).

If we form the sum of these numbers $\xi$, it follows that:
\[ (g-1) \Sigma \xi \equiv g^e - 1 \, (\mathrm{mod}\ m), \]
and this number is divisible or not divisible by $q$, according as $e$ is divisible or not divisible by $q-1$. Since, moreover, $g-1$ is not divisible by $q$, $\Sigma\xi$ is divisible or not divisible by $q$, according as $p-1$ is divisible or not divisible by $q$.

2. We turn to the case in which $m = 2^\varkappa$ is a power of 2, and here we have again to distinguish whether $p-1$ is divisible by 4 or only by 2.

a) If $p-1$ is divisible by 4, then (\S.\,15):
\[ p \equiv 5^\beta \, (\mathrm{mod}\ m), \]
and $f$ is the least positive number that satisfies the condition
\[ \beta f \equiv 0 \, [\mathrm{mod}\ \tfrac{1}{2}\varphi(m)]. \]
If $f=1$, that is, if $\beta$ is divisible by $\tfrac{1}{2}\varphi(m) = 2^{\varkappa-2}$, then $e = \varphi(m)$, and $\xi$ runs through the whole group $\mathfrak{N}$, i.\,e. all odd numbers between 0 and $m$. This is the case in which $p-1$ is divisible by $m$.

If $2^{\varkappa_1-2}$ is the highest power of 2 that divides $\beta$, and $2 \leqq \varkappa_1 < \varkappa$, then, if we set $m_1 = 2^{\varkappa_1}$ and $m = m_1 m_2$,
\[ f = 2^{\varkappa - \varkappa_1}, \quad e = 2^{\varkappa_1 - 1} = \varphi(m_1), \]
and, if $s$ is an odd number, we have
\[ p - 1 \equiv 5^{s\tfrac{1}{2}\varphi(m_1)} - 1 \ldots \, (\mathrm{mod}\ m). \]
% [sic] Source prints three low dots after -1 in this congruence; no continuation term is specified.
By \S.\,15, $5^{\tfrac{1}{2}\varphi(m_1)} - 1$ is divisible by $m_1$, but by no higher power of 2. Since, moreover, one has by the same theorems
\[ 5^{s\tfrac{1}{2}\varphi(m_1)} \equiv 5^{\tfrac{1}{2}\varphi(m_1)} \, (\mathrm{mod}\ 2m_1) \]
it follows that $m_1$ is the greatest common divisor of $p-1$ and $m$.

If, then, as above, we call in the field $\Omega_{m_1}$, it results that, when $m_1 \geqq 4$ is the greatest common divisor of $p-1$ and $m$, $\xi$ runs through a complete system of odd residues of $m_1$, and that $\mathfrak{p}_\xi$ are at the same time the prime factors of $p$ in the field $\Omega_{m_1}$.


% ===== scan leaf 648, printed 627 =====
% !PRINTED 627
\weberpage{627}
b) If $p-1$ is not divisible by 4, and hence
\weq{11}{p \equiv -5^\beta \, (\mathrm{mod}\ m),}
then $f$ is the least positive number that satisfies the congruences
\[ f \equiv 0 \, (\mathrm{mod}\ 2), \quad 2\beta f \equiv 0 \, [\mathrm{mod}\ \varphi(m)] \]
and consequently $e$ is the greatest common divisor of $\varphi(m)$ and $2\beta$, or, if $2\beta$ is divisible by $\varphi(m)$, equal to $\tfrac{1}{2}\varphi(m)$. (An exception is formed here by the case $m=4$, in which $e=1$. This case, which has indeed already been completely settled in earlier chapters, we therefore now set aside.)

It then follows from (11) that a number of the form $5^\lambda$ can be congruent to a power of $p$ only if $\lambda$ is a multiple of $2\beta$, and hence a multiple of $e$.

Now here a number can be found in each of the subsidiary groups $\mathfrak{A}\xi$ that is $\equiv 1 \, (\mathrm{mod}\ 4)$, since $p \equiv -1 \, (\mathrm{mod}\ 4)$ is a number in $\mathfrak{A}$, and we accordingly obtain a complete system of the $\xi$ in the form:
\[ 1, 5, 5^2, \ldots, 5^{e-1}, \]
and in it $e$ is at most $= \tfrac{1}{2}\varphi(m) = 2^{\varkappa-2}$.

We collect what has been proved in the following theorem:

III. \gesperrt{If one decomposes the prime $p$ different from $q$ in the field $\Omega_m$ into its prime factors:}
\[ p = \mathfrak{p}_{\xi_1} \mathfrak{p}_{\xi_2} \cdots \mathfrak{p}_{\xi_e}, \]
\gesperrt{then, if $p-1$ is divisible by $q$, or in the case of an even $m$ by 4, one may let the $\xi$ run through a complete system of residues not divisible by $q$ with respect to the modulus $m_1$, and take the $\mathfrak{p}_\xi$ to be prime functionals in the field $\Omega_{m_1}$, if $m_1$ is the greatest common divisor of $p-1$ and $m$.}

\gesperrt{But if $p-1$ is not divisible by $q$, or for an even $m$ not by 4, and if $g$ denotes in the first case a primitive root of $m$, in the second case the number 5, then $\xi$ runs through the sequence of numbers}
\[ 1, g, g^2, \ldots, g^{e-1}, \]
\gesperrt{and for an odd $m$, $e$ is not divisible by $q-1$, and for an even $m$ it is at least $=2$ and at most $=\tfrac{1}{2}\varphi(m)$. Excluded from this is the case $m=4$.}


% ===== scan leaf 649, printed 628 =====
% !PRINTED 628
\weberpage{628}
\pdfbookmark[2]{§ 173. Representation of the prime factors of p}{bm198}
\begin{center}
\S.\,173.

\gesperrt{Representation of the prime factors of $p$.}
\end{center}

For the representation of the prime functionals of the field $\Omega_m$ we can apply a procedure that we became acquainted with in \S.\,157, and which, in consequence of the circumstance that
\[ 1, r, r^2, \ldots, r^{\mu-1} \]
is a basis of $\mathfrak{o}$, is here essentially simplified.

The natural prime $q$ is, as we have already seen, the $m$th power of a prime $\sigma$ existing in the field $\Omega_m$; with this we need not concern ourselves further, and we therefore consider here only the primes $p$ different from $q$.
% [sic] Source prints the m-th power of sigma here; the ramification index of the prime above q is mu = phi(m). Preserve the printed m.
Let $\mathfrak{p}$ be a prime factor of $p$, and let $\mathfrak{A}$ have the same meaning as above. If $p$ belongs to the exponent $f$ and
\[ ef = \varphi(m) = \mu \]
holds, then the group $\mathfrak{A}$ consists of the numbers
\[ 1, p, p^2, \ldots, p^{f-1}. \]
If now
\[ f(t) = N(t-r) = \prod^n (t - r^n) \]
denotes the polynomial in $t$ of degree $\mu$ whose roots are the $\mu$ quantities $r^n$, then we can decompose it into factors in such a way that each factor corresponds to one of the subsidiary groups of $\mathfrak{A}$. Namely, if we set
\weq{1}{F_1(t) = \prod_{0,\,f-1}^{h} (t - r^{p^h}),}
and in general, for every $n$ not divisible by $q$,
\weq{2}{F_n(t) = \prod_{0,\,f-1}^{h} (t - r^{np^h}),}
then $F_n(t)$ is identical with $F_{n'}(t)$ when $n$ and $n'$ belong to the same subsidiary group $\mathfrak{A}\xi_i$. If, then, we understand by $\xi_1, \xi_2, \ldots, \xi_e$ the system of numbers applied in \S.\,172, (5), we have
\weq{3}{f(t) = F_{\xi_1}(t)\, F_{\xi_2}(t) \ldots F_{\xi_e}(t).}

Now, however, by the binomial theorem:
\[ (t - r^{p^h})^p \equiv (t^p - r^{p^{h+1}}) \, (\mathrm{mod}\ p), \]
%CONT>>


% ===== scan leaf 650, printed 629 =====
% !PRINTED 629
\weberpage{629}
%<<CONT
and if we apply this to every factor of $F_n(t)$, and observe
that $p^{h+1}$ runs through the same sequence of numbers with
respect to the modulus $m$ as $p^h$ does, it follows:
\[
[F_n(t)]^p \equiv F_n(t^p) \ (\mathrm{mod}\ p),
\]
and this congruence then naturally holds also for the modulus $\mathfrak p$.

But this is the criterion for $F_n(t)$ being congruent, with
respect to the modulus $\mathfrak p$, to a polynomial $P_n(t)$ with integral
rational numerical coefficients (\S.\,150, 4.); thus
\weq{4}{F_n(t) \equiv P_n(t) \ (\mathrm{mod}\ \mathfrak p).}

If we substitute this in (3), there results in the first place
a congruence with respect to the modulus $\mathfrak p$, which however,
since it is a congruence between \gesperrt{rational} functions, must also
hold for the modulus $p$; thus
\weq{5}{f(t) \equiv P_{\xi_1}(t)\ P_{\xi_2}(t)\ \ldots\ P_{\xi_e}(t) \ (\mathrm{mod}\ p).}

By the definition (1), (4) we have
\weq{6}{P_1(r) \equiv 0 \ (\mathrm{mod}\ \mathfrak p),}
and all the roots of this congruence are
\weq{7}{r,\ r^p,\ \ldots,\ r^{p^{f-1}}.}

In the same way the congruence
\[
P_n(t) \equiv 0 \ (\mathrm{mod}\ \mathfrak p)
\]
has the roots
\[
r^n,\ r^{np},\ \ldots,\ r^{np^{f-1}},
\]
and these roots, if one lets $n$ run through the system of numbers
$\xi_1, \xi_2, \ldots, \xi_e$, are all incongruent with respect to the
modulus $\mathfrak p$.

If, in the congruence
\[
F_1(t) \equiv P_1(t) \ (\mathrm{mod}\ \mathfrak p)
\]
one makes the substitution $(r, r^n)$, then $F_1(t)$ passes into $F_n(t)$,
and $\mathfrak p$ into the conjugate ideal $\mathfrak p_n$, while $P_1(t)$, as a rational
form, remains unchanged. We therefore have
\weq{8}{F_n(t) \equiv P_1(t) \ (\mathrm{mod}\ \mathfrak p_n),}
and if we set $t=r$ in this, it follows:
\weq{9}{F_n(r) \equiv P_1(r) \ (\mathrm{mod}\ \mathfrak p_n),}
$F_n(r)$ is, however, if $n$ is not contained in $\mathfrak A$, relatively
prime to $p$, and hence not divisible by $\mathfrak p_n$.

By (9), then, $P_1(r)$ is divisible by $\mathfrak p_1$, but by none of the
%CONT>>


% ===== scan leaf 651, printed 630 =====
% !PRINTED 630
\weberpage{630}
%<<CONT
prime factors conjugate to $\mathfrak p_1$, and it follows,
since $p$ is not divisible by $\mathfrak p^2$:

1. \gesperrt{The prime factor $\mathfrak p_1$ is the greatest common
divisor of $p$ and $P_1(r)$; in the same way it follows that
$\mathfrak p_n$ is the greatest common divisor of $p$ and $P_1(r^n)$,
or also the greatest common divisor of $p$ and
$P_{n'}(r)$, when $n n' \equiv 1 \ (\mathrm{mod}\ m)$.}

The latter follows from the congruence which results from (4)
by the substitution $(r, r^{n'})$,
\[
F_1(t) \equiv P_n(t) \ (\mathrm{mod}\ \mathfrak p_{n'}).
\]

One thus obtains all the prime factors of $p$ if one
seeks out the greatest common divisors of $p$ with each
of the rational functions
\weq{10}{P_{\xi_1}(r),\ P_{\xi_2}(r),\ \ldots,\ P_{\xi_e}(r)}
and indeed $\mathfrak p_{\xi_i}$ is the greatest common divisor
of $p$ and $P_{\xi_i^{-1}}$.
% [sic] Source omits the argument (r) after P_{xi_i^{-1}} here; (10) and theorem 1 give the intended evaluation.

We shall consider somewhat more closely the case where $f=1$,
and thus
\weq{11}{p \equiv 1 \ (\mathrm{mod}\ m).}

In this case $e = \varphi(m)$; the group $\mathfrak A$ reduces
to the unit, and all the $\mu$ conjugate factors $\mathfrak p$
of $p$ are different from one another. The forms
\[
P_{\xi_1}(t),\ P_{\xi_2}(t),\ \ldots,\ P_{\xi_e}(t)
\]
all become linear, that is, $r$ is congruent to a rational number
with respect to each of the moduli $\mathfrak p_i$.

This follows also from the fact that here the number of
numbers incongruent with respect to the modulus $\mathfrak p$ is equal to $N(\mathfrak p) = p$,
and that therefore $0, 1, 2, \ldots, p-1$ is a complete system of residues.

If accordingly, say, $r \equiv c \ (\mathrm{mod}\ \mathfrak p)$, then, since $r^m = 1$,
we must have $c^m \equiv 1 \ (\mathrm{mod}\ \mathfrak p)$ and hence also $(\mathrm{mod}\ p)$, and if
now $g$ is a primitive root of the prime $p$, then we must have
\weq{12}{c \equiv g^{-n\frac{p-1}{m}} \ (\mathrm{mod}\ p)}
in which $n$ is relatively prime to $m$. If we let $\mathfrak p$ run through
the system of conjugate ideals, then $n$ in (12) must run through the
group $\mathfrak N$, and it is therefore only a matter of notation
if we stipulate:
\weq{13}{r \equiv g^{-\frac{p-1}{m}} \ (\mathrm{mod}\ \mathfrak p).}


% ===== scan leaf 652, printed 631 =====
% !PRINTED 631
\weberpage{631}
If we make the substitution $(r, r^n)$ in it, this becomes
\weq{14}{r^n \equiv g^{-\frac{p-1}{m}} \ (\mathrm{mod}\ \mathfrak p_n),}
or, if we define $n'$ by the congruence
\weq{15}{n n' \equiv 1 \ (\mathrm{mod}\ m),}
then:
\weq{16}{r \equiv g^{-n'\frac{p-1}{m}} \ (\mathrm{mod}\ \mathfrak p_n).}

In this case, therefore, in agreement with the general
 rule, $\mathfrak p_n$ is the greatest common divisor of
\[
p \ \text{and} \ \left(r - g^{-n'\frac{p-1}{m}}\right).
\]

By this determination the individual prime forms $\mathfrak p_n$
are exactly characterised. We thus obtain the theorem:

2. \gesperrt{A prime $p$ which is congruent to 1 with respect
to the modulus $m$ splits in the field $\Omega_m$ into $\varphi(m)$
prime factors $\mathfrak p_n$ of the first degree, distinct
from one another, which one obtains as the greatest
 common divisors of $p$ with the various
numbers $r - g^{-n'\frac{p-1}{m}}$, if $g$ is a primitive
root of $p$ and $n'$ is determined by the congruence
$n n' \equiv 1 \ (\mathrm{mod}\ m)$.}

According to \S.\,141 we can decompose every integer or fractional
number $\omega$, and every functional of the field $\Omega_m$, into prime
 factors in such a way that numerator and denominator contain
 no common prime factor, and this decomposition is,
apart from unit factors, completely determined, so that, if
$\mathfrak p$ is a prime ideal, there exists a perfectly determinate positive
or negative exponent $k$ such that the functional $\omega \mathfrak p^{-k}$ contains
 the prime factor $k$ neither in the numerator nor in the denominator. % [sic] source prints $k$; context requires $\mathfrak p$
We then say,
for the sake of brevity, that $\mathfrak p^k$ is \gesperrt{the power of $\mathfrak p$ contained
in $\omega$}.

If $k=0$, we say that $\mathfrak p$ is \gesperrt{not contained} in $\omega$.

Two numbers (or also functionals) $\omega, \omega'$ of a field, in
which one and the same prime functional $\mathfrak p$ is contained, \gesperrt{we
shall call divisor-related, or related for short}; if, on the other hand,
no prime functional is contained in both $\omega$ and $\omega'$, then
$\omega$ and $\omega'$ are called \gesperrt{coprime, or foreign}. This expression
will be particularly useful to us in what follows, when
%CONT>>


% ===== scan leaf 653, printed 632 =====
% !PRINTED 632
\weberpage{632}
\pdfbookmark[2]{§ 174. Kummer's theorem}{bm199}
%<<CONT
the place of $\omega'$ is taken by the natural prime $p$, which is divisible
by $\mathfrak p$; and we speak accordingly of \gesperrt{the primes related
to a number $\omega$ or foreign to it}.

\gesperrt{A number $\omega$ that has no related primes at all is
a unit.}

In the field $\Omega_m$ (as indeed in every normal field) all
the conjugate numbers $\omega_n$ are related to the same primes.

\begin{center}
\S.\ 174.

\gesperrt{Kummer's Theorem.}
\end{center}

In the theory of the division of the circle (Bd.\,I, \S.\,169)
we have already become acquainted with certain decompositions of the
natural primes into factors composed of roots of unity,
and it is now of importance to investigate how these
factors stand to the prime factors in the field $\Omega_m$.

We consider a prime $p$ which belongs to the exponent 1,
so that $m$ is thus a divisor of $p-1$, and set
\[
p - 1 = m m'.
\]

From the $p$th roots of unity $\varrho$, periods of $m'$ terms each
can be formed according to Bd.\,I, \S.\,167, which, if $g$ is a primitive
root of $p$ (not, as above, of $m$), have the expression:
\weq{1}{
\begin{aligned}
\eta &= \varrho \phantom{^g} + \varrho^{g^m} + \varrho^{g^{2m}} + \cdots + \varrho^{g^{(m'-1)m}} \\
\eta_1 &= \varrho^g + \varrho^{g^{m+1}} + \varrho^{g^{2m+1}} + \cdots + \varrho^{g^{(m'-1)m+1}} \\
&\ \, \cdots\cdots\cdots\cdots\cdots\cdots\cdots\cdots\cdots\cdots \\
\eta_{m-1} &= \varrho^{g^{m-1}} + \varrho^{g^{2m-1}} + \varrho^{g^{3m-1}} + \cdots + \varrho^{g^{p-2}}.
\end{aligned}
}

If, as in \S.\,169 of the first volume, we denote by $\varepsilon$ a
primitive $(p-1)$th root of unity, so that
\weq{2}{\varepsilon^{m'} = r,}
as hitherto, is a primitive $m$th root of unity, then the
resolvents are
\begin{equation*}\tag*{3)}
(\varepsilon^\lambda, \varrho) = \varrho + \varepsilon^\lambda \varrho^g + \varepsilon^{2\lambda}\varrho^{g^2} + \cdots + \varepsilon^{(p-2)\lambda}\varrho^{g^{p-2}},
\end{equation*}
% [sic] printed "3)" without an opening parenthesis, unlike the other display numbers on this page
in which $\lambda$ may be any integral rational number, and in
particular for $\lambda = m'$:
\weq{4}{(r, \eta) = \eta + r\eta_1 + r^2\eta_2 + \cdots + r^{m-1}\eta_{m-1}.}


% ===== scan leaf 654, printed 633 =====
% !PRINTED 633
\weberpage{633}
These resolvents are integral algebraic numbers in certain
cyclotomic fields, and we have, when $\varepsilon^\lambda$ is not $=1$,
\[
(\varepsilon^\lambda,\varrho)\,(\varepsilon^{-\lambda},\varrho) = (-1)^\lambda p,
\]
and for $\lambda = m'n$
\weq{5}{(r^n,\eta)\,(r^{-n},\eta) = \pm\, p.}

All these resolvents are therefore divisors of the natural
prime $p$.

Besides these resolvents we have derived, at the place cited,
still other integers of the cyclotomic fields, namely
\weq{6}{\psi_{\lambda,\mu}(\varepsilon) = \sum\limits_{1,\,p-2}^{t} \varepsilon^{\mu\,\mathrm{ind}\,t - (\lambda+\mu)\,\mathrm{ind}\,(t+1)},}
(in which the indices refer to the base $g$), which in general belong to the field $\Omega_{p-1}$, and in the special case where $\lambda$ and $\mu$ are divisible by $m'$, to the field $\Omega_m$. Of particular importance among them are the numbers of the field $\Omega_m$:
\weq{7}{\psi_{sm',\,m'}(\varepsilon) = \psi_s(r) = \sum\limits_{1,\,p-2}^{t} r^{\mathrm{ind}\,t - (s+1)\,\mathrm{ind}\,(t+1)}.}

These numbers $\psi_{\lambda,\mu}(\varepsilon)$ too, when none of the numbers $\lambda$, $\mu$, $\lambda+\mu$ is divisible by $p-1$, are divisors of $p$ by the formula
\[
\psi_{\lambda,\mu}(\varepsilon)\,\psi_{\lambda,\mu}(\varepsilon^{-1}) = p,
\]
from which again, by (7), there results:
\weq{8}{\psi_s(r)\,\psi_s(r^{-1}) = p,}
provided that neither $s$ nor $s+1$ is divisible by $m$. Certain combinations of the resolvents (4), among them the $m$th power, we have expressed by means of the functions $\psi$ [Bd.\ I, \S.\,169, (14), (15)]:
\weq{9}{(r,\eta)^n\,(r^n,\eta)^{-1} = \psi_1(r)\,\psi_2(r)\cdots\psi_{n-1}(r),}
so long as $n<m$, and
\weq{10}{(r,\eta)\,(r^{-1},\eta) = \pm\,p}
\weq{11}{(r,\eta)^m = (-1)^{\frac{p-1}{m}}\, p\, \psi_1(r)\,\psi_2(r)\cdots\psi_{m-2}(r),}
from which it appears that these $m$th powers belong to the field $\Omega_m$.

We now denote, as before, by $\mathfrak{p}_n$ the various
prime divisors of $p$ in the field $\Omega_m$. The numbers $\psi_s(r)$, which likewise belong to this field and are factors of $p$, must
%CONT>>


% ===== scan leaf 655, printed 634 =====
% !PRINTED 634
\weberpage{634}
%<<CONT
therefore be divisible by some of these prime factors, and they can, no more than $p$ itself, be divisible by the square of a $\mathfrak{p}_n$. We must establish by which prime divisors $\mathfrak{p}_n$ each of these numbers $\psi_s(r)$ is divisible.

To this end, however, we are led on the one hand by the congruences of the preceding section:
\weq{12}{r \equiv g^{-n\frac{p-1}{m}} \ (\mathrm{mod}\ \mathfrak{p}_{n'}), \qquad n n' \equiv 1\ (\mathrm{mod}\ m),}
and on the other hand by the congruences (17) derived in \S.\,170 of the first volume:
\weq{13}{\begin{aligned}
&\psi_{\lambda,\mu}(g) \equiv 0, && \lambda+\mu<p-1\\[4pt]
&\psi_{\lambda,\mu}(g) \equiv -\dfrac{\Pi(2p-\lambda-\mu-2)}{\Pi(p-\lambda-1)\,\Pi(p-\mu-1)}, && (\mathrm{mod}\ p)\\
&&& p-1<\lambda+\mu<2p-2,
\end{aligned}}
in which $\lambda$ and $\mu$ are taken between $0$ and $p-1$, so that the binomial coefficient % [sic] German source prints "Binominalcoëfficient"; "Binomialcoëfficient" is meant
\[
\frac{\Pi(2p-\lambda-\mu-2)}{\Pi(p-\lambda-1)\,\Pi(p-\mu-1)}
\]
is an integer not divisible by $p$.

By the congruence (12) we have
\[
\psi_s(r) \equiv \psi_s(g^{-nm'}) \equiv \sum\limits^{t} g^{-nm'\,[\mathrm{ind}\,t - (s+1)\,\mathrm{ind}\,(t+1)]},
\]
or, by (6):
\[
\psi_s(r) \equiv \psi_{\lambda,\mu}(g) \ (\mathrm{mod}\ \mathfrak{p}_{n'}),
\]
if
\weq{14}{\begin{aligned}
\lambda &\ \text{denotes the least positive residue of } {-}nm's\\
\mu &\ \text{\textquotedbl\ \textquotedbl\ \textquotedbl\ \textquotedbl\ \textquotedbl}\ {-}nm'
\end{aligned}\quad (\mathrm{mod}\ p-1)}
and the congruences (13) then show that $\psi_s(r)$ is divisible by $\mathfrak{p}_{n'}$ if $\lambda+\mu<p-1$, and not divisible if $\lambda+\mu>p-1$.

Since the number $n$, coprime to $m$, is determined only with respect to the modulus $m$, we now take it \gesperrt{positive and less than $m$}. Then
\[
\mu = m'(m-n).
\]

The number $\lambda$ determined from (14) is a multiple of $m'$, and if we set $\lambda = m'\varkappa$, then
\weq{15}{\varkappa \ \text{is the least positive residue of } {-}ns \ (\mathrm{mod}\ m).}

According to this, $\psi_s(r)$ will be divisible by $\mathfrak{p}_{\varkappa'}$ if
% [sic] Source prints the index varkappa' in this sentence; the preceding derivation and following factor count indicate mathfrak p_{n'}. Preserve the printed index.
\[
m'(\varkappa+m-n) < m m'
\]
or $\varkappa < n$, and not divisible if $\varkappa > n$.


% ===== scan leaf 656, printed 635 =====
% !PRINTED 635
\weberpage{635}
Since we are now in a position to determine the prime factors of $\psi_s(r)$ for every $s$, we proceed to decompose the number $(r,\eta)^m$ into its prime factors according to formula (11).

For this purpose we form, according to (15), the least positive residues $\varkappa$ of the numbers, $-n, -2n, \cdots -(m-2)n$ with respect to the modulus $m$, and find, apart from the order, the numbers
\[
\varkappa = 1, 2, \ldots, n-1, n+1, \ldots, m-1;
\]
the residues $0$ and $n$ do not occur among them, because $s$ receives no value for which $s$ or $s+1$ becomes divisible by $m$. Of these numbers, however, exactly $n-1$ are smaller than $n$, and so the factor $\mathfrak{p}_{n'}$ occurs that many times in the product $\psi_1\psi_2\ldots\psi_{m-2}$. Then $p$ contains this factor once more, and consequently, by (11), it occurs exactly $n$ times in $(r,\eta)^m$. Since no prime factors other than the $\mathfrak{p}_{n'}$ can divide $(r,\eta)^m$, the decomposition is hereby completely carried out:
\weq{16}{(r,\eta)^m = \prod^{n} \mathfrak{p}_{n'}^{\,n};}
here $n$ runs through the sequence of positive numbers that are smaller than $m$ and relatively prime to $m$, and from the congruence
\weq{17}{n n' \equiv 1\ (\mathrm{mod}\ m)}
$n'$ is in each case to be determined\footnote{This theorem is due to \gesperrt{Kummer}. Cf. Theorie der idealen Primfactoren der complexen Zahlen etc. Abhandlungen der Berliner Akademie, 1856.}.

Formula (16) holds for an even as well as for an odd $m$, and is still valid also in the case $m=2$, for which it gives a result already known, since then $(r,\eta) = (-1,\eta)$ agrees with a \gesperrt{Gauss} sum (Bd.\ I, \S.\,171).

But now, excluding the case $m=2$ from this point on, we shall give formula (16) a somewhat more general shape.

Let $m=q^{\varkappa}$ be an arbitrary prime power, only greater than $2$, and let $p$ be a prime different from $q$, with the property that $p-1$ is divisible by $q$ or, if $q=2$, by $4$.

Further, let $m_1$ be the greatest common divisor of $p-1$ and $m$, and
\[
m = m_1 m_2,
\]
%CONT>>


% ===== scan leaf 657, printed 636 =====
% !PRINTED 636
\weberpage{636}
%<<CONT
then by \S.\,172, III. the decomposition of $p$ into prime factors in the field $\Omega_m$ is the same as in the field $\Omega_{m_1}$, and we have
\[
\mathfrak{p}_n = \mathfrak{p}_{n+sm_1},
\]
if $s$ is an arbitrary integer. We now set
\weq{18}{t = n + s m_1}
and let $n$ run through the sequence of the positive integers $<m_1$ not divisible by $q$, and $s$ through the sequence of numbers $0,1,2,\ldots,m_2-1$; then $t$ runs through a complete system of residues of $m$ not divisible by $q$, and indeed the completely determinate system whose elements are all $<m$. We then have by \S.\,172
\weq{19}{p = \prod^{n} \mathfrak{p}_n.}

If, further, $t'$ is a number which satisfies the congruence
\weq{20}{t\,t' \equiv 1\ (\mathrm{mod}\ m),}
then, if $n'$ is determined by the condition $n n' \equiv 1\ (\mathrm{mod}\ m_1)$,
\weq{21}{t' \equiv n'\ (\mathrm{mod}\ m_1).}

According to this the product
\[
\prod^{t} \mathfrak{p}_{t'}^{\,t}
\]
can be determined by the formula (16), which, applied to the field $\Omega_{m_1}$, consisting of all rational functions of the $m_1$th root of unity $r_1 = r^{m_2}$, yields the formula:
\weq{22}{(r_1,\eta)^{m_1} = \prod^{n} \mathfrak{p}_{n'}^{\,n}.}

But by (18) and (21)
\[
\prod^{t} \mathfrak{p}_{t'}^{\,t} = \prod^{n} \mathfrak{p}_{n'}^{\,n m_2} \prod^{n'}\prod^{s} \mathfrak{p}_{n'}^{\,s m_1},
\]
and since $n'$, according to the modulus $m_1$, runs through the same sequence of numbers as $n$, we have by (19):
\[
\prod^{s}\prod^{n'} \mathfrak{p}_{n'}^{\,s m_1} = p^{m_1 \Sigma s} = p^{\frac{1}{2}m(m_2-1)},
\]
and there results:
\[
\prod^{t} \mathfrak{p}_{t'}^{\,t} = \left[p^{\frac{1}{2}(m_2-1)}\,(r^{m_2},\eta)\right]^m.
\]

If we set
\weq{23}{\varrho = p^{\frac{1}{2}(m_2-1)}\,(r_1,\eta), \qquad \varrho_n = p^{\frac{1}{2}(m_2-1)}\,(r_1^{\,n},\eta),}
then $\varrho$ is a \gesperrt{cyclotomic number}, since indeed $\sqrt{p}$ too, as a value
%CONT>>


% ===== scan leaf 658, printed 637 =====
% !PRINTED 637
\weberpage{637}
%<<CONT
of a \gesperrt{Gauss} sum belongs to the cyclotomic numbers,
and we have
\weq{24}{\varrho^m = \Pi^{t}\, \mathfrak{p}^t_{t'}.}

But by (9) and (11), if $m$ in them is replaced by $m_1$,
\weq{25}{\varrho_n^{m_1} = \pm\, p^{\frac{1}{2}\,m_1(m_2-1)+1}\, \psi_1(r_1^n)\,\dots\,\psi_{m_1-2}(r_1^n),}
\weq{26}{\varrho^n \varrho_n^{-1} = p^{\frac{1}{2}(n-1)(m_2-1)}\, \psi_1(r_1)\,\psi_2(r_1)\,\dots\,\psi_{n-1}(r_1)}
and thus $\varrho^n \varrho_n^{-1}$, $\varrho_n^{m_1}$, and consequently also $\varrho_n^m$, are numbers
of the field $\Omega_m$, and by (25) the combination $\varrho^n \varrho_n^{-1}$ occurring
in (26) remains contained in $\Omega_m$ when $n$ is increased by a
multiple of $m_1$. Accordingly $\varrho^n \varrho_n^{-1}$ is still a number in $\Omega_m$
even when $n$ does not have precisely the restricted meaning it
has had hitherto, but denotes an arbitrary number not divisible
by $q$.

The purpose of this consideration is to make \gesperrt{Kummer}'s
theorem (16) applicable simultaneously to several different primes
which belong to different values of $m_1$, and the result is then
expressed in the following theorem:

3. \gesperrt{Let $\varphi_n$ be a system of conjugate functionals of
the field $\Omega_m$ which is related to no other
natural primes than such as are congruent to 1
with respect to the modulus $q$ (or, for $q = 2$,
with respect to the modulus 4); further, let $t$
run through those numbers of the sequence
$1, 2, \dots, m - 1$ that are not divisible by $q$, and
let $t'$ be determined from the congruence
$t\,t' \equiv 1 \pmod{m}$.}

\gesperrt{Then the functional of the field $\Omega_m$:}
\[ \text{a)}\qquad \Phi_n = \Pi^{t}\, \varphi^t_{nt'} = \varepsilon\,\vartheta_n^m \]
\gesperrt{is associated with the $m$th power of a cyclotomic
number $\vartheta_n$, which is indeed itself contained in a
higher field, but whose $m$th power belongs to the
field $\Omega_m$, and}
\[ \text{b)}\qquad \vartheta_n^{-1}\vartheta_1^n = \alpha \]
\gesperrt{is a number of the field $\Omega_m$.}

The proof is contained in the foregoing discussions and
follows from the formulae (25), (26), if one
%CONT>>


% ===== scan leaf 659, printed 638 =====
% !PRINTED 638
\weberpage{638}
\pdfbookmark[2]{§ 175. The roots of unity in the field Ωm}{bm200}
%<<CONT
understands by $\vartheta$ a product of powers of the numbers $\varrho$ which
are derived according to the formulae (24) from the prime factors $\mathfrak{p}$
contained in $\varphi$.

\begin{center}
§. 175.

\gesperrt{The Roots of Unity in the Field $\Omega_m$.}
\end{center}

For the applications to be made further on, a more precise
knowledge of the units of the field $\Omega_m$ is required, which indeed
is also of great interest in itself. We concern ourselves
here not with the functional but only with the numerical
units, by which, as we recall, \gesperrt{are to be understood integers of the
field $\Omega_m$ whose norm is $\pm 1$,
or an integer whose reciprocal value is likewise integral.}

Among the units are certainly all the roots of unity contained
in $\Omega_m$; but we can easily prove that these
roots of unity can only be powers of $r$, with positive
and negative sign.

Namely, let $\varrho$ be a root of unity contained in $\Omega_m$, and
\weq{1}{\varrho = \varphi(r).}

If $q^h$ is the highest power of $q$ dividing the degree of $\varrho$,
then $\varrho^{q^h}$ is a root of unity whose degree is not divisible by $q$,
and
\weq{2}{\varrho^{q^h} = [\varphi(r)]^{q^h}.}

On account of the irreducibility of the cyclotomic equation in
the wider sense\footnote{Cf. Bd.~I, §.\,134 and the supplement to this volume.} all the substitutions
 $(r, r^n)$ can therefore be made in (2) without the left side
changing, and consequently $\varrho^{q^h}$ is a rational number, and since it
is a unit, it must be $= \pm 1$. Hence $\varrho$, apart from sign,
is equal to a root of unity of degree $q^h$ or, when $q = 2$,
of degree $q^{h+1}$. It is to be shown that $q^h$ or
$q^{h+1}$ cannot be greater than $m$. But this results
from (1). For if $\varrho$ were a root of unity of degree higher than $m$,
then $r$ would be a power of $\varrho$, and $\varrho$ would have to satisfy an irreducible
 equation of higher degree than $r$, and hence also of higher
degree than $\varphi(r)$, which by (1) is a contradiction. Thus
we have the theorem:
%CONT>>


% ===== scan leaf 660, printed 639 =====
% !PRINTED 639
\weberpage{639}
%<<CONT
1. \gesperrt{The only roots of unity contained in $\Omega_m$
are the powers of $r$ taken with positive and negative
sign.}

To this attaches a quite general theorem on the roots of
unity, valid, that is, in every algebraic number field, which
we shall now prove.

The theorem reads:

2. \gesperrt{If $\alpha$ is an integer of an algebraic number field
 of degree $n$, and if all the numbers $\alpha_1, \alpha_2, \dots, \alpha_n$
 conjugate to $\alpha$ have the absolute
value 1, then $\alpha$ is a root of unity}\footnote{\gesperrt{Kronecker}, ``Zwei Sätze über Gleichungen mit ganzzahligen
Co\"efficienten''. Crelle's Journal, Bd. 53 (1857). \gesperrt{Minkowski}, ``Geometrie
der Zahlen'', Art. 43.}.

For the proof it is first to be remarked that the absolute
value of the norm of an integer $\omega$ is equal to the product of the
absolute values of the conjugate numbers $\omega_1, \omega_2, \dots, \omega_n$, and,
being an integral rational number, is in any case greater than or
equal to 1. It is therefore not possible that all these absolute
values are less than 1. But if they are all equal to 1, as for the
number $\alpha$ assumed in our theorem, then the norm must be $= \pm 1$,
and $\alpha$ is a unit. If $\beta$ is a second unit which, together with its
conjugates, has the absolute value 1, then the quotient $\alpha : \beta$ has
the same property. If therefore, among the numbers conjugate to
this quotient, even a single one is real, we must have
\[ \frac{\alpha}{\beta} = \pm 1 \]
or $\alpha = \pm \beta$, and this then holds also when $\alpha, \beta$
are replaced by the corresponding numbers of a conjugate
field.

If then $\alpha$ and $\beta$ are neither equal nor opposite, their
ratio is not real, and there accordingly subsists between the
absolute values the inequality:
\[ |\alpha \pm \beta| < |\alpha| + |\beta|. \]
(Cf. the Introduction to the first Band, S. 19.)

We have therefore
\[ |\alpha \pm \beta| < 2, \]
%CONT>>


% ===== scan leaf 661, printed 640 =====
% !PRINTED 640
\weberpage{640}
%<<CONT
and consequently $\frac{1}{2}(\alpha \pm \beta)$ cannot be an integer, because the
absolute value of the norm of this number is less than 1. Now
if $\omega_1, \omega_2, \dots, \omega_n$ is a basis of $\mathfrak{o}$, and
\[
\begin{aligned}
\alpha &= a_1\omega_1 + a_2\omega_2 + \dots + a_n\omega_n \\
\beta &= b_1\omega_1 + b_2\omega_2 + \dots + b_n\omega_n,
\end{aligned}
\]
then the integral rational numbers $a, b$ cannot satisfy the congruences
\[ a_1 \equiv b_1,\ a_2 \equiv b_2,\ \dots,\ a_n \equiv b_n \pmod{2} \]
since otherwise $\frac{1}{2}(\alpha - \beta)$ would be an integer.

If, then, we distribute all the numbers $\alpha$ into $2^n$ compartments,
by throwing into one compartment all the numbers in which
$a_1, a_2, \dots, a_n$ leave the same residues (0 or 1) with respect to the modulus 2,
then in each of these compartments at most 2 numbers,
namely $\alpha$ and $-\alpha$, can occur; and if we further exclude
$\alpha = 0$, then in one of these compartments, namely that in which
the $a_1, a_2, \dots, a_n$ are all even, there is in fact no number $\alpha$ at all. The number
of all possible numbers $\alpha$ is therefore finite, and at most
$= 2^{n+1} - 2$.

If now the absolute value of $\alpha = 1$, then the same holds
of all the powers of $\alpha$. Consequently, in the unbounded
sequence
\[ 1,\ \alpha,\ \alpha^2,\ \alpha^3,\ \dots \]
the same number must necessarily recur a second time, and hence
$\alpha^h = \alpha^k$ and $k > h$. But then
\[ \alpha^{k-h} = 1, \]
that is, $\alpha$ is a root of unity, as was to be proved.

We further add the remark that if a number $\alpha$
of a field $\Omega$ is a root of unity of degree $m$, then all
the numbers conjugate to $\alpha$ are also roots of unity of the same degree.
For in the rational equation $\alpha^m - 1 = 0$, $\alpha$ can be
replaced by each of the conjugate values.

The number of roots of unity that are contained in a field $\Omega$
can only ever be a finite one, because every
number in $\Omega$ satisfies a rational equation whose degree is
at most equal to the degree of the field. The degree of the roots of unity
in $\Omega$ therefore cannot exceed a finite value.


% ===== scan leaf 662, printed 641 =====
% !PRINTED 641
\weberpage{641}
\pdfbookmark[2]{§ 176. The real field Hm contained in Ωm}{bm201}
\begin{center}
\S.\,176.\\[4pt]
\gesperrt{The real field $H_m$ contained in $\Omega_m$.}
\end{center}

Every number of the field $\Omega_m$ passes by the substitution $(r, r^{-1})$ into the conjugate imaginary number, which is also obtained by the interchange $(i, -i)$. The product of two such conjugate imaginary numbers is the \gesperrt{square of the absolute value} of each one of them.

A number in $\Omega_m$ is real when it remains unchanged under the interchange $(r, r^{-1})$, and only under this assumption. Every such number can be expressed rationally in terms of the \gesperrt{two-term period} $r + r^{-1}$, and therefore belongs to a real field of degree $\frac12\varphi(m)$, which is a divisor of $\Omega_m$. This field we will call the \gesperrt{real field contained in $\Omega_m$} and denote by $H_m$ (although still other real fields, namely all divisors of $H_m$, are contained in $\Omega_m$ as well).

The $\frac12\varphi(m) = \frac12\mu$ numbers
\weq{1}{1,\ r+r^{-1},\ r^2+r^{-2},\ \dots,\ r^{\frac12\mu-1}+r^{-\frac12\mu+1}}
form a basis of the integers of the field $H_m$. For according to \S.\,170, (7),
\weq{2}{\omega = x_0+x_1r+x_2r^2+\cdots+x_{\frac12\mu-1}r^{\frac12\mu-1}+x_{\frac12\mu}r^{\frac12\mu} +x_{-1}r^{-1}+x_{-2}r^{-2}+\cdots+x_{-\frac12\mu+1}r^{-\frac12\mu+1}}
is then and only then an integer of the field $\Omega_m$, when the coefficients $x_0, x_1, \dots$ are integral rational numbers.

The condition for this number to belong to the field $H_m$ is obtained by carrying out the substitution $(r, r^{-1})$ and setting the difference $=0$. If we further divide by $r-r^{-1}$, there results in this way:
\[
(x_1-x_{-1})+(x_2-x_{-2})(r+r^{-1})+\cdots
+(x_{\frac12\mu-1}-x_{-\frac12\mu+1})\frac{r^{\frac12\mu-1}-r^{-\frac12\mu+1}}{r-r^{-1}}
+x_{\frac12\mu}\frac{r^{\frac12\mu}-r^{-\frac12\mu}}{r-r^{-1}}=0.
\]
The divisions can here be carried out, and if we then multiply by $r^{\frac12\mu-1}$, there results for $r$ a rational equation of degree lower than $\mu$, which can only
%CONT>>


% ===== scan leaf 663, printed 642 =====
% !PRINTED 642
\weberpage{642}
%<<CONT
be satisfied when all of its coefficients vanish. This leads to the equations
\[
x_{\frac12\mu}=0,\ x_{\frac12\mu-1}=x_{-\frac12\mu+1},\ \dots,\ x_1=x_{-1},
\]
and every integer $\omega$ of the field $H_m$ is therefore contained in the form:
\[
\omega = x_0+x_1(r+r^{-1})+x_2(r^2+r^{-2})+\cdots
+x_{\frac12\mu-1}(r^{\frac12\mu-1}+r^{-\frac12\mu+1}).
\]

Instead of the basis (1) one may also, if one sets
\[
\varrho = r+r^{-1}
\]
choose the powers of $\varrho$:
\weq{3}{1,\ \varrho,\ \varrho^2,\ \dots,\ \varrho^{\frac12\mu-1}}
as a basis of the integers of $H_m$. For the quantities (3) can be expressed with integral coefficients in terms of (1), and conversely.

The number $\varrho$ is the root of an irreducible equation of degree $\frac12\mu$, which we shall denote by $\psi(\varrho)=0$. If $f(r)=0$ is the equation for $r$ (\S.\,169), then there subsists the identical relation
\weq{4}{f(x) = x^{\frac12\mu}\,\psi\Bigl(x+\frac1x\Bigr),}
from which, by formation of the derivative,
\weq{5}{f'(r) = r^{\frac12\mu}\,\psi'(\varrho)\,(1-r^{-2}).}

From this we can form the fundamental number $\Delta_1$ of the field $H_m$, which, since $H_m$ contains only real numbers, must be positive. This fundamental number, if we denote the norm with respect to $H_m$ by $N_1$, is:
\[
\Delta_1 = \pm\,N_1\,\psi'(\varrho).
\]

But if $N$ is the norm formed with respect to $\Omega_m$, then
\[
N\,\psi'(\varrho) = [N_1\,\psi'(\varrho)]^2 = \Delta_1^2,
\]
and accordingly formula (5), having regard to \S.\,169, yields:
\[
\Delta = N(1-r^{-2})\,\Delta_1^2.
\]

By \S.\,169, (7) and (10) we have:
\weq{6}{
\begin{array}{rll}
N(1-r^{-2}) &{}=q & \text{for odd } q\\
&{}=4 & \text{for } q=2,\\[4pt]
\pm\,\Delta &{}=q\,\Delta_1^2 & \text{for odd } q\\
&{}=4\,\Delta_1^2 & \text{for } q=2.
\end{array}
}


% ===== scan leaf 664, printed 643 =====
% !PRINTED 643
\weberpage{643}
\pdfbookmark[2]{§ 177. The prime ideals of the field Hm}{bm202}
If one sets therein
\[
\Delta = \pm\,q^{\,q^{\varkappa-1}[\varkappa(q-1)-1]}
\]
it follows that
\weq{7}{
\begin{array}{rl}
\Delta_1 = q^{\varkappa\, q^{\varkappa-1}\frac{q-1}{2} - \frac{q^{\varkappa-1}+1}{2}} & \text{for odd } q\\
= 2^{(\varkappa-1)2^{\varkappa-2}-1} & \text{for } q=2,
\end{array}
}
so that $\Delta_1$ is always a power of $q$.

\begin{center}
\S.\,177.\\[4pt]
\gesperrt{The Prime Ideals in the Field $H_m$.}
\end{center}

We have now to investigate the question in what relation the prime ideals of the field $H_m$ stand to those of the field $\Omega_m$.

If $\mathfrak{P}$ denotes any prime ideal in $H_m$ of degree $f_1$, then $\mathfrak{P}$ must be a divisor of a natural prime $p$, and $\mathfrak{P}$ can therefore (by \S.\,171) be divisible by the second or a higher power of a prime factor $\mathfrak{p}$ in $\Omega_m$ only when $p=q$.

But if $\mathfrak{P}$ is a divisor of $q$, it must be a power of $\sigma=1-r$. We set, say,
\[
\mathfrak{P}=\sigma^\lambda.
\]
If we take the norm of this with respect to $\Omega_m$, which is the square of the norm with respect to $H_m$, it follows (by \S.\,169) that
\[
q^{2f_1}=q^\lambda,
\]
hence $f_1=\frac12\lambda$. Since $f_1$ is an integer, $\lambda$ must be at least $=2$. But $\lambda$ also cannot be greater than $2$, since $\sigma^2$ is associated with the integer $(1-r)(1-r^{-1})$ contained in $H_m$, and hence $\sigma^2$ must be divisible by $\mathfrak{P}$. Accordingly $f_1=1$, and we obtain the theorem:

1. \gesperrt{The prime $q$ is the $\frac12\mu$th power of a prime of the first degree existing in $H_m$.}

Now let $p$ be different from $q$, and $\mathfrak{p}$ a prime factor of $\mathfrak{P}$ in the field $\Omega_m$ of degree $f$, which goes over into $\mathfrak{p}'$ under the substitution $(r,r^{-1})$. Then $\mathfrak{P}$ is divisible both by $\mathfrak{p}$ and by $\mathfrak{p}'$; on the other hand, $\mathfrak{p}\mathfrak{p}'$ is contained in $H_m$ and hence divisible by $\mathfrak{P}$, and we have to distinguish whether $\mathfrak{p}, \mathfrak{p}'$ are different or not.

If $\mathfrak{p}$ is different from $\mathfrak{p}'$, then $\mathfrak{P}$ is divisible by $\mathfrak{p}\mathfrak{p}'$, and it follows that $\mathfrak{P}=\mathfrak{p}\mathfrak{p}'$, and by formation of the norm $f_1=f$.


% ===== scan leaf 665, printed 644 =====
% !PRINTED 644
\weberpage{644}
But if $\mathfrak{p}=\mathfrak{p}'$, then $\mathfrak{P}=\mathfrak{p}$, and the formation of the norm yields $2f_1=f$.

Which of these two cases occurs therefore depends on whether the prime ideal $\mathfrak{p}$ admits the substitution $(r,r^{-1})$ or not, i.e. on whether $(r,r^{-1})$ is or is not contained in the group to which the prime ideal belongs. According to \S.\,172 this distinction is conditioned by whether there is any exponent $h$ for which the congruence
\[
p^h \equiv -1 \ (\mathrm{mod}\ m)
\]
is fulfilled, or whether there is no such exponent.

We summarise the result as follows:

2. \gesperrt{The natural primes $p$ different from $q$ split into two kinds $p_1, p_2$.}

\gesperrt{The first kind $p_1$ is characterised by the fact that for some exponent $h$}
\[
p_1^h \equiv -1 \ (\mathrm{mod}\ m),
\]
\gesperrt{and these primes, if they belong to the exponent $f$ for the modulus $m$ and $\mu=ef$ is set, split in the field $H_m$ into $e$ prime factors of degree $\frac12 f$.}

\gesperrt{Their prime factors are not further decomposable in $\Omega_m$ either.}

\gesperrt{The primes of the second kind $p_2$ are determined by the fact that none of their powers becomes congruent to $-1$ with respect to the modulus $m$, and they split in the field $H_m$ into $\frac12 e$ prime factors of degree $f$. Their prime factors are still decomposable in $\Omega_m$ into two factors each.}

For the primes of the first kind $f$ must certainly be an even number. If $m$ is odd, it also suffices, in order that $p$ be a prime of the first kind, that it belong to an even exponent; for then
\[
(p^{\frac12 f}-1)(p^{\frac12 f}+1)\equiv 0\ (\mathrm{mod}\ m).
\]
But the first of these factors $p^{\frac12 f}-1$ is not divisible by $m$, because otherwise $p$ would belong to the exponent $\frac12 f$. Consequently $p^{\frac12 f}+1$ must be divisible by $q$. But here the two factors $p^{\frac12 f}-1$ and $p^{\frac12 f}+1$ cannot both be divisible by $q$, because otherwise their difference, which is $=2$, would also be divisible by $q$; and consequently $p^{\frac12 f}+1$ is divisible by $m$.


% ===== scan leaf 666, printed 645 =====
% !PRINTED 645
\weberpage{645}
\pdfbookmark[2]{§ 178. The units of the field Hm}{bm203}

But if $m$ is a power of 2, the congruence
\[
p^h \equiv -1 \ (\mathrm{mod}\ m)
\]
is possible only when $p \equiv -1\ (\mathrm{mod}\ m)$. For every even power of an odd number is $\equiv + 1\ (\mathrm{mod}\ 8)$ and therefore cannot be $\equiv -1\ (\mathrm{mod}\ m)$. But if $h$ is odd, then
\[
\frac{p^h + 1}{p + 1} = p^{h-1} - p^{h-2} + \cdots + 1
\]
is itself odd, and consequently $p^h + 1$ is divisible by no higher power of 2 than $p + 1$. Accordingly we may add to Theorem 2. the following closer determination:

3. \gesperrt{If $m$ is odd, the primes $p_1$ belong to an even exponent, the primes $p_2$ to an odd exponent.}

\gesperrt{If $m$ is even, the primes of the form $km - 1$ belong to the first kind, and all other odd primes to the second.}

\begin{center}
\S.\,178.\\[4pt]
\gesperrt{The Units of the Field $H_m$.}
\end{center}

The units belonging to the field $H_m$ are of particular importance for the applications. The units of the field $\Omega_m$ can also be reduced to them according to the following theorem:

1. \gesperrt{Every unit of the field $\Omega_m$ is the product of a root of unity and a real unit of the field $\Omega_m$.}

If we denote any unit of the field $\Omega_m$ by $\mathfrak{E}(r)$, then $\mathfrak{E}(r^{-1})$ is the conjugate imaginary value, which likewise represents a unit, and the quotient $\mathfrak{E}(r) : \mathfrak{E}(r^{-1})$ is again a unit, whose absolute value
\weq{1}{\sqrt{\dfrac{\mathfrak{E}(r)}{\mathfrak{E}(r^{-1})}\ \dfrac{\mathfrak{E}(r^{-1})}{\mathfrak{E}(r)}}}
is equal to 1, and consequently this quotient is a root of unity and hence $= \pm r^h$, where $h$ is any integer exponent (§. 175, Theorem 1., 2.). We therefore have
\weq{2}{\mathfrak{E}(r) = \pm r^h \mathfrak{E}(r^{-1}).}


% ===== scan leaf 667, printed 646 =====
% !PRINTED 646
\weberpage{646}

In the further course of the proof a small distinction must now be made, according as $m$ is odd or even.

If $m$ is, to begin with, \gesperrt{odd}, then in (2) we can assume the exponent $h$ to be even, since we can if necessary replace it by $h + m$. Further, in formula (2) only the upper sign is admissible. For if the lower sign stood there, it would follow that:
\[
\mathfrak{E}(r) \equiv \mathfrak{E}(1) \equiv -\mathfrak{E}(1), \quad 2\,\mathfrak{E}(1) \equiv 0\ [\mathrm{mod}\ (1 - r)].
\]

By §. 169, however, $1 - r$ is a divisor of $q$, hence relatively prime to 2, and therefore $\mathfrak{E}(1)$, and consequently also $\mathfrak{E}(r)$, is divisible by $(1 - r)$. But this contradicts the assumption that $\mathfrak{E}(r)$ is to be a unit.

Accordingly, from (2) we obtain:
\[
r^{-1/2\,h}\, \mathfrak{E}(r) = r^{1/2\,h}\, \mathfrak{E}(r^{-1}),
\]
from which it follows that $r^{-1/2\,h}\, \mathfrak{E}(r)$ is a real unit. If we denote it by $e(r)$, we thus have
\weq{3}{\mathfrak{E}(r) = r^{1/2\,h}\, e(r),}
by which Theorem 1. is proved for this case.

But if, secondly, $m$ \gesperrt{is a power of 2}, then $\mu = \tfrac{1}{2} m$ and $r^\mu = -1$, and together with $r$, $-r$ is also an $m$th root of unity. If we replace $h$ if necessary by $h + \tfrac{1}{2} m$, we can assume the upper sign in (2) and set:
\weq{4}{\mathfrak{E}(r) = r^h\, \mathfrak{E}(r^{-1}).}

It is now a matter of proving that $h$ must be even.

If we represent $\mathfrak{E}(r)$ by the basis $1, r, r^2, \ldots r^{\mu-1}$, we obtain
\weq{5}{\mathfrak{E}(r) = \sum_{0,\,\mu-1}^{s} x_s\, r^s,}
in which the $x_s$ are integral rational numbers. But by (4) we then have
\[
\sum_{0,\,\mu-1}^{s} x_s\, r^s = \sum_{1,\,\mu-1}^{s} x_s\, r^{h-s}, % [sic] source prints lower limit 1 on the right; substitution in (4) requires 0
\]
and from this it follows that $x_s = \pm x_{s'}$ when $s + s' \equiv h\ (\mathrm{mod}\ \mu)$. Now if $h$ is \gesperrt{odd}, then of two numbers so connected the one is even, the other odd (since $\mu$ is a power of 2). The expression (5) therefore yields for $\mathfrak{E}(r)$ an aggregate of terms of the form
\[
x_s\,(r^s \pm r^{s'}).
\]
%CONT>>


% ===== scan leaf 668, printed 647 =====
% !PRINTED 647
\weberpage{647}
%<<CONT
But $r^s \pm r^{s'} = r^s (1 \pm r^{h-2s}) = r^s (1 \mp r^{\mu+h-2s})$ is always divisible by $1 - r$, and from this it would follow that $\mathfrak{E}(r)$ would be divisible by $1 - r$, which is not a unit, and this contradicts the concept of the unit. Hence the assumption that in formula (4) the exponent $h$ is odd is inadmissible, and, as above, we can set
\[
\mathfrak{E}(r) = r^{\tfrac12 h}\, e(r)
\]
where $e(r)$ is a real unit. With this, Theorem 1 is proved.

We wish to single out a system of real units: According to §. 169 the $\mu$ quantities $1 - r^n$ are all associated with one another. Accordingly the quotient
\[
\frac{1 - r^n}{1 - r^{n'}},
\]
in which $n, n'$ denote two numbers relatively prime to $m$, is a unit. But from this one derives, when
\[
r = e^{\frac{2\pi i}{m}}
\]
is set, the real unit
\weq{6}{\dfrac{\sin \dfrac{n\pi}{m}}{\sin \dfrac{n'\pi}{m}} = \dfrac{r^{\frac{n}{2}} - r^{-\frac{n}{2}}}{r^{\frac{n'}{2}} - r^{-\frac{n'}{2}}} = r^{\frac{n'-n}{2}}\, \dfrac{1 - r^n}{1 - r^{n'}}}
If $m$ is even, one can set $n' = n + \dfrac{m}{2}$, and obtains for this case the real units
\weq{7}{\tau_n = \mathrm{tang}\,\dfrac{n\pi}{m},}
in which $n$ can denote any odd number.

If one replaces $n$ by $\tfrac{1}{2} m - n$, then $\tau_n$ goes over into the reciprocal value. If one lets $n$ run through the sequence of numbers
\[
1,\ 3,\ 5,\ \cdots,\ \frac{m}{4} - 1
\]
then $\tau_n$ takes nothing but positive properly fractional values.

\begin{center}
\rule{0.2\linewidth}{0.4pt}
\end{center}


% ===== scan leaf 669, printed 648 =====
\weberpage{648}% page number not printed on this page
\pdfbookmark[1]{Twenty-second Section. Abelian Fields and Cyclotomic Fields.}{bm204}
\pdfbookmark[2]{§ 179. Simple Abelian fields}{bm205}
% !PRINTED none

\begin{center}
\gesperrt{Twenty-Second Chapter.}

\gesperrt{Abelian Fields and Cyclotomic Fields.}

\rule{3cm}{0.4pt}

\S.\,179.

\gesperrt{Simple Abelian Fields.}
\end{center}

We turn now to one of the most interesting applications of the theory of algebraic numbers.

\begingroup
\renewcommand{\thefootnote}{\arabic{footnote})}
It concerns the proof of the general theorem\footnote{\gesperrt{Kronecker} first stated this theorem, but published no complete proof of it. The first proof was made known by the author of the present work in Bd. 8 of the Acta Mathematica (1886). Most recently a second proof by \gesperrt{Hilbert} has been published, which rests upon the theorems developed in the nineteenth chapter (Nachrichten d. Gesellschaft d. Wissenschaften zu Göttingen, 1896, Heft 1).}:
\endgroup

I. \gesperrt{All number fields that are Abelian in the absolute domain of rationality are cyclotomic fields.}

Once we have proved this theorem, the investigations of the third chapter on cyclotomic fields gain a heightened interest, because it is thereby shown that on the path indicated there not only all cyclotomic fields, but all \gesperrt{Abelian} fields whatever, are found.

Let $\Omega = R(x)$ be an \gesperrt{Abelian} field of degree $m$, that is, a field which consists of all rational functions of a root $x$ of an irreducible \gesperrt{Abelian} equation of degree $m$, when the field $R$ of the rational numbers is taken as the domain of rationality. The \gesperrt{Galois} group $\mathfrak{G}$ of this equation, which is thus an \gesperrt{Abelian} group and has the same degree $m$ as the field $\Omega$, is also the group of the field $\Omega$.

If $x_1$ is a non-primitive number from $\Omega$, then $\Omega_1 = R(x_1)$ is likewise an \gesperrt{Abelian} field, which we call a proper divisor of $\Omega$, and whose degree is a proper divisor of the degree of $\Omega$. For if $\mathfrak{A}$ is the group to which the number $x_1$ belongs,
%CONT>>


% ===== scan leaf 670, printed 649 =====
% !PRINTED 649
\weberpage{649}
%<<CONT
then $\mathfrak{G}\,|\,\mathfrak{A}$ (\S.\,4 of this volume) is the group of the field $\Omega_1$, and,
like $\mathfrak{G}$ itself, this too is an \gesperrt{Abelian} group. A divisor
of $\Omega$ which is of the same degree as $\Omega$ is identical with $\Omega$
(cf. Bd. I, \S.\,144, 149, 156).

If now there are two non-primitive numbers $x_1, x_2$ in $\Omega$ of
such a kind that we can set
\[
\Omega = R(x_1, x_2)
\]
then the fields $\Omega_1 = R(x_1)$, $\Omega_2 = R(x_2)$ are
proper divisors of $\Omega$, and $\Omega$ is said to be \gesperrt{composed} of the two fields
$\Omega_1, \Omega_2$.

If the field $\Omega$ admits no such representation, it is called
\gesperrt{simple}.

If, then, it is assumed as proved that $\Omega_1, \Omega_2$ are
cyclotomic fields, i.\,e. that all their numbers can be represented
rationally by roots of unity, the same follows for $\Omega$.

If one of the fields $\Omega_1, \Omega_2$ is decomposable, we can
repeat the decomposition, and since the degrees are decreasing
positive integers, we must finally arrive at simple
fields.

\gesperrt{Theorem I thus needs to be proved only for simple Abelian
fields.}

It is to be remembered, however, that the decomposability of a
field $\Omega$ by no means follows merely from the existence of a
divisor, and that for a decomposable field the decomposition into
simple fields can take place in quite different ways, so
that for fields the laws of divisibility do not hold
as in the region of numbers.

A criterion for the simplicity of a field we can
derive from its group.

If the group $\mathfrak{G}$ has two proper divisors $\mathfrak{A}$ and $\mathfrak{B}$ of such a
kind that every element can be composed once and only once from an
element of $\mathfrak{A}$ and an element of $\mathfrak{B}$, we call the group
$\mathfrak{G}$ decomposable into the two
components $\mathfrak{A}, \mathfrak{B}$; $\mathfrak{G}$ is thus decomposable if, in the product
\weq{1}{\mathfrak{G} = \mathfrak{A}\mathfrak{B}}
formed according to the composition of the parts (\S.\,4), every
element of $\mathfrak{G}$ appears once and only once. The degree
of $\mathfrak{G}$ is then equal to the product of the degrees of $\mathfrak{A}$ and $\mathfrak{B}$.

If there are no such divisors, the group $\mathfrak{G}$ is called
indecomposable.


% ===== scan leaf 671, printed 650 =====
% !PRINTED 650
\weberpage{650}

We can then prove the theorem:

1. \gesperrt{If the group of an Abelian field is decomposable, then the field is also composite.}

To prove this, let us assume that the group $\mathfrak{G}$ of the
field $\Omega$ is decomposable according to the formula (1), and that $m, a, b$
are the degrees of $\mathfrak{G}, \mathfrak{A}, \mathfrak{B}$. We take a number $\xi$ of the field $\Omega$
\gesperrt{belonging} to the group $\mathfrak{B}$, which therefore passes into $a$ different
conjugate values under the substitutions of $\mathfrak{G}$, and likewise
a $b$-valued number $\eta$ belonging to $\mathfrak{A}$. We can
then form a number
\[
x = \alpha \xi + \beta \eta
\]
with \gesperrt{rational} numerical coefficients $\alpha, \beta$, which has $m$ different
conjugate values, and then $\Omega = R(x)$. Hence
$\Omega$ is composed of $R(\xi)$ and $R(\eta)$ (Bd. I, \S.\,143).

If we now recall the representation of an Abelian group by
a basis derived in \S.\,9, 10 of this volume, then repeated
application of what has just been proved yields:

2. \gesperrt{Every Abelian field $\Omega$ can be composed of
fields whose degree is a prime power
and whose group is cyclic. The degrees
of these composing fields are the invariants
of the group of $\Omega$.}

These theorems are supplemented by the following one:

3. \gesperrt{An Abelian field of prime-power degree
with cyclic group is simple.}

If the group $\mathfrak{G}$ is cyclic, then its elements
can be represented by the powers of a basis, $G^\gamma$, when $\gamma$
runs through a complete system of residues with respect to the modulus $m$. We
obtain all divisors $\mathfrak{A}$ of $\mathfrak{G}$ represented by
\[
\mathfrak{A} = G^{b\gamma},
\]
when $b$ is a divisor of $m$ and $\gamma$ runs through a complete system of residues
with respect to the modulus $a = m : b$.

If
\[
\mathfrak{A}' = G^{b'\gamma}
\]
is a second divisor of $\mathfrak{G}$ of degree $a'$, and if
\[
b' \leqq b, \quad a' \geqq a,
\]
then, if we now assume that $m$, and hence also
$a, b, a', b'$, are powers of a prime $q$, $b'$ is a divisor of $b$,
and consequently $\mathfrak{A}$ is a divisor of $\mathfrak{A}'$.


% ===== scan leaf 672, printed 651 =====
% !PRINTED 651
\weberpage{651}
\pdfbookmark[2]{§ 180. The resolvents}{bm206}

If, then, $\xi$ is a number of the field $\Omega$ that belongs to the
group $\mathfrak{A}$, then $\xi$ is a primitive number of the field
$R(\xi)$ of degree $b$, and this field also contains every number $\xi'$
that admits the substitutions of the group $\mathfrak{A}'$. If now
$\mathfrak{A}, \mathfrak{B}$ and $\mathfrak{A}', \mathfrak{B}'$ are proper divisors of $\mathfrak{G}$, then $R(\xi)$ is different from
$\Omega$, and $R(\xi)$ is not extended by the adjunction of $\xi'$.
Hence $\Omega$ also cannot be composed from $R(\xi)$ and $R(\xi')$,
by which 3. is proved.

\begin{center}
\S.\,180.

\gesperrt{The Resolvents.}
\end{center}

After what has now been proved, we may restrict ourselves, in
the further considerations that have the proof of Theorem I. as
their goal, to those \gesperrt{Abelian} fields whose degree is a prime power
and whose group is cyclic. For this reason it also sufficed for
the intended application that in the preceding chapter we confined
ourselves to the consideration of the cyclotomic fields $\Omega_m$,
in which $m$ was a prime power. We retain, as far as possible,
the notation used there, and set
\weq{1}{m = q^{\varkappa},}
where $q$ is a prime and $\varkappa$ a positive exponent, and in the
case that $q = 2$ we assume $\varkappa > 1$. By $n$ we denote every
number not divisible by $q$. Here $r$ is a primitive $m$th root of
unity and $\Omega_m$ is the field of the rational functions of $r$.

Now let $\Omega = R(x)$ be a simple \gesperrt{Abelian} field of degree $m$,
so that $x$ is a root of a rational irreducible cyclic equation of
degree $m$. If we denote the roots of this equation, in a suitable
order, by
\weq{2}{x_0, x_1, x_2, \ldots, x_{m-1},}
and take $x_m = x_0$, $x_{m+1} = x_1, \ldots$, then
\[
x_1 = \Phi(x_0), \quad x_2 = \Phi(x_1), \quad \ldots, \quad x_0 = \Phi(x_{m-1}),
\]
where $\Phi$ denotes an integral function with rational coefficients,
and the group of the field $\Omega$ consists of the powers of the
substitution $(x_0, x_1)$ or of the cyclic permutations
%CONT>>


% ===== scan leaf 673, printed 652 =====
% !PRINTED 652
\weberpage{652}
%<<CONT
of the roots (2). Every cyclic function of these quantities is
a rational number. What we have to prove is that the $x$ can be
expressed rationally in terms of roots of unity. If we adjoin
to the field $\Omega$ in addition the $m$th root of unity $r$, there
arises a field $R(x, r)$, which contains the two fields $\Omega = R(x)$ and
$R(r) = \Omega_m$ as divisors.

\gesperrt{If a number} $\omega = \Phi(x, r)$ of the field $R(x, r)$ admits
all the substitutions $(x_0, x_1), (x_0, x_2), \ldots, (x_0, x_{m-1})$,
then it is contained in the field $\Omega_m$; for then
\[
m\,\omega = \Phi(x_0, r) + \Phi(x_1, r) + \cdots + \Phi(x_{m-1}, r),
\]
and in this the coefficients of the powers of $r$ are not only
cyclic, but even symmetric functions of the roots
(2), and hence rational numbers. This holds even in the
case that the equation for $x$ becomes reducible upon the adjunction
of $r$.

We apply this to the \gesperrt{resolvents}
\weq{3}{(r, x_0) = x_0 + r x_1 + r^2 x_2 + \cdots + r^{m-1} x_{m-1},}
and more generally, if $s$ denotes an arbitrary integral exponent:
\weq{4}{(r^s, x_0) = x_0 + r^s x_1 + r^{2s} x_2 + \cdots + r^{(m-1)s} x_{m-1}.}

By means of these resolvents $x_0$ itself can in turn be expressed
rationally, namely:
\weq{5}{m x_0 = \sum\limits_{0,\,m-1}^{s} (r^s, x_0),}
and if therefore we can show that all these resolvents
$(r^s, x_0)$ are cyclotomic numbers, we have attained our
goal. If we set
\weq{6}{(r^s, x_{\varkappa}) = x_{\varkappa} + r^s x_{\varkappa+1} + r^{2s} x_{\varkappa+2} + \cdots + r^{(m-1)s} x_{\varkappa+m-1}}
then $(r^s, x_{\varkappa})$ arises from $(r^s, x_0)$ by the substitution
$(x_0, x_{\varkappa})$, and there subsists the fundamental relation
\weq{7}{r^{\varkappa s} (r^s, x_{\varkappa}) = (r^s, x_0).}

All these resolvents are numbers of the field $R(x, r)$;
we consider by preference the two following combinations:
\weq{8}{(r^s, x_0)^m = \omega_s,}
\weq{9}{(r^s, x_0)^{-1} (r, x_0)^s = \alpha.}


% ===== scan leaf 674, printed 653 =====
% !PRINTED 653
\weberpage{653}

These two numbers admit, as follows immediately from (7), the substitutions $(x_0, x_\varkappa)$, and are therefore numbers of the field $\Omega_m$.

If $n$ is not divisible by $q$, then on account of the irreducibility of the cyclotomic equation none of the resolvents $(r^n, x_0)$ can vanish unless they all vanish; for the equation
\[
(r, x_0)^m = 0
\]
would be an equation for $r$ with rational coefficients which admits all the substitutions $(r, r^n)$.

In the same way none of the resolvents $(r^s, x_0)$ can vanish unless the whole family vanishes in which $s$ has a definite greatest common divisor with $m$.

If all the $(r, x_0)$ vanish, formula (9) is not applicable. But then we can consider directly the resolvents $(r^q, x_0)$, which belong to a field of degree $m q^{-1}$ arising from
\[
y_0 = x_0 + x_q + x_{2q} + \cdots + x_{m-q}
\]
and if these quantities also vanish, we pass on to the resolvents $(r^{q^2}, x_0)$, and so on.

But if $(r, x_0)$ does not vanish, then by formula (9) all the $(r^s, x_0)$ can be represented rationally in terms of $(r, x_0)$ and of cyclotomic numbers.

Accordingly it is sufficient for our purpose if we can prove that the \gesperrt{non-vanishing resolvents} $(r, x_0)$ \gesperrt{are cyclotomic numbers}.

If we let $n$ run through a complete system of incongruent numbers not divisible by $q$, then
\[
\sum n \equiv 0 \ (\mathrm{mod}\ m),
\]
for the numbers $n$ split into pairs of numbers completing one another to $m$. From this it follows by (7) that the product
\weq{10}{\prod^n (r^n, x_0) = a}
is a number in $\Omega_m$. But since this number, moreover, does not change when any one of the substitutions $(r, r^n)$ is carried out in it, $a$ is a \gesperrt{rational number}.


% ===== scan leaf 675, printed 654 =====
% !PRINTED 654
\weberpage{654}
\pdfbookmark[2]{§ 181. Preparation for the Proof}{bm207}

\begin{center}
\S.\,181.\\
\gesperrt{Preparation for the Proof.}
\end{center}

After the expositions of the preceding section, all that is still required for the proof of the great Theorem I is the demonstration that the number contained in $\Omega_m$
\[
\omega = (r, x_0)^m
\]
is the $m$th power of a cyclotomic number.

1. \gesperrt{The essential property of the numbers $\omega_n$ conjugate to $\omega$, on which the proof rests, is that}
\weq{1}{\omega_n^{-1} \omega_1^n = \alpha^m}
\gesperrt{is the $m$th power of a number in $\Omega_m$.}

Here $n$ may be any number not divisible by $q$. If we form the norms of the right and left sides of (1) in the field $\Omega_m$, then, denoting the norm of the number $\alpha$, which is a rational number, by $a$, we obtain, since every $\omega_n$ has the same norm as $\omega$,
\[
N(\omega)^{n-1} = a^m.
\]
Now if $m$ is odd, we may here take $n = 2$, and obtain as a \gesperrt{conclusion from 1.}:

2. \gesperrt{The norm of $\omega$ is the $m$th power of a rational number}
\weq{2}{N(\omega) = a^m.}

% [sic] The German source prints the anomalous punctuation ``Potenz von 2.,''; the translation normalizes it without concealing the source reading.
If $m$ is a power of 2, this conclusion is no longer admissible. It then follows from 1. only that $N(\omega)^2$ is the $m$th power of a rational number. Nevertheless the numbers $\omega$, as is to be seen from (10) of the preceding section, must in this case too still satisfy condition 2., which then, however, is no longer of itself already contained in 1., but is added as a new property.

\gesperrt{What we have to prove is that from 1. and 2. it follows that $\omega$ itself is also the $m$th power of a cyclotomic number, though it may be in a higher field.}

For this, however, it is necessary to investigate the decomposition of the numbers $\omega$ into their prime factors in the field $\Omega_m$.

The prime number $q$ is, as demonstrated in \S.\,169, with the $\varphi(m)$th power of a prime number $\sigma = 1 - r$ existing in $\Omega_m$
%CONT>>


% ===== scan leaf 676, printed 655 =====
% !PRINTED 655
\weberpage{655}
%<<CONT
associated. If therefore $\sigma^k$ is the power of $\sigma$ contained in $\omega$, we set
\[
\omega = \sigma^k \omega',
\]
where $\omega'$ is coprime to the prime $q$. The formation of the norm yields, by 2.,
\[
a^m = q^k b,
\]
where $b$ is a rational number which, in its simplest shape, contains the prime $q$ neither in the numerator nor in the denominator. From this it follows, however, that $k$ must be divisible by $m$, and from this the first result:

3. \gesperrt{The exponent of the prime $\sigma$ contained in $\omega$ is always divisible by $m$.}

Now let $p$ further be a prime different from $q$, and $\mathfrak{p}_\xi$ a prime factor of $p$ in the field $\Omega_m$. We let $\xi$ run through the sequence of numbers $\xi_1, \xi_2, \ldots, \xi_e$ determined more closely in the various cases in \S.\,172, so that
\weq{3}{p = \mathfrak{p}_{\xi_1}\mathfrak{p}_{\xi_2}\ldots\mathfrak{p}_{\xi_e}}
results, and assume that $\mathfrak{p}_\xi^{k_\xi}$ is the power of $\mathfrak{p}_\xi$ contained in $\omega$. The power of $\mathfrak{p}_{n\xi}$ contained in $\omega_n$ is then $\mathfrak{p}_{n\xi}^{k_\xi}$, and if we determine $n'$ from the congruence
\weq{4}{n n' \equiv 1 \ (\mathrm{mod}\ m)}
then $\mathfrak{p}_\xi^{k_{n'\xi}}$ is the power of $\mathfrak{p}_\xi$ contained in $\omega_n$. Accordingly, in
\weq{5}{\omega_n^{-1} \omega_1^n \text{ the power } \mathfrak{p}_\xi^{n k_\xi - k_{n'\xi}}}
is contained, and on account of 1. the exponent of this power must be divisible by $m$.

We thus obtain, for every $n$ not divisible by $q$, the congruence
\weq{6}{n k_\xi \equiv k_{n'\xi} \ (\mathrm{mod}\ m),}
in which we can now also replace $\xi$ by $n\xi$, whereby
\weq{7}{n k_{n\xi} \equiv k_\xi \ (\mathrm{mod}\ m)}
results. If we here take $\xi = 1$ and then put $\xi, \xi'$ for $n, n'$ and $k$ for $k_1$, it follows from (7) that
\weq{8}{k_\xi \equiv \xi' k \ (\mathrm{mod}\ m),}
where $k$ now has the same value for all $\xi$, and between $\xi, \xi'$ the congruence
\[
\xi \xi' \equiv 1 \ (\mathrm{mod}\ m)
\]
subsists.


% ===== scan leaf 677, printed 656 =====
% !PRINTED 656
\weberpage{656}

If we assume $n = p$, then $\mathfrak{p}_{p\xi} = \mathfrak{p}_\xi$ (\S.\,172), and hence also $k_{p\xi} = k_\xi$, and from (7) and (8) it follows that:
\weq{9}{k(p-1) \equiv 0 \ (\mathrm{mod}\ m).}

From this there now follow important consequences: If, in the first place, $p - 1$ is not divisible by $q$, it follows that $k$ must be divisible by $m$, and we therefore have the theorem:

4. \gesperrt{The prime factors of a prime $p$ which is not congruent to 1 with respect to the modulus $q$ are contained in $\omega$ only in such powers whose exponents are divisible by $m$.}

If, in the case where $m$ is a power of 2, $p - 1$ is divisible by 2 but not by 4, that is, if $p \equiv -1 \ (\mathrm{mod}\ 4)$, then it follows from (9) only that $k$ is divisible by $\tfrac{1}{2} m$, but not that $k$ is divisible by $m$. The exponent of the power of $\mathfrak{p}_\xi$ contained in $\omega$ is, apart from multiples of $m$, equal to $k$, and if $k$ is divisible by $m$, then everything is as in case 4. But if $k$ is divisible only by $\tfrac{1}{2} m$ and not by $m$, so that
\[
k \equiv \tfrac{1}{2} m \ (\mathrm{mod}\ m),
\]
then $k - \tfrac{1}{2} m$ is divisible by $m$, and, having regard to the decomposition (3), we can formulate the theorem thus:

5. \gesperrt{If $m$ is a power of 2 and $p$ a prime congruent to $-1$ with respect to the modulus 4, then the exponent $h$ can be so determined that all the prime factors of $p$ are contained in}
\[
\omega \sqrt{p}^{-hm}
\]
\gesperrt{only to such powers whose exponent is divisible by $m$.}

Finally, let $p - 1$ be divisible by $q$, or, if $m$ is even, by 4. If $m_1$ is the greatest common divisor of $m$ and $p - 1$, and if
\[
m = m_1 m_2,
\]
then by (9) $k$ must be divisible by $m_2$, and if we set
\[
k = h m_2
\]
then by (8) the following product is contained in $\omega$:
\weq{10}{\left(\prod^\xi \mathfrak{p}_\xi^{\xi'}\right)^{h m_2}.}
Moreover, the prime factors $\mathfrak{p}_\xi$ occur in $\omega$ only in such powers whose exponents are divisible by $m$.


% ===== scan leaf 678, printed 657 =====
% !PRINTED 657
\weberpage{657}

In this case $\xi$ now runs, by \S.\,172, III, through a complete
system of residues not divisible by $q$ with respect to the modulus
$m_1$, and in (10) we may understand by $\xi'$ the least positive
solution of the congruence
\[
\xi\xi' \equiv 1\ (\mathrm{mod}\ m_1)
\]
because, if we alter the exponents $\xi'$ in (10) with respect to the
modulus $m_1$, only $m$th powers of the prime factors $\mathfrak{p}$
are added.

We can then apply \gesperrt{Kummer's} theorem [\S.\,174, (16)] to
the product (10), replacing $m$ there by $m_1$, and, if we understand
by $r_1 = r^{m_2}$ an $m_1$th root of unity, and by the $\eta$ % [source sic: German prints a full stop after ``weil''; the translation supplies the grammatically required comma]
certain periods of the $p$th roots of unity, we obtain:
\[
\prod^{\xi} \mathfrak{p}_\xi^{\xi'} = (r_1, \eta)^{m_1},
\]
and consequently
\weq{11}{(\prod^{\xi} \mathfrak{p}_\xi^{\xi'})^{hm_2} = (r_1, \eta)^{hm}.}
Accordingly we have the theorem:

6. \gesperrt{If $p$ is a prime and $p-1$ is divisible by $q$, or,
when $q=2$, by 4, then the positive exponent $h$ can be so
determined that the number contained in $\Omega_m$}
\[
\omega (r_1, \eta)^{-hm}
\]
\gesperrt{contains the prime factors of $p$ only to such powers
whose exponent is divisible by $m$.}

It is now also no longer necessary to distinguish Theorem 5. from
Theorem 6.; for when $m_1 = 2$, 6. passes over into 5., because we
then have $r_1 = -1$, and by Bd. I, \S.\,171
\[
(r_1, \eta)^m = (-1, \eta)^m = \sqrt{p}^{\,m}
\]
holds $^1$).

If we combine the result of Theorems 3. to 6. into one formula, we
obtain, if $\varepsilon$ denotes a unit functional, $\varphi$ any
functional of the field $\Omega_m$, and $p$ runs through a sequence
of primes that are congruent to 1 with respect to the modulus $q$,
among which the same prime $p$ may also occur several times,
\weq{12}{\omega = (r, x_0)^m = \varepsilon \varphi^m \prod^{p} (r^{m_2}, \eta)^m.}

\noindent
$^1$) In the author's memoir in Bd. 8 of the Acta mathematica it was
neglected to give special emphasis to the case of Theorem 5.


% ===== scan leaf 679, printed 658 =====
% !PRINTED 658
\weberpage{658}

The resolvents of the division of the circle $(r^{m_2}, \eta)$
occurring in these formulae are cyclotomic numbers in a higher
field. But their $m$th powers are numbers contained in the field
$\Omega_m$, which under the substitution $(r, r^n)$ pass into
\[
(r^{n m_2}, \eta)^m.
\]
Likewise the combinations
\weq{13}{(r^{n m_2}, \eta)^{-1} (r^{m_2}, \eta)^n}
are contained in the field $\Omega_m$ [\S.\,174, (9), (11)]. These
resolvents are special cases of the general resolvent
$(r, x_0)$. From (12) we obtain, if we denote by $\varphi_n,
\varepsilon_n$ the conjugate functionals of $\varphi$ and
$\varepsilon$,
\weq{14}{\omega_n = \varepsilon_n \varphi_n^m \prod^{p} (r^{m_2 n}, \eta)^m.}
This formula acquaints us with the most important properties of the
functionals $\varphi_n$ and $\varepsilon_n$; first of all:

7. \gesperrt{The functionals $\varphi_n^m$ are associated with
numbers of the field $\Omega_m$, and hence belong to the
principal class (\S.\,152).}

By 1., $\omega_n^{-1} \omega_1^n = \alpha^m$ is the % [sic] German source omits "in" before $\Omega_m$
$m$th power of a number $\Omega_m$. Since the
combinations (13) are also numbers in $\Omega_m$, we can set
\[
\prod^{p} (r^{m_2 n}, \eta)^{-1} (r^{m_2}, \eta)^n = \frac{\alpha}{\beta}
\]
and obtain from (14)
\[
\varepsilon_n^{-1} \varepsilon_1^n (\varphi_n^{-1} \varphi_1^n)^m = \beta^m,
\]
in which $\beta$ is a number in $\Omega_m$.

From this it follows that the product
\[
\beta \varphi_n \varphi_1^{-n} = \varepsilon
\]
is a unit functional of the field $\Omega_m$, and from this there
result for the functionals $\varphi_n$ and the
unit functionals $\varepsilon_n$ the following two theorems:

8. \gesperrt{The conjugate functionals $\varphi_n$ have the
property that all the products}
\[
\varphi_n \varphi_1^{-n}
\]
\gesperrt{belong to the principal class.}

9. \gesperrt{The unit functionals $\varepsilon_n$ have the
property that the products}
\[
\varepsilon_n \varepsilon_1^{-n} = \varepsilon^m
\]
\gesperrt{are $m$th powers of unit functionals.}


% ===== scan leaf 680, printed 659 =====
% !PRINTED 659
\weberpage{659}
\pdfbookmark[2]{§ 182. Proof of the first lemma for an odd m}{bm208}

For the proof of the Principal Theorem I. everything now depends on
the proof of the two theorems that are to be inferred from the
properties of the functionals $\varphi$ and $\varepsilon$:

A. \gesperrt{The functionals $\varphi_n$ belong to the principal class.}

If we may assume this lemma, then $\varphi_n$ is associated with a
number $\alpha_n$, and we can represent the formula (14), with a
somewhat altered meaning of $\varepsilon_n$, in this
way:
\weq{15}{\omega_n = \varepsilon_n \alpha_n^m \prod^{p} (r^{m_2 n}, \eta)^m.}
In this formula, however, $\varepsilon_n$ is a \gesperrt{numerical unit}
of the field $\Omega_m$ which satisfies Theorem 9.

If, then, we can prove from Theorem 9. the second lemma as
well:

B. \gesperrt{A numerical unit $\varepsilon_n$ of the field
$\Omega_m$ which satisfies Theorem 9. is the $m$th
power of a unit in $\Omega_m$, multiplied by a power
of $r$;}

then we can extract the $m$th root in (15), and obtain, if by
$\varrho$ a root of unity of possibly higher degree than $m$,
and by $\alpha$ a number of the field $\Omega_m$, is denoted:
\weq{16}{(r, x_0) = \varrho \alpha \prod^{p} (r^{m_2}, \eta),}
where now nothing but cyclotomic numbers stand on the right, and
Theorem I. is therefore proved.

\begin{center}
\S.\,182.

\gesperrt{Proof of the first lemma for an odd $m$.}
\end{center}

In the preceding section we made the proof of Theorem I. depend
on two lemmas, whose proof is now still to be sought.

We formulate the first of these lemmas thus:

a) \gesperrt{If $\varphi_n$ is a system of conjugate
functionals of the field $\Omega_m$, different from zero, of
which it is known that}
\[
\varphi_n^m \text{ \gesperrt{and} } \varphi_n^{-1} \varphi_1^n
\]
\gesperrt{belong to the principal class, then the
functionals $\varphi$ themselves belong to the principal class.}


% ===== scan leaf 681, printed 660 =====
% !PRINTED 660
\weberpage{660}

The theorem will be proved if it can be shown of any power
$\varphi^k$ of $\varphi$ whose exponent is not divisible by $q$
that it belongs to the principal class. For if $\alpha, \beta$
are numbers in $\Omega_m$ and
\[
\varphi^m = \varepsilon' \alpha, \quad \varphi^k = \varepsilon'' \beta,
\]
then we can determine the integral rational numbers $x, y$ in such
a way that
\[
mx + ky = 1
\]
holds, and then
\[
\varphi = \varphi^{mx} \varphi^{ky} = \varepsilon''' \alpha^x \beta^y,
\]
hence, if $\varepsilon', \varepsilon'', \varepsilon'''$ are units,
$\varphi$ itself is also associated with a number.

In the present state of the matter, the proof for an even $m$ must
be conducted along an entirely different path than for an odd one,
and so we occupy ourselves first with the easier case of odd
$m$ $^1$).

The assumptions made in a) about $\varphi_n$ can be expressed, if
$\alpha, \beta$ again denote any numbers of the field $\Omega_m$
and $\varepsilon$ unit functionals, by the two formulae:
\weq{1}{\varphi_n^m = \varepsilon \beta,}
\weq{2}{\varphi_n = \varepsilon \alpha \varphi_1^n.}
We shall now solve the problem step by step.

If first, as before, $\sigma$ is the prime factor contained in
$q$, which is a number existing in $\Omega_m$, and $\sigma^k$ the
power of $\sigma$ contained in $\varphi$, then we set
\weq{3}{\varphi = \sigma^k \psi,}
in which $\psi$ denotes a functional coprime to $q$ that satisfies
the same conditions (1), (2) as $\varphi$; and if one of the two
belongs to the principal class, the same holds of the other. The
theorem a) therefore need only be proved under the assumption that
$\varphi$ is coprime to $q$.

Now let $p$ be a prime different from $q$ which, decomposed into
prime factors, has the expression:
\weq{4}{p = \prod^{\xi} \mathfrak{p}_\xi,}
%CONT>>

\noindent
$^1$) Cf., moreover, the new proof by \gesperrt{Hilbert}, in which
the difference between even and odd $m$ recedes far more
(S. 648).


% ===== scan leaf 682, printed 661 =====
% !PRINTED 661
\weberpage{661}
%<<CONT
in which $\xi$ runs through the sequence of numbers determined more closely in \S.\,172.
Let $\mathfrak{p}_\xi^{k_\xi}$ be the power of $\mathfrak{p}_\xi$ contained in $\varphi$,
so that, if we set
\[
\varphi = \chi \prod^{\xi} \mathfrak{p}_\xi^{k_\xi}
\]
then $\chi$ is coprime to $p$ and is related to no primes other than
those to which $\varphi$ itself is related (see the conclusion of
\S.\,173). From this it follows that
\weq{5}{\varphi_n = \chi_n \prod^{\xi} \mathfrak{p}_{n\xi}^{k_\xi},}
and here we again let $n$ run through the sequence of numbers $\xi$.
If we further set
\[
\psi = \prod^{\xi} \chi_\xi,
\]
then $\psi$ too is coprime to $p$ and is related to no primes other
than those to which $\varphi$ is related, and from (5), using (4), we obtain:
\weq{6}{\prod^{\xi} \varphi_\xi = \psi\, p^{\Sigma k_\xi},}
and from this it follows that the functional $\psi$ satisfies the same two
conditions (1), (2) as $\varphi$.

It appears further that, if the $\varphi_\xi$ lie in the principal class,
then $\psi$ too is contained in the principal class.

If, however, $p-1$ is not divisible by $q$, then we can also
infer the converse.

For by (2) $\varphi_\xi$ is equivalent to $\varphi^\xi$, and consequently:
\[
\begin{gathered}
\prod^{\xi} \varphi_\xi \ \text{equivalent}\ \varphi^{\Sigma\xi}, \\
\text{equivalent}\ \psi,
\end{gathered}
\]
and if $\psi$ is associated with a number, the same holds
of $\varphi^{\Sigma\xi}$. But by the theorem of \S.\,172, if $p-1$ is not
divisible by $q$, then $\Sigma \xi$ is also not divisible by $q$, and consequently
 $\varphi$ itself is likewise contained in the principal class.

This consideration we can now repeat, if $\psi$ is related to
still further primes that are not congruent to 1 with respect to the
 modulus $q$, and we thus arrive at the
theorem:

\gesperrt{Theorem a) is proved in general if it can be
proved under the assumption
that $\varphi$ is related only to such primes as
are congruent to 1 with respect to the modulus $q$.}


% ===== scan leaf 683, printed 662 =====
% !PRINTED 662
\weberpage{662}
If we make this assumption about the functionals $\varphi_n$, then
we can apply \gesperrt{Kummer's} theorem in the form given in
\S.\,174, 3.

In what follows we denote by $\varepsilon$ unit functionals, of
which no closer knowledge is required, and set
\weq{7}{\Phi_n = \prod^{t} \varphi_{nt'}^{t} = \varepsilon\, \vartheta_n^m,}
where $t$ runs through the residues of $m$ that are not divisible by
$q$ and are less than $m$, and $t'$ is determined by
$t\,t' \equiv 1\ (\mathrm{mod}\ m)$. Then $\vartheta_n$ is a
cyclotomic number in a higher field, whose $m$th power is contained
in $\Omega_m$ and which satisfies the second condition, that
\weq{8}{\vartheta_n^{-1} \vartheta_1^n = \alpha}
is a number of the field $\Omega_m$.

We now introduce a symbol that is used in many places, by
understanding by $E(x)$ the greatest integer contained in the real
number $x$, and by $(x)$ the residue, which is always a proper
fraction and which, should $x$ itself be an integer, is to be set
equal to zero. Then
\weq{9}{x = E(x) + (x).}

If $x$ is an integer, then according to this mode of notation
\[
m \left(\frac{x}{m}\right)
\]
is likewise an integer, and indeed the least positive residue of the
division of $x$ by $m$.

In formula (7) the index of the functionals $\varphi$ is determined
only with respect to the modulus $m$. The exponent, however, is
reduced to its least residue with respect to the modulus $m$. If,
then, we replace the index $t'$ in (7) by $n' t'$, then $t$ must be
replaced by the residue of the division of $n t$ by $m$, that is
[by (9)], by
\[
m \left(\frac{n t}{m}\right) = n t - m E\left(\frac{n t}{m}\right)
\]
and we can set
\[
\Phi_n = \prod^{t} \varphi_{t'}^{\,n t - m E\left(\frac{n t}{m}\right)} = \varepsilon\, \vartheta_n^m,
\]
on the other hand
\[
\Phi_1 = \prod^{t} \varphi_{t'}^{t} = \varepsilon\, \vartheta_1^m,
\]
%CONT>>


% ===== scan leaf 684, printed 663 =====
% !PRINTED 663
\weberpage{663}
%<<CONT
hence by (8)
\weq{10}{\Phi_1^n \Phi_n^{-1} = \prod^{t} \varphi_{t'}^{\,m E\left(\frac{n t}{m}\right)} = \varepsilon\, \alpha^m.}

According to this, the quotient
\[
\left(\frac{\prod\limits^{t} \varphi_{t'}^{\,E\left(\frac{n t}{m}\right)}}{\alpha}\right)^m
\]
is a unit functional, and since it is at the same time the $m$th power of a
functional in $\Omega_m$, the base of this power is also a unit, that is,
\[
\prod^{t} \varphi_{t'}^{\,E\left(\frac{n t}{m}\right)} = \varepsilon \alpha,
\]
and this product belongs to the principal class. But now, by our
assumption,
\[
\varphi_{t'}\ \text{equivalent to}\ \varphi^{t'},
\]
and therefore
\[
\prod^{t} \varphi_{t'}^{\,E\left(\frac{n t}{m}\right)}\ \text{equivalent to}\ \varphi^{\sum^{t} t' E\left(\frac{n t}{m}\right)},
\]
and this holds for every arbitrary $n$ not divisible by $q$.

If, then, we can show that $n$ can be so determined that the sum
\[
S_n = \sum^{t} t' E\left(\frac{n t}{m}\right)
\]
is likewise not divisible by $q$, we shall have proved Theorem a).

We shall first undertake a transformation of the sum $S_n$ by which
the question is reduced to the far simpler case in which $m$ is a
prime.

If $m$ is a higher than the first power of $q$, we set
\[
m = q m', \quad t = t_1 + q t_2, \quad
\begin{aligned}
t_1 &= 1, 2, \dots, q-1 \\
t_2 &= 0, 1, \dots, m'-1.
\end{aligned}
\]

Here $m'$ is again a power of $q$; the summation letter $t_1$ runs
through the sequence of numbers $1, 2, \dots, q-1$, and $t_2$ the
sequence of numbers $0, 1, 2, \dots, m'-1$. We determine $t_1'$ from the
congruence
\[
t_1 t_1' \equiv 1\ (\mathrm{mod}\ q)
\]
and obtain
\[
t' \equiv t_1'\ (\mathrm{mod}\ q).
\]


% ===== scan leaf 685, printed 664 =====
% !PRINTED 664
\weberpage{664}
It is then
\weq{11}{S_n = \sum^{t} t' E\left(\frac{t n}{m}\right) \equiv \sum^{t_1} t_1' \sum^{t_2} E\left(\frac{n t_2}{m'} + \frac{n t_1}{m}\right)\ (\mathrm{mod}\ q).}

We now assume $n < q$, and hence also $n t_1 < m$, so that
$n t_1 : m$ is a proper fraction. Then
\[
\begin{aligned}
\text{a)} \quad E\left(\frac{n t_2}{m'} + \frac{n t_1}{m}\right)
  &\quad \text{is either}\quad = E\left(\frac{n t_2}{m'}\right) \\
\text{b)} \quad
  &\quad \text{or}\quad = E\left(\frac{n t_2}{m'}\right) + 1,
\end{aligned}
\]
and case b) occurs only when between
\[
\frac{n t_2}{m'}\ \text{and}\ \frac{n t_2}{m'} + \frac{n t_1}{m}
\]
there lies an integer. But this is the case only when
the difference between $n t_2 : m'$ and the next greater
integer is smaller than $n t_1 : m$, that is, when
\weq{12}{E\left(\frac{n t_2}{m'}\right) + 1 - \frac{n t_2}{m'} < \frac{n t_1}{m}.}

If in the expression on the left-hand side of (12)
\weq{13}{E\left(\frac{n t_2}{m'}\right) + 1 - \frac{n t_2}{m'}}
we let $t_2$ run through all its values, we obtain nothing but positive
rational fractions, not exceeding unity, with the
denominator $m'$, no two of which are equal to one another. For if
two of the values that arise from (13) by substitution of
$t_2', t_2''$ for $t_2$ were equal to one another, then $\dfrac{n(t_2' - t_2'')}{m'}$ would have to be an
integer, which, since $n$ is relatively prime to $m'$ and $t_2' - t_2''$ is smaller
than $m'$, cannot be. The expressions (13) therefore run through,
in some order, the numerical values
\weq{14}{\frac{1}{m'}, \frac{2}{m'}, \dots, \frac{m'-1}{m'}, 1,}
and if $a$ denotes an arbitrary positive fraction, then
$E(m' a)$ is the number of the numbers of the sequence (14) that are not greater
than $a$. By (12) the number of values of $t_2$ for
which case b) occurs is therefore equal to $E\left(\dfrac{n t_1}{q}\right)$, and there results
\[
\sum^{t_2} E\left(\frac{n t_2}{m'} + \frac{n t_1}{m}\right) = \sum^{t_2} E\left(\frac{n t_2}{m}\right) + E\left(\frac{n t_1}{q}\right). % [source sic: the first denominator on the right is printed m, not the mathematically required m']
\]


% ===== scan leaf 686, printed 665 =====
% !PRINTED 665
\weberpage{665}

%% [running head, not body text -- commented out 2026-08-07] Page furniture (running
%% head, folio, signature) is excluded from this edition by standing policy.
%% \S.\,182. \hfill \gesperrt{Proof of the Lemmas.} \hfill 665

In order to form $S_n$, we multiply this equation by $t_1'$ and sum.
This gives

\[
S_n \equiv \sum^{t_1'} t_1' \sum^{t_2} E\left(\frac{n t_2}{m}\right)
+ \sum^{t_1} t_1' E\left(\frac{n t_1}{q}\right) \pmod q.
% [source sic: the denominator m repeats the source's m-for-m' error from the preceding leaf]
\]

Now here $\sum t_1' = \sum t_1 \equiv 0 \pmod q$, because the $t_1$
yield in pairs the sum $q$, and it follows:

\weq{15}{S_n \equiv \sum^{t_1} t_1' E\left(\frac{n t_1}{q}\right) \pmod q.}

The sum that stands here on the right side is formed in just the same
way as $S_n$ itself in the case where $m$ is a prime. For this case,
however, we obtain

\[
\frac{n t_1}{q} = E\left(\frac{n t_1}{q}\right) + \left(\frac{n t_1}{q}\right)
\]
\[
\frac{(q-n) t_1}{q} = E\left(\frac{(q-n) t_1}{q}\right) + \left(\frac{(q-n) t_1}{q}\right),
\]

and from this, by addition,

\[
t_1 = E\left(\frac{n t_1}{q}\right) + E\left(\frac{(q-n) t_1}{q}\right)
+ \left(\frac{n t_1}{q}\right) + \left(\frac{(q-n) t_1}{q}\right),
\]

and here the sum of the two positive proper fractions, which must after
all be an integer, can only have the value 1. Hence it follows

\[
E\left(\frac{n t_1}{q}\right) + E\left(\frac{(q-n) t_1}{q}\right) = t_1 - 1,
\]

and from this, by multiplication by $t_1'$ and summation over all
values of $t_1$ from 1 to $q - 1$,

\[
\sum^{t_1} t_1' E\left(\frac{n t_1}{q}\right)
+ \sum^{t_1} t_1' E\left(\frac{(q-n) t_1}{q}\right) \equiv -1 \pmod q.
\]

Consequently the two sums on the left side can certainly not both be
divisible by $q$. The expression (15) then shows that of the two sums
$S_n$ and $S_{q-n}$ one is certainly not divisible by $q$.

But this was what was to be proved.


% ===== scan leaf 687, printed 666 =====
% !PRINTED 666
\weberpage{666}
\pdfbookmark[2]{§ 183. Proof of the second lemma for an odd m}{bm209}

%% [running head, not body text -- commented out 2026-08-07] Page furniture (running
%% head, folio, signature) is excluded from this edition by standing policy.
%% 666 \hfill \gesperrt{Twenty-second Chapter.} \hfill \S.\,183.

\begin{center}
\S.\,183.

\gesperrt{Proof of the second lemma for an odd $m$.}
\end{center}

The second lemma, which by \S.\,181 still remains to be proved, is the
following:

2. \gesperrt{If $\varepsilon_n$ is a system of conjugate numerical units
of the field $\Omega_m$ which satisfies the condition that}

\weq{1}{\varepsilon_n \varepsilon_1^{-n} = \mathfrak{E}^m}

\gesperrt{is the $m$th power of a unit in $\Omega_m$, then the
$\varepsilon_n$ are themselves $m$th powers of units in $\Omega_m$,
multiplied by a root of unity of the form $r^{nk}$.}

This proof too is entirely elementary in the case of an odd $m$. We
therefore consider this case here first.

By \S.\,178, 1. we can represent every unit of the field $\Omega_m$ in
the following form:

\weq{2}{\varepsilon = \pm r^k \mathfrak{E}(r), \quad \varepsilon_n = \pm r^{nk} \mathfrak{E}(r^n),}

where $\mathfrak{E}(r)$ is a real unit of the field $\Omega_m$.
Accordingly $\mathfrak{E}(r)$ satisfies the condition

\weq{3}{\mathfrak{E}(r) = \mathfrak{E}(r^{-1}),}

from which, by (2), there results:

\weq{4}{\varepsilon_1 \varepsilon_{-1} = \mathfrak{E}(r)^2;}

but from (1) one finds, on putting $n = -1$,

\[
\varepsilon_1 \varepsilon_{-1} = \mathfrak{E}^m,
\]

that is, the left side of (4) is the $m$th power of a unit in
$\Omega_m$, and therefore

\weq{5}{\mathfrak{E}(r)^2 = \mathfrak{E}^m}

is likewise the $m$th power of such a unit.

Since $m$ is now odd, the integral rational number $x$ can be so
determined that

\weq{6}{2 x + m = 1}

holds, and from this it follows that

\[
\mathfrak{E}(r) = \mathfrak{E}(r)^{2x} \mathfrak{E}(r)^m.
\]
%CONT>>


% ===== scan leaf 688, printed 667 =====
% !PRINTED 667
%% [running head, not body text -- commented out 2026-08-07] Page furniture (running
%% head, folio, signature) is excluded from this edition by standing policy.
%% \S.\,184. \hfill \gesperrt{Case of an even $m$.} \hfill 667
\weberpage{667}
\pdfbookmark[2]{§ 184. Preliminary remarks on the case of an even m}{bm210}
%<<CONT
If therefore we put, in accordance with (5),

\[
\mathfrak{E}^x \mathfrak{E}(r) = e(r)
\]

then it follows that

\weq{7}{\mathfrak{E}(r) = e(r)^m.}

But by (7) Theorem 2. is proved, and with it at the same time the
\gesperrt{whole Theorem I. for an odd $m$.}

This inference too fails for an even $m$, because the equation (6) is
then no longer solvable.

\begin{center}
\S.\,184.

\gesperrt{Preliminaries concerning the case of an even $m$.}
\end{center}

For the case in which $m$ is a power of 2 we take an entirely
different route, on which a few preliminary remarks may find their
place here first.

We have already seen in \S.\,182 that the first lemma is proved if we
can show that some power of $\varphi$ whose exponent is not divisible
by $q$ belongs to the principal class. In doing so, of the two
properties of the functional $\varphi$ only the first is used, namely
that the $m$th power of $\varphi$ is contained in the principal
class.

Now, by the general theorems in \S.\,154, we know in every field an
exponent $h$, namely the number of ideal classes, for which
$\varphi^h$ belongs to the principal class, whenever $\varphi$ is an
arbitrary functional in this field.

If, therefore, we can prove that the \gesperrt{class number $h$ in the field
$\Omega_m$, when $m$ is a power of 2, is an odd number,} then Lemma 1.
is thereby proved without anything further.

This is to be the aim of our considerations in the next chapters.

With regard to the second lemma the case is similar. Here one can
derive from the assumptions in \S.\,183, 2., quite as above, the
formula \S.\,183, (5); but from it one can conclude only that
$\mathfrak{E}(r)$ is the $\frac{1}{2} m$th power of a real unit
$e(r)$.

If, then, we assume the first lemma to be proved, the formula
\S.\,181, (14) yields:

\weq{1}{(r^n, x_0) = \varrho_n \sqrt{e(r^n)}\ \alpha_n\ \Pi(r^{m_2 n}, \eta),}
%CONT>>


% ===== scan leaf 689, printed 668 =====
% !PRINTED 668
%% [running head, not body text -- commented out 2026-08-07] Page furniture (running
%% head, folio, signature) is excluded from this edition by standing policy.
%% 668 \hfill \gesperrt{Twenty-second Chapter.} \hfill \S.\,184.
\weberpage{668}
%<<CONT
where $\varrho_n$ is a root of unity of degree $m^2$, and $\alpha_n$ is a
number in $\Omega_m$, and it would still remain to be proved that
$e(r^n)$ is the \gesperrt{square of a unit} in $\Omega_m$.

Since $e(r)$ is a real unit, we have

\weq{2}{e(r^n) = e(r^{-n}),}

and we shall further stipulate that the sign of the root is so
determined (which, upon a suitable assumption about $\varrho_n$, we
may always take to be the case) that

\weq{3}{\sqrt{e(r^n)} = \sqrt{e(r^{-n})}.}

By hypothesis the product $(r^n, x_0)\,(r^{-n}, x_0)$ is a number of
the field $\Omega_m$ which is not altered by the substitution
$(r, r^{-1})$, i.e. a real number. Further, by \S.\,174, (5),

\[
(r^{m_2 n}, \eta)\,(r^{-m_2 n}, \eta) = \pm p,
\]

and hence likewise real. In the same way $\alpha_n \alpha_{-n}$ is
real, and from this it follows by (1) that

\[
\varrho_n \varrho_{-n} = \frac{(r^n, x_0)\,(r^{-n}, x_0)}
{e(r^n)\ \alpha_n \alpha_{-n}\ \Pi(r^{m_2 n}, \eta)\,(r^{-m_2 n}, \eta)}
\]

is a real number, which, since it is a root of unity, can only be
$= \pm 1$.

Now in the equation

\[
(r, x_0)\,(r^{-1}, x_0) = \pm\ e(r)\ \alpha_1 \alpha_{-1}\ \Pi(\pm p)
\]

which results from this for $n = 1$, we can carry out every
substitution $(r, r^n)$, and from this it appears that the sign of
$\varrho_n \varrho_{-n}$ is independent of $n$, so that

\weq{4}{\varrho_n \varrho_{-n} = \varrho_1 \varrho_{-1} = \pm 1}

holds.

Now we derive further from (1):

\weq{5}{\varrho_n \varrho_1^{-1} \sqrt{e(r^n)}\ \sqrt{e(r)}^{-n}
= \frac{(r^n, x_0)\,(r, x_0)^{-n}}
{\alpha_n \alpha_1^{-1}\ \Pi(r^{m_2 n}, \eta)\,(r^{m_2}, \eta)^{-n}},}
% [source sic: (5) prints both \varrho_1^{-1} and \alpha_1^{-1}; formation from (1) and the displayed -n powers requires \varrho_1^{-n} and \alpha_1^{-n}]

and the right side is a number of the field $\Omega_m$. The left
side shows, however, that it is a unit, and accordingly, if
$\mathfrak{E}(r)$ denotes a real unit, we may set

\weq{6}{\varrho_n \varrho_1^{-1} \sqrt{e(r^n)}\ \sqrt{e(r)}^{-n} = r^k \mathfrak{E}(r).}
% [source sic: (6) repeats the inherited \varrho_1^{-1}; the derivation requires exponent -n]

If we substitute this value in formula (5) and make the substitution
$(r, r^{-1})$, there results

\weq{7}{\varrho_{-n} \varrho_{-1}^{-1} \sqrt{e(r^n)}\ \sqrt{e(r)}^{-n} = r^{-k} \mathfrak{E}(r),}
% [source sic: (7) prints \varrho_{-1}^{-1}; the substituted form of the preceding relation requires exponent -n]
%CONT>>


% ===== scan leaf 690, printed 669 =====
% !PRINTED 669
\weberpage{669}
%<<CONT
and by multiplication of (6) and (7), with regard to (4)

\weq{8}{e(r^n)\, e(r)^{-n} = \mathfrak{E}(r)^2.}

In our considerations on the class number the following theorem will yet appear:

\gesperrt{A unit $e(r)$ of the real field $H_m$, which is positive together with all its conjugates, is the square of a unit in the field $H_m$.}

But formula (8) shows that the unit $\pm e(r)$ has this property, and hence that $e(r)$ (since $-1 = i^2$ is likewise the square of a unit) is the square of a unit of the field $\Omega_m$.

With this the question is reduced to the proof of the two theorems mentioned. That proof appears as the keystone of a long chain of considerations, extremely rich in interesting relations, to which the next chapters are to be devoted.

\begin{center}
\rule{6em}{0.4pt}
\end{center}


% ===== scan leaf 691, printed 670 =====
\weberpage{670}% page number not printed on this page
\pdfbookmark[1]{Twenty-third Section. Class Number.}{bm211}
\pdfbookmark[2]{§ 185. Dirichlet's theorem on the units}{bm212}
% !PRINTED none
\weberpage{none}
% [source note: this Abschnitt opener has no printed folio; sequence position corresponds to inferred page 670.]

\begin{center}
\gesperrt{Twenty-Third Chapter.}

\gesperrt{Class Number.}

\rule{4em}{0.4pt}

\S.\,185.

\gesperrt{Dirichlet's Theorem on the Units.}
\end{center}

The investigations concerning the number of ideal classes in an algebraic field, which are now to give us the necessary completions of the proofs of the preceding chapter, make use of the methods that \gesperrt{Dirichlet} introduced into number theory\begingroup\renewcommand{\thefootnote}{\arabic{footnote})}\footnote{\gesperrt{Dirichlet's}, Recherches sur diverses applications de l'analyse infinitésimale à la théorie des nombres. \gesperrt{Crelle's} Journal, Bd. 19, 21. \gesperrt{Dirichlet's} Werke, Bd. I (1839, 1840). (Here belong, moreover, the posthumous memoirs of \gesperrt{Gauss} as well, \glqq De nexu inter multitudinem classium etc.\grqq{} \gesperrt{Gauss'} Werke, Bd. II.) The application of these principles to the cyclotomic numbers was first made by \gesperrt{Kummer} [\gesperrt{Crelle's} Journal, Bd. 35, 40 (1847, 1850); \gesperrt{Liouville's} Journal, Bd. 16 (1851)]. The general theory of the units was likewise founded by \gesperrt{Dirichlet} [\gesperrt{Dirichlet's} Werke, Bd. I, S. 622, 633, 639 (1840, 1842, 1846)]. The general theory of the units and the determination of the class numbers was given by \gesperrt{Dedekind}: \gesperrt{Dirichlet-Dedekind}, Vorlesungen über Zahlentheorie, 4. Auflage, \S.\,183 f. \gesperrt{Minkowski}, Geometrie der Zahlen.}\endgroup.

They are no longer purely algebraic, inasmuch as transcendental functions and numbers occur in them as well, such as the natural logarithm and the number $\pi$. Use is also made of passages to the limit, as in the integral calculus, from which we have already drawn a beautiful result in \S.\,156.

The investigations we have first to deal with are entirely general; that is, they hold for every algebraic number field, and it offers no advantage in simplicity to restrict them to special fields, such as the cyclotomic fields.


% ===== scan leaf 692, printed 671 =====
% !PRINTED 671
\weberpage{671}
Let $\Omega$ then be a field of degree $n$, and

\weq{1}{\Omega_1,\ \Omega_2,\ \ldots,\ \Omega_n}

the conjugate fields. The irreducible equation of degree $n$ whose roots determine these conjugate fields for us will in general have both real and imaginary roots; the latter, however, always occur in pairs, so that to every imaginary field $\Omega^{(i)}$ there belongs another whose numbers are conjugate imaginary to those of $\Omega^{(i)}$, and which therefore arises from $\Omega^{(i)}$ by the substitution $(i, -i)$.

Such imaginary pairs we combine into a unit, and we denote \gesperrt{the number of real fields and of pairs of imaginary fields of the sequence (1)} by $\nu$. If all the conjugate fields are real, then $n = \nu$; if they are all imaginary, then $n = 2\nu$.

Apart from the case of the rational field, $\nu$ is $= 1$ only in the one further case of an imaginary quadratic field. Otherwise $\nu$ is always greater than 1.

The number of imaginary pairs is $n - \nu$, and the number of real fields $2\nu - n$.

In the determinant by which the square root of the fundamental number, $\sqrt{\Delta}$, was defined in \S.\,145, the substitution $(i, -i)$ interchanges each two conjugate imaginary rows, and $\sqrt{\Delta}$ therefore changes its sign whenever this number is odd, but not otherwise. In the first case $\sqrt{\Delta}$ is imaginary, in the second real. From this it follows that the \gesperrt{fundamental number $\Delta$ of the field $\Omega$ is positive or negative according as $n - \nu$ is even or odd.}

If of two conjugate imaginary fields we retain only the one, there results the sequence of conjugate fields

\[ \Omega_1,\ \Omega_2,\ \ldots,\ \Omega_\nu, \]

to which we assign a system of symbols

\[ \delta_1,\ \delta_2,\ \ldots,\ \delta_\nu \]

with the provision that $\delta_s = 1$ is to hold when $\Omega_s$ is real, and $= 2$ when $\Omega_s$ is imaginary, so that $\Sigma \delta_s = n$.

Let $\eta$ denote a number of the field $\Omega$ different from zero, and

\[ \eta_1,\ \eta_2,\ \ldots,\ \eta_n \]

the conjugate numbers. To this system of numbers we assign another system of numbers

\[ \lambda_1,\ \lambda_2,\ \ldots,\ \lambda_n, \]
%CONT>>


% ===== scan leaf 693, printed 672 =====
% !PRINTED 672
\weberpage{672}
%<<CONT
which we define by the equation

\weq{2}{\lambda_s = \delta_s \log|\eta_s|,}

where $|\eta_s|$ denotes the absolute value of $\eta_s$, and $\log|\eta_s|$ the real natural logarithm of the positive number $|\eta_s|$. If $\Omega_s$ is real and the sign $\pm$ is so chosen that $\pm \eta_s$ is positive, then

\[ \lambda_s = \log(\pm \eta_s). \]

But if $\eta_s, \eta_s'$ form an imaginary pair, then

\[ \lambda_s = \log \eta_s \eta_s', \]

and to $\eta_s$ and $\eta_s'$ there belongs the same $\lambda_s$, so that the number of distinct $\lambda_s$ is equal to $\nu$. From this determination there follows further

\weq{3}{\lambda_1 + \lambda_2 + \cdots + \lambda_\nu = \log N_a(\eta),}

where, as before, $N_a$ denotes the absolute norm.

The numbers $\lambda_1, \lambda_2, \ldots, \lambda_\nu$ we shall call the \gesperrt{conjugate logarithms} of the number $\eta$.

If $\eta$ is an \gesperrt{integer} of the field $\Omega$, then the absolute norm is a \gesperrt{natural number} which is $= 1$ only when $\eta$ is a unit, and from this it follows by (3):

1. \gesperrt{The sum of the conjugate logarithms of an integer $\eta$ is positive, and equal to zero only when $\eta$ is a unit.}

If $\omega_1, \omega_2, \ldots, \omega_n$ denotes a basis of the integers in $\Omega$ (a minimal basis of $\Omega$, \S.\,145), and $\omega_{1,s}, \omega_{2,s}, \ldots, \omega_{n,s}$ for $s = 1, 2, \ldots, n$ the conjugate numbers, then all integers of the field $\Omega_s$ can be represented in the form:

\weq{4}{\eta_s = x_1 \omega_{1,s} + x_2 \omega_{2,s} + \cdots + x_n \omega_{n,s},}

with integral rational coefficients $x_1, x_2, \ldots, x_n$. The determinant of this system

\[ \Sigma \pm \omega_{1,1} \omega_{2,2} \ldots \omega_{n,n} = \sqrt{\Delta} \]

is the square root of the discriminant $\Delta$ of the field, and therefore always different from zero. Accordingly the system (4) can be solved for the unknowns $x_1, x_2, \ldots, x_n$, and if we now stipulate that the absolute values of the numbers $\eta_s$ are not to exceed a certain limit, then from these solutions there result limits for the integral rational numbers $x_i$.


% ===== scan leaf 694, printed 673 =====
% !PRINTED 673
\weberpage{673}
This simple remark yields the following important theorem:

2. \gesperrt{In an algebraic field $\Omega$ there is only a finite number of integers $\eta$ of such a nature that the absolute values of the conjugate numbers $\eta_s$ lie below a given limit.}

We now make use of the general theorem \S.\,155, (25), according to which, if $f(x_1,x_2,\dots,x_n)$ is any \gesperrt{positive quadratic form} in $n$ variables with the determinant $D$, the variables $x_i$ can be determined as integral rational numbers, not all vanishing, in such a way that
\[
f(x_1,x_2,\dots,x_n) < n\sqrt[n]{0{,}404\,D},
\]
i.\,e.\ so that it becomes smaller than a number depending on the determinant $D$ alone.

We apply this theorem, as in \S.\,156, to the form:
\[
f(x_1,x_2,\dots,x_n) = \frac{|\eta_1|^2}{c_1^2} + \frac{|\eta_2|^2}{c_2^2} + \cdots + \frac{|\eta_n|^2}{c_n^2},
\]
in which the $\eta_s$ are the linear forms (4), and, in the case of a real $\eta_s$,
\[
|\eta_s|^2 = (\eta_s)^2,
\]
while in the case of an imaginary pair $\eta_s, \eta_s'$
\[
|\eta_s|^2 = |\eta_s'|^2 = \eta_s\eta_s'.
\]
In the latter case $\eta_s\eta_s'$ is the sum of two real squares, and accordingly, if the $c_s$ are assumed to be real, $f$ is a positive form. Concerning the real quantities $c_s$, which have so far been arbitrary, we further stipulate that, if $\eta_s, \eta_s'$ form a conjugate imaginary pair,
\weq{5}{c_s = c_s'}
shall hold; besides this we shall subject the $c_s$ to the further condition:
\weq{6}{c_1 c_2 \dots c_n = 1.}
We form the determinant of the function $f$, as in \S.\,156, by first regarding it as a form in the variables $\eta_1, \eta_2, \dots, \eta_n$, and the determinant of this form is, on account of (6), equal to $\pm 1$. The determinant of $f$ with respect to the variables $x_i$ is therefore (Bd.\ I, \S.\,56), apart from the sign, equal to the
%CONT>>


% ===== scan leaf 695, printed 674 =====
% !PRINTED 674
\weberpage{674}
%<<CONT
square of the determinant of the linear forms $\eta_s$, that is, of the fundamental number of the field $\Omega$.

From this it follows that one can find a number $A$ which is independent of the $c_s$ and determined by the field $\Omega$ alone, and which is such that, whatever the $c_s$ may be, $f(x_1,x_2,\dots,x_n)$ sinks below $A$ for integer non-vanishing $x$. Since $f$ is the sum of positive summands, the same holds of each of these summands, and with that the following theorem is proved:

3. \gesperrt{If $c_s$ is any system of real positive numbers satisfying the conditions (5), (6), then there is a real number $A$ determined by the field $\Omega$ alone which is such that one can always determine an integer $\eta$ of the field $\Omega$ that, together with its conjugate numbers $\eta_s$, satisfies the conditions}
\weq{7}{|\eta_s| < c_s A.}
Incidentally it follows from this that in every algebraic field $\Omega$, with the exception of the rational one and the imaginary quadratic one, one can find integers different from zero whose absolute value sinks below every given limit.

We wish to recast the statement of theorem 3. somewhat by introducing the conjugate logarithms of $\eta$. We put $\log A = k$, so that $k$ is a real number determined by $\Omega$; further
\weq{8}{c_s = e^{\frac{\gamma_s u}{\delta_s}},}
in which $u$ may be any real quantity, while the numbers $\gamma_s$, in virtue of (5) and (6), satisfy the condition:
\weq{9}{\gamma_1 + \gamma_2 + \cdots + \gamma_\nu = \sum_{1,\nu}^{i} \gamma_i = 0.}
Since moreover
\[
|\eta_s| = e^{\frac{\lambda_s}{\delta_s}}
\]
holds, it follows from (7) and (8):
\weq{10}{\lambda_s - \gamma'_s u < k\delta_s.}
% [sic] The printed prime has no referent; the unprimed \gamma_s is required by (8), (9), and the following sum.
The sum of the $\nu$ expressions on the left side of this inequality is, in virtue of (9), equal to $\Sigma\lambda_s$, and therefore, by theorem 1., not negative, and we find
\[
0 \leqq \sum^{s} (\lambda_s - \gamma_s u) < k n.
\]


% ===== scan leaf 696, printed 675 =====
% !PRINTED 675
\weberpage{675}
From this it follows, on account of (10), that a term of this sum, $\lambda_s - \gamma_s u$, cannot be smaller than $-(n - \delta_s) k$, since otherwise the total sum would be negative; and if for simplicity we take the lower limit somewhat smaller still
\weq{11}{-n k \delta_s < \lambda_s - \gamma_s u < k \delta_s.}
It is thereby proved that for every given system of the numbers $u, \gamma_s$, provided only that condition (9) is satisfied, there exists an integer $\eta$ of the field $\Omega$ which, together with its conjugate numbers, satisfies the limiting conditions (11).

Let there further be given a system $g_s$ of real quantities, of which we require only that
\[
\gamma_1 g_1 + \gamma_2 g_2 + \cdots + \gamma_\nu g_\nu = \alpha
\]
be different from zero. It is thereby excluded [by (9)] that all the $g_s$ are equal to one another; but if they are not, then for every given system of the $g_s$ infinitely many systems $\gamma_s$ can be determined which satisfy requirement (9).

If we now set $k \Sigma \delta_s g_s = g$, then from (11), by multiplication with $g_s$ and summation, there results
\weq{12}{-n g + \alpha u < \Sigma g_s \lambda_s < g + \alpha u.}
% [source note] Multiplying the inequalities (11) in this direction requires nonnegative g_s, not arbitrary real g_s as stated. One may first add a sufficiently large common constant to all g_s: this leaves alpha unchanged since sum gamma_s=0 and leaves the final linear form on unit logarithms unchanged since sum l_s=0. The printed proof omits this reduction.
In accordance with this limiting condition we now determine, with $g_s$ and $\alpha$ held fixed, a sequence of integers
\weq{13}{\eta, \eta', \eta'', \eta''', \dots,}
by making for $u$ a sequence of assumptions:
\[
u, u', u'', u''', \dots,
\]
which are more precisely determined as follows:

We first choose $u$ so that
\[
-n g + \alpha u = a, \quad g + \alpha u = a'
\]
become positive; then we set
\[
\delta = \frac{a' - a}{\alpha} = \frac{(n+1)g}{\alpha},
\]
and set
\[
\begin{gathered}
u' = u + \delta, \quad u'' = u' + \delta, \quad u''' = u'' + \delta, \dots,\\
a + \alpha\delta = a', \quad a' + \alpha\delta = a'', \quad a'' + \alpha\delta = a''', \dots,
\end{gathered}
\]
%CONT>>


% ===== scan leaf 697, printed 676 =====
% !PRINTED 676
\weberpage{676}
%<<CONT
and thus obtain from (12), if $\lambda_s', \lambda_s'', \dots$ are the conjugate logarithms of $\eta', \eta'', \dots$:
\weq{14}{
\begin{aligned}
a &< \Sigma g_s \lambda_s < a'\\
a' &< \Sigma g_s \lambda_s' < a''\\
a'' &< \Sigma g_s \lambda_s'' < a'''\\
&\cdots\cdots\cdots\cdots,
\end{aligned}
}
and the sequence of these inequalities can be continued without limit.

All these numbers $\eta, \eta' \eta'', \dots$ have moreover, by (6) and (7), the property that their absolute norms $N, N', N'', \dots$ remain below a definite finite limit $e^{nk}$.

% [sic] The source omits the comma between \eta' and \eta'' in the preceding list.

The number of possible values of $N, N', N'', \dots$ is therefore finite, namely certainly not greater than the greatest integer contained in $e^{nk}$. Likewise the number of all possible residues of the numbers $x_1, x_2, \dots, x_n$ [in (4)] with respect to all the moduli $N, N', N'', \dots$ is only a finite one, and from this it follows that, if we continue the sequence of the numbers (13) only far enough, there must occur in the sequence \gesperrt{two distinct numbers $\eta^{(h)}, \eta^{(k)}$ in which $N^{(h)} = N^{(k)}$ and the residues of $x_1, x_2, \dots, x_n$ with respect to the modulus $N^{(h)}$ are exactly the same.}

From (14), however, it follows, if $h > k$ is assumed:
\weq{15}{\Sigma g_s (\lambda_s^{(h)} - \lambda_s^{(k)}) > 0.}
If $N$ is the absolute norm of these two numbers $\eta^{(h)}, \eta^{(k)}$, then, because of the agreement of the residues of the $x$, the difference $\eta^{(h)} - \eta^{(k)}$ is divisible by $N$. On the other hand, by \S.\,138\,(13) the number $N$ is divisible by $\eta^{(h)}$, and consequently $\eta^{(h)} - \eta^{(k)}$ is also divisible by $\eta^{(h)}$. From this it results that $\eta^{(k)}$ is also divisible by $\eta^{(h)}$, and likewise $\eta^{(h)}$ by $\eta^{(k)}$. Both numbers are therefore associated, and if we set
\[
\frac{\eta^{(h)}}{\eta^{(k)}} = \varepsilon
\]
then $\varepsilon$ is a \gesperrt{unit} of the field $\Omega$.

If we further set, denoting by $|\varepsilon_s|$ the absolute value of $\varepsilon_s$,
\weq{16}{\delta_s \log|\varepsilon_s| = l_s,}
then from (15) there results:
\weq{17}{g_1 l_1 + g_2 l_2 + \cdots + g_\nu l_\nu > 0.}


% ===== scan leaf 698, printed 677 =====
% !PRINTED 677
\weberpage{677}
\pdfbookmark[2]{§ 186. Systems of independent units and systems of exponents of the units}{bm213}
This now proves the theorem:
% [source note] The source prints der Satz; den Satz is apparently intended, since the preceding argument establishes the theorem which follows. The English renders that intended direction; the German preserves the printed wording.

4. \gesperrt{If $g_1, g_2, \ldots, g_\nu$ is an arbitrary system of real numbers that are not all equal to one another, then in every algebraic field $\Omega$ a unit $\varepsilon$ can be so determined that the sum}
\[
g_1 l_1 + g_2 l_2 + \cdots + g_\nu l_\nu
\]
\gesperrt{is different from zero, when $l_s$ is the system of the conjugate logarithms of $\varepsilon$.}

This important theorem, which is the foundation for what follows, goes back, though in a somewhat altered version, to \gesperrt{Dirichlet}. This formulation was given by \gesperrt{Minkowski}.

\begin{center}
\S.\,186.\\
\gesperrt{Systems of independent units and systems of exponents of the units.}
\end{center}

If in the theorem last proved we set $g_1 = 1, g_2 = 0, \ldots, g_\nu = 0$, which is permissible once we assume from now on that $\nu > 1$, it follows that there is in $\Omega$ a unit $\varepsilon$ for which one of the conjugate logarithms, $l_1$, is different from zero. Since every power of $\varepsilon$ has the same property, there are also infinitely many such units.

This theorem now admits a very important generalisation:

5. \gesperrt{If $s \leqq \nu - 1$, then in the field $\Omega$ a system of units with the corresponding conjugate logarithms}
\weq{1}{
\begin{array}{l}
\varepsilon_1 : l_{1,1}, l_{1,2}, \ldots, l_{1,\nu} \\
\varepsilon_2 : l_{2,1}, l_{2,2}, \ldots, l_{2,\nu} \\
\cdots\cdots\cdots\cdots\cdots\cdots \\
\varepsilon_s : l_{s,1}, l_{s,2}, \ldots, l_{s,\nu}
\end{array}
}
\gesperrt{can be so determined that an arbitrary one of the $s$-rowed determinants of the matrix of the $l_{h,k}$, say}
\[
L_s = \Sigma \pm l_{1,1} l_{2,2} \ldots l_{s,s}
\]
\gesperrt{is different from zero.}

For $s = 1$ the theorem as stated is correct by the remark made at the outset. We therefore assume that it
%CONT>>


% ===== scan leaf 699, printed 678 =====
% !PRINTED 678
\weberpage{678}
%<<CONT
be proved for $s - 1$ and derive it in general by complete induction. To this end we arrange the determinant $L_s$ according to the elements of the last row, and write it thus:
\weq{2}{
L_s = g_1 l_{s,1} + g_2 l_{s,2} + \cdots + g_s l_{s,s},
}
in which $g_s = L_{s-1}$, which can therefore, by the assumption made, be taken to be different from zero. If we substitute these values of $g$ into the formula of Theorem 4., \S.\,185, assuming $g_{s+1} = 0, \ldots, g_\nu = 0$ while $g_s$ is not $= 0$, then, so long as $s < \nu$, the quantities $g_1, g_2, \ldots, g_\nu$ are certainly not all equal to one another, and it follows that $\varepsilon_s$ can be so chosen that $L_s$ is different from zero, as was to be proved.

This mode of reasoning cannot be extended to $s = \nu$, and the determinant $L_\nu$ must indeed always be zero, because for every unit $\sum\limits_{1,\nu}^{s} l_s$ vanishes.

We now wish, for the case that $s = \nu - 1$, to denote the determinant $L_s$ thus:
\weq{3}{
L_{\nu-1} = L(\varepsilon_1, \varepsilon_2, \ldots, \varepsilon_{\nu-1}),
}
and to call a system of units $\varepsilon_1, \varepsilon_2, \ldots, \varepsilon_{\nu-1}$ for which this determinant is different from zero a \gesperrt{system of independent units}.

The absolute value $L$ of the determinant $L(\varepsilon_1, \varepsilon_2, \ldots, \varepsilon_{\nu-1})$ is called the \gesperrt{regulator} of the system $\varepsilon_1, \varepsilon_2, \ldots, \varepsilon_{\nu-1}$\footnote{\gesperrt{Dirichlet-Dedekind}, Vorlesungen \"uber Zahlentheorie 4. Auflage, \S.\,183.}.

If in the determinant
\weq{4}{
\begin{vmatrix}
l_{1,1}, & l_{1,2}, & \ldots, & l_{1,\nu} \\
\cdots & \cdots & \cdots & \cdots \\
l_{\nu-1,1}, & l_{\nu-1,2}, & \ldots, & l_{\nu-1,\nu} \\
x_1, & x_2, & \ldots, & x_\nu
\end{vmatrix},
}
in which the $x_1, x_2, \ldots, x_\nu$ are arbitrary quantities, we add all the columns to the last, then, having regard to the relations $\sum\limits^{i} l_{s,i} = 0$
we obtain its value equal to
\[
\pm (x_1 + x_2 + \cdots + x_\nu) L.
\]
If therefore we let all the $x$ with the exception of one arbitrary
%CONT>>


% ===== scan leaf 700, printed 679 =====
% !PRINTED 679
\weberpage{679}
%<<CONT
% [sic] The source cites (3), but the determinant specialized by the x_i choices is (4) on printed p. 678.
$= 0$, and one of them $= 1$, then from (3) there results some one of the $(\nu-1)$-rowed determinants of the matrix
\[
\begin{matrix}
l_{1,1}, & l_{1,2}, & \ldots, & l_{1,\nu} \\
\cdots & \cdots & \cdots & \cdots \\
l_{\nu-1,1}, & l_{\nu-1,2}, & \ldots, & l_{\nu-1,\nu}
\end{matrix}
\]
and all these determinants therefore have (apart from the sign) the same value. It is therefore a matter of indifference which among the $\nu$ conjugate values is omitted in the formation of the regulator $L$.

If $\varepsilon_1, \varepsilon_2, \ldots, \varepsilon_{\nu-1}$ is an independent system of units, and if $\lambda_1, \lambda_2, \ldots, \lambda_\nu$ are the conjugate logarithms of any unit $\varepsilon$ in $\Omega$, then the numbers $\xi_1, \xi_2, \ldots, \xi_{\nu-1}$ can always be determined, and in only one way, so that
\weq{5}{
\begin{array}{l}
\xi_1 l_{1,1} + \xi_2 l_{2,1} + \cdots + \xi_{\nu-1} l_{\nu-1,1} = \lambda_1, \\
\xi_1 l_{1,2} + \xi_2 l_{2,2} + \cdots + \xi_{\nu-1} l_{\nu-1,2} = \lambda_2, \\
\cdots\cdots\cdots\cdots\cdots\cdots\cdots \\
\xi_1 l_{1,\nu} + \xi_2 l_{2,\nu} + \cdots + \xi_{\nu-1} l_{\nu-1,\nu} = \lambda_\nu
\end{array}
}
holds. For of these equations the first $\nu - 1$ form a system of linear equations with non-vanishing determinant, and the $\nu$th equation follows from the rest, because the sum of the conjugate logarithms of every unit vanishes.

The system of numbers $\xi_1, \xi_2, \ldots, \xi_{\nu-1}$ we call the \gesperrt{system of exponents} of the unit $\varepsilon$ with respect to the system $\varepsilon_1, \varepsilon_2, \ldots, \varepsilon_{\nu-1}$, and it is now to be proved:

6. \gesperrt{That the exponents of every unit $\varepsilon$ are rational numbers whose denominators do not exceed a certain finite limit, and which can therefore all be assumed to be affected with a common denominator depending only on the system $\varepsilon_i$.}

In order to see this, one has to consider the following:

1) If we regard the quantities $\xi_1, \xi_2, \ldots, \xi_{\nu-1}$ as variables, and let each of them, independently of the others, run from 0 to 1, then the left sides of (5) remain enclosed within finite limits depending only on the units $\varepsilon_1, \varepsilon_2, \ldots, \varepsilon_{\nu-1}$. Within these limits, then, the conjugate logarithms $\lambda$ of the unit $\varepsilon$ must also remain, so long as
%CONT>>


% ===== scan leaf 701, printed 680 =====
% !PRINTED 680
\weberpage{680}
%<<CONT
the exponents $\xi_i$ remain restricted to non-negative proper fractional values, and from this there results, by the definition of the conjugate logarithms, an upper limit for $\varepsilon$ and its conjugates.

But by Theorem 2., \S.\,185 there is in $\Omega$ only a finite number of integers, and hence all the more only a finite number of units which, in absolute value, lie together with all their conjugates below a finite limit.

If we now call a unit whose exponents lie within the limits 0 and 1, with inclusion of the lower and with exclusion of the upper limit, a unit \gesperrt{reduced with respect to the system $\varepsilon_1, \varepsilon_2, \ldots, \varepsilon_{\nu-1}$}, then we have the theorem:

\gesperrt{There is only a finite number of units in $\Omega$ which are reduced with respect to an independent system $\varepsilon_1, \varepsilon_2, \ldots, \varepsilon_{\nu-1}$.}

2) The units reproduce themselves by multiplication and division. If
\[
\begin{array}{l}
\xi_1', \xi_2', \ldots, \xi_{\nu-1}' \\
\xi_1'', \xi_2'', \ldots, \xi_{\nu-1}''
\end{array}
\]
are the exponents of two units $\varepsilon', \varepsilon''$, then the sums and the differences
\[
\xi_1' \pm \xi_1'', \xi_2' \pm \xi_2'', \ldots, \xi_{\nu-1}' \pm \xi_{\nu-1}''
\]
are the exponents of the units $\varepsilon' \varepsilon''$ and $\varepsilon' : \varepsilon''$.

This follows immediately from the theorem that the logarithm of a product or of a quotient is equal to the sum or to the difference of the logarithms of the two parts.

3) The units $\varepsilon_1, \varepsilon_2, \ldots$ themselves have the exponents $1, 0, \ldots, 0; 0, 1, \ldots, 0; \ldots$, and from this it follows by 2) that every system of integers, put for $\xi_1, \xi_2, \ldots, \xi_{\nu-1}$, gives the system of exponents of a unit which can be formed by multiplication of powers of the $\varepsilon_1, \varepsilon_2, \ldots$.

4) From 2) and 3) it results that a system of exponents $\xi_1, \xi_2, \ldots, \xi_{\nu-1}$ of a unit remains a system of exponents when all its elements are multiplied by one and the same integer, and that it still retains this character even when its elements $\xi_i$ are increased or diminished by arbitrary integers. If therefore we use the symbol $(x)$ in the sense of \S.\,182, namely that it denotes the ex-
%CONT>>


% ===== scan leaf 702, printed 681 =====
% !PRINTED 681
\weberpage{681}
\pdfbookmark[2]{§ 187. Fundamental systems of units}{bm214}
%<<CONT
cess of the number $x$ over the greatest integer contained in $x$,
the following results:

If $\xi_1, \xi_2, \ldots, \xi_{\nu-1}$ is the system of exponents of a unit, and
$m$ an integer, then

\weq{6}{(m\,\xi_1), (m\,\xi_2), \ldots, (m\,\xi_{\nu-1})}

is likewise the system of exponents of a unit.

Now the quantities $(m\,\xi_i)$, when they are not zero, are
positive proper fractions, and by 1) it must therefore happen, if the
sequence of integers is continued far enough, that for two
different values $m'$, $m''$ of $m$ the sequence of numbers (6)
coincides. Since $m'\xi_i$, $m''\xi_i$ accordingly have the same
excess over an integer, $(m' - m'')\,\xi_i$ is itself an
integer, and there is therefore an integral rational
number $m$, at most equal to the number of systems of exponents
of reduced units, with the property that

\[
m\,\xi_1, m\,\xi_2, \ldots, m\,\xi_{\nu-1}
\]

become integral rational numbers.

This number $m$ may change when the unit $\varepsilon$ is
changed. But since all possible values of this number lie below a
finite limit, there is a finite number in which
all these values of $m$ are contained and which can itself be
taken for $m$. There is therefore a certain positive
integer $m$ that can be taken as the denominator of all those rational
 numbers which can occur as exponents of
 units. Theorem 6. is thereby proved.

\begin{center}
\S.\,187.\\[4pt]
\gesperrt{Fundamental Systems of Units.}
\end{center}

In place of the system of independent units $\varepsilon_i$ we now
 choose another, whose conjugate logarithms shall be

\weq{1}{\lambda_{i,1}, \lambda_{i,2}, \ldots, \lambda_{i,\nu} \qquad i = 1, 2, \ldots, \nu-1}

By \S.\,186, (5) we have

\weq{2}{\begin{aligned}
&\lambda_{i,k} = \xi_{1,i}\, l_{1,k} + \xi_{2,i}\, l_{2,k} + \cdots + \xi_{\nu-1,i}\, l_{\nu-1,k} \\
&i = 1, 2, \ldots, \nu-1; \quad k = 1, 2, \ldots, \nu,
\end{aligned}}

where $\xi_{1,i}, \xi_{2,i}, \ldots, \xi_{\nu-1,i}$ denote the exponents of one of the new units
 with respect to the system $\varepsilon_1, \varepsilon_2, \ldots, \varepsilon_{\nu-1}$, which
we may assume to be rational numbers with the common denominator $m$.


% ===== scan leaf 703, printed 682 =====
% !PRINTED 682
\weberpage{682}
But by the multiplication theorem of determinants

\weq{3}{\begin{vmatrix} \lambda_{1,1}, & \cdots, & \lambda_{\nu-1,1} \\ \cdot & \cdot & \cdot \\ \lambda_{1,\nu-1}, & \cdots, & \lambda_{\nu-1,\nu-1} \end{vmatrix}
=
\begin{vmatrix} \xi_{1,1}, & \cdots, & \xi_{1,\nu-1} \\ \cdot & \cdot & \cdot \\ \xi_{\nu-1,1}, & \cdots, & \xi_{\nu-1,\nu-1} \end{vmatrix}
\begin{vmatrix} l_{1,1}, & \cdots, & l_{\nu-1,1} \\ \cdot & \cdot & \cdot \\ l_{1,\nu-1}, & \cdots, & l_{\nu-1,\nu-1} \end{vmatrix},}

or, if we set

\[
\Lambda_{\nu-1} = \Sigma \pm \lambda_{1,1}\lambda_{2,2} \cdots \lambda_{\nu-1,\nu-1}
\]

and observe that the determinant

\[
\Sigma \pm \xi_{1,1}\xi_{2,2} \cdots \xi_{\nu-1,\nu-1}
\]

is a rational number with the denominator $m^{\nu-1}$:

\weq{4}{\Lambda_{\nu-1} = \frac{a}{m^{\nu-1}} \, L(\varepsilon_1, \varepsilon_2, \ldots, \varepsilon_{\nu-1}),}

in which $a$ is an integral rational number.

From this it follows that the newly introduced units form an
independent system always and only then, when the
determinant of the $\xi_{i,k}$, that is, the number $a$, is different from zero.

Now among a finite or infinite number of
non-vanishing rational fractions with the same denominator
there is always a least in absolute value, and consequently
we can, by (4), choose the new system of independent units
so that its regulator $\pm \Lambda_{\nu-1}$ becomes \gesperrt{as small as possible}.
Such a system we call a \gesperrt{fundamental system
of units}, and the minimal value of the regulator itself,
that is, the regulator of a fundamental system, we call the
\gesperrt{regulator of the field}. We now assume that our
system $\varepsilon_1, \varepsilon_2, \ldots, \varepsilon_{\nu-1}$ is itself a fundamental system.

Under this assumption it can be proved that
all exponents of units must be \gesperrt{integers}.

For if we assume that there exists a unit $\varepsilon$
whose exponents, determined by the formulae \S.\,186, (5), are not
all integers, then by \S.\,186, 4) there is also a
unit $\varepsilon_0$ whose exponents are proper fractions which do not all
vanish.

If by $\lambda_1, \lambda_2, \ldots, \lambda_{\nu-1}$ in the formulae (5),
\S.\,186 we understand the system of conjugate logarithms of $\varepsilon_0$, then there
 results

\[
L(\varepsilon_0, \varepsilon_2, \ldots, \varepsilon_{\nu-1}) = \xi_1\, L(\varepsilon_1, \varepsilon_2, \ldots, \varepsilon_{\nu-1}),
\]
%CONT>>


% ===== scan leaf 704, printed 683 =====
% !PRINTED 683
\weberpage{683}
%<<CONT
and if we therefore assume, which is plainly no restriction of
generality, that $\xi_1$ is a non-vanishing proper fraction, then this
formula contradicts the assumption that

\[
\varepsilon_1, \varepsilon_2, \ldots, \varepsilon_{\nu-1}
\]

is a fundamental system, because the determinant $L_{\nu-1}$ would be
 diminished if $\varepsilon_1$ were replaced by $\varepsilon_0$.

If we have a fundamental system, we may take an arbitrary
 system of integral rational numbers $\xi_1, \xi_2, \ldots, \xi_{\nu-1}$ as
system of exponents and form a unit $\varepsilon$ with these exponents:

\[
\varepsilon = \varepsilon_1^{\xi_1} \varepsilon_2^{\xi_2} \ldots \varepsilon_{\nu-1}^{\xi_{\nu-1}}.
\]

There thus remains the question to be answered, how far
a unit is determined by its system of exponents.

Let us assume that $\varepsilon', \varepsilon''$ are two units with the same
system of exponents; then the system of exponents of the
unit $\varepsilon' : \varepsilon'' = \varrho$ consists of nothing but zeros, and consequently the conjugate
 logarithms of $\varrho$ are all equal to zero. Hence $\varrho$ is an
integer with the property that the absolute values of all
numbers conjugate to $\varrho$ are equal to 1, and therefore, as was proved in
\S.\,175, 2., a \gesperrt{root of unity}. With this the
following theorem is proved:

I. \gesperrt{In the field $\Omega$ there is a system of $\nu - 1$ fundamental
 units}

\weq{5}{\varepsilon_1, \varepsilon_2, \ldots, \varepsilon_{\nu-1},}

\gesperrt{which has the property that in the form}

\weq{6}{\varepsilon = \varrho\, \varepsilon_1^{\xi_1} \varepsilon_2^{\xi_2} \ldots \varepsilon_{\nu-1}^{\xi_{\nu-1}}}

\gesperrt{all units of the field are contained, each only once,
when $\xi_1, \xi_2, \ldots, \xi_{\nu-1}$ run through all integral rational numbers and $\varrho$ through all
roots of unity present in $\Omega$.}

It is now easy to derive all the other fundamental systems
from one of them, by putting in (6) for the exponents $\nu - 1$ different
 systems of integers $\xi_{i,k}$ whose determinant
$= \pm 1$. For if one sets

\[
\varepsilon_i' = \varrho_i\, \varepsilon_1^{\xi_{1,i}} \varepsilon_2^{\xi_{2,i}} \ldots \varepsilon_{\nu-1}^{\xi_{\nu-1,i}},
\]

then it follows from (3):

\[
L(\varepsilon_1', \varepsilon_2', \ldots, \varepsilon_{\nu-1}') = \pm L(\varepsilon_1, \varepsilon_2, \ldots, \varepsilon_{\nu-1}).
\]


% ===== scan leaf 705, printed 684 =====
% !PRINTED 684
\weberpage{684}
\pdfbookmark[2]{§ 188. Reduced numbers}{bm215}
% [source note] The same-parity restriction on x,y below applies only to odd m. For a fundamental discriminant -m the general condition is x congruent to m*y modulo 2; for m=4, x=0,y=1 represents i. The printed blanket restriction is retained.
Both determinants therefore have the minimal value, and
both systems are fundamental systems.

The case $\nu = 1$, which we excluded above, occurs, apart
from the case of the rational field, in which only the two
units $\pm 1$ exist, in the case of the \gesperrt{imaginary quadratic
 field}. In this case we have represented, in the twentieth
 chapter, the integers of the field in the form

\[
\frac{x + y\sqrt{-m}}{2}
\]

in which $-m$ denotes a fundamental discriminant, so
that $m \equiv 0$ or $\equiv 3\ (\mathrm{mod}\ 4)$, and $x, y$ are both even or
both odd. We obtain the units when we solve the
equation

\[
x^2 + m y^2 = 4
\]

in all possible ways in integral rational numbers. But this
equation has, when $m > 4$, only the two solutions
$x = \pm 2$, $y = 0$, and so in these cases, as in the
field of rational numbers, there are only the two units $\pm 1$. If
$m = 4$, one finds the four units $\pm 1, \pm i$, and if finally
 $m = 3$, the six units

\[
\pm 1, \; \pm \frac{-1 + i\sqrt{3}}{2}, \; \pm \frac{-1 - i\sqrt{3}}{2}.
\]

In the next simplest case, $n = 2$, $\nu = 2$, that is, in the
real quadratic field, the theory of units coincides
with the theory of the Pell equation treated in \S.\,127 of the
first volume.

\begin{center}
\S.\,188.\\[4pt]
\gesperrt{Reduced Numbers.}
\end{center}

If $\alpha$ is any integral or fractional number of the field $\Omega$
 different from zero, then we consider, as with the
units, the \gesperrt{system of conjugate logarithms} of $\alpha$,
by which we understand, as in \S.\,185, the \gesperrt{real numbers}

\weq{1}{\lambda_1 = \delta_1 \log|\alpha_1|, \; \lambda_2 = \delta_2 \log|\alpha_2|, \; \ldots, \; \lambda_\nu = \delta_\nu \log|\alpha_\nu|}

which we also denote by $\lambda_s(\alpha)$. If we consult the
fundamental system $\varepsilon_1, \varepsilon_2, \ldots, \varepsilon_{\nu-1}$ of units, then
%CONT>>


% ===== scan leaf 706, printed 685 =====
% !PRINTED 685
\weberpage{685}
%<<CONT
we can form the linear equations, in analogy with \S.\,186, (5):
\weq{2}{
\begin{aligned}
\xi_1 l_{1,1} + \xi_2 l_{2,1} + \cdots + \xi_{\nu-1} l_{\nu-1,1} + \delta_1 \xi_\nu &= \lambda_1, \\
\xi_1 l_{1,2} + \xi_2 l_{2,2} + \cdots + \xi_{\nu-1} l_{\nu-1,2} + \delta_2 \xi_\nu &= \lambda_2, \\
\cdot\ \cdot\ \cdot\ \cdot\ \cdot\ \cdot\ \cdot\ \cdot\ \cdot\ \cdot\ \cdot\ \cdot & \\
\xi_1 l_{1,\nu} + \xi_2 l_{2,\nu} + \cdots + \xi_{\nu-1} l_{\nu-1,\nu} + \delta_\nu \xi_\nu &= \lambda_\nu,
\end{aligned}
}
whose determinant [\S.\,186, (4)] has the value, different from zero:
\weq{3}{
\left|
\begin{array}{ccccc}
l_{1,1}, & l_{2,1}, & \ldots, & l_{\nu-1,1}, & \delta_1\\
l_{1,2}, & l_{2,2}, & \ldots, & l_{\nu-1,2}, & \delta_2\\
\cdot & \cdot & \cdot & \cdot & \cdot\\
l_{1,\nu}, & l_{2,\nu}, & \ldots, & l_{\nu-1,\nu}, & \delta_\nu
\end{array}
\right| = n\,L\,(\varepsilon_1, \varepsilon_2, \ldots, \varepsilon_{\nu-1})
}
Accordingly the unknowns $\xi_1, \xi_2, \ldots, \xi_{\nu-1}, \xi_\nu$ are
\gesperrt{uniquely} determined from (2).

By addition of the equations (2) one finds first of all,
very simply,
\weq{4}{
n\,\xi_\nu = \log N_a(\alpha),
}
and $\xi_\nu$ therefore remains the same for all \gesperrt{numbers $\alpha$ that are
 associated with one another.}

The numbers $\xi_1, \xi_2, \ldots, \xi_{\nu-1}$ are called the exponents of the
number $\alpha$ [with respect to the fundamental system $(\varepsilon_1, \varepsilon_2, \ldots, \varepsilon_{\nu-1})$].

From this definition it follows at once that the
exponents of a product or of a quotient are equal to the
sum or the difference of the corresponding exponents of the
individual parts. The exponents of a unit are
integral rational numbers, and every system of integral rational numbers
$\xi_1, \xi_2, \ldots, \xi_{\nu-1}$ is also the system of exponents of a unit.

The exponents of associated numbers therefore differ
from one another by integers, and to every number $\alpha$
one can find an associated number $\alpha_0$ whose exponents lie between 0
and 1. Such numbers are called \gesperrt{reduced numbers} (with
respect to the fundamental system $\varepsilon_1, \ldots, \varepsilon_{\nu-1}$). Reduced units
are all and only the \gesperrt{roots of unity} present in $\Omega$,
\gesperrt{whose number we shall denote by $w$.} Since
the two roots of unity $\pm 1$ are contained in every field,
$w$ is at least $= 2$ and always an even number.

If two associated numbers have the same system of exponents,
then they differ only by a factor which
%CONT>>


% ===== scan leaf 707, printed 686 =====
% !PRINTED 686
\weberpage{686}
\pdfbookmark[2]{§ 189. Bounds for the number of integers of the field Ω divisible by an ideal}{bm216}
%<<CONT
is a root of unity. The system of exponents of the reduced number
$\alpha_0$ derived from $\alpha$ is completely determined by $\alpha$ itself.
This proves the theorem:

II. \gesperrt{To every number $\alpha$ (different from zero) of the
field $\Omega$ there are $w$, and not more, reduced
numbers $\varrho\,\alpha_0$.}

If in place of the fundamental system $(\varepsilon_1, \varepsilon_2, \ldots, \varepsilon_{\nu-1})$ another
one is used, the exponents $\xi_1, \xi_2, \ldots, \xi_{\nu-1}$ undergo
a linear substitution with integral coefficients of determinant $\pm 1$.
A reduced number may then cease to be reduced for the new fundamental
system. The concept of the reduced number
therefore depends on the choice of the fundamental system. But the
number of the reduced numbers derived from a given number
is always the same.

By this theorem we have attained our purpose, namely to single
out, from the whole infinite system of numbers associated with one
another, a definite finite number of representatives.

\begin{center}
\S.\,189.\\[4pt]
\gesperrt{Limits of the number of the integers of the field $\Omega$
divisible by an ideal.}
\end{center}
% [source note] The following finiteness assertion and definition of T omit counting modulo association. Infinitely many units already have absolute norm 1. The intended finite count is of non-associated integers (equivalently the corresponding principal ideals), explicitly specified on printed p. 690. The source also changes the endpoint from norm at most t here to natural numbers less than t below; this can alter a boundary term but not the limiting quotient. The wording is preserved.

Our earlier considerations have shown that every
integer of a field $\Omega$ has only a finite number of ideals
as divisors. Since every integral ideal is a divisor of its norm,
there are therefore also only a finite number of integers
in $\Omega$ whose absolute norm does not exceed a given
limit.

From these numbers we now single out a part once more,
and ask for the number $T$ of all the \gesperrt{integers}
of the field $\Omega$ whose absolute norm does not exceed a positive quantity $t$
and which are divisible by a given ideal $\mathfrak{a}$.

Let us take the rational field as a preliminary example;
there $T$ is the number of the natural numbers that are less than
$t$ and divisible by a given integer $m$, that is, in
the notation used earlier (\S.\,182), $T = E\left(\frac{t}{m}\right)$.


% ===== scan leaf 708, printed 687 =====
% !PRINTED 687
\weberpage{687}
The number $T$ is a function of $t$ which in general cannot be
 determined more precisely; it grows to infinity together with $t$,
but, since it takes only integral values, it changes only
stepwise, and therefore discontinuously. What matters here is
\gesperrt{the limiting value which the ratio $T:t$ approaches as
$t$ grows infinitely} (which in the example cited
 above is $1:m$).

This limiting value can be found by considerations of the kind
we have already used in \S.\,155.

We first form a basis form of the given ideal $\mathfrak{a}$:
\weq{1}{
\alpha = \alpha_1 x_1 + \alpha_2 x_2 + \cdots + \alpha_n x_n,
}
that is, a linear function of the $n$ variables $x_i$ with the property
that, by substituting all integral rational numbers for
$x_i$, there arise from $\alpha$ all the integers of the field $\Omega$ divisible by $\mathfrak{a}$
 (\S.\,146). By $\alpha_{i,1}, \alpha_{i,2}, \ldots, \alpha_{i,n}$ we denote the values conjugate
to $\alpha_i$, and we form the linear substitution for $x_i$:
\weq{2}{
\begin{aligned}
y_1 &= \alpha_{1,1} x_1 + \alpha_{2,1} x_2 + \cdots + \alpha_{n,1} x_n\\
y_2 &= \alpha_{1,2} x_1 + \alpha_{2,2} x_2 + \cdots + \alpha_{n,2} x_n\\
&\ \cdot\ \cdot\ \cdot\ \cdot\ \cdot\ \cdot\ \cdot\ \cdot\ \cdot\ \cdot\ \cdot\ \cdot\\
y_n &= \alpha_{1,n} x_1 + \alpha_{2,n} x_2 + \cdots + \alpha_{n,n} x_n.
\end{aligned}
}
If $A$ is the determinant of this system of equations
\[
A = \Sigma \pm \alpha_{1,1} \alpha_{2,2} \ldots \alpha_{n,n},
\]
then $A^2$ is the discriminant of the system of functions $(\alpha_1, \alpha_2, \ldots, \alpha_n)$,
and therefore, by \S.\,147, equal to the product of the square
of the absolute norm of $\alpha$ and the field discriminant $\Delta$.
Since the absolute norm of $\alpha$ is at the same time the norm of the ideal $\mathfrak{a}$
(\S.\,151), we set
\weq{3}{
A^2 = N(\mathfrak{a})^2 \, \Delta.
}
Now, as in \S.\,155, we may interpret $x_1, x_2, \ldots, x_n$ as
coordinates of a point $(x)$ in a space of $n$ dimensions,
 and we let a point $(x')$ correspond to every point $(x)$,
whose coordinates are determined by
\weq{4}{
x_1' = t^{\frac{1}{n}} x_1, \quad x_2' = t^{\frac{1}{n}} x_2, \quad \ldots, \quad x_n' = t^{\frac{1}{n}} x_n
}
where $t$ denotes an arbitrary positive quantity.
To every region $S$ of points $(x)$ there corresponds a region $S'$ of
points $(x')$. All figures in $S$ are similar to the corresponding figures
in $S'$, and the linear dimensions in $S$ are to
those in $S'$ as $1 : t^{\frac{1}{n}}$.


% ===== scan leaf 709, printed 688 =====
% !PRINTED 688
\weberpage{688}
If we think of the region $S$ as given, the corresponding
 region $S'$ will depend on $t$ and will grow with $t$.
The points with \gesperrt{integer} coordinates $x_i$ we call,
as before, the \gesperrt{lattice points}. The number of lattice points
lying in a finite region $S'$ is always finite,
but grows with $t$, and shall be denoted by $Z_t$. Then, by
the theorems of \S.\,155, (14), the volume of the region $S$, that is, the
$n$-fold integral extended over all points of $S$
\[
V = \int\!\!\int \ldots \int dx_1\, dx_2 \ldots dx_n
\]
is equal to the limiting value of the ratio $Z_t : t$ for infinitely
increasing $t$:
\weq{5}{
V = \mathop{\mathrm{Lim}}\limits_{t=\infty} \frac{Z_t}{t}.
}
Herein lies at the same time the more exact specification that we must
make concerning the region $S$, if our consideration is to be
 applicable. $S$ must be so constituted that the
$n$-fold integral $V$ has a definite meaning. Every lattice point
 is now, by (1), the image of an integer of the field $\Omega$ divisible
 by $\mathfrak{a}$, and we can therefore also say
that $Z_t$ is the number of integers in $\Omega$ divisible by $\mathfrak{a}$
whose images lie in the region $S'$.

In order to delimit the region $S'$ in a suitable manner, we choose
a system of fundamental units of the field $\Omega$
\[
\varepsilon_1, \varepsilon_2, \ldots, \varepsilon_{\nu-1}
\]
with the system of conjugate logarithms
\[
l_{i,1}, l_{i,2}, \ldots, l_{i,\nu}
\]
and the regulator
\weq{6}{
\pm L(\varepsilon_1, \varepsilon_2, \ldots, \varepsilon_{\nu-1}) = \pm \Sigma \pm l_{1,1} l_{2,2} \ldots l_{\nu-1,\nu-1},
}
which may be denoted by $L$.

We now define a system of functions $\xi_1, \xi_2, \ldots, \xi_\nu$
of the variables $x_i$ by the following conditions:

Let, when the $y_i$ have the meaning (2),
\weq{7}{
z_i = \delta_i \log |y_i|,
}
so that, if a system of integral rational numbers is
substituted for the $x_i$, $z_i$ passes, by \S.\,185, into the system of
conjugate logarithms of the number $\alpha$.

Then, if $y_k$ belongs to a real field,
\weq{8}{
z_k = \tfrac{1}{2} \log y_k^2,
}
%CONT>>


% ===== scan leaf 710, printed 689 =====
% !PRINTED 689
\weberpage{689}
%<<CONT
and if $y_k, y_k'$ belong to an imaginary pair,
\weq{9}{z_k = \log y_k y_k',}
by which the definition becomes completely unambiguous.

Now we determine the variables $\xi_1, \xi_2, \ldots, \xi_{\nu-1}, \xi_\nu$ from the linear equations
\weq{10}{\begin{array}{l}
\xi_1 l_{1,1} + \cdots + \xi_{\nu-1} l_{\nu-1,1} + \delta_1 \xi_\nu = z_1\\
\xi_1 l_{1,2} + \cdots + \xi_{\nu-1} l_{\nu-1,2} + \delta_2 \xi_\nu = z_2\\
. \; . \; . \; . \; . \; . \; . \; . \; . \; . \; . \; . \; . \; .\\
\xi_1 l_{1,\nu} + \cdots + \xi_{\nu-1} l_{\nu-1,\nu} + \delta_\nu \xi_\nu = z_\nu,
\end{array}}
whose determinant does not vanish.

If one gives the variables $x_i$ integer values, then by \S.\,188, (2), (3) the $\xi_1, \xi_2, \ldots, \xi_{\nu-1}$ pass over into the system of exponents of the number $\alpha$, and by addition of the equations (10) [\S.\,185, (3)] there results:
\weq{11}{\xi_\nu = \frac{1}{n} \log N_a(\alpha).}

We now pass from the point $(x)$ to the point $(x')$ by setting $x_i' = t^{\frac{1}{n}} x_i$. Thereby $y_i$ passes into
\[ y_i' = t^{\frac{1}{n}} y_i \]
and $z_i$ into
\[ z_i' = z_i + \frac{1}{n} \delta_i \log t. \]

If, then, we put $z_i'$ in place of $z_i$, and denote the values of $\xi_1, \xi_2, \ldots, \xi_\nu$ that then result from (10) by $\xi_1', \xi_2', \ldots, \xi_\nu'$, we have
\weq{12}{\xi_1' = \xi_1, \xi_2' = \xi_2, \ldots, \xi_{\nu-1}' = \xi_{\nu-1}, \xi_\nu' = \xi_\nu + \frac{1}{n} \log t.}

We now propose to delimit the region $S'$ by requiring that
\weq{13}{\begin{array}{c}
0 \leqq \xi_1' < 1, \ldots, 0 \leqq \xi_{\nu-1}' < 1, \\
\xi_\nu' < \frac{1}{n} \log t
\end{array}}
shall hold. By this we achieve that the lattice points $(x')$ lying in the region $S'$ represent only reduced integers $\alpha$, and all and only those whose absolute norm is less than $t$.

% German ends 'die Anzahl der in $\Omega$ enthaltenen |' with the noun on printed 690;
% English word order forces the noun onto this page so that 'contained in $\Omega$' can
% follow it. i0713 opens correspondingly at 'and if we therefore'.
But among the reduced numbers there are always $w$ at a time associated with one another, if $w$ denotes the number of the roots of unity contained in $\Omega$,
%CONT>>


% ===== scan leaf 713, printed 690 =====
% !PRINTED 690
\weberpage{690}
\pdfbookmark[2]{§ 190. Determination of the volume}{bm217}
%<<CONT
and if we therefore combine these $w$ associated numbers into one complex, then the number of these complexes is equal to the number $T$ of all non-associated integers in $\Omega$ divisible by $\mathfrak{a}$ whose absolute norm is less than $t$. Accordingly the number of lattice points lying in $S'$ is
\weq{14}{Z_t = w T.}

For the boundary of the region $S$ one obtains from (12) and (13):
\weq{15}{0 \leqq \xi_1 < 1, \ldots, 0 \leqq \xi_{\nu-1} < 1, \xi_\nu < 0,}
and from (5) we obtain
\weq{16}{\mathrm{Lim}\, \frac{T}{t} = \frac{1}{w} \int\!\int \ldots \int dx_1\, dx_2 \ldots dx_n = \frac{1}{w} V,}
in which the $n$-fold integral is to be extended over the region $S$ determined by (15).

\begin{center}
\S.\,190.

\gesperrt{Determination of the Volume.}
\end{center}

The volume $V$ can be determined by the transformation formula for multiple integrals already cited in \S.\,155, according to which, if $x_1, x_2, \ldots, x_n$ are functions of the variables $u_1, u_2, \ldots, u_n$,
\weq{1}{\begin{aligned}
V &= \int\!\int \ldots \int dx_1\, dx_2 \ldots dx_n\\
&= \int\!\int \ldots \int \left( \Sigma \pm \frac{\partial x_1}{\partial u_1} \frac{\partial x_2}{\partial u_2} \cdots \frac{\partial x_n}{\partial u_n} \right) du_1\, du_2 \ldots du_n
\end{aligned}}
holds, where the functional determinant in the integral taken with respect to the variables $u$ is to be taken with its absolute value.

This functional determinant is also equal to the reciprocal value of the determinant
\[ \Sigma \pm \frac{\partial u_1}{\partial x_1} \frac{\partial u_2}{\partial x_2} \cdots \frac{\partial u_n}{\partial x_n}. \]

We first introduce the variables $u$ by a linear substitution, and remark that among the fields conjugate with $\Omega$ there are, according to our assumption, $n - \nu$ conjugate imaginary pairs
%CONT>>


% ===== scan leaf 714, printed 691 =====
% !PRINTED 691
\weberpage{691}
%<<CONT
and $2\nu - n$ real ones occur. Accordingly, of the $n$ functions $y$ [§. 189, (2)]

\weq{2}{y_k = \alpha_{1,k}x_1 + \alpha_{2,k}x_2 + \cdots + \alpha_{n,k}x_n.}

$2\nu - n$ are real and $2n-2\nu$ conjugate in pairs. By the $n$ variables $u$ we shall understand real linear combinations of the $x$, namely the $2\nu-n$ real $y$ and the $2n-2\nu$ parts of the conjugate imaginary pairs of the $y$, so that, when $y, y'$ is an imaginary pair, one sets
\[
2u_1 = y+y', \quad 2iu_2 = y-y'
\]
or

\weq{3}{y = u_1+iu_2, \quad y' = u_1-iu_2.}

If, in the functional determinant of the $y$, one applies to each pair of conjugate imaginary rows twice the elementary determinant theorem on the addition of rows (Bd. I, §. 22), there results

\[
\begin{aligned}
\Sigma \pm \frac{\partial u_1}{\partial x_1}\frac{\partial u_2}{\partial x_2}\cdots\frac{\partial u_n}{\partial x_n}
&= \frac{1}{(-2i)^{n-\nu}}\Sigma\pm\frac{\partial y_1}{\partial x_1}\frac{\partial y_2}{\partial x_2}\cdots\frac{\partial y_n}{\partial x_n} \\
&= \frac{A}{(-2i)^{n-\nu}},
\end{aligned}
\]

in which $A$ has the meaning fixed in §. 189, (3).

The field discriminant is positive or negative according as the number of conjugate imaginary pairs, that is $n-\nu$, is even or odd; and accordingly, if the absolute value of the field discriminant is denoted by $\pm\varDelta$, one obtains from §. 189, (4), (16):

\weq{4}{\mathrm{Lim}\,\frac{T}{t} = \frac{2^{n-\nu}}{w\,N(\mathfrak{a})\sqrt{\pm\varDelta}}\int\int\cdots\int du_1\,du_2\ldots du_n.}

Thus the integral

\[
U = \int\int\cdots\int du_1\,du_2\ldots du_n
\]

still remains to be treated.

From the equations §. 189, (10) one obtains, for every system of values of the variables $\xi_1, \xi_2, \ldots, \xi_\nu$ that satisfies the conditions (15), a definite finite system of values for $z_1, z_2, \ldots, z_\nu$, and from these, by §. 189, (7), an absolute value different from zero for every $y$. If, in the case of an imaginary pair (3), one sets

\weq{5}{u_1 = e^{1/2\,z}\cos\vartheta, \quad u_2 = e^{1/2\,z}\sin\vartheta,}

\weq{6}{y_1y_1' = u_1^2+u_2^2 = e^z,}
%CONT>>


% ===== scan leaf 715, printed 692 =====
% !PRINTED 692
\weberpage{692}
%<<CONT
it follows that, if $\vartheta$ is taken within the limits
\[
0 \leqq \vartheta < 2\pi
\]
then, for a fixed $z$, there belongs to every value of $\vartheta$ a corresponding pair of values $u_1, u_2$. Such an angle $\vartheta$ belongs to every imaginary pair.

To every point of the region $S$ there then belongs one and only one system of values of the variables $\xi_1, \xi_2, \ldots, \xi_\nu, \vartheta \ldots$ within the bounds

\weq{7}{\begin{aligned}
&0 \leqq \xi_1 < 1, \ldots, 0 \leqq \xi_{\nu-1} < 1 \\
&\xi_\nu < 0, \qquad 0 \leqq \vartheta < 2\pi, \ldots
\end{aligned}}
% [source notation] Both first-row less-than-or-equal signs place the equality bars above the angle; other occurrences on this page place them below. Both variants express the same relation and are represented by \leqq. No damaged-type or reversed-inequality claim is made.

If conversely any system of values of the variables $\xi, \vartheta$ within the limits (7) is given, one obtains from it the corresponding parts $u_1, u_2$ of an imaginary $y$ uniquely, but of the real $y$ only the absolute values, and so each of the real $y$ may still be given either of the two signs. The number of these possible determinations is $2^{2\nu-n}$. If one of these determinations has been singled out, the variables $x$ are uniquely determined, and they are finite for every finite system of values of the $\xi, \vartheta$. For a $\xi_\nu$ decreasing towards $-\infty$ the values of the $y$ will in part converge to zero.

The region $S$ accordingly splits into $2^{2\nu-n}$ partial regions $S_1, S_2, \ldots$, each of which is characterised by a definite sign-combination of the real $y_1, y_2, \ldots$, and accordingly the integral $U$ splits into just as many partial integrals $U_1, U_2, \ldots$

\[
U = U_1+U_2+\cdots
\]

In carrying out the integration in one of these parts, all of which prove to be of equal value, we begin with an imaginary pair, if such a one is present.

If we regard $u_1, u_2$ as rectangular coordinates in a plane, then by (5) we may conceive $e^{1/2\,z}, \vartheta$ as polar coordinates, and obtain

\[
\int\int du_1\,du_2 = \frac12\int\int e^z\,dz\,d\vartheta,
\]

in which the integration with respect to $\vartheta$, which extends from $0$ to $2\pi$, can be carried out. One thus finds

\[
\int\int du_1\,du_2 = \pi\int e^z\,dz.
\]


% ===== scan leaf 716, printed 693 =====
% !PRINTED 693
\weberpage{693}
Now one changes the order of integration and puts first a second imaginary pair, if such a pair is still present, which is then treated again in the same way.

The various parts $U_1, U_2, \ldots$ differ in the signs of the $u$ corresponding to the real $y$. If we take the sign so that $\pm u$ becomes positive, and set

\weq{8}{z = \log(\pm u), \quad \pm du = e^z\,dz,}

then for the various $U_1, U_2, \ldots$ we obtain the same expression by means of a $\nu$-fold integral

\weq{9}{U_1 = \pi^{n-\nu}\int\int\cdots\int e^{z_1+z_2+\cdots+z_\nu}\,dz_1\,dz_2\ldots dz_\nu,}

provided we assume $dz_1, dz_2, \ldots, dz_\nu$ to be positive.

But by (6) and (8) the variables $z_k$ are here none other than those in §. 189, (8), (9), (10), and accordingly

\[
z_1+z_2+\cdots+z_\nu = n\xi_\nu.
\]

If we now apply the transformation formula (1) to the $\nu$-fold integral (9) in order to introduce the variables $\xi_1, \xi_2, \ldots, \xi_\nu$, then we have to form the determinant

\[
\Sigma\pm\frac{\partial z_1}{\partial\xi_1}\frac{\partial z_2}{\partial\xi_2}\cdots\frac{\partial z_\nu}{\partial\xi_\nu}
\]

which by §. 188, (3) agrees in absolute value with $nL$, that is, with the $n$-fold value of the \gesperrt{field regulator} $L$, and accordingly

\[
U_1 = \pi^{n-\nu}nL\int\limits_0^1 d\xi_1\ldots\int\limits_0^1 d\xi_{\nu-1}\int\limits_{-\infty}^0 e^{n\xi_\nu}\,d\xi_\nu = \pi^{n-\nu}L
\]

and

\[
U = 2^{2\nu-n}\pi^{n-\nu}L.
\]

From this there follows at last, by (4), the theorem:

1. \gesperrt{If $T$ denotes the number of non-associated integers divisible by $\mathfrak{a}$ whose absolute norm is less than $t$, then}

\weq{10}{\mathop{\mathrm{Lim}}\limits_{t=\infty}\frac{T}{t} = \frac{2^\nu\pi^{n-\nu}L}{w\,N(\mathfrak{a})\sqrt{\pm\varDelta}} = \frac{g}{N(\mathfrak{a})},}

\gesperrt{where $g$ is a positive number completely determined by the nature of the field $\Omega$, namely:}

\weq{11}{g = \frac{2^\nu\pi^{n-\nu}L}{w\sqrt{\pm\varDelta}}.}


% ===== scan leaf 717, printed 694 =====
% !PRINTED 694
\weberpage{694}
\pdfbookmark[2]{§ 191. Theorems from the theory of series}{bm218}
\begin{center}
§. 191.

\gesperrt{Theorems from the theory of series.}
\end{center}

For the further applications of the results obtained hitherto, certain theorems from the theory of infinite series are required, and these shall first be derived here.

1. \gesperrt{If $a_1, a_2, a_3, \ldots, a_n\ldots$ is an unlimited system of real (positive, negative, or also vanishing) quantities, of which we assume that the sum}

\weq{1}{\sigma_n = a_1+a_2+\cdots+a_n}

\gesperrt{remains, for every arbitrary $n$, in absolute value below a finite limit $C$, then the series}

\weq{2}{S = \frac{a_1}{1^s}+\frac{a_2}{2^s}+\frac{a_3}{3^s}+\cdots}

\gesperrt{is, for every positive $s$, not only convergent, but also a continuous function of $s$.}

To prove this theorem, we decompose $S$ in the following way:
\[
S = S_m+R_m,
\]
in which
\[
S_m = \frac{a_1}{1^s}+\frac{a_2}{2^s}+\cdots+\frac{a_{m-1}}{(m-1)^s},
\]
\[
R_m = \frac{a_m}{m^s}+\frac{a_{m+1}}{(m+1)^s}+\frac{a_{m+2}}{(m+2)^s}+\cdots
\]
Now by (1):
\[
a_m = \sigma_m-\sigma_{m-1}, \quad a_{m+1}=\sigma_{m+1}-\sigma_m, \quad a_{m+2}=\sigma_{m+2}-\sigma_{m+1}, \ldots
\]
and therefore
\[
\begin{aligned}
R_m+\frac{\sigma_{m-1}}{m^s}
&= \sigma_m\left(\frac{1}{m^s}-\frac{1}{(m+1)^s}\right) \\
&\quad +\sigma_{m+1}\left(\frac{1}{(m+1)^s}-\frac{1}{(m+2)^s}\right)+\cdots
\end{aligned}
\]
Since now, by the assumption, $\sigma_{m-1}, \sigma_m, \sigma_{m+1}, \ldots$ are enclosed between finite limits $\pm C$, it follows from this that $R_m+\dfrac{\sigma_{m-1}}{m^s}$ is contained between the two limits
\[
\pm C\left(\frac{1}{m^s}-\frac{1}{(m+1)^s}+\frac{1}{(m+1)^s}-\frac{1}{(m+2)^s}+\cdots\right) = \pm\frac{C}{m^s}
\]
%CONT>>


% ===== scan leaf 718, printed 695 =====
% !PRINTED 695
\weberpage{695}
%<<CONT
enclosed, and in absolute value
\[
R_m < \frac{2\,C}{m^s},
\]
or, if $s > c$,
\[
R_m < \frac{2\,C}{m^c}.
\]

This upper limit for $R_m$, which is independent of $s$, can however, if $c$ is positive, be pressed down below any value however small by taking $m$ sufficiently large.

But from this there follows not only the convergence, but also the continuity of $S$. For if one takes $m$ sufficiently large, then not only $R_m$, but also the fluctuation of $R_m$ under varying $s$ becomes infinitely small, and $S_m$ is, for a fixed $m$, a continuous function of $s$. Hence $S$ too is continuous.
% [sic] The German source has a singular subject (die Schwankung) with the plural verb werden.

It is also to be remarked here that the assumption about $\sigma_m$ by no means presupposes the convergence of the infinite series $\Sigma a_k$. It is not even necessary that the $a_k$ be real; for if they are imaginary, one need only apply the theorem to the real and to the imaginary part. Finally, it is also not necessary to assume the $a_1, a_2, a_3, \dots$ to be constants. Everything remains valid if they are \gesperrt{continuous functions} of $s$.

2. \gesperrt{If $\mu_1, \mu_2, \dots, \mu_n, \dots$ denotes an infinite set of positive numbers of such a nature that two finite numbers $\alpha, \beta$ can be assigned in such a way that for every $n$, however large,}
\weq{3}{\alpha < \dfrac{n}{\mu_n} < \beta,}
\gesperrt{then the infinite series}
\[
S = \frac{1}{\mu_1^s} + \frac{1}{\mu_2^s} + \frac{1}{\mu_3^s} + \cdots
\]
\gesperrt{is convergent so long as the exponent $s$ is greater than 1, and the product $(s-1)\,S$ likewise remains between the limits $\alpha$ and $\beta$ when $s$ approaches the value 1 sufficiently closely, or at least approaches one of these values without limit.}


% ===== scan leaf 719, printed 696 =====
% !PRINTED 696
\weberpage{696}
The proof is very simply carried out by reduction to an integral. For evidently
\[
\frac{1}{(n+1)^s} < \int\limits_n^{n+1} \frac{dx}{x^s} < \frac{1}{n^s},
\]
because the first or the second value, $(n+1)^{-s}$ or $n^{-s}$, results when, under the integral sign, the too small or the too large value $(n+1)^{-s}$ or $n^{-s}$ is put in place of $x^{-s}$.

Accordingly there results from (3), for every positive $s$,
\[
\alpha^s \int\limits_n^{n+1} \frac{dx}{x^s} < \frac{1}{\mu_n^s} < \beta^s \int\limits_{n-1}^n \frac{dx}{x^s}.
\]
If we put
\[
R_m = \frac{1}{\mu_{m+1}^s} + \frac{1}{\mu_{m+2}^s} + \frac{1}{\mu_{m+3}^s} + \cdots,
\]
then there results from this
\[
\alpha^s \int\limits_{m+1}^{\infty} \frac{dx}{x^s} < R_m < \beta^s \int\limits_m^{\infty} \frac{dx}{x^s}
\]
or
\weq{4}{\dfrac{\alpha^s}{(s-1)(m+1)^{s-1}} < R_m < \dfrac{\beta^s}{(s-1)\,m^{s-1}}.}

Thus $R_m$ remains finite, however many terms may be summed, provided $s > 1$ holds; and for series with positive terms this is a sufficient condition for convergence.

If one further puts (4) into the form
\[
\frac{\alpha^s}{(m+1)^{s-1}} < (s-1)\,R_m < \frac{\beta^s}{m^{s-1}},
\]
then one sees that the two limits can be brought arbitrarily near to $\alpha, \beta$ if $s-1$ is assumed small enough; and from this it follows that $(s-1)\,R_m$, for sufficiently small $s-1$, can no longer lie outside the limits $\alpha, \beta$, unless it approaches one of these values without limit as $s-1$ decreases. Now since $R_m$ differs from $S$ only by a part that is finite for every positive $s$, the theorem is herewith proved.

But the theorem still holds even when the inequality (3) does not hold for all $n$, but subsists only as soon as $n$ has exceeded an arbitrarily given limit $m$. On
%CONT>>


% ===== scan leaf 720, printed 697 =====
% !PRINTED 697
\weberpage{697}
%<<CONT
of this remark one may, if $n : \mu_n$ approaches a finite limiting value $\gamma$, take $\alpha$ and $\beta$ as near to this limiting value as one pleases, and so obtains the somewhat more precise formulation of the theorem:

3. \gesperrt{If the positive numbers $\mu_1, \mu_2, \mu_3, \dots$ are such that}
\[
\mathop{\mathrm{Lim}}\limits_{n=\infty} \frac{n}{\mu_n} = \gamma
\]
\gesperrt{is a finite limiting value, then the infinite series}
\[
S = \frac{1}{\mu_1^s} + \frac{1}{\mu_2^s} + \frac{1}{\mu_3^s} + \cdots
\]
\gesperrt{is convergent for every $s$ that is greater than 1, and we have}
\[
\mathop{\mathrm{Lim}}\limits_{s=1} (s-1)\,S = \gamma.
\]

Let there now be given any system $M$ of \gesperrt{infinitely many} positive numbers $\mu_n$, which we order thus
\weq{5}{\mu_1 \leqq \mu_2 \leqq \mu_3 \leqq \dots \qquad (M),}
so that in the sequence (5) a smaller term never follows a greater one, and let $T$ denote the number of terms in $M$ that are not greater than an arbitrarily given positive quantity $t$. The following theorem then holds:

4. \gesperrt{If one of the two limiting values}
\weq{6}{\mathop{\mathrm{Lim}}\limits_{n=\infty} \frac{n}{\mu_n}, \quad \mathop{\mathrm{Lim}}\limits_{t=\infty} \frac{T}{t}}
\gesperrt{is finite, then the other has the same finite value.}

Let first
\weq{7}{\mathop{\mathrm{Lim}}\limits_{n=\infty} \frac{n}{\mu_n} = \gamma}
be a finite limiting value. Then the $\mu_n$ must necessarily grow to infinity as $n$ increases without limit, and for every positive $t$ there is a value $m$ such that
\weq{8}{\mu_m \leqq t < \mu_{m+1};}
$m$ increases to infinity together with $t$, and one has $T = m$ [because of (5)]. Hence by (8)
\weq{9}{\frac{m}{\mu_m} \geqq \frac{T}{t} > \frac{m}{\mu_{m+1}}.}


% ===== scan leaf 721, printed 698 =====
% !PRINTED 698
\weberpage{698}
If now $\gamma = 0$, it follows from this, on letting $t$ and $m$ grow to infinity, that $T : t$ too has the limiting value 0. But if $\gamma$ is different from 0, then by (7)
\[
\operatorname{Lim} \frac{\mu_m}{\mu_{m+1}} = \operatorname{Lim} \frac{m}{m+1} = 1,
\]
and therefore the two limiting values in (9) approach the value $\gamma$, and accordingly there follows:
\weq{10}{\mathop{\mathrm{Lim}}\limits_{t=\infty} \frac{T}{t} = \gamma.}

If, secondly, we conversely assume the limit equation (10), then here too $\mu_n$ will have to grow to infinity with $n$; for otherwise, since the total number of all $\mu$ is infinite, $T$ would already be infinite for a finite $t$, and $\gamma$ could not be finite.

Let us now take a value $\mu_n$ which occurs in the sequence of the mutually equal
\[
\mu_{m+1}, \; \mu_{m+2}, \; \dots, \; \mu_{m+l}
\]
so that
\[
m + 1 \leqq n \leqq m + l,
\]
then $T$, when $t$ passes through the value $\mu_n$, will suddenly increase by $l$ units, and the ratio $T : t$ increases by $l : \mu_n$. But since $T : t$ is to have a finite limiting value, we must have
\weq{11}{\operatorname{Lim} \frac{l}{\mu_n} = 0.}
Now if $t = \mu_n$, then $T = m + l$, and consequently
\weq{12}{\frac{T}{t} = \frac{m+l}{\mu_n}, \quad \operatorname{Lim} \frac{m+l}{\mu_n} = \gamma.}

Further
\[
\frac{m}{\mu_n} < \frac{n}{\mu_n} \leqq \frac{m+l}{\mu_n},
\]
and consequently, by (11) and (12),
\[
\operatorname{Lim} \frac{n}{\mu_n} = \operatorname{Lim} \frac{m}{\mu_n} = \gamma.
\]
Thus equation (7) has been deduced from equation (10) \textsuperscript{1}).

By using theorem 4. we can now also give theorem 3. the form:

\begingroup
\footnotesize
\textsuperscript{1}) This theorem is proved in essentially the same way in \gesperrt{Dirichlet-Dedekind}, Supplement II. Cf. also \gesperrt{Riemann's} mathematische Werke, 2. Auflage, Nr. XXX.
\endgroup


% ===== scan leaf 722, printed 699 =====
% !PRINTED 699
\weberpage{699}
5. \gesperrt{If $T$ is the number of the quantities $\mu_n$ that are not greater than $t$, then}
\weq{13}{\mathop{\mathrm{Lim}}\limits_{s=1}\sum_{1,\infty}^{n}\frac{s-1}{\mu_n^{s-1}} = \mathop{\mathrm{Lim}}\limits_{t=\infty}\frac{T}{t},}
% [sic] The source prints exponent s-1 in the denominator. Theorem 3 and the application in section 192,(1) require exponent s; for mu_n=n the printed series diverges for 1<s<=2. Printed bytes are retained.
\gesperrt{if the limiting value of $T:t$ is finite.}

Let us assume that the quantities $\mu_1, \mu_2, \mu_3, \ldots$ satisfy not only the conditions of theorem 3., but also the further condition that, when
\weq{14}{\mu_n = \frac{n+c_n}{\gamma}}
is set, the $c_n$ do not become infinite as $n$ increases indefinitely (i. e. for every $n$, however large, they are enclosed between finite limits); then we can replace theorem 3. by the still sharper one:

6. \gesperrt{If the numbers $\mu_n$ have the property that a finite $\gamma$ can be so determined that}
\weq{15}{\gamma \mu_n - n = c_n}
\gesperrt{does not become infinite as $n$ increases indefinitely, then the function of $s$ defined for every $s>1$}
\weq{16}{S = \frac{1}{\mu_1^s} + \frac{1}{\mu_2^s} + \cdots}
\gesperrt{has the property that the difference}
\weq{17}{S - \frac{\gamma}{s-1} = C_s}
\gesperrt{approaches a finite limit as $s-1$ decreases indefinitely.}

To prove this, we set, according to (14) and (16):
\weq{18}{\sum^{n} \frac{\gamma^s}{n^s} - S = \gamma^s \sum^{n} \frac{1}{n^s}\left[1-\left(1+\frac{c_n}{n}\right)^{-s}\right].}

Now if $\varepsilon$ is a positive proper fraction, then, if we put
\[
a_n = \frac{1}{n^\varepsilon}\left[1-\left(1+\frac{c_n}{n}\right)^{-s}\right]
\]
$a_n n^{1+\varepsilon}$ does not become infinite for an infinitely increasing $n$, and consequently, by a well-known elementary theorem of the theory of series, $a_1 + a_2 + \cdots$ is an unconditionally convergent series.


% ===== scan leaf 723, printed 700 =====
% !PRINTED 700
\weberpage{700}
If we now set
\[
s = s_1 + \varepsilon,
\]
then from (18) there results:
\[
\sum \frac{\gamma^s}{n^s} - S = \gamma^s \sum \frac{a_s}{n^{s_1}},
\]
% [sic] The German source prints a_s; a_n is required by the definition on printed p. 699 and by the summation index.
and by 1. this is a function of $s_1$ that is finite and continuous for positive $s_1$. But for $s_1 = 1-\varepsilon$ we have $s=1$, and consequently
\weq{19}{\sum \frac{\gamma^s}{n^s} - S = D_s}
is finite for $s=1$, namely equal to
\[
D_1 = \sum \frac{\gamma c_n}{n(n+c_n)}.
\]
By applying simple theorems from the theory of the $\Gamma$-functions, however, the sum $\sum \dfrac{1}{n^s}$ can be expressed by a definite integral. We have
\[
\frac{1}{n^s} = \frac{1}{\Gamma(s)}\int_0^{\infty} e^{-nx} x^{s-1}\, dx,
\]
and consequently
\[
\sum \frac{1}{n^s} = \frac{1}{\Gamma(s)}\int_0^{\infty} \frac{x^{s-1}e^{-x}\,dx}{1-e^{-x}}.
\]
Further we have [by the theorem $(s-1)\Gamma(s-1) = \Gamma(s)$]:
\[
\frac{1}{s-1} = \frac{1}{\Gamma(s)}\int_0^{\infty} e^{-x} x^{s-2}\, dx,
\]
and from this
\[
\sum \frac{1}{n^s} - \frac{1}{s-1} = \frac{1}{\Gamma(s)}\int_0^{\infty} x^{s-1}e^{-x}\left(\frac{1}{1-e^{-x}} - \frac{1}{x}\right)dx.
\]
Here the right-hand side is a function continuous for all positive values of $s$, which for $s=1$ passes over into \gesperrt{Euler}'s constant
\[
\int_0^{\infty} e^{-x}\left(\frac{1}{1-e^{-x}} - \frac{1}{x}\right)dx = -\Gamma'(1) = 0{,}57721566\ldots
\]
Now by (19):
\[
D_s = \gamma^s\left(\sum \frac{1}{n^s} - \frac{1}{s-1}\right) - \left(S - \frac{\gamma_s}{s-1}\right),
\]
% [sic] The source has a subscript in the final gamma_s, unlike the leading gamma^s. Equation (19) requires gamma^s in both places; the printed subscript is preserved.
%CONT>>


% ===== scan leaf 724, printed 701 =====
% !PRINTED 701
\weberpage{701}
\pdfbookmark[2]{§ 192. Application to the determination of the class number}{bm219}
%<<CONT
and therefore
\weq{20}{\mathop{\mathrm{Lim}}\limits_{s=1}\left(S - \frac{\gamma_s}{s-1}\right) = -\gamma\,\Gamma'(1) + D_1 = C_1,}
% [sic] The source prints gamma_s, +D_1, and =C_1. With the intended gamma^s, (19) gives the limit -gamma*Gamma'(1)-D_1. Definition (17), which subtracts gamma/(s-1), instead gives C_1=-gamma*Gamma'(1)+gamma*log(gamma)-D_1. The printed subscript, sign, and identification are retained and distinguished here.
is finite, as was to be proved\footnote{The theorems on $\Gamma$-functions used here are to be found in the more detailed textbooks of the integral calculus, e.\,g. \gesperrt{Serret-Harnack}, Lehrbuch der Differential- und Integralrechnung, Bd.\ II, S.\,170\,f. (Leipzig 1885).}.

\begin{center}
\S.\,192.\\[4pt]
\gesperrt{Application to the determination of the class number.}
\end{center}

We now make application of these theorems to the theory of the field $\Omega$. By the numbers $\mu_n$ of Theorem 5. we understand the norms of all the ideals of the field $\Omega$, and hence by $T$ the number of ideals in $\Omega$ whose norm is not greater than a given positive quantity $t$, and we obtain
\weq{1}{\mathop{\mathrm{Lim}}\limits_{s=1}\sum \frac{s-1}{N(\mathfrak{a})^s} = \mathop{\mathrm{Lim}}\limits_{t=\infty}\frac{T}{t},}
and in this the sum on the left side extends over \gesperrt{all} ideals $\mathfrak{a}$ of the field. The ideals now split, according to \S.\,153, into a finite number of classes
\[
A_1, A_2, \ldots, A_h,
\]
and the great goal is the determination of this number $h$, the \gesperrt{class number}.

The ideals $\mathfrak{a}_1$ of one of these classes $A_1$ are characterised by this, that their products with one and the same integral functional $\varphi_1$, which belongs to the class $A_1^{-1}$, are \gesperrt{principal ideals}, so that therefore
\[
\varphi_1 \mathfrak{a}_1 = \alpha
\]
is an integer of the field $\Omega$.

If the norm of $\mathfrak{a}_1$ is not greater than $t$, then
\[
N_a(\alpha) \leqq N_a(\varphi_1)\, t = t_1.
\]
If $T_1$ is the number of ideals of the class $A_1$ whose norms are not greater than $t$, then $T_1$ is at the same time the number of the non-
%CONT>>


% ===== scan leaf 725, printed 702 =====
% !PRINTED 702
\weberpage{702}
%<<CONT
associated integers divisible by $\varphi_1$ whose norms are not greater than $t_1$, and by the theorem \S.\,190,\,1. we have
\weq{2}{\mathop{\mathrm{Lim}}\limits_{t_1=\infty}\frac{T_1}{t_1} = \frac{g}{N_a(\varphi_1)},}
where $g$ has the meaning given at the place referred to, and is accordingly a number \gesperrt{different from zero}, determined by the nature of the field $\Omega$. If now $T_2, t_2, \ldots, T_h, t_h$ have the corresponding meaning for the classes $A_2, \ldots, A_h$ that $T_1, t_1$ have for $A_1$, then
\[
T = T_1 + T_2 + \cdots + T_h,
\]
\[
\frac{T}{t} = \frac{T_1}{t_1}\, N_a(\varphi_1) + \frac{T_2}{t_2}\, N_a(\varphi_2) + \cdots + \frac{T_h}{t_h}\, N_a(\varphi_h),
\]
and from (1) and (2) the fundamental formula results:
\weq{3}{\mathop{\mathrm{Lim}}\limits_{s=1}\sum \frac{s-1}{N(\mathfrak{a})^s} = gh.}
% [sic] The following source citation is section 185,(11); the formula defining g is section 190,(11) on printed p. 693.
The number $g$ we may regard as known by \S.\,185,\,(11), even though its actual calculation still rests on the assumption that a fundamental system of units is known, and even though the ascertaining of such a system in higher fields has always been regarded as one of the greatest difficulties. The calculation of the class number then still depends on the determination of the limiting value on the left side of (3), which of course likewise succeeds only in special cases, but which often permits important conclusions about the nature of the class number. We wish to develop one further transformation, preparatory to the calculation, of the sum
\weq{4}{\sum \frac{1}{N(\mathfrak{a})^s} = \Phi(s)}
into an infinite product.

Let $\mathfrak{p}$ be any prime ideal of degree $f$ in the field $\Omega$, so that
\[
N(\mathfrak{p}) = p^f.
\]
Then the infinite geometric series is
\[
1 + \frac{1}{N(\mathfrak{p})^s} + \frac{1}{N(\mathfrak{p})^{2s}} + \cdots = \frac{1}{1-p^{-sf}};
\]
if one multiplies these series for all the different prime ideals $\mathfrak{p}$ with one another, then by known theorems one obtains from
%CONT>>


% ===== scan leaf 726, printed 703 =====
% !PRINTED 703
\weberpage{703}
%<<CONT
the theory of infinite series a sum of terms of the
form
\[
\left(\frac{1}{N(\mathfrak{p})^k}\,\frac{1}{N(\mathfrak{p}')^{k'}}\,\frac{1}{N(\mathfrak{p}'')^{k''}}\cdots\right)^s = \left(\frac{1}{N(\mathfrak{p}^k\mathfrak{p}'^{k'}\mathfrak{p}''^{k''}\dots)}\right)^s,
\]
all of which are contained in the form $N(\mathfrak{a})^{-s}$, and each such
term arises once and only once. Accordingly we have
\weq{5}{\varPhi(s) = \prod \frac{1}{1-N(\mathfrak{p})^{-s}}.}

If therefore $\mathfrak{p}_1, \mathfrak{p}_2, \dots, \mathfrak{p}_e$ are the prime factors of $p$ that are
distinct from one another, and $f_1, f_2, \dots, f_e$ their degrees (\S.\,143), then
we obtain
\weq{6}{\varPhi(s) = \overset{p}{\prod} \frac{1}{(1-p^{-sf_1})(1-p^{-sf_2})\dots(1-p^{-sf_e})},}
in which the infinite product $\prod$ is to be extended over all natural
 primes $p$, and by (3), (4)
\weq{7}{gh = \mathop{\mathrm{Lim}}\limits_{s=1}\,(s-1)\varPhi(s).}

From this a conclusion may be drawn that is important for
manifold applications:

If from the development
\[
1 + \frac{1}{p^s} + \frac{1}{p^{2s}} + \cdots = \frac{1}{1-p^{-s}}
\]
we form the product for all primes $p$, then we obtain:
\weq{8}{\overset{p}{\prod}\frac{1}{1-p^{-s}} = \overset{n}{\sum}\frac{1}{n^s},}
in which the sum on the right side extends over all natural
numbers $n$; and this product therefore has a finite value for all
values of $s$ that are greater than 1. The reciprocal
value, namely the product
\weq{9}{P = \overset{p}{\prod}(1-p^{-s}),}
whose factors are all less than 1, has therefore, as
long as $s>1$, a value different from zero. This
property is now preserved if, in forming the
product $P$, we take account of only a part of all the primes $p$,
because by the omission of arbitrary factors the product is only
increased, without ever exceeding unity.


% ===== scan leaf 727, printed 704 =====
% !PRINTED 704
\weberpage{704}
From this it follows that in the product (6) the partial products
\[
\prod \frac{1}{1-p^{-sf}},
\]
which extend over all primes $p$ for which $f>1$, still remain finite
for $s=1$ as well. Since on the other hand $g, h$ have definite
values different from zero, there follows from (7) the important
theorem:

I. \gesperrt{If $\mathfrak{p}$ runs through the totality of prime ideals
of the first degree of any algebraic field $\Omega$, then the product}
\weq{10}{(s-1)\prod\frac{1}{1-N(\mathfrak{p})^{-s}}}
\gesperrt{has for $s=1$ a finite limiting value different from
zero.}

This property is retained even when, in the formation of the
product, an arbitrary \gesperrt{finite} number of prime ideals of the first
degree is omitted.

The norms $N(\mathfrak{p})$ in the formula (10) are equal to natural
primes $p$. But a prime $p$ occurs in it as often as it contains
prime factors of the first degree in $\Omega$, hence at most
$n$ times. If $\Omega$ is a \gesperrt{normal field}, then every prime $p$ that
is decomposable into prime factors of the first degree does in
fact occur in it exactly $n$ times, and for the product (10) we
may also put:
\[
(s-1)\prod\frac{1}{(1-p^{-s})^n},
\]
in which the product $\prod$ extends over all primes $p$ that are
\gesperrt{decomposable into prime factors of the first degree} (\S.\,160).
% [source note] The assertion of exactly n distinct degree-one prime factors requires an unramified prime. For Q(i), 2 ramifies as a unit times (1+i)^2 and has only one distinct degree-one prime. Omitting the finitely many ramified Euler factors preserves the finite nonzero limiting-value conclusion, but the displayed products are not literally identical without that qualification.

An immediate consequence of Theorem I. is this:

\gesperrt{In every algebraic field there exist infinitely
many prime ideals of the first degree.}

\begin{center}
\rule{6em}{0.4pt}
\end{center}


% ===== scan leaf 728, printed 705 =====
\weberpage{705}% page number not printed on this page
\pdfbookmark[1]{Twenty-fourth Section. Class Number of the Cyclotomic Fields.}{bm220}
\pdfbookmark[2]{§ 193. Representation of the class number in the cyclotomic field Ωm}{bm221}
% !PRINTED none
\weberpage{none}
% [source note: this Abschnitt opener has no printed folio; sequence position corresponds to inferred page 705.]

\begin{center}
\gesperrt{Twenty-fourth Chapter.}

\gesperrt{Class number of the cyclotomic fields.}

\rule{4em}{0.4pt}
\end{center}

\begin{center}
\S.\,193.

\gesperrt{Class number representation in the cyclotomic field $\Omega_m$.}
\end{center}

The general results concerning the class number $h$ are now to
be applied to the full cyclotomic field $\Omega_m$, under the
assumption made hitherto that $m$ is a power of a prime
$q$, and that accordingly
\weq{1}{m = q^{\varkappa}, \quad \varphi(m) = q^{\varkappa-1}(q-1) = \mu}
has been set.

In the twenty-first chapter (\S.\,169, 171) we have seen
that only a single prime ideal of degree $1$ divides $q$, and
that a prime $p$ \gesperrt{belonging} to the exponent $f$, and
different from $q$, splits into $e$ distinct prime ideals of degree $f$.
Here $f$ denotes the least positive exponent for which
\weq{2}{p^f \equiv 1\ (\mathrm{mod}\ m)}
holds, and the equation
\weq{3}{\mu = ef}
determines $e$.

The function $\varPhi(s)$ determined at the end of the preceding
section by formula (6) therefore takes here the expression:
\weq{4}{\varPhi(s) = \frac{1}{1-q^{-s}}\prod\frac{1}{(1-p^{-sf})^e},}
and in it the product sign $\varPi$ extends over all primes
 $p$ that are different from $q$. Here $s$ is a variable
that is always greater than 1, and that is finally to approach
the limit 1.


% ===== scan leaf 729, printed 706 =====
% !PRINTED 706
\weberpage{706}
In order now to transform this expression for the function $\varPhi$ further
and to prepare it for computation, one has to use the theorems
on \gesperrt{Abel}ian groups and group characters
which we became acquainted with in sections 14 to 16 of this
volume.

The totality of the numbers $n$ not divisible by $q$
forms, taken with respect to the modulus $m$, an \gesperrt{Abel}ian group $\mathfrak{N}$
of degree $\mu$, whose characters $\chi(n)$ are determined in \S.\,16. There
was in this a small difference, according as $q$ is odd or
$q=2$.

1) If $q$ is odd, there is a primitive root $c$ of $m$,
and for every number $n$ there is an index $\gamma$ such that
\weq{5}{n \equiv c^{\gamma}\ (\mathrm{mod}\ m).}
If then $\Theta$ is a $\mu$th root of unity, then
\weq{6}{\chi(n) = \Theta^{\gamma},}
and one obtains all $\mu$ characters if for $\Theta$
one puts the different $\mu$th roots of unity.

If $n$ belongs to the exponent $f$, then $e$ is the greatest common
divisor of $\gamma$ and $\mu$, and $\chi(n)$ is an $f$th root of unity.

2) For $q=2$, $m>4$ (we do not take the case $m=4$ into account
here) every odd number $n$ has two indices $\alpha, \beta$, which are determined
with respect to the moduli $2, \frac{1}{2}\mu$, for which
\weq{7}{n \equiv (-1)^{\alpha}\,5^{\beta}\ (\mathrm{mod}\ m)}
holds. As characters one obtains, when $\varepsilon = \pm 1$ and $\Theta$ is set equal
to a $\frac{1}{2}\mu$th root of unity:
\weq{8}{\chi(n) = \varepsilon^{\alpha}\Theta^{\beta};}
$f$ is the least positive solution of the two congruences:
\[
\alpha f \equiv 0\ (\mathrm{mod}\ 2), \quad \beta f \equiv 0\ (\mathrm{mod}\ \tfrac{1}{2}\mu),
\]
and here too $\chi(n)$ is an $f$th root of unity.

For both cases, however, the following remark holds:

3) If $n$ belongs to the exponent $f$, then, if one lets $\chi$
run through the entire group of the $\mu$ characters, one obtains from
$\chi(n)$ every $f$th root of unity $e$ times.

For since the $\chi$ form an \gesperrt{Abel}ian group, one obtains
in the first place every $f$th root of unity that occurs at all
equally often, namely as often as the value 1, because, if $\chi_1(n)=\chi_2(n)$,
then $\chi_1\chi_2^{-1}(n) = \chi_0(n) = 1$. On the other hand one obtains
%CONT>>


% ===== scan leaf 730, printed 707 =====
% !PRINTED 707
\weberpage{707}
%<<CONT
every $f$th root of unity; for if in case 1) one takes for $\Theta$ a
primitive $\mu$th root of unity, in case 2) a primitive $\frac{1}{2}\mu$th
root of unity, and in the last case sets at the same time $\varepsilon = -1$,
then $\chi(n)$ is a primitive $f$th root of unity, and from its powers
all the $f$th roots of unity can be derived.

If therefore $x$ now denotes a variable, then, when $n$ belongs to the
exponent $f$, the product extended over all the characters $\chi$ is
\weq{9}{\overset{\chi}{\Pi} [1 - \chi(n)\,x] = (1-x^f)^e,}
and if one therefore sets $x = p^{-s}$, there results from (4):
\weq{10}{\Phi(s) = \frac{1}{1-q^{-s}} \overset{\chi}{\Pi} \overset{p}{\Pi} \frac{1}{1-\chi(p)\,p^{-s}},}
where the first product sign extends over all characters $\chi$, the second
over all primes $p$. Thus $\Phi(s)$ is a product of a finite number, $\mu$,
of factors, each of which is an infinite product.

We now propose to develop each factor of one of these infinite products in
ascending powers of $p^{-s}$. This gives
\[
\frac{1}{1-\chi(p)\,p^{-s}} = 1 + \chi(p)\,p^{-s} + \chi(p^2)\,p^{-2s} + \chi(p^3)\,p^{-3s} + \cdots
\]
The product of all the factors that one obtains from this, when $\chi$ is
held fast while $p$ runs through all the primes different from $q$, is an
aggregate of terms of the form
\[
\chi(p^k)\,\chi(p'^{k'})\,\chi(p''^{k''}) \cdots p^{-sk}\,p'^{-sk'}\,p''^{-sk''} \cdots = \chi(n)\,n^{-s},
\]
and every term of this form occurs in the product once and only once
[cf. \S.\,192, (8)]. Accordingly there results the transformation of $\Phi(s)$
into a finite product of infinite series:
\weq{11}{\Phi(s) = \frac{1}{1-q^{-s}} \overset{\chi}{\Pi} \overset{n}{\Sigma} \frac{\chi(n)}{n^s},}
in which the sum extends over all numbers $n$ not divisible by $q$,
and the product over all $\mu$ characters $\chi$.

This expression is suited to determining the limiting value of
$(s-1)\,\Phi(s)$ for $s=1$.

Let us first consider the sum corresponding to the principal character
$\chi = 1$,
\[
\overset{n}{\Sigma} \frac{1}{n^s}.
\]


% ===== scan leaf 731, printed 708 =====
% !PRINTED 708
\weberpage{708}
\pdfbookmark[2]{§ 194. Determination of the sums X}{bm222}
One obtains the numbers $n$ if one removes from the totality of all natural
numbers $k$ the multiples of $q$, that is, the totality of the numbers $qk$.
Accordingly
\[
\Sigma \frac{1}{n^s} = \Sigma \frac{1}{k^s} - \Sigma \frac{1}{q^s k^s} = (1-q^{-s}) \Sigma \frac{1}{k^s}.
\]
If, however, in Theorem 3., \S.\,191, one puts for the $\mu_n$ the sequence of
the natural numbers $k$, it follows that:
\[
\underset{s=1}{\mathrm{Lim}} (s-1) \Sigma \frac{1}{k^s} = 1,
\]
and we therefore obtain
\weq{12}{\underset{s=1}{\mathrm{Lim}} \frac{s-1}{1-q^{-s}} \Sigma \frac{1}{n^s} = 1.}
The other factors of the product $\Phi(s)$, however, are continuous functions of
$s$ by Theorem 1., \S.\,191, if the infinite series are ordered according to
increasing values of $n$.

For the sum
\weq{13}{\Sigma \chi(n)}
ordered according to increasing values of $n$ is indeed not convergent, but it
cannot for all that exceed a finite limit in absolute value. For as often as $n$
runs through a complete system of residues with respect to the modulus $m$,
there is added to the sum (13) a contribution which by \S.\,11, 6. has the value
0. Accordingly, if $\chi$ is not the principal character,
\weq{14}{\underset{s=1}{\mathrm{Lim}} \overset{n}{\Sigma} \frac{\chi(n)}{n^s} = \overset{n}{\Sigma} \frac{\chi(n)}{n},}
and accordingly we now obtain for the class number $h$ in the cyclotomic field
$\Omega_m$ the expression
\weq{15}{g\,h = \overset{\chi}{\Pi} \overset{n}{\Sigma} \frac{\chi(n)}{n},}
in which the product sign $\Pi$ refers only to those characters $\chi$ which are
\gesperrt{different from the principal character}.

\begin{center}
\S.\,194.

\gesperrt{Determination of the Sums X.}
\end{center}

The infinite series
\weq{1}{X = \overset{n}{\Sigma} \frac{\chi(n)}{n},}
%CONT>>


% ===== scan leaf 732, printed 709 =====
% !PRINTED 709
\weberpage{709}
%<<CONT
on which the determination of the class number still depends, can be
represented by finite expressions. For we have
\[
\frac{1}{n} = \int_0^1 x^{n-1}\,dx,
\]
and accordingly
\[
X = \int_0^1 \Sigma \chi(n)\,x^{n-1}\,dx.
\]

Now $\chi(n)$ remains unchanged when $n$ increases by a multiple of $m$.
But if by $t$ we understand the least positive residue of $n$ and set
\[
n = t + ml,
\]
then $\chi(n) = \chi(t)$, and it follows that:
\[
X = \int_0^1 \overset{t}{\Sigma} \chi(t)\,x^{t-1} \underset{0,\,\infty}{\overset{l}{\Sigma}} x^{ml}\,dx.
\]

If we now set
\weq{2}{f(x) = \overset{t}{\Sigma} \chi(t)\,x^t,}
then $f(x)$ is an integral function of $x$ of degree $m-1$, and at the same
time, in consequence of the relation $\Sigma \chi(t) = 0$:
\weq{3}{f(0) = 0,\quad f(1) = 0,}
and therefore $f(x)$ is divisible by $x(1-x)$. For $X$ there then results,
by the summation formula for the geometric series $\Sigma x^{ml}$:
\weq{4}{X = \int_0^1 \frac{f(x)\,dx}{x(1-x^m)}.}

To find such an integral, the integral calculus prescribes that the
rational fraction under the integral sign be decomposed into partial fractions.
This decomposition into partial fractions gives, however, in consequence of (3), when
\[
r = e^{\frac{2\pi i}{m}}
\]
is set:
\[
\frac{f(x)}{x(1-x^m)} = -\frac{1}{m} \sum\limits_{1,\,m-1}^{s} \frac{f(r^s)}{x-r^s}
\]
[Bd. I, \S.\,29, (11)], and now we have further, by known
%CONT>>


% ===== scan leaf 733, printed 710 =====
% !PRINTED 710
\weberpage{710}
%<<CONT
elementary theorems of the integral calculus, \gesperrt{if $\alpha$ denotes any angle
between 0 and $2\pi$,}
\[
\int_0^1 \frac{dx}{x-e^{i\alpha}} = \tfrac{1}{2}\log\left(4\sin^2\frac{\alpha}{2}\right) + i\,\frac{\pi-\alpha}{2}.
\]
From this there results, by (4):
\weq{5}{X = -\frac{1}{m}\underset{1,\,m-1}{\overset{s}{\Sigma}} f(r^s)\left[\tfrac{1}{2}\log 4\left(\sin\frac{s\pi}{m}\right)^2 - \frac{i\pi(m-2s)}{2m}\right].}
In this, by (2),
\weq{6}{f(r^s) = \overset{t}{\Sigma}\chi(t)\,r^{ts},}
and from this there results [because of (3)]:
\weq{7}{\underset{1,\,m-1}{\overset{s}{\Sigma}} f(r^s) = \underset{0,\,m-1}{\overset{s}{\Sigma}} f(r^s) = 0,}
since
\[
\underset{0,\,m-1}{\overset{s}{\Sigma}} r^{st} = 0
\]
holds. Accordingly formula (5) simplifies somewhat further and gives
\weq{8}{X = -\frac{1}{m}\underset{1,\,m-1}{\overset{s}{\Sigma}} f(r^s)\left[\tfrac{1}{2}\log\left(\sin\frac{s\pi}{m}\right)^2 + \frac{i\pi s}{m}\right],}
for which one may also, if one replaces $s$ by $m-s$ and again uses (7), put
\weq{9}{X = -\frac{1}{m}\underset{1,\,m-1}{\overset{s}{\Sigma}} f(r^{-s})\left[\tfrac{1}{2}\log\left(\sin\frac{s\pi}{m}\right)^2 - \frac{i\pi s}{m}\right]}
in its place.

But from (6) it follows that
\[
f(r^{-s}) = \overset{t}{\Sigma}\chi(t)\,r^{-ts},
\]
and the sum taken over $t$ does not change when $t$ is replaced by $m-t$.
From this it follows:
\[
f(r^{-s}) = \chi(-1)\,f(r^s).
\]
The factor $\chi(-1)$ has the value $+1$ or $-1$, since its square $=+1$.
We therefore now distinguish two kinds of characters $\chi_1(n)$,
$\chi_2(n)$, so that
\weq{10}{\chi_1(-1) = 1,\quad \chi_2(-1) = -1}
%CONT>>


% ===== scan leaf 734, printed 711 =====
% !PRINTED 711
\weberpage{711}
\pdfbookmark[2]{§ 195. On the class number in the real field contained in Ωm}{bm223}
%<<CONT
is, and accordingly the functions $f$ are also to be distinguished
into $f_1$ and $f_2$, so that

\weq{11}{f_1(r^{-s}) = f_1(r^s), \quad f_2(r^{-s}) = -\,f_2(r^s)}
holds. Likewise we distinguish the expressions (8), (9) as $X_1$
and $X_2$, and if both expressions are added, there results, by
(11):
\weq{12}{\begin{aligned}
X_1 &= -\frac{1}{2m}\sum_{1,\,m-1}^{s} f_1(r^s)\,\log\left(\sin\frac{s\pi}{m}\right)^{2}, \\
X_2 &= -\frac{i\pi}{m^2}\sum_{1,\,m-1}^{s} s\,f_2(r^s).
\end{aligned}}

\begin{center}
\S.\,195.\\[4pt]
\gesperrt{On the class number in the real field contained in
$\Omega_m$.}
\end{center}

Our general principles can also be applied to determine the class
numbers in the divisors of the field $\Omega_m$.

In particular the formulae (6), (7), \S.\,192 are applicable if we
know the degrees of the prime ideals in such a divisor. The question
of the ratio of the class numbers in the divisors of a field
$\Omega_m$ to the class number in $\Omega_m$ itself is of great
interest, and it shall be treated here for the case that the field
concerned is the real field $H_m$ contained in $\Omega_m$, which we
have indeed already investigated in \S.\,176 f.\footnote{Cf.
\gesperrt{Kummer}, \glqq Bestimmung der Anzahl u.\,s.\,f.\grqq{}.
\gesperrt{Crelle's Journ.}, Bd.\,40 (1849).} It was found there that
$q$ in the field $H_m$ is the $\tfrac12\mu$th power of a prime of
the first degree, that the other primes $p$ split into $e_1$ prime
factors of degree $f_1$, so that

\weq{1}{f_1 e_1 = \tfrac{1}{2}\mu,}
and that two kinds of primes $p_1$, $p_2$ are to be distinguished,
so that for the primes of the first kind $f_1 = \tfrac12 f$,
$e_1 = e$, and for those of the second kind $f_1 = f$,
$e_1 = \tfrac12 e$.

For the field $H_m$ we accordingly obtain from \S.\,192, (6), (7)
the class number $h_1$ by the formula:
%CONT>>


% ===== scan leaf 735, printed 712 =====
% !PRINTED 712
\weberpage{712}
%<<CONT
\weq{2}{g_1 h_1 = \mathop{\mathrm{Lim}}\limits_{s=1} (s-1)\,\Phi_1(s),}
\weq{3}{\Phi_1(s) = \frac{1}{1-q^{-s}}\;\prod^{p_1}\frac{1}{(1-p_1^{-1/2\,sf})^{e}}\;\prod^{p_2}\frac{1}{(1-p_2^{-sf})^{1/2\,e}},}
in which the two product symbols refer to all the primes $p_1$ and
$p_2$. With the use of the symbols $f_1$, $e_1$ explained above one
can also write for this:
\weq{4}{\Phi_1(s) = \frac{1}{1-q^{-s}}\;\prod^{p}\frac{1}{(1-p^{-sf_1})^{e_1}},}
and extend the product symbol over \gesperrt{all} primes $p$.

The meaning of $g_1$ follows from \S.\,190, (11).

Now from the group $C$ of the characters $\chi$ we will single out a
divisor $C_1$, which shall consist of all those characters $\chi_1$
for which
\weq{5}{\chi_1(-1) = +1}
holds. It is clear at the outset that these characters $\chi_1$ form
a group $C_1$. But this group is not identical with $C$, for there
certainly is a character $\chi_0$ for which $\chi_0(-1) = -1$
holds. To obtain such a one we need only put, in \S.\,193,
(6), a primitive $\mu$th root of unity for $\Theta$,
or set $\varepsilon = -1$ in \S.\,193, (8). Then the
totality of the characters $\chi_0\chi_1$, none of which belong to
$C_1$, forms a subsidiary group $C_2$, whose elements may be
denoted by $\chi_2$. Every character $\chi$ is then
contained either in $C_1$ or in $C_2$; for if $\chi$ does not
occur in $C_1$, then $\chi_0^{-1}\chi$ occurs in it, and
consequently $\chi$ in $C_2$. Therefore $C_1$ is a divisor of $C$
of index 2.

The characters $\chi_1$ can be formed in the following way:

1) If $m$ is odd, we set, as in \S.\,193,
\[ n \equiv c^{\gamma} \pmod{m}, \]
and in
\weq{6}{\chi_1(n) = \Theta^{\gamma}}
we let $\Theta$ run through the totality of the $\tfrac12\mu$th
roots of unity.

2) If $m$ is even, we set
\[ n \equiv (-1)^{\alpha}5^{\beta}, \]
and obtain
\weq{7}{\chi_1(n) = \Theta^{\beta},}
in which $\Theta$ again runs through all the $\tfrac12\mu$th
roots of unity.


% ===== scan leaf 736, printed 713 =====
% !PRINTED 713
\weberpage{713}
Both follow easily from the fact that for $n = -1$ in the first
case $\gamma = \tfrac12\mu$, and in the second $\alpha=1$,
$\beta=0$, so that in both cases $\chi_1(-1)=1$, and from the fact
that each of the formulae (6), (7) represents exactly
$\tfrac12\mu$ different characters.

What now matters essentially is the following:

\begin{quote}
\gesperrt{If $p$ is any prime different from $q$, then in the form
$\chi_1(p)$, as $\chi_1$ runs through the group $C_1$, one obtains
every $f_1$th root of unity, and each equally often, that is,
$e_1$ times.}
\end{quote}

That all the $\chi_1(p)$ are roots of unity of degree $f_1$ is
clear, for it is always the case (cf. \S.\,177) that
$p^{f_1}\equiv\pm1\pmod{m}$, and therefore
\[ [\chi_1(p)]^{f_1} = \chi_1(\pm1) = 1. \]
But if a primitive $f_1$th root of unity occurs among the
$\chi_1(p)$, then all the other $f_1$th roots of unity occur among
them as well, and equally often, namely as often as the root of
unity 1.

This follows immediately from the group nature of the $\chi_1$. It
therefore remains only to show that \gesperrt{a primitive} $f_1$th root of
unity occurs among the $\chi_1(p)$.

But this possibility follows very simply from the determination of
$\chi_1$. Suppose first that, for \gesperrt{odd $m$, $\Theta$ in (6) is a
primitive root of unity of degree} $\tfrac12\mu$, and that $\gamma$
is the index of $p$; then a power of $\Theta^{\gamma}$,
$\Theta^{\gamma\lambda}$, is $=1$ then and only then, when
\[ 2\gamma\lambda \equiv 0 \pmod{\mu}. \]
Now $\mu=ef$, and $e$ is the greatest common divisor of $\gamma$
and $\mu$; and therefore this condition becomes
$2\lambda\equiv0\pmod{f}$. But by \S.\,177, 3. $f_1=f$
or $=\tfrac12 f$, according as $f$ is odd or even, and
consequently $\lambda\equiv0\pmod{f_1}$ must hold. Thus
$\Theta^{\gamma}$ is a primitive $f_1$th root of unity.

But if $m$ is \gesperrt{even}, then in (7) we likewise take for $\Theta$
a primitive $\tfrac12\mu$th root of unity and put
\[ p \equiv (-1)^{\alpha}5^{\beta} \pmod{m}, \]
so that $f$ is the least positive number for which the congruences
\[ \alpha f \equiv 0 \pmod{2}, \quad \beta f \equiv 0 \pmod{\tfrac12\mu} \]
are satisfied. Now $\Theta^{\beta\lambda}$ is $=1$ only when
$\beta\lambda\equiv0\pmod{\tfrac12\mu}$. Unless $\lambda=1$
(which forces $\beta=0$), the least possible value of $\lambda$ is
even and hence equals $f$.


% ===== scan leaf 737, printed 714 =====
% !PRINTED 714
\weberpage{714}
But if $\lambda=1$, then $\beta=0$ must hold, and two further
cases are still possible:
\[
\begin{aligned}
&\text{either}\quad f=1,\ \alpha=0,\ p\equiv\phantom{-}1\pmod{m},\ \lambda=f\\
&\text{or}\qquad\quad\, f=2,\ \alpha=1,\ p\equiv -1\pmod{m},\ \lambda=\tfrac12 f.
\end{aligned}
\]
Only in this last case does $p$ belong to the first kind, and it
follows therefore that here too $\Theta^{\beta}$ is in general a
primitive $f_1$th root of unity. With this the theorem is then
proved.

This theorem permits us the application of the formula \S.\,193, (9),
which here gives:
\weq{8}{(1-p^{-sf_1})^{e_1} = \prod^{\chi_1}\bigl[1-\chi_1(p)p^{-s}\bigr],}
in which the product extends over all characters $\chi_1$ of the
group $C_1$, and now the function $\Phi_1(s)$ can be further
developed just as $\Phi(s)$ was in \S.\,193. One obtains, in the
same way as there, the formula (11):
\weq{9}{\Phi_1(s) = \frac{1}{1-q^{-s}}\;\prod^{\chi_1}\sum^{n}\frac{\chi_1(n)}{n^s},}
and when one passes to the limit $s=1$,
\weq{10}{g_1 h_1 = \prod^{\chi_1}\sum^{n}\frac{\chi_1(n)}{n},}
in which, however, the product is now to be extended over all
characters of the group $C_1$ \gesperrt{with the exception of the principal
character}. This product is therefore a part of the product by which
the number $gh$ is expressed in \S.\,193.

The numbers $g$, $g_1$ follow from the general expression given
in \S.\,190, (11):
\weq{11}{\frac{2^{\nu}\pi^{n-\nu}L}{w\sqrt{\pm\Delta}}.}
Since, as was shown in \S.\,178, the units in $\Omega_m$ and
$H_m$, apart from the roots of unity, are the same, a fundamental
system of units in $H_m$ is at the same time a fundamental system
in $\Omega_m$.

If therefore $\varepsilon_1,\varepsilon_2,\dots,\varepsilon_{\nu-1}$
denotes a fundamental system of real units, and if
$\varepsilon_{i,1},\varepsilon_{i,2},\dots,\varepsilon_{i,\nu}$ are
the conjugate units to $\varepsilon_i$, then the conjugate
logarithms in the field $H_m$ are:
\[ \log|\varepsilon_{i,1}|,\ \log|\varepsilon_{i,2}|,\ \dots,\ \log|\varepsilon_{i,\nu}|, \]
whereas in the field $\Omega_m$, in which only conjugate imaginary
pairs occur (\S.\,185),
\[ 2\log|\varepsilon_{i,1}|,\ 2\log|\varepsilon_{i,2}|,\ \dots,\ 2\log|\varepsilon_{i,\nu}|. \]


% ===== scan leaf 738, printed 715 =====
% !PRINTED 715
\weberpage{715}

If we therefore denote by
\[
(12)\quad E = \pm
\begin{vmatrix}
\log|\varepsilon_{1,1}|, & \log|\varepsilon_{1,2}|, & \ldots, & \log|\varepsilon_{1,\nu-1}| \\
\log|\varepsilon_{2,1}|, & \log|\varepsilon_{2,2}|, & \ldots, & \log|\varepsilon_{2,\nu-1}| \\
\cdots & \cdots & \cdots & \cdots \\
\log|\varepsilon_{\nu-1,1}|, & \log|\varepsilon_{\nu-1,2}|, & \ldots, & \log|\varepsilon_{\nu-1,\nu-1}|
\end{vmatrix}
\]
the regulator of the field $H_m$ (\S.\,187), then in applying the expression (11) we have to put, in the field $\Omega_m$,
\[
L = 2^{\nu-1} E,
\]
and in the field $H_m$
\[
L = E.
\]
Further we have
\[
\begin{aligned}
\text{in } \Omega_m: & \quad n = \mu, \quad \nu = \tfrac{1}{2}\mu, \\
\text{in } H_m: & \quad n = \tfrac{1}{2}\mu, \quad \nu = \tfrac{1}{2}\mu.
\end{aligned}
\]
In the field $H_m$ only the two roots of unity $\pm 1$ are contained. In $\Omega_m$ the powers of $r$ taken with positive and with negative sign are contained, but no other roots of unity (\S.\,175), and in the counting it must further be observed that for even $m$ the number $-1$ is also contained among the powers of $r$, while for odd $m$ it is not. Hence we have
\[
\begin{aligned}
\text{in } \Omega_m: & \quad w = 2m, \quad m \text{ odd} \\
& \quad = m, \quad m \text{ even,} \\
\text{in } H_m: & \quad w = 2.
\end{aligned}
\]
Finally we have also, if $\varDelta$, $\varDelta_1$ are the fundamental numbers of the fields $\Omega_m$, $H_m$ (\S.\,176):
\[
\begin{aligned}
\pm \varDelta & = q\varDelta_1^2 \quad m \text{ odd} \\
& = 4\varDelta_1^2 \quad m \text{ even,}
\end{aligned}
\]
and it follows, if for simplicity we put $\nu$ for $\tfrac{1}{2}\mu$:
\[
\begin{aligned}
g & = \frac{2^{2(\nu-1)}\pi^\nu E}{m\sqrt{q}\,\varDelta_1} \quad m \text{ odd} \\
g & = \frac{2^{2(\nu-1)}\pi^\nu E}{m\varDelta_1} \quad m \text{ even,}
\end{aligned}
\]
\weq{13}{g_1 = \frac{2^{\nu-1}E}{\sqrt{\varDelta_1}},}
from which
\weq{14}{\begin{aligned}
g & = \frac{2^{\nu-1}\pi^\nu g_1}{m\sqrt{q\varDelta_1}} \quad m \text{ odd} \\
& = \frac{2^{\nu-1}\pi^\nu g_1}{m\sqrt{\varDelta_1}} \quad m \text{ even.}
\end{aligned}}


% ===== scan leaf 739, printed 716 =====
% !PRINTED 716
\weberpage{716}

If we combine formula (10) with formula \S.\,193, (15) and employ the notation introduced in \S.\,194
\[
X_1 = \sum{}^n \frac{\chi_1(n)}{n}, \quad X_2 = \sum{}^n \frac{\chi_2(n)}{n}
\]
then there results
\[
g\,h = g_1 h_1 \prod X_2.
\]
If we replace $g$ by its value (14), then for the class number it follows that
\weq{15}{h = k h_1,}
where
\weq{16}{
\begin{aligned}
k & = \frac{m\sqrt{q\varDelta_1}}{2^{\nu-1}\pi^\nu} \prod X_2 \quad \text{for odd } m \\
& = \frac{m\sqrt{\varDelta_1}}{2^{\nu-1}\pi^\nu} \prod X_2 \quad \text{for even } m,
\end{aligned}
}
and in which [by (10), (13)]:
\weq{17}{h_1 = \frac{\sqrt{\varDelta_1}}{2^{\nu-1}E} \prod X_1}
is the class number of the field $H_m$. In the latter expression the product extends over all characters $\chi_1$ of the group $C_1$ \gesperrt{with the exception of the principal character}.

That $k$ is an integer, and hence $h$ a multiple of $h_1$, will soon be shown. The numbers $k$ and $h_1$ are called the \gesperrt{first and second factor of the class number}.

In the case in which
\[
m = 2^\varkappa
\]
is a power of 2, we have by \S.\,176, (7):
\[
\varDelta_1 = 2^{(\varkappa-1)2^{\varkappa-2}-1},
\]
and if we substitute this value, then in this case we can write for the foregoing formulae:
\weq{18}{
\begin{aligned}
k & = \frac{\sqrt{2}\,2^{2^{\varkappa-3}(\varkappa-3)}\,2^\varkappa}{\pi^\nu} \prod X_2 \\
h_1 & = \frac{\sqrt{2}\,2^{2^{\varkappa-3}(\varkappa-3)}}{E} \prod X_1.
\end{aligned}
}


% ===== scan leaf 740, printed 717 =====
% !PRINTED 717
\weberpage{717}
\pdfbookmark[2]{§ 196. Class number in the field of the eighth roots of unity}{bm224}

\begin{center}
\S.\,196.

\gesperrt{Class number in the field of the eighth roots of unity.}
\end{center}

We wish to treat first the case $m = 8$, $\mu = 4$, $\nu = 2$, partly in order to illustrate the results obtained so far, and partly in order to gain a basis for the complete induction to be applied later.

Here, besides the principal character, we have only three characters, namely, if
\[
n \equiv (-1)^\alpha 5^\beta \pmod{8}
\]
then
\[
\begin{aligned}
\chi_1(n) & = (-1)^\beta \\
\chi_2(n) & = (-1)^\alpha \\
\chi_3(n) & = (-1)^{\alpha+\beta}.
\end{aligned}
\]
For $n = -1$ we have $\alpha = 1$, $\beta = 0$, and consequently $\chi_1$ belongs to the first kind, $\chi_2$, $\chi_3$ to the second [\S.\,195, (5)]; and for the four numbers $n = 1, 3, 5, 7$ the following values of these characters result:
\[
\begin{aligned}
\chi_1: & \quad +1, \; -1, \; -1, \; +1 \\
\chi_2: & \quad +1, \; -1, \; +1, \; -1 \\
\chi_3: & \quad +1, \; +1, \; -1, \; -1.
\end{aligned}
\]
From this one obtains the three functions $f_1, f_2, f_3$ [\S.\,194, (2)]:
\[
\begin{aligned}
f_1(x) & = x - x^3 - x^5 + x^7 = x(1-x^2)(1-x^4) \\
f_2(x) & = x - x^3 + x^5 - x^7 = x(1-x^2)(1+x^4) \\
f_3(x) & = x + x^3 - x^5 - x^7 = x(1+x^2)(1-x^4).
\end{aligned}
\]
For $r$ we may set:
\[
r = \frac{1+i}{\sqrt{2}}, \quad r^2 = i, \quad r^3 = -\frac{1-i}{\sqrt{2}}, \quad r^4 = -1,
\]
and accordingly there results:
\[
\begin{aligned}
f_1(r) & = 2\sqrt{2}, & f_2(r) & = 0, & f_3(r) & = 2\sqrt{2}\,i, \\
f_1(r^2) & = 0, & f_2(r^2) & = 4i, & f_3(r^2) & = 0, \\
f_1(r^3) & = -2\sqrt{2}, & f_2(r^3) & = 0, & f_3(r^3) & = 2\sqrt{2}\,i.
\end{aligned}
\]
For $x = r^4$ all three functions $f(x)$ vanish, and for $x = r^5, r^6, r^7$ the function $f_1(x)$ takes the same values, while $f_2(x)$ and $f_3(x)$ take the opposite values, as for $x = r^3, r^2, r$.


% ===== scan leaf 741, printed 718 =====
% !PRINTED 718
\weberpage{718}

Accordingly there results from \S.\,194, (12), having regard to the trigonometric formula
\[
\sin\frac{3\pi}{8} = \cos\frac{\pi}{8}
\]
\[
X_1 = -\frac{1}{\sqrt{2}} \log \operatorname{tg}\frac{\pi}{8}, \quad X_2 = \frac{-\pi}{4}, \quad X_3 = \frac{-\pi}{2\sqrt{2}}.
\]
Now
\[
\tau = \frac{r-1}{r^2(r+1)} = \operatorname{tg}\frac{\pi}{8} = \sqrt{2} - 1
\]
is a unit of the field $H_8$, which here is nothing other than the quadratic field arising from $\sqrt{2}$, and the units in it are therefore obtained from the solutions of the \gesperrt{Pell} equation
\[
u^2 - 2v^2 = \pm 1
\]
in the form $u + v\sqrt{2}$.

But the least positive solution of this \gesperrt{Pell} equation is $u = 1$, $v = 1$, and consequently one finds, by Bd.\,I, \S.\,128, all units in $H_8$ in the form $\pm(1+\sqrt{2})^a$, in which $a$ is a positive or negative integer exponent; $\tau$ results from this for $a = -1$, and hence all units are contained in the form
\[
\pm \tau^a
\]
as well. Thus $\tau$ is a \gesperrt{fundamental unit}, and the determinant $E$ occurring in the formulae of the preceding section [\S.\,195, (12)] here becomes (since $\tau$ is a proper fraction):
\[
E = -\log \tau.
\]
One thus obtains from \S.\,195, (15), (18):
\[
h_1 = 1, \quad k = 1, \quad h = 1.
\]
The field $\Omega_8$ is therefore a \gesperrt{one-class field}, i.e. one in which the decomposition of the integers into prime factors is possible according to the same laws as in the field of rational numbers.

At the same time it has emerged that all units in the field $H_8$ are representable in the form $\pm \tau^a$. But the conjugate values of the unit $\tau$ are
\[
\tau_1 = \operatorname{tg}\frac{\pi}{8}, \quad \tau_3 = -\operatorname{tg}\frac{3\pi}{8} = -\tau_1^{-1},
\]
and so have different signs. From this we conclude, for this case, the theorem:
%CONT>>


% ===== scan leaf 742, printed 719 =====
% !PRINTED 719
\weberpage{719}
\pdfbookmark[2]{§ 197. Recursive computation of the class number in the field Ωm when m is a power of 2}{bm225}
%<<CONT

\gesperrt{A unit in $H_8$ which is of one and the same sign as its conjugates is,
apart from the sign, the square of a unit.}

\begin{center}
\S.\,197.

\gesperrt{Recurrent calculation of the class number in the field $\Omega_m$,
when $m$ is a power of 2.}
\end{center}

For the general determination of the class number in the field $\Omega_m$,
under the assumption that $m$ is some power of 2 and greater than 8, we
apply a recurrent procedure. Let then
\weq{1}{m = 2^\varkappa, \quad \mu = 2^{\varkappa-1}, \quad
\nu = 2^{\varkappa-2}, \quad \nu' = 2^{\varkappa-3}.}
Beside the field $\Omega_m$ we consider the field $\Omega_\mu$, which is a
divisor of $\Omega_m$, and which we shall here also denote by $\Omega_m'$.
Let $h'$ be the class number in the field $\Omega_m'$, that is, the number
which results from $h$ when $\varkappa$ is changed into $\varkappa - 1$;
and we set
\weq{2}{h = h' H:}
If we assume $h'$ to be already known, then everything turns only on the
calculation of $H$, which, as will be shown, is an integer. Decomposing
$h$ and $h'$ as in \S.\,195, we set
\[
\begin{aligned}
(3)\quad h &= k\,h_1, & \qquad h' &= k'\,h_1', \\
(4)\quad k &= k'A, & \qquad h_1 &= h_1'B,
\end{aligned}
\qquad H = AB,
\]
in which $k'$ arises from $k$, and $h_1'$ from $h_1$, by the interchange of
$\varkappa$ with $\varkappa - 1$. In general we shall now indicate by
accents on the letters that $\varkappa - 1$ is to be put in place of
$\varkappa$.

According to \S.\,176, (7) we then have
\[
\begin{aligned}
\Delta_1 &= 2^{(\varkappa-1)2^{\varkappa-2}-1}, \\
\Delta_1' &= 2^{(\varkappa-2)2^{\varkappa-3}-1},
\end{aligned}
\]
and consequently
\weq{5}{\Delta_1 = 2^{\varkappa\,2^{\varkappa-3}} \Delta_1'.}
If $E$ has the meaning of \S.\,195, (12), and if $E'$ arises from $E$ when
$\varkappa$ is replaced by $\varkappa - 1$, we further set
\weq{6}{E = E'D,}
%CONT>>


% ===== scan leaf 743, printed 720 =====
% !PRINTED 720
\weberpage{720}

%<<CONT
whereby $D$ is defined as a positive number different from zero.
As regards the sums
\[
X = \overset{n}{\Sigma}\, \frac{\chi(n)}{n}
\]
it is to be recalled that $\chi(n)$ [by \S.\,193, (8)] is defined by
\weq{7}{n \equiv (-1)^\alpha 5^\beta \pmod{m}, \qquad
\chi(n) = (\pm 1)^\alpha \Theta^\beta,}
in which $\Theta$ denotes a $\nu$th root of unity.

But among the $\nu$th roots of unity the $\nu'$th roots of unity are
also contained, and among the characters $\chi$ also the characters for
the field $\Omega_m'$. We accordingly distinguish \gesperrt{primitive and
imprimitive characters of the field $\Omega_m$}, in that we designate as
primitive \gesperrt{those characters $\chi$ peculiar to the field $\Omega_m$,
i.\,e. those in which $\Theta$ is a primitive $\nu$th root of unity}.

Among the sums $X$ there occur also all the sums belonging to the field
$\Omega_m'$:
\[
X' = \overset{n}{\Sigma}\, \frac{\chi'(n)}{n}.
\]

If one now substitutes in (4) for $k, k'$ and $h_1, h_1'$ the
expressions formed in \S.\,195, (16), (17), and then cancels out the
sums $X'$, there results:
\weq{8}{\begin{aligned}
\pi^{\nu'} A &= 2 \cdot 2^{2^{\varkappa-4}(\varkappa-2)} \prod X_2 \\
D B &= 2^{2^{\varkappa-4}(\varkappa-2)} \prod X_1,
\end{aligned}}
\gesperrt{in which, however, the products $\prod$ now extend only to the sums
$X_1, X_2$ formed with the primitive characters $\chi_1, \chi_2$.}

In the expressions for the $X$ given in \S.\,194, (12):
\weq{9}{\begin{aligned}
X_1 &= -\frac{1}{2m} \sum_{1,\,m-1}^{s} f_1(r^s)
\log\left(\sin \frac{s\pi}{m}\right)^2 \\
X_2 &= -\frac{i\pi}{m^2} \sum_{1,\,m-1}^{s} s\, f_2(r^s)
\end{aligned}}
considerable simplifications now occur, under the assumption that $m$ is
a power of 2 and that $\chi_1, \chi_2$ are primitive characters.


% ===== scan leaf 744, printed 721 =====
% !PRINTED 721
\weberpage{721}

The number $1 + \mu$ has the indices $\alpha = 0$, $\beta = \nu'$, and
consequently, by (7),
\[
\chi(1+\mu) = \Theta^{\nu'} = -1,
\]
if $\chi$ is a primitive character.

Further, for an odd $n$:
\[
n(1+\mu) \equiv n + \mu \pmod{m},
\]
and consequently
\[
\chi(n+\mu) = -\chi(n).
\]
If we now consider the function $f(r^s)$:
\weq{10}{f(r^s) = \overset{n}{\Sigma} \chi(n)\, r^{sn},}
in which $n$ runs through the odd numbers of a complete system of
residues for the modulus $m$, then, since in (10) we may replace $n$ by
$n+\mu$ under the summation sign, we obtain:
\[
f(r^s) = -r^{\mu s} \sum \chi(n)\, r^{sn},
\]
and since $r^\mu = -1$:
\[
f(r^s) = -(-1)^s f(r^s).
\]
Thus
\weq{11}{f(r^s) = 0, \quad \text{\gesperrt{if $s$ is even,}}}
\gesperrt{and accordingly we may, in the formulae (9), restrict the summation
letter $s$ to the odd numbers that are less than $m$.}

But by (10) we have:
\[
f(r^s) = \chi(s)^{-1} \overset{n}{\Sigma} \chi(ns)\, r^{sn},
\]
and if one replaces $sn$ by $n$, which is permitted when $s$ is odd:
\weq{12}{f(r^s) = \chi(s)^{-1} f(r), \quad \text{\gesperrt{if $s$ is odd.}}}
The function $f(r)$ has, when $\alpha, \beta$ are the indices of $n$,
the expression
\weq{13}{f(r) = \overset{n}{\Sigma} (\pm 1)^\alpha \Theta^\beta r^n,}
and is therefore identical in meaning with the cyclotomic resolvent
considered in \S.\,17 of this volume
\[
(\pm 1, \Theta, r)
\]
for which the relation was set up in \S.\,17, (17):
\weq{14}{(\pm 1, \Theta, r)(\pm 1, \Theta^{-1}, r) = \pm m.}
Here the upper signs hold throughout for the sums $X_1$, the lower for
the sums $X_2$ [\S.\,195, (7)].

We now divide the sequence of the numbers $s$ occurring in (9) into
four parts. Let $t$ be a symbol which
%CONT>>


% ===== scan leaf 745, printed 722 =====
% !PRINTED 722
\weberpage{722}
\pdfbookmark[2]{§ 198. The class number factor A}{bm226}

%<<CONT
runs through the sequence of the positive odd numbers from 1 to $\nu - 1$;
then the odd $s$ runs through the four sequences of numbers:
\[
t, \quad t+\mu, \quad m-t, \quad m-t-\mu,
\]
and we have:
\[
\begin{aligned}
\chi(t+\mu) &= -\chi(t), \\
\chi_1(m-t) &= \chi_1(t), \qquad \chi_1(m-\mu-t) = -\chi_1(t), \\
\chi_2(m-t) &= -\chi_2(t), \qquad \chi_2(m-\mu-t) = \chi_2(t).
\end{aligned}
\]
By (12) there results from this:
\[
\begin{aligned}
f(r^{t+\mu}) &= -f(r^t), \\
f_1(r^{m-t}) &= f_1(r^t), \qquad f_1(r^{m-\mu-t}) = -f_1(r^t), \\
f_2(r^{m-t}) &= -f_2(r^t), \qquad f_2(r^{m-\mu-t}) = f_2(r^t).
\end{aligned}
\]
If accordingly we divide the sums (9) each into four parts, and use
besides the trigonometric equation
\[
\sin \frac{(t+\mu)\pi}{m} = \cos \frac{t\pi}{m},
\]
then we obtain
\[
\begin{aligned}
X_1 &= \frac{-1}{m} \sum^{t} f_1(r^t)
\log\left(\operatorname{tg} \frac{t\pi}{m}\right)^2, \\
X_2 &= \frac{i\pi}{m} \sum^{t} f_2(r^t).
\end{aligned}
\]
In this sum we have
\[
\frac{t\pi}{m} < \frac{\pi}{4}, \qquad
\operatorname{tg} \frac{t\pi}{m} > 0,
\]
and consequently, having regard to (12), (13), we obtain, when we put
$\Theta^{-1}$ for $\Theta$, whereby $\chi(t)^{-1}$ passes into
$\chi(t)$:
\weq{15}{X_1 = -\frac{2}{m} (+1, \Theta^{-1}, r) \sum^{t} \chi_1(t)
\log\left(\operatorname{tg} \frac{t\pi}{m}\right),}
\weq{16}{X_2 = \frac{i\pi}{m} (-1, \Theta^{-1}, r) \sum^{t} \chi_2(t),
\qquad 1 \leqq t \leqq \nu - 1.}

\begin{center}
\S.\,198.

\gesperrt{The Class Number Factor $A$.}
\end{center}

In order now to calculate $A$ first of all, we have to substitute the
expression (16) for $X_2$ into the first formula (8) of the preceding
section. We can also represent this formula thus:
\weq{1}{\pi^{\nu'} A = 2\, m^{1/4\,\nu} 2^{-\nu'} \prod X_2.}


% ===== scan leaf 746, printed 723 =====
% !PRINTED 723
\weberpage{723}

If we denote by $\alpha$, $\beta$ the indices of $t$, then $\chi_2(t) = (-1)^\alpha\Theta^\beta;$

\weq{2}{\overset{t}{\Sigma}\,\chi_2(t) = \Sigma\,(-1)^\alpha\Theta^\beta = \varphi(\Theta)}

is an \gesperrt{integer of the field} $\Omega_\nu$, and we obtain:
\[
X_2 = \frac{i\pi}{m}\,(-1,\Theta^{-1},r)\,\varphi(\Theta).
\]

Now one has to put for $\Theta$ all the roots of the equation
\weq{3}{\Theta^{\nu'} + 1 = 0}
and obtains for $A$, with regard to \S.\,197, (14):
\weq{4}{A = 2^{1-\nu'}\overset{\Theta}{\prod}\,\varphi(\Theta).}

The sum by which the number $\varphi(\Theta)$ is defined in (2)
contains $\nu'$ terms of the form $(-1)^\alpha\Theta^\beta$, in which $\alpha$, $\beta$ are to be
determined so that
\weq{5}{(-1)^\alpha\,5^\beta \equiv t \pmod{m}}
and $t$ lies between 0 and $\nu$. Now it is at once evident that
among the exponents $\beta$ no two congruent with respect to the modulus $\nu'$ can occur. For if $\beta \equiv \beta' \pmod{\nu'}$, then
$5^\beta \equiv 5^{\beta'} \pmod{\mu}$, and from
\[
(-1)^\alpha\,5^\beta \equiv t, \quad (-1)^{\alpha'}\,5^{\beta'} \equiv t' \pmod{m}
\]
there follows:
\[
(-1)^\alpha\,t - (-1)^{\alpha'}\,t' \equiv 0 \pmod{\mu}.
\]

But since $t$ and $t'$ are absolutely smaller than $\nu = \tfrac{1}{2}\mu$, this
is only possible if $t = t'$, $\alpha = \alpha'$. Accordingly
$\beta$ runs in (2) through a complete system of residues with respect to the modulus $\nu'$.

But since $\beta$ is to be taken with respect to the modulus $\nu$, the
$\beta$ determined by (5) will turn out partly smaller, partly greater than $\nu'$; and so, if we let $\beta_1$ run through the sequence of numbers $0, 1, \dots, \nu' - 1$,
we obtain in (2) two kinds of
terms:
\weq{6}{\begin{aligned}
1.\ &(-1)^\alpha\Theta^\beta = (-1)^\alpha\Theta^{\beta_1}, & \beta = \beta_1 < \nu', \\
2.\ &(-1)^\nu\Theta^\beta = -(-1)^\alpha\Theta^{\beta_1}, & \beta = \beta_1 + \nu' \geqq \nu'.
\end{aligned}}
% [sic] in row 2 the page prints the exponent \nu where \alpha is required; reproduced as printed

Now by (5) the absolutely least residue of $(-1)^\alpha\,5^\beta$
with respect to the modulus $m$ is positive, and if we therefore determine a number
$c_\beta = \pm 1$ so that the absolutely least residue of $c_\beta\,5^\beta$
becomes positive, then in the case (6), 1. we have
\[
c_{\beta_1} = (-1)^\alpha.
\]


% ===== scan leaf 747, printed 724 =====
% !PRINTED 724
\weberpage{724}

Further, $5^{\nu'} \equiv 1 \pmod{\mu}$, and since $5^{\nu'}$ is not also congruent
with 1 with respect to the modulus $m$, it follows that:
\weq{7}{5^{\nu'} \equiv 1 + \mu \pmod{m}.}
Accordingly, in case (6), 2. we have
\[
(-1)^\alpha\,5^\beta = (-1)^\alpha\,5^{\beta_1}\,5^{\nu'} \equiv (-1)^\alpha\,5^{\beta_1} + \mu \pmod{m},
\]
and consequently the absolutely least residue of $(-1)^\alpha\,5^{\beta_1}$
with respect to the modulus $m$ is here negative. Thus in case (6), 2. the value
\[
c_{\beta_1} = -(-1)^\alpha
\]
is to be set.

From this there follows, by (2) (if we again write $\beta$ for $\beta_1$):
\weq{8}{\varphi(\Theta) = \sum_{0,\,\nu'-1}^{\beta} c_\beta\,\Theta^\beta = c_0 + c_1\Theta + c_2\Theta^2 + \cdots + c_{\nu'-1}\Theta^{\nu'-1},}
and here we have
\weq{9}{c_\beta = +1 \ \text{or} \ = -1,}
according as the absolutely least residue of $5^\beta$ is positive or
negative. Accordingly $\varphi(\Theta)$ is an \gesperrt{integer of the
field} $\Omega_\nu$.

It still remains to determine the behaviour of this number with
respect to the number 2, or rather with respect to the prime factor
$(1 - \Theta)$ of 2 contained in $\Omega_\nu$. If for this purpose we form
\weq{10}{(1 - \Theta)\,\varphi(\Theta) = 2\,\psi(\Theta),}
there results
\weq{11}{\psi(\Theta) = \frac{c_0 + c_{\nu'-1}}{2} + \frac{c_1 - c_0}{2}\,\Theta + \frac{c_2 - c_1}{2}\,\Theta^2 + \cdots + \frac{c_{\nu'-1} - c_{\nu'-2}}{2}\,\Theta^{\nu'-1}.}

The coefficients in this expression are all $= 0$ or
$= 1$, and therefore $\psi(\Theta)$ is likewise an integer of the field $\Omega_\nu$.
For this number, however, if one sets $\Theta \equiv 1$, there results:
\weq{12}{\psi(\Theta) \equiv c_{\nu'-1} \equiv \pm 1\ [\mathrm{mod}\ (1-\Theta)],}
and consequently $\psi(\Theta)$ is relatively prime to 2.

The product occurring in (4) is nothing other than the norm of the
number $\varphi(\Theta)$ taken with respect to the field $\Omega_\nu$,
and if we denote this norm by $N$, then by \S.\,169
\[
N(1-\Theta) = 2,
\]
%CONT>>


% ===== scan leaf 748, printed 725 =====
% !PRINTED 725
\weberpage{725}
%<<CONT
and consequently, by (10),
\[
N\varphi(\Theta) = 2^{\nu'-1}\,N\psi(\Theta).
\]
From this there results
\weq{13}{A = N\psi(\Theta),}
and from this it follows, having regard to (12), that $A$ is an \gesperrt{odd
integer.}

If we denote the number $A$ belonging to $m = 2^\varkappa$ by $A_\varkappa$, then
from \S.\,197, (4) there results, by complete induction, the
expression for the first factor of the class number:
\weq{14}{k = A_4\,A_5\,\dots\,A_\varkappa,}
and from this the theorem:

1. \gesperrt{The first factor of the class number is an odd integer.}

The number $A_\varkappa$ is, by formula (10), not hard to compute for the
first values of $\varkappa$.

For $m = 16$ one finds $\nu' = 2$, $\beta = 0, 1$, $c_0 = c_1 = 1$
consequently
\[
\psi(\Theta) = 1, \quad A_4 = 1.
\]

For $m = 32$, $\nu' = 4$, $\beta = 0, 1, 2, 3$.
\[
c_0 = 1, \quad c_1 = 1, \quad c_2 = -1, \quad c_3 = -1,
\]
\[
\psi(\Theta) = -\Theta^2,
\]
hence also $A_5 = 1$.

For $m = 64$, $\nu' = 8$, $\beta = 0, 1, 2, 3, 4, 5, 6, 7$ the
absolutely least residues of $5^\beta$ are
\[
1,\ 5,\ 25,\ -3,\ -15,\ -11,\ 9,\ -19;
\]
so the $c_\beta$ are:
\[
+1, +1, +1, -1, -1, -1, +1, -1
\]
and
\[
\psi(\Theta) = -\Theta^3 + \Theta^6 - \Theta^7, \quad \psi(\Theta)\,\psi(-\Theta) = \Theta^6(\Theta^6 - 2i),
\]
from which it easily results:
\[
A_6 = 17.
\]
A somewhat longer calculation gives in addition
\[
A_7 = 21121.
\]


% ===== scan leaf 749, printed 726 =====
% !PRINTED 726
\weberpage{726}
\pdfbookmark[2]{§ 199. The class number factor B}{bm227}

\begin{center}
\S.\,199.

\gesperrt{The class number factor} $B$.
\end{center}

The class number factor $B$ must be calculated in quite another way. Here one has to set, if $\Theta$ is again a primitive root
of the equation \S.\,198, (3) ($\Theta^{\nu'} + 1 = 0$),
\weq{1}{t \equiv (-1)^\alpha\,5^\beta \pmod{m},}
\weq{2}{\chi_1(t) = \Theta^\beta}
and formula \S.\,197, (15) then gives:
\weq{3}{X_1 = -\frac{2}{m}\,(1, \Theta^{-1}, r)\,\overset{t}{\Sigma}\,\Theta^\beta\,\log\,\mathrm{tg}\,\frac{t\pi}{m}.}

Now for every odd $t$, if [with regard to (1)]
\[
r = e^{\frac{2\pi i}{m}}, \quad r^\nu = i, \quad r^{\nu t} = (-1)^\alpha i
\]
is set, we have
\weq{4}{\mathrm{tg}\,\frac{t\pi}{m} = (-1)^\alpha\,r^{t\nu}\,\frac{1 - r^t}{1 + r^t},}
and this, as we have already seen in \S.\,178, is a unit of the real field $H_m$. One obtains the conjugate values
of this unit with respect to this field by setting $t \equiv 5^\beta \pmod{m}$ and letting $\beta$ run through a complete system of residues with respect to the modulus $\mu$.
We accordingly set:
\weq{5}{\tau_\beta = \mathrm{tg}\,\frac{5^\beta\pi}{m}.}

By \S.\,198, (7) we have $5^{\beta+\nu'} \equiv 5^\beta + \mu \pmod{m}$, and from this
there follows, by a well-known trigonometric formula:
\weq{6}{\tau_{\beta+\nu'} = \frac{-1}{\tau_\beta}.}

If we therefore now set
\weq{7}{\lambda_\beta = \tfrac{1}{2}\,\log\,\tau_\beta^2,}
then by (6)
\weq{8}{\lambda_{\beta+\nu'} = -\lambda_\beta,}
and accordingly
\weq{9}{\Theta^{\beta+\nu'}\,\lambda_{\beta+\nu'} = \Theta^\beta\,\lambda_\beta;}
among the values of $\beta$ occurring in the sum (3)
no two are, as we have already seen in the preceding section,
congruent with respect to the modulus $\nu'$, and if we therefore by means of
%CONT>>


% ===== scan leaf 750, printed 727 =====
% !PRINTED 727
\weberpage{727}

%<<CONT
formula (9) reduce all the values of $\beta$ to the interval from $0$ to
$\nu' - 1$, then we obtain all the numbers of this interval.

Accordingly formula (3) can be represented thus:
\weq{10}{
X_1 = -\frac{2}{m}\,(1,\,\Theta^{-1},\,r)\sum_{0,\,\nu'-1}^{\beta}\Theta^\beta\lambda_\beta.
}
This we have to substitute in the expression given in \S.\,197, (8)
\[
D B = 2^{2^{\varkappa-4}(\varkappa-2)}\prod X_1 = m^{\nu''}2^{-\nu'}\prod X_1
\]
and then, if we set
\weq{11}{
U_\Theta = \sum_{0,\,\nu'-1}^{\beta}\Theta^\beta\lambda_\beta
}
there results, with regard to \S.\,197, (14):
\weq{12}{
D B = \prod^{\Theta} U_\Theta.
}
This product can now be represented in the following way by a
determinant which contains only the conjugate logarithms of the
units $\tau$.

If we omit the index $\Theta$ on $U$, we can set up the system of
equations:
\[
\begin{array}{l}
U \;=\; \lambda_0 \;+\; \lambda_1\Theta \;+\; \lambda_2\Theta^2 \;+\; \cdots \;+\; \lambda_{\nu'-1}\Theta^{\nu'-1} \\[4pt]
\Theta U \;=\; -\lambda_{\nu'-1} \;+\; \lambda_0\Theta \;+\; \lambda_1\Theta^2 \;+\; \cdots \;+\; \lambda_{\nu'-2}\Theta^{\nu'-1} \\[4pt]
\text{\normalfont .\;.\;.\;.\;.\;.\;.\;.\;.\;.\;.\;.\;.\;.\;.\;.\;.\;.\;.\;.\;.\;.\;.\;.\;.\;.\;.\;.\;.\;.} \\[4pt]
\Theta^{\nu'-1} U \;=\; -\lambda_1 \;-\; \lambda_2\Theta \;-\; \lambda_3\Theta^2 \;+\; \cdots \;+\; \lambda_0\Theta^{\nu'-1},
\end{array}
\]
from which, by elimination of the powers of $\Theta$, it follows that:
\weq{13}{
\begin{vmatrix}
\lambda_0 - U, & \lambda_1, & \lambda_2, & \ldots, & \lambda_{\nu'-1} \\
-\lambda_{\nu'-1}, & \lambda_0 - U, & \lambda_1, & \ldots, & \lambda_{\nu'-2} \\
\multicolumn{5}{c}{\text{\normalfont .\;.\;.\;.\;.\;.\;.\;.\;.\;.\;.\;.\;.\;.\;.\;.}} \\
-\lambda_1, & -\lambda_2, & -\lambda_3, & \ldots, & \lambda_0 - U
\end{vmatrix} = 0.
}
% [source note: the first row prints a semicolon-shaped mark after \lambda_0-U; the mathematically required list punctuation is normalized to a comma.]
This is, for $U$, an equation of degree $\nu'$, whose roots are the
quantities $U_\Theta$. If we arrange it according to powers of $U$,
then, since $\nu'$ is even, the highest power $U^{\nu'}$ of $U$ receives
the coefficient $1$, and one therefore finds the product of the roots
occurring in (12) from the determinant (13) by setting $U = 0$.

The determinant so obtained we now wish to arrange somewhat
differently, by leaving the first row in its place, while the
%CONT>>


% ===== scan leaf 751, printed 728 =====
% !PRINTED 728
\weberpage{728}
\pdfbookmark[2]{§ 200. Normal units in Hm}{bm228}
%<<CONT
remaining ones, however, in reverse order and with altered sign.
In doing this $\frac{1}{2}\nu'$ changes of sign are necessary,
and we therefore obtain for the product (12) the determinant
defined by the following equation, whose absolute value we
denote by $T$:
\weq{14}{(-1)^{\frac{1}{2}\nu'}
\begin{vmatrix}
\lambda_0, & \lambda_1, & \ldots, & \lambda_{\nu'-2}, & \lambda_{\nu'-1}\\
\lambda_1, & \lambda_2, & \ldots, & \lambda_{\nu'-1}, & -\lambda_0\\
\lambda_2, & \lambda_3, & \ldots, & -\lambda_0, & -\lambda_1\\
\cdots & \cdots & \cdots & \cdots & \cdots\\
\lambda_{\nu'-1}, & -\lambda_0, & \ldots, & -\lambda_{\nu'-3}, & -\lambda_{\nu'-2}
\end{vmatrix} = \pm T,}
in which the negative signs stand to the right of the diagonal row
running from lower left to upper right.

The determinant $T$ is here so arranged that each row
arises from the preceding one by the substitution $(r, r^{5^\beta})$.
Each column therefore contains a system of conjugate logarithms.

By (12) we now have
\weq{15}{B = \frac{T}{D}.}
In order to gain clarity about the nature of this result,
it is necessary to consider somewhat more closely the number $D$, which was
defined by \S.\,195, (12), \S.\,197, (6); and this requires us to enter
 into the theory of the units of the field $H_m$.

\begin{center}
\S.\,200.\\
\gesperrt{Normal units in $H_m$.}
\end{center}

We now introduce a simplifying notation, by setting
\weq{1}{r^{5^\beta} = r_\beta}
and letting $\beta$ run through a complete system of residues with respect to the modulus $\nu$.
 From this it follows that
\weq{2}{r_{\beta+\nu'} = -r_\beta,}
and in the form $(r, r_\beta)$ not all the substitutions
of the group of the field $\Omega_m$ are then contained; rather, there
must still be added the substitutions $(r, r_\beta^{-1})$; but
$(r, r_\beta)$ does furnish us all the substitutions of the group of $H_m$.

We now denote by $\mathfrak{E}(r)$ a unit of the field $H_m$.
Among these are also the units of the field $H_\mu$ (as
%CONT>>


% ===== scan leaf 752, printed 729 =====
% !PRINTED 729
\weberpage{729}
%<<CONT
imprimitive numbers) and is characterised by the equation $\mathfrak{E}(r)
= \mathfrak{E}(-r)$.

By a \gesperrt{normal unit} of the field $H_m$ we shall now understand a
unit which is not $= \pm 1$, but for which the equation
\weq{3}{\mathfrak{E}(r)\mathfrak{E}(-r) = \pm 1}
is satisfied.

It is clear that a unit of the field $H_\mu$ can in no case satisfy
this condition. But not every unit in $H_m$ which does not belong
to the field $H_\mu$ will be a normal unit $^{1)}$.

The units $\tau_\beta$ [\S.\,199, (5)], however, do belong to the normal units.
For we have
\weq{4}{\tau_\beta = \mathrm{tg}\,\frac{5^\beta\pi}{m} = \frac{1-r_\beta}{1+r_\beta}r_\beta^{\nu},}
and if in this the substitution $(r,-r) = (r_\beta, r_{\beta+\nu'})$ is carried
out, then $\tau_\beta$ passes into
\weq{5}{\tau_{\beta+\nu'} = -\frac{1}{\tau_\beta}}
which corresponds to the relation (3).

A system of $\nu'$ normal units
\[
\mathfrak{E}_0(r), \mathfrak{E}_1(r), \ldots, \mathfrak{E}_{\nu'-1}(r)
\]
shall be called an \gesperrt{independent system} if the determinant
\weq{6}{\pm\Sigma\pm\log|\mathfrak{E}_0(r_0)|\log|\mathfrak{E}_1(r_1)|\ldots\log|\mathfrak{E}_{\nu'-1}(r_{\nu'-1})|
= L(\mathfrak{E}_0, \mathfrak{E}_1, \ldots, \mathfrak{E}_{\nu'-1})}
has a value different from zero.

The absolute value of this determinant, $L$, we call here too
the \gesperrt{regulator} of the system $\mathfrak{E}_0, \mathfrak{E}_1, \ldots, \mathfrak{E}_{\nu'-1}$.

The units $\tau_0, \tau_1, \ldots, \tau_{\nu'-1}$ form an independent system
of normal units, for their regulator
\[
L(\tau_0, \tau_1, \ldots, \tau_{\nu'-1})
\]
is precisely the determinant $T$ defined in the preceding section
[\S.\,199, (14), (15)]. And that this cannot be zero follows
immediately from the fact that it occurs as a factor in the expression for the
%CONT>>

\par\noindent{\footnotesize $^{1)}$ In the paper mentioned above (Acta
mathematica, Bd.\,8) the author designated the normal units as ``primitive
units''. This name is rejected here, because it could give occasion to the
error of supposing that every unit of the field $H_m$ which does not at the
same time belong to the field $H_\mu$ is to be reckoned among them.}


% ===== scan leaf 753, printed 730 =====
% !PRINTED 730
\weberpage{730}
%<<CONT
class number. But the class number, since there
certainly always exists at least one class, the principal class, is
different from zero, and therefore $T$ also cannot be equal to zero.
We therefore have the theorem:

1. \gesperrt{The numbers $\tau_0, \tau_1, \ldots, \tau_{\nu'-1}$ are an independent
system of normal units of the field $H_m$.}

If $\mathfrak{E}(r)$ is a normal unit, then, if we set
\weq{7}{L_\beta = \log|\mathfrak{E}(r_\beta)|}
we have, because of (3),
\weq{8}{L_\beta = -L_{\beta+\nu'}.}
Now if $\mathfrak{E}_0, \mathfrak{E}_1, \ldots, \mathfrak{E}_{\nu'-1}$ is any system of independent
% [source note: the separator after \mathfrak{E}_0 is printed as a period-shaped or damaged comma; the mathematically required list punctuation is normalized to a comma.]
normal units, we set:
\weq{9}{L_{i,k} = \log|\mathfrak{E}_i(r_k)|, \quad L_{i,k+\nu'} = -L_{i,k},}
so that the regulator of the system of the $\mathfrak{E}_i$ takes the expression:
\weq{10}{\pm\Sigma\pm L_{0,0}L_{1,1}\cdots L_{\nu'-1,\nu'-1} = L(\mathfrak{E}_0, \mathfrak{E}_1, \ldots, \mathfrak{E}_{\nu'-1}),}
and is a number different from zero, which we denote by $T_0$.

If $\mathfrak{E}(r)$ denotes any other normal unit, we can determine
the unknowns $\xi_0, \xi_1, \ldots, \xi_{\nu'-1}$ from the $\nu'$ linear equations
 that we obtain when in
\weq{11}{L_\beta = \xi_0 L_{0,\beta} + \xi_1 L_{1,\beta} + \cdots + \xi_{\nu'-1}L_{\nu'-1,\beta}}
we let the index $\beta$ run from 0 to $\nu'-1$; and because of (8) and (9)
this equation is correct for the remaining values of $\beta$ as well,
so that from it the formula valid for all $r$ can be derived:
\weq{12}{\mathfrak{E}(r) = \pm\mathfrak{E}_0(r)^{\xi_0}\mathfrak{E}_1(r)^{\xi_1}\cdots\mathfrak{E}_{\nu'-1}(r)^{\xi_{\nu'-1}}.}
If on the other hand $m_0, m_1, \ldots, m_{\nu'-1}$ denote any integral
rational numbers whatever, then all the numbers contained in the form
\[
\pm\mathfrak{E}_0(r)^{m_0}\mathfrak{E}_1(r)^{m_1}\cdots\mathfrak{E}_{\nu'-1}(r)^{m_{\nu'-1}}
\]
are normal units of the field $H_m$. Accordingly, the considerations
that we carried out in detail in \S.\,186 can be applied to the
system of equations (10), from which it follows:

2. \gesperrt{The exponents $\xi_0, \xi_1, \ldots, \xi_{\nu'-1}$ determined
from (10) or (11) are rational numbers.}


% ===== scan leaf 754, printed 731 =====
% !PRINTED 731
\weberpage{731}

From this it can be concluded, as in \S.\,187, that there are systems of $\nu'$ independent normal units whose regulator $T_0$ becomes \gesperrt{as small as possible}, and such systems are called \gesperrt{fundamental systems of normal units}.

If $\mathcal{E}_0, \mathcal{E}_1, \ldots, \mathcal{E}_{\nu'-1}$ is such a fundamental system, then the exponents $\xi_0, \xi_1, \ldots, \xi_{\nu'-1}$ in (10) and (11) come out as \gesperrt{integral rational numbers}. This too can be proved in the same way as in \S.\,187.

If, under this assumption about the $\mathcal{E}_i$, we set up the equations (11) for any other system of independent units $\mathcal{E}_0^{(1)}, \mathcal{E}_1^{(1)}, \ldots, \mathcal{E}_{\nu'-1}^{(1)}$, then we obtain expressions of the form:

\[
L_{i,k}^{(1)} = \xi_{i,0} L_{0,k} + \xi_{i,1} L_{1,k} + \cdots + \xi_{i,\nu'-1} L_{\nu'-1,k}.
\]

If, then, $T_1$ is the regulator of the system of the $\mathcal{E}_i^{(1)}$, it follows by the multiplication theorem for determinants, when by $C$ is denoted the absolute value of the determinant of the integer exponents $\xi_{i,k}$, and hence an \gesperrt{integral rational positive number}:

\weq{13}{T_1 = C\,T_0.}

It cannot be proved, and for the higher values of $m$ it probably does not even hold, that $\tau_0, \tau_1, \ldots, \tau_{\nu'-1}$ form a fundamental system\footnote{It is here, as with all questions that concern the units, very difficult to arrive at definite numerical results. According to the results contained in the comprehensive work of \gesperrt{Reuschle}, \glqq Tafeln complexer Primzahlen, welche aus Wurzeln der Einheiten gebildet sind\grqq{} (Berlin 1875), the $\tau$ still form a fundamental system for $m=16$. For here the fundamental number of the field $H_m$ is equal to $2^{11}$, the degree of the field is equal to 4, and hence, by the remark to \S.\,156, there is in every ideal class an ideal whose norm is less than 6. But according to the tables of \gesperrt{Reuschle} every natural prime below 100 can be decomposed in the field $H_m$ into actually existing prime factors. Consequently only principal ideals certainly divide all numbers below 6, and consequently there is here only the one class, the principal class. Since, as we have seen, the first class-number factor is also $=1$, the cyclotomic field $\Omega_{16}$ is of one class. In it the same laws of primes hold as in the field of rational numbers. This is doubtful for the field $\Omega_{32}$, and certainly no longer true in the field $\Omega_{64}$, in which the class number is a multiple of 17.}.

For our purpose, however, it is of importance that these units, with respect to the denominator 2 in the exponents
%CONT>>


% ===== scan leaf 9732, printed 732 =====
% !PRINTED 732
\weberpage{732}
%<<CONT
behave in a certain sense like a fundamental system, which is expressed
more definitely by the following theorem:

2. \gesperrt{If a normal unit of the field $H_m$
of the form}
\[
\mathfrak{E}(r) = \tau_0^{e_0}\tau_1^{e_1}\ldots\tau_{\nu'-1}^{e_{\nu'-1}},
\]
\gesperrt{in which the exponents $e_0, e_1, \ldots, e_{\nu'-1}$ are integral rational
 numbers, has one and the same sign as all its conjugate
values, then all of the $e_0, e_1, \ldots, e_{\nu'-1}$ must be even numbers,
and $\mathfrak{E}(r)$ is the square of a normal unit.}

We can again carry out the proof by complete induction, but wish first
to transform the expression for $\tau_\beta$ somewhat.

We obtain the conjugate values to the unit $\mathfrak{E}(r)$
in the form
\[
\mathfrak{E}(r_\beta) = \tau_\beta^{e_0}\tau_{\beta+1}^{e_1}\ldots\tau_{\beta+\nu'-1}^{e_{\nu'-1}},
\]
and these expressions are therefore to have the same sign for all
values of $\beta$.

By §.\,199, (5) we have
\[
\tau_\beta = \mathrm{tg}\,\frac{5^\beta\pi}{m} = \frac{1}{2}\frac{\sin\left(\dfrac{2\pi}{m}5^\beta\right)}{\left[\cos\left(\dfrac{5^\beta\pi}{m}\right)\right]^2}.
\]

In this we put, for abbreviation,
\weq{14}{\sigma_\beta = \sin\left(\frac{2\pi}{m}5^\beta\right),}
so that $\sigma_\beta$ always has the same sign as $\tau_\beta$.

From this we form the product
\weq{15}{S_\beta = \sigma_\beta^{e_0}\sigma_{\beta+1}^{e_1}\ldots\sigma_{\beta+\nu'-1}^{e_{\nu'-1}},}
and our assumption is therefore this, that $S_\beta$ has the same sign
for all values of $\beta$.

It is to be shown that the integers
$e_0, e_1, \ldots, e_{\nu'-1}$ must all be even. The theorem is
correct for $m=8$; for in this case we have
\[
\sigma_0 = \sin\frac{2\pi}{8}, \sigma_1 = \sin\frac{10\pi}{8} = -\sin\frac{6\pi}{8},
\]
%CONT>>


% ===== scan leaf 9733, printed 733 =====
% !PRINTED 733
\weberpage{733}
%<<CONT

therefore $\sigma_0$ positive, $\sigma_1$ negative. Here now $S_\beta = \sigma_\beta^{e_0}$,
and hence $S_0$ and $S_1$ can have the same sign only when $e_0$
is even.

Accordingly we apply complete induction. For simplification
we set $\nu' = 2\nu''$ and recall the notation

\weq{16}{m = 2^{\varkappa}, \quad \mu = \tfrac{1}{2}m, \quad \nu = \tfrac{1}{2}\mu, \quad \nu' = \tfrac{1}{2}\nu, \quad \nu'' = \tfrac{1}{2}\nu',}
and assume that our theorem has been proved correct when
$\mu$ takes the place of $m$.

We then have
\[
\begin{gathered}
5^{\nu'} \equiv 1 \pmod{m}, \quad 5^{\nu'} \equiv 1+\mu \pmod{m}, \quad 5^{\nu''} \equiv 1+\nu \pmod{\mu} \\
5^{\nu''} \equiv 1+\nu+\mu \pmod{m}
\end{gathered}
\]
and from this there results

\[
\sigma_{\beta+\nu'} = -\sigma_\beta,
\]
\[
\sigma_{\beta+\nu''} = -\cos\left(\frac{2\pi}{m}5^\beta\right),
\]
\weq{17}{\sigma_\beta \sigma_{\beta+\nu''} = -\tfrac{1}{2}\sin\left(\frac{2\pi}{\mu}5^\beta\right) = -\tfrac{1}{2}\sigma'_\beta,}
\weq{18}{\sigma'_{\beta+\nu''} = -\sigma'_\beta,}
if the $\sigma'_\beta$ arise from the $\sigma_\beta$ through the transition from $\varkappa$ to
$\varkappa - 1$. Now, by (17), (18), we form the product
$S_\beta S_{\beta+\nu''}$, and obtain, when in addition, for abbreviation,
\[
e' = e_0 + e_1 + \cdots + e_{\nu'-1}, \quad e'' = e_0 + e_1 + \cdots + e_{\nu''-1}
\]
is set,
\[
S_\beta S_{\beta+\nu''} = (-1)^{e''}\left(\tfrac{1}{2}\right)^{e'}\sigma_\beta'^{\,e_0+e_{\nu''}}\sigma_{\beta+1}'^{\,e_1+e_{\nu''+1}}\cdots\sigma_{\beta+\nu''-1}'^{\,e_{\nu''-1}+e_{\nu'-1}}.
\]

But this is a product of the same form as $S_\beta$, in
which $\mu$ has taken the place of $m$, and which likewise has
the same sign as all its conjugates. Therefore, by
the assumption made,
\weq{19}{e_0 \equiv e_{\nu''+1}, \; e_1 \equiv e_{\nu''+1}, \; \ldots, \; e_{\nu''-1} \equiv e_{\nu'-1} \pmod{2}.} % [sic: the source prints the same right-hand subscript in the first two congruences; the preceding product indicates that the first should be e_{\nu''}.]
If, then, we set
\[
S'_\beta = \sigma_\beta'^{\,e_0}\sigma_{\beta+1}'^{\,e_1}\cdots\sigma_{\beta+\nu''-1}'^{\,e_{\nu''-1}},
\]
\[
R_\beta = \sigma_{\beta+\nu''}^{\frac{1}{2}(e_0-e_{\nu''})}\sigma_{\beta+\nu''+1}^{\frac{1}{2}(e_1-e_{\nu''+1})}\cdots\sigma_{\beta+\nu'-1}^{\frac{1}{2}(e_{\nu''-1}-e_{\nu'-1})},
\]
then from (15) and (17) it follows that
\[
S_\beta R_\beta^2 = \left(-\tfrac{1}{2}\right)^{e''}S'_\beta,
\]
%CONT>>


% ===== scan leaf 755, printed 734 =====
% !PRINTED 734
\weberpage{734}

%<<CONT
and from this it follows that $S'_\beta$, just like $S_\beta$, has the same sign as all its conjugate values. But $S'_\beta$ arises from $S_\beta$ by putting $\mu$ in place of $m$, and by our assumption we therefore have

\[
e_0 \equiv 0,\ e_1 \equiv 0,\ \ldots,\ e_{\nu''-1} \equiv 0 \ (\mathrm{mod}\ 2),
\]

and, with the help of (19),

\[
e_{\nu''} \equiv 0,\ e_{\nu''+1} \equiv 0,\ \ldots,\ e_{\nu'-1} \equiv 0 \ (\mathrm{mod}\ 2),
\]

whereby the proof of Theorem 2. is carried through.

If we now choose in formula (12) the system $\tau_0, \tau_1, \ldots, \tau_{\nu'-1}$ for $\mathcal{E}_0, \mathcal{E}_1, \ldots, \mathcal{E}_{\nu'-1}$, and put the exponents $\xi$ into the form:

\[
\xi_0 = \frac{e_0}{e},\ \xi_1 = \frac{e_1}{e},\ \cdots,\ \xi_{\nu'-1} = \frac{e_{\nu'-1}}{e},
\]

in which $e, e_0, e_1, \ldots, e_{\nu'-1}$ are integral rational numbers without a common divisor, then we obtain

\[
\mathcal{E}(r)^e = \pm\, \tau_0^{e_0}\tau_1^{e_1} \cdots \tau_{\nu'-1}^{e_{\nu'-1}},
\]

and it follows from our theorem that $e$ must be \gesperrt{odd}; for if it were even, then $\mathcal{E}(r)^e$ would be the square of a real quantity, and consequently $\pm\, \mathcal{E}(r)^e$ would have the same sign as its conjugates. But then, by our theorem, all the $e_0, e_1, \ldots, e_{\nu'-1}$ would have to be divisible by 2, which again runs counter to the assumption that the exponents $e, e_0, \ldots, e_{\nu'-1}$ are to have no common divisor. We therefore have the theorem:

3. \gesperrt{If $\mathcal{E}(r)$ is any normal unit of the field $H_m$, then the odd number $e$ and the integers $e_0, e_1, \ldots, e_{\nu'-1}$ can be so determined that}

\weq{20}{\mathcal{E}(r_\beta)^e = \pm\, \tau_\beta^{e_0}\tau_{\beta+1}^{e_1} \cdots \tau_{\beta+\nu'-1}^{e_{\nu'-1}}}

\gesperrt{holds.}

This result leads to several important consequences:

4. \gesperrt{If $T$ is the regulator of the system of the $\tau_\beta$, and $T_0$ the regulator of a fundamental system of normal units in $H_m$, i.\,e. the minimal value which the regulator of any system of independent normal units can assume, then}

\weq{21}{T = C\,T_0, \quad C \equiv 1\ (\mathrm{mod}\ 2),}

\gesperrt{in which $C$ denotes an odd natural number.}


% ===== scan leaf 756, printed 735 =====
% !PRINTED 735
\weberpage{735}
\pdfbookmark[2]{§ 201. Fundamental system of units of the field Hm}{bm229}

That $T:T_0 = C$ is an integer follows, namely, from formula (13).

If, however, we apply theorem 3. to all elements of the fundamental system $\mathcal{E}_i(r)$, there results a system of integers $e_{i,k}$ and an odd number $e$, such that

\[
\mathcal{E}_i(r)^e = \pm\, \tau_0^{e_{i,0}}\tau_1^{e_{i,1}} \cdots \tau_{\nu'-1}^{e_{i,\nu'-1}}.
\]

If we take the logarithms of this, and then form the regulator $T_0$, there results

\[
e^{\nu'} T_0 = \pm\, \Sigma \pm e_{0,0}\, e_{1,1} \cdots e_{\nu'-1,\nu'-1}\, T,
\]

consequently

\[
e^{\nu'} = \pm\, C\, \Sigma \pm e_{0,0}\, e_{1,1} \cdots e_{\nu'-1,\nu'-1},
\]

and since the left-hand side here is odd, $C$ cannot be even, which was to be proved.

If all conjugate values of the normal unit $\mathcal{E}(r_\beta)$ have one and the same sign, then the same holds of $\mathcal{E}(r_\beta)^e$ as well, and consequently in this case the exponents $e_0, e_1, \ldots, e_{\nu'-1}$ in formula (20) must, by 2., be even numbers. If one then writes formula (20) thus:

\[
\mathcal{E}(r) = \pm\, \tau_0^{e_0}\tau_1^{e_1} \cdots \tau_{\nu'-1}^{e_{\nu'-1}}\, \mathcal{E}(r)^{1-e},
\]

one obtains from it the theorem:

5. \gesperrt{A normal unit of the field $H_m$ which has the same sign as its conjugates is, apart from sign, a square of a normal unit.}

\begin{center}
\S.\,201.

\gesperrt{Fundamental system of units of the field $H_m$.}
\end{center}

It still remains to determine the denominator $D$ in the formula for $B$ [\S.\,197, (6)], which was defined by the equation

\[
E = E'\, D
\]

where $E$ and $E'$ denoted the regulators of the fields $H_m, H_\mu$.

Here, then, it is no longer a question of the normal units, but of the units of the field $H_m$ in general.


% ===== scan leaf 757, printed 736 =====
% !PRINTED 736
\weberpage{736}

We assume a fundamental system of units in the field $H_m$:

\weq{1}{\varepsilon_1(r_i),\ \varepsilon_2(r_i),\ \ldots,\ \varepsilon_{\nu-1}(r_i),}

and let the conjugate logarithms of this system be:

\weq{2}{l_{1,k},\ l_{2,k},\ \ldots,\ l_{\nu-1,k} \quad k=0,1,2,\ldots,\nu-1.}

Beside this we consider a fundamental system in the field $H_\mu$

\weq{3}{\varepsilon_1',\ \varepsilon_2',\ \ldots,\ \varepsilon_{\nu'-1}'}

with the conjugate logarithms

\weq{4}{l_{1,k}',\ l_{2,k}',\ \ldots,\ l_{\nu'-1,k}' \quad k=0,1,2,\ldots,\nu'-1.}

In the formation of the regulators one of the conjugate values, say $k=0$, is omitted (cf. \S.\,186), and we therefore obtain

\weq{5}{\begin{aligned} E &= \pm\, \Sigma \pm l_{1,1} l_{2,2} \ldots l_{\nu-1,\nu-1} \\ E' &= \pm\, \Sigma \pm l_{1,1}' l_{2,2}' \ldots l_{\nu'-1,\nu'-1}'. \end{aligned}}

Now we add to the system (3) a fundamental system of normal units of the field $H_m$

\weq{6}{\mathcal{E}_0,\ \mathcal{E}_1,\ \ldots,\ \mathcal{E}_{\nu'-1},}

and we assert that the system

\weq{7}{\varepsilon_1',\ \varepsilon_2',\ \ldots,\ \varepsilon_{\nu'-1}',\ \mathcal{E}_0,\ \mathcal{E}_1,\ \ldots,\ \mathcal{E}_{\nu'-1}}

is an independent system of units of the field $H_m$ (though not a fundamental system).

If we denote by

\weq{8}{L_{0,k},\ L_{1,k},\ \ldots,\ L_{\nu'-1,k} \quad k=0,1,\ldots,\nu'-1}

the conjugate logarithms of the system (6), then by the preceding section

\weq{9}{\pm\, \Sigma \pm L_{0,0} L_{1,1} \ldots L_{\nu'-1,\nu'-1} = T_0}

is the regulator of this system, and with regard to the definition of the normal units [\S.\,200, (9)] we have:

\weq{10}{L_{i,k+\nu'} = -L_{i,k},}

and since the $\varepsilon_i'$, as numbers of the field $H_\mu$, remain unchanged under the substitution $(r,-r)$:

\weq{11}{l_{i,k+\nu'}' = l_{i,k}'.}


% ===== scan leaf 758, printed 737 =====
% !PRINTED 737
\weberpage{737}

From this the regulator $R$ of the system (7) is obtained from
the determinant:

\[
\left|
\begin{array}{cccccccc}
l'_{1,1}, & l'_{2,1}, & \ldots, & l'_{\nu'-1,1}, & L_{0,1}, & L_{1,1}, & \ldots, & L_{\nu'-1,1} \\
\cdots & \cdots & \cdots & \cdots & \cdots & \cdots & \cdots & \cdots \\
l'_{1,\nu'-1}, & l'_{2,\nu'-1}, & \ldots, & l'_{\nu'-1,\nu'-1}, & L_{0,\nu'-1}, & L_{1,\nu'-1}, & \ldots, & L_{\nu'-1,\nu'-1} \\
l'_{1,0}, & l'_{2,0}, & \ldots, & l'_{\nu'-1,0}, & -L_{0,0}, & -L_{1,0}, & \ldots, & -L_{\nu'-1,0} \\
l'_{1,1}, & l'_{2,1}, & \ldots, & l'_{\nu'-1,1}, & -L_{0,1}, & -L_{1,1}, & \ldots, & -L_{\nu'-1,1} \\
\cdots & \cdots & \cdots & \cdots & \cdots & \cdots & \cdots & \cdots \\
l'_{1,\nu'-1}, & l'_{2,\nu'-1}, & \ldots, & l'_{\nu'-1,\nu'-1}, & -L_{0,\nu'-1}, & -L_{1,\nu'-1}, & \ldots, & -L_{\nu'-1,\nu'-1}
\end{array}
\right|.
\]

This determinant has $\nu-1$ rows; we add the first to the
$(\nu'+1)$th, the second to the $(\nu'+2)$th, and so on, the $(\nu'-1)$th to
the last. In this way the $L$ cancel out in the last $(\nu'-1)$ rows
and the $l$ acquire the factor 2 [by (10), (11)], and from this
there follows, by a theorem on determinants (Bd.\,I,
\S.\,23, V.):
\weq{12}{R = 2^{\nu'-1} E' T_0.}

Since $R$ accordingly does not vanish, and there is therefore
present in (7) a system of $\nu-1$ independent units of the field $H_m$,
we can [by \S.\,186, (5)] determine the rational (integral
or fractional) numbers $m_{i,k}$, $M_{i,k}$ from the equations
\weq{13}{
\begin{aligned}
2\, l_{i,k} & = m_{1,i}\, l'_{1,k} + \cdots + m_{\nu'-1,i}\, l'_{\nu'-1,k} \\
& \quad + M_{0,i}\, L_{0,k} + M_{1,i}\, L_{1,k} + \cdots + M_{\nu'-1,i}\, L_{\nu'-1,k}.
\end{aligned}
}

But it can easily be shown by means of the relations (10), (11)
that the numbers $m_{i,k}$, $M_{i,k}$ here must be \gesperrt{integral}.
For we have
\[
\begin{aligned}
l_{i,k} + l_{i,k+\nu'} & = \log \lvert \varepsilon_i(r_k)\, \varepsilon_i(-r_k) \rvert = \\
& \quad m_{1,i}\, l'_{1,k} + \cdots + m_{\nu'-1,i}\, l'_{\nu'-1,k} \\
l_{i,k} - l_{i,k+\nu'} & = \log \lvert \varepsilon_i(r_k)\, \varepsilon_i(-r_k)^{-1} \rvert = \\
& \quad M_{0,i}\, L_{0,k} + M_{1,i}\, L_{1,k} + \cdots + M_{\nu'-1,i}\, L_{\nu'-1,k}.
\end{aligned}
\]

But the units $\varepsilon_i(r)\, \varepsilon_i(-r)$ belong to the field $H_\mu$,
and since the $\varepsilon'_i$ are by hypothesis a fundamental system
of this field, the $m_{i,k}$ must, by \S.\,187, be integral rational
 numbers.

The units $\varepsilon_i(r)\, \varepsilon_i(-r)^{-1}$ are normal units of the field
 $H_m$, and since the $\mathcal{E}_i$ were to be a fundamental system of normal
 units, the $M_{i,k}$ are, by \S.\,200, integers.


% ===== scan leaf 759, printed 738 =====
% !PRINTED 738
\weberpage{738}

For a better overview we set down the equations (13) here once
more without the index $k$, which distinguishes the conjugate values
from one another:
\weq{14}{
\begin{aligned}
2\, l_i & = m_{1,i}\, l'_1 + m_{2,i}\, l'_2 + \cdots + m_{\nu'-1,i}\, l'_{\nu'-1} \\
& \quad + M_{0,i}\, L_0 + M_{1,i}\, L_1 + M_{2,i}\, L_2 + \cdots + M_{\nu'-1,i}\, L_{\nu'-1}.
\end{aligned}
}

If we now form the regulator $E$ according to (13), and denote
by $M$ the absolute value of the determinant of the $\nu-1$ rows
of the numbers $m_{k,i}$, $M_{k,i}$:
\[
M = \pm \Sigma \pm m_{1,1} \ldots m_{\nu'-1,\nu'-1}\, M_{0,0}\, M_{1,1} \ldots M_{\nu'-1,\nu'-1},
\]
then it follows:
\weq{15}{2^{\nu-1} E = M R,}
and from this, having regard to (12):
\weq{16}{E = 2^{-\nu'} M E' T_0.}

It can now be shown further that $M$ is a power of 2 and is
divisible by $2^{\nu'}$.

For since the system (1) was assumed to be a fundamental system
in $H_m$, there exist \gesperrt{integral rational numbers}
$n_{k,i}$, $N_{k,i}$ which satisfy the equations:
\weq{17}{
\begin{aligned}
l'_i & = n_{1,i}\, l_1 + n_{2,i}\, l_2 + \cdots + n_{\nu-1,i}\, l_{\nu-1}, \quad i = 1, 2, \ldots, \nu'-1 \\
L_i & = N_{1,i}\, l_1 + N_{2,i}\, l_2 + \cdots + N_{\nu-1,i}\, l_{\nu-1}, \quad i = 0, 1, 2, \ldots, \nu'-1.
\end{aligned}
}

If we form from these the determinant $R$ and denote by $N$
the absolute value of the determinant of the $n_{k,i}$, $N_{k,i}$, then it follows:
\[
R = N E,
\]
and from this, by (15):
\[
M N = 2^{\nu-1};
\]
from this it follows that each of the two integers $M$, $N$ must
be a power of 2.

In order to decide by which power of 2 the determinant
$M$ is divisible, we first determine a system of integral
rational numbers $a_i$ without a common divisor, which satisfies
the linear equations
\weq{18}{\sum\limits_{1,\, \nu-1}^{i} a_i\, m_{k,i} = 0 \quad (k = 1, 2, \ldots, \nu'-1).}
If we then put
\[
\sum\limits_{1,\, \nu-1}^{i} a_i\, M_{k,i} = \xi_k,
\]
%CONT>>


% ===== scan leaf 760, printed 739 =====
% !PRINTED 739
\weberpage{739}

%<<CONT
then from equations (13) we obtain
\[
2 \sum\limits_{1,\, \nu-1}^{i} a_i\, l_{i,k} = \xi_0 L_{0,k} + \xi_1 L_{1,k} + \cdots + \xi_{\nu'-1} L_{\nu'-1,k}.
\]
From this there results, by (10),
\[
\sum^{i} a_i\, l_{i,k+\nu'} = -\sum^{i} a_i\, l_{i,k},
\]
and this means that
\[
\delta = \varepsilon_1^{a_1}\, \varepsilon_2^{a_2} \ldots \varepsilon_{\nu-1}^{a_{\nu-1}}
\]
is a normal unit of the field $H_m$. Since $\mathcal{E}_0, \mathcal{E}_1, \ldots, \mathcal{E}_{\nu'-1}$
is a fundamental system of normal units, it then follows further that
 the numbers $\xi_0, \xi_1, \ldots, \xi_{\nu'-1}$ must be even. From this
there results:

1. \gesperrt{The system of equations (18) has the congruences
}
\weq{19}{\begin{gathered}
\sum\limits_{1,\, \nu-1}^{i} a_i\, M_{k,i} \equiv 0 \pmod{2} \\
k = 0, 1, \ldots, \nu'-1
\end{gathered}}
\gesperrt{as a consequence.}

But equations (18), whose number amounts to only $\nu'-1$, contain
$\nu-1$ unknowns $a_i$. Hence there are several systems of mutually
independent solutions, which we can make somewhat more precise
as follows:

2. \gesperrt{There are $\nu'$ integral solutions of the system
(18):}
\weq{20}{
\begin{aligned}
& a_{1,1}, a_{2,1}, \ldots, a_{\nu-1,1} \\
& a_{1,2}, a_{2,2}, \ldots, a_{\nu-1,2} \\
& \ldots\ldots\ldots\ldots \\
& a_{1,\nu'}, a_{2,\nu'}, \ldots, a_{\nu-1,\nu'}
\end{aligned}
}
\gesperrt{of such a kind that from the matrix (20) a
determinant of $\nu'$ rows can be formed
which is an odd number.}

For there is in the first place a solution of (18):
\[
a_{1,1}, a_{2,1}, a_{3,1}, \ldots, a_{\nu-1,1},
\]
in which an odd number occurs. Let us assume---as is permissible,
since here nothing depends on the arrangement of the indices---that
%CONT>>


% ===== scan leaf 761, printed 740 =====
% !PRINTED 740
\weberpage{740}

%<<CONT
$a_{1,1}$ is this odd number. Then we determine a
second solution:
\[
0, a_{2,2}, a_{3,2}, \ldots, a_{\nu-1,2},
\]
in which again an odd number, say $a_{2,2}$, occurs. Then we
determine a third solution:
\[
0, 0, a_{3,3}, \ldots, a_{\nu-1,3},
\]
and we continue setting to zero as long as the number of the
unknowns still remaining is greater than the number of the
equations.

In the last system, therefore, $\nu'-1$ zeros will occur, so
that $\nu'$ unknowns remain.

We may take these rows of numbers $a_{\nu,k}$ as the system (20),
for in it the determinant
\weq{21}{\Sigma \pm a_{1,1}\, a_{2,2} \ldots a_{\nu',\nu'}}
reduces to the diagonal term $a_{1,1}\, a_{2,2} \ldots a_{\nu',\nu'}$, whose factors are all
 odd.

If now the condition in statement 2 is satisfied, the matrix
(20) can be completed, by the addition of $\nu'-1$ rows, to a
determinant of $\nu-1$ rows which likewise has an odd integral
value:
\weq{22}{a = \Sigma \pm a_{1,1}\, a_{2,2} \ldots a_{\nu',\nu'}\, a_{\nu'+1,\nu'+1} \ldots a_{\nu-1,\nu-1}.}

One need only, when the determinant (21) is assumed to be odd,
set nothing but zeros for the elements in the added rows, except
those on the diagonal,
 $a_{\nu'+1,\nu'+1} \ldots a_{\nu-1,\nu-1}$, which may be taken equal to 1. Then $a$
receives the same value as the determinant
 (21).

But if we now form the product $a M$ according to the
multiplication theorem for determinants, there results a
determinant in which the $\nu-1$ sums
\[
\begin{aligned}
& \sum\limits_{1,\, \nu-1}^{i} a_{i,s}\, m_{k,i}, \qquad \sum\limits_{1,\, \nu-1}^{i} a_{i,s}\, M_{k,i} \\
& k = 1, 2, \ldots \nu'-1, \qquad k = 0, 1, 2 \ldots \nu'-1
\end{aligned}
\]
form a row for each value $s = 1, 2, \ldots, \nu'$, and in which
therefore $\nu'$ rows occur whose elements are all divisible by
 2.


% ===== scan leaf 762, printed 741 =====
% !PRINTED 741
\weberpage{741}
Accordingly $aM$, and consequently $M$ as well, is divisible by $2^{\nu'}$,
and if we set
\weq{23}{M = 2^{\nu'+\sigma}, \quad N = 2^{\nu'-\sigma-1}}
then $\sigma$ is a non-negative integral rational number.

If we substitute this in (16), it follows that:
\[ E = 2^\sigma E' T_0. \]
Now in \S.\,197, (6) we had put
\[ E = E' D \]
and there results therefore:
\[ D = 2^\sigma T_0. \]
From this there follows for the class number factor $B$, by \S.\,199, (15)
and \S.\,200, (21):
\weq{24}{B = 2^{-\sigma} \frac{T}{T_0} = 2^{-\sigma} C,}
in which $C$ is an odd integer.

Accordingly there results, by \S.\,197, (4), for the second factor
of the class number $h_1$, which, as we have seen, is an
\gesperrt{integer}:
\weq{25}{h_1 = 2^{-\sigma} h_1' C.}
If we now assume it proved that $h_1'$ is odd, then, since
$C$ is odd as well, it follows that $\sigma = 0$ and that $h_1$
is odd; and both are therewith proved by complete
induction.

With this the results derived are simplified as
follows:

3. \gesperrt{We have}
\weq{26}{M = 2^{\nu'}, \quad D = T_0, \quad B = \frac{T}{T_0};}
\gesperrt{the class number factor $B$ is equal to the ratio
of the regulator $T$ of the units $\tau$ to the regulator $T_0$
of a fundamental system of normal units, and is an
odd integer.}

If we denote the number $B$ belonging to $m = 2^\varkappa$ by
$B_\varkappa$, then the class number $h_1$ is, as follows by
complete induction,
\weq{27}{h_1 = B_4 B_5 \ldots B_\varkappa,}
and is therefore likewise an odd number. With this, then, having
regard to \S.\,198, 1., the main theorem is proved:

A. \gesperrt{The class number in the field $\Omega_m$ is odd when
$m$ is a power of 2.}


% ===== scan leaf 763, printed 742 =====
% !PRINTED 742
\weberpage{742}
\pdfbookmark[2]{§ 202. Positive units}{bm230}

\S.\,202.

\gesperrt{Positive Units.}

The results last obtained show that the system of independent
units \S.\,201, (7)
\[ \varepsilon'_1, \varepsilon'_2, \ldots, \varepsilon'_{\nu'-1}, \; \mathcal{E}_0, \mathcal{E}_1, \ldots, \mathcal{E}_{\nu'-1} \]
forms no \gesperrt{fundamental system} in the field $H_m$. For otherwise the
% [sic] -- the source omits a separator between the coefficient families $m_{k,i}$ and $M_{k,i}$; the printed form is preserved.
numbers $m_{k,i} M_{k,i}$ in the formulae \S.\,201, (13) would all have
to be divisible by 2, and the determinant $M$, whose value we
found $= 2^{\nu'}$, would be divisible by $2^{\nu-1}$.

In place of the system of equations (18), \S.\,201, we now consider
the following congruences:
\weq{1}{
\begin{aligned}
\sum_{1,\,\nu-1}^{i} b_i m_{1,i} &\equiv 1 \\
\sum_{1,\,\nu-1}^{i} b_i m_{2,i} &\equiv 0 \pmod{2}, \\
& \cdots\cdots\cdots\cdots \\
\sum_{1,\,\nu-1}^{i} b_i m_{\nu'-1,i} &\equiv 0
\end{aligned}
}
and we prove that it is possible to satisfy them by integral
values of the $b_i$.

If it were not possible, then from the $\nu'-2$ congruences
\[ \sum_{1,\,\nu-1}^{i} b_i m_{k,i} \equiv 0 \pmod{2} \quad k = 2, 3, \ldots, \nu'-1 \]
the last one
\[ \sum_{1,\,\nu-1}^{i} b_i m_{1,i} \equiv 0 \pmod{2} \]
would have to follow as a \gesperrt{consequence}, and then, as in the
preceding section, it would follow further:
\[ \sum_{1,\,\nu-1}^{i} b_i M_{k,i} \equiv 0 \pmod{2} \quad k = 0, 1, \ldots, \nu'-1. \]


% ===== scan leaf 764, printed 743 =====
% !PRINTED 743
\weberpage{743}
In place of the matrix \S.\,201, (20) one could then set a matrix of
the $b_{k,i}$ with $\nu'+1$ rows, and along the path taken in the
preceding section it could be concluded that $M$ would have to be
divisible by $2^{\nu'+1}$, which is not the case.

But if we apply the system of multipliers $b_i$, which satisfies the
conditions (1), to the equations (14) of the preceding section,
multiplying by $b_i$ and then summing, we obtain, when we omit
multiples of 2 and express this here too (where irrational numbers,
namely the logarithms, occur) by the symbol of congruence:
\weq{2}{
\begin{aligned}
l'_1 \equiv{} & L_0 \sum M_{0,i} b_i + L_1 \sum M_{1,i} b_i + \cdots \\
& + L_{\nu'-1} \sum M_{\nu'-1,i} b_i \pmod{2}.
\end{aligned}
}
If we wish to transform the congruence into an equation, there is
added a sum of quantities
\[ l_1, l_2, \ldots, l_{\nu-1}, \quad l'_1, l'_2, \ldots, l'_{\nu'-1} \]
with \gesperrt{even integers} as coefficients.

The same consideration can, however, be applied to $l'_2, l'_3,
\ldots, l'_{\nu'-1}$, and accordingly, if one passes back from the
logarithms to the numbers, formula (2) can be put into words as
follows:

1. \gesperrt{Every unit of the field $H_\mu$ can be represented as the
product of a normal unit of the field $H_m$ and a square of a unit
in $H_m$.}

If we then apply theorem 2., \S.\,200, it follows that every unit
$\mathcal{E}(r)$ of the field $H_\mu$ which has the same sign as all
its conjugates must, apart from the sign, be the square of a unit in
$H_m$. But if
\[ \pm\, \mathcal{E}(r) = \Delta(r)^2, \]
where $\mathcal{E}(r)$ is a unit in $H_\mu$ and $\Delta(r)$ a unit in
$H_m$, then $\Delta(r)$ must also be contained in $H_\mu$.

For if we set
\[ \Delta(r) = \Delta_1(r^2) + r \Delta_2(r^2), \]
then the square of it can be contained in $H_\mu$, and so remain
unchanged under the change of sign of $r$, only if either $\Delta_1$
or $\Delta_2 = 0$. But we cannot have $\Delta(r) = r \Delta_2(r^2)$,
for otherwise $\Delta_2(r^2)$ would be a unit in $\Omega_\mu$ and
would have to
%CONT>>


% ===== scan leaf 765, printed 744 =====
% !PRINTED 744
\weberpage{744}
%<<CONT
be, by \S.\,178, 1., the product of a real unit and a root of
unity $r^{2\lambda}$; and since $\Delta(r)$ is real as well,
$r^{2\lambda+1}$ would have to be real, which is not
possible. Accordingly $\Delta(r)$ is contained in $H_\mu$.

Since now $\mu$, just like $m$, can be any power of 2 whatever,
there results from this the second main theorem:

B. \gesperrt{A unit of the field $H_m$ which, together with all
its conjugates, is positive, is the square of a unit of the
same field.}

But in \S.\,184 we made the proof of the principal theorem on
the \gesperrt{Abel}ian number fields (\S.\,179, I.) depend upon the two
theorems A. and B.

\begin{center}
\rule{4cm}{0.4pt}
\end{center}


% ===== scan leaf 766, printed 745 =====
\weberpage{745}% page number not printed on this page
\pdfbookmark[1]{Twenty-fifth Section. Transcendental Numbers.}{bm231}
\pdfbookmark[2]{§ 203. Countable sets}{bm232}
% !PRINTED none
\weberpage{none}

\begin{center}
\gesperrt{Twenty-fifth Chapter.}

\gesperrt{Transcendental Numbers.}
\end{center}

\begin{center}
\rule{4cm}{0.4pt}
\end{center}

\begin{center}
\S.\,203.

\gesperrt{Enumerable Sets.}
\end{center}

In the introduction to our work the general concept of number was defined, and it was with this general concept of number that we first operated, for example in the proof of the existence of roots. In the further course of our considerations we then occupied ourselves only with \gesperrt{algebraic numbers}, without giving ourselves any account of whether the extent of the domain of numbers is thereby exhausted, or whether there are also non-algebraic numbers. The existence of non-algebraic numbers, which are also called transcendental numbers, was first proved by \gesperrt{Liouville}. Other proofs of this theorem are due to \gesperrt{G. Cantor}\footnote{\gesperrt{G. Cantor}, Crelle's Journal, Bd. 77 (1873). Ueber eine Eigenschaft des Inbegriffes aller reellen algebraischen Zahlen.}.

Here we take up the concept of a \gesperrt{set} or \gesperrt{manifold} already set forth in the introduction, by which is to be understood a system of elements of any kind whatever, so precisely delimited that of every arbitrary object it is completely decided whether it belongs to the system or not.

We distinguish finite and infinite sets, and adduce as the first and most important example of an infinite set the totality of the \gesperrt{natural numbers} $1, 2, 3, \ldots$. Then the following definition holds:

1. \textbf{Definition.} \gesperrt{A set is called enumerable when its elements can be placed in a mutually}
%CONT>>


% ===== scan leaf 767, printed 746 =====
% !PRINTED 746
\weberpage{746}
%<<CONT
\gesperrt{one-to-one relation with the natural sequence of numbers or with a part of it}\footnote{\gesperrt{Cantor}, l.\,c. The concept of enumerable sets coincides with the concept of \gesperrt{simply infinite systems}, which \gesperrt{Dedekind} defined without the \gesperrt{assumption of the system of numbers}. (\gesperrt{Dedekind}, Was sind und was sollen die Zahlen? \S.\,6. Braunschweig 1887. Zweite unveränderte Auflage 1893.)}.

Accordingly every finite set is enumerable, and in the enumeration the natural sequence of numbers is used only up to a certain highest number. In what follows we shall consider chiefly infinite sets.

We may also give the definition of an infinite enumerable set the following expression: it is a set in which to every element a definite number of the natural sequence of numbers can be assigned as a name, in such a way that every number of the sequence of numbers is thereby used as well; that is, a set in which there is a first, a second, a third \ldots a hundredth \ldots element.

Finally we may also say that an infinite enumerable set is one that can be ordered in a sequence in such a way that a \gesperrt{first} element is present, and that upon every element there follows a \gesperrt{definite other element of the set}, and that every element, with the exception of the first, is preceded by a \gesperrt{definite other element}. Such a sequence we may call a countable arrangement.

It is clear that an enumerable set is enumerable not merely in one way, but in infinitely many different ways.

Besides the natural sequence of numbers itself, which is certainly enumerable, we may adduce as a second example the system of rational positive proper fractions, which, among other ways, can be enumerated in the following manner:
\[
\frac{1}{2}, \frac{1}{3}, \frac{2}{3}, \frac{1}{4}, \frac{3}{4}, \frac{1}{5}, \frac{2}{5}, \frac{3}{5}, \frac{4}{5}, \frac{1}{6}, \frac{5}{6}, \ldots
\]
i.\,e. in such a way that every larger denominator follows the smaller denominator, and that the fractions with the same denominator are ordered according to the size of the numerator. If, however, one wished to order the fractions according to the size of their numerical value, one would obtain no countable arrangement.

2. \gesperrt{The totality of all algebraic numbers is an enumerable set.}


% ===== scan leaf 768, printed 747 =====
% !PRINTED 747
\weberpage{747}
In order to prove this important theorem, we recall that every algebraic number $\Theta$ is the root of one and only one irreducible equation
\weq{1}{f(\Theta) = a_0 \Theta^n + a_1 \Theta^{n-1} + \cdots + a_n = 0}
in which $a_0, a_1, \ldots, a_n$ are integral rational numbers without a common divisor, and $a_0$ is different from zero and positive. The degree $n$ of equation (1) is a positive integer, and therefore at least $=1$. The numbers $a_1, \ldots, a_{n-1}$ may in part be zero.

We now wish to determine the signs $\pm$ in such a way that $\pm a_1, \pm a_2, \pm a_3, \ldots, \pm a_n$ are not negative, and we call the sum
\weq{2}{N = (n-1) + a_0 \pm a_1 \pm a_2 \cdots \pm a_n}
the \gesperrt{height} of the algebraic number $\Theta$. The height is then always a \gesperrt{positive} integer.

Now it is easy to see that for a given value of the height $N$ there can always be only a finite number of algebraic numbers. For, in the first place, by (2) $n$ can never be greater than $N$, and for each given $N$ and $n$ the numbers $a_0, a_1, \ldots, a_n$ can be determined in only a finite number of ways. One has then to retain, among the functions $f(\Theta)$ so determined, only the irreducible ones. If we now order the algebraic numbers in such a way that we place the numbers of smaller height before those of greater height, that among numbers of equal height we place first those whose real part is smaller, and that among numbers of equal height and equal real part we let those of smaller imaginary part precede, then we have a countable arrangement of the algebraic numbers, and it is proved that the totality of all algebraic numbers is an enumerable set.

We obtain, for example:
\[
\begin{array}{llll}
N=1, & n=1, & a_0=1, & a_1=0\\
N=2, & n=1, & a_0=1, & a_1=\pm 1\\
N=3, & n=1, & a_0=1, & a_1=\pm 2\\
     &      & a_0=2, & a_1=\pm 1\\
     & n=2, & a_0=1, & a_1=0, \quad a_2=1,
\end{array}
\]
and the beginning of the ordered sequence of the algebraic numbers therefore becomes:
\[
0, -1, +1, -2, -\tfrac{1}{2}, -\sqrt{-1}, \sqrt{-1}, \tfrac{1}{2}, 2, \ldots
\]


% ===== scan leaf 769, printed 748 =====
% !PRINTED 748
\weberpage{748}
\pdfbookmark[2]{§ 204. Uncountable sets}{bm233}
Every part of the natural sequence of numbers is an enumerable set; for one need only order the numbers of such a part according to their magnitude in order to obtain a countable arrangement.

From this, however, it follows that every part of an enumerable set is itself an enumerable set. It follows from this, among other things, that the set of the \gesperrt{real algebraic} numbers is also enumerable.

\begin{center}
\S.\,204.

\gesperrt{Uncountable Sets.}
\end{center}

We now come, secondly, to the proof of the theorem that among the sets of numbers there are also non-enumerable ones, and we shall show specifically:

\begin{center}
\gesperrt{That the totality of all real numbers, even if we restrict ourselves to a finite interval, is not enumerable.}
\end{center}

For this purpose we consider any enumerable set of real numbers different from one another, which, set into a countable arrangement $\Omega$, we shall denote as follows:
\weq{$\Omega$}{\omega_1, \omega_2, \omega_3, \omega_4, \ldots}
(where it is to be recalled that this is not meant to signify a succession according to magnitude).

For brevity we shall call, of two elements of the sequence $\Omega$, the one with the smaller index the \gesperrt{earlier}, the one with the greater the \gesperrt{later}.

We now assume any two real numbers $\alpha, \beta$ such that $\alpha < \beta$, and show that in the interval $\delta = (\alpha, \beta)$ there is at least one number which does not occur in the sequence $\Omega$. Once we have demonstrated such a number for every interval, there are also infinitely many of them, since the same mode of reasoning can be applied to every part of the interval $(\alpha, \beta)$.

First of all it is clear that our assertion is correct if in any finite part of the interval $\delta$ only a finite number of numbers of the sequence $\Omega$ is contained, and we may therefore pass at once to the case in which in every part of the interval $\delta$, however small, there lies an infinite set of numbers of the sequence $\Omega$.


% ===== scan leaf 770, printed 749 =====
% !PRINTED 749
\weberpage{749}

% [sic] -- the source uses $\delta_1$ and $\delta_2$ both for scalar lengths and for intervals; the printed symbols are preserved.
We denote by $\alpha_1$, $\beta_1$ the two \gesperrt{earliest} numbers of the sequence $\Omega$ that lie in the interval $\delta$, we assume $\alpha_1 < \beta_1$, and we set $\delta_1 = \beta_1 - \alpha_1$, so that $\delta_1 = (\alpha_1, \beta_1)$ is a part of the interval $\delta$.

Now in the same way we denote by $\alpha_2$, $\beta_2$ the two earliest numbers of $\Omega$ that lie \gesperrt{in the interior} of the interval $\delta_1$ (with the exclusion of the limits), we presuppose $\alpha_2 < \beta_2$, and we set $\delta_2 = \beta_2 - \alpha_2$.

In this way we can continue, and we obtain an unlimited sequence of intervals:
\[
\delta,\ \delta_1,\ \delta_2,\ \delta_3,\ \dots,
\]
each of which includes all the following ones, and two sequences of numbers:
\[
\begin{gathered}
\alpha,\ \alpha_1,\ \alpha_2,\ \dots \\
\beta,\ \beta_1,\ \beta_2,\ \dots
\end{gathered}
\]
which, with the possible exception of the first, $\alpha$ and $\beta$, all belong to the sequence $\Omega$. The $\alpha, \alpha_1, \alpha_2, \dots$ form an increasing, the $\beta, \beta_1, \beta_2, \dots$ a decreasing sequence of numbers, and at the same time every $\alpha$ is smaller than every $\beta$.

1. \gesperrt{From this it follows that the numbers $\alpha_\nu$ have an upper limit $a$, and the numbers $\beta_\nu$ a lower limit $b$, and that $a$ is in any case not greater than $b$.} (Cf. Bd. I, \S.\,38, 1.)

But it is possible that $a = b$.

From the mode of formation of the intervals $\delta, \delta_1, \delta_2, \dots$ the following further results:

2. \gesperrt{If any number $\omega$ of the sequence $\Omega$ lies in the interval $\delta_\nu$, with the exclusion of its limits $\alpha_\nu, \beta_\nu$, then $\omega$ stands later in the sequence $\Omega$ than the pair of numbers $\alpha_\nu, \beta_\nu$ belonging to that same sequence.}

For $\alpha_\nu, \beta_\nu$ were indeed the two \gesperrt{earliest} numbers of $\Omega$ lying in the interval $\delta_{\nu-1}$. From this it follows:

3. \gesperrt{The pairs of numbers $\alpha_\nu, \beta_\nu$ belonging to the sequence $\Omega$ are the later terms of the sequence $\Omega$ the greater the index $\nu$ is, and since we have assumed that the sequence of intervals $\delta_\nu$ does not break off, we can, for a sufficiently large $\nu$, let $\alpha_\nu, \beta_\nu$ move out arbitrarily far in the sequence $\Omega$.}


% ===== scan leaf 771, printed 750 =====
% !PRINTED 750
\weberpage{750}

Now it follows very simply that no number $g$ that coincides with one of the numbers $a, b$, or that lies between them when $a$ and $b$ are different, can belong to the sequence $\Omega$.

% [sic] -- the source prints singular ``kommt'' with the pair $\alpha_\nu,\beta_\nu$; English renders the collective meaning grammatically.
For the number $g$ lies in the interior of every one of the intervals $\delta_\nu$. Suppose that $g$ occurs in $\Omega$; if, following 3., we go so far with $\nu$ that the pair $\alpha_\nu, \beta_\nu$ comes later than $g$ in $\Omega$, then by 2. $g$ cannot lie in the interval $\delta_\nu$, and with that the impossibility of our assumption is proved.

\gesperrt{From this it follows, then, that the totality of the numbers of an interval $\alpha, \beta$ does not form an enumerable set.}

\begin{center}
\rule{4cm}{0.4pt}
\end{center}

This fact can be proved in yet another way, which is in a certain respect still simpler, and which we will set out in a few words. We do not restrict the generality if we take the interval from 0 to 1 as our basis. We think of all the numbers of this interval as represented by infinite decimal fractions. The finite decimal fractions are contained among them as well, if we set all figures, from a certain one on, equal to zero. In order to make the representation by decimal fractions a unique one, let it further be stipulated that for a finite decimal fraction this representation shall \gesperrt{always} be chosen, so that, e.\,g., 0,4999\dots\ shall not be put for 0,5000\dots.

We will now assume that these decimal fractions form an enumerable set; they can therefore be arranged in a countable sequence, which we represent thus:
\weq{$\Omega$}{
\begin{aligned}
\omega_1 &= 0,\ \alpha_1^{(1)}\, \alpha_2^{(1)}\, \alpha_3^{(1)}\, \dots\\
\omega_2 &= 0,\ \alpha_1^{(2)}\, \alpha_2^{(2)}\, \alpha_3^{(2)}\, \dots\\
\omega_3 &= 0,\ \alpha_1^{(3)}\, \alpha_2^{(3)}\, \alpha_3^{(3)}\, \dots\\
&\quad\cdots\cdots\cdots\cdots\cdots
\end{aligned}
}
where the $\alpha_\mu^{(\nu)}$ denote figures of the decadic system.

But it is now very easy to exhibit one decimal fraction (or indeed arbitrarily many) not contained in the sequence $\Omega$. We need only form
\[
\eta = 0,\ \beta_1\ \beta_2\ \beta_3\ \dots
\]
% German splits 'ver-|schieden' across the leaf but keeps $\alpha_\nu^{(\nu)}$ on this
% page. English cannot split 'dif-|ferent' and still carry 'from alpha' here, so the whole
% clause is absorbed onto this page; i0772 opens at the next sentence.
% [sic] -- the printed condition permits a recurring-9 expansion equivalent to an earlier terminating decimal; one repair is to choose $\beta_\nu\in\{1,2\}$ unequal to $\alpha_\nu^{(\nu)}$.
where the $\beta_\nu$ are figures of the decadic system that satisfy the single condition that $\beta_\nu$ shall for every $\nu$ differ from $\alpha_\nu^{(\nu)}$.
%CONT>>


% ===== scan leaf 772, printed 751 =====
% !PRINTED 751
\weberpage{751}
\pdfbookmark[2]{§ 205. Transcendence of the number e}{bm234}

%<<CONT
This number $\eta$, which after all likewise belongs to the interval $(0, 1)$, can agree with no number of the sequence $\Omega$.

The formation of $\eta$ can be generalised still further by choosing the initial $\beta_\nu$ arbitrarily through any prescribed finite index, and imposing only thereafter the law that $\beta_\nu$ differ from $\alpha_\nu^{(\nu)}$.

Since it has now been proved that the real algebraic numbers form an enumerable set, and that in every interval there are numbers which do not belong to a given enumerable set, there follows from this quite immediately:

\begin{center}
\gesperrt{In every real interval there are transcendental numbers.}
\end{center}

\begin{center}
\S.\,205.

\gesperrt{Transcendence of the number $e$.}
\end{center}

It is now a far more difficult task to decide, of a definitely given number, whether it is algebraic or transcendental. Here the interest has concentrated chiefly on the two numbers that occur so frequently in analysis: $e$, that is, the base of the natural system of logarithms, and the \gesperrt{Ludolphian} number $\pi$, the ratio of the circumference of the circle to the diameter.

For the number $e$ the question was decided by \gesperrt{Hermite} in a celebrated memoir, which became the foundation for the later investigations on the number $\pi$$^{1)}$. The decision for the number $\pi$, however, which was quite particularly interesting on account of its connection with the long-celebrated problem of the quadrature of the circle, still presented insuperable difficulties for a long time. Finally the demonstration was furnished by \gesperrt{Lindemann} that $\pi$ too belongs to the transcendental numbers. But the proof given by \gesperrt{Lindemann} at first still presented great difficulties to the understanding, which through later investigations by \gesperrt{Weierstrass}, \gesperrt{Hilbert}, \gesperrt{Hurwitz} and \gesperrt{Gordan}$^{2)}$

\noindent\rule{4cm}{0.4pt}

$^{1)}$ \gesperrt{Hermite}, Sur la fonction exponentielle. Comptes rendus, T.\,LXXVII, 1873.

$^{2)}$ \gesperrt{Lindemann}, Ueber die Zahl $\pi$. Mathem. Annalen, Bd. 20, 1882. \gesperrt{Weierstrass}, Zu \gesperrt{Lindemann}'s Abhandlung \glqq Ueber die \gesperrt{Ludolph}'sche
%CONT>>


% ===== scan leaf 773, printed 752 =====
% !PRINTED 752
\weberpage{752}

%<<CONT
have gradually been so completely removed that the proof can now be carried out by wholly elementary means and in the simplest way.

If, as already before, we denote by $\Pi(n)$ the product
\[
\Pi(n) = 1.2.3\dots n
\]
then, if $x$ is an arbitrary quantity,
\weq{1}{\frac{x^n}{\Pi(n)}}
will approach the limit zero as $n$ increases indefinitely, and indeed more strongly than the terms of a decreasing geometric series.

For if $k$ is an integer greater than the absolute value of $x$, then the absolute value of $x : h$ is a proper fraction whenever $h$ is equal to or greater than $k$. Consequently
\[
\left|\frac{x^n}{\Pi(n)}\right| = \left|\frac{x^k}{\Pi(k)}\, \frac{x}{k+1}\, \frac{x}{k+2} \cdots \frac{x}{n}\right| < \left|\frac{x^k}{\Pi(k)}\right| \left|\frac{x}{k}\right|^{n-k}.
\]
% [sic] -- the strict inequality tacitly assumes $n>k$ and $x\ne0$; $x=0$ is the trivial equality and convergence case.
From this it follows that the infinite series
\weq{2}{e^x = 1 + x + \frac{x^2}{\Pi(2)} + \frac{x^3}{\Pi(3)} + \cdots}
converges for all values of $x$, and by it we define the exponential function $e^x$, whose simplest properties we here assume from analysis. In particular, for any two values $x, y$ the relation
\[
e^{x+y} = e^x e^y,
\]
holds, from which it then follows that, for a positive integer $n$, $e^n$ is the $n$th power of the number
\weq{3}{e = 2 + \frac{1}{\Pi(2)} + \frac{1}{\Pi(3)} + \frac{1}{\Pi(4)} + \cdots = 2{,}718281828459\cdots\,.}

If now $r$ denotes any positive integer, we can represent formula (2) in the following way:
\weq{4}{\Pi(r)\, e^x = \Pi(r) + \frac{\Pi(r)}{\Pi(1)} x + \frac{\Pi(r)}{\Pi(2)} x^2 + \cdots + x^r + x^r\, U_r,}

\noindent\rule{4cm}{0.4pt}

Zahl\grqq{}. Sitzungsbericht der Berliner Akademie, 3. December 1885. The papers by \gesperrt{Hilbert}, \gesperrt{Hurwitz} and \gesperrt{Gordan} are all three to be found in volume 43 of Mathematische Annalen (1893), the first two also in the G\"ottinger Nachrichten of 1893.
%CONT>>


% ===== scan leaf 774, printed 753 =====
% !PRINTED 753
\weberpage{753}
%<<CONT
in which

\weq{5}{U_r = \frac{x}{r+1}+\frac{x^2}{(r+1)(r+2)}+\frac{x^3}{(r+1)(r+2)(r+3)}+\cdots,}

and this formula holds also for $r=0$, if one assumes $\Pi(0)=1$.

If we denote the absolute value of $x$ by $\xi$, then, by the theorem proved in the Introduction that the absolute value of a sum is never greater than the sum of the absolute values of its parts, the absolute value of $U_r$ is certainly not greater than
\[
\frac{\xi}{r+1}+\frac{\xi^2}{(r+1)(r+2)}+\frac{\xi^3}{(r+1)(r+2)(r+3)}+\cdots,
\]
and hence smaller than
\[
1+\frac{\xi}{1}+\frac{\xi^2}{1.2}+\frac{\xi^3}{1.2.3}+\cdots=e^{\xi}.
\]
If we therefore set
\[
U_r = q_r\,e^{\xi}
\]
then $q_r$ is a function of $x$, of which we establish only this much, that its absolute value is smaller than 1.

Accordingly we set up formula (4) for $r=0,1,\ldots,n$, and write the terms, arranged in reverse order, in a clear survey thus:
\[
\begin{aligned}
\Pi(n)\,e^x &= q_n\,x^n e^{\xi} + x^n + n\,x^{n-1} + n(n-1)x^{n-2}+\cdots\\
&\qquad{}+\Pi(n)\\[4pt]
(6)\qquad \Pi(n-1)\,e^x &= q_{n-1}\,x^{n-1}e^{\xi}+x^{n-1}+(n-1)x^{n-2}\\
&\qquad{}+(n-1)(n-2)x^{n-3}+\cdots\\[4pt]
&\hphantom{\Pi(n)\,e^x{}}\;\cdots\cdots\cdots\cdots\cdots\cdots\cdots\cdots\cdots\cdots\cdots\cdots\cdots\\[4pt]
e^x &= q_0\,e^{\xi}+1.
\end{aligned}
\]

Now we take an arbitrary integral function of degree $n$ in $x$:

\weq{7}{f(x) = c_n x^n + c_{n-1}x^{n-1} + c_{n-2}x^{n-2}+\cdots+c_0,}

and its derivatives
\[
\begin{aligned}
f'(x) &= n\,c_n x^{n-1} + (n-1)c_{n-1}x^{n-2}+\cdots\\[4pt]
(8)\qquad f''(x) &= n(n-1)c_n x^{n-2}+(n-1)(n-2)c_{n-1}x^{n-3}+\cdots\\[4pt]
&\hphantom{f''(x){}}\;\cdots\cdots\cdots\cdots\cdots\cdots\cdots\cdots\cdots\cdots\cdots\cdots\cdots\\[4pt]
f^{(n)}(x) &= \Pi(n)\,c_n,
\end{aligned}
\]

and we set further, for abbreviation,

\weq{9}{F(x) = f(x)+f'(x)+f''(x)+\cdots+f^{(n)}(x).}
% signature "Weber, Algebra. II." and gathering-mark "48" at foot -- not text, omitted


% ===== scan leaf 775, printed 754 =====
% !PRINTED 754
\weberpage{754}
Then $F(x)$ is an integral function of $x$ of degree $n$. If finally we also set

\weq{10}{Q(x) = c_n q_n x^n + c_{n-1}q_{n-1}x^{n-1}+\cdots+c_0 q_0,}

\weq{11}{P = c_n\Pi(n)+c_{n-1}\Pi(n-1)+\cdots+c_0,}

then $Q(x)$ depends on $x$, though it is not expressible rationally in terms of $x$; $P$, however, is independent of $x$.

If we now multiply equations (6) in succession by $c_n, c_{n-1},\ldots, c_0$ and add them, it follows that

\weq{12}{e^x P = F(x) + e^{\xi}Q(x).}

This formula is now the starting point of the further conclusions.

Let us assume that $e$ is an algebraic number; then an equation must subsist, whose degree let us call $m$:

\weq{13}{C_0 + C_1 e + C_2 e^2 + \cdots + C_m e^m = 0,}

whose coefficients $C_0, C_1, \ldots, C_m$ are \gesperrt{integral rational numbers}, among which $C_0$ and $C_m$ are different from zero. The matter is to demonstrate the impossibility of this assumption.

For this purpose we substitute in equation (12) for $x$ in succession the integers $0,1,2,\ldots,m$, so that $\xi$ becomes identical with $x$, multiply by $C_0, C_1, \ldots, C_m$ and add. Then, by (13), since $P$ is independent of $x$, there results:

\weq{14}{0 = C_0 F(0) + C_1 F(1) + \cdots + C_m F(m) + C_0 Q(0) + C_1 e\,Q(1) + \cdots + C_m e^m Q(m),}

and now it is to be shown, from a suitable assumption about the still arbitrary function $f(x)$, that equation (14) is impossible.

We choose a prime $p$ which is greater than $m$\,\footnote{That there is always a prime $p$ greater than an arbitrary number $\mu$ was already proved in \gesperrt{Euclid} (Elemente, Buch IX, Nr. XX, Bd. 2 of the \gesperrt{Heiberg} edition). The proof is simply this, that the integer $\Pi(\mu)+1$, which is evidently greater than $\mu$, is divisible by no prime that is not greater than $\mu$, because on division by each of the numbers $2,3,\ldots,\mu$ this number gives the residue 1. It is therefore impossible that no prime lies above $\mu$.}, and set

\weq{15}{f(x) = \frac{x^{p-1}(1-x)^p(2-x)^p\ldots(m-x)^p}{\Pi(p-1)},}
%CONT>>


% ===== scan leaf 776, printed 755 =====
% !PRINTED 755
\weberpage{755}
%<<CONT
so that the degree $n$ of $f(x)$ is equal to $(m+1)p-1$, and we now prove two things:

1) $C_0 F(0) + C_1 F(1) + \cdots + C_m F(m)$ is an integer different from zero, and therefore, apart from the sign, at least equal to 1,

2) $C_0 Q(0) + C_1 e\,Q(1) + \cdots + C_m e^m Q(m)$ is less than 1,

both under the assumption that a suitable choice is made for $p$. If both of these are proved, one recognises the impossibility of equation (14), and hence that of equation (13), from which (14) had been inferred.

If we arrange the numerator of $f(x)$ according to powers of $x$, an expression of the following form results:

\weq{16}{f(x) = \frac{A_{p-1}x^{p-1}+A_p x^p + A_{p+1}x^{p+1}+\cdots}{\Pi(p-1)},}

in which $A_{p-1}, A_p, A_{p+1}, \ldots$ are integers, and $A_{p-1}=[\Pi(m)]^p$, hence certainly not divisible by $p$. Accordingly, if one compares (16) with the \gesperrt{Taylor} expansion
\[
f(x) = f(0) + xf'(0) + \frac{x^2}{\Pi(2)}f''(0)+\cdots
\]
one finds
\[
f(0)=f'(0)=f''(0)=\cdots=f^{(p-2)}(0)=0,
\]
\[
f^{(p-1)}(0)=A_{p-1},\quad f^{(p)}(0)=pA_p,\quad f^{(p+1)}(0)=p(p+1)A_{p+1}\cdots,
\]
and therefore
\[
F(0) = A_{p-1}+pA_{p+1}+p(p+1)A_{p+1}+\cdots, % [sic] -- printed with A_{p+1} repeated where the pattern would expect A_p, A_{p+1}
\]
an \gesperrt{integer not divisible by $p$}. If, however, we arrange $f(x)$ according to powers of $x-1$, there follows:
\[
f(x) = \frac{B_p(x-1)^p+B_{p+1}(x-1)^{p+1}+\cdots}{\Pi(p-1)},
\]
in which the $B_p, B_{p+1}, \ldots$ are again integers. From this there follows, as above, by comparison with
\[
f(x) = f(1) + (x-1)f'(1) + \frac{(x-1)^2}{\Pi(2)}f''(1)+\cdots
\]
\[
F(1) = pB_p + p(p+1)B_{p+1}+\cdots,
\]
consequently $F(1)$ is an \gesperrt{integer divisible by $p$}, and in exactly the same way it follows that $F(2), F(3), \ldots, F(m)$ are integers divisible by $p$. Since $p$ can now also be chosen so large that $C_0$ is not divisible by $p$, it follows
%CONT>>
% signature "48*" at foot -- not text, omitted


% ===== scan leaf 777, printed 756 =====
% !PRINTED 756
\weberpage{756}
\pdfbookmark[2]{§ 206. Transcendence of the number π}{bm235}
%<<CONT
that $C_0 F(0) + C_1 F(1) + \cdots + C_m F(m)$ is an integer not divisible by $p$, and therefore also not vanishing, and with this 1) is proved.

If we can now show further that $Q(x)$ can be made arbitrarily small for every positive $x$ by increasing $p$, it follows that for sufficiently large $p$ the expression 2) is certainly less than 1, and our proof is complete.

But if we go back to the expression (10) for $Q(x)$ and denote by $\gamma_n, \gamma_{n-1}, \ldots, \gamma_0$ the absolute values of the coefficients $c_n, c_{n-1}, \ldots, c_0$ of $f(x)$, then, since the $q_n, q_{n-1}, \ldots$ are less than 1 in absolute value, we see that for every positive $x$ the absolute value of $Q(x)$ is less than

\weq{17}{\psi(x) = \gamma_n x^n + \gamma_{n-1}x^{n-1}+\cdots+\gamma_0.}

But the coefficients $c_n, c_{n-1}, \ldots, c_0$ of the function $f(x)$ in (15) differ from the $\gamma_n, \gamma_{n-1}, \ldots, \gamma_0$ only in sign. If we therefore replace $x$ by $-x$, and so form the function
\[
f(-x) = \frac{x^{p-1}(1+x)^p(2+x)^p\ldots(m+x)^p}{\Pi(p-1)},
\]
then this function has the same coefficients as $f(x)$, but all with positive signs. It is therefore none other than the function $\psi(x)$.

If we then also set
\[
X = x(1+x)(2+x)\ldots(m+x),
\]
we obtain

\weq{18}{\psi(x) = \frac{X}{x}\cdot\frac{X^{p-1}}{\Pi(p-1)},}

and this, as was shown at the beginning of this section, approaches the limit \gesperrt{zero} as $p$ increases without bound.

\gesperrt{Thus $e$ is a transcendental number.}

\begin{center}
§. 206.

\gesperrt{Transcendence of the number $\pi$.}
\end{center}

By the same means the transcendence of the number $\pi$ can now also be proved. As the definition of this number we take the fact that it is the least positive number which satisfies the equation

\weq{1}{e^{i\pi} = -1}

when $e^x$ is defined by the formula \S.\,205, (2).


% ===== scan leaf 778, printed 757 =====
% !PRINTED 757
\weberpage{757}

Let us therefore assume that $\pi$, and consequently also $i\pi$, is an algebraic number; then $i\pi$ is one of the roots of an irreducible rational equation $\chi(x)=0$ whose coefficients are integral rational numbers.

If we denote all the roots of this algebraic equation by $\beta_1,\beta_2,\ldots,\beta_\nu$, and denote the coefficient of $x^\nu$ in $\chi$ by $a$, then
\weq{2}{\chi(x) = a(x-\beta_1)(x-\beta_2)\ldots(x-\beta_\nu),}
and the products $a\beta_1,a\beta_2,\ldots,a\beta_\nu$ are \gesperrt{integral} algebraic numbers, and hence their integral symmetric functions are \gesperrt{integral rational numbers} (\S.\,133).

Since the number $i\pi$ is to occur among the $\beta$, we have by (1):
\[
(1+e^{\beta_1})(1+e^{\beta_2})\ldots(1+e^{\beta_\nu}) = 0,
\]
and if we carry out the multiplication, an equation results:
\[
1 + \sum e^{\beta_i} + \sum e^{\beta_i+\beta_h} + \sum e^{\beta_i+\beta_h+\beta_k} + \cdots = 0.
\]
Among the exponents in this sum the number zero may occur several times; let us assume $(C-1)$ times, so that $C$ is a positive integer, at least $=1$. The remaining exponents $\beta_i,\beta_i+\beta_h,\beta_i+\beta_h+\beta_k,\ldots$, some of which may also be equal to one another, we shall denote by $\alpha_1,\alpha_2,\ldots,\alpha_\mu$, so that the equation holds:
\weq{3}{C + e^{\alpha_1} + e^{\alpha_2} + \cdots + e^{\alpha_\mu} = 0.}
Here the $\alpha_1,\alpha_2,\ldots,\alpha_\mu$ are algebraic numbers which, multiplied by the integral rational number $a$, pass over into integral algebraic numbers.

But the symmetric functions of all the sums $\beta_i,\beta_i+\beta_h,\beta_i+\beta_h+\beta_k,\ldots$ are at the same time symmetric functions of the $\beta_1,\ldots,\beta_\nu$, and consequently rational numbers. These sums are therefore also the roots of a rational equation, and since the root $0$ can be separated off as often as it is present, the $\alpha_1,\alpha_2,\ldots,\alpha_\mu$ are likewise the roots of a rational equation; the symmetric fundamental functions of the $a\alpha_1,a\alpha_2,\ldots,a\alpha_\mu$ are \gesperrt{integral rational numbers}.

The absolute values of the numbers
\[
\alpha_1,\alpha_2,\ldots,\alpha_\mu
\]
shall now be denoted by
\[
a_1,a_2,\ldots,a_\mu
\]
respectively.


% ===== scan leaf 779, printed 758 =====
% !PRINTED 758
\weberpage{758}

If we now set $x=0,\alpha_1,\alpha_2,\ldots,\alpha_\mu$ in equation (12) of the preceding section and apply equation (3), we obtain:
\weq{4}{0 = CF(0) + F(\alpha_1) + F(\alpha_2) + \cdots + F(\alpha_\mu) + CQ(0) + e^{a_1}Q(\alpha_1) + e^{a_2}Q(\alpha_2) + \cdots + e^{a_\mu}Q(\alpha_\mu),}
and we have to demonstrate the impossibility of such an equation for a suitable choice of $f(x)$.

Let $p$ again be a prime determined to increase without limit, and
\weq{5}{f(x) = \dfrac{a^{\mu p+p-1}x^{p-1}(x-\alpha_1)^p(x-\alpha_2)^p\ldots(x-\alpha_\mu)^p}{\Pi(p-1)},}
so that $f(x)$ is a function with rational coefficients.

We now form, as in the preceding section, by arranging according to powers of $x$:
\[
f(x) = \dfrac{A_{p-1}x^{p-1} + A_px^p + A_{p+1}x^{p+1}\ldots}{\Pi(p-1)},
\]
in which the $A_{p-1},A_p,\ldots$ are integral rational numbers, and
\[
A_{p-1} = (-1)^\mu a^{\mu p+p-1}\alpha_1^p\alpha_2^p\ldots\alpha_\mu^p.
\]
If we therefore assume $p$ to be greater than each of the integers
\[
a, a^\mu\alpha_1\alpha_2\ldots\alpha_\mu
\]
then $A_{p-1}$ is not divisible by $p$, and we have
\[
F(0) = A_{p-1} + pA_p + p(p+1)A_{p+1} + \cdots,
\]
i.\,e. an integer not divisible by $p$. Likewise, if $p$ is taken large enough, $C$ is not divisible by $p$.

On the other hand, we arrange the numerator of $f(x)$ according to powers of $a(x-\alpha_1)$ and obtain
\[
f(x) = \dfrac{B_pa^p(x-\alpha_1)^p + B_{p+1}a^{p+1}(x-\alpha_1)^{p+1} + \cdots}{\Pi(p-1)},
\]
in which the $B_p,B_{p+1},\ldots$ are indeed no longer rational, but are nevertheless \gesperrt{integral algebraic numbers}, since the numerator of $f(x)$ is a function with integral coefficients in $ax,a\alpha_1,\ldots,a\alpha_\mu$.

From this it now follows, as in the preceding section:
\[
F(\alpha_1) = pB_pa^p + p(p+1)B_{p+1}a^{p+1} + \cdots
\]
If one forms $F(\alpha_2),\ldots,F(\alpha_\mu)$ in the same way, and observes that the sum $B_p(\alpha_1)+B_p(\alpha_2)+\cdots+B_p(\alpha_\mu)$ is an integral rational number, it follows that
\[
F(\alpha_1) + F(\alpha_2) + \cdots + F(\alpha_\mu)
\]
%CONT>>


% ===== scan leaf 780, printed 759 =====
% !PRINTED 759
\weberpage{759}
%<<CONT
is an integral rational number divisible by $p$. From this it follows finally that
\[
CF(0) + F(\alpha_1) + F(\alpha_2) + \cdots + F(\alpha_\mu)
\]
is an integral rational number not divisible by $p$, and therefore also not vanishing, which consequently, apart from its sign, is at least $=1$.

If we can now finally prove that, under the present assumption about $f$ as well, the function $Q(x)$ can be made arbitrarily small for every finite $x$ by making $p$ sufficiently large, then the impossibility of equation (4) follows exactly as above.

But this may be seen in the following way: instead of the function
\weq{6}{f(x) = c_nx^n + c_{n-1}x^{n-1} + \cdots + c_0}
we consider the function
\[
\begin{aligned}
\psi(x) &= \dfrac{a^{\mu p+p-1}x^{p-1}(x+a_1)^p(x+a_2)^p\ldots(x+a_\mu)^p}{\Pi(p-1)} \\
&= \gamma_nx^n + \gamma_{n-1}x^{n-1} + \cdots + \gamma_0,
\end{aligned}
\]
which now has none but positive (though not necessarily rational) coefficients. The coefficients $c_n,c_{n-1}\ldots$ are formed by multiplication and addition from the numbers $a,-\alpha_1,-\alpha_2,\ldots,-\alpha_\mu$, and the corresponding coefficients $\gamma_n,\gamma_{n-1}\ldots$ are obtained from them if the $-\alpha_1,-\alpha_2,\ldots,-\alpha_\mu$ are replaced by their absolute values $a_1,a_2,\ldots,a_\mu$; from which it follows, by the theorem of the Introduction already mentioned above, that the coefficients $\gamma_n,\gamma_{n-1},\ldots$ are certainly not smaller than the absolute values of the corresponding $c_n,c_{n-1}\ldots$

Now for every finite $x$ whose absolute value is $\xi$, the absolute value of $Q(x)$ is, by \S.\,205, (10), smaller than, or at least not greater than,
\[
\gamma_n\xi^n + \gamma_{n-1}\xi^{n-1} + \cdots + \gamma_0 = \psi(\xi),
\]
and that $\psi(\xi)$ sinks below every limit when $p$ becomes large enough is inferred from the representation:
\[
\psi(x) = \dfrac{X}{ax}\dfrac{X^{p-1}}{\Pi(p-1)},
\]
if
\[
% [sic] -- the source prints Latin $u+1$; the derivation from the exponent $\mu p+p-1$ requires $\mu+1$.
X = a^{u+1}x(x+a_1)(x+a_2)\ldots(x+a_\mu)
\]
is put.


% ===== scan leaf 781, printed 760 =====
% !PRINTED 760
\weberpage{760}
\pdfbookmark[2]{§ 207. The general theorem of Lindemann on the exponential function}{bm236}

With this it is proved:

\gesperrt{The number $\pi$ is a transcendental number. The ``quadrature of the circle'' cannot be solved by a geometric construction in which only algebraic curves and surfaces are employed.}

\begin{center}
\S.\,207.

\gesperrt{The general theorem of Lindemann on the exponential function.}
\end{center}

The transcendence of the numbers $e$ and $\pi$, which is hereby proved, is contained as a special case in a very general theorem on the exponential function that \gesperrt{Lindemann} announced in the memoir mentioned above, and of which a fully worked-out proof is contained in the cited memoir of \gesperrt{Weierstrass}. We will prove this theorem here in conclusion, using the same means that we have employed for the two special cases. The theorem runs as follows:

I. \gesperrt{There exists no equation of the form} \weq{1}{C_1e^{z_1} + C_2e^{z_2} + \cdots + C_me^{z_m} = 0,} \gesperrt{in which the coefficients $C_1,C_2,\ldots,C_m$ are algebraic numbers and the exponents $z_1,z_2,\ldots,z_m$ are algebraic numbers distinct from one another, unless all the coefficients $C_1,C_2,\ldots,C_m$ are equal to zero.}

In order to prove it, we first derive a lemma:

Let
\[
\alpha_1,\alpha_2,\ldots,\alpha_r
\]
be arbitrary real or imaginary quantities, \gesperrt{distinct however from one another}, and let
\[
A_1,A_2,\ldots,A_r
\]
be likewise arbitrary quantities that \gesperrt{do not all vanish}.

The same shall hold of the two sequences of quantities
\[
\begin{gathered}
\beta_1,\beta_2,\ldots,\beta_s \\
B_1,B_2,\ldots,B_s
\end{gathered}
\]
We denote by
\[
\gamma_1,\gamma_2,\ldots,\gamma_t
\]
%CONT>>


% ===== scan leaf 782, printed 761 =====
% !PRINTED 761
\weberpage{761}
%<<CONT
\gesperrt{those distinct from one another} among the $rs$ sums
$\alpha_i+\beta_k$, and set

\weq{2}{\begin{aligned} A &= A_1e^{\alpha_1}+A_2e^{\alpha_2}+\cdots+A_re^{\alpha_r}\\ B &= B_1e^{\beta_1}+B_2e^{\beta_2}+\cdots+B_se^{\beta_s}, \end{aligned}}

\weq{3}{AB = C_1e^{\gamma_1}+C_2e^{\gamma_2}+\cdots+C_te^{\gamma_r}.} % [sic] printed exponent index $\gamma_r$, inconsistent with $C_t$

The lemma to be proved then runs:

1. \gesperrt{The coefficients $C_1, C_2, \ldots, C_t$ cannot all vanish.}

In proving this theorem we may evidently assume that none of the coefficients $A_1, A_2, \ldots, A_r, B_1, B_2, \ldots, B_s$ vanishes (since we can simply leave out any that happen to vanish).

For the sake of a shorter mode of expression we shall, for the moment, when $a, b$ are two different complex numbers, call $a$ smaller than $b$, $(a<b)$, if the real part of $a$ is smaller than the real part of $b$, or if the real parts are equal and the imaginary part of $a$ is smaller than the imaginary part of $b$.

If then, in this sense, $a<b$, $b<c$, then also $a<c$, and if $a<b$, $c<d$, then $a+c<b+d$.

Among any finite sequence of complex numbers different from one another there is then a definite smallest one, and so, if $\alpha_1$ is the smallest among the numbers $\alpha$, and $\beta_1$ the smallest among the numbers $\beta$, then $\alpha_1+\beta_1$ is the smallest among the numbers $\gamma$, and this sum is equal to none of the other sums $\alpha_i+\beta_k$. Hence $C_1=A_1B_1$, and $C_1$ is different from zero, which was to be proved.

This theorem can now be generalised at once by complete induction.

2. \gesperrt{If}

\[
\begin{aligned}
A' &= A_1'e^{\alpha_1'}+A_2'e^{\alpha_2'}+\cdots\\
A'' &= A_1''e^{\alpha_1''}+A_2''e^{\alpha_2''}+\cdots\\
A''' &= A_1'''e^{\alpha_1'''}+A_2'''e^{\alpha_2'''}+\cdots\\
&\ \cdots\cdots\cdots\cdots\cdots\cdots
\end{aligned}
\]

\gesperrt{are sums of the form (2) in arbitrary number, and $\gamma_1,\gamma_2,\ldots$ those distinct from one another among the sums $\alpha_h'+\alpha_i''+\alpha_k'''+\cdots$, then in the product}

\weq{4}{A'A''A'''\cdots = C_1e^{\gamma_1}+C_2e^{\gamma_2}+\cdots}

\gesperrt{not all the coefficients $C_1, C_2, \ldots$ are equal to zero.}


% ===== scan leaf 783, printed 762 =====
% !PRINTED 762
\weberpage{762}
This lemma allows us first of all to make, in the proof of Theorem (1), the simplifying assumption that the coefficients $C_1, C_2, \ldots, C_m$ are integral rational numbers.

For suppose that an equation

\weq{5}{X_1e^{x_1}+X_2e^{x_2}+\cdots = 0}

holds, with algebraic coefficients $X_1, X_2, \ldots$ and with exponents $x_1, x_2, \ldots$ different from one another. Then we may always assume that the $X_1, X_2, \ldots$ belong to one and the same algebraic field $\Omega$.

We denote variables by $u_1, u_2, \ldots$ and form the norm of the linear form $X_1u_1+X_2u_2+\cdots$

\weq{6}{N(X_1u_1+X_2u_2+\cdots) = \Phi(u_1,u_2,\ldots),}

which is an integral homogeneous function of the variables $u$ with \gesperrt{rational coefficients}, and whose degree is equal to the degree of the field $\Omega$. If in (6) we set

\[
u_1 = e^{x_1},\ u_2 = e^{x_2},\ \ldots,
\]

then every product of variables $u$ passes over into a quantity of the form $e^z$, in which $z$ is a sum of numbers $x$, and $\Phi(u_1,u_2,\ldots)$ becomes an expression of the form $C_1e^{z_1}+C_2e^{z_2}+\cdots$ whose coefficients are rational numbers. By 2., these coefficients cannot all be zero, provided the $X_1, X_2, \ldots$ do not all vanish, and they cannot all be zero even when those terms of the function $\Phi(e^{x_1},e^{x_2},\ldots)$ which are equal to one another are collected into a single term $Ce^z$. But if now equation (5) holds, then $\Phi(e^{x_1},e^{x_2},\ldots)=0$ as well, and consequently

\[
C_1e^{z_1}+C_2e^{z_2}+\cdots = 0.
\]

If there are fractional numbers among the $C_1, C_2, \ldots$, we can convert them into integers by multiplying the whole equation by the common denominator.

3. \gesperrt{It therefore only remains to be proved that equation (1) can hold for no system of integral rational numbers $C_1, C_2, \ldots, C_m$ which do not all vanish.}

Accordingly we now assume that an equation

\weq{7}{A_1e^{x_1}+A_2e^{x_2}+\cdots = 0,}

holds whose coefficients $A_1, A_2, \ldots$ are integral rational numbers,
%CONT>>


% ===== scan leaf 784, printed 763 =====
% !PRINTED 763
\weberpage{763}
%<<CONT
which are not all $=0$, and in which the exponents $x_1, x_2, \ldots$ are algebraic numbers distinct from one another. We denote by $\Omega$ a \gesperrt{normal field} to which all these numbers $x_1, x_2, \ldots$ belong, and by $\Theta$ a primitive element of this field, and further by

\[
\sigma' = (\Theta,\Theta'),\ \sigma'' = (\Theta,\Theta''),\ \ldots
\]

the substitutions of the field $\Omega$. If by one of these substitutions, $\sigma'$, the numbers $x_1, x_2, \ldots$ pass into $x_1', x_2', \ldots$, then these too must be distinct from one another.

Let $u$ now be a variable, and

\weq{8}{U = A_1e^{ux_1}+A_2e^{ux_2}+\cdots}

In this we carry out all the substitutions $\sigma', \sigma'', \ldots$ and denote the functions so arising by $U', U'', \ldots$ The product of all these functions we may call the \gesperrt{norm} of $U$, although we are not dealing merely with algebraic numbers. This product has the form

\weq{9}{N(U) = C_1e^{uz_1}+C_2e^{uz_2}+\cdots,}

in which the $C_1, C_2, \ldots$ are likewise integral rational numbers. The $z_1, z_2, \ldots$ are numbers of the field $\Omega$, and if we combine the terms with equal exponents into a single term, we may assume them to be \gesperrt{distinct from one another}. At the same time, on account of (7),

\weq{10}{C_1e^{z_1}+C_2e^{z_2}+\cdots = 0.}

But the expression on the right-hand side of (9) has, as its origin as a norm shows, the further property of remaining unchanged when any one of the substitutions $\sigma', \sigma'', \ldots$ is carried out. If, however, (9) is expanded in powers of $u$, this invariance must hold for every term of the power-series expansion; therefore, if $h$ is any positive integer, it holds for

\[
C_1z_1^h+C_2z_2^h+\cdots;
\]

but this sum is a number of the field $\Omega$, and therefore a \gesperrt{rational number}. This we may summarise as follows:

If $g(z)$ denotes any integral function of $z$ with rational coefficients, then

\[
C_1g(z_1)+C_2g(z_2)+\cdots
\]

is a \gesperrt{rational number}.


% ===== scan leaf 785, printed 764 =====
% !PRINTED 764
\weberpage{764}
4. \gesperrt{The proof of Theorem I. is hereby reduced to the proof of the following special case: If there subsists an equation}

\weq{11}{C_1e^{z_1}+C_2e^{z_2}+\cdots+C_me^{z_m} = 0,}

\gesperrt{in which the $z_1, z_2, \ldots, z_m$ are algebraic numbers different from one another, in which $C_1, C_2, \ldots, C_m$ and the combinations}

\weq{12}{C_1g(z_1)+C_2g(z_2)+\cdots+C_mg(z_m)}

\gesperrt{are rational numbers, when $g(z)$ is any integral function with rational coefficients, then the $C_1, C_2, \ldots, C_m$ must all vanish.}

In order to prove this theorem, we first determine a positive integral rational number $c$ in such a way that

\weq{13}{x_1 = cz_1,\ x_2 = cz_2,\ \ldots,\ x_m = cz_m}

become integral algebraic numbers (§.\,133, 5.). We assume the coefficients $C_1, C_2, \ldots, C_m$ to be \gesperrt{integral} rational numbers, and form an integral function

\weq{14}{\varphi(x) = ax^r+a_1x^{r-1}+\cdots+a_r}

with integral rational numerical coefficients, which has the following properties:

1) The numbers $x_1, x_2, \ldots, x_m$ occur among the roots of $\varphi(x)=0$.

2) The sum

\[
C_1\varphi'(x_1)+C_2\varphi'(x_2)+\cdots+C_m\varphi'(x_m) = k,
\]

which by the assumptions (12), (13) is an integral rational number, is different from zero.

In order to see that such a function $\varphi(x)$ always exists, let one first take a function $\chi(x)$ which satisfies only the condition 1), and the second condition that none of the numbers $x_1, x_2, \ldots, x_m$ is a double root of $\chi(x)$; such a function evidently exists [in order to form a function $\chi$ one may, for example, free the norm of the product $(x-x_1)(x-x_2)\cdots(x-x_m)$ from common divisors with its derivatives]. If one then sets

\[
\varphi(x) = x^h\chi(x),\ \varphi'(x_1) = x_1^h\chi'(x_1),\ \varphi'(x_2) = x_2^h\chi'(x_2)\ldots,
\]
%CONT>>


% ===== scan leaf 786, printed 765 =====
% !PRINTED 765
\weberpage{765}
%<<CONT
then the sum

\[
C_1\,\varphi'(x_1)+\cdots+C_m\,\varphi'(x_m) = C_1\,\chi'(x_1)\,x_1^h+\cdots+C_m\,\chi'(x_m)\,x_m^h
\]

can certainly not vanish for every exponent $h$, since otherwise, contrary to the assumption, $C_1\,\chi'(x_1)\ldots C_m\,\chi'(x_m)$ would all have to be equal to zero.

Now that the existence of a function $\varphi(x)$ of the kind required has been established in this way, we apply the formula \S.\,205, (12):

% [sic] -- the source prints $P(x)$, although $P$ in \S.\,205, (12) is constant and the deduction below requires it to be independent of $x$.
\[
e^x\,P(x) = F(x) + e^{\xi}\,Q(x).
\]

In it we set $x=z_1, z_2, \ldots, z_m$, and denote the absolute values of $z_1, z_2, \ldots, z_m$ by $\xi_1, \xi_2, \ldots, \xi_m$.

Then it follows from (11):

\weq{15}{0 = C_1\,F(z_1)+C_2\,F(z_2)+\cdots+C_m\,F(z_m)+C_1\,e^{\xi_1}\,Q(z_1)+C_2\,e^{\xi_2}\,Q(z_2)+\cdots+C_m\,e^{\xi_m}\,Q(z_m),}

and, if $f(x)$ is any integral function of degree $n$, we have

\weq{16}{F(z) = f(z)+f'(z)+\cdots+f^{(n)}(z).}

Denoting by $p$ a rational prime that may be taken arbitrarily large, we set $x=cz$ [by (13)] and

\weq{17}{f(z) = \dfrac{\varphi(x)^{p-1}\,\varphi'(x)}{\Pi(p-1)},}

where $\varphi(x)$ is the function satisfying the conditions 1), 2).

Arranging by powers of $(x-x_1)$, we may obtain:

\[
\varphi(x)^{p-1}\,\varphi'(x) = \varphi'(x_1)^p\,(x-x_1)^{p-1}+A_p(x_1)\,(x-x_1)^p+A_{p+1}(x_1)\,(x-x_1)^{p+1}+\cdots,
\]

and in this the $A_p(x_1), A_{p+1}(x_1), \ldots$ are integral algebraic numbers, and indeed rational functions of $x_1$. On the other hand, by \gesperrt{Taylor}'s theorem:

\[
f(z) = f(z_1)+(z-z_1)\,f'(z_1)+\dfrac{(z-z_1)^2}{1\,.\,2}\,f''(z_1)+\cdots
\]

and $c(z-z_1)=x-x_1$. The comparison then gives

\[
f(z_1)=0,\ f'(z_1)=0,\ \ldots,\ f^{(p-2)}(z_1)=0,\ f^{(p-1)}(z_1)=c^{p-1}\,\varphi'(x_1)^p,
\]
\[
f^{(p)}(z_1)=p\,c^p\,A_p(x_1),\ f^{(p+1)}(z_1)=p(p+1)\,c^{p+1}\,A_{p+1}(x_1),\ \ldots,
\]

and consequently

\[
F(z_1) = c^{p-1}\,\varphi'(x_1)^p+p\,c^p\,A_p(x_1)+p(p+1)\,c^{p+1}\,A_{p+1}(x_1)+\cdots;
\]

in this, $x_1, z_1$ may be replaced by $x_2, z_2$ or by $x_3, z_3$, and so on.


% ===== scan leaf 787, printed 766 =====
% !PRINTED 766
\weberpage{766}
Accordingly the integral rational number

\[
C_1\,F(z_1)+C_2\,F(z_2)+\cdots+C_m\,F(z_m)
\]

is congruent modulo $p$ to

\[
c^{p-1}\left[C_1\,\varphi'(x_1)^p+C_2\,\varphi'(x_2)^p+\cdots+C_m\,\varphi'(x_m)^p\right].
\]

Now since $C_1, C_2, \ldots, C_p$ are integral rational numbers, % [sic] -- printed $C_p$; the surrounding system has $m$ constants
we have $C_1^p\equiv C_1, C_2^p\equiv C_2, \ldots \pmod{p}$, and by applying the polynomial theorem it follows that:

\[
\begin{aligned}
C_1\,\varphi'(x_1)^p+C_2\,\varphi'(x_2)^p+\cdots+C_m\,\varphi'(x_m)^p &\equiv \\
\left[C_1\,\varphi'(x_1)+C_2\,\varphi'(x_2)+\cdots+C_m\,\varphi'(x_m)\right]^p &\equiv k^p \pmod{p}.
\end{aligned}
\]

From this we therefore obtain the congruence

\weq{18}{C_1\,F(z_1)+C_2\,F(z_2)+\cdots+C_m\,F(z_m) \equiv c^{p-1}\,k^p \pmod{p}.}

Now the numbers $c, k$ are independent of $p$, and we can therefore assume $p$ so large that it divides neither $c$ nor $k$.

5. \gesperrt{Then the sum $C_1\,F(z_1)+C_2\,F(z_2)+\cdots+C_m\,F(z_m)$ is an integral rational number different from zero, and therefore in absolute value at least equal to 1.}

We now denote by $\varphi_1(x)$ the integral function which arises from $\varphi(x)$ when the negative ones among the coefficients $a, a_1, a_2, \ldots$ are replaced by their positive values. The function

\[
f_1(z) = \dfrac{\varphi_1(x)^{p-1}\,\varphi_1'(x)}{\Pi(p-1)},
\]

which arises according to (17), then has only positive coefficients, and these coefficients are in absolute value certainly not smaller than the corresponding coefficients of $f(z)$; for the coefficients of $f_1(z)$ arise by addition from the same numbers which, in the formation of the coefficients of $f(z)$, are partly added and partly subtracted.

From this, however, there follows by the formula \S.\,205, (10), if we denote by $\xi$ the absolute value of $x$, for the absolute value of $Q(z)$:

\[
|Q(z)| \leqq \dfrac{\varphi_1(\xi)^{p-1}\,\varphi_1'(\xi)}{\Pi(p-1)},
\]

6. \gesperrt{and therefore $Q(z)$ can, for every finite $z$, be made arbitrarily small by a sufficient increase of $p$.}


% ===== scan leaf 788, printed 767 =====
% !PRINTED 767
\weberpage{767}
But by 5. and 6. it has been shown that, if $p$ is large enough, equation (15) cannot hold. Thereby 4., and with it the whole of Theorem I., is proved.

This theorem now admits manifold applications. It gives us first of all the transcendence of $e$, if we take $C_1, C_2, \ldots, z_1, z_2, \ldots$ to be integral rational numbers. It further gives the transcendence of $\pi$. For from the equation $1+e^{i\pi}=0$ it follows by I. that $i\pi$, and consequently $\pi$ also, cannot be algebraic.

From this there follows further:

\gesperrt{For every algebraic number $x$, with the exception of $x=0$, $X=e^x$ is a transcendental number.}

\gesperrt{For every algebraic $X$, with the exception of $X=1$, every natural logarithm $x=\log X$ is a transcendental number.}

\gesperrt{For every arc that stands to the radius in a ratio $x$ expressible algebraically, with the exception of $x=0$, $X=\sin x$ is a transcendental number.}

This follows by I. from $2i\,X=e^{ix}-e^{-ix}$.

% [sic] -- the source prints the chord as $\frac12\sin(x/2)$; for unit radius the chord is $2\sin(x/2)$.
% [sic] -- the source's claim for $\mathrm{tg}\,x=\alpha x$ omits the algebraic root $x=0$; the intended claim concerns nonzero roots.
The same holds for the other trigonometric lines, $\cos x$, $\mathrm{tg}\,x$, and for the chord $\dfrac{1}{2}\sin\dfrac{x}{2}$. To adduce one thing more: the transcendental equation $\mathrm{tg}\,x=\alpha x$ has, for an algebraic $\alpha$ other than 0, only transcendental numbers for roots.

\begin{center}
\rule{4cm}{0.4pt}
\end{center}


% ===== scan leaf 789, printed 768 =====
\weberpage{768}% page number not printed on this page
% !PRINTED none
% part-title page (no running head, no printed page number): centred heading
% "Nachträge." with a horizontal rule beneath it, then subsection "I." and
% "Irreducibilität der Kreistheilungsgleichung.", followed by body text.
\begin{center}
{\Large \gesperrt{Supplements.}}

\rule{4em}{0.4pt}

I.

\textbf{\gesperrt{Irreducibility of the Cyclotomic Equation.}}
\end{center}

\S.\,134 of the first volume stands in need, in its last part, which treats of the irreducibility of the cyclotomic equation in the general case in which the order of the root of unity is a number composed of several distinct primes, of a correction, which shall be supplied here; and we ask the reader to complete the passage named in accordance with what is set out here.

If $n=ab$ is put, $a$ and $b$ being relatively prime, and if $\varrho, \alpha, \beta$ are primitive $n$th, $a$th, $b$th roots of unity, then $\varrho=\alpha\beta$, and it is proved in the passage cited that, if $\varrho$ is a root of a rational equation $\Phi(x)=0$, then among the roots of this equation $\varrho=\alpha\beta$ all primitive $a$th roots of unity $\alpha$ and all primitive $b$th roots of unity $\beta$ must occur. But it does not follow from this, as was there erroneously assumed, that among them \gesperrt{all} primitive $n$th roots of unity occur, because it has not been shown that every $\alpha$ occurs combined with every $\beta$.

To fill this gap, let an entirely elementary proof first find a place here.\footnote{\gesperrt{Arndt}, \gesperrt{Crelle}'s Journ., Bd. 56, S. 178 (1859). \gesperrt{Lebesgue}, \gesperrt{Liouville}'s Journ., 2. Ser., Bd. 4, S. 105 (1859).}

If $p$ is a prime, then all the binomial coefficients for the $p$th power

\[
B_h^{(p)} = \dfrac{p(p-1)\ldots(p-h+1)}{1\,.\,2\ldots h},
\]

are, with the exception of $B_0^{(p)}$ and $B_p^{(p)}$, integers divisible by $p$,
%CONT>>


% ===== scan leaf 790, printed 769 =====
% !PRINTED 769
\weberpage{769}
\pdfbookmark[0]{ Addenda.}{bm237}
\pdfbookmark[1]{I. Irreducibility of the cyclotomic equation}{bm238}
%<<CONT

and consequently, if $u$, $v$ are indeterminate quantities and the
\gesperrt{congruence is referred only to the coefficients},
\[
(u+v)^p \equiv u^p+v^p \ (\mathrm{mod}\ p).
\]

If we again replace $v$ in it by a sum $v+w+\cdots$,
we obtain by complete induction (or else from the
polynomial theorem):
\[
(u+v+w+\cdots)^p \equiv u^p+v^p+w^p+\cdots \ (\mathrm{mod}\ p),
\]
and if we apply this formula repeatedly, and understand by $k$
an arbitrary positive exponent,
\weq{1}{(u+v+w+\cdots)^{p^k} \equiv u^{p^k}+v^{p^k}+w^{p^k}+\cdots \ (\mathrm{mod}\ p).}

Let $x$ now be a variable, and let $a_1, a_2, \dots, a_\nu$ be integral rational
numbers,
\[
f(x) = x^\nu + a_1 x^{\nu-1} + a_2 x^{\nu-2} + \cdots + a_\nu
\]
an integral function of $x$, whose roots we may take to be $\alpha, \beta, \gamma, \dots$,
so that
\[
-a_1 = \Sigma\alpha, \quad a_2 = \Sigma\alpha\beta, \quad -a_3 = \Sigma\alpha\beta\gamma \dots
\]
are the symmetric fundamental functions of the $\alpha, \beta, \gamma, \dots$.

We form the function
\[
F(x) = x^\nu + A_1 x^{\nu-1} + A_2 x^{\nu-2} + \cdots + A_\nu,
\]
whose roots are the $p^{k}$th powers of the roots of $f$, that is,
\[
\alpha^{p^k}, \ \beta^{p^k}, \ \gamma^{p^k}, \dots
\]
The coefficients of $F$:
\[
A_1 = -\Sigma\alpha^{p^k}, \quad A_2 = \Sigma\alpha^{p^k}\beta^{p^k}, \quad A_3 = -\Sigma\alpha^{p^k}\beta^{p^k}\gamma^{p^k} \dots
\]
are, as symmetric functions of the $\alpha, \beta, \gamma, \dots$ with integral
coefficients, representable rationally and with integral coefficients by means of the $a_1, a_2, \dots$, and are
therefore integers. The quotients
\[
\frac{A_1 - a_1^{p^k}}{p}, \quad \frac{A_2 - a_2^{p^k}}{p}, \quad \frac{A_3 - a_3^{p^k}}{p}, \dots
\]
are, however, by formula (1) likewise integers, and thus by
\gesperrt{Fermat's theorem}
\weq{2}{F(x) \equiv f(x) \ (\mathrm{mod}\ p).}

These theorems are now to be applied to the function $f_n(x)$,
whose roots are all the \gesperrt{primitive} $n$th roots of unity $\alpha, \beta, \gamma, \dots$, and which, as was proved in \S.\,133 of the first
volume, has integral coefficients $a_1, a_2, \dots$.


% ===== scan leaf 791, printed 770 =====
% !PRINTED 770
\weberpage{770}

If $p$ is a prime that divides $n$, then
\weq{3}{\frac{x^n-1}{x^{\frac{n}{p}}-1} = x^{\frac{n}{p}(p-1)} + x^{\frac{n}{p}(p-2)} + \cdots + x^{\frac{n}{p}} + 1,}
and this function vanishes for $x=\alpha, \beta, \gamma, \dots$ and is
therefore divisible by $f_n(x)$. We will set it $= f_n(x)\, g(x)$,
where $g(x)$ (by Bd. I, \S.\,2) is then an integral function of $x$
with integral coefficients.

Now if
\[
n = p^k n',
\]
where $n'$ is no longer divisible by $p$, and $\alpha'$ is a primitive $n'$th root of unity, then
\[
\alpha'^{\frac{n}{p}} = 1,
\]
and consequently, if one puts $x = \alpha'$ [by (3)]:
\weq{4}{f_n(\alpha')\, g(\alpha') = p.}

From this we first derive a proof of the irreducibility of $f_n(x)$ under the hypothesis that $n$ is a
\gesperrt{prime power}, even though for this case the proof given in Bd. I,
\S.\,134 is unobjectionable.

Suppose that $f_n(x)$ splits into two rational factors:
\weq{5}{f_n(x) = \varphi(x)\, \psi(x).}

Then all the roots, both of $\varphi(x)$ and of $\psi(x)$, are
$n$th roots of unity, and their $n$th powers are therefore equal to 1. If
now $n = p^k$, then from $\varphi(x)$ and $\psi(x)$ we can derive two integral
functions $\Phi(x)$, $\Psi(x)$ whose roots are the $n$th powers of the roots of $\varphi(x)$ and $\psi(x)$ and which are therefore all
equal to 1, and from (2) there follows:
\[
\varphi(x) \equiv \Phi(x), \quad \psi(x) \equiv \Psi(x) \ (\mathrm{mod}\ p),
\]
from which, for $x = 1$, it follows that:
\[
\varphi(1) \equiv 0, \quad \psi(1) \equiv 0 \ (\mathrm{mod}\ p).
\]
Accordingly there follows from (5):
\weq{6}{f_n(1) \equiv 0 \ (\mathrm{mod}\ p^2).}

But in this case $f_n(1) = p$ (Bd. I, \S.\,133, V.),
which contradicts formula (6). Our assumption was therefore
inadmissible, and $f_n(x)$ is irreducible.


% ===== scan leaf 792, printed 771 =====
% !PRINTED 771
\weberpage{771}

In order to prove the irreducibility for an arbitrarily composite $n$,
we apply complete induction. We set
\[
n = p^k n',
\]
and understand by $p$ a prime which divides $n$, but not $n'$, and
assume the \gesperrt{irreducibility of} $f_{n'}(x)$ \gesperrt{to be already proved}. If
then $f_n(x)$ is decomposable into two rational factors
$\varphi(x)$, $\psi(x)$, neither of which is constant:
\weq{7}{f_n(x) = \varphi(x)\, \psi(x),}
then from $\varphi(x)$ we derive the function $\Phi(x)$, whose roots
are the $p^{k}$th powers of the roots of $\varphi(x)$, and which satisfies the condition
\[
\varphi(x) \equiv \Phi(x) \ (\mathrm{mod}\ p)
\]
or also, if we denote by $\chi(x)$ a second function of $x$ with
integral coefficients:
\[
\varphi(x) = \Phi(x) + p\, \chi(x).
\]

But the roots of $\Phi(x)$ are primitive roots of unity of degree
$n'$, and therefore there must be among them at least one,
$\alpha'$, for which $\Phi(\alpha')$ vanishes, so that
\[
\varphi(\alpha') = p\, \chi(\alpha')
\]
holds. On account of the presupposed irreducibility of $f_{n'}(x)$
this equation must, however, hold for \gesperrt{all} roots of unity $\alpha'$, and
for the same reason we have, if $\vartheta(x)$ is a function with
integral coefficients, for all $\alpha'$:
\[
\psi(\alpha') = p\, \vartheta(\alpha').
\]

Accordingly there follows from (7):
\[
f_n(\alpha') = p^2\, \chi(\alpha')\, \vartheta(\alpha'),
\]
and by (4):
\[
1 = p\, \chi(\alpha')\, \vartheta(\alpha')\, g(\alpha').
\]

But the product $\chi(x)\, \vartheta(x)\, g(x)$ is an integral
function of $x$ with integral coefficients, whose remainder on division by $f_{n'}(x)$ likewise has integral coefficients and may be
denoted by $\Theta(x)$. This remainder is of lower degree than $f_{n'}(x)$, and the
function
\[
p\, \Theta(x) - 1
\]
vanishes for \gesperrt{all roots} of $f_{n'}(x)$. Consequently the
equation
\[
p\, \Theta(x) = 1
\]
must be identical, which however is obviously impossible, since $\Theta(x)$ for an
%CONT>>


% ===== scan leaf 793, printed 772 =====
% !PRINTED 772
\weberpage{772}

%<<CONT
integer $x$ passes into an integer which, multiplied by $p$, cannot be $= 1$.

With this the irreducibility of $f_n(x)$ is proved in general.

\begin{center}
\rule{4cm}{0.4pt}
\end{center}

Now that the irreducibility of the cyclotomic equation
$f_n(x) = 0$ has been established in general, it follows, as in \S.\,134 of the
first volume, that it also remains irreducible when an arbitrary root of unity is adjoined whose degree is relatively
prime to $n$.

\gesperrt{Kronecker} takes, in a certain sense, the opposite path\footnote{Mém. sur les facteurs irréd. de l'expression $(x^n-1)$. \gesperrt{Liouville}'s
Journal, Bd. 19 (1854). A simple proof of the irreducibility of the general cyclotomic equation, resting on the theory of the higher congruences,
has been given by \gesperrt{Dedekind} [\gesperrt{Crelle}'s Journal für Mathematik, Bd. 58 (1859)].}.

He first proves, under the assumption that $n$ is a
power of a prime $p$, the irreducibility of $f_n(x)$ after
the adjunction of a root of unity whose degree is not divisible by $p$,
and from this he derives the general proof of irreducibility, as we shall now show briefly.

\begin{center}
\rule{4cm}{0.4pt}
\end{center}

Let now
\[
n = p^k
\]
be a power of a prime $p$, and let
\weq{8}{f_n(x) = x^{p^{k-1}(p-1)} + x^{p^{k-1}(p-2)} + \cdots + x^{p^{k-1}} + 1}
be the function whose roots are the primitive $n$th roots of unity.
Let further $\alpha$ be any root of unity whose degree $m$
is not divisible by $p$, and let $\nu$ be the degree of the field $R(\alpha)$.
\gesperrt{It is to be proved that $f_n(x)$ is irreducible in the field $R(\alpha)$.}

Every number of the field $R(\alpha)$ can be represented in one and only
one way in the form:
\weq{9}{\varphi(\alpha) = m_0 + m_1 \alpha + \cdots + m_{\nu-1} \alpha^{\nu-1},}
where $m_0, m_1, \dots, m_{\nu-1}$ are rational numerical coefficients.

If $\chi(x) = 0$ is the rational equation of lowest degree
having $\alpha$ as its root, and
\weq{10}{\chi(x) = x^\nu + c_1 x^{\nu-1} + \cdots + c_{\nu-1} x + c_\nu,}
%CONT>>


% ===== scan leaf 794, printed 773 =====
% !PRINTED 773
\weberpage{773}
%<<CONT
then the $c_1, c_2, \ldots$ are \gesperrt{integers}, because $\chi(x)$ must be a divisor of $x^m - 1$ (by \gesperrt{Gauss}'s theorem, Bd.~I, \S.\,2).

From this it follows that, if $\varphi(x)$ denotes an integral function of $x$ with integral rational coefficients, the coefficients $m_0, m_1, \ldots, m_{\nu-1}$ in (9) become integers, even when the degree of $\varphi(x)$ is originally higher than $\nu$; for one obtains these coefficients by seeking out the residue of the division of $\varphi(x)$ by $\chi(x)$.

We now assume that $f_n(x)$ is decomposable in the field $R(\alpha)$, and that accordingly:

\weq{11}{f_n(x) = \varphi(x,\alpha)\ \psi(x,\alpha),}

where $\varphi(x,\alpha)$ and $\psi(x,\alpha)$ are two integral functions of $x$, of which neither is constant, and whose coefficients are contained in $R(\alpha)$. From (8) and (11) it follows, if we set $x=1$,

\weq{12}{p = \varphi(1,\alpha)\ \psi(1,\alpha),}

where $\varphi(1,\alpha)$, $\psi(1,\alpha)$ can be represented, as numbers of the field $R(\alpha)$, in the form (9). If we denote the least common denominators of these two representations by $a$ and $b$, we obtain:

\weq{13}{\begin{aligned}
a\,\varphi(1,\alpha) &= a_0 + a_1\alpha + a_2\alpha^2 + \cdots + a_{\nu-1}\alpha^{\nu-1} = A(\alpha)\\
b\,\psi(1,\alpha) &= b_0 + b_1\alpha + b_2\alpha^2 + \cdots + b_{\nu-1}\alpha^{\nu-1} = B(\alpha),
\end{aligned}}

where $a, a_0, a_1, \ldots, a_{\nu-1}$ are integers without a common divisor and $A$ is a function symbol. The same holds for $b, b_0, b_1, \ldots, b_{\nu-1}$ and $B$.

The roots of the equation $\varphi(x,\alpha) = 0$ are to be found among the primitive $n$th roots of unity, and these are all powers of one among them, $r$. Accordingly there are certain powers $r^h, r^{h'}, r^{h''}, \ldots$ such that

\[
\varphi(x,\alpha) = (x - r^h)(x - r^{h'})(x - r^{h''}) \ldots
\]

holds, and consequently, by (13):

\[
A(\alpha) = a(1 - r^h)(1 - r^{h'})(1 - r^{h''}) \ldots
\]

If we take the $n$th power of this, it follows by (1), having regard to the equation $r^n = 1$:

\[
[A(\alpha)]^n = p\,\chi(r),
\]

where $\chi(x)$ denotes an integral function of the variable $x$ with integral coefficients.


% ===== scan leaf 795, printed 774 =====
% !PRINTED 774
\weberpage{774}
Now, understanding by $t$ a variable, we set:

\[
\begin{aligned}
{}[t - \chi(1)][t - \chi(r)] \ldots [t - \chi(r^{n-1})] &= \\
t^n + A_1t^{n-1} + A_2t^{n-2} + \cdots + A_n, &
\end{aligned}
\]

in which the coefficients $A_1, A_2, \ldots, A_n$, being integral symmetric functions of the $n$ quantities $\chi(1), \chi(r), \ldots, \chi(r^{n-1})$, are integral rational numbers. If, however, we set in it

\[
t = \chi(r) = \dfrac{[A(\alpha)]^n}{p}
\]

then the left side vanishes, and multiplication by $p^n$ gives:

\weq{14}{[A(\alpha)]^{n^2} = p\,C(\alpha),}

where $C(\alpha)$ denotes a number of the form (9) with \gesperrt{integral coefficients}.

If formula (14) is repeatedly raised to the $p$th power and theorem (1) is applied, there follows from it, for every power that is not less than $n^2$:

\weq{15}{A(\alpha^{p^\lambda}) = p\,\xi(\alpha),}

in which the integral function $\xi(\alpha)$ likewise has integral coefficients.

Here we may assume $\lambda$ to be divisible by $\varphi(m)$, and then, since $m$ is not divisible by $p$, we have by the generalised theorem of \gesperrt{Fermat} (Bd.~II, \S.\,14):

\[
p^\lambda \equiv 1 \ (\mathrm{mod}\ m), \quad \alpha^{p^\lambda} = \alpha,
\]

so that from (14) there follows: % [sic] in the source; equation (16) follows mathematically from (15), not (14)

\weq{16}{A(\alpha) = p\,\xi(\alpha).}

Hence in (13) the coefficients $a_0, a_1, \ldots, a_{\nu-1}$ are all divisible by $p$, and $a$ is consequently not divisible by $p$. In the same way it is shown that $b_0, b_1, \ldots, b_{\nu-1}$ are divisible by $p$, while $b$ is indivisible by $p$.

Accordingly there follows from (12), if we set $A(\alpha)B(\alpha) = p^2\Theta(\alpha)$:

\weq{17}{ab = p\,\Theta(\alpha),}

where $\Theta(\alpha)$ denotes a function of $\alpha$ of the form (9) with integral coefficients. But equation (17), since it is at most of degree $(\nu-1)$, would have to be identical with respect to $\alpha$, and $ab$ would have to be divisible by $p$, which is not possible.


% ===== scan leaf 796, printed 775 =====
% !PRINTED 775
\weberpage{775}
From this follows the inadmissibility of our assumption that $f_n(x)$, in accordance with formula (11), is decomposable into two factors in the field $R(\alpha)$.

\begin{center}
\rule{4cm}{0.4pt}
\end{center}

With this it can now be proved anew and in general that the equation whose roots are all the primitive $m$th roots of unity, and whose degree is $\varphi(m)$, is irreducible even when $m$ is an arbitrary composite number. This is equivalent to the following theorem:

\gesperrt{If an integral function $\varphi(x)$ with rational coefficients vanishes when a primitive $m$th root of unity is put for $x$, then $\varphi(x)$ vanishes also when any other primitive $m$th root of unity is put for $x$.}

When $m$ is a power of a prime, the theorem has already been proved in what precedes. We therefore apply complete induction: we assume it established for $m$th roots of unity and derive it for $mn$th roots of unity, where $n = p^k$ and $p$ is a prime that does not divide $m$.

Every $mn$th root of unity $\beta$ can be represented in the form

\[
\beta = \alpha r
\]

where $\alpha$ is an $m$th and $r$ an $n$th root of unity. What is to be proved can be stated thus: if $\beta$ is a root of a rational equation $\Phi(x) = 1$, % [sic]
then all the other primitive roots of unity of degree $mn$ must be roots of this equation as well.

If, as above, $f_n(x) = 0$ is the equation whose roots are the quantities $r$, then, as we have proved, $f_n(x)$ is not decomposable into rational factors, nor into such as depend rationally on $\alpha$. The function $f_n(\alpha^{-1}x)$, which vanishes for $x = \beta$, can therefore likewise not be decomposed into rational factors, because from a decomposition of the form

\[
f_n(\alpha^{-1}x) = \varphi(x,\alpha)\ \psi(x,\alpha)
\]

there would follow:

\[
f_n(x) = \varphi(\alpha x,\alpha)\ \psi(\alpha x,\alpha),
\]

which cannot hold. Since, accordingly, $f_n(\alpha^{-1}x)$ has a root
%CONT>>


% ===== scan leaf 797, printed 776 =====
% !PRINTED 776
\weberpage{776}
%<<CONT
in common with $\Phi(x)$, $\Phi(x)$ must be divisible by $f_n(\alpha^{-1}x)$, and we can accordingly set, for an indeterminate $x$,

\weq{18}{\Phi(x) = f_n(\alpha^{-1}x)\ \Psi(x,\alpha)}

where $\Psi(x,\alpha)$ is an integral rational function of $x$.

If $\alpha, \alpha_1, \alpha_2, \ldots$ are all the primitive $m$th roots of unity, which, as we have assumed, are roots of an \gesperrt{irreducible} rational equation

\[
F_m(x) = 0
\]

then, by this assumption, $\alpha$ can be replaced in the identity (18) by any other one of $\alpha_1, \alpha_2, \ldots$, and it follows that $\Phi(x)$ is divisible by each of the factors

\[
f_n(\alpha^{-1}x),\ f_n(\alpha_1^{-1}x),\ f_n(\alpha_2^{-1}x) \ldots
\]

But two of these factors can have no common divisor; for if, say,

\[
f_n(\alpha^{-1}x) = 0,\ f_n(\alpha_1^{-1}x) = 0
\]

had a common root $x$, then two primitive $n$th roots of unity $r, r_1$ would have to exist such that

\[
x = \alpha r = \alpha_1 r_1
\]

held; so we should have $\alpha^n = \alpha_1^n$, and consequently, since $n$ is relatively prime to $m$, $\alpha = \alpha_1$.

Accordingly $\Phi(x)$ is divisible by the product

\[
F_{mn}(x) = f_n(\alpha^{-1}x)\ f_n(\alpha_1^{-1}x)\ f_n(\alpha_2^{-1}x) \ldots,
\]

whose roots are all the $\beta$, and this function is therefore also indecomposable, as was to be proved.

\begin{center}
\rule{4cm}{0.4pt}
\end{center}


% ===== scan leaf 798, printed 777 =====
% !PRINTED 777
\weberpage{777}
\pdfbookmark[1]{II. The irreducibility of the cyclotomic equation and the theorem on the primes contained in a linear form}{bm239}

\begin{center}
II.

\gesperrt{The Irreducibility of the Cyclotomic Equation
and the Theorem on the Primes contained in a
Linear Form.}
\end{center}

In the first supplement, which is to be regarded as a completion of the first Band,
no use was made, of set purpose, of the theory of integral algebraic numbers. In
conclusion we therefore give one more proof of the irreducibility of the cyclotomic
equation, resting on other foundations, namely on the theorems concerning the class
number (\S.\,192); it is remarkable on account of the generalisations which it admits.

Let $m$ be an arbitrary natural number, and let $n$ be the system of the positive
numbers coprime to $m$; among these there are
\weq{1}{\varphi(m) = \mu,}
which are less than $m$, and these we denote by $a$.

If we unite into a single class $A$ all the numbers $n$ that are congruent to one
and the same $a$ with respect to the modulus $m$, we obtain $\mu$ classes of numbers:
\weq{2}{A_1, A_2, \ldots, A_\mu,}
which, under composition by multiplication, form an \gesperrt{Abelian} group of
degree $\mu$, $\mathfrak{N}$. This group has already been considered in \S.\,16 of
this Band.

The $\mu$ characters of this group we denote by
\weq{3}{\chi_1, \chi_2, \ldots, \chi_\mu,}
and we put, when $\chi$ is one of these characters and $n$ denotes a number from
the class $A$,
\weq{4}{\chi(A) = \chi(n).}
These characters, which are more closely determined in \S.\,16, are all
$\mu$th roots of unity.


% ===== scan leaf 799, printed 778 =====
% !PRINTED 778
\weberpage{778}
% running head "II. Nachtrag." at the top of this page is not transcribed as body text.

Here, however, we have no need to make use of the special expressions for these
characters. We rely instead on the following theorem, which holds for all
\gesperrt{Abel}ian groups:

1. \gesperrt{If} $A$ \gesperrt{is an element of degree} $f$ \gesperrt{of an}
\gesperrt{Abel}\gesperrt{ian group of degree} $\mu$, \gesperrt{and if} $\mu = ef$,
\gesperrt{then all the} $\chi_i(A)$ \gesperrt{are} $f$\gesperrt{th roots of unity,
and among them every} $f$\gesperrt{th root of unity occurs exactly} $e$
\gesperrt{times.}

The theorem is contained implicitly in theorem 7., \S.\,12 of this volume. For if we
apply that theorem to the group
\[
T = 1, A, A^2, \ldots, A^{f-1}
\]
whose index with respect to $\mathfrak{N}$ is equal to $e$, it follows that there is
a group $\Xi$ of $e$ characters $\xi$ for which the
conditions
\[
\xi(A) = 1
\]
are satisfied.

If, then, one decomposes the group $X$ of the characters $\chi$ into the subsidiary
groups
\[
\Xi_1, \Xi_2, \ldots, \Xi_f,
\]
then $\chi_1(A)$ and $\chi_2(A)$ are equal to one another or different from one
another according as $\chi_1$ and $\chi_2$ occur in the same or in different ones of
these subsidiary groups. Since, moreover, because
$\chi(A)^f = \chi(A^f) = 1$ all the $\chi(A)$ are roots of unity of degree $f$, and
there are only $f$ different such roots, our theorem is thereby proved.

If $n$ is contained in the class $A$, then the degree $f$ of $A$ is the exponent to
which $n$ \gesperrt{belongs} with respect to the modulus $m$, that is, the least positive
exponent for which
\weq{5}{n^f \equiv 1 \pmod{m}}
holds. If then $n$ belongs to the exponent $f$, all the $\chi_i(n)$ are $f$th roots
of unity, and every $f$th root of unity occurs $e$ times among them.

If $a$ denotes any one of the $\varphi(m)$ residues of the numbers $n$, then in the
form
\weq{6}{\mu_x = mx + a - m}
every number of the number class (2) represented by $a$ is contained; if by $a$ we
understand the least positive residue, then we obtain the totality of the numbers of
this class when we let $x$ run through all positive integral values.


% ===== scan leaf 800, printed 779 =====
% !PRINTED 779
\weberpage{779}
% running head "Primzahlen in Linearformen." / "779" at the top of this page is not
% transcribed as body text.

Now
\[
\frac{\mu_x}{m} - x = \frac{a-m}{m},
\]
and consequently we can apply the theorem \S.\,191, 6. to the quantities
$\mu_x$, in which we have to set $n = x$, $\gamma = 1:m$, $c_n = (a-m):m$.
Accordingly we have
\weq{7}{A(s) = \overset{x}{\underset{0,\,\infty}{\Sigma}} \frac{1}{(mx+a)^s}
= \overset{x}{\underset{1,\,\infty}{\Sigma}} \frac{1}{(mx+a-m)^s},}
which, so long as $s > 1$, is a continuous function of $s$, and we have
\weq{8}{A(s) = \frac{1}{m(s-1)} + C(s),}
where $C(s)$ remains finite for $s=1$\footnote{By the use of known
theorems on $\Gamma$-functions one finds, from \S.\,191, (20):
\[
C(1) = \frac{1}{m} \frac{\Gamma'\left(\frac{a}{m}\right)}{\Gamma\left(\frac{a}{m}\right)}.
\]}.

We now let $A$ run through all the classes (2), denote by
$\chi$ any one of the characters of the group $\mathfrak{N}$, and set
\weq{9}{Q(s) = \overset{A}{\Sigma} \chi(A)\, A(s),}
or, what is the same thing,
\weq{10}{Q(s) = \overset{n}{\Sigma} \frac{\chi(n)}{n^s},}
in which the sum extends over all the numbers $n$. The number of these sums
$Q$ is as great as the number of the characters, that is, $= \varphi(m)$. We
denote them, when a distinction is necessary, by
\[
Q_1, Q_2, \ldots, Q_\mu.
\]

One of them, $Q_1$, corresponds to the principal character and has the expression
\[
Q_1(s) = \overset{n}{\Sigma} \frac{1}{n^s}.
\]

Now [\S.\,11, (21)]:
\[
\overset{A}{\Sigma} \chi(A) = \mu \ \text{or} = 0,
\]
according as $\chi$ is the principal character or not. Furthermore, there
results from (8)
%CONT>>


% ===== scan leaf 801, printed 780 =====
% !PRINTED 780
\weberpage{780}
% running head "II. Nachtrag." at the top of this page is not transcribed as body text.
%<<CONT
\[
Q(s) = \frac{1}{m(s-1)} \overset{A}{\Sigma} \chi(A) + \overset{A}{\Sigma} \chi(A)\, C(s),
\]
and from this the theorem:

2. \gesperrt{The sums $Q_2, \ldots, Q_\mu$ take finite values for $s=1$,}
$Q_1$ \gesperrt{becomes infinite, and indeed we have}
\[
\mathrm{Lim}_{s=1} (s-1)\, Q_1(s) = \frac{\mu}{m}.
\]

In order to transform the sums $Q(s)$, we denote by $p$ every prime that does
not divide $m$, and multiply all the infinite series of the form
\[
1 + \frac{\chi(p)}{p^s} + \frac{\chi(p^2)}{p^{2s}} + \cdots
= \frac{1}{1 - \chi(p)\, p^{-s}}
\]
with one another. In this way we obtain [cf. \S.\,192, (8)]:
\weq{11}{Q(s) = \overset{p}{\Pi} \frac{1}{1 - \chi(p)\, p^{-s}}.}
Now if $p$ belongs to the exponent $f$, that is, if $f$ is the least positive
number satisfying the congruence
\weq{12}{p^f \equiv 1 \pmod{m}}
and if $\mu = ef$, then among the $\chi(p)$ every $f$th root of unity occurs
exactly $e$ times, and we have therefore (for a variable $x$):
\[
\overset{\chi}{\Pi} [1 - \chi(p)\, x] = (1 - x^f)^e,
\]
and accordingly there follows from (11):
\weq{13}{Q_1 Q_2 \ldots Q_\mu = \overset{p}{\Pi} \frac{1}{(1-p^{-fs})^e}.}

We now introduce, without assuming the irreducibility of the cyclotomic
equation, the cyclotomic field $\Omega_m$, whose still unknown degree we shall
denote by $\nu$. But since the primitive $m$th root of unity $r$ satisfies a
rational equation of degree $\mu$, we have
\weq{14}{\nu \leqq \mu.}

The numbers $\omega$ of the field $\Omega_m$ can all be represented in the form
\weq{15}{a\omega = a_0 + a_1 r + \cdots + a_{\nu-1} r^{\nu-1}}
with integral rational coefficients $a, a_0, a_1, \ldots, u_{\nu-1}$; % [sic] -- printed $u_{\nu-1}$; display (15) has $a_{\nu-1}$
and if
$\omega$ is an \gesperrt{integer}, we can set for $a$ a fixed integer, a closer
knowledge of which is not neces-
%CONT>>


% ===== scan leaf 802, printed 781 =====
% !PRINTED 781
\weberpage{781}
%<<CONT
sary. Its value is equal to the square root of the quotient
of the discriminants
\[
\Delta (1, r, r^2, \ldots, r^{\nu-1}) : \Delta (\omega_1, \omega_2, \ldots, \omega_\nu),
\]
in which $\omega_1, \omega_2, \ldots, \omega_\nu$ is a minimal basis (\S.\,144, 145).

From the system of the $p$ we exclude all primes which divide the discriminant of the cyclotomic equation and all which divide $m$ and $a$, whose number is in any case finite, and we denote by $\mathfrak{p}$ a prime ideal dividing $p$. Then the congruence
\[
r^p \equiv r \ (\mathrm{mod}\ \mathfrak{p})
\]
is satisfied always, but also only, when

\weq{16}{p \equiv 1 \ (\mathrm{mod}\ m)}

holds, that is, when $p$ belongs to the exponent 1, and hence when $r^p = r$; and under the same condition we therefore have also, by (15), for every integer $\omega$ in $\Omega_m$:
\[
\omega^p \equiv \omega \ (\mathrm{mod}\ \mathfrak{p}).
\]
But by \S.\,162 this is the necessary and sufficient condition for the prime $p$, which does not divide the field discriminant of $\Omega_m$, to split into nothing but prime ideals of the first degree.

If, then, we denote by
\[
P_f(s) = \prod (1 - p^{-sf})
\]
the product extended over all primes $p$ which belong to the exponent $f$, it follows from \S.\,192 that, if $f>1$, $P_f(1)$ is finite and different from zero, and that

\weq{17}{\dfrac{s-1}{P_1(s)^{\nu}}}

\gesperrt{is, for} $s=1$, \gesperrt{finite and different from zero.}

Now by (13), if $S$ denotes the \gesperrt{finite} product $\Pi(1-p^{-sf})^{-e}$ extended over the \gesperrt{excluded} primes $p$,
\[
Q_1 Q_2 \ldots Q_\mu = S \overset{f}{\Pi} \frac{1}{P_f(s)^{e}},
\]
and since $e = \mu$ for $f = 1$, it follows from (17) that

\weq{18}{(s-1)^{\frac{\mu}{\nu}} Q_1 Q_2 \ldots Q_\mu}

\gesperrt{is, for} $s = 1$, \gesperrt{finite and different from zero.}


% ===== scan leaf 803, printed 782 =====
% !PRINTED 782
\weberpage{782}
On the other hand, however, by 2.:

\weq{19}{(s-1) Q_1 Q_2 \ldots Q_\mu}

is finite for $s = 1$, and so if we divide (19) by (18), it follows that:
\[
(s-1)^{1-\frac{\mu}{\nu}} \text{ finite for } s = 1 \text{,}
\]
i.e.

\weq{20}{\nu \geqq \mu.}

Taken together with (14), however, this yields $\nu = \mu$, whereby the following theorem is proved:

3. \gesperrt{The degree of the cyclotomic field $\Omega_m$ is equal to $\varphi(m)$, and hence the cyclotomic equation of degree $\varphi(m)$, whose roots are the primitive $m$th roots of unity, is irreducible.}

But this consideration renders us yet another important service, in that it furnishes us with the proof of the famous theorem \gesperrt{that in every arithmetic progression whose difference and initial term are natural numbers without a common divisor there are contained infinitely many primes}\textsuperscript{1)}.

For it now follows from (18) that the product
\[
(s-1) Q_1 Q_2 \ldots Q_\mu
\]
has a finite value different from zero for $s = 1$, and since by 2. none of the factors is infinite, it follows that:

4. \gesperrt{The sums}
\[
(s-1) Q_1(s), Q_2(s), \ldots, Q_\mu(s)
\]
\gesperrt{have finite values different from zero for $s = 1$.}

If we apply to all the factors of (11) the formula:
\[
\log \left[1 - \chi(p) p^{-s}\right] = \frac{\chi(p)}{p^s} + \frac{1}{2} \frac{\chi(p^2)}{p^{2s}} + \frac{1}{3} \frac{\chi(p^3)}{p^{3s}} + \cdots % [sic] the source omits the mathematically required minus sign before \log
\]

\noindent\rule{4cm}{0.4pt}

\textsuperscript{1)} The first complete and general proof of this theorem was given by \gesperrt{Dirichlet} (Abhandlungen der Berl. Akademie vom Jahre 1837. Gesammelte Werke Nr.\ XXI). \gesperrt{Dedekind} has drawn attention to the essential simplification of this proof by the use of the \gesperrt{Kummer} formulae for the class number in the cyclotomic fields (Vorlesungen über Zahlentheorie, 3. Auflage, S. 596).
%CONT>>


% ===== scan leaf 804, printed 783 =====
% !PRINTED 783
\weberpage{783}
%<<CONT
(in which the imaginary part of the logarithm is to be taken between $\pm \pi i$), then, if the imaginary part of $\log Q(s)$ is suitably determined, it follows that

\weq{21}{\log Q(s) = \Sigma \frac{\chi(p)}{p^s} + \frac{1}{2}\Sigma \frac{\chi(p^2)}{p^{2s}} + \frac{1}{3}\Sigma \frac{\chi(p^3)}{p^{3s}} + \cdots}

If now $A$ denotes any one of the classes (2) and $a$ a number from $A$, then by \S.\,11, (6) of this volume

\weq{22}{\overset{\chi}{\Sigma} \chi(A) \chi(n) = \Sigma \chi(an) = \mu \text{ or } = 0,}

according as $an \equiv 1$ or not $\equiv 1$ to the modulus $m$, that is, according as $n$ is contained in the class $A^{-1}$ or is not contained in it. If, then, we multiply formula (21) by $\chi(A)$ and sum with respect to $\chi$, it follows that

\weq{23}{\frac{1}{\mu} \overset{\chi}{\Sigma} \chi(A) \log Q(s) = \Sigma \frac{1}{p^s} + \frac{1}{2}\Sigma \frac{1}{p^{2s}} + \frac{1}{3}\Sigma \frac{1}{p^{3s}} + \cdots,}

where on the right the first, second, third, \ldots\ sum refers to all primes $p$ whose first, second, third, \ldots\ power is contained in $A^{-1}$.

The second part of this sum
\[
R(s) = \frac{1}{2}\Sigma \frac{1}{p^{2s}} + \frac{1}{3}\Sigma \frac{1}{p^{3s}} + \cdots
\]
is increased if each of its partial sums is extended to \gesperrt{all} primes $p$, and accordingly
\[
R(s) < \frac{1}{2}\Sigma \left(\frac{1}{p^{2s}} + \frac{1}{p^{3s}} + \cdots\right) = \frac{1}{2}\Sigma \frac{1}{p^s (p^s - 1)},
\]
and so all the more:
\[
R(s) < \frac{1}{2} \sum_{2,\infty}^{n} \frac{1}{n^{2s} (1 - n^{-s})},
\]
or, since $1 - n^{-s} > \frac{1}{2}$:
\[
R(s) < \sum_{2,\infty}^{n} \frac{1}{n^{2s}},
\]
and from this one concludes that $R$ remains finite for $s = 1$ (\S.\,191, 2.).

If we now, in formula (23),
\[
\frac{1}{\mu} \overset{\chi}{\Sigma} \chi(A) \log Q(s) = \Sigma \frac{1}{p^s} + R(s),
\]
%CONT>>


% ===== scan leaf 805, printed 784 =====
% !PRINTED 784
\weberpage{784}
%<<CONT
$s$ pass over into 1, then $\log Q_1(s)$ becomes infinite by 4., while the remaining terms of the left side, $\log Q_2(s), \ldots \log Q_\mu(s)$, remain finite. Since $R(s)$ likewise remains finite, the sum $\Sigma p^{-s}$ must become infinite, and this is certainly possible only if the sum consists of infinitely many terms. With this it is proved:

5. \gesperrt{In each of the number classes $A$ (with respect to the modulus $m$) infinitely many primes are contained.}

One can also express this by saying that the linear form $mx + a$, in which $m$ and $a$ are integers without a common divisor, represents a prime for infinitely many integral values of $x$.

\begin{center}
\rule{4cm}{0.4pt}
\end{center}

% [The remainder of this leaf and the following leaf are blank in the source.]


% ===== scan leaf 806, printed 785 =====
\weberpage{785}% page number not printed on this page
\pdfbookmark[1]{Index to Volumes I and II}{bm240}
% !PRINTED none
% Register (index) opening page: no folio number appears in the running head on
% this page image (title page of the alphabetical index). Footer at the foot of
% the page reads "Weber, Algebra. II." followed by the gathering signature "50".
\begin{center}
\gesperrt{INDEX.}
\end{center}

\begin{center}
(The Roman figures denote the volume, the Arabic figures the page.)
\end{center}

\begin{center}
\rule{3cm}{0.5pt}
\end{center}

\noindent
\gesperrt{\textbf{A}bel}ian equations I, 533.\\
--- --- cubic II, 107.\\
--- --- biquadratic II, 110.\\
--- groups I, 476, 536; II, 6, 32.\\
--- fields II, 588, 648.\\
Derived functions I, 52.\\
Absolute norm II, 502, 536.\\
Absolute value of an imaginary number I, 18.\\
--- --- of a functional II, 501.\\
Enumerable sets II, 745.\\
Addition I, 14.\\
Adjunction I, 451.\\
Similar substitutions II, 154.\\
Equivalent functionals II, 552.\\
--- numbers I, 371.\\
--- --- and functionals in the quadratic field II, 610.\\
Affect I, 480.\\
Algebraic functions II, 526.\\
--- fields I, 455; II, 527.\\
--- numbers II, 489.\\
Algorithm of the greatest common divisor for numbers I, 2.\\
--- for integral functions I, 34.\\
Alternating group I, 496, 600.\\
--- --- of five figures II, 138, 230.\\
--- function I, 496.\\
Arithmetic series I, 47.\\
\gesperrt{Aronhold} systems of double tangents II, 367.\\
Associated functionals II, 510.\\

Associated numbers I, 586.\\
Solution of equations by root signs I, 595.\\
Axes of a ternary substitution group II, 438.\\

\medskip
\noindent
\textbf{B}ases of the functionals II, 532, 581.\\
--- of the ideals II, 550.\\
Basis of an \gesperrt{Abel}ian group II, 33.\\
--- of an algebraic field II, 528.\\
Base of the natural system of logarithms II, 752.\\
Basis form of a functional II, 535.\\
--- --- --- of $\mathfrak{o}$ II, 535.\\
--- --- --- in the quadratic field II, 604, 608.\\
Removal of the second term from an equation I, 114.\\
\gesperrt{Bernoulli} method of approximation for the solution of equations I, 341.\\
\gesperrt{Bertrand}'s theorem on the limits of the index of permutation groups II, 143.\\
\gesperrt{Bezout}'s theorem I, 161.\\
Bezoutiant I, 225, 265.\\
--- and the reality of the roots I, 252.\\
Binomial theorem I, 42.\\
Binomial coefficients I, 42.\\
Binary forms I, 189.\\
--- linear substitution II, 191.\\
Biquadratic \gesperrt{Abel}ian equations II, 110.\\

\bigskip
\noindent
%% [page furniture, not body text -- commented out 2026-08-10] Page furniture (running
%% head, folio, signature) is excluded from this edition by standing policy.
%% Weber, Algebra. II. \hfill 50


% ===== scan leaf 807, printed 786 =====
% !PRINTED 786
\weberpage{786}
\noindent
Biquadratic forms I, 199.\\
--- equations I, 119, 201, 528; II, 319.\\
--- cyclotomic fields II, 101.\\
\gesperrt{Bring-Jerrard} form of the equation of the fifth degree I, 176, 233.\\
\gesperrt{Brioschi} normal form of the equation of the fifth degree I, 236.\\
Fractions I, 10.\\
Literal calculation I, 20.\\
\gesperrt{Budan-Fourier} theorem I, 299.\\

\medskip
\noindent
\textbf{C}anonical form of the ternary cubic form II, 332.\\
Cardan formula I, 116.\\
Cartesian (\gesperrt{Harriot}) theorem on the number of real roots I, 308.\\
Casus irreducibilis of the cubic equations I, 608.\\
\gesperrt{Cayley} expression of the Cardan formula I, 118.\\
Characters of an \gesperrt{Abel}ian group II, 44.\\
Characteristic of a system of functions I, 285.\\
Classes of functionals II, 553.\\
--- --- ideals II, 553.\\
--- --- quadratic forms I, 400.\\
--- --- quadratic irrational numbers I, 383.\\
Class number of an algebraic field II, 554.\\
--- in the field of the eighth roots of unity II, 717.\\
Class number formula II, 702.\\
--- in the cyclotomic field II, 705.\\
--- in the real cyclotomic field II, 711.\\
--- in the field of the $2^{\varkappa}$th roots of unity II, 719.\\
Class number factors II, 716.\\
Class number factor $A$ II, 722.\\
--- $B$ II, 726.\\
Commutative groups I, 476; II, 6.\\
Commutative normal divisors of a metacyclic group II, 27.\\
Complexes of double tangents II, 357.\\

Complex roots (their bounding) I, 292.\\
--- numbers I, 18.\\
--- --- of \gesperrt{Gauss} I, 585.\\
Complex pairs, complex triples II, 363.\\
Composition in a group II, 3.\\
--- of the parts II, 12.\\
--- of quadratic irrational numbers II, 613.\\
--- of permutations I, 474.\\
--- of substitutions I, 470.\\
--- of linear substitutions II, 153.\\
Composition series of a group II, 17.\\
Configuration of the ternary substitution group of degree 168 II, 451.\\
Congruence of numbers I, 358, 368; II, 539.\\
Congruences of the first degree I, 368.\\
--- --- --- in algebraic fields II, 542.\\
Congruence group II, 250.\\
Congruence field II, 245.\\
--- of the second degree II, 259.\\
Congruence roots I, 426.\\
Conjugate functions I, 505.\\
--- groups I, 506; II. 10.\\ % [sic] printed full stop after the volume numeral
--- fields I, 460; II, 493.\\
--- logarithms II, 672.\\
--- poles of a linear substitution group II, 201.\\
Constancy of the index series II, 18.\\
Contragredient groups II, 480.\\
--- transformations II, 157, 478.\\
Contravariants II, 480.\\
Covariants I, 187.\\
--- of the ternary forms II, 327.\\
--- of the ternary cubic forms II, 333.\\
Cubic \gesperrt{Abel}ian equations I, 107.\\
--- forms, binary I, 192.\\
--- --- ternary II, 331.\\
--- equations I, 116, 522.\\
--- --- trigonometric solution I, 349.\\
--- cyclotomic fields II, 93.\\
Curves, algebraic II, 324.\\
Cyclic functions I, 531, 538.\\
--- equations I, 538.\\
--- groups I, 476.\\
--- --- of linear substitutions II, 209.\\


% ===== scan leaf 808, printed 787 =====
% !PRINTED 787
\weberpage{787}
\noindent
Cyclic groups in the congruence field II, 269, 275.\\
--- permutations I, 476, 493.\\

\medskip
\noindent
\gesperrt{\textbf{D}edekind} ideals II, 547.\\
\gesperrt{Dedekind}'s theorem on the field discriminant II, 578.\\
Derivative functions I, 51.\\
Derivative of a product I, 54.\\
--- of functions of several variables I, 59.\\
Determinants I, 68.\\
--- bordered I, 81.\\
--- from subdeterminants I, 97.\\
Determinant of a system of linear equations I, 84.\\
--- of a quadratic form I, 180.\\
Dense sets I, 4.\\
Dihedral group II, 210.\\
--- in the congruence field II, 276.\\
Differential quotients I, 54, 60.\\
Differences I, 46.\\
Difference product of the $m$th roots of unity I, 423.\\
\gesperrt{Dirichlet}'s theorem on the units II, 670.\\
Discrete sets I, 4.\\
Discriminant I, 150.\\
--- of the cubic form I, 37, 153, 193.\\
--- of the biquadratic form I, 155, 200.\\
--- of the quadratic irrational numbers I, 378; II, 602.\\
--- of the cyclotomic equation for a prime degree I, 421.\\
--- of an algebraic curve II, 325.\\
--- of an algebraic field II, 529.\\
--- of the cyclotomic field II, 619. 643.\\ % [sic] printed full stop instead of a comma after 619
Discriminants in an algebraic field II, 528.\\
Discriminant surface I, 249.\\
Division of numbers I, 1.\\
--- of integral functions I, 28.\\
Divisors of an \gesperrt{Abel}ian group II, 48.\\
--- of a group I, 476, 501; II, 7.\\
Rotations II, 186.\\
Rotation group II, 189.\\

Triangular coordinates II, 322.\\
% [sic] In print, Double points and Double tangents follow Rotation group and Triangular coordinates.
Double points of a curve II, 326.\\
Double tangents II, 327.\\
--- of a curve of the fourth order II, 351.\\
Intersection of groups I, 510; II, 9.\\

\medskip
\noindent
\textbf{P}roper and improper equivalence I, 371.\\
Proper and improper orthogonal substitutions II, 195.\\
Simple groups I, 511; II, 136.\\
--- --- of degree 60 II, 138.\\
--- \gesperrt{Abel}ian fields II, 649.\\
Simplicity of the alternating group I, 600.\\
--- --- congruence groups II, 254.\\
Unit of a group II, 5.\\
Units I, 399.\\
--- in the field $R(i)$ I, 586.\\
--- in algebraic fields II, 510.\\
--- in the real cyclotomic field II, 645.\\
--- in the field of the eighth roots of unity II, 718.\\
--- independent II, 677.\\
--- fundamental system II, 682.\\
Roots of unity I, 114, 408.\\
--- third I, 118.\\
--- primitive I, 411.\\
--- in the cyclotomic fields II, 638.\\
--- in the congruence field II, 260.\\
Elimination from linear equations I, 87.\\
--- from higher equations I, 157, 159.\\
--- from several equations I, 161.\\
Final equation I, 157.\\
Finite groups II, 4.\\
--- --- of linear substitutions II, 195.\\
Finiteness of the system of invariants of a linear group II, 169.\\
Opposite elements of a group II, 5.\\
\gesperrt{Euler}'s theorem on homogeneous functions I, 62.\\
Exponent system of units II, 679.\\

% Printer's signature “50*” excluded as page furniture.


% ===== scan leaf 809, printed 788 =====
% !PRINTED 788
\weberpage{788}
\noindent
System of exponents of numbers II, 685.\\
Exponential functions II, 752.\\

\medskip
\noindent
\textbf{F}actors of integral functions I, 39, 103.\\
--- of the natural primes in algebraic fields II, 576.\\
\gesperrt{Fermat}'s theorem I, 427; II, 55.\\
--- --- in the congruence field II, 247.\\
--- --- in algebraic fields II, 544.\\
Forms I, 58.\\
Form problem of the linear substitution group II, 175.\\
System of forms of the cubic form I, 196.\\
--- of the biquadratic form I, 203, 206.\\
Functions, integral I, 23.\\
--- --- decomposition into linear factors I, 103.\\
--- --- decomposition into prime functions I, 164, 452; II, 495.\\
--- homogeneous I, 56.\\
--- in a field I, 452; II, 520.\\
Function congruences II, 243.\\
Functional determinants I, 185.\\
Functionals II, 499.\\
Functional classes II, 552.\\
\gesperrt{Frobenius}'s theorems on groups II, 129, 134.\\
Fundamental theorem of algebra I, 121, 127, 296.\\
Fundamental systems of units II, 681.\\
--- of normal units II, 731.\\

\medskip
\noindent
\gesperrt{\textbf{G}alois} group I, 477, 511; II, 588.\\
--- resolvents I, 467.\\
--- field I, 465; II, 588.\\
Integral functions of one variable I, 23.\\
--- --- of several variables I, 24.\\
--- --- in a field II, 520.\\
--- algebraic numbers, II, 491.\\ % [sic] Comma before the volume numeral as printed.
--- functionals II, 504.\\
\gesperrt{Gauss} sums I, 575.\\
Fractional functions I, 32.\\
Ordered sets I, 4.\\

Weight of an invariant I, 187.\\
Lattice II, 561.\\
Equations I, 101.\\
--- linear homogeneous I, 81.\\
--- --- inhomogeneous I, 89.\\
--- of the third degree I, 116, 118, 522.\\
--- of the fourth degree I, 119, 201, 528.\\
--- of the fifth degree I, 233, 621; II, 404.\\
--- of the seventh degree II, 476, 481.\\
Degree of a group I, 472, 475; II, 4.\\
--- of an element of a group II, 10.\\
--- of a permutation I, 500.\\
--- of a prime ideal and prime functional II, 516.\\
\gesperrt{Gr\"affe} method of the approximate solution of an equation I, 344.\\
Limits I, 121.\\
Greatest common divisor of numbers I, 2.\\
--- --- --- of integral functions I, 34.\\
--- --- --- of groups I, 510.\\
--- --- --- of functionals II, 512, 519.\\
Greatest normal divisor of a group II, 17.\\
Fundamental curve of the group $G_{168}$ II, 456.\\
Fundamental forms of the polyhedral groups II, 205.\\
--- --- cyclic groups II, 210.\\
--- --- dihedral groups II, 210.\\
--- --- tetrahedral group II, 214.\\
--- --- octahedral group II, 216, 218.\\
--- --- icosahedral group II, 221, 230.\\
Fundamental ideal II, 579, 593.\\
Fundamental number of a field II, 529.\\
Group (in general) II, 3.\\
--- of the double-tangent equation II, 380.\\
--- of the triple equations II, 344.\\
--- of the tetrahedron II, 212.\\
--- of the octahedron II, 216.\\
--- of the icosahedron II, 220.\\
--- of an equation I, 477.\\
--- of a normal field II, 588.\\
--- of the cyclotomic fields II, 68.\\
--- of an ideal in the normal field II, 589.\\
--- of the ideal classes II, 557.\\
--- of degree 168 II, 282, 433.\\


% ===== scan leaf 810, printed 789 =====
% !PRINTED 789
\weberpage{789}

\noindent
Groups, \gesperrt{Abel}ian I, 476, 536; II, 6.\\
--- finite II, 4.\\
--- of permutations I, 475.\\
--- of substitutions I, 472.\\
--- of degree $p^{\alpha}$ II, 127.\\
--- of degree $p\,q$ II, 141.\\
--- of linear substitutions II, 152.\\
--- of orthogon. substitutions II, 184.\\
--- of linear fractional substitutions II, 189, 195.\\
Group characters II, 44.\\
Group table II, 115.

\medskip
\noindent
\textbf{S}emimetacyclic groups I, 616.\\
Principal axes of a ternary linear substitution II, 441.\\
--- of the group $G_{168}$ II, 444.\\
Principal equation of the fifth degree I, 175, 230.\\
Principal series II, 24.\\
Principal subdeterminants I, 256.\\
\gesperrt{Hermite}'s solution of \gesperrt{Sturm}'s problem I, 280.\\
\gesperrt{Hesse} curve II, 330.\\
--- determinant I, 186.\\
\gesperrt{Hesse-Cayley} notation for the double tangents II, 369.\\
\gesperrt{Hilbert}'s theorem II, 165.\\
Homogeneous functions I, 56.

\medskip
\noindent
\textbf{I}deals II, 547.\\
Ideal numbers after \gesperrt{Kummer} II, 551.\\
Identical group I, 496.\\
--- substitution II, 152.\\
Icosahedral equation II, 232, 418.\\
Icosahedral group II, 220, 279.\\
Icosahedral invariants II, 232.\\
Icosahedral resolvents II, 422, 426.\\
Imaginary numbers I, 17.\\
Imaginary form of the linear congruence group II, 266.\\
Imaginary congruence field of the second degree II, 259.\\
Imaginary quadratic irrational numbers I, 379.\\
Imprimitive groups I, 483, 516.\\
--- fields I, 460, 483.\\
Index of an invariant II, 163.\\
--- of a divisor of a group I, 502; II, 8.

Index moduli II, 61.\\
Index series of a group II, 17.\\
Indices I, 431.\\
--- of the elements of a group II, 44.\\
--- with respect to an odd prime power II, 54.\\
--- with respect to a power of 2 II, 58.\\
--- with respect to a composite modulus II, 60.\\
Inflexion points II, 326.\\
--- of a curve of the third order II, 331.\\
Interpolation I, 44, 98, 331.\\
Intransitive groups I, 482.\\
--- normal divisors I, 520.\\
Invariants I, 186.\\
--- of the binary cubic form I, 193.\\
--- --- --- biquadratic form I, 200, 206.\\
--- of an \gesperrt{Abel}ian group II, 42.\\
--- of finite groups of linear substitutions II, 161.\\
--- absolute and relative II, 162, 177.\\
--- in the wider sense II, 179.\\
--- of the ternary forms II, 328.\\
--- of the curves of the third order II, 337.\\
--- of the group $G_{168}$ II, 453, 461.\\
Invariant curves II, 454.\\
Inverse substitution II, 155.\\
Irrational ratios I, 10.\\
--- numbers I, 12, 370.\\
Irreducibility I, 453.\\
--- of pure equations I, 607.\\
--- of the cyclotomic equation I, 417; II, 768.\\
Isomorphic groups I, 477; II, 6.\\
Isomorphism (multiple) II, 15.

\medskip
\noindent
\gesperrt{\textbf{J}acobi}'s estimation of the number of roots between two limits I, 311.\\
\gesperrt{Jordan}'s theorem on composite groups II, 18.

\medskip
\noindent
\textbf{C}ontinued fractions I, 358.\\
--- for equivalent numbers I, 374.\\
Chains of principal subdeterminants I, 257.\\
--- \gesperrt{Sturm} I, 272.


% ===== scan leaf 811, printed 790 =====
% !PRINTED 790
\weberpage{790}

\noindent
Least common multiple of numbers I, 3.\\
--- --- --- of functionals II, 519.\\
--- --- --- of groups II, 29.\\
Field I, 449; II, 493.\\
Spherical functions I, 273.\\
\gesperrt{Kummer}'s theorem on the resolvents II, 632.\\
Cyclotomic equation I, 422, 554.\\
Cyclotomic field II, 67.\\
--- of given group II, 79.\\
--- of prime power degree II, 617.\\
Theory of the division of the circle I, 408, 554.\\
Crystallographic groups II, 241.

\medskip
\noindent
\gesperrt{\textbf{L}agrange}'s theorem on functions that admit the permutations of a group I, 508.\\
\gesperrt{Laguerre}'s theorems on equations with only real roots I, 322.\\
\gesperrt{Legendre} symbol I, 441.\\
Linear congruence groups I, 611; II, 250, 302.\\
--- --- homogeneous II, 303.\\
--- --- real II, 261.\\
--- --- imaginary form II, 266.\\
--- --- of degree 168 II, 282.\\
--- --- binary for the modulus 3 II, 305.\\
--- --- ternary for the modulus 2 II, 306.\\
--- fractional substitutions II, 189.\\
--- functions I, 30.\\
--- equations I, 81.\\
--- transformation I, 178.\\
--- substitutions I, 178, 371; II, 151.\\
\gesperrt{Lüroth}'s theorem II, 404.

\medskip
\noindent
\textbf{M}anifold I, 4.\\
Matrix I, 82.\\
Set I, 4.\\
Measurable set I, 8.\\
Metacyclic equations I, 597; II, 292.\\
--- --- of prime degree I, 609.\\
--- --- of the fifth degree I, 621, 648.\\
--- --- of the sixth degree II, 296.\\
--- --- of the eighth degree II, 314.\\
--- --- of the ninth degree II, 305.

Metacyclic groups I, 598; II, 27.\\
--- functions I, 617, 635.\\
Minimal basis of a field II, 529.\\
--- of the cyclotomic field II, 620.\\
--- of a quadratic field II, 603.\\
Minimum I, 122.\\
--- of a quadratic form II, 559.\\
--- of the fundamental number of a field II, 569.\\
Modulus I, 358.\\
Multiplicative substitutions II, 152.\\
Multiplication of numbers I, 14.\\
--- of determinants I, 92.\\
--- of trigonometric functions I, 433.

\medskip
\noindent
% the German alphabet division here is "N" (Näherungsbrüche); the bold section initial is
% mirrored onto the initial letter of the corresponding English entry
\textbf{C}onvergents of a continued fraction I, 362.\\
Method of approximation for the calculation of roots of equations by the regula falsi I, 331.\\
--- for the calculation of roots of equations by \gesperrt{Daniell Bernoulli} I, 341.\\ % [sic] printed "Daniell" rather than "Daniel"
--- \gesperrt{Gräffe} I, 344.\\
--- \gesperrt{Newton} I, 335.\\
--- by continued fractions I, 401.\\
Method of approximation for the solution of trinomial equations I, 352.\\
Natural irrationalities I, 516.\\
Subsidiary groups I, 502; II, 8.\\
Negative numbers I, 16.\\
Ninth roots of unity I, 582.\\
\gesperrt{Newton} formulae for the power sums I, 142.\\
\gesperrt{Newton} rule for the estimation of the roots I, 305.\\
Non-metacyclic equations I, 603.\\
Norm I, 461; II, 494, 498.\\
--- absolute II, 502, 538.\\
--- of the \gesperrt{Gauss} complex numbers I, 586.\\
--- of an ideal II, 551.\\
--- of a field I, 465.\\
--- of a quadratic irrational number I, 377.\\
Normal units in the cyclotomic field II, 729.


% ===== scan leaf 812, printed 791 =====
% !PRINTED 791
\weberpage{791}

\noindent
Normal forms of the linear substitution groups II, 181.\\
Normal equation I, 465.\\
Normal field I, 464; II, 588.\\
Normal divisor of a group I, 511; II, 11.\\
Zero I, 16.

\medskip
\noindent
% German alphabetical block "O" (Obere Grenze ...); the bold initial mirrors the printed device
\textbf{U}pper limit for the roots I, 316.\\
Octahedral group II. 216.\\ % [sic] printed full stop after the volume numeral
--- in the congruence field II, 277.\\
Orthogonal substitution II, 184.

\medskip
\noindent
% German alphabetical block "P" (Partialdiscriminante ...)
\textbf{P}artial discriminant II, 597.\\
Partial fundamental ideal II, 585, 597.\\
Partial norm II, 584.\\
Partial trace II, 585.\\
Partial resolvents I, 512.\\
\gesperrt{Pell} equation I, 395.\\
Period of a permutation I, 500.\\
Periods of the division of the circle I, 557, 579; II, 74.\\
--- of the roots of a cyclic equation I, 545.\\
--- of the reduced numbers I, 388.\\
Periodic continued fractions I, 387.\\
Permutations I, 64, 473, 500.\\
--- their analytic representation I, 613; II, 299.\\
--- of the first and second kind I, 65, 496.\\
Permutation groups I, 475, 489; II, 117.\\
--- of degree $168$ of seven figures II, 473.\\
Polars I, 188.\\
Poles of linear fractional substitutions II, 196.\\
--- of a group II, 198.\\
--- --- --- in the congruence field II, 274.\\
--- of a ternary group II, 438.\\
Polyhedral groups II, 202.\\
--- of the second kind II, 234.\\
--- in the congruence field II, 273.\\
Polynomial coefficients I, 51.\\
Polynomial theorem I, 50.\\
Positive units in the cyclotomic field II, 742.\\
Power residues I, 430.\\
Power sums I, 140.

Primary divisors of a cyclotomic field II, 71.\\
--- --- of an \gesperrt{Abel}ian group II, 72.\\
Prime factors of the numbers of an algebraic field II, 523.\\
--- of the natural primes II, 576.\\
--- of the cyclotomic resolvents II, 635.\\
Prime functionals II, 515, 576.\\
--- relative II, 514.\\
Prime ideals II, 550.\\
--- in the quadratic field II, 604.\\
--- in the relatively normal field II, 588.\\
--- in the divisors of a normal field II, 599.\\
--- in the cyclotomic fields II, 619, 622.\\
--- in the real cyclotomic field II, 643.\\
Primitive congruence roots I, 428; II, 546.\\
--- roots of unity I, 410.\\
--- functions I, 25; II, 497.\\
--- roots of primes I, 429.\\
--- --- of squares of primes II, 57.\\
--- --- of prime functionals II, 546.\\
--- --- of prime ideals II, 591.\\
--- --- of a congruence field II, 248.\\
--- and imprimitive characters in cyclotomic fields II, 720.\\
--- --- --- groups I, 484.\\
--- --- --- fields I, 463.\\
--- --- --- form problems II, 180.\\
Primes I, 1.\\
--- relative I, 359.\\
--- in the field $R(i)$ I, 587.\\
--- in the field of the third roots of unity I, 593.

\medskip
\noindent
% German alphabetical block "Q" (Quadrate im Congruenzkörper ...)
\textbf{S}quares in the congruence field II, 248.\\
Quadratic forms I, 179.\\
--- equations I, 116.\\
--- irrational numbers I, 377.\\
--- residues I, 444.\\
--- fields II, 601.\\
Quadrature of the circle II, 760.

\medskip
\noindent
% German alphabetical block "R" (Rationale Wurzeln ...)
\textbf{R}ational roots of an equation I, 403.


% ===== scan leaf 813, printed 792 =====
% !PRINTED 792
\weberpage{792}

\noindent
Rational numbers I, 12.\\
Rational ratio I, 10.\\
Domain of rationality I, 452.\\
Space of $n$ dimensions II, 560.\\
Reality of the roots I, 241.\\
--- --- --- for quadratic and cubic equations I, 243.\\
--- --- --- for biquadratic equations I, 246.\\
--- --- --- for cyclic equations I, 551.\\
--- --- --- for metacycl. equations I, 647.\\
--- --- --- of the double tangents of a curve of the fourth order II, 391.\\
Reality conditions of the orthogonal group II, 193.\\
Reality relations of the points of inflection of a curve of the third order II, 340.\\
--- for triple equations II, 348.\\
Operations of arithmetic I, 1.\\
Calculation with integral functions I, 24.\\
--- with numbers I, 14.\\
Law of reciprocity of the quadratic residues I, 443.\\
Reducible and irreducible functions I, 453.\\
--- --- --- with respect to a prime number modulus II, 243.\\
--- --- --- equations I, 404, 482.\\
Reduced numbers I, 380, 384.\\
--- units II, 680.\\
Real cyclotomic fields II, 641.\\
--- radicals I, 606.\\
Regula falsi I, 334.\\
Regular bodies II, 188.\\
Regulator of a system of units II, 678, 729.\\
--- of the field II, 682.\\
Pure equations I, 109.\\
Series II, 694.\\
Relatively prime numbers I, 2.\\
--- prime functionals II, 514.\\
Relative norm II, 584.\\
Relative trace II, 585.\\
Resolvents I, 511.\\
--- of the biquadratic equation I, 120.

Resolvents of the composition series II, 289.\\
--- of equations of the seventh degree II, 476.\\
--- of equations of the eighth degree II, 308.\\
--- of the group $G_{168}$ II, 466, 473.\\
--- of the icosahedral equation II, 419.\\
--- in the division of the circle I, 564; II, 63, 651.\\
--- of \gesperrt{Lagrange} I, 542.\\
--- --- for metacyclic. equations I, 631.\\
--- with one parameter II, 404.\\
Residue of the division of integral functions I, 28.\\
System of residues I, 359; II, 541.\\
Resultants I, 156.\\
Resultant of two quadratic functions I, 36, 158.\\
\gesperrt{Rolle}'s theorem on the number of real roots I, 319.

\medskip
\noindent
\textbf{S}ecular equation I, 276.\\
Closing numbers of a continued fraction I, 361.\\
Cut I, 5.\\
Seven-systems of double tangents II, 367.\\
Seventeen-gon I, 568.\\
Singular points of a curve II, 325.\\
Trace I, 461; II, 494.\\
\gesperrt{Steiner} complexes of double tangents II, 357.\\
Continuity I, 5.\\
--- of integral functions I, 105.\\
--- of the roots I, 132.\\
\gesperrt{Sturm} functions I, 279.\\
--- chains I, 271.\\
\gesperrt{Sturm}'s problem I, 270.\\
Substitution, linear I, 92, 372; II, 151.\\
--- of the ratios II, 158.\\
Substitutions of a normal field I, 468; II, 588.\\
Substitution determinant I, 92, 178.\\
\gesperrt{Sylow}'s theorems II, 121, 125.\\
Symmetric determinants I, 69.\\
--- functions I, 138, 144, 147.\\
--- fundamental functions I, 139.\\
--- group I, 492.


% ===== scan leaf 814, printed 793 =====
% !PRINTED 793
\weberpage{793}
\begin{flushleft}
Symmetric group, its normal divisors I, 602. \\
Syzygetic and azygetic systems of double tangents II, 362. \\
— — — complexes of double tangents II, 366. \\[6pt]
\textbf{T}angents of a curve II, 326. \\
\gesperrt{Taylor} development I, 59. \\
Ternary forms II, 322. \\
— linear groups of degree 168 II, 433. \\
— — congruence group of degree 168 II, 475. \\
Tetrahedral group II, 212. \\
— in the congruence field II, 276. \\
Divisibility of integral functions I, 32. \\
— of integral functionals II, 509. \\
— of integers II, 510. \\
Divisors of integers I, 1. \\
— of integers, greatest common I, 359. \\
— of integral functions I, 27; II, 497. \\
— of integral functions, greatest common I, 34. \\
— of a group I, 476, 501; II, 7. \\
— of the icosahedral group II, 227. \\
— of a field I, 450. \\
— of the linear congruence group II, 271. \\
Coprime numbers I, 2, 359. \\
— functions I, 35. \\
— functionals II, 514. \\
Partial denominator of a continued fraction I, 361. \\
Division of the angle I, 435. \\
Total resolvents I, 512. \\
Law of inertia of quadratic forms I, 183, 255. \\
Transcendental functions for the solution of the icosahedral equation II, 429. \\
— numbers II, 745. \\
Transcendence of the number $e$ II, 751. \\
— of the number $\pi$ II, 756. \\
Transformation of a group I, 506. \\
— of the quadratic form into a sum of squares I, 181. \\
— of forms of degree $n$ I, 185. \\
— of the cubic equation I, 219. \\
Transformation of the equation of the fifth degree I, 230. \\
Transformed substitutions II, 156. \\
Transitive and intransitive groups I, 482, 506. \\
— permutation group of index 6 of six figures I, 628. \\
Transposed substitutions II, 156. \\
Transpositions I, 65, 492. \\
Trigonometric solution of pure equations I, 111. \\
— — of cubic equations I, 349. \\
Trinomial equations I, 352. \\
Triple equations II, 342. \\
Triple systems II, 310. \\
\gesperrt{Tschirnhausen} transformation I, 170, 210. \\
— — of the cubic equation I, 219. \\[6pt]
\textbf{R}eversible periods I, 391. \\
Unknowns I, 20. \\
Indeterminate equations I, 365. \\
Infinite roots of an equation I, 136. \\
Infinity of integral rational functions I, 104. \\
Unicursal curves II, 416. \\
Subdeterminants I, 72. \\
— higher I, 77. \\
— complementary I, 79. \\
Uncountable sets II, 748. \\
Original functions I, 25. \\[6pt]
\textbf{V}ariable in field theory II, 499. \\
Ratios I, 10. \\
Related numbers and functionals II, 631. \\
Complete system of residues I, 359; II, 541. \\
Complete systems of double tangents II, 367. \\
Volume II, 561, 690. \\
Change of sign of integral functions I, 107. \\[6pt]
\textbf{P}oints of inflection II, 326. \\
Inflectional tangents II, 326. \\
— of a curve of the third order II, 331. \\
\gesperrt{Wilson}'s theorem I, 428. \\
Division of the angle I, 551, 609.
\end{flushleft}


% ===== scan leaf 815, printed 794 =====
% !PRINTED 794
\weberpage{794}
\begin{flushleft}
Duplication of the cube I, 609. \\
Roots of equations I, 101. \\
— of equations of odd degree I, 109. \\
— of pure equations I, 109. % [sic] printed \textquotedbl{}Gleichnngen\textquotedbl{} for \textquotedbl{}Gleichungen\textquotedbl{}
\\
— of metacyclic equations I, 638. \\
— rational I, 403. \\[6pt]
\textbf{N}umber I, 1, 12. \\
— of the terms of a homogeneous function I, 58. \\
— of the permutations I, 64. \\
Number classes with respect to a composite modulus II, 60.
\\
Sequences of numbers I, 15. \\
Number field I, 449; II, 527.
\end{flushleft}

\begin{flushleft}
Change of sign and succession of signs I, 262. \\
Decomposable and indecomposable functions of several variables I, 164. \\
Decomposition of integral functions into linear factors I, 102. % [sic] printed \textquotedbl{}Functioaen\textquotedbl{} for \textquotedbl{}Functionen\textquotedbl{}
\\
— of groups into subsidiary groups I, 502; II, 8. \\
— — — with respect to two divisors II, 120. \\
Composed Abelian fields II, 649. \\
Composition of linear substitutions I, 372; II, 153. \\
— of substitutions of a normal field I, 470. \\
— of permutations I, 473. \\
Two-sided numbers I, 390.
\end{flushleft}


% ===== scan leaf 816, printed 795 =====
\weberpage{795}% page number not printed on this page
\pdfbookmark[1]{Corrections}{bm241}
% !PRINTED none
\weberpage{none}
\begin{center}
\textbf{\gesperrt{Corrections to the First Volume.}}
\end{center}

\begin{center}
\rule{4cm}{0.4pt}
\end{center}

\begin{flushleft}
Page 93, line 6 from below, read $B_h^{(i)}$ instead of $B^{(i)}$. \\
„ 117, formula (10) read $\sqrt[3]{\dfrac{-b}{2}+\sqrt{R}}$ instead of $\sqrt{\dfrac{-b}{2}+\sqrt{R}}$. \\
„ 158, formulae (14), on the right $m$ and $n$ are to be interchanged. \\
„ 172, formula (14) read $n-1$ instead of $n+1$. \\
„ 185, line 9 from above „ $\Phi'$ instead of $\Phi$. \\
„ 190, „ 11 from below „ $a_1'$ „ $(n-1)a_1'$. \\
„ 194, „ 7 from below „ $+9$ „ $-9$. \\
„ 200, „ 7 from below „ \S.\,61 „ \S.\,67. \\
„ 202, „ 3 from below „ $u$ „ $v_1$ (in both formulae). \\
„ 213, „ 6 from below „ $\Phi(\tau,\xi)$ instead of $\Phi(\tau,k)$. \\
„ 221, „ 12 from below „ $\dfrac{2}{9}$ instead of $\dfrac{2}{3}$. \\
„ 222, „ 4,5 from below „ $+q_0, -q_0$ instead of $-q_0, +q_0$. \\
„ 233, „ 4,7 from above „ $\tfrac{1}{5}D$ instead of $5D$. \\
„ 236, „ 12 from above „ $5c^2$ „ $-5c^2$. \\
„ 251, „ 10 from above „ $-\psi_3$ „ $\psi_3$. \\
„ 268, „ 5 from below „ $\tfrac{1}{5}D$ „ $D$. \\
„ 282, formula (2) „ $f_s$ „ $f^s$. \\
„ 318, line 2 from below „ $F(z)$ instead of $F(x)$. \\
„ 332, formula (5) „ $f(\beta)-$ instead of $f(\beta)+$. \\
„ 334, line 6 from below „ 0,00164 „ 0,0164. \\
„ 344, „ 6 from below, formula, read $\alpha_n\varphi(\alpha_n)$ and $\varphi(\alpha_n)$ instead of $\alpha_m\varphi(\alpha_m)$
in $\varphi(\alpha_m)$. \\
„ 357, „ 3 from below read 6 instead of 5. \\
„ 378, 379, formulae (13), (14), $a_1, b_1, c_1$ are to be replaced by $\pm a_1, \pm b_1, \pm c_1$,
and the sign is to be so determined that $c_1$ becomes positive. \\
„ 402, line 6,7 from below, $a_4=1$ and $a_3=2$ are to be interchanged. \\
„ — „ 4 from below read 9083 instead of 9183. \\
„ 419, on this compare the I. Supplement to the second volume. \\
„ 423, line 15 from above read $\dfrac{(n-1)(n-2)}{2}$ instead of $\dfrac{n\cdot n-1}{2}$. \\
„ — „ 24 from above there is to be read
$\dfrac{1}{2}\left(\dfrac{(n-1)(n-2)}{2}-\dfrac{n-1}{2}\right)
=\dfrac{(n-1)(n-3)}{4}\equiv 0\ (\mathrm{mod}\ 2)$.
\end{flushleft}


% ===== scan leaf 817, printed 796 =====
% !PRINTED 796
\weberpage{796}
\begin{flushleft}
Page 423, line 3 from below, read $(-1)^{\frac{(n-1)(n-3)}{4}}=+1$ instead of $(-1)^{\frac{n-1}{2}}$. \\
„ 424, Formula (14) „ $P=i^{\frac{n-1}{2}}n^{\frac{n-3}{2}}\sqrt{n}$. \\
„ 436, line 3, 22 from above, read $\dfrac{\pi}{n}$ instead of $\pi$. \\
„ 482, „ 8 from below, transitive and intransitive are to be interchanged. \\
„ 502, „ 18 from above, read $\pi_1$ instead of $\pi$. \\
„ 504, „ 9 from below, „ $k_{q-1}$ instead of $k_{\nu-1}$. \\
„ 514, „ 5 from below, „ $\Theta_1$ instead of $\Theta_i$. \\
„ 531, „ 11 from above, „ $y_1^2$ „ $y^2$. \\
„ 540, Formula (10) „ $\gamma_{e-1}$ instead of $\gamma_{c-1}$. \\
„ 544, line 15 from above, „ $mB_k=\overset{\varepsilon}{\Sigma}\varepsilon^{-k}(\varepsilon,\alpha)^\nu
(\varepsilon^{\lambda_1},\alpha)^{\nu_1}(\varepsilon^{\lambda_2},\alpha)^{\nu_2}\ldots$ \\
„ — „ 23 from above, „ $\varepsilon^{\lambda_2}$ instead of $\varepsilon^2$. \\
„ 551, „ 15 from above, „ $\alpha_k$ „ $\alpha^k$. \\
„ 592, „ 8 from below, „ which is satisfied only for $b=0$, $a=\pm 1$ or
$a=0$, $b=\pm 1$ or $a=b=\pm 1$. \\
„ 612, „ 12 from above, „ $1-a$ instead of $a-b$. \\
„ 618, „ 1 from below, „ $c_1^h$ instead of $c_1$. \\
„ 623, „ 6 from below, „ $a_6$ „ $\alpha_6$.
\\
„ 636, „ 18, 21 from above, read $-g^{n-1}+1$ instead of $g^{n-1}-1$. \\
„ 639, „ 10, 27 from above, „ $n-2$, $n-3$ instead of $n-1$, $n-2$. \\
„ 641, „ 1 from above, read (6) instead of (5). \\
„ 646, „ 11 from above, „ rational in $\mathfrak{K}(f_0)$ instead of rational. \\
„ — „ 14, 19, 22, 26 from above, read $\mathfrak{K}(f_0)$ instead of $\mathfrak{K}$. \\
„ — „ 17, 18 from above, from ``All quantities'' on to be deleted. \\
„ — „ 1, 2 from below, delete from ``the others'' onward. \\
„ 653, „ 4 from above, read $-A_4 r\varrho$ instead of $+A_4 r\varrho$. \\
„ — „ 5 from above, „ $+A_2\varrho'$ „ $-A_2\varrho'$.
\end{flushleft}

\begin{center}
\rule{6em}{0.4pt}
\end{center}

\begin{center}
\textbf{\gesperrt{Corrections to the second volume.}}
\end{center}

\begin{center}
\rule{6em}{0.4pt}
\end{center}

\begin{flushleft}
Page 15, line 19 from above, read $N$ instead of $P$. \\
„ 16, „ 12, 13 from above, $P$ and $Q$ are to be interchanged. \\
„ 21, „ 19 from above, read (17) instead of (13). \\
„ 26, „ 17 from above, „ \S.\,6, 4 instead of \S.\,6, 2. \\
„ 569, „ 5 from below, in the exponent read $n-\nu$ instead of $\mu-\nu$. \\
„ 592, „ 6, 3 from below, read $\gamma^{P^{f-1}}$ instead of $\gamma^{P^f-1}$.
\end{flushleft}

\begin{center}
\rule{6em}{0.4pt}
\end{center}


% ===== scan leaf 818, printed None =====
% !PRINTED none
% Publisher's advertisement page of Friedrich Vieweg & Sohn in Braunschweig; no folio. The
% publisher's masthead and a short centred rule are followed by five complete advertisements,
% separated by short centred rules:
% 1. “Elliptic Functions and Algebraic Numbers: Academic Lectures,” by H. Weber, Professor
% of Mathematics at the University of Strasbourg; royal octavo; sewn; price 13 marks.
% 2. “Electrodynamics, with Consideration of Thermoelectricity, Electrolysis, and
% Thermochemistry,” by Dr Heinrich Weber, Professor at the Ducal Technical College in
% Braunschweig; with wood engravings; royal octavo; sewn; price 6 marks.
% 3. “Five Popular Scientific Lectures,” delivered in the auditorium of the Ducal Technical
% College in Braunschweig, by Dr Heinrich Weber, Professor of Physics at the Ducal Technical
% College; with 84 illustrations; octavo; sewn; price 2 marks 50 pfennigs. Contents: an
% overview of the origin and effects of galvanic current; telegraphy and telephony;
% electromagnetic, magnetoelectric, and dynamo machines; electroplating and electric light;
% on perpetual motion.
% 4. “Continuity and Irrational Numbers,” by Richard Dedekind, Professor of Mathematics at
% the Technical College in Braunschweig; second unaltered edition; royal octavo; sewn;
% price 1 mark. The edition statement is a separate capitalised display line.
% 5. “A Treatise on Differential Equations,” by Dr Andrew Russell Forsyth, Professor at
% Trinity College, Cambridge; with an appendix containing the results of the exercises cited
% in the textbook; edited by H. Maser; authorised translation; royal octavo; sewn;
% price 14 marks.


% ===== scan leaf 819, printed None =====
% !PRINTED none
% Publisher's advertisement page of Friedrich Vieweg & Sohn in Braunschweig; no folio. The
% publisher's masthead and a short centred rule are followed by six complete advertisements,
% separated by short centred rules:
% 1. “What Are Numbers and What Should They Be?” by Richard Dedekind, Professor of
% Mathematics at the Technical College in Braunschweig; second unaltered edition; octavo;
% sewn; price 1 mark 60 pfennigs.
% 2. “Lectures on Number Theory,” by P. G. Lejeune-Dirichlet; edited and supplemented by
% R. Dedekind, Professor at the Carolo-Wilhelmina Technical College in Braunschweig; fourth
% revised and enlarged edition; royal octavo; sewn; price 14 marks.
% 3. “Textbook of Political Arithmetic,” prepared for higher commercial schools (commercial
% academies) and self-instruction by F. S. Holzinger, Professor at the Public Commercial
% Academy in Linz; royal octavo; sewn; price 3 marks 50 pfennigs.
% 4. “Guide to Surveying the Heavens: Astronomical Objects for Ordinary Telescopes. A
% Handbook and Aid for All Friends of Astronomy, Especially Owners of Telescopes,” by
% Dr Hermann J. Klein; second improved edition; with 75 wood engravings, 5 plates, partly
% colour-printed, 4 star charts, and 1 frontispiece; octavo; sewn; price 24 marks.
% 5. “Collection of Formulas of Pure and Applied Mathematics,” by Dr W. Láska; with three
% plates; royal octavo; sewn, 26 marks; half-morocco bound, 28 marks.
% 6. “Stargazing: Past and Present,” by J. Norman Lockyer, Fellow of the Royal Society and
% corresponding member of the Institute of France; authorised German edition; translated
% by G. Siebert; with 217 wood engravings; octavo; sewn; price 18 marks.


% ===== scan leaf 820, printed None =====
% !PRINTED none
% Publisher's advertisement page of Friedrich Vieweg & Sohn in Braunschweig; no folio. The
% publisher's masthead and a short centred rule are followed by five complete advertisements,
% separated by short centred rules:
% 1. “Practical Guide to Making Astronomical Observations, with Particular Regard to
% Astrophysics, Together with a Modern Account of Instruments,” by Nicolaus von Konkoly,
% PhD, member of the Academy of Sciences in Budapest, the Royal Astronomical Society in
% London, etc.; with 345 wood engravings; royal octavo; sewn; price 24 marks.
% 2. “Lectures and Addresses,” by Hermann von Helmholtz; fourth edition; with a portrait of
% the author and numerous wood engravings printed in the text; two volumes; royal octavo;
% per volume, sewn 8 marks, bound 9 marks 50 pfennigs.
% 3. “The Moon and the Condition and Configurations of Its Surface,” by Edmund Neison,
% member of the Royal Astronomical Society in London, etc.; authorised German original
% edition; second edition, enlarged by an appendix containing the author's investigations
% into the new formation Hyginus N on the Moon; with a quarto atlas (4°) of 26 maps and
% 5 colour-printed plates; royal octavo; sewn; 18 marks together with the atlas.
% 4. “Photogrammetry and International Cloud Measurement,” by Dr Carl Koppe, Professor at
% the Ducal Technical College in Braunschweig; with illustrations and five plates; royal
% octavo; sewn; price 7 marks.
% 5. “Five-Place Logarithmic and Trigonometric Tables,” edited by Dr O. Schlömilch, retired
% Royal Saxon privy councillor and member of the Royal Saxon Society of Sciences in Leipzig,
% the Royal Swedish Academy in Stockholm, the Imperial Leopoldine Academy, etc.; electrotype
% stereotyping; inexpensive school edition; twelfth edition; octavo; sewn; price 1 mark.
% This advertisement is complete on the present page.


% ===== scan leaf 821, printed None =====
% !PRINTED none
% Publisher's advertisement page of Friedrich Vieweg & Sohn in Braunschweig; no folio. The
% publisher's masthead and a short centred rule are followed by five complete advertisements:
% 1. “Partial Differential Equations and Their Application to Physical Questions: Lectures,”
% by Bernhard Riemann; prepared for the press and edited by Karl Hattendorff; third edition;
% with wood engravings; royal octavo; sewn; price 8 marks.
% 2. “Polydimensional Quantities and Perfect Numbers,” by Dr Hermann Scheffler, member of
% the Imperial Society of Naturalists in Moscow and corresponding member of the Royal
% Academy of Sciences in Padua; with 42 wood engravings; royal octavo; sewn;
% price 5 marks 60 pfennigs.
% 3. “Compendium of Higher Analysis,” by Dr Oskar Schlömilch, retired Royal Saxon privy
% councillor and member of the Royal Saxon Society of Sciences in Leipzig, the Royal Swedish
% Academy in Stockholm, the Imperial Leopoldine Academy, etc.; in two volumes; with wood
% engravings; royal octavo; sewn; first volume, fifth improved edition, 9 marks; second
% volume, fourth edition, 9 marks.
% 4. “The Book of Nature, Embracing the Teachings of Physics, Astronomy, Chemistry,
% Mineralogy, Geology, Botany, Zoology, and Physiology,” dedicated to all friends of natural
% science, especially grammar schools, secondary schools, and higher civic schools, by
% Dr Friedrich Schoedler, Director of the Grand-Ducal Hessian First-Order Secondary School
% in Mainz. Part I: Physics, Astronomy, and Chemistry; twenty-second improved edition with
% the author's portrait; 404 wood engravings, a colour-printed spectral plate, star charts,
% and a lunar chart; royal octavo; sewn; 4 marks 80 pfennigs. Part II: Mineralogy, Geology,
% Botany, Zoology, and Physiology; twenty-second improved edition; 683 wood engravings and a
% colour-printed geognostic plate; royal octavo; sewn; 4 marks 80 pfennigs. The separator
% before the next advertisement shares the second part's price line.
% 5. “The Fundamental Theories of Modern Geometry and the Elements of the Theory of Conic
% Sections, Prepared for School Instruction,” by H. Seeger; with 60 wood engravings; royal
% octavo; sewn; price 2 marks 80 pfennigs. The advertisement is complete and followed by
% blank space.


% ===== scan leaf 822, printed None =====
% !PRINTED none
% Scan shows only faint speckling/noise on all three bands; no printed text, no running head, no legible content. Appears to be a blank or near-blank leaf.


% ===== scan leaf 823, printed None =====
% !PRINTED none
% A library "NON-CIRCULATING BOOK" ownership stamp lies across the upper part of this otherwise blank leaf. Faint verso show-through is visible but not legible. A pencil call-number annotation, best read as "QH 11-80", appears near the lower-left corner. No printed page number or Weber body text is present.


% ===== scan leaf 824, printed None =====
% !PRINTED none
% The leaf is blank apart from faint speckling, tonal variation, and the dark scanned gutter at the left edge. No handwriting, printed text, running head, page number, or Weber body content is present.


% ===== scan leaf 825, printed None =====
% !PRINTED none
% This scan is the library's rear pastedown and due-date slip, set against a marbled endpaper. The heading reads "RETURN TO -> Astronomy/Mathematics/Statistics/Computer Science Library, 100 Evans Hall, 642-3381". In the six-cell loan-period table, the first cell's "7 DAYS" is struck out and the second is stamped "1 MONTH". Below "ALL BOOKS MAY BE RECALLED AFTER 7 DAYS" and "DUE AS STAMPED BELOW" are a cancelled "MAR - 5 1988", heavily overprinted and only partly legible 1989 stamps, a cancelled "OCT 14 1998", and "OCT 05 2002". The footer reads "FORM NO. DD3, 1/83" and "UNIVERSITY OF CALIFORNIA, BERKELEY / BERKELEY, CA 94720", with a small circled-P device. No Weber text is present.

\end{document}
